\documentclass[11pt,oneside]{book}
\usepackage[margin=1in]{geometry}
\usepackage{amsmath,amssymb,amsthm,mathtools,mathrsfs}
\usepackage{fontspec}
\setmonofont[Path=fonts/,Scale=0.85]{JetBrainsMono.ttf}
\usepackage{newunicodechar}
\usepackage{booktabs,longtable,array,enumitem}
\usepackage{xcolor}
\usepackage[colorlinks=true,linkcolor=blue!50!black,urlcolor=blue!50!black,citecolor=blue!50!black]{hyperref}
\setlist{itemsep=2pt,topsep=4pt}
\sloppy\emergencystretch=5em

\newtheorem{theorem}{Theorem}[chapter]
\newtheorem{proposition}[theorem]{Proposition}
\newtheorem{lemma}[theorem]{Lemma}
\newtheorem{corollary}[theorem]{Corollary}
\theoremstyle{definition}
\newtheorem{definition}[theorem]{Definition}
\newtheorem{example}[theorem]{Example}
\theoremstyle{remark}
\newtheorem{remark}[theorem]{Remark}

\newcommand{\Lean}[1]{\texttt{#1}}
\newcommand{\leanref}[3]{\par{\footnotesize\raggedright\noindent\textsf{Lean:}\ \texttt{#1}\ \ (\texttt{#2}, lines #3)\par}}
\newcommand{\Q}{\mathbf{Q}}\newcommand{\Z}{\mathbf{Z}}\newcommand{\C}{\mathbf{C}}
\newcommand{\R}{\mathbf{R}}\newcommand{\N}{\mathbf{N}}\newcommand{\F}{\mathbf{F}}
\newcommand{\Sha}{\mbox{\fontencoding{U}\fontfamily{wncy}\selectfont Sh}}
\renewcommand{\Sha}{\mathrm{Sha}}
\renewcommand{\Re}{\operatorname{Re}}\renewcommand{\Im}{\operatorname{Im}}
\DeclareMathOperator{\Gal}{Gal}\DeclareMathOperator{\GL}{GL}\DeclareMathOperator{\U}{U}
\DeclareMathOperator{\Frob}{Frob}
\newcommand{\ASK}{\mathrm{ASK}}\newcommand{\Main}{\mathrm{Main}}
\newcommand{\Ga}{\mathrm{Ga}}\newcommand{\Fo}{\mathrm{Fo}}\newcommand{\Ee}{\mathrm{Ee}}
\newcommand{\sg}{\mathrm{sg}}\newcommand{\gainD}{\mathrm{gainD}}

% common Unicode characters that may appear in running text
\newunicodechar{ψ}{\ensuremath{\psi}}\newunicodechar{χ}{\ensuremath{\chi}}\newunicodechar{ε}{\ensuremath{\varepsilon}}
\newunicodechar{ℓ}{\ensuremath{\ell}}\newunicodechar{≤}{\ensuremath{\le}}\newunicodechar{≥}{\ensuremath{\ge}}
\newunicodechar{→}{\ensuremath{\to}}\newunicodechar{∈}{\ensuremath{\in}}\newunicodechar{×}{\ensuremath{\times}}
\newunicodechar{θ}{\ensuremath{\theta}}\newunicodechar{ρ}{\ensuremath{\rho}}\newunicodechar{σ}{\ensuremath{\sigma}}
\newunicodechar{μ}{\ensuremath{\mu}}\newunicodechar{λ}{\ensuremath{\lambda}}\newunicodechar{π}{\ensuremath{\pi}}
\newunicodechar{Φ}{\ensuremath{\Phi}}\newunicodechar{Λ}{\ensuremath{\Lambda}}\newunicodechar{ζ}{\ensuremath{\zeta}}
\newunicodechar{γ}{\ensuremath{\gamma}}\newunicodechar{δ}{\ensuremath{\delta}}\newunicodechar{η}{\ensuremath{\eta}}
\newunicodechar{α}{\ensuremath{\alpha}}\newunicodechar{β}{\ensuremath{\beta}}\newunicodechar{τ}{\ensuremath{\tau}}
\newunicodechar{ω}{\ensuremath{\omega}}\newunicodechar{ξ}{\ensuremath{\xi}}\newunicodechar{κ}{\ensuremath{\kappa}}
\newunicodechar{≠}{\ensuremath{\ne}}\newunicodechar{∞}{\ensuremath{\infty}}\newunicodechar{⊂}{\ensuremath{\subset}}
\newunicodechar{⊆}{\ensuremath{\subseteq}}\newunicodechar{∑}{\ensuremath{\sum}}\newunicodechar{√}{\ensuremath{\surd}}\newunicodechar{ℝ}{\ensuremath{\mathbb R}}


\begin{document}

\begin{titlepage}\centering\vspace*{3cm}{\LARGE\bfseries The Riemann hypothesis and GRH for Dirichlet $L$-functions\par}\vspace{1em}{\Large A readable account of the Lean proof in Levent Alpöge's archive\par}\vspace{2em}{\large rh-lean-roadmap\par}\vspace{1em}{November 2026\par}\end{titlepage}

\tableofcontents

\chapter*{Introduction}\addcontentsline{toc}{chapter}{Introduction}\label{ch:intro}

On 20 November 2026 Levent Alpöge announced a proof, found with Claude, of the Riemann hypothesis and of the generalized Riemann hypothesis for Dirichlet $L$-functions. The proof is a Lean~4 archive of 175{,}187 files and 17.4 million lines, without comments, checked by several independent kernels \cite{sutherland-2026}; on 21 November 2026 the Lean and Mathlib communities, with the endorsement of the Lean FRO and the Mathlib Initiative, certified that its statement and proof are correct \cite{leanfro-2026}. This book gives the argument of the archive as mathematics: every statement the proof needs, in the order in which the proof uses it, with a proof written for a reader and the location of the Lean declaration that proves it.

\section*{The theorem}

The archive proves, for Mathlib's definitions of $\zeta$ and of $L(s,\chi)$:
\begin{itemize}
\item if $\zeta(s)=0$, $s\neq1$ and $s$ is not a negative even integer, then $\Re s=\tfrac12$;
\item for every Dirichlet character $\chi$ of any modulus, $L(s,\chi)\ne0$ when $\Re s>\tfrac12$ (for the principal character, $s=1$ is excluded, where Mathlib assigns a value to the pole);
\item every zero of $L(s,\chi)$ with $0<\Re s<1$ lies on $\Re s=\tfrac12$, for all $\chi$ and for primitive $\chi$.
\end{itemize}
All five statements follow from one theorem: no Dirichlet $L$-function vanishes to the right of the critical line.

\section*{The argument in six steps}

\begin{enumerate}
\item \emph{The critical line} (Part~I). For primitive $\chi$ the functional equation reflects zeros in $0<\Re s<\tfrac12$ to $\tfrac12<\Re s<1$, and imprimitive $L$-functions differ from primitive ones by Euler factors without zeros in $\Re s>0$. It suffices to show $L(s,\chi)\ne0$ for $\Re s>\tfrac12$ and primitive $\chi$.
\item \emph{A window criterion} (Part~I). Fix a primitive character $\psi$ of order $o$ and let $\mathcal X$ be the $o$ primitive characters induced by its powers. A Landau-type argument shows: if a weighted, centred prime sum $G_m(y)$ satisfies $G_m(y)\le y^{\theta}$ on whole intervals $[Y,Y^{3/2}]$ for arbitrarily large $Y$, then no $L(s,\chi)$, $\chi\in\mathcal X$, vanishes for $\Re s>\theta$. Only a one-sided bound is needed, and only on a sequence of windows.
\item \emph{The prime sum over Frobenius classes} (Part~II). For each point $x$ of a window the prime sum is distributed over the Frobenius classes of a finite abelian Galois group. Orthogonality of characters turns the count in each class into a sum over characters; averaging over a cyclic subgroup leaves the target family $\mathcal X$ together with three families of companion characters twisted by the quadratic characters of two auxiliary primes $q,d$. The auxiliary primes are chosen for each window by a mean-square estimate over candidate primes and a zero-density estimate, so that every twisted companion obeys a square-root bound on the whole window.
\item \emph{Multiplicity from theta lifts} (Parts~III and~IV). The proof follows Stepanov's method for curves over finite fields, in which an auxiliary function vanishes to high order at many points. Over the integers the archive builds the multiplicity from a chain of about $\sqrt x$ theta lifts between definite unitary groups. Each lift applies a differential operator, given by the weights, to a local theta kernel; the factorials it produces, minus what the denominators of a formal-group expansion can cancel, give a computable lower bound for the divisibility of the periods by each prime, valid on a whole lattice. The input lattice is then divided by these prime powers, and the pairing stays integral against a global lattice glued across Hecke characters.
\item \emph{Two bounds on one slope} (Part~V). Division raises the average slope of the lattices by a sum of logarithms of primes, which is the prime sum of step~3 up to a small loss. The size of the period maps at the complex places, through a slope inequality for the whole family of maps, bounds the same slope from above. This is Stepanov's inequality ``multiplicity times number of points is at most the degree''. The main terms on the arithmetic and geometric sides are the same number, and the chain length $\sqrt x$ puts the error at square-root size.
\item \emph{Extra divisibility} (Part~VI). Divisibility of special values that does not come from the chain is bounded by an estimate for the total length of Tate--Shafarevich-type groups over a family of CM subfields of cyclotomic fields.
\end{enumerate}

\section*{The three inputs}

In Lean, steps~3 to~6 start from three statements, proved elsewhere in the archive:
\begin{itemize}
\item \textsf{Green}: the classification of the cuspidal characters of $\GL_d(\F_p)$ by irreducible polynomials (Part~VII);
\item \textsf{VWin}: the window bound for Tate--Shafarevich-type groups (Part~VI);
\item \textsf{HypsAlt}: a structure of 46 fields asserting that one family of data satisfies all the conditions of the theta chain at once (Part~VIII).
\end{itemize}
The proofs of \textsf{VWin} and \textsf{HypsAlt} use Galois representations attached to polarized cuspidal automorphic representations and a local Langlands correspondence; Part~IX describes the forms in which the archive proves them.

\section*{How the book is written}

\begin{itemize}
\item The arguments made for this proof are given with proofs. Results from the literature (prime number theorems in progressions, zero-density estimates, Siegel's mass formula, Galois representations, local Langlands, Kato's Euler system, foundations of étale cohomology) are stated in the form the archive uses, with references, and a sentence records how the archive's version compares with the published one.
\item Statements are numbered by node identifiers such as \ref{B3.12}, which group them by part. After each statement, a line of the form \textsf{Lean:} \Lean{name} gives the archive declarations that prove it; the Index of Lean declarations at the end of the appendix gives, for each statement, the module and line numbers of each declaration. The archive often contains several versions of a result; the one cited is the version on the path of the final proof. The archive version is the one with SHA-256 \Lean{9631b4f0\ldots c7fc}.
\item Lean engineering (coercions, bookkeeping, routine inequalities) is absorbed into the proofs.
\item The main text keeps every statement and the proofs of moderate length. Longer proofs are given in the appendix; in the main text an outline of each says how the proof goes and which results it uses.
\item Chapter~\ref{ch:p3c0} collects the standard facts on the Weil representation and theta lifts that Parts~III, IV and~VIII use.
\end{itemize}

\section*{Reading order}

Part~I and the first chapters of Parts~II and~V give the frame of the argument. Parts~III, IV and~V contain the new mechanism: the local divisibility of one theta lift, its accumulation along the chain and the slope inequality that turns it into a bound for prime sums. Parts~VI to~IX treat the inputs.

\section*{About this text}

The text was written by several AI models, including Claude, working from the Lean source. An independent author, who is not a specialist in this area, set the goals and the editorial rules (plain mathematical English, no hedging, every statement tied to the Lean declaration that proves it) and organised the review: models working in isolation drafted and revised the chapters, and other models read anonymised copies of the whole book as referees, in five rounds; every point they raised was settled against the Lean source. No human expert has reviewed the text. Where the text and the Lean archive differ, the archive is right; the proofs of a few deep lemmas in Parts V, VI and VIII are still given in outline, and the Lean declarations cited after these statements prove them. The aim is a readable account for mathematicians; shorter versions of the same argument, down to one for readers without mathematics, are listed in the repository. Alpöge works at Anthropic; this text is independent of him and of Anthropic. The mathematics is Alpöge's archive; the exposition is ours.


\part{The five statements and the window criterion}

% ---- p1c1 ----
\chapter{Five statements and the critical line}\label{ch:p1c1}

Let $\mathbf1_N$ be the principal character of positive modulus $N$.
Write $L(s,\chi)$ for the continued Dirichlet $L$-function and
$\chi^\ast$ for the primitive character inducing $\chi$.
The assertion $H$ is defined in \ref{B1.10}.

\section{The five statements}

\begin{theorem}[Riemann hypothesis]\label{B1.01}
Every zero $s$ of $\zeta$ with $s\ne1$ and
$s\notin\{-2,-4,-6,\ldots\}$ satisfies $\Re s=1/2$.
\end{theorem}
\leanref{riemannHypothesis}{Main}{4--4}
\begin{proof}
Apply Lemma~\ref{B1.06} to Theorem~\ref{B1.09}.
\end{proof}

\begin{theorem}[Explicit zeta statement]\label{B1.02}
For every $s\in\C$,
\[
 \zeta(s)=0,\qquad
 \neg\exists n\in\N\;\;s=-2(n+1),\qquad s\ne1
 \quad\Longrightarrow\quad \Re s=\tfrac12 .
\]
Here $\N$ includes zero.
\end{theorem}
\leanref{riemannZeta\_\allowbreak{}eq\_\allowbreak{}zero\_\allowbreak{}re}{Main}{6--6}
\leanref{RiemannZetaZerosOnLine}{Statements}{3--4}
\begin{proof}
Apply Theorem~\ref{B1.01} with the negative even integers indexed by $n\in\N$.
\end{proof}

\begin{theorem}[Nonvanishing for every character]\label{B1.03}
Let $\chi$ be a complex Dirichlet character modulo $N\ge1$,
and let $s\in\C$. Then
\[
 \Re s>\tfrac12,\qquad
 (\chi\ne\mathbf1_N\ \text{or}\ s\ne1)
 \quad\Longrightarrow\quad L(s,\chi)\ne0 .
\]
\end{theorem}
\leanref{LFunction\_\allowbreak{}ne\_\allowbreak{}zero\_\allowbreak{}of\_\allowbreak{}half\_\allowbreak{}lt\_\allowbreak{}re}{Main}{8--8}
\leanref{DirichletLNonvanishingRightOfHalf}{Statements}{6--8}
\begin{proof}
This is Theorem~\ref{B1.09}, with $H$ expanded.
\end{proof}

\begin{theorem}[Zeros in the strip for every character]\label{B1.04}
Let $\chi$ be a complex Dirichlet character modulo $N\ge1$,
and let $s\in\C$. Then
\[
 0<\Re s<1,\qquad L(s,\chi)=0
 \quad\Longrightarrow\quad \Re s=\tfrac12 .
\]
\end{theorem}
\leanref{LFunction\_\allowbreak{}eq\_\allowbreak{}zero\_\allowbreak{}re}{Main}{10--10}
\leanref{DirichletLZerosOnLine}{Statements}{10--12}
\begin{proof}
Apply Lemma~\ref{B1.07} to Theorem~\ref{B1.09}.
\end{proof}

\begin{theorem}[Zeros in the strip for primitive characters]\label{B1.05}
Let $\chi$ be a primitive complex Dirichlet character modulo $N\ge1$,
and let $s\in\C$. Then
\[
 0<\Re s<1,\qquad L(s,\chi)=0
 \quad\Longrightarrow\quad \Re s=\tfrac12 .
\]
\end{theorem}
\leanref{LFunction\_\allowbreak{}eq\_\allowbreak{}zero\_\allowbreak{}re\_\allowbreak{}of\_\allowbreak{}isPrimitive}{Main}{12--12}
\leanref{PrimitiveDirichletLZerosOnLine}{Statements}{14--16}
\begin{proof}
Apply Lemma~\ref{B1.08} to Theorem~\ref{B1.09}.
\end{proof}

\section{Reduction to a right half-plane}

\begin{lemma}[The character of modulus one]\label{B1.06}
The assertion $H$ implies the Riemann hypothesis.
\end{lemma}
\leanref{Formalism.\allowbreak{}VW.\allowbreak{}rh\_\allowbreak{}of\_\allowbreak{}every}{Formalism.\allowbreak{}GammaFactorNonvanishing.\allowbreak{}CritLine}{29--32}
\begin{proof}
The unique character modulo $1$ satisfies $L(s,\mathbf1_1)=\zeta(s)$.
Apply $H$ to $\mathbf1_1$ for $\Re s>1/2$ and $s\ne1$, then use
Lemma~\ref{B1.11}.
\end{proof}

\begin{lemma}[Reflection for an arbitrary character]\label{B1.07}
Assume $H$. For every character $\chi$ modulo $N\ge1$, a zero
$s$ of $L(\cdot,\chi)$ with $0<\Re s<1$ satisfies $\Re s=1/2$.
\end{lemma}
\leanref{Claim.\allowbreak{}strip}{Development.\allowbreak{}AutomorphicArithmetic.\allowbreak{}CharacterBoxes.\allowbreak{}ClaimStrip}{59--73}
\begin{proof}
By Lemma~\ref{B1.12}, $L(s,\chi^\ast)=0$. Apply Lemma~\ref{B1.08}
to $\chi^\ast$.
\end{proof}

\begin{lemma}[Reflection for a primitive character]\label{B1.08}
Assume $H$. For every primitive character $\chi$ modulo $N\ge1$,
a zero $s$ of $L(\cdot,\chi)$ with $0<\Re s<1$ satisfies $\Re s=1/2$.
\end{lemma}
\leanref{Formalism.\allowbreak{}VW.\allowbreak{}critLine}{Formalism.\allowbreak{}GammaFactorNonvanishing.\allowbreak{}CritLine}{11--27}
\begin{proof}
In the strip, $H$ excludes zeros with $\Re s>1/2$. If $\Re s<1/2$,
Lemma~\ref{B1.13} gives a zero at $1-s$ for $\chi^{-1}$, which $H$
also excludes.
\end{proof}

\begin{theorem}[Common input to the five statements]\label{B1.09}
The assertion $H$ holds.
\end{theorem}
\leanref{allChi}{Supplements.\allowbreak{}Kernels.\allowbreak{}MainTheorems}{7--7}
\begin{proof}
Apply Theorem~\ref{B1.15}.
\end{proof}

\begin{definition}[Two nonvanishing assertions]\label{B1.10}
Let $\mathcal P$ consist of the pairs $(N,\psi)$ with $N\ge1$ and
$\psi$ a primitive complex Dirichlet character modulo $N$.
For $(N,\psi)\in\mathcal P$, say that $\psi$ is zero-free above
$\theta\in\R$ if $L(s,\psi)\ne0$ whenever $\Re s>\theta$. Define
\[
\begin{split}
 H:\;&\ \forall N\ge1\ \forall\chi\bmod N\ \forall s\in\C,\quad
 \Re s>\tfrac12\ \Longrightarrow\\[-2pt]
 &\hspace{23mm}
 (\chi\ne\mathbf1_N\ \text{or}\ s\ne1)
 \ \Longrightarrow\ L(s,\chi)\ne0;\\
 H_{\rm prim}:\;&\ \forall(N,\psi)\in\mathcal P\ \forall s\in\C,\quad
 \Re s>\tfrac12\ \Longrightarrow\ L(s,\psi)\ne0 .
\end{split}
\]
\end{definition}
\leanref{UpperTheory.\allowbreak{}EveryCharacter}{Development.\allowbreak{}AutomorphicArithmetic.\allowbreak{}CharacterBoxes.\allowbreak{}EveryCharacter}{13--14}
\leanref{Formalism.HC.PChar}{Formalism.\allowbreak{}SharpRestriction.\allowbreak{}RamificationFrobenius.\allowbreak{}TotientGaloisClasses.\allowbreak{}PrimePairSizeCondition}{7--11}
\leanref{Formalism.HC.PChar.ZeroFreeAbove}{Formalism.\allowbreak{}SharpRestriction.\allowbreak{}RamificationFrobenius.\allowbreak{}TotientGaloisClasses.\allowbreak{}ExplicitFormulaErrorOfHostConditions}{11--12}
\leanref{Formalism.HC.ZeroFreeAll}{Formalism.\allowbreak{}SharpRestriction.\allowbreak{}RamificationFrobenius.\allowbreak{}TotientGaloisClasses.\allowbreak{}ExplicitFormulaErrorOfHostConditions}{18--18}

At $s=1$, the formalized value is
$L(1,\mathbf1_1)=\zeta(1)=(\gamma-\log(4\pi))/2\ne0$,
where $\gamma$ is Euler's constant.
\leanref{riemannZeta\_\allowbreak{}one}{Mathlib.\allowbreak{}NumberTheory.\allowbreak{}Harmonic.\allowbreak{}ZetaAsymp}{418--422}
\leanref{riemannZeta\_\allowbreak{}one\_\allowbreak{}ne\_\allowbreak{}zero}{Mathlib.\allowbreak{}NumberTheory.\allowbreak{}Harmonic.\allowbreak{}ZetaAsymp}{441--450}

\begin{lemma}[Zeta reflection from the right]\label{B1.11}
Suppose that $\zeta(s)\ne0$ for every $s$ satisfying
$\Re s>1/2$ and $s\ne1$. Then the Riemann hypothesis holds.
\end{lemma}
\leanref{Formalism.\allowbreak{}VW.\allowbreak{}rh\_\allowbreak{}of\_\allowbreak{}right}{Formalism.\allowbreak{}GammaFactorNonvanishing.\allowbreak{}RhOfRight}{5--62}
The classical reflection formula
\cite[\S12.8, p.~259, (12)]{apostol-1976}\cite[(25.4.1)]{olver-2010}
gives
\[
 \zeta(1-t)=2(2\pi)^{-t}\Gamma(t)
                 \cos(\pi t/2)\zeta(t),
 \qquad t\notin\{0,-1,-2,\ldots\},\quad t\ne1 .
\]
If $\zeta(s)=0$ and $\Re s<1/2$, put $t=1-s$. Then $\Re t>1/2$
and $t\ne1$ since $\zeta(0)=-1/2$. Both $\zeta(t)$ and $\Gamma(t)$
are nonzero, so the formula gives $\cos(\pi t/2)=0$. Thus $t$ is an
odd integer greater than $1$, and $s=1-t$ is negative and even.

\begin{lemma}[Euler factors and the conductor]\label{B1.12}
Let $\chi$ be a character modulo $N\ge1$, and let
$f$ be its conductor. If $\chi^\ast\ne\mathbf1_f$ or $s\ne1$, then
\[
 L(s,\chi)=L(s,\chi^\ast)
              \prod_{p\mid N}\bigl(1-\chi^\ast(p)p^{-s}\bigr).
\]
Let $\varphi$ be a character of positive modulus.
For a prime $p$ and $\Re s>0$, one has
\[
 |\varphi(p)p^{-s}|\le p^{-\Re s}<1,\qquad
 1-\varphi(p)p^{-s}\ne0 .
\]
Consequently, if $0<\Re s<1$ and $L(s,\chi)=0$, then
$L(s,\chi^\ast)=0$.
\end{lemma}
\leanref{DirichletCharacter.\allowbreak{}LFunction\_\allowbreak{}changeLevel}{Mathlib.\allowbreak{}NumberTheory.\allowbreak{}LSeries.\allowbreak{}DirichletContinuation}{150--162}
\leanref{Claim.\allowbreak{}factor\_\allowbreak{}ne\_\allowbreak{}zero}{Development.\allowbreak{}AutomorphicArithmetic.\allowbreak{}CharacterBoxes.\allowbreak{}ClaimStrip}{7--21}
\leanref{Claim.\allowbreak{}prim\_\allowbreak{}zero}{Development.\allowbreak{}AutomorphicArithmetic.\allowbreak{}CharacterBoxes.\allowbreak{}ClaimStrip}{23--33}
If $M\mid N$ and $\varphi$ modulo $M$ induces $\chi$, the identity
holds with $\varphi$ in place of $\chi^\ast$ provided
$\varphi\ne\mathbf1_M$ or $s\ne1$. Taking $M=f$ and $\varphi=\chi^\ast$
gives the conductor formula
\cite[p.~262]{apostol-1976}\cite[(25.15.4)]{olver-2010}.
For $p\mid f$ the factor is $1$ since $\chi^\ast(p)=0$.

\begin{lemma}[Reflection of primitive zeros]\label{B1.13}
Let $\psi$ be primitive modulo $N\ge1$. Put $a_\psi=0$ if
$\psi(-1)=1$, and $a_\psi=1$ if $\psi(-1)=-1$. Write
$\Gamma_{\R}(z)=\pi^{-z/2}\Gamma(z/2)$ for $\Re z>0$.
Let $\mathcal L_\psi$ denote the archive's completed $L$-function.
For $0<\Re s<1$, it satisfies
\[
 \mathcal L_\psi(s)=\Gamma_{\R}(s+a_\psi)L(s,\psi).
\]
The gamma factor is finite and nonzero for $\Re s>0$
by \cite[\S5.2(i)]{olver-2010}.
If $0<\Re s<1$ and $L(s,\psi)=0$, then
\[
 L(1-s,\psi^{-1})=0 .
\]
\end{lemma}
\leanref{Claim.\allowbreak{}prim\_\allowbreak{}reflect}{Development.\allowbreak{}AutomorphicArithmetic.\allowbreak{}CharacterBoxes.\allowbreak{}ClaimStrip}{42--57}
\leanref{Claim.\allowbreak{}gamma\_\allowbreak{}ne\_\allowbreak{}zero}{Development.\allowbreak{}AutomorphicArithmetic.\allowbreak{}CharacterBoxes.\allowbreak{}ClaimStrip}{35--40}
\leanref{Formalism.\allowbreak{}VW.\allowbreak{}gammaFactor\_\allowbreak{}ne\_\allowbreak{}zero}{Formalism.\allowbreak{}GammaFactorNonvanishing.\allowbreak{}CritLine}{6--9}
\leanref{DirichletCharacter.\allowbreak{}LFunction\_\allowbreak{}eq\_\allowbreak{}completed\_\allowbreak{}div\_\allowbreak{}gammaFactor}{Mathlib.\allowbreak{}NumberTheory.\allowbreak{}LSeries.\allowbreak{}DirichletContinuation}{259--264}
The primitive functional equation and its root number are
\[
 \begin{gathered}
 \mathcal L_{\psi^{-1}}(1-s)
   =N^{s-1/2}W(\psi^{-1})\,\mathcal L_\psi(s),\\
 W(\chi)=i^{-a_\chi}N^{-1/2}\sum_{u\bmod N}\chi(u)e^{2\pi iu/N}.
 \end{gathered}
\]
For $N>1$ this is Apostol's equation with the parity gamma factors
\cite[\S12.10, p.~263, (16)]{apostol-1976}. For $N=1$,
$W(\mathbf1_1)=1$ and this is the completed zeta equation.
A zero in the strip gives $\mathcal L_\psi(s)=0$. Apply the functional
equation to the primitive inverse character to get
$\mathcal L_{\psi^{-1}}(1-s)=0$. Since $\Re(1-s)>0$, the gamma factor
there is nonzero, so $L(1-s,\psi^{-1})=0$.
\leanref{DirichletCharacter.\allowbreak{}IsPrimitive.\allowbreak{}completedLFunction\_\allowbreak{}one\_\allowbreak{}sub}{Mathlib.\allowbreak{}NumberTheory.\allowbreak{}LSeries.\allowbreak{}DirichletContinuation}{284--313}

\begin{lemma}[Passage to all characters]\label{B1.14}
The assertion $H_{\rm prim}$ implies $H$.
\end{lemma}
\leanref{UpperTheory.\allowbreak{}Top.\allowbreak{}imprimitive}{Development.\allowbreak{}AutomorphicArithmetic.\allowbreak{}CharacterBoxes.\allowbreak{}Imprimitive}{5--35}
In the domain of $H$, the pole proviso of Lemma~\ref{B1.12} holds
because a principal inducing character induces the principal character.
Use $H_{\rm prim}$ for the primitive factor and Lemma~\ref{B1.12}
for the Euler factors.

\section{The common arithmetic input}

Let $\mathcal G$ denote Green's polynomial form in
Definition~\ref{B7.06}. For each prime $p$ and compatible label datum, it gives a bijection for each $r\ge1$ from monic irreducible polynomials
$P\in\F_p[X]$ of degree $r$ with $P(0)\ne0$ to the cuspidal irreducible
characters of $\GL_r(\F_p)$, with their degrees and values on elements
with irreducible characteristic polynomial.

Let $\mathcal V$ denote the cohomological window assertion of
Definition~\ref{B6.03}. Fix a primitive $\psi$ and $0<\varepsilon<1/52$.
Let $a$ be a page function as in Definition~\ref{B6.02}. There are
real numbers $J\ge1$ and $X_0\ge1$ such that for every $X\ge X_0$
and every admissible pair $q,d$ at scale $X^{1/J}$, there is a real
$Y$ with $X^{1/J}\le Y$ and $Y^2\le X$. For each $x\in[Y,Y^2]$ put
$t=\lfloor x^{401}\rfloor$. For every field $E$ specified by
$(\psi,a(x),q,d)$ and every algebraic closure $\Omega/E$, choose any
family of local data over $\Omega$ as in Definition~\ref{B6.01}.
Then the lengths $s_{\ell,e}(E)$ are finite for every prime $\ell$
and $1\le e<t$, and every finite set $S$ of primes satisfies
\[
 \sum_{\ell\in S}\log\ell\sum_{1\le e<t}s_{\ell,e}(E)
 \le \exp(x^\varepsilon/2).
\]

Write $\mathcal D=\Lean{GlobalTheory.\allowbreak{}Fc.\allowbreak{}DefiniteSideCondition}$ and set
\[
\begin{aligned}
 \mathcal N(\mathcal C)
   &=\Lean{Closed}\bigl(\mathcal C,
              (F_m,C_H)\mapsto\Lean{LinkN}(F_m)\bigr),\\
 \mathcal H&=\Lean{GlobalTheory.\allowbreak{}Fc.\allowbreak{}HypsAlt}(\mathcal N,\mathcal D).
\end{aligned}
\]
The closure quantifiers and local error are in
Definitions~\ref{B8.39} and \ref{B8.40}. A record in $\mathcal H$
gives the family class $\mathcal A$ and selection $\sigma$ of
Definition~\ref{B8.02}. In the closure range on $D_3(3,\mathcal A)$,
the field \Lean{rRows} supplies a common $C\ge0$ for each qualifying
family $F_m$ and closure parameter $C_H$. At every sufficiently large
selected point, the bundle chosen by $\sigma$ satisfies the indefinite
bounds and the lattice condition $\mathcal D$ of Definition~\ref{B8.35}
with these constants. Each period estimate has its own closure
threshold on $D_3(3,\mathcal A)$. For each qualifying family there is
a threshold $x_0$ such that the estimate holds for every $x\ge x_0$
with $x>F_m.x_1$ and every $i\in F_m.Ix(x)$. The last field is
$\mathcal N(D_3(3,\mathcal A))$.

\begin{theorem}[Construction of the common input]\label{B1.15}
Let $\mathcal R$ be the Galois representation assertion of
Definitions~\ref{B9.04}--\ref{B9.06}, with the local parameter laws of
Definition~\ref{B9.47}. Then $\mathcal R$ implies $H$.
Theorem~\ref{B9.02} supplies $\mathcal R$, so $H$ holds unconditionally.
\end{theorem}
\leanref{UpperTheory.\allowbreak{}FinalTheorems.\allowbreak{}everyCharacter\_\allowbreak{}of\_\allowbreak{}row}{Development.\allowbreak{}AutomorphicArithmetic.\allowbreak{}CharacterBoxes.\allowbreak{}MainConclusion}{7--8}
\leanref{UpperTheory.\allowbreak{}FinalTheorems.\allowbreak{}everyCharacter}{Development.\allowbreak{}AutomorphicArithmetic.\allowbreak{}CharacterBoxes.\allowbreak{}MainConclusion}{10--11}
\begin{proof}
Given $h:\mathcal R$, form $O=(h,\Lean{recTExists},h)$ as in \ref{B9.01},
with local-parameter existence from \ref{B9.46}. Use \ref{B6.04} for
$\mathcal V$.
Let $P,F$ be the unitary-group and local theta inputs of \ref{B8.01}
from $O$. Let $h_L$ assert that for every number field $E$ and every algebraic
closure $\Omega/E$, a correspondence as in \ref{B9.47} exists for every
valuation subring of $\Omega$ over a finite place of $E$.
The regular-character assertion in $P$ gives $\mathcal G$ by \ref{B7.37}.

The doubling continuation and definite Siegel--Weil formula from $O$
(\ref{B8.12}) give local matching and proportionality of
matched theta integrals by \ref{B8.11}.
The Hilbert--Samuel input of \ref{B8.25} gives the lower successive-minima
bound. The fixed-variety inequality of \ref{B8.27} gives the upper
successive-minima and lower arithmetic-volume bounds. Apply thirteen
lemmas to $O$ for the remaining hypotheses. Proposition~\ref{B8.26c}
bounds the mean component slope from below with $C_T$ uniform in the
field and tower. Proposition~\ref{B8.26} uses $C_{\rm deep}=C_T+120$.
Write $h_X$ for the normalized log covolume upper bound of \ref{B8.24}.

For $c\in\R$, apply \ref{B8.41} to $P.\Lean{nLink}$ at the parameters
set by $h_L,c$. The bound applies at twisted local indices above $2$
and above $3$ if $3\notin D.\mathrm{host.small}$ for the family datum $D$.
These estimates, $P,F$, the theta conclusions and $h_X,h_L$ give a
witness in $\mathcal H$ by \ref{B8.01}. Apply \ref{B1.16}.
\end{proof}
\leanref{UpperTheory.LiteratureTheorems.published\_results\_eigen}{Development.AutomorphicArithmetic.ThetaCorrespondence.\allowbreak{}ThetaPairLifts.ThirdPlaceSeries.EveryCharacterOfLiteratureOfVWin}{666--726}
\leanref{UpperTheory.LiteratureTheorems.published\_results\_first}{Development.AutomorphicArithmetic.ThetaCorrespondence.\allowbreak{}ThetaPairLifts.ThirdPlaceSeries.EveryCharacterOfLiteratureOfVWin}{727--780}
\leanref{UpperTheory.LiteratureTheorems.local\_langlands\_gl\_recT\_exists\_of\_literature}{Development.AutomorphicArithmetic.ThetaCorrespondence.\allowbreak{}ThetaPairLifts.ThirdPlaceSeries.EveryCharacterOfLiteratureOfVWin}{313--314}
\leanref{UpperTheory.LiteratureTheorems.every\_character\_of\_literature\_of\_vWin}{Development.\allowbreak{}AutomorphicArithmetic.\allowbreak{}ThetaCorrespondence.\allowbreak{}ThetaPairLifts.\allowbreak{}ThirdPlaceSeries.\allowbreak{}EveryCharacterOfLiteratureOfVWin}{810--813}
\leanref{UpperTheory.Lit.every\_character\_of\_theta\_integrals\_local\_matching}{Development.\allowbreak{}AutomorphicArithmetic.\allowbreak{}ThetaCorrespondence.\allowbreak{}ThetaPairLifts.\allowbreak{}ThirdPlaceSeries.\allowbreak{}SiegelEisensteinSecondTermHypotheses}{69--73}

\begin{theorem}[Three inputs give all characters]\label{B1.16}
Assume $\mathcal G$ and $\mathcal V$, and suppose that
$\mathcal H$ has an element. Then $H$ holds.
\end{theorem}
\leanref{UpperTheory.Top.every\_character\_of\_green\_v\_win}{Development.\allowbreak{}AutomorphicArithmetic.\allowbreak{}ThetaCorrespondence.\allowbreak{}ThetaPairLifts.\allowbreak{}ThirdPlaceSeries.\allowbreak{}TopPlug}{16--18}
\begin{proof}
Apply Theorem~\ref{B1.17}, then Lemma~\ref{B1.14}.
\end{proof}

\begin{theorem}[Three inputs give primitive nonvanishing]\label{B1.17}
Assume $\mathcal G$ and $\mathcal V$, and suppose that
$\mathcal H$ has an element. Then $H_{\rm prim}$ holds.
\end{theorem}
\leanref{UpperTheory.Top.zero\_free\_all\_of\_green\_v\_win}{Development.\allowbreak{}AutomorphicArithmetic.\allowbreak{}ThetaCorrespondence.\allowbreak{}ThetaPairLifts.\allowbreak{}ThirdPlaceSeries.\allowbreak{}TopPlug}{12--15}
\begin{proof}
Choose $h\in\mathcal H$. By Theorem~\ref{B7.54}, $\mathcal G$
gives the family implication $\Lean{TopShapeAlt}$. Apply it to the family class and selection in $h$ with the estimates from the same record. Its two remaining
propositional hypotheses are $\Lean{True}$. The conclusion is
$\Lean{TopG}$, the connectedness assertion used in
Theorem~\ref{B1.18} for requests satisfying host existence and
the familywise cohomological hypotheses.
\leanref{GlobalTheory.Fc.forall\_goal\_of\_shape\_of\_hyps}{Development.\allowbreak{}AutomorphicArithmetic.\allowbreak{}ThetaCorrespondence.\allowbreak{}ThetaPairLifts.\allowbreak{}ThirdPlaceSeries.\allowbreak{}ShapeOf}{31--32}

For each primitive $\psi$, $\Lean{supplyChi}$ gives the host condition
at $\lambda=3/2$ from the window assertion. For $0<\varepsilon<1/52$,
choose $J\ge1$ and real constants $C_{T2},x_0,C_8,C_5,c_6,C_b$ with
$C_8\ge1$ and $C_5,c_6,C_b>0$. For every real $Y_0$, some
$Y\ge\max\{Y_0,x_0,2\}$ gives rank data $R$ and a host $H_R$ with
constant $C_{T2}$ at each scale $s=(x,\varepsilon,J)$ with $x\in[Y,Y^{3/2}]$.
It also gives an auxiliary prime $q_0$ at each such scale. The rank condition and
host bounds hold, and $(C_5,s,R,H_R,q_0)$ satisfies the request for $\psi$.
\leanref{UpperTheory.\allowbreak{}Top.\allowbreak{}supplyChi}{Development.\allowbreak{}AutomorphicArithmetic.\allowbreak{}CharacterBoxes.\allowbreak{}SupplyChi}{8--16}
Apply Theorem~\ref{B1.18} with $\lambda=3/2>1$ to these host conditions
and $\Lean{TopG}$.
\end{proof}


% ---- p1c2 ----
\chapter{How arithmetic windows enter the zero-free criterion}\label{ch:arithmetic-windows}

Use the $L$-value convention of \ref{B1.10}. Write
$\mathcal H_\lambda(S)$ if, for every $Y_0\in\R$, some $Y\geq Y_0$
satisfies $[Y,Y^\lambda]\subset S$. In each window condition we may take
$Y\geq\max(Y_0,1)$, since $Y_0$ is arbitrary.

\section{The arithmetic interface}

A family $\mathcal F$ at $\varepsilon$ has scales $W(\mathcal F)$ and
objects $D_{x,i}$, $i\in I_{\mathcal F}(x)$.
Write $x_1(\mathcal F)>1$ for the threshold in the family data.
Write $E_{\mathcal L}(\mathcal F)$ for
$\exists C,x_0\ \forall x\geq x_0\ \forall i\in I_{\mathcal F}(x),\,
EAt(\mathcal L,D_{x,i},x,\varepsilon,C)$, where $EAt$ evaluates
$EAtH$ of \ref{B2.58} at the raw host and top set of $D_{x,i}$.
For a relation $T$ on families and the real parameters below,
write $\mathcal T_{e_D,e_L}(\mathcal F)$ for
\[
\exists K,c_{HH},e_K,C_L,x_0\in\R\
\forall x\geq x_0,\quad x\in W(\mathcal F)
\Longrightarrow T(\mathcal F;e_D,e_L,K,c_{HH},e_K,C_L,x).
\]
The comparison $\mathcal C_{\mathcal L}$ asserts, for every
$e_D,e_L\in\R$, $\varepsilon>0$, and $\mathcal F$,
\[
\begin{gathered}
E_{\mathcal L}(\mathcal F)\land\mathcal T_{e_D,e_L}(\mathcal F)
\Longrightarrow\exists x_1\ \forall x\geq x_1,\\
x\in W(\mathcal F)\Longrightarrow
G_m^{\mathcal L}(x)\leq x^{1/2+(e_D+e_L+7)\varepsilon},
\end{gathered} \tag{1}
\]
with $G_m^{\mathcal L}$ from \ref{B1.21}.
Write $\mathcal W_{\mathcal L}(\lambda)$ for
\[
\begin{gathered}
\exists e_D,e_L\in\R\ \exists\varepsilon_1>0\
\forall\,0<\varepsilon<\varepsilon_1\ \exists\mathcal F,\\
E_{\mathcal L}(\mathcal F),\qquad
\mathcal H_\lambda(W(\mathcal F)),\qquad
\mathcal T_{e_D,e_L}(\mathcal F).
\end{gathered} \tag{2}
\]

\begin{theorem}[Arithmetic input to the window criterion]\label{B1.18}
Fix $\lambda>1$. Assume $\Lean{TopG}$ and
$\Lean{HostExW}(\lambda,(\Lean{GChi}\,\psi).\Lean{req})$ for every
primitive character $\psi$ of positive modulus. Then
$L(s,\psi)\neq0$ for every such $\psi$ and $\Re s>1/2$.

For each $0<\varepsilon<1/52$, host existence fixes $J\geq1$,
a threshold, common host bounds, and a set of $x$ with $\mathcal H_\lambda$.
At each scale $s=(x,\varepsilon,J)$ in this set above the threshold,
there are rank data $R$ satisfying $\Lean{RankVal}(s,R)$, a host $H$
on $R$, and an auxiliary prime $q_0$. With the common bounds, $H$ and
$q_0$ satisfy $\Lean{HostGood}$, and the request of
\ref{B2.58}--\ref{B2.59} holds for these data.

Use the family type $\mathrm{FamP}(\varepsilon)$ and
$T=\Lean{ChainTied}$, the common chain realisation of \ref{B5.33}.
For each $\lambda'>1$ and request $A$, let
$\mathcal C_{\lambda',A}$ be the families with $\lambda'$-windows
satisfying $A$ for every $x>x_1(\mathcal F)$ and every
$i\in I_{\mathcal F}(x)$. Then $\Lean{TopG}$ says that
host existence for $A$, together with $\Lean{HypV}$ for every
family in this class, yields fixed $e_D,e_L\in\R,\varepsilon_1>0$
such that each $0<\varepsilon<\varepsilon_1$ admits a family in
the class satisfying $\mathcal T_{e_D,e_L}$.
\end{theorem}
\leanref{UpperTheory.\allowbreak{}Top.\allowbreak{}zeroFreeAll\_\allowbreak{}of\_\allowbreak{}topG}{Development.\allowbreak{}AutomorphicArithmetic.\allowbreak{}CharacterBoxes.\allowbreak{}ZeroFreeAllOfTopG}{26--35}
\leanref{UpperTheory.\allowbreak{}Top.\allowbreak{}TopG}{Development.\allowbreak{}AutomorphicArithmetic.\allowbreak{}CharacterBoxes.\allowbreak{}ZeroFreeAllOfTopG}{12--15}

\begin{proof}
At scale $(x,\varepsilon,J)$, let $R$ be the rank data and $H$ the host.
Put $t=R.t$ and $TT=H.TT$. Write $H_{\rm raw}$ for the raw host
underlying $H$.
The request $A=\Lean{GChi}\,\psi$ requires
$EAtH(\mathcal L_\psi,H_{\rm raw},TT,x,\varepsilon,4)$ of \ref{B2.58}.
Put $E=H.E$ and $b=e^{x^\varepsilon}$. For every algebraic closure
of $E$ and every compatible family of local data from \ref{B6.01},
the request also requires finite
$\Z_\ell$-lengths of $\Sha_{\ell,e}(E)$ for every prime $\ell$
and $1\leq e<t$, with
$\sum_{\ell\in S}\log\ell\sum_{1\leq e<t}
\operatorname{length}_{\Z_\ell}\Sha_{\ell,e}(E)\leq\sqrt b$
for every finite set $S$ of rational primes.

For a family in $\mathcal C_{\lambda,A}$, each indexed scale has
coordinates $x,\varepsilon$ when $x\geq x_1(\mathcal F)+1$.
The two clauses give $E_{\mathcal L_\psi}(\mathcal F)$ with $C=4$
and $\Lean{HypV}(\mathcal F)$, the eventual length bound at every index.
Hence $\Lean{TopG}$ gives (2).

To prove (1), fix $e_D,e_L\in\R$, $\varepsilon>0$, and any
$\mathcal F\in\mathrm{FamP}(\varepsilon)$ satisfying its two hypotheses.
For all sufficiently large $x\in W(\mathcal F)$, the tied chain gives
an index $i$ with the chain hypothesis of \ref{B2.20} at $D_{x,i}$.
The condition $E_{\mathcal L_\psi}(\mathcal F)$ gives the arithmetic
hypothesis at the same index. Thus \ref{B2.20} gives (1).
Apply \ref{B1.19} and \ref{B1.22}.
\end{proof}
\leanref{UpperTheory.\allowbreak{}Comparison₂.\allowbreak{}level\_\allowbreak{}gm\_\allowbreak{}le\_\allowbreak{}on\_\allowbreak{}family\_\allowbreak{}scales}{Development.\allowbreak{}AutomorphicArithmetic.\allowbreak{}PrimitiveCosets.\allowbreak{}LevelGmLeOnFamilyScales}{28--31}

\begin{lemma}[Covering primitive characters by cyclic levels]\label{B1.19}
Fix $\lambda>1$ and a property $Q$ of primitive characters.
For every $\psi$ satisfying $Q$, assume (1), the supply (2) at
$\lambda$, and \ref{B1.22} for every window exponent $>1$ at
$\mathcal L_\psi$. Then $L(s,\psi)\neq0$ whenever
$Q(\psi)$ and $\Re s>1/2$.
Each $\psi$ may use a containing level with these properties and its own
$\lambda>1$.
\end{lemma}
\leanref{Formalism.\allowbreak{}HC.\allowbreak{}zeroFreeOn\_\allowbreak{}of\_\allowbreak{}tiedWin\_\allowbreak{}cyc}{Formalism.\allowbreak{}SharpRestriction.\allowbreak{}RamificationFrobenius.\allowbreak{}TotientGaloisClasses.\allowbreak{}TopWin}{90--94}
\leanref{Formalism.\allowbreak{}HC.\allowbreak{}zeroFreeOn\_\allowbreak{}of\_\allowbreak{}tiedWin}{Formalism.\allowbreak{}SharpRestriction.\allowbreak{}RamificationFrobenius.\allowbreak{}TotientGaloisClasses.\allowbreak{}TopWin}{83--89}

\begin{proof}
Apply \ref{B1.20} to the assigned level.
The cyclic level contains $\psi$, its own primitive reduction.
\end{proof}
\leanref{Formalism.\allowbreak{}HC.\allowbreak{}Level.\allowbreak{}mem\_\allowbreak{}cyc}{Formalism.\allowbreak{}SharpRestriction.\allowbreak{}RamificationFrobenius.\allowbreak{}TotientGaloisClasses.\allowbreak{}\allowbreak ExplicitFormulaErrorOfHostConditions}{123--125}

\begin{lemma}[Removing the exponent loss]\label{B1.20}
Fix $\lambda\in\R$, a family type with its map to the underlying
family data, and a relation $T$.
Assume (1), (2), and $\mathrm{FormBAt}(\mathcal L,\lambda)$ as defined
in \ref{B1.22}. Then every kept character has nonvanishing $L$-values
throughout $\Re s>1/2$.
\end{lemma}
\leanref{Formalism.\allowbreak{}HC.\allowbreak{}level\_\allowbreak{}of\_\allowbreak{}tiedWin}{Formalism.\allowbreak{}SharpRestriction.\allowbreak{}RamificationFrobenius.\allowbreak{}TotientGaloisClasses.\allowbreak{}TopWin}{68--82}

\begin{proof}
Choose $e_D,e_L,\varepsilon_1$ from (2).
Fix a kept character and $s$ with $\sigma=\Re s>1/2$. Set
\[
\begin{gathered}
b_0=\tfrac12\min(\sigma-\tfrac12,\tfrac12),\qquad
\theta=\tfrac12+b_0,\qquad D=e_D+e_L+7,\\
\varepsilon=\min\left(\tfrac{\varepsilon_1}{2},
                      \tfrac{b_0}{|D|+1}\right).
\end{gathered}
\]
Then $0<b_0\leq1/4$, $1/2<\theta<\min(\sigma,1)$,
$0<\varepsilon<\varepsilon_1$, and $D\varepsilon\leq b_0$.
Take the family from (2) and the threshold $x_1$ in (1).
On its allowed scales above $\max(x_1,1)$,
\[
G_m^{\mathcal L}(y)\leq y^{1/2+D\varepsilon}\leq y^\theta.
\]
Choosing each family window beyond $\max(Y_0,x_1,1)$ gives
$\mathcal H_\lambda(\{y\in[1,\infty):G_m^{\mathcal L}(y)\leq y^\theta\})$.
Apply $\mathrm{FormBAt}(\mathcal L,\lambda)$, since $\theta<\sigma$.
\end{proof}

\section{Levels and positive weighted families}

\begin{definition}[Levels and cyclic levels]\label{B1.21}
A level has a finite kept set $\mathcal X_{\mathcal L}$ of primitive
characters with positive moduli. The set contains the modulus-one
character and is closed under inversion. The level also specifies
$o_{\mathcal L}\in\N$ and may designate a real odd character
$\xi\in\mathcal X_{\mathcal L}$. Put
\[
\begin{gathered}
w_{\mathcal L}(\chi)=
\begin{cases}
1,&\chi=\xi\text{ when }\xi\text{ is specified},\\
(3o_{\mathcal L}+1)^{-1},&\text{otherwise},
\end{cases}\\
G_m^{\mathcal L}(y)=G_z(y)+
\sum_{\chi\in\mathcal X_{\mathcal L}\setminus\{1\}}
w_{\mathcal L}(\chi)\Re g_\chi(y).
\end{gathered}
\]
Use $G_z,g_\chi$ from \ref{B1.26}. Write
$\mathrm{ZF}_{\mathcal L}(\theta)$ if $L(s,\chi)\neq0$ for every
$\chi\in\mathcal X_{\mathcal L}$ and every $s\in\C$ with $\Re s>\theta$.

A subgroup $\Xi$ of characters modulo $M\geq1$ gives a level
by primitive reduction, with $o_{\mathcal L}=|\Xi|$ and no distinguished
character. Reduction preserves inversion and sends the principal
character to modulus one, so the axioms hold.
For $\Xi=\langle\psi\rangle$, write $\mathcal L_\psi$. If $\psi$
has order $o$, each nonprincipal weight is $1/(3o+1)$.
\end{definition}
\leanref{Formalism.\allowbreak{}HC.\allowbreak{}KSet}{Formalism.\allowbreak{}SharpRestriction.\allowbreak{}RamificationFrobenius.\allowbreak{}TotientGaloisClasses.\allowbreak{}PrimePairSizeCondition}{39--46}
\leanref{Formalism.\allowbreak{}HC.\allowbreak{}Level}{Formalism.\allowbreak{}SharpRestriction.\allowbreak{}RamificationFrobenius.\allowbreak{}TotientGaloisClasses.\allowbreak{}\allowbreak ExplicitFormulaErrorOfHostConditions}{8--10}
\leanref{Formalism.\allowbreak{}HC.\allowbreak{}Level.\allowbreak{}wt}{Formalism.\allowbreak{}SharpRestriction.\allowbreak{}RamificationFrobenius.\allowbreak{}TotientGaloisClasses.\allowbreak{}\allowbreak ExplicitFormulaErrorOfHostConditions}{36--37}
\leanref{Formalism.\allowbreak{}HC.\allowbreak{}Level.\allowbreak{}ofSub}{Formalism.\allowbreak{}SharpRestriction.\allowbreak{}RamificationFrobenius.\allowbreak{}TotientGaloisClasses.\allowbreak{}\allowbreak ExplicitFormulaErrorOfHostConditions}{105--115}
\leanref{Formalism.\allowbreak{}HC.\allowbreak{}Level.\allowbreak{}Gm}{Formalism.\allowbreak{}SharpRestriction.\allowbreak{}RamificationFrobenius.\allowbreak{}TotientGaloisClasses.\allowbreak{}\allowbreak ExplicitFormulaErrorOfHostConditions}{38--39}
\leanref{Formalism.\allowbreak{}HC.\allowbreak{}Level.\allowbreak{}cyc}{Formalism.\allowbreak{}SharpRestriction.\allowbreak{}RamificationFrobenius.\allowbreak{}TotientGaloisClasses.\allowbreak{}\allowbreak ExplicitFormulaErrorOfHostConditions}{121--122}

\begin{theorem}[The level window criterion]\label{B1.22}
For a level $\mathcal L$, $\lambda>1$ and $0<\theta<1$,
\[
\mathcal H_\lambda(\{y\in[1,\infty):G_m^{\mathcal L}(y)\leq y^\theta\})
\Longrightarrow\mathrm{ZF}_{\mathcal L}(\theta).
\]
For $\lambda\in\R$, write $\mathrm{FormBAt}(\mathcal L,\lambda)$ for
this implication for all $0<\theta<1$.
Then $\mathrm{FormB}(\mathcal L)$ requires it for every $\lambda>1$.
\end{theorem}
\leanref{UpperTheory.\allowbreak{}Comparison₂.\allowbreak{}formB}{Development.\allowbreak{}AutomorphicArithmetic.\allowbreak{}PrimitiveCosets.\allowbreak{}Imports}{8--20}
\leanref{Formalism.\allowbreak{}HC.\allowbreak{}Level.\allowbreak{}FormBAt}{Formalism.\allowbreak{}SharpRestriction.\allowbreak{}RamificationFrobenius.\allowbreak{}TotientGaloisClasses.\allowbreak{}Window}{30--32}
\leanref{Formalism.\allowbreak{}HC.\allowbreak{}Level.\allowbreak{}FormB}{Formalism.\allowbreak{}SharpRestriction.\allowbreak{}RamificationFrobenius.\allowbreak{}TotientGaloisClasses.\allowbreak{}Window}{33--34}

\begin{proof}
Apply \ref{B1.24} with $\delta=\lambda-1$ and $C=1$.
The identities in \ref{B1.23} transfer the window bound and
the nonvanishing conclusion.
\end{proof}

\begin{lemma}[Realising a level as a positive weighted family]\label{B1.23}
Every level $\mathcal L$ admits a family $F$ of \ref{B1.25} with
$G_m(F,y)=G_m^{\mathcal L}(y)$ for all $y\in\R$ and
$\mathrm{ZF}(F;\theta)\Longleftrightarrow
\mathrm{ZF}_{\mathcal L}(\theta)$ for all $\theta\in\R$.
\end{lemma}
\leanref{UpperTheory.\allowbreak{}Comparison₂.\allowbreak{}level\_\allowbreak{}family}{Development.\allowbreak{}AutomorphicArithmetic.\allowbreak{}PrimitiveCosets.\allowbreak{}LevelFamily}{41--53}

\begin{proof}
Index $F$ by $\mathcal X_{\mathcal L}$, with its characters
and moduli, weight $1$ at modulus one and $w_{\mathcal L}(\chi)>0$
elsewhere. A primitive principal character has modulus $1$, so
admissibility holds.

The $n=0$ term vanishes since $\Lambda(0)=0$.
On the remaining support, $1\leq n\leq\lfloor y\rfloor_+$ gives
$n\leq y$, hence $\max(y/n-1,0)=y/n-1$. For $y<1$ the support
is empty. At modulus one, $\chi(n)=1$ for all $n$, and the centred
term is $G_z(y)$. For each nonprincipal $\chi$, both definitions give
$S_\chi(y)+yL'(1,\chi)/L(1,\chi)$, with $S_\chi$ as in \ref{B1.26}.
The nonvanishing assertions have identical indices.
\end{proof}

\begin{theorem}[Pointwise power-window criterion]\label{B1.24}
Let $F$ be a positive weighted character family,
$0<\theta<1$, $\delta>0$, and $C\in\R$. If
\[
\forall Y_0\in\R\ \exists Y\geq Y_0\
\forall y\in[Y,Y^{1+\delta}],\quad G_m(F,y)\leq Cy^\theta,
\tag{3}
\]
then $\mathrm{ZF}(F;\theta)$.
\end{theorem}
\leanref{Deduction.\allowbreak{}zeroFree\_\allowbreak{}of\_\allowbreak{}rpow\_\allowbreak{}window}{Development.\allowbreak{}AutomorphicArithmetic.\allowbreak{}DedekindLFunctions.\allowbreak{}LandauTheorems.\allowbreak{}ZeroFreeOfRpowWindow}{11--15}

\begin{proof}
By \ref{B1.27}, (3) gives the integral window condition.
Apply \ref{B1.28}.
\end{proof}

\begin{definition}[Positive weighted character families]\label{B1.25}
A family $F$ has a finite index set $I$, positive integer moduli $N_i$,
complex Dirichlet characters $\chi_i$ modulo $N_i$, and weights $a_i>0$,
with $N_i=1$ or $\chi_i\neq1$ for each $i$.
Characters may be imprimitive.
For $\theta\in\R$, write $\mathrm{ZF}(F;\theta)$ if
$L(s,\chi_i)\neq0$ for every $i\in I$ and every
$s\in\C$ with $\Re s>\theta$.
\end{definition}
\leanref{Deduction.\allowbreak{}CharFamily}{Development.\allowbreak{}AutomorphicArithmetic.\allowbreak{}DedekindLFunctions.\allowbreak{}CharFamily}{24--33}
\leanref{Deduction.\allowbreak{}ZeroFree}{Development.\allowbreak{}AutomorphicArithmetic.\allowbreak{}DedekindLFunctions.\allowbreak{}CharFamily}{38--40}

\begin{definition}[Centred smoothed prime sums]\label{B1.26}
Let $\Lambda$ be the von Mangoldt function and $\gamma$ Euler's
constant. For $y\in\R$, set
\[
\begin{gathered}
S_\chi(y)=\sum_{1\leq n\leq\lfloor y\rfloor_+}
\Lambda(n)\chi(n)\max(y/n-1,0),\\
G_\chi(y)=
\begin{cases}
S_\chi(y)-y\log y+(1+\gamma)y,&\chi=1,\\
S_\chi(y)+yL'(1,\chi)/L(1,\chi),&\chi\neq1.
\end{cases}
\end{gathered}
\]
Put $G_m(F,y)=\sum_{i\in I}a_i\Re G_{\chi_i}(y)$.
Here $\lfloor y\rfloor_+=\max(\lfloor y\rfloor,0)$.
Use $\log y=\log|y|$ for $y<0$ and $\log0=0$.
Write $G_z=G_1$ at modulus one and $g_\chi=G_\chi$ for nontrivial
primitive $\chi$. Put $g_F(u)=G_m(F,e^u)$, $u\in\R$.
\end{definition}
\leanref{Deduction.\allowbreak{}ramp}{Development.\allowbreak{}AutomorphicArithmetic.\allowbreak{}DedekindLFunctions.\allowbreak{}CharFamily}{10--12}
\leanref{Deduction.\allowbreak{}Gchi}{Development.\allowbreak{}AutomorphicArithmetic.\allowbreak{}DedekindLFunctions.\allowbreak{}CharFamily}{15--20}
\leanref{Deduction.\allowbreak{}Gm}{Development.\allowbreak{}AutomorphicArithmetic.\allowbreak{}DedekindLFunctions.\allowbreak{}CharFamily}{36--37}
\leanref{Deduction.\allowbreak{}gm}{Development.\allowbreak{}AutomorphicArithmetic.\allowbreak{}DedekindLFunctions.\allowbreak{}LandauZeroFreeCriterion}{9--10}

\section{From pointwise bounds to integral windows}

\begin{lemma}[An integrable upper envelope]\label{B1.27}
For $\theta,\delta\in\R$, let $\mathrm{WH}(F;\theta,\delta)$
mean that every $\eta>0$ and $Y_0\in\R$ admit $Y\geq Y_0$ and
$R:\R\to\R$ with $y\mapsto R(y)y^{-\theta-1}\in L^1(J_Y)$,
where $J_Y=[Y,Y^{1+\delta}]$, and
\[
\forall y\in J_Y,\quad 0\leq R(y),\quad G_m(F,y)\leq R(y),
\qquad
\int_{J_Y}R(y)y^{-\theta-1}\,dy\leq Y^\eta. \tag{4}
\]
For every $\theta,C\in\R$ and $\delta>0$, (3) implies
$\mathrm{WH}(F;\theta,\delta)$.
\end{lemma}
\leanref{Deduction.\allowbreak{}WindowHyp}{Development.\allowbreak{}AutomorphicArithmetic.\allowbreak{}DedekindLFunctions.\allowbreak{}CharFamily}{41--46}
\leanref{Deduction.\allowbreak{}windowHyp\_\allowbreak{}of\_\allowbreak{}rpow}{Development.\allowbreak{}AutomorphicArithmetic.\allowbreak{}DedekindLFunctions.\allowbreak{}LandauTheorems.\allowbreak{}WindowHypOfRpow}{27--58}

\begin{proof}
Fix $\eta>0,Y_0$. Choose a window in (3) with $Y\geq\max(Y_0,1)$
and $|C|\delta\log Y\leq Y^\eta$, possible since
$(\log Y)/Y^\eta\to0$.
Set $R(y)=|C|y^\theta$ for $y>0$ and $R(y)=0$ for $y\leq0$.
On $J_Y$, $R$ is nonnegative and dominates $G_m(F,y)$.
The integrand $R(y)y^{-\theta-1}=|C|/y$ is continuous there, and
\[
\int_{J_Y}R(y)y^{-\theta-1}\,dy
=|C|\log\frac{Y^{1+\delta}}Y
=|C|\delta\log Y\leq Y^\eta.
\]
\end{proof}

\begin{theorem}[Integral window criterion for a character family]\label{B1.28}
For a positive weighted character family $F$, $0<\theta<1$, and
$\delta>0$, one has
$\mathrm{WH}(F;\theta,\delta)\Longrightarrow\mathrm{ZF}(F;\theta)$.
\end{theorem}
\leanref{Deduction.\allowbreak{}lemmaL}{Development.\allowbreak{}AutomorphicArithmetic.\allowbreak{}DedekindLFunctions.\allowbreak{}LandauTheorems.\allowbreak{}lemmaL}{13--23}

\begin{proof}
By \ref{B1.30}, the Laplace transform of $g_F$ on $u>0$ continues
meromorphically to $\Re s>0$ as $\Phi$, with the pole data of \ref{B1.33}.
Its pole set $Z_F$ contains all zeros of the family $L$-functions
with positive real part.
Apply \ref{B1.29} and then \ref{B1.31} to obtain
$\Re\rho\leq\theta$ for $\rho\in Z_F$.
A zero $s$ with $\Re s>\theta$ contradicts \ref{B1.32} if
$\Re s\geq1$, and the bound on $Z_F$ if $\Re s<1$.
\end{proof}

\begin{lemma}[Logarithmic windows]\label{B1.29}
For $g:\R\to\R$ and $\theta,\delta\in\R$, let
$\mathrm{WH}_U(g;\theta,\delta)$ mean that every $\eta>0$ and $U_0\in\R$
admit $U\geq U_0$ and $R_U:\R\to\R$ with
$u\mapsto R_U(u)e^{-\theta u}\in L^1(I_U)$,
where $I_U=[U,(1+\delta)U]$, and
\[
\forall u\in I_U,\quad 0\leq R_U(u),\quad g(u)\leq R_U(u),
\qquad
\int_{I_U}R_U(u)e^{-\theta u}\,du\leq e^{\eta U}. \tag{5}
\]
For every character family $F$, $\theta\in\R$, and $\delta>0$,
$\mathrm{WH}(F;\theta,\delta)\Longrightarrow
\mathrm{WH}_U(g_F;\theta,\delta)$.
\end{lemma}
\leanref{Deduction.\allowbreak{}WindowHypU}{Development.\allowbreak{}AutomorphicArithmetic.\allowbreak{}DedekindLFunctions.\allowbreak{}LandauZeroFreeCriterion}{25--30}
\leanref{Deduction.\allowbreak{}windowHypU\_\allowbreak{}of\_\allowbreak{}windowHyp}{Development.\allowbreak{}AutomorphicArithmetic.\allowbreak{}DedekindLFunctions.\allowbreak{}LandauTheorems.\allowbreak{}WindowHypUOfWindowHyp}{29--57}

\begin{proof}
Given $\eta>0$ and $U_0\in\R$, use $\mathrm{WH}(F;\theta,\delta)$
to choose $Y\geq\max(e^{U_0},1)$ and its envelope $R$.
Put $U=\log Y$ and $R_U(u)=R(e^u)$.
Then $U\geq U_0$, and $\exp$ maps $I_U$ bijectively onto $J_Y$,
since $e^{(1+\delta)U}=Y^{1+\delta}$.
The envelope inequalities follow.
Apply Folland's change-of-variables theorem
\cite[Theorem~2.47(a)]{folland-1999} to $u\mapsto e^u$ and
$y\mapsto R(y)y^{-\theta-1}$ on $J_Y$, extended by zero outside
$J_Y$. This specialization gives integrability and
\[
\begin{aligned}
\int_{I_U}R(e^u)e^{-\theta u}\,du
&=\int_{I_U}R(e^u)(e^u)^{-\theta-1}e^u\,du\\
&=\int_{J_Y}R(y)y^{-\theta-1}\,dy
\leq Y^\eta=e^{\eta U}.
\end{aligned}
\]
\end{proof}


% ---- p1c3 ----
\chapter{Zeros become poles}\label{ch:zeros-become-poles}

Throughout, characters have positive integer modulus. Fix a family
\(F=(\chi_i,a_i)_{i\in I}\) as in Definition~\ref{B1.25} and use the sums of
Definition~\ref{B1.26}. Put
\[
 g(u)=G_m(F,e^u)=\sum_{i\in I}a_i\Re G_{\chi_i}(e^u),
 \qquad
 \mathcal M H(s)=\int_1^\infty H(y)y^{-s-1}\,dy.
\]
Put \(D_\chi(s)=L'(s,\chi)/L(s,\chi)\) where defined.
Use real logarithms in powers.

\section{The analytic interface}

\begin{theorem}[Pole data for a character family]\label{B1.30}
Choose the functions of \ref{B1.34} for \(\chi_i,\chi_i^{-1}\).
The sum
\[
 \Phi(s)=\sum_{i\in I}\frac{a_i}{2}
 \bigl(\Phi_{\chi_i}(s)+\Phi_{\chi_i^{-1}}(s)\bigr)
\]
gives pole data for \(g,Z_F\) as in \ref{B1.33}.
\end{theorem}
\leanref{Deduction.\allowbreak{}poleData}{Development.\allowbreak{}AutomorphicArithmetic.\allowbreak{}DedekindLFunctions.\allowbreak{}LandauTheorems.\allowbreak{}poleData}{248--401}
\begin{proof}
By \ref{B1.38}, \(g\) is Borel measurable and
\(|g(u)|\le(\sum_i a_iC_{\chi_i})e^{2u}\) for \(u\ge0\).
By \ref{B1.45}, \(2\Re G_{\chi_i}=G_{\chi_i}+G_{\chi_i^{-1}}\).
Multiply by \(e^{-su}\) and integrate termwise for \(\Re s>2\)
by \ref{B1.37}.
Both characters satisfy \ref{B1.34}, since
\(\chi_i^{-1}=1\iff\chi_i=1\).
Fix \(z\) with \(\Re z>0\). By \ref{B1.34} and finiteness of \(I\), there
is a punctured neighbourhood of \(z\) on which
\(L(w,\chi_i)L(w,\chi_i^{-1})\ne0\) for every \(i\in I\).
By \ref{B1.45}, this neighbourhood misses \(Z_F\).
If \(\Re s>0\) and \(s\notin Z_F\), each summand is analytic at \(s\)
by \ref{B1.34} and \ref{B1.45}. By definition \(Z_F\) lies in
\(\Re s>0\), and \ref{B1.46} gives the residues.
\end{proof}

\begin{theorem}[Integral windows exclude poles]\label{B1.31}
Let \(g:\R\to\R\), \(\Phi:\C\to\C\), and \(Z\subset\C\)
satisfy \ref{B1.33}. Fix \(0<\theta<1\) and \(\delta>0\).
Assume the logarithmic window condition
\(\mathrm{WH}_U(g;\theta,\delta)\) of Lemma~\ref{B1.29}.
Then \(\Re\rho\le\theta\) for every \(\rho\in Z\).
\end{theorem}
\leanref{Deduction.\allowbreak{}landau\_\allowbreak{}window}{Development.\allowbreak{}AutomorphicArithmetic.\allowbreak{}DedekindLFunctions.\allowbreak{}LandauTheorems.\allowbreak{}landauWindow}{83--167}
\begin{proof}
Suppose \(\rho_0\in Z\), \(\Re\rho_0>\theta\).
Take \(\sigma>2\) from \ref{B1.47}. All poles lie to its left
by \ref{B1.54}. In \ref{B1.56}, the first alternative contradicts
\ref{B1.61}, and the second contradicts \ref{B1.62} by positivity
of the residue at its real pole.
\end{proof}

\begin{lemma}[Nonvanishing on the closed half-plane]\label{B1.32}
For \(i\in I\) and \(\Re s\ge1\), \(L(s,\chi_i)\ne0\),
including the assigned modulus-one value at \(1\).
\end{lemma}
\leanref{Deduction.\allowbreak{}charFamily\_\allowbreak{}LFunction\_\allowbreak{}ne\_\allowbreak{}zero\_\allowbreak{}of\_\allowbreak{}one\_\allowbreak{}le\_\allowbreak{}re}{Development.\allowbreak{}AutomorphicArithmetic.\allowbreak{}DedekindLFunctions.\allowbreak{}LandauTheorems.\allowbreak{}\allowbreak CharFamilyLFunctionNeZeroOfOneLeRe}{9--21}
Davenport gives nonvanishing at regular points
\cite[Sections~4, 13--14]{davenport-2000}.
At modulus one and \(s=1\), \ref{B1.10} gives
\(\zeta(1)=(\gamma-\log(4\pi))/2\ne0\).
% [Springer edition record](https://link.springer.com/book/9780387950976)

\begin{definition}[Zero set and pole data]\label{B1.33}
Define
\[
 Z_F=\{\rho\in\C:\Re\rho>0,\ 
 \exists i\in I,\ L(\rho,\chi_i)=0
 \ \hbox{or}\ L(\bar\rho,\chi_i)=0\}.
\]
For \(g:\R\to\R\), \(\Phi:\C\to\C\), and \(Z\subset\C\),
\emph{pole data} require Borel measurability of \(g\) and some
\(C_0\in\R\) with
\[
 |g(u)|\le C_0e^{2u}\quad(u\ge0),\qquad
 \Phi(s)=\int_0^\infty g(u)e^{-su}\,du\quad(\Re s>2).
\]
Every point of \(Z\) has positive real part.
\(\Phi\) is holomorphic on \(\{s:\Re s>0\}\setminus Z\).
The set \(Z\) has no accumulation in \(\Re s>0\).
Each \(\rho\in Z\) admits \(c_\rho\ne0\) and \(h_\rho\)
analytic near \(\rho\) with
\[
 \Phi(s)=\frac{c_\rho}{s-\rho}+h_\rho(s)
 \quad(0<|s-\rho|\ \hbox{sufficiently small}).
\]
If \(\rho\) is real, \(c_\rho\) is a positive real number.
\end{definition}
\leanref{Deduction.\allowbreak{}zeroSet}{Development.\allowbreak{}AutomorphicArithmetic.\allowbreak{}DedekindLFunctions.\allowbreak{}LandauZeroFreeCriterion}{11--14}
\leanref{Deduction.\allowbreak{}PoleData}{Development.\allowbreak{}AutomorphicArithmetic.\allowbreak{}DedekindLFunctions.\allowbreak{}LandauZeroFreeCriterion}{15--24}

\section{One character and the point \(1\)}

\begin{lemma}[Continuation for one character]\label{B1.34}
Let \(\chi\) have modulus \(N\ge1\), with \(N=1\) or
\(\chi\ne1\). There exists \(\Phi_\chi:\C\to\C\) such that
\[
 \Phi_\chi(s)=\int_0^\infty G_\chi(e^u)e^{-su}\,du
 \quad(\Re s>2).
\]
For each \(\rho\ne0\), \(\Phi_\chi\) is analytic at \(\rho\) if
\(L(\rho,\chi)\ne0\), and there are \(m\in\N\) and \(h\) analytic near
\(\rho\) with
\[
 m=0\ \Longleftrightarrow\ L(\rho,\chi)\ne0,\qquad
 \Phi_\chi(s)=\frac{c_\chi(\rho)}{s-\rho}+h(s)
\]
on a punctured neighbourhood of \(\rho\).
Here \(c_\chi(\rho)=m/[\rho(1-\rho)]\) if \(\rho\ne1\),
and \(c_\chi(1)=0\). At a zero, \(m\) is its multiplicity.
The zeros of \(L(\,\cdot\,,\chi)\) have no accumulation point in \(\C\).
\end{lemma}
\leanref{Deduction.\allowbreak{}poleData\_\allowbreak{}char}{Development.\allowbreak{}AutomorphicArithmetic.\allowbreak{}DedekindLFunctions.\allowbreak{}LandauTheorems.\allowbreak{}poleData}{137--240}
\begin{proof}
For \(\chi\ne1\), apply \ref{B1.36} with \(b=0\) to the entire
function \(f=L(\,\cdot\,,\chi)\), which has \(f(1)\ne0\). Use
\ref{B1.41} and \ref{B1.37} for the integral.

If \(\chi=1\), then \(N=1\). The entire function \(f(s)=(s-1)\zeta(s)\) for
\(s\ne1\) has \(f(1)=1\) and \(f'(1)=\gamma\) by the Laurent expansion
\cite[Section~25.2]{olver-2010}. On \(\Re s>2\), the Euler product gives
\(\zeta(s)\ne0\), so \(f'/f=\zeta'/\zeta+1/(s-1)\).
With \(b=1\), \ref{B1.36} gives the numerator
\(\zeta'/\zeta+1/(s-1)-s\gamma+1-s\) of \ref{B1.43}.
Apply \ref{B1.37} for the integral. The factor \(s-1\) preserves
zero orders away from \(1\). At \(1\), \(f(1)\) and the assigned \(\zeta(1)\) are
both nonzero.
\end{proof}

\begin{lemma}[Local logarithmic derivative]\label{B1.35}
Let \(f\) be analytic near \(\rho\) with a nonzero germ.
Then \(m=\operatorname{ord}_\rho f\in\N\), with
\(m=0\iff f(\rho)\ne0\), and some \(h\) analytic at \(\rho\) satisfies
\(f'(s)/f(s)=m/(s-\rho)+h(s)\) on a punctured neighbourhood of \(\rho\).
\end{lemma}
\leanref{Deduction.\allowbreak{}logDeriv\_\allowbreak{}local}{Development.\allowbreak{}AutomorphicArithmetic.\allowbreak{}DedekindLFunctions.\allowbreak{}LandauTheorems.\allowbreak{}LogDerivLocal}{9--40}
This is the logarithmic derivative form of the standard local factorization
theorem \cite[Chapter~4]{ahlfors-1979}.

\begin{lemma}[Removing the apparent singularity]\label{B1.36}
Let \(f:\C\to\C\) be entire, \(f(1)\ne0\), and \(b\in\C\).
Write \(D=f'/f\), with division by zero assigned the value zero, and put
\[
 N(s)=D(s)-sD(1)+b(1-s),\qquad
 \Phi_f(s)=-\frac{N(s)}{s(s-1)}\quad(s\ne1).
\]
It extends analytically at \(1\), with
\(\Phi_f(1)=D(1)+b-D'(1)\).
It is analytic where \(s\ne0\) and \(f(s)\ne0\).
At a nonzero zero \(\rho\) of order \(m\),
\[
 \Phi_f(s)=\frac{m}{\rho(1-\rho)(s-\rho)}+h_\rho(s)
\]
on a punctured neighbourhood of \(\rho\), with \(h_\rho\) analytic at
\(\rho\).
The zeros of \(f\) are isolated.
\end{lemma}
\leanref{Deduction.\allowbreak{}poleData\_\allowbreak{}abstract}{Development.\allowbreak{}AutomorphicArithmetic.\allowbreak{}DedekindLFunctions.\allowbreak{}LandauTheorems.\allowbreak{}poleData}{22--90}
\leanref{Deduction.\allowbreak{}poleData\_\allowbreak{}dslope}{Development.\allowbreak{}AutomorphicArithmetic.\allowbreak{}DedekindLFunctions.\allowbreak{}LandauTheorems.\allowbreak{}poleData}{17--21}
\begin{proof}
Since \(N(1)=0\), \(T(s)=N(s)/(s-1)\) extends analytically with
\(T(1)=N'(1)\). Near \(1\) use \(\Phi_f(s)=-T(s)/s\).
Away from \(0,1\) use the stated quotient.
The identity theorem and \(f(1)\ne0\) give isolated zeros.
At a zero \(\rho\ne0,1\), \ref{B1.35} gives
\(D(s)=m/(s-\rho)+h(s)\), with \(m=\operatorname{ord}_\rho f\).
For \(Q(s)=-1/[s(s-1)]\),
\[
 \begin{split}
 \Phi_f(s)={}&\frac{mQ(\rho)}{s-\rho}
 +m\frac{Q(s)-Q(\rho)}{s-\rho}\\
 &+Q(s)\bigl(h(s)-sD(1)+b(1-s)\bigr).
 \end{split}
\]
The second term extends with value \(mQ'(\rho)\), and the third is
analytic. The residue is \(mQ(\rho)=m/[\rho(1-\rho)]\).
If \(\rho\ne0\) and \(f(\rho)\ne0\), take \(m=0\) and \(h_\rho=\Phi_f\).
\end{proof}

\section{The Mellin calculations}

\begin{lemma}[Exponential substitution]\label{B1.37}
For \(G:\R\to\C\) and \(s\in\C\),
\[
 \int_0^\infty G(e^u)e^{-su}\,du
 =\int_1^\infty G(y)y^{-s-1}\,dy,
\]
where the Bochner integral has value zero for a nonintegrable integrand.
Both integrands are absolutely integrable for \(G=G_\chi\) and \(\Re s>2\).
\end{lemma}
\leanref{Deduction.\allowbreak{}poleData\_\allowbreak{}change}{Development.\allowbreak{}AutomorphicArithmetic.\allowbreak{}DedekindLFunctions.\allowbreak{}LandauTheorems.\allowbreak{}poleData}{114--136}
\leanref{Deduction.\allowbreak{}poleData\_\allowbreak{}integrable}{Development.\allowbreak{}AutomorphicArithmetic.\allowbreak{}DedekindLFunctions.\allowbreak{}LandauTheorems.\allowbreak{}poleData}{91--113}
The smooth inverse maps \(y=e^u\) and \(u=\log y\) preserve Lebesgue
measurability. Change of variables for the absolute values
\cite[Theorem~2.47(a)]{folland-1999} gives equivalence of integrability.
In the integrable case the same substitution gives the identity.
For \(G_\chi\), \ref{B1.38} gives the majorant
\(C_\chi e^{-(\Re s-2)u}\) when \(\Re s>2\).

\begin{lemma}[Measurability and quadratic growth]\label{B1.38}
For every \(\chi\), \(G_\chi:\R\to\C\) is Borel measurable,
and some \(C_\chi\in\R\) satisfies
\[
 |G_\chi(y)|\le C_\chi y^2\qquad(y\ge1).
\]
For \(f:\N\to\C\) with \(|f(n)|\le1\),
\[
 \left|\sum_{1\le n\le y}\Lambda(n)f(n)(y/n-1)\right|
 \le2y^2\qquad(y\ge1).
\]
\end{lemma}
\leanref{Deduction.\allowbreak{}measurable\_\allowbreak{}Gchi}{Development.\allowbreak{}AutomorphicArithmetic.\allowbreak{}DedekindLFunctions.\allowbreak{}LandauTheorems.\allowbreak{}measurableGchi}{24--29}
\leanref{Deduction.\allowbreak{}Gchi\_\allowbreak{}growth\_\allowbreak{}ramp}{Development.\allowbreak{}AutomorphicArithmetic.\allowbreak{}DedekindLFunctions.\allowbreak{}LandauTheorems.\allowbreak{}GchiGrowth}{9--36}
\leanref{Deduction.\allowbreak{}Gchi\_\allowbreak{}growth}{Development.\allowbreak{}AutomorphicArithmetic.\allowbreak{}DedekindLFunctions.\allowbreak{}LandauTheorems.\allowbreak{}GchiGrowth}{37--67}
\begin{proof}
For \(t\in\R\), write \(t_+=\max(t,0)\).
Since \(0\le\Lambda(n)\le n\) and \((y/n-1)_+\le y/n\), each
summand has norm at most \(y\), and there are at most \(2y\) terms.
Use \(0\le\log y\le y\) to take \(C_\chi=3+|1+\gamma|\)
for principal \(\chi\) and \(C_\chi=2+|D_\chi(1)|\) otherwise.
On each Borel level set of \(\lfloor y\rfloor_+\), the ramp is a
finite sum of Borel functions, and the centring terms are Borel.
\end{proof}

\begin{lemma}[Transform of a single ramp]\label{B1.39}
For \(n\in\N\), \(n\ge1\), and \(s\in\C\) with
\(\sigma=\Re s>1\),
\[
 \int_1^\infty (y/n-1)_+y^{-s-1}\,dy
 =\frac{n^{-s}}{s(s-1)},\qquad
 \int_1^\infty |(y/n-1)_+y^{-s-1}|\,dy
 =\frac{n^{-\sigma}}{\sigma(\sigma-1)}.
\]
\end{lemma}
\leanref{Deduction.\allowbreak{}mellin\_\allowbreak{}ramp}{Development.\allowbreak{}AutomorphicArithmetic.\allowbreak{}DedekindLFunctions.\allowbreak{}LandauTheorems.\allowbreak{}MellinRamp}{10--52}
\leanref{Deduction.\allowbreak{}mellin\_\allowbreak{}Gchi\_\allowbreak{}of\_\allowbreak{}ne\_\allowbreak{}one\_\allowbreak{}norm}{Development.\allowbreak{}AutomorphicArithmetic.\allowbreak{}DedekindLFunctions.\allowbreak{}LandauTheorems.\allowbreak{}mellinGchiOfNeOne}{27--42}
By \cite[Equation~(5.9.1), \(\mu=\nu=1\)]{olver-2010},
the substitution \(y=ne^u\) gives
\(n^{-s}((s-1)^{-1}-s^{-1})\) for the first integral.
Taking the modulus replaces \(s\) by \(\sigma\) in the integrand.
% [DLMF (5.9.1)](https://dlmf.nist.gov/5.9.E1)

\begin{lemma}[Termwise transformation]\label{B1.40}
Let \(f:\N\to\C\), \(\Re s=\sigma>1\), and suppose
\(\sum_{n\ge1}|f(n)|\Lambda(n)n^{-\sigma}<\infty\).
If \(E:(1,\infty)\to\C\) is integrable, then
\[
 \begin{split}
 &\int_1^\infty
 \left(\sum_{1\le n\le y}\Lambda(n)f(n)(y/n-1)y^{-s-1}
 +E(y)\right)\,dy\\
 &\hspace{1cm}=
 \frac{\sum_{n\ge1}f(n)\Lambda(n)n^{-s}}{s(s-1)}
 +\int_1^\infty E(y)\,dy .
 \end{split}
\]
\end{lemma}
\leanref{Deduction.\allowbreak{}mellin\_\allowbreak{}Gchi\_\allowbreak{}of\_\allowbreak{}ne\_\allowbreak{}one\_\allowbreak{}ramp}{Development.\allowbreak{}AutomorphicArithmetic.\allowbreak{}DedekindLFunctions.\allowbreak{}LandauTheorems.\allowbreak{}mellinGchiOfNeOne}{65--115}
\leanref{Deduction.\allowbreak{}mellin\_\allowbreak{}Gchi\_\allowbreak{}modOne\_\allowbreak{}sum}{Development.\allowbreak{}AutomorphicArithmetic.\allowbreak{}DedekindLFunctions.\allowbreak{}LandauTheorems.\allowbreak{}mellinGchiModOne}{10--26}
By \ref{B1.39}, the sum of the \(L^1\)-norms of the ramp terms is
\(\sum_{n\ge1}|f(n)|\Lambda(n)n^{-\sigma}/[\sigma(\sigma-1)]\).
Folland's Theorem~2.25 \cite{folland-1999} gives the stated interchange.

\begin{lemma}[Transform for a nonprincipal character]\label{B1.41}
For nonprincipal \(\chi\) and \(\Re s>1\),
\[
 \mathcal M G_\chi(s)
 =-\frac{D_\chi(s)-sD_\chi(1)}{s(s-1)}.
\]
\end{lemma}
\leanref{Deduction.\allowbreak{}mellin\_\allowbreak{}Gchi\_\allowbreak{}of\_\allowbreak{}ne\_\allowbreak{}one}{Development.\allowbreak{}AutomorphicArithmetic.\allowbreak{}DedekindLFunctions.\allowbreak{}LandauTheorems.\allowbreak{}mellinGchiOfNeOne}{116--143}
\begin{proof}
Logarithmic differentiation of the Euler product
\cite[Section~11.5, p.~231]{apostol-1976} gives
\[
 \sum_{n\ge1}\chi(n)\Lambda(n)n^{-s}=-D_\chi(s).
\]
In \ref{B1.40} take \(E(y)=D_\chi(1)y^{-s}\) and use
\(\int_1^\infty y^{-s}\,dy=(s-1)^{-1}\).
\end{proof}
% [Euler product](https://dlmf.nist.gov/25.15.E2)

\begin{lemma}[The logarithmic centring term]\label{B1.42}
For \(\Re s>1\), the function \((\log y)y^{-s}\) is
absolutely integrable on \((1,\infty)\), and
\[
 \lim_{y\to\infty}(\log y)y^{1-s}=0,\qquad
 \int_1^\infty(\log y)y^{-s}\,dy=\frac1{(s-1)^2}.
\]
\end{lemma}
\leanref{Deduction.\allowbreak{}mellin\_\allowbreak{}Gchi\_\allowbreak{}modOne\_\allowbreak{}log\_\allowbreak{}int}{Development.\allowbreak{}AutomorphicArithmetic.\allowbreak{}DedekindLFunctions.\allowbreak{}LandauTheorems.\allowbreak{}mellinGchiModOne}{57--84}
\leanref{Deduction.\allowbreak{}mellin\_\allowbreak{}Gchi\_\allowbreak{}modOne\_\allowbreak{}log}{Development.\allowbreak{}AutomorphicArithmetic.\allowbreak{}DedekindLFunctions.\allowbreak{}LandauTheorems.\allowbreak{}mellinGchiModOne}{85--114}
With \(y=e^u\), equation~(5.9.1) of \cite{olver-2010} with
\(\mu=1,\nu=2\) gives \(\int_0^\infty ue^{-(s-1)u}\,du=(s-1)^{-2}\).
Since \(\Re(s-1)>0\), the integral converges absolutely and
\(ue^{-(s-1)u}\to0\).

\begin{lemma}[Transform for the modulus-one character]\label{B1.43}
For \(1\) modulo \(1\) and \(\Re s>1\),
\[
 \mathcal M G_1(s)=
 -\frac{\zeta'(s)/\zeta(s)+1/(s-1)+1-(1+\gamma)s}{s(s-1)}.
\]
\end{lemma}
\leanref{Deduction.\allowbreak{}mellin\_\allowbreak{}Gchi\_\allowbreak{}modOne}{Development.\allowbreak{}AutomorphicArithmetic.\allowbreak{}DedekindLFunctions.\allowbreak{}LandauTheorems.\allowbreak{}mellinGchiModOne}{115--200}
\begin{proof}
Here
\(G_1(y)=\sum_{1\le n\le y}\Lambda(n)(y/n-1)
-y\log y+(1+\gamma)y\).
Use \(\sum_{n\ge1}\Lambda(n)n^{-s}=-\zeta'(s)/\zeta(s)\).
Applying \ref{B1.40} to the ramp and \ref{B1.42} to the
logarithmic term gives
\[
 \mathcal M G_1(s)=
 -\frac{\zeta'(s)/\zeta(s)}{s(s-1)}
 -\frac1{(s-1)^2}+\frac{1+\gamma}{s-1}.
\]
The identity \(s/(s-1)=1+1/(s-1)\) gives the stated numerator.
\end{proof}

\section{Conjugation and residues}

\begin{lemma}[Conjugation of the Hurwitz function]\label{B1.44}
Let \(a\in\R/\Z\). Write \(H_a\) for the archive's Hurwitz
zeta function with parameter modulo \(1\).
For every \(s\in\C\),
\[
 H_a(\bar s)=\overline{H_a(s)},\qquad
 \Gamma_{\R}(\bar s)=\overline{\Gamma_{\R}(s)},\qquad
 \Gamma_{\R}(s)=\pi^{-s/2}\Gamma(s/2).
\]
The difference \(H_a(s)-\Gamma_{\R}(s)^{-1}/(s-1)\) is entire
with value \(H_a(1)\) at \(1\).
\end{lemma}
\leanref{Deduction.\allowbreak{}LFunctionConjugation.\allowbreak{}hurwitz\_\allowbreak{}conj}{Development.\allowbreak{}AutomorphicArithmetic.\allowbreak{}DedekindLFunctions.\allowbreak{}LandauTheorems.\allowbreak{}LFunctionConj}{76--85}
\begingroup
\fontspec[Scale=0.85]{DejaVu Sans Mono}
\makeatletter\let\ttdefault\f@family\makeatother
\leanref{Deduction.\allowbreak{}LFunctionConjugation.\allowbreak{}Gammaℝ\_\allowbreak{}conj}{Development.\allowbreak{}AutomorphicArithmetic.\allowbreak{}DedekindLFunctions.\allowbreak{}LandauTheorems.\allowbreak{}LFunctionConj}{19--25}
\endgroup
\leanref{Deduction.\allowbreak{}LFunctionConjugation.\allowbreak{}differentiable\_\allowbreak{}reg}{Development.\allowbreak{}AutomorphicArithmetic.\allowbreak{}DedekindLFunctions.\allowbreak{}LandauTheorems.\allowbreak{}LFunctionConj}{38--49}
\leanref{HurwitzZeta.\allowbreak{}tendsto\_\allowbreak{}hurwitzZeta\_\allowbreak{}sub\_\allowbreak{}one\_\allowbreak{}div\_\allowbreak{}nhds\_\allowbreak{}one}{Mathlib.\allowbreak{}NumberTheory.\allowbreak{}LSeries.\allowbreak{}HurwitzZeta}{96--100}
\begin{proof}
Choose \(a_*\in(0,1]\) representing \(a\), so \(H_a(s)=\zeta(s,a_*)\)
for \(s\ne1\) \cite[Section~25.11]{olver-2010}.
The identity \(\Gamma(\bar z)=\overline{\Gamma(z)}\) and \(\pi>0\)
give conjugation for \(\Gamma_{\R}\) and its entire reciprocal
\(J=1/\Gamma_{\R}\) \cite[Section~5.2]{olver-2010}.
Since \(J(1)=1\) and \(H_a\) has residue \(1\),
\(R_a(s)=H_a(s)-J(s)/(s-1)\) extends across \(1\).
In Mathlib's normalization,
\(\lim_{s\to1,\,s\ne1}(H_a(s)-J(s)/(s-1))=H_a(1)\).
Hence the extension has \(R_a(1)=H_a(1)\).
The Hurwitz series gives \(R_a(\bar s)=\overline{R_a(s)}\) on
\(\Re s>1\). The identity theorem extends this equality to \(\C\).
Add the polar term away from \(1\) and use the extension value at \(1\).
At \(s=-2n\) with \(n\in\Z_{\ge0}\), both sides of the gamma
conjugation identity have the assigned value zero.
\end{proof}
% [DLMF 25.11](https://dlmf.nist.gov/25.11)

\begin{lemma}[Inverse characters and conjugation]\label{B1.45}
For every \(\chi\), \(s\in\C\), and \(y\in\R\),
\[
 L(\bar s,\chi^{-1})=\overline{L(s,\chi)},\qquad
 G_{\chi^{-1}}(y)=\overline{G_\chi(y)}.
\]
Consequently \(L(\rho,\chi^{-1})=0\) if and only if
\(L(\bar\rho,\chi)=0\).
For nonprincipal \(\chi\), also
\(D_{\chi^{-1}}(1)=\overline{D_\chi(1)}\).
\end{lemma}
\leanref{Deduction.\allowbreak{}LFunction\_\allowbreak{}conj}{Development.\allowbreak{}AutomorphicArithmetic.\allowbreak{}DedekindLFunctions.\allowbreak{}LandauTheorems.\allowbreak{}LFunctionConj}{88--97}
\leanref{Deduction.\allowbreak{}Gchi\_\allowbreak{}conj}{Development.\allowbreak{}AutomorphicArithmetic.\allowbreak{}DedekindLFunctions.\allowbreak{}LandauTheorems.\allowbreak{}GchiConj}{38--52}
\leanref{Deduction.\allowbreak{}poleData\_\allowbreak{}inv\_\allowbreak{}zero}{Development.\allowbreak{}AutomorphicArithmetic.\allowbreak{}DedekindLFunctions.\allowbreak{}LandauTheorems.\allowbreak{}poleData}{241--247}
\begin{proof}
For modulus \(N\), the definition gives
\[
 L(s,\chi)=N^{-s}\sum_{j\in\Z/N\Z}\chi(j)H_{j/N}(s).
\]
This is Apostol's finite Hurwitz expansion away from \(1\)
\cite[Chapter~12, (8), p.~255]{apostol-1976},
extended to \(1\) by the values in \ref{B1.44}.
For \(N=1\), the sole class \(j=0\) contributes \(H_0=\zeta\).
Conjugate the sum using \ref{B1.44} and
\(\overline{\chi(j)}=\chi^{-1}(j)\).
For nonprincipal \(\chi\), differentiate
\(L(s,\chi^{-1})=\overline{L(\bar s,\chi)}\) at \(1\)
to obtain the identity for \(D(1)\).

The real weights \(\Lambda(n)(y/n-1)_+\) and the identity
\(\overline{\chi(n)}=\chi^{-1}(n)\) give conjugation for the ramps.
The principal centring term is real. The identity for \(D(1)\) gives
conjugation for the nonprincipal centring terms.
\end{proof}
% [Hurwitz expansion](https://dlmf.nist.gov/25.15.E3)

\begin{lemma}[Positive weighted multiplicities]\label{B1.46}
For \(\rho\in Z_F\), let \(m_i(\rho),m_i^-(\rho)\) be the
multiplicities at \(\rho\) for \(\chi_i,\chi_i^{-1}\), respectively,
with multiplicity zero at nonzeros. The function in \ref{B1.30}
has residue
\[
 c_\rho=\frac{S_\rho}{\rho(1-\rho)},\qquad
 S_\rho=\sum_{i\in I}\frac{a_i}{2}
 \bigl(m_i(\rho)+m_i^-(\rho)\bigr)>0.
\]
Thus \(c_\rho\ne0\). If \(\rho=\beta\in\R\),
then \(0<\beta<1\) and \(c_\beta>0\).
\end{lemma}
\leanref{Deduction.\allowbreak{}poleData}{Development.\allowbreak{}AutomorphicArithmetic.\allowbreak{}DedekindLFunctions.\allowbreak{}LandauTheorems.\allowbreak{}poleData}{336--400}
\begin{proof}
By \ref{B1.32} and \ref{B1.45}, \(0<\Re\rho<1\).
Sum the expansions of \ref{B1.34} on a common punctured
neighbourhood. By \ref{B1.34} and \ref{B1.45}, an index \(i_0\)
witnessing \(\rho\in Z_F\) contributes at least \(a_{i_0}/2>0\)
to \(S_\rho\). The other terms are nonnegative. The denominator is
nonzero, and for real \(\rho=\beta\) it is positive.
\end{proof}


% ---- p1c4 ----
\chapter{Windows control high derivatives}\label{ch:p1c4}

Let \(g:\R\to\R\) be Borel measurable, with
\(|g(u)|\le C_0e^{2u}\) for \(u\ge0\), where \(C_0\ge0\).
For \(k\in\N\), \(a>0\), and \(u\ge0\), put
\[
 K_{k,a}(u)=\frac{u^ke^{-au}}{k!},\qquad
 B_k(z)=\int_0^\infty g(u)\frac{u^ke^{-zu}}{k!}\,du
       \quad(\Re z>2).
\]
When \((g,\Phi,Z)\) satisfies the pole data of \ref{B1.33}, set
\(A_k(z)=(-1)^k\Phi^{(k)}(z)/k!\).
For a character family \(F\), Definition~\ref{B1.26} gives the
arithmetic example \(g(u)=G_m(F,e^u)\). All integrals are Lebesgue
integrals.

\begin{theorem}[Derivative bound from integral windows]\label{B1.47}
Suppose that \((g,\Phi,Z)\) satisfies \ref{B1.33}.
Fix \(0<\theta<1\) and \(\delta>0\). Assume the logarithmic window
condition \(\mathrm{WH}_U(g;\theta,\delta)\) of Lemma~\ref{B1.29}.
Then there exist \(\sigma>2\) and \(C\in\R\) such that, for every
\(\eta>0\) and \(U_0\in\R\), one can find \(U\ge U_0\) and
\(k\in\N\), \(k\ge U\), for which
\[
 |A_k(\sigma+it)|+\Re A_k(\sigma)
       \le \frac{Ce^{\eta U}}{(\sigma-\theta)^k}
       \qquad\text{for every }t\in\R.
\]
\end{theorem}
\leanref{Deduction.\allowbreak{}landau\_\allowbreak{}window\_\allowbreak{}key}{Development.\allowbreak{}AutomorphicArithmetic.\allowbreak{}DedekindLFunctions.\allowbreak{}LandauTheorems.\allowbreak{}landauWindow}{40--74}
\begin{proof}
Take \(\sigma\) and the positive witness \(D\) displayed in
Lemma~\ref{B1.52}. Apply the window condition with lower bound
\(\max(U_0,1)\), and choose \(k\ge U\) from that lemma.
By \ref{B1.49} and \ref{B1.53}, for every real \(t\),
\[
 |A_k(\sigma+it)|+\Re A_k(\sigma)
 \le \frac{2e^{\eta U}}{(\sigma-\theta)^k}
                         +\frac{2C_0D}{\sigma^k}.
\]
Since \(0<\sigma-\theta<\sigma\) and \(e^{\eta U}\ge1\), take
\(C=2+2C_0D\).
\end{proof}

\section{The differentiated Laplace kernel}

\begin{lemma}[Differentiation under the Laplace integral]\label{B1.48}
For \(g\) as above and \(j\in\N\), define
\[
 H_j(z)=\int_0^\infty g(u)(-u)^je^{-zu}\,du
                  \qquad(\Re z>2).
\]
These integrals converge absolutely. For every \(j\in\N\) and
\(s\in\C\) with \(\Re s>2\),
\[
 H_j'(s)=H_{j+1}(s),\qquad H_0^{(j)}(s)=H_j(s).
\]
\end{lemma}
\leanref{Deduction.\allowbreak{}iteratedDeriv\_\allowbreak{}laplace\_\allowbreak{}hasDerivAt}{Development.\allowbreak{}AutomorphicArithmetic.\allowbreak{}DedekindLFunctions.\allowbreak{}LandauTheorems.\allowbreak{}IteratedDerivLaplace}{48--83}
\leanref{Deduction.\allowbreak{}iteratedDeriv\_\allowbreak{}laplace}{Development.\allowbreak{}AutomorphicArithmetic.\allowbreak{}DedekindLFunctions.\allowbreak{}LandauTheorems.\allowbreak{}IteratedDerivLaplace}{84--106}
Apply Widder's Laplace--Stieltjes differentiation formula
\cite[Chapter II, \S5]{widder-1941} to the locally absolutely
continuous function \(\alpha(u)=\int_0^u g(v)\,dv\) on \([0,\infty)\).
% [Widder, title page, contents and introduction](https://api.pageplace.de/preview/DT0400.9781400876457_A26058229/preview-9781400876457_A26058229.pdf).

\begin{lemma}[Normalized coefficients]\label{B1.49}
If \((g,\Phi,Z)\) satisfies \ref{B1.33}, then
\(A_k(z)=B_k(z)\) for every \(k\in\N\) and \(\Re z>2\).
In particular, for \(\sigma>2\) and \(t\in\R\),
\[
 \begin{aligned}
 A_k(\sigma+it)&=\int_0^\infty g(u)K_{k,\sigma}(u)e^{-itu}\,du,\\
 \Re A_k(\sigma)&=\int_0^\infty g(u)K_{k,\sigma}(u)\,du.
 \end{aligned}
\]
\end{lemma}
\leanref{Deduction.\allowbreak{}landau\_\allowbreak{}coeff}{Development.\allowbreak{}AutomorphicArithmetic.\allowbreak{}DedekindLFunctions.\allowbreak{}LandauTheorems.\allowbreak{}LandauCoeff}{10--29}
On \(\Re z>2\), the pole data identify \(\Phi\) with \(H_0\), so
\(A_k=B_k\) is Widder's differentiation formula
\cite[Chapter II, \S5]{widder-1941} in Lemma~\ref{B1.48} with the
normalization \((-1)^k/k!\).

\begin{lemma}[Mass and pointwise bound for the kernel]\label{B1.50}
For \(a>0\) and \(k\in\N\), the kernel \(K_{k,a}\) is integrable
on \((0,\infty)\), and
\[
 \int_0^\infty K_{k,a}(u)\,du=a^{-k-1},\qquad
 0\le K_{k,a}(u)\le a^{-k}\quad(u\ge0).
\]
\end{lemma}
\leanref{Deduction.\allowbreak{}kernel\_\allowbreak{}integral}{Development.\allowbreak{}AutomorphicArithmetic.\allowbreak{}DedekindLFunctions.\allowbreak{}LandauTheorems.\allowbreak{}kernelIntegral}{9--21}
\leanref{Deduction.\allowbreak{}landau\_\allowbreak{}kernel\_\allowbreak{}max}{Development.\allowbreak{}AutomorphicArithmetic.\allowbreak{}DedekindLFunctions.\allowbreak{}LandauTheorems.\allowbreak{}LandauKernelMax}{9--20}
\leanref{Deduction.\allowbreak{}landau\_\allowbreak{}tails\_\allowbreak{}integrable}{Development.\allowbreak{}AutomorphicArithmetic.\allowbreak{}DedekindLFunctions.\allowbreak{}LandauTheorems.\allowbreak{}landauTails}{10--17}
The integral is the specialization \(\mu=1\), \(\nu=k+1\), \(z=a\)
of Euler's gamma integral
\cite[Section 5.9(i), equation (5.9.1)]{olver-2010}, divided by \(k!\).
The pointwise bound follows from \((au)^k/k!\le e^{au}\) after
multiplication by \(a^{-k}e^{-au}\).
% [DLMF, gamma integral (5.9.1)](https://dlmf.nist.gov/5.9.E1).

\section{Both tails at a common derivative order}

\begin{lemma}[Exponential bounds for the two tails]\label{B1.51}
For \(k\in\N\), \(0\le\ell<a\), and \(U,V\in\R\),
\[
 \int_{(0,\infty)\setminus[U,V]}K_{k,a}(u)\,du
 \le \frac{e^{-\ell V}}{(a-\ell)^{k+1}}
                +\frac{e^{\ell U}}{(a+\ell)^{k+1}}.
\]
Furthermore, for \(b,\sigma>0\) and \(x,c\in\R\),
\[
 x+k\frac{\sigma-b}{b}\le c
 \quad\Longrightarrow\quad
 \frac{e^x}{b^{k+1}}\le\frac{e^c}{b\sigma^k}.
\]
\end{lemma}
\leanref{Deduction.\allowbreak{}landau\_\allowbreak{}tails\_\allowbreak{}chernoff}{Development.\allowbreak{}AutomorphicArithmetic.\allowbreak{}DedekindLFunctions.\allowbreak{}LandauTheorems.\allowbreak{}landauTails}{18--61}
\leanref{Deduction.\allowbreak{}landau\_\allowbreak{}tails\_\allowbreak{}pow}{Development.\allowbreak{}AutomorphicArithmetic.\allowbreak{}DedekindLFunctions.\allowbreak{}LandauTheorems.\allowbreak{}landauTails}{62--79}
\begin{proof}
If \(u>0\) lies outside \([U,V]\), then \(u<U\) or \(u>V\). Hence
\[
 K_{k,a}(u)\le
 e^{\ell U}K_{k,a+\ell}(u)+e^{-\ell V}K_{k,a-\ell}(u).
\]
By nonnegativity, the integral of this majorant over the complement
is at most its integral over \((0,\infty)\).
Lemma~\ref{B1.50} evaluates the latter.

For the second assertion, \(1+y\le e^y\), with
\(y=(\sigma-b)/b\), gives
\[
 \frac{e^x}{b^{k+1}}
 =\frac{e^x}{b\sigma^k}\left(\frac{\sigma}{b}\right)^k
 \le\frac{\exp\!\left(x+k(\sigma-b)/b\right)}{b\sigma^k}
 \le\frac{e^c}{b\sigma^k}.
\]
\end{proof}

\begin{lemma}[Choosing the order from the window]\label{B1.52}
For every \(\delta>0\), there exist \(\sigma>2\) and \(D\in\R\)
such that, for every real \(U\ge1\), there is \(k\in\N\) with
\(k\ge U\) and
\[
 \int_{(0,\infty)\setminus[U,(1+\delta)U]}
             K_{k,\sigma-2}(u)\,du\le\frac{D}{\sigma^k}.
\]
One may take \(\ell=2+6/\delta\), \(M=\ell(1+\delta)\), and
\(\sigma=M(\ell+2)\), with
\[
 k=\left\lceil(1+\delta/2)\sigma U\right\rceil,\qquad
 D=\frac{e}{\sigma-2-\ell}+\frac1{\sigma-2+\ell}.
\]
\end{lemma}
\leanref{Deduction.\allowbreak{}landau\_\allowbreak{}tails}{Development.\allowbreak{}AutomorphicArithmetic.\allowbreak{}DedekindLFunctions.\allowbreak{}LandauTheorems.\allowbreak{}landauTails}{80--132}
\begin{proof}
Write \(p=1+\delta/2\) and \(b_\pm=\sigma-2\pm\ell\).
Since \(M>\ell>2\), we have
\(b_-=(M-1)(\ell+2)>0\). Thus \(\sigma>2\) and \(D>0\).
Direct computation gives
\[
 p(\ell+2)=M-1,\qquad p(\ell-2)=\ell+1,\qquad
 p\sigma=M(M-1).
\]
The ceiling satisfies \(p\sigma U\le k<p\sigma U+1\), so \(k\ge U\).
By \ref{B1.51}, with \(a=\sigma-2\), the tail integral is at most
\(T_++T_-\), where
\[
 T_+=\frac{e^{-MU}}{b_-^{k+1}},\qquad
 T_-=\frac{e^{\ell U}}{b_+^{k+1}}.
\]
For the upper tail, \((\sigma-b_-)/b_-=1/(M-1)\). The upper
ceiling bound gives
\[
 \begin{aligned}
 -MU+k\frac{\sigma-b_-}{b_-}
 &=-MU+\frac{k}{M-1}\\
 &<-MU+\frac{p\sigma U+1}{M-1}
   =\frac1{M-1}\le1.
 \end{aligned}
\]
The second bound of \ref{B1.51}, with \(x=-MU\) and \(c=1\),
gives \(T_+\le e/(b_-\sigma^k)\).

For the lower tail, \((\sigma-b_+)/b_+=-(\ell-2)/b_+\).
Since \(\sigma\ge\ell(\ell-2)\), we have
\(\ell b_+\le(\ell+1)\sigma\). The lower ceiling bound gives
\[
 \frac{k(\ell-2)}{b_+}
 \ge\frac{p\sigma U(\ell-2)}{b_+}
 =\frac{(\ell+1)\sigma U}{b_+}\ge\ell U.
\]
Thus \(\ell U+k(\sigma-b_+)/b_+\le0\).
Apply the same bound with \(x=\ell U\) and \(c=0\) to obtain
\(T_-\le1/(b_+\sigma^k)\). Adding proves the claim.
\end{proof}

\section{The one-sided estimate}

\begin{lemma}[Combining the window with its complement]\label{B1.53}
Let \(g\) satisfy the growth and measurability assumptions above.
Let \(\sigma>2\), \(\theta<\sigma\), \(U>0\), and \(V\in\R\),
and put \(I=[U,V]\). Suppose that \(R:\R\to\R\) satisfies
\(R(u)\ge\max(g(u),0)\) on \(I\), and that
\(R(u)e^{-\theta u}\) is integrable on \(I\).
Then, for every \(k\in\N\) and \(t\in\R\),
\[
 \begin{split}
 |B_k(\sigma+it)|+\Re B_k(\sigma)
 \le{}&\frac{2}{(\sigma-\theta)^k}
                 \int_I R(u)e^{-\theta u}\,du\\
 &+2C_0\int_{(0,\infty)\setminus I}K_{k,\sigma-2}(u)\,du.
 \end{split}
\]
\end{lemma}
\leanref{Deduction.\allowbreak{}landau\_\allowbreak{}window\_\allowbreak{}ineq}{Development.\allowbreak{}AutomorphicArithmetic.\allowbreak{}DedekindLFunctions.\allowbreak{}LandauTheorems.\allowbreak{}landauWindowIneq}{27--112}
\leanref{Deduction.\allowbreak{}landau\_\allowbreak{}window\_\allowbreak{}ineq\_\allowbreak{}int}{Development.\allowbreak{}AutomorphicArithmetic.\allowbreak{}DedekindLFunctions.\allowbreak{}LandauTheorems.\allowbreak{}landauWindowIneq}{10--26}
\begin{proof}
Set \(w(u)=K_{k,\sigma}(u)\).
The bound \(|g(u)|w(u)\le C_0K_{k,\sigma-2}(u)\) for \(u>0\),
together with \ref{B1.50}, gives absolute convergence of the
integrals below. Since \(|e^{-itu}|=1\), the triangle inequality gives
\[
 \begin{aligned}
 |B_k(\sigma+it)|+\Re B_k(\sigma)
 &\le\int_0^\infty |g(u)|w(u)\,du+\int_0^\infty g(u)w(u)\,du\\
 &=\int_0^\infty (|g(u)|+g(u))w(u)\,du.
 \end{aligned}
\]
On \(I\), the envelope assumption and \ref{B1.50} give
\[
 \begin{aligned}
 |g(u)|+g(u)&=2\max(g(u),0)\le2R(u),\\
 w(u)&=K_{k,\sigma-\theta}(u)e^{-\theta u}
       \le(\sigma-\theta)^{-k}e^{-\theta u}.
 \end{aligned}
\]
Since \(R\ge0\) and \(w\ge0\) on \(I\), multiplying these
inequalities and integrating gives the first term.
On the complement,
\((|g(u)|+g(u))w(u)\le2|g(u)|w(u)\le2C_0K_{k,\sigma-2}(u)\).
Adding the integral over \((0,\infty)\setminus I\) proves the inequality.
\end{proof}


% ---- p1c5 ----
\chapter{Excluding zeros to the right of theta}\label{ch:p1c5}

We use the pole data $(g,\Phi,Z)$ of Definition~\ref{B1.33}. Write
\[
 H=\{s\in\C:\Re s>0\},\qquad
 B(z,r)=\{w\in\C:|w-z|<r\},\qquad
 \overline B(z,r)=\{w\in\C:|w-z|\le r\}.
\]
At $\rho\in Z$, denote the coefficient of $(s-\rho)^{-1}$ by
$c_\rho$. It is nonzero, and it is positive real when $\rho$ is real.

\begin{lemma}[The absolute-convergence region]\label{B1.54}
For every $\sigma_0>2$ there is $M\in\R$ such that
$|\Phi(s)|\le M$ whenever $\Re s\ge\sigma_0$.
More precisely, a growth bound $|g(u)|\le C_0e^{2u}$ for $u\ge0$
gives $M=C_0/(\sigma_0-2)$. Every $\rho\in Z$ satisfies $\Re\rho\le2$.
\end{lemma}
\leanref{Deduction.\allowbreak{}landau\_\allowbreak{}pole\_\allowbreak{}re\_\allowbreak{}le\_\allowbreak{}bound}{Development.\allowbreak{}AutomorphicArithmetic.\allowbreak{}DedekindLFunctions.\allowbreak{}LandauTheorems.\allowbreak{}LandauPoleReLe}{9--28}
\leanref{Deduction.\allowbreak{}landau\_\allowbreak{}pole\_\allowbreak{}re\_\allowbreak{}le}{Development.\allowbreak{}AutomorphicArithmetic.\allowbreak{}DedekindLFunctions.\allowbreak{}LandauTheorems.\allowbreak{}LandauPoleReLe}{29--54}

The Laplace majorant is $C_0e^{-(\sigma_0-2)u}$
\cite[Chapter~II, \S3]{widder-1941}.
If $\Re\rho>2$, take $\sigma_0=(2+\Re\rho)/2$.
Local boundedness gives $(s-\rho)\Phi(s)\to0$. The pole expansion
gives the limit $c_\rho\ne0$.
This is the bounded case of the removable-singularity criterion
\cite[Chapter~IV, \S3.1, Theorem~7]{ahlfors-1979}.

\begin{lemma}[Local finiteness]\label{B1.55}
Let $Z\subset\C$. Suppose every $z\in H$ has a punctured
neighbourhood disjoint from $Z$. For every $z\in\C$ and $r\in\R$
with $r<\Re z$, the set $Z\cap\overline B(z,r)$ is finite.
\end{lemma}
\leanref{Deduction.\allowbreak{}landau\_\allowbreak{}finite}{Development.\allowbreak{}AutomorphicArithmetic.\allowbreak{}DedekindLFunctions.\allowbreak{}LandauTheorems.\allowbreak{}LandauFinite}{9--27}

This is the closed-disc case of the Bolzano--Weierstrass theorem
\cite[Chapter~III, \S1.4, Theorem~7]{ahlfors-1979}.

\begin{lemma}[Dichotomy at a real centre]\label{B1.56}
Suppose $Z$ satisfies the local condition of Lemma~\ref{B1.55}.
Fix $0<\theta<\sigma$ with $Z\subset\{w:\Re w<\sigma\}$.
If $\Re\rho_0>\theta$ for some $\rho_0\in Z$, at least one of the
following alternatives holds.
\begin{enumerate}
\item There are $D\le\sigma-\theta$ and $\rho_1\in Z$ such that
$Z\cap B(\sigma,D)=\varnothing$ and $\sigma-\Re\rho_1<D$.
\item There are $\beta\in\R$ and $r\in\R$ such that
$\beta\in Z$, $\beta>\theta$, $\sigma-\beta<r\le\sigma-\theta$,
and $Z\cap B(\sigma,r)=\{\beta\}$.
\end{enumerate}
\end{lemma}
\leanref{Deduction.\allowbreak{}landau\_\allowbreak{}dichotomy}{Development.\allowbreak{}AutomorphicArithmetic.\allowbreak{}DedekindLFunctions.\allowbreak{}LandauTheorems.\allowbreak{}LandauDichotomy}{24--88}

\begin{proof}
Set $R=\sigma-\theta$ and $F=Z\cap\overline B(\sigma,R)$.
Lemma~\ref{B1.55} makes $F$ finite. If $F$ is empty, take
$D=R$ and $\rho_1=\rho_0$.

Otherwise choose $\rho_*\in F$ minimising distance to $\sigma$,
and put $D=|\sigma-\rho_*|>0$. Then $D\le R$ and
$Z\cap B(\sigma,D)=\varnothing$. If some $\rho_1\in Z$ has
$\sigma-\Re\rho_1<D$, the first alternative holds.
In the remaining case every $w\in Z$ satisfies
$D\le\sigma-\Re w$. For every minimiser $w\in F$,
\[
 D\le\sigma-\Re w\le|\sigma-w|=D,\qquad
 |\sigma-w|^2=(\sigma-\Re w)^2+(\Im w)^2.
\]
Hence every minimiser equals the real point $\beta=\sigma-D$.
Applying $D\le\sigma-\Re\rho_0$ gives $D<R$, so $\beta>\theta$.
Every point of $F\setminus\{\beta\}$ has distance strictly greater
than $D$. Choose $r$ with $D<r\le R$ no larger than any of these
finitely many distances. Points of $Z$ outside $F$ have distance
greater than $R$. Thus $Z\cap B(\sigma,r)=\{\beta\}$.
\end{proof}
\leanref{Deduction.\allowbreak{}eq\_\allowbreak{}ofReal\_\allowbreak{}sub\_\allowbreak{}of\_\allowbreak{}dist\_\allowbreak{}le\_\allowbreak{}of\_\allowbreak{}le\_\allowbreak{}sub\_\allowbreak{}re}{Development.\allowbreak{}AutomorphicArithmetic.\allowbreak{}DedekindLFunctions.\allowbreak{}LandauTheorems.\allowbreak{}LandauDichotomy}{10--23}

\begin{lemma}[An isolated pole after moving the centre]\label{B1.57}
Suppose $Z$ satisfies the local condition of Lemma~\ref{B1.55},
$\Re w<\sigma$ for all $w\in Z$, and $D<\sigma$.
If $\rho_1\in Z$ satisfies $\sigma-\Re\rho_1<D$, then there exist
$t\in\R$, $\rho\in Z$, and $r\in\R$ such that, for $z=\sigma+it$,
\[
 |z-\rho|<r\le D,\qquad Z\cap B(z,r)=\{\rho\}.
\]
\end{lemma}
\leanref{Deduction.\allowbreak{}landau\_\allowbreak{}nearest}{Development.\allowbreak{}AutomorphicArithmetic.\allowbreak{}DedekindLFunctions.\allowbreak{}LandauTheorems.\allowbreak{}LandauNearest}{36--110}

\begin{proof}
Put $a=\sigma-\Re\rho_1$, so $0<a<D$, and choose
\[
 \xi=\tfrac12\min\{D-a,\sigma-D\}>0,\qquad
 z_0=\sigma+i\Im\rho_1.
\]
Since $D+\xi<\sigma$, the candidate set
$F=Z\cap\overline B(z_0,D+\xi)$ is finite and contains
$\rho_1$. For distinct $v,w\in F$, equality of distances to
$\sigma+it$ is equivalent to
\[
 2t(\Im v-\Im w)
 =(\sigma-\Re v)^2+(\Im v)^2
  -(\sigma-\Re w)^2-(\Im w)^2.
\]
If $\Im v\ne\Im w$, the displayed equation determines a unique
height at which the distances tie.
If $\Im v=\Im w$, equal distances would imply
$(\sigma-\Re v)^2=(\sigma-\Re w)^2$. Both bases are positive,
so $v=w$. Thus each distinct pair excludes at most one height.
Choose $t\in(\Im\rho_1-\xi,\Im\rho_1+\xi)$ outside this finite
set of heights and put $z=\sigma+it$.
All distances from $z$ to points of $F$ are distinct, and
\[
 |z-\rho_1|\le |z-z_0|+|z_0-\rho_1|
 <\xi+a<D.
\]
Let $\rho$ be the unique nearest point in $F$ and set
\[
 r=\min\bigl(\{D\}\cup\{|z-w|:w\in F\setminus\{\rho\}\}\bigr).
\]
Since $|z-\rho|\le|z-\rho_1|<D$, the choice of $\rho$ gives
$|z-\rho|<r\le D$. If $w\in Z\cap B(z,r)$, then
$|w-z_0|<D+\xi$, so $w\in F$ and $w=\rho$.
\end{proof}
\leanref{Deduction.\allowbreak{}landau\_\allowbreak{}nearest\_\allowbreak{}equi}{Development.\allowbreak{}AutomorphicArithmetic.\allowbreak{}DedekindLFunctions.\allowbreak{}LandauTheorems.\allowbreak{}LandauNearest}{10--35}

For any $\Phi:\C\to\C$, write $A_k(z)=(-1)^k\Phi^{(k)}(z)/k!$ for
$k\in\N$. In Lemmas~\ref{B1.58} and~\ref{B1.59}, $\Phi$ is arbitrary.

\begin{lemma}[Cauchy's coefficient bound]\label{B1.58}
If $\Phi$ is holomorphic on $B(z,r)$ and $0<r'<r$, then some
$M\in\R$ satisfies
\[
 \frac{|\Phi^{(k)}(z)|}{k!}\le\frac{M}{(r')^k}
 \qquad(k\in\N).
\]
\end{lemma}
\leanref{Deduction.\allowbreak{}landau\_\allowbreak{}cauchy}{Development.\allowbreak{}AutomorphicArithmetic.\allowbreak{}DedekindLFunctions.\allowbreak{}LandauTheorems.\allowbreak{}LandauCauchy}{9--34}

Ahlfors gives Cauchy's estimate with the circle maximum.
Taking absolute values inside his integral formula gives $M=(2\pi)^{-1}\int_0^{2\pi}|\Phi(z+r'e^{it})|\,dt$
\cite[Chapter~IV, \S2.3, formulas~(24)--(25)]{ahlfors-1979}.

\begin{lemma}[Coefficients of an isolated simple pole]\label{B1.59}
Let $\rho\in B(z,r)$, $\rho\ne z$, and $c\in\C$.
Suppose $\Phi$ is holomorphic on $B(z,r)\setminus\{\rho\}$ and
agrees with $c/(s-\rho)+h(s)$ on a punctured neighbourhood of
$\rho$, where $h$ is holomorphic at $\rho$.
For every $0<r'<r$ there is $M\in\R$ such that
\[
 \left|A_k(z)-\frac{c}{(z-\rho)^{k+1}}\right|
 \le\frac{M}{(r')^k}\qquad(k\in\N).
\]
\end{lemma}
\leanref{Deduction.\allowbreak{}landau\_\allowbreak{}pole\_\allowbreak{}asym}{Development.\allowbreak{}AutomorphicArithmetic.\allowbreak{}DedekindLFunctions.\allowbreak{}LandauTheorems.\allowbreak{}LandauPoleAsym}{19--80}

After subtraction of the simple principal part, the remainder extends
across $\rho$ with value $h(\rho)$. Lemma~\ref{B1.58} applies on
the whole disc \cite[Chapter~IV, \S\S2.3, 3.2]{ahlfors-1979}.
The normalisation follows from
\[
 \frac{d^k}{ds^k}\frac{c}{s-\rho}
   =\frac{(-1)^k k!\,c}{(s-\rho)^{k+1}}.
\]
\leanref{Deduction.\allowbreak{}landau\_\allowbreak{}pole\_\allowbreak{}asym\_\allowbreak{}deriv}{Development.\allowbreak{}AutomorphicArithmetic.\allowbreak{}DedekindLFunctions.\allowbreak{}LandauTheorems.\allowbreak{}LandauPoleAsym}{10--18}

\begin{lemma}[Incompatible exponential rates]\label{B1.60}
Let $c>0$ and $0<d<e$. For every $C\in\R$ there are
$\eta>0$ and $U_*\in\R$ such that, for every $U\ge U_*$
and every $k\in\N$ with $k\ge U$,
\[
 C\,\frac{\exp(\eta U)}{e^k}<\frac{c}{d^k}.
\]
Consequently, for $e\le E$ and fixed $M,C\in\R$, the assertion
\[
 \forall\eta>0\ \forall U_0\in\R\ \exists U\ge U_0\
 \exists k\in\N,\quad k\ge U,\quad
 \frac{c}{d^k}\le\frac{M}{e^k}
                    +C\frac{\exp(\eta U)}{E^k}
 \tag{*}
\]
is false.
\end{lemma}
\leanref{Deduction.\allowbreak{}landau\_\allowbreak{}growth\_\allowbreak{}contra}{Development.\allowbreak{}AutomorphicArithmetic.\allowbreak{}DedekindLFunctions.\allowbreak{}LandauTheorems.\allowbreak{}LandauGrowthContra}{9--41}
\leanref{Deduction.\allowbreak{}landau\_\allowbreak{}window\_\allowbreak{}contra}{Development.\allowbreak{}AutomorphicArithmetic.\allowbreak{}DedekindLFunctions.\allowbreak{}LandauTheorems.\allowbreak{}landauWindow}{20--39}

\begin{proof}
Write $L=\log(e/d)>0$ and $\eta=L/2$.
After multiplication by $d^k$, the first left-hand side is at most
\[
 |C|\exp(\eta U-kL)\le |C|\exp(-\eta U),
\]
which is less than $c$ for every sufficiently large $U$.
The first assertion with coefficient $|M|+|C|$ gives
$\eta>0$ and $U_*$. Apply $(*)$ with this $\eta$ and
$U_0=\max\{U_*,0\}$. For the resulting $U,k$, the right-hand
side of $(*)$ is at most
$(|M|+|C|)\exp(\eta U)/e^k<c/d^k$, a contradiction.
\end{proof}

Assume the integral window hypothesis of~\ref{B1.29}, with
$0<\theta<1$ and $\delta>0$. By~\ref{B1.47}, there are fixed
$\sigma>2$ and $C\in\R$ such that, for every $\eta>0$ and
$U_0\in\R$, some $U\ge U_0$ and $k\in\N$, $k\ge U$, satisfy
\begin{equation}\label{eq:p1c5-window}
 |A_k(\sigma+it)|+\Re A_k(\sigma)
 \le C\frac{\exp(\eta U)}{(\sigma-\theta)^k}
 \qquad\text{for every }t\in\R.
\end{equation}
Lemma~\ref{B1.54} gives $\Re\rho<\sigma$ for every $\rho\in Z$.

Theorem~\ref{B1.31} follows if both alternatives of
Lemma~\ref{B1.56} are excluded for the fixed $\sigma,C$ in
\eqref{eq:p1c5-window}.

\begin{lemma}[The moving-centre case]\label{B1.61}
Under these hypotheses, the first alternative of
Lemma~\ref{B1.56} is impossible.
\end{lemma}
\leanref{Deduction.\allowbreak{}landau\_\allowbreak{}window}{Development.\allowbreak{}AutomorphicArithmetic.\allowbreak{}DedekindLFunctions.\allowbreak{}LandauTheorems.\allowbreak{}landauWindow}{92--133}

\begin{proof}
Let $D$ and $\rho_1$ be supplied by that alternative.
Since $D\le\sigma-\theta<\sigma$, Lemma~\ref{B1.57} supplies
$z=\sigma+it$, $\rho\in Z$, and $r\le D$ with
$d=|z-\rho|<r$ and $Z\cap B(z,r)=\{\rho\}$.
Here $d>0$, since $\Re\rho<\sigma$. Set
\[
 e=\frac{d+r}{2},\qquad E=\sigma-\theta,
 \qquad 0<d<e<r\le D\le E.
\]
The discs $B(z,r)$ and $B(\sigma,D)$ lie in $H$ because their centres
have real part $\sigma$ and their radii are at most $\sigma-\theta$.
By Definition~\ref{B1.33}, $\Phi$ is holomorphic on
$B(z,r)\setminus\{\rho\}$ and on $B(\sigma,D)$.
Lemmas~\ref{B1.59} and~\ref{B1.58} give constants
$M_1,M_2$, independent of $k$, with
\[
 \left|A_k(z)-\frac{c_\rho}{(z-\rho)^{k+1}}\right|
 \le\frac{M_1}{e^k},\qquad |A_k(\sigma)|\le\frac{M_2}{e^k}.
\]
The triangle inequality and $\Re A_k(\sigma)\ge-|A_k(\sigma)|$
therefore give
\[
 \frac{|c_\rho|/d}{d^k}
 \le |A_k(z)|+\Re A_k(\sigma)+\frac{M_1+M_2}{e^k}.
\]
For the $U,k$ supplied by~\eqref{eq:p1c5-window}, this is $(*)$
with $c=|c_\rho|/d>0$. Lemma~\ref{B1.60} gives the contradiction.
\end{proof}

\begin{lemma}[The real-pole case]\label{B1.62}
Under the same hypotheses, the second alternative of
Lemma~\ref{B1.56} is impossible. Hence every $\rho\in Z$ satisfies
$\Re\rho\le\theta$.
\end{lemma}
\leanref{Deduction.\allowbreak{}landau\_\allowbreak{}window}{Development.\allowbreak{}AutomorphicArithmetic.\allowbreak{}DedekindLFunctions.\allowbreak{}LandauTheorems.\allowbreak{}landauWindow}{134--166}

\begin{proof}
Let $\beta\in Z\cap\R$, $\beta>\theta$, and $r$ be supplied by the
second alternative. Put
\[
 d=\sigma-\beta>0,\qquad e=\frac{d+r}{2},\qquad
 E=\sigma-\theta.
\]
Then $d<e<r\le E$, and $B(\sigma,r)\subset H$ contains the single
pole $\beta$. Write $c_\beta>0$ for its real residue.
Lemma~\ref{B1.59} gives a fixed $M$ such that
\[
 \left|A_k(\sigma)-\frac{c_\beta}{d^{k+1}}\right|
 \le\frac{M}{e^k}.
\]
Since $c_\beta/d^{k+1}$ is real, the bound gives
$c_\beta/d^{k+1}-\Re A_k(\sigma)\le M/e^k$.
Adding $|A_k(\sigma)|$ to the right gives
\[
 \frac{c_\beta/d}{d^k}
 \le \Re A_k(\sigma)+|A_k(\sigma)|+\frac{M}{e^k}.
\]
Use~\eqref{eq:p1c5-window} at $t=0$.
For every $\eta>0$ and every lower bound on $U$, the resulting
inequality is $(*)$ with $c=c_\beta/d>0$.
Lemma~\ref{B1.60} gives a contradiction.

If any $\rho_0\in Z$ had $\Re\rho_0>\theta$,
Lemma~\ref{B1.56} would give one of the two excluded alternatives.
This proves Theorem~\ref{B1.31}. For a family $F$ satisfying the
window hypothesis of Theorem~\ref{B1.28}, Lemma~\ref{B1.29} gives
the logarithmic window condition for $g_F$. Apply
Theorem~\ref{B1.31} to the pole data of Theorem~\ref{B1.30}.
A zero with $\theta<\Re s<1$ would lie in $Z_F$ and contradict
the pole bound. Lemma~\ref{B1.32} excludes zeros with $\Re s\ge1$.
\end{proof}


\part{The arithmetic side: prime sums and auxiliary primes}

% ---- p2c1 ----
\chapter{Counting primes class by class}\label{ch:p2c1}

Fix a frame with finite abelian group $A$ and distinguished element
$a_0$, and put $A_0=\langle a_0\rangle$. The arithmetic data in
Alp\"oge's archive \cite{alpoge-2026} attach a Frobenius class and an
inertia subgroup to each rational prime. Character orthogonality
expresses the weighted prime-power count as a finite sum. After replacing
the weights on thick positions by a common weight, we average over
$A_0$. This bounds the real part of a weighted sum over the surviving
characters.

\section{The host and its local classes}

\begin{definition}[Host data used in the prime count]\label{B2.01}
Let $H$ be a host for the frame $\phi$ in the sense of
\Lean{Formalism.\allowbreak{}HC.\allowbreak{}HostAx}.
The frame has $e\in\N$ and a surjective homomorphism
$h:A\to(\Z/2^{e+2}\Z)^\times/\{\pm1\}$ whose restriction to
$A_0$ is injective. Write
$\mathcal G_\phi=(\{\pm1\}\times\{\pm1\})\times A$,
with the last factor written multiplicatively in host coordinates.
The host supplies primes $q,d$ and a positive integer $Q_S$ such that
\[
 \begin{gathered}
 e\ge2,\qquad q\equiv3\pmod8,\qquad d\equiv7\pmod{16},
 \qquad \left(\frac{-d}{q}\right)=1,\\
 2^{e+2}\mid Q_S,\qquad (Q_S,qd)=1,\qquad
 \left(\frac{-d}{p}\right)=1
 \quad\text{for every odd prime }p\mid Q_S.
 \end{gathered}
\]
Set $\mathfrak m=Q_Sqd$. The field data of $H$ include a Galois number
field $E$, an embedding $E\hookrightarrow\Q(\zeta_{\mathfrak m})$,
an isomorphism
$\operatorname{Gal}(E/\Q)\simeq\mathcal G_\phi$, and an element
$s_d\in E$ with
\[
 s_d^2=-d,\qquad
 \sigma(s_d)=\sigma_d s_d
 \quad\text{if }\sigma\longleftrightarrow((\sigma_q,\sigma_d),a).
\]
There is an Artin homomorphism
$\operatorname{art}:(\Z/\mathfrak m\Z)^\times\to\mathcal G_\phi$.
If $\operatorname{art}(u)=((\sigma_q,\sigma_d),a)$, then
$h(a)=[u\bmod 2^{e+2}]$.
It is compatible with restriction of the cyclotomic automorphism
$\zeta_{\mathfrak m}\mapsto\zeta_{\mathfrak m}^{\,u}$ and sends $-1$
to $((-1,-1),1)$. It also satisfies
\[
 \begin{split}
 u\equiv1\pmod{Q_Sq}&\ \Longrightarrow\
 \operatorname{art}(u)\in\{1,\operatorname{iod}\},
 \qquad \operatorname{iod}=((1,-1),1),\\
 u\equiv1\pmod{Q_Sd}&\ \Longrightarrow\
 \operatorname{art}(u)^2=1 .
 \end{split}
\]
The extension $E/\Q$ has degree $4|A|$. The Artin map is surjective,
with $d$-coordinate $(u/d)$ at $u$. For each prime $p\nmid\mathfrak m$,
arithmetic Frobenius corresponds to $\operatorname{art}(p)$ under
the Galois isomorphism. More generally, let
$\mathfrak m^{(p)}$ be the largest divisor of $\mathfrak m$ prime
to $p$. The inertia subgroup at $p$ is the Artin image of the
units congruent to $1$ modulo $\mathfrak m^{(p)}$. The Frobenius
coset is the Artin image of the units congruent to $p$ modulo
$\mathfrak m^{(p)}$.
\end{definition}
\leanref{Formalism.\allowbreak{}HC.\allowbreak{}Frame}{Formalism.\allowbreak{}SharpRestriction.\allowbreak{}RamificationFrobenius.\allowbreak{}TotientGaloisClasses.\allowbreak{}ZsmulGeneratorEqZeroOfProjection}{15--24}
\leanref{Formalism.\allowbreak{}HC.\allowbreak{}HostRaw}{Formalism.\allowbreak{}SharpRestriction.\allowbreak{}RamificationFrobenius.\allowbreak{}TotientGaloisClasses.\allowbreak{}ConjEMemHostFIff}{12--47}
\leanref{Formalism.\allowbreak{}HC.\allowbreak{}HostDer}{Formalism.\allowbreak{}SharpRestriction.\allowbreak{}RamificationFrobenius.\allowbreak{}TotientGaloisClasses.\allowbreak{}ConjEMemHostFIff}{87--171}
\leanref{Formalism.\allowbreak{}HC.\allowbreak{}HostAx}{Formalism.\allowbreak{}SharpRestriction.\allowbreak{}RamificationFrobenius.\allowbreak{}TotientGaloisClasses.\allowbreak{}ConjEMemHostFIff}{174--176}

For the sums we use additive coordinates. Put
\[
 E_0=(\Z/2\Z)^2,\qquad G=E_0\times A,\qquad
 \Phi=\{((v_d,v_q),h):v_d=0\},\qquad f=|\Phi|=2|A|.
\]
Write $\epsilon_d=((1,0),0)$, $c=((1,1),0)$ and
$\tau(t)=((0,0),t)$.
The character group $\widehat G$ consists of homomorphisms
$G\to\C^\times$. Write $\mathbf1$ for the trivial character.
Let
$\chi_d((v_d,v_q),h)=(-1)^{v_d}$ and
$\chi_q((v_d,v_q),h)=(-1)^{v_q}$.

\begin{lemma}[Changing coordinates and characters]\label{B2.02}
For every frame $\phi$ there is a group isomorphism
\[
 T:\mathcal G_\phi\longrightarrow G,\qquad
 T((\sigma_q,\sigma_d),a)
 =((\overline{\sigma_d},\overline{\sigma_q}),a),
 \qquad \overline{1}=0,\quad\overline{-1}=1.
\]
Here the last coordinate is converted to additive notation.
Put $\mathrm{cc}=((-1,-1),1)$. For a host $H$, its Galois
isomorphism sends complex conjugation to $\mathrm{cc}$.
We have $T(\operatorname{iod})=\epsilon_d$ and $T(\mathrm{cc})=c$.
For $t\in A$, the host element $\tau_1(t)=((1,1),t)$ satisfies
$T(\tau_1(t))=\tau(t)$.
The subgroup $\phi.\mathrm{Psi}=\{((\sigma_q,1),a)\}$ maps onto $\Phi$.
Composition with $T$ gives a bijection
$\widehat G\to\widehat{\mathcal G_\phi}$,
$\alpha\mapsto\theta_\alpha=\alpha\circ T$.
Pullback sends the trivial character to the trivial character.
At $((\sigma_q,\sigma_d),a)$, the pullbacks of $\chi_d$ and
$\chi_q$ take the values $\sigma_d$ and $\sigma_q$, respectively.
Then $\theta_\alpha$ is trivial on $\tau_1(A_0)$ exactly when
$\alpha$ is trivial on $\tau(A_0)$. The set $X_+$ of such host
characters with value $1$ at $\operatorname{iod}$ therefore corresponds
to the $\alpha$ trivial on $\tau(A_0)$ with $\alpha(\epsilon_d)=1$.
\end{lemma}
\leanref{UpperTheory.\allowbreak{}Comparison₂.\allowbreak{}toE\_\allowbreak{}bij}{Development.\allowbreak{}AutomorphicArithmetic.\allowbreak{}PrimitiveCosets.\allowbreak{}ToEBij}{37--85}
\leanref{UpperTheory.\allowbreak{}Comparison₂.\allowbreak{}chr\_\allowbreak{}bij}{Development.\allowbreak{}AutomorphicArithmetic.\allowbreak{}PrimitiveCosets.\allowbreak{}ChrBij}{21--91}

\begin{proof}
The map $\{\pm1\}\to\Z/2\Z$ displayed above is an isomorphism.
Read the two signs in the order $d,q$ and write the last coordinate
additively. The inverse reads the signs in the order $q,d$.
The formulas for the distinguished elements follow by substitution.
The defining condition for $\phi.\mathrm{Psi}$ is $\sigma_d=1$,
which becomes $v_d=0$.
Composition with an isomorphism is a bijection of character groups,
with inverse composition with $T^{-1}$.
Evaluation at $\tau_1(s)$ and at $\operatorname{iod}$ gives the
assertions about trivial restrictions and $X_+$.
\end{proof}

Let $I_p\subseteq G$ be the inertia subgroup, and choose a Frobenius
representative $F_p\in G$. At a ramified prime its class is
$F_p+I_p$.

\begin{lemma}[The three local cases]\label{B2.03}
For any host, let
$S=\{q\}\cup\{p:p\text{ is prime and }p\mid Q_S\}$.
Then $2,q\in S$, $d\notin S$, and
\[
 \begin{array}{ll}
 p\text{ prime},\ p\notin S,\ p\ne d:& I_p=\{0\},\\
 p\in S:& I_p\subseteq\Phi,\quad F_p\in\Phi,\\
 p=d:& I_d=\{0,\epsilon_d\}.
 \end{array}
\]
Consequently every coset $iF_d+I_d$ contains exactly one element
outside $\Phi$.
\end{lemma}
\leanref{UpperTheory.\allowbreak{}Comparison₂.\allowbreak{}hostFactsS}{Development.\allowbreak{}AutomorphicArithmetic.\allowbreak{}PrimitiveCosets.\allowbreak{}hostFactsS}{127--137}
\leanref{Deduction.\allowbreak{}CH.\allowbreak{}HostFactsS}{Development.\allowbreak{}AutomorphicArithmetic.\allowbreak{}DedekindLFunctions.\allowbreak{}SplitInertia}{14--23}

\begin{proof}
The divisibility $2^{e+2}\mid Q_S$ gives $2\in S$.
Also $q\in S$ and $d\nmid Q_S$, while $q\ne d$ because their
residues modulo $8$ differ. Thus $d\notin S$.
If $p$ is prime and $p\notin S\cup\{d\}$, then $p\nmid\mathfrak m$.
The local unit subgroup defining
inertia is trivial, so $I_p=\{0\}$.

Suppose $p\in S$. Since $d\mid\mathfrak m^{(p)}$, a unit defining
inertia is $1$ modulo $d$. Its $d$-sign is therefore $1$.
A unit defining Frobenius is $p$ modulo $d$, and its $d$-sign
is $(p/d)$. For odd $p$, quadratic reciprocity and
$d\equiv3\pmod4$ give
\[
 \left(\frac{p}{d}\right)=\left(\frac{-d}{p}\right)=1.
\]
This uses the splitting condition on $Q_S$ when $p\mid Q_S$,
and the host condition when $p=q$. For $p=2$ the same conclusion
follows from $d\equiv7\pmod8$. Thus $F_p\in\Phi$ and
$I_p\subseteq\Phi$.

At $d$, the local units are those congruent to $1$ modulo $Q_Sq$,
so their Artin images lie in $\{1,\operatorname{iod}\}$.
The Chinese remainder theorem supplies such a unit which is
$-1$ modulo $d$. Its $d$-sign is $(-1/d)=-1$, so its image is
$\operatorname{iod}$. Lemma~\ref{B2.02} now gives
$I_d=\{0,\epsilon_d\}$. Addition by $\epsilon_d$ exchanges
$\Phi$ and its complement.
\end{proof}

\section{The exact prime-power identity}

Let $A$ be any finite abelian group. An arithmetic datum $D$ consists of $q,d\in\N$ and, for each
$n\in\N$, a subgroup $I_n\le G$ and an element $F_n\in G$.
The local hypotheses mean that $d$ is prime and there is a finite
$S\subseteq\N$ with $2,q\in S$, $d\notin S$, and the three
local conditions displayed in Lemma~\ref{B2.03}.
Write $[z]_+=\max(z,0)$ and use $a/0=0$ throughout.

\begin{definition}[Local values, sums, and centring]\label{B2.04}
For $\alpha\in\widehat G$, a prime $p$, and $i\ge1$, set
\[
 V_\alpha(p^i)=\frac1{|I_p|}
 \sum_{h\in I_p}\alpha(iF_p+h),\qquad
 \operatorname{primVal}(\alpha,n)
 =\prod_{p\mid n}V_\alpha(p^{v_p(n)})\quad(n\ge1).
\]
At $n=1$ the product is empty. At $n=0$ the archive also uses the
empty product, so $\operatorname{primVal}(\alpha,0)=\operatorname{primVal}(\alpha,1)=1$.
Let $\ell:\widehat G\to\C$ be arbitrary. For every $y\in\R$ define
\[
 G_\alpha(y)=
 \begin{cases}
 G_z(y),&\alpha=\mathbf1,\\[2pt]
 \displaystyle\sum_{1\le n\le y}\Lambda(n)
 \operatorname{primVal}(\alpha,n)(y/n-1)+y\ell_\alpha,
 &\alpha\ne\mathbf1.
 \end{cases}
\]
Here $G_z$ and the total real logarithm are as in
Definition~\ref{B1.26}. A sum over $1\le n\le y$ is empty for $y<1$.
For $t\in A$, define the partner map $r_t(g)=c+g+\tau(t)$.
It maps $G\setminus\Phi$ bijectively onto $\Phi$.
Given $a\in\R$, $w:G\to\R$, and $N\in\N$, put
\[
\begin{split}
 \ASK_N(a,t,w)={}&
 \sum_{\substack{p\le N,\ p\ {\rm prime}\\
 p\ne2,q,d,\ F_p\notin\Phi}}
 2\log p
 \sum_{\substack{1\le i\le N\\iF_p\notin\Phi}}
 [a/p^i-w(r_t(iF_p))]_+\\
 &+\log d\sum_{i=1}^N
 \sum_{\substack{g\in iF_d+I_d\\g\notin\Phi}}
 [a/d^i-w(r_t(g))]_+ .
\end{split}
\]
For the same $a,t,w$, define the complex quantities
\[
 \mathrm{MAIN}(a,w)
 =a\left(\log a-1-\gamma+\ell_{\chi_d}
 -\frac1f\sum_{\rho\in\Phi}\log w(\rho)\right)
\]
and
\[
 \mathrm{De}(a,t,w)
 =\frac1f\sum_{g\notin\Phi}\sum_{\alpha\in\widehat G}
 \overline{\alpha(g)}\,w(r_t(g))\,
 G_\alpha\left(\frac{a}{w(r_t(g))}\right).
\]
\end{definition}
\leanref{Formalism.\allowbreak{}LFunctionChainParameters.\allowbreak{}characterSumData.\allowbreak{}primVal}{Formalism.\allowbreak{}SizeWindowBrackets.\allowbreak{}FirstPrimeNeSecondPrime}{65--67}
\leanref{Formalism.\allowbreak{}LFunctionChainParameters.\allowbreak{}characterSumData.\allowbreak{}GPsi}{Formalism.\allowbreak{}SizeWindowBrackets.\allowbreak{}FirstPrimeNeSecondPrime}{68--71}
\leanref{Formalism.\allowbreak{}LFunctionChainParameters.\allowbreak{}characterSumData.\allowbreak{}ASK}{Formalism.\allowbreak{}SizeWindowBrackets.\allowbreak{}FirstPrimeNeSecondPrime}{84--88}
\leanref{Formalism.\allowbreak{}LFunctionChainParameters.\allowbreak{}characterSumData.\allowbreak{}MAIN}{Formalism.\allowbreak{}SizeWindowBrackets.\allowbreak{}FirstPrimeNeSecondPrime}{89--92}
\leanref{Formalism.\allowbreak{}LFunctionChainParameters.\allowbreak{}characterSumData.\allowbreak{}De}{Formalism.\allowbreak{}SizeWindowBrackets.\allowbreak{}FirstPrimeNeSecondPrime}{93--97}

The value $V_\alpha(p^i)$ is unchanged when $F_p$ is replaced
by another representative of its inertia coset. Character
orthogonality on $I_p$ also gives
\[
 V_\alpha(p^i)=
 \begin{cases}
 \alpha(iF_p),&\alpha|_{I_p}=\mathbf1,\\
 0,&\alpha|_{I_p}\ne\mathbf1.
 \end{cases}
\]
In particular $|\operatorname{primVal}(\alpha,n)|\le1$.

\begin{lemma}[Orthogonality on an inertia coset]\label{B2.05}
For every arithmetic datum $D$, every prime $p$, every integer
$i\ge1$, and every $g\in G$,
\[
 \frac1f\sum_{\alpha\in\widehat G}
 \overline{\alpha(g)}\operatorname{primVal}(\alpha,p^i)
 =\frac{2}{|I_p|}\,\mathbf1_{\,g\in iF_p+I_p}.
\]
Under the local conditions of Lemma~\ref{B2.03}, the coefficient
on the indicated coset is $2$ at an unramified prime and $1$ at $d$.

\end{lemma}
\leanref{Formalism.\allowbreak{}LFunctionChainParameters.\allowbreak{}characterSumData.\allowbreak{}orthogonality\_\allowbreak{}ramified}{Formalism.\allowbreak{}SizeWindowBrackets.\allowbreak{}ChainSlopeWindows.\allowbreak{}StronglyIrreducible}{311--317}
\leanref{Formalism.\allowbreak{}LFunctionChainParameters.\allowbreak{}characterSumData.\allowbreak{}OrthoRam}{Formalism.\allowbreak{}SizeWindowBrackets.\allowbreak{}FirstPrimeNeSecondPrime}{116--122}

\begin{proof}
At $p^i$ the product defining $\operatorname{primVal}$ is
$V_\alpha(p^i)$. Since $|\alpha(g)|=1$, the left side is
\[
 \frac1{f|I_p|}\sum_{h\in I_p}
 \sum_{\alpha\in\widehat G}\alpha(iF_p+h-g).
\]
The inner sum is $|G|$ when $iF_p+h=g$ and is $0$ otherwise.
There is one such $h$ precisely when $g\in iF_p+I_p$.
Finally $|G|=2f$.
\end{proof}

\begin{lemma}[Expansion of one layer]\label{B2.06}
Assume the local conditions of Lemma~\ref{B2.03}.
Let $a\in\R$, $t\in A$, and $w:G\to\R$ satisfy
$w(\rho)>0$ for $\rho\in\Phi$. Let $N\in\N$ satisfy
\[
 a/w(\rho)\le N\qquad(\rho\in\Phi).
\]
Then
\[
 \ASK_N(a,t,w)
 =\frac1f\sum_{g\notin\Phi}\sum_{\alpha\in\widehat G}
 \overline{\alpha(g)}
 \sum_{n=1}^N\Lambda(n)\operatorname{primVal}(\alpha,n)
 [a/n-w(r_t(g))]_+ .
\]
The equality is in $\C$, with the real left side embedded in $\C$.
\end{lemma}
\leanref{Deduction.\allowbreak{}CharacterBalance.\allowbreak{}askExp}{Development.\allowbreak{}AutomorphicArithmetic.\allowbreak{}DedekindLFunctions.\allowbreak{}LandauTheorems.\allowbreak{}ChBalance}{122--197}

\begin{proof}
For $n>N$ and $\rho\in\Phi$ the positive part
$[a/n-w(\rho)]_+$ is zero. Expand the von Mangoldt sum
into prime powers, using $\Lambda(p^i)=\log p$, and interchange
the finite sums. Lemma~\ref{B2.05} makes the contribution
of a prime power $p^i$ equal to
\[
 \frac{2\log p}{|I_p|}
 \sum_{\substack{g\in iF_p+I_p\\g\notin\Phi}}
 [a/p^i-w(r_t(g))]_+ .
\]
For $p\in S$, the coset $iF_p+I_p$ is contained in $\Phi$,
so the inner sum is empty. For $p\notin S\cup\{d\}$, inertia
is trivial and the contribution is
\[
 2\log p\,
 \mathbf1_{\,iF_p\notin\Phi}
 [a/p^i-w(r_t(iF_p))]_+ .
\]
If $F_p\in\Phi$, all its multiples lie in $\Phi$.
If $F_p\notin\Phi$, its first binary coordinate is $1$,
so $iF_p\notin\Phi$ precisely when $i$ is odd.
At $d$, inertia has order $2$, which gives the second
line of the definition of $\ASK_N$.

Since $p^i\le N$ implies $p,i\le N$, the rectangular cutoff
retains every nonzero term. When $d>N$, every $d$-term vanishes.
\end{proof}

\begin{theorem}[Exact balance for one layer]\label{B2.07}
Under the hypotheses of Lemma~\ref{B2.06}, assume further that
$a>0$. For every choice of $\ell:\widehat G\to\C$,
\[
 \ASK_N(a,t,w)-\mathrm{MAIN}(a,w)
 =\mathrm{De}(a,t,w).
\]
\end{theorem}
\leanref{Deduction.\allowbreak{}ch\_\allowbreak{}balance}{Development.\allowbreak{}AutomorphicArithmetic.\allowbreak{}DedekindLFunctions.\allowbreak{}LandauTheorems.\allowbreak{}ChBalance}{200--261}

\begin{proof}
For the trivial character every inertia average equals $1$, so
$\operatorname{primVal}(\mathbf1,n)=1$.
Fix $g\notin\Phi$ and put $W=w(r_t(g))$. Since $r_t(g)\in\Phi$,
we have $W>0$ and $a/W\le N$. Then
\[
 \sum_{n=1}^N\Lambda(n)\operatorname{primVal}(\alpha,n)
 [a/n-W]_+
 = WG_\alpha(a/W)+R_\alpha(W),
\]
where
\[
 R_\alpha(W)=
 \begin{cases}
 a\log a-a\log W-(1+\gamma)a,&\alpha=\mathbf1,\\
 -a\ell_\alpha,&\alpha\ne\mathbf1.
 \end{cases}
\]
Indeed the terms with $n>a/W$ vanish, and multiplication of
the defining sum for $G_\alpha(a/W)$ by $W$ gives the displayed
formula.

The half-group character sum is
\[
 \sum_{g\notin\Phi}\overline{\alpha(g)}
 =\begin{cases}
 f,&\alpha=\mathbf1,\\
 -f,&\alpha=\chi_d,\\
 0,&\alpha\notin\{\mathbf1,\chi_d\}.
 \end{cases}
\]
Since $G\setminus\Phi=\epsilon_d+\Phi$, the character sum is
$\overline{\alpha(\epsilon_d)}\sum_{h\in\Phi}\overline{\alpha(h)}$,
which vanishes when $\alpha|_\Phi$ is nontrivial.
The two characters trivial on $\Phi$ are $\mathbf1$ and
$\chi_d$, whose values on the other coset are $1$ and $-1$.

Substitution in Lemma~\ref{B2.06} gives $\mathrm{De}$ and the
remainders. The coset sum leaves $a\ell_{\chi_d}$ from the
nontrivial characters. In the trivial remainder, change variables by
$\rho=r_t(g)$ to obtain
$a\log a-(1+\gamma)a-(a/f)\sum_{\rho\in\Phi}\log w(\rho)$.
Together these give $\mathrm{MAIN}(a,w)$.
\end{proof}

\section{Averaging the thick contribution}

Let $A_0$ be any subgroup of $A$. Define
\[
 P=\{\alpha\in\widehat G:\alpha(\tau(s))=1
 \text{ for every }s\in A_0\},
 \qquad e_1(t)=((1,0),t).
\]
For a finite subset $T_2\subseteq A$, put
$\delta=|T_2|/|A|$.
For $b,y\in\R$ and $u,t\in A$, define the uniform thick sum
\[
 \mathrm{Th}(u,t;y)
 =\frac1f\sum_{s\in T_2}\sum_{\alpha\in\widehat G}
 \alpha(-e_1(s+u-t))\,bG_\alpha(y).
\]

\begin{lemma}[Averaging over the subgroup]\label{B2.08}
Let $s\in A_0$, $t\in A$, and $\alpha\in\widehat G$. Then
\[
 \sum_{u\in A_0}\alpha(-e_1(s+u-t))
 =\begin{cases}
 |A_0|\alpha(e_1(t)),&\alpha\in P,\\
 0,&\alpha\notin P.
 \end{cases}
\]
Consequently, if $T_2\subseteq A_0$, then for all $b,y\in\R$
and $t\in A$,
\[
 \sum_{u\in A_0}\mathrm{Th}(u,t;y)
 =|A_0|\,b\,\frac{\delta}{2}
 \sum_{\alpha\in P}\alpha(e_1(t))G_\alpha(y).
\]
\end{lemma}
\leanref{Deduction.\allowbreak{}CharacterBalanceMoves.\allowbreak{}sum\_\allowbreak{}A}{Development.\allowbreak{}AutomorphicArithmetic.\allowbreak{}DedekindLFunctions.\allowbreak{}LandauTheorems.\allowbreak{}ChCore}{28--55}
\leanref{Deduction.\allowbreak{}CharacterBalanceMoves.\allowbreak{}sum\_\allowbreak{}Th}{Development.\allowbreak{}AutomorphicArithmetic.\allowbreak{}DedekindLFunctions.\allowbreak{}LandauTheorems.\allowbreak{}ChCore}{56--74}

\begin{proof}
The identity $-e_1(s+u-t)=e_1(t)+\tau(-s)+\tau(-u)$ gives
\[
 \sum_{u\in A_0}\alpha(-e_1(s+u-t))
 =\alpha(e_1(t))\alpha(\tau(-s))
 \sum_{u\in A_0}\alpha(\tau(-u)).
\]
For $\alpha\in P$ the last two factors are $1$ and $|A_0|$.
Otherwise the restriction to $A_0$ is a nontrivial character.
Translation by an element on which it differs from $1$ shows
that its sum is zero. Multiply the identity for each $\alpha$ by
$bG_\alpha(y)/f$, then sum over $s\in T_2$ and $\alpha\in\widehat G$.
The surviving coefficient is $|A_0|b|T_2|/f=|A_0|b\delta/2$.
\end{proof}

For $u\in A$, the thick positions form
$\mathcal T_u=\{((0,1),s+u):s\in T_2\}\subseteq\Phi$.
Fix a scale $X'>0$ and an error bound $\theta\ge0$. The weight comparison uses the estimate
\begin{equation}\label{eq:b2-thin}
 |G_\alpha(y)|\le y\theta
 \qquad(\alpha\in\widehat G,\ y\ge X'/3),
\end{equation}
and a bound for changing a positive weight,
\begin{equation}\label{eq:b2-lip}
 \left|vG_\alpha(a/v)-v'G_\alpha(a/v')\right|
 \le \frac{3a}{\min(v,v')}|v-v'|
 \qquad(a,v,v'>0,\ \alpha\in\widehat G).
\end{equation}
These are the archive's conditions \Lean{ThinScales} and \Lean{LipW}.

\begin{lemma}[Replacing the weights on thick positions]\label{B2.09}
Let $T_2\subseteq A$ and $u,t\in A$. Suppose
\[
 \begin{gathered}
 x\ge3,\quad b>0,\quad b'>0,\quad \eta\ge0,\quad
 X'>0,\quad\theta\ge0,\qquad b'X'\le xb,\\
 |a/b-x|\le1.
 \end{gathered}
\]
Assume \eqref{eq:b2-thin} and \eqref{eq:b2-lip}, and suppose
$w:G\to\R$ satisfies
\[
 \begin{array}{ll}
 \rho\in\Phi:&w(\rho)>0,\\
 \rho\in\mathcal T_u:&b\le w(\rho)\le b+\eta+1,\\
 \rho\in\Phi\setminus\mathcal T_u:&w(\rho)\le2b'.
 \end{array}
\]
Then
\[
 \Re\mathrm{Th}(u,t;a/b)
 \le\Re\mathrm{De}(a,t,w)
 +2f\left(3(x+1)(\eta+1)+(x+1)b\theta\right).
\]
\end{lemma}
\leanref{Deduction.\allowbreak{}CharacterBalanceMoves.\allowbreak{}th\_\allowbreak{}le}{Development.\allowbreak{}AutomorphicArithmetic.\allowbreak{}DedekindLFunctions.\allowbreak{}LandauTheorems.\allowbreak{}ChCore}{75--153}

\begin{proof}
The condition on $a/b$ gives $0<(x-1)b\le a\le(x+1)b$.
The inverse image of the thick position $((0,1),s+u)$
under $r_t$ is $e_1(s+u-t)$.
Since $\overline{\alpha(g)}=\alpha(-g)$, replacing
$w(r_t(g))G_\alpha(a/w(r_t(g)))$ by $bG_\alpha(a/b)$
at those positions, and by zero at the other positions,
gives exactly $\mathrm{Th}(u,t;a/b)$.

At a thick position, \eqref{eq:b2-lip} bounds the change by
\[
 \frac{3a}{b}|w(\rho)-b|
 \le3(x+1)(\eta+1).
\]
At a position $\rho\in\Phi\setminus\mathcal T_u$, we have
\[
 \frac{X'}3w(\rho)
 \le\frac23 b'X'
 \le\frac23xb
 \le(x-1)b
 \le a.
\]
Thus $a/w(\rho)\ge X'/3$, and \eqref{eq:b2-thin} gives
\[
 \left|w(\rho)G_\alpha(a/w(\rho))\right|
 \le a\theta\le(x+1)b\theta.
\]
There are $f$ positions and $|\widehat G|=|G|=2f$
characters. Divide the sum of these bounds by $f$ and
use $|\alpha(g)|=1$. This bounds the absolute value of
$\mathrm{De}-\mathrm{Th}$ by the stated error, and hence
gives the inequality of real parts.
\end{proof}

\begin{theorem}[From bounds on translates to an average bound]\label{B2.10}
Assume the local conditions of Lemma~\ref{B2.03}.
Let $A_0\subseteq A$ be a subgroup, let $TT,T_2\subseteq A$,
and let $m,N\in\N$. Choose $a:A\to\R$ and write $a_t=a(t)$.
Choose weights $w_{u,j}:G\to\R$ for $u\in A$ and $0\le j<m$.
Suppose $T_2\subseteq A_0$, $x\ge3$, $b,b',X'>0$,
$\eta,\theta\ge0$, and $b'X'\le xb$.
Assume $|a_t/b-x|\le1$ for every $t\in TT$.
For each $u\in A_0$ and $0\le j<m$, assume the position bounds
of Lemma~\ref{B2.09} for $w_{u,j}$, and require
$a_t/w_{u,j}(\rho)\le N$ for every $t\in TT$ and $\rho\in\Phi$.
Assume \eqref{eq:b2-thin} and the two cost inequalities
$2f(x+1)\theta\le1/x$ and $6f(x+1)(\eta+1)\le b/x$.
Define
\[
 \operatorname{gap}(u)=
 \sum_{t\in TT}\sum_{j=0}^{m-1}
 \left(\ASK_N(a_t,t,w_{u,j})
 -\Re\mathrm{MAIN}(a_t,w_{u,j})\right),
 \qquad Q=|TT|mb .
\]
For any $\mathscr E,Y_*\in\R$, if
$\operatorname{gap}(u)\le\mathscr E$ for every $u\in A_0$ and
\[
 |TT|Y_*
 \le \sum_{t\in TT}\Re
 \sum_{\alpha\in P}
 \alpha(e_1(t))G_\alpha(a_t/b),
\]
then
\[
 Q\,\frac{\delta}{2}\,Y_*
 \le\mathscr E+\frac{2Q}{x}.
\]
\end{theorem}
\leanref{Deduction.\allowbreak{}CharacterBalanceMoves.\allowbreak{}core}{Development.\allowbreak{}AutomorphicArithmetic.\allowbreak{}DedekindLFunctions.\allowbreak{}LandauTheorems.\allowbreak{}ChCore}{283--339}

\begin{proof}
Write $\Psi(z)=\sum_{1\le n\le z}\Lambda(n)$.
Rosser and Schoenfeld prove $\Psi(z)<1.03883z$ for $z>0$
\cite[Theorem~12]{rosser-schoenfeld-1962}.
For $z>0$ the archive uses the weaker estimate $\Psi(z)\le3z$,
which also holds at $z=0$ since $\Psi(0)=0$.
This bound gives \eqref{eq:b2-lip} for every $D$ and $\ell$.

Fix $a>0$. For $\alpha\ne\mathbf1$ and $v>0$, multiplication of the
definition of $G_\alpha$ by $v$ gives
\[
 vG_\alpha(a/v)=
 \sum_{n\ge1}\Lambda(n)\operatorname{primVal}(\alpha,n)
 [a/n-v]_+ +a\ell_\alpha.
\]
The constant $a\ell_\alpha$ cancels
in the difference for $v,v'$. For $n\le a/\min(v,v')$, each
positive part changes by at most $|v-v'|$. All other summands vanish.
The Chebyshev bound and $|\operatorname{primVal}|\le1$
prove \eqref{eq:b2-lip}.

For the trivial character, take $0<v\le v'$ and let $B$
be the decrease of $\sum_n\Lambda(n)[a/n-v]_+$ when $v$
is replaced by $v'$. Then
\[
 0\le B\le3(a/v)(v'-v),\qquad
 0\le a\log(v'/v)\le(a/v)(v'-v).
\]
The difference $vG_z(a/v)-v'G_z(a/v')$ equals
$B-a\log(v'/v)$. Both nonnegative quantities are bounded
by $3(a/v)(v'-v)$, so their difference has absolute
value at most that bound. Interchanging $v,v'$ covers
the other order.
\leanref{Formalism.\allowbreak{}LFunctionChainParameters.\allowbreak{}ThickSetMeanBounds.\allowbreak{}AN.\allowbreak{}psi\_\allowbreak{}le\_\allowbreak{}three}{Formalism.\allowbreak{}SizeWindowBrackets.\allowbreak{}ChainSlopeWindows.\allowbreak{}LipschitzWeightAndMinimum}{668--691}
\leanref{Formalism.\allowbreak{}LFunctionChainParameters.\allowbreak{}ThickSetMeanBounds.\allowbreak{}lipschitz\_\allowbreak{}weight\_\allowbreak{}all}{Formalism.\allowbreak{}SizeWindowBrackets.\allowbreak{}ChainSlopeWindows.\allowbreak{}LipschitzWeightAndMinimum}{957--965}

For each $t\in TT$ the hypotheses imply $a_t>0$.
Theorem~\ref{B2.07} and Lemma~\ref{B2.09} therefore give,
for every $u\in A_0$ and $0\le j<m$,
\[
 \Re\mathrm{Th}(u,t;a_t/b)
 \le \ASK_N(a_t,t,w_{u,j})
 -\Re\mathrm{MAIN}(a_t,w_{u,j})+\frac{2b}{x}.
\]
The two cost inequalities bound each error contribution by $b/x$.
After summation over $t,j$, the assumed bound on the gap yields
\[
 \sum_{t\in TT}\sum_{j=0}^{m-1}
 \Re\mathrm{Th}(u,t;a_t/b)
 \le\mathscr E+\frac{2Q}{x}.
\]
Sum over $u\in A_0$. By Lemma~\ref{B2.08}, the left side becomes
\[
 |A_0|\,m b\,\frac{\delta}{2}
 \sum_{t\in TT}\Re\sum_{\alpha\in P}
 \alpha(e_1(t))G_\alpha(a_t/b).
\]
Divide by $|A_0|$. Multiplying the assumed lower bound for the sum
over $t\in TT$ by $mb\delta/2\ge0$ gives
$Q(\delta/2)Y_*\le\mathscr E+2Q/x$.
\end{proof}


% ---- p2c2 ----
\chapter{Reading off the target sum}\label{ch:p2c2}

Use the arithmetic data and functions \(G_\alpha\) of \ref{B2.04}. Write \(1\) for the trivial character of \(G=E_0\times A\), so \(G_1=G_z\). Let \(P\) be the surviving set of \ref{B2.08}. For \(K_p\subseteq\widehat G\) and \(c\in\R\), set \(w_c(\alpha)=1\) when \(\alpha|_{\{0\}\times A}=1\), and \(w_c(\alpha)=c\) otherwise. Put
\[
 G_m^A(y)=\sum_{\alpha\in K_p}w_c(\alpha)\Re G_\alpha(y),
 \qquad
 \Sigma_A(y)=\sum_{\alpha\in P\setminus K_p}|G_\alpha(y)|.
\]

\section{A common point of evaluation}

The condition \Lean{DerivInt} with \(C_d\ge0\) means that, for \(y,z\ge1\),
\begin{align}
 |G_\alpha(y)-G_\alpha(z)|
 &\le |y-z|\bigl(\log\max(y,z)+C_d+|\ell_\alpha|\bigr)
 &&(\alpha\ne1),\label{eq:p2c2-deriv}\\
 |G_z(y)-G_z(z)|&\le C_d|y-z|.\label{eq:p2c2-zeta-lip}
\end{align}

\begin{lemma}[Moving one evaluation point]\label{B2.11}
Assume \Lean{DerivInt} with \(C_d\ge0\). For \(x\ge3\), \(|y-x|\le1\), and \(\alpha\ne1\),
\[
 |G_\alpha(y)-G_\alpha(x)|\le\log(x+1)+C_d+|\ell_\alpha|,
 \qquad
 |G_\alpha(x)|\le x(\log x+C_d+|\ell_\alpha|).
\]
For the trivial character, \(|G_z(y)-G_z(x)|\le C_d\).
\end{lemma}
\leanref{Deduction.\allowbreak{}CharacterBalanceMoves.\allowbreak{}move}{Development.\allowbreak{}AutomorphicArithmetic.\allowbreak{}DedekindLFunctions.\allowbreak{}LandauTheorems.\allowbreak{}ChCore}{154--166}
\leanref{Deduction.\allowbreak{}CharacterBalanceMoves.\allowbreak{}triv}{Development.\allowbreak{}AutomorphicArithmetic.\allowbreak{}DedekindLFunctions.\allowbreak{}LandauTheorems.\allowbreak{}ChCore}{167--180}
\leanref{Formalism.\allowbreak{}LFunctionChainParameters.\allowbreak{}ThickSetMeanBounds.\allowbreak{}DerivInt}{Formalism.\allowbreak{}SizeWindowBrackets.\allowbreak{}GrowthBoundGivesRiemannHypothesis}{62--65}
\begin{proof}
Since \(y\ge2\) and \(\max(y,x)\le x+1\), \eqref{eq:p2c2-deriv} gives the first bound. The defining sum gives \(G_\alpha(1)=\ell_\alpha\), so
\[
 |G_\alpha(x)|\le(x-1)(\log x+C_d+|\ell_\alpha|)+|\ell_\alpha|
 \le x(\log x+C_d+|\ell_\alpha|).
\]
The trivial-character bound follows from \eqref{eq:p2c2-zeta-lip}.
\end{proof}

\begin{lemma}[Comparison for each character]\label{B2.12}
Let \(TT\subseteq A\), \(n_T=|TT|\), \(x\ge3\), and \(|y_t-x|\le1\) for \(t\in TT\). Assume \Lean{DerivInt} with \(C_d\ge0\). Let \(K_p\subseteq\widehat G\), \(c,L,R\in\R\), and \(\rho,K_o\ge0\). Suppose \(1\in K_p\), \(\alpha(\epsilon_d)=1\) for \(\alpha\in K_p\), and
\begin{align*}
 \left|\sum_{t\in TT}\alpha(0,t)-cn_T\right|&\le\rho n_T
 &&(\alpha\in K_p,\ \alpha|_{\{0\}\times A}\ne1),\\
 |\ell_\alpha|&\le K_o&&(\alpha\in K_p,\ \alpha\ne1),\\
 |\ell_\alpha|&\le L&&(\alpha\in P\setminus K_p).
\end{align*}
Require \(C_d\le R\) and
\[
 \begin{aligned}
 \log(x+1)+C_d+K_o+\rho x(\log x+C_d+K_o)&\le R,\\
 \log(x+1)+C_d+L&\le R.
 \end{aligned}
\]
Set \(v_\alpha=w_c(\alpha)\Re G_\alpha(x)\) for \(\alpha\in K_p\), and \(v_\alpha=-|G_\alpha(x)|\) otherwise. Then
\[
 n_T(v_\alpha-R)\le
 \sum_{t\in TT}\Re\bigl(\alpha(e_1(t))G_\alpha(y_t)\bigr)
 \qquad(\alpha\in P).
\]
If \(K_p\subseteq P\), then \(\sum_{\alpha\in P}v_\alpha=G_m^A(x)-\Sigma_A(x)\).
\end{lemma}
\leanref{Deduction.\allowbreak{}CharacterBalanceMoves.\allowbreak{}one\_\allowbreak{}char}{Development.\allowbreak{}AutomorphicArithmetic.\allowbreak{}DedekindLFunctions.\allowbreak{}LandauTheorems.\allowbreak{}ChCore}{202--282}
\leanref{Deduction.\allowbreak{}CharacterBalanceMoves.\allowbreak{}sum\_\allowbreak{}vv}{Development.\allowbreak{}AutomorphicArithmetic.\allowbreak{}DedekindLFunctions.\allowbreak{}LandauTheorems.\allowbreak{}ChCore}{187--200}
\begin{proof}
For \(\alpha=1\), sum \(G_z(y_t)\ge G_z(x)-C_d\).
For \(\alpha\ne1\) in \(K_p\), the identities \(e_1(t)=\epsilon_d+(0,t)\) and \(\alpha(\epsilon_d)=1\) give \(\alpha(e_1(t))=\alpha(0,t)\).
Put \(w=w_c(\alpha)\) and \(S_\alpha=\sum_{t\in TT}\alpha(0,t)\).
Then \(|S_\alpha-wn_T|\le\rho n_T\), with left side zero when \(\alpha|_{\{0\}\times A}=1\), and
\[
 \begin{split}
 \sum_{t\in TT}\alpha(0,t)G_\alpha(y_t)-wn_TG_\alpha(x)
 ={}&\sum_{t\in TT}\alpha(0,t)(G_\alpha(y_t)-G_\alpha(x))\\
 &+(S_\alpha-wn_T)G_\alpha(x).
 \end{split}
\]
Lemma~\ref{B2.11} bounds its absolute value by
\(n_T\{\log(x+1)+C_d+K_o+\rho x(\log x+C_d+K_o)\}\le n_TR\).
Take real parts.

For \(\alpha\in P\setminus K_p\), we have \(\alpha\ne1\), and Lemma~\ref{B2.11} gives
\(|G_\alpha(y_t)|\le|G_\alpha(x)|+\log(x+1)+C_d+L\le|G_\alpha(x)|+R\).
Since \(|\alpha(e_1(t))|=1\), summation proves the required lower bound.
For \(K_p\subseteq P\), summing the definition of \(v_\alpha\) gives the stated identity.
\end{proof}

\section{The geometric bound and the arithmetic comparison}

Write \(n=|A|\) and \(n_2=|T_2|\), with \(Q,\delta,\operatorname{gap}\) as in \ref{B2.10}. Let \(k>0\) be the rank and \(a_E\ge0\) the top-weight scale. Define
\begin{equation}\label{eq:p2c2-bracket}
 \mathcal B=n_Ta_E+n_Ta_E\frac{m^2}{k}
 +m(2nmbx^\varepsilon)
 +m\frac{n_2kbx^\varepsilon}{n(nx^{-\varepsilon})}.
\end{equation}
Write \(\ASK_D(N,a,t,w)\) for the sum \(\ASK_N(a,t,w)\) of Definition~\ref{B2.04} with datum \(D\). The geometric main term \(\Main_{\mathscr D}\) is defined in Definition~\ref{B5.34}.

\begin{theorem}[The bound at each translation]\label{B2.13}
Let \(\mathscr D\) be a datum of \ref{B5.43}, with finite translation and top sets \(U,\mathcal T\), finite nonempty word sets, and \(m,f\in\N\). Choose a top-weight function \(a:A\to\R\), written \(a_t=a(t)\), and maps
\[
 \operatorname{tr}:U\to A,\qquad e_T:\mathcal T\simeq TT,\qquad
 e_P:\{0,\ldots,f-1\}\simeq\Phi .
\]
Assume \(\mathscr D.a_T(\tau)=a_{e_T(\tau)}\) for every \(\tau\in\mathcal T\) and
\[
 \mathscr D.w(u,j,r)=w_{\operatorname{tr}(u),j}(e_P(r))
 \qquad(u\in U,\ 0\le j<m,\ 0\le r<f).
\]
Assume \(\mathscr D.\gamma=\gamma\).
Require the real parameter \(L_pL\) of Definition~\ref{B5.34} to equal \(\Re\ell_{\chi_d}\), for the sign character \(\chi_d\) of \ref{B2.02}. Assume also that
\[
 \ASK_{\mathscr D}(u)=
 \sum_{t\in TT}\sum_{j=0}^{m-1}\ASK_D(N,a_t,t,w_{\operatorname{tr}(u),j})
 \qquad(u\in U).
\]
Let \(\mathcal C\) be the chain constants and \(X'',e_K,K\in\R\). Assume \(\mathrm{H}1\) through \(\mathrm{H}11\) of Theorem~\ref{B5.43}. For the same data, impose the window hypotheses of Theorem~\ref{B5.45} through its bound on \((\bar\eta+4\sqrt{\bar e})x^{c_{HH}}b\). Set \(C_{\rm err}=e_{\rm Div}+e_{\rm Loss}+3\), and identify \(x,\varepsilon,b,k\) with the parameters of \(\mathscr D\). Put \(a_E=\mathscr D.a\). Suppose
\[
 \begin{gathered}
 B_0\le2nmbx^\varepsilon,\qquad
 N_c\le\frac{n_2kbx^\varepsilon}{n(nx^{-\varepsilon})},\\
 \log x\le x^\varepsilon,\qquad
 200(1+K)^2(1+c_{HH})(1+e_K)\le x^\varepsilon,
 \end{gathered}
\]
where \(c_{HH}\) is the component of \(\mathcal C\). Then, for every \(u\in U\),
\[
 \operatorname{gap}(\operatorname{tr}(u))
 =\ASK_{\mathscr D}(u)-\Main_{\mathscr D}(u)
 \le x^{C_{\rm err}\varepsilon}\mathcal B.
\]
\end{theorem}
\leanref{Deduction.\allowbreak{}ch\_\allowbreak{}top\_\allowbreak{}main\_\allowbreak{}eq}{Development.\allowbreak{}AutomorphicArithmetic.\allowbreak{}DedekindLFunctions.\allowbreak{}LandauTheorems.\allowbreak{}ChTop}{11--35}
\leanref{Deduction.\allowbreak{}ch\_\allowbreak{}top}{Development.\allowbreak{}AutomorphicArithmetic.\allowbreak{}DedekindLFunctions.\allowbreak{}LandauTheorems.\allowbreak{}ChTop}{37--72}
\begin{proof}
Apply Theorem~\ref{B5.47}. Its bracket is \(\mathcal B_x(\mathscr D.a)=\mathcal B\) by \eqref{eq:p2c2-bracket}.
\end{proof}

For fixed real \(M',K_o,C_E,c\), the condition \Lean{CharsA} is
\begin{equation}\label{eq:p2c2-chars}
 \begin{gathered}
 1\in K_p,\quad K_p\subseteq P,\quad
 \alpha(\epsilon_d)=1\ (\alpha\in K_p),\quad |P|\le M',\\
 \left|\sum_{t\in TT}\alpha(0,t)-cn_T\right|
 \le M'x^{-1+\varepsilon}n_T
 \quad(\alpha\in K_p,\ \alpha|_{\{0\}\times A}\ne1),\\
 |\ell_\alpha|\le K_o\quad(\alpha\in K_p,\ \alpha\ne1),\qquad
 |\ell_\alpha|\le C_E\log x\quad(\alpha\in P\setminus K_p).
 \end{gathered}
\end{equation}

\begin{theorem}[Arithmetic comparison]\label{B2.14}
Fix \(C,C_E,K_o,c,M'\in\R\) and \(\varepsilon>0\), with \(C\ge0\) and \(|c|\le1\).
There exists \(x_1\), depending on these constants, such that the following holds for \(x\ge x_1\).

Let \(A\) be a finite abelian group, \(A_0\le A\), \(K_p\subseteq\widehat G\), and \(S\subset\N\) finite. Let \(m,N\in\N\) and \(D,\ell,TT,T_2,a_t,w_{s,j}\) be arithmetic data as above. Assume \(d\) is prime, \(2,q\in S\), and \(d\notin S\). For primes \(p\notin S\cup\{d\}\), require \(I_p=\{0\}\). For every \(p\in S\), require \(I_p\subseteq\Phi\) and \(F_p\in\Phi\). Require \(I_d=\{0,\epsilon_d\}\).

Let \(k>0\) and \(a_E\ge0\). For real \(b,b',\eta,X',\theta\), require \(b,b',X'>0\), \(\eta,\theta\ge0\), and \(b'X'\le xb\). Impose the weight and cost hypotheses of \ref{B2.10}. Its thin-value hypotheses and the cutoff \(a_t/w_{s,j}(\rho)\le N\) also apply.
Require \(T_2\subseteq A_0\), \eqref{eq:p2c2-chars}, and
\begin{equation}\label{eq:p2c2-sizes}
 \begin{gathered}
 a_E\le2xb,\qquad \tfrac12nx^{-\varepsilon/8}\le n_T,\qquad
 x^{-\varepsilon}\le n_2/n,\\
 x^{1/2-\varepsilon}\le m\le x^{1/2+\varepsilon/2},\qquad
 xm\le k\le n^2x^{1/2-\varepsilon}.
 \end{gathered}
\end{equation}
If \(\operatorname{gap}(s)\le x^{C\varepsilon}\mathcal B\) for every \(s\in A_0\), then
\(G_m^A(x)\le x^{1/2+(C+3)\varepsilon}+\Sigma_A(x)\).
\end{theorem}
\leanref{Deduction.\allowbreak{}ch\_\allowbreak{}core}{Development.\allowbreak{}AutomorphicArithmetic.\allowbreak{}DedekindLFunctions.\allowbreak{}LandauTheorems.\allowbreak{}ChCore}{343--442}
\begin{proof}
To obtain \Lean{DerivInt} with an absolute \(C_d\ge0\), use \(|\operatorname{primVal}(\alpha,n)|\le1\) and
\(|[y/n-1]_+-[z/n-1]_+|\le|y-z|/n\) to bound the nontrivial-character difference by
\[
 |y-z|\left(\sum_{n\le\max(y,z)}\frac{\Lambda(n)}n+|\ell_\alpha|\right).
\]
Montgomery and Vaughan give \(\sum_{n\le u}\Lambda(n)/n=\log u+O(1)\) for \(u\ge2\) \cite[\S2.2, Theorem~2.7(a)]{montgomery-vaughan-2007}. For integers \(r\ge1\), the explicit bounds are
\[
 \log r-1\le\sum_{n\le r}\frac{\Lambda(n)}n\le\log r+\log4+4.
\]
Taking \(r=\lfloor\max(y,z)\rfloor\) gives \eqref{eq:p2c2-deriv}. The same bounds and \(0\le\log y-\log\lfloor y\rfloor\le\log2\) for \(y\ge1\) bound
\(G_z'(y)=\sum_{n\le y}\Lambda(n)/n-\log y+\gamma\)
uniformly between integers. Integrate there and use continuity at the integers to obtain \eqref{eq:p2c2-zeta-lip}.
\leanref{Formalism.\allowbreak{}LogFactorialMertensBounds.\allowbreak{}mertens\_\allowbreak{}lower}{Formalism.\allowbreak{}ZetaNonvanishing.\allowbreak{}GzLipschitz}{86--105}
\leanref{Formalism.\allowbreak{}LogFactorialMertensBounds.\allowbreak{}mertens\_\allowbreak{}upper}{Formalism.\allowbreak{}ZetaNonvanishing.\allowbreak{}GzLipschitz}{55--84}
\leanref{Formalism.\allowbreak{}LFunctionChainParameters.\allowbreak{}ThickSetMeanBounds.\allowbreak{}derivative\_\allowbreak{}integral\_\allowbreak{}all}{Formalism.\allowbreak{}SizeWindowBrackets.\allowbreak{}ChainSlopeWindows.\allowbreak{}LipschitzWeightAndMinimum}{966--982}

Put \(A_1=5+C_d+|K_o|+|M'|(4+C_d+|K_o|)+4|C_E|\) and \(R=A_1x^{1/4+\varepsilon}\).
Choose \(x_1\ge3\) so that \(26\le x^{3\varepsilon/8}\) and \(|M'|A_1\le x^{1/4}\) for \(x\ge x_1\).
Apply Lemma~\ref{B2.12} with \(y_t=a_t/b\), \(\rho=|M'|x^{-1+\varepsilon}\), and bounds \(|K_o|\), \(L=|C_E|\log x\).
Use \(x^\varepsilon\ge1\) together with \(\log x\le4x^{1/4}\) and \(\log(x+1)\le5x^{1/4}\). These inequalities give
\[
 \rho x(\log x+C_d+|K_o|)
 \le |M'|(4+C_d+|K_o|)x^{1/4+\varepsilon}.
\]
Summing Lemma~\ref{B2.12} over \(P\) gives the lower bound
\(n_T(G_*-|P|R)\), where \(G_*=G_m^A(x)-\Sigma_A(x)\).
Since \(|P|R\le |M'|A_1x^{1/4+\varepsilon}\le x^{1/2+\varepsilon}\), we obtain
\begin{equation}\label{eq:p2c2-top-lower}
 n_T(G_*-4x^{1/2+\varepsilon})
 \le\sum_{t\in TT}\Re\sum_{\alpha\in P}\alpha(e_1(t))G_\alpha(a_t/b).
\end{equation}
Apply \ref{B2.10}, with the stated local conditions and error \(x^{C\varepsilon}\mathcal B\), to obtain
\begin{equation}\label{eq:p2c2-balance}
 Q\frac\delta2(G_*-4x^{1/2+\varepsilon})
 \le x^{C\varepsilon}\mathcal B+\frac{2Q}{x}.
\end{equation}
Apply Theorem~\ref{B5.55} to \eqref{eq:p2c2-balance} with \(G_x=G_*\), \(E=x^{C\varepsilon}\mathcal B\), and \(\delta_2=\delta=n_2/n\).
Since \(T_2\subseteq A_0\subseteq A\), we have \(n_2\le n\). The other size bounds and the choice of \(x_1\) give the remaining hypotheses.
\end{proof}

\section{Reading the characters on the host}

Let \(\phi\) be a frame and \(H\) a host satisfying \ref{B2.01}.
Let \(D\) be the arithmetic datum obtained from the Artin map of \(H\) in the coordinates of \ref{B2.02}.
Write \(\theta_\alpha=\operatorname{chr}_\phi(\alpha)\) for the host character corresponding to \(\alpha\), and \(\chi_\alpha=H.\operatorname{prim}(\theta_\alpha)\) for its primitive Dirichlet representative with its conductor.
Compatibility requires, for every \(\alpha\ne1\), that each primitive Dirichlet character of positive modulus with values \(\operatorname{primVal}(\alpha,n)\) for all \(n\ge1\) have logarithmic derivative \(\ell_\alpha\) at \(1\).

\begin{lemma}[Primitive representatives]\label{B2.15}
The map \(H.\operatorname{prim}\) is injective and sends \(1\) to the primitive trivial character. If \(\ell\) is compatible, then
\[
 \ell_\alpha=\frac{L'(1,\chi_\alpha)}{L(1,\chi_\alpha)}\quad(\alpha\ne1),
 \qquad
 G_\alpha(y)=g_{\chi_\alpha}(y)\quad(\alpha\in\widehat G,\ y\in\R),
\]
where \(g_1=G_z\).
\end{lemma}
\leanref{UpperTheory.\allowbreak{}Comparison₂.\allowbreak{}prim\_\allowbreak{}read}{Development.\allowbreak{}AutomorphicArithmetic.\allowbreak{}PrimitiveCosets.\allowbreak{}PrimRead}{48--97}
\leanref{UpperTheory.\allowbreak{}Comparison₂.\allowbreak{}prim\_\allowbreak{}inj}{Development.\allowbreak{}AutomorphicArithmetic.\allowbreak{}PrimitiveCosets.\allowbreak{}PrimInj}{6--27}
\leanref{Formalism.\allowbreak{}LFunctionChainParameters.\allowbreak{}LHost.\allowbreak{}primitive\_\allowbreak{}value\_\allowbreak{}primitive}{Formalism.\allowbreak{}SizeWindowBrackets.\allowbreak{}ChainSlopeWindows.\allowbreak{}GaloisMapComplexConjugation}{233--245}
\begin{proof}
Put \(M_H=Q_Sqd\) and \(r_H=T\circ\operatorname{art}\), with \(T\) as in Lemma~\ref{B2.02}. On units modulo \(M_H\), we have \(\theta_\alpha\circ\operatorname{art}=\alpha\circ r_H\). The Dirichlet character extending this common character has primitive representative \(\chi_\alpha\).
Let \(f_\alpha\mid M_H\) be its conductor. For a prime \(p\), inertia \(I_p\) is the image under \(r_H\) of the kernel of reduction of units modulo the prime-to-\(p\) part \(M_H^{(p)}\). Thus \(\alpha|_{I_p}=1\) precisely when the pullback factors through \(M_H^{(p)}\), equivalently when \(p\nmid f_\alpha\).
In that case the Frobenius unit is congruent to \(p\) modulo \(f_\alpha\), and for \(i\ge1\)
\[
 \frac1{|I_p|}\sum_{h\in I_p}\alpha(iF_p+h)=\chi_\alpha(p)^i.
\]
When \(p\mid f_\alpha\), the restriction to \(I_p\) is nontrivial. Translation by an \(h_0\in I_p\) with \(\alpha(h_0)\ne1\) multiplies this sum by \(\alpha(h_0)\), so the sum is zero. In this case \(\chi_\alpha(p)=0\) as well.
Multiply over \(p^i\Vert n\) to obtain
\(\operatorname{primVal}(\alpha,n)=\chi_\alpha(n)\) for every \(n\ge1\).
Compatibility supplies the logarithmic derivative.

Equal primitive representatives induce equal characters on units modulo \(M_H\). Artin surjectivity then gives equal host characters. The trivial character pulls back and reduces to the primitive trivial character, so \(\chi_\alpha=1\) precisely when \(\alpha=1\).

For \(\alpha\ne1\), substituting these identities into \(G_\alpha(y)\) gives \(g_{\chi_\alpha}(y)\), since \([y/n-1]_+\) equals \(y/n-1\) for \(1\le n\le y\) and vanishes for \(n>y\). At \(\alpha=1\) both functions equal \(G_z\).
\end{proof}

\begin{lemma}[The number of surviving characters]\label{B2.16}
For a frame \(\phi\), let \(A=\phi.A\), \(A_0=\langle a_0\rangle\), and \(o=\phi.o\). Then \([A:A_0]=o\) and \(|P|=4o\).
For any subgroup \(A_0\) of a finite abelian group \(A\),
\[
 P\simeq\widehat{E_0\times(A/A_0)},\qquad |P|=4[A:A_0].
\]
\end{lemma}
\leanref{UpperTheory.\allowbreak{}Comparison₂.\allowbreak{}card\_\allowbreak{}PsiA}{Development.\allowbreak{}AutomorphicArithmetic.\allowbreak{}PrimitiveCosets.\allowbreak{}CardPsiA}{60--63}
\leanref{card\_\allowbreak{}PsiA\_\allowbreak{}index}{Development.\allowbreak{}AutomorphicArithmetic.\allowbreak{}PrimitiveCosets.\allowbreak{}CardPsiA}{26--58}
\leanref{Formalism.\allowbreak{}HC.\allowbreak{}Frame.\allowbreak{}index\_\allowbreak{}A0\_\allowbreak{}eq\_\allowbreak{}o}{Formalism.\allowbreak{}SharpRestriction.\allowbreak{}RamificationFrobenius.\allowbreak{}TotientGaloisClasses.\allowbreak{}FrameFacts}{214--218}
\begin{proof}
Pullback along \(E_0\times A\to E_0\times(A/A_0)\) identifies the quotient's character group with \(P\), since these are exactly the characters constant on its fibres.

For a finite abelian group \(B\), Serre gives \(|\widehat B|=|B|\) \cite[Theorems~7 and~9]{serre-1977}. Apply this identity with \(B=E_0\times(A/A_0)\). The frame gives \(\phi.o=|A|/\operatorname{ord}(a_0)=[A:A_0]\).
\end{proof}

A level \(\Lambda\) has an inversion-stable finite family \(\mathcal K=\Lambda.\mathrm{KEPT}\) of primitive characters containing \(1\), and \(o=\Lambda.o\in\N\). Its designated \(\chi_{\rm sp}\), when present, is a real odd member of \(\mathcal K\).
For \(\chi\ne1\), put \(\kappa_\chi=L'(1,\chi)/L(1,\chi)\).
Put \(\kappa_1=0\), since Lean's total derivative vanishes at the pole of \(\zeta\). Set
\[
 c=(3o+1)^{-1},\qquad K_H=\{\alpha:\chi_\alpha\in\mathcal K\},
 \qquad K_o=\sum_{\chi\in\mathcal K}|\kappa_\chi|.
\]
Let \(\mathrm{Kept}_H(\mathcal K)=\{\theta\in\widehat{\mathcal G_\phi}:H.\operatorname{prim}(\theta)\in\mathcal K\}\).
The bounds in Definition~\ref{B2.58} apply to
\(\mathrm{Bounded}=(X_+\setminus\mathrm{Kept}_H(\mathcal K))\cup X_-\).
Here \(X_\pm\) consists of host characters trivial on \(A_0\) with value \(\pm1\) on \(\mathrm{iod}\), whose additive coordinate is \(\epsilon_d\).

\begin{lemma}[Coset shares give the required weights]\label{B2.17}
Assume \(\mathrm{EAtH}(\Lambda,H,TT,x,\varepsilon,C)\), with \(x>0\), \(\varepsilon,C\in\R\), and compatible \(\ell\). Then \Lean{CharsA} holds with
\(K_p=K_H\), \(c=(3o+1)^{-1}\), \(M'=4o\), \(C_E=C\), and \(K_o\) as above.
\end{lemma}
\leanref{UpperTheory.\allowbreak{}Comparison₂.\allowbreak{}charsA}{Development.\allowbreak{}AutomorphicArithmetic.\allowbreak{}PrimitiveCosets.\allowbreak{}charsA}{80--125}
\leanref{Formalism.HC.Frame.Shares}{Formalism.SharpRestriction.RamificationFrobenius.TotientGaloisClasses.\allowbreak{}ExplicitFormulaErrorOfHostConditions}{139--141}
\begin{proof}
Each \(\chi\in\mathcal K\) has a representative in \(X_+\), unique by Lemma~\ref{B2.15}. Thus every retained character lies in \(P\) and takes value \(1\) on \(\epsilon_d\). The trivial character belongs to \(K_H\) because its primitive representative belongs to \(\mathcal K\). By \ref{B2.16} and \(\mathrm{EAtH}.o_{\rm eq}\), \(|P|=4\phi.o=4\Lambda.o\).

Put \(N_\gamma=|TT\cap\gamma|\) for \(\gamma\in A/A_0\), and \(h=x^{-1+\varepsilon}\).
The share condition gives
\[
 |N_\gamma-\lambda_\gamma n_T|\le hn_T,\qquad
 \lambda_0=\frac4{3o+1},\quad
 \lambda_\gamma=\frac3{3o+1}\quad(\gamma\ne0).
\]
For \(\alpha\in P\) nontrivial on \(A\), put \(\bar\alpha(t+A_0)=\alpha(0,t)\). Membership in \(P\) makes this well defined, and \(\bar\alpha\) is nontrivial. Orthogonality yields
\[
 \sum_\gamma\bar\alpha(\gamma)=0,\qquad
 \sum_\gamma\bar\alpha(\gamma)\lambda_\gamma
 =\frac{3\sum_\gamma\bar\alpha(\gamma)+1}{3o+1}=c.
\]
Since \([A:A_0]=\phi.o=\Lambda.o=o\),
\[
 \left|\sum_{t\in TT}\alpha(0,t)-cn_T\right|
 =\left|\sum_\gamma\bar\alpha(\gamma)(N_\gamma-\lambda_\gamma n_T)\right|
 \le ohn_T\le4ohn_T.
\]
For retained \(\alpha\ne1\), Lemma~\ref{B2.15} identifies \(|\ell_\alpha|\) with a summand of \(K_o\).

For \(\alpha\in P\setminus K_H\), the relation \(\mathrm{iod}^2=1\) gives \(\theta_\alpha(\mathrm{iod})=\pm1\). The plus sign puts \(\theta_\alpha\) in \(X_+\setminus\mathrm{Kept}_H(\mathcal K)\), and the minus sign puts it in \(X_-\).
Thus \(\theta_\alpha\in\mathrm{Bounded}\), and its logarithmic derivative bound gives \(|\ell_\alpha|\le C\log x\).
\end{proof}

\begin{theorem}[The retained sum is the target sum]\label{B2.18}
For any level \(\Lambda\), assume \(\mathrm{EAtH}(\Lambda,H,TT,x,\varepsilon,C)\) for real \(x,\varepsilon,C\), with compatible \(\ell\).
With \(K_p=K_H\) and \(c=(3o+1)^{-1}\), for every \(y\in\R\),
\[
 G_m^A(y)=\Lambda.G_m(y)
 =G_z(y)+\sum_{\chi\in\mathcal K\setminus\{1\}}
 w_\Lambda(\chi)\Re g_\chi(y).
\]
Here \(w_\Lambda(\chi)=1\) for the designated \(\chi_{\rm sp}\) when present, and \(w_\Lambda(\chi)=c\) otherwise.
For the cyclic level of \(\psi\), this gives
\[
 G_m(y)=G_z(y)+\frac1{3o+1}\sum_{\chi\in\mathcal X\setminus\{1\}}\Re g_\chi(y).
\]
\end{theorem}
\leanref{UpperTheory.\allowbreak{}Comparison₂.\allowbreak{}gmA\_\allowbreak{}eq}{Development.\allowbreak{}AutomorphicArithmetic.\allowbreak{}PrimitiveCosets.\allowbreak{}GmAEq}{61--149}
\leanref{Formalism.HC.Level.ofSub}{Formalism.SharpRestriction.RamificationFrobenius.TotientGaloisClasses.\allowbreak{}ExplicitFormulaErrorOfHostConditions}{105--114}
\leanref{Formalism.HC.Level.cyc}{Formalism.SharpRestriction.RamificationFrobenius.TotientGaloisClasses.\allowbreak{}ExplicitFormulaErrorOfHostConditions}{121--121}
\begin{proof}
By \ref{B2.15}, \ref{B2.02}, and \(\mathrm{EAtH}\), the map \(\alpha\mapsto\chi_\alpha\) gives a bijection \(K_H\simeq\mathcal K\).

The character \(\chi_q\) of Lemma~\ref{B2.02} corresponds to \(\phi.\mathrm{psC}\).
For \(\alpha\in K_H\), the representative condition in \(\mathrm{EAtH}\) gives \(\alpha(\epsilon_d)=1\), so
\[
 \alpha|_{\{0\}\times A}=1\quad\Longleftrightarrow\quad
 \alpha=1\ \text{or}\ \alpha=\chi_q.
\]
A character trivial on \(A\) factors through \(E_0\), and its first sign is \(+1\).
The trivial character contributes \(G_z(y)\) with weight \(1\).
When a character \(\chi_{\rm sp}\) is designated, the exceptional-character clause of \(\mathrm{EAtH}\) gives \(H.\operatorname{prim}(\phi.\mathrm{psC})=\chi_{\rm sp}\). Its retained term has weight \(1\). In the absence of a designated character, the same clause excludes \(\chi_q\) from \(K_H\).
Injectivity assigns every remaining term to a character of target weight \(c\), which equals \(w_c(\alpha)\).
Use \ref{B2.15} and reindex the sum by \(K_H\simeq\mathcal K\). By definition the cyclic level has \(\mathcal K=\mathcal X\) and all its nontrivial terms have weight \(c\).
\end{proof}

\section{The remaining characters and the final bound}

\begin{theorem}[The sum over the remaining characters]\label{B2.19}
Under the hypotheses of \ref{B2.17},
\[
 \Sigma_A(x)\le4o\,x^{1/2+\varepsilon},\qquad o=\Lambda.o,\quad K_p=K_H.
\]
\end{theorem}
\leanref{UpperTheory.\allowbreak{}Comparison₂.\allowbreak{}sigA\_\allowbreak{}le}{Development.\allowbreak{}AutomorphicArithmetic.\allowbreak{}PrimitiveCosets.\allowbreak{}SigALe}{38--60}
\begin{proof}
For \(\alpha\in P\setminus K_H\), the sign argument in \ref{B2.17} puts \(\theta_\alpha\) in \(\mathrm{Bounded}\).
The clause \(\mathrm{EAtH.bdG}\) and Lemma~\ref{B2.15} give
\(|G_\alpha(x)|=|g_{\chi_\alpha}(x)|\le x^{1/2+\varepsilon}\).
By \ref{B2.16} and \(\mathrm{EAtH}.o_{\rm eq}\), \(|P|=4o\). Hence
\[
 \Sigma_A(x)\le|P\setminus K_H|x^{1/2+\varepsilon}
 \le4o\,x^{1/2+\varepsilon}.
\]
\end{proof}

The condition \(\mathrm{CoreAtH}(H,TT,e_D,e_L,K,c_{HH},e_K,\varepsilon,x)\) asserts the existence of a chain datum and compatible \(\ell\), with parameters \(x,\varepsilon\) and
\(e_{\rm Div}\le e_D\), \(e_{\rm Loss}\le e_L\), \(c_{HH}(\mathcal C)\le c_{HH}\).
The witness supplies the identifications and the chain and window hypotheses of \ref{B2.13}, including the bounds on \(B_0,N_c\). Choose \(x_1\) so that the two large-\(x\) inequalities hold for every \(x\ge x_1\). The witness satisfies the weight and cost hypotheses of \ref{B2.14}. It also satisfies the thin-value hypotheses. Its translation map covers \(A_0\), and \(T_2\subseteq A_0\). In \eqref{eq:p2c2-sizes}, the bounds \(m\le x^{1/2+\varepsilon/2}\) and \(xm\le k\) come from \Lean{ErrWindows}, and the other five from \Lean{RestSize}.
\leanref{UpperTheory.\allowbreak{}Comparison₂.\allowbreak{}CoreAtH}{Development.\allowbreak{}AutomorphicArithmetic.\allowbreak{}PrimitiveCosets.\allowbreak{}CoreAtH}{29--41}

\begin{theorem}[The target prime-sum bound]\label{B2.20}
For every level \(\Lambda\), fixed \(e_D,e_L,K,c_{HH},e_K,C\in\R\), and \(\varepsilon>0\), there exists \(x_1\in\R\) such that for all \(x\ge x_1\), every frame \(\phi\), every host \(H\) satisfying \ref{B2.01}, and every \(TT\subseteq\phi.A\),
\[
 \mathrm{EAtH}(\Lambda,H,TT,x,\varepsilon,C),\qquad
 \mathrm{CoreAtH}(H,TT,e_D,e_L,K,c_{HH},e_K,\varepsilon,x)
\]
imply \(\Lambda.G_m(x)\le x^{1/2+(e_D+e_L+7)\varepsilon}\).
\end{theorem}
\leanref{UpperTheory.\allowbreak{}Comparison₂.\allowbreak{}core\_\allowbreak{}gm}{Development.\allowbreak{}AutomorphicArithmetic.\allowbreak{}PrimitiveCosets.\allowbreak{}CoreGm}{12--74}
\begin{proof}
Put \(o=\Lambda.o\). Take the threshold from \ref{B2.14} with
\(C'=\max(e_D+e_L,0)+3\), \(c=(3o+1)^{-1}\), \(M'=4o\), \(C_E=C\), and \(K_o=\sum_{\chi\in\mathcal K}|\kappa_\chi|\).
Here \(C'\ge0\) and \(|c|\le1\).
Choose \(x_1\ge3\) large enough also to ensure
\[
 \log x\le x^\varepsilon,\qquad
 200(1+K)^2(1+c_{HH})(1+e_K)\le x^\varepsilon,\qquad
 1+4o\le x^\varepsilon.
\]
Take the chain datum and compatible \(\ell\) supplied by \(\mathrm{CoreAtH}\).
The window conditions give \(e_{\rm Div},e_{\rm Loss},c_{HH}(\mathcal C),e_K\ge0\).
Thus \(e_D,e_L\ge0\) and \(C'=e_D+e_L+3\). The inequality \(c_{HH}(\mathcal C)\le c_{HH}\) gives the large-\(x\) condition with \(c_{HH}(\mathcal C)\) in place of \(c_{HH}\). By \ref{B2.13} and \(\mathcal B\ge0\), for every \(v\in U\) we have
\[
 \operatorname{gap}(\operatorname{tr}(v))
 \le x^{(e_{\rm Div}+e_{\rm Loss}+3)\varepsilon}\mathcal B
 \le x^{C'\varepsilon}\mathcal B .
\]
For each \(u\in A_0\), choose \(v\in U\) with \(\operatorname{tr}(v)=u\). The same \(a:A\to\R\) from \(\mathrm{CoreAtH}\) occurs in every bound, so for fixed \(t\) and \(\alpha\) the factor \(G_\alpha(a(t)/b)\) is independent of \(u\) in \ref{B2.08}.
Lemma~\ref{B2.17} supplies \Lean{CharsA}. The remaining assumptions of \ref{B2.14} follow from \(\mathrm{CoreAtH}\), with the finite set \(S\) and local conditions furnished by \ref{B2.03}. Consequently
\(G_m^A(x)\le x^{1/2+(e_D+e_L+6)\varepsilon}+\Sigma_A(x)\).
Apply \ref{B2.18} and \ref{B2.19}. Since \(x\ge3\) and \(e_D+e_L\ge0\),
\[
 \begin{aligned}
 \Lambda.G_m(x)
 &\le x^{1/2+(e_D+e_L+6)\varepsilon}+4o\,x^{1/2+\varepsilon}\\
 &\le(1+4o)x^{1/2+(e_D+e_L+6)\varepsilon}
 \le x^{1/2+(e_D+e_L+7)\varepsilon}.
 \end{aligned}
\]
\end{proof}


% ---- p2c3 ----
\chapter{Choosing the auxiliary primes \(q\) and \(d\)}
\label{ch:p2c3}

Fix a positive integer \(M\), \(0<\varepsilon<1/52\), and \(A\geq1\).
For every sufficiently large \(Y\), we choose primes \(q,d\) such that the bounds of Definition~\ref{B2.42} hold throughout \([Y,Y^A]\) for every character modulo \(M\) and every \((i,j)\in\{0,1\}^2\setminus\{(0,0)\}\).
We first choose \(q\equiv3\pmod8\).
For this \(q\), we choose \(d\equiv7\pmod{16}\) subject to congruences involving \(q\) and all small odd primes.
At each stage we discard candidates with a large finite character sum, and candidates whose twisted \(L\)-function has a zero in a specified rectangle.
A mean-square estimate controls the first set, and a zero-density estimate controls the second.

\section{Candidate primes and the analytic inputs}

\begin{definition}[Candidate sets and finite twists]
\label{B2.21}
For \(U>1\), \(m_0\in\N_{>0}\), and \(a_0\in\N\), put
\[
 P(m_0,a_0,U)
 =
 \{p\text{ prime}:\sqrt U<p\leq U,\ p\equiv a_0\pmod{m_0}\}.
\]
The candidate estimates apply at the positive scales used in the selections.
The archive also defines \(P\) for real \(U\), with \(P=\varnothing\) for \(U\leq0\).

For \(y,Z\in\R\), \(q\in\N\), and \(\kappa:\N\to\C\), define
\(c_y(n)=\Lambda(n)\min\{1,y/n\}\) for \(n\geq1\), and
\[
 S(Z,y,q,\kappa)
 =
 \sum_{1\leq n\leq\lfloor Z\rfloor_+}
 c_y(n)\kappa(n)\left(\frac nq\right).
\]
Here \(\lfloor Z\rfloor_+=\max\{0,\lfloor Z\rfloor\}\). All real upper limits in integer sums have this meaning.
For general \(q\), the symbol is the product of the Legendre symbols over its prime factors, counted with multiplicity, with value \(1\) at \(q=0\).
The factor at \(2\) is \(0\) for even \(n\) and \(1\) for odd \(n\).
The selections use odd prime denominators.
% Source: [Mathlib's Jacobi symbol](https://leanprover-community.github.io/mathlib4_docs/Mathlib/NumberTheory/LegendreSymbol/JacobiSymbol.html).

For \(\alpha>0\), \(Y>1\), and \(U>1\), call
\((m_0,a_0,U)\) \emph{admitted at \((\alpha,Y)\)} if
\[
 \begin{gathered}
 m_0>0,\qquad 8\mid m_0,\qquad (a_0,m_0)=1,\\
 \frac14(\log Y)^2\log m_0\leq\log U,\qquad
 Y^{\alpha/6}\leq\log U
 \leq Y^{\alpha/5}(\log Y)^2.
 \end{gathered}
\]
For a finite set \(P\subset\N\) and any set \(I\subseteq\R\), write
\(\mathrm{MeanSq}(C,Z,P,I)\) for the assertion that
\[
 \sum_{q\in P}|S(Z,y,q,\kappa)|^2
 \leq Cy\log y\,|P|
\]
for every \(y\in I\) and every \(\kappa:\N\to\C\) with
\(|\kappa(n)|\leq1\) for all \(n\).
For arbitrary \(I\), the archive uses \(\log 0=0\) and
\(\log y=\log|y|\) for \(y<0\).

For \(Q>0\), let
\(\mathcal R(Q)=\{s\in\C:\Re s>399/400,\ |\Im s|\leq Q\}\).
For a Dirichlet character \(\chi\) of positive modulus \(N\), write
\(\mathrm{GoodC}(Q,\chi)\) if
\[
 \chi\ne1,\qquad N\leq Q,\qquad
 L(s,\chi)\ne0\quad(s\in\mathcal R(Q)).
\]
For such a nonprincipal character, use the centred sum
\[
 g_\chi(y)
 =
 \sum_{1\leq n\leq y}
 \Lambda(n)\chi(n)\left(\frac yn-1\right)
 +y\frac{L'}{L}(1,\chi).
\]
The predicate \(\mathrm{Box}(Q,C,\tau,Y,X,\chi)\) requires
\(\mathrm{GoodC}(Q,\chi)\) and \(\left|(L'/L)(1,\chi)\right|\leq C\log Y\).
It also requires \(|g_\chi(y)|\leq y^{1/2+\tau}\) for every \(y\in[Y,X]\).

A primitive representative of a function \(f:\N\to\C\) at a positive level \(L\) means a primitive character \(\chi\), of some positive conductor, such that
\(\chi(n)=f(n)\) whenever \((n,L)=1\).
The assertion \(\mathrm{BoxAll}(Q,C,\tau,Y,X,f,L)\) requires
\(\mathrm{Box}(Q,C,\tau,Y,X,\chi)\) for every such representative.
\end{definition}
\leanref{QH.\allowbreak{}S}{Development.\allowbreak{}AutomorphicArithmetic.\allowbreak{}MeanSquareBounds.\allowbreak{}PrimeSetTwoBounds}{13--15}
\leanref{QH.\allowbreak{}admitted}{Development.\allowbreak{}AutomorphicArithmetic.\allowbreak{}MeanSquareBounds.\allowbreak{}PrimeSetTwoBounds}{18--25}
\leanref{QH.\allowbreak{}MeanSq}{Development.\allowbreak{}AutomorphicArithmetic.\allowbreak{}MeanSquareBounds.\allowbreak{}PrimeSetTwoBounds}{26--28}
\leanref{QH.\allowbreak{}BoxAll}{Development.\allowbreak{}AutomorphicArithmetic.\allowbreak{}MeanSquareBounds.\allowbreak{}PrimeSetTwoBounds}{38--40}
\leanref{Formalism.\allowbreak{}HermitianLattices.\allowbreak{}ZeroFreeRegionGL2.\allowbreak{}PPset}{Formalism.\allowbreak{}GL3MeanSquares.\allowbreak{}SecondEstimateOfIntervalHypothesis}{50--52}

Set \(z=Y^{\varepsilon/5}\) and \(z'_0=z(\log Y)^2\).
When a window exponent \(K\geq1\) has been fixed, the finite sums will have upper limit \(Z=(Y^K)^{400}\).

\begin{theorem}[Deuring--Heilbronn repulsion]
\label{B2.22}
There exist absolute constants \(c_0,c_1,c_2>0\) and an integer
\(B\geq0\) with the following property.
Let \(M_0\geq3\) be an integer, let \(T\geq2\), and put
\(\mathcal L=\log(M_0T)\).
Let \(\chi_1,\chi\) be characters modulo \(M_0\).
Suppose
\[
 \chi_1^2=1,\quad \chi_1\ne1,\quad
 L(\beta_1,\chi_1)=0,\quad
 1-\frac{c_1}{\mathcal L}<\beta_1<1.
\]
If \(L(\rho,\chi)=0\), with the pole pair \((1,1)\) excluded, and
\[
 \Re\rho\geq1-c_0,\qquad
 |\Im\rho|\leq T,\qquad
 (\chi,\rho)\ne(\chi_1,\beta_1),
\]
then
\[
 1-\Re\rho
 \geq
 \frac{c_2}{\mathcal L}
 \log\!\left(
 \frac{c_1}{(1-\beta_1)\mathcal L^B}
 \right).
\]
\end{theorem}
\leanref{Formalism.\allowbreak{}HermitianLattices.\allowbreak{}Ext.\allowbreak{}PrimesInProgressions.\allowbreak{}DeuringHeilbronnMB}{Formalism.\allowbreak{}GL3MeanSquares.\allowbreak{}GenusMeanAsymptotics.\allowbreak{}JacobiCharacterSums.\allowbreak{}PNTDHBSt}{11--18}
\leanref{Formalism.\allowbreak{}HermitianLattices.\allowbreak{}Ext.\allowbreak{}PrimesInProgressions.\allowbreak{}deuringHeilbronnMB\_\allowbreak{}proof'}{Formalism.\allowbreak{}GL3MeanSquares.\allowbreak{}GenusMeanAsymptotics.\allowbreak{}JacobiCharacterSums.\allowbreak{}PNTFinal}{20--21}

Corollary~1.1 of the revised paper of Benli, Goel, Twiss, and Zaman \cite{benli-2026} gives explicit constants for \(M_0>400{,}000\), \(T\geq4\), and \(\Re\rho>1/2\).
Take \(T'=\max\{T,4\}\) and \(c_0<1/2\).
Choose \(c_1\leq1/10\). Then
\(1-\beta_1<c_1/\mathcal L\leq1/(10\log M_0)\), as required by the corollary.
In this interval \(\beta_1\) is a simple zero of \(\prod_{\chi\bmod M_0}L(s,\chi)\) by the zero-free region in \cite[(1.2)]{benli-2026}.
All factors are holomorphic at \(\beta_1\), so \(\chi_1\) is the unique character whose \(L\)-function vanishes there.
Nonvanishing at \(1\) for nonprincipal characters and the exclusion of the pole pair give \(\rho\ne1\).
Thus every pair \((\chi,\rho)\) in the statement gives a zero \(\rho\ne\beta_1\) of the product.
Its denominator \(10\log M_0+\log T'+107\) is comparable to \(\mathcal L\), so the corollary gives the bound with \(B=1\) after decreasing \(c_1\) and \(c_2\).
For the finitely many smaller moduli, nonvanishing at \(1\) for nonprincipal real characters lets us decrease \(c_1\).
The archive allows a general nonnegative integer exponent \(B\).
% Source: [Benli--Goel--Twiss--Zaman, explicit Deuring--Heilbronn repulsion](https://arxiv.org/html/2410.06082v3).

\begin{theorem}[Relative prime number theorem in progressions]
\label{B2.23}
There are absolute real constants \(c,A_0,C_0\), with \(c>0\), such that, for every positive integer \(m_0\) divisible by \(8\), one can choose either no exceptional pair or a pair \((\chi_1,\beta_1)\), where \(\chi_1\) is a character modulo \(m_0\), satisfying
\[
 \chi_1^2=1,\qquad \chi_1\ne1,\qquad
 L(\beta_1,\chi_1)=0,\qquad
 1-\frac1{8\log m_0}<\beta_1<1.
\]
The choice of no exceptional pair means that no pair with these properties exists.

Put
\[
 \operatorname{li}_2(t)=\int_2^t\frac{du}{\log u}.
\]
Let \(\pi(t;m_0,a_0)\) count primes \(p\leq t\) with
\(p\equiv a_0\pmod{m_0}\).
For \((a_0,m_0)=1\), define
\[
 H_{a_0}(t)=
 \begin{cases}
 \operatorname{li}_2(t)/\varphi(m_0),
       &\text{if there is no exceptional pair},\\[2mm]
 \bigl(\operatorname{li}_2(t)
       -\chi_1(a_0)\operatorname{li}_2(t^{\beta_1})\bigr)/
       \varphi(m_0),
       &\text{otherwise}.
 \end{cases}
\]
Then, whenever \(t\geq2\) and
\(\log t\geq C_0\log m_0\log(\log m_0+2)\), one has
\[
 \left|\pi(t;m_0,a_0)-H_{a_0}(t)\right|
 \leq
 A_0\left(
 e^{-c\log t/\log m_0}
 +e^{-\sqrt{c\log t}}
 \right)H_{a_0}(t).
\]
\end{theorem}
\leanref{Formalism.\allowbreak{}HermitianLattices.\allowbreak{}Ext.\allowbreak{}LogarithmicIntegralBounds.\allowbreak{}PNTAPRelW}{Formalism.\allowbreak{}GL3MeanSquares.\allowbreak{}GenusMeanAsymptotics.\allowbreak{}JacobiCharacterSums.\allowbreak{}MeanSquareHypWIff}{38--47}
\leanref{Formalism.\allowbreak{}HermitianLattices.\allowbreak{}Ext.\allowbreak{}PrimesInProgressions.\allowbreak{}pnt\_\allowbreak{}progressions\_\allowbreak{}relative\_\allowbreak{}weak\_\allowbreak{}holds}{Formalism.\allowbreak{}GL3MeanSquares.\allowbreak{}GenusMeanAsymptotics.\allowbreak{}JacobiCharacterSums.\allowbreak{}PNTFinal}{22--25}

Take \(L=\Q(\zeta_{m_0})\) and \(H=G=\operatorname{Gal}(L/\Q)\) in Thorner--Zaman \cite[Theorem~1.4]{thorner-2019}.
The fixed field is \(\Q\), and the Frobenius class specifies \(a_0\bmod m_0\).
The maximum conductor \(Q\) in their equation~(1-11) equals \(m_0\).
Since \(8\mid m_0\), each prime-power factor of \(m_0\) admits a primitive character.
Their product has conductor \(m_0\).
The exceptional interval is therefore \(1-1/(8\log m_0)<\beta_1<1\).
The cyclotomic estimate has the displayed relative error for \(t\geq m_0^{c_2}\), where \(c_2\) is the absolute constant of that theorem.
The stated estimate uses a larger lower threshold.
% Source: [Thorner--Zaman, Theorem 1.4](https://msp.org/ant/2019/13-5/ant-v13-n5-p02-p.pdf#page=6).

The archive uses Theorem~\ref{B2.22} and a zero-density estimate to bound the sum over zeros in the explicit formula for the weighted prime count.
It combines the estimates from two ranges and applies partial summation to obtain the stated bound for \(\pi(t;m_0,a_0)\).

\begin{theorem}[Exceptional characters across moduli]
\label{B2.24}
There is an absolute \(c_{\mathrm P}>0\) such that the following holds.
Let \(D\geq3\), and let \(\chi_i\) be nonprincipal real characters modulo positive integers \(m_i\leq D\), for \(i=1,2\).
If
\(L(\beta_i,\chi_i)=0\) and
\(1-c_{\mathrm P}/\log D<\beta_i<1\), then
\[
 \chi_1(n)=\chi_2(n)\qquad((n,m_1m_2)=1).
\]
Moreover, for every \(\nu>0\), there is \(c_\nu>0\) such that every real zero \(\beta\) of a nonprincipal real character modulo a positive integer \(m_0\) satisfies
\(1-\beta\geq c_\nu m_0^{-\nu}\).
\end{theorem}
\leanref{Formalism.\allowbreak{}HermitianLattices.\allowbreak{}Ext.\allowbreak{}LogarithmicIntegralBounds.\allowbreak{}page\_\allowbreak{}across}{Formalism.\allowbreak{}GL3MeanSquares.\allowbreak{}GenusMeanAsymptotics.\allowbreak{}JacobiCharacterSums.\allowbreak{}MSqPage}{37--66}
\leanref{Formalism.\allowbreak{}ZeroFreeDirichlet.\allowbreak{}excBound}{Formalism.\allowbreak{}LandauZeroFreeBounds.\allowbreak{}SiegelWalfiszAPOfSiegelLower}{32--33}

The first assertion is Page's exceptional-character theorem
\cite{page-1935}, expressed as equality of induced character values.
Equality on integers coprime to \(m_1m_2\) gives equality at that common level, hence equality of the primitive representatives.
Conversely, a common primitive representative gives the displayed equality.
The archive obtains the first assertion by applying \Lean{page\_\allowbreak{}across} to
\Lean{Formalism.\allowbreak{}ZeroFreeDirichlet.\allowbreak{}zFRDirichlet}.
\leanref{Formalism.\allowbreak{}ZeroFreeDirichlet.\allowbreak{}zFRDirichlet}{Formalism.\allowbreak{}LandauZeroFreeBounds.\allowbreak{}NoZeroComplex}{32--34}

Davenport \cite[Chapter~21]{davenport-2000} gives Siegel's theorem.
Thorner--Zaman \cite[Remark~1.2, (1.2)]{thorner-2024} state the bound for the exceptional zero with an effectively computable \(c_\nu\) for \(\nu\geq1/2\).
Zeros outside the exceptional interval satisfy the same bound after decreasing \(c_\nu\), since
\(1/\log(3m_0)\gg_\nu m_0^{-\nu}\).
% Sources: [Page, publication record](https://academic.oup.com/plms/article-abstract/s2-39/1/116/1475562); [Thorner--Zaman, equation (1.2)](https://arxiv.org/pdf/2108.10878).

\begin{lemma}[Size of every candidate set]
\label{B2.25}
For every \(\alpha>0\), there is \(Y_0=Y_0(\alpha)\) such that, for \(Y\geq Y_0\), every progression admitted at \((\alpha,Y)\) satisfies
\(|P(m_0,a_0,U)|\geq U^{9/10}\).
\end{lemma}
\leanref{Formalism.\allowbreak{}HermitianLattices.\allowbreak{}Ext.\allowbreak{}LogarithmicIntegralBounds.\allowbreak{}count\_\allowbreak{}a}{Formalism.\allowbreak{}GL3MeanSquares.\allowbreak{}GenusMeanAsymptotics.\allowbreak{}JacobiCharacterSums.\allowbreak{}MSqCount}{346--398}

\begin{proof}
Apply the second assertion of Theorem~\ref{B2.24} with \(\nu=1\).
Decrease its constant to obtain \(0<c_*\leq1\) such that
\(1-\beta_1\geq c_*/m_0\) for every exceptional pair modulo \(m_0\).

For \(\log U\geq4\), the corrected main term from Theorem~\ref{B2.23} satisfies
\[
 H_{a_0}(U)\geq
 \frac{c_*U}{2m_0^2\log U}.
 \tag{1}\label{eq:aux-count-main}
\]
Integration over \([U/2,U]\) gives
\(\operatorname{li}_2(U)\geq U/(2\log U)\).
For a selected exceptional zero, \(\beta_1>1/2\) gives \(U^{\beta_1}\geq\sqrt U>2\).
This proves \eqref{eq:aux-count-main} when the pair is absent or \(\chi_1(a_0)=-1\).
If \(\chi_1(a_0)=1\), the gap
\(\operatorname{li}_2(U)-\operatorname{li}_2(U^{\beta_1})\)
is bounded below by \((U-U^{\beta_1})/\log U\).
The inequality
\(1-e^{-v}\geq\tfrac12\min\{v,1\}\) for \(v\geq0\), together with
\(1-\beta_1\geq c_*/m_0\), gives
\(U-U^{\beta_1}\geq c_*U/(2m_0)\).
Finally \(\varphi(m_0)\leq m_0\), which proves
\eqref{eq:aux-count-main}.

Write \(L=\log Y\) and \(u=\log U\).
Admissibility gives
\[
 \log m_0\leq\frac{4u}{L^2},
 \qquad
 Y^{\alpha/6}\leq u\leq Y^{\alpha/5}L^2.
 \tag{2}\label{eq:aux-count-scales}
\]
Thus
\(\log(\log m_0+2)=O_\alpha(L)\) and
\(u/\log m_0\geq L^2/4\).
For sufficiently large \(Y\), these inequalities place \(U\) in the range of Theorem~\ref{B2.23} and make its relative error at most \(1/2\).
Consequently \(\pi(U;m_0,a_0)\geq\tfrac12H_{a_0}(U)\).

Uniformly in the allowed \(u\),
\[
 \frac1u
 \log\!\left(\frac{c_*U}{2m_0^2\log U}\right)
 \geq
 1-\frac8{L^2}-\frac{\log u+O(1)}u
 \longrightarrow1.
\]
Increase \(Y_0\) so that \(H_{a_0}(U)\geq4U^{9/10}\).
Removing primes at most \(\sqrt U\) removes at most \(\sqrt U\) integers.
Hence
\[
 |P(m_0,a_0,U)|
 \geq2U^{9/10}-U^{1/2}
 \geq U^{9/10}.
\]
\end{proof}

\section{Quadratic cancellation and the mean square}

\begin{lemma}[One exceptional squarefree kernel]
\label{B2.26}
Fix \(\alpha,A,B>0\).
For all sufficiently large \(Y\), uniformly over progressions
\((m_0,a_0,U)\) admitted at \((\alpha,Y)\), there is \(n_*\in\N\) such that
\[
 \left|\sum_{q\in P(m_0,a_0,U)}
             \left(\frac qs\right)\right|
 \leq Y^{-B}|P(m_0,a_0,U)|
 \tag{3}\label{eq:aux-kernel-cancel}
\]
for every odd squarefree integer \(s\) satisfying
\(1<s\leq Y^A\), \((s,m_0)=1\), and \(s\ne n_*\).
\end{lemma}
\leanref{QH.\allowbreak{}twistK}{Development.\allowbreak{}AutomorphicArithmetic.\allowbreak{}MeanSquareBounds.\allowbreak{}twistK}{148--354}

\begin{proof}
Put \(L=\log Y\), \(u=\log U\), \(D=m_0Y^A\), and \(\ell=\log D\).
Choose the optional exceptional pair of Theorem~\ref{B2.23} for every modulus \(m_0s\).
For an allowed \(s\), denote by
\(\lambda_s(v)=(v/s)\) the quadratic character on
\((\Z/s\Z)^\times\).

Call \(s\) exceptional if the selected pair
\((\chi_s,\beta_s)\) modulo \(m_0s\) satisfies
\[
 1-\beta_s<\frac{c_{\mathrm P}}{\ell},
 \qquad
 \chi_s(v)=\lambda_s(v)
 \quad
 \bigl(v\equiv1\pmod{m_0},\ (v,s)=1\bigr).
 \tag{4}\label{eq:aux-aligned-exception}
\]
There is at most one such \(s\).
Indeed, Theorem~\ref{B2.24} makes two corresponding exceptional characters agree on integers coprime to the product of their moduli.
Their restrictions in \eqref{eq:aux-aligned-exception} would therefore make
\(\lambda_{s_1}\) and \(\lambda_{s_2}\) agree on integers congruent to \(1\) modulo \(m_0\) and coprime to \(s_1s_2\).
If \(s_1\ne s_2\), choose a prime dividing exactly one of them.
At that prime choose a quadratic nonresidue, and at every other prime dividing \(s_1s_2\) choose \(1\).
The Chinese remainder theorem, also imposing \(v\equiv1\pmod{m_0}\), gives an integer on which the two Jacobi symbols differ.
Take \(n_*\) to be that \(s\), or \(0\) if none exists.

Fix \(s\ne n_*\), and apply Theorem~\ref{B2.23} modulo \(m_0s\).
Let
\[
 \mathcal B_s=
 \{b\in(\Z/m_0s\Z)^\times:b\equiv a_0\pmod{m_0}\}.
\]
The Chinese remainder theorem identifies \(\mathcal B_s\) with
\((\Z/s\Z)^\times\).
Therefore
\[
 \sum_{b\in\mathcal B_s}\lambda_s(b)=0.
 \tag{5}\label{eq:aux-primary-cancel}
\]
For the exceptional contribution, character orthogonality on the kernel of reduction modulo \(m_0\) gives
\[
 \sum_{b\in\mathcal B_s}\lambda_s(b)\chi_s(b)=0
\]
unless the restriction of \(\chi_s\) to that kernel is \(\lambda_s\).
In this aligned case, \(s\ne n_*\) gives
\(\beta_s\leq1-c_{\mathrm P}/\ell\).
Since \(2\leq U^{\beta_s}\leq U\), the bounds
\(\operatorname{li}_2(U^{\beta_s})\leq2U^{\beta_s}\) and
\(\operatorname{li}_2(U)\geq U/(2u)\) give
\[
 \frac{|\operatorname{li}_2(U^{\beta_s})|}
      {\operatorname{li}_2(U)}
 \leq r,\qquad
 r=4u\exp(-c_{\mathrm P}u/\ell).
 \tag{6}\label{eq:aux-exception-ratio}
\]

Put
\(k=A_0(e^{-cu/\log(m_0s)}+e^{-\sqrt{cu}})\).
Write \(H_b\) for the corrected main term in class \(b\), and put
\(A_U=\operatorname{li}_2(U)/\varphi(m_0s)\).
These main terms are nonnegative.
In the aligned case, \(r\leq1/2\) gives \(H_b\geq(1-r)A_U\), while the exceptional term in each class has magnitude at most \(rA_U\).
Its signed sum is therefore bounded by
\(r(1-r)^{-1}\sum_bH_b\leq2r\sum_bH_b\).
The signed sum vanishes when there is no exceptional pair and, by orthogonality, in the unaligned case.
For \(|k|\leq1/2\), the relative estimate gives
\[
 \sum_{b\in\mathcal B_s}H_b
 \leq2\sum_{b\in\mathcal B_s}\pi(U;m_0s,b)
 \leq2(|P|+\sqrt U),
\]
where \(P=P(m_0,a_0,U)\).
The summed errors contribute at most \(|k|\sum_bH_b\), and the primary terms cancel by \eqref{eq:aux-primary-cancel}.
Primes dividing \(s\), and the passage from primes at most \(U\) to \(P\), are covered by the primes at most \(\sqrt U\).
Here \(s\leq Y^A<\sqrt U\) for large \(Y\).
It follows that
\[
 \left|\sum_{q\in P}\lambda_s(q)\right|
 \leq
 (4r+2|k|)(|P|+\sqrt U)+\sqrt U.
 \tag{7}\label{eq:aux-signed-relative}
\]

Since \(s\leq Y^A\), we have \(\log(m_0s)\leq\ell=\log m_0+AL\).
For large \(Y\), admissibility gives \(\log m_0\leq4u/L^2\) and \(AL\leq4u/L^2\), so \(u/\ell\geq L^2/8\).
Also \(\log u\leq(\alpha/5+2)L\) and \(u\geq Y^{\alpha/6}\).
These bounds place every \(m_0s\) in the range of the progression theorem and yield, after increasing \(Y_0\),
\[
 4r+2|k|\leq\tfrac14Y^{-B},\qquad
 U^{-2/5}\leq\tfrac14Y^{-B}.
\]
Lemma~\ref{B2.25} gives \(\sqrt U\leq U^{-2/5}|P|\).
Substitution in \eqref{eq:aux-signed-relative} proves
\eqref{eq:aux-kernel-cancel}.
\end{proof}

For \(N\in\N\), define
\[
 s_p(y,N)=
 \sum_{\substack{a\geq1\\p^a\leq N}}
 \Lambda(p^a)\min\{1,y/p^a\}
 \quad(p\text{ prime}),
\]
and put \(s_p(y,N)=0\) for nonprime \(p\).
Let
\[
 R_0(y,m_0,N)=
 \sum_{\substack{p^a\leq N\\
                  a\text{ even or }p\mid m_0}}
 \Lambda(p^a)\min\{1,y/p^a\}.
\]
Each prime power in this last sum is counted once.

\begin{lemma}[Expansion of the quadratic moment]
\label{B2.27}
Let \(P\) be a finite set of primes, let \(8\mid m_0\) and
\((a_0,m_0)=1\), and set \(N=\lfloor Z\rfloor_+\).
Suppose \(y,\eta\geq0\), \(|\kappa(n)|\leq1\) for all \(n\), and
\(q\equiv a_0\pmod{m_0}\), \(q>N^2\) for every \(q\in P\).
Suppose \(s_p(y,N)\leq B\) for every \(p\), and suppose that for some
\(n_*\in\N\),
\[
 \left|\sum_{q\in P}\left(\frac qs\right)\right|
 \leq\eta|P|
\]
whenever
\[
 1<s\leq N^2,\qquad
 s\text{ odd and squarefree},\qquad
 (s,m_0)=1,\qquad s\ne n_*.
\]
Then
\[
 \begin{split}
 \sum_{q\in P}|S(Z,y,q,\kappa)|^2
 \leq |P|\bigg\{&
 \sum_{\substack{p\leq N\\p\text{ prime}}}s_p(y,N)^2
 +R_0(y,m_0,N)^2\\
 &+2B^2+2BR_0(y,m_0,N)
 +\eta\left(\sum_{n=1}^Nc_y(n)\right)^2
 \bigg\}.
 \end{split}
\]
\end{lemma}
\leanref{QH.\allowbreak{}meanSq\_\allowbreak{}abs}{Development.\allowbreak{}AutomorphicArithmetic.\allowbreak{}MeanSquareBounds.\allowbreak{}MeanSqAbs}{8--59}

\begin{proof}
For a positive integer \(v\), let
\[
 r(v)=\prod_{\substack{p\mid v\\
                        v_p(v)\text{ odd}\\p\nmid m_0}}p.
\]
The product contains each prime at most once and excludes every divisor of \(m_0\), including \(2\).
Quadratic reciprocity gives, for \(1\leq n,n'\leq N\),
\[
 \left|
 \sum_{q\in P}\left(\frac nq\right)\left(\frac{n'}q\right)
 \right|
 =
 \left|\sum_{q\in P}\left(\frac q{r(nn')}\right)\right|.
 \tag{8}\label{eq:aux-reciprocity-pair}
\]
The reciprocity sign and the factors contributed by primes dividing \(m_0\) are constant on the progression.
The factor at \(2\) is constant because \(8\mid m_0\).
The condition \(q>N^2\) prevents a candidate prime from dividing \(nn'\).

Expand the square and use \(|\kappa(n)|\leq1\).
The resulting upper bound is
\[
 \sum_{n,n'\leq N}c_y(n)c_y(n')
 \left|
 \sum_{q\in P}\left(\frac nq\right)\left(\frac{n'}q\right)
 \right|.
 \tag{9}\label{eq:aux-moment-expansion}
\]
The coefficients vanish outside prime powers.
If \(n=p^a\), its kernel after removing primes dividing \(m_0\) is \(p\) when \(a\) is odd and \(p\nmid m_0\), and is \(1\) otherwise.

Pairs with \(r(nn')=1\) either have the same prime base or have both entries in the set counted by \(R_0\).
Their total coefficient is at most
\[
 \sum_p s_p(y,N)^2+R_0(y,m_0,N)^2.
\]
For a fixed kernel \(r>1\), two distinct nontrivial kernels \(p,p'\) can produce \(r\) only when \(r=pp'\).
These pairs have total weight at most \(2B^2\).
For prime \(r\), the two kernel orders are \((r,1)\) and \((1,r)\).
Together they contribute at most \(2s_r(y,N)R_0(y,m_0,N)\leq2BR_0(y,m_0,N)\).
Thus the exceptional nontrivial kernel has total coefficient at most
\(2B^2+2BR_0\).
If \(n_*\leq1\), there is no exceptional nontrivial kernel.

The total coefficient of all remaining pairs is at most
\((\sum_{n\leq N}c_y(n))^2\).
Their character sums have magnitude at most \(\eta|P|\), so their contribution to \eqref{eq:aux-moment-expansion} is at most
\(\eta|P|(\sum_{n\leq N}c_y(n))^2\).
Combining these bounds with \eqref{eq:aux-moment-expansion} proves the assertion.
\end{proof}

\begin{lemma}[Prime-power coefficient estimates]
\label{B2.28}
There are absolute positive constants \(C_1,C_2,C_3,C_4\) such that, for every \(y\geq3\) and \(N\in\N\),
\[
 s_p(y,N)\leq C_2\log y\quad(p\in\N),
 \qquad
 \sum_{\substack{p\leq N\\p\text{ prime}}}s_p(y,N)^2
 \leq C_3y\log y.
 \tag{10}\label{eq:aux-prime-power-bounds}
\]
If \(w\geq1\) and every prime \(p\mid m_0\) with \(p\leq N\) satisfies \(p\leq w\), then
\[
 R_0(y,m_0,N)\leq C_1(\sqrt y+w\log y).
 \tag{11}\label{eq:aux-R0-bound}
\]
For every \(y\geq0\) and \(Z\geq2\),
\[
 \sum_{1\leq n\leq Z}c_y(n)\leq C_4y\log Z.
 \tag{12}\label{eq:aux-total-coefficient}
\]
One may take \(C_2=3\).
\end{lemma}
\leanref{QH.\allowbreak{}RR\_\allowbreak{}le}{Development.\allowbreak{}AutomorphicArithmetic.\allowbreak{}MeanSquareBounds.\allowbreak{}RRLe}{9--100}
\leanref{QH.\allowbreak{}sPsq\_\allowbreak{}le}{Development.\allowbreak{}AutomorphicArithmetic.\allowbreak{}MeanSquareBounds.\allowbreak{}SPsqLe}{161--172}
\leanref{QH.\allowbreak{}exists\_\allowbreak{}primePowerSum\_\allowbreak{}le\_\allowbreak{}mul\_\allowbreak{}log}{Development.\allowbreak{}AutomorphicArithmetic.\allowbreak{}MeanSquareBounds.\allowbreak{}SumMinLeDiv}{71--102}
\leanref{QH.\allowbreak{}sum\_\allowbreak{}cN\_\allowbreak{}le}{Development.\allowbreak{}AutomorphicArithmetic.\allowbreak{}MeanSquareBounds.\allowbreak{}SumCNLe}{8--28}

\begin{proof}
For a prime \(p\),
\[
 s_p(y,N)\leq
 \log p\sum_{a\geq1}\min\{1,y/p^a\}.
\]
If \(p\leq y\), split the sum at \(p^a\leq y\).
The initial part contributes at most \(\log y\), and the geometric tail contributes at most \(2\log p\).
If \(p>y\), the whole sum is at most
\(2y\log p/p\leq2\log y\), since \((\log t)/t\) decreases for \(t\geq3\).
This proves the first bound.

To sum the squares, first consider \(p\leq\sqrt y\).
Their contribution is \(O(\sqrt y(\log y)^2)=O(y\log y)\).
For \(\sqrt y<p\leq y\), the bound \(s_p\leq3\log y\) and
\(\log y\leq2\log p\) give \(s_p^2\ll(\log y)\log p\).
The elementary Chebyshev bound
\(\sum_{p\leq t}\log p\leq t\log4\) handles this range.
For \(p>y\), we have \(s_p^2\leq4y^2(\log p)^2/p^2\).
Partial summation of the same Chebyshev bound gives
\[
 \sum_{p>y}\frac{(\log p)^2}{p^2}\ll\frac{\log y}{y}.
\]
This proves the second bound in \eqref{eq:aux-prime-power-bounds}.

For the even powers, write \(p^{2a}=(p^a)^2\).
Their contribution is bounded by
\[
 \sum_{v\geq1}\Lambda(v)\min\{1,y/v^2\}\ll\sqrt y.
\]
The last estimate follows by splitting at \(v=\sqrt y\) and applying
\(\sum_{v\leq t}\Lambda(v)\ll t\).
There are at most \(w\) prime bases \(p\mid m_0\) with \(p\leq N\).
Each contributes at most \(3\log y\).
This proves \eqref{eq:aux-R0-bound}.

Finally, \(c_y(n)\leq y\Lambda(n)/n\), and the elementary Mertens upper bound
\(\sum_{n\leq Z}\Lambda(n)/n\ll\log Z\) proves
\eqref{eq:aux-total-coefficient}.
\end{proof}

\begin{theorem}[Uniform mean square over candidate primes]
\label{B2.29}
There is an absolute \(C_{\mathrm m}>0\) with the following property.
Fix \(0<\alpha\leq1\) and \(K\geq1\).
For all sufficiently large \(Y\), let
\((m_0,a_0,U)\) be admitted at \((\alpha,Y)\), and assume
\[
 p\mid m_0,\ p\text{ prime}
 \quad\Longrightarrow\quad
 p\leq Y^{\alpha/5}\ \text{or}\ p>(Y^K)^{400}.
 \tag{13}\label{eq:aux-modulus-gap}
\]
Then, with \(Z=(Y^K)^{400}\),
\[
 \sum_{q\in P(m_0,a_0,U)}
 |S(Z,y,q,\kappa)|^2
 \leq C_{\mathrm m}y\log y\,|P(m_0,a_0,U)|
\]
for every \(Y\leq y\leq Y^K\) and every
\(\kappa:\N\to\C\) satisfying \(|\kappa(n)|\leq1\) for all \(n\).
\end{theorem}
\leanref{QH.\allowbreak{}meanSqK}{Development.\allowbreak{}AutomorphicArithmetic.\allowbreak{}MeanSquareBounds.\allowbreak{}meanSqK}{51--222}

\begin{proof}
Apply Lemma~\ref{B2.26} with its parameters \(A=800K\) and \(B=K+1\).
It supplies the cancellation hypothesis of Lemma~\ref{B2.27} with
\(\eta=Y^{-(K+1)}\), since \(\lfloor Z\rfloor_+^2\leq Y^{800K}\).
For large \(Y\), admissibility gives
\(\log U\geq Y^{\alpha/6}\geq1600K\log Y\).
Therefore \(\sqrt U\geq Z^2\), and every candidate \(q\) exceeds
\(\lfloor Z\rfloor_+^2\).

Put \(w=Y^{\alpha/5}\).
By \eqref{eq:aux-modulus-gap}, every prime divisor \(p\mid m_0\) with
\(p\leq\lfloor Z\rfloor_+\) is at most \(w\).
Since \(\alpha\leq1\) and \(y\geq Y\), we have \(w\leq y^{1/5}\).
For large \(Y\), the estimates of Lemma~\ref{B2.28} consequently give
\(R_0\leq5C_1\sqrt y\) and \(B=C_2\log y\).
Thus \(R_0^2+2B^2+2BR_0\ll y\log y\) with an absolute implied constant.

For the remaining term, use
\(\log Z=400K\log Y\) and \((400K\log Y)^2\leq Y\).
Then
\[
 \begin{split}
 \eta\left(\sum_{n\leq Z}c_y(n)\right)^2
 &\leq
 C_4^2Y^{-(K+1)}y^2(400K\log Y)^2\\
 &\leq C_4^2y^2Y^{-K}
 \leq C_4^2y
 \leq C_4^2y\log y.
 \end{split}
\]
Insert these estimates and
\(\sum_p s_p^2\leq C_3y\log y\) into Lemma~\ref{B2.27}.
Take \(C_{\mathrm m}\) to be the sum of the absolute constants in the resulting bounds.
\end{proof}

\section{Bounds throughout a window and zero-density exclusions}

\begin{theorem}[Candidates exceeding the bound somewhere in a window]
\label{B2.30}
Fix \(C\geq0\), \(\tau>0\), and \(\eta>0\).
There is \(Y_0=Y_0(C,\tau,\eta)\) such that, for \(Y\geq Y_0\), the following holds for every real \(X,Z\), finite sets \(B\subseteq P\subset\N\), and function
\(\kappa:\N\to\C\) with \(|\kappa(n)|\leq1\) for all \(n\).
Assume \(\mathrm{MeanSq}(C,Z,P,[Y,X])\).
If every \(q\in B\) has a point \(y\in[Y,X]\) such that
\[
 |S(Z,y,q,\kappa)|>\tfrac12y^{1/2+\tau},
\]
then \(|B|\leq\eta|P|\).
\end{theorem}
\leanref{QH.\allowbreak{}bad\_\allowbreak{}share}{Development.\allowbreak{}AutomorphicArithmetic.\allowbreak{}MeanSquareBounds.\allowbreak{}BadShare}{251--319}

\begin{proof}
For fixed \(Z,q,\kappa\), the function
\(S_q(y)=S(Z,y,q,\kappa)\) is continuous and piecewise linear.
Its right derivative is
\[
 D_q(v)=
 \sum_{\substack{1\leq n\leq Z\\n>v}}
 \frac{\Lambda(n)\kappa(n)}n\left(\frac nq\right).
\]
For \(v>0\),
\[
 vD_q(v)=
 S\bigl(Z,v,q,\kappa\,1_{\{n>v\}}\bigr).
 \tag{14}\label{eq:aux-derivative-twist}
\]
The coefficient function on the right has modulus at most \(1\).
The full quantifier over \(\kappa\) in the mean-square hypothesis therefore gives
\[
 \sum_{q\in P}|D_q(v)|^2
 \leq C\frac{\log v}{v}|P|
 \qquad(Y\leq v\leq X).
 \tag{15}\label{eq:aux-derivative-moment}
\]

Consider \(T\leq W\leq2T\) with \([T,W]\subseteq[Y,X]\).
Let \(B_T\) consist of candidates with a witness in \([T,W]\).
For every \(q\in P\), choose \(y_q\in[T,W]\) and take it to be a witness when \(q\in B_T\).
The fundamental theorem of calculus and Cauchy--Schwarz give
\[
 |S_q(y_q)|^2
 \leq2|S_q(T)|^2+
 2T\int_T^W|D_q(v)|^2\,dv.
\]
After summing and using \eqref{eq:aux-derivative-moment},
\[
 \sum_{q\in P}|S_q(y_q)|^2
 \leq4CT\log(2T)|P|.
 \tag{16}\label{eq:aux-block-maximal}
\]
For \(q\in B_T\), the chosen witness satisfies \(|S_q(y_q)|^2>T^{1+2\tau}/4\).
Hence
\[
 |B_T|
 \leq16C\log(2T)T^{-2\tau}|P|
 \leq\frac{32C}{\tau}T^{-\tau}|P|,
 \tag{17}\label{eq:aux-block-bad}
\]
where \(T\geq2\) and \(\log(2T)\leq2T^\tau/\tau\).

If \(X<Y\), the witness hypothesis gives \(B=\varnothing\).
Otherwise, for each integer \(i\geq0\) with \(T_i=2^iY\leq X\), put
\(W_i=\min\{2T_i,X\}\).
Let \(B_i\) consist of the \(q\in B\) with a witness in \([T_i,W_i]\).
These intervals cover \([Y,X]\), so \(|B|\leq\sum_i|B_i|\).
Summing \eqref{eq:aux-block-bad} gives
\[
 |B|
 \leq
 \frac{32C}{\tau}
 \frac{Y^{-\tau}}{1-2^{-\tau}}|P|.
\]
Choose \(Y\) so that the coefficient is at most \(\eta\).
\end{proof}

\begin{theorem}[Primitive characters with a zero near \(1\)]
\label{B2.31}
There is an absolute \(C_{\mathrm z}>0\) such that, for every \(Q\geq3\), every finite set \(\mathcal E\) of primitive Dirichlet characters satisfying
\[
 \operatorname{cond}(\chi)\leq Q,\qquad
 L(s,\chi)=0\text{ for some }s\in\mathcal R(Q)
\]
has \(|\mathcal E|\leq C_{\mathrm z}Q^{1/4}\).
A character is counted once, regardless of how many zeros it has in the rectangle.
\end{theorem}
\leanref{Formalism.\allowbreak{}HermitianLattices.\allowbreak{}Ext.\allowbreak{}PrimitiveCharacterCasts.\allowbreak{}ZeroDensityBox}{Formalism.\allowbreak{}GL3MeanSquares.\allowbreak{}GenusMeanAsymptotics.\allowbreak{}JacobiCharacterSums.\allowbreak{}CountSt}{7--12}
\leanref{Formalism.\allowbreak{}HermitianLattices.\allowbreak{}Ext.\allowbreak{}PrimitiveCharacterCasts.\allowbreak{}ZD.\allowbreak{}zero\_\allowbreak{}density\_\allowbreak{}box}{Formalism.\allowbreak{}GL3MeanSquares.\allowbreak{}GenusMeanAsymptotics.\allowbreak{}JacobiCharacterSums.\allowbreak{}CountZAll}{13--16}

For \(T\geq2\), let \(N(\sigma,T,\chi)\) count the nontrivial zeros
\(\rho\) with \(\Re\rho\geq\sigma\) and \(|\Im\rho|\leq T\), with multiplicity.
For \(\sigma=399/400\) and \(\delta=1-\sigma\), summing
Montgomery's Theorem~12.2 \cite{montgomery-1971} over dyadic conductor ranges gives
\[
 \sum_{\substack{\chi\ \mathrm{primitive}\\
                  \operatorname{cond}(\chi)\leq Q}}
 N(\sigma,T,\chi)
 \ll (Q^2T)^{3\delta/(1+\delta)}(\log QT)^9.
\]
At \(T=Q\), this is \(O(Q^{9/401}(\log Q)^9)\).
Every \(\chi\in\mathcal E\) of conductor greater than \(1\) contributes at least one to the zero count.
Allowing one further character of conductor \(1\) gives the archive's weaker bound \(|\mathcal E|\ll Q^{1/4}\).
% Source for the displayed form of Montgomery's theorem: [Goldston--Pintz--Yildirim, (12.23)](https://arxiv.org/pdf/0710.2728v1#page=25).
% Source: [Montgomery, Topics in Multiplicative Number Theory](https://link.springer.com/book/10.1007/BFb0060851).

\begin{lemma}[From bad characters to bad candidate primes]
\label{B2.32}
There is an absolute \(C_{\mathrm b}\geq0\) with the following property.
Fix real numbers \(Q\geq3\) and \(F\), and let \(B\) be a finite set of primes.
Suppose that every \(q\in B\) satisfies \(q>F\) and admits a primitive character \(\chi_q\) of conductor
\(N_q=f_qq\), where \(f_q\in\N_{>0}\), \(f_q\leq F\), and \(N_q\leq Q\), whose \(L\)-function has a zero in \(\mathcal R(Q)\).
Then \(|B|\leq C_{\mathrm b}Q^{1/4}\).
\end{lemma}
\leanref{QH.\allowbreak{}count\_\allowbreak{}bad\_\allowbreak{}core}{Development.\allowbreak{}AutomorphicArithmetic.\allowbreak{}MeanSquareBounds.\allowbreak{}CountBadCore}{18--60}
\leanref{QH.\allowbreak{}count\_\allowbreak{}bad}{Development.\allowbreak{}AutomorphicArithmetic.\allowbreak{}MeanSquareBounds.\allowbreak{}CountBad}{9--17}

\begin{proof}
Choose one witness \((N_q,\chi_q)\) for each \(q\).
The resulting map to primitive characters with their conductors is injective.
Indeed, equality of two witnesses gives \(f_qq=f_{q'}q'\).
The prime \(q\) divides \(f_{q'}q'\).
If \(q\mid f_{q'}\), then \(q\leq f_{q'}\leq F\), contrary to \(q>F\).
Thus \(q\mid q'\), so \(q=q'\).
Apply Theorem~\ref{B2.31} to the image.
\end{proof}

\begin{lemma}[Primitive twists and uniqueness of representatives]
\label{B2.33}
Let \(\xi\) be a primitive character of positive conductor \(f\), and let \(q\) be an odd prime with \(q\nmid f\).
Then there is a nonprincipal primitive character \(\chi\) of conductor \(fq\) such that
\[
 \chi(n)=\xi(n)\left(\frac nq\right)
 \qquad(n\in\N).
 \tag{18}\label{eq:aux-primitive-twist}
\]
Also, if primitive characters \(\chi,\chi'\) of positive conductors
\(N,N'\) satisfy \(\chi(n)=\chi'(n)\) whenever \((n,L)=1\), for some positive integer \(L\), then their conductors and characters agree.
\end{lemma}
\leanref{QH.\allowbreak{}twist\_\allowbreak{}prim}{Development.\allowbreak{}AutomorphicArithmetic.\allowbreak{}MeanSquareBounds.\allowbreak{}TwistPrim}{9--19}
\leanref{QH.\allowbreak{}ind\_\allowbreak{}unique}{Development.\allowbreak{}AutomorphicArithmetic.\allowbreak{}MeanSquareBounds.\allowbreak{}IndUnique}{7--52}

\begin{proof}
Under the Chinese remainder isomorphism
\[
 (\Z/fq\Z)^\times
 \simeq(\Z/f\Z)^\times\times(\Z/q\Z)^\times,
\]
take the product of \(\xi\) with the nontrivial quadratic character modulo \(q\).
A character of coprime product level has conductor equal to the product of the conductors of its two factors.
Both factors here are primitive, so the conductor is \(fq\).
The quadratic factor also makes the product nonprincipal.
Extending the characters by zero gives \eqref{eq:aux-primitive-twist} on all integers.

For uniqueness, induce both characters to level \(NN'L\).
They agree on every unit of that ring by the assumed equality.
Their induced characters at level \(NN'\) therefore agree as well.
A character induced from both levels \(N\) and \(N'\) factors through
\(\gcd(N,N')\).
Primitivity gives \(N\leq\gcd(N,N')\) and \(N'\leq\gcd(N,N')\).
Hence \(N=N'\), and injectivity of induction gives \(\chi=\chi'\).
\end{proof}

\section{Recovering the centred prime sum}

For a nonprincipal character \(\chi\) of positive modulus, define
\[
 D_\chi=-\frac{L'}{L}(1,\chi),\qquad
 T_\chi(V)=D_\chi-\sum_{1\leq n\leq V}
                         \frac{\Lambda(n)\chi(n)}n.
\]
Thus \(T_\chi(V)\) is the archive's \(\Lean{tailL}\).

\begin{lemma}[Tail estimate from a zero-free rectangle]
\label{B2.34}
Fix \(0<\alpha\leq1\) and \(K\geq1\).
There is \(Y_0=Y_0(\alpha,K)\) such that, for \(Y\geq Y_0\), every character \(\chi\) of positive modulus \(N\) satisfying
\[
 Y^{\alpha/6}\leq\log Q\leq Y^{\alpha/4},
 \qquad \mathrm{GoodC}(Q,\chi)
\]
obeys
\(\left|Y^KT_\chi((Y^K)^{400})\right|\leq Y^\alpha\).
\end{lemma}
\leanref{QH.\allowbreak{}tailK}{Development.\allowbreak{}AutomorphicArithmetic.\allowbreak{}MeanSquareBounds.\allowbreak{}TailKExp}{92--146}

\begin{proof}
Write
\[
 x=Y^K,\quad V=x^{400},\quad
 \theta=x^{-2},\quad h=V\theta,\quad
 \delta=\frac{\alpha}{1600K}.
\]
In particular \(0<\delta\leq1/800\).
Introduce
\[
 \mathcal S_\chi(t)
 =
 \sum_{1\leq n\leq t}
 \Lambda(n)\chi(n)\left(\frac tn-1\right).
\]

For every \(0<\delta\leq1/800\), there is \(A_\delta\geq0\) such that a character satisfying \(\mathrm{GoodC}(Q,\chi)\), with \(Q\geq2\), obeys
\[
 \begin{split}
 \bigl|\mathcal S_\chi(V+V\theta)
       -\mathcal S_\chi(V)-V\theta D_\chi\bigr|
 \leq A_\delta\log Q\bigl(
    (\log Q)\theta V^{1-1/400+\delta}
    +V^2/Q\bigr)
 \end{split}
 \tag{19}\label{eq:aux-contour-difference}
\]
for \(V\geq1\) and \(0<\theta\leq1\).

On \(\Re s=2\), the integrated Perron kernel is
\[
 \frac{-L'(s,\chi)}{L(s,\chi)}
       \frac{t^s}{s(s-1)}.
\]
Move the segment of height \(Q/2\) to
\(\sigma_0=1-1/400+\delta\).
The residue at \(s=1\) is \(tD_\chi\).
Suppose \(\sigma_0\leq\Re s<11/10\) and \(|\Im s|\leq Q-1\).
The local logarithmic-derivative formula writes \(L'/L(s,\chi)\), up to \(O(\log Q)\), as a sum of \(1/(s-\rho)\) over the zeros in the closed disk of radius \(11/60\) centred at \(21/20+i\Im s\), counted with multiplicity.
Their total multiplicity is \(O(\log Q)\).
Every such zero satisfies \(|\Im\rho|\leq Q-1+11/60<Q\), so the zero-free rectangle gives \(\Re\rho\leq399/400\).
Hence \(|s-\rho|\geq\delta\).
For \(\Re s\geq11/10\), the absolutely convergent Dirichlet series gives an absolute bound.
Thus
\[
 \left|\frac{L'}L(s,\chi)\right|
 \ll_\delta\log Q
 \quad
 \bigl(\Re s\geq\sigma_0,\ |\Im s|\leq Q-1\bigr).
\]
For the difference of the two left-edge integrals, use
\[
 \bigl|(V+V\theta)^s-V^s\bigr|
 \ll \theta|s|V^{\sigma_0}
 \qquad(\Re s=\sigma_0).
\]
After \(|s|\) cancels, use
\(\int_{-Q/2}^{Q/2}|\sigma_0-1+it|^{-1}\,dt\ll\log Q\), since
\(1-\sigma_0\geq1/800\).
The horizontal segments, and the tails of the original line at real part \(2\), contribute \(O_\delta((\log Q)V^2/Q)\).
These estimates give \eqref{eq:aux-contour-difference}.
\leanref{Formalism.\allowbreak{}HermitianLattices.\allowbreak{}Ext.\allowbreak{}BoxTail.\allowbreak{}tail\_\allowbreak{}box\_\allowbreak{}negLD}{Formalism.\allowbreak{}GL3MeanSquares.\allowbreak{}GenusMeanAsymptotics.\allowbreak{}JacobiCharacterSums.\allowbreak{}TailBox}{35--120}
\leanref{Formalism.\allowbreak{}HermitianLattices.\allowbreak{}Ext.\allowbreak{}BoxTail.\allowbreak{}tail\_\allowbreak{}contour}{Formalism.\allowbreak{}GL3MeanSquares.\allowbreak{}GenusMeanAsymptotics.\allowbreak{}JacobiCharacterSums.\allowbreak{}TailContour}{14--136}

There is also a direct finite-sum estimate
\[
 \left|
 \mathcal S_\chi(V+h)-\mathcal S_\chi(V)
 -h\sum_{n\leq V}\frac{\Lambda(n)\chi(n)}n
 \right|
 \leq
 h(h+1)\frac{\log(V+h)}V.
 \tag{20}\label{eq:aux-finite-difference}
\]
Indeed, terms with \(n\leq V\) cancel exactly.
There are at most \(h+1\) integers in \(V<n\leq V+h\).
Each corresponding term satisfies
\[
 \left|\Lambda(n)\chi(n)
       \left(\frac{V+h}{n}-1\right)\right|
 \leq h\,\frac{\log(V+h)}V.
\]
Put \(\Delta=\mathcal S_\chi(V+h)-\mathcal S_\chi(V)\) and
\(A_V=\sum_{n\leq V}\Lambda(n)\chi(n)/n\).
Since \(hT_\chi(V)=(\Delta-hA_V)-(\Delta-hD_\chi)\), the triangle inequality and
\eqref{eq:aux-contour-difference}--\eqref{eq:aux-finite-difference}, followed by division by \(h\), give
\[
 |T_\chi(V)|
 \leq
 A_\delta(\log Q)^2V^{-1/400+\delta}
 +A_\delta(\log Q)\frac{V}{Q\theta}
 +(h+1)\frac{\log(V+h)}V.
 \tag{21}\label{eq:aux-tail-general}
\]

Multiply by \(x\).
Since \(V=x^{400}\), the first term becomes
\[
 A_\delta(\log Q)^2x^{400\delta}
 \leq A_\delta Y^{\alpha/2}Y^{\alpha/4}
 =A_\delta Y^{3\alpha/4}.
\]
The second becomes \(A_\delta(\log Q)x^{403}/Q=o(1)\), uniformly because
\(Q\geq\exp(Y^{\alpha/6})\).
The third is \(O(x^{-1}\log x+x^{-399}\log x)\).
Their sum is at most \(Y^\alpha\) for sufficiently large \(Y\).
\end{proof}

\begin{lemma}[Logarithmic derivatives, including the constant \(4\)]
\label{B2.35}
There is an absolute \(C_{\mathrm L}>0\) such that, for every
\(0<\alpha\leq1\), there is \(Y_0=Y_0(\alpha)\) such that, for every \(Y\geq Y_0\) and every character \(\chi\) of positive modulus,
\[
 \left|\frac{L'}L(1,\chi)\right|\leq C_{\mathrm L}\log Y
 \tag{22}\label{eq:aux-kappa-C}
\]
whenever
\(Y^{\alpha/6}\leq\log Q\leq Y^{\alpha/4}\) and
\(\mathrm{GoodC}(Q,\chi)\).

More precisely, fix \(0<\varepsilon<1/52\) and \(A_1>0\).
There is \(Y_1=Y_1(\varepsilon,A_1)\) such that, for \(Y\geq Y_1\),
\[
 Y^{\varepsilon/6}\leq\log Q
 \leq A_1Y^{\varepsilon/5}(\log Y)^2,
 \qquad \mathrm{GoodC}(Q,\chi)
\]
imply
\[
 \left|\frac{L'}L(1,\chi)\right|\leq4\log Y.
 \tag{23}\label{eq:aux-kappa-four}
\]
\end{lemma}
\leanref{QH.\allowbreak{}kappaC}{Development.\allowbreak{}AutomorphicArithmetic.\allowbreak{}MeanSquareBounds.\allowbreak{}kappaC}{9--63}
\leanref{QH.\allowbreak{}kappa4}{Development.\allowbreak{}AutomorphicArithmetic.\allowbreak{}MeanSquareBounds.\allowbreak{}Kappa4Loglog}{126--154}

\begin{proof}
For the first assertion, apply Lemma~\ref{B2.34} with \(K=1\).
It gives \(|T_\chi(Y^{400})|\leq1\), since \(\alpha\leq1\).
Consequently
\[
 \left|\frac{L'}L(1,\chi)\right|
 \leq1+\sum_{n\leq Y^{400}}\frac{\Lambda(n)}n
 \leq C_{\mathrm L}\log Y
\]
with an absolute constant.

For the explicit coefficient, there are absolute constants
\(C_0,\ell_0\), with \(\ell_0\geq1\), such that
\[
 \left|\frac{L'}L(1,\chi)\right|
 \leq1000\log\log Q+C_0
 \qquad(\log Q\geq\ell_0,\ \mathrm{GoodC}(Q,\chi)).
 \tag{24}\label{eq:aux-kappa-loglog}
\]
Put \(\ell=\log Q\) and choose
\[
 x=\ell^{5/2},\qquad V=x^{400}=\ell^{1000},
 \qquad\theta=x^{-2},\qquad\delta=1/4000.
\]
Use \eqref{eq:aux-tail-general} with this fixed \(\delta\).
Its first two terms are bounded respectively by constant multiples of
\[
 \ell^2V^{-1/400+1/4000}=\ell^{-1/4},
 \qquad
 \ell\,\frac{Vx^2}{Q}=\frac{\ell^{1006}}{e^\ell}.
\]
The finite-difference term is bounded as well.
Thus \(|T_\chi(V)|\leq C_1\), for an absolute \(C_1\) and all sufficiently large \(\ell\).

Mertens's upper bound gives
\[
 \sum_{n\leq V}\frac{\Lambda(n)}n\leq\log V+2
 \qquad(V\geq1).
\]
It follows that
\(|D_\chi|\leq C_1+\log V+2=1000\log\ell+C_1+2\), which proves
\eqref{eq:aux-kappa-loglog}.
\leanref{Deduction.\allowbreak{}mertens\_\allowbreak{}upper}{Development.\allowbreak{}AutomorphicArithmetic.\allowbreak{}DedekindLFunctions.\allowbreak{}LandauTheorems.\allowbreak{}mertensUpper}{103--113}

Finally,
\[
 \log\log Q
 \leq
 \log A_1+\frac{\varepsilon}{5}\log Y+2\log\log Y.
\]
The coefficient of \(\log Y\) obtained from
\eqref{eq:aux-kappa-loglog} is \(200\varepsilon<50/13\).
Since \(4-50/13=2/13>0\), the terms
\(2000\log\log Y+1000\log A_1+C_0\) are absorbed by the remaining multiple of \(\log Y\).
This gives \eqref{eq:aux-kappa-four}.
\end{proof}

\begin{lemma}[From a finite sum to the analytic box]
\label{B2.36}
There is an absolute \(C>0\) such that the following holds.
Fix \(0<\alpha\leq1/4\), \(\tau>0\), and \(K\geq1\).
For all sufficiently large \(Y\), suppose \(\chi\) is a character of positive modulus satisfying
\[
 Y^{\alpha/6}\leq\log Q\leq Y^{\alpha/4},
 \qquad \mathrm{GoodC}(Q,\chi),
\]
and suppose that, for some \(q\in\N\) and \(\kappa:\N\to\C\),
\(\chi(n)=\kappa(n)(n/q)\) for every \(n\in\N\).
If
\[
 |S((Y^K)^{400},y,q,\kappa)|
 \leq\tfrac12y^{1/2+\tau}
 \qquad(Y\leq y\leq Y^K),
\]
then \(\mathrm{Box}(Q,C,\tau,Y,Y^K,\chi)\) holds.
\end{lemma}
\leanref{QH.\allowbreak{}box\_\allowbreak{}of\_\allowbreak{}S}{Development.\allowbreak{}AutomorphicArithmetic.\allowbreak{}MeanSquareBounds.\allowbreak{}BoxOfS}{28--61}

\begin{proof}
Put \(Z=(Y^K)^{400}\).
For \(0<y\leq Z\), splitting at \(n=y\) gives
\[
 g_\chi(y)
 =
 -\sum_{n\leq Z}\Lambda(n)\chi(n)\min\{1,y/n\}
 -yT_\chi(Z).
 \tag{25}\label{eq:aux-finite-splitting}
\]
Indeed, for \(n\leq y\),
\(y/n-1=y/n-\min\{1,y/n\}\), whereas for \(y<n\leq Z\), the two terms on the right cancel.

The first sum in \eqref{eq:aux-finite-splitting} is \(S(Z,y,q,\kappa)\).
Lemma~\ref{B2.34} gives
\(|yT_\chi(Z)|\leq|Y^KT_\chi(Z)|\leq Y^\alpha\).
For \(Y\geq16\), \(\alpha\leq1/4\), and \(y\geq Y\), we have
\(Y^\alpha\leq\tfrac12y^{1/2+\tau}\).
The triangle inequality proves the desired bound for \(g_\chi(y)\) at every point of the window.
The absolute constant in \eqref{eq:aux-kappa-C} gives the logarithmic derivative bound.
\end{proof}

\section{The two selections}

\begin{theorem}[One selection stage]
\label{B2.37}
There is an absolute \(C>0\) with the following property.
Fix \(0<\alpha\leq1/4\), \(\tau,\eta>0\), \(K\geq1\), and \(R\in\N\).
There is \(Y_0=Y_0(\alpha,\tau,K,\eta,R)\) such that the following holds for \(Y\geq Y_0\).
Let \((m_0,a_0,U)\) be admitted at \((\alpha,Y)\), and assume
\eqref{eq:aux-modulus-gap}.
Let \(1\leq F\leq\sqrt U\), and let \(\Xi\) be a finite set of primitive characters such that
\(|\Xi|\leq R\) and \(\operatorname{cond}(\xi)\leq F\) for \(\xi\in\Xi\).
Then \(P=P(m_0,a_0,U)\) is nonempty.

Moreover, if \(B\subseteq P\) and every \(q\in B\) has a
\(\xi\in\Xi\) and a positive integer \(L\) such that
\[
 \mathrm{BoxAll}
 \left(FU,C,\tau,Y,Y^K,
       n\mapsto\xi(n)\left(\frac nq\right),L\right)
\]
fails, then \(|B|\leq\eta|P|\).
\end{theorem}
\leanref{QH.\allowbreak{}stage}{Development.\allowbreak{}AutomorphicArithmetic.\allowbreak{}MeanSquareBounds.\allowbreak{}StagePow}{41--171}

\begin{proof}
Set \(Q=FU\) and \(Z=(Y^K)^{400}\).
The inequalities \(1\leq F\leq\sqrt U\) give
\[
 Y^{\alpha/6}\leq\log Q
 \leq\tfrac32\log U
 \leq\tfrac32Y^{\alpha/5}(\log Y)^2
 \leq Y^{\alpha/4}
\]
for sufficiently large \(Y\).
Lemma~\ref{B2.25} gives \(|P|\geq U^{9/10}>0\).

Define \(B_1\subseteq B\) to consist of primes admitting a primitive character of conductor \(fq\), with \(f\leq F\) and \(fq\leq Q\), whose \(L\)-function has a zero in \(\mathcal R(Q)\).
By Lemma~\ref{B2.32},
\(|B_1|\leq C_{\mathrm b}Q^{1/4}\leq C_{\mathrm b}U^{3/8}\).
Since \(3/8<9/10\), increasing \(Y_0\) gives
\[
 |B_1|\leq\tfrac12\eta|P|.
 \tag{26}\label{eq:aux-stage-zero-bad}
\]

For each \(\xi\in\Xi\), let \(B_2(\xi)\subseteq B\) consist of primes for which
\(|S(Z,y,q,\xi)|>\tfrac12y^{1/2+\tau}\) at some \(y\in[Y,Y^K]\).
Theorem~\ref{B2.29} supplies its mean-square hypothesis.
Apply Theorem~\ref{B2.30} with proportion \(\eta/(2(R+1))\).
It gives
\[
 |B_2(\xi)|\leq
 \frac{\eta}{2(R+1)}|P|,
 \qquad
 \left|\bigcup_{\xi\in\Xi}B_2(\xi)\right|
 \leq\tfrac12\eta|P|.
 \tag{27}\label{eq:aux-stage-sum-bad}
\]

We prove
\[
 B\subseteq B_1\cup\bigcup_{\xi\in\Xi}B_2(\xi).
 \tag{28}\label{eq:aux-stage-cover}
\]
Take \(q\in B\) outside the right-hand side, and let \(\xi,L\) witness its defining failure.
With \(f=\operatorname{cond}(\xi)\), we have \(q>\sqrt U\geq F\geq f\), so \(q\nmid f\).
Also \(q\equiv a_0\pmod{m_0}\), where \(8\mid m_0\) and \((a_0,m_0)=1\), so \(q\) is odd.
Lemma~\ref{B2.33} constructs a nonprincipal primitive
\(\chi\) of conductor \(fq\) with
\(\chi(n)=\xi(n)(n/q)\) for every \(n\).
Its conductor is at most \(FU\).
Since \(q\notin B_1\), it satisfies \(\mathrm{GoodC}(Q,\chi)\).
Since \(q\notin B_2(\xi)\), its finite sum obeys the required bound throughout the window.
Lemma~\ref{B2.36} gives its box.

Any primitive representative of the same function at level \(L\) agrees with \(\chi\) on all integers coprime to \(L\).
The uniqueness assertion of Lemma~\ref{B2.33} identifies it with \(\chi\).
This contradicts the chosen failure of \(\mathrm{BoxAll}\) and proves
\eqref{eq:aux-stage-cover}.
Combine \eqref{eq:aux-stage-zero-bad} and
\eqref{eq:aux-stage-sum-bad}.
\end{proof}

Put
\[
 P_1(z)=P(8,3,e^z),\qquad
 \mathfrak m(z,q)=q\prod_{\substack{3\leq p\leq z\\p\text{ prime}}}p.
\]
Define
\[
 \begin{split}
 P_2(z,z',q)=\{d\text{ prime}:\
 &e^{z'/2}<d\leq e^{z'},\\
 &d\equiv7\pmod{16},\
 \mathfrak m(z,q)\mid d+1\}.
 \end{split}
 \tag{29}\label{eq:aux-second-set}
\]
For a positive integer \(N_0\), let
\(\mathrm{GoodQ}(N_0,C,\tau,Y,X,z,q)\) mean
\[
 \mathrm{BoxAll}
 \left(N_0e^z,C,\tau,Y,X,
       n\mapsto\xi(n)(n/q),N_0q\right)
\]
for every character \(\xi\) modulo \(N_0\).
Let \(\mathrm{GoodD}(N_0,C,\tau,Y,X,z,z',q,d)\) mean
\[
 \mathrm{BoxAll}
 \left(N_0e^{z+z'},C,\tau,Y,X,
       n\mapsto\xi(n)(n/q)^a(n/d),N_0qd\right)
\]
for every such \(\xi\) and \(a\in\{0,1\}\).
These are the predicates \(\Lean{QH.\allowbreak{}GoodQ}\) and \(\Lean{QH.\allowbreak{}GoodD}\).

\begin{lemma}[First selection]
\label{B2.38}
There is an absolute \(C_1>0\) with the following property.
Fix \(0<\varepsilon<1/52\), \(\tau,\eta>0\), \(K\geq1\), and \(N_0\in\N_{>0}\).
There is \(Y_0=Y_0(\varepsilon,\tau,K,\eta,N_0)\) such that, for \(Y\geq Y_0\), with \(z=Y^{\varepsilon/5}\), the set \(P_1(z)\) is nonempty.
Every subset \(B\subseteq P_1(z)\) whose members fail
\(\mathrm{GoodQ}(N_0,C_1,\tau,Y,Y^K,z,q)\) satisfies
\(|B|\leq\eta|P_1(z)|\).
\end{lemma}
\leanref{QH.\allowbreak{}stage1}{Development.\allowbreak{}AutomorphicArithmetic.\allowbreak{}MeanSquareBounds.\allowbreak{}Stage1Succ}{36--96}

\begin{proof}
For each character \(\xi\) modulo \(N_0\), choose its primitive representative \(\xi^*\).
Its conductor divides \(N_0\), and
\(\xi^*(n)=\xi(n)\) whenever \((n,N_0)=1\).
The set \(\Xi\) of these representatives has at most
\(R=\#\{\text{Dirichlet characters modulo }N_0\}\) members.

Use \(\alpha=\varepsilon\), \(U=e^z\), \(m_0=8\), \(a_0=3\), and \(F=N_0\) in Theorem~\ref{B2.37}.
For large \(Y\), we have
\(\tfrac14(\log Y)^2\log8\leq z\),
\(Y^{\varepsilon/6}\leq z\leq z(\log Y)^2\), and
\(N_0\leq e^{z/2}\).
Thus the progression is admitted and the conductor bound holds.
The sole prime divisor of \(m_0=8\) is \(2\leq z\), so
\eqref{eq:aux-modulus-gap} holds.

A failure of \(\mathrm{GoodQ}\) supplies a character \(\xi\) modulo \(N_0\) and a primitive representative of its \(q\)-twist at level \(N_0q\).
On integers coprime to this level the values of \(\xi\) and \(\xi^*\) agree.
The same representative therefore witnesses a failure for the member
\(\xi^*\in\Xi\) in Theorem~\ref{B2.37}.
Apply Theorem~\ref{B2.37} to this bad set.
\end{proof}

\begin{lemma}[Second selection, uniformly in \(q\)]
\label{B2.39}
There is an absolute \(C_2>0\) with the following property.
Fix \(0<\varepsilon<1/52\), \(\tau,\eta>0\), \(K\geq1\), and \(N_0\in\N_{>0}\).
There is \(Y_0=Y_0(\varepsilon,\tau,K,\eta,N_0)\) such that the following holds for \(Y\geq Y_0\).
Put \(z=Y^{\varepsilon/5}\).
For every \(z(\log Y)^2\leq z'\leq K^2z(\log Y)^2\) and every
\(q\in P_1(z)\), the set \(P_2(z,z',q)\) is nonempty.
Every subset \(B\subseteq P_2(z,z',q)\) whose members fail
\(\mathrm{GoodD}(N_0,C_2,\tau,Y,Y^K,z,z',q,d)\) satisfies
\(|B|\leq\eta|P_2(z,z',q)|\).
\end{lemma}
\leanref{QH.\allowbreak{}stage2}{Development.\allowbreak{}AutomorphicArithmetic.\allowbreak{}MeanSquareBounds.\allowbreak{}CharOf}{61--218}

\begin{proof}
Write \(m_0=16\mathfrak m(z,q)\).
The integer \(\mathfrak m(z,q)\) is odd.
By the Chinese remainder theorem there is a reduced residue class \(a_0\) modulo \(m_0\) specified by
\(a_0\equiv7\pmod{16}\) and
\(a_0\equiv-1\pmod{\mathfrak m(z,q)}\).
Hence
\[
 P_2(z,z',q)=P(m_0,a_0,e^{z'}).
 \tag{30}\label{eq:aux-second-progression}
\]
Use the parameters
\[
 \alpha=\frac65\varepsilon,\qquad
 U=e^{z'},\qquad F=N_0e^z
 \tag{31}\label{eq:aux-second-parameters}
\]
in Theorem~\ref{B2.37}.

The elementary primorial bound gives
\[
 \prod_{\substack{3\leq p\leq z\\p\text{ prime}}}p
 \leq4^{\lfloor z\rfloor_+}.
\]
Since \(q\leq e^z\),
\[
 \log m_0\leq\log16+z+z\log4\leq4z
 \tag{32}\label{eq:aux-second-log-modulus}
\]
for large \(Y\).
It follows that
\(\tfrac14(\log Y)^2\log m_0\leq z(\log Y)^2\leq z'=\log U\).
For the other two admission bounds, observe that
\(Y^{\alpha/6}=Y^{\varepsilon/5}=z\) and
\(Y^{\alpha/5}=zY^{\varepsilon/25}\).
The lower bound \(z'\geq z(\log Y)^2\) gives
\(Y^{\alpha/6}\leq\log U\).
The upper bound on \(z'\), together with \(K^2\leq Y^{\varepsilon/25}\), gives
\(\log U\leq Y^{\alpha/5}(\log Y)^2\).
Thus the progression in \eqref{eq:aux-second-progression} is admitted.

Every prime factor of \(m_0\) other than \(q\) is at most \(z\) for large \(Y\).
The small factors are at most \(Y^{\alpha/5}\).
For \(q\), use \(q>e^{z/2}\) and \(z\geq800K\log Y\) to obtain
\[
 q>(Y^K)^{400}.
 \tag{33}\label{eq:aux-q-beyond-truncation}
\]
This proves the prime-divisor condition \eqref{eq:aux-modulus-gap}.

For large \(Y\), every \(q\in P_1(z)\) exceeds \(N_0\).
For each character \(\xi\) modulo \(N_0\) and each \(a\in\{0,1\}\), choose a primitive character \(\xi_{q,a}^*\) representing
\(n\mapsto\xi(n)(n/q)^a\) on integers coprime to \(N_0\).
For \(a=0\), use the primitive representative of \(\xi\).
For \(a=1\), Lemma~\ref{B2.33} gives the twist of that representative by \(q\).
In both cases the conductor is at most \(N_0q\leq F\).
The resulting family has at most
\(R=2\#\{\text{Dirichlet characters modulo }N_0\}\) members, uniformly in \(q\).

Finally, \(\log F=\log N_0+z\leq z'/2\) for large \(Y\), because
\(z'\geq z(\log Y)^2\).
Hence \(1\leq F\leq\sqrt U\).
The parameter \(\alpha=6\varepsilon/5\) is at most \(1/4\).

The box height in Theorem~\ref{B2.37} is \(FU=N_0e^{z+z'}\).
Agreement of \(\xi_{q,a}^*\) with \(\xi(n)(n/q)^a\) on integers coprime to \(N_0\) transfers every failure of \(\mathrm{GoodD}\) into the bad set of that theorem.
Apply Theorem~\ref{B2.37} to this bad set.
\end{proof}

\begin{theorem}[Proportions in the two selections]
\label{B2.40}
There is an absolute \(C>0\) such that, for fixed
\(0<\varepsilon<1/52\), \(\tau,\eta>0\), \(K\geq1\), and \(N_0\in\N_{>0}\), the following holds for all sufficiently large \(Y\).
Put \(z=Y^{\varepsilon/5}\).
For every real \(X\leq Y^K\) and
\(z(\log Y)^2\leq z'\leq K^2z(\log Y)^2\), there is
\(G_1\subseteq P_1(z)\) such that
\[
 |P_1(z)|>0,\qquad |G_1|\geq(1-\eta)|P_1(z)|.
\]
Every \(q\in G_1\) satisfies
\(\mathrm{GoodQ}(N_0,C,\tau,Y,X,z,q)\).
For each such \(q\), there is \(G_2(q)\subseteq P_2(z,z',q)\) such that
\[
 |P_2(z,z',q)|>0,\qquad
 |G_2(q)|\geq(1-\eta)|P_2(z,z',q)|,
\]
and every \(d\in G_2(q)\) satisfies
\(\mathrm{GoodD}(N_0,C,\tau,Y,X,z,z',q,d)\).
\end{theorem}
\leanref{QH.\allowbreak{}pair\_\allowbreak{}share}{Development.\allowbreak{}AutomorphicArithmetic.\allowbreak{}MeanSquareBounds.\allowbreak{}PairShare}{26--53}

\begin{proof}
Apply Lemmas~\ref{B2.38} and \ref{B2.39}, and take the larger of their constants.
A larger constant preserves the logarithmic derivative bound because
\(\log Y\geq0\).
Restriction from \([Y,Y^K]\) to \([Y,X]\) preserves every box.

Let \(G_1\) be the complement in \(P_1(z)\) of the candidates failing the first box.
For each \(q\in G_1\), let \(G_2(q)\) be the complement in its own candidate set \(P_2(z,z',q)\) of the candidates failing the second box.
The two bad-set estimates give the stated cardinality inequalities.
The candidate sets are nonempty by the same lemmas.
\end{proof}

\begin{theorem}[A single pair satisfying both stages]
\label{B2.41}
There is an absolute \(C>0\) such that the following holds.
Fix \(0<\varepsilon<1/52\), \(\tau>0\), \(K\geq1\), and \(N_0\in\N_{>0}\).
For every sufficiently large \(Y\), every real \(X\leq Y^K\), and every
\(z'_0\leq z'\leq K^2z'_0\), where
\(z=Y^{\varepsilon/5}\) and \(z'_0=z(\log Y)^2\), there are primes \(q,d\) with \(N_0<q<d\) satisfying
\[
 \begin{gathered}
 q\equiv3\pmod8,\qquad
 z/2<\log q\leq z,\\
 d\equiv7\pmod{16},\qquad
 z'/2<\log d\leq z',\\
 p\mid d+1
 \quad\text{for every odd prime }p
       \text{ with }p\leq z\text{ or }p=q,
 \end{gathered}
 \tag{34}\label{eq:aux-pair-rows}
\]
and both
\(\mathrm{GoodQ}(N_0,C,\tau,Y,X,z,q)\) and
\(\mathrm{GoodD}(N_0,C,\tau,Y,X,z,z',q,d)\).
\end{theorem}
\leanref{QH.\allowbreak{}pair\_\allowbreak{}exists}{Development.\allowbreak{}AutomorphicArithmetic.\allowbreak{}MeanSquareBounds.\allowbreak{}PairExists}{9--50}

\begin{proof}
Use Theorem~\ref{B2.40} with \(\eta=1/2\).
Its first good set has cardinality at least half the positive cardinality of \(P_1(z)\), so it is nonempty.
Choose \(q\) from it.
Choose \(d\) from the second good set associated to \(q\), which is nonempty by the same cardinality bound.
Membership in \(P_1(z)\) and \(P_2(z,z',q)\) gives \eqref{eq:aux-pair-rows}.
Membership in the corresponding good sets gives both box bounds.

For large \(Y\), we have \(e^{z/2}>N_0\) and \(z'\geq2z\).
Hence \(q>N_0\) and \(d>e^{z'/2}\geq e^z\geq q\).
\end{proof}

\section{The pair required by the arithmetic construction}

\begin{definition}[Window pairs and good pairs]
\label{B2.42}
In the parameter range of Theorem~\ref{B2.43}, write
\(z=Y^{\varepsilon/5}\) and \(z'_0=z(\log Y)^2\).
A \emph{window pair} \((q,d)\) consists of primes satisfying
\eqref{eq:aux-pair-rows} with \(z'=z'_0\).
The archive defines this predicate for arbitrary real \(\varepsilon,Y\), and
\(\mathrm{PairGood}\) for arbitrary real window endpoints.
For \(i,j\in\{0,1\}\), put
\[
 f_{\xi,i,j}(n)=
 \xi(n)\left(\frac nq\right)^i
        \left(\frac nd\right)^j,\qquad
 H_j=8Me^ze^{jz'_0}.
\]

The assertion \(\mathrm{PairGood}(M,\varepsilon,Y,X,q,d)\) means that
\((q,d)\) is a window pair, \((M,qd)=1\), and the following two conditions hold.

For every character \(\xi\) modulo \(M\), every
\((i,j)\in\{0,1\}^2\setminus\{(0,0)\}\), and every primitive representative
\(\chi\) of \(f_{\xi,i,j}\) at level \(2Mqd\),
\[
 |g_\chi(y)|\leq y^{1/2+\varepsilon/2},
 \qquad
 \left|\frac{L'}L(1,\chi)\right|\leq4\log y
 \quad(Y\leq y\leq X).
 \tag{35}\label{eq:aux-pair-good-bounds}
\]
For every real primitive character \(\rho\) whose conductor divides \(8M\), every
\((i,j)\in\{0,1\}^2\setminus\{(0,0)\}\), and every primitive representative
\(\chi\) of \(f_{\rho,i,j}\) at level \(2Mqd\),
\[
 L(s,\chi)\ne0
 \quad
 \bigl(\Re s>399/400,\ |\Im s|\leq H_j\bigr).
 \tag{36}\label{eq:aux-pair-good-boxes}
\]
\end{definition}
\leanref{UpperTheory.\allowbreak{}Top.\allowbreak{}WinPair}{Development.\allowbreak{}AutomorphicArithmetic.\allowbreak{}CharacterBoxes.\allowbreak{}ExistsPageN}{27--38}
\leanref{UpperTheory.\allowbreak{}Top.\allowbreak{}PairGood}{Development.\allowbreak{}AutomorphicArithmetic.\allowbreak{}CharacterBoxes.\allowbreak{}PairGood}{18--27}

\begin{theorem}[Good auxiliary primes for every sufficiently large window]
\label{B2.43}
For every positive integer \(M\), every \(0<\varepsilon<1/52\), and every
\(A\geq1\), there is \(Y_1\) such that, for every \(Y\geq Y_1\), there exist primes \(q,d\) satisfying
\(\mathrm{PairGood}(M,\varepsilon,Y,Y^A,q,d)\).
\end{theorem}
\leanref{QH.\allowbreak{}pairEx}{Development.\allowbreak{}AutomorphicArithmetic.\allowbreak{}MeanSquareBounds.\allowbreak{}pairEx}{112--170}
\leanref{UpperTheory.\allowbreak{}Top.\allowbreak{}PairEx}{Development.\allowbreak{}AutomorphicArithmetic.\allowbreak{}CharacterBoxes.\allowbreak{}PairGood}{28--31}

\begin{proof}
Apply Theorem~\ref{B2.41} with \(N_0=8M\), \(\tau=\varepsilon/2\),
\(K=A\), \(X=Y^A\), and \(z'=z'_0\).
It gives primes \(q,d\) satisfying the window conditions and \(8M<q<d\).
Thus neither prime divides \(M\), and \((M,qd)=1\).

Let \(f:\N\to\C\) agree, on integers coprime to \(8M\), with a character
\(\xi_8\) modulo \(8M\).
Let \(\chi\) be a primitive representative at level \(2Mqd\) of
\(n\mapsto f(n)(n/q)^i(n/d)^j\), where
\((i,j)\in\{0,1\}^2\setminus\{(0,0)\}\).
We claim that
\[
 \mathrm{Box}(H_j,C,\varepsilon/2,Y,Y^A,\chi)
 \tag{37}\label{eq:aux-final-box}
\]
holds, with the constant \(C\) from Theorem~\ref{B2.41}.

If \(j=1\), restrict the agreement defining \(\chi\) to integers coprime to
\(8Mqd\).
On these integers \(f=\xi_8\), so the \(\mathrm{GoodD}\) predicate applies directly.
Its height is \(8Me^{z+z'_0}=H_1\).

If \(j=0\), then \(i=1\).
Choose the primitive representative \(\xi_8^*\) of \(\xi_8\), with conductor
\(f_0\mid8M\).
Since \(q>8M\), Lemma~\ref{B2.33} gives a primitive character of conductor
\(f_0q\) whose values are \(\xi_8^*(n)(n/q)\).
It agrees with \(\chi\) on integers coprime to \(2Mqd\).
The uniqueness part of Lemma~\ref{B2.33} therefore identifies the characters.
In particular, \(\chi\) represents \(\xi_8(n)(n/q)\) at level \(8Mq\).
The \(\mathrm{GoodQ}\) predicate now gives \eqref{eq:aux-final-box} with height
\(8Me^z=H_0\).

For the first family in Definition~\ref{B2.42}, induce the given character
\(\xi\) modulo \(M\) to level \(8M\).
Induction preserves its values on integers coprime to \(8M\).
The claim supplies \eqref{eq:aux-final-box}.
Its \(\mathrm{GoodC}\) component gives \(\chi\ne1\), for which the archive's centred sum equals \(g_\chi\).
\leanref{QH.\allowbreak{}pairEx\_\allowbreak{}g}{Development.\allowbreak{}AutomorphicArithmetic.\allowbreak{}MeanSquareBounds.\allowbreak{}pairEx}{61--78}
Its box yields
\(|g_\chi(y)|\leq y^{1/2+\varepsilon/2}\) for \(Y\leq y\leq Y^A\).

For \(j=0,1\), we have \(\log H_j=\log(8M)+z+jz'_0\).
For \(Y\geq\exp(1+\log(8M))\),
\[
 Y^{\varepsilon/6}\leq\log H_j
 \leq3Y^{\varepsilon/5}(\log Y)^2.
\]
Apply \eqref{eq:aux-kappa-four} with \(A_1=3\) and
\(\mathrm{GoodC}(H_j,\chi)\), which is part of
\eqref{eq:aux-final-box}.
It gives
\(\left|L'/L(1,\chi)\right|\leq4\log Y\leq4\log y\).
This proves \eqref{eq:aux-pair-good-bounds}.

Finally, let \(\rho\) be a real primitive character of conductor dividing \(8M\).
Induce it to level \(8M\) and apply the same claim with \(f=\rho\).
The \(\mathrm{GoodC}\) component of its box gives nonvanishing on
\(\mathcal R(H_j)\), as required by \eqref{eq:aux-pair-good-boxes}.
Take \(Y_1=\max\{Y_{\rm pair},Y_4,\exp(1+\log(8M))\}\), where
\(Y_{\rm pair}\) is the threshold in Theorem~\ref{B2.41} for the parameters
used above and \(Y_4\) is the threshold in \eqref{eq:aux-kappa-four} for \(A_1=3\).
\end{proof}


% ---- p2c4 ----
\chapter{Building the arithmetic host of a window}\label{ch:arithmetic-window-host}

Fix a primitive Dirichlet character $\psi$ of conductor $M\geq1$.
Write $\Xi=\langle\psi\rangle$, $o_\psi=|\Xi|$, and
$\mathcal X=\{\xi^{\mathrm{prim}}:\xi\in\Xi\}$ for the retained family
of $\Lambda=\operatorname{Level.cyc}(\psi)$.

In the frame of \ref{B2.01}, put $n=|A|$, $n_0=|A_0|$, and
$o_\phi=[A:A_0]$, where $A_0=\langle a_0\rangle$. For a scale
$s=(x,\varepsilon,J)$ with $x>1$, $0<\varepsilon<1/52$ and real $J\geq1$,
write $b=\exp(x^\varepsilon)$, $\delta_1=x^{-\varepsilon/8}$ and
$\varepsilon_z=\varepsilon/J$.
The rank condition $\operatorname{RankVal}(s,R)$ includes
\begin{equation}\label{eq:host-rank-ranges}
 \begin{gathered}
 x^{2+8\varepsilon}\leq n\leq2x^{2+8\varepsilon},\qquad
 \sqrt{x\delta_1}\leq m\leq2\sqrt{x\delta_1},\\
 x^{9/2+10\varepsilon}\leq k\leq2x^{9/2+10\varepsilon}.
 \end{gathered}
\end{equation}
Here $f=2n$, $r=k+m-1$, and the auxiliary rank $t$ satisfies
$r+t=8fr+4+T_On_W$, with $T_O,n_W$ as in \ref{B8.59}.

\section{The field and its characters}

For a subgroup $\Xi$ of the characters modulo $M$, let
\[
 B=(\Z/M\Z)^\times/\operatorname{ann}(\Xi),\qquad
 \operatorname{ann}(\Xi)=\{u:\xi(u)=1\text{ for every }\xi\in\Xi\}.
\]
Finite character duality gives $\widehat B=\Xi$ and $|B|=|\Xi|$.
For $e\geq0$, put $H_e=(\Z/2^{e+2}\Z)^\times/\{\pm1\}$ and
$D_e=\operatorname{lcm}(2^{e+2},M)$. Let $A_e$ be the image of
$(\Z/D_e\Z)^\times$ in $H_e\times B$ under reduction to the two factors.
Choose $Q_S\mid D_e$ minimal among multiples of $2^{e+2}$ through
which the joint map factors. Let $A_0$ be the kernel of $A_e\to B$
and $a_0$ a generator of this cyclic group.

\begin{lemma}[The order of the frame and the Artin kernel]\label{B2.44}
For every $M\geq1$ and subgroup $\Xi$ of the Dirichlet characters modulo
$M$, there are integers $1\leq c\leq|\Xi|$ and $e_0$ such that
$|A_e|=c2^e$ for $e\geq e_0$. For every $e$, $[A_e:A_0]=|\Xi|$.

Take $\Xi=\langle\psi\rangle$. Suppose $e\geq2$, the primes $q,d$
satisfy
\begin{equation}\label{eq:host-prime-hyp}
 q\equiv3\pmod8,\qquad d\equiv7\pmod{16},\qquad
 \left(\frac{-d}{q}\right)=1,
\end{equation}
and $(M,qd)=1$ and $(-d/\ell)=1$ for every odd prime $\ell\mid M$.
In the resulting cyclic host, $M\mid Q_S$. For every integer $a\geq0$
coprime to $Q_Sqd$,
\begin{equation}\label{eq:cyclic-artin-kernel}
 \begin{split}
 \operatorname{Art}(a)=1\quad\Longleftrightarrow\quad&
 a\equiv1\text{ or }-1\pmod{2^{e+2}},\\
 &\psi(a)=1,\qquad
 \left(\frac aq\right)=\left(\frac ad\right)=1.
 \end{split}
\end{equation}
\end{lemma}
\leanref{UpperTheory.\allowbreak{}Top.\allowbreak{}ordAny}{Development.\allowbreak{}AutomorphicArithmetic.\allowbreak{}CharacterBoxes.\allowbreak{}OrdAny}{86--108}
\leanref{Formalism.\allowbreak{}HC.\allowbreak{}AnyGrp.\allowbreak{}o\_\allowbreak{}eq}{Formalism.\allowbreak{}SharpRestriction.\allowbreak{}RamificationFrobenius.\allowbreak{}TotientGaloisClasses.\allowbreak{}InertiaGaloisSymmetry.\allowbreak{}AnyLevel}{29--35}
\leanref{UpperTheory.\allowbreak{}Top.\allowbreak{}kerChiAny}{Development.\allowbreak{}AutomorphicArithmetic.\allowbreak{}CharacterBoxes.\allowbreak{}KerChiAny}{104--142}

\begin{proof}
The surjection $A_e\to H_e$ has kernel $\{1\}\times C_e$, where
$C_e=\{b\in B:(1,b)\in A_e\}$, so $|A_e|=2^e|C_e|$.
Reduction gives $C_{e+1}\subset C_e$, hence $|C_e|$ stabilizes to $c$.
The projection to $B$ is onto, with kernel $A_0$ embedded in $H_e$.

Let $\eta(a)\in A_e$ be the joint image, and let
$w:\{\pm1\}\to A_e$ take a sign to the joint image of that integer.
The Artin map is
\begin{equation}\label{eq:artin-with-sign}
 \operatorname{Art}(a)=
 \left(\left(\frac aq\right),\left(\frac ad\right),
       \eta(a)w\left(\left(\frac aq\right)\right)\right).
\end{equation}
Its kernel consists of units with both signs $1$ and $\eta(a)=1$.
The two projections of $\eta$ give the congruence and $\psi(a)=1$.
Since $\psi$ is primitive and factors through $Q_S$, we have $M\mid Q_S$.
\end{proof}

\begin{lemma}[Characters trivial on the averaging subgroup]\label{B2.45}
Under the hypotheses of the second part of \ref{B2.44}, every host
character $\theta$ trivial on $A_0$ has a representation
\begin{equation}\label{eq:host-twist-form}
 \theta(\operatorname{Art}(a))
   =\xi(a)\left(\frac aq\right)^i\left(\frac ad\right)^j,
 \qquad (a,Q_Sqd)=1,
\end{equation}
where $\xi\in\Xi$ and $i,j\in\{0,1\}$. Its value on the distinguished
inertia element at $d$ is $\theta(\mathrm{iod})=(-1)^j$.
These assertions constitute $\operatorname{TwistForm}(\psi,H)$.
\end{lemma}
\leanref{UpperTheory.\allowbreak{}Top.\allowbreak{}twistAny}{Development.\allowbreak{}AutomorphicArithmetic.\allowbreak{}CharacterBoxes.\allowbreak{}TwistAny}{115--169}
\leanref{UpperTheory.\allowbreak{}Top.\allowbreak{}TwistForm}{Development.\allowbreak{}AutomorphicArithmetic.\allowbreak{}CharacterBoxes.\allowbreak{}TwistForm}{14--18}

\begin{proof}
The restriction to $A$ factors through $A/A_0\simeq B$ and hence gives
$\xi\in\widehat B=\Xi$. Factor
\[
 ((s_q,s_d),z)=
 ((s_q,1),w(s_q))\,((1,s_d),1)\,((1,1),z\,w(s_q)^{-1}).
\]
The first two factors give the characters $s_q\mapsto s_q^i$ and
$s_d\mapsto s_d^j$. Substitution into \eqref{eq:artin-with-sign} gives
\eqref{eq:host-twist-form}. At $\mathrm{iod}=((1,-1),1)$ the value is
$(-1)^j$.
\end{proof}

For $B_0\in\N$, $\operatorname{FieldRows}(H,B_0)$ requires the first
Artin coordinate to be $(a/q)$ and $E$ to contain $\sqrt{-q}$ with
that action. For each prime $\ell\mid Q_S$ it requires a unit $u$
modulo $Q_Sqd$ with $u\equiv1\bmod Q_Sqd/\ell$ and
$\operatorname{Art}(u)\ne1$. It also requires $Q_S\leq n^2$,
$o_\phi^{64}\leq2^e$, $q\ne3$, $(-d/p)=1$ for odd primes $p\leq B_0$,
and $4\ell\leq B_0$ for odd primes $\ell\mid Q_S$.

\begin{theorem}[Constructing the field]\label{B2.46}
For every primitive $\psi$, there are integers $c\geq1$ and $e_0$
with the following property. Let $e\geq e_0$ and let $q,d$ be primes
satisfying \eqref{eq:host-prime-hyp}, with $q\ne3$ and $(M,qd)=1$.
Suppose $B_0\geq4M$ and $(-d/p)=1$ for every odd prime $p\leq B_0$.
There are a frame $\phi$ and a host arithmetic datum $H$ on $\phi$
with $\phi.e=e$, $n=c2^e$, $o_\phi=o_\psi$, $H.q=q$ and $H.d=d$,
and
\begin{equation}\label{eq:field-conductor-bounds}
 Q_S\mid\operatorname{lcm}(2^{e+2},M),\qquad Q_S\leq4Mn.
\end{equation}
They satisfy $\operatorname{FieldRows}(H,B_0)$, $M\mid Q_S$,
\eqref{eq:cyclic-artin-kernel}, and $\operatorname{TwistForm}(\psi,H)$.
Every member of $\mathcal X$ is represented by a host character
trivial on $A_0$ with value $1$ on $\mathrm{iod}$. The distinguished
$q$-sign character $\mathrm{psC}$ has its primitive representative
outside $\mathcal X$.
\end{theorem}
\leanref{UpperTheory.\allowbreak{}Top.\allowbreak{}fieldChi}{Development.\allowbreak{}AutomorphicArithmetic.\allowbreak{}CharacterBoxes.\allowbreak{}FieldChi}{20--60}

\begin{proof}
Apply \ref{B2.44} with $e_0\geq\max(2,4M,64o_\psi)$.
Take $E$ to be the fixed field of the kernel of \eqref{eq:artin-with-sign}
inside $\Q(\zeta_{Q_Sqd})$. For each unit residue modulo $Q_S$,
the Chinese remainder theorem realizes all four sign pairs, so the
Artin map surjects onto $\{\pm1\}^2\times A_e$.

The third coordinate projects to $[a]\in H_e$ because $w(\pm1)$ projects
to $1$. For $a\equiv1\bmod Q_Sq$, the Artin image lies in
$\{1,\mathrm{iod}\}$. For $a\equiv1\bmod Q_Sd$ the third coordinate is
$w((a/q))$ and the image squares to $1$. Restriction to $E$ identifies
the Artin map with the cyclotomic action (\Lean{art\_\allowbreak{}spec}).
The two sign characters cut out $\Q(\sqrt{-q})$ and $K=\Q(\sqrt{-d})$.
Complex conjugation maps to $((-1,-1),1)$. These give the raw host
of \ref{B2.01} and its arithmetic datum $\mathrm{HostAx.ofRaw}(H)$.

Fix a prime $\ell\mid Q_S$. If $2^{e+2}\mid Q_S/\ell$, minimality of
$Q_S$ gives a unit $1\bmod Q_S/\ell$ with nontrivial joint image.
Otherwise $\ell=2$ and $v_2(Q_S)=e+2$. The unit $1+Q_S/2$ is
$1\bmod Q_S/2$ and $1+2^{e+1}\bmod 2^{e+2}$, so its $H_e$-coordinate
is nontrivial. Lift the witness to a unit equal to $1\bmod qd$.
The lift is $1\bmod Q_Sqd/\ell$ and its Artin image is $((1,1),z)$
with $z\ne1$. The bounds follow from
$Q_S\leq2^{e+2}M\leq4Mn\leq n^2$ and
$o_\phi^{64}=o_\psi^{64}\leq2^{64o_\psi}\leq2^e$.
Every odd prime $\ell\mid Q_S$ divides $M$, so $4\ell\leq B_0$.

For $\xi\in\Xi$ let $\xi_A$ be its pullback to $A_e$. The character
$\theta_\xi((s_q,s_d),z)=\xi_A(z\,w(s_q)^{-1})$ is trivial on $A_0$
and takes value $1$ on $\mathrm{iod}$. Substitution into
\eqref{eq:artin-with-sign} gives the primitive representative
$\xi^{\mathrm{prim}}$. Choose a unit equal to $1\bmod Q_Sd$ and
nonsquare modulo $q$. The retained characters take value $1$ there
and $\mathrm{psC}$ takes value $-1$. The kernel and twist descriptions
follow from \ref{B2.44} and \ref{B2.45}.
\end{proof}

Let $X_+$ and $X_-$ denote the host characters trivial on $A_0$ with
values $1$ and $-1$ on $\mathrm{iod}$, respectively. A character is
\emph{retained} if its primitive Dirichlet representative belongs to
$\mathcal X$. Put
$\operatorname{Bounded}=(X_+\setminus\operatorname{Retained})\cup X_-$.

\begin{theorem}[A nonzero twist for every remaining character]\label{B2.47}
Let $\varepsilon>0$, $x\geq1$, and $Y\leq x$. Suppose a host $H$
satisfies $\operatorname{TwistForm}(\psi,H)$, $M\mid Q_S$, and every
prime divisor of $Q_S$ divides $2M$. Assume
$\operatorname{PairGood}(M,\varepsilon,Y,x,q,d)$.
For every $\theta\in\operatorname{Bounded}$ and every primitive
Dirichlet representative $\chi$ of $\theta$,
\[
 |g_\chi(x)|\leq x^{1/2+\varepsilon},
 \qquad \left|\frac{L'}{L}(1,\chi)\right|\leq4\log x.
\]
\end{theorem}
\leanref{UpperTheory.\allowbreak{}Top.\allowbreak{}boundedTw}{Development.\allowbreak{}AutomorphicArithmetic.\allowbreak{}CharacterBoxes.\allowbreak{}BoundedTw}{45--57}

\begin{proof}
In \eqref{eq:host-twist-form}, $(i,j)=(0,0)$ would give
$\theta\in X_+$ and $\chi=\xi^{\mathrm{prim}}\in\mathcal X$.
The assumptions on $Q_S$ identify $\chi$ with the primitive
representative of the nonzero twist at level $2Mqd$.
Apply $\operatorname{PairGood}$ at $y=x$ and use
$x^{1/2+\varepsilon/2}\leq x^{1/2+\varepsilon}$.
\end{proof}

\section{Matching the two windows}

The condition $\operatorname{VVle}(E,t,B)$ requires finite lengths
$\operatorname{shaLen}_{p,j}$ for $1\leq j<t$ in every algebraic closure
of $E$ and for all local data of \ref{B6.01}, with
\begin{equation}\label{eq:vv-budget}
 \sum_{p\in S}\log p\sum_{j=1}^{t-1}
          \operatorname{shaLen}_{p,j}\leq B
 \quad\text{for every finite set of primes }S.
\end{equation}
The predicate $\operatorname{IsChiField}(\psi,e,q,d,E)$ requires an
embedding $j:E\hookrightarrow\Q(\zeta_N)$ for some positive common
multiple $N$ of $2^{e+2},M,q,d$. The automorphism
$\zeta_N\mapsto\zeta_N^a$ fixes $j(E)$ pointwise exactly for the
unit residues in \eqref{eq:cyclic-artin-kernel}. Call this the
prescribed $\psi$-kernel.

For each $0<\varepsilon<1/52$ and page function $e$ of \ref{B6.02},
$\operatorname{VWinAt}(\psi,401)$ in \ref{B6.03} gives $J\geq1$ and $X_0$.
If $X\geq X_0$ and $\operatorname{WinPair}(\varepsilon,X^{1/J},q,d)$
holds, some $W\geq X^{1/J}$ satisfies $W^2\leq X$. For every
$u\in[W,W^2]$, a field with the prescribed kernel at $(e(u),q,d)$
satisfies $\operatorname{VVle}(E,\lfloor u^{401}\rfloor,\sqrt{\exp(u^\varepsilon)})$.

\begin{theorem}[A common analytic and cohomological window]\label{B2.48}
Let $0<\varepsilon<1/52$ and $c\geq1$. Assume
$\operatorname{VWinAt}(\psi,401)$ and the pair-existence theorem
\ref{B2.43}. There is $J\geq1$ such that, for every $y\in\R$, there
is $Z\geq y$ for which every $x\in[Z,Z^{3/2}]$ admits
$Y\in\R$ and $q,d,e\in\N$ satisfying
\begin{equation}\label{eq:window-datum}
 \begin{gathered}
 x^{1/J}\leq Y\leq x,\qquad
 \operatorname{PairGood}(M,\varepsilon,Y,x,q,d),\\
 x^{2+8\varepsilon}\leq c2^e\leq2x^{2+8\varepsilon}.
 \end{gathered}
\end{equation}
Every field with the prescribed $\psi$-kernel at $(e,q,d)$ then satisfies
\begin{equation}\label{eq:window-vv}
 \operatorname{VVle}
 \left(E,\lfloor x^{400}\rfloor,\sqrt{\exp(x^\varepsilon)}\right).
\end{equation}
Write $\operatorname{WinDatum}(\psi,\varepsilon,J,c,x,Y,q,d,e)$ for
the conjunction of \eqref{eq:window-datum} and \eqref{eq:window-vv}.
\end{theorem}
\leanref{UpperTheory.\allowbreak{}Top.\allowbreak{}winPairs}{Development.\allowbreak{}AutomorphicArithmetic.\allowbreak{}CharacterBoxes.\allowbreak{}WinPairs}{5--82}
\leanref{UpperTheory.\allowbreak{}Top.\allowbreak{}WinDatum}{Development.\allowbreak{}AutomorphicArithmetic.\allowbreak{}CharacterBoxes.\allowbreak{}TwistForm}{19--27}

\begin{proof}
Put $e(u)=\lceil\log_2(u^{2+8\varepsilon})\rceil$ for $u\geq1$ and
choose it arbitrarily elsewhere. Take $J,X_0$ from the cohomological
hypothesis and $Y_1$ from \ref{B2.43} for this $J$.
Put $\kappa=c^{1/(2+8\varepsilon)}\geq1$. Choose $X\geq X_0$ with
$Y=X^{1/J}\geq\max(1,y,Y_1,\kappa^3,\kappa^{400})$.
Take a good pair $q,d$ on $[Y,X]$ and apply the cohomological hypothesis
to obtain $W\geq Y$ with $W^2\leq X$.

Set $Z=\kappa W$. For $x\in[Z,Z^{3/2}]$, write $x=\kappa u$.
Since $W\geq\kappa^3$, we have $W\leq u\leq x\leq W^2$.
Take $e=e(u)$ and multiply the bounds on $2^e$ by $c$.
Since $x^{1/J}\leq Y\leq x\leq X$, restrict the pair to $[Y,x]$.
The bound at $u$ gives \eqref{eq:window-vv}, since
$u\geq\kappa^{400}$ gives $x^{400}\leq u^{401}$ and $u\leq x$
gives $\exp(u^\varepsilon)\leq\exp(x^\varepsilon)$.
\end{proof}

\section{Primes for the level structure}

\begin{theorem}[A polynomial bound for a progression prime]\label{B2.49}
There is an absolute constant $A\geq1$ such that, for every integer
$N\geq2$, there is a prime $p$ with $p\equiv1\pmod N$ and $p\leq N^A$.
\end{theorem}
\leanref{UpperTheory.\allowbreak{}T3.\allowbreak{}smallOne}{Development.\allowbreak{}AutomorphicArithmetic.\allowbreak{}ShortNorms.\allowbreak{}Primes}{7--8}

This is Linnik's theorem restricted to the class $1\bmod N$
\cite[\S1.2]{thorner-2024}. The relative prime-counting estimate for progressions gives this case.
\leanref{GlobalTheory.\allowbreak{}T.\allowbreak{}primesInProgressions\_\allowbreak{}relativeError}{Development.\allowbreak{}AbelianVarietyTheory.\allowbreak{}Proofs.\allowbreak{}GlobalInterpolation.\allowbreak{}PrimesInProgressionsRelativeErrorStatement}{8--10}
\leanref{GlobalTheory.\allowbreak{}Lin.\allowbreak{}PrimesRel}{Development.\allowbreak{}AbelianVarietyTheory.\allowbreak{}Definitions.\allowbreak{}GlobalTwisting.\allowbreak{}ZerosDisc}{35--44}

A coefficient family has a finite additive abelian group $G$ and parameters
$D\geq\max(4,|G|)$ and $D^{-1}\leq\kappa\leq4$. Its coefficients
$a_g(n),\lambda_g(n)$ are nonnegative for $g\in G$, $n\in\N$.
Require $a_g(0)=0$, $a_g(1)=\mathbf1_{g=0}$ and
\[
 \left|\sum_{1\leq n\leq y}a_g(n)-\kappa y\right|\leq10\sqrt y\quad(y\geq1),
 \qquad
 \left|\sum_n a_g(n)e^{-nt}-\kappa/t\right|\leq D^{10}t\quad(0<t\leq1).
\]
With $\Lambda(n)$ the von Mangoldt function, require
$\sum_g\lambda_g(n)\leq2\Lambda(n)$ and
\[
 a_g(n)\log n=\sum_{uv=n}\sum_{j\in G}\lambda_j(u)a_{g-j}(v)\qquad(n\geq1).
\]
For $\chi\in\widehat G$, put $a_\chi(n)=\sum_g\chi(g)a_g(n)$ and
denote the trivial character by $1$. For $\chi^2=1$, require the
Dirichlet convolution $a_1*a_\chi$ to be nonnegative and
$(a_1*a_\chi)(n^2)\geq a_1(n)$. These are the hypotheses of an
\Lean{LFam}. Write $\widehat\lambda_\chi(n)=\sum_g\chi(g)\lambda_g(n)$,
$\Psi_\chi(y)=\sum_{1\leq n\leq y}\widehat\lambda_\chi(n)$ and
$\ell(t)=\log(D(|t|+2))$.
\leanref{GlobalTheory.\allowbreak{}Hk.\allowbreak{}LFam}{Development.\allowbreak{}AbelianVarietyTheory.\allowbreak{}Definitions.\allowbreak{}GlobalTwisting.\allowbreak{}LamIdealOfNotPrimePow}{12--41}

\begin{lemma}[The possible exceptional zero]\label{B2.49a}
For a family as above let $F_\chi(s)$ denote the continuation to
$\Re s>1/2$ of $(s-1)^{\mathbf1_{\chi=1}}\sum_{n\geq1}a_\chi(n)n^{-s}$.
There are absolute $0<c_0\leq1/20$, $a>0$ and $0<c_9\leq1$ such that,
in every family, the zeros of all the $F_\chi$ in
$\Re s\geq1-c_0/\ell(\Im s)$ consist of at most one pair $(\chi,\beta)$.
For this pair $\chi^2=1$ and $\beta$ is a simple real zero with
$3/4\leq\beta\leq1-c_9D^{-a}$.
\end{lemma}
\begin{proof}
Write $L_\chi(s)=\sum_{n\geq1}a_\chi(n)n^{-s}$ and
$P_\chi(s)=\sum_{n\geq1}\widehat\lambda_\chi(n)n^{-s}$.
For $\Re s>1$ the convolution gives $L'_\chi=-P_\chi L_\chi$.
Integrate from $+\infty$, using $a_\chi(1)=1$ and
$|\widehat\lambda_\chi|\leq2\Lambda$, to get
$|L_\chi(\sigma+it)|\geq\zeta(\sigma)^{-2}$.
Partial summation continues $F_\chi$ to $\Re s>1/2$ and gives
$|F_\chi(s)|\leq300D(1+|s|)^2$ for $\Re s\geq11/20$.
On the disc of radius $7/10$ about $z_t=5/4+it$, the logarithm of
the ratio of the maximum modulus to $|F_\chi(z_t)|$ is $O(\ell(t))$.
Apply \cite[Lemmas~6.1 and~6.3]{montgomery-vaughan-2007} to
$z\mapsto F_\chi(z_t+7z/10)$ with $r=5/7$ and $R=6/7$.
For $|s-z_t|\leq1/2$ away from the zeros this gives
\begin{equation}\label{eq:host-local-logderivative}
 \frac{F'_\chi}{F_\chi}(s)=
 \sum_{|\rho-z_t|\leq3/5}\frac{m_\rho}{s-\rho}+O(\ell(t)),
 \qquad \sum_{|\rho-z_t|\leq3/5}m_\rho\ll\ell(t),
\end{equation}
where $m_\rho$ is the multiplicity. Since
$P_\chi(s)=\mathbf1_{\chi=1}/(s-1)-F'_\chi(s)/F_\chi(s)$ and all zeros
have real part at most $1$, each zero contributes a nonpositive real
part to $P_\chi(s)$ when $\Re s>1$.

Write a zero as $1-b+it$, with $0\leq b\leq v\leq1/400$ and $v>0$.
At $\sigma=1+6v$, the inequality $3+4\cos u+\cos 2u\geq0$ gives
$3\Re P_1(\sigma)+4\Re P_\chi(\sigma+it)+\Re P_{\chi^2}(\sigma+2it)\geq0$.
If $\chi^2\ne1$, both shifted pole terms vanish. If $|t|\geq30v$,
their real parts before weighting are at most $1/(156v)$. The
indicated zero then gives
\[
 4/(7v)\leq3/(6v)+5/(156v)+O(\ell(t)).
\]
For $\chi^2=1$ and $0<|t|<30v$, sum $P_\xi(\sigma)$ over
$\xi\in\langle\chi\rangle$ at $\sigma=1+60v$. Its coefficients are
nonnegative by orthogonality. The conjugate pair contributes
$2(60v+b)/((60v+b)^2+t^2)\geq120/(4621v)$ against one pole of size
$1/(60v)$. Both comparisons give $1/v\ll\ell(t)$, which excludes
these zeros for $v=c_0/\ell(t)$ and sufficiently small $c_0$.

If the strip contains two real zeros counted with multiplicity, sum
$P_\xi(\sigma)$ at $\sigma=1+2c_0/\ell(0)$ over the subgroup generated
by their characters. It has at most four elements and the sum has
nonnegative coefficients. By \eqref{eq:host-local-logderivative} the
zeros contribute at least $2\ell(0)/(3c_0)$ against one pole
$\ell(0)/(2c_0)$ and an error $O(\ell(0))$. This is impossible for
small $c_0$.

Put $r=\kappa|G|$ and $H_\chi(t)=\sum_n a_\chi(n)e^{-nt}$.
Partial summation gives $F_1(1)=r\geq D^{-1}$.
For nontrivial real $\chi$, the theta estimate and $a_g(n)\ll n$ give
\[
 |H_1(t)-r/t|+|H_\chi(t)|\ll D^{11}\min(t,1/t).
\]
For $t\geq1$, summing $a_g(n)\ll n$ against $e^{-nt}$ gives
$|H_\chi(t)|\ll De^{-t}$. This bound and the $O(D^{11}t)$ bound
at zero extend
$\int_0^\infty H_\chi(t)t^{s-1}\,dt=\Gamma(s)F_\chi(s)$
from $\Re s>1$ to $s=1$. Hence $F_\chi(1)=\int_0^\infty H_\chi(t)\,dt$.

Put $J(y)=\int_0^\infty e^{-y/w-w}\,dw/w$. For $0<y\leq1$,
expand the product and use $a_1*a_\chi\geq0$ and
$(a_1*a_\chi)(1)=1$ to get
\[
 \begin{aligned}
 J(1)\leq J(y)&\leq\sum_{n\geq1}(a_1*a_\chi)(n)J(ny)\\
 &=\int_0^\infty H_1(y/w)H_\chi(w)\,\frac{dw}{w}
 =\frac r y F_\chi(1)+O(D^{22}\sqrt y).
 \end{aligned}
\]
Absolute convergence follows from $J(z)\ll z^{-3}$ and
$\sum_n a_1(n)n^{-3}<\infty$. The error uses
$\min(y/w,w/y)\leq\sqrt{y/w}$ and integrability of
$\min(w,1/w)w^{-3/2}$. Take $\sqrt y=\delta D^{-22}$ for a sufficiently
small absolute $\delta>0$. Since $r\leq4D$, this gives
$|F_\chi(1)|\gg D^{-45}$. Cauchy's estimate with radius $1/10$ gives
$|F'_\chi|\leq18000D$ on $|s-1|\leq1/4$ for every $\chi$.
Thus $1-\beta\gg D^{-46}$.
\end{proof}
\leanref{GlobalTheory.\allowbreak{}HkLocal.\allowbreak{}logDeriv\_\allowbreak{}local}{Development.\allowbreak{}AbelianVarietyTheory.\allowbreak{}Proofs.\allowbreak{}GlobalInterpolation.\allowbreak{}RayClassLFunctionLogDerivLocalFormula}{21--157}
\leanref{GlobalTheory.\allowbreak{}Hk.\allowbreak{}ZeroFree.\allowbreak{}zero\_\allowbreak{}free\_\allowbreak{}region}{Development.\allowbreak{}AbelianVarietyTheory.\allowbreak{}Proofs.\allowbreak{}GlobalInterpolation.\allowbreak{}RayClassLFunctionZeroFreeRegion}{286--312}
\leanref{GlobalTheory.\allowbreak{}Hk.\allowbreak{}RealZero.\allowbreak{}main}{Development.\allowbreak{}AbelianVarietyTheory.\allowbreak{}Proofs.\allowbreak{}GlobalInterpolation.\allowbreak{}RayClassLFunctionUniqueSimpleExceptionalZero}{135--259}
\leanref{GlobalTheory.\allowbreak{}Hk.\allowbreak{}RealChar.\allowbreak{}norm\_\allowbreak{}F\_\allowbreak{}one\_\allowbreak{}ge}{Development.\allowbreak{}AbelianVarietyTheory.\allowbreak{}Proofs.\allowbreak{}GlobalInterpolation.\allowbreak{}NormRayClassFamilyFOneGeOfRealCharacter}{343--423}
\leanref{GlobalTheory.\allowbreak{}Hk.\allowbreak{}LFam.\allowbreak{}real\_\allowbreak{}zero\_\allowbreak{}far\_\allowbreak{}from\_\allowbreak{}one}{Development.\allowbreak{}AbelianVarietyTheory.\allowbreak{}Proofs.\allowbreak{}GlobalInterpolation.\allowbreak{}ExistsRayClassLFunctionRealZeroFarFromOne}{47--77}
\leanref{GlobalTheory.\allowbreak{}Hk.\allowbreak{}ShortInt.\allowbreak{}zeros}{Development.\allowbreak{}AbelianVarietyTheory.\allowbreak{}Proofs.\allowbreak{}GlobalInterpolation.\allowbreak{}TrivialClassShortInterval}{20--102}

\begin{lemma}[A weighted estimate at one endpoint]\label{B2.49b}
Fix $0<c_0\leq1/20$. For any family and character $\chi$, suppose
the set $Z_\chi$ of zeros of $F_\chi$ in $\Re s\geq1-c_0/\ell(\Im s)$
is empty or consists of a simple real zero $\beta<1$. For $u\geq2$
and $\log D\leq\sqrt{\log u}$,
\[
 \left|\Psi_\chi(u)-\mathbf1_{\chi=1}u+\sum_{\beta\in Z_\chi}\frac{u^\beta}{\beta}\right|
 \leq C u e^{-c_7\sqrt{\log u}},
\]
where $C,c_7>0$ depend only on $c_0$.
\end{lemma}
\begin{proof}
Put $w=\sqrt{\log u}$ and $\alpha=c_0/32$. For $\alpha w\leq1$ use
$|\Psi_\chi(u)|\leq2\sum_{n\leq u}\Lambda(n)\ll u$ and
$\beta\geq1-c_0/\ell(0)>1/2$. Otherwise take $T=e^{\alpha w}$, with
$2\leq T\leq u$ and $\ell(T)\leq3w$. Truncated Perron inversion on
$\Re s=1+1/\log u$ costs $O(u\log^2u/T+\log u)$. For $u/2<n<2u$,
sum the kernel error $O(\log u)\min(1,u/(T|n-u|))$, with the
minimum taken as $1$ at $n=u$. For the other $n$ use
$|\log(u/n)|\geq\log2$ and
$\sum_n\Lambda(n)n^{-1-1/\log u}\ll\log u$.

In $1-c_0/(2\ell(t))\leq\Re s\leq3$ at height $t$, away from
$Z_\chi$, we have
\[
 \frac{F'_\chi}{F_\chi}(s)-\sum_{\beta\in Z_\chi}\frac1{s-\beta}
 = O_{c_0}(\ell(t)^2).
\]
In \eqref{eq:host-local-logderivative} there are $O(\ell(t))$ zeros
counted with multiplicity. For each nonexceptional zero in that sum,
$\ell(\Im\rho)\leq4\ell(t)/3$ gives a horizontal distance
$\gg c_0/\ell(t)$ from $s$. If $|\beta-z_t|>3/5$ and
$|s-z_t|\leq1/2$, then $|s-\beta|>1/10$. For $\Re s\geq7/4$ the
absolutely convergent series for $P_\chi$ gives the bound directly.

Take $b=c_0/(8\ell(T))$. Shift the integral of $P_\chi(s)u^s/s$ to
the line $\sigma=1-2b$ or $1-4b$ at distance at least $b$ from
$Z_\chi$. The residues give $\mathbf1_{\chi=1}u$ and $-u^\beta/\beta$
for $\beta>\sigma$. A zero to the left contributes at most
$2u^{1-2b}$ in absolute value. On the new sides
$|P_\chi|\ll_{c_0}\ell(T)^2+1/b$. The horizontal and vertical integrals
cost $O_{c_0}(w^2u/T)$ and $O_{c_0}(Tw^2u^{1-2b})$, respectively.
Since $2b\log u\geq c_0w/12\geq2\log T$, the total error is
$O_{c_0}(uw^4e^{-\alpha w})=O_{c_0}(ue^{-\alpha w/2})$.
\end{proof}
\leanref{GlobalTheory.\allowbreak{}PerronCoeff.\allowbreak{}truncated\_\allowbreak{}perron\_\allowbreak{}of\_\allowbreak{}le\_\allowbreak{}two\_\allowbreak{}vonMangoldt}{Development.\allowbreak{}AbelianVarietyTheory.\allowbreak{}Proofs.\allowbreak{}GlobalInterpolation.\allowbreak{}TruncatedPerronOfLeTwoVonMangoldt}{9--158}
\leanref{GlobalTheory.\allowbreak{}Hk.\allowbreak{}LogDerivHalf.\allowbreak{}logDeriv\_\allowbreak{}sub\_\allowbreak{}exceptional\_\allowbreak{}le}{Development.\allowbreak{}AbelianVarietyTheory.\allowbreak{}Proofs.\allowbreak{}GlobalInterpolation.\allowbreak{}LogDerivSubExceptionalLe}{92--261}
\leanref{GlobalTheory.\allowbreak{}Hk.\allowbreak{}PsiExc.\allowbreak{}psi\_\allowbreak{}exceptional}{Development.\allowbreak{}AbelianVarietyTheory.\allowbreak{}Proofs.\allowbreak{}GlobalInterpolation.\allowbreak{}RayClassFamilyPsiWithExceptionalZero}{438--503}

\begin{lemma}[Principal-class weight in a short interval]\label{B2.49c}
There are absolute $A,c,\eta>0$ such that for every family, $X\geq1$,
$\log D\leq c\sqrt{\log X}$ and $e^{-c\sqrt{\log X}}\leq h\leq1$ imply
\[
 \sum_{X<n\leq X(1+h)}\lambda_0(n)\geq\eta hX/D^A .
\]
\end{lemma}
\begin{proof}
Apply Lemma~\ref{B2.49b} with the $c_0$ of Lemma~\ref{B2.49a}.
Put $L=\sqrt{\log X}$, $U=X(1+h)$ and $\Delta=c_9D^{-a}$. Choose
$c=\min\{1,c_7/(2(a+2)),c_7c_9/(24C)\}$. Then $cL\geq\log D\geq1$,
so $X\geq\mathrm e$ and both endpoints $u=X,U$ satisfy
$\log D\leq\sqrt{\log u}$.
Orthogonality gives $\sum_\chi\Psi_\chi(u)=|G|\sum_{1\leq n\leq u}\lambda_0(n)$.
Subtract the estimates at the endpoints and put $W=\sum_{X<n\leq U}\lambda_0(n)$.
Since $U\leq2X$,
\[
 |G|W=hX-B+R,\qquad |R|\leq3CDX e^{-c_7L}.
\]
Here $B=0$ in the absence of an exceptional pair. If the pair is present,
\[
 B=\int_X^U t^{\beta-1}\,dt\leq hX e^{-\Delta\log X}\leq hX(1-\Delta/2)
\]
since $0<\Delta\leq1$. The endpoint errors satisfy
\[
 \frac{|R|}{hX\Delta}\leq\frac{3C}{c_9}e^{((a+2)c-c_7)L}
 \leq\frac{3C}{c_9}e^{-c_7L/2}\leq\frac14 .
\]
The last step uses $e^{c_7L/2}\geq c_7L/2\geq c_7/(2c)\geq12C/c_9$.
Hence $|G|W\geq hX\Delta/4$. Use $|G|\leq D$ and take
$A=a+1$, $\eta=c_9/4$.
\end{proof}
\leanref{GlobalTheory.\allowbreak{}Hk.\allowbreak{}ShortInt.\allowbreak{}trivialClass\_\allowbreak{}shortInterval}{Development.\allowbreak{}AbelianVarietyTheory.\allowbreak{}Proofs.\allowbreak{}GlobalInterpolation.\allowbreak{}TrivialClassShortInterval}{218--282}
\leanref{GlobalTheory.\allowbreak{}T.\allowbreak{}exists\_\allowbreak{}lam\_\allowbreak{}zero\_\allowbreak{}sum\_\allowbreak{}shortInterval\_\allowbreak{}ge}{Development.\allowbreak{}AbelianVarietyTheory.\allowbreak{}Proofs.\allowbreak{}GlobalInterpolation.\allowbreak{}TrivialClassShortIntervalStatement}{8--9}

For $d\equiv3\pmod4$, define $\operatorname{NormOne}(d,N,p)$ by the
existence of integers $u,v$ such that
\begin{equation}\label{eq:norm-one-coordinates}
 4p=u^2+dv^2,\qquad N\mid v,\qquad 2N\mid u-2-v.
\end{equation}
Write $\mathcal P(d,N;X,X')$ for the primes in $(X,X']$ satisfying
$\operatorname{NormOne}(d,N,p)$.

\begin{theorem}[Norm primes in a short interval]\label{B2.50}
There are absolute $A,c,c'>0$ as follows.
Let $d\geq100$ be prime with $d\equiv3\pmod4$, $N\in\N$ with $N\geq3$,
and $X,h\in\R$ with
\[
 X\geq1,\qquad
 \log(4dN^2)\leq c\sqrt{\log X},\qquad
 e^{-c\sqrt{\log X}}\leq h\leq1.
\]
Then
\begin{equation}\label{eq:short-norm-primes}
 |\mathcal P(d,N;X,X(1+h))|
 \geq\frac{c'hX}{(4dN^2)^A\log X}.
\end{equation}
\end{theorem}
\leanref{UpperTheory.\allowbreak{}T3.\allowbreak{}shortNorm}{Development.\allowbreak{}AutomorphicArithmetic.\allowbreak{}ShortNorms.\allowbreak{}Primes}{5--6}
\leanref{GlobalTheory.\allowbreak{}T.\allowbreak{}normPrimesImagQuad\_\allowbreak{}short\_\allowbreak{}interval\_\allowbreak{}lower\_\allowbreak{}bound}{Development.\allowbreak{}AbelianVarietyTheory.\allowbreak{}Proofs.\allowbreak{}GlobalInterpolation.\allowbreak{}NormPrimesImagQuadShortIntervalLowerBoundStatement}{8--10}
\leanref{GlobalTheory.\allowbreak{}Mem.\allowbreak{}An.\allowbreak{}NormOne}{Development.\allowbreak{}AbelianVarietyTheory.\allowbreak{}Definitions.\allowbreak{}Foundations.\allowbreak{}RayPrimes}{40--42}

\begin{proof}
Let $K=\Q(\sqrt{-d})$, $G$ its ray class group modulo $N\mathcal O_K$,
and $D=4dN^2$. For ideals prime to $N$ of norm $n$ in class $g$,
let $a_g(n)$ count the ideals and $\lambda_g(n)$ sum their ideal von
Mangoldt weights. The ray class number is at most $h_KN^2\leq D$.
Minkowski's bound gives an ideal of norm at most $\sqrt d$
in each ordinary class. At norm $n$ there are at most $\tau(n)\leq n$
ideals, where $\tau$ counts positive divisors, so $h_K\leq d$.

Choose an integral ideal $J$ prime to $N$ in class $-g$.
Ideals $I$ in class $g$ correspond to generators $z$ of $IJ$ with
$z\equiv1\bmod N$. Such a generator is unique since $N\geq3$.
The possible $z$ form a translate of $NJ$ in $J$. In a reduced basis
of $NJ$, write
\[
 \frac{|z|^2}{\mathrm{N}J}
 =\alpha(m+\theta_1+\ell(k+\theta_2))^2+\gamma(k+\theta_2)^2
 \qquad(m,k\in\Z),
\]
where $\theta_1,\theta_2$ describe the translate. For the first basis
vector $w_1$, $\alpha=|w_1|^2/\mathrm{N}J\geq N^2$.
Reduction gives $|\ell|\leq1/2$ and $\gamma\geq3\alpha/4$.
The covolume gives $\alpha\gamma=dN^4/4$, so
$\kappa=\pi/\sqrt{\alpha\gamma}=2\pi/(\sqrt d\,N^2)$ lies in $[D^{-1},4]$.
Each horizontal chord costs at most one in rounding. The sum of chord
lengths differs from its integral by at most the total variation
$4\sqrt{y/\alpha}$. The area is $\kappa y$ and the total count error
is at most $2\sqrt{y/\gamma}+3+4\sqrt{y/\alpha}\leq10\sqrt y$.
For the Gaussian sum, one-dimensional Poisson summation twice gives
relative error at most $(\alpha t)^2+(\gamma t)^2+(\alpha\gamma t^2)^2$.
At scale $s$ the shifted one-dimensional relative error is at most
$2\sum_{j\geq1}e^{-\pi^2j^2/s}\leq s^2$, by $e^{-z}\leq2/z^2$.
For $0<t\leq1$ the absolute error is at most
$3\pi(\alpha\gamma)^2t\leq D^{10}t$.

The counts give $a_g(0)=0$ and $a_g(1)=\mathbf1_{g=0}$. Factorization
gives the logarithmic convolution and $\sum_g\lambda_g(n)\leq2\Lambda(n)$
since $[K:\Q]=2$. For real $\chi$, the coefficient at an ideal $I$
prime to $N$ in $L_1L_\chi$ is
$\prod_{\mathfrak p^e\parallel I}\sum_{j=0}^e\chi(\mathfrak p)^j$.
It is nonnegative and at least one on squares. Summing over ideals
of norm $n$ gives $(a_1*a_\chi)(n)\geq0$. The $a_1(n)$ ideals of
norm $n$ prime to $N$ have distinct squares of norm $n^2$, each with
coefficient at least one. Hence $(a_1*a_\chi)(n^2)\geq a_1(n)$.
Lemma~\ref{B2.49c} gives $W=\sum_{X<n\leq U}\lambda_0(n)\geq\eta hX/D^A$,
where $U=X(1+h)$.
\leanref{GlobalTheory.\allowbreak{}Hk.\allowbreak{}rayCountTheta}{Development.\allowbreak{}AbelianVarietyTheory.\allowbreak{}Proofs.\allowbreak{}GlobalInterpolation.\allowbreak{}RayCountThetaFibreCount}{53--122}
\leanref{GlobalTheory.\allowbreak{}Hk.\allowbreak{}exists\_\allowbreak{}lFam\_\allowbreak{}ray}{Development.\allowbreak{}AbelianVarietyTheory.\allowbreak{}Proofs.\allowbreak{}GlobalInterpolation.\allowbreak{}ExistsLFamRay}{73--129}

If $p$ is prime and $\lambda_0(p)>0$, an ideal of norm $p$ has a
generator $z=(u-v)/2+v(1+\sqrt{-d})/2\equiv1\bmod N$. Its norm
and congruence give \eqref{eq:norm-one-coordinates}. Each such prime
contributes at most $2\log U$, and the other positive weights lie
on higher prime powers. Thus
\[
 W\leq2\log U\,|\mathcal P(d,N;X,U)|+2\sum_{j\geq2}\vartheta(U^{1/j}),
\]
where $\vartheta(y)=\sum_{p\leq y}\log p$. By Chebyshev's bound the
last term is $O(\sqrt U+U^{1/3}\log U)=O(\sqrt X)$. For
$v=\sqrt{\log X}$ we have
$hX/(D^A\sqrt X)\geq\exp(v^2/2-(A+1)cv)$ and $v\geq\log4/c$.
Shrink $c$ so that the higher powers cost at most $\eta hX/(2D^A)$.
Since $\log U\leq2\log X$, take $c'=\eta/8$.
\end{proof}
\leanref{GlobalTheory.\allowbreak{}T.\allowbreak{}exists\_\allowbreak{}lFam\_\allowbreak{}normOne\_\allowbreak{}of\_\allowbreak{}mod\_\allowbreak{}four\_\allowbreak{}eq\_\allowbreak{}three}{Development.\allowbreak{}AbelianVarietyTheory.\allowbreak{}Proofs.\allowbreak{}GlobalInterpolation.\allowbreak{}ExistsLFamRayStatement}{8--9}
\leanref{GlobalTheory.\allowbreak{}Hk.\allowbreak{}normPrimes\_\allowbreak{}short\_\allowbreak{}interval\_\allowbreak{}lower}{Development.\allowbreak{}AbelianVarietyTheory.\allowbreak{}Proofs.\allowbreak{}GlobalInterpolation.\allowbreak{}NormPrimesImagQuadShortIntervalLowerBound}{98--148}

For the separation of $g_0,g_1$, define
\[
 \begin{split}
 \operatorname{Exc}(g,j,\ell)
 &\Longleftrightarrow
 \ell\text{ is prime and }g^j<\ell^{\,4v_\ell(g^j-1)},\\
 \mathcal B_r(g)
 &=\{g\}\cup\{\ell:\operatorname{Exc}(g,j,\ell)
                  \text{ for some }1\leq j\leq r\}.
 \end{split}
\]

\begin{theorem}[Choosing the five level primes]\label{B2.51}
There is an absolute $C_8\geq1$ such that the following holds for every
$0<\varepsilon<1/52$ and all sufficiently large scales $s$.
Suppose $\operatorname{RankVal}(s,R)$ holds, $N,d\in\N$, and
$d\geq100$ is prime with $d\equiv3\pmod4$. Assume $N>0$, $d\mid N$
and $\log N\leq3x^{\varepsilon/5}(\log x)^2$.
There are primes $g_0,g_1,p,p',q_0$ with
\[
 (kb)^8\leq g_i\leq2(kb)^8,\quad
 \operatorname{NormOne}(d,N,g_i)\quad(i=0,1),\qquad
 g_0\ne g_1,\quad \mathcal B_r(g_0)\cap\mathcal B_r(g_1)=\varnothing.
\]
The primes $p,p',q_0$ are pairwise distinct and at most $N^{C_8}$.
They satisfy $p\equiv p'\equiv1\pmod{15N}$, $q_0\equiv1\pmod N$,
and $p,p'\ne g_i$ for $i=0,1$.
\end{theorem}
\leanref{UpperTheory.\allowbreak{}T3.\allowbreak{}levelData}{Development.\allowbreak{}AutomorphicArithmetic.\allowbreak{}ShortNorms.\allowbreak{}levelData}{36--78}
\leanref{UpperTheory.\allowbreak{}T3.\allowbreak{}typePrimesN}{Development.\allowbreak{}AutomorphicArithmetic.\allowbreak{}ShortNorms.\allowbreak{}typePrimesN}{78--130}

\begin{proof}
Take $V=(kb)^8$ and $D=4dN^2$. Since $d\mid N$,
$\log D\leq\log4+9x^{\varepsilon/5}(\log x)^2$ and
$\log V\geq8x^\varepsilon$. Thus $\log D=o(\sqrt{\log V})$ and
$A\log D+\log\log V-\log c'\leq(\log V)/10$ for large $x$.
Apply \ref{B2.50} with $h=1$ to obtain
\begin{equation}\label{eq:norm-candidate-count}
 |\mathcal P(d,N;V,2V)|\geq V^{9/10}.
\end{equation}
\leanref{UpperTheory.\allowbreak{}T3.\allowbreak{}normPrimesN}{Development.\allowbreak{}AutomorphicArithmetic.\allowbreak{}ShortNorms.\allowbreak{}normPrimesN}{108--158}
The rank ranges give $r\leq x^5$. Fix $g_*\in[V,2V]\cap\Z$.
For each $j$, four distinct exceptional prime powers would divide
$g_*^j-1$ and have product greater than $g_*^j$.
Hence $|\mathcal B_r(g_*)|\leq3r+1$.

For composite $\ell$, membership $\ell\in\mathcal B_r(g)$ forces
$g=\ell$. Fix a prime $\ell$ and $1\leq j\leq r$. Take the
least $a$ with $\ell^a>V^{j/4}$. Then $\operatorname{Exc}(g,j,\ell)$
implies $g^j\equiv1\bmod\ell^a$. For odd $\ell$ there are at most
$j$ roots by cyclicity. For $\ell=2$, the cyclic subgroup
$\langle5\rangle$ has index at most two, so there are at most $2j$.
Each residue class contains at most $1+V/\ell^a\leq1+V^{3/4}$
candidates. For each $\ell$, sum over $1\leq j\leq r$ and include $g=\ell$.
Then sum over $\ell\in\mathcal B_r(g_*)$. At most
\begin{equation}\label{eq:bb-exclusion-count}
 (3r+1)\bigl(1+2r^2(V^{3/4}+1)\bigr)
\end{equation}
integers $g\in[V,2V]$ have $\mathcal B_r(g)\cap\mathcal B_r(g_*)\ne\varnothing$.
Since $20r^3<V^{3/20}$ for large $x$, the bound is less than $V^{9/10}$.
Choose $g_0$ among the norm primes and $g_1$ outside its exclusion set.
Disjointness gives $g_0\ne g_1$.

Let $A\geq1$ be the exponent of \ref{B2.49}. Choose $p,p',q_0$
successively in the classes $1$ modulo $15N$, $15Np$, and $Npp'$.
They are pairwise distinct. Since $N\geq100$,
\[
 p\leq N^{2A},\qquad p'\leq N^{(2A+2)A},\qquad
 q_0\leq N^{(1+2A+(2A+2)A)A}.
\]
Take $C_8=(1+2A+(2A+2)A)A$. For large $x$,
$C_8\log N\leq3C_8x^{\varepsilon/5}(\log x)^2<x^\varepsilon$.
Thus $p,p',q_0<b\leq(kb)^8\leq g_i$.
\leanref{GlobalTheory.\allowbreak{}choice\_\allowbreak{}count}{Development.\allowbreak{}AbelianVarietyTheory.\allowbreak{}Proofs.\allowbreak{}FoundationProofs.\allowbreak{}DivLeRpowThreeQuarters}{221--222}
\end{proof}

\section{Real characters and the thin positions}

Extend the window-pair parameters of \ref{B2.42} to real $Y$.
Put $a=\varepsilon/5$. For $Y>0$, let $z=Y^a$. For $Y<0$, let
$z=|Y|^a\cos(\pi a)$. For $Y\ne0$, put $z'=z(\log|Y|)^2$.
At $Y=0$, take $z=\mathbf1_{a=0}$ and $z'=0$.
Set $X_1(Y)=\exp z$, $X_2(Y)=\exp z'$ and
$H_M(j)=8MX_1(Y)X_2(Y)^j$ for $j\in\{0,1\}$.
Write $\operatorname{ZeroBox}(T,\chi)$ for nonvanishing of $L(s,\chi)$
when $\Re s>399/400$ and $|\Im s|\leq T$. Let $Q_{\mathrm{odd}}$
be the odd part of $Q_S$.

\begin{theorem}[Real-character boxes and the thin-position bound]\label{B2.52}
Fix a raw host $H$ on a frame $\phi$, with $q=H.q$ and $d=H.d$.
Let $M\geq1$ be an integer and $\varepsilon,Y,X\in\R$. Suppose
$Q_{\mathrm{odd}}\mid M$ and every prime divisor of $Q_S$ divides $2M$.
If $\operatorname{PairGood}(M,\varepsilon,Y,X,q,d)$ holds, then every
primitive real character $\chi$ of the host field whose conductor
$N_\chi$ is divisible by $q$ or $d$ satisfies
\[
 \operatorname{ZeroBox}(H_M(j),\chi),\qquad
 N_\chi\leq8Q_{\mathrm{odd}}X_1(Y)X_2(Y)^j,\qquad
 j=\mathbf1_{d\mid N_\chi}.
\]
In particular, these characters have no real zero $\beta>399/400$.

Take $c_6=1$. For every $B\in\R$ and $0<\varepsilon<1/52$ there is
$x_0$ such that for every $x\geq x_0$ the following holds.
Suppose $Q_{\mathrm{odd}}\leq B$ and
$\log(Q_Sqd)\leq2x^{\varepsilon/5}(\log x)^2$, and suppose every
primitive real host character with conductor divisible by $q$ or $d$
has no real zero greater than $399/400$. Then every primitive host
character $\chi$ satisfies
\begin{equation}\label{eq:thin-character-bound}
 |g_\chi(y)|\leq y\exp(-c_6x^{\varepsilon/8})
 \qquad\left(y\geq\frac{\exp(x^{\varepsilon/3})}{3}\right).
\end{equation}
\end{theorem}
\leanref{UpperTheory.\allowbreak{}Top.\allowbreak{}realBox}{Development.\allowbreak{}AutomorphicArithmetic.\allowbreak{}CharacterBoxes.\allowbreak{}RealBox}{99--113}
\leanref{QH.\allowbreak{}thinH}{Development.\allowbreak{}AutomorphicArithmetic.\allowbreak{}MeanSquareBounds.\allowbreak{}GaussSumFunctionEq}{55--108}
\leanref{Formalism.\allowbreak{}HC.\allowbreak{}HostRaw.\allowbreak{}ChebyshevCharacterBound}{Formalism.\allowbreak{}SharpRestriction.\allowbreak{}RamificationFrobenius.\allowbreak{}TotientGaloisClasses.\allowbreak{}Analytic}{145--148}

\begin{proof}
By CRT and the conductor classification of primitive real characters,
$\chi(a)=\rho(a)(a/q)^i(a/d)^j$ on units, with $\rho$ primitive real,
$N_\rho\mid8Q_{\mathrm{odd}}$, $q,d\nmid N_\rho$, and
$N_\chi=N_\rho q^id^j$ for $i,j\in\{0,1\}$. If $q$ or $d$ divides
$N_\chi$, the twist is nonzero and $j=\mathbf1_{d\mid N_\chi}$.
The prime-support hypothesis makes the integers tested by
$\operatorname{PairGood}$ coprime to $N_\chi$. Since $N_\rho\mid8M$,
its real-character boxes apply. The pair sizes give the conductor bound.

The nontrivial real host characters of conductor prime to $qd$ have
conductor at most $8\lfloor B\rfloor$. Nonvanishing at $1$ gives
this finite family a common zero-free interval to the left of $1$.
Take the smaller of its width and $1/400$ to cover the other real
host characters. The classical zero-free region then gives, for
$\lambda=\max(1,\log N_\chi)$ and every nontrivial zero $\rho$ of
$L(s,\chi)$,
\[
 \Re\rho\leq1-\frac{c}{\lambda+\log(|\Im\rho|+3)}
\]
with $c>0$ depending on $B$. Apply the same bound to $1-\rho$ for
$\overline\chi$ by the functional equation. For $|\Im\rho|\leq2$
this gives $|\rho|\gg_B1/\lambda$. The zero count in a unit interval
at height $h$ is $O(\lambda+\log(|h|+3))$ with multiplicity. Thus the
low zeros contribute $O_B(\lambda^2)$ to the sum below, and summing
the other unit-height counts with weights $1/j$ gives
\[
 \sum_{|\Im\rho|<T}\frac{m_\rho}{|\rho|}
 \ll_B(\lambda+\log T)^2\qquad(T\geq1),
\]
where $m_\rho$ denotes multiplicity.

For $u\geq\mathrm e$, put $v=\log u$, $T=\exp(\sqrt v)$ and
$D_v=\exp(-c'v/(\lambda+\sqrt v))$, with $c'=\min(c/3,1)$.
Then $|u^\rho|\leq uD_v$ for $|\Im\rho|<T$ and $T^{-1}\leq D_v$.
In the truncated explicit formula for
$\psi_0(\chi,u)=\sum_{n\leq u}\Lambda(n)\chi(n)-\mathbf1_{\chi=1}u$,
the zero sum is $O_B(u(v+\lambda)^2D_v)$. The error is bounded by
a constant times
$u\log^2(N_\chi uT)/T+\log(N_\chi u)+\sum_{|\Im\rho|<1}m_\rho/|\rho|$.
Since $uD_v\geq1$, this has the same bound.

Use $g_\chi(y)=-y\int_y^\infty\psi_0(\chi,u)\,du/u^2$.
Since $N_\chi\mid Q_Sqd$, we have
$\lambda\leq1+2x^{\varepsilon/5}(\log x)^2$.
For $v\geq V=\log y\geq x^{\varepsilon/3}/2$ and large $x$, this gives
$\lambda\leq v^{61/100}$. Thus
$(v+\lambda)^2D_v\leq4v^2e^{-(c'/2)v^{39/100}}$.
Substituting $u=\mathrm e^v$ bounds $|g_\chi(y)|/y$ by
\[
 4C e^{-(c'/4)V^{39/100}}\int_0^\infty v^2e^{-(c'/4)v^{39/100}}\,dv
\]
with $C$ depending on $B$. The integral is finite and
$V^{39/100}\gg x^{13\varepsilon/100}$, with $13/100>1/8$.
This gives \eqref{eq:thin-character-bound} for large $x$.
\leanref{QH.\allowbreak{}psi\_\allowbreak{}thin}{Development.\allowbreak{}AutomorphicArithmetic.\allowbreak{}MeanSquareBounds.\allowbreak{}PsiThin}{385--422}
\leanref{QH.\allowbreak{}g\_\allowbreak{}integral}{Development.\allowbreak{}AutomorphicArithmetic.\allowbreak{}MeanSquareBounds.\allowbreak{}GIntegral}{229--234}
\leanref{QH.\allowbreak{}int\_\allowbreak{}thin}{Development.\allowbreak{}AutomorphicArithmetic.\allowbreak{}MeanSquareBounds.\allowbreak{}IntThin}{169--263}
\end{proof}

\section{The thick set and the top set}

The thick places are $\tau_2(u_0+sa_0)$ for $s\in T_2$, with $u_0$
the host's offset in $A$.

\begin{theorem}[A thick set without short relations]\label{B2.53}
There is an absolute $C_T>0$. Fix $0<\varepsilon<1/52$ and real
$C_o,C_L$. For all sufficiently large $x$, let a frame $\phi$ and
an integer $L\geq2$ satisfy $L\leq C_L\log x$, $o_\phi\leq C_o$,
and $x^{2+8\varepsilon}/2\leq n\leq C_o x^{2+8\varepsilon}$.
There is a finite set $T_2\subset\N$ with $4Ls<n_0$ for $s\in T_2$,
$4L|T_2|\leq n_0$, and
\begin{equation}\label{eq:thick-size-bounds}
 |T_2|\geq
 \frac{n_0}{L}\exp\left(-C_T\sqrt{\log n_0\log L}\right),
 \qquad |T_2|\geq x^{-\varepsilon}n.
\end{equation}
For distinct $s_1,s_2,s_3\in T_2$ and integers $i,j$ with
$0<|i|,|j|\leq L$ and $i\ne j$,
\begin{equation}\label{eq:no-short-relation}
 j(s_2-s_1)\ne i(s_3-s_1).
\end{equation}
\end{theorem}
\leanref{UpperTheory.\allowbreak{}BernoulliSumsFrobeniusOrbits.\allowbreak{}thickFrame2}{Development.\allowbreak{}AutomorphicArithmetic.\allowbreak{}ThickFrames.\allowbreak{}ThickFrame2Large}{29--199}
\leanref{Formalism.\allowbreak{}FrobeniusOrbits.\allowbreak{}NoRel}{Formalism.\allowbreak{}FrobeniusOrbits.\allowbreak{}HostThickExists}{13--17}

\begin{proof}
First let $N\geq1$ and $L\geq2$. One sphere of constant squared norm
in $\{0,\ldots,d-1\}^S$ contains at least $d^S/(Sd^2)$ points.
Map them to integers by $v\mapsto\sum_{\nu=0}^{S-1}v_\nu(4Ld)^\nu$.
In a putative relation \eqref{eq:no-short-relation}, each coefficient
has absolute value less than half the base. Successive reduction
modulo the base makes every coefficient zero. The three vectors
would lie on a line, which meets the sphere in at most two points.

Choose the parameters so that $4Ls<N$ and the cardinality is at least
\begin{equation}\label{eq:integer-sphere-bound}
 \frac NL\exp(-32\sqrt{\log N\log L}).
\end{equation}
If this is at most $1$, take $\{0\}$. Otherwise put $v=\log N$,
$w=\log L$, $\sigma=\sqrt{vw}$, $S=\lfloor\sigma/w\rfloor$,
$a=(v-w-3)/S$ and $d=\lfloor e^a/(4L)\rfloor$.
Then $S\geq32$, $a\geq7+w$ and $e^a/(8L)\leq d\leq e^a/(4L)$.
Since $N=e^{aS}Le^3$, we have $4L(4Ld)^S\leq N$.
The sphere count gives \eqref{eq:integer-sphere-bound} because
$e^3(8L)^SS e^{2a}\leq e^{32\sigma}$. This follows from
$Sw\leq\sigma$, $a\leq2\sigma$ and $3+S(7+w)+S+2a\leq32\sigma$.

Take $N=n_0$, $E=\exp(32\sqrt{\log n_0\log L})$, and $Q=n_0/(LE)$.
The bounds $\log n_0=O(\log x)$ and $\log L=O(\log\log x)$ give
$4C_oLE\leq x^\varepsilon$ eventually, hence $Q\geq4x^{-\varepsilon}n$.
Since $x^{-\varepsilon}n\geq1$ and $E\geq2$, deleting $0$ leaves at
least $Q-1\geq\max(x^{-\varepsilon}n,Q/E)$ points. This gives
\eqref{eq:thick-size-bounds} with $C_T=64$. The remaining $s$ are
positive and $4Ls<n_0$, so $4L|T_2|\leq n_0$.
\leanref{Formalism.\allowbreak{}FrobeniusOrbits.\allowbreak{}hostThickExists}{Formalism.\allowbreak{}FrobeniusOrbits.\allowbreak{}FrobeniusOrbitsProofs.\allowbreak{}ThickExists}{181--335}
\end{proof}

For $\gamma\in A$, take $\lambda_\gamma=4/(3o_\phi+1)$ on $A_0$ and
$3/(3o_\phi+1)$ elsewhere. The assertion $\operatorname{Shares}(TT,\rho)$
means
\begin{equation}\label{eq:top-shares}
 \left||TT\cap(\gamma+A_0)|-\lambda_\gamma|TT|\right|
       \leq\rho|TT|\qquad(\gamma\in A).
\end{equation}

Put $X'=\exp(x^{\varepsilon/3})$, and let $h_p\in A$ be the host
Frobenius coordinate. Let $\mathcal P_s$ be the inert-class primes
$p<\lceil X'\rceil$ with $\operatorname{ord}(h_p)\geq x^{\varepsilon/4}$.
Let $b_p(TT)$ count the members of $TT$ whose $h_p$-orbit fails the
good-orbit condition below. Then $\operatorname{BadOK}(s,C_b,H)$
is the bound
\begin{equation}\label{eq:bad-top-budget}
 \sum_{p\in\mathcal P_s}\frac{\log p}{p}\,b_p(TT)
 \leq4C_bx^{-2}\delta_1 n
          \sum_{p\in\mathcal P_s}\frac{\log p}{p}
\end{equation}
with $TT=H.TT$.

For a selected point $t$, let $t_0$ be the first element of
$t+\langle h_p\rangle$ in the frame's enumeration of $A$. Mark
coordinate $i$ when $t_0+ih_p\in TT$. Define
\[
 \begin{gathered}
 H_\varepsilon=2+8\varepsilon,\qquad
 K_\varepsilon=\left\lceil8(H_\varepsilon+2)/\varepsilon\right\rceil,
 \qquad R_\varepsilon=K_\varepsilon+1,\qquad u=\lceil1/\delta_1\rceil,\\
 \ell=(6R_\varepsilon+1)u,\qquad
 M_0=\left\lceil(H_\varepsilon+3)\log x\,10^{6R_\varepsilon+1}\right\rceil,
 \qquad G'=(M_0+5)\ell,\\
 N_1=1+\left\lceil\max(6\mathrm e G'\delta_1,(H_\varepsilon+3)\log_2x)\right\rceil,
 \qquad
 L_2=2\sup_{\substack{p\ {\rm prime}\\p\ {\rm inert\ class}}}
              \left\lceil\log_p((kx^{26})^4)\right\rceil .
 \end{gathered}
\]
On an orbit of length $N_{\rm orb}$, the marked set must have at most
$K_\varepsilon$ points in every cyclic window of length $L_2$ and at
most $N_1$ in every window of length $G'$. Divide the orbit into
consecutive full tiles of length $\ell$. A tile is marked when its
odd-numbered cells of length $u$ are empty and each even-numbered cell
contains exactly one marked point, of even orbit coordinate. Every
run of $M_0+1$ tiles, cyclically among the full tiles, must contain
a marked tile.

\begin{theorem}[Top positions with the required shares]\label{B2.54}
Fix $0<\varepsilon<1/52$ and $C_8\in\R$. There are $C_b>0$ and $x_0$
such that if $s=(x,\varepsilon,J)$, $R$, $H$ and $Y$ satisfy
\[
 \begin{gathered}
 x\geq x_0,\quad \operatorname{RankVal}(s,R),\quad
 J,o_\phi\leq C_8,\quad Q_S\leq C_8n,\\
 x^{1/J}\leq Y\leq x,\qquad \operatorname{PairRow}(Y,\varepsilon,H),
 \end{gathered}
\]
then there is $TT\subset A$ with $\delta_1n/2\leq|TT|\leq2\delta_1n$,
$\operatorname{Shares}(TT,x^{-1+\varepsilon})$, and \eqref{eq:bad-top-budget}.
\end{theorem}
\leanref{UpperTheory.\allowbreak{}BernoulliSumsFrobeniusOrbits.\allowbreak{}topsChoiceH}{Development.\allowbreak{}AutomorphicArithmetic.\allowbreak{}ThickFrames.\allowbreak{}topsChoiceH}{30--68}

\begin{proof}
Give $1_{a\in T}$ the product Bernoulli law with parameter
$p_a=\delta_1$ on $A_0$ and $p_a=3\delta_1/4$ elsewhere. For each
coset $C=\gamma+A_0$, put $Y_C=|T\cap C|-\mathbb E|T\cap C|$
and $D=\sqrt{2\delta_1n}$. Then
\[
 \begin{gathered}
 \mu=\mathbb E|T|=\frac{3o_\phi+1}{4o_\phi}\delta_1n,\qquad
 \mathbb E|T\cap C|=\lambda_\gamma\mu,\\
 \mathbb E\sum_CY_C^2=\sum_ap_a(1-p_a)\leq D^2/2.
 \end{gathered}
\]
Condition on $t\in T$. The other indicators retain their product law
with parameters in $[3\delta_1/4,\delta_1]$.
For large $x$, $(M_0+1)\ell=O_\varepsilon(x^{\varepsilon/8}\log x)
\leq x^{\varepsilon/4}\leq N_{\rm orb}$, so a run contains $M_0+1$
distinct full tiles. A tile away from $t$ is marked with probability
at least
\[
 q_*=((1-\delta_1)^u)^{3R_\varepsilon+1}
 \left(\lfloor u/2\rfloor\frac{3\delta_1}{4}
                 (1-\delta_1)^{u-1}\right)^{3R_\varepsilon}.
\]
At least $M_0$ tiles in a run have disjoint coordinates away from $t$,
so the run fails with probability at most $(1-q_*)^{M_0}$. A
short-window failure needs at least $K_\varepsilon$ further selected
points and a long-window failure needs at least $N_1-1$. The
conditional orbit failure probability is at most
\begin{equation}\label{eq:orbit-failure-probability}
 N_{\rm orb}\left[
 \binom{L_2}{K_\varepsilon}\delta_1^{K_\varepsilon}
 +(1-q_*)^{M_0}
 +\operatorname{BinTail}(G',\delta_1,N_1-1)\right],
\end{equation}
where $\operatorname{BinTail}(G',\delta_1,a)$ is the probability
that $\operatorname{Bin}(G',\delta_1)\geq a$.

Write $\delta=\delta_1\leq1/100$. Since $u=\lceil1/\delta\rceil$,
\[
 (1-\delta)^u\geq\exp\left(-\frac{1+\delta}{1-\delta}\right)>\frac3{10},
 \qquad \delta\lfloor u/2\rfloor\geq\frac{1-\delta}{2}.
\]
An occupied-cell factor is at least $(3/4)(99/200)(3/10)>1/10$.
Thus $q_*\geq10^{-(6R_\varepsilon+1)}$ and
$(1-q_*)^{M_0}\leq x^{-(H_\varepsilon+3)}$.
The binomial tail is at most $(\mathrm e G'\delta_1/a)^a$ at
$a=N_1-1$, which gives the same estimate. Multiplying either bound
by $N_{\rm orb}\leq2x^{H_\varepsilon}$ gives $2x^{-3}$.

Put $C=\max(C_8,2)$. The pair row gives $p>x^{\varepsilon/(5C)}$ for
every inert-class prime. For $x\geq16$, \eqref{eq:host-rank-ranges}
gives $\log_p((kx^{26})^4)\leq620C/\varepsilon$, so
$L_2\leq2(620C/\varepsilon+1)$. Since
$\delta_1^{K_\varepsilon}\leq x^{-(H_\varepsilon+2)}$, the orbit length times the first term in brackets in
\eqref{eq:orbit-failure-probability} is $O_{\varepsilon,C_8}(x^{-2})$. Hence that probability
is at most $\beta=C_bx^{-2}$ for suitable $C_b$.

Let $B(T)$ be the left side of \eqref{eq:bad-top-budget} with $TT=T$,
and put $m_B=\beta\delta_1n\sum_{p\in\mathcal P_s}\log p/p$.
The conditional bound gives $\mathbb EB(T)\leq m_B$. If $m_B>0$,
then $\mathbb E(\sum_CY_C^2/D^2+B(T)/(4m_B))\leq3/4$.
Some $T$ has all $|Y_C|\leq D$ and $B(T)\leq4m_B$. If $m_B=0$,
then $B(T)=0$ almost surely and the second-moment bound gives the
same choice. For this choice,
\[
 \bigl||T|-\mu\bigr|\leq o_\phi D,\qquad
 \left||T\cap C|-\lambda_\gamma|T|\right|
 \leq (1+\lambda_\gamma o_\phi)D.
\]
Since $\delta_1n\geq x^{2+63\varepsilon/8}$ and $o_\phi$ is bounded,
these give the size and share bounds for $TT=T$ at large $x$.
\leanref{UpperTheory.\allowbreak{}BernoulliSumsFrobeniusOrbits.\allowbreak{}not\_\allowbreak{}good\_\allowbreak{}le}{Development.\allowbreak{}AutomorphicArithmetic.\allowbreak{}ThickFrames.\allowbreak{}NotGoodLe}{84--139}
\leanref{UpperTheory.\allowbreak{}BernoulliSumsFrobeniusOrbits.\allowbreak{}select}{Development.\allowbreak{}AutomorphicArithmetic.\allowbreak{}ThickFrames.\allowbreak{}select}{70--175}
\end{proof}

\section{Completing a host over the chosen field}

Write $\operatorname{PairWin}(s,Y)$ for $x^{1/J}\leq Y\leq x$.
The condition $\operatorname{PairRow}(Y,\varepsilon,H)$ requires prime
$q,d$ satisfying \eqref{eq:host-prime-hyp}, together with
\[
 \begin{gathered}
 \sqrt{X_1(Y)}<q\leq X_1(Y),\qquad
 \sqrt{X_2(Y)}<d\leq X_2(Y),\\
 \left(\frac{-d}{p}\right)_{\!K}=1
       \quad(p\text{ prime},\ p\leq Y^{\varepsilon/5}),\qquad
 \log(Q_Sqd)\leq2Y^{\varepsilon/5}(\log Y)^2.
 \end{gathered}
\]
The subscript $K$ denotes the Kronecker symbol, so the splitting
condition includes $p=2$. The condition
$\operatorname{QDq}(x,\varepsilon,C,q,d)$ bounds $|g_\chi(x)|$ by
$x^{1/2+\varepsilon}$ and $|(L'/L)(1,\chi)|$ by $C\log x$ for the
primitive representatives of the two sign characters
$(\,\cdot\,/q)$ and $(\,\cdot\,/d)$ and their product.
The condition $\operatorname{Boxq}(Y,\varepsilon,q,d)$ requires
$\operatorname{ZeroBox}(T,\chi)$ and $N_\chi\leq T$ for these
characters and their products with $\chi_8(n)=(2/n)_K$ of conductor
$8$. Take $T=8X_1(Y)$ for the two twists involving $q$ alone and
$T=8X_1(Y)X_2(Y)$ for the four involving $d$.

\begin{lemma}[The numerical conditions supplied by a good pair]\label{B2.55}
Fix $0<\varepsilon<1/52$ and $C_0\in\R$.
For all sufficiently large $x$, suppose $\operatorname{RankVal}(s,R)$
holds and $H$ is a host with $J,o_\phi\leq C_0$ and $Q_S\leq C_0n$.
If $M\geq1$, $\operatorname{PairWin}(s,Y)$ and
$\operatorname{PairGood}(M,\varepsilon,Y,x,q,d)$ hold, then so do
$\operatorname{PairRow}(Y,\varepsilon,H)$,
$\operatorname{QDq}(x,\varepsilon,4,q,d)$ and
$\operatorname{Boxq}(Y,\varepsilon,q,d)$. The host also satisfies
$q,d>4tr+1$.
\end{lemma}
\leanref{UpperTheory.\allowbreak{}Top.\allowbreak{}pairRows}{Development.\allowbreak{}AutomorphicArithmetic.\allowbreak{}CharacterBoxes.\allowbreak{}PairRows}{334--388}

\begin{proof}
The pair congruences give the prime ranges and $-d\equiv1$ modulo $q$
and every odd prime $p\leq Y^{\varepsilon/5}$. The congruence
$d\equiv7\pmod{16}$ gives Kronecker value $1$ at $2$.
Put $z=Y^{\varepsilon/5}$. Since $n\leq2x^{9/4}$,
$\log Q_S\leq2C_0+\frac94\log x$. Using $\log x\leq C_0\log Y$,
for large $x$ we have $\log Y\geq2$, $z\geq820C_0\log Y$ and
$\log Q_S\leq z((\log Y)^2-1)$. Add $\log q\leq z$ and
$\log d\leq z(\log Y)^2$ for the pair-row bound.

The twists of the trivial character in $\operatorname{PairGood}$
give $\operatorname{QDq}$ since $x^{1/2+\varepsilon/2}\leq x^{1/2+\varepsilon}$.
For $\operatorname{Boxq}$ use base characters $1,\chi_8$ with conductors
$f=1,8$. Coprimality gives primitive twist conductor
$fq^id^j\leq8X_1(Y)X_2(Y)^j$. Restrict the box of height
$8MX_1(Y)X_2(Y)^j$ to this height. On odd test integers, quadratic
reciprocity identifies $(-q/\,\cdot\,)$ and $(-d/\,\cdot\,)$ with the
two sign characters. The host data include $q,d>4tr+1$.
\end{proof}

For $\operatorname{HostGood}(s,C_8,C_5,c_6,C_b,H,q_0)$ put $N=Q_Sqd$.
The five level primes have the size bounds and $\mathcal B_r$-disjointness
of \ref{B2.51}. Use the thick bounds of \ref{B2.53} at
$L_w=10\lceil405\log x\rceil+2\lceil30/\varepsilon\rceil+22$, and
the top size and $\operatorname{BadOK}(s,C_b,H)$ bounds of \ref{B2.54}.
Require \eqref{eq:thin-character-bound} with constant $c_6$ and
$\operatorname{QDq}$ with constant $C_5$.

The primes $p,p'$ split completely in $E$, and all level primes split
in $K=\Q(\sqrt{-d})$. There are $\pi_i\in\mathcal O_K$ with
\begin{equation}\label{eq:level-norm-and-congruence}
 \pi_i\overline{\pi_i}=g_i,\qquad
 \pi_i\equiv1\pmod{Q_Sqd\,\mathcal O_E}\quad(i=0,1).
\end{equation}
The congruences of \ref{B2.51} use $\operatorname{lcm}(15,N)$ as the
modulus for $p,p'$. Require $|d\log H|\leq x^\varepsilon$, where
$d\log H$ is the real logarithmic derivative at $1$ of the
$d$-quadratic character. Take $\mathfrak c_K=\mathcal O_K$, splitting
bound at least $\lfloor x^{\varepsilon_z/5}\rfloor$, $J,o_\phi\leq C_8$,
and $Q_S\leq C_8n$. At a common $Y$, require $\operatorname{PairWin}$,
$\operatorname{PairRow}$, $\operatorname{Boxq}$, and the boxes of
\ref{B2.52} at height $H_{Q_{\mathrm{odd}}}(j)$ for the primitive real
host characters with conductor divisible by $q$ or $d$. As in
\ref{B2.52}, $j=\mathbf1_{d\mid N_\chi}$ and the conductor is at most
this height.

Write $\operatorname{admitted}(R,H_0)$ if for every $B_0$ and every
record of level primes, auxiliary ideal and positions satisfying the
choice conditions, $\operatorname{FieldRows}(H_0,B_0)$ and $q,d>4tr+1$
give an extension with that record and splitting bound $B_0$.

\begin{theorem}[Completing a fixed arithmetic datum]\label{B2.56}
Fix $0<\varepsilon<1/52$ and $C_0\geq1$. There are real constants
$C_{T_2},C_8,c_6,C_b,x_0$, with $C_8\geq C_0$, $c_6>0$ and $C_b>0$,
with the following property.
Let $s=(x,\varepsilon,J)$, $R$, an arithmetic datum $H_0$ on its
frame, an integer $M\geq1$, and $Y\in\R$ satisfy
\[
 \begin{gathered}
 x\geq x_0,\quad\operatorname{RankVal}(s,R),\quad
 J,o_\phi\leq C_0,\quad Q_S\leq C_0n,\quad M\leq C_0,\\
 Q_S\mid\operatorname{lcm}(2^{e+2},M),\quad
 \operatorname{PairWin}(s,Y),\quad
 \operatorname{PairGood}(M,\varepsilon,Y,x,q,d),\\
 \operatorname{FieldRows}(H_0,\lfloor Y^{\varepsilon/5}\rfloor),\qquad
 \operatorname{admitted}(R,H_0).
 \end{gathered}
\]
There are a host $H$ extending exactly $H_0$ with thickness constant
$C_{T_2}$ and a prime $q_0$ such that
\[
 \operatorname{HostGood}(s,C_8,4,c_6,C_b,H,q_0),\qquad
 \operatorname{Shares}(TT,x^{-1+\varepsilon}).
\]
\end{theorem}
\leanref{UpperTheory.\allowbreak{}Top.\allowbreak{}hostOfField}{Development.\allowbreak{}AutomorphicArithmetic.\allowbreak{}CharacterBoxes.\allowbreak{}HostOfField}{8--16}
\leanref{UpperTheory.\allowbreak{}T3.\allowbreak{}hostOfFieldOf}{Development.\allowbreak{}AutomorphicArithmetic.\allowbreak{}ShortNorms.\allowbreak{}hostOfFieldOf}{159--256}

\begin{proof}
Use the constants of \ref{B2.51}--\ref{B2.55}, with $C_{T_2}=64$.
Take $C_8\geq C_0$ and a common threshold.

Put $N=Q_Sqd$ and $z=Y^{\varepsilon/5}$. The ranks bound $r,t$ polynomially in $x$.
The lower bounds are $\log q>z/2$ and $z\geq x^{\varepsilon/(5C_0)}$. Together with the lower bound for $d$
these give $d\geq100$ and $q,d>4tr+1$, as in \Lean{UpperTheory.\allowbreak{}T3.\allowbreak{}qdBig}.
For an initial host over $H_0$, use admissibility with two primes
greater than $Q_S$, two larger primes $1\bmod15$, the unit ideal,
empty position sets and thickness length $2$. Its pair row from
\ref{B2.55} gives $\log N\leq2z(\log Y)^2\leq2x^{\varepsilon/5}(\log x)^2$.
Choose the level primes by \ref{B2.51}.

For $\log x\geq1$, the bound $L_w\leq(4084+60/\varepsilon)\log x$
allows us to choose $T_2$ by \ref{B2.53} with $C_o=\max(C_0,2)$ and $L=L_w$.
By \eqref{eq:norm-one-coordinates}, the generators satisfy
\eqref{eq:level-norm-and-congruence} and $g_i\equiv1\pmod N$.
The five primes are coprime to $N$. Apply admissibility to the
choice record with $\mathfrak c_K=\mathcal O_K$ and splitting bound
$\lfloor z\rfloor\geq\lfloor x^{\varepsilon_z/5}\rfloor$ to obtain $D$.
Since each prime is $1\bmod N$, it splits completely in $\Q(\zeta_N)$,
hence in $E$ and $K$. Apply \ref{B2.55} and \ref{B2.54} to $D$ and
use the resulting $TT$.

The modulus conditions give $Q_{\mathrm{odd}}\mid M$ and ensure that
every prime divisor of $Q_S$ divides $2M$. For the primitive real
characters with $q\mid N_\chi$ or $d\mid N_\chi$, \ref{B2.52} gives
$N_\chi\leq H_{Q_{\mathrm{odd}}}(j)$. Restrict the box at $H_M(j)$
to this height at the same $Y$. Since $Q_{\mathrm{odd}}\leq C_0$,
\ref{B2.52} and the bound on $\log N$ give the thin estimate.
$\operatorname{QDq}$ gives $|d\log H|\leq4\log x\leq x^\varepsilon$,
and $\operatorname{lcm}(15,N)\mid 15N$ gives the required congruences.
\end{proof}

\begin{theorem}[Admissible ranks on a prescribed frame]\label{B2.57}
For every $0<\varepsilon<1/52$ there is $x_0$ such that for
$s=(x,\varepsilon,J)$ with $x\geq x_0$, a frame $\phi$ satisfying
$x^{2+8\varepsilon}\leq|\phi.A|\leq2x^{2+8\varepsilon}$, and any
arithmetic datum $H_0$ on $\phi$, there is a rank datum $R$ on
$\phi$ satisfying $\operatorname{RankVal}(s,R)$ and
$\operatorname{admitted}(R,H_0)$.
\end{theorem}
\leanref{UpperTheory.\allowbreak{}Top.\allowbreak{}admissibleClassesTwistedChar}{Development.\allowbreak{}AutomorphicArithmetic.\allowbreak{}CoprimeRanks.\allowbreak{}AdmChi}{13--21}

\begin{proof}
Choose ranks on $\phi$ by \ref{B8.59}. Fix a choice record and $B_0$
as in admissibility. The field rows give the square root of $-q$
with its Galois action. The arithmetic datum supplies an embedding
in the chosen CM type and an ideal $\mathfrak b$ stable under
$\Gal(E/K)$ that corresponds to $\phi.\mathrm{Psi}$ in the chosen
coordinates, with
$\mathfrak D_{E/K}=\mathfrak b\,\overline{\mathfrak b}$ and
$\mathfrak b+\overline{\mathfrak b}=\mathcal O_E$.
Adjoin these choices and set the splitting bound to $B_0$.
\leanref{HostParam.\allowbreak{}Seam.\allowbreak{}exists\_\allowbreak{}host}{Development.\allowbreak{}AutomorphicArithmetic.\allowbreak{}CoprimeRanks.\allowbreak{}SecondSetNoRelation}{100--105}
\end{proof}

\section{The arithmetic conditions and the window supply}

\begin{definition}[The arithmetic requirement at a host]\label{B2.58}
Let $\Lambda$ be a level, $\phi$ a frame, and $H$ a raw host on
$\phi$. Take $TT\subset A$ and $x,\varepsilon,C\in\R$.
Put $S=\Lambda.\mathrm{KEPT}$ and let $\mathrm{Kept}_H(S)$ consist
of host characters whose primitive representative belongs to $S$.
Set $\mathcal B_H(S)=(X_+\setminus\mathrm{Kept}_H(S))\cup X_-$.
The predicate $\operatorname{EAtH}(\Lambda,H,TT,x,\varepsilon,C)$
requires $[A:A_0]=\Lambda.o$, a representative in $X_+$ for each
$\chi\in S$, and $\operatorname{Shares}(TT,x^{-1+\varepsilon})$.
If $\Lambda.\mathrm{psK}$ designates a character, it must represent
$\mathrm{psC}$. Otherwise $\mathrm{psC}$ lies outside the retained
family. For every $\theta\in\mathcal B_H(S)$ and every primitive
representative $\chi$ of $\theta$, require
$|g_\chi(x)|\leq x^{1/2+\varepsilon}$ and $|(L'/L)(1,\chi)|\leq C\log x$.

For an arithmetic datum $H_0$, define the cyclic host requirement
\[
 \operatorname{GChi}(\psi)(s,H_0,TT,t)
 \Longleftrightarrow
 \operatorname{EAtH}(\operatorname{Level.cyc}(\psi),
                  H_0.\mathrm{toHostRaw},TT,x,\varepsilon,4)
 \ \text{and}\ 
 \operatorname{VVle}(E,t,\sqrt b).
\]
\end{definition}
\leanref{Formalism.\allowbreak{}HC.\allowbreak{}EAtH}{Formalism.\allowbreak{}SharpRestriction.\allowbreak{}RamificationFrobenius.\allowbreak{}TotientGaloisClasses.\allowbreak{}ExplicitFormulaErrorOfHostConditions}{144--152}
\leanref{UpperTheory.\allowbreak{}Top.\allowbreak{}GChi}{Development.\allowbreak{}AutomorphicArithmetic.\allowbreak{}CharacterBoxes.\allowbreak{}Supply}{42--44}

The assertion $\operatorname{HostExW}(\lambda,P)$ means that for
each $0<\varepsilon<1/52$ there are constants $C_{T_2},C_8,C_5,c_6,C_b,x_0$,
with $C_8\geq1$ and $C_5,c_6,C_b>0$, a number $J\geq1$, and a set
$\mathcal W\subset\R$ such that
\begin{equation}\label{eq:window-set-property}
 \forall y\in\R\ \exists Z\geq y,\qquad
 [Z,Z^\lambda]\subset\mathcal W.
\end{equation}
For every scale $s=(x,\varepsilon,J)$ with $x\geq x_0$ in $\mathcal W$,
there must be $R,H,q_0$ satisfying $\operatorname{RankVal}(s,R)$,
$H.C_{T_2}=C_{T_2}$, $\operatorname{HostGood}(s,C_8,C_5,c_6,C_b,H,q_0)$,
and $P(C_5,s,R,H,q_0)$.
\leanref{GlobalTheory.\allowbreak{}EX.\allowbreak{}HostExW}{Development.\allowbreak{}AbelianVarietyTheory.\allowbreak{}Definitions.\allowbreak{}GlobalTwisting.\allowbreak{}ThinScalesIffHostBound}{120--124}

For a predicate $F$ on scales, frames, arithmetic data, top sets and
auxiliary ranks, put
\[
 F.\mathrm{req}(C_5,s,R,H,q_0)=F(s,R.\mathrm{toFrame},H.\mathrm{toHostAx},H.TT,R.t).
\]

\begin{theorem}[Assembling the data on one window]\label{B2.59}\label{B2.60}
Fix a primitive $\psi$ and assume $\operatorname{VWinAt}(\psi,401)$.
For every $0<\varepsilon<1/52$ assume the completion assertion of
\ref{B2.56} for every $C_0\geq1$. Suppose also that for each
$0<\varepsilon<1/52$ there is $x_0(\varepsilon)$ such that, at
every scale $s=(x,\varepsilon,J)$ with $x\geq x_0(\varepsilon)$,
every datum $H_0$ with $M\mid Q_S$ and \eqref{eq:cyclic-artin-kernel},
on a frame with
\[
 x^{2+8\varepsilon}\leq n\leq2x^{2+8\varepsilon},\qquad o_\phi=o_\psi,
\]
admits ranks on that frame satisfying $\operatorname{RankVal}(s,R)$
and $\operatorname{admitted}(R,H_0)$. Then
\[
 \operatorname{HostExW}\left(\frac32,
                    (\operatorname{GChi}(\psi)).\mathrm{req}\right)
\]
holds with $C_5=4$.
\end{theorem}
\leanref{UpperTheory.\allowbreak{}T3.\allowbreak{}supplyOfAdmOf}{Development.\allowbreak{}AutomorphicArithmetic.\allowbreak{}ShortNorms.\allowbreak{}supplyOfAdmOf}{13--136}
\leanref{UpperTheory.\allowbreak{}Top.\allowbreak{}supplyOfAdm}{Development.\allowbreak{}AutomorphicArithmetic.\allowbreak{}CharacterBoxes.\allowbreak{}HostOfField}{17--24}

\begin{proof}
Fix $\varepsilon$. Take $c,e_0$ from \ref{B2.46}, and $J$ and its
windows from \ref{B2.48}. Choose $C_0\geq\max(1,J,o_\psi,8M)$.
Take a common threshold $x_1$ for completion, rank supply and
\ref{B6.61}. Increase $x_1$ so that $x>c2^{e_0}$ and
$Y^{\varepsilon/5}\geq\max(5,4M)$ for $x\geq x_1$ and $Y\geq x^{1/J}$.

Let $\mathcal W$ consist of $x$ admitting a
$\operatorname{WinDatum}(\psi,\varepsilon,J,c,x,Y,q,d,e)$.
By \ref{B2.48} it satisfies \eqref{eq:window-set-property} at $\lambda=3/2$.
Fix such an $x\geq x_1$ and its datum. The pair gives
\eqref{eq:host-prime-hyp}, $(M,qd)=1$, and $(-d/p)=1$ for odd primes
$p\leq Y^{\varepsilon/5}$. Its size gives $q\ne3$, and
$c2^e\geq x^{2+8\varepsilon}>c2^{e_0}$ gives $e\geq e_0$.

Apply \ref{B2.46} with $B_0=\lfloor Y^{\varepsilon/5}\rfloor$ to get
$H_0$. The rank supply gives $\operatorname{RankVal}(s,R)$ and
$\operatorname{admitted}(R,H_0)$ on its frame. Complete with $Y$ and
$C_0$. The host has the required shares and inherits the character
data of \ref{B2.46}. Each prime divisor of $Q_S$ divides $2M$, so
\ref{B2.47} gives the bounds in $\operatorname{EAtH}$ with constant $4$.

Restriction from $\Q(\zeta_{Q_Sqd})$ is the Artin map. Its kernel
gives $\operatorname{IsChiField}(\psi,e,q,d,E)$. The window datum
gives $\operatorname{VVle}$ with cutoff $\lfloor x^{400}\rfloor$
and budget $\sqrt b$. By \ref{B6.61}, $t\leq x^{400}$, so
$\operatorname{VVle}(E,t,\sqrt b)$ follows. Take $C_5=4$ and $x_0=x_1$.
\end{proof}

\begin{theorem}[Arithmetic hosts on arbitrarily distant windows]\label{B2.61}
For every primitive Dirichlet character $\psi$,
\[
 \operatorname{VWinAt}(\psi,401)
 \ \Longrightarrow\
 \operatorname{HostExW}\left(
       \frac32,(\operatorname{GChi}(\psi)).\mathrm{req}\right).
\]
For each $0<\varepsilon<1/52$, one $J\geq1$ works on arbitrarily
distant intervals $[Z,Z^{3/2}]$. Every $x$ in those intervals admits
a host at scale $(x,\varepsilon,J)$ with $C_5=4$ and both conditions
of Definition~\ref{B2.58}.
\end{theorem}
\leanref{UpperTheory.\allowbreak{}Top.\allowbreak{}supplyChi}{Development.\allowbreak{}AutomorphicArithmetic.\allowbreak{}CharacterBoxes.\allowbreak{}SupplyChi}{8--17}

\begin{proof}
Apply \ref{B2.59} with completion from \ref{B2.56} and rank supply
from \ref{B2.57}.
\end{proof}


\part{Divisibility from one theta lift}

% ---- p3c1 ----
\chapter{The local kernel and the differential operator}\label{ch:p3c1}

Fix a prime $p$. For a finitely supported multi-index $n$, write $|n|=\sum_xn_x$ and $n!=\prod_xn_x!$. Use the coefficient notation $[Z^n]f$. Substitutions have zero constant term.

\section{Local data and weight diagrams}

\begin{definition}[Local systems and lattice bounds]\label{B3.01}
Let $V$ be a finite set of places and $O$ a discrete valuation ring with structure map $\iota:O\to K$ to a field. Take $O$-submodules $\Lambda_j$ in $K$-vector spaces $\mathcal M_j$ and $K$-linear lifts $K_j:\mathcal M_j\to\mathcal M_{j+1}$.
Fix $k_1,t,c,m,r,e_O\in\N$ and $d_G:\N\to\N$. Require $p$ to be associated in $O$ to $\varpi^{e_O}$ for an irreducible $\varpi\in O$. The stabilizer $\Gamma_{j,\xi,L}$ of each lattice tuple $L$ in level component $\xi$ is finite for $0\le j\le m$ and satisfies $v_p(|\Gamma_{j,\xi,L}|)\le d_G(j)$. Define
\[
 v_j(f)\ge g\quad\Longleftrightarrow\quad
 f\in\iota(\varpi)^{\lceil e_Og\rceil}\Lambda_j
 \qquad(f\in\mathcal M_j,\ g\in\R).
\]
At each $v\in V$, require $p\ge5$, $m,r\ge1$ and $2r\le k_1+1$. Use positions $s=0,\ldots,n'_v-1$ modulo an even $n'_v\ge2$. Choose marked positions $T_v$ and integer weights $\lambda_{s,1}\ge\cdots\ge\lambda_{s,k_1+1}$. For $s\in T_v$,
\[
 \lambda_{s,1}=a_v(s),\quad \lambda_{s,k_1+1}=-b_v(s),
 \quad |\lambda_{s,i}|\le c\quad(1<i<k_1+1),
\]
where $a_v,b_v$ are integer functions. For $s\notin T_v$, all $|\lambda_{s,i}|\le c$. Require the integer functions $w_j$ to satisfy
\[
 w_1(s)\ge\max\{c,-\lambda_{s,k_1+1}\},\qquad
 w_l(s)\le w_{l+1}(s)\quad(1\le l<m).
\]
Choose the dual bases of \ref{B3.08}.
\end{definition}
\leanref{Formalism.\allowbreak{}MainNotions.\allowbreak{}FormalGroupSystem.\allowbreak{}Std}{Formalism.\allowbreak{}SharpRestriction.\allowbreak{}GrowthBounds.\allowbreak{}FormalGroupSystem}{89--102}
\leanref{Formalism.\allowbreak{}MainNotions.\allowbreak{}FormalGroupSystem.\allowbreak{}Sys}{Formalism.\allowbreak{}SharpRestriction.\allowbreak{}GrowthBounds.\allowbreak{}FormalGroupSystem}{264--278}
\leanref{Formalism.\allowbreak{}MainNotions.\allowbreak{}FormalGroupSystem.\allowbreak{}Sys.\allowbreak{}vsc}{Formalism.\allowbreak{}SharpRestriction.\allowbreak{}GrowthBounds.\allowbreak{}FormalGroupSystem}{309--310}
\leanref{Formalism.\allowbreak{}MainNotions.\allowbreak{}FormalGroupSystem.\allowbreak{}Sys.\allowbreak{}vGe}{Formalism.\allowbreak{}SharpRestriction.\allowbreak{}GrowthBounds.\allowbreak{}FormalGroupSystem}{311--312}
\leanref{Formalism.\allowbreak{}MainNotions.\allowbreak{}FormalGroupSystem.\allowbreak{}Pend}{Formalism.\allowbreak{}SharpRestriction.\allowbreak{}GrowthBounds.\allowbreak{}FormalGroupSystem}{447--494}

Suppress $v$ below.

\begin{definition}[Layer weights and remaining columns]\label{B3.02}
Put $N_j=k_1+1+j(t+1)$. For $1\le i\le N_j$, set
\[
 \nu_{j,s,i}=
 \begin{cases}
 \lambda_{s,i}+w_j(s),&i\le k_1+1,\\
 w_j(s)-w_b(s),&
 i=k_1+1+(b-1)(t+1)+h,\quad 1\le b\le j,\ 1\le h\le t+1.
 \end{cases}
\]
Set $\nu_{j,s,i}=0$ outside this range. For $1\le j\le m$ the first block decreases to at least $w_j-w_1$, and the appended blocks decrease to zero. Thus the rows form a Young diagram. Put
\[
 \alpha_s=\begin{cases}\lambda_{s,1}-\lambda_{s,2},&s\in T,\\0,&s\notin T,\end{cases}
 \qquad
 C_{j,s}=\begin{cases}\nu_{j,s,2},&s\in T,\\\nu_{j,s,1},&s\notin T,\end{cases}
 \qquad
 L_{j,s}=\sum_{i=1}^{N_j}\nu_{j,s,i}-\alpha_s.
\]
Removing the $\alpha_s$ height-one columns leaves $C_{j,s}$ columns and $L_{j,s}$ boxes.
\end{definition}
\leanref{Formalism.\allowbreak{}MainNotions.\allowbreak{}FormalGroupSystem.\allowbreak{}NN}{Formalism.\allowbreak{}SharpRestriction.\allowbreak{}GrowthBounds.\allowbreak{}FormalGroupSystem}{26--27}
\leanref{Formalism.\allowbreak{}MainNotions.\allowbreak{}FormalGroupSystem.\allowbreak{}nu}{Formalism.\allowbreak{}SharpRestriction.\allowbreak{}GrowthBounds.\allowbreak{}FormalGroupSystem}{28--32}
\leanref{Formalism.\allowbreak{}MainNotions.\allowbreak{}FormalGroupSystem.\allowbreak{}al}{Formalism.\allowbreak{}SharpRestriction.\allowbreak{}GrowthBounds.\allowbreak{}FormalGroupSystem}{33--35}
\leanref{Formalism.\allowbreak{}MainNotions.\allowbreak{}FormalGroupSystem.\allowbreak{}cap}{Formalism.\allowbreak{}SharpRestriction.\allowbreak{}GrowthBounds.\allowbreak{}FormalGroupSystem}{36--39}
\leanref{Formalism.\allowbreak{}MainNotions.\allowbreak{}FormalGroupSystem.\allowbreak{}lsz}{Formalism.\allowbreak{}SharpRestriction.\allowbreak{}GrowthBounds.\allowbreak{}FormalGroupSystem}{40--43}

\begin{definition}[Digit loss]\label{B3.03}
For $z\in\Z$, set $\ell_p(z)=0$ if $z\le0$, and $\ell_p(z)=1+\lfloor\log_pz\rfloor$ otherwise. For $1\le j\le m$, put $r_0=N_{j-1}$, $r_1=N_j$, and
\[
 \delta_j=\sum_{s\in\Z/n'\Z}
 \bigl(r_1\ell_p(\alpha_s)+r_0r_1\ell_p(L_{j,s})\bigr)
 +\sum_{s\in T}\ell_p(\alpha_s)^2.
\]
When $\alpha$ and $L_{j,\bullet}$ are natural-valued, $\delta_j=\operatorname{deLoss}(p,r_0,r_1,T,\alpha,L_{j,\bullet})$.
\end{definition}
\leanref{Formalism.\allowbreak{}MainNotions.\allowbreak{}FormalGroupSystem.\allowbreak{}lgp}{Formalism.\allowbreak{}SharpRestriction.\allowbreak{}GrowthBounds.\allowbreak{}FormalGroupSystem}{44--45}
\leanref{Formalism.\allowbreak{}MainNotions.\allowbreak{}FormalGroupSystem.\allowbreak{}de}{Formalism.\allowbreak{}SharpRestriction.\allowbreak{}GrowthBounds.\allowbreak{}FormalGroupSystem}{46--50}
\leanref{Formalism.\allowbreak{}InvariantFactorPaths.\allowbreak{}G7.\allowbreak{}deLoss}{Formalism.\allowbreak{}ThetaLogCoordinates.\allowbreak{}UniformThetaSubstitution.\allowbreak{}StatementAB}{100--102}

\section{Formal parameters and admissible kernels}

\begin{definition}[Parameters and phase equivalence]\label{B3.04}
Let $W=W(\overline{\F}_p)$ and $K_0=\operatorname{Frac}(W)$. Let $\sigma$ be Witt Frobenius extended to $K_0$. Put
\[
 \ell(T)=\sum_{b\ge0}\frac{T^{p^{2b}}}{p^b},
 \qquad f^\ell(X)=f((\ell(X_i))_i).
\]
For $f(0)=0$, define
\[
 [X^n]\exp(f)=\sum_{k=0}^{|n|}\frac{[X^n]f^k}{k!}.
\]
For $f,g\in K_0[[X]]$ in finitely many variables, set
\[
 f\simeq g\quad\Longleftrightarrow\quad
 (f-g)(0)=0,\qquad \exp((f-g)^\ell)\in1+(X)W[[X]].
\]
Use the same substitution over field extensions $K_1/K_0$.
\end{definition}
\leanref{Formalism.\allowbreak{}MainNotions.\allowbreak{}FormalGroupSystem.\allowbreak{}logF}{Formalism.\allowbreak{}SharpRestriction.\allowbreak{}GrowthBounds.\allowbreak{}FormalGroupSystem}{126--128}
\leanref{Formalism.\allowbreak{}MainNotions.\allowbreak{}FormalGroupSystem.\allowbreak{}inParam}{Formalism.\allowbreak{}SharpRestriction.\allowbreak{}GrowthBounds.\allowbreak{}FormalGroupSystem}{129--131}
\leanref{Formalism.\allowbreak{}MainNotions.\allowbreak{}FormalGroupSystem.\allowbreak{}inParamK}{Formalism.\allowbreak{}SharpRestriction.\allowbreak{}GrowthBounds.\allowbreak{}FormalGroupSystem}{132--135}
\leanref{Formalism.\allowbreak{}MainNotions.\allowbreak{}FormalGroupSystem.\allowbreak{}expS}{Formalism.\allowbreak{}SharpRestriction.\allowbreak{}GrowthBounds.\allowbreak{}FormalGroupSystem}{136--138}
\leanref{Formalism.\allowbreak{}MainNotions.\allowbreak{}FormalGroupSystem.\allowbreak{}Simeq}{Formalism.\allowbreak{}SharpRestriction.\allowbreak{}GrowthBounds.\allowbreak{}FormalGroupSystem}{139--143}

\begin{definition}[The local theta condition]\label{B3.05}
For even $n'$ and $r_0,r_1\ge0$, put
\[
 E=\{0,\ldots,n'/2-1\}\times\{0,\ldots,r_0-1\}\times\{0,\ldots,r_1-1\}.
\]
A pairing $b\in XYK_0[[X,Y]]$ satisfies
\[
 \begin{split}
 b(X+X',Y)-b(X,Y)-b(X',Y)&\simeq0,\\
 b(X,Y+Y')-b(X,Y)-b(X,Y')&\simeq0,\\
 b(\eta X,Y)&\simeq b(X,\sigma(\eta)Y)
 \quad(\eta\in W,\ \sigma^2\eta=\eta).
 \end{split}
\]
For fixed $b$ and $R_1\subset K_1$ a subring, the local condition on $\Theta\in K_1[[X_e,Y_e]]$ requires $\Phi\in K_0[[X_e,Y_e]]$ and $U\in K_1[[X_e,Y_e]]$ such that
\[
 \Theta=\exp(-\Phi)U,\quad U^\ell\in R_1[[X,Y]],\quad
 \Phi(-X,-Y)\simeq\Phi(X,Y),\quad
 \Phi\simeq\sum_{e\in E}b(X_e,Y_e).
\]
For a framed lattice pair at levels $j,j+1$, let $\Theta_A$ be its generator kernel and $\Theta_v$ its restriction with variables outside $v$ set to zero. Use the displayed factorization for $\Theta_v$ on $E_v$ with $(r_0,r_1)=(N_j,N_{j+1})$, and for $\Theta_A$ on $\coprod_v E_v$. Both use the system's $b$ and $(K_1,R_1)=(K_c,O_c)$, the kernel coefficient field and its valuation subring.
\end{definition}
\leanref{Formalism.\allowbreak{}MainNotions.\allowbreak{}FormalGroupSystem.\allowbreak{}BeProp}{Formalism.\allowbreak{}SharpRestriction.\allowbreak{}GrowthBounds.\allowbreak{}FormalGroupSystem}{147--159}
\leanref{Formalism.\allowbreak{}MainNotions.\allowbreak{}FormalGroupSystem.\allowbreak{}ExtensionLocalCondition}{Formalism.\allowbreak{}SharpRestriction.\allowbreak{}GrowthBounds.\allowbreak{}FormalGroupSystem}{162--171}
\leanref{Formalism.\allowbreak{}MainNotions.\allowbreak{}FormalGroupSystem.\allowbreak{}ExtensionGeneratorCondition}{Formalism.\allowbreak{}SharpRestriction.\allowbreak{}GrowthBounds.\allowbreak{}FormalGroupSystem}{237--244}
\leanref{Formalism.\allowbreak{}MainNotions.\allowbreak{}FormalGroupSystem.\allowbreak{}adm2}{Formalism.\allowbreak{}SharpRestriction.\allowbreak{}GrowthBounds.\allowbreak{}FormalGroupSystem}{548--551}
\leanref{Formalism.\allowbreak{}MainNotions.\allowbreak{}FormalGroupSystem.\allowbreak{}adm2b}{Formalism.\allowbreak{}SharpRestriction.\allowbreak{}GrowthBounds.\allowbreak{}FormalGroupSystem}{552--555}

\begin{definition}[Admissibility of the kernel and lifts]\label{B3.05a}
Let $X$ be the system of \ref{B3.01} and $Y$ its kernel data. The pair is admissible if it satisfies the conditions below. For every $0\le j<m$, the tests \(\Lean{adm2}\), \(\Lean{adm2b}\) and \(\Lean{adm3}\) apply to every pair of component--lattice data at levels $j$ and $j+1$ and every frame between them.
\begin{enumerate}
\item \(\Lean{adm1}\) requires Lemma~\ref{B3.06} for every odd prime. The tests \(\Lean{adm2}\) and \(\Lean{adm2b}\) require the respective factorizations of \ref{B3.05} for each one-place restriction and for the generator on the full disjoint union. In \(\Lean{adm2c}\) the chosen \(b\) satisfies the pairing laws of \ref{B3.05}.
\item In \(\Lean{adm3}\) each coefficient in (3.29.1) of \ref{B3.29} satisfies $\operatorname{ge}_p(-,0)$ of \ref{B3.12} for all column choices at zero neighbour indices. In \(\Lean{adm3b}\), \(K_j(K\Lambda_j)\subseteq K\Lambda_{j+1}\).
\item \(\Lean{ArchimedeanSignatureMatching}\) requires (3.30.4) for \(j+1<m\) and (3.30.5) for \(j+1=m\) from \ref{B3.30}, on every form in \(K\Lambda_j\) at every target lattice tuple. For each $j$, choose \(c_K\) in advance for all such forms and targets, with \(c_K,c_K^{-1}\in O_c\) and \(0^{-1}=0\). The terminal denominators $w_x$ are positive with $v_p(w_x)\le d_G(j)$. At $j+1=m$, require the second identity of Lemma~\ref{B3.35c} for every place and every $i\ge0$, with the same $L_{xy},w_x$ as in (3.30.5) and weights $1/w_x$.
\item In \(\Lean{adm5}\), for each \(v\in V\) and \(0\le j<m\), choose \(a,a'\in K^\times\) of valuation zero. For every lattice pair \(L,L'\), the fixed-coordinate series of Lemma~\ref{B3.35a} satisfies
\[
 \Theta^F_{j,L,L'}=a^{l}(a')^{l'}p^{-(n'_v/2)(N_jl'-N_{j+1}l)}\,\Xi_{v,j,\omega}(\mathbf1_{\{xL_v\subseteq L'_v\}}).
\]
Here $l,l'$ are the local lengths and $\Xi_{v,j,\omega}$ is the additive map for the pair's datum $\omega$. In \(\Lean{MinkowskiCounting}\) the source level group preserves the kernel value under $(L,G)\mapsto(\gamma L,\chi(\gamma)\gamma G)$ for source weight vectors $G$.
\item For \(0\le j\le m\), \(\Lean{adm6}\) requires the first identity of Lemma~\ref{B3.35c} to hold at each place for every \(i\ge0\) and every diagonally invariant \(K_c\)-valued edge function. The left sum fixes the source representative and the right sum fixes the target representative.

In \(\Lean{adm7}\), let \(0\le j<m\) and fix lattice data \(L,L'\) at levels \(j,j+1\). Take \(0\le i_v\le r\) and a target \(L''\) in the component of \(L'\), equal to \(L'\) at places where \(i_v=0\) and a lower \(i_v\)-neighbour elsewhere. For each such \(L,L',L'',\mathbf i\), there exists a frame \(\phi\) from \(L\) to \(L''\) that works for every series \(f\in K_c[[Z]]\) and every real \(g\). If \(\operatorname{ge}_p([\tau^n]\mathcal D_{P_{I,\mathbf i}}(f),g)\) holds for all \(I,n\), then for every \(G\) in the source weight lattice \(H_j(L)\), the value
\[
 \left(\prod_v p^{N'_{v,j+1}(i_v)}\right) D(L,L'')\,
 \mathcal D_j(G,f(B_\phi Z))
\]
lies in the target weight lattice at \(L'\) with coefficient bound \(g\) as in (3.30.7). Here \(B_\phi\) is the frame matrix, and \(D\) and \(\mathcal D_j\) are as in Lemma~\ref{B3.35b}.
\item In \(\Lean{adm8}\) the source Hecke operators $y_v$ commute and preserve $\Lambda_0$. Condition \(\Lean{adm11}\) requires \(\operatorname{LowAt}(j,v,0)\) of Definition~\ref{B3.12} for every \(v\in V\) and \(0\le j<m\).
\end{enumerate}
\end{definition}
\leanref{Formalism.\allowbreak{}MainNotions.\allowbreak{}FormalGroupSystem.\allowbreak{}admitted}{Formalism.\allowbreak{}SharpRestriction.\allowbreak{}GrowthBounds.\allowbreak{}FormalGroupSystem}{771--773}

\begin{lemma}[Dieudonn\'e--Dwork criterion]\label{B3.06}
Let $p$ be odd, $d\ge0$, and $\Psi\in K_0[[X_1,\ldots,X_d]]$, with $\Psi(0)=0$. Then
\[
 \exp(\Psi)\in W[[X]]
 \quad\Longleftrightarrow\quad
 (\sigma\Psi)(X_1^p,\ldots,X_d^p)-p\Psi(X)\in pW[[X]].
\]
\end{lemma}
\leanref{Formalism.\allowbreak{}MainNotions.\allowbreak{}FormalGroupSystem.\allowbreak{}DworkLemma}{Formalism.\allowbreak{}SharpRestriction.\allowbreak{}GrowthBounds.\allowbreak{}FormalGroupSystem}{225--230}
\leanref{Formalism.\allowbreak{}InvariantFactorPaths.\allowbreak{}G7.\allowbreak{}dwork\_\allowbreak{}multivariable}{Formalism.\allowbreak{}ThetaLogCoordinates.\allowbreak{}UniformThetaSubstitution.\allowbreak{}FrobeniusLiftNorms.\allowbreak{}LocalComponentGL}{448--480}

The one-variable criterion over $\Q_p$ holds for all primes \cite[Lemma 1]{krattenthaler-2015}. See \cite[Lemma 2.3]{hazewinkel-1979} for the functional equation approach.

\begin{proof}
Put $F=\exp\Psi$ and $\varphi Q=(\sigma Q)(X_1^p,\ldots,X_d^p)$. Then $\varphi F/F^p=\exp(\varphi\Psi-p\Psi)$. Since $p^k/k!$ and $p^k/k$ belong to $pW$ for $k\ge1$, exponential and logarithm identify $p(X)W[[X]]$ with $1+p(X)W[[X]]$. For integral $F$, Frobenius gives $\varphi F\equiv F^p\pmod p$. Conversely assume $\varphi F/F^p\in1+p(X)W[[X]]$. Induct on total degree $h\ge1$. Divide $F$ by the integral polynomial formed by its terms of degree $<h$ to obtain $H=1+H_h+\cdots$. By the forward direction the divisor's quotient is in $1+p(X)W[[X]]$, so the quotient for $H$ is too. Its degree-$h$ part is $-pH_h$, since $\varphi(H-1)$ starts in degree $ph$. Thus $H_h$ and the degree-$h$ part of $F$ are integral.
\end{proof}
\leanref{Formalism.\allowbreak{}MainNotions.\allowbreak{}FormalGroupSystem.\allowbreak{}Pf.\allowbreak{}sx\_\allowbreak{}dwork}{Formalism.\allowbreak{}SharpRestriction.\allowbreak{}GrowthBounds.\allowbreak{}MultisetValuation.\allowbreak{}PositivityConditionHolds}{107--113}

\begin{definition}[An expansion with controlled denominators]\label{B3.07}
Let $K'/K_0$ be a field extension and $R'\subset K'$ a subring. Take $Z_{s,a,b}$ with $s\in\Z/n'\Z$, $a<r_0$, $b<r_1$ and $a,b\ge0$. Given $u:\Z/n'\Z\to K'$, put
\[
 S_{s,a,b}=\sum_{l\ge0}(-1)^lZ_{s-2l,a,b}^{p^{2l}},
 \qquad H_u=\sum_su_s\sum_{a,b}S_{s,a,b}^pS_{s+1,a,b}.
\]
For $f,H\in K'[[Z]]$, write $f\in\mathcal C_{R'}(H)$ if some $U_i\in R'[[Z]]$ satisfy
\[
 [Z^n]f=\sum_{i=0}^{|n|}p^{-i}[Z^n](H^iU_i)\quad\text{for every }n.
\]
The normed-field version requires $\|[Z^n]U_i\|\le1$ for all $i,n$.
\end{definition}
\leanref{Formalism.\allowbreak{}MainNotions.\allowbreak{}FormalGroupSystem.\allowbreak{}sSer}{Formalism.\allowbreak{}SharpRestriction.\allowbreak{}GrowthBounds.\allowbreak{}FormalGroupSystem}{200--204}
\leanref{Formalism.\allowbreak{}MainNotions.\allowbreak{}FormalGroupSystem.\allowbreak{}hhS}{Formalism.\allowbreak{}SharpRestriction.\allowbreak{}GrowthBounds.\allowbreak{}FormalGroupSystem}{205--209}
\leanref{Formalism.\allowbreak{}MainNotions.\allowbreak{}FormalGroupSystem.\allowbreak{}InClass}{Formalism.\allowbreak{}SharpRestriction.\allowbreak{}GrowthBounds.\allowbreak{}FormalGroupSystem}{210--215}
\leanref{Formalism.\allowbreak{}InvariantFactorPaths.\allowbreak{}G7.\allowbreak{}InClass}{Formalism.\allowbreak{}ThetaLogCoordinates.\allowbreak{}UniformThetaSubstitution.\allowbreak{}OddLogRelation}{76--80}

\section{Reduction of the phase}

\begin{theorem}[The explicit odd-power phase]\label{B3.08}
Let $p\ge5$ and let $b$ satisfy \ref{B3.05}. There is $c\in W$ fixed by $\sigma^2$ such that
\[
 \exp\bigl(b(\ell(X),\ell(Y))-G_c(X,Y)\bigr)\in1+(X,Y)W[[X,Y]],
\]
where
\[
 A(X,Y)=\sum_{0\le i\le l}\frac{X^{p^{2i}}Y^{p^{2l+1}}}{p^{i+l+1}},
 \qquad G_c(X,Y)=cA(X,Y)-\sigma(c)A(Y,X).
\]
For even $n'\ge2$, choose $\omega_r,\omega_r^*\in W$, $0\le r<n'/2$, fixed by $\sigma^{n'}$, such that
\[
 \sum_{h=0}^{n'/2-1}\sigma^{2h}(\omega_r\omega_{r'}^*)=\delta_{r,r'}.
\]
Put
\[
 L^X_{r,a,b}=\sum_{s\ {\rm even}}\sigma^s(\omega_r^*)Z_{s,a,b},
 \qquad
 L^Y_{r,a,b}=\sum_{s\ {\rm odd}}\sigma^s(\omega_r)Z_{s,a,b}.
\]
Let $K'/K_0$ contain $\pi$ with $\pi^{p^2-1}=p$, and put $R_\pi=W[\pi]\subset K'$. If $\Phi\simeq\sum_e b(X_e,Y_e)$, then
\[
 \exp(-\Phi(\pi L^X,\pi L^Y))
 =\exp\left(-\frac{\pi^{p+1}}pH_u\right)V,\quad
 V\in1+(Z)R_\pi[[Z]],\quad
 u_s=\begin{cases}-\sigma(c),&s\ {\rm even},\\c,&s\ {\rm odd}.\end{cases}
\]
Also $\pi^{(p+1)i}/i!\in R_\pi$ for every $i\ge0$.
\end{theorem}
\leanref{Formalism.\allowbreak{}MainNotions.\allowbreak{}FormalGroupSystem.\allowbreak{}Pf.\allowbreak{}restBeta}{Formalism.\allowbreak{}SharpRestriction.\allowbreak{}GrowthBounds.\allowbreak{}MultisetValuation.\allowbreak{}BiadditiveSecondVariable}{137--154}
\leanref{Formalism.\allowbreak{}MainNotions.\allowbreak{}FormalGroupSystem.\allowbreak{}Pf.\allowbreak{}restGc}{Formalism.\allowbreak{}SharpRestriction.\allowbreak{}GrowthBounds.\allowbreak{}MultisetValuation.\allowbreak{}CommonEigenvectorThird}{21--23}
\leanref{Formalism.\allowbreak{}MainNotions.\allowbreak{}FormalGroupSystem.\allowbreak{}Pf.\allowbreak{}restSExp\_\allowbreak{}of}{Formalism.\allowbreak{}SharpRestriction.\allowbreak{}GrowthBounds.\allowbreak{}MultisetValuation.\allowbreak{}DworkExponentialRestrictions}{95--106}
\leanref{Formalism.\allowbreak{}MainNotions.\allowbreak{}FormalGroupSystem.\allowbreak{}Pf.\allowbreak{}restTwoSums\_\allowbreak{}of}{Formalism.\allowbreak{}SharpRestriction.\allowbreak{}GrowthBounds.\allowbreak{}MultisetValuation.\allowbreak{}TripleSumComm}{205--361}
\leanref{Formalism.\allowbreak{}MainNotions.\allowbreak{}FormalGroupSystem.\allowbreak{}Pf.\allowbreak{}piFact}{Formalism.\allowbreak{}SharpRestriction.\allowbreak{}GrowthBounds.\allowbreak{}MultisetValuation.\allowbreak{}NatCastInvMemRp}{52--71}

\begin{proof}
Write $P\equiv Q$ if $(P-Q)(0)=0$ and $\exp(P-Q)$ has coefficients in $W$. Put
\[
 \beta(X,Y)=b(\ell(X),\ell(Y)),\qquad
 h(T)=\sum_{i\ge0}\frac{T^{p^{2i+1}}}{p^{i+1}},\qquad
 F(X,T)=\ell^{-1}(\ell(X)+\ell(T)).
\]
Since $\ell(T)-p^{-1}\ell(T^{p^2})=T$, Hazewinkel's functional equation lemma with $q=p^2$ and $s_1=p^{-1}$ gives $F\in\Z_p[[X,T]]$ \cite[Lemma 2.3, (2.5)]{hazewinkel-1979}. Put
\[
 \delta f=f(F(X,T))-f(X)-f(T),\qquad
 d_1Q=Q(F(X,T),Y)-Q(X,Y)-Q(T,Y).
\]
Both $d_1\beta\equiv0$ and $d_1(\beta(Y,X))\equiv0$.

For a primitive $(p^2-1)$-st root $\zeta\in W$, set
\[
 B(X,Y)=\frac1{p^2-1}\sum_{k=0}^{p^2-2}\beta(\zeta^kX,\zeta^{-pk}Y).
\]
The scalar relation and \ref{B3.06} give $B\equiv\beta$, with support
\[
 B_{ij}\ne0\ \Longrightarrow\ i,j\ge1,\quad i\equiv pj\pmod{p^2-1}.
\]
The average preserves both additivity relations since $F(\eta X,\eta T)=\eta F(X,T)$ for $\eta^{p^2-1}=1$.

Suppose $f(0)=0$ and $\delta f\in W[[X,T]]$. Then
\[
 f=a_\infty h+b_\infty\ell+R,\qquad
 a_\infty\in W,\ b_\infty\in K_0,\ R\in W[[T]].
 \tag{1}
\]
Use the identities $\ell(T^p)=ph(T)$ and $h(T^p)=\ell(T)-T$. $F(X,T)^p-F(X^p,T^p)\in pW[[X,T]]$. These imply
\[
 \delta h=\frac{\ell(F(X,T)^p)-\ell(F(X^p,T^p))}{p}\in W[[X,T]].
\]
For integral $P',D$ of zero constant term and a series $C$ with integral derivative $C'$, the quotient $(C(P'+pD)-C(P'))/p$ is integral by $\binom nk=(n/k)\binom{n-1}{k-1}$ and $p^{k-1}/k\in\Z_p$ for $k\ge1$. Here $\ell'(T)=\sum_{b\ge0}p^bT^{p^{2b}-1}$ is integral, so the estimate applies to $C=\ell$.

For $f_0=f-ah-b\ell$ with $a\in W$, $b\in K_0$, we have $\delta f_0=\delta f-a\delta h\in W[[X,T]]$. If $n\ge2$ and $f_0$ has integral coefficients below degree $n$, take $[X^iT^{n-i}]\delta f_0$ to get
\[
 \binom ni[T^n]f_0\in W\qquad(0<i<n).
\]
If $n$ has a nonzero base-$p$ digit in position $k$ with $p^k<n$, then $\binom n{p^k}$ is a unit. For $n=p^{k+1}$, $\binom n{p^k}$ has valuation one.
Start with $a_0=0,\ b_0=[T]f$ and set
\[
 \begin{aligned}
 \mu_j&=p[T^{p^{2j+1}}](f-a_jh-b_j\ell),
 &a_{j+1}&=a_j+p^j\mu_j,\\
 \nu_j&=p[T^{p^{2j+2}}](f-a_{j+1}h-b_j\ell),
 &b_{j+1}&=b_j+p^j\nu_j.
 \end{aligned}
\]
Thus $\mu_j,\nu_j\in W$. Each correction cancels the exceptional coefficient and preserves integrality below it. The unit case covers intervening degrees. The increments lie in $p^jW$, so $a_j,b_j$ converge and the remainder is integral by closedness.

Logarithmic derivatives of integral units are integral. Use $\partial_2F(X,0)=\ell'(X)^{-1}$ to differentiate $d_1B\equiv0$ at $T=0$. This gives
\[
 \partial_XB-\ell'(X)f(Y)\in W[[X,Y]],\qquad f(Y)=[X]B.
\]
The other relation gives $\delta f\in W[[X,T]]$. In (1), the support condition at degrees $p^{2k}$ gives $b_\infty/p^k\in W$ for every $k$, hence $b_\infty=0$. Thus, for $c,c'\in W$,
\[
 \partial_XB-c\ell'(X)h(Y)\in W[[X,Y]],\qquad
 \partial_YB-c'\ell'(Y)h(X)\in W[[X,Y]].
\]
Put
\[
 D_\sigma Q(\mathbf X)=\frac{(\sigma Q)(\mathbf X^p)-pQ(\mathbf X)}p,\qquad
 J=D_\sigma B-\sigma(c)h(X)\ell(Y)-\sigma(c')\ell(X)h(Y).
\]
Here $\mathbf X^p$ raises each variable to the $p$-th power. Then
\[
 d_1(D_\sigma B)-\sigma(c)\delta h(X,T)\ell(Y)\in W[[X,T,Y]].
 \tag{2}
\]
Set $P=F(X,T)^p$, $P'=F(X^p,T^p)$, and
\[
 C(U,Y)=(\sigma B)(U,Y)-\sigma(c)\ell(U)h(Y).
\]
Then $P-P'\in pW[[X,T]]$ and $\partial_UC$ is integral. The Taylor calculation above gives
\[
 \frac{C(P,Y^p)-C(P',Y^p)}p\in W[[X,T,Y]].
\]
Since $(\ell(P)-\ell(P'))/p=\delta h$ and $h(Y^p)=\ell(Y)-Y$, the difference
$d_1(D_\sigma B)-D_\sigma(d_1B)$ equals
$\sigma(c)\delta h\,\ell(Y)$ modulo integral series. By \ref{B3.06}, $D_\sigma(d_1B)$ is integral. This proves (2). Swapping variables proves that both partial coboundaries of $J$ are integral.

Let $Q$ be divisible by $XY$ with both partial coboundaries integral. Then
\[
 Q-\gamma\ell(X)\ell(Y)\in W[[X,Y]]
 \quad\text{for some }\gamma\in K_0.
 \tag{3}
\]
By (1), write $[Y^m]Q=A_mh(X)+B_m\ell(X)$ modulo $W[[X]]$, and $[X^i]Q=A'_ih(Y)+B'_i\ell(Y)$ modulo $W[[Y]]$. Comparing $[X^iY^m]Q$, with $h_i=[T^i]h$ and $\ell_i=[T^i]\ell$, gives
\[
 A_mh_i+B_m\ell_i-A'_ih_m-B'_i\ell_m\in W,\qquad A_m,A'_i\in W.
\]
Fix $i=p^{2k+1}$ and vary $m=p^{2l}$. Then $B'_i/p^l$ stays bounded, so $B'_i=0$. With $m$ fixed and $k$ variable, the relation bounds $A_m/p^{k+1}$, so $A_m=0$. With $i=p^{2k}$, it bounds $B_m/p^k$ when $\ell_m=0$, so $B_m=0$ there. For $m=p^{2l}$, subtract $p^{-l}$ times the relation for $m=1$. This bounds
\[
 \frac{B_{p^{2l}}-p^{-l}B_1}{p^k}
\]
uniformly in $k$, so $B_{p^{2l}}=p^{-l}B_1$. Taking $\gamma=B_1$ proves (3).
The support condition makes $J_{p^{2k},p^{2l}}=0$. Equation (3) forces $\gamma/p^{k+l}\in W$ for all $k,l$, hence $\gamma=0$ and $J$ is integral.

Put $m=p^{2l+1}$. Take $[X^0Y^m]$ in the first derivative relation and $[XY^m]$ in $J$ to get
\[
 B_{1,p^{2l+1}}-\frac{c}{p^{l+1}}\in W,\qquad
 -B_{1,p^{2l+1}}-\frac{\sigma(c')}{p^{l+1}}\in W.
\]
Thus $c+\sigma(c')\in\bigcap_l p^{l+1}W=0$. Swapping variables gives $c'+\sigma(c)=0$, hence $\sigma^2c=c$. Reindexing yields
\[
 A(X^p,Y^p)=p(h(X)\ell(Y)-A(Y,X)),\quad
 A(Y^p,X^p)=p(\ell(X)h(Y)-A(X,Y)).
\]
Therefore $D_\sigma G_c=\sigma(c)h(X)\ell(Y)-c\ell(X)h(Y)$, so
$D_\sigma(B-G_c)=J$ is integral. By \ref{B3.06}, $\beta\equiv B\equiv G_c$.

Set $\lambda_\pi(T)=\pi^{-1}\ell(\pi T)$, $\mu=\lambda_\pi^{-1}$. Then
\[
 \lambda_\pi(T)\equiv T+T^{p^2}\pmod p,\qquad
 \mu(T)\equiv\sum_{l\ge0}(-1)^lT^{p^{2l}}\pmod p.
\]
Since $\lambda_\pi(T)\in T+T^2\Z[[T]]$, coefficient recursion gives $\mu(T)\in T\Z[[T]]$. The identity $\ell(\pi\mu(T))=\pi T$ transports the phase equivalences to $R_\pi$.

Separating $i=l=0$ gives
\[
 G_c(\pi X,\pi Y)=\frac{\pi^{p+1}}p(cXY^p-\sigma(c)X^pY)
 +\pi^{p+1}R(X,Y),\quad R\in(X,Y)R_\pi[[X,Y]].
 \tag{4}
\]
For every other term, $l\ge1$, and
\[
 \frac{p^{2i}+p^{2l+1}-(p+1)}{p^2-1}-(i+l+1)
 =\sum_{a<i}(p^{2a}-1)+\sum_{a<l}(p^{2a+1}-1)-1\ge p-2.
\]
Modulo $p$, the $(t,l)$ term of $\mu(L)$ is $(-1)^l\sigma^{t+2l}(w)Z_{t,a,b}^{p^{2l}}$ for the basis element $w$, by Frobenius. Set $s=t+2l\bmod n'$ and use $\sigma^{n'}w=w$ to obtain
\[
 \mu(L^X_{r,a,b})\equiv\sum_{s\ {\rm even}}\sigma^s(\omega_r^*)S_{s,a,b},
 \qquad
 \mu(L^Y_{r,a,b})\equiv\sum_{s\ {\rm odd}}\sigma^s(\omega_r)S_{s,a,b}.
\]
Set $M_{rh}=\sigma^{2h}(\omega_r)$,
$N_{hr'}=\sigma^{2h}(\omega_{r'}^*)$. Since $MN=I$, also $NM=I$. Apply $\sigma$ for odd indices and use $\sigma^{n'}$-invariance at the wraparound. For even $s$, odd $t$,
\[
 \sum_r\sigma^s(\omega_r^*)\sigma^{t+1}(\omega_r)
 =\begin{cases}1,&s=t+1\text{ in }\Z/n'\Z,\\0,&\text{otherwise}.\end{cases}
\]
Use the transposed identity for the other parity. Since $(\sigma^t(w))^p\equiv\sigma^{t+1}(w)\pmod p$, the two identities give
\[
 \sum_e\bigl(c\mu(L_e^X)\mu(L_e^Y)^p-\sigma(c)\mu(L_e^X)^p\mu(L_e^Y)\bigr)
 =H_u+pY_0,\qquad Y_0\in(Z)R_\pi[[Z]].
\]
Substituting $(\mu(L_e^X),\mu(L_e^Y))$ in (4) and summing gives $\pi^{p+1}H_u/p$ modulo $\pi^{p+1}(Z)R_\pi[[Z]]$.
Write $i!=p^am_i$, $p\nmid m_i$. Since $(p-1)a\le i$,
\[
 \frac{\pi^{(p+1)i}}{i!}=m_i^{-1}\pi^{(p+1)i-(p^2-1)a}\in R_\pi.
\]
Multiply the resulting integral exponential of the error by the phase-equivalence unit to get $V$ with $V(0)=1$.
\end{proof}

\section{Kernel expansions}

\begin{theorem}[Single-place and simultaneous expansions]\label{B3.09}
Assume admissibility in Definition~\ref{B3.05a}. For $j<m$, fix two lattice tuples and a frame. Let $K'/K_c$ contain $\pi$ with $\pi^{p^2-1}=p$. Let $R'\subset K'$ be a subring containing $W[\pi]$ and $O_c$. Write $Z_v=(Z_{v,s,a,b})_{s,a,b}$ and $Z=(Z_v)_{v\in V}$.

For a choice of $u_{v,s}\in R'$, set
\[
 H'_v=H_{u_v}+pG_v,\qquad G_v(0)=0.
\]
At each $v$ the coefficients $u_{v,s}$ can be chosen so that $\Theta_v(\pi L_v^X,\pi L_v^Y)\in\mathcal C_{R'}(H'_v)$ for every $G_v\in R'[[Z_v]]$. For the generator one can choose $u_{v,s}$ so that, for every family $G_v\in R'[[Z]]$, there are $U_{\boldsymbol i}\in R'[[Z]]$, $\boldsymbol i:V\to\N$, satisfying, for every $n$,
\[
 [Z^n]\Theta_A(\pi L^X,\pi L^Y)
 =\sum_{\substack{\boldsymbol i:V\to\N\\i_v\le|n|\ (\forall v)}}
 [Z^n]\left(\left(\prod_v p^{-i_v}(H'_v)^{i_v}\right)U_{\boldsymbol i}\right).
\]
\end{theorem}
\leanref{Formalism.\allowbreak{}MainNotions.\allowbreak{}FormalGroupSystem.\allowbreak{}Pf.\allowbreak{}theta\_\allowbreak{}kernel\_\allowbreak{}expansion}{Formalism.\allowbreak{}SharpRestriction.\allowbreak{}GrowthBounds.\allowbreak{}MultisetValuation.\allowbreak{}CommonEigenvectorThird}{32--34}
\leanref{Formalism.\allowbreak{}MainNotions.\allowbreak{}FormalGroupSystem.\allowbreak{}Pf.\allowbreak{}expansions\_\allowbreak{}multiply}{Formalism.\allowbreak{}SharpRestriction.\allowbreak{}GrowthBounds.\allowbreak{}MultisetValuation.\allowbreak{}CommonEigenvectorThird}{35--37}

\begin{proof}
Proposition~\ref{B8.56c} supplies the local and generator factorizations for the selected arithmetic kernel under its joint comparison hypotheses. In either factorization of Definition~\ref{B3.05}, $U(\pi L)=U^\ell(\pi\mu(L))$ is integral, with $\mu$ as in \ref{B3.08}. Write $q=\pi^{p+1}$. At one place, \ref{B3.08} gives the kernel as $\exp(-qH_u/p)V_0$ with $V_0$ integral. Replacing $H_u$ by $H'$ changes the integral factor to $V_1=\exp(qG)V_0$, since $G(0)=0$ and $q^i/i!\in R'$. The factors in \ref{B3.07} are $U_i=(-q)^iV_1/i!$.

For the generator put $S=\sum_{v,e}b(X_{v,e},Y_{v,e})$. Since $\Phi_A\simeq S$, $\exp((S-\Phi_A)(\pi L))$ is integral by the same substitution. Apply \ref{B3.08} to the summands of $S$ grouped by place and absorb this unit into $V_A$ to get
\[
 \Theta_A(\pi L^X,\pi L^Y)=\left(\prod_v\exp(-qH_{u_v}/p)\right)V_A,
 \qquad V_A\in R'[[Z]].
\]
For $V=\varnothing$ the phase is zero. Take
\[
 U_{\boldsymbol i}=V_A\prod_v\frac{(-q)^{i_v}}{i_v!}\exp(qG_v).
\]
In both expansions, the $i$-th power of a zero-constant series has order at least $i$, which gives the stated coefficient ranges.
\end{proof}

\section{The polynomial and its coefficients}

\begin{definition}[Products of column minors]\label{B3.10}
Fix $0\le j<m$, and put $r_0=N_j,r_1=N_{j+1}$. Use $Z_{v,s,a,b}$ with $0\le a<r_0,0\le b<r_1$, and frame variables $\tau_{v,s,h,a}$, $h\in\N$.
A column choice assigns maps $I_{v,s,c}:\N\to\{0,\ldots,r_1-1\}$. For $0\le c<\nu_{j+1,s,1}$, set
\[
 h_{v,s,c}=\#\{1\le i\le r_1:c<\nu_{j+1,s,i}\},\qquad
 D_{v,s,c}=\det\left(\sum_{a<r_0}\tau_{v,s,h,a}Z_{v,s,a,I_{v,s,c}(h')}\right)_{h,h'<h_{v,s,c}}.
\]
Define
\[
 P_{I,v}=\prod_{s,c}D_{v,s,c}\in\Z[\tau][Z],
\]
where $c$ ranges over the columns at position $(v,s)$. For $i_0:V\to\N$, put
\[
 q_{v,s,c}=\#\{h'<h_{v,s,c}:I_{v,s,c}(h')<i_0(v)\},\quad
 e_{v,s,c}=\begin{cases}
 q_{v,s,c},&s\ {\rm odd},\\
 \max\{\min(h_{v,s,c},i_0(v))-q_{v,s,c},0\},&s\ {\rm even}.
 \end{cases}
\]
Put $P_{I,i_0}=\prod_{v,s,c}p^{e_{v,s,c}}D_{v,s,c}$, so $P_{I,0}=\prod_vP_{I,v}$.
\end{definition}
\leanref{Formalism.\allowbreak{}MainNotions.\allowbreak{}FormalGroupSystem.\allowbreak{}Sys.\allowbreak{}colHt}{Formalism.\allowbreak{}SharpRestriction.\allowbreak{}GrowthBounds.\allowbreak{}FormalGroupSystem}{398--400}
\leanref{Formalism.\allowbreak{}MainNotions.\allowbreak{}FormalGroupSystem.\allowbreak{}Sys.\allowbreak{}colFac}{Formalism.\allowbreak{}SharpRestriction.\allowbreak{}GrowthBounds.\allowbreak{}FormalGroupSystem}{405--411}
\leanref{Formalism.\allowbreak{}MainNotions.\allowbreak{}FormalGroupSystem.\allowbreak{}Sys.\allowbreak{}colExp}{Formalism.\allowbreak{}SharpRestriction.\allowbreak{}GrowthBounds.\allowbreak{}FormalGroupSystem}{412--416}
\leanref{Formalism.\allowbreak{}MainNotions.\allowbreak{}FormalGroupSystem.\allowbreak{}Sys.\allowbreak{}PI}{Formalism.\allowbreak{}SharpRestriction.\allowbreak{}GrowthBounds.\allowbreak{}FormalGroupSystem}{417--420}
\leanref{Formalism.\allowbreak{}MainNotions.\allowbreak{}FormalGroupSystem.\allowbreak{}Sys.\allowbreak{}PIv}{Formalism.\allowbreak{}SharpRestriction.\allowbreak{}GrowthBounds.\allowbreak{}FormalGroupSystem}{421--423}

Each determinant has bidegree $(h_{v,s,c},h_{v,s,c})$. Thus $P_{I,v}$ is homogeneous of degree
\[
 D_v=\sum_s\sum_{i=1}^{N_{j+1}}\nu_{j+1,s,i}
 =\sum_s(\alpha_s+L_{j+1,s})
\]
in each set of variables.

\begin{definition}[Differential evaluation]\label{B3.11}
For $P\in\Z[\tau][Z]$, $f\in K_c[[Z]]$, define
\[
 \mathcal D_P(f)=\sum_{n\in\operatorname{supp}_Z(P)}
 n!\,[Z^n]f\,[Z^n]P\in K_c[\tau].
\]
This is $(P(\partial_Z)f)(0)$ since $\partial^nf(0)=n![Z^n]f$.
\end{definition}
\leanref{Formalism.\allowbreak{}MainNotions.\allowbreak{}FormalGroupSystem.\allowbreak{}applyD}{Formalism.\allowbreak{}SharpRestriction.\allowbreak{}GrowthBounds.\allowbreak{}FormalGroupSystem}{256--260}

\begin{definition}[The single-place coefficient bound]\label{B3.12}
By \Lean{Pend.\allowbreak{}witt\_\allowbreak{}mem\_\allowbreak{}integers\_\allowbreak{}Cp} and \Lean{Pend.\allowbreak{}Oc\_\allowbreak{}K}, $W\subseteq O_c$ and $K\cap O_c=\iota(O)$ inside $K_c$. Define
\[
 \operatorname{ge}_p(z,g)\quad\Longleftrightarrow\quad
 \forall a\in\Z\ \forall b\in\N_{>0},\
 a\le bg\ \Longrightarrow\ z^bp^{-a}\in O_c
 \qquad(z\in K_c,\ g\in\R).
\]
For a real valuation $v$ with ring $O_c$ and $v(p)=1$, this says $v(z)\ge g$.

For $j\in\N$ and $v\in V$, put $\Phi_{j,v}=\sum_{e\in E_v}b(X_{v,e},Y_{v,e})$ in the full variable set.
Use $\mathcal T_j$ for substitution of $(L^X,L^Y)$ at all places. Then $\operatorname{LowAt}(j,v,g)$ asserts that for all $U\in K_c[[X,Y]]$ with $U^\ell\in O_c[[X,Y]]$, column choices $I$, and frame multi-indices $n$,
\[
 \operatorname{ge}_p\left(
 [\tau^n]\mathcal D_{P_{I,v}}\bigl(\mathcal T_j(\exp(-\Phi_{j,v})U)\bigr),g\right).
\]
\end{definition}
\leanref{Formalism.\allowbreak{}MainNotions.\allowbreak{}FormalGroupSystem.\allowbreak{}Pend.\allowbreak{}PhAt}{Formalism.\allowbreak{}SharpRestriction.\allowbreak{}GrowthBounds.\allowbreak{}FormalGroupSystem}{506--509}
\leanref{Formalism.\allowbreak{}MainNotions.\allowbreak{}FormalGroupSystem.\allowbreak{}Pend.\allowbreak{}geC}{Formalism.\allowbreak{}SharpRestriction.\allowbreak{}GrowthBounds.\allowbreak{}FormalGroupSystem}{522--524}
\leanref{Formalism.\allowbreak{}MainNotions.\allowbreak{}FormalGroupSystem.\allowbreak{}Pend.\allowbreak{}LowAt}{Formalism.\allowbreak{}SharpRestriction.\allowbreak{}GrowthBounds.\allowbreak{}FormalGroupSystem}{529--534}

For $g'\le g$, $\operatorname{LowAt}(j,v,g)$ implies $\operatorname{LowAt}(j,v,g')$.


% ---- p3c2 ----
\chapter{Entry estimates: factorials against denominators}\label{ch:b3-entry-estimates}

The preceding chapter associates a differential polynomial to the
target weight diagram and controls the denominators of the scaled
theta kernel. The entry estimates bound their differential pairing at one step of
the theta chain.

Fix a prime \(p>2\) and a positive integer \(\nu\). Put
\[
 S=\Z/(2\nu)\Z,\qquad
 \rho_p=\frac1{p^2-1},\qquad
 \gamma_p=\frac{p}{p^2-1}.
\]
The parity of a position is the parity of its representative in
\(\{0,\ldots,2\nu-1\}\). Addition to a position takes place in \(S\).
The constants satisfy
\begin{equation}\label{eq:b3-entry-geometric}
 \sum_{l\ge0}p^{-(2l+1)}=\gamma_p,
 \qquad
 \frac1{p-1}=\gamma_p+\rho_p.
\end{equation}

For a nonarchimedean normed field \(K\), put
\(\mathcal O_K=\{a\in K:|a|\le1\}\). A polynomial or power series is
integral when its coefficients lie in \(\mathcal O_K\). When
\(|p|=p^{-1}\), write \(v_p(a)=-\log_p|a|\) for \(a\ne0\) and
\(v_p(0)=+\infty\). For \(a\in K\) (with \(a=0\) allowed) and \(x\in\R\), the bound
\(|a|\le p^{-x}\) is equivalent to \(v_p(a)\ge x\). All power series operations are formal.

\section{The two local exponents}

For the weight diagram at a place \(v\) in \ref{B3.02}, take
\(2\nu=n'_v\), the number of positions. Identify each position with
its residue class in \(S\) and let \(T\) be the image of the marked set
\(T_v\). Fix a layer \(1\le j\le m\), where \(m\) is the chain length.
At \(s\), let \(\alpha_s\) count the distinguished height-one columns.
The remaining columns number \(C_s\) and contain \(L_s\) boxes.
Thus \(\alpha_s=0\) for \(s\notin T\). Take \(r_0=N_{j-1}\) and
\(r_1=N_j\), with \(N_j\) as in \ref{B3.02}. In the entry estimates below, the
marked set \(T\subseteq S\) and the ranks \(r_0,r_1\in\N\) are arbitrary.

For \(a\in\N\), set
\[
 \ell_p(a)=
 \begin{cases}
  0,&a=0,\\
  1+\lfloor\log_p a\rfloor,&a>0.
 \end{cases}
\]
We use the error from \ref{B3.03} in the form
\begin{equation}\label{eq:b3-entry-error}
 \delta=
 \sum_{s\in S}
 \bigl(r_1\ell_p(\alpha_s)+r_0r_1\ell_p(L_s)\bigr)
 +\sum_{s\in T}\ell_p(\alpha_s)^2.
\end{equation}

\begin{definition}[The ordinary local exponent]\label{B3.13}
For \(s\in S\), let
\[
 D_s=
 \bigl\{(s'-s)_{\mathrm{rep}}:
       s'\in T,\ s'\text{ has parity opposite to }s\bigr\},
\]
where \((s'-s)_{\mathrm{rep}}\) denotes the representative in
\(\{0,\ldots,2\nu-1\}\). Define
\begin{equation}\label{eq:b3-entry-ga}
 g_A=
 \sum_{s\in T}
 \sum_{\substack{d\ge1,\ d\ {\rm odd}\\d\notin D_s}}
       [\alpha_sp^{-d}-C_{s+d}]_+
 -\delta,
 \qquad [x]_+=\max(x,0).
\end{equation}
For these weight data, \(g_A=g_{0,v}(j)\).
\end{definition}
\leanref{Formalism.\allowbreak{}MainNotions.\allowbreak{}FormalGroupSystem.\allowbreak{}dSet}{Formalism.\allowbreak{}SharpRestriction.\allowbreak{}GrowthBounds.\allowbreak{}FormalGroupSystem}{63--64}
\leanref{Formalism.\allowbreak{}MainNotions.\allowbreak{}FormalGroupSystem.\allowbreak{}gA}{Formalism.\allowbreak{}SharpRestriction.\allowbreak{}GrowthBounds.\allowbreak{}FormalGroupSystem}{66--70}
\leanref{Formalism.\allowbreak{}MainNotions.\allowbreak{}FormalGroupSystem.\allowbreak{}Sys.\allowbreak{}g0v}{Formalism.\allowbreak{}SharpRestriction.\allowbreak{}GrowthBounds.\allowbreak{}FormalGroupSystem}{325--325}

Every \(d\in D_s\) is the first positive distance from \(s\) to its
endpoint. The sum omits this odd integer \(d\) and includes
\(d+2\nu\). The bound \([\alpha_sp^{-d}-C_{s+d}]_+\le\alpha_sp^{-d}\)
gives convergence.

The identification with the local system follows by substituting
\[
 \alpha_s=\lambda_{s,1}-\lambda_{s,2},\quad C_s=\nu_{j,s,2}
 \quad(s\in T),
 \qquad
 \alpha_s=0,\quad C_s=\nu_{j,s,1}\quad(s\notin T).
\]
The error \eqref{eq:b3-entry-error} is then \(\delta_j\).

\begin{definition}[The enhanced local exponent]\label{B3.14}
Write
\[
 T_{\mathrm{odd}}=\{s\in T:s\text{ is odd}\},\qquad
 T_{\mathrm{ev}}=T\setminus T_{\mathrm{odd}},\qquad
 A_{\mathrm{ev}}=\sum_{s\in T_{\mathrm{ev}}}\alpha_s.
\]
Define
\begin{equation}\label{eq:b3-entry-gb}
 g_B=
 \gamma_p\sum_{s\in T}\alpha_s
 -\sum_{s\in T_{\mathrm{odd}}}
      \sum_{\substack{d\ge1\\d\ {\rm odd}}}
         \min(\alpha_sp^{-d},C_{s+d})
 -\delta.
\end{equation}
For these weight data, \(g_B=g_v(j)\) and
\(A_{\mathrm{ev}}=\mathrm{alEven}(v)\).
\end{definition}
\leanref{Formalism.\allowbreak{}MainNotions.\allowbreak{}FormalGroupSystem.\allowbreak{}gB}{Formalism.\allowbreak{}SharpRestriction.\allowbreak{}GrowthBounds.\allowbreak{}FormalGroupSystem}{57--61}
\leanref{Formalism.\allowbreak{}MainNotions.\allowbreak{}FormalGroupSystem.\allowbreak{}Sys.\allowbreak{}gv}{Formalism.\allowbreak{}SharpRestriction.\allowbreak{}GrowthBounds.\allowbreak{}FormalGroupSystem}{323--323}
\leanref{Formalism.\allowbreak{}MainNotions.\allowbreak{}FormalGroupSystem.\allowbreak{}Sys.\allowbreak{}alEven}{Formalism.\allowbreak{}SharpRestriction.\allowbreak{}GrowthBounds.\allowbreak{}FormalGroupSystem}{355--355}

The series converges by
\(0\le\min(\alpha_sp^{-d},C_{s+d})\le\alpha_sp^{-d}\).

\section{Differential pairings and column bounds}

For a finite variable set \(\mathcal V\), a multi-index
\(n\in\N^{\mathcal V}\), a polynomial \(P\), and a power series \(F\), put
\[
 |n|=\sum_{x\in\mathcal V}n_x,\qquad
 n!=\prod_{x\in\mathcal V}n_x!,\qquad
 \langle P,F\rangle
   =\sum_n [z^n]P\,n!\,[z^n]F.
\]
The sum in this differential evaluation (\ref{B3.11}) runs over
the monomials of \(P\). Write \(\mathbf1_x\) for the unit multi-index
at \(x\).

\begin{lemma}[Linear factors and Hasse derivatives]\label{B3.15}
Let \(K\) be a field, and let \(\mathcal V,\mathcal K\) be finite sets.
For directions \(b_k\in K^{\mathcal V}\), define
\[
 L_k(z)=\sum_{x\in\mathcal V}b_{k,x}z_x,\qquad
 D_b^{[\beta]}P
  =[t^\beta]P\left(z+\sum_{k\in\mathcal K}t_kb_k\right)
 \quad(\beta\in\N^{\mathcal K}).
\]
For every polynomial \(P\) and power series \(F\),
\begin{equation}\label{eq:b3-hasse-pair}
 \left\langle P,\left(\prod_kL_k^{\beta_k}\right)F\right\rangle
   =\left(\prod_k\beta_k!\right)
       \left\langle D_b^{[\beta]}P,F\right\rangle.
\end{equation}
For the integrality and degree assertions, \(\mathcal V\) may be
arbitrary. Let \(K\) be an ultrametric normed field. Integral \(P\) and
integral directions give integral \(D_b^{[\beta]}P\). If \(P\) is
homogeneous of degree \(d\), then the derivative is homogeneous of
degree \(d-|\beta|\) for \(|\beta|\le d\), and is zero for \(|\beta|>d\).
\end{lemma}
\leanref{Formalism.\allowbreak{}InvariantFactorPaths.\allowbreak{}G7.\allowbreak{}Top.\allowbreak{}e\_\allowbreak{}one\_\allowbreak{}identity}{Formalism.\allowbreak{}ThetaLogCoordinates.\allowbreak{}UniformThetaSubstitution.\allowbreak{}KOneConsequence}{25--25}
\leanref{Formalism.\allowbreak{}InvariantFactorPaths.\allowbreak{}G7.\allowbreak{}Top.\allowbreak{}e\_\allowbreak{}one\_\allowbreak{}integral}{Formalism.\allowbreak{}ThetaLogCoordinates.\allowbreak{}UniformThetaSubstitution.\allowbreak{}KOneConsequence}{26--26}

\begin{proof}
For a monomial \(z^m\), the contributions from multiplying by \(z_x\)
on the series side and differentiating on the polynomial side both
vanish when \(m_x=0\). When \(m_x>0\), their equality is
\[
 m!\,[z^{m-\mathbf1_x}]F
 =m_x(m-\mathbf1_x)!\,[z^{m-\mathbf1_x}]F.
\]
By linearity, multiplication by \(L_k\) is adjoint to
\(\partial_{b_k}=\sum_x b_{k,x}\partial/\partial z_x\).
Expanding the monomials of \(P(z+\sum_kt_kb_k)\) gives
\[
 \prod_k\partial_{b_k}^{\beta_k}P
   =\left(\prod_k\beta_k!\right)D_b^{[\beta]}P.
\]
The multinomial coefficients are integers, so this identity holds in
every characteristic. Apply the adjoint relation repeatedly to
obtain \eqref{eq:b3-hasse-pair}.

The expansion takes place over \(\mathcal O_K\) when \(P\) and the
directions are integral. For homogeneous \(P\), every term has total
\((z,t)\)-degree \(d\). If \(|\beta|\le d\), every monomial in the
coefficient of \(t^\beta\) has \(z\)-degree \(d-|\beta|\).
If \(|\beta|>d\), the coefficient is zero.
\end{proof}

Let \(\Sigma\) be a set of positions and use variables
\[
 z_{s,a,b},\qquad
 s\in\Sigma,\quad 0\le a<r_0,\quad 0\le b<r_1.
\]
For \(c=1,\ldots,C\), choose a position \(s_c\), a height \(h_c\in\N\),
a matrix \(M_c=(M_{c,i,a})_{0\le i<h_c,\ 0\le a<r_0}\), and a map
\[
 I_c:\{0,\ldots,h_c-1\}\longrightarrow\{0,\ldots,r_1-1\}.
\]
Put
\begin{equation}\label{eq:b3-low-minors}
 \Delta_c(z)=
 \det\left(\sum_{a<r_0}M_{c,i,a}z_{s_c,a,I_c(j)}
                  \right)_{0\le i,j<h_c},
 \qquad Q_0(z)=\prod_{c=1}^C\Delta_c(z).
\end{equation}
A height-zero determinant is \(1\). The maps \(I_c\) may repeat column
labels. A direction at \(s\) with row vector \(\xi\in K^{r_0}\)
and column label \(b\) is
\[
 B(s,\xi,b)_{t,a,b'}=
 \begin{cases}
  \xi_a,&t=s,\ b'=b,\\
  0,&\text{otherwise}.
 \end{cases}
\]

\begin{lemma}[Vanishing when columns are exhausted]\label{B3.16}
Let \(K\) be any field. For a finite family of directions
\(b_k=B(\operatorname{pos}(k),\xi_k,\operatorname{col}(k))\)
indexed by \(\mathcal K\), let \(\beta\in\N^{\mathcal K}\).
A direction is used when its multiplicity is positive.
Each of the following conditions forces \(D_b^{[\beta]}\Delta_c=0\).
\begin{enumerate}
 \item Some \(\beta_k\ge2\).
 \item Distinct used directions \(k,k'\) have the same position and
       either \(\xi_k=\xi_{k'}\) or
       \(\operatorname{col}(k)=\operatorname{col}(k')\).
 \item A used direction has position different from \(s_c\), or has
       column outside \(\operatorname{im}I_c\).
\end{enumerate}
Let \(\mathcal K_0\subseteq\mathcal K\) consist of directions at a
single position \(s_0\). If \(D_b^{[\beta]}Q_0\ne0\),
the following bounds hold.
\begin{equation}\label{eq:b3-hits-cap}
 \sum_{k\in\mathcal K_0}\beta_k
 \le \#\{c:s_c=s_0\}
 \quad\text{if all }\xi_k\ (k\in\mathcal K_0)\text{ are equal},
\end{equation}
and, for every \(J\subseteq\{0,\ldots,r_1-1\}\),
\begin{equation}\label{eq:b3-hits-columns}
 \sum_{k\in\mathcal K_0}\beta_k
 \le \sum_{c:s_c=s_0}\#(\operatorname{im}I_c\cap J)
 \quad\text{if }\operatorname{col}(k)\in J
       \text{ for }k\in\mathcal K_0.
\end{equation}
\end{lemma}
\leanref{Formalism.\allowbreak{}InvariantFactorPaths.\allowbreak{}G7.\allowbreak{}Top.\allowbreak{}minorHits}{Formalism.\allowbreak{}ThetaLogCoordinates.\allowbreak{}UniformThetaSubstitution.\allowbreak{}KOneConsequence}{28--28}
\leanref{Formalism.\allowbreak{}InvariantFactorPaths.\allowbreak{}G7.\allowbreak{}Top.\allowbreak{}hitsCapPos}{Formalism.\allowbreak{}ThetaLogCoordinates.\allowbreak{}UniformThetaSubstitution.\allowbreak{}KOneConsequence}{33--33}
\leanref{Formalism.\allowbreak{}InvariantFactorPaths.\allowbreak{}G7.\allowbreak{}Top.\allowbreak{}hitsColumnsPos}{Formalism.\allowbreak{}ThetaLogCoordinates.\allowbreak{}UniformThetaSubstitution.\allowbreak{}KOneConsequence}{34--34}

\begin{proof}
If \(I_c\) repeats a label, \(\Delta_c=0\). For injective \(I_c\), a
direction at \(s_c\) with label \(b\in\operatorname{im}I_c\) adds
\(t_kM_c\xi_k\) to the unique column labeled \(b\). Every other
direction leaves the matrix unchanged. The determinant has degree at
most one in the parameters attached to any single column. If two
directions have the same row vector, a term using both replacement
columns vanishes by alternation.

The Hasse product formula is
\begin{equation}\label{eq:b3-hasse-product}
 D_b^{[\beta]}\prod_c\Delta_c
   =\sum_{\beta^{(1)}+\cdots+\beta^{(C)}=\beta}
        \prod_cD_b^{[\beta^{(c)}]}\Delta_c.
\end{equation}
In a nonzero summand, directions from \(\mathcal K_0\) go to factors
at \(s_0\). With a common row vector, each such factor receives at most
one direction. With labels in \(J\), it receives at most one per label
in \(\operatorname{im}I_c\cap J\). Summing the respective bounds over
these factors proves \eqref{eq:b3-hits-cap} and
\eqref{eq:b3-hits-columns}.
\end{proof}

\section{Four parts of the substituted kernel}

On the cyclic position set \(S\), the series from \ref{B3.07} is
\begin{equation}\label{eq:b3-hh}
 S_{s,a,b}(z)=\sum_{l\ge0}(-1)^lz_{s-2l,a,b}^{p^{2l}},
 \qquad
 H(z)=\sum_{s\in S}\sum_{a<r_0,\ b<r_1}
               u_sS_{s,a,b}(z)^pS_{s+1,a,b}(z),
\end{equation}
where \(u_s\in\mathcal O_K\). Introduce top variables \(t_{s,b}\) and remaining variables \(y_{s,a,b}\).
Let \(\xi_s=(\xi_{s,a})_{a<r_0}\) be integral vectors. Make the substitution
\begin{equation}\label{eq:b3-top-low-sub}
 z_{s,a,b}\longmapsto
 y_{s,a,b}+\mathbf1_{s\in T}\xi_{s,a}t_{s,b}.
\end{equation}
Write \(h(t,y)\) for the resulting series. The remaining degree of
a monomial is its total degree in \(y\).

For \(E\in\N\), index the pure-top terms by
\[
 \phi=(v,l,l',a,b),\qquad
 v\in S,\quad 0\le l,l'<E,\quad a<r_0,\quad b<r_1.
\]
Set \(s_\phi=v-2l\), \(t_\phi=v+1-2l'\), and define
\begin{align}
 d^A_\phi
   &=p^{2l+1}\mathbf1_{(s_\phi,b)}
       +p^{2l'}\mathbf1_{(t_\phi,b)},\label{eq:b3-da}\\
 a^A_\phi
   &=\mathbf1_{s_\phi,t_\phi\in T}\,
      u_v(-1)^{l+l'}
      \xi_{s_\phi,a}^{p^{2l+1}}\xi_{t_\phi,a}^{p^{2l'}}.
       \label{eq:b3-aa}
\end{align}
For the second family, let
\(\kappa=(s,l,b)\) with \(s\in S\), \(0\le l<E\), \(b<r_1\), and put
\begin{align}
 d^B_\kappa&=p^{2l+1}\mathbf1_{(s,b)},\label{eq:b3-db}\\
 a^B_\kappa&=\mathbf1_{s\in T}(-1)^lu_{s+2l},\label{eq:b3-ab}\\
 L_\kappa(y)&=
   \sum_{a<r_0}\xi_{s,a}^{p^{2l+1}}y_{s+2l+1,a,b}.
       \label{eq:b3-b-direction}
\end{align}
The exponents \(d^A_\phi,d^B_\kappa\) index top monomials. The direction
\(b_\kappa\) of \(L_\kappa\) is \(B(s+2l+1,(\xi_{s,a}^{p^{2l+1}})_a,b)\).

\begin{lemma}[The four-part decomposition]\label{B3.17}
Assume \(K\) is nonarchimedean, \(p>2\) is prime, \(|p|=p^{-1}\),
and all \(u_s,\xi_{s,a}\) are integral. Let \(D,E\in\N\) satisfy
\(D<p^{2E}\). There are integral power series \(Z_{\ge p}\) and \(Y\)
such that every monomial of \(Z_{\ge p}\) with nonzero coefficient has
remaining degree at least \(p\), and
\begin{equation}\label{eq:b3-four-parts}
 h(t,y)\equiv
  \sum_\phi a^A_\phi t^{d^A_\phi}
  +\sum_\kappa a^B_\kappa t^{d^B_\kappa}L_\kappa(y)
  +Z_{\ge p}(t,y)+pY(t,y)
       \pmod{\text{total degree }>D}.
\end{equation}
The families and coefficients are given by
\eqref{eq:b3-da}--\eqref{eq:b3-b-direction}.
\end{lemma}
\leanref{Formalism.\allowbreak{}InvariantFactorPaths.\allowbreak{}G7.\allowbreak{}Top.\allowbreak{}hhDecomp}{Formalism.\allowbreak{}ThetaLogCoordinates.\allowbreak{}UniformThetaSubstitution.\allowbreak{}KOneConsequence}{38--38}
\leanref{Formalism.\allowbreak{}InvariantFactorPaths.\allowbreak{}G7.\allowbreak{}decomposition\_\allowbreak{}of\_\allowbreak{}hh}{Formalism.\allowbreak{}ThetaLogCoordinates.\allowbreak{}UniformThetaSubstitution.\allowbreak{}FrobeniusLiftNorms.\allowbreak{}EDecomp}{498--546}

\begin{proof}
If \(E=0\), then \(D=0\). Both families are empty and \(h\) has
zero constant term, so \(Z_{\ge p}=Y=0\) suffices.
Assume \(E\ge1\) and replace each \(S_{s,a,b}\) by its sum over
\(0\le l<E\). The omitted terms have degree at least \(p^{2E}\).
The linear substitution \eqref{eq:b3-top-low-sub} preserves this
degree bound, so the coefficients of \(h\) through degree \(D\)
agree with those of the finite substitution.

Work modulo \(p\mathcal O_K\), where Frobenius is additive. Since
\(p\) is odd, the \(p\)-th power preserves the signs \((-1)^l\).
For fixed \(v,a,b\), the substituted product
\(S_{v,a,b}^pS_{v+1,a,b}\) becomes \((Z_1+T_1)(Z_2+T_2)\), where
\begin{align*}
 Z_1&=\sum_{l<E}(-1)^ly_{v-2l,a,b}^{p^{2l+1}},\\
 T_1&=\sum_{l<E}(-1)^l
       \bigl(\mathbf1_{v-2l\in T}\xi_{v-2l,a}t_{v-2l,b}
                    \bigr)^{p^{2l+1}},\\
 Z_2&=\sum_{l'<E}(-1)^{l'}y_{v+1-2l',a,b}^{p^{2l'}},\\
 T_2&=\sum_{l'<E}(-1)^{l'}
       \bigl(\mathbf1_{v+1-2l'\in T}
                   \xi_{v+1-2l',a}t_{v+1-2l',b}
                    \bigr)^{p^{2l'}}.
\end{align*}
Write \(Z_2=y_{v+1,a,b}+Z_2'\), where every monomial of \(Z_2'\)
has remaining degree at least \(p^2\). Then
\begin{equation}\label{eq:b3-sort-four}
 (Z_1+T_1)(Z_2+T_2)
   =T_1T_2+T_1y_{v+1,a,b}
      +Z_1(Z_2+T_2)+T_1Z_2'.
\end{equation}

Multiply by \(u_v\) and sum over \(v,a,b\). Expanding \(T_1T_2\)
gives the first sum of \eqref{eq:b3-four-parts}, with the exponents
and coefficients in \eqref{eq:b3-da} and \eqref{eq:b3-aa}.
In the second term of \eqref{eq:b3-sort-four}, change variables
from \(v\) to \(s=v-2l\). Summing over \(a\) gives
\[
 (-1)^lu_{s+2l}\,
 t_{s,b}^{p^{2l+1}}
 \sum_a\xi_{s,a}^{p^{2l+1}}y_{s+2l+1,a,b},
\]
when \(s\in T\), and zero otherwise. Summing over \(s\in S\),
\(0\le l<E\) and \(0\le b<r_1\) gives the sum indexed by \(\kappa=(s,l,b)\).

Define \(Z_{\ge p}\) by summing
\(u_v[Z_1(Z_2+T_2)+T_1Z_2']\) over \(v,a,b\). Its monomials have
\(y\)-degree at least \(p\) by the bounds for \(Z_1\) and \(Z_2'\).
Subtract the two indexed sums and \(Z_{\ge p}\) from the finite
substitution and divide by \(p\). The congruence modulo
\(p\mathcal O_K\) gives an integral quotient \(Y\), which completes
\eqref{eq:b3-four-parts}.
\end{proof}

For a series \(h\) with zero constant term, write \(g\in\mathcal C_p(h)\)
if there are integral series \(U_i\), \(i\in\N\), with
\begin{equation}\label{eq:b3-class-recall}
 [z^n]g=\sum_{i=0}^{|n|}p^{-i}[z^n](h^iU_i)
 \quad\text{for every multi-index }n.
\end{equation}
This is the expansion class of \ref{B3.07}. For finitely many source
variables, let \(\phi\) be a substitution by integral series with zero
constant term. If \(g\in\mathcal C_p(h)\), then
\(g\circ\phi\in\mathcal C_p(h\circ\phi)\). The witnesses are
\(U_i\circ\phi\), since substitution preserves integrality and lower
bounds on total degree.

\begin{lemma}[The factorial bound for a nonzero pairing]\label{B3.18}
Let \(\mathcal T,\mathcal L,\mathcal A,\mathcal K\) be finite sets,
and let \(K\) be a nonarchimedean normed field with \(|p|=p^{-1}\),
where \(p\) is prime. Use variables \(t_\tau\), \(\tau\in\mathcal T\),
and \(y_x\), \(x\in\mathcal L\).

Fix a total degree \(D\in\N\).
Suppose \(h\) has zero constant term, \(g\in\mathcal C_p(h)\), and,
in total degrees at most \(D\),
\[
 h=
 \sum_{\phi\in\mathcal A}a^A_\phi t^{d^A_\phi}
 +\sum_{\kappa\in\mathcal K}
          a^B_\kappa t^{d^B_\kappa}L_\kappa(y)
 +Z_{\ge p}+pY.
\]
Here \(d^A_\phi,d^B_\kappa\in\N^{\mathcal T}\) and
\(L_\kappa(y)=\sum_x b_{\kappa,x}y_x\). The scalars \(a^A_\phi,a^B_\kappa\),
the coordinates \(b_{\kappa,x}\), and the series \(Z_{\ge p},Y\) are
integral. Every nonzero monomial of \(Z_{\ge p}\) has \(y\)-degree
at least \(p\).

Let \(Q_0(y)\) be an integral homogeneous polynomial of degree \(d_0\),
and let \(e\in\N^{\mathcal T}\) satisfy \(|e|+d_0=D\).
If \(J=\langle t^eQ_0(y),g\rangle\ne0\), there exist multiplicities
\(m_\phi,\beta_\kappa\in\N\), an integer \(i_2\ge0\), and a multi-index
\(n\in\operatorname{supp}D_b^{[\beta]}Q_0\) such that
\begin{gather}
 m_\phi>0\Longrightarrow a^A_\phi\ne0,\qquad
 \beta_\kappa>0\Longrightarrow a^B_\kappa\ne0,
       \label{eq:b3-term-active}\\
 \sum_\phi m_\phi d^A_\phi
       +\sum_\kappa\beta_\kappa d^B_\kappa\le e,
 \qquad pi_2\le|n|,\label{eq:b3-term-consumption}\\
 |J|\le
 p^{\,\sum_\phi m_\phi+\sum_\kappa\beta_\kappa+i_2}
       |e!|\,
       \left|\left(\sum_\kappa\beta_\kappa\right)!\right|\,
       |n!|.\label{eq:b3-term-factorials}
\end{gather}
The inequality between top multi-indices is coordinatewise.
\end{lemma}
\leanref{Formalism.\allowbreak{}InvariantFactorPaths.\allowbreak{}G7.\allowbreak{}Top.\allowbreak{}termBound}{Formalism.\allowbreak{}ThetaLogCoordinates.\allowbreak{}UniformThetaSubstitution.\allowbreak{}KOneConsequence}{37--37}
\leanref{Formalism.\allowbreak{}InvariantFactorPaths.\allowbreak{}G7.\allowbreak{}TermBound\_\allowbreak{}of}{Formalism.\allowbreak{}ThetaLogCoordinates.\allowbreak{}UniformThetaSubstitution.\allowbreak{}FrobeniusLiftNorms.\allowbreak{}NormNatCastLeOfDvd}{495--543}

\begin{proof}
The polynomial \(t^eQ_0\) is homogeneous of degree \(D\).
Thus \eqref{eq:b3-class-recall} expresses its pairing with \(g\) as
a finite sum of pairings with \(p^{-i}h^iU_i\), \(0\le i\le D\).
In these pairings we may use the displayed expression for \(h\),
because multiplication by power series preserves agreement through
degree \(D\).

Expand \(h^i\) by the multinomial formula. A summand chooses
\(m_\phi\) occurrences of each pure-top term,
\(\beta_\kappa\) occurrences of each term with a linear remaining
factor, \(i_2\) occurrences of \(Z_{\ge p}\), and \(i_3\) occurrences
of \(pY\). Put
\[
 A=\sum_\phi m_\phi,\qquad B=\sum_\kappa\beta_\kappa,
 \qquad i=A+B+i_2+i_3.
\]
Its multinomial coefficient and the exponent of its extracted top
monomial are
\[
 M=\frac{i!}
        {(\prod_\phi m_\phi!)(\prod_\kappa\beta_\kappa!)i_2!i_3!},
 \qquad
 q=\sum_\phi m_\phi d^A_\phi+\sum_\kappa\beta_\kappa d^B_\kappa.
\]
Use \ref{B3.15} to move the product of linear factors to the polynomial
side. This contributes \(\prod_\kappa\beta_\kappa!\) and replaces
\(Q_0\) by \(D_b^{[\beta]}Q_0\), with \(t^e\) unchanged. The resulting
pairing is a sum over \(n\in\operatorname{supp}D_b^{[\beta]}Q_0\)
with factorial factor \(e!n!\). A nonzero coefficient requires \(q\le e\).

Apart from integral scalar factors and a coefficient of the integral
Hasse derivative, such a summand has the form
\begin{equation}\label{eq:b3-term-shape}
 p^{-i+i_3}M
 \left(\prod_\kappa\beta_\kappa!\right)e!n!\,
 [t^{e-q}y^n](Z_{\ge p}^{\,i_2}Y^{i_3}U_i).
\end{equation}
If it is nonzero, every scalar factor selected with positive
multiplicity is nonzero. The minimum remaining degree of
\(Z_{\ge p}^{\,i_2}\) gives \(pi_2\le|n|\).

The multinomial product is divisible by the factorial of the total
number \(B\) of linear factors. Indeed,
\[
 M\prod_\kappa\beta_\kappa!
   =B!\,
      \frac{i!}{(\prod_\phi m_\phi!)B!i_2!i_3!},
\]
where the quotient is an integer multinomial coefficient. Hence
\(\left|M\prod_\kappa\beta_\kappa!\right|\le|B!|\).
The coefficient in \eqref{eq:b3-term-shape} has norm at most one
and \(|p^{-i+i_3}|=p^{A+B+i_2}\). For each nonzero summand \(w\)
this gives \(|w|\le p^{A+B+i_2}|e!|\,|B!|\,|n!|\) with its own data.
Choose a summand of maximal norm. The ultrametric inequality gives
\(|J|\le |w|\), so its data satisfy \eqref{eq:b3-term-factorials}.
\end{proof}

\section{The arithmetic contribution of the factorials}

\begin{lemma}[A factorial valuation bound]\label{B3.19}
For every prime \(p\) and every \(a\in\N\), the arithmetic valuation
satisfies
\begin{equation}\label{eq:b3-factorial-digit}
 v_p(a!)\ge\frac{a}{p-1}-\ell_p(a).
\end{equation}
\end{lemma}
\leanref{Formalism.\allowbreak{}InvariantFactorPaths.\allowbreak{}G7.\allowbreak{}factorialValuationLowerBound}{Formalism.\allowbreak{}ThetaLogCoordinates.\allowbreak{}UniformThetaSubstitution.\allowbreak{}FrobeniusLiftNorms.\allowbreak{}ABArith}{35--53}

Legendre's formula gives \(v_p(a!)=(a-s_p(a))/(p-1)\), where \(s_p(a)\)
is the base-\(p\) digit sum \cite[Ch.~I, \S2, Ex.~14, p.~7]{koblitz-1984}.
The archive uses the weaker lower bound obtained by substituting
\(s_p(a)\le(p-1)\ell_p(a)\) in Legendre's formula.
% Formula source: [Mathlib, sub_one_mul_padicValNat_factorial](https://leanprover-community.github.io/mathlib4_docs/Mathlib/NumberTheory/Padics/PadicVal/Basic.html#sub_one_mul_padicValNat_factorial).
% Bibliographic source: [Koblitz, second edition](https://link.springer.com/book/10.1007/978-1-4612-1112-9).

\begin{lemma}[The net contribution of three factorial families]\label{B3.20}
Let \(S\) be a finite set, \(T\subseteq S\), \(p\) prime, and
\(r_0,r_1\in\N\). Let \(\alpha,L:S\to\N\), with \(\alpha_s=0\)
outside \(T\). Suppose nonnegative integer arrays
\[
 e_{s,b}\ (b<r_1),\qquad
 a_{s,a,b}\ (a<r_0,\ b<r_1),\qquad
 i_{s,l}\ (l\in\N)
\]
and integers \(i_0,i_2\ge0\) satisfy
\begin{align}
 &\sum_{b<r_1}e_{s,b}=\alpha_s,\qquad
   \sum_{a,b}a_{s,a,b}\le L_s,\label{eq:b3-star-local}\\
 &i_{s,l}=0\quad\text{if }s\notin T
                   \text{ or }l\notin E,\label{eq:b3-star-support}\\
 &p^{2l+1}i_{s,l}\le\alpha_s,\label{eq:b3-star-source}\\
 &\sum_{s,a,b}a_{s,a,b}+\sum_{s\in S}\sum_{l\in E}i_{s,l}
        =\sum_sL_s,\qquad
   pi_2\le\sum_sL_s,\label{eq:b3-star-total}
\end{align}
where \(E\subseteq\N\) is finite. Define \(\delta\) by
\eqref{eq:b3-entry-error}, and set
\[
 \mathcal F=
 \sum_{s,b}v_p(e_{s,b}!)
 +\sum_{s\in S}\sum_{l\in E}v_p(i_{s,l}!)
 +\sum_{s,a,b}v_p(a_{s,a,b}!).
\]
Then
\begin{equation}\label{eq:b3-star-arith}
\begin{aligned}
 &\mathcal F-\rho_p\sum_{s\in S}(\alpha_s+L_s)
       -\left(i_0+\sum_{s\in S}\sum_{l\in E}i_{s,l}+i_2\right)\\
 &\qquad\ge \gamma_p\sum_{s\in T}\alpha_s
       -i_0-\sum_{s\in T}\sum_{l\in E}i_{s,l}-\delta.
\end{aligned}
\end{equation}
\end{lemma}
\leanref{Formalism.\allowbreak{}InvariantFactorPaths.\allowbreak{}G7.\allowbreak{}gainFactorialInequality}{Formalism.\allowbreak{}ThetaLogCoordinates.\allowbreak{}UniformThetaSubstitution.\allowbreak{}FrobeniusLiftNorms.\allowbreak{}ABArith}{103--183}

\begin{proof}
Put
\[
 c_p=\frac1{p-1},\qquad
 A=\sum_s\alpha_s,\qquad L=\sum_sL_s,\qquad
 I=\sum_{s\in S}\sum_{l\in E}i_{s,l}.
\]
By \ref{B3.19}, monotonicity of \(\ell_p\), and
\eqref{eq:b3-star-local},
\begin{align}
 \sum_bv_p(e_{s,b}!)
   &\ge c_p\alpha_s-r_1\ell_p(\alpha_s),
       \label{eq:b3-star-top}\\
 \sum_{a,b}v_p(a_{s,a,b}!)
   &\ge c_p\sum_{a,b}a_{s,a,b}
             -r_0r_1\ell_p(L_s).
       \label{eq:b3-star-rem}
\end{align}
For the middle family,
\begin{equation}\label{eq:b3-star-logsum}
 \sum_{l\in E}\ell_p(i_{s,l})\le\ell_p(\alpha_s)^2.
\end{equation}
For \(\alpha_s=0\), all the \(i_{s,l}\) vanish. A positive \(i_{s,l}\)
implies \(p^{2l+1}\le\alpha_s\), hence \(l<\ell_p(\alpha_s)\).
There are at most \(\ell_p(\alpha_s)\) such indices, and each
\(i_{s,l}\le\alpha_s\). Another application of \ref{B3.19} gives
\[
 \sum_{l\in E}v_p(i_{s,l}!)
   \ge c_p\sum_{l\in E}i_{s,l}-\ell_p(\alpha_s)^2.
\]
Outside \(T\), this inequality has zero on both sides. Summing the
three bounds over \(s\) and using \eqref{eq:b3-star-total} gives
\begin{equation}\label{eq:b3-three-factorial-sum}
 \mathcal F\ge c_p(A+L)-\delta.
\end{equation}
Since \(c_p-\rho_p=\gamma_p\), it follows that
\[
 \mathcal F-\rho_p(A+L)-(i_0+I+i_2)
 \ge \gamma_p A-i_0-I-\delta+\gamma_p L-i_2.
\]
By \eqref{eq:b3-star-total} and \(1/p\le\gamma_p\), we have
\(i_2\le L/p\le\gamma_p L\).
\end{proof}

\section{A common estimate for every nonzero entry}

Let \(K\) be an ultrametric normed field with \(|p|=p^{-1}\), where
\(p>2\) is prime, and let \(\nu\ge1\). Use variables \(z_{s,a,b}\)
with \(s\in S\), \(a<r_0\), \(b<r_1\). Take the determinant product
\(Q_0\) of \eqref{eq:b3-low-minors} with positions \(s_c\in S\)
and integral frame entries \(M_{c,i,a}\).

Choose integers \(e_{s,b}\ge0\), supported on \(T\), and integral
vectors \(\xi_s\in\mathcal O_K^{r_0}\). Define
\begin{align}
 P_{\mathrm{ent}}(z)
  &=\prod_{s\in S,\ b<r_1}
       \left(\sum_{a<r_0}\xi_{s,a}z_{s,a,b}\right)^{e_{s,b}}
       Q_0(z),\label{eq:b3-entry-polynomial}\\
 \alpha_s&=\sum_{b<r_1}e_{s,b},\qquad
 C_s=\#\{c:s_c=s\},\qquad
 L_s=\sum_{c:s_c=s}h_c,\label{eq:b3-entry-counts}\\
 D&=\sum_s(\alpha_s+L_s).\label{eq:b3-entry-degree}
\end{align}
Let \(H\) be \eqref{eq:b3-hh}, with integral \(u_s\), and take
\(F\in\mathcal C_p(H)\). Under these entry hypotheses, the entry is
\(\mathcal J=\langle P_{\mathrm{ent}},F\rangle\).

\begin{theorem}[The common entry estimate]\label{B3.21}
Under the entry hypotheses, either \(\mathcal J=0\), or there are
an integer \(E\ge0\) and nonnegative multiplicities \(m_\phi,\beta_\kappa\)
for the families of \ref{B3.17} such that
\begin{gather}
 m_\phi>0\Longrightarrow s_\phi,t_\phi\in T,\qquad
 \beta_{(s,l,b)}>0\Longrightarrow s\in T,
       \label{eq:b3-entry-support}\\
 \sum_\phi m_\phi d^A_\phi
       +\sum_\kappa\beta_\kappa d^B_\kappa\le e,
       \label{eq:b3-entry-budget}\\
 D_b^{[\beta]}Q_0\ne0,\label{eq:b3-entry-hasse-nonzero}
\end{gather}
and
\begin{equation}\label{eq:b3-entry-star}
 v_p(\mathcal J)\ge
 \gamma_p\sum_{s\in T}\alpha_s
     -\sum_\phi m_\phi-\sum_\kappa\beta_\kappa
     -\delta+\rho_pD.
\end{equation}
Here \(b_\kappa\) is the direction in \eqref{eq:b3-b-direction}.
\end{theorem}
\leanref{Formalism.\allowbreak{}InvariantFactorPaths.\allowbreak{}G7.\allowbreak{}EntryHyp}{Formalism.\allowbreak{}ThetaLogCoordinates.\allowbreak{}UniformThetaSubstitution.\allowbreak{}StatementE}{170--174}
\leanref{Formalism.\allowbreak{}InvariantFactorPaths.\allowbreak{}G7.\allowbreak{}Top.\allowbreak{}entryStar}{Formalism.\allowbreak{}ThetaLogCoordinates.\allowbreak{}UniformThetaSubstitution.\allowbreak{}KOneConsequence}{39--41}

\begin{proof}
Transposing the linear substitution \eqref{eq:b3-top-low-sub}
in the differential pairing gives
\begin{equation}\label{eq:b3-entry-transpose}
 \mathcal J
  =\left\langle t^eQ_0(y),
       F\bigl(y+\mathbf1_T\xi t\bigr)\right\rangle.
\end{equation}
A derivative in \(t_{s,b}\) is carried by the chain rule to
\(\sum_a\xi_{s,a}\partial/\partial z_{s,a,b}\) when \(s\in T\).
Top exponents outside \(T\) vanish. A derivative in \(y_{s,a,b}\)
becomes \(\partial/\partial z_{s,a,b}\). Iterating and evaluating
at the origin gives \eqref{eq:b3-entry-transpose}.

The substituted \(F\) belongs to \(\mathcal C_p(h)\).
Choose \(E\) with \(D<p^{2E}\) and apply \ref{B3.17}.
Each determinant contributes its height to the degree of \(Q_0\),
so \(Q_0\) is homogeneous of degree \(\sum_sL_s\).
Apply \ref{B3.18} to \eqref{eq:b3-entry-transpose}.

For a nonzero entry, the resulting \(m,\beta,i_2,n\) satisfy
\eqref{eq:b3-entry-support} by the nonzero scalar conditions.
The other conclusions give \eqref{eq:b3-entry-budget},
\eqref{eq:b3-entry-hasse-nonzero}, and
\begin{equation}\label{eq:b3-entry-factorial-value}
 v_p(\mathcal J)\ge
 v_p(e!)+v_p(B!)+v_p(n!)
       -\left(\sum_\phi m_\phi+B+i_2\right),
 \qquad B=\sum_\kappa\beta_\kappa.
\end{equation}
Every integer has norm at most one. For an integer \(u\) with
\(p\nmid u\), a Bézout identity \(au+bp=1\) gives \(|u|=1\) since
\(|p|<1\). Writing each factorial factor \(m\) in
\eqref{eq:b3-entry-factorial-value} as \(p^ku\) with \(p\nmid u\)
gives \(v_p(m)=k=\operatorname{ord}_p(m)\), where \(\operatorname{ord}_p\)
denotes the arithmetic valuation.

For each source \(s\), set \(i_{s,l}=\sum_{b<r_1}\beta_{(s,l,b)}\)
for \(l<E\) and \(i_{s,l}=0\) for \(l\ge E\). The top budget gives
\(p^{2l+1}i_{s,l}\le\alpha_s\). Since \(\sum_{s,l}i_{s,l}=B\),
\[
 \prod_{s,l}i_{s,l}!\ \mid\ B!,
 \qquad
 v_p(B!)\ge\sum_{s,l}v_p(i_{s,l}!).
\]
A monomial \(y^n\) of \(D_b^{[\beta]}Q_0\) comes from a monomial
\(y^{n_0}\) of \(Q_0\) with \(n\le n_0\) and \(|n|+B=|n_0|\).
Every monomial of \(Q_0\) has degree \(L_s\) in the variables at \(s\),
so
\[
 \sum_{a,b}n_{s,a,b}\le L_s,\qquad
 |n|+B=\sum_sL_s.
\]
Also \(pi_2\le|n|\le\sum_sL_s\). Apply \ref{B3.20} with
\(a_{s,a,b}=n_{s,a,b}\), \(i_0=\sum_\phi m_\phi\), and the grouped
array \(i_{s,l}\). Its conclusion and
\eqref{eq:b3-entry-factorial-value} give \eqref{eq:b3-entry-star}.
\end{proof}

\section{First distances and the ordinary gain}

\begin{lemma}[Charging losses to first distances]\label{B3.22}
Fix \(n,a,a',d\in\N\) and \(s,s'\in\Z/n\Z\).
For \(n=0\), interpret \(\Z/0\Z\) as \(\Z\).
Suppose \(a'<a\), \(s+a=s'+a'\), and \(d\) is the least positive
integer with \(s+d=s'\). Then \(d\le a\).

For the scalar estimate, take a finite subset \(T\) of \(S=\Z/n\Z\),
where \(n\in\N\). Let \(p>1\) and \(\delta\) be real and let
\(\alpha,C:S\to\R_{\ge0}\). Give each \(s\) a finite
set \(D_s\subseteq\N\). Suppose nonnegative real arrays
\(i_{s,l},\operatorname{ch}_{s,l}\) are supported in a common finite
set of indices \(l\in\N\) and satisfy
\begin{align}
 &\operatorname{ch}_{s,l}=0
       &&\text{if }2l+1\notin D_s,\label{eq:b3-charge-support}\\
 &\operatorname{ch}_{s,l}+i_{s,l}
        \le\alpha_sp^{-(2l+1)}
       &&\text{if }s\in T,\ 2l+1\in D_s,\label{eq:b3-charge-excluded}\\
 &i_{s,l}\le
       \min(\alpha_sp^{-(2l+1)},C_{s+2l+1})
       &&\text{if }s\in T,\ 2l+1\notin D_s.
       \label{eq:b3-charge-free}
\end{align}
If a real number \(V\) satisfies
\[
 V\ge
 \gamma_p\sum_{s\in T}\alpha_s
 -\sum_{s\in T,l}\operatorname{ch}_{s,l}
 -\sum_{s\in T,l}i_{s,l}-\delta,
\]
then
\begin{equation}\label{eq:b3-charge-conclusion}
 V\ge
 \sum_{s\in T}
   \sum_{\substack{d\ge1,\ d\ {\rm odd}\\d\notin D_s}}
        [\alpha_sp^{-d}-C_{s+d}]_+-\delta.
\end{equation}
\end{lemma}
\leanref{Formalism.\allowbreak{}InvariantFactorPaths.\allowbreak{}G7.\allowbreak{}cyclicShiftLeast}{Formalism.\allowbreak{}ThetaLogCoordinates.\allowbreak{}UniformThetaSubstitution.\allowbreak{}FrobeniusLiftNorms.\allowbreak{}ABArith}{185--191}
\leanref{Formalism.\allowbreak{}InvariantFactorPaths.\allowbreak{}G7.\allowbreak{}excessGainBound}{Formalism.\allowbreak{}ThetaLogCoordinates.\allowbreak{}UniformThetaSubstitution.\allowbreak{}FrobeniusLiftNorms.\allowbreak{}ABArith}{193--238}

\begin{proof}
The positive integer \(a-a'\) satisfies \(s+(a-a')=s'\), so
minimality gives \(d\le a-a'\le a\).

For \(s\in T\), write \(x_l=\alpha_sp^{-(2l+1)}\), with \(d=2l+1\).
At an excluded distance, \eqref{eq:b3-charge-excluded} gives
\(0\le x_l-\operatorname{ch}_{s,l}-i_{s,l}\). At an unexcluded
distance, \eqref{eq:b3-charge-support} and \eqref{eq:b3-charge-free} give
\[
 [x_l-C_{s+d}]_+
 \le x_l-i_{s,l}
 =x_l-\operatorname{ch}_{s,l}-i_{s,l}.
\]
The series of positive parts is bounded termwise by the convergent
series \(\sum_lx_l=\gamma_p\alpha_s\). The other two arrays have finite
support. Summation therefore gives
\[
 \sum_{\substack{d\ge1,\ d\ {\rm odd}\\d\notin D_s}}
      [\alpha_sp^{-d}-C_{s+d}]_+
 \le\gamma_p\alpha_s
      -\sum_l\operatorname{ch}_{s,l}-\sum_li_{s,l}.
\]
Sum over \(s\in T\) and subtract \(\delta\).
\end{proof}

\begin{theorem}[The ordinary entry bound]\label{B3.23}
Under the entry hypotheses of \ref{B3.21}, with \(D_s\) as in
\ref{B3.13},
\begin{equation}\label{eq:b3-entry-a1}
 |\mathcal J|\le p^{-(g_A+\rho_pD)}.
\end{equation}
\end{theorem}
\leanref{Formalism.\allowbreak{}InvariantFactorPaths.\allowbreak{}G7.\allowbreak{}Top.\allowbreak{}entry\_\allowbreak{}a\_\allowbreak{}one}{Formalism.\allowbreak{}ThetaLogCoordinates.\allowbreak{}UniformThetaSubstitution.\allowbreak{}KOneConsequence}{43--43}
\leanref{Formalism.\allowbreak{}InvariantFactorPaths.\allowbreak{}G7.\allowbreak{}entryExcessBound\_\allowbreak{}of}{Formalism.\allowbreak{}ThetaLogCoordinates.\allowbreak{}UniformThetaSubstitution.\allowbreak{}FrobeniusLiftNorms.\allowbreak{}EntryAOneParity}{177--312}

\begin{proof}
The bound holds for \(\mathcal J=0\). For a nonzero entry take the
multiplicities of \ref{B3.21}. Set \(i_{s,l}=\sum_b\beta_{(s,l,b)}\)
for \(l<E\) and \(i_{s,l}=0\) for \(l\ge E\), as in that proof.
For fixed \(s,l\) with \(l<E\), the directions share the target
\(s+2l+1\) and the row vector \((\xi_{s,a}^{p^{2l+1}})_a\).
By \eqref{eq:b3-entry-hasse-nonzero} and \eqref{eq:b3-hits-cap},
\(i_{s,l}\le C_{s+2l+1}\). The top budget gives
\(p^{2l+1}i_{s,l}\le\alpha_s\), so for every \(l\ge0\)
\begin{equation}\label{eq:b3-arrival-source-min}
 i_{s,l}\le\min(\alpha_sp^{-(2l+1)},C_{s+2l+1}).
\end{equation}

Assign each occurrence of a pure-top term to an excluded distance.
For \(\phi=(v,l,l',a,b)\), put
\[
 s_1=v-2l,\qquad s_2=v+1-2l',\qquad
 a_1=2l+1,\qquad a_2=2l'.
\]
Then \(s_1+a_1=s_2+a_2\). The endpoints have opposite parity, and
\(a_1,a_2\) are unequal. Choose as source the endpoint with larger
\(a_i\), and as target the other endpoint. Let \(d_\phi\) be the
first positive distance from the source to the target.
By \ref{B3.22}, \(d_\phi\le\max(a_1,a_2)\). The distance is odd.
If \(m_\phi>0\), both endpoints belong to \(T\), so
\(d_\phi\in D_{\operatorname{source}(\phi)}\). The pure-top monomial
consumes at least \(p^{d_\phi}\) units of top degree at its chosen
source in column \(b\).

For \(d=2l+1\), let
\[
 \operatorname{ch}_{s,l}
  =\sum_{\substack{\phi:\ \operatorname{source}(\phi)=s\\
                             d_\phi=2l+1}}m_\phi.
\]
These counts have finite support, and
\[
 \sum_{s\in T,l}\operatorname{ch}_{s,l}=\sum_\phi m_\phi,
 \qquad
 \operatorname{ch}_{s,l}=0\quad\text{if }2l+1\notin D_s.
\]
For fixed \(s,l\), sum the coordinatewise top budget
\eqref{eq:b3-entry-budget} over the columns at \(s\).
The charged \(A\)-terms and the \(B\)-terms of distance \(2l+1\)
consume at least \(p^{2l+1}(\operatorname{ch}_{s,l}+i_{s,l})\)
units of top degree. Consequently
\begin{equation}\label{eq:b3-combined-charge-budget}
 \operatorname{ch}_{s,l}+i_{s,l}
       \le\alpha_sp^{-(2l+1)}.
\end{equation}
Apply the scalar part of \ref{B3.22} with
\[
 V=
 \gamma_p\sum_{s\in T}\alpha_s
       -\sum_\phi m_\phi-\sum_\kappa\beta_\kappa-\delta.
\]
By \eqref{eq:b3-combined-charge-budget} and
\eqref{eq:b3-arrival-source-min}, this gives \(V\ge g_A\).
The common estimate \eqref{eq:b3-entry-star} proves the theorem.
\end{proof}

\section{Parity and the enhanced entry bounds}

\begin{theorem}[The second unmodified entry bound]\label{B3.24}
Under the entry hypotheses,
\begin{equation}\label{eq:b3-entry-a2}
 |\mathcal J|\le p^{-(g_B-A_{\mathrm{ev}}+\rho_pD)}.
\end{equation}
\end{theorem}
\leanref{Formalism.\allowbreak{}InvariantFactorPaths.\allowbreak{}G7.\allowbreak{}Top.\allowbreak{}entry\_\allowbreak{}a\_\allowbreak{}two}{Formalism.\allowbreak{}ThetaLogCoordinates.\allowbreak{}UniformThetaSubstitution.\allowbreak{}KOneConsequence}{44--44}
\leanref{Formalism.\allowbreak{}InvariantFactorPaths.\allowbreak{}G7.\allowbreak{}entryCappedBound\_\allowbreak{}of}{Formalism.\allowbreak{}ThetaLogCoordinates.\allowbreak{}UniformThetaSubstitution.\allowbreak{}FrobeniusLiftNorms.\allowbreak{}EntryATwoAndB}{220--266}

\begin{proof}
For a nonzero entry take the data of \ref{B3.21}, with \(i_{s,l}\)
extended by zero as above, and put
\[
 A=\sum_\phi m_\phi,\qquad
 B_{\mathrm{ev}}=\sum_{s\in T_{\mathrm{ev}}}\sum_{l\ge0}i_{s,l},
 \qquad
 B_{\mathrm{odd}}=\sum_{s\in T_{\mathrm{odd}}}\sum_{l\ge0}i_{s,l}.
\]
Each pure-top term consumes at least one unit of top degree at its
even endpoint. Each even-source linear-factor term also consumes at
least one unit at an even position. Summing the top budget
\eqref{eq:b3-entry-budget} over the even coordinates gives
\begin{equation}\label{eq:b3-even-payment}
 A+B_{\mathrm{ev}}\le A_{\mathrm{ev}}.
\end{equation}
For odd sources, \eqref{eq:b3-arrival-source-min} gives
\[
 B_{\mathrm{odd}}
 \le\sum_{s\in T_{\mathrm{odd}}}
       \sum_{l\ge0}\min(\alpha_sp^{-(2l+1)},C_{s+2l+1}).
\]
Substitute both inequalities in the common lower bound
\[
 v_p(\mathcal J)\ge
 \gamma_p\sum_{s\in T}\alpha_s
       -A-B_{\mathrm{ev}}-B_{\mathrm{odd}}-\delta+\rho_pD.
\]
This gives \eqref{eq:b3-entry-a2}. The norm bound also holds for
a zero entry.
\end{proof}

\begin{theorem}[The prescribed-column entry bound]\label{B3.25}
Under the entry hypotheses, let
\(J\subseteq\{0,\ldots,r_1-1\}\) be arbitrary. Define the nonnegative
integer
\begin{align}
 N_J={}&
 \sum_{\substack{s\ {\rm even}\\b\notin J}}e_{s,b}
 +\sum_{\substack{s\ {\rm odd}\\b\in J}}e_{s,b}
 +\sum_{c:\ s_c\ {\rm odd}}\#(\operatorname{im}I_c\cap J).
       \label{eq:b3-entry-column-payment}
\end{align}
Then
\begin{equation}\label{eq:b3-entry-b}
 |p^{N_J}\mathcal J|\le p^{-(g_B+\rho_pD)}.
\end{equation}
\end{theorem}
\leanref{Formalism.\allowbreak{}InvariantFactorPaths.\allowbreak{}G7.\allowbreak{}Top.\allowbreak{}bCore}{Formalism.\allowbreak{}ThetaLogCoordinates.\allowbreak{}UniformThetaSubstitution.\allowbreak{}KOneConsequence}{32--32}
\leanref{Formalism.\allowbreak{}InvariantFactorPaths.\allowbreak{}G7.\allowbreak{}Top.\allowbreak{}entryB}{Formalism.\allowbreak{}ThetaLogCoordinates.\allowbreak{}UniformThetaSubstitution.\allowbreak{}KOneConsequence}{45--45}
\leanref{Formalism.\allowbreak{}InvariantFactorPaths.\allowbreak{}G7.\allowbreak{}EntryB\_\allowbreak{}of}{Formalism.\allowbreak{}ThetaLogCoordinates.\allowbreak{}UniformThetaSubstitution.\allowbreak{}FrobeniusLiftNorms.\allowbreak{}EntryATwoAndB}{268--345}

\begin{proof}
For a nonzero entry take the multiplicities of \ref{B3.21}.
Let \(\mathcal P\) consist of the even-source coordinates \((s,b)\)
with \(b\notin J\) and the odd-source coordinates with \(b\in J\).
The endpoints of a pure-top term have opposite parity and share a
column label, so exactly one belongs to \(\mathcal P\). Each occurrence
consumes at least one unit of top degree there, as does a \(B\)-term
whose source coordinate belongs to \(\mathcal P\). Summing the top
budget over \(\mathcal P\) gives
\begin{equation}\label{eq:b3-top-column-payment}
 \sum_\phi m_\phi+
 \sum_{\kappa:\ (\operatorname{source}(\kappa),
                         \operatorname{col}(\kappa))\in\mathcal P}
                     \beta_\kappa
 \le\sum_{(s,b)\in\mathcal P}e_{s,b}.
\end{equation}

The even-source \(B\)-terms left outside this payment have labels
in \(J\) and odd targets. Group their directions by target.
By the nonvanishing in \eqref{eq:b3-entry-hasse-nonzero},
\eqref{eq:b3-hits-columns} applies at each such target with label set \(J\). Summing the bounds gives
\[
 \sum_{\substack{\kappa:\ \operatorname{source}(\kappa)\ {\rm even}\\
                              \operatorname{col}(\kappa)\in J}}
       \beta_\kappa
 \le
 \sum_{c:\ s_c\ {\rm odd}}\#(\operatorname{im}I_c\cap J).
\]
Together with \eqref{eq:b3-top-column-payment}, this gives
\begin{equation}\label{eq:b3-full-column-payment}
 \sum_\phi m_\phi+
 \sum_{\kappa:\ \operatorname{source}(\kappa)\ {\rm even}}
       \beta_\kappa
 \le N_J.
\end{equation}

For odd-source \(B\)-terms use \eqref{eq:b3-arrival-source-min}.
Add \(N_J\) to the common estimate \eqref{eq:b3-entry-star} and use
\eqref{eq:b3-full-column-payment} to obtain
\[
 N_J+v_p(\mathcal J)\ge
 \gamma_p\sum_{s\in T}\alpha_s
 -\sum_{s\in T_{\mathrm{odd}}}
       \sum_{l\ge0}\min(\alpha_sp^{-(2l+1)},C_{s+2l+1})
 -\delta+\rho_pD.
\]
This is \(v_p(p^{N_J}\mathcal J)\ge g_B+\rho_pD\). The norm
bound also holds for a zero entry.
\end{proof}

\section{From entries to coefficients of the frame jet}

At each \(s\) take \(c_s\in\N\) columns of heights
\(h_{s,c}\in\N\), \(h_{s,c}\le r_0\) for \(0\le c<c_s\).
Choose maps
\[
 I_{s,c}:\{0,\ldots,h_{s,c}-1\}
                  \longrightarrow\{0,\ldots,r_1-1\}.
\]
Use independent frame variables
\(\tau_{s,i,a}\), \(s\in S\), \(i,a<r_0\), and define
\begin{equation}\label{eq:b3-frame-polynomial}
 \mathcal P(\tau,z)=
 \prod_{s\in S}\prod_{c<c_s}
  \det\left(
       \sum_{a<r_0}\tau_{s,i,a}z_{s,a,I_{s,c}(j)}
      \right)_{0\le i,j<h_{s,c}}.
\end{equation}
For a power series \(\Theta(z)\), its frame jet is the polynomial
\begin{equation}\label{eq:b3-frame-jet}
 \mathcal J_\Theta(\tau)
   =\sum_n n!\,[z^n]\Theta\,[z^n]\mathcal P(\tau,z).
\end{equation}
This is the coefficientwise version of the operator in
\ref{B3.10} and \ref{B3.11}. For these diagrams put
\begin{align}
 \alpha_s&=
   \begin{cases}
    \#\{c<c_s:h_{s,c}=1\},&s\in T,\\
    0,&s\notin T,
   \end{cases}\label{eq:b3-gauss-alpha}\\
 C_s&=c_s-\alpha_s,\qquad
 L_s=\sum_{c<c_s}h_{s,c}-\alpha_s,\qquad
 D=\sum_s(\alpha_s+L_s)
   =\sum_s\sum_{c<c_s}h_{s,c}.
       \label{eq:b3-gauss-counts}
\end{align}
Use these counts in \(\delta,g_A,g_B,A_{\mathrm{ev}}\).

\begin{theorem}[Frame-jet coefficients of the unscaled kernel]\label{B3.26}
Let \(K\) be a characteristic-zero nonarchimedean normed field.
Fix a prime \(p>2\) with \(\lvert p\rvert=p^{-1}\).
Take \(\pi\in K\) satisfying \(\pi^{p^2-1}=p\). Take \(\nu\ge1\), \(r_0,r_1\in\N\),
\(u_s\in\mathcal O_K\), and \(H\) as in \eqref{eq:b3-hh}. Suppose
\(\Theta(\pi z)\in\mathcal C_p(H)\). For the diagrams and maps in
\eqref{eq:b3-frame-polynomial}, each coefficient
\(a_\eta=[\tau^\eta]\mathcal J_\Theta\) satisfies
\begin{align}
 |a_\eta|&\le p^{-g_A},\label{eq:b3-kernel-a1}\\
 |a_\eta|&\le p^{-(g_B-A_{\mathrm{ev}})}.\label{eq:b3-kernel-a2}
\end{align}
For every integer \(i\ge1\), set
\[
 J_i=\{b:0\le b<r_1,\ b<i\},\qquad
 q_{s,c}(i)=\#(\operatorname{im}I_{s,c}\cap J_i),
\]
and
\begin{equation}\label{eq:b3-kernel-edge-exponent}
 M_i=
 \sum_{\substack{s\in S\\s\ {\rm even}}}
        \sum_{c<c_s}\bigl(\min(h_{s,c},i)-q_{s,c}(i)\bigr)
 +\sum_{\substack{s\in S\\s\ {\rm odd}}}
        \sum_{c<c_s}q_{s,c}(i).
\end{equation}
Then
\begin{equation}\label{eq:b3-kernel-b}
 |p^{M_i}a_\eta|\le p^{-g_B}.
\end{equation}
The differences in \eqref{eq:b3-gauss-counts} and
\eqref{eq:b3-kernel-edge-exponent} are nonnegative integers.
\end{theorem}
\leanref{Formalism.\allowbreak{}InvariantFactorPaths.\allowbreak{}G7.\allowbreak{}kernel\_\allowbreak{}a\_\allowbreak{}one\_\allowbreak{}generic}{Formalism.\allowbreak{}ThetaLogCoordinates.\allowbreak{}UniformThetaSubstitution.\allowbreak{}FrobeniusLiftNorms.\allowbreak{}EGaussK}{213--215}
\leanref{Formalism.\allowbreak{}InvariantFactorPaths.\allowbreak{}G7.\allowbreak{}kernel\_\allowbreak{}a\_\allowbreak{}two\_\allowbreak{}generic}{Formalism.\allowbreak{}ThetaLogCoordinates.\allowbreak{}UniformThetaSubstitution.\allowbreak{}FrobeniusLiftNorms.\allowbreak{}EGaussK}{216--218}
\leanref{Formalism.\allowbreak{}InvariantFactorPaths.\allowbreak{}G7.\allowbreak{}kernel\_\allowbreak{}b\_\allowbreak{}generic}{Formalism.\allowbreak{}ThetaLogCoordinates.\allowbreak{}UniformThetaSubstitution.\allowbreak{}FrobeniusLiftNorms.\allowbreak{}EGaussK}{219--220}

\begin{proof}
Let \(F\in\mathcal C_p(H)\). For \(A(\tau)\in K[\tau]\), put
\[
 \|A\|_{\mathrm G}=\max_\eta|[\tau^\eta]A|,
 \qquad \|0\|_{\mathrm G}=0.
\]
To prove multiplicativity, divide each of two nonzero polynomials
by a coefficient of maximal norm. Each normalized polynomial is
integral and has nonzero reduction modulo the maximal ideal of
\(\mathcal O_K\). Their reduced product is nonzero because the residue
polynomial ring is a domain. The normalized product therefore has
Gauss norm one. Rescaling proves multiplicativity.

The Gauss norm extends to the rational function field \(K(\tau)\).
Write \(Y_{s,i,a}\) for its coordinate functions. For every
\(A\in K[\tau]\), evaluation at \(Y\) gives
\begin{equation}\label{eq:b3-gauss-generic-eval}
 |A(Y)|=\|A\|_{\mathrm G}.
\end{equation}
Taking \(A\) constant gives the isometric inclusion of \(K\).
Taking \(A=\tau_{s,i,a}\) gives \(|Y_{s,i,a}|=1\).

Evaluate the frame variables in \(\mathcal P\) at \(Y\).
Separate the height-one factors at marked positions. Each has
the form \(\sum_{a<r_0}Y_{s,0,a}z_{s,a,b}\). Choose
\(\xi_{s,a}=Y_{s,0,a}\) and let \(e_{s,b}\) count the separated
factors at \(s\) with label \(b\). If \(r_0=0\), the vectors are
empty and there are no height-one factors. Leave the remaining
determinants in \(Q_0\). Thus \(\mathcal P(Y,z)=P_{\mathrm{ent}}(z)\)
of \eqref{eq:b3-entry-polynomial}, with integral frame and direction
coordinates and counts \eqref{eq:b3-gauss-alpha} and
\eqref{eq:b3-gauss-counts}. The isometric extension preserves the
expansion-class identities for \(F\).

Evaluation in the finite frame-jet sum gives
\(\mathcal J_F(Y)=\langle P_{\mathrm{ent}},F\rangle\).
By \eqref{eq:b3-gauss-generic-eval},
\begin{equation}\label{eq:b3-coeff-to-entry}
 |[\tau^\eta]\mathcal J_F|
 \le\|\mathcal J_F\|_{\mathrm G}
 =|\langle P_{\mathrm{ent}},F\rangle|.
\end{equation}
Apply \ref{B3.23} and \ref{B3.24} to this entry to get
\begin{align}
 |[\tau^\eta]\mathcal J_F|
       &\le p^{-(g_A+\rho_pD)},\label{eq:b3-scaled-a1}\\
 |[\tau^\eta]\mathcal J_F|
       &\le p^{-(g_B-A_{\mathrm{ev}}+\rho_pD)}.
          \label{eq:b3-scaled-a2}
\end{align}
\leanref{Formalism.\allowbreak{}InvariantFactorPaths.\allowbreak{}G7.\allowbreak{}EGaussK.\allowbreak{}reduce}{Formalism.\allowbreak{}ThetaLogCoordinates.\allowbreak{}UniformThetaSubstitution.\allowbreak{}FrobeniusLiftNorms.\allowbreak{}EGaussK}{38--101}

For the prescribed-column bound, apply \ref{B3.25} with \(J=J_i\).
A separated height-one factor contributes \(1\) to \(N_{J_i}\) if
its position is even and its label lies outside \(J_i\), or if its
position is odd and its label lies in \(J_i\). Since \(i\ge1\), these
are exactly its contributions to \(M_i\). Each remaining determinant
at an odd position contributes \(q_{s,c}(i)\) to both. At an even
position, a remaining determinant adds
\(\min(h_{s,c},i)-q_{s,c}(i)\) to \(M_i\). This is nonnegative
because \(q_{s,c}(i)\le h_{s,c}\) and \(q_{s,c}(i)\le |J_i|\le i\).
Hence
\begin{equation}\label{eq:b3-edge-payment-comparison}
 N_{J_i}\le M_i.
\end{equation}
Since \(|p|\le1\), \ref{B3.25} and \eqref{eq:b3-coeff-to-entry} give
\begin{equation}\label{eq:b3-scaled-b}
 |p^{M_i}[\tau^\eta]\mathcal J_F|
 \le p^{-(g_B+\rho_pD)}.
\end{equation}

Take \(F(z)=\Theta(\pi z)\). The product \(\mathcal P(\tau,z)\)
is homogeneous of \(z\)-degree \(D\), since each determinant
contributes its height. Substitution multiplies \([z^n]\Theta\)
by \(\pi^{|n|}\), and the frame pairing uses terms of degree \(D\).
Together with \(\pi^{p^2-1}=p\) this gives
\begin{align}
 \mathcal J_{\Theta(\pi z)}&=\pi^D\mathcal J_\Theta,
       \label{eq:b3-frame-scaling}\\
 |\pi^D|&=p^{-\rho_pD}.\label{eq:b3-pi-degree}
\end{align}
Dividing each scaled bound by \(|\pi^D|=p^{-\rho_pD}\) proves
the three asserted estimates.
\leanref{Formalism.\allowbreak{}InvariantFactorPaths.\allowbreak{}G7.\allowbreak{}EGaussK.\allowbreak{}descend}{Formalism.\allowbreak{}ThetaLogCoordinates.\allowbreak{}UniformThetaSubstitution.\allowbreak{}FrobeniusLiftNorms.\allowbreak{}EGaussK}{132--144}
\end{proof}

For the target diagrams, \ref{B3.27} uses the ordinary coefficient
bound. The bounds with exponents \(g_B-A_{\mathrm{ev}}\) and \(g_B\)
give the two estimates of \ref{B3.28}.


% ---- p3c3 ----
\chapter{One lift on a lattice and the Hecke commutation}\label{ch:b3-one-lift}

Fix a local system as in Definition~\ref{B3.01}, satisfying Definition~\ref{B3.05a} \cite{alpoge-2026}. Let \(\iota:O\to K\) be the coefficient map and let \(\varpi\) be a uniformizer of \(O\), with \(p\sim\varpi^{e_O}\). Set \(k=k_1+1\) and \(N_j=k+j(t+1)\). At level \(0\le j\le m\), an element of \(\mathcal H_j\) assigns a polynomial over \(K\) in \(X_{v,s,a,b}\) to each component and submodule tuple. The indices \(v,s\) range over the places and their embeddings, and \(1\le a,b\le N_j\). A form in \(\Lambda_j\) is level-equivariant and vanishes off lattice tuples. Its value at each lattice tuple belongs to the integral weight lattice there. Put \(V_j=\operatorname{span}_K\Lambda_j\) and, for \(g\in\R\), write
\[
 v_j(f)\ge g\quad\Longleftrightarrow\quad
 f\in\iota(\varpi)^{\lceil e_Og\rceil}\Lambda_j .
 \tag{3.27.1}
\]
For the lifts \(K_j:\mathcal H_j\to\mathcal H_{j+1}\), put \(K^{[0]}=1\) and \(K^{[b]}=K_{b-1}\cdots K_0\) for \(1\le b\le m\).

Let \(\mathcal P\) be the finite set of local positions. At \(v\in\mathcal P\), the embeddings are indexed cyclically by \(s=0,\ldots,n'_v-1\), with \(n'_v\ge2\) even. The discrete valuation ring \(W_v\) has uniformizer \(p\) and residue field of cardinality \(p^{n'_v}\). Set \(\eta_v=n'_v/2\). Write the initial weight as \(\lambda_{v,s,1}\ge\cdots\ge\lambda_{v,s,k}\) and the level weights as \(w_{v,b}(s)\). When \(\mathcal P\ne\varnothing\), standardness gives \(p\ge5\) and \(1\le r\le k/2\). The weights satisfy the lower bounds and monotonicity in Definition~\ref{B3.01}. At step \(j\), abbreviate
\[
 A_v=g_{A,v}(j+1),\quad B_v=g_{B,v}(j+1),\quad
 E_v=\sum_{\substack{s\in T_v\\s\ {\rm even}}}\alpha_{v,s},\quad
 R_j(v)=\sum_{\substack{v'\in\mathcal P\\v'\ne v}}A_{v'}-d_G(j).
 \tag{3.27.2}
\]
The marked set \(T_v\) and the integers \(\alpha_{v,s}\) are as in Definitions~\ref{B3.01} and~\ref{B3.02}. The functions \(g_A,g_B\) are given in Definitions~\ref{B3.13} and~\ref{B3.14}, and \(d_G(j)\) bounds the \(p\)-valuation of the stabilizer orders.

\section{From the abstract kernel to the actual weights}

Let \(K_c\) and \(O_c\) be the kernel coefficient field and ring. Identify \(K\) with its image in \(K_c\). In the convention of Definition~\ref{B3.12},
\[
 \nu_c(z)\ge g\quad\Longleftrightarrow\quad
 z^bp^{-a}\in O_c
 \quad\text{for all }a\in\Z,\ b\in\Z_{\ge1}\text{ with }a\le bg .
 \tag{3.27.3}
\]
For an ultrametric norm with \(\|p\|=p^{-1}\) and unit ball \(O_c\), this means \(\|z\|\le p^{-g}\). Call \(U\) integral after logarithmic substitution if \(U(\ell(X))\in O_c[[X]]\), with the formal group logarithm \(\ell\) applied coordinatewise. For the pairing \(b\) of Definition~\ref{B3.05}, put \(\Phi_v=\sum_{e\in E_v^{\rm pair}}b(X_{v,e},Y_{v,e})\) in the full variable set, where \(E_v^{\rm pair}\) is the index set of that definition at \(v\). Let \(\mathcal C_j\) denote the change to the chosen \(Z\)-coordinates and \(P_{I,v}\) the determinant polynomial for column choice \(I\) in Definition~\ref{B3.10}. The differential pairing of Definition~\ref{B3.11} gives
\[
 \mathcal J_{I,v}(U)=
 \left.P_{I,v}(\partial_Z)\mathcal C_j(\exp(-\Phi_v)U)\right|_{Z=0}.
 \tag{3.27.4}
\]

\begin{theorem}[The ordinary coefficient bound for the actual weight diagram]\label{B3.27}
Fix \(v\in\mathcal P\) and \(0\le j<m\). For every \(U\) integral after logarithmic substitution, all column choices \(I\) and frame multi-indices \(n\) satisfy
\[
 \nu_c([\tau^n]\mathcal J_{I,v}(U))\ge A_v .
 \tag{3.27.5}
\]
Write \(\mathrm{LowAt}(j,v,g)\) for this condition with \(g\) in place of \(A_v\).
\end{theorem}
\leanref{Formalism.\allowbreak{}MainNotions.\allowbreak{}FormalGroupSystem.\allowbreak{}Pf.\allowbreak{}first\_\allowbreak{}theorem\_\allowbreak{}place\_\allowbreak{}bound}{Formalism.\allowbreak{}SharpRestriction.\allowbreak{}GrowthBounds.\allowbreak{}MultisetValuation.\allowbreak{}FirstTheoremPlaceBound}{20--23}
\leanref{Formalism.\allowbreak{}MainNotions.\allowbreak{}FormalGroupSystem.\allowbreak{}Pf.\allowbreak{}eb\_\allowbreak{}jet}{Formalism.\allowbreak{}SharpRestriction.\allowbreak{}GrowthBounds.\allowbreak{}MultisetValuation.\allowbreak{}EbBrJet}{35--39}
\leanref{Formalism.\allowbreak{}MainNotions.\allowbreak{}FormalGroupSystem.\allowbreak{}Pf.\allowbreak{}ea\_\allowbreak{}brNum}{Formalism.\allowbreak{}SharpRestriction.\allowbreak{}GrowthBounds.\allowbreak{}MultisetValuation.\allowbreak{}GainOfCapacityLoss}{252--256}

\begin{proof}
Index embeddings by \(\Z/n'_v\Z\) and write \(\nu_{s,a}=\nu_{j+1,s,a}\) for the target row lengths and \(h_{s,c}=\#\{a:\nu_{s,a}>c\}\) for \(0\le c<\nu_{s,1}\). The last \(t+1\) rows vanish, so \(h_{s,c}\le N_j\). At a marked position delete the \(\nu_{s,1}-\nu_{s,2}=\lambda_{v,s,1}-\lambda_{v,s,2}=\alpha_{v,s}\) height-one columns. The remaining \(C_s=\nu_{s,2}\) columns form the reduced diagram. At an unmarked position \(C_s=\nu_{s,1}\) and \(\alpha_{v,s}=0\). The count \(L_s=\sum_a\nu_{s,a}-\alpha_{v,s}\) agrees with Definition~\ref{B3.02}. Cyclic indexing preserves parity and distances, so the exponent of Theorem~\ref{B3.26} is \(A_v\).

Definition~\ref{B3.10} identifies the determinants of Theorem~\ref{B3.26} with \(P_{I,v}\) by injective renamings of the local \(Z\) and frame variables. Restriction to \(v\) commutes with the pairing, so every coefficient of (3.27.4) is zero or a coefficient of that pairing.

Since \(O\) is finite over \(W_v\), \(p\) lies in its maximal ideal, so \(e_O>0\). By Lemma~\ref{B3.30}, \(O_c\cap K=\iota(O)\) and \(\iota\) is injective, hence \(p^{-1}\notin O_c\). Extend the rank-one valuation normalized by \(\nu_c(p)=1\) to an algebraic closure of \(K_c\) and take \(\pi^{p^2-1}=p\). Write \(U^\ell(X)=U(\ell(X))\) and let \(L\) be the tuple of integral linear coordinate forms. As in Theorem~\ref{B3.09}, \(\lambda_\pi(T)=\pi^{-1}\ell(\pi T)\) has an integral compositional inverse \(\mu\). Thus \(\ell(\pi\mu(L))=\pi L\) and
\[
 U(\pi L)=U^\ell(\pi\mu(L)).
\]
This series is integral. By Theorem~\ref{B3.08} the scaled phase exponential is \(\exp(-\pi^{p+1}H_u/p)V\) with \(V\) integral and \((-\pi^{p+1})^d/d!\) integral for every \(d\). The \(d\)-th term of the scaled product is \(p^{-d}H_u^d[(-\pi^{p+1})^d/d!]VU(\pi L)\). Restricting to \(v\) gives the expansion in Theorem~\ref{B3.26}. That theorem bounds the norm of each paired coefficient by \(p^{-A_v}\), so (3.27.3) gives (3.27.5).
\end{proof}

Fix \(1\le i\le r\). For a determinant of height \(h=h_{s,c}\), let \(d_{s,c}(I,i)=\#\{a:0\le a<h,\ I(s,c,a)<i\}\). The target labels start at zero. Put
\[
 P_{I,v}^{(i)}=p^{\sum_{s,c}\epsilon_{s,c}(I,i)}P_{I,v},\qquad
 \epsilon_{s,c}(I,i)=
 \begin{cases}
 d_{s,c}(I,i),&s\ {\rm odd},\\
 [\min(h_{s,c},i)-d_{s,c}(I,i)]_+,&s\ {\rm even}.
 \end{cases}
 \tag{3.28.1}
\]

\begin{theorem}[The shifted ordinary bound and the selected-column bound]\label{B3.28}
For \(v\in\mathcal P\) and \(0\le j<m\),
\[
 \nu_c([\tau^n]\mathcal J_{I,v}(U))\ge B_v-E_v .
 \tag{3.28.2}
\]
For \(1\le i\le r\), replacing \(P_{I,v}\) by \(P_{I,v}^{(i)}\) gives
\[
 \nu_c\!\left([\tau^n]\left.
 P_{I,v}^{(i)}(\partial_Z)\mathcal C_j(\exp(-\Phi_v)U)\right|_{Z=0}\right)
 \ge B_v .
 \tag{3.28.3}
\]
Both bounds hold for all \(I,n\) and all \(U\) integral after logarithmic substitution.
\end{theorem}
\leanref{Formalism.\allowbreak{}MainNotions.\allowbreak{}FormalGroupSystem.\allowbreak{}Pf.\allowbreak{}second\_\allowbreak{}lemma\_\allowbreak{}place\_\allowbreak{}bound}{Formalism.\allowbreak{}SharpRestriction.\allowbreak{}GrowthBounds.\allowbreak{}MultisetValuation.\allowbreak{}FirstTheoremPlaceBound}{24--27}
\leanref{Formalism.\allowbreak{}MainNotions.\allowbreak{}FormalGroupSystem.\allowbreak{}Pf.\allowbreak{}hecke\_\allowbreak{}place\_\allowbreak{}bound}{Formalism.\allowbreak{}SharpRestriction.\allowbreak{}GrowthBounds.\allowbreak{}MultisetValuation.\allowbreak{}FirstTheoremPlaceBound}{28--31}

\begin{proof}
Use the diagram and scaled expansion of Theorem~\ref{B3.27}. The second ordinary estimate of Theorem~\ref{B3.26} gives \(g_B-A_{\rm ev}=B_v-E_v\). For the selected-column estimate take \(J_0=\{0,\ldots,i-1\}\). A determinant with a repeated target label vanishes. Otherwise \(d_{s,c}(I,i)=|\operatorname{im}I_{s,c}\cap J_0|\le\min(h_{s,c},i)\), where \(I_{s,c}\) is restricted to the labels of the determinant. Thus (3.28.1) gives the exponent \(M_i\) of Theorem~\ref{B3.26}. The selected-column estimate there and (3.27.3) give (3.28.3).
\end{proof}

\section{Products of positions and the class expansion}

For \(0\le i_v\le r\), put \(\mathbf i=(i_v)_v\) and \(P_{I,\mathbf i}=\prod_vP_{I,v}^{(i_v)}\), with \(P_{I,v}^{(0)}=P_{I,v}\). Write \(L=(c,\mathbf L)\) and \(L'=(c',\mathbf L')\) for the source and target component--lattice pairs. For a frame \(\phi\) for this pair, write \(\Theta=\Theta_{j,L,L',\phi}\) and set
\[
 \mathcal J_{I,\mathbf i}(\Theta)
 =\left.P_{I,\mathbf i}(\partial_Z)\mathcal C_j(\Theta)\right|_{Z=0}.
 \tag{3.29.1}
\]

\begin{theorem}[Addition of the coefficient exponents]\label{B3.29}
Fix \(0\le j<m\) and source and target component--lattice pairs \(L,L'\). Choose a frame \(\phi\) for this pair. There is a series \(U_1\), integral after logarithmic substitution, with
\[
 \Theta_{j,L,L',\phi}=\left(\prod_{v\in\mathcal P}\exp(-\Phi_v)\right)U_1.
 \tag{3.29.2}
\]
Choose \(q_v\in\Q\) at each place under the conditions
\[
 \begin{aligned}
 i_v=0&\quad\Longrightarrow\quad
 \mathrm{LowAt}(j,v,q_v)\ \text{or}\ q_v\le B_v-E_v,\\
 1\le i_v\le r&\quad\Longrightarrow\quad q_v\le B_v .
 \end{aligned}
 \tag{3.29.3}
\]
Then every coefficient satisfies
\[
 \nu_c([\tau^n]\mathcal J_{I,\mathbf i}(\Theta))\ge\sum_{v\in\mathcal P}q_v.
 \tag{3.29.4}
\]
\end{theorem}
\leanref{Formalism.\allowbreak{}MainNotions.\allowbreak{}FormalGroupSystem.\allowbreak{}Pf.\allowbreak{}kernel\_\allowbreak{}product\_\allowbreak{}form}{Formalism.\allowbreak{}SharpRestriction.\allowbreak{}GrowthBounds.\allowbreak{}MultisetValuation.\allowbreak{}RestEntOf}{20--59}
\leanref{Formalism.\allowbreak{}MainNotions.\allowbreak{}FormalGroupSystem.\allowbreak{}Pf.\allowbreak{}restEnt\_\allowbreak{}of}{Formalism.\allowbreak{}SharpRestriction.\allowbreak{}GrowthBounds.\allowbreak{}MultisetValuation.\allowbreak{}RestEntOf}{62--64}
\leanref{Formalism.\allowbreak{}MainNotions.\allowbreak{}FormalGroupSystem.\allowbreak{}Pf.\allowbreak{}en\_\allowbreak{}places\_\allowbreak{}add}{Formalism.\allowbreak{}SharpRestriction.\allowbreak{}GrowthBounds.\allowbreak{}MultisetValuation.\allowbreak{}KernelProductForm}{51--110}

\begin{proof}
By the generator condition of Definition~\ref{B3.05a}, \(\Theta=\exp(-\Phi)U_0\) with \(U_0\) integral after logarithmic substitution. For \(\Delta=\Phi-\sum_v\Phi_v\), phase equivalence makes \(\exp(\Delta(\ell(X)))\) an integral unit of constant term one. Thus \(U_1=\exp(-\Delta)U_0\) gives (3.29.2).

If \(\mathcal P\) is empty, the expression is the constant coefficient of \(U_1\), which belongs to \(O_c\) and satisfies (3.29.4) with bound zero. Otherwise use the normalized coefficient norm of Theorem~\ref{B3.27}. Let \(\mathfrak e\) be the coordinatewise compositional inverse of \(\ell\). Writing \(U_1(\ell(X))=\sum_\beta a_\beta X^\beta\), with \(a_\beta\in O_c\), gives
\[
 U_1(Y)=\sum_\beta a_\beta\prod_v U_{\beta,v}(Y_v),\qquad
 U_{\beta,v}(Y_v)=\mathfrak e(Y_v)^{\beta_v}.
 \tag{3.29.5}
\]
The coordinate change respects positions, so
\[
 \mathcal J_{I,\mathbf i}(\Theta)=
 \sum_\beta a_\beta\prod_v\left.
 P_{I,v}^{(i_v)}(\partial_Z)\mathcal C_j(\exp(-\Phi_v)U_{\beta,v})
 \right|_{Z=0}.
 \tag{3.29.6}
\]
Since \(\mathfrak e(0)=0\), a contributing term has \(|\beta|\le\deg_Z P_{I,\mathbf i}\). Each \(U_{\beta,v}\) is integral after logarithmic substitution. Thus the local factors have coefficient bounds \(q_v\) by (3.29.3) and Theorem~\ref{B3.28}. Products add the bounds, and the ultrametric inequality preserves that sum over the finite expansion with \(a_\beta\in O_c\).
\end{proof}

Write \(\mathrm{Ent}(j,\mathbf i,g)\) for the uniform coefficient bound \(g\) in (3.29.1). For a list \(v_1,\ldots,v_a\) of distinct places containing the support of \(\mathbf i\), use the operators of Definition~\ref{B3.35} to form
\[
 \widetilde T_{\mathbf i}^{(j+1)}
 =\widetilde T_{v_1,i_{v_1}}^{(j+1)}\cdots
   \widetilde T_{v_a,i_{v_a}}^{(j+1)} .
 \tag{3.30.1}
\]

\begin{lemma}[The loss in summing lattice classes]\label{B3.30}
Let \(0\le j<m\), \(0\le i_v\le r\), and use a list as in (3.30.1). If \(\mathrm{Ent}(j,\mathbf i,g)\) holds for \(g\in\R\), then every \(f\in\Lambda_j\) satisfies
\[
 v_{j+1}(\widetilde T_{\mathbf i}^{(j+1)}K_jf)\ge g-d_G(j).
 \tag{3.30.2}
\]
The class expansion has a common factor \(c_K\in O_c\), which is a unit if nonzero. The nonterminal coefficients are \(c_K/|\Gamma_{j,L}|\), where \(\Gamma_{j,L}\) is the stabilizer of \(L\) and \(v_p(|\Gamma_{j,L}|)\le d_G(j)\). At the terminal step the positive integer denominators \(w_x\) satisfy \(v_p(w_x)\le d_G(j)\).
\end{lemma}
\leanref{Formalism.\allowbreak{}MainNotions.\allowbreak{}FormalGroupSystem.\allowbreak{}ExpansionEquality}{Formalism.\allowbreak{}SharpRestriction.\allowbreak{}GrowthBounds.\allowbreak{}FormalGroupSystem}{566--571}
\leanref{Formalism.\allowbreak{}MainNotions.\allowbreak{}FormalGroupSystem.\allowbreak{}ExpansionDecomposition}{Formalism.\allowbreak{}SharpRestriction.\allowbreak{}GrowthBounds.\allowbreak{}FormalGroupSystem}{572--586}
\leanref{Formalism.\allowbreak{}MainNotions.\allowbreak{}FormalGroupSystem.\allowbreak{}Pf.\allowbreak{}formB\_\allowbreak{}of}{Formalism.\allowbreak{}SharpRestriction.\allowbreak{}GrowthBounds.\allowbreak{}MultisetValuation.\allowbreak{}KernelSumOfClassData}{255--278}

\begin{proof}
The class data for fixed \(j\) apply to every \(f\in V_j\). At a target lattice pair \(L'\) extend coefficients to \(K_c\). Let \(\operatorname{Rep}_j\) contain one representative of each orbit of lattice tuples under the level group in each component. Before the terminal step,
\[
 (K_jf)(L')=\sum_{L\in\operatorname{Rep}_j}
 \frac{c_K}{|\Gamma_{j,L}|}\operatorname{ker}_{j,L,L'}(f(L)).
 \tag{3.30.4}
\]
At the terminal step,
\[
 (K_jf)(L')=\sum_x\frac{c_K}{w_x}\sum_y
 \operatorname{ker}_{j,L_{x,y},L'}(f(L_{x,y})).
 \tag{3.30.5}
\]
Both sums are finite.

For \(f\in\Lambda_j\), (3.30.1) takes one lower neighbour at each place with \(i_v>0\). For each source \(L\) and resulting target \(L''\), Definition~\ref{B3.05a} supplies a frame where the entry bound applies. The frame identity and the weight-lattice condition put \(\bigl(\prod_v p^{N'_{v,j+1}(i_v)}\bigr)\operatorname{ker}_{j,L,L''}(f(L))\) in the \(O_c\)-span of \(az\) with \(\nu_c(a)\ge g\) and \(z\in H_{j+1}(L')\), the integral weight lattice at \(L'\). For each denominator \(w\), \(p^{d_G(j)}/w\in O_c\) since the prime-to-\(p\) part is a unit. Thus the weighted sum has such an expansion with \(h=g-d_G(j)\) in place of \(g\).

Let \(y\) be defined over \(K\) with an expansion
\[
 y=\sum_\alpha a_\alpha z_\alpha,\qquad
 z_\alpha\in H_{j+1}(L'),\qquad \nu_c(a_\alpha)\ge h .
 \tag{3.30.7}
\]
Put \(R=\iota(O)=O_c\cap K\). Since \(R\) is a valuation ring of \(K\), in a nonzero relation \(\sum_\alpha b_\alpha z_\alpha=0\) over \(K\) we can choose \(b_\gamma\ne0\) with all \(b_\alpha/b_\gamma\in R\). Eliminating \(z_\gamma\) preserves the bounds on the \(a_\alpha\). Repeat until the remaining \(z_\alpha\) are \(K\)-linearly independent. An invertible coordinate minor gives \(a_\alpha\in K\).

If \(e_O>0\), any nonzero kernel ideal of \(\iota\) would contain a power of \(p\), which is nonzero in \(K_c\). Thus \(\iota\) is injective and \(K=\operatorname{Frac}(O)\), since \(R\) is a valuation ring of \(K\). For \(a_\alpha\ne0\), (3.27.3) gives \(b\operatorname{ord}_\varpi(a_\alpha)\ge e_Oa\) for every rational \(a/b\le h\). Hence \(\operatorname{ord}_\varpi(a_\alpha)\ge\lceil e_Oh\rceil\). If \(e_O=0\), \(p\) is a unit and (3.27.3) with \(a=\lfloor h\rfloor\), \(b=1\) gives \(a_\alpha\in R\). Thus \(y\in\iota(\varpi)^{\lceil e_Oh\rceil}H_{j+1}(L')\) in either case. The lift takes \(V_j\) into \(V_{j+1}\), and the neighbour operators preserve equivariance and vanishing off lattice tuples, so the fibre bounds give (3.30.2).
\end{proof}

\begin{lemma}[Real exponents and the lattice endpoint]\label{B3.31}
Let \(0\le j<m\), \(0\le i_v\le r\), and let \(q_v\in\R\) satisfy (3.29.3). For a list as in (3.30.1) and every \(f\in\Lambda_j\),
\[
 v_{j+1}(\widetilde T_{\mathbf i}^{(j+1)}K_jf)\ge\sum_vq_v-d_G(j).
 \tag{3.31.1}
\]
\end{lemma}
\leanref{Formalism.\allowbreak{}MainNotions.\allowbreak{}FormalGroupSystem.\allowbreak{}Pf.\allowbreak{}ent\_\allowbreak{}strict}{Formalism.\allowbreak{}SharpRestriction.\allowbreak{}GrowthBounds.\allowbreak{}MultisetValuation.\allowbreak{}RamifiedLatticeOfChainEstimates}{46--64}
\leanref{Formalism.\allowbreak{}MainNotions.\allowbreak{}FormalGroupSystem.\allowbreak{}Pf.\allowbreak{}formB\_\allowbreak{}lim}{Formalism.\allowbreak{}SharpRestriction.\allowbreak{}GrowthBounds.\allowbreak{}MultisetValuation.\allowbreak{}RamifiedLatticeOfChainEstimates}{65--89}
\leanref{Formalism.\allowbreak{}MainNotions.\allowbreak{}FormalGroupSystem.\allowbreak{}Pf.\allowbreak{}form\_\allowbreak{}of\_\allowbreak{}place}{Formalism.\allowbreak{}SharpRestriction.\allowbreak{}GrowthBounds.\allowbreak{}MultisetValuation.\allowbreak{}RamifiedLatticeOfChainEstimates}{90--95}

\begin{proof}
Put \(s=\sum_vq_v\). For \(s'<s\), choose \(\epsilon>0\) with \(|\mathcal P|\epsilon<s-s'\) and rationals \(q_v-\epsilon<q'_v<q_v\). Then \(\sum_vq'_v\ge s'\), and decreasing the local bounds preserves (3.29.3). Theorem~\ref{B3.29} gives \(\mathrm{Ent}(j,\mathbf i,s')\), and Lemma~\ref{B3.30} gives the lattice bound \(s'-d_G(j)\).

Set \(x=e_O(s-d_G(j))\) and \(\delta=(x-\lceil x\rceil+1)/(e_O+1)>0\). Then \(\lceil x\rceil-1<x-e_O\delta\le\lceil x\rceil\), so
\[
 \lceil e_O(s-\delta-d_G(j))\rceil=\lceil e_O(s-d_G(j))\rceil .
 \tag{3.31.2}
\]
The bound with \(s'=s-\delta\) gives the desired lattice.
\end{proof}

\section{The three single-step estimates}

\begin{theorem}[The ordinary lift of an integral input]\label{B3.32}
For \(0\le j<m\) and \(f\in\Lambda_j\),
\[
 v_{j+1}(K_jf)\ge\sum_{v\in\mathcal P}A_v-d_G(j).
 \tag{3.32.1}
\]
Before the class sum, every coefficient of the ordinary full differential kernel has bound \(\sum_v A_v\). The integral-kernel hypothesis also gives
\[
 v_{j+1}(K_jf)\ge-d_G(j).
 \tag{3.32.2}
\]
\end{theorem}
\leanref{Formalism.\allowbreak{}MainNotions.\allowbreak{}FormalGroupSystem.\allowbreak{}Pf.\allowbreak{}theorem\_\allowbreak{}A\_\allowbreak{}entries\_\allowbreak{}of\_\allowbreak{}corollary}{Formalism.\allowbreak{}SharpRestriction.\allowbreak{}GrowthBounds.\allowbreak{}MultisetValuation.\allowbreak{}RamifiedLatticeOfChainEstimates}{104--115}
\leanref{Formalism.\allowbreak{}MainNotions.\allowbreak{}FormalGroupSystem.\allowbreak{}Pf.\allowbreak{}slopeEstimate\_\allowbreak{}of\_\allowbreak{}lemmaB\_\allowbreak{}slopeLine}{Formalism.\allowbreak{}SharpRestriction.\allowbreak{}GrowthBounds.\allowbreak{}MultisetValuation.\allowbreak{}RamifiedLatticeOfChainEstimates}{116--121}
\leanref{Formalism.\allowbreak{}MainNotions.\allowbreak{}FormalGroupSystem.\allowbreak{}Pf.\allowbreak{}form\_\allowbreak{}valuation\_\allowbreak{}of\_\allowbreak{}bound}{Formalism.\allowbreak{}SharpRestriction.\allowbreak{}GrowthBounds.\allowbreak{}MultisetValuation.\allowbreak{}RamifiedLatticeOfChainEstimates}{96--102}

\begin{proof}
Use \(i_v=0\) and \(q_v=A_v\) from Theorem~\ref{B3.27}. Lemma~\ref{B3.31} with the empty list gives (3.32.1). For the coefficient endpoint, the digit error is integral and the excluded distance set is finite. For large \(e\), \(\alpha_{v,s}p^{-(2e+1)}<1\). The positive-part tail vanishes on residues modulo \(n'_v\) with positive capacity and is geometric with rational ratio on the zero-capacity residues. Thus \(A_v\in\Q\), and Theorem~\ref{B3.29} gives the coefficient bound. Definition~\ref{B3.05a} gives \(\mathrm{Ent}(j,\mathbf 0,0)\), so Lemma~\ref{B3.30} with the empty list gives (3.32.2).
\end{proof}

\begin{theorem}[The distinguished-position bound for an ordinary lift]\label{B3.33}
For \(v\in\mathcal P\), \(0\le j<m\), and \(f\in\Lambda_j\),
\[
 v_{j+1}(K_jf)\ge B_v-E_v+R_j(v).
 \tag{3.33.1}
\]
\end{theorem}
\leanref{Formalism.\allowbreak{}MainNotions.\allowbreak{}FormalGroupSystem.\allowbreak{}Pf.\allowbreak{}hecke\_\allowbreak{}integral\_\allowbreak{}of\_\allowbreak{}forms}{Formalism.\allowbreak{}SharpRestriction.\allowbreak{}GrowthBounds.\allowbreak{}MultisetValuation.\allowbreak{}RamifiedLatticeOfChainEstimates}{122--141}

\begin{proof}
Take \(q_v=B_v-E_v\) from (3.28.2) and \(q_{v'}=A_{v'}\) for \(v'\ne v\) from Theorem~\ref{B3.27}. Apply Lemma~\ref{B3.31} with all indices zero and the empty operator list.
\end{proof}

\begin{theorem}[An output Hecke operator]\label{B3.34}
For \(v\in\mathcal P\), \(0\le j<m\), and \(1\le i\le r\), the normalized operator of Definition~\ref{B3.35} satisfies
\[
 \widetilde T_{v,i}^{(j+1)}\Lambda_{j+1}\subseteq\Lambda_{j+1}.
 \tag{3.34.1}
\]
Every \(f\in\Lambda_j\) satisfies
\[
 v_{j+1}(\widetilde T_{v,i}^{(j+1)}K_jf)\ge B_v+R_j(v).
 \tag{3.34.2}
\]
\end{theorem}
\leanref{Formalism.\allowbreak{}MainNotions.\allowbreak{}FormalGroupSystem.\allowbreak{}Pf.\allowbreak{}latticeChain\_\allowbreak{}succ\_\allowbreak{}of\_\allowbreak{}partial}{Formalism.\allowbreak{}SharpRestriction.\allowbreak{}GrowthBounds.\allowbreak{}MultisetValuation.\allowbreak{}RamifiedLatticeOfChainEstimates}{165--168}
\leanref{Formalism.\allowbreak{}MainNotions.\allowbreak{}FormalGroupSystem.\allowbreak{}Pf.\allowbreak{}restricted\_\allowbreak{}hecke\_\allowbreak{}integral}{Formalism.\allowbreak{}SharpRestriction.\allowbreak{}GrowthBounds.\allowbreak{}MultisetValuation.\allowbreak{}RestrictedHeckeIntegralOfEdges}{289--290}

\begin{proof}
For (3.34.2), use the list \([v]\) in Lemma~\ref{B3.31}, with \(i_v=i\) and \(i_{v'}=0\) for \(v'\ne v\). Use \(q_v=B_v\) from (3.28.3) and \(q_{v'}=A_{v'}\) from Theorem~\ref{B3.27}.

For an edge \(L''\subseteq L_v\) of index \(i\), choose an adapted basis with matrix \(g\) in the fixed coordinates. Then
\[
 L''=gW_v^N,\qquad L_v=gDW_v^N,\qquad
 D=\operatorname{diag}(p^{-1}I_i,I_{N-i}),\qquad N=N_{j+1}.
 \tag{3.34.3}
\]
For one embedding with dominant weight \(\kappa\), put \(e=\kappa_N\), \(\nu_a=\kappa_a-e\), and \(M_i=\sum_{a=1}^i\nu_a\). The integral Weyl module before twisting is spanned by divided-power words on the highest-weight polynomial. Each elementary divided power shifts the row multidegree by a fixed vector, so each word is homogeneous of some degree \(d\) in the first \(i\) rows. The dual Weyl module, which contains the Weyl module, is spanned by products of column minors. A column of height \(h\) contributes at most \(\min(i,h)\) to \(d\), so \(d\le M_i\). Thus \(p^{M_i}\operatorname{rowAct}(D)\) multiplies the word by \(p^{M_i-d}\).

Let \(G\) be the image of \(g\) at this embedding and \(F=(\det G)^{-e}F_0\) a vector in the fibre at \(L''\). Since \(\operatorname{rowAct}(GD)=\operatorname{rowAct}(D)\circ\operatorname{rowAct}(G)\), \(\operatorname{rowAct}(GD)(p^{M_i}F_0)\) is in the integral Weyl module. The identity \((\det(GD))^{-e}p^{M_i}F_0=p^{M_i+ie}F\), with \(M_i+ie=\sum_{a=1}^i\kappa_a\), gives the inclusion at \(L_v\). Integral basis changes preserve the Weyl module and have unit determinant.

At an odd embedding the weight is the negative reverse of the displayed level weight. Write \(\kappa_{s,a}\) for the displayed entries. The tensor product over embeddings therefore has normalization exponent
\[
 \mathcal A_i(\kappa)=
 \sum_{s\ {\rm even}}\sum_{a=1}^{i}\kappa_{s,a}
 -\sum_{s\ {\rm odd}}\sum_{a=N-i+1}^{N}\kappa_{s,a}.
 \tag{3.34.4}
\]
At level \(j+1\) the even terms use the initial entries. The initial lower bound and monotonicity of the appended weights give \(\kappa_{s,a}\ge-w_{v,j+1}(s)\) at odd \(s\). Hence
\[
 \mathcal A_i(\kappa)
 \le\sum_{s\ {\rm even}}\sum_{a=1}^{i}\lambda_{v,s,a}
    +i\sum_{s\ {\rm odd}}w_{v,j+1}(s)
 =N'_{v,j+1}(i).
 \tag{3.34.5}
\]
The normalized sum has integral values in the target fibres and preserves equivariance and vanishing off lattice tuples.
\end{proof}

\section{The operators and the neighbour count}

For full \(W_v\)-lattices an index-\(i\) edge satisfies
\[
 pB\subseteq A\subseteq B,\qquad |B/A|=(p^{n'_v})^i .
 \tag{3.35.1}
\]
The unnormalized neighbour operator at level \(b\) is
\[
 (T_{v,i}^{(b)}f)(c,\mathbf L)=
 \sum_{L''\subseteq L_v\ {\rm of\ index}\ i} f(c,\mathbf L[v\leftarrow L'']),
 \tag{3.35.2}
\]
where \(L_v\) is the \(v\)-entry of \(\mathbf L\).

\begin{definition}[The input operator and the output polynomial]\label{B3.35}
Put
\[
 \mathcal A_v=\sum_{s\ {\rm even}}\sum_{a=1}^{r}\lambda_{v,s,a}
 -\sum_{s\ {\rm odd}}\sum_{a=k-r+1}^{k}\lambda_{v,s,a},
 \qquad y_v=p^{\mathcal A_v}T_{v,r}^{(0)} .
 \tag{3.35.3}
\]
For \(b\ge0\) and \(0\le i\le r\), define
\[
 \begin{split}
 N'_{v,b}(i)&=\sum_{s\ {\rm even}}\sum_{a=1}^{i}\lambda_{v,s,a}
                 +i\sum_{s\ {\rm odd}}w_{v,b}(s),\\
 \widetilde T_{v,0}^{(b)}&=1,\qquad
 \widetilde T_{v,i}^{(b)}=p^{N'_{v,b}(i)}T_{v,i}^{(b)}\quad(i>0).
 \end{split}
 \tag{3.35.4}
\]
For \(u_l,\mathfrak m_l\in K\) with \(u_l\ne0\) for \(1\le l\le b\), set
\[
 U_b=\prod_{l=1}^bu_l,\quad
 \beta_{v,a}=p^{\eta_v(t-2a)}\ (0\le a\le t),\quad
 \mathcal F_{v,b}=(\mathfrak m_l\beta_{v,a})_{1\le l\le b,\ 0\le a\le t}.
 \tag{3.35.5}
\]
Let \(h_\ell\) be the complete homogeneous symmetric polynomial of degree \(\ell\), with \(h_0=1\). For \(0\le\ell\le r\), define \(c_{v,b,\ell}\) and the associated operators by
\[
 \begin{split}
 c_{v,b,\ell}&=
 p^{N'_{v,b}(r)-N'_{v,b}(r-\ell)+\eta_v\ell(N_b-2r+\ell)}
 h_\ell(\mathcal F_{v,b}),\\
 \widetilde T_{v,i}^{\,u,(b)}&=U_b^{-i}\widetilde T_{v,i}^{(b)},\\
 Z_{v,b}&=\sum_{\ell=0}^r(-1)^\ell c_{v,b,\ell}
                       \widetilde T_{v,r-\ell}^{\,u,(b)}.
 \end{split}
 \tag{3.35.6}
\]
\end{definition}
\leanref{Formalism.\allowbreak{}MainNotions.\allowbreak{}G.\allowbreak{}Level.\allowbreak{}yv}{Formalism.\allowbreak{}SharpRestriction.\allowbreak{}GrowthBounds.\allowbreak{}DataG}{148--152}
\leanref{Formalism.\allowbreak{}MainNotions.\allowbreak{}FormalGroupSystem.\allowbreak{}Sys.\allowbreak{}Tt}{Formalism.\allowbreak{}SharpRestriction.\allowbreak{}GrowthBounds.\allowbreak{}FormalGroupSystem}{331--333}
\leanref{Formalism.\allowbreak{}MainNotions.\allowbreak{}FormalGroupSystem.\allowbreak{}Sys.\allowbreak{}cl}{Formalism.\allowbreak{}SharpRestriction.\allowbreak{}GrowthBounds.\allowbreak{}FormalGroupSystem}{344--348}
\leanref{Formalism.\allowbreak{}MainNotions.\allowbreak{}FormalGroupSystem.\allowbreak{}Sys.\allowbreak{}Ttu}{Formalism.\allowbreak{}SharpRestriction.\allowbreak{}GrowthBounds.\allowbreak{}FormalGroupSystem}{349--351}
\leanref{Formalism.\allowbreak{}MainNotions.\allowbreak{}FormalGroupSystem.\allowbreak{}Sys.\allowbreak{}Zop}{Formalism.\allowbreak{}SharpRestriction.\allowbreak{}GrowthBounds.\allowbreak{}FormalGroupSystem}{352--354}

\begin{lemma}[The theta series as a function of a lattice indicator]\label{B3.35a}
Fix \(v\in\mathcal P\) and \(0\le j<m\). Write \(N=N_j\), \(N'=N_{j+1}\), \(\eta=\eta_v\), and \(M=\operatorname{Mat}_{N'\times N}(\operatorname{Frac}W_v)\). For component--lattice pairs \(L,L'\) at these levels put
\[
 I_{L,L'}(x)={\bf1}_{xL_v\subseteq L'_v},\qquad
 s=\operatorname{len}_v(L),\quad s'=\operatorname{len}_v(L'),\quad d=Ns'-N's .
\]
The kernel data assign a parameter \(\omega(L,L')\) to each pair and an additive map \(\Xi_\omega:(M\to\Q)\to K_c[[Z_j]]\) to each parameter, where \(Z_j\) denotes the fixed coordinate variables at level \(j\). Replacing either \(v\)-entry preserves \(\omega\). There are \(a,a'\in K^\times\) of valuation zero, the same for all pairs at \(v,j\), such that the theta series in the fixed coordinates satisfies
\[
 \Theta^F_{L,L'}=a^s(a')^{s'}p^{-\eta d}\,\Xi_{\omega(L,L')}(I_{L,L'}).
\]
\end{lemma}
\begin{proof}
The theta identity is part of Definition~\ref{B3.05a}. At the places of \ref{B8.56} used in \ref{B3.56}, \(K_c=\C_p\) and the length of \(BW_v^n\) is \(-\operatorname{ord}_p\det B\). At these places \(\omega\) is the pair with both \(v\)-entries replaced by the zero submodule.

Take the split prime \(P\) of the CM extension corresponding to \(v\). The level is prime to \(p\), and the finite test support avoids \(P,\bar P\). Choose a norm-one idele \(\pi\) supported over this split place with order one in the chosen factor. Write \(M_\pi,C_\pi\) for the \(p\)-adic images of the step character \(\mu_{j+1}(\pi)\) and the determinant character \(\chi_I(\pi)\). Put \(b_l=\lfloor l(t+1)/2\rfloor\) and \(\Delta=\Delta_v(j+1)\), with \(\Delta_w(l)\) as in (3.36.2). When both component sets are nonempty, the coefficient-field comparison gives \(M_\pi,C_\pi^{b_j},C_\pi^{b_{j+1}}\in K\). Take
\[
 a=p^{-\Delta}M_\pi C_\pi^{b_j},\qquad
 a'=p^\Delta M_\pi^{-1}C_\pi^{-b_{j+1}}.
\]
The step character has valuation \(\sum_w\Delta_w(j+1)\operatorname{ord}_w(\pi)=\Delta\). A positive power of \(\pi\) is principal times an integral norm-one idele. The integral factor has finite-order character value. On the principal factor the character type gives the signed sum of its embedding exponents, and the common shift cancels between even and odd embeddings. The determinant character has valuation zero. A positive power of its norm-one ideal is generated by \(z\) with \(z\bar z=1\), and the product of the CM-type embeddings of \(z\) has norm one because the odd prime \(p\) is inert in the auxiliary quadratic field. Thus both scalars are units. If either component set is empty, take \(a=a'=1\).
\leanref{GlobalTheory.\allowbreak{}TwV.\allowbreak{}Xc.\allowbreak{}two\_\allowbreak{}units}{Development.\allowbreak{}AbelianVarietyTheory.\allowbreak{}Proofs.\allowbreak{}SatakeParameters.\allowbreak{}TwoUnitsAtPlaceUniformizingIdele}{38--82}

On a fibre of \(\omega\), put \(k_i=a^{s_i}(a')^{s'_i}p^{-\eta d_i}\) for pair \(i\). Take a finite rational relation \(\sum_i c_i I_i=0\) and choose a reference pair \(0\) in its nonempty support. The completed global Hom lattice has indicator \(I_i\) after the same change of frames \(x\mapsto B'^{-1}xB\) for every \(i\). Its volume relative to pair \(0\) is \(p^{2\eta(d_i-d_0)}\). Left multiplication by the target basis matrix contributes its determinant once for each of the \(N\) columns. Right multiplication by the inverse source basis matrix contributes its determinant once for each of the \(N'\) rows.

Let \(H_i\) be the complex theta kernel and \(\chi_i\) its nonzero determinant-ideal factor. In the split theta model the move from pair \(0\) to pair \(i\) sends the reference theta vector to the image of \(I_i\) under a common linear map, divided by the square root of the relative volume. The source and target moves compose. The remaining finite theta operator is common because the lattices agree away from \(P\) and the translated test datum agrees. Applying both maps to the indicator relation gives
\[
 \sum_i c_i p^{\eta(d_i-d_0)}\chi_i^{-1}H_i=0 .
\]
Take Taylor coefficients and apply the fixed embedding of the algebraic coefficients into \(\C_p\). The chart substitution and restriction of variables are common linear maps. Write \(s_w,s'_w\) for the two lengths at \(w\). When the lengths at \(v\) increase by \(h,h'\), the chart factor \(p^{\sum_w\Delta_w(j+1)(s'_w-s_w)}\) is multiplied by \(p^{\Delta(h'-h)}\). The adelic determinants have orders minus the lengths, so their step and determinant characters contribute \(M_\pi^{h-h'}C_\pi^{b_jh-b_{j+1}h'}\). The product is \(a^h(a')^{h'}\). Thus the scalar relating the chart series of \(\chi_i^{-1}H_i\) to \(\Theta_i^F\) is a common constant times \(a^{s_i}(a')^{s'_i}\). Multiplying the Taylor relation by this constant and \(p^{\eta d_0}\) gives \(\sum_i c_i k_i^{-1}\Theta_i^F=0\).

Thus \(I_i\mapsto k_i^{-1}\Theta_i^F\) defines a rational linear map on the indicators' span. Extend it by zero on a basis of a complement in the space of all functions \(M\to\Q\). The kernel data use the underlying additive map. Its representative formula agrees with the chart formula on lattice pairs.
\end{proof}
\leanref{Formalism.\allowbreak{}MainNotions.\allowbreak{}FormalGroupSystem.\allowbreak{}Pend}{Formalism.\allowbreak{}SharpRestriction.\allowbreak{}GrowthBounds.\allowbreak{}FormalGroupSystem}{447--494}
\leanref{Formalism.\allowbreak{}MainNotions.\allowbreak{}FormalGroupSystem.\allowbreak{}adm5}{Formalism.\allowbreak{}SharpRestriction.\allowbreak{}GrowthBounds.\allowbreak{}FormalGroupSystem}{590--596}
\leanref{GlobalTheory.\allowbreak{}TwV.\allowbreak{}Xc.\allowbreak{}AwayId.\allowbreak{}scal\_\allowbreak{}away}{Development.\allowbreak{}AbelianVarietyTheory.\allowbreak{}Proofs.\allowbreak{}SatakeParameters.\allowbreak{}KernelSiteScalarsAgreeAwayFromPlace}{66--160}
\leanref{GlobalTheory.\allowbreak{}KerRel.\allowbreak{}kerTh\_\allowbreak{}relation\_\allowbreak{}of\_\allowbreak{}indicator}{Development.\allowbreak{}AbelianVarietyTheory.\allowbreak{}Proofs.\allowbreak{}SatakeParameters.\allowbreak{}KernelThetaRelationOfIndicatorRelationSplit}{76--122}
\leanref{GlobalTheory.\allowbreak{}XiGen.\allowbreak{}xiEx\_\allowbreak{}theta\_\allowbreak{}general}{Development.\allowbreak{}AbelianVarietyTheory.\allowbreak{}Proofs.\allowbreak{}SatakeParameters.\allowbreak{}XiExThetaGeneral}{86--304}
\leanref{GlobalTheory.\allowbreak{}TwV.\allowbreak{}site\_\allowbreak{}law\_\allowbreak{}xiIn}{Development.\allowbreak{}AbelianVarietyTheory.\allowbreak{}Proofs.\allowbreak{}SatakeParameters.\allowbreak{}KernelSiteLawAndFixedCoordinateStabiliserLaw}{20--49}
\leanref{Formalism.\allowbreak{}L2.\allowbreak{}G.\allowbreak{}Pend.\allowbreak{}xiOf\_\allowbreak{}spec}{Formalism.\allowbreak{}LocalToGlobal.\allowbreak{}GroupLattices.\allowbreak{}PendingLemmas.\allowbreak{}XiOfSpec}{31--34}

\begin{lemma}[Evaluation at a fixed lattice pair]\label{B3.35b}
Keep the notation of Lemma~\ref{B3.35a}, and let \(s_w,s'_w\) denote the lengths at \(w\). Put
\[
 D(L,L')=\prod_w p^{\Delta_w(j+1)(s_w-s'_w)} .
\]
Let \(\mathcal D_j(G,S)\) be the differential evaluation on a source polynomial \(G\) and a series \(S\) in the fixed coordinates, with values in target polynomials over \(K_c\). It is \(K\)-linear in each variable. For fixed \(L,L',G\) and a function \(b:M\to\Q\), define
\[
 J_{L,L'}(b;G)=D(L,L')\,a^s(a')^{s'}p^{-\eta d}
                    \mathcal D_j(G,\Xi_{\omega(L,L')}(b)).
\]
Then \(\operatorname{ker}_{j,L,L'}(G)=J_{L,L'}(I_{L,L'};G)\). For any finite family of functions \(b_i:M\to\Q\) and \(n_i\in\N\),
\[
 J_{L,L'}\Bigl(\sum_i n_i b_i;G\Bigr)=\sum_i n_i J_{L,L'}(b_i;G).
\]
\end{lemma}
\begin{proof}
The kernel is \(e_S e_T^{-1}\mathcal D_j(G,\Theta^F_{L,L'})\), where \(e_S=\prod_w p^{\Delta_w(j+1)s_w}\) and \(e_T=\prod_w p^{\Delta_w(j+1)s'_w}\). Substitute Lemma~\ref{B3.35a} and use linearity in the second variable. Both \(\Xi_\omega\) and evaluation preserve addition, so repeated addition gives the formula for multiplicities.
\end{proof}
\leanref{Formalism.\allowbreak{}MainNotions.\allowbreak{}FormalGroupSystem.\allowbreak{}Pf.\allowbreak{}kernel\_\allowbreak{}eq\_\allowbreak{}theta}{Formalism.\allowbreak{}SharpRestriction.\allowbreak{}GrowthBounds.\allowbreak{}MultisetValuation.\allowbreak{}CoupledSumOfEntries}{23--29}
\leanref{Formalism.\allowbreak{}MainNotions.\allowbreak{}FormalGroupSystem.\allowbreak{}Pf.\allowbreak{}integral\_\allowbreak{}span\_\allowbreak{}coordinates}{Formalism.\allowbreak{}SharpRestriction.\allowbreak{}GrowthBounds.\allowbreak{}MultisetValuation.\allowbreak{}IntegralSpanCoordinates}{62--67}
\leanref{Formalism.\allowbreak{}MainNotions.\allowbreak{}FormalGroupSystem.\allowbreak{}Pf.\allowbreak{}kernel\_\allowbreak{}nsmul}{Formalism.\allowbreak{}SharpRestriction.\allowbreak{}GrowthBounds.\allowbreak{}MultisetValuation.\allowbreak{}IntegralSpanCoordinates}{68--73}

\begin{lemma}[Changing the endpoint in a weighted edge sum]\label{B3.35c}
Fix a place \(v\) and an integer \(h\ge0\). At level \(0\le j\le m\) write \(A\to_h B\) when the component and the entries away from \(v\) agree and \(A_v\subseteq B_v\) is an edge of index \(h\). Put \(w(A)=|\Gamma_A|\), the order of the stabilizer in the level group of that component. For every \(K_c\)-valued function \(\phi(A,B)\) invariant under simultaneous translation in each component,
\[
 \sum_{A\in\operatorname{Rep}_j}\frac1{w(A)}\sum_{B:\,A\to_h B}\phi(A,B)
 =\sum_{B\in\operatorname{Rep}_j}\frac1{w(B)}\sum_{A:\,A\to_h B}\phi(A,B).
\]
For \(j<m\), \(f\in V_j\) and a target lattice pair \(L'\), the class data of (3.30.4) or (3.30.5) satisfy
\[
 \sum_x\frac1{w_x}\sum_y\sum_{L_{xy}\to_h B}
          \operatorname{ker}_{j,B,L'}(f(L_{xy}))
 =\sum_x\frac1{w_x}\sum_y
          \operatorname{ker}_{j,L_{xy},L'}((T_{v,h}^{(j)}f)(L_{xy})).
\]
In (3.30.4), \(x\) runs over representatives with \(w_x=w(L_x)\). The index \(y\) has one value.
\end{lemma}
\begin{proof}
Fix representatives \(A,B\) in one component, with level group \(\Gamma\). Among the elements \(\gamma\) with \(A\to_h\gamma B\), each outgoing endpoint in the orbit of \(B\) occurs \(w(B)\) times. The contribution from these two orbits is
\[
 \frac{1}{w(A)w(B)}\sum_{\gamma:\,A\to_h\gamma B}\phi(A,\gamma B).
\]
The neighbour set and stabilizer are finite. Inversion identifies this sum with the one over \(\gamma^{-1}A\to_h B\), with the same summands by invariance. Sum over the representatives and components.

Apply the first identity to each monomial coefficient of \(\operatorname{ker}_{j,B,L'}(f(A))\), extended by zero off lattice tuples. Write \(\rho(\gamma)G=\chi(\gamma)\operatorname{act}(\gamma)G\). Equivariance gives \(f(\gamma A)=\rho(\gamma)f(A)\), and \(f\in V_j\) puts \(f(A)\) in the source weight space. The covariance in Definition~\ref{B3.05a} gives \(\operatorname{ker}_{j,\gamma B,L'}(\rho(\gamma)G)=\operatorname{ker}_{j,B,L'}(G)\) on lattice tuples. The level action preserves lattice tuples, so the zero extension is invariant. By linearity in Lemma~\ref{B3.35b}, the sum at \(B\) has input \(\sum_{A\to_h B}f(A)=(T_{v,h}^{(j)}f)(B)\). Equality for each coefficient gives the polynomial identity for (3.30.4).

At \(j+1=m\) the identity is the neighbour clause of Definition~\ref{B3.05a}. In the arithmetic construction \(y\) has one value and the terminal source classes \(x\) correspond bijectively to representatives \(L_x\) with \(w_x=|\Gamma_{L_x}|\). Reindex the two sums by this bijection.
\end{proof}
\leanref{Formalism.\allowbreak{}L2.\allowbreak{}G.\allowbreak{}Pend.\allowbreak{}up\_\allowbreak{}down\_\allowbreak{}adjoint}{Formalism.\allowbreak{}LocalToGlobal.\allowbreak{}GroupLattices.\allowbreak{}PendingLemmas.\allowbreak{}OrbitSumStabilizer}{257--299}
\leanref{Formalism.\allowbreak{}MainNotions.\allowbreak{}FormalGroupSystem.\allowbreak{}MinkowskiCounting}{Formalism.\allowbreak{}SharpRestriction.\allowbreak{}GrowthBounds.\allowbreak{}FormalGroupSystem}{597--602}
\leanref{Formalism.\allowbreak{}MainNotions.\allowbreak{}FormalGroupSystem.\allowbreak{}ExpansionDecomposition}{Formalism.\allowbreak{}SharpRestriction.\allowbreak{}GrowthBounds.\allowbreak{}FormalGroupSystem}{572--586}
\leanref{Formalism.\allowbreak{}MainNotions.\allowbreak{}FormalGroupSystem.\allowbreak{}Pf.\allowbreak{}hk\_\allowbreak{}adj}{Formalism.\allowbreak{}SharpRestriction.\allowbreak{}GrowthBounds.\allowbreak{}MultisetValuation.\allowbreak{}CoupledSumOfEntries}{65--110}
\leanref{Formalism.\allowbreak{}MainNotions.\allowbreak{}FormalGroupSystem.\allowbreak{}Pf.\allowbreak{}ksa\_\allowbreak{}of\_\allowbreak{}adm}{Formalism.\allowbreak{}SharpRestriction.\allowbreak{}GrowthBounds.\allowbreak{}MultisetValuation.\allowbreak{}CoupledSumOfEntries}{140--151}
\leanref{Formalism.\allowbreak{}L2.\allowbreak{}KerInert.\allowbreak{}SecondKernelCoupled.\allowbreak{}hecke\_\allowbreak{}coupled\_\allowbreak{}entry}{Formalism.\allowbreak{}LocalToGlobal.\allowbreak{}InertKernelTheory.\allowbreak{}InertCoprimeLemmas.\allowbreak{}SecondKernelCoupled}{104--137}

\begin{lemma}[Neighbour counting and the one-step Hecke relation]\label{B3.36}
Let \(W\) be a discrete valuation ring with fraction field \(Q_W\) and uniformizer \(p\). Its residue field has cardinality \(q=p^{n'}\). Let \(N,g,i\) be nonnegative integers, \(A\subset Q_W^N\) and \(B\subset Q_W^{N+g}\) full \(W\)-lattices, and \(x:Q_W^N\to Q_W^{N+g}\) linear. Then
\[
 \begin{split}
 &\#\{B_1:pB\subseteq B_1\subseteq B,\ |B/B_1|=q^i,\ xA\subseteq B_1\}\\
 &\quad=\sum_{\ell=0}^{\min(i,g)}q^{(i-\ell)(g-\ell)}{g\brack\ell}_q\,
 \#\{A_1:A\subseteq A_1\subseteq p^{-1}A,\ |A_1/A|=q^{i-\ell},\
                                  xA_1\subseteq B\}.
 \end{split}
 \tag{3.36.1}
\]
Here \({g\brack\ell}_q\) is the Gaussian binomial coefficient. For a place \(v\) in the admissible local system, put
\[
 \Delta_v(l)=\sum_{s\ {\rm even}}w_{v,l}(s)-\sum_{s\ {\rm odd}}w_{v,l}(s),
 \qquad \mathcal E_{v,i}^{(b)}=p^{-\eta_vi(N_b-i)}T_{v,i}^{(b)} .
 \tag{3.36.2}
\]
There exist \(u_l,\mathfrak a_l\in K^\times\), \(1\le l\le m\), with
\[
 v_p(u_l)=0,\qquad v_p(\mathfrak a_l)=\Delta_v(l),
 \tag{3.36.3}
\]
where \(v_p\) is the coefficient valuation on \(K\) normalized by \(v_p(p)=1\). For every \(0\le j<m\), \(i\ge0\), and \(f\in V_j\), these parameters satisfy
\[
 \mathcal E_{v,i}^{(j+1)}K_jf
 =u_{j+1}^{\,i}K_j\left(
   \sum_{\ell=0}^i e_\ell((\mathfrak a_{j+1}\beta_{v,a})_{a=0}^{t})
                     \mathcal E_{v,i-\ell}^{(j)}f\right).
 \tag{3.36.4}
\]
The elementary symmetric polynomial \(e_\ell\) here vanishes for \(\ell>t+1\).
\end{lemma}
\leanref{Formalism.\allowbreak{}MainNotions.\allowbreak{}FormalGroupSystem.\allowbreak{}Pf.\allowbreak{}countLoc}{Formalism.\allowbreak{}SharpRestriction.\allowbreak{}GrowthBounds.\allowbreak{}MultisetValuation.\allowbreak{}SmulPsiGUp}{276--289}
\leanref{Formalism.\allowbreak{}MainNotions.\allowbreak{}FormalGroupSystem.\allowbreak{}Pf.\allowbreak{}heckeCondition\_\allowbreak{}of\_\allowbreak{}adm}{Formalism.\allowbreak{}SharpRestriction.\allowbreak{}GrowthBounds.\allowbreak{}MultisetValuation.\allowbreak{}HeckeConditionOfAdm}{75--147}

\begin{proof}
If \(xA\not\subseteq B\), both sides of (3.36.1) vanish. For \(xA\subseteq B\), let \(\bar x:A/pA\to B/pB\) be the reduction of \(x\) and put \(d=\dim_{\F_q}\ker\bar x\). The quotient \((B/pB)/\operatorname{im}\bar x\) has dimension \(d+g\). Its codimension-\(i\) subspaces correspond to the lower neighbours on the left. Multiplication by \(p\) identifies \(p^{-1}A/A\) with \(A/pA\), so the upper neighbours on the right correspond to \((i-\ell)\)-dimensional subspaces of \(\ker\bar x\). Thus it suffices to show
\[
 {d+g\brack i}_q=\sum_{\ell=0}^{\min(i,g)}
 q^{(i-\ell)(g-\ell)}{g\brack\ell}_q{d\brack i-\ell}_q.
 \tag{3.36.5}
\]
Count \(i\)-dimensional subspaces \(S\subset D\oplus G\) with \(\dim D=d\) and \(\dim G=g\). Fix \(J=S\cap G\) of dimension \(\ell\) and \(H=\operatorname{pr}_D S\) of dimension \(i-\ell\). Then \(S/J\) is the graph in \(H\oplus(G/J)\) of a unique map \(H\to G/J\). Conversely, the inverse image of each graph has the prescribed intersection and image. For each \((J,H)\) there are \(q^{(i-\ell)(g-\ell)}\) maps, which gives (3.36.5).

Use \(g=t+1\), \(N=N_j\), \(N'=N+g\), \(\eta=\eta_v\) and \(\Delta=\Delta_v(j+1)\), with the units of Lemma~\ref{B3.35a}. A target lower edge of index \(i\) decreases \(s'\) by \(i\) and \(d\) by \(Ni\). The lower-edge length rule in reverse gives an increase of \(h\) in \(s\) on a source upper edge of index \(h\), so \(d\) decreases by \(N'h\). The other lengths and \(\omega\) are fixed. Thus \(D(L,L')\) gains factors \(p^{\Delta i}\) and \(p^{\Delta h}\), respectively. With \(b,G\) fixed, the corresponding multipliers of \(J(b;G)\) in Lemma~\ref{B3.35b} are
\[
 C_A(i)=p^{\Delta i}(a')^{-i}p^{\eta Ni},\qquad
 C_C(h)=p^{\Delta h}a^hp^{\eta(N+g)h}.
 \tag{3.36.6}
\]
\leanref{Formalism.\allowbreak{}MainNotions.\allowbreak{}FormalGroupSystem.\allowbreak{}Pf.\allowbreak{}kernel\_\allowbreak{}target\_\allowbreak{}identity}{Formalism.\allowbreak{}SharpRestriction.\allowbreak{}GrowthBounds.\allowbreak{}MultisetValuation.\allowbreak{}IntegralSpanCoordinates}{80--126}
\leanref{Formalism.\allowbreak{}MainNotions.\allowbreak{}FormalGroupSystem.\allowbreak{}Pf.\allowbreak{}kernel\_\allowbreak{}source\_\allowbreak{}identity}{Formalism.\allowbreak{}SharpRestriction.\allowbreak{}GrowthBounds.\allowbreak{}MultisetValuation.\allowbreak{}IntegralSpanCoordinates}{127--174}
View (3.36.1) as an identity of functions of \(x\) and apply \(J_{L,L'}(\,\cdot\,;G)\) at the original pair. Substituting the moved indicators by Lemma~\ref{B3.35b} and using (3.36.6) gives
\[
 \sum_{B\to_i L'}\operatorname{ker}_{j,L,B}(G)
 =\sum_{\ell=0}^{\min(i,g)}
 \frac{C_A(i)}{C_C(i-\ell)}q^{(i-\ell)(g-\ell)}{g\brack\ell}_q
       \sum_{L\to_{i-\ell}A}\operatorname{ker}_{j,A,L'}(G).
\]
Here \(q=p^{n'_v}\) and \(C_C(i-\ell)\ne0\) since \(p,a\ne0\).
\leanref{Formalism.\allowbreak{}MainNotions.\allowbreak{}FormalGroupSystem.\allowbreak{}Pf.\allowbreak{}ker\_\allowbreak{}dn\_\allowbreak{}sum}{Formalism.\allowbreak{}SharpRestriction.\allowbreak{}GrowthBounds.\allowbreak{}MultisetValuation.\allowbreak{}CoupledSumOfEntries}{38--64}

At each level \(j\), the Weyl-lattice calculation in Theorem~\ref{B3.34} gives \(p^{\mathcal A_i(\kappa)}T_{v,i}^{(j)}\Lambda_j\subseteq\Lambda_j\) for \(0\le i\le N_j\), with \(\kappa\) the displayed weight at that level. For \(i>N_j\) the sum is empty. Taking \(K\)-spans gives \(T_{v,i}^{(j)}V_j\subseteq V_j\).

Take \(L=L_{xy}\) and \(G=f(L_{xy})\) in the kernel identity. Multiply by \(c_K/w_x\) and sum over \(x,y\). The left side is \(T_{v,i}^{(j+1)}K_jf\) at \(L'\) by the class expansion. On the right Lemma~\ref{B3.35c} moves the source neighbour sum onto \(f\). Since \(T_{v,i-\ell}^{(j)}f\in V_j\), the class expansion gives \(K_j(T_{v,i-\ell}^{(j)}f)\) at \(L'\) with the displayed coefficient. The equality descends by the injectivity of \(K\to K_c\).
\leanref{Formalism.\allowbreak{}MainNotions.\allowbreak{}FormalGroupSystem.\allowbreak{}Pf.\allowbreak{}hk\_\allowbreak{}point}{Formalism.\allowbreak{}SharpRestriction.\allowbreak{}GrowthBounds.\allowbreak{}MultisetValuation.\allowbreak{}HeckeConditionOfAdm}{15--74}

Write \(Q=p^\eta\). For \(0\le\ell\le g\),
\[
 {g\brack\ell}_{Q^2}=Q^{\ell(g-\ell)}
 e_\ell(Q^{g-1},Q^{g-3},\ldots,Q^{1-g}).
 \tag{3.36.7}
\]
The generating polynomial for the elementary symmetric functions at dimension \(g\) is the one at \(g-1\) with \(z\) replaced by \(Qz\), multiplied by \(1+Q^{1-g}z\). Multiplying the coefficient of \(z^\ell\) by \(Q^{\ell(g-\ell)}\) gives the recursion with coefficients \(Q^{2\ell}\) and \(1\) in (3.36.5) for \(d=g-1\) and a one-dimensional summand. Both sides start at \(1\) for \(g=\ell=0\). At \(g=t+1\) the arguments are the \(\beta_{v,a}\).

Here \(q=p^{2\eta}\). After target normalization by (3.36.2), the coefficient of the transferred sum at \(\ell\le\min(i,g)\) is
\[
 \begin{split}
 &p^{-\eta i(N+g-i)}
 \frac{C_A(i)}{C_C(i-\ell)}p^{2\eta(i-\ell)(g-\ell)}{g\brack\ell}_{p^{2\eta}}\\
 &\qquad=((aa')^{-1})^i(ap^\Delta)^\ell
 e_\ell((\beta_{v,a})_{a=0}^{t})\,
 p^{-\eta(i-\ell)(N-(i-\ell))}.
 \end{split}
 \tag{3.36.8}
\]
The quadratic exponents agree since
\(-i(N+g-i)+Ni-(N+g)(i-\ell)+2(i-\ell)(g-\ell)+\ell(g-\ell)=-(i-\ell)(N-i+\ell)\).
Set \(u_{j+1}=(aa')^{-1}\) and \(\mathfrak a_{j+1}=ap^\Delta\). This gives (3.36.3), and (3.36.8) gives (3.36.4) on lattice tuples. Since \(K_j(V_j)\subseteq V_{j+1}\) and the Hecke operators preserve these spaces, both sides vanish off lattice tuples.
\end{proof}

\begin{definition}[The exact commutation exponent]\label{B3.37}
For \(v\in\mathcal P\) and \(b\ge0\), put
\[
 P_v^{\Z}(b)=
 \sum_{s\ {\rm odd}}\left(rw_{v,b}(s)+\sum_{a=k-r+1}^{k}\lambda_{v,s,a}\right)
 +\eta_vrb(t+1).
 \tag{3.37.1}
\]
Define \(P_v(b)\) by the same odd-position sum in \(\R\), with last term \(n'_v rb(t+1)/2\). Then \(P_v(b)\) is the real image of \(P_v^{\Z}(b)\).
\end{definition}
\leanref{Formalism.\allowbreak{}MainNotions.\allowbreak{}FormalGroupSystem.\allowbreak{}PrZ}{Formalism.\allowbreak{}SharpRestriction.\allowbreak{}GrowthBounds.\allowbreak{}FormalGroupSystem}{76--80}
\leanref{Formalism.\allowbreak{}MainNotions.\allowbreak{}FormalGroupSystem.\allowbreak{}Pr}{Formalism.\allowbreak{}SharpRestriction.\allowbreak{}GrowthBounds.\allowbreak{}FormalGroupSystem}{81--85}
\leanref{Formalism.\allowbreak{}MainNotions.\allowbreak{}FormalGroupSystem.\allowbreak{}Pf.\allowbreak{}PrZ\_\allowbreak{}cast}{Formalism.\allowbreak{}SharpRestriction.\allowbreak{}GrowthBounds.\allowbreak{}MultisetValuation.\allowbreak{}LiftBoundOfIntegrality}{91--101}

\begin{proof}
Since \(n'_v\) is even, its integer half maps to \(n'_v/2\) in \(\R\).
\end{proof}

\section{Commutation along a prefix}

\begin{theorem}[The prefix Hecke commutation and its integral coefficients]\label{B3.38}
For each \(v\in\mathcal P\), there are sequences \(u_l,\mathfrak m_l\in K^\times\), \(1\le l\le m\), with
\[
 v_p(u_l)=0,\qquad v_p(\mathfrak m_l)=\Delta_v(l),
 \tag{3.38.1}
\]
such that, for \(1\le b\le m\) and every \(f\in V_0\),
\[
 Z_{v,b}K^{[b]}f=p^{P_v^{\Z}(b)}K^{[b]}(y_vf).
 \tag{3.38.2}
\]
The coefficients satisfy
\[
 c_{v,b,\ell}\in O\quad(0\le\ell\le r),\qquad v_p(c_{v,b,r})\ge E_v .
 \tag{3.38.3}
\]
For \(b=j+1\) and every \(f\in\Lambda_j\),
\[
 v_{j+1}(Z_{v,j+1}K_jf)\ge B_v+R_j(v).
 \tag{3.38.4}
\]
\end{theorem}
\leanref{Formalism.\allowbreak{}MainNotions.\allowbreak{}FormalGroupSystem.\allowbreak{}Conclusion4}{Formalism.\allowbreak{}SharpRestriction.\allowbreak{}GrowthBounds.\allowbreak{}FormalGroupSystem}{671--680}
\leanref{Formalism.\allowbreak{}MainNotions.\allowbreak{}FormalGroupSystem.\allowbreak{}Pf.\allowbreak{}Ci}{Formalism.\allowbreak{}SharpRestriction.\allowbreak{}GrowthBounds.\allowbreak{}MultisetValuation.\allowbreak{}HeckeValueEstimateOfRecordValue}{195--270}
\leanref{Formalism.\allowbreak{}MainNotions.\allowbreak{}FormalGroupSystem.\allowbreak{}Pf.\allowbreak{}Cii}{Formalism.\allowbreak{}SharpRestriction.\allowbreak{}GrowthBounds.\allowbreak{}MultisetValuation.\allowbreak{}TheoremOfAlternateForm}{538--549}
\leanref{Formalism.\allowbreak{}MainNotions.\allowbreak{}FormalGroupSystem.\allowbreak{}Pf.\allowbreak{}Ciii}{Formalism.\allowbreak{}SharpRestriction.\allowbreak{}GrowthBounds.\allowbreak{}MultisetValuation.\allowbreak{}TheoremOfAlternateForm}{608--617}

\begin{proof}
Fix \(v\) and suppress its subscript where possible. The input operator \(y_v\) preserves \(\Lambda_0\) by Definition~\ref{B3.05a}. Choose \(u_l,\mathfrak a_l\) from Lemma~\ref{B3.36}. With \(U_b\) as in (3.35.5) and \(U_0=1\), put
\[
 \mathfrak m_l=\mathfrak a_l/U_{l-1}.
 \tag{3.38.5}
\]
Each \(U_b\) is a unit, so (3.38.1) holds. Put \(\widehat{\mathcal E}_i^{(b)}=U_b^{-i}\mathcal E_i^{(b)}\) and \(\widehat{\mathcal E}^{(b)}(z)=\sum_{i\ge0}\widehat{\mathcal E}_i^{(b)}z^i\). Divide (3.36.4) by \(U_{j+1}^i\) and use \(U_{j+1}^{-i}u_{j+1}^i\mathfrak a_{j+1}^\ell U_j^{i-\ell}=\mathfrak m_{j+1}^\ell\). On \(V_j\) this gives
\[
 \widehat{\mathcal E}^{(j+1)}(z)K_j
 =K_j\widehat{\mathcal E}^{(j)}(z)
   \prod_{a=0}^t(1+\mathfrak m_{j+1}\beta_a z).
 \tag{3.38.6}
\]
Each \(K_j\) carries \(V_j\) into \(V_{j+1}\), so successive application gives
\[
 \widehat{\mathcal E}^{(b)}(z)K^{[b]}
 =K^{[b]}\widehat{\mathcal E}^{(0)}(z)
   \prod_{1\le l\le b,\ 0\le a\le t}(1+\mathfrak m_l\beta_a z).
 \tag{3.38.7}
\]
The reciprocal product is \(\sum_{\ell\ge0}(-1)^\ell h_\ell(\mathcal F_b)z^\ell\) by geometric series. Multiply (3.38.7) by it and take the coefficient of \(z^r\) to obtain
\[
 \left(\sum_{\ell=0}^r(-1)^\ell h_\ell(\mathcal F_b)
          \widehat{\mathcal E}_{r-\ell}^{(b)}\right)K^{[b]}
 =K^{[b]}\mathcal E_r^{(0)} .
 \tag{3.38.8}
\]

On a lattice tuple the sole index-zero neighbour is the tuple itself. Forms in \(V_b\) vanish off lattice tuples, so \(T_0^{(b)}=\mathcal E_0^{(b)}=1\) on \(V_b\). Hence
\(\widetilde T_i^{\,u,(b)}=p^{N'_b(i)+\eta i(N_b-i)}\widehat{\mathcal E}_i^{(b)}\)
also holds at \(i=0\). In (3.35.6), use
\(\ell(N_b-2r+\ell)+(r-\ell)(N_b-r+\ell)=r(N_b-r)\)
to obtain on \(V_b\)
\[
 Z_b=p^{N'_b(r)+\eta r(N_b-r)}
    \sum_{\ell=0}^r(-1)^\ell h_\ell(\mathcal F_b)
          \widehat{\mathcal E}_{r-\ell}^{(b)}.
 \tag{3.38.9}
\]
By (3.38.8) and \(\mathcal E_r^{(0)}=p^{-\eta r(k-r)-\mathcal A_v}y_v\), the exponent in (3.38.2) is
\[
 \begin{split}
 &N'_b(r)+\eta r(N_b-r)-\eta r(k-r)-\mathcal A_v\\
 &\qquad=N'_b(r)+\eta rb(t+1)-\mathcal A_v\\
 &\qquad=\sum_{s\ {\rm odd}}
    \left(rw_b(s)+\sum_{a=k-r+1}^{k}\lambda_{s,a}\right)
    +\eta rb(t+1)
 =P_v^{\Z}(b).
 \end{split}
 \tag{3.38.10}
\]

Put
\[
 \mu_b=\sum_{s\ {\rm even}}w_1(s)-\sum_{s\ {\rm odd}}w_b(s)-\eta t.
 \tag{3.38.11}
\]
By monotonicity, \(w_1(s)\le w_l(s)\le w_b(s)\) for \(1\le l\le b\). Take the lower endpoint for even \(s\) and the upper endpoint for odd \(s\). With \(a\le t\), this gives \(v_p(\mathfrak m_l\beta_a)\ge\mu_b\). Each degree-\(\ell\) monomial in \(h_\ell(\mathcal F_b)\) has valuation at least \(\ell\mu_b\), as does the sum. After the odd weights cancel in (3.35.6), the bound is
\[
 v_p(c_{b,\ell})\ge
 \sum_{s\ {\rm even}}\sum_{a=r-\ell+1}^{r}(\lambda_{s,a}+w_1(s))
 +\eta\ell(N_b-2r+\ell-t).
 \tag{3.38.14}
\]
Every summand of the double sum is nonnegative by dominance and \(\lambda_{s,k}+w_1(s)\ge0\), and
\(N_b-2r+\ell-t=(k-2r)+(b-1)t+b+\ell\ge0\).
This gives \(c_{b,\ell}\in O\). When \(\ell=r\), each marked even position contributes at least \(\alpha_{v,s}\), since \(\lambda_{s,1}+w_1(s)\ge\lambda_{s,1}-\lambda_{s,2}=\alpha_{v,s}\). Summing over these positions gives \(v_p(c_{b,r})\ge E_v\).

For \(b=j+1\),
\[
 Z_bK_jf=\sum_{\ell=0}^r(-1)^\ell c_{b,\ell}U_b^{-(r-\ell)}
                        \widetilde T_{r-\ell}^{(b)}K_jf .
 \tag{3.38.15}
\]
For \(\ell<r\), Theorem~\ref{B3.34} and the integral scalar give the bound \(B_v+R_j(v)\). For \(\ell=r\), the Hecke factor is the identity and \((-1)^rc_{b,r}\) adds at least \(E_v\) to the bound of Theorem~\ref{B3.33}. The fractional lattice is closed under the sum.
\end{proof}


% ---- p3c4 ----
\chapter{Induction along the chain}\label{ch:chain-induction}

Fix a prime \(p\) and a local system \(X\) as in \ref{B3.01}, with finite place set \(\mathcal P\) and chain length \(m\in\N\). The standard weight conditions at any \(v\in\mathcal P\) give \(m\ge1\). Let \(K\) be the coefficient field and \(O\) the coefficient DVR with uniformizer \(\varpi\). Write \((p)=(\varpi^e)\) in \(O\), with \(e=e_O\ge0\).

Let \(\Lambda_j\subset M_j\) be the integral lattice in the space of forms at level \(j\), and write \(K\Lambda_j\) for its \(K\)-span. For the lifts \(K_j\), put
\[
 K^{[0]}=1,\qquad K^{[b]}=K_{b-1}\cdots K_0:M_0\longrightarrow M_b.
\]
Write \(\nu_j(F)\ge h\) for \(F\in\varpi^{\lceil eh\rceil}\Lambda_j\), and \(\nu_K(a)\ge h\) for \(a\in\varpi^{\lceil eh\rceil}O\), with scalars mapped into \(K\). If \(e>0\) and \(K=\operatorname{Frac}(O)\), the scalar relation uses the valuation normalized by \(\nu_K(p)=1\). Write \(d_j=d_G(j)\) for the stabilizer loss and \([h]_+=\max(h,0)\) for the positive part of a real number.

For \(v\in\mathcal P\) and \(1\le b\le m\), write \(A_v(b)=g_{A,v}(b)\) and \(B_v(b)=g_{B,v}(b)\) as in \ref{B3.13}--\ref{B3.14}. Let \(P_v(b)\) be the real image of the integer exponent \(P_v^{\Z}(b)\) of \ref{B3.37}. For \(s\in T_v\), put \(\alpha_{v,s}=\lambda_{v,s,1}-\lambda_{v,s,2}\) as in \ref{B3.02}, and set \(E_v=\sum_{s\in T_v,\ s\ {\rm even}}\alpha_{v,s}\).

Fix kernel data \(Y\) for \(X\). The input operators \(y_v\) commute. The lifts, the \(y_v\), and the normalized output operators \(\widetilde T_{v,i,b}\) of \ref{B3.35} are \(K\)-linear. For \ref{B3.39}--\ref{B3.42}, assume \(K=\operatorname{Frac}(O)\), \(p\ne0\) in \(K\), and \(y_v\Lambda_0\subseteq\Lambda_0\) for every \(v\). Thus \(y_v\) preserves \(K\Lambda_0\). Assume the joint coefficient estimate of \ref{B3.29} for the local bounds (3.29.3) and the coefficient-to-lattice implication of \ref{B3.30}. Write \(\mathcal L_{v,j}(h)\) for the local coefficient bound \Lean{Y.\allowbreak{}LowAt} of \ref{B3.12} at \((j,v,h)\).

At each place assume given sequences \((u_{v,a})_a\) and \((\mathfrak m_{v,a})_a\) as in \ref{B3.38}. Use the coefficients \(c_{v,b,\ell}\) of (3.35.6) and put \(U_v(b)=\prod_{a=1}^b u_{v,a}\). For \(1\le a,b\le m\), we have
\[
 u_{v,a}\in O^\times,\qquad
 \nu_K(c_{v,b,\ell})\ge0\quad(0\le\ell\le r),\qquad
 \nu_K(c_{v,b,r})\ge E_v,
\]
where \(r\) is the system parameter of \ref{B3.35}. Define
\[
 Z_{v,b}=\sum_{\ell=0}^r(-1)^\ell c_{v,b,\ell}
 U_v(b)^{-(r-\ell)}\widetilde T_{v,r-\ell,b},
 \qquad \widetilde T_{v,0,b}=1.
\]
The prefix relation is
\begin{equation}\label{eq:chain-prefix-relation}
 Z_{v,b}K^{[b]}F
 =p^{P_v^{\Z}(b)}K^{[b]}y_vF
 \qquad(F\in K\Lambda_0,\ 1\le b\le m).
\end{equation}

\section{Several places and the input lattice}

For an ordered list \(L=(v_1,\ldots,v_a)\), put
\[
 Y_L=y_{v_1}\cdots y_{v_a},\qquad
 Z_{L,b}=Z_{v_1,b}\cdots Z_{v_a,b}.
\]
The rightmost operator acts first, and empty products are identities. For a subset \(C\) of the places in \(L\), let \(L|_C\) be the sublist in its original order and write \(Z_{C,b}=Z_{L|_C,b}\). For any finite set \(C\) of places, write \(Y_C\) for the product obtained from an enumeration of \(C\). The \(y_v\) commute, so this product is independent of the enumeration.

\begin{lemma}[Simultaneous operators at distinct places]\label{B3.39}
Let \(L\) be a list of distinct places. For \(1\le b\le m\) and \(F\in K\Lambda_0\),
\begin{equation}\label{eq:chain-multiple-commutation}
 Z_{L,b}K^{[b]}F=p^{\sum_{v\in L}P_v^{\Z}(b)}K^{[b]}Y_LF.
\end{equation}
Fix \(0\le j<m\) and let \(q_v\in\R\) for every \(v\in\mathcal P\). Suppose \(q_v\le B_v(j+1)\) for \(v\in L\) and \(\mathcal L_{v,j}(q_v)\) holds for \(v\notin L\). Then every \(F\in\Lambda_j\) satisfies
\begin{equation}\label{eq:chain-multiple-bound}
 \nu_{j+1}\bigl(Z_{L,j+1}K_jF\bigr)\ge\sum_{v\in\mathcal P}q_v-d_j.
\end{equation}
\end{lemma}
\leanref{Formalism.\allowbreak{}MainNotions.\allowbreak{}FormalGroupSystem.\allowbreak{}Pf.\allowbreak{}ZL\_\allowbreak{}chain}{Formalism.\allowbreak{}SharpRestriction.\allowbreak{}GrowthBounds.\allowbreak{}MultisetValuation.\allowbreak{}YEPEComm}{110--123}
\leanref{Formalism.\allowbreak{}MainNotions.\allowbreak{}FormalGroupSystem.\allowbreak{}Pf.\allowbreak{}ZL\_\allowbreak{}bound}{Formalism.\allowbreak{}SharpRestriction.\allowbreak{}GrowthBounds.\allowbreak{}MultisetValuation.\allowbreak{}AofUpdate}{172--217}

\begin{proof}
Induct on \(L\), with the empty product as the base case. For \(L=(v,L')\), pass the scalar from the induction hypothesis through \(Z_{v,b}\) and apply \eqref{eq:chain-prefix-relation} to \(Y_{L'}F\in K\Lambda_0\). The exponents of \(p\) add.

For the bound, expand the product of the \(Z_{v,j+1}\). A summand is a scalar \(a\) times the ordered product of \(\widetilde T_{v,i_v,j+1}\) on \(K_jF\), where \(0\le i_v\le r\) on \(L\) and \(i_v=0\) elsewhere. The inverse powers of \(U_v(j+1)\) are units. At each selected zero index the coefficient \(c_{v,j+1,r}\) contributes \(E_v\) to the bound for \(a\). Thus \(\nu_K(a)\ge H\), where \(H=\sum_{v\in L,\ i_v=0}E_v\).

Put \(q'_v=q_v-E_v\) if \(v\in L\) and \(i_v=0\), and \(q'_v=q_v\) otherwise. For \(v\in L\), the inequality \(q_v\le B_v(j+1)\) gives \(q'_v\le B_v(j+1)-E_v\) when \(i_v=0\) and \(q'_v\le B_v(j+1)\) when \(i_v>0\). Outside \(L\), use \(\mathcal L_{v,j}(q_v)\). Thus (3.29.3) holds for the \(q'_v\), and Lemma~\ref{B3.31} gives the operator bound \(\sum_v q'_v-d_j=\sum_v q_v-H-d_j\). Multiplication by \(a\) restores \(H\), since
\[
 \lceil e(x+y)\rceil\le\lceil ex\rceil+\lceil ey\rceil.
\]
The target lattice is closed under addition.
\end{proof}

\begin{definition}[The lattice allowing enhanced estimates]\label{B3.40}
For \(V\subseteq\mathcal P\) and \(S:\mathcal P\to\R\), let \(I(S,V)\subseteq M_0\) be the \(O\)-span of the vectors
\begin{equation}\label{eq:chain-generator}
 \left(\prod_{v\in C}y_v^m\right)
 \varpi^{\,\sum_{v\in V\setminus C}\lceil eS_v\rceil}F_0,
 \qquad C\subseteq V,\quad F_0\in\Lambda_0.
\end{equation}
\end{definition}
\leanref{Formalism.\allowbreak{}MainNotions.\allowbreak{}FormalGroupSystem.\allowbreak{}Sys.\allowbreak{}Iprod}{Formalism.\allowbreak{}SharpRestriction.\allowbreak{}GrowthBounds.\allowbreak{}FormalGroupSystem}{361--364}

Commutativity gives \(\prod_{v\in C}y_v^m=Y_C^m\) for any enumeration of \(C\).

\section{Prefix induction and the maximum at each layer}

\begin{lemma}[Prefix estimate for an arbitrary selection rule]\label{B3.41}
Let \(L\) be a list of distinct places. Let \(G^+,G^0:\mathcal P\times\N\to\R\) satisfy, for \(0\le j<m\),
\[
 G^+_v(j+1)\le B_v(j+1)\quad(v\in L),\qquad
 \mathcal L_{v,j}\bigl(G^0_v(j+1)\bigr)\quad(v\in\mathcal P).
\]
For each \(0\le j<m\), choose any subset \(C_j\) of the places in \(L\) and put
\begin{equation}\label{eq:chain-step-budget}
 \beta_j=
 \sum_{v\in C_j}G^+_v(j+1)+\sum_{v\notin C_j}G^0_v(j+1)
 -d_j-\sum_{v\in C_j}P_v(j+1).
\end{equation}
Then, for every \(0\le b\le m\) and \(F\in\Lambda_0\),
\begin{equation}\label{eq:chain-prefix-bound}
 \nu_b\bigl(K^{[b]}Y_L^bF\bigr)\ge\sum_{j=0}^{b-1}\beta_j.
\end{equation}
\end{lemma}
\leanref{Formalism.\allowbreak{}MainNotions.\allowbreak{}FormalGroupSystem.\allowbreak{}Pf.\allowbreak{}bbJ}{Formalism.\allowbreak{}SharpRestriction.\allowbreak{}GrowthBounds.\allowbreak{}MultisetValuation.\allowbreak{}YEPEComm}{127--129}
\leanref{Formalism.\allowbreak{}MainNotions.\allowbreak{}FormalGroupSystem.\allowbreak{}Pf.\allowbreak{}main\_\allowbreak{}ind}{Formalism.\allowbreak{}SharpRestriction.\allowbreak{}GrowthBounds.\allowbreak{}MultisetValuation.\allowbreak{}YEPEComm}{131--185}

\begin{proof}
At \(b=0\) the assertion is \(F\in\Lambda_0\). Suppose it holds at \(b<m\) for every integral input. Put \(C=C_b\) and let \(D\) be its complement in \(L\). Write \(H_b=\sum_{j<b}\beta_j\). Commutativity gives \(Y_L^{b+1}=Y_CY_L^bY_D\). Since \(Y_DF\in\Lambda_0\), the induction hypothesis gives \(K^{[b]}W=\varpi^{\lceil eH_b\rceil}F'\) with \(W=Y_L^bY_DF\) and \(F'\in\Lambda_b\). Each \(y_v\) preserves \(\Lambda_0\), so \(W\in\Lambda_0\).

Set \(N=\sum_{v\in C}P_v^{\Z}(b+1)\). By \ref{B3.39} and linearity,
\[
 K^{[b+1]}Y_L^{b+1}F
 =p^{-N} Z_{C,b+1}K^{[b+1]}W
 =p^{-N}\varpi^{\lceil eH_b\rceil} Z_{C,b+1}K_bF'.
\]
Use the second assertion of \ref{B3.39} with \(q_v=G^+_v(b+1)\) on \(C\) and \(q_v=G^0_v(b+1)\) elsewhere. This gives \(\sum_v q_v-d_b\) for the operator term. The scalar factors contribute \(-N\) and \(H_b\) to the bound. Since the real image of \(N\) is \(\sum_{v\in C}P_v(b+1)\), the resulting bound is \(H_b+\beta_b\).
\end{proof}

\begin{theorem}[Chain bound with external numerical estimates]\label{B3.42}
Fix \(V\subseteq\mathcal P\) and \(S:\mathcal P\to\R\). Let \(G^+,G^0,P^*:\mathcal P\times\N\to\R\). Suppose
\begin{align*}
 G^+_v(b)&\le B_v(b),&
 P_v(b)&\le P^*_v(b)
 &&(v\in V,\ 1\le b\le m),\\
 \mathcal L_{v,j}&\bigl(G^0_v(j+1)\bigr)
 &&&&(v\in\mathcal P,\ 0\le j<m),\\
 \sum_{b=1}^m[G^+_v(b)-P^*_v(b)-G^0_v(b)]_+&\le S_v
 &&&&(v\in V).
\end{align*}
Then every \(F\in I(S,V)\) satisfies
\begin{equation}\label{eq:chain-final-bound}
 \nu_m(K^{[m]}F)\ge T,
\end{equation}
where
\begin{equation}\label{eq:chain-final-budget}
 T=\sum_{b=1}^m\left(
 \sum_{v\in V}\max\{G^+_v(b)-P^*_v(b),G^0_v(b)\}
 +\sum_{v\notin V}G^0_v(b)\right)-\sum_{j=0}^{m-1}d_j.
\end{equation}
\end{theorem}
\leanref{Formalism.\allowbreak{}MainNotions.\allowbreak{}FormalGroupSystem.\allowbreak{}LowTermBound}{Formalism.\allowbreak{}SharpRestriction.\allowbreak{}GrowthBounds.\allowbreak{}FormalGroupSystem}{708--716}
\leanref{Formalism.\allowbreak{}MainNotions.\allowbreak{}FormalGroupSystem.\allowbreak{}Pf.\allowbreak{}lift\_\allowbreak{}bound\_\allowbreak{}of\_\allowbreak{}integrality}{Formalism.\allowbreak{}SharpRestriction.\allowbreak{}GrowthBounds.\allowbreak{}MultisetValuation.\allowbreak{}LiftBoundOfIntegrality}{103--174}

\begin{proof}
Take a generator in \ref{B3.40} with \(C\subseteq V\) and \(F_0\in\Lambda_0\). Enumerate \(C\) without repetition and apply \ref{B3.41} with
\[
 C_j=\{v\in C:G^0_v(j+1)\le G^+_v(j+1)-P^*_v(j+1)\}.
\]
Let \(g=\sum_{j=0}^{m-1}\beta_j\) be the resulting lower bound for \(K^{[m]}Y_C^mF_0\).
On \(C_j\), the contribution to \(\beta_j\) is at least the maximum in \eqref{eq:chain-final-budget} because \(P_v\le P^*_v\). On \(C\setminus C_j\) and outside \(V\), both contributions equal \(G^0_v(j+1)\). At each \(v\in V\setminus C\), the excess in \eqref{eq:chain-final-budget} over all layers is the sum of \([G^+_v(b)-P^*_v(b)-G^0_v(b)]_+\) over \(b\). The bound on this sum gives
\begin{equation}\label{eq:chain-real-comparison}
 T\le g+\sum_{v\in V\setminus C}S_v.
\end{equation}

Put \(n=\sum_{v\in V\setminus C}\lceil eS_v\rceil\). Since \(e\ge0\), the comparison gives \(eT\le\lceil eg\rceil+n\). The right side is integral, hence \(\lceil eT\rceil\le\lceil eg\rceil+n\). By \ref{B3.41} and linearity, the lifted generator lies in \(\varpi^{\lceil eg\rceil+n}\Lambda_m\subseteq\varpi^{\lceil eT\rceil}\Lambda_m\). Linearity extends the inclusion to the \(O\)-span.
\end{proof}

\section{The scalar budget and the local theorem}

\begin{lemma}[Budget for enhancement]\label{B3.43}
Fix \(p,n',k_1,t,m,c,r\in\N\) and \(T\subseteq\{0,\ldots,n'-1\}\). Let the functions
\[
 a,b:\{0,\ldots,n'-1\}\to\Z,\quad
 \lambda:\{0,\ldots,n'-1\}\times\{1,\ldots,k_1+1\}\to\Z,\quad
 w:\N\times\{0,\ldots,n'-1\}\to\Z
\]
satisfy the standard weight conditions of \ref{B3.01}. In particular \(p\ge5\) is prime, \(n'\ge2\) is even, \(r\ge1\), \(2r\le k_1+1\), and \(m\ge1\). With the gains \(g_A(j),g_B(j)\) of \ref{B3.13}--\ref{B3.14} and the price \(P(j)\) of \ref{B3.37}, set
\[
 S=\sum_{j=1}^m[g_B(j)-P(j)-g_A(j)]_+,\qquad
 \gamma_p=\frac{p}{p^2-1}.
\]
Then \(0\le S\le m\gamma_p\sum_{s\in T}\alpha_s\).
\end{lemma}
\leanref{Formalism.\allowbreak{}MainNotions.\allowbreak{}FormalGroupSystem.\allowbreak{}Sv}{Formalism.\allowbreak{}SharpRestriction.\allowbreak{}GrowthBounds.\allowbreak{}FormalGroupSystem}{86--87}
\leanref{Formalism.\allowbreak{}MainNotions.\allowbreak{}FormalGroupSystem.\allowbreak{}RelativePresentationLiterature}{Formalism.\allowbreak{}SharpRestriction.\allowbreak{}GrowthBounds.\allowbreak{}FormalGroupSystem}{703--706}
\leanref{Formalism.\allowbreak{}MainNotions.\allowbreak{}FormalGroupSystem.\allowbreak{}Pf.\allowbreak{}relative\_\allowbreak{}presentation\_\allowbreak{}of\_\allowbreak{}d7}{Formalism.\allowbreak{}SharpRestriction.\allowbreak{}GrowthBounds.\allowbreak{}MultisetValuation.\allowbreak{}RelativePresentationOfD7}{110--126}

\begin{proof}
The summands of \(S\) are nonnegative. Dominance gives \(\alpha_s\ge0\), and the lower bound on \(w(1,s)\) and monotonicity in \(j\) give \(\lambda_{s,i}+w(j,s)\ge0\) for \(1\le i\le k_1+1\) and \(1\le j\le m\). Thus the capacities in the gain formulas are nonnegative. With \(\delta_j\) the digit loss of \ref{B3.03}, the terms added to \(-\delta_j\) in \(g_A(j)\) and the minima subtracted in \(g_B(j)\) give
\[
 g_A(j)\ge-\delta_j,\qquad
 g_B(j)\le\gamma_p\sum_{s\in T}\alpha_s-\delta_j.
\]
The inner sum in the price formula of \ref{B3.37} has \(r\) terms, so
\[
 P(j)=\sum_{s\ {\rm odd}}\sum_{i=k_1+2-r}^{k_1+1}
 (w(j,s)+\lambda_{s,i})+\frac{n'rj(t+1)}2\ge0.
\]
Hence \([g_B(j)-P(j)-g_A(j)]_+\le\gamma_p\sum_{s\in T}\alpha_s\). Sum over \(1\le j\le m\).
\end{proof}

\begin{theorem}[Chain bound for the original local parameters]\label{B3.44}
Assume \(\mathcal L_{v,j}(A_v(j+1))\) for every \(v\in\mathcal P\) and \(0\le j<m\). Assume also that the inequality \eqref{eq:chain-final-bound} holds for every choice of external data satisfying the four displayed conditions on \(G^+,G^0,P^*,S,V\) in Theorem~\ref{B3.42}. For each \(v\), set
\[
 S_v=\sum_{b=1}^m[B_v(b)-P_v(b)-A_v(b)]_+.
\]
Then \(S_v\le m\gamma_p\sum_{s\in T_v}\alpha_{v,s}\), and for every \(V\subseteq\mathcal P\) and \(F\in I(S,V)\),
\begin{equation}\label{eq:chain-original-bound}
 \nu_m(K^{[m]}F)\ge
 \sum_{b=1}^m\left(
 \sum_{v\in V}\max\{B_v(b)-P_v(b),A_v(b)\}
 +\sum_{v\notin V}A_v(b)\right)-\sum_{j=0}^{m-1}d_j.
\end{equation}
\end{theorem}
\leanref{Formalism.\allowbreak{}MainNotions.\allowbreak{}FormalGroupSystem.\allowbreak{}Conclusion5}{Formalism.\allowbreak{}SharpRestriction.\allowbreak{}GrowthBounds.\allowbreak{}FormalGroupSystem}{697--701}
\leanref{Formalism.\allowbreak{}MainNotions.\allowbreak{}FormalGroupSystem.\allowbreak{}Pf.\allowbreak{}ramifiedLattice\_\allowbreak{}of\_\allowbreak{}chain\_\allowbreak{}estimates}{Formalism.\allowbreak{}SharpRestriction.\allowbreak{}GrowthBounds.\allowbreak{}MultisetValuation.\allowbreak{}RamifiedLatticeOfChainEstimates}{169--173}

\begin{proof}
Apply \ref{B3.43} at each place. Apply the assumed external bound with \(G^+_v=B_v\), \(G^0_v=A_v\), and \(P^*_v=P_v\). The local coefficient condition holds by hypothesis. The definition of \(S_v\) gives equality in the budget condition.
\end{proof}

\begin{theorem}[The local-system theorem]\label{B3.45}
Fix a datum \(D\) of type \Lean{Formalism.\allowbreak{}BasicNotions.\allowbreak{}D7}. For any prime \(p\), an admissible system \((X,Y)\) at \(p\) satisfies \(\Lean{Conclusion₂}(D,X,Y)\). In Definition~\ref{B3.05a}, the tests \(\Lean{adm2}\), \(\Lean{adm2b}\), and \(\Lean{adm3}\) apply to every frame for each source and target lattice pair at levels \(j,j+1\), \(0\le j<m\). In \(\Lean{adm7}\), each admissible lattice and neighbour configuration admits a frame that works for every series and real exponent.

The conjunction has eleven assertions. The first is the basic lift bound \(-d_j\) of \ref{B3.32}. The next two are the single-place and all-place product expansions of \ref{B3.09}. The fourth and fifth are the ordinary lift bound \(\sum_v A_v(j+1)-d_j\) and the full differentiated-kernel coefficient bound \(\sum_v A_v(j+1)\) of \ref{B3.32}. The sixth is the family of ordinary local coefficient bounds of \ref{B3.27}.

The seventh is the distinguished-place bound of \ref{B3.33} for the unmodified lift. The eighth asserts that, for \(1\le i\le r\), \(\widetilde T_{v,i,j+1}\) preserves \(\Lambda_{j+1}\) and satisfies the bound of \ref{B3.34} on \(K_j\Lambda_j\). For each \(v\), the ninth assertion supplies sequences \(u_{v,a},\mathfrak m_{v,a}\) such that, for \(1\le a\le m\),
\[
 \nu_K(u_{v,a})=0,\qquad
 \nu_K(\mathfrak m_{v,a})
 =\sum_{s\ {\rm even}}w_{v,a}(s)-\sum_{s\ {\rm odd}}w_{v,a}(s),
\]
where the sums run over the local positions \(0\le s<n'_v\). For these same sequences, it gives the prefix identity and coefficient estimates of \ref{B3.38} and the bound
\[
 \nu_{j+1}(Z_{v,j+1}K_jF)\ge
 B_v(j+1)+\sum_{v'\ne v}A_{v'}(j+1)-d_j
 \qquad(v\in\mathcal P,\ 0\le j<m,\ F\in\Lambda_j).
\]
The tenth assertion gives the two bounds in \ref{B3.44}. The eleventh gives the inequality in \ref{B3.42} for every choice of \(G^+,G^0,P^*,S,V\) satisfying the four displayed conditions in Theorem~\ref{B3.42}. The system assertion and the upper bound of \ref{B3.43} for all standard local weight data together form \(\Lean{Stmt}(D)\).

Fix a word \(u\) of \(D\) and a coefficient-field place \(w\) with residue field \(k(w)\). Put \(T_u=(D.\Lean{mem}\,u).A.T\). Suppose \(p=\operatorname{char}k(w)\ge5\) lies in the host's inert class and outside its level primes, and \(p\nmid T_u\). At each prime of the host field above \(p\), require that the second component of every arithmetic Frobenius under the host's Galois identification have even order. Write \(h\) for these \Lean{GIdx} conditions. For a system with \(X.G=D.\Lean{gat}(u,w,h)\), admissibility gives the same eleven conclusions at \((u,w)\).
\end{theorem}
\leanref{Formalism.\allowbreak{}MainNotions.\allowbreak{}FormalGroupSystem.\allowbreak{}admitted}{Formalism.\allowbreak{}SharpRestriction.\allowbreak{}GrowthBounds.\allowbreak{}FormalGroupSystem}{771--773}
\leanref{Formalism.\allowbreak{}MainNotions.\allowbreak{}FormalGroupSystem.\allowbreak{}Conclusion₂}{Formalism.\allowbreak{}SharpRestriction.\allowbreak{}GrowthBounds.\allowbreak{}FormalGroupSystem}{775--777}
\leanref{Formalism.\allowbreak{}MainNotions.\allowbreak{}FormalGroupSystem.\allowbreak{}Stmt}{Formalism.\allowbreak{}SharpRestriction.\allowbreak{}GrowthBounds.\allowbreak{}FormalGroupSystem}{779--780}
\leanref{Formalism.\allowbreak{}MainNotions.\allowbreak{}FormalGroupSystem.\allowbreak{}Pf.\allowbreak{}stmt\_\allowbreak{}proved}{Formalism.\allowbreak{}SharpRestriction.\allowbreak{}GrowthBounds.\allowbreak{}MultisetValuation.\allowbreak{}PrvAndKernelCommutationUpToUnitsOfRests}{107--108}
\leanref{Formalism.\allowbreak{}MainNotions.\allowbreak{}FormalGroupSystem.\allowbreak{}Pf.\allowbreak{}stmtAt\_\allowbreak{}proved}{Formalism.\allowbreak{}SharpRestriction.\allowbreak{}GrowthBounds.\allowbreak{}MultisetValuation.\allowbreak{}PrvAndKernelCommutationUpToUnitsOfRests}{110--111}

\begin{proof}
Use Definition~\ref{B3.05a}. Theorem~\ref{B3.32} gives the basic lift bound and both ordinary bounds.

If \(v\in\mathcal P\), the finite map \(W_v\to O\) from the DVR with maximal ideal \((p)\) has kernel contained in \((p)\). Lying over for the induced integral extension gives a prime of \(O\) containing \(p\) \cite{cohen-1946}. Hence \(e>0\).

The kernel field \(K_c\) contains \(\operatorname{Frac}W(\overline{\F}_p)\), where \(W\) denotes Witt vectors, so \(p\ne0\) in \(K_c\) and in \(K\). For \(e>0\), the map \(O\to K\) is injective, since a nonzero kernel would contain a power of \(\varpi\) and therefore kill \(p\). The valuation ring \(O_c\subseteq K_c\) contains the image of \(z\) or \(z^{-1}\) for every \(z\in K^\times\). Its inverse image in \(K\) is the image of \(O\), which therefore has fraction field \(K\).

Theorem~\ref{B3.09} gives the one-place and full-variable expansions. Theorem~\ref{B3.27} gives the ordinary local coefficient bounds.

The seventh and eighth assertions follow from \ref{B3.33} and \ref{B3.34}, respectively. For the ninth, use \Lean{adm3b} for the chain compatibility in \ref{B3.38} and take its sequences with their one-step bound (3.38.4). For \(e>0\), apply \ref{B3.42} with these sequences, the joint coefficient estimate of \ref{B3.29}, and \Lean{FormB}. The operators \(y_v\) preserve \(\Lambda_0\) by \Lean{adm8}.

If \(e=0\), then \(\mathcal P=\varnothing\) and \(I(S,\varnothing)=\Lambda_0\). Since \(\lceil eh\rceil=0\) for every \(h\), the basic lift bound gives \(K_j\Lambda_j\subseteq\Lambda_{j+1}\) at each step. Thus \(K^{[m]}\Lambda_0\subseteq\Lambda_m\), which is the external-parameter bound at \(e=0\).

The local bounds of \ref{B3.27} and the external-parameter bound give \ref{B3.44}. Lemma~\ref{B3.43} supplies the numerical conjunct of \(\Lean{Stmt}(D)\). Apply the system assertion at \(p=\operatorname{char}k(w)\) to the given system at \((u,w)\).
\end{proof}

% ---- p3c5 ----
\chapter{Global Frobenius-orbit parameters}\label{ch:global-frobenius-orbit-parameters}

Frobenius orbits bound the gain from below and the Hecke exchange cost from above \cite{alpoge-2026}.

Fix arithmetic data \(D\) and size data \(L\). Choose \(X\in\mathrm{Inst}(D,L)\). Write the finite abelian group \(H=D.R.A\) additively and put \((n,k,t,m)=(D.R.n,D.R.k,D.R.t,D.R.m)\). The rank after \(j\) lifts is \(N_j=k+j(t+1)\). Extend the integer weights \(a_s,b_s\) on \(TT\subseteq H\) by zero.

Put \(x=L.\mathrm{sc}.x\) and \(\varepsilon=L.\mathrm{sc}.\mathrm{eps}\), so \(x>1\) and \(0<\varepsilon<1/52\). Set
\[
 b=\exp(x^\varepsilon),\qquad a_\ast=xb,\qquad
 X'=\exp(x^{\varepsilon/3}),\qquad b'=\frac{a_\ast}{X'},\qquad
 w''=\frac{b'}{xn},\qquad X''=\frac{a_\ast}{w''}.
\]
Put \(c=4L.\mathrm{Sc}\) and \(\eta=L.\eta_{\mathrm{Sz}}\). Set \(p_0=k^2x^{26}\). Write \([y]_+=\max(y,0)\).

\section{Orbits and global weight diagrams}

\begin{definition}[Frobenius orbits and marked positions]\label{B3.46}
For a rational prime \(p\), let \(\mathcal F_p\subseteq D.R.GG\) be the set of arithmetic Frobenius images under \(D.\mathrm{host}.\mathrm{galE}\) at primes of \(D.\mathrm{host}.E\) above \(p\). Put \(h=h_p=D.R.\iota^{-1}((g_p)_2)\), where \(g_p\) is the Hilbert choice for the predicate \(g\in\mathcal F_p\).

Write \(q,d,Q_S\) for the host parameters. Among primes \(p\nmid Q_S\) with \(p\ne q,d\) and Jacobi symbol \(\bigl(\frac{-d}{p}\bigr)\ne1\), the signal class has \(D.\mathrm{host}.\mathrm{psAt}(p)=1\). Here \(\mathrm{psAt}(p)\) is the first sign of the host Artin image of \(p\bmod Q_Sqd\). The inert class has sign \(-1\) in this range.

With \(\operatorname{ord}\) denoting additive order, set
\[
 \Delta_p=\begin{cases}2h,&p\text{ of signal class},\\ h,&\text{otherwise},\end{cases}
 \qquad
 n_p=\begin{cases}2\operatorname{ord}(2h),&p\text{ of signal class},\\
                  \operatorname{ord}(h),&\text{otherwise}.\end{cases}
\]
For each \(a\in H\), the instance \(X\) supplies a map \(\beta_a:H\to H\) with
\(s\in\beta_a(s)+\langle a\rangle\) and \(\beta_a(s+a)=\beta_a(s)\) for all \(s\in H\). Put
\[
 \mathcal R_p=\beta_{\Delta_p}(H),\qquad
 T(t_0)=\{s\in TT:\beta_{\Delta_p}(s)=t_0\},\qquad N_v(t_0)=|T(t_0)|.
\]
For \(s\in H\) and \(d\in\N\), the labelled partner is
\[
 z_d(s)=\bigl(\mathbf 1_{\{p\text{ of signal class},\ d\text{ odd}\}},s+dh\bigr)
 \in\{0,1\}\times H.
\]
Write \(\operatorname{IsV}_p(s)\) if \(\beta_{\Delta_p}(s)+ih=s\) for some even index \(0\le i<n_p\).
\end{definition}
\leanref{Formalism.\allowbreak{}MainNotions.\allowbreak{}G.\allowbreak{}frobSet}{Formalism.\allowbreak{}SharpRestriction.\allowbreak{}GrowthBounds.\allowbreak{}IsResidueClassZero}{93--95}
\leanref{Formalism.\allowbreak{}MainNotions.\allowbreak{}G.\allowbreak{}hP}{Formalism.\allowbreak{}SharpRestriction.\allowbreak{}GrowthBounds.\allowbreak{}IsResidueClassZero}{97--98}
\leanref{Formalism.\allowbreak{}MainNotions.\allowbreak{}G.\allowbreak{}partner}{Formalism.\allowbreak{}SharpRestriction.\allowbreak{}GrowthBounds.\allowbreak{}IsResidueClassZero}{131--132}
\leanref{Formalism.\allowbreak{}MainNotions.\allowbreak{}G.\allowbreak{}stepP}{Formalism.\allowbreak{}SharpRestriction.\allowbreak{}GrowthBounds.\allowbreak{}IsResidueClassZero}{161--161}
\leanref{Formalism.\allowbreak{}MainNotions.\allowbreak{}G.\allowbreak{}nP}{Formalism.\allowbreak{}SharpRestriction.\allowbreak{}GrowthBounds.\allowbreak{}IsResidueClassZero}{163--164}
\leanref{Formalism.\allowbreak{}MainNotions.\allowbreak{}G.\allowbreak{}reps}{Formalism.\allowbreak{}SharpRestriction.\allowbreak{}GrowthBounds.\allowbreak{}IsResidueClassZero}{166--166}
\leanref{Formalism.\allowbreak{}MainNotions.\allowbreak{}G.\allowbreak{}topsOn}{Formalism.\allowbreak{}SharpRestriction.\allowbreak{}GrowthBounds.\allowbreak{}IsResidueClassZero}{168--169}
\leanref{Formalism.\allowbreak{}MainNotions.\allowbreak{}G.\allowbreak{}IsV}{Formalism.\allowbreak{}SharpRestriction.\allowbreak{}GrowthBounds.\allowbreak{}IsResidueClassZero}{179--180}
\leanref{Formalism.\allowbreak{}BasicNotions.\allowbreak{}primeClass}{Formalism.\allowbreak{}LocalRingCardinality.\allowbreak{}DataHost}{33--37}
\leanref{Formalism.\allowbreak{}HC.\allowbreak{}HostRaw.\allowbreak{}psAt}{Formalism.\allowbreak{}SharpRestriction.\allowbreak{}RamificationFrobenius.\allowbreak{}TotientGaloisClasses.\allowbreak{}ConjEMemHostFIff}{68--69}

\begin{definition}[Global capacities and logarithmic error]\label{B3.47}
A position \((0,s)\) is marked if \(s\in TT\). A position \((1,s)\) is thick if \(s=D.R.\operatorname{dig}(i)+D.\mathrm{host}.u_0\) for some \(i\in D.\mathrm{host}.T2\). The remaining positions are thin. For each \(j\in\N\), \(X\) supplies offsets \(\operatorname{off}_j(z)\) in the intervals
\[
 [b'/x,b'/x+\eta],\quad [b,b+\eta+1],\quad [w'',w''+\eta]
\]
at marked, thick, and thin positions, respectively. For \(z=(e,s)\), define
\[
 \operatorname{fill}_j(z)=\mathbf 1_{\{e=0,\ s\in TT\}}b_s+
 \operatorname{off}_j(z),\qquad
 \kappa_j=\frac{k+(2j-1)(t+1)}2,\qquad B_j(z)=\operatorname{fill}_j(z)-\kappa_j.
\]
For \(s\in H\) and \(d\in\N\), set
\[
 \bar\alpha_s=a_s-2c,\qquad \bar C_{s,d}=B_j(z_d(s))+c.
\]
Let \(K\) be the global constants record and write
\(K_{\delta,1}=K.\mathrm{dl1}\), \(K_{\delta,2}=K.\mathrm{dl2}\), and \(K_{\mathrm{Pr}}=K.\mathrm{Pr1}\). Put
\[
 z_p=\log_p(2a_\ast),\qquad
 \bar\delta(t_0)=K_{\delta,1}n_pN_m^2z_p+K_{\delta,2}N_v(t_0)z_p^2.
\]
\end{definition}
\leanref{Formalism.\allowbreak{}MainNotions.\allowbreak{}G.\allowbreak{}Inst}{Formalism.\allowbreak{}SharpRestriction.\allowbreak{}GrowthBounds.\allowbreak{}IsResidueClassZero}{141--150}
\leanref{Formalism.\allowbreak{}MainNotions.\allowbreak{}G.\allowbreak{}fill}{Formalism.\allowbreak{}SharpRestriction.\allowbreak{}GrowthBounds.\allowbreak{}IsResidueClassZero}{154--155}
\leanref{Formalism.\allowbreak{}MainNotions.\allowbreak{}G.\allowbreak{}fillG}{Formalism.\allowbreak{}SharpRestriction.\allowbreak{}GrowthBounds.\allowbreak{}IsResidueClassZero}{159--159}
\leanref{Formalism.\allowbreak{}MainNotions.\allowbreak{}G.\allowbreak{}al}{Formalism.\allowbreak{}SharpRestriction.\allowbreak{}GrowthBounds.\allowbreak{}IsResidueClassZero}{230--230}
\leanref{Formalism.\allowbreak{}MainNotions.\allowbreak{}G.\allowbreak{}capU}{Formalism.\allowbreak{}SharpRestriction.\allowbreak{}GrowthBounds.\allowbreak{}IsResidueClassZero}{232--233}
\leanref{Formalism.\allowbreak{}MainNotions.\allowbreak{}G.\allowbreak{}dl}{Formalism.\allowbreak{}SharpRestriction.\allowbreak{}GrowthBounds.\allowbreak{}IsResidueClassZero}{241--243}

\begin{definition}[Ordinary global gain]\label{B3.48}
For \(s,s'\in H\) and \(d\in\N\), let \(\operatorname{IsDist}_p(s,s',d)\) mean that \(d>0\), \(z_d(s)=(0,s')\), and \(z_e(s)\ne(0,s')\) for every \(0<e<d\). Let \(\operatorname{Excl}_p(s,d)\) mean that \(d\) is odd and \(\operatorname{IsDist}_p(s,s',d)\) holds for some \(s'\in TT\setminus\{s\}\). The ordinary gain is
\[
 g_0(t_0,j)=
 \sum_{s\in T(t_0)}
 \sum_{\substack{d\ge1,\ d\ {\rm odd}\\\neg\operatorname{Excl}_p(s,d)}}
 \left[\frac{\bar\alpha_s}{p^d}-\bar C_{s,d}\right]_+-\bar\delta(t_0).
\]
\end{definition}
\leanref{Formalism.\allowbreak{}MainNotions.\allowbreak{}G.\allowbreak{}IsDist}{Formalism.\allowbreak{}SharpRestriction.\allowbreak{}GrowthBounds.\allowbreak{}IsResidueClassZero}{235--236}
\leanref{Formalism.\allowbreak{}MainNotions.\allowbreak{}G.\allowbreak{}Excl}{Formalism.\allowbreak{}SharpRestriction.\allowbreak{}GrowthBounds.\allowbreak{}IsResidueClassZero}{238--239}
\leanref{Formalism.\allowbreak{}MainNotions.\allowbreak{}G.\allowbreak{}g0}{Formalism.\allowbreak{}SharpRestriction.\allowbreak{}GrowthBounds.\allowbreak{}IsResidueClassZero}{247--250}

\begin{definition}[Enhanced global gain]\label{B3.49}
Put \(T^V(t_0)=\{s\in T(t_0):\operatorname{IsV}_p(s)\}\) and
\(\gamma_p=p/(p^2-1)=\sum_{d\ge1,\ d\ {\rm odd}}p^{-d}\).
The enhanced gain is
\[
 g_U(t_0,j)=\gamma_p\sum_{s\in T(t_0)}\bar\alpha_s
 -\sum_{s\in T(t_0)\setminus T^V(t_0)}
   \sum_{\substack{d\ge1\\d\ {\rm odd}}}
   \min\left(\frac{\bar\alpha_s}{p^d},\bar C_{s,d}\right)
 -\bar\delta(t_0).
\]
\end{definition}
\leanref{Formalism.\allowbreak{}MainNotions.\allowbreak{}G.\allowbreak{}gU}{Formalism.\allowbreak{}SharpRestriction.\allowbreak{}GrowthBounds.\allowbreak{}IsResidueClassZero}{254--258}

The gains are defined for \(t_0\in H\) and \(j\in\N\), with nonsummable sums set to zero. For nonnegative capacities use
\[
 [u/p^d-C]_+\le [u]_+/p^d,\qquad
 |\min(u/p^d,C)|\le |u|/p^d.
\]
Hence the sums converge and \(u-\min(u,C)=[u-C]_+\) gives
\[
 g_U(t_0,j)+\bar\delta(t_0)
 =\gamma_p\sum_{s\in T^V(t_0)}\bar\alpha_s
 +\sum_{s\in T(t_0)\setminus T^V(t_0)}
   \sum_{\substack{d\ge1\\d\ {\rm odd}}}
   \left[\frac{\bar\alpha_s}{p^d}-\bar C_{s,d}\right]_+ .
 \tag{1}
\]

\section{Prices and selected orbits}

Put
\[
 h_G=2+8\varepsilon,\qquad
 K_G=\left\lceil\frac{8(h_G+2)}{\varepsilon}\right\rceil,\qquad
 r_G=K_G+1,\qquad M=w''+\eta+c+N_m.
\]

\begin{proposition}[The global price and its two bounds]\label{B3.50}
Define
\[
 \mathrm{Pr}(t_0,j)=
 \begin{cases}
 \displaystyle\sum_{s\in T(t_0)\setminus T^V(t_0)}(B_j(0,s)-b_s)
           +K_{\mathrm{Pr}}n_pM,&p\ge p_0,\\[4pt]
 r_G\bigl(N_v(t_0)(2b'+c)+n_pM\bigr),&p<p_0 .
 \end{cases} \tag{2}
\]
Assume \(\Lean{SizesAt}(L)\) and the tying conditions \(\Lean{GatTies}\) and \(\Lean{LatTied}\). Fix a word \(u\) and a coefficient-field place \(w\), with \(p=\operatorname{char}k(w)\). Assume \(\Lean{GIdx}\) at \((u,w)\). Thus \(p\ge5\) is an inert-class prime outside the level primes and \(p\nmid T_u\). Every arithmetic Frobenius above \(p\) has second component of even order under \(D.\mathrm{host}.\mathrm{galE}\).

For \(t_0\in\mathcal R_p\), let \(v=v(t_0)\) be its local position and \(P_v(j)\) the exact exchange exponent of \ref{B3.37}. If \(p\ge p_0\) and \(K_{\mathrm{Pr}}\ge1\), the local Hecke rank is \(1\) and \(P_v(j)\le\mathrm{Pr}(t_0,j)\) for \(1\le j\le m\). If \(p<p_0\), \(r_G<k\), and \(x\ge2\), the rank is \(r_G\) and the same inequality holds.
\end{proposition}
\leanref{Formalism.\allowbreak{}MainNotions.\allowbreak{}G.\allowbreak{}Pr}{Formalism.\allowbreak{}SharpRestriction.\allowbreak{}GrowthBounds.\allowbreak{}IsResidueClassZero}{262--267}
\leanref{Formalism.\allowbreak{}MainNotions.\allowbreak{}FormalGroupSystem.\allowbreak{}Pf.\allowbreak{}price\_\allowbreak{}bound\_\allowbreak{}large\_\allowbreak{}primes}{Formalism.\allowbreak{}SharpRestriction.\allowbreak{}GrowthBounds.\allowbreak{}MultisetValuation.\allowbreak{}PriceBoundLargePrimes}{117--180}
\leanref{Formalism.\allowbreak{}MainNotions.\allowbreak{}FormalGroupSystem.\allowbreak{}Pf.\allowbreak{}price\_\allowbreak{}bound\_\allowbreak{}small\_\allowbreak{}primes}{Formalism.\allowbreak{}SharpRestriction.\allowbreak{}GrowthBounds.\allowbreak{}MultisetValuation.\allowbreak{}PriceLinearCombinationLe}{89--142}

\begin{proof}
For \(0\le i<n_p\), put \(s=t_0+ih\). This identifies the marks with \(T(t_0)\) and the odd marks with \(T(t_0)\setminus T^V(t_0)\). The ties give \(w_j(i)=B_j(0,s)\) and \(c_{\rm loc}\le c\). They fix the rank \(r\) in each branch. In one-based row notation,
\[
 P_v(j)=\sum_{\substack{0\le i<n_p\\ i\ {\rm odd}}}
 \left(rw_j(i)+\sum_{\ell=k-r+1}^{k}\lambda_{i,\ell}\right)
 +\frac{n_prj(t+1)}2. \tag{3}
\]
For \(r=1\), \(\lambda_{i,k}=-b_s\) at marks, so their summands are \(B_j(0,s)-b_s\). Elsewhere \(w_j(i)\le w''+\eta\) and \(\lambda_{i,k}\le c\). Thus \(j(t+1)/2\le N_m\) and \(K_{\mathrm{Pr}}\ge1\) give the first branch of (2).

For \(r=r_G<k\), the last \(r\) rows omit the first, so their weights are at most \(c_{\rm loc}\). This includes the weight \(-b_s\le0\) at a mark. At marks \(\Lean{SizesAt}\) gives \(b_s\le5b'/4\), so
\[
 w_j(i)\le b_s+b'/x+\eta\le7b'/4+\eta\le2b'+\eta.
\]
Each summand is at most \(r(w''+\eta+c)+\mathbf 1_{\{i\ {\rm marked}\}}\,2rb'\). Summing and using \(n_prj(t+1)/2\le rn_pN_m\) gives \(r(2N_v(t_0)b'+n_pM)\le\mathrm{Pr}(t_0,j)\).
\end{proof}

Put \(n'=\operatorname{ord}(h)\) and \(S_{t_0}=\{0\le i<n':t_0+ih\in TT\}\). Use the constants
\[
 \begin{gathered}
 \delta_1=x^{-\varepsilon/8},\quad u_1=\lceil1/\delta_1\rceil,\quad
 \ell=(6r_G+1)u_1,\quad M_0=\lceil(h_G+3)(\log x)10^{6r_G+1}\rceil,\\
 G'=(M_0+5)\ell,\qquad
 N_{\mathrm{win}}=\left\lceil\max\{6eG'\delta_1,(h_G+3)\log_2x\}\right\rceil+1 .
 \end{gathered} \tag{4}
\]
Let \(L_1(p)=\lceil\log_p((kx^{26})^4)\rceil\) and \(L_2=2\sup_{p\ \text{of inert class}}L_1(p)\).
Call \(t_0\) good when \(S_{t_0}\subseteq\Z/n'\Z\) satisfies the window and tile conditions of \ref{B2.54}, with \(K_\varepsilon=K_G\), \(R_\varepsilon=r_G\), \(u=u_1\), and \(N_1=N_{\mathrm{win}}\).
\leanref{Formalism.\allowbreak{}MainNotions.\allowbreak{}G.\allowbreak{}IsGoodOrbit}{Formalism.\allowbreak{}SharpRestriction.\allowbreak{}GrowthBounds.\allowbreak{}UnramifiedQuadraticSetting}{250--252}
\leanref{Formalism.\allowbreak{}G0.\allowbreak{}GoodOrbitIntervalTileCondition}{Formalism.\allowbreak{}SubspaceCounting.\allowbreak{}GoodOrbitIntervalTileCondition}{24--25}

\begin{definition}[Selection and counted gain]\label{B3.51}
Let \(\mathcal G_p\subseteq\mathcal R_p\) be the good representatives and set \(\Sigma_p=m\sum_{t_0\in\mathcal G_p}\kappa_v(t_0)\), where
\[
 \kappa_v(t_0)=
 \begin{cases}
 \min\{1,K.\mathrm{kII}\,|S_{t_0}|r_G^2 kx^{-26}\},&p<p_0,\\
 X.\mathrm{kapI}\bigl(p,n',S_{t_0},(a_{t_0+ih})_i,(b_{t_0+ih})_i\bigr),&p\ge p_0 .
 \end{cases}
\]
At an inert-class prime, \(N_v(t_0)=|S_{t_0}|\) for \(t_0\in\mathcal R_p\).

A candidate is a prime of inert class with \(p<X'\) and \(n'\ge x^{\varepsilon/4}\), subject also to \(\neg\operatorname{InBand}(p)\) when \(p\ge p_0\). The predicate \(\operatorname{InBand}\) means \(p_0\le p<X'\) and
\[
 X'(1-\rho_B)\le p^d\le X'(1+\rho_B)\quad\text{for some integer }d\ge1,
\]
where \(N_\ast=2n\delta_1\), \(\mathcal L=K.\mathrm{La}\,N_{\mathrm{win}}^2\log a_\ast\), and
\[
 \rho_B=\frac{4N_\ast x^\varepsilon\mathcal L}{k}
         +\frac{K.\rho\,K.\mathrm{CN}(\log x)\mathcal L}{k}.
\]
Set
\[
 V_p=
 \begin{cases}
 \mathcal G_p,&p\text{ a candidate and }\Sigma_p\le x^{-\varepsilon},\\
 \varnothing,&\text{otherwise}.
 \end{cases}\tag{5}
\]
For \(1\le j\le m\), the counted gain is
\[
 \operatorname{gainP}(p,j)=
 \begin{cases}
 0,&(p\text{ of inert class and }h=0)\text{ or }p>X'',\\[2pt]
 \displaystyle\sum_{t_0\in V_p}
 \max\{g_U(t_0,j)-\mathrm{Pr}(t_0,j),g_0(t_0,j),0\}
 +\sum_{t_0\in\mathcal R_p\setminus V_p}[g_0(t_0,j)]_+,
 &\text{otherwise}.
 \end{cases}\tag{6}
\]
\end{definition}
\leanref{Formalism.\allowbreak{}MainNotions.\allowbreak{}G.\allowbreak{}kapv}{Formalism.\allowbreak{}SharpRestriction.\allowbreak{}GrowthBounds.\allowbreak{}IsResidueClassZero}{276--277}
\leanref{Formalism.\allowbreak{}MainNotions.\allowbreak{}G.\allowbreak{}Cand}{Formalism.\allowbreak{}SharpRestriction.\allowbreak{}GrowthBounds.\allowbreak{}IsResidueClassZero}{279--281}
\leanref{Formalism.\allowbreak{}MainNotions.\allowbreak{}G.\allowbreak{}VV}{Formalism.\allowbreak{}SharpRestriction.\allowbreak{}GrowthBounds.\allowbreak{}IsResidueClassZero}{292--293}
\leanref{Formalism.\allowbreak{}MainNotions.\allowbreak{}G.\allowbreak{}gainP}{Formalism.\allowbreak{}SharpRestriction.\allowbreak{}GrowthBounds.\allowbreak{}IsResidueClassZero}{355--358}

\section{Identifying the global data}

\begin{definition}[Filling compatibility]\label{B3.52}
Let \(w\) be a coefficient-field place with \(p=\operatorname{char}k(w)\). The condition \(\Lean{ParametersFill}(X,K,w)\) imposes
\[
 \operatorname{fill}_j(z)=w_j(\tau)
 \quad(1\le j\le m,\ D.eTau(\tau)=\operatorname{plGG}(z)) \tag{7}
\]
for all embedding indices \(\tau\) and labelled positions \(z\). Here \(w_j(\tau)=(D.\mathrm{hc}(\tau)).e_{j-1}\) and \(\operatorname{plGG}(e,s)=D.R.\tau_{e+1}(s)\) for \(e\in\{0,1\}\). Also require
\[
 L.wPart(j,s,p,d)=\operatorname{fill}_j(z_d(s))
 \qquad(j\in\N,\ s\in TT,\ d\in\N). \tag{8}
\]
Under \(TT\hookrightarrow H\), the image of \(L.\mathrm{TTw}(w)\) must equal
\[
 \operatorname{TTP}(p)=
 \begin{cases}
 TT,&p\text{ of signal class or }p\ge X',\\
 \{s\in TT:\beta_{\Delta_p}(s)\in V_p\},&\text{otherwise}.
 \end{cases}\tag{9}
\]
For \(w\in L.\mathrm{supp}\), \(p\) of inert or signal class, and \(1\le j\le m\), impose
\[
 L.\mathrm{gain}(j,w)=\operatorname{gainP}(p,j). \tag{10}
\]
The remaining equalities are
\[
 L.\tau(w)=\tau_P(p),\qquad L.\mathrm{io}(w)=\mathrm{io}_P(p),\qquad
 L.\mathrm{LaG}=\mathcal L. \tag{11}
\]
Put
\[
 S_v(t_0)=\left\lceil\sum_{j=1}^m
   \bigl[g_U(t_0,j)-\mathrm{Pr}(t_0,j)-[g_0(t_0,j)]_+\bigr]_+
 \right\rceil .
\]
Both budgets are zero for \(V_p=\varnothing\). For nonempty \(V_p\),
\[
 \tau_P(p)=X.\mathrm{ep}(p)+
 m\sum_{t_0\in V_p}\bigl([\lceil t_v(t_0)\rceil]_++D_p+1\bigr),
 \qquad \mathrm{io}_P(p)=2\Sigma_p\max_{t_0\in V_p}S_v(t_0),
\]
where \((t_v(t_0),D_p)=(t_{II},D_{II})\) for \(p<p_0\) and \((t_I,D_I)\) for \(p\ge p_0\). Use the cycle of length \(n'=n_p\) with marks \(S_{t_0}\) and weights \(\mathbf a=(a_{t_0+ih})_i,\ \mathbf b=(b_{t_0+ih})_i\). The formulas and the family choice of \(X.\mathrm{kapI},X.\mathrm{ep}\) are given in \ref{B3.52a}.
\end{definition}
\leanref{Formalism.\allowbreak{}MainNotions.\allowbreak{}G.\allowbreak{}ParametersFill}{Formalism.\allowbreak{}SharpRestriction.\allowbreak{}GrowthBounds.\allowbreak{}IsResidueClassZero}{401--410}
\leanref{Formalism.\allowbreak{}MainNotions.\allowbreak{}G.\allowbreak{}TTP}{Formalism.\allowbreak{}SharpRestriction.\allowbreak{}GrowthBounds.\allowbreak{}IsResidueClassZero}{362--364}
\leanref{Formalism.\allowbreak{}MainNotions.\allowbreak{}G.\allowbreak{}Sv}{Formalism.\allowbreak{}SharpRestriction.\allowbreak{}GrowthBounds.\allowbreak{}IsResidueClassZero}{336--337}
\leanref{Formalism.\allowbreak{}MainNotions.\allowbreak{}G.\allowbreak{}tauP}{Formalism.\allowbreak{}SharpRestriction.\allowbreak{}GrowthBounds.\allowbreak{}IsResidueClassZero}{341--344}
\leanref{Formalism.\allowbreak{}MainNotions.\allowbreak{}G.\allowbreak{}ioP}{Formalism.\allowbreak{}SharpRestriction.\allowbreak{}GrowthBounds.\allowbreak{}IsResidueClassZero}{348--351}
\leanref{Formalism.\allowbreak{}MainNotions.\allowbreak{}G.\allowbreak{}tIIf}{Formalism.\allowbreak{}SharpRestriction.\allowbreak{}GrowthBounds.\allowbreak{}IsResidueClassZero}{197--200}
\leanref{Formalism.\allowbreak{}MainNotions.\allowbreak{}G.\allowbreak{}DIIf}{Formalism.\allowbreak{}SharpRestriction.\allowbreak{}GrowthBounds.\allowbreak{}IsResidueClassZero}{188--191}
\leanref{Formalism.\allowbreak{}MainNotions.\allowbreak{}G.\allowbreak{}tIf}{Formalism.\allowbreak{}SharpRestriction.\allowbreak{}GrowthBounds.\allowbreak{}IsResidueClassZero}{217--225}
\leanref{Formalism.\allowbreak{}MainNotions.\allowbreak{}G.\allowbreak{}DIf}{Formalism.\allowbreak{}SharpRestriction.\allowbreak{}GrowthBounds.\allowbreak{}IsResidueClassZero}{193--195}

\begin{lemma}[The two local presentations]\label{B3.52a}
Put \(A=a_\ast\), \(r=r_G\), \(z=\log_p A\), \(q_p=p^{G'}\), \(c_0=2b'+2c\), and \(U=40n'\delta_1\). Write \(K_1,K_2\) for the fields \Lean{kappaThresholdOne}, \Lean{kappaThresholdTwo} of \(K\). The other subscripts on \(K\) give its field names. For \(N\) marks on a cycle of length \(n'\), the denominator bounds and costs of \ref{B3.52} are
\[
\begin{aligned}
 D_I&=\lceil K_{\rm DI}Ux^{26}(N_{\rm win}^5+k^3)z\rceil_+,\\
 D_{II}&=\lceil K_{\rm DII}U(N_{\rm win}^3+r^5+k^2)z+2rn'c+rUc_0\rceil_+,\\
 t_{II}&=K_{\rm tII}rN(\mathcal LA/(pk)+Ap^{-L_1(p)}+c_0+kq_p)+2rn'c,\\
 c_6&=8(D_I+1)+K_2Ukz+4Uc+2n'c,\\
 t_I&=J+K_1N N_{\rm win}(A/p^2+\mathcal LA/(pk))
      +K_2N(q_pc_6+kq_p+kz+c)+2n'c .
\end{aligned}
\]
Here \(\lceil y\rceil_+=\max(0,\lceil y\rceil)\). Each odd mark \(y\) contributes \(\min(b_y,a_{y-d_y}/p^{d_y})\) to \(J\), where \(d_y\) is the least positive distance back to an even mark within one period. With no even mark the contribution is zero. Parity uses cyclic representatives \(0,\ldots,n'-1\).

Assume the family instance \(X\) satisfies \(\Lean{EpDef}(X,C_e)\) and \(\Lean{KapIDef}(X,K)\). Then \(X.\mathrm{ep}(p)=\lceil C_e f k^2z\rceil_+\), where \(f=D.R.f\). Fix \(p,n',S,\mathbf a,\mathbf b\), with mark set \(S\) and integer weights \(\mathbf a,\mathbf b\) on the cycle of length \(n'\). The share \(X.\mathrm{kapI}(p,n',S,\mathbf a,\mathbf b)\) is the least \(\kappa\ge0\) for which every local model with these parameters passes the following rank-one test. Set \(r'=1\). Use the budgets \(t_\ast=\lceil t_I\rceil_+\) and \(D_0=D_I\).

Range over complete DVRs \(W\) with uniformizer \(p\ge5\) prime, residue order \(p^{n'}\), \(n'\ge2\) even, and a Frobenius lift \(\sigma\) with \(\sigma^{n'}=1\). Take every finite extension \(W\to O\) of complete DVRs with compatible inclusion of fraction fields, uniformizer \(\omega\) satisfying \(\omega^2=p\), residue order \(p^{2n'}\), and square roots of all rational integers in \(O\). Use every dominant base weight with \(k\) entries at each position subject to the following conditions. The extremes at marks are \(a_i,-b_i\), and all unspecified entries are bounded in absolute value by \(\max(0,\lfloor c\rfloor)\). For odd \(n'\) the class of local models is empty and the share is zero.

For a full \(W\)-lattice \(\Lambda\subset\operatorname{Frac}(W)^k\), let \(H_\Lambda\) be the integral dual Weyl tensor lattice on its Frobenius translates over \(O\), with the negative reversed weight at odd positions. Its \(\operatorname{Frac}(O)\)-span is the tensor representation \(V\). For a rank-\(r'\) edge \(e:\Lambda'\to\Lambda\) as in \ref{B3.35}, form \(H_e\) on the intermediate lattice \(\Lambda'_O+\omega\Lambda_O\) at each translate. Put
\[
 R_e=\omega^{-A_{r'}}H_e,\qquad
 R'_e=\omega^{-A_{r'}}H_e^\vee,
\]
where \(A_{r'}=\mathcal A_{r'}\) of (3.34.4) at level zero and \({}^\vee\) denotes the \(O\)-linear dual lattice. Each target \(\Lambda\) must admit a finite list \((e_i,F_i)_{i=1}^M\) of incoming edges \(e_i:\Lambda_i\to\Lambda\) and vectors \(F_i\) satisfying the conditions below. An edge may occur more than once.
\[
 F_i\in\omega^{-t_\ast}H_{\Lambda_i}\cap\bigcap_{e:\Lambda_i\to *}R_e,
 \qquad
 F_i\in\omega^{8(D_0+1)}R_e\quad(e:\Lambda_i\to *,\ e\ne e_i).
\]
For some \(h_0\in\N\) with \(h_0\ge(1-\kappa)\dim V\), there must be a matrix \(B\in O^{h_0\times M}\) and an \(O\)-linear map
\[
 U_0:O^{h_0}\longrightarrow
 \omega^{-t_\ast}H_\Lambda^\vee\cap\bigcap_{e:*\to\Lambda}R'_e
\]
such that
\[
 \sum_{i=1}^{M} B_{ji}\,(U_0\xi)(p^{A_{r'}}F_i)
   =p^{D_0}\xi_j
   \qquad(\xi\in O^{h_0},\ 1\le j\le h_0).
\]
\end{lemma}
\begin{proof}
Substitution of (4) in the local budgets gives the displayed formulas. The family defines the share as the infimum \(s\) of the passing fractions. The empty list with \(h_0=0\) passes at one, so \(0\le s\le1\). Fix a local model and choose a weight for it. For a fixed target, set \(q=(1-s)\dim V\). If \(q>0\), choose a passing fraction \(\kappa\) close enough to \(s\) that \((1-\kappa)\dim V>\lceil q\rceil-1\). Its witness has integer rank at least \(\lceil q\rceil\) and passes at \(s\) for this target. If \(q=0\), take the empty list.
\end{proof}
\leanref{Formalism.\allowbreak{}MainNotions.\allowbreak{}G.\allowbreak{}EpDef}{Formalism.\allowbreak{}SharpRestriction.\allowbreak{}GrowthBounds.\allowbreak{}DefsInst}{18--19}
\leanref{Formalism.\allowbreak{}MainNotions.\allowbreak{}G.\allowbreak{}kapISet}{Formalism.\allowbreak{}SharpRestriction.\allowbreak{}GrowthBounds.\allowbreak{}DefsInst}{38--46}
\leanref{Formalism.\allowbreak{}MainNotions.\allowbreak{}G.\allowbreak{}sInf\_\allowbreak{}mem\_\allowbreak{}kapISet}{Formalism.\allowbreak{}SharpRestriction.\allowbreak{}GrowthBounds.\allowbreak{}DefsInst}{68--115}
\leanref{Formalism.\allowbreak{}MainNotions.\allowbreak{}G.\allowbreak{}KapIDef}{Formalism.\allowbreak{}SharpRestriction.\allowbreak{}GrowthBounds.\allowbreak{}DefsInst}{54--55}
\leanref{Formalism.\allowbreak{}MainNotions.\allowbreak{}G.\allowbreak{}exists\_\allowbreak{}inst\_\allowbreak{}defs}{Formalism.\allowbreak{}SharpRestriction.\allowbreak{}GrowthBounds.\allowbreak{}DefsInst}{120--124}

\section{The numerical comparison}

For tied data \(Y\) at a \(\Lean{GIdx}\) place and each \(t_0\in\mathcal R_p\), use the orbit coordinate \(\iota(i)=t_0+ih\) of \ref{B3.50}. The ties make \(t_0\mapsto v(t_0)\) injective on \(\mathcal R_p\).

Write \(\alpha_i,C_i,\delta_v\) for the local column count, capacity, and digit error of \ref{B3.02}--\ref{B3.03}, and \(g_{A,v}(j),g_{B,v}(j)\) for the gains of \ref{B3.13}--\ref{B3.14}.

\begin{theorem}[Comparison when \(z_p\ge1\)]\label{B3.53}
Assume \(K_{\delta,1}\ge5\), \(K_{\delta,2}\ge4\), \(K_{\mathrm{Pr}}\ge1\), \(r_G<k\), \(x\ge2\), and \(\Lean{SizesAt}(L)\). For every choice of tied data, every \((u,w)\) satisfying \(\Lean{GIdx}\), every \(t_0\in\mathcal R_p\), and every \(1\le j\le m\), the condition \(z_p\ge1\) implies
\[
 g_U(t_0,j)\le g_{B,v}(j),\qquad
 g_0(t_0,j)\le g_{A,v}(j),\qquad
 P_v(j)\le\mathrm{Pr}(t_0,j).
 \tag{12}
\]
\end{theorem}
\leanref{Formalism.\allowbreak{}MainNotions.\allowbreak{}FormalGroupSystem.\allowbreak{}SmallResidueCharComparison}{Formalism.\allowbreak{}SharpRestriction.\allowbreak{}GrowthBounds.\allowbreak{}FormalGroupSystem}{741--754}
\leanref{Formalism.\allowbreak{}MainNotions.\allowbreak{}FormalGroupSystem.\allowbreak{}Pf.\allowbreak{}numeric\_\allowbreak{}comparison\_\allowbreak{}of\_\allowbreak{}constants}{Formalism.\allowbreak{}SharpRestriction.\allowbreak{}GrowthBounds.\allowbreak{}ProofLW}{207--219}

\begin{proof}
At a marked position, \(\alpha_i=a_{\iota(i)}-\lambda_{i,2}\) and \(C_i=w_j(i)+\lambda_{i,2}\). Elsewhere, \(\alpha_i=0\) and \(C_i=w_j(i)+\lambda_{i,1}\). For this row, \(|\lambda_{i,\ell}|\le c_{\rm loc}\le c\) and \(w_j(i)\ge c_{\rm loc}\), so
\[
 \bar\alpha_{\iota(i)}\le\alpha_i\quad(i\text{ marked}),\qquad
 0\le C_i\le B_j(0,\iota(i))+c. \tag{13}
\]
Even local indices correspond to \(\operatorname{IsV}_p\). If \(D_i\) is the local exclusion set of \ref{B3.13}, then
\[
 d\in D_i,\quad d\text{ odd}\quad\Longrightarrow\quad
 \operatorname{Excl}_p(\iota(i),d). \tag{14}
\]
The set \(D_i\) consists of the least positive representatives of \(i'-i\bmod n_p\) as \(i'\) runs over marks of parity opposite to \(i\). On an inert orbit \(\iota\) is injective on a period and the partner label is \(0\).

Put \(A=2a_\ast\) and \(z=\log_p A\). The size conditions give \(c_{\rm loc}\le4w''/5\) and \(\eta\le w''/2\). They also give \(N_m\le a_\ast\). Since \(w'',b'/x\le b'\) and \(b_s\le5b'/4\), the offsets give \(w_j(i)\le11b'/4\). At marks \(a_{\iota(i)}\le6a_\ast/5\), and after removing the single-box columns each row has at most \(w_j(i)+c_{\rm loc}\) boxes. Thus the remaining box count \(L_i\) and the other size parameters satisfy
\[
 \alpha_i\le A,\qquad L_i\le N_j(w_j(i)+c_{\rm loc})\le a_\ast(4a_\ast)=A^2,
 \qquad N_{j-1},N_j\le N_m. \tag{15}
\]
For the base-\(p\) digit length this gives \(\ell_p(\alpha_i)\le1+z\le2z\) and \(\ell_p(L_i)\le1+2z\le3z\). The error formula of \ref{B3.03} yields
\[
 \begin{split}
 \delta_v
 &\le\sum_{i\bmod n_p}\bigl(2N_mz+3N_m^2z\bigr)
       +\sum_{i\ {\rm marked}}4z^2\\
 &\le5n_pN_m^2z+4N_v(t_0)z^2\le\bar\delta(t_0).
 \end{split}\tag{16}
\]
For a marked \(i\), (13) gives
\[
 \left[\frac{\bar\alpha_{\iota(i)}}{p^d}-\bar C_{\iota(i),d}\right]_+
 \le \left[\frac{\alpha_i}{p^d}-C_{i+d}\right]_+ .
\]
By (14), the local ordinary gain retains every index retained in the global ordinary gain. For (1), sum the positive-part inequality at odd marks and use \(\gamma_p\bar\alpha_{\iota(i)}\le\gamma_p\alpha_i\) at even marks. Convergence follows from the geometric bounds after \ref{B3.49}. Subtract the errors by (16) and use \ref{B3.50} for the price.
\end{proof}
\leanref{Formalism.\allowbreak{}MainNotions.\allowbreak{}FormalGroupSystem.\allowbreak{}Pf.\allowbreak{}nc\_\allowbreak{}loc}{Formalism.\allowbreak{}SharpRestriction.\allowbreak{}GrowthBounds.\allowbreak{}MultisetValuation.\allowbreak{}NumericalClosedOfNeed}{12--19}
\leanref{Formalism.\allowbreak{}MainNotions.\allowbreak{}FormalGroupSystem.\allowbreak{}Pf.\allowbreak{}lw\_\allowbreak{}size}{Formalism.\allowbreak{}SharpRestriction.\allowbreak{}GrowthBounds.\allowbreak{}ProofLW}{135--203}
\leanref{Formalism.\allowbreak{}MainNotions.\allowbreak{}FormalGroupSystem.\allowbreak{}Pf.\allowbreak{}lw\_\allowbreak{}de}{Formalism.\allowbreak{}SharpRestriction.\allowbreak{}GrowthBounds.\allowbreak{}ProofLW}{62--104}
\leanref{Formalism.\allowbreak{}MainNotions.\allowbreak{}FormalGroupSystem.\allowbreak{}Pf.\allowbreak{}nc\_\allowbreak{}gB\_\allowbreak{}ge}{Formalism.\allowbreak{}SharpRestriction.\allowbreak{}GrowthBounds.\allowbreak{}MultisetValuation.\allowbreak{}NcSummableMin}{120--146}
\leanref{Formalism.\allowbreak{}MainNotions.\allowbreak{}FormalGroupSystem.\allowbreak{}Pf.\allowbreak{}nc\_\allowbreak{}gA\_\allowbreak{}ge}{Formalism.\allowbreak{}SharpRestriction.\allowbreak{}GrowthBounds.\allowbreak{}MultisetValuation.\allowbreak{}NcSummableMin}{147--166}

\begin{theorem}[Zero counted gain when \(z_p<1\)]\label{B3.54}
Assume \(\Lean{ParameterChoice'}(D,L)\), that every component of \(K\) is nonnegative, and that \(K_{\mathrm{Pr}}\ge1\). For every coefficient-field place \(w\), every \(t_0\in\mathcal R_p\), and every \(1\le j\le m\), where \(p=\operatorname{char}k(w)\), one has
\[
 z_p<1\quad\Longrightarrow\quad
 g_U(t_0,j)-\mathrm{Pr}(t_0,j)\le0,\qquad g_0(t_0,j)\le0. \tag{17}
\]
\end{theorem}
\leanref{Formalism.\allowbreak{}MainNotions.\allowbreak{}G.\allowbreak{}TopChoiceBounds.\allowbreak{}hLarge\_\allowbreak{}proved}{Formalism.\allowbreak{}SharpRestriction.\allowbreak{}GrowthBounds.\allowbreak{}TopChoiceBounds}{4226--4373}

\begin{proof}
The hypothesis gives \(p>2a_\ast\). Write \(S=L.\mathrm{Sc}\). The parameter bounds give \(k\ge5\), \(S\ge100\), \(\kappa_j\le S\), and
\[
 kS\le w''\le b'/x\le b'\le a_\ast/2 .
\]
At a mark the weight window gives \(a_s\ge a_\ast\ge2kS\ge2c\) and \(a_s\le a_\ast+a_\ast/(kx^{13})\le6a_\ast/5\). The filling is nonnegative, so \(\bar C_{s,d}\ge4S-\kappa_j\ge3S\ge1\). Since \(z_p\ge0\), we have
\[
 0\le\bar\alpha_s\le6a_\ast/5\quad(s\in T(t_0)),\qquad
 \bar C_{s,d}\ge1,\qquad \bar\delta(t_0)\ge0. \tag{18}
\]
For each positive odd \(d\), \(\bar\alpha_s/p^d\le6a_\ast/(5p)<1\). Thus \(g_0=-\bar\delta\le0\). All minimum terms in \(g_U\) are nonnegative. Since \(a_\ast<p/2\) and \(3p^2/5\le p^2-1\) for \(p\ge2\), one has \(\gamma_p(6a_\ast/5)\le1\). Therefore
\[
 g_U(t_0,j)\le\gamma_p\sum_{s\in T(t_0)}\bar\alpha_s\le N_v(t_0).
 \tag{19}
\]
The fiber lies in a single \(\Delta_p\)-coset, so \(N_v(t_0)\le n_p\). The scale bounds give \(M\ge1\). At a mark,
\(B_j(0,s)-b_s=\operatorname{off}_j(0,s)-\kappa_j\ge b'/x-S\ge0\).
The large-prime branch gives \(\mathrm{Pr}\ge K_{\mathrm{Pr}}n_pM\ge N_v(t_0)\). The small-prime branch gives the same bound since \(r_G\ge1\) and \(2b'+c\ge1\).
\end{proof}

\section{Substitution into the integral theta chain}

Let \(\mathcal O\) be the local coefficient discrete valuation ring with fraction field \(K_f\). Choose a uniformizer \(\varpi\). Put \(e_{\mathcal O}=v_\varpi(p)\) and write \(v_\Lambda(f)\ge q\) for \(f\in\varpi^{\lceil e_{\mathcal O}q\rceil}\Lambda\), as in \ref{B3.01}.

\begin{theorem}[Global exponents in the local chain]\label{B3.55}
Fix tied data \(D,L,X,Y,K\), a word \(u\), and a place \(w\) satisfying \(\Lean{GIdx}\). Assume (12) of \ref{B3.53} when \(z_p\ge1\) and (17) of \ref{B3.54} otherwise. Suppose the local system satisfies the ordinary coefficient bound of \ref{B3.27} and the zero coefficient bound \(\Lean{adm11}\). Assume the chain bound of \ref{B3.42} for the same system. Assume \(\mathcal O\to K_f\) is injective.

At each local position \(v\notin v(\mathcal R_p)\), assume real bounds \(F_n(v,j)\) with \(\operatorname{LowAt}(j-1,v,F_n(v,j))\) for \(1\le j\le m\). Let \(\mathrm{low}\in\R\) satisfy
\[
 \mathrm{low}\le\sum_{j=1}^m\sum_{v\notin v(\mathcal R_p)}F_n(v,j)
       -\sum_{j=0}^{m-1}d_G(j). \tag{20}
\]
These assumptions give \(\Lean{SysUAt}\).

For \(V'\subseteq V_s\subseteq\mathcal R_p\) and \(f\) in the initial integral lattice \(M\), let \(y_{V'}f=(\prod_{t_0\in V'}y_{v(t_0)}^{\,m})f\), in the prescribed finite-set order. Then
\[
 v_{\Lambda_m}\bigl(K^{[m]}y_{V'}f\bigr)
 \ge \mathrm{low}+\sum_{j=1}^m
 \left(\sum_{t_0\in V'}\max\{g_U(t_0,j)-\mathrm{Pr}(t_0,j),g_0(t_0,j),0\}
       +\sum_{t_0\in\mathcal R_p\setminus V'}[g_0(t_0,j)]_+\right).
 \tag{21}
\]
Here \(\Lambda_m\) is the terminal integral lattice and \(K^{[m]}=K_{m-1}\cdots K_0\).
\end{theorem}
\leanref{Formalism.\allowbreak{}MainNotions.\allowbreak{}G.\allowbreak{}yProd}{Formalism.\allowbreak{}SharpRestriction.\allowbreak{}GrowthBounds.\allowbreak{}IsResidueClassZero}{453--457}
\leanref{Formalism.\allowbreak{}MainNotions.\allowbreak{}G.\allowbreak{}Joints.\allowbreak{}UAt\_\allowbreak{}of}{Formalism.\allowbreak{}SharpRestriction.\allowbreak{}GrowthBounds.\allowbreak{}JointsTw}{23--215}
\leanref{Formalism.\allowbreak{}MainNotions.\allowbreak{}WeilJacquetShalikaLaw.\allowbreak{}UAt}{Formalism.\allowbreak{}SharpRestriction.\allowbreak{}GrowthBounds.\allowbreak{}WeilJacquetShalikaLaw}{161--170}

\begin{proof}
In \(\Lean{UAt\_\allowbreak{}of}\), (12) ranges over all tied \(\Lean{GIdx}\) pairs and (17) over all coefficient-field places with \(z_p<1\). Apply the corresponding bound at \((u,w)\).

The ordinary and zero coefficient bounds give \(\operatorname{LowAt}(j-1,v,q)\) for \(q\le\max\{g_{A,v}(j),0\}\).
\leanref{Formalism.\allowbreak{}MainNotions.\allowbreak{}FormalGroupSystem.\allowbreak{}lowAt\_\allowbreak{}max}{Formalism.\allowbreak{}SharpRestriction.\allowbreak{}GrowthBounds.\allowbreak{}FormalGroupSystem}{718--725}

Put \(B=v(V')\) and set \(G_n=P_n=0\) off \(B\). If \(z_p\ge1\), choose on \(B\)
\[
 G_n(v(t_0),j)=g_U(t_0,j),\qquad P_n(v(t_0),j)=\mathrm{Pr}(t_0,j),
\]
and set \(G_{0n}(v(t_0),j)=[g_0(t_0,j)]_+\) for every \(t_0\in\mathcal R_p\). Equation (12) gives the selected-position bounds for \(G_n,P_n\) and gives \(G_{0n}(v(t_0),j)\le\max\{g_{A,v(t_0)}(j),0\}\).

If \(z_p<1\), choose on \(B\)
\[
 G_n(v,j)=\min\{g_{B,v}(j),P_v(j)\},\qquad P_n(v,j)=P_v(j),
\]
and set \(G_{0n}=0\) on \(v(\mathcal R_p)\). On \(B\), \(G_n\le g_B\) and \(G_n-P_n\le0\), with price bound \(P_v=P_n\). Equation (17) makes the global gain terms in (21) zero.

Thus \(G_{0n}\) satisfies the required \(\operatorname{LowAt}\) bound on \(v(\mathcal R_p)\). Elsewhere set \(G_{0n}(v,j)=F_n(v,j)\).

Set \(S(v)=\sum_{j=1}^m[G_n(v,j)-P_n(v,j)-G_{0n}(v,j)]_+\). The element \(y_{V'}f\) is a generator of \(\operatorname{Iprod}(S,B)\) in \ref{B3.40}, with distinct positions by injectivity and complementary scalar \(1\).

Apply \ref{B3.42}. Its terminal exponent is
\[
 \sum_{j=1}^m
 \left(\sum_{v\in B}\max\{G_n(v,j)-P_n(v,j),G_{0n}(v,j)\}
       +\sum_{v\notin B}G_{0n}(v,j)\right)-\sum_{j=0}^{m-1}d_G(j).
 \tag{22}
\]
On \(B\) the first sum dominates the selected gain in (21). The represented positions off \(B\) provide the ordinary positive parts. By (20) the remaining positions and stabilizer terms contribute at least \(\mathrm{low}\). In \(\Lean{UAt}\) the ordinary sum splits over \(V_s\setminus V'\) and \(\mathcal R_p\setminus V_s\).

Multiplication by \(e_{\mathcal O}\) and taking ceilings preserve the inequality. For the resulting integers \(a\ge b\), \(\mathcal O\)-stability gives \(\varpi^a\Lambda_m\subseteq\varpi^b\Lambda_m\).
\end{proof}

\section{Passage to the global service condition}

\begin{theorem}[The local theorem supplies the global service]\label{B3.56}
Let \(\mathcal C\) be a class of admissible parameter families. Assume \(\Lean{DGamSize}\) on \(\mathcal C\) with the \(\Lean{Closed}\) quantifiers and assume \(\Lean{SumOverPrimesHypotheses}(\mathcal C)\). For sufficiently large members, the first clause supplies levels and chains \(LC\) for which the given map satisfies \(\Lean{CmpV}\) and the local systems satisfy (20) at \(\mathrm{low}=\operatorname{lowS}(u,w)\) under the hypotheses below. The second clause supplies the kernel data listed in \ref{B4.16}. For the selected family these data come from \ref{B8.56}. Then \(\Lean{ServF}\) holds with the \(\Lean{Closed}\) family quantifiers.

For each resulting family \(F_m\) there are \(C\ge0\) and \(x_0\) such that \(\Lean{Serv}\) holds for every \(x\ge x_0\) with \(x>x_1\) and every member \(i\in I_x\). At service parameter \(1\), \(\Lean{Serv}\) includes \(\Lean{TypePrimeSizeConds}\). For each member let \(D_x\) be its divisibility data and put \(Q'(\mu)=D_x.\mathrm{Qp}(\mu)\). Fix a word \(u\) and a used twist \(\mu\). At every finite place \(w\) with \(p\ne T_u,g_0,g_1\) and \(p\notin Q'(\mu)\), the bounds are
\[
 \begin{array}{ll}
 \operatorname{ColGe}(u,\mu,w,f,\operatorname{totS}(u,w)),&
 f\text{ satisfies }\Lean{InNP},\\[2pt]
 \operatorname{ColGe}(u,\mu,w,f,\operatorname{lowS}(u,w)),&
 f\in S^{\rm ram}_{u,w}.
 \end{array} \tag{23}
\]
The first bound assumes \(\Lean{Ties9}\) and \(\Lean{ErrDefs}\). The second assumes \(\Lean{FDef}\) and \(\Lean{ErrDefs}\). By \ref{B4.07}, \(\Lean{InNP}\) means \(c_{u,w,h_g}(f)\in NP_{u,w,h_g}\) for some \(\Lean{GIdx}\) witness \(h_g\), or \(V_p=\varnothing\) and \(f\in S^{\rm ram}_{u,w}\).
\end{theorem}
\leanref{Formalism.\allowbreak{}MainNotions.\allowbreak{}Pk.\allowbreak{}sum\_\allowbreak{}over\_\allowbreak{}primes\_\allowbreak{}service}{Formalism.\allowbreak{}SharpRestriction.\allowbreak{}SmoothDesignExistence.\allowbreak{}SumOverPrimesHypotheses}{187--221}
\leanref{Formalism.\allowbreak{}MainNotions.\allowbreak{}WeilJacquetShalikaLaw.\allowbreak{}Serv}{Formalism.\allowbreak{}SharpRestriction.\allowbreak{}GrowthBounds.\allowbreak{}WeilJacquetShalikaLaw}{241--243}
\leanref{Formalism.\allowbreak{}MainNotions.\allowbreak{}WeilJacquetShalikaLaw.\allowbreak{}ServF}{Formalism.\allowbreak{}SharpRestriction.\allowbreak{}GrowthBounds.\allowbreak{}WeilJacquetShalikaLaw}{252--255}

\begin{proof}
Fix a used twist \(\mu\) and assume \(\Lean{ErrDefs}\). Impose the tying conditions \(\Lean{GatTies}\) and \(\Lean{LatTied}\). At each \(\Lean{GIdx}\) place with \(p\notin Q'(\mu)\), the first clause supplies an admitted system. By \ref{B3.45} it satisfies the ordinary coefficient and chain bounds, and admissibility gives the zero coefficient bound. Use \ref{B3.53} and \ref{B3.54} for the numerical comparisons. With (20), \ref{B3.55} gives \(\Lean{UAt}\) at \(\mathrm{low}=\operatorname{lowS}(u,w)\).
\leanref{Formalism.\allowbreak{}MainNotions.\allowbreak{}G.\allowbreak{}Joints.\allowbreak{}unit\_\allowbreak{}bound\_\allowbreak{}of\_\allowbreak{}sum\_\allowbreak{}over\_\allowbreak{}primes}{Formalism.\allowbreak{}SharpRestriction.\allowbreak{}GrowthBounds.\allowbreak{}JointsTw}{243--270}

By \ref{B4.10}, each generator's scalar exponent covers the omitted orbit gains, and linearity extends the bound to \(NP\). At every \(\Lean{GIdx}\) place for the used twist, including \(p\in Q'(\mu)\), the place-specific form of \ref{B3.45} gives \(\Lean{LowTermBound}\) for each matching admitted system.
\leanref{Formalism.\allowbreak{}MainNotions.\allowbreak{}WeilJacquetShalikaLaw.\allowbreak{}local\_\allowbreak{}theorem\_\allowbreak{}of\_\allowbreak{}statement}{Formalism.\allowbreak{}SharpRestriction.\allowbreak{}GrowthBounds.\allowbreak{}WeilJacquetShalikaLaw}{371--375}

At a nonexceptional signal prime \(p\le X''\), (5) gives \(V_p=\varnothing\), and (4.21) of \ref{B4.11} supplies the ordinary gain. For \(p>X''\), (6) gives zero counted gain. With \((C_6,C)=(11,1)\), \ref{B4.11a} gives the orbit estimates used in \ref{B4.11}(i)--(iii). For (iv), jet integrality gives entry exponent zero at the level-free places away from \(d\), and (4.20) of \ref{B4.10a} gives the lift bound. Apply \ref{B4.16} with the common threshold (4.49) to obtain (23) for the original filling.
\leanref{Formalism.\allowbreak{}MainNotions.\allowbreak{}WeilJacquetShalikaLaw.\allowbreak{}serv\_\allowbreak{}of}{Formalism.\allowbreak{}SharpRestriction.\allowbreak{}GrowthBounds.\allowbreak{}WeilJacquetShalikaLaw}{376--379}
\leanref{Formalism.\allowbreak{}MainNotions.\allowbreak{}Pk.\allowbreak{}serv\_\allowbreak{}re}{Formalism.\allowbreak{}SharpRestriction.\allowbreak{}SmoothDesignExistence.\allowbreak{}SumOverPrimesHypotheses}{98--107}
\end{proof}

Suppose \(\Lean{ParametersFill}\) holds at every finite place and \(L.\mathrm{drop}\) contains every inert-class prime for which \(h_p\) is nonzero and has odd order. Then \ref{B4.25} gives \(\gainD_j(w)\le\operatorname{gainS}(p,j)\) for \(p=\operatorname{char}k(w)\) and \(1\le j\le m\), with \(\gainD\) as in \ref{B4.04}. The period-divisibility theorem \ref{B4.29} uses this inequality with \(\Lean{Serv}\).


\part{Divisibility along the chain and division of lattices}

% ---- p4c1 ----
\chapter{What each layer records}\label{ch:layer-records}

Fix chain data $D$ of initial rank $k$ and length $m$ with size data $L$
\cite{alpoge-2026}. At a finite place $w$ of $K$ put $p=\operatorname{resChar}(w)$
and $e_w=v_w(p)$. Use the first applicable case in each formula.
Let $\mathcal C,\mathcal S$ be the inert and signal classes of \ref{B3.46}.
For $z=(\epsilon,t)\in Z=\{0,1\}\times A$, let $\tau(z)$ be the
embedding index for $\tau_{\epsilon+1}(t)$ under $D.eTau$.
For $1\le j\le m$ let $w_j(z)=w_j(\tau(z))$ be extension weight $j-1$.

\begin{definition}[The orbit base value]\label{B4.01}
For a prime $p$, an integer $r\ge1$ and weights $d_s\ge0$ on $\Z/r\Z$,
let $\mathcal P$ consist of arrays $x_{s,i}\ge0$ for odd $1\le i<r$
with $\sum_i(p^ix_{s,i}+x_{s-i,i})\le d_s$. Put
\[
 \operatorname{lpF}(p,r,\mathbf d)=\frac p{p^2-1}\sum_s d_s
 -\sup_{x\in\mathcal P}\sum_{s,i}x_{s,i}.
\]
Use $h_p,n_p$ from \ref{B3.46}. Set
$P_p(\epsilon,t)=(\epsilon+\mathbf1_{p\in\mathcal S},t+h_p)$
with first coordinate modulo $2$, and define
\[
 F_{\rm orb}(j,p,z)=\operatorname{lpF}
 \bigl(p,n_p,((k-1)w_j(P_p^sz))_{s\in\Z/n_p\Z}\bigr).
\]
\end{definition}
\leanref{Formalism.\allowbreak{}BasicNotions.\allowbreak{}lpF}{Formalism.\allowbreak{}LocalRingCardinality.\allowbreak{}DataParams₂}{58--60}
\leanref{Formalism.\allowbreak{}BasicNotions.\allowbreak{}FOrb}{Formalism.\allowbreak{}LocalRingCardinality.\allowbreak{}DataForced}{99--100}
\begin{proof}
The period is $\operatorname{ord}(h_p)$ off $\mathcal S$ and
$2\operatorname{ord}(2h_p)$ on $\mathcal S$, hence $n_p$. The demands are
nonnegative by the size bounds and $k\ge1$. The set $\mathcal P$ is
compact and contains zero. Translating $s$ preserves the sums and shifts
the demands, so $F_{\rm orb}$ is constant on each orbit.
\end{proof}
\leanref{Formalism.\allowbreak{}BasicNotions.\allowbreak{}FOrb\_\allowbreak{}step}{Formalism.\allowbreak{}LocalRingCardinality.\allowbreak{}DataForced}{102--108}

\begin{lemma}[Changing the demands on an orbit]\label{B4.01a}
Write $F(D)=\operatorname{lpF}(p,r,D)$ and $\gamma=p/(p^2-1)$.
Suppose $p\ge5$ is prime, $r\ge2$ is even, $T\subseteq\Z/r\Z$,
$k\ge3$, $w_s>0$ and $0\le e_s\le4w_s$. Set $W_s=(k-1)w_s$ and
$D_s=W_s+e_s-2w_s\mathbf1_T(s)$. Then
\[
 F(D)\ge F(W)-\frac{11}{p}\sum_s w_s.
\]
\end{lemma}
\begin{proof}
The dual minimizes $\sum_s W_sy_s$ for $y_s\ge0$ with
$p^iy_s+y_{s+i}\ge1$ for odd $1\le i<r$. Take optimizers $x,y$ with
$y_s\le1$, so $F(W)=\sum_s W_s(\gamma-y_s)$. For $y_s\ge2/p$,
complementary slackness gives tight capacity and zero outgoing flow,
hence $W_s=\sum_i x_{s-i,i}$ and
\[
 \sum_{y_s\ge2/p}W_s\le\sum_{s,i}x_{s,i}\le p^{-1}\sum_s W_s .
\]
The displayed bound on incoming demand holds for odd periods too.
Put $S=\sum_s w_s$. Since $\gamma\ge1/p$, testing $y$ at $W+e$ loses
at most $\sum_s e_s(y_s-\gamma)_+\le4S/p+4\sum_{y_s\ge2/p}w_s\le8S/p$.
Primal monotonicity gives $F(W+e)\le F(D)+2\gamma\sum_{s\in T}w_s
\le F(D)+3S/p$.
\end{proof}
\leanref{Formalism.\allowbreak{}DworkPowerSeriesWeights.\allowbreak{}OddLevelGainBounds.\allowbreak{}proposition\_\allowbreak{}six}{Formalism.\allowbreak{}GaussJetBounds.\allowbreak{}OneLetterGaussJets.\allowbreak{}OneLetterGaussProofs.\allowbreak{}CompanionTheory}{38--147}
\leanref{Formalism.\allowbreak{}MainNotions.\allowbreak{}DIV.\allowbreak{}Bridge.\allowbreak{}DworkPowerSeriesWeights.\allowbreak{}local\_\allowbreak{}bound\_\allowbreak{}eleven\_\allowbreak{}comparison}{Formalism.\allowbreak{}SharpRestriction.\allowbreak{}Divisibility.\allowbreak{}JetPolynomialLogs.\allowbreak{}LocalBoundElevenComparison}{33--55}

\begin{proposition}[Cyclic and one-letter bounds]\label{B4.01b}
Let $p\ge5$ be prime and $N_0,N_1\in\N$ with $N_0<N_1$. On the even
cycle $\Z/r\Z$, $r\ge2$, take Young diagrams $\nu_s$ with at most $N_0$
rows. Set $\ell_{s,ib}=\sum_{h\ge0}Z_{s-2h,ib}^{p^{2h}}/p^h$ and suppose
\[
 \Theta(\xi\ell)=U\exp\!\left(\epsilon\sum_{i,b}
 \beta(\ell_{\bullet,ib})\right),
 \qquad |\xi|_p=1,\quad\epsilon=\pm1,\quad |[Z^n]U|_p\le1
\]
for all $n$, with $(i,b)$ over the $N_0N_1$ entries. Take $\beta$ over
the completed maximal unramified extension of $\Q_p$, invariant under
$Z_s\mapsto u^{(-1)^sp^s}Z_s$ for $u^{p^r-1}=1$. Each monomial has
positive degree in both parity blocks and bidegree different from $(1,1)$.
Set $\Delta_X\beta=\beta(X+X',Y)-\beta(X,Y)-\beta(X',Y)$, and similarly for $Y$.
Assume both exponentials of these differences are integral with constant
one after letter-logarithm substitution.

Let $A_s,B_s$ be the first two row lengths, with missing rows of length
zero. For $T\subseteq\Z/r\Z$ set $D_s=|\nu_s|-A_s\mathbf1_T(s)$,
where $|\nu_s|$ counts boxes. Put $c_s=B_s$ on $T$ and $c_s=A_s$
elsewhere. Use positive odd distances $e$ below. For $s\in T$ let
$h_s=\inf\{e:s+e\in T\}$ and $M_s=\max_{p^e\le|\nu_{s-e}|}c_{s-e}$,
with empty values $\infty$ and zero. Set
\[
 \begin{aligned}
 H_s&=\sum_{e<h_s}(A_sp^{-e}-c_{s+e})_++\sum_{e>h_s}A_sp^{-e},\\
 J_s&=\sum_{e\ne h_s}((A_s-B_s)p^{-e}-c_{s+e})_+ .
 \end{aligned}
\]
For each column of height $h$, choose a strictly increasing map
$\{1,\ldots,h\}\to\{1,\ldots,N_1\}$ and form $P_I$ by \ref{B3.10}.
If $p,|\nu_s|\le a$ for all $s$ and $2\log_p a+2\le N_0-1$, every
coefficient $[\tau^\mu]\mathcal D_{P_I}(\Theta)$ of \ref{B3.11} has
valuation at least
\[
 F(D)+\sum_{s\in T}(H_s-\gamma M_s)-rN_0N_1(\log_p a+1)
\]
with $F,\gamma$ from \ref{B4.01a}. Under its hypotheses
$F(D)+\sum_{s\in T}H_s$ may be replaced by
$F(W)-11\sum_s w_s/p+\sum_{s\in T}J_s$.

For one letter let $\nu$ have at most $N_0$ rows and set
$\ell=\sum_{h\ge0}Z^{p^{2h}}/p^h$. Use $(\epsilon/2)\beta(\ell_{ib},\xi'\ell_{ib})$
at entry $(i,b)$, with $|\xi'|_p=1$ and the two-letter hypotheses on
$\beta$. For $p,|\nu|\le a$ the bound is
$\rho|\nu|-N_0N_1(\log_p a+1)$ with $\rho=1/(p^2-1)$.
\end{proposition}
\begin{proof}
Give $cZ^n$ weight $v_p(c)+\rho|n|$, with $v_p(p)=1$, and put
$q_*=1/(p-1)$. From $p^{2h}\rho-h\ge\rho$, degree recursion gives
the inverse logarithm weight at least $\rho$. Substitution gives the
integral coboundary exponentials positive nonconstant weights.
Invariance gives $p+1\mid|n|$ since $(-1)^sp^s\equiv1\bmod(p+1)$.
With integral valuations, positive weights are at least $q_*$.
The logarithm preserves this by $j q_*-v_p(j)\ge q_*$.

A split extracts the monomial coefficient times a binomial factor
from the coboundary. A block with two variables or a sole exponent
of digit sum greater than one admits a unit split. If neither block
does, the monomial is $Z_s^{p^u}Z_t^{p^v}$ with $u>0$ or $v>0$.
Since $v_p\binom{p^h}{p^{h-1}}=1$ for $h>0$, its weight is at least
$q_*-1$. Invariance gives $t=s+u-v\bmod r$, with $u-v$ odd.

In the Gauss frame of \ref{B3.26}, reserve an ordinary row at $s+j$
for each marked $s$ and signed odd $j$ satisfying $p^j\le A_s$ for $j>0$
and $p^{-j}\le|\nu_{s+j}|$ for $j<0$. The rank bound suffices since
$p^{|j|}\le a$. Raise marked first rows to $p^M$-th powers, $M\ge|j|$
for all labels. At $s+j$ use the form whose coefficients are the
$p^{M+j}$-th powers of the original first row at $s$. On reserved rows
the matrix of the forms has unit determinant by the distinct diagonal
residue monomial. Invert it and subtract labelled components from other
ordinary rows. The forms are coordinate projections on the ordinary
span. The marked first rows stay fixed. Expressing the old rows in the
new ones rewrites the minors integrally. By alternation the partner degree
is at most $c_{s+j}$, since one-row minors at a marked target use its first row.

Fix $\kappa$ in the $Z$-support of a transformed minor product.
Contributing phase exponents satisfy $n\le\kappa$ coordinatewise.
For a mixed pair with marked source $s$ and $j=u-v$, this gives
$p^u\le A_s$ and $p^v\le|\nu_{s+j}|$, so the label exists. Modulo $p$
the contraction is the $p^v$-th power of the form on the target row.
Off the labelled row divisibility by $p$ raises the weight to at least $q_*$.

On ordinary rows add the optimal dual weight $y_s$ for $D$. Reduction
modulo $r$ extends the constraint to every odd distance $e>0$. On
marked first rows add $p^{-b_s}$, where $b_s$ is the first odd $e<h_s$
with $A_sp^{-e}\le c_{s+e}$, or $h_s$ if none exists. Set $p^{-\infty}=0$.
Raise the labelled increment to at least $(1-p^{j-b_s})_+$ for $j>0$
and $p^{-|j|}$ for $j<0$. The two increments times $p^u,p^v$ sum to
at least one. For two marked rows use $b_s\le h_s$ at the forward
source, and for two ordinary rows use the dual constraint. Thus phase
terms with $n\le\kappa$ have weight at least $q_*$. The bound
$v_p(h!)\le hq_*$ gives nonnegative exponential weights through
$\kappa$. The integral factor has nonnegative weight after inverse
logarithm substitution and the frame change.

With $\kappa!=\prod_i\kappa_i!$ and $\kappa_i\le a$, \ref{B3.19}
gives $v_p(\kappa!)\ge|\kappa|q_*-rN_0N_1(\log_p a+1)$. Since
$q_*-\rho=\gamma$, ordinary degrees give $F(D)$ before label costs.
From $\gamma A_s$ subtract $A_sp^{-b_s}$ and the forward costs
$c_{s+e}(1-p^{e-b_s})_+$. These cost at most $c_{s+e}$ before $b_s$
and vanish thereafter. The balance is at least $H_s$ since
$\gamma=\sum_e p^{-e}$. Each reverse label has $c_{s-e}\le M_s$,
so reverse costs total at most $M_s\sum_e p^{-e}=\gamma M_s$.
The Gauss norm bounds each frame coefficient. Powering the marked
variables preserves it by injectivity on monomials. Apply
\ref{B4.01a} and $J_s\le H_s$.

For one letter use weight $\gamma$. Exceptional pairs have degree at
least $p+1$ and $(p+1)(\gamma-\rho)=1$. The unit specialization and
$\epsilon/2$ preserve the bound. The factorial gives $q_*-\gamma=\rho$
per box before digit loss.

The minor product $P_I$ is homogeneous in $Z$ of degree $d_\nu$, the
box count. Thus $\mathcal D_{P_I}(\Theta(\xi Z))=
\xi^{d_\nu}\mathcal D_{P_I}(\Theta)$, which preserves valuations.
\end{proof}
\leanref{Formalism.\allowbreak{}MainNotions.\allowbreak{}DIV.\allowbreak{}Bridge.\allowbreak{}DworkPowerSeriesWeights.\allowbreak{}split\_\allowbreak{}bound\_\allowbreak{}eleven\_\allowbreak{}one}{Formalism.\allowbreak{}SharpRestriction.\allowbreak{}Divisibility.\allowbreak{}JetPolynomialLogs.\allowbreak{}LocalBoundElevenComparison}{64--68}

\begin{definition}[The base term of a layer]\label{B4.02}
Put $n=|A|$, $X'=e^{x^{\varepsilon/3}}$ and $X''=nxX'$.
For the integer row tuples $\kappa^{\rm big}_{j,\tau}$,
$\kappa^{\rm small}_{j,\tau}$ of lengths $k+j-1$, $m-j$, set
\[
 \begin{aligned}
 F_d(j)&=-\tfrac{k-1}{2}\sum_{t\in A}\min\{w_j(0,t),w_j(1,t)\},\\
 D_j^{\rm Fock}&=\sum_\tau\left(\sum_a|\kappa^{\rm big}_{j,\tau,a}|
 +\sum_a|\kappa^{\rm small}_{j,\tau,a}|\right),\qquad F_g(j)=-D_j^{\rm Fock}/2.
 \end{aligned}
\]
Under \Lean{FDef},
\[
 F_j(w)=\begin{cases}
 F_g(j),&p\in\{g_0,g_1\},\\
 F_d(j),&p=d,\\
 \sum_\tau(k-1)w_j(\tau)/(p^2-1),&p\in\mathcal C,\ h_p=0,\ p\le X'',\\
 n_p^{-1}\sum_{z\in Z}F_{\rm orb}(j,p,z),&p\in\mathcal C\cup\mathcal S,\ p\le X'',\\
 0,&\text{otherwise}.
 \end{cases}
\]
\end{definition}
\leanref{Formalism.\allowbreak{}BasicNotions.\allowbreak{}SizeData₂.\allowbreak{}FDef}{Formalism.\allowbreak{}LocalRingCardinality.\allowbreak{}DataForced}{134--135}
\begin{proof}
By \ref{B4.01}, division by $n_p$ counts each orbit once.
As $d\nmid Q_S$ and $d\ne q$, $d$ has class \Lean{ramD}.
\end{proof}

\begin{definition}[The loss charged to a column]\label{B4.03}
Let $M_u>0$ be the stabilizer denominator of column $u$ and put
$\delta_{u,w}=v_p(M_u)$. Under \Lean{DGamDef} all columns and twists
have denominator $M_{u_0}$, which by \Lean{GamMult} is divisible by
each class stabilizer order at stages $0\le i\le m$.

Set $t_{\rm ch}=D.R.t$ (odd), $N_i=k+i(t_{\rm ch}+1)$, $f=2|A|$,
$\eta=L.\mathrm{etaSz}$, $\mathrm{Sc}=L.\mathrm{Sc}$ and $a_*=xb$.
Use the auxiliary error datum $X$ of \Lean{ErrDefs}, with constants
$c_g\ge10$, $c_{\rm dig}\ge3$ and $c_6\ge12$. For the host field $E$,
$\pi_i=X.\mathrm{piG}(i)\in\mathcal O_E\cap\Q(\sqrt{-d})$ satisfies
$\pi_i\bar\pi_i=g_i$ for $i=0,1$.

For $0\le i<N_{j-1}$ let $s_{u,\tau,i}$ be initial weight $i$ at $\tau$
if $i<k$. For $i=k+\ell(t_{\rm ch}+1)+r$, $0\le r\le t_{\rm ch}$, set
$s_{u,\tau,i}=-w_{\ell+1}(\tau)+(t_{\rm ch}-2r)/2$. Put
\[
 \nu_{u,j,\tau,i}=s_{u,\tau,i}+w_j(\tau)-\frac{N_{j-1}+N_j-1-2i}{2},
 \qquad A_{u,j,\tau}=\sum_{i<N_{j-1}}\nu_{u,j,\tau,i}.
\]
For fixed $u,j$, index $A,\nu$ by $(\epsilon,t)$ via $\tau(\epsilon,t)$.
Put
\[
\begin{aligned}
 G_{i,u,j,w}&=F_g(j)+e_w^{-1}\sum_\tau A_{u,j,\tau}v_w(\tau(\bar\pi_i))
   +c_gf(N_m^2+m^2\eta),\\
 c_{u,j,t}&=A_{0,t}+\mathbf1_{t\in TT}(\nu_{0,t,1}-\nu_{0,t,0}+\nu_{1,t,0}),\\
 D_j^*&=F_d(j)+\tfrac12\max_u\sum_{t\in A}c_{u,j,t},\\
 C_{j,p}^*&=c_{\rm dig}fN_m^2\log_p a_*+\frac{\sum_\tau w_j(\tau)}{p^2-1}.
\end{aligned}
\]
Set $Z_A=\{(\epsilon,t)\in Z:p\notin\mathcal S,\ \epsilon=0\}$,
$Z_{BC}=Z\setminus Z_A$ and $Z_{\rm top}=Z_{BC}\cap(\{0\}\times TT)$.
Let $R_{BC}(j,p)$ be (4.26) of \ref{B4.11} with $BC=Z_{BC}$,
$TT_{BC}=Z_{\rm top}$, $\gamma(p)=p/(p^2-1)$, $c_{e,6}=c_6$ and
$c_{\mathrm{dl},i}=K_{\delta,i}$ from \ref{B3.47}. Put
\[
 I_{j,p}^*=\max\{0,n_p^{-1}\sum_{z\in Z_A}F_{\rm orb}(j,p,z)\}
       +R_{BC}(j,p)+X_{j,p}^*,
\]
where $X_{j,p}^*=(k-1)p^{-3}\sum_{z\in Z_{BC}}w_j(z)+c_6fN_m^2\log_p a_*$
if $p\in\mathcal C$ and $h_p$ has odd order greater than $1$, and zero
otherwise.

Let $T_u$ be the level prime and $\mathrm{Top}=\{\tau(0,t):t\in TT\}$.
For integral highest weights $\lambda^u_{i,\tau,r}$ put
$W_{u,\tau}=\max_{i,r}|\lambda^u_{i,\tau,r}-\lfloor i(t_{\rm ch}+1)/2\rfloor|$
over $0\le i\le m$, $0\le r<N_i$, omitting $r=0$ on $\mathrm{Top}$.
Use zero for an empty maximum. Put $\pi_{T,u}=X.\mathrm{piT}(u)\in
\mathcal O_E\cap\Q(\sqrt{-d})$, with $\pi_{T,u}\bar\pi_{T,u}=T_u$.
At $p=T_u$, \Lean{SeenDef} sets
$\epsilon_u(w)=\mathbf1_{\{\forall\tau,\ v_w(\tau(\pi_{T,u}))>0\}}$.
Then
\[
 T_{u,w}^*=\epsilon_u(w)\bigl(2\sum_\tau W_{u,\tau}+fN_m+
        |\mathrm{Top}|(t_{\rm ch}+1)(m+1)/4\bigr).
\]
Under \Lean{EDef},
\[
 E_{u,j}(w)=\begin{cases}
 T_{u,w}^*,&p=T_u,\\
 \delta_{u,w}+G_{i,u,j,w},&p=g_i,\ i=0,1,\\
 \delta_{u,w}+D_j^*,&p=d,\\
 \delta_{u,w}+C_{j,p}^*,&p\in\mathcal C,\ h_p=0,\ p\le X'',\\
 \delta_{u,w}+I_{j,p}^*,&p\in\mathcal C\cup\mathcal S,\ p\le X'',\\
 \delta_{u,w},&\text{otherwise}.
 \end{cases}
\]
\end{definition}
\leanref{Formalism.\allowbreak{}BasicNotions.\allowbreak{}SizeData₂.\allowbreak{}EDef}{Formalism.\allowbreak{}LocalRingCardinality.\allowbreak{}DataErr}{171--172}
\leanref{Formalism.\allowbreak{}BasicNotions.\allowbreak{}DGamDef}{Formalism.\allowbreak{}LocalRingCardinality.\allowbreak{}DataErr}{74--75}
\leanref{Formalism.\allowbreak{}BasicNotions.\allowbreak{}Eval}{Formalism.\allowbreak{}LocalRingCardinality.\allowbreak{}DataErr}{159--169}
\begin{proof}
The trivial twist has denominator $M_u$, so \Lean{DGamDef} applies.
\end{proof}
See \ref{B4.11a} for level-free bounds and \ref{B4.23}--\ref{B4.24} at $g_i$.

\begin{definition}[The gain retained after selection]\label{B4.04}
Use $\rho_B,\Sigma_p$ and candidacy from \ref{B3.51}. Let $C_r\in\R$.
The test $B_p$ holds if $\operatorname{ord}(h_p)<x^{\varepsilon/4}$,
if $p\ge k^2x^{26}$ and $|p^i-X'|\le\rho_BX'$ for some $i\in\N_{>0}$,
or if $p$ is a candidate with $\Sigma_p>x^{-\varepsilon}$.
The predicate $\mathrm{NoGain}(p)$ is the disjunction of $p>X''$,
$p\notin\mathcal C\cup\mathcal S$, and $p\in\mathcal C$ with either
$\operatorname{ord}(h_p)>1$ odd or both $p<X'(1+C_r\rho_B)$ and $B_p$.
For prime $p$, \Lean{DropDef} gives $p\in\mathcal D=L.\mathrm{drop}$
exactly when $\mathrm{NoGain}(p)$ holds and $p\ne T_u$ for all $u$.
Set $\gainD_j(w)=0$ for $p\in\mathcal D$ and $\mathrm{gain}_j(w)$ otherwise.
Assume filling compatibility at $w$ as in \ref{B3.52}.
For $w\in L.\mathrm{supp}$ and $p\in\mathcal C\cup\mathcal S$,
$\mathrm{gain}_j(w)=\mathrm{gainP}(p,j)$ for $1\le j\le m$.
\end{definition}
\leanref{Formalism.\allowbreak{}BasicNotions.\allowbreak{}SizeData₂.\allowbreak{}gainD}{Formalism.\allowbreak{}LocalRingCardinality.\allowbreak{}DataParams₂}{212--213}
\leanref{Formalism.\allowbreak{}MainNotions.\allowbreak{}LocalQuadraticParts.\allowbreak{}DropDef}{Formalism.\allowbreak{}SharpRestriction.\allowbreak{}Divisibility.\allowbreak{}LocalQuadraticParts}{517--518}
\leanref{Formalism.\allowbreak{}MainNotions.\allowbreak{}NoGain}{Formalism.\allowbreak{}SharpRestriction.\allowbreak{}ParConsts}{24--29}
\leanref{Formalism.\allowbreak{}MainNotions.\allowbreak{}G.\allowbreak{}ParametersFill}{Formalism.\allowbreak{}SharpRestriction.\allowbreak{}GrowthBounds.\allowbreak{}IsResidueClassZero}{401--411}
\begin{proof}
The gain is retained at $p=T_u$. At $p=g_0,g_1$ the size data give
zero gain.
\end{proof}

\begin{definition}[The net budget of a column]\label{B4.05}
Set $\sg(u,w)=\sum_{j=1}^m(\gainD_j(w)+F_j(w)-E_{u,j}(w))$.
It vanishes outside $L.\mathrm{supp}$.
\end{definition}
\leanref{Formalism.\allowbreak{}BasicNotions.\allowbreak{}SizeData₂.\allowbreak{}sg}{Formalism.\allowbreak{}LocalRingCardinality.\allowbreak{}DataParams₂}{215--217}
\begin{proof}
Each summand vanishes there by the support conditions.
\end{proof}


% ---- p4c2 ----
\chapter{From local lifts to bounds on whole columns}\label{ch:local-to-column}

Fix the data of Part IV. At a finite place \(w\) of the coefficient field \(K\), normalise \(v_w\) to have value group \(\Z\) and put
\[
 p=\operatorname{resChar}(w),\qquad e_w=v_w(p),\qquad
 \delta_{u,w}=v_p(D\Gamma_u).
\]
Here \(D\Gamma_u\) is the stabiliser denominator and \(T_u\) the level prime of column \(u\). The \(m\ge2\) lifts connect spaces \(V_{u,i}\), \(0\le i\le m\), with lattices \(\Lambda_{u,i}\) and complex embeddings \(\iota_{u,i}\). Superscripts \(\mu\) denote twisted data. The initial isomorphism \(e^\mu_0:V_{u,0}\to V^\mu_{u,0}\) preserves the lattices.

The used twists are \(\mu_0,\mu_\ast,\mu'_\ast\) and the indices \(\mu_w\) for finite \(w\), with \(\mu_0\) the reference index. Let \(Q'(\mu)\subseteq\N\) be the finite set of excluded residue characteristics. With \(F_j(w)\) as in \ref{B4.02} and \(X_e,c_e\) as in \ref{B4.03}, put
\[
 E^{\mathrm{ev}}_{u,j}(w)
 =\operatorname{Eval}(L,X_e,\mathbf c,c_e,u,j,w),
\]
where \(\mathbf c\) denotes the archive's fixed constants \Lean{Consts} \cite{alpoge-2026}. For \(1\le j\le m\), \Lean{ErrDefs} gives \(E^{\mathrm{ev}}_{u,j}(w)=E_{u,j}(w)\). It also supplies the stabiliser denominator condition and the prescribed level and local-place data.

\section{Valuation bounds and the input lattice}

\begin{definition}[Bounds for vectors and columns]\label{B4.06}
Let \(\Lambda\) be an \(\mathcal O_K\)-submodule of a \(K\)-vector space \(V\) and put \(\Lambda_{(w)}=\mathcal O_{K,(w)}\Lambda\). For \(y\in V\) and \(s\in\R\), write \(\operatorname{VGe}_w(\Lambda,y;s)\) if there are \(a\in K\) and \(z\in\Lambda_{(w)}\) with
\[
 y=az,\qquad a=0\ \text{or}\ v_w(a)\ge e_ws.
\]

Let \(C_u,C^\mu_u\) be the reference and twisted column maps. For \(x\in V_{u,0}\), the assertion \(\operatorname{ColGe}_{u,\mu,w}(x;s)\) at \(\mu=\mu_0\) requires \(y\in V_{u,m}\) with
\[
 \iota_{u,m}(y)=C_u(\iota_{u,0}(x)),\qquad
 \operatorname{VGe}_w(\Lambda_{u,m},y;s).
\]
For \(\mu\ne\mu_0\), it requires \(y\in V^\mu_{u,m}\) with
\[
 \iota^\mu_{u,m}(y)=C^\mu_u\bigl(\iota^\mu_{u,0}(e^\mu_0x)\bigr),\qquad
 \operatorname{VGe}_w(\Lambda^\mu_{u,m},y;s).
\]
Here \(\iota^\mu_{u,0}(e^\mu_0x)=D.e0Tw(u,\mu)(\iota_{u,0}(x))\).
\end{definition}
\leanref{Formalism.\allowbreak{}MainNotions.\allowbreak{}WeilJacquetShalikaLaw.\allowbreak{}VGe}{Formalism.\allowbreak{}SharpRestriction.\allowbreak{}GrowthBounds.\allowbreak{}WeilJacquetShalikaLaw}{19--23}
\leanref{Formalism.\allowbreak{}MainNotions.\allowbreak{}WeilJacquetShalikaLaw.\allowbreak{}ColGe0}{Formalism.\allowbreak{}SharpRestriction.\allowbreak{}GrowthBounds.\allowbreak{}WeilJacquetShalikaLaw}{24--26}
\leanref{Formalism.\allowbreak{}MainNotions.\allowbreak{}WeilJacquetShalikaLaw.\allowbreak{}ColGeTw}{Formalism.\allowbreak{}SharpRestriction.\allowbreak{}GrowthBounds.\allowbreak{}WeilJacquetShalikaLaw}{28--31}
\leanref{Formalism.\allowbreak{}MainNotions.\allowbreak{}WeilJacquetShalikaLaw.\allowbreak{}ColGe}{Formalism.\allowbreak{}SharpRestriction.\allowbreak{}GrowthBounds.\allowbreak{}WeilJacquetShalikaLaw}{33--35}

A permitted index \(h_g:\Lean{GIdx}(D.\mathrm{host},T_u,p)\) requires an inert-class prime \(p\ge5\) outside the host level primes, with \(p\nmid T_u\). At every host-field place above \(p\) the second Frobenius component has even order. Write \(\mathcal O,\mathcal K,M\) for the model's coefficient DVR, fraction field and integral input module.

Let \(R_p\) be the representatives of \ref{B3.46} and \(W_p\subseteq R_p\) the selected set of \ref{B3.51}. With top weights \(a_L,b_L\), put
\[
 A_{t,j}=\max\{g_{0,t}(j),0\},\qquad
 H_{t,j}=\max\{g_{U,t}(j)-\operatorname{Pr}_t(j),A_{t,j}\}
 \quad(t\in R_p,\ 1\le j\le m).
\]
In a permitted model let \(\iota(t)\) be the local position of \(t\). For \(U\subseteq R_p\) put \(Y_U=\prod_{t\in U}y_{\iota(t)}^{\,m}\) in the model's order. Set
\[
 S_t=\left\lceil\sum_{j=1}^m(H_{t,j}-A_{t,j})\right\rceil,\qquad
 n(W')=\sum_{t\in W_p\setminus W'}S_t\quad(W'\subseteq W_p).
\]

\begin{definition}[The lattice supporting the enhanced estimate]\label{B4.07}
Let \(V^{\mathrm{ram}}\) be the prescribed ramified submodule in a permitted model. Define
\[
 N_P=
 \operatorname{span}_{\mathcal O}
 \left\{Y_{W'}(p^{n(W')}F):W'\subseteq W_p,\ F\in M\right\}
 \cap p^{\max\{0,\lceil\tau_p\rceil\}}M\cap V^{\mathrm{ram}},
 \tag{4.6}
\]
where \(\tau_p=\tau_P(p)\) is the presentation budget of \ref{B3.52}.

Let \(c_{u,w,h_g}\) compare \(V_{u,0}\) with the initial local function space, and set \(S^{\mathrm{ram}}_{u,w}=(\Lambda_{u,0})_{(w)}\cap V^{\mathrm{ram}}_{u,0}\). The assertion \(\operatorname{InNP}_{u,w}(x)\) holds if \(c_{u,w,h_g}(x)\in N_P\) for some permitted \(h_g\), or if \(W_p=\varnothing\) and \(x\in S^{\mathrm{ram}}_{u,w}\).
\end{definition}
\leanref{Formalism.\allowbreak{}MainNotions.\allowbreak{}G.\allowbreak{}NP}{Formalism.\allowbreak{}SharpRestriction.\allowbreak{}GrowthBounds.\allowbreak{}IsResidueClassZero}{464--473}
\leanref{Formalism.\allowbreak{}MainNotions.\allowbreak{}WeilJacquetShalikaLaw.\allowbreak{}InNP}{Formalism.\allowbreak{}SharpRestriction.\allowbreak{}GrowthBounds.\allowbreak{}WeilJacquetShalikaLaw}{60--64}

\begin{proposition}[Colength of the enhanced lattice]\label{B4.07a}
In a permitted model write \(R=M\cap V^{\mathrm{ram}}\), \(e=v_{\mathcal O}(p)\), \(d_M=\operatorname{rk}_{\mathcal O}M\), and \(W=W_p\). Assume finite level classes with trivial stabilisers and commuting neighbours that preserve the naive and integral lattices. Require
\[
 p^{X.\mathrm{ep}(p)}M_{\mathrm{na}}\subseteq M\subseteq M_{\mathrm{na}},
 \qquad R\text{ saturated in }M .
\]
Assume \(W\) maps to distinct local positions. At each selected position \(v\), require the test of \ref{B3.52a} with \(t=\lceil t_v\rceil_+\), \(D_0=D_p\) and \(\kappa=\kappa_v\), at rank \(r_G\) for \(p<p_0\) and rank one for \(p\ge p_0\). If \(m\ge1\) and \(m\sum_{v\in W}\kappa_v\le1/2\), then
\[
 N_P\subseteq p^{\tau_p}R,\qquad
 \operatorname{length}_{\mathcal O}(R/N_P)
 \le e d_M(\tau_p+\mathrm{io}_p).
\]
\end{proposition}
\begin{proof}
For \(W=\varnothing\), both budgets vanish and \(N_P=R\). Otherwise put \(s=m\sum_{v\in W}\kappa_v=\Sigma_p\), \(S=\max_{v\in W}S_v\), and \(q_v=\lceil t_v\rceil_+\). Let \(I(A)\) be the span in (4.6) with \(A\) in place of \(M\). The empty monomial gives \(p^{\sum_v S_v}A\subseteq I(A)\).

At each class representative, the naive lattice is the tensor product of the base lattices \(H_\Lambda\) of \ref{B3.52a}. At a local factor, the ball lattice is the source test lattice \(N_{\rm src}\), and the dual lattice is
\[
 N_{\rm tgt}^{\circ}
 =\{x\in V:\phi(x)\in\mathcal O\ \text{for every }\phi\in N_{\rm tgt}\},
\]
where \(N_{\rm tgt}\) is the target test lattice in \(V^\vee\) and \(V\) is the local tensor representation. The ball contains the base, which contains the dual. Multiplication by \(\omega^{q_v}\), where \(\omega^2=p\), sends the ball into the base and the base into the dual. The neighbour sends the ball into the dual. For nonempty \(B\subseteq W\), let \(L_B\) use dual factors at \(B\) and base factors elsewhere. Let \(L_{B,v}^{\mathrm{ball}}\) replace the factor at \(v\in B\) by the ball.

Fix a selected factor and put \(q=p^{D_p}\). The normalised edge vectors \(p^{A_{r'}}F_i\) of \ref{B3.52a}, with \(r'\) the test rank, belong to \(N_{\rm tgt}^{\circ}\) by the integral pairing of the two edge lattices. Write \(a_jb_j\) for a Smith basis of their span in a basis \((b_j)\) of \(N_{\rm tgt}^{\circ}\). The rank-\(h_0\) test gives
\[
 C\operatorname{diag}(a_j)X=qI_{h_0},
\]
where the rows of \(C\) express the tested combinations of edge vectors in the Smith basis, and \(X_{jk}=(U_0e_k)(b_j)\) for the standard vectors \(e_k\) of \(\mathcal O^{h_0}\). Both matrices are integral. If \(a_j\nmid q\), then \(a_j/q\) belongs to the maximal ideal. Move these terms to the right to obtain \(q(I_{h_0}-Y)\) with \(Y\) zero modulo that ideal. This matrix has rank \(h_0\), so at least \(h_0\ge(1-\kappa_v)\dim V\) of the \(a_j\) divide \(q\). At each such coordinate, \(qb_j\) is an integral combination of the normalised edge vectors.

For a tensor with such a coordinate at \(v\in B\), the test realises \(q\) times the tensor by a neighbour image modulo \(p^{D_p+1}L_B\), by the off-edge divisibility. At a class representative the fraction of basis tensors bad at every \(v\in B\) is at most \(\prod_{v\in B}\kappa_v\). Trivial stabilisers extend these tensors to a set \(C_B\subset L_B\) of at most \(d_M\prod_{v\in B}\kappa_v\) exceptional vectors. The sandwich identifies the naive and integral ranks. In the quotient by the neighbour images and \(\operatorname{span}_{\mathcal O}C_B\), the image of \(p^{D_p}L_B\) equals its multiple by \(p\). Nakayama's lemma gives
\[
 p^{D_p+1}L_B\subseteq
 \sum_{v\in B} y_{\iota(v)} L_{B,v}^{\mathrm{ball}}
 +\operatorname{span}_{\mathcal O} C_B .
\]

Order \(W\) and start with its first \(\beta+1\) positions active, \(0\le\beta<|W|\). In each neighbour term advance the chosen exponent by one. Below \(m\), return the ball input to the dual at cost \(\omega^{2q_v}\). At \(m\), return it to the base at cost \(\omega^{q_v}\) and activate the next unused position. Completed positions keep exponent \(m\). If all positions have been used when one reaches \(m\), omit the other \(\beta\) active positions at cost at most \(p^{\beta S}\). Their partial powers preserve \(M_{\mathrm{na}}\), so they can be included in the input of the resulting monomial in \(I(M_{\mathrm{na}})\).

The exponent sum increases at every continuation. The cost per position on each path, including entry and exit, is at most \(2q_v+2(m-1)q_v+2m(D_p+1)\) powers of \(\omega\). Thus the uniform cost is at most \(\tau_p-X.\mathrm{ep}(p)+\beta S\) powers of \(p\). An exceptional term is a neighbour product applied to \(C_B\). The completed positions are those preceding the largest member of \(B\) outside \(B\). By commutation the exponent lists for a given \(B\) number at most \(m^{|B|}\). Consequently \(p^{\tau_p-X.\mathrm{ep}(p)+\beta S}M_{\mathrm{na}}\) lies in \(I(M_{\mathrm{na}})\) plus the span of at most
\[
 d_M m^{\beta+1}
 \sum_{\substack{B\subseteq W\\|B|=\beta+1}}\prod_{v\in B}\kappa_v
 \le d_M\frac{s^{\beta+1}}{(\beta+1)!}
\]
vectors. The last inequality follows by expanding \((\sum_v\kappa_v)^{\beta+1}\). Multiply by \(p^{X.\mathrm{ep}(p)}\) to send the span into \(M\) and \(I(M_{\mathrm{na}})\) into \(I(M)\).

Let \(a_\ell\) be the Smith exponents of \(I(M)\subseteq M\) in a uniformiser of \(\mathcal O\). The image of \(p^{\tau_p+\beta S}M\) in \(M/I(M)\) lies in a submodule generated by the exceptional vectors. Over a DVR a submodule of a module with \(q'\) generators has at most \(q'\) generators, by taking the inverse image in \(\mathcal O^{q'}\). The image has one nonzero cyclic factor per \(a_\ell>e(\tau_p+\beta S)\), so their number is bounded by the last display. For \(\beta\ge|W|\) the empty monomial gives \(a_\ell\le e\sum_v S_v\le e\beta S\).

Intersecting with \(p^{\tau_p}M\) replaces \(a_\ell\) by \(\max(a_\ell,e\tau_p)\). Sum the excess above \(e\tau_p\) in intervals of width \(eS\) to bound it by \(eSd_Ms/(1-s)\le2eSd_Ms\). For \(S=0\) the empty monomial gives zero excess. The quotient \(R/N_P\) injects into \(M/(I(M)\cap p^{\tau_p}M)\). Saturation gives \(p^{\tau_p}M\cap R=p^{\tau_p}R\), and \(\mathrm{io}_p=2Ss\) by \ref{B3.52}.
\end{proof}
\leanref{Formalism.\allowbreak{}MainNotions.\allowbreak{}G.\allowbreak{}LatticeSystems.\allowbreak{}walk\_\allowbreak{}allBut\_\allowbreak{}bound}{Formalism.\allowbreak{}SharpRestriction.\allowbreak{}GrowthBounds.\allowbreak{}LatticeSystems}{217--313}
\leanref{Formalism.\allowbreak{}MainNotions.\allowbreak{}FormalGroupSystem.\allowbreak{}Pf.\allowbreak{}stmt6G}{Formalism.\allowbreak{}SharpRestriction.\allowbreak{}GrowthBounds.\allowbreak{}MultisetValuation.\allowbreak{}StabilizerCount}{384--491}

\begin{definition}[The two column budgets]\label{B4.08}
For \(\sigma\in\R\), put
\[
 \begin{split}
 \operatorname{lowS}_{u,w}(\sigma)
 &=\sum_{j=1}^m\bigl(F_j(w)-E^{\mathrm{ev}}_{u,j}(w)\bigr)
 +\sigma\delta_{u,w},\\
 \operatorname{totS}_{u,w}(\sigma)
 &=\operatorname{lowS}_{u,w}(\sigma)+\sum_{j=1}^m\operatorname{gainS}_{p,j}.
 \end{split} \tag{4.7}
\]
Let \(\operatorname{Exc}(p)\) mean \(p\mid\operatorname{condN}\,g_0g_1\), and \(\operatorname{IsCs}(p)\) mean that \(p\) is inert-class with \(h_P\ne0\) of odd additive order. The served gain is
\[
 \operatorname{gainS}_{p,j}=
 \begin{cases}
 \operatorname{gainP}_{p,j},
 &(\operatorname{inertClass}(p)\text{ or }\operatorname{signal}(p)),
 \ \neg\operatorname{Exc}(p),\ \neg\operatorname{IsCs}(p),\\
 0,&\text{otherwise}.
 \end{cases}
 \tag{4.8}
\]
The gain \(\operatorname{gainP}_{p,j}\) is zero if \(p>X''\) or if \(p\) is inert-class with \(h_P=0\). Otherwise it equals
\[
 \sum_{t\in W_p}H_{t,j}+\sum_{t\in R_p\setminus W_p}A_{t,j}.
 \tag{4.9}
\]
\end{definition}
\leanref{Formalism.\allowbreak{}MainNotions.\allowbreak{}WeilJacquetShalikaLaw.\allowbreak{}lowS}{Formalism.\allowbreak{}SharpRestriction.\allowbreak{}GrowthBounds.\allowbreak{}WeilJacquetShalikaLaw}{73--75}
\leanref{Formalism.\allowbreak{}MainNotions.\allowbreak{}WeilJacquetShalikaLaw.\allowbreak{}totS}{Formalism.\allowbreak{}SharpRestriction.\allowbreak{}GrowthBounds.\allowbreak{}WeilJacquetShalikaLaw}{76--78}
\leanref{Formalism.\allowbreak{}MainNotions.\allowbreak{}WeilJacquetShalikaLaw.\allowbreak{}gainS}{Formalism.\allowbreak{}SharpRestriction.\allowbreak{}GrowthBounds.\allowbreak{}WeilJacquetShalikaLaw}{49--51}
\leanref{Formalism.\allowbreak{}MainNotions.\allowbreak{}G.\allowbreak{}gainP}{Formalism.\allowbreak{}SharpRestriction.\allowbreak{}GrowthBounds.\allowbreak{}IsResidueClassZero}{355--358}

\section{Hecke images and their integral span}

Let \(K_\mu^{[m]}:\mathrm{Fm}_\mu(0)\to\mathrm{Fm}_\mu(m)\) be the local chain operator, with initial space \(G.\mathrm{Fn}\) and integral input \(M\). For a uniformiser \(\varpi\) of \(\mathcal O\), the bound \(\operatorname{vKge}_\mu(F;s)\) means
\[
 K_\mu^{[m]}F\in
 \varpi^{\lceil v_{\mathcal O}(p)s\rceil}\Lambda^{\mathrm{loc}}_{\mu,m},
 \tag{4.10}
\]
where \(\Lambda^{\mathrm{loc}}_{\mu,m}\) is the terminal integral module.

\begin{theorem}[The whole-chain estimate on Hecke images]\label{B4.09}
Fix local data \(Y\) and \(W\in\Lean{TwLoc}(D,Y)\). Assume \Lean{ErrDefs}. For each used \(\mu\), each tying datum \(h_T\) for \(W.Yt(\mu)\), and each permitted \(u,w,h_g\) with \(p\notin Q'(\mu)\), assume a matching formal-group system and kernel datum admissible as in \ref{B3.45}. Let \(d_G(i)\) be its stabiliser losses and \(\iota:R_p\to\mathcal V\) the injection into local positions. Require \(f_{v,j}\in\R\) with \(\operatorname{LowAt}(j-1,v,f_{v,j})\) as in \ref{B3.12} for \(v\notin\iota(R_p)\), \(1\le j\le m\), and
\[
 \operatorname{lowS}_{u,w}(\sigma)
 \le \sum_{j=1}^m\sum_{v\notin\iota(R_p)} f_{v,j}
 -\sum_{i=0}^{m-1}d_G(i).
 \tag{4.11}
\]
Write \(a_\ast=L.\mathrm{sc}.a\). For each tying datum for \(Y\) and each permitted \(u,w,h_g\) with \(1\le\log_p(2a_\ast)\), assume the comparisons of \ref{B3.53} for \(t\in R_p\), \(1\le j\le m\),
\[
 g_{U,t}(j)\le g^{\mathrm{loc}}_{B,\iota(t)}(j),\qquad
 g_{0,t}(j)\le g^{\mathrm{loc}}_{A,\iota(t)}(j),\qquad
 P^{\mathrm{loc}}_{\iota(t)}(j)\le\operatorname{Pr}_t(j).
 \tag{4.12}
\]
At every coefficient-field place with \(\log_p(2a_\ast)<1\), assume instead for the same \(t,j\) that
\[
 g_{U,t}(j)-\operatorname{Pr}_t(j)\le0,\qquad g_{0,t}(j)\le0.
 \tag{4.13}
\]
Then for each \(\mu,h_T,u,w,h_g\) as above, every \(V'\subseteq V\subseteq R_p\) and \(F\in M\) satisfy
\[
 \operatorname{vKge}_\mu\left(Y_{V'}F;\
 \operatorname{lowS}_{u,w}(\sigma)
 +\sum_{j=1}^m\left(
 \sum_{t\in V'}H_{t,j}
 +\sum_{t\in V\setminus V'}A_{t,j}
 +\sum_{t\in R_p\setminus V}A_{t,j}\right)\right).
 \tag{4.14}
\]
\end{theorem}
\leanref{Formalism.\allowbreak{}MainNotions.\allowbreak{}G.\allowbreak{}Joints.\allowbreak{}unit\_\allowbreak{}bound\_\allowbreak{}of\_\allowbreak{}sum\_\allowbreak{}over\_\allowbreak{}primes}{Formalism.\allowbreak{}SharpRestriction.\allowbreak{}GrowthBounds.\allowbreak{}JointsTw}{243--271}
\leanref{Formalism.\allowbreak{}MainNotions.\allowbreak{}WeilJacquetShalikaLaw.\allowbreak{}UAt}{Formalism.\allowbreak{}SharpRestriction.\allowbreak{}GrowthBounds.\allowbreak{}WeilJacquetShalikaLaw}{161--171}

\begin{proof}
Use \ref{B3.45} for the ordinary coefficient bound. Use the chain bound \ref{B3.42} and the zero coefficient bound \Lean{adm11} as well. The twist shares the filling and representative map of \(Y\), so the tie and (4.12) carry over. Apply \ref{B3.55} with \(\mathcal O\hookrightarrow\mathcal K\) and \(\mathrm{low}=\operatorname{lowS}_{u,w}(\sigma)\) by (4.11)--(4.13). Split \(R_p\setminus V'\) at \(V\).
\end{proof}

\begin{lemma}[Extension from generators to \(N_P\)]\label{B4.10}
Assume \ref{B4.09} and the error identities. For every used \(\mu\), every tying datum for \(W.Yt(\mu)\), and every permitted \(u,w,h_g\) with \(p\notin Q'(\mu)\), every \(F\in N_P\) satisfies
\[
 \operatorname{vKge}_\mu\left(F;\
 \operatorname{lowS}_{u,w}(\sigma)+\sum_{j=1}^m\operatorname{gainP}_{p,j}\right).
 \tag{4.17}
\]
For fixed \(u,w,h_g\), \(N_P\) is independent of \(\mu\).
\end{lemma}
\leanref{Formalism.\allowbreak{}MainNotions.\allowbreak{}WeilJacquetShalikaLaw.\allowbreak{}Proof.\allowbreak{}first\_\allowbreak{}in\_\allowbreak{}local\_\allowbreak{}models}{Formalism.\allowbreak{}SharpRestriction.\allowbreak{}GrowthBounds.\allowbreak{}AdmissibleIntegers.\allowbreak{}StatementsOfRemainingBound}{111--152}

\begin{proof}
For \(t\in W_p\), the definition gives
\[
 \sum_{j=1}^m H_{t,j}\le\sum_{j=1}^m A_{t,j}+S_t.
 \tag{4.18}
\]
On a generator \(Y_{W'}(p^{n(W')}F_0)\), apply (4.14) with \(V=W_p\), \(V'=W'\) and \(F=F_0\). The scalar adds \(n(W')\) to the exponent. By (4.18) for the omitted representatives,
\[
 \begin{split}
 \sum_{j=1}^m\operatorname{gainP}_{p,j}\le
 &\sum_{j=1}^m\left(\sum_{t\in W'}H_{t,j}
 +\sum_{t\in W_p\setminus W'}A_{t,j}
 +\sum_{t\in R_p\setminus W_p}A_{t,j}\right)+n(W').
 \end{split} \tag{4.19}
\]
This also holds when \(\operatorname{gainP}\) vanishes, since the right side is nonnegative. Linearity and (4.10) give (4.17) on the span. By \Lean{NP\_\allowbreak{}tw}, the shared filling and representative map together with \(V^{\mathrm{ram}}\) give the same (4.6) after twisting.
\end{proof}

\begin{lemma}[From kernel entries to a lift]\label{B4.10a}
Fix \(u\), a used twist \(\mu\) and a level-free place \(w\), so \(p\notin\{T_u,g_0,g_1\}\) and \(p\notin Q'(\mu)\). For \(0\le i<m\), let \(\nu_z=\nu_{u,i,z}\) be the entry diagrams and \(\Theta\) the series for a source and target class pair. A choice \(I\) assigns to every column of height \(h\) in \(\nu_z\) an increasing injection \(\{1,\dots,h\}\hookrightarrow\{1,\dots,N_{i+1}\}\). Let \(G_{\nu,I}\) be the product of the minors of \(X_zZ_z\) on the first \(h\) rows and the columns given by \(I\). The bound \(\Lean{EntGe}(B)\) requires every monomial coefficient of
\[
 \left.G_{\nu,I}(X,\partial_Z)\Theta(Z)\right|_{Z=0}
\]
to have norm at most \(p^{-B}\) for every class pair and every \(I\).

Suppose the diagrams and series agree with the kernel datum. Assume the unit and rationality conditions of \ref{B4.16} as well as its \(\mathrm{Fi}\)-preservation condition. Assume also the stabiliser denominator condition. Take \(\sigma=0\), or \(\sigma=1\) with trivial initial stabilisers. Then \(\Lean{EntGe}(B)\) gives the lift bound
\[
 B-\bigl(1-\sigma\mathbf1_{\{i=0\}}\bigr)\delta_{u,w}.
 \tag{4.20}
\]
Thus each integral input has a rational output representative of this bound in the target lattice, in the sense of \ref{B4.06}. At \(\mu_0\) the bound holds for the reference lift.
\end{lemma}
\begin{proof}
By Cauchy--Binet the entry polynomial belongs to the span of products of source minors. For the prescribed shapes, the straightened tableaux have distinct diagonal leading monomials of coefficient one. Successive elimination of those monomials gives an integral projection onto the span. Divide a nonzero entry polynomial by a coefficient of largest norm. Apply the projection over the valuation ring and multiply back. The resulting expansion in products of source minors has coefficient norms at most \(p^{-B}\).

For integral input each class-pair target coordinate \(c_I\) evaluates this expansion on integral source coordinates, up to a unit, so \(|c_I|_p\le p^{-B}\). Choose a \(K\)-rational representative of the class-pair output by rationality. The kernel datum gives \(b_I\in K\) with \(|c_I|_p=|b_I|_w\), where \(|b|_w=p^{-v_w(b)/e_w}\) for \(b\ne0\). If \(c_I\ne0\), then \(v_w(b_I)\ge\lceil e_wB\rceil\). Choose \(a\in K^\times\) of valuation \(\lceil e_wB\rceil\). Every class-pair coordinate divided by the image of \(a\) has norm at most one.

Let \(n>0\) be a common multiple of the source stabiliser orders and \(b_0=v_p(n)\). The lift weights the class-pair maps by \(1/|\operatorname{Stab}_i(c)|\), so \(p^{b_0}/a\) times the lift has integral coordinates in the rational class representatives. It lies in \(\mathrm{Fi}_{i+1}\) by hypothesis and hence has a representative in the localised target lattice. Scaling back gives \(B-b_0\). Use \(n=D\Gamma_u\), or \(n=1\) when \(i=0\) and \(\sigma=1\). The lattice isomorphisms at the reference twist intertwine the lifts and transfer the bound.
\end{proof}
\leanref{Formalism.\allowbreak{}MainNotions.\allowbreak{}DIV.\allowbreak{}Lemmas.\allowbreak{}theta\_\allowbreak{}isometry\_\allowbreak{}facts\_\allowbreak{}hold}{Formalism.\allowbreak{}SharpRestriction.\allowbreak{}Divisibility.\allowbreak{}ChainOperators.\allowbreak{}GenericFramesLetters}{10--13}
\leanref{Formalism.\allowbreak{}MainNotions.\allowbreak{}OrbitPlaces.\allowbreak{}kform\_\allowbreak{}of\_\allowbreak{}ker}{Formalism.\allowbreak{}SharpRestriction.\allowbreak{}Divisibility.\allowbreak{}OrbitPlaces}{1621--1657}

\section{Orbit bounds and ordinary places}

Write \(\mathscr E\) for the orbit entry data. In the cyclic expansion the positions split into orbits of length \(n_P\), each with the expansion of \ref{B4.01b}. The full phase is the sum of the local phases and there is a single integral multiplier. For \(h_P=0\) the blocks are single positions with the one-letter phases of \ref{B4.01b}.

For a nonzero odd-order inert prime let \(P\) denote orbit shift, and write \(\ell_{z,ab}=\sum_{h\ge0} Z_{P^{-2h}z,ab}^{p^{2h}}/p^h\) for the letter logarithm. The conjugate theta product expansion is
\[
 \Theta(\xi\ell)=U\exp\!\left(
 \sum_{r\in R}\frac{\epsilon_r}{2}
 \sum_{a,b}\beta_r\bigl((\ell_{P^sr,ab})_{s\in\Z/(2n_P)}\bigr)\right).
\]
Here \(R\) is a finite set of positions, \(|\xi|_p=1\), \(\epsilon_r=\pm1\), and \(U\) has integral coefficients. The sum over \(a,b\) runs over the \(N_iN_{i+1}\) matrix entries. The \(\beta_r\) satisfy the letter conditions at period \(2n_P\) with a common unramified coefficient subfield and a common primitive \((p^{2n_P}-1)\)-st root of unity. Their arguments use the projection of the \(2n_P\) letters onto the orbit of \(r\). The selected kernel has the shape and expansions of \ref{B8.56}.

\begin{lemma}[Four consequences of the orbit estimates]\label{B4.11}
Suppose
\[
 \begin{gathered}
 \Lean{OrbitPlaces.\allowbreak{}admitted}(D,L,X,\mathbf c,D_x,T,\mathscr E,c_e,\sigma,C_6,C),\\
 \Lean{OrbitPlaces.\allowbreak{}Core}(D,L,X,\mathbf c,D_x,\mathscr E,c_e,C),
 \end{gathered}
\]
where \(T\) is the datum of twisted forms. Assume \Lean{FDef} and \Lean{ErrDefs} in (i)--(iii), which hold for each used \(\mu\) with \(p\notin Q'(\mu)\). Put \(g_{0,S}(p,j)=0\) if \(\operatorname{IsCs}(p)\), and \(g_{0,S}(p,j)=\sum_{t\in R_p}A_{t,j}\) otherwise.

\begin{enumerate}
\item[(i)] If \(p\ne T_u\), \(\neg\operatorname{Exc}(p)\), \(p\le X''\), and \(p\) is signal or inert-class with \(h_P\ne0\), every \(x\in S^{\mathrm{ram}}_{u,w}\) satisfies
\[
 \operatorname{ColGe}_{u,\mu,w}
 \left(x;\operatorname{lowS}_{u,w}(\sigma)+\sum_{j=1}^m g_{0,S}(p,j)\right).
 \tag{4.21}
\]
\item[(ii)] If \(p\ne T_u\), \(p\) is inert-class with \(h_P=0\), \(\neg\operatorname{Exc}(p)\), and \(p\le X''\), every \(x\in S^{\mathrm{ram}}_{u,w}\) has column bound
\[
 \sum_{j=1}^m\bigl(F_j(w)-\operatorname{EatC}_j(p)-\delta_{u,w}\bigr)
 +\sigma\delta_{u,w}.
 \tag{4.22}
\]
\item[(iii)] If \(p\ne T_u\) and \(p\) has class \(\mathrm{ramD}\), every \(x\in S^{\mathrm{ram}}_{u,w}\) has column bound
\[
 \sum_{j=1}^m\bigl(F_j(w)-\operatorname{EatD}_j-\delta_{u,w}\bigr)
 +\sigma\delta_{u,w}.
 \tag{4.23}
\]
\item[(iv)] Under the stabiliser denominator condition, if \(p\notin\{T_u,d,g_0,g_1\}\), each step \(0\le i<m\) has bound
\[
 -\bigl(1-\sigma\mathbf1_{\{i=0\}}\bigr)\delta_{u,w}
 \tag{4.24}
\]
on every vector of its integral lattice. This holds throughout the stated range for the reference chain and outside \(Q'(\mu)\) for each nonreference used twist.
\end{enumerate}
\end{lemma}
\leanref{Formalism.\allowbreak{}MainNotions.\allowbreak{}OrbitPlaces.\allowbreak{}lift\_\allowbreak{}bounds\_\allowbreak{}of\_\allowbreak{}essential}{Formalism.\allowbreak{}SharpRestriction.\allowbreak{}Divisibility.\allowbreak{}OrbitPlaces}{1340--1344}
\leanref{Formalism.\allowbreak{}MainNotions.\allowbreak{}J.\allowbreak{}four\_\allowbreak{}orbit\_\allowbreak{}conditions\_\allowbreak{}of\_\allowbreak{}proof}{Formalism.\allowbreak{}SharpRestriction.\allowbreak{}AdmissibleParameterJoins.\allowbreak{}GandSab}{11--17}

\begin{proof}
The even \(k\) and odd \(t_{\rm ch}=D.R.t\) give even \(N_i=k+i(t_{\rm ch}+1)\). In (i), at a signal or even-order inert prime use respectively \(\mathcal R_p=\{0\}\times R_p\) or \(\{0,1\}\times R_p\). Put \(g_{0,(0,t),j}=A_{t,j}\) and \(g_{0,(1,t),j}=0\). Use the no-operator bound \(\beta_{r,j}=g_{0,r,j}\) on kind \(A\), and clause (2) of \ref{B4.11a} on kind \(BC\), with \(\beta_{r,j}=g_{0,r,j}+F_{\mathrm{Orb},r,j}-\mathrm{err}_{r,j}\). Clause (1) adds these bounds.

Here \(F_j(w)=\sum_{r\in\mathcal R_p}F_{\mathrm{Orb},r,j}\) and \(\operatorname{EatX}_j(p)=0\). The kind-\(A\) representative sum is \(n_P^{-1}\sum_{z\in\mathrm{kindA}(p)}F_{\mathrm{Orb},j}(p,z)\). Subtract the positive part of the kind-\(A\) sum and the kind-\(BC\) error from the gain plus the full base term. The entry exponent is at least
\[
 \sum_{r\in\mathcal R_p}\beta_{r,j}
 \ge \sum_{t\in R_p}A_{t,j}+F_j(w)-\operatorname{EatI}_j(p).
 \tag{4.25}
\]
Let \(BC\) be the kind-\(BC\) positions, \(TT_{BC}\) their top positions, \(\gamma(p)=p/(p^2-1)\), and \(w_j(z)\) the position weights. Summing \(\mathrm{err}_{r,j}\) over the kind-\(BC\) orbits gives
\[
 \begin{split}
 &\gamma(p)\,\#TT_{BC}(b+\eta+1+3\mathrm{Sc})
 +\frac{c_{e,6}}p\sum_{z\in BC}w_j(z)\\
 &\quad+\#BC\,N_{j-1}N_j\bigl(\log_p(4n_PN_ma_\ast)+1\bigr)\\
 &\quad+c_{\mathrm{dl},1}\#BC\,N_m^2\log_p(2a_\ast)
 +c_{\mathrm{dl},2}\#TT_{BC}\log_p(2a_\ast)^2.
 \end{split} \tag{4.26}
\]

On the odd-order branch of (i), \(p\ge5\) by orbit admissibility. Since \(N_iN_{i+1}\) is even, its odd-order clause supplies the conjugate theta product expansion for every class pair. Clause (5) of \ref{B4.11a} gives entry bound \(L_{\rm odd}\), and the large-error condition in \Lean{ErrDefs} gives \(F_{i+1}(w)-\operatorname{EatI}_{i+1}(p)\le L_{\rm odd}\). Here \(g_{0,S}(p,i+1)=0\). In both branches the prime exclusions give a level-free place and \(E^{\mathrm{ev}}_{u,j}(w)=\delta_{u,w}+\operatorname{EatI}_j(p)\). Apply (4.20) to obtain the step exponent
\[
 g_{0,S}(p,i+1)+F_{i+1}(w)-E^{\mathrm{ev}}_{u,i+1}(w)
 +\sigma\mathbf1_{\{i=0\}}\delta_{u,w}.
 \tag{4.27}
\]

In (ii), admissibility gives \(p\ge5\) and the one-position expansion. The one-letter estimate of \ref{B4.01b} and clause (1) of \ref{B4.11a} give the entry bound \(B_C\) of clause (4). Apply (4.20) with \(F_{i+1}(w)-\operatorname{EatC}_{i+1}(p)\).

In (iii), \(p=d\) and \(g_0,g_1\ne d\). By the discriminant error definition,
\[
 F_j(w)-\operatorname{EatD}_j
 =-\tfrac12\max_{u'}\sum_t\mathrm{ct}_{u',j,t}
 \le-\tfrac12\sum_t\mathrm{ct}_{u,j,t}.
\]
Use clause (3) and (4.20). In (i)--(iii), \ref{B4.12} adds the step exponents with the correction \(\sigma\delta_{u,w}\) at \(i=0\).

In (iv), even rank product and the jet-integrality clause give entry bound zero. Apply (4.20). For the reference step, \(Q'(\mu_0)=\varnothing\).
\end{proof}

\begin{proposition}[The orbit entry estimates]\label{B4.11a}
Assume the orbit-letter estimates of \ref{B4.01b} with weight-loss coefficient \(C_6\) and digit-loss coefficient \(C\). Assume also the weight-shape inequalities of \ref{B8.56} and the quadratic expansion (8.56.2) at level-free discriminant places. Suppose
\[
 C\ge0,\qquad C_6\le c_{e,6},\qquad c_{\mathrm{dl},1},c_{\mathrm{dl},2}\ge0,
 \qquad 3C\le\min(1+c_{\mathrm{dl},1},c_{\rm dig}).
\]
Then \(\Lean{OrbitPlaces.\allowbreak{}Core}\) consists of the following bounds and restrictions, for each column \(u\). In clauses (2)--(5), \(0\le i<m\) and \(j=i+1\).

\begin{enumerate}
\item For every \(i\in\N\) and prime \(p\) satisfying \Lean{OrbOK}, meaning \(p\ge5\) and \(n_P\ge2\) even, the cyclic expansion restricts to each orbit for every choice of representatives and admissible letter data. For every \(i\in\N\), the one-position expansion restricts to each position when \(p\ge5\) and \(n_P=1\). For \(i<m\), uniform local exponent bounds add to a full kernel bound in either case.
\item For \(p\le X''\) satisfying \Lean{OrbOK}, with \(p\) signal or inert-class with \(h_P\ne0\), each kind-\(BC\) representative \(r\) has exponent at least \(g_{0,r,j}+F_{\mathrm{Orb},r,j}-\mathrm{err}_{r,j}\) for every local expansion with the prescribed unit and sign as well as the admissible series. The error \(\mathrm{err}_{r,j}\) is (4.26) restricted to the orbit of \(r\).
\item At a level-free place of class \(\mathrm{ramD}\), every class pair has entry exponent at least \(-\frac12\sum_t\mathrm{ct}_{u,j,t}\).
\item At an inert prime \(5\le p\le X''\) with \(h_P=0\), one has \(F_{\rm inert}-\operatorname{EatC}_j(p)\le B_C\), where \(F_{\rm inert}\) is the base term of \ref{B4.02} at \(j,p\) and
\[
 B_C=\rho(p)\sum_z|\nu_{u,i,z}|
 -CfN_iN_{i+1}\bigl(\log_p(a_\ast^2)+1\bigr),\qquad
 \rho(p)=\frac1{p^2-1}.
\]
Here \(f\) counts positions and \(|\nu|\) counts boxes.
\item At a nonzero odd-order inert prime \(5\le p\le X''\), every kernel with the conjugate theta product expansion has exponent at least
\[
 L_{\rm odd}=\frac1{n_P}\sum_{z\in BC}F_{\mathrm{Orb},j}(p,z)
 -\left(\frac{k-1}{p^3}+\frac8p\right)\sum_{z\in BC}w_j(z)
 -3fN_m^2\log_p a_\ast .
\]
Under the large-error condition of \ref{B4.03}, \(F_{\rm inert}-\operatorname{EatI}_j(p)\le L_{\rm odd}\).
\end{enumerate}
\end{proposition}
\begin{proof}
For restriction, set the variables outside the block to zero. The discarded phases have zero constant term. For addition, expand \(U\) into monomials, each a product of monomials on the blocks. The inverse logarithmic substitution respects the blocks and has zero constant term. Each monomial times the product of exponentials pulls back to a product of local kernels, with integral monomial multipliers in logarithmic coordinates. Minor differential evaluation factors over the blocks. At each differential order the sum is finite, so the local exponents add.

On a kind-\(BC\) orbit use the refined cyclic estimate of \ref{B4.01b} at \(a=a_\ast^2\). The shape inequalities give \(2\log_p(a_\ast^2)+2\le N_i-1\) and the demand bounds of \ref{B4.01a} with \(W_s=(k-1)w_j(s)\). The base value is \(F_{\mathrm{Orb},r,j}\). At a signal prime an odd step from a top position changes sheets, so the first positive odd distance to a top position is infinite. The first-row difference in \ref{B4.01b} dominates the numerator of \ref{B3.48}, and the global partner cap dominates the local cap. Hence the sum of the marked terms \(J_s\) is at least \(g_{0,r,j}\). At inert kind-\(BC\) positions \(g_{0,r,j}=0\).

Write \(TT_r\) for the top positions on this orbit and \(l=\log_p a_\ast\ge1\). The shape bound for the reverse-partner maximum gives
\[
 \gamma(p)\sum_{s\in TT_r}M_s
 \le\gamma(p)\,\#TT_r\,(b+\eta+1+3\mathrm{Sc}).
\]
This is the top term of the orbit error. Its weight term bounds \(C_6\sum_s w_j(s)/p\). Both logarithms are at least \(l\). Since \(C(2l+1)\le(1+c_{\mathrm{dl},1})l\) and \(N_iN_{i+1}\le N_m^2\), the rank and digit terms together cover \(Cn_PN_iN_{i+1}(2l+1)\).

At \(h_P=0\), the shape inequalities give \(|\nu_z|\ge(k-1)w_j(z)\), since on top positions the first row restores at least \(2w_j(z)\) to the off-first-row demand. Hence \(\rho(p)\sum_z|\nu_z|\ge F_{\rm inert}\). The digit term of \(\operatorname{EatC}_j(p)\) covers \(CfN_iN_{i+1}(2l+1)\) by \(C(2l+1)\le c_{\rm dig}l\).

For odd order choose optimal nonnegative dual coordinates \(y_s\le1\) in \ref{B4.01} on each kind-\(BC\) orbit with demands \(W_s=(k-1)w_j(s)\). Set \(y'_s=\max(y_s,p^{-3})\). The retained distances obey \(p^e y_s+y_{s+e}\ge1\). An omitted positive odd \(e\) is at least three, so \(p^e y'_s+y'_{s+e}\ge1\) for every positive odd \(e\). The change costs at most \(p^{-3}\sum_s W_s\) in \(\sum_s W_s(\gamma(p)-y_s)\). Kind-\(BC\) orbits here have no top positions, and \(|\nu_s|\in[W_s,W_s+4w_j(s)]\). By the tight-capacity estimate of \ref{B4.01a}, \(\{s:y_s\ge2/p\}=\{s:y'_s\ge2/p\}\) has demand at most \(p^{-1}\sum_s W_s\). Outside this set, \(y'_s<2/p\), and inside it \(y'_s\le1\). The loss from replacing \(W_s\) by \(|\nu_s|\) is at most \(8\sum_s w_j(s)/p\) by the same estimate. On kind \(A\) the constant choice \(y'_s=1/(p+1)\) satisfies the odd-distance inequalities and \(\gamma(p)-y'_s\ge0\).

Give each variable at \(z\) weight \(\rho(p)+y'_z\) in addition to the valuation of its coefficient. In the \(2n_P\)-letter phase the factor \(\epsilon_r/2\) is a unit. In the coefficient calculation of \ref{B4.01b}, an exceptional pair of degrees \(p^\alpha,p^\beta\) with \(\alpha>\beta\) projects to an odd displacement \(\alpha-\beta\). Its weight increases by
\[
 p^\beta\bigl(p^{\alpha-\beta}y'_z+y'_{P^{\alpha-\beta}z}\bigr)\ge1.
\]
Thus every phase monomial has weight at least \(1/(p-1)\). Exponentiation and the integral factor preserve nonnegative weight. The inverse logarithm also preserves it, since its linear term preserves the variable and its higher terms have weight at least \(\rho(p)+1\) at the base weights \(\rho(p)\). Legendre's factorial bound gives entry exponent at least
\[
 \sum_z|\nu_z|(\gamma(p)-y'_z)-fN_iN_{i+1}(2l+1)\ \ge\ L_{\rm odd}.
\]
Use orbit invariance for the kind-\(BC\) base values and \(2l+1\le3l\). For the comparison with \(F_{\rm inert}\), the positive-part term in \(\operatorname{EatI}_j(p)\) covers the kind-\(A\) average. The cubic term in \(\operatorname{EatX}_j(p)\) covers the change from \(y_s\) to \(y'_s\). The remaining weight and logarithmic deductions fit the terms with coefficient \(c_{e,6}\ge12\).

At \(p=d\), let \(\alpha=\nu_{(0,t)}\) and \(\lambda=\nu_{(1,t)}\), with first-row lengths \(a_1,b_1\). The phase in (8.56.2) pairs variables on opposite sheets with scalar \(-Su_t\), where \(|Su_t|_p\le p^{1/2}\). The factorial-weighted coefficient for paired exponents \(h,h\) is \(h!(-Su_t)^h\), and unequal exponents give zero. Thus the order-\(n_t\) pairing acts by \((-Su_t)^{n_t}\Delta_t^{n_t}/n_t!\), where
\[
 \Delta_t=\sum_{v,h}\partial_{(0,t),vh}\partial_{(1,t),vh}.
\]
The multinomial expansion pairs divided derivatives on one side with ordinary derivatives on the other. This makes the contraction integral after removal of the scalar. At most \(|\alpha|-a_1\) contractions use rows below the first on the \(\alpha\) side. A first-row contraction differentiates the other side in the direction of the first frame row. Each of the \(b_1\) minors is affine under this rank-one translation, since a term with two columns replaced by the same vector vanishes. At most \(\min(a_1,b_1)\) first-row contractions survive. The total is at most \(|\alpha|\) off the top and at most \(|\alpha|-a_1+\alpha_2+b_1\) on the top, where \(\alpha_2\) is the second-row length. These are the bounds \(\mathrm{ct}_{u,j,t}\). The residual jet has integral factorial-weighted coefficients. Each surviving contraction lowers the valuation by at most \(1/2\), so the entry exponent is at least \(-\tfrac12\sum_t\mathrm{ct}_{u,j,t}\).
\end{proof}
\leanref{PI.\allowbreak{}OrbBoundF}{Formalism.\allowbreak{}PeriodIntegrals.\allowbreak{}OrbitalBoundsConclusion.\allowbreak{}OrbBoundF}{3--5}
\leanref{UpperTheory.\allowbreak{}kvalGe\_\allowbreak{}theta\_\allowbreak{}of\_\allowbreak{}shape}{Development.\allowbreak{}AutomorphicArithmetic.\allowbreak{}OrbitLabellings.\allowbreak{}KvalGeThetaOfShape}{155--205}
\leanref{jet\_\allowbreak{}bound\_\allowbreak{}ramified\_\allowbreak{}prime}{Formalism.\allowbreak{}PeriodIntegrals.\allowbreak{}OrbitalBoundsConclusion.\allowbreak{}EntAtD}{11--14}
\leanref{Formalism.\allowbreak{}MainNotions.\allowbreak{}DIV.\allowbreak{}Lemmas.\allowbreak{}OrbitEntryBounds.\allowbreak{}orbit\_\allowbreak{}entry\_\allowbreak{}bounds\_\allowbreak{}hold}{Formalism.\allowbreak{}SharpRestriction.\allowbreak{}Divisibility.\allowbreak{}ChainOperators.\allowbreak{}OrbitEntryBounds}{49--53}

\begin{lemma}[Adding the bounds at the remaining primes]\label{B4.12}
Let \(n\in\N\) and \(b_i\in\R\) for \(0\le i<n\). In a fixed reference or twisted chain, let \(V_i,\Lambda_i,\iota_i\) be the spaces, lattices and embeddings, and let \(K^{[r]}\) compose the first \(r\) lifts. Suppose the lift at step \(i\) has bound \(b_i\) on every vector of \(\Lambda_i\). For \(x\in(\Lambda_0)_{(w)}\) and \(0\le r\le n\) there is \(y_r\in V_r\) with
\[
 \iota_r(y_r)=K^{[r]}(\iota_0(x)),\qquad
 \operatorname{VGe}_w\left(\Lambda_r,y_r;\sum_{i<r}b_i\right).
 \tag{4.28}
\]
In particular, under (4.24) and the stabiliser denominator condition, for \(p\notin\{T_u,d,g_0,g_1\}\) and each used \(\mu\) with \(p\notin Q'(\mu)\), every \(x\in S^{\mathrm{ram}}_{u,w}\) satisfies
\[
 \operatorname{ColGe}_{u,\mu,w}\bigl(x;-(m-\sigma)\delta_{u,w}\bigr).
 \tag{4.29}
\]
\end{lemma}
\leanref{Formalism.\allowbreak{}MainNotions.\allowbreak{}WeilJacquetShalikaLaw.\allowbreak{}Proof.\allowbreak{}chain\_\allowbreak{}ge}{Formalism.\allowbreak{}SharpRestriction.\allowbreak{}GrowthBounds.\allowbreak{}AdmissibleIntegers.\allowbreak{}RemainingPrimesOfAdmitted}{144--162}
\leanref{Formalism.\allowbreak{}MainNotions.\allowbreak{}WeilJacquetShalikaLaw.\allowbreak{}Proof.\allowbreak{}remaining\_\allowbreak{}primes\_\allowbreak{}of\_\allowbreak{}admitted}{Formalism.\allowbreak{}SharpRestriction.\allowbreak{}GrowthBounds.\allowbreak{}AdmissibleIntegers.\allowbreak{}RemainingPrimesOfAdmitted}{193--209}

\begin{proof}
The vectors of bound \(s\) form an \(\mathcal O_{K,(w)}\)-submodule. In a sum factor out the nonzero scalar of smaller valuation, so the ratio multiplying the other lattice vector is integral. The step bounds therefore extend to the localised lattices.

Start with \(y_0=x\). If \(y_r=az\) has bound \(\sum_{i<r}b_i\), apply the localised step bound to \(z\) and multiply the representative by \(a\). A nonzero resulting scalar has valuation at least \(e_w(\sum_{i<r}b_i+b_r)\). By linearity the embedded image is \(K^{[r+1]}(\iota_0(x))\).

By \ref{B4.06}, apply (4.28) to \(x\) in the reference chain or \(e^\mu_0x\) in the twisted chain. The sum of (4.24) is \(-(m-\sigma)\delta_{u,w}\).
\end{proof}

\begin{lemma}[The base column bound at an ordinary place]\label{B4.13}
Assume the \(F\)-identities, error identities, the first three conclusions of \ref{B4.11} and (4.29). For \(p\notin\{T_u,g_0,g_1\}\), each used \(\mu\) with \(p\notin Q'(\mu)\), and \(x\in S^{\mathrm{ram}}_{u,w}\), one has
\[
 \operatorname{ColGe}_{u,\mu,w}\bigl(x;\operatorname{lowS}_{u,w}(\sigma)\bigr).
 \tag{4.30}
\]
\end{lemma}
\leanref{Formalism.\allowbreak{}MainNotions.\allowbreak{}WeilJacquetShalikaLaw.\allowbreak{}Proof.\allowbreak{}first\_\allowbreak{}clause\_\allowbreak{}cases}{Formalism.\allowbreak{}SharpRestriction.\allowbreak{}GrowthBounds.\allowbreak{}AdmissibleIntegers.\allowbreak{}StatementsOfRemainingBound}{25--96}

\begin{proof}
In class \(\mathrm{ramD}\) the error is \(\delta_{u,w}+\operatorname{EatD}_j\), so use (4.23). For inert-class or signal primes \(p\le X''\), coprimality to the conductor and \(p\ne g_0,g_1\) give \(\neg\operatorname{Exc}(p)\). If the prime is inert with \(h_P=0\), use (4.22) with error \(\delta_{u,w}+\operatorname{EatC}_j(p)\). For the other primes in this range, drop the nonnegative gain from (4.21).

For the remaining primes, \(p\ne d\). By \ref{B4.02} and \ref{B4.03}, \(F_j(w)=0\) and \(E^{\mathrm{ev}}_{u,j}(w)=\delta_{u,w}\). Then (4.29) gives (4.30) with the stabiliser condition from \Lean{ErrDefs.\allowbreak{}dgam}.
\end{proof}

For each permitted \(h_g\), each used \(\mu\) with \(p\notin Q'(\mu)\), \(x\in V_{u,0}\) and \(s\in\R\), the condition \Lean{CmpV} is
\[
 \operatorname{vKge}_\mu(c_{u,w,h_g}(x);s)
 \quad\Longrightarrow\quad \operatorname{ColGe}_{u,\mu,w}(x;s).
 \tag{4.31}
\]
The condition \Lean{Ties9} ties the local positions and weights. It also supplies the ramified comparison \Lean{VramDef} and the base-term identity \Lean{FDef}.

\begin{theorem}[The complete column estimates at ordinary places]\label{B4.14}
Fix \(\sigma\in\R\). In \(\Lean{Stmt}\) the parameter range is \(\sigma\in\{0,1\}\). Fix \(W\) satisfying (4.31). Assume \ref{B4.09} and the conclusions of \ref{B4.11}--\ref{B4.12}. For every used twist and permitted \(u,w,h_g\), assume \ref{B3.42} for each matching admissible local system with its kernel datum.

For \(p\notin\{T_u,g_0,g_1\}\) and each used \(\mu\) with \(p\notin Q'(\mu)\), \(\Lean{Ties9}\) and \(\Lean{ErrDefs}\) give, for every \(x\in V_{u,0}\),
\[
 \operatorname{InNP}_{u,w}(x)\quad\Longrightarrow\quad
 \operatorname{ColGe}_{u,\mu,w}\bigl(x;\operatorname{totS}_{u,w}(\sigma)\bigr).
 \tag{4.32}
\]
For the same primes and indices, \(\Lean{FDef}\) and \(\Lean{ErrDefs}\) give, for every \(x\in S^{\mathrm{ram}}_{u,w}\),
\[
 \operatorname{ColGe}_{u,\mu,w}\bigl(x;\operatorname{lowS}_{u,w}(\sigma)\bigr).
 \tag{4.33}
\]
The served gain vanishes under the condition
\[
 \begin{split}
 \operatorname{Exc}(p)\quad&\text{or}\quad
 p\text{ belongs to neither the inert nor signal class}\\
 &\text{or}\quad(\operatorname{inertClass}(p)\text{ and }h_P=0)
 \quad\text{or}\quad p>X''
 \quad\text{or}\quad\operatorname{IsCs}(p).
 \end{split} \tag{4.34}
\]
For each permitted model at \(u,w,h_g\), each used \(\mu\) with \(p\notin Q'(\mu)\), and each tying datum for \(W.Yt(\mu)\), \(\Lean{ErrDefs}\) gives (4.17).
\end{theorem}
\leanref{Formalism.\allowbreak{}MainNotions.\allowbreak{}WeilJacquetShalikaLaw.\allowbreak{}Proof.\allowbreak{}statements\_\allowbreak{}of\_\allowbreak{}remaining\_\allowbreak{}bound}{Formalism.\allowbreak{}SharpRestriction.\allowbreak{}GrowthBounds.\allowbreak{}AdmissibleIntegers.\allowbreak{}StatementsOfRemainingBound}{194--200}
\leanref{Formalism.\allowbreak{}MainNotions.\allowbreak{}WeilJacquetShalikaLaw.\allowbreak{}Proof.\allowbreak{}baseChange\_\allowbreak{}preserves\_\allowbreak{}unitary\_\allowbreak{}of}{Formalism.\allowbreak{}SharpRestriction.\allowbreak{}GrowthBounds.\allowbreak{}AdmissibleIntegers.\allowbreak{}StatementsOfRemainingBound}{162--193}

\begin{proof}
Equation (4.33) follows from \ref{B4.13}. Equations (4.34) and (4.17) follow from (4.8)--(4.9) and \ref{B4.10}, respectively.

If \(\operatorname{InNP}\) holds through a permitted \(h_g\), the lattice \(N_P\) agrees on \(W.Yt(\mu)\) and \Lean{LatTied\_\allowbreak{}tw} supplies the tie. Apply (4.17) and (4.31), then use \(\operatorname{gainS}_{p,j}\le\operatorname{gainP}_{p,j}\).

Otherwise \(W_p=\varnothing\) and \(x\in S^{\mathrm{ram}}_{u,w}\). Under (4.34) use (4.33), with \Lean{FDef} from \Lean{Ties9}. In the remaining range \(p\) is a nonexceptional signal or nonzero even-order inert prime with \(p\le X''\). The complementary sum in (4.9) runs over \(R_p\), so
\[
 \operatorname{gainS}_{p,j}=g_{0,S}(p,j).
 \tag{4.35}
\]
This also holds when \(\operatorname{IsCs}(p)\), since both sides vanish. In the remaining range (4.21) gives (4.32).
\end{proof}

\section{Uniform constants and service for a family}

Let \(\mathcal F\in\mathrm{FamP}(\varepsilon)\) have data \(D_{x,i},L_{x,i}\) for \(x\in\R\), \(i\in I_x\), with parameter and size conditions for \(x>x_1\). For a member put \(f=D.R.f\), \(n=D.R.n\), \(Y=X''(n)\), \(N_m=D.R.N(m)\), and
\[
 F_{\mathrm{kindA},j}(p)=
 \frac1{n_P}\sum_{z\in\mathrm{kindA}(p)}F_{\mathrm{Orb},j}(p,z).
 \tag{4.36}
\]
The condition \(\Lean{DGamSize}(\mathcal F)\) gives \(C_D,x_D\in\R\) such that for \(x\ge x_D\), \(i\in I_x\) and all columns \(u\),
\[
 \log D\Gamma_u\le C_D fN_m\log(fN_m).
 \tag{4.37}
\]

\begin{lemma}[Uniform size constants]\label{B4.15}
Suppose \(\mathcal F\) satisfies \Lean{DGamSize}. There are \(C\ge0\) and \(x_0\in\R\) such that for \(x\ge x_0\), \(x>x_1\), \(i\in I_x\), the condition \(\Lean{TypePrimeSizeConds}(D,L,\mathbf c,c_e,C)\) holds. For every column \(u\) and \(1\le j\le m\) it gives
\[
 \sum_{t\in HH}\mathrm{ct}_{u,j,t}\le2\,\#TT\,b+2fkb',
 \tag{4.38}
\]
\[
 \log D\Gamma_u\le CfN_m\log(fN_m).
 \tag{4.39}
\]
For \(1\le j\le m\) and each finite set \(S\) of inert-class or signal primes \(p\le Y\) excluding inert-class primes with \(h_P=0\), it gives
\[
 \sum_{p\in S}\log p\,
 \bigl(\operatorname{EatI}_j(p)-\max\{0,F_{\mathrm{kindA},j}(p)\}\bigr)
 \le Cfb\log Y.
 \tag{4.40}
\]
\end{lemma}
\leanref{Formalism.\allowbreak{}MainNotions.\allowbreak{}Pk.\allowbreak{}admissible\_\allowbreak{}size\_\allowbreak{}family}{Formalism.\allowbreak{}SharpRestriction.\allowbreak{}SmoothDesignExistence.\allowbreak{}SumOverPrimesHypotheses}{57--80}
\leanref{Formalism.\allowbreak{}MainNotions.\allowbreak{}Pk.\allowbreak{}szEatI\_\allowbreak{}fam}{Formalism.\allowbreak{}SharpRestriction.\allowbreak{}SmoothDesignExistence.\allowbreak{}SumOverPrimesHypotheses}{38--48}

\begin{proof}
For \(x\ge3\), sum \(\mathrm{ct}_{u,j,t}\le2b\mathbf1_{TT}(t)+4kb'\) from \Lean{ct\_\allowbreak{}le} over the \(n=f/2\) positions to get (4.38). For (4.40), put
\[
 \begin{split}
 \mathcal R={}&N_m^2\bigl(\log(8nN_ma_\ast)+\log Y\bigr)
 +|c_{\mathrm{dl},1}|N_m^2\log(2a_\ast)
 +2|c_{\mathrm{dl},2}|\log(2a_\ast)^2,\\
 c_1={}&12+5|c_{e,6}|.
 \end{split} \tag{4.41}
\]
The scales are \(b=e^{x^\varepsilon}\), \(a_\ast=xb\), \(X'=e^{x^{\varepsilon/3}}\), \(b'=a_\ast/X'\), \(Y=xnX'\). With \(n\le x^3\) and \(N_m^2\le e^{x^{\varepsilon/3}}\), one has \(Y\mathcal R\le C_\varepsilon x^{4+2\varepsilon}e^{2x^{\varepsilon/3}}\le b\) and \(X'\ge x\) for large \(x\). In particular,
\[
 b'\le b,\qquad \mathcal R\ge0,\qquad Y\mathcal R\le b,\qquad
 Y>1,\qquad \log Y\ge1.
 \tag{4.42}
\]

Subtract the positive part of (4.36) from the error to obtain (4.26) plus \(\operatorname{EatX}_j(p)\). With \(1,\eta,\mathrm{Sc}\le b\) and \(\#TT_{BC}\le f\), use \(\gamma(p)\le2/p\) to bound the first term by \(12fb/p\). The weight term costs at most \(5|c_{e,6}|fb/p\) by \(0\le w_j(z)\le5b\) and \(\#BC\le f\). Multiply the remaining terms by \(\log p\) and use \(n_P\le2n\), \(p\le Y\), \(N_{j-1},N_j\le N_m\), and \(1/\log p\le2\) to bound them by \(f\mathcal R\). Hence
\[
 \log p\,\bigl(\operatorname{EatI}_j(p)-\max\{0,F_{\mathrm{kindA},j}(p)\}\bigr)
 \le c_1fb\frac{\log p}{p}+f\mathcal R+\log p\,\operatorname{EatX}_j(p).
 \tag{4.43}
\]

The supported primes of \(\operatorname{EatX}\) are inert with \(h_P\ne0\) of odd order. The image under \(h_H\) of the second Frobenius component must be trivial, since nontrivial image forces even order. Put \(e=D.R.e\) and \(q_e=2^e\). Frobenius gives \(p\equiv\pm1\bmod 4q_e\), and the size conditions give \(e\ge1\), \(q_e\ge x^2\), \(k\le x^5\). For any finite set \(S'\) of positive integers in these classes,
\[
 \sum_{p\in S'}\frac{\log p}{p^3}\le\frac{\log q_e}{q_e^3}.
 \tag{4.44}
\]
The integer \(1\) contributes zero. For \(l=4q_er\pm1\), \(r\ge1\), use \(\log l\le2(l/q_e)\log q_e\) and \(l\ge3q_er\) to get \(\log l/l^3\le\log q_e/(4q_e^3r^2)\). Sum over both signs with \(\sum_{r\ge1}r^{-2}\le2\). As \(k\le q_e^{5/2}\) and \(\log q_e\le2\sqrt{q_e}\), \(k\) times the right side of (4.44) is at most \(2\).

By (4.44) and the weight bound \(5fb\), the cubic term weighted by \(\log p\) sums to at most \(10fb\). The other term sums to at most \(|c_{e,6}|fb\) by \(\#S\le Y\) and \(YN_m^2\log a_\ast\le b\). Therefore
\[
 \sum_{p\in S}\log p\,\operatorname{EatX}_j(p)
 \le (10+|c_{e,6}|)fb\log Y.
 \tag{4.45}
\]
Rosser and Schoenfeld give \(\sum_{p\le Y}\log(p)/p<\log Y\) for \(Y>1\) \cite[(3.24)]{rosser-schoenfeld-1962}. Sum (4.43) using this inequality. For the remaining terms, use \(\#S\,\mathcal R\le b\) and (4.45). This gives (4.40) with \(C_E=c_1+1+10+|c_{e,6}|\). Take \(C=\max\{C_E,C_D,0\}\) and \(x_0\) above the size thresholds and \(x_D\).
\end{proof}
\leanref{UpperTheory.\allowbreak{}G2.\allowbreak{}szEatI\_\allowbreak{}of\_\allowbreak{}stepA}{Development.\allowbreak{}AutomorphicArithmetic.\allowbreak{}OrbitLabellings.\allowbreak{}SzEatIFree}{11--68}
\leanref{PI.\allowbreak{}SzEatI.\allowbreak{}W.\allowbreak{}stepA}{Formalism.\allowbreak{}PeriodIntegrals.\allowbreak{}UnitaryGroupMassFormula.\allowbreak{}SizeAbsorption.\allowbreak{}SizeDel}{121--344}
\leanref{UpperTheory.\allowbreak{}sum\_\allowbreak{}log\_\allowbreak{}mul\_\allowbreak{}max\_\allowbreak{}zero\_\allowbreak{}le\_\allowbreak{}of\_\allowbreak{}odd\_\allowbreak{}addOrderOf}{Development.\allowbreak{}AutomorphicArithmetic.\allowbreak{}OrbitLabellings.\allowbreak{}SumLogMulMaxZeroLeOfOddAddOrderOf}{101--170}

A family class \(\mathscr C\) assigns a predicate to each \(\Lean{FamP}(\varepsilon)\). For a family property \(P\), the assertion \(\operatorname{Closed}_{\mathscr C}(P)\) means
\[
 \begin{gathered}
 \exists\varepsilon_1>0\quad
 \forall\,0<\varepsilon\le\varepsilon_1\quad
 \exists\mathcal K_{\mathrm{size}}\quad
 \forall\,\mathcal F\in\mathrm{FamP}(\varepsilon),\ C_H\in\R,\\
 \mathscr C(\mathcal F)\ \land\
 \operatorname{Large}(\mathcal F,\mathcal K_{\mathrm{size}},C_H)
 \ \Longrightarrow\ P(\mathcal F,C_H).
 \end{gathered} \tag{4.46}
\]
The datum \(\mathcal K_{\mathrm{size}}\) sets staged lower thresholds for the family constants. To combine finitely many closed conditions take the smallest \(\varepsilon_1\) and componentwise largest thresholds.

For a member, \(\mathrm{Serv}(C)\) consists of \(\mathrm{TypePrimeSizeConds}(C)\) and (4.32)--(4.33) with their stated hypotheses at \(\sigma=1\). The family property \(\Lean{ServF}(\mathcal F)\) is
\[
 \exists C\ge0\ \exists x_0\in\R\ \forall x\in\R,\quad
 (x\ge x_0\land x>x_1)\ \Longrightarrow\
 \forall i\in I_x,\ \Lean{Serv}_{x,i}(C).
 \tag{4.47}
\]

\begin{theorem}[Uniform service for the family]\label{B4.16}
Let \(\mathscr C\) be a family class. Assume \(\operatorname{Closed}_{\mathscr C}(\Lean{DGamSize})\) and \(\Lean{SumOverPrimesHypotheses}(\mathscr C)\). The latter has two clauses, each under (4.46).

First, beyond a family threshold each member admits levels and chains \(LC(\mu)\) on the given filling. Use \(LC(\mu_w)\) at \(w\) to form \(Y'\) from \(Y\). The resulting \(W\in\Lean{TwLoc}(D,Y')\) satisfies (4.31) and, under \Lean{ErrDefs}, the local-system hypothesis of \ref{B4.09} with \(\sigma=1\).

Second, beyond a family threshold each member admits kernel data \(E\). At every level-free \((u,\mu,w)\) and \(0\le i<m\) the class-pair scalar has \(p\)-adic norm one. For used \(\mu\) the class-pair map \(K_{i,c,c'}\) takes the evaluated \(K\)-rational source space into the evaluated \(K\)-rational target space. Each reference lift and each used twisted lift takes the complex span \(\mathrm{Fi}_i\) of its embedded lattice into \(\mathrm{Fi}_{i+1}\). Their initial stabilisers are trivial.

For this \(E\), choose \(ep_w\in\mathbb C_p\) and, at each position \(z\), a scalar \(sg_{w,z}\in\mathbb C_p\) and a series \(be_{w,z}\) over \(\mathbb C_p\) in variables indexed by \(\Z/n_P\Z\). Require the even-orbit and conjugate-place expansions. These include the odd-order theta product. Require also the no-operator bound, jet integrality, the discriminant kernel condition and the weight-shape bounds.

Then
\[
 \operatorname{Closed}_{\mathscr C}(\Lean{ServF}).
 \tag{4.48}
\]
The conclusion also holds under the closed stabiliser size condition and \(\Lean{SelectionServed}(\mathscr C,\mathrm{sel})\) for a twist selection \(\mathrm{sel}\).
\end{theorem}
\leanref{Formalism.\allowbreak{}MainNotions.\allowbreak{}Pk.\allowbreak{}sum\_\allowbreak{}over\_\allowbreak{}primes\_\allowbreak{}service}{Formalism.\allowbreak{}SharpRestriction.\allowbreak{}SmoothDesignExistence.\allowbreak{}SumOverPrimesHypotheses}{187--222}
\leanref{Formalism.\allowbreak{}L2.\allowbreak{}G.\allowbreak{}ServT.\allowbreak{}closed\_\allowbreak{}servF\_\allowbreak{}of\_\allowbreak{}servT\_\allowbreak{}of\_\allowbreak{}dGamSize}{Development.\allowbreak{}AutomorphicArithmetic.\allowbreak{}ThetaCorrespondence.\allowbreak{}ThetaPairLifts.\allowbreak{}ThirdPlaceSeries.\allowbreak{}GServT}{31--33}

\begin{proof}
Combine the closed hypotheses with \Lean{G.\allowbreak{}num\_\allowbreak{}closed}. For each resulting family, \Lean{G.\allowbreak{}numCmpF\_\allowbreak{}fixed} gives \ref{B3.53} beyond a threshold \(x_n\) with \(x\ge2\) and \(r_G<k\). Let \(x_b,x_c\) be the thresholds of the two clauses and choose \(C_s\ge0,x_a\) by \ref{B4.15}. With splitting exponent \(\mathrm{ez}>0\), take
\[
 x_0=\max\{x_a,x_b,x_c,x_n,3^{5/\mathrm{ez}}\}.
 \tag{4.49}
\]
For \(x\ge x_0\), \(x>x_1\), \(i\in I_x\), the family tie gives \(\mathrm{splitBound}\ge\lfloor x^{\mathrm{ez}/5}\rfloor\ge3\). The host modulus and splitting condition exclude \(2,3\) from the inert and signal classes. At a signal prime \(n_P=2\operatorname{ord}(2h_P)\), and at an even-order inert prime \(n_P=\operatorname{ord}(h_P)\). Thus both satisfy \Lean{OrbOK}.

Choose the chains from the first clause. As in \ref{B3.56}, the numerical comparisons and \ref{B3.54} give the Hecke-image estimate by \ref{B4.09} for the chosen systems outside \(Q'(\mu)\). At every permitted place, including those with \(p\in Q'(\mu)\), \ref{B3.45} gives \ref{B3.42} for every matching system and kernel datum satisfying its admissibility hypotheses.

For the kernel from the second clause, \ref{B4.10a} gives (4.20) with \(\sigma=1\) at every level-free place. The estimates of \ref{B4.01b} give \Lean{OrbitLetterBoundsAll} with \((C_6,C)=(11,1)\). The fixed values \(c_{e,6}=12\), \(c_{\rm dig}=3\), \(c_{\mathrm{dl},1}=5\) and \(c_{\mathrm{dl},2}\ge0\) meet the inequalities of \ref{B4.11a}. That proposition gives \Lean{Core}. With orbit parity and \(g_0,g_1\ne d\), the remaining kernel data give \Lean{OrbitPlaces.\allowbreak{}admitted}.

Apply \ref{B4.14} to \(Y'\). The ties for \(Y\) also hold for \(Y'\), so \Lean{serv\_\allowbreak{}re} gives \(\Lean{Serv}(C_s)\) for \(Y\), with threshold (4.49). For a twist selection, \Lean{SelectionServed} supplies the first clause and the kernel of \ref{B8.56} as witness to the second.
\end{proof}

\begin{lemma}[The orbit and discriminant entry bounds]\label{B4.16a}
Under the hypotheses of \ref{B4.16}, at \((C_6,C)=(11,1)\), each representative eligible for clause (2) of \ref{B4.11a} has entry exponent at least \(g_{0,r,j}+F_{\mathrm{Orb},r,j}-\mathrm{err}_{r,j}\). At every level-free \(w\mid d\) each class pair has exponent at least \(-\frac12\sum_t\mathrm{ct}_{u,j,t}\).
\end{lemma}
\begin{proof}
Use clause (2) for the orbit bound and clause (3) for the discriminant bound.
\end{proof}
\leanref{PI.\allowbreak{}OrbBoundFW.\allowbreak{}prvF\_\allowbreak{}of\_\allowbreak{}orbitLetterBounds}{Formalism.\allowbreak{}PeriodIntegrals.\allowbreak{}UnitaryGroupMassFormula.\allowbreak{}OrbitValuationBounds.\allowbreak{}PrvFOfOrbitLetterBounds}{17--59}
\leanref{jet\_\allowbreak{}bound\_\allowbreak{}ramified\_\allowbreak{}prime}{Formalism.\allowbreak{}PeriodIntegrals.\allowbreak{}OrbitalBoundsConclusion.\allowbreak{}EntAtD}{11--13}

The period estimate of \ref{B4.19} gives exponent \(e_w(s-\delta_{u,w})\) from a column bound \(s\). Theorems \ref{B4.26} and \ref{B4.27} use it with (4.32) and (4.33), respectively.


% ---- p4c3 ----
\chapter{From columns to periods}\label{ch:columns-to-periods}

Fix \(D,L\) and the divisibility datum \(D_x\) of Alpöge's
archive~\cite{alpoge-2026}. Write \(K=D.\mathrm{what}_2.\mathrm{LL}\),
with specified embedding \(\iota:K\hookrightarrow\C\), and
\(\mathcal O=\mathcal O_K\). At a finite place \(w\), let \(p_w\),
\(e_w\) and \(\mathfrak p_w\) denote its residue characteristic,
ramification index and prime ideal. Normalize \(v_w:K^\times\to\Z\)
by \(v_w(p_w)=e_w\).

For a column \(u\), put
\[
 T_u=(D.\mathrm{mem}\,u).A.T,\qquad
 \Delta_u=(D.\mathrm{mem}\,u).C.\mathrm{DGam},\qquad
 \delta_{u,w}=v_{p_w}(\Delta_u),
\]
where \(v_p\) is the exponent of \(p\) in a positive integer. Let
\(\mu_0\) be the reference twist index and \(\mu_w\) the index selected
at \(w\). The used tuples form the set
\[
 \mathcal U_T=\{\mu_0\}\cup\{\mu_w:w\text{ a finite place of }K\}.
\]
For \(Q'(\mu)=D_x.\mathrm{Qp}(\mu)\), the selection gives
\(p_w\notin Q'(\mu_w)\).

Let \(V_u=(D.\mathrm{div}.\mathrm{Nu}\,u).V\) and let
\(\mathrm{nuIn}_u:V_u\hookrightarrow V^0_{u,0}\) be the inclusion
into the initial reference space. With initial reference lattice
\(\Lambda^0_{u,0}\), put
\[
 N_u=\mathrm{NuLat}_u,\qquad
 M_u^{\mathrm{ram}}=\mathrm{nuIn}_u^{-1}(\Lambda^0_{u,0}).
\]
The complex realization \(jN_u=D.jN\,u\) of \(V_u\) is
\(\iota\)-semilinear. Define the column period by
\[
 P_u^\mu(\phi,H)=D.\mathrm{what}_2.\mathrm{perCol}\,u\,\mu\,(jN_u\phi)\,H.
\]
Write \(\widetilde M_\mu\) for the ambient test lattice and \(V_0(\mu)\)
for the image of the self-adjoint projector \(eOp_\mu\). The admissible
test vectors belong to \(\widetilde M_\mu\cap V_0(\mu)\).

Write \(S_{u,w}=\sg(u,w)\) as in \ref{B4.05}, and set
\[
 B_{u,w}=\sum_{j=1}^{m}\bigl(F_j(w)-E_{u,j}(w)\bigr).
\]

\section{Localizing the terminal pairing}

For an \(\mathcal O\)-submodule \(\Lambda\) of a \(K\)-vector space,
write \(\Lambda_{(w)}=\operatorname{span}_{\mathcal O_{K,\mathfrak p_w}}\Lambda\).

\begin{lemma}[A pairing bound survives localization]\label{B4.17}
Let \(V\) be a \(K\)-vector space, \(\Lambda\subset V\) an
\(\mathcal O\)-submodule and \(g:V\to\C\) an \(\iota\)-semilinear map.
Fix \(b\in\R\) and a finite place \(w\). Suppose every \(x\in\Lambda\)
has a representative \(a_x\in K\) with
\[
 \iota(a_x)=g(x),\qquad a_x=0\ \text{or}\ v_w(a_x)\ge b.
\]
Then the same assertion holds on \(\Lambda_{(w)}\).
\end{lemma}
\leanref{Formalism.\allowbreak{}MainNotions.\allowbreak{}WeightDegreeSizes.\allowbreak{}pair\_\allowbreak{}core}{Formalism.\allowbreak{}SharpRestriction.\allowbreak{}Divisibility.\allowbreak{}WeightDegreeSizes}{783--801}

\begin{proof}
Write \(z=\sum_\nu c_\nu x_\nu\) with \(x_\nu\in\Lambda\) and
\(c_\nu\in\mathcal O_{K,\mathfrak p_w}\). By semilinearity,
\(a=\sum_\nu c_\nu a_{x_\nu}\) represents \(g(z)\). Every nonzero
summand has valuation at least \(b\), since \(v_w(c_\nu)\ge0\)
for \(c_\nu\ne0\). If \(a\ne0\), the nonarchimedean inequality
gives \(v_w(a)\ge b\).
\end{proof}

Fix \(T:\Lean{TwForms}(D)\). Write \(\Lambda^\mu_{u,m}\) and
\(\jmath^\mu_{u,m}\) for the terminal lattice and its complex embedding,
from the reference chain at \(\mu_0\) and from \(T.\mathrm{fm}(u,\mu)\)
otherwise. The terminal pairing is
\[
 \mathcal T_u^\mu(z,H)=P^{\mu,\mathrm{til}}_u
 \bigl(\jmath^\mu_{u,m}z,(e^\mu_{M,u})^{-1}H\bigr),
\]
where \(P^{\mu,\mathrm{til}}_u\) pairs after evaluation and lowering, and
\(e^\mu_{M,u}\) identifies its test space with the ambient one.

The point hypotheses give a rational terminal point map \(q_u^\mu\)
taking \(\Lambda^\mu_{u,m}\) into the point lattice. In the reference
case, \Lean{MemberChainIntegral} supplies \Lean{Conclusion6}. For the
other used tuples the hypothesis is \Lean{TwistedPointIntegral}.
Under \Lean{DGamDef}, the period hypothesis \Lean{PeriodClassSumWeyl}
gives one \(K\)-representative for the period of each reference
point-lattice vector against an admissible test vector. Its valuation,
if it is nonzero, is at least \(-e_w\delta_{u,w}\) at every \(w\).
For the other used tuples, \Lean{TwistedPeriodAdmissible} gives the
same bound at the places with \(p_w\notin Q'(\mu)\).

\begin{lemma}[The terminal denominator]\label{B4.18}
Assume \Lean{DGamDef} and the reference and twisted point and period
hypotheses. For every \(u,w\), take \(\mu=\mu_0\) or
\(\mu\in\mathcal U_T\setminus\{\mu_0\}\) with \(p_w\notin Q'(\mu)\).
For every \(z\in(\Lambda^\mu_{u,m})_{(w)}\) and
\(H\in\widetilde M_\mu\cap V_0(\mu)\), there is \(a\in K\) with
\[
 \iota(a)=\mathcal T_u^\mu(z,H),\qquad
 a=0\ \text{or}\ v_w(a)\ge-e_w\delta_{u,w}.
\]
\end{lemma}
\leanref{Formalism.\allowbreak{}L2.\allowbreak{}G.\allowbreak{}Adm2InU.\allowbreak{}pair0U}{Formalism.\allowbreak{}LocalToGlobal.\allowbreak{}GroupLattices.\allowbreak{}AdmissibleAtGeneratorPlaces}{39--54}
\leanref{Formalism.\allowbreak{}L2.\allowbreak{}G.\allowbreak{}Adm2InU.\allowbreak{}pairTwU}{Formalism.\allowbreak{}LocalToGlobal.\allowbreak{}GroupLattices.\allowbreak{}AdmissibleAtGeneratorPlaces}{55--71}

\begin{proof}
For \(z\in\Lambda^\mu_{u,m}\), the point hypothesis puts \(q_u^\mu(z)\)
in the point lattice, with
\[
 \operatorname{ptEmb}(q_u^\mu(z))
 =\operatorname{lowerMap}
 \bigl(\operatorname{ev}_{u,\mu}(\jmath^\mu_{u,m}z)\bigr).
\]
Here \(\operatorname{ev}_{u,\mu}\) is the chain evaluation map.
Pairing with \((e^\mu_{M,u})^{-1}H\) gives \(\mathcal T_u^\mu(z,H)\).
The class denominator of \ref{B4.18a} supplies the period hypothesis
on the constructed columns. For fixed data its bound is
\(b=-e_w\delta_{u,w}\). Apply \ref{B4.17} to the semilinear pairing in
\(z\).
\end{proof}

\begin{lemma}[The class denominators on the constructed columns]\label{B4.18a}
Take the constructed columns with the reference and twisted
evaluations of \ref{B8.09}, and assume \Lean{DGamDef}. For every
\(u\) and every \(\mu\in\mathcal U_T\), evaluation and lowering take
the terminal lattice into the point lattice. Fix \(x\) in that point
lattice and \(H\in\widetilde M_\mu\cap V_0(\mu)\). Let \(h_c\) be the
order of the point stabilizer at class \(c\). The point period has
one representative
\[
 y=\sum_c b_c/h_c\in K.
\]
For \(\mu=\mu_0\) at every finite \(w\), and for \(\mu\ne\mu_0\)
at every \(w\) with \(p_w\notin Q'(\mu)\), one has
\(b_c\in\mathcal O_{K,\mathfrak p_w}\) for all \(c\) and
\(\Delta_u y\in\mathcal O_{K,\mathfrak p_w}\).
\end{lemma}
\leanref{GlobalTheory.\allowbreak{}RowIntQ.\allowbreak{}rowInt\_\allowbreak{}over}{Development.\allowbreak{}AbelianVarietyTheory.\allowbreak{}Proofs.\allowbreak{}GlobalInterpolation.\allowbreak{}PeriodRowLatticesIntegralProvisionalRows}{47--64}
\leanref{GlobalTheory.\allowbreak{}PairAtClass.\allowbreak{}pair\_\allowbreak{}integral}{Development.\allowbreak{}AbelianVarietyTheory.\allowbreak{}Proofs.\allowbreak{}GlobalInterpolation.\allowbreak{}ClassPairingIntegralFormRowIsIntegral}{49--70}
\leanref{GlobalTheory.\allowbreak{}PerRed.\allowbreak{}admissiblePeriodZero\_\allowbreak{}of\_\allowbreak{}classPeriodDecomposition}{Development.\allowbreak{}AbelianVarietyTheory.\allowbreak{}Proofs.\allowbreak{}GlobalInterpolation.\allowbreak{}AdmissiblePeriodZeroOfClassPeriodDecompositionAl}{196--240}
\leanref{GlobalTheory.\allowbreak{}PerRedU.\allowbreak{}perTwAt\_\allowbreak{}of\_\allowbreak{}perDecTw}{Development.\allowbreak{}AbelianVarietyTheory.\allowbreak{}Proofs.\allowbreak{}GlobalInterpolation.\allowbreak{}AdmissiblePeriodTwistsReductionPlaceTuples}{12--40}
\leanref{GlobalTheory.\allowbreak{}Mem.\allowbreak{}Col2.\allowbreak{}TC.\allowbreak{}formsMap\_\allowbreak{}at\_\allowbreak{}every\_\allowbreak{}trueColumnMember}{Development.\allowbreak{}AbelianVarietyTheory.\allowbreak{}Proofs.\allowbreak{}GlobalInterpolation.\allowbreak{}FormsMapAtEveryTrueColumnMember}{15--27}
\leanref{GlobalTheory.\allowbreak{}PerQ.\allowbreak{}twistedPointIntegral\_\allowbreak{}of\_\allowbreak{}TC\_\allowbreak{}ok}{Development.\allowbreak{}AbelianVarietyTheory.\allowbreak{}Proofs.\allowbreak{}GlobalInterpolation.\allowbreak{}TwistedPointRowAdmissibleTrueColumns}{12--16}
\leanref{GlobalTheory.\allowbreak{}RedPt0Pend.\allowbreak{}exists\_\allowbreak{}integralPendingObject\_\allowbreak{}symmetricFormsLevel}{Development.\allowbreak{}AbelianVarietyTheory.\allowbreak{}Proofs.\allowbreak{}GlobalInterpolation.\allowbreak{}ExistsIntegralPendingObjectSymmetricFormsLevel}{291--333}

\begin{proof}
In each tensor factor, evaluation followed by lowering uses the
interlacing branching map. This map takes integral polynomial
generators to integral linear combinations of target generators.
In reference class \(c\), the change of frame contributes
\(\zeta_c^{-1}\), with \(\iota(\zeta_c)\) a finite-order level
character value. Thus \(\zeta_c^n=1\) for some \(n>0\) and
\(v_w(\zeta_c)=0\) at every finite place. The coordinates stay integral.
The coordinate criterion for the point lattice gives a local point
representative at each \(w\). Since the terminal lattice localized at
one place spans the rational space, the evaluation descends by
linearity. Injectivity of the point embedding identifies the local
representatives with this rational evaluation. The intersection of
the local point lattices is the global point lattice.

Write \(H=\sum_v a_v v\) on the integral row generators. The same
linear combination of their rational lifts gives a rational row
with the prescribed class readings. At each relevant place its
polynomial \(r_c\) at \(c\) belongs to the dual of the point lattice
for the Weyl pairing. Thus \(b_c=\langle x_c,r_c\rangle\) is integral
there, where \(x_c\) is the point coordinate. The restriction to \(c\)
leaves one term in the Petersson sum, with weight \(1/h_c\). Summing
gives \(\iota(y)\) equal to the point period.

The point stabilizer injects into the terminal class stabilizer.
Its order therefore divides \(\Delta_u\) by the reference chain
data and, on a twist, by \Lean{DGamDef}. Hence
\[
 \Delta_u y=\sum_c(\Delta_u/h_c)b_c
 \in\mathcal O_{K,\mathfrak p_w}.
\]
Since \(v_w(\Delta_u)=e_wv_{p_w}(\Delta_u)\), this is the
period bound of \ref{B4.18}.
\end{proof}

By \ref{B4.06}, a column has bound \(s\in\R\) at \(w\) if its
terminal image has a representative \(az\) with
\[
 z\in(\Lambda^\mu_{u,m})_{(w)},\qquad a\in K,\qquad
 a=0\ \text{or}\ v_w(a)\ge e_ws.
\]
The reference column starts from \(\mathrm{nuIn}_u\phi\).
On a twist, \Lean{ColGe} uses the corresponding complex input
under the twist comparison.

\begin{theorem}[From a column bound to a period bound]\label{B4.19}
Assume \Lean{DGamDef} and the point and period hypotheses of
\ref{B4.18}. Fix a column \(u\), a finite place \(w\), a vector
\(\phi\in V_u\), a real number \(s\) and a tuple \(\mu\in\mathcal U_T\)
with \(p_w\notin Q'(\mu)\).
If the \(\mu\)-column from \(\mathrm{nuIn}_u\phi\) has bound \(s\)
at \(w\), then for every \(H\in\widetilde M_\mu\cap V_0(\mu)\)
there is \(y\in K\) with
\[
 \iota(y)=P_u^\mu(\phi,H),\qquad
 y=0\ \text{or}\ v_w(y)\ge e_w(s-\delta_{u,w}).
\]
\end{theorem}
\leanref{Formalism.\allowbreak{}L2.\allowbreak{}G.\allowbreak{}Adm2InU.\allowbreak{}period\_\allowbreak{}geTU}{Formalism.\allowbreak{}LocalToGlobal.\allowbreak{}GroupLattices.\allowbreak{}AdmissibleAtGeneratorPlaces}{72--91}

\begin{proof}
Choose a terminal representative \(az\) and apply \ref{B4.18} to
\(\mathcal T_u^\mu(z,H)\), with representative \(y_0\).
The column-period comparison and semilinearity give
\[
 P_u^\mu(\phi,H)=\iota(a)\mathcal T_u^\mu(z,H)=\iota(ay_0).
\]
If \(ay_0\ne0\), then \(v_w(ay_0)=v_w(a)+v_w(y_0)\ge
e_ws-e_w\delta_{u,w}\).
\end{proof}

\section{The level prime}

Write \(\tau(w)=L.\mathrm{tau}(w)\) and
\(\mathrm{EatT}(u,w)=T^*_{u,w}\) as in \ref{B4.03}.

\begin{lemma}[The budget above \(T_u\)]\label{B4.20}
Under \Lean{ParameterChoice'} and the formulas \Lean{FDef}
and \Lean{EDef}, every \(u,w\) with \(p_w=T_u\) satisfies
\[
 \tau(w)=0,\qquad
 \mathrm{gain}_j(w)=F_j(w)=0,\qquad
 E_{u,j}(w)=\mathrm{EatT}(u,w)
 \quad(1\le j\le m).
\]
Whenever these identities hold, \(S_{u,w}=-m\,\mathrm{EatT}(u,w)\).
\end{lemma}
\leanref{Formalism.\allowbreak{}MainNotions.\allowbreak{}J.\allowbreak{}sizeData\_\allowbreak{}at\_\allowbreak{}level\_\allowbreak{}primes}{Formalism.\allowbreak{}SharpRestriction.\allowbreak{}AdmissibleParameterJoins.\allowbreak{}DIVfromGLs}{25--28}
\leanref{Formalism.\allowbreak{}MainNotions.\allowbreak{}WeightDegreeSizes.\allowbreak{}sg\_\allowbreak{}atT}{Formalism.\allowbreak{}SharpRestriction.\allowbreak{}Divisibility.\allowbreak{}WeightDegreeSizes}{887--902}

\begin{proof}
The parameter choice puts \(T_u\) in \(\mathrm{splitK}\), where
\(\tau(w)\) and the raw gain vanish. The same class and
\(T_u\notin\{g_0,g_1\}\) give \(F_j(w)=0\).
The first branch of \(E\) gives \(\mathrm{EatT}(u,w)\).
With zero raw gain the screened gain is zero, so the budget sums
\(m\) copies of \(-\mathrm{EatT}(u,w)\).
\end{proof}

\begin{lemma}[Eventual coprimality of the level prime]\label{B4.21}
Fix a family class \(\mathcal C\) and \(C_0,\varepsilon\in\R\).
Let \(\mathfrak F\in\mathrm{FamP}(\varepsilon)\) belong to
\(\Lean{Des}\,\mathcal C\,C_0\). There is \(x_0\in\R\) such that
for every \(x\ge x_0\) with \(x>\mathfrak F.x_1\), every
\(i\in\mathfrak F.I_x\), and every column \(u\) of \(D_{x,i}\),
\[
 T_u\nmid\Delta_u.
\]
The eventual assertion is closed on \(\Lean{Des}\,\mathcal C\,C_0\).
\end{lemma}
\leanref{Formalism.\allowbreak{}MainNotions.\allowbreak{}eventually\_\allowbreak{}level\_\allowbreak{}prime\_\allowbreak{}coprime\_\allowbreak{}stabiliser\_\allowbreak{}of\_\allowbreak{}design}{Formalism.\allowbreak{}SharpRestriction.\allowbreak{}GlobalSchurBounds.\allowbreak{}PlugDgt}{10--101}
\leanref{Formalism.\allowbreak{}MainNotions.\allowbreak{}closed\_\allowbreak{}eventually\_\allowbreak{}levelPrimeCoprimeStabiliser}{Formalism.\allowbreak{}SharpRestriction.\allowbreak{}GlobalSchurBounds.\allowbreak{}PlugDgt}{102--107}

\begin{proof}
Take the maximum of the design thresholds for stabilizer
smoothness and \Lean{TuSize}. In any member beyond it, a prime
\(T_u\mid\Delta_u\) would satisfy \(T_u\le2fN_m+1\).
For the chain parameter \(t\) and ranks \(N_i=k+i(t+1)\),
the rank conditions give \(t\ge7\) and \(f,k,m\le t\). Hence
\[
 2fN_m+1\le 2t^3+4t^2+1\le7t^3<49t^3\le t^5.
\]
Coprimality of \(T_u\) and \(Q_S\) gives \(Q_Sqd\ge1\). The size
condition gives \(t^{150}\le(Q_Sqd)^{36}t^{150}\le\sqrt b\).
Squaring and applying \Lean{TuSize} with \(m\ge2\) gives
\[
 t^5\le t^{300}\le b\le b^{m/2}\le T_u.
\]
This contradicts the bound for a prime divisor of \(\Delta_u\).
For closedness take the cutoff \(1\) for \(\varepsilon\) and zero
size thresholds.
\end{proof}

Put \(V^{\mathrm{ram}}_{u,0}=\operatorname{im}(\mathrm{nuIn}_u)\) and
\(S^{\mathrm{ram}}_{u,w}=(\Lambda^0_{u,0})_{(w)}\cap V^{\mathrm{ram}}_{u,0}\).
Above \(T_u\), let \(VVT_{u,w}\) consist of the level indices \(v\)
whose labels \(\tau_w(v)\) are top positions at a seen place. At an
unseen place set \(VVT_{u,w}=\varnothing\). Let \(U\) be the ordered
product of the initial level operators at these indices. With
highest-weight coordinates numbered from zero, put
\[
 b_{\rm top}(\tau)=-\lambda^u_{0,\tau,k-1},\qquad
 a=m\sum_{v\in VVT_{u,w}}b_{\rm top}(\tau_w(v)).
\]
The submodule \(N_{\lambda,u,w}\) consists of the \(x\in V^0_{u,0}\)
whose image under the initial embedding \(\jmath^0_{u,0}\) is
\(T_u^aU^m\jmath^0_{u,0}(z)\) for some \(z\in S^{\mathrm{ram}}_{u,w}\).
The condition \Lean{NuAtTDef} identifies \((N_u)_{(w)}\) with the
inverse image of \(N_{\lambda,u,w}\) under \(\mathrm{nuIn}_u\).

The reference bound \Lean{Conclusion5Sharp} puts the column of
every \(x\in N_{\lambda,u,w}\), scaled by \(\iota(c)\), in the image
of \((\Lambda^0_{u,m})_{(w)}\) whenever \(c\in K^\times\) and
\(v_w(c)\ge e_wmT^*_{u,w}\). The hypothesis
\Lean{TwistedColumnBoundAtT} gives bound \(-mT^*_{u,w}\) on the same
inputs for each nonreference used tuple \(\mu\) with
\(p_w\notin Q'(\mu)\).

\begin{theorem}[Periods above the level prime]\label{B4.22}
Assume \Lean{DGamDef}, the point and period hypotheses of
\ref{B4.19}, \Lean{LevelPrimeSizes} and \Lean{NuAtTDef}.
Assume the reference and twisted sharp bounds and \(T_u\nmid\Delta_u\)
for every \(u\). Then for every \(u,w\) with \(p_w=T_u\),
every \(\phi\in N_u\), and every
\(H\in\widetilde M_{\mu_w}\cap V_0(\mu_w)\), there is \(y\in K\) with
\[
 \iota(y)=P_u^{\mu_w}(\phi,H),\qquad
 y=0\ \text{or}\ v_w(y)\ge e_wS_{u,w}.
\]
\end{theorem}
\leanref{Formalism.\allowbreak{}MainNotions.\allowbreak{}J.\allowbreak{}ColGe0\_\allowbreak{}atT}{Formalism.\allowbreak{}SharpRestriction.\allowbreak{}AdmissibleParameterJoins.\allowbreak{}DIVfromGLi}{34--68}
\leanref{Formalism.\allowbreak{}L2.\allowbreak{}G.\allowbreak{}Adm2InU.\allowbreak{}admissible\_\allowbreak{}at\_\allowbreak{}twist\_\allowbreak{}places}{Formalism.\allowbreak{}LocalToGlobal.\allowbreak{}GroupLattices.\allowbreak{}AdmissibleAtGeneratorPlaces}{92--112}

\begin{proof}
By \Lean{NuAtTDef}, \(x=\mathrm{nuIn}_u\phi\) lies in
\(N_{\lambda,u,w}\). The integer \(n_0=mT^*_{u,w}\) is nonnegative.
At a seen place, \(m\) even and \(t\) odd give \(4\mid m(t+1)\)
in the formula of \ref{B4.03}. At an unseen place \(n_0=0\).
For \(c=T_u^{n_0}\) the reference sharp bound represents the
reference column by \(c^{-1}z'\), with
\(z'\in(\Lambda^0_{u,m})_{(w)}\) and \(v_w(c^{-1})=-e_wn_0\).
At a nonreference selected tuple apply the twisted column bound.
Proposition~\ref{B4.22a} supplies both bounds on the constructed chains.

Apply \ref{B4.19} to \(\mu_w\) with \(s=-n_0\). Coprimality gives
\(\delta_{u,w}=0\), and \Lean{LevelPrimeSizes} gives \(S_{u,w}=s\)
by \ref{B4.20}.
\end{proof}

\begin{proposition}[The sharp column exponent at the level]\label{B4.22a}
Take the constructed columns with their rational jets, weight signs
and compatible level-zero lattices as in \ref{B8.09}. Fix \(u,w\)
with \(p_w=T_u\). For every \(x\in N_{\lambda,u,w}\) and
\(c\in K^\times\) with \(v_w(c)\ge e_wmT^*_{u,w}\), the reference
column of \(x\), scaled by \(\iota(c)\), belongs to the complex
image of its terminal local lattice. For every
\(\mu\in\mathcal U_T\setminus\{\mu_0\}\) with \(p_w\notin Q'(\mu)\),
the twisted column of \(x\) has bound \(-mT^*_{u,w}\).
\end{proposition}
\leanref{GlobalTheory.\allowbreak{}LvT.\allowbreak{}Built.\allowbreak{}TrueLaws.\allowbreak{}laws\_\allowbreak{}true\_\allowbreak{}fam}{Development.\allowbreak{}AbelianVarietyTheory.\allowbreak{}Proofs.\allowbreak{}GlobalInterpolation.\allowbreak{}FivePointLawsTrueChains}{46--63}
\leanref{GlobalTheory.\allowbreak{}LvT.\allowbreak{}UKS.\allowbreak{}ukSharp}{Development.\allowbreak{}AbelianVarietyTheory.\allowbreak{}Proofs.\allowbreak{}GlobalInterpolation.\allowbreak{}ExtInOfValIn}{155--199}
\leanref{GlobalTheory.\allowbreak{}LvT.\allowbreak{}levelChainSharpColumnValuation\_\allowbreak{}of\_\allowbreak{}stability}{Development.\allowbreak{}AbelianVarietyTheory.\allowbreak{}Proofs.\allowbreak{}GlobalInterpolation.\allowbreak{}LevelChainSharpColumnValuationOfStability}{174--194}
\leanref{GlobalTheory.\allowbreak{}LvT.\allowbreak{}Built.\allowbreak{}TrueCols.\allowbreak{}row\_\allowbreak{}true}{Development.\allowbreak{}AbelianVarietyTheory.\allowbreak{}Proofs.\allowbreak{}GlobalInterpolation.\allowbreak{}TrueColumnsRowSharpBounds}{36--51}
\leanref{GlobalTheory.\allowbreak{}LvT.\allowbreak{}StabOrder.\allowbreak{}stabUnit\_\allowbreak{}of\_\allowbreak{}embed}{Development.\allowbreak{}AbelianVarietyTheory.\allowbreak{}Proofs.\allowbreak{}GlobalInterpolation.\allowbreak{}NotDvdCardStabilizerForm}{157--164}

\begin{proof}
Put \(p=T_u\). In one chain let \(K_i\) be the lift from level \(i\)
to level \(i+1\), \(C_i=K_{i-1}\cdots K_0\) with \(C_0=1\), and
\(U_i\) the product of the level operators at the
\(r=|VVT_{u,w}|\) selected indices. Identify each local lattice \(S_i\)
with its image in the space of forms. Set
\[
 d_0=
 \begin{cases}
 2\sum_\tau W_{u,\tau}+fN_m,& w\text{ seen},\\
 0,& w\text{ unseen},
 \end{cases}
\]
with \(W_{u,\tau}\) as in \ref{B4.03}.

At a seen place, the integral moments of the point kernel in Weyl
coordinates give exponent \(fN_i+\sum_\tau s_{i,\tau}\) on the
target lattice with all level indices present. Here \(s_{i,\tau}\)
is the negative determinant shift of the shifted target weight.
Each rank-\(N_i\) class stabilizer acts faithfully on \(E^{N_i}\)
over the host field \(E\). For any prime divisor \(q\) of its order,
an element of order \(q\) has the \(q\)-th cyclotomic factor in
its rational minimal polynomial. Thus \(q-1\le[E:\Q]N_i\le2fN_m\).
The bound \(p>2fN_m+1\) makes the stabilizer orders units at \(w\).
Each selected operator sums incident values in
the lattice with that index removed. For each remaining index \(v\),
let \(\ell_{i+1,\tau_w(v)}\) be the first coordinate of the shifted
target weight at \(\tau_w(v)\). Multiplying the form by
\(p^{\ell_{i+1,\tau_w(v)}}\) removes \(v\) from the index set.
The resulting exponent is
\[
 fN_i+\sum_\tau s_{i,\tau}
 +\sum_{v\notin VVT_{u,w}}\ell_{i+1,\tau_w(v)}
 \ \le\ fN_m+2\sum_\tau W_{u,\tau}.
\]
Indeed \(N_i\le N_m\), the level labels are bijective with the
embeddings, and both \(s_{i,\tau}\) and each remaining
\(\ell_{i+1,\tau}\) are at most \(W_{u,\tau}\).
At an unseen place the conjugate-prime kernel bound gives an
integral value and \(U_{i+1}=1\). Thus in either case
\[
 p^{d_0}U_{i+1}K_i S_i\subseteq S_{i+1}.
\]

Set \(s=(t+1)/2\). On forms at level \(i\), each selected Hecke
operator contributes a factor \(p^s\). The lifts and level operators
preserve the form spaces, so iteration gives the identities on forms
\[
 U_{i+1}K_i=p^{sr}K_iU_i,\qquad U_i C_i=p^{sri}C_iU_0.
\]
For \(z\in S_0\), induction gives \(p^{b_i}C_i U_0^i z\in S_i\),
where
\[
 b_0=0,\qquad b_{i+1}=b_i+d_0+sr(i+1).
\]
Indeed applying \(p^{d_0}U_{i+1}K_i\) to the level-\(i\) vector
gives \(p^{b_{i+1}}C_{i+1}U_0^{i+1}z\). At a seen place the labels
identify the selected indices with the top positions, so \ref{B4.03}
gives
\[
 b_m=md_0+sr\,m(m+1)/2=mT^*_{u,w}.
\]
At an unseen place \(d_0=r=T^*_{u,w}=0\).

By definition, \(x\) is represented by \(p^aU_0^m z\) with
\(z\in S^{\mathrm{ram}}_{u,w}\). The top weight signs give \(a\ge0\).
The scalar \(cp^{a-b_m}\) has valuation
\(v_w(c)+e_w(a-b_m)\ge0\). Multiplying the level-\(m\) vector of
the induction by this scalar proves the reference assertion.

The twisted comparison carries \(x\) into the corresponding
initial submodule of the twisted chain. It preserves the seen
condition and identifies the weights, so the calculation gives
the same \(b_m\). With \(c=p^{b_m}\), it gives a representative
\(c^{-1}z'\) in which \(z'\) lies in the twisted terminal local
lattice and \(v_w(c^{-1})=-e_wmT^*_{u,w}\).
\end{proof}

\section{The two fixed primes}

Fix the local error datum \(X_e\). With the notation of \ref{B4.03},
including \(\eta=L.\mathrm{etaSz}\), put
\[
 \nu_{i,u,j,w}=e_w^{-1}\sum_\tau A_{u,j,\tau}\,v_w\bigl(\tau(\bar\pi_i)\bigr),
 \qquad
 \mathcal E_G=f\bigl(mN_m(N_m+t)+m^3\eta\bigr).
\]
These are \Lean{vNN} and \Lean{LiftErrorLocal.\allowbreak{}errC}.
At \(p_w=g_i\ne T_u\), the tuple hypothesis
\Lean{UsedTuplesBoundByConstant} with \(C_1\in\R\) gives, for each
\(\mu\in\mathcal U_T\), \(\phi\in M_u^{\mathrm{ram}}\) and
\(H\in\widetilde M_\mu\), a representative \(y\in K\) of
\(P_u^\mu(\phi,H)\) with
\[
 y=0\quad\text{or}\quad
 \frac{v_w(y)}{e_w}\ge-\sum_{j=1}^{m}\nu_{i,u,j,w}-C_1\mathcal E_G.
\]
The numerical hypothesis \Lean{SgAtG} is
\[
 B_{u,w}\le-\sum_{j=1}^{m}\nu_{i,u,j,w}-C_1\mathcal E_G.
\]

\begin{lemma}[The bound on the original lattice at \(g_i\)]\label{B4.23}
If \(p_w=g_i\) for \(i\in\{0,1\}\), then \(S_{u,w}=B_{u,w}\).
Suppose also that \(p_w\ne T_u\) and that the two fixed-prime
hypotheses hold with the same \(C_1\). For every
\(\phi\in M_u^{\mathrm{ram}}\) and \(H\in\widetilde M_{\mu_w}\)
there is \(y\in K\) with
\[
 \iota(y)=P_u^{\mu_w}(\phi,H),\qquad
 y=0\ \text{or}\ v_w(y)\ge e_wB_{u,w}.
\]
\end{lemma}
\leanref{Formalism.\allowbreak{}MainNotions.\allowbreak{}WeightDegreeSizes.\allowbreak{}sg\_\allowbreak{}atG}{Formalism.\allowbreak{}SharpRestriction.\allowbreak{}Divisibility.\allowbreak{}WeightDegreeSizes}{1007--1017}
\leanref{Formalism.\allowbreak{}L2.\allowbreak{}G.\allowbreak{}Adm2InU.\allowbreak{}atG\_\allowbreak{}coreU}{Formalism.\allowbreak{}LocalToGlobal.\allowbreak{}GroupLattices.\allowbreak{}AdmissibleAtGeneratorPlaces}{113--133}

\begin{proof}
The raw gain, hence the screened gain, vanishes above either
fixed prime. This gives \(S_{u,w}=B_{u,w}\).
Proposition~\ref{B4.23a} supplies the two hypotheses with a common
constant on the selected family. Apply the tuple bound to \(\mu_w\).
For \(y\ne0\), the numerical inequality gives \(v_w(y)/e_w\ge B_{u,w}\).
\end{proof}

\begin{proposition}[The common constant at the two fixed primes]\label{B4.23a}
Use the constructed rows whose rational forms agree with the
complex period maps as in \ref{B8.06}. Put
\(D_{u,j}=\sum_\tau A_{u,j,\tau}\). Suppose for \(p=g_0,g_1\) that
\[
 2fN_m+1<p,\qquad 2D_{u,j}+2<p,\qquad D_{u,j}+N_m<p
 \quad\text{for all }u,\ 1\le j\le m.
\]
Then \Lean{UsedTuplesBoundByConstant} holds with \(C_1=1\).
For size data satisfying \Lean{FDef} and \Lean{ErrDefs},
\Lean{SgAtG} holds for every \(C_1\ge0\) with \(2C_1\le c_g\),
where \(c_g\) is the fixed-prime error coefficient of \ref{B4.03}.
For \(0<\varepsilon\le1/100\), every selected family satisfying
the middle size threshold \(2\) has a cutoff beyond which both
bounds hold with \(C_1=1\), uniformly in the member index.
\end{proposition}
\leanref{GlobalTheory.\allowbreak{}JetsFromTables.\allowbreak{}jetLossAll\_\allowbreak{}of\_\allowbreak{}jetSplitAmb}{Development.\allowbreak{}AbelianVarietyTheory.\allowbreak{}Proofs.\allowbreak{}SatakeParameters.\allowbreak{}ThetaJetIntegralityTypePrimeOfSplitTableBound}{221--293}
\leanref{GlobalTheory.\allowbreak{}ColPeriodTies.\allowbreak{}boundAt\_\allowbreak{}of\_\allowbreak{}chainTie}{Development.\allowbreak{}AbelianVarietyTheory.\allowbreak{}Proofs.\allowbreak{}SatakeParameters.\allowbreak{}PeriodValuationBoundTypePrimePlaceOfChainAgreeme}{19--95}
\leanref{GlobalTheory.\allowbreak{}ColPeriod.\allowbreak{}vG\_\allowbreak{}loss\_\allowbreak{}le}{Development.\allowbreak{}AbelianVarietyTheory.\allowbreak{}Proofs.\allowbreak{}SatakeParameters.\allowbreak{}PeriodBoundOfJetLoss}{87--142}
\leanref{Formalism.\allowbreak{}MainNotions.\allowbreak{}J.\allowbreak{}sumFE\_\allowbreak{}atG\_\allowbreak{}le}{Formalism.\allowbreak{}SharpRestriction.\allowbreak{}AdmissibleParameterJoins.\allowbreak{}GandSx}{12--85}
\leanref{GlobalTheory.\allowbreak{}AtG.\allowbreak{}row\_\allowbreak{}atG\_\allowbreak{}of\_\allowbreak{}jetSplitAmb}{Development.\allowbreak{}AbelianVarietyTheory.\allowbreak{}Proofs.\allowbreak{}SatakeParameters.\allowbreak{}TypePrimePeriodRowOfConstructedPointChains}{15--29}
\leanref{GlobalTheory.\allowbreak{}AtG.\allowbreak{}jetSplitAmb\_\allowbreak{}holds}{Development.\allowbreak{}AbelianVarietyTheory.\allowbreak{}Proofs.\allowbreak{}SatakeParameters.\allowbreak{}JetTableBoundSplitPrime}{7--9}
\leanref{GlobalTheory.\allowbreak{}JetsFromTables.\allowbreak{}datum\_\allowbreak{}bound}{Development.\allowbreak{}AbelianVarietyTheory.\allowbreak{}Proofs.\allowbreak{}SatakeParameters.\allowbreak{}ThetaJetIntegralityTypePrimeOfSplitTableBound}{108--165}

\begin{proof}
Fix \(w\mid g_i\) with \(g_i\ne T_u\), a used tuple \(\mu\),
\(\phi\in M_u^{\mathrm{ram}}\) and \(H\in\widetilde M_\mu\).
Put
\[
 d_j=g_i^{\,fN_{j-1}N_j}\prod_\tau
 \tau(\bar\pi_i)^{A_{u,j,\tau}},\qquad
 d_*=\prod_{j=1}^m d_j .
\]
The factors are nonzero since \(\pi_i\bar\pi_i=g_i\).
At step \(j\), put \(p=g_i\). The source and target lattices are
self-dual and standard at \(p\). The split theta estimate clears
the \(p\)-part \(p^{fN_{j-1}N_j}\) of the averaging index. Its
diagonal substitution contributes the factor
\(\prod_\tau\tau(\bar\pi_i)^{A_{u,j,\tau}}\) on a monomial of the
prescribed multidegree. The multi-index factorial for passing from
coefficients to derivatives is integral.

Write \(c'\) for the prime-to-\(p\) part of the product of the two
level factors. Refinement identifies the rational jet sum with the
sum at level \(pc'\). At this level each nonzero test weight is
\(w\)-integral and its argument belongs to the localized homomorphism
lattice divided by \(c'\). The split estimate puts each such jet
multiplied by \(d_j\) in \(\mathcal O_{K,\mathfrak p_w}\), so the
weighted sum satisfies the same bound. The degree bound gives
\(D_{u,j}+1<p\). The integral Weyl operator therefore takes the
source lattice into the target lattice after multiplication by
\(d_j\). Equivariance gives the inclusion in the class lattices.
Composition gives the factor \(d_*\) for the column.

Lift the row generators as in \ref{B4.18a}. The pairing of the
rational column with the lift of \(H\) gives \(y\in K\), and
agreement of the row forms gives \(\iota(y)=P_u^\mu(\phi,H)\).
In the class pairing for \(d_*y\), branching preserves the local
lattices. The level factor is a unit since \(g_i\ne T_u\), and the
argument of \ref{B4.22a} makes the stabilizer orders units since
\(g_i>2fN_m+1\). Thus \(d_*y\) is integral at \(w\). If \(y\ne0\), then
\[
 \frac{v_w(y)}{e_w}
 \ge-\frac{v_w(d_*)}{e_w}
 =-\sum_{j=1}^m\bigl(\nu_{i,u,j,w}+fN_{j-1}N_j\bigr)
 \ge-\sum_{j=1}^m\nu_{i,u,j,w}-\mathcal E_G .
\]
The last step uses \(N_{j-1},N_j\le N_m\) and \(\eta\ge0\).

For the numerical inequality set
\[
 A=f(N_m^2+m^2\eta),\qquad
 B=f\bigl(N_m(N_m+t)+m^2\eta\bigr).
\]
Above \(g_i\ne T_u\), the base term in \Lean{FDef} cancels the
base term of the loss in \ref{B4.03}. Hence
\(F_j(w)-E_{u,j}(w)=-\nu_{i,u,j,w}-c_gA-\delta_{u,w}\).
Since \(t\le N_m\), we have \(B\le2A\) and \(\mathcal E_G=mB\).
Thus
\[
 mc_gA+m\delta_{u,w}\ge 2mC_1A\ge C_1\mathcal E_G.
\]
Summation proves \Lean{SgAtG}.

For a selected family with \(0<\varepsilon\le1/100\), the middle size
threshold \(2\) gives \(c_g\ge2\). With \(A'=kb\ge10\)
and \(N=N_m\ge2\), the size conditions give
\[
 D_{u,j}\le3(A')^2N,\qquad
 g_i\ge\max\{(A')^8,N^{16}\}.
\]
For the degree bound, sum the row-entry bound \(3xb\) over at most
\(N\) entries at each of the \(f\le k\) embeddings. Here
\(x\le k\) and \(b\ge1\). If \(N\le(A')^5\), then
\(6(A')^2N+2\le6(A')^7+2<(A')^8\).
If \(N>(A')^5\), then \((A')^2<N\) and
\(6(A')^2N+2<6N^2+2<N^{16}\). Each of the three size expressions
in the statement is at most \(6(A')^2N+2\), since \(f\le k\).
The bounds therefore hold for both primes and all words beyond
the same threshold. The two-sided estimate for the log-weighted
sum of losses is \ref{B5.13}.
\end{proof}

Away from \(T_u\), \Lean{NuInM} gives, for each \(\phi\in N_u\),
a multiplier \(\rho\in\mathcal O\setminus\mathfrak p_w\) with
\(\rho\phi\in M_u^{\mathrm{ram}}\). By \Lean{PerColSmul} the complex
period is linear under scaling its column input. Semilinearity of
\(jN_u\) therefore gives
\[
 P_u^\mu(\rho\phi,H)=\iota(\rho)P_u^\mu(\phi,H)
 \qquad(\rho\in K).
\]

\begin{theorem}[Periods above the fixed primes]\label{B4.24}
Assume the two hypotheses of \ref{B4.23}, \Lean{NuInM} and
\Lean{PerColSmul}. For every \(u,w,i\) with \(p_w=g_i\ne T_u\),
every \(\phi\in N_u\), and every \(H\in\widetilde M_{\mu_w}\),
there is \(y\in K\) with
\[
 \iota(y)=P_u^{\mu_w}(\phi,H),\qquad
 y=0\ \text{or}\ v_w(y)\ge e_wS_{u,w}.
\]
\end{theorem}
\leanref{Formalism.\allowbreak{}L2.\allowbreak{}G.\allowbreak{}Adm2InU.\allowbreak{}admissible\_\allowbreak{}at\_\allowbreak{}generator\_\allowbreak{}places}{Formalism.\allowbreak{}LocalToGlobal.\allowbreak{}GroupLattices.\allowbreak{}AdmissibleAtGeneratorPlaces}{134--170}

\begin{proof}
Choose \(\rho\) as in \Lean{NuInM} and apply \ref{B4.23} to
\(\rho\phi\), with representative \(a\). Then \(y=\rho^{-1}a\)
represents the period of \(\phi\). Since \(v_w(\rho)=0\) and
\(B_{u,w}=S_{u,w}\), the valuation bound is unchanged.
\end{proof}

\section{All remaining finite places}

Use the local data \(X,Y,\boldsymbol\kappa,c\) of
\ref{B4.07}--\ref{B4.14}, with \(\boldsymbol\kappa\) the
\Lean{Consts} record. Put \(G_{p,j}=\operatorname{gainS}_{p,j}\)
as in \ref{B4.08}.

\begin{lemma}[Comparison of the two gains]\label{B4.25}
Suppose \Lean{ParametersFill} holds at every finite place and every
prime satisfying \(\mathrm{IsCs}\) belongs to \(\mathrm{drop}\).
Then \(\gainD_j(w)\le G_{p_w,j}\) for every \(w\) and
\(1\le j\le m\).
\end{lemma}
\leanref{Formalism.\allowbreak{}MainNotions.\allowbreak{}WeilJacquetShalikaLaw.\allowbreak{}gainD\_\allowbreak{}le\_\allowbreak{}gainS}{Formalism.\allowbreak{}SharpRestriction.\allowbreak{}GrowthBounds.\allowbreak{}WeilJacquetShalikaLaw}{552--576}
\leanref{Formalism.\allowbreak{}MainNotions.\allowbreak{}WeightDegreeSizes.\allowbreak{}gainLe\_\allowbreak{}of\_\allowbreak{}fills}{Formalism.\allowbreak{}SharpRestriction.\allowbreak{}Divisibility.\allowbreak{}WeightDegreeSizes}{931--934}

\begin{proof}
The orbit contributions to \(\mathrm{gainP}_{p,j}\) are nonnegative,
so \(G_{p,j}\ge0\). The screened gain is zero on \(\mathrm{drop}\).
The raw gain vanishes off \(L.\mathrm{supp}\) and outside the inert
and signal classes. In the remaining cases \Lean{ParametersFill}
identifies the raw gain with \(\mathrm{gainP}_{p_w,j}\).
A prime in either class that divides \(\mathrm{condN}\,g_0g_1\)
must be a fixed prime, where the raw gain vanishes. For
\(p_w\nmid\mathrm{condN}\,g_0g_1\) outside \(\mathrm{drop}\), the
screening assumption excludes \(\mathrm{IsCs}\). Both gains then
equal \(\mathrm{gainP}_{p_w,j}\).
\end{proof}

At normalization \(1\), \Lean{ErrDefs.\allowbreak{}edef} gives the column bounds
\[
 \mathrm{lowS}_{u,w}=B_{u,w}+\delta_{u,w},\qquad
 \mathrm{totS}_{u,w}=\sum_{j=1}^{m}
 \bigl(G_{p_w,j}+F_j(w)-E_{u,j}(w)\bigr)+\delta_{u,w}.
\]
By definition \(\mathrm{nuIn}_u(M_u^{\mathrm{ram}})\subseteq
S^{\mathrm{ram}}_{u,w}\), and \(\mu_w\) belongs to the column family
of \ref{B4.14}.

\begin{theorem}[The enhanced period bound away from \(T_u,g_0,g_1\)]\label{B4.26}
Assume the point and period hypotheses of \ref{B4.19},
\Lean{ErrDefs}, \Lean{Ties9} and the enhanced column bound of
\ref{B4.14} at normalization \(1\).
Assume \Lean{NuInM}, the gain comparison of \ref{B4.25}, and the
following local model conditions. If \(W_{p_w}\), the selected
orbit set, is nonempty and \(p_w\ne T_u\), a permitted index
\(h_g\) as in \ref{B4.07} exists. For every such index, with lattice
\(NP_{u,w,h_g}\) and comparison map \(c_{u,w,h_g}\), and every
\(\phi\in V_u\), assume
\[
 \phi\in(N_u)_{(w)}
 \quad\Longleftrightarrow\quad
 c_{u,w,h_g}(\mathrm{nuIn}_u\phi)\in NP_{u,w,h_g}.
\]
These conditions are \Lean{ServedHasModel} and \Lean{NuServedDef}.

Then for every \(u,w\) with \(p_w\notin\{T_u,g_0,g_1\}\), every
\(\phi\in N_u\), and every \(H\in\widetilde M_{\mu_w}\cap V_0(\mu_w)\),
there is \(y\in K\) with
\[
 \iota(y)=P_u^{\mu_w}(\phi,H),\qquad
 y=0\ \text{or}\ v_w(y)\ge e_wS_{u,w}.
\]
\end{theorem}
\leanref{Formalism.\allowbreak{}L2.\allowbreak{}G.\allowbreak{}Adm2InU.\allowbreak{}admissible\_\allowbreak{}off\_\allowbreak{}twist\_\allowbreak{}places}{Formalism.\allowbreak{}LocalToGlobal.\allowbreak{}GroupLattices.\allowbreak{}AdmissibleAtGeneratorPlaces}{171--197}

\begin{proof}
Put \(x=\mathrm{nuIn}_u\phi\). If \(W_{p_w}\ne\varnothing\), choose
a permitted \(h_g\). Since \(\phi\in(N_u)_{(w)}\), its image
\(c_{u,w,h_g}(x)\) belongs to \(NP_{u,w,h_g}\).
If \(W_{p_w}=\varnothing\), choose
\(\rho\in\mathcal O\setminus\mathfrak p_w\) by \Lean{NuInM}.
Then \(\mathrm{nuIn}_u(\rho\phi)\in S^{\mathrm{ram}}_{u,w}\)
and multiplication by the local unit \(\rho^{-1}\) gives
\(x\in S^{\mathrm{ram}}_{u,w}\). These are the two clauses of
\Lean{InNP} in \ref{B4.14}.

The enhanced theorem gives bound \(\mathrm{totS}_{u,w}\) for the
selected column. By \ref{B4.25}, this is at least
\(s=S_{u,w}+\delta_{u,w}\). Apply \ref{B4.19} with
\Lean{ErrDefs.\allowbreak{}dgam} to obtain the lower bound
\(e_w(s-\delta_{u,w})=e_wS_{u,w}\).
\end{proof}

\begin{theorem}[The complementary bound on the original lattice]\label{B4.27}
Assume the point and period hypotheses of \ref{B4.19},
\Lean{ErrDefs}, \Lean{FDef} and the base column bound
\Lean{ColumnValuationAwayFromTypePrimes} at normalization \(1\),
as in \ref{B4.13}--\ref{B4.14}. For every \(u,w\) with
\(p_w\notin\{T_u,g_0,g_1\}\), every \(\phi\in M_u^{\mathrm{ram}}\),
and every \(H\in\widetilde M_{\mu_w}\cap V_0(\mu_w)\), there is
\(y\in K\) with
\[
 \iota(y)=P_u^{\mu_w}(\phi,H),\qquad
 y=0\ \text{or}\ v_w(y)\ge e_wB_{u,w}.
\]
\end{theorem}
\leanref{Formalism.\allowbreak{}L2.\allowbreak{}G.\allowbreak{}Adm2InU.\allowbreak{}admissible\_\allowbreak{}away\_\allowbreak{}from\_\allowbreak{}twist}{Formalism.\allowbreak{}LocalToGlobal.\allowbreak{}GroupLattices.\allowbreak{}AdmissibleAtGeneratorPlaces}{198--213}

\begin{proof}
Since \(\mathrm{nuIn}_u\phi\in S^{\mathrm{ram}}_{u,w}\), the base
theorem gives column bound \(B_{u,w}+\delta_{u,w}\) at \(\mu_w\).
Apply \ref{B4.19} with \Lean{ErrDefs.\allowbreak{}dgam} to subtract
\(\delta_{u,w}\) from the normalized bound.
\end{proof}

\section{The all-place conclusion}

\begin{lemma}[Rounding the valuation budget]\label{B4.28}
Suppose that for every column \(u\), every finite place \(w\),
every \(\phi\in N_u\) and every
\(H\in\widetilde M_{\mu_w}\cap V_0(\mu_w)\), there is \(y\in K\) with
\[
 \iota(y)=P_u^{\mu_w}(\phi,H),\qquad
 y=0\ \text{or}\ v_w(y)\ge e_wS_{u,w}.
\]
Then the same representative is zero or has valuation at least
\(s_{u,w}=\lfloor e_wS_{u,w}\rfloor\).
The hypothesis is \Lean{UnroundedValuationBound} and the conclusion
is \Lean{adm2}.
\end{lemma}
\leanref{Formalism.\allowbreak{}MainNotions.\allowbreak{}LocalQuadraticParts.\allowbreak{}UnroundedValuationBound}{Formalism.\allowbreak{}SharpRestriction.\allowbreak{}Divisibility.\allowbreak{}LocalQuadraticParts}{212--217}
\leanref{Formalism.\allowbreak{}MainNotions.\allowbreak{}LocalQuadraticParts.\allowbreak{}adm2}{Formalism.\allowbreak{}SharpRestriction.\allowbreak{}Divisibility.\allowbreak{}LocalQuadraticParts}{206--211}
\leanref{Formalism.\allowbreak{}MainNotions.\allowbreak{}LocalQuadraticParts.\allowbreak{}valuation\_\allowbreak{}bound\_\allowbreak{}of\_\allowbreak{}unrounded}{Formalism.\allowbreak{}SharpRestriction.\allowbreak{}Divisibility.\allowbreak{}LocalQuadraticParts}{218--222}

\begin{proof}
The floor is at most \(e_wS_{u,w}\).
\end{proof}

Write \Lean{adm12} for the assertion that, for every \(u,w\) with
\(p_w\ne T_u\), \(\phi\in M_u^{\mathrm{ram}}\) and
\(H\in\widetilde M_{\mu_w}\cap V_0(\mu_w)\), there is a representative
\(y\in K\) of \(P_u^{\mu_w}(\phi,H)\) with \(y=0\) or
\(v_w(y)\ge e_wB_{u,w}\).

\begin{theorem}[Divisibility for fixed arithmetic data]\label{B4.29}
Fix \(D,L\) and local data
\(X,Y,\boldsymbol\kappa,c,D_x,T,\mathrm{lev},X_e,c_e\).
Assume the point and period hypotheses of \ref{B4.18} and the
following conditions.
\begin{enumerate}
 \item The data satisfy \Lean{ErrDefs} and \Lean{Ties9}, and
 \Lean{Serv} of \ref{B4.16} holds for some \(C\in\R\).
 It includes the enhanced and base column bounds at normalization \(1\).
 \item The level-prime hypotheses of \ref{B4.22} hold for every column.
 \item \Lean{NuInM}, \Lean{NuServedDef} and \Lean{ServedHasModel}
 hold. Also \Lean{ParametersFill} holds at every \(w\), and every
 \(\mathrm{IsCs}\) prime belongs to \(\mathrm{drop}\).
 \item The two fixed-prime hypotheses of \ref{B4.23} hold with
 one \(C_1\in\R\), and \Lean{PerColSmul} holds.
\end{enumerate}
Then \(\Lean{adm2}\) and \(\Lean{adm12}\) hold.
\end{theorem}
\leanref{Formalism.\allowbreak{}L2.\allowbreak{}G.\allowbreak{}Adm2InU.\allowbreak{}admissible\_\allowbreak{}level\_\allowbreak{}two}{Formalism.\allowbreak{}LocalToGlobal.\allowbreak{}GroupLattices.\allowbreak{}AdmissibleAtGeneratorPlaces}{214--246}
\leanref{Formalism.\allowbreak{}MainNotions.\allowbreak{}LocalQuadraticParts.\allowbreak{}adm12}{Formalism.\allowbreak{}SharpRestriction.\allowbreak{}Divisibility.\allowbreak{}LocalQuadraticParts}{323--329}

\begin{proof}
For \(\phi\in N_u\) and an admissible test vector at \(\mu_w\),
apply \ref{B4.22} with \Lean{ErrDefs.\allowbreak{}dgam} above \(T_u\).
For \(p_w\ne T_u\), use \ref{B4.24} above a fixed prime and
\ref{B4.26} elsewhere, with the gain comparison from \ref{B4.25}.
These bounds give \(\Lean{adm2}\) by \ref{B4.28}.

For \(\Lean{adm12}\), take \(\phi\in M_u^{\mathrm{ram}}\) and
\(p_w\ne T_u\). Above the fixed primes use \ref{B4.23}.
Elsewhere \Lean{Ties9.\allowbreak{}fdef} gives \Lean{FDef} and \Lean{Serv}
gives the base column bound, so \ref{B4.27} applies.
\end{proof}

\section{Uniformity in the selected family}

A parameter family \(\mathfrak F:\Lean{FamP}\,\varepsilon\) has members
\(D_{x,i},L_{x,i}\) for \(x>\mathfrak F.x_1\) and
\(i\in\mathfrak F.I_x\). Use closedness as in (4.46).
For a threshold record \(\mathcal K:\Lean{TransportedAlt}\),
the condition \(\mathfrak F.\mathrm{Large}(\mathcal K,C_H)\)
bounds the entries of \(c_{\rm Pre},c_{\rm Mid}(C_H),c_{\rm Last}\)
below by \(\mathcal K.\mathrm{pre},\mathcal K.\mathrm{mid}(c_{\rm Out}),
\mathcal K.\mathrm{last}(c_{\rm Out}++c_{\rm Mid}(C_H))\), respectively.
Here \(++\) denotes list concatenation.

\begin{theorem}[A common threshold for the family]\label{B4.30}
Let \(\varepsilon>0\) and \(\mathfrak F:\Lean{FamP}\,\varepsilon\).
Here \emph{eventually} means for all sufficiently large
\(x>\mathfrak F.x_1\) and all \(i\in\mathfrak F.I_x\).
Assume the following conditions.
\begin{enumerate}
 \item Eventual reference member-chain integrality.
 \item Eventual twisted point integrality for every used nonreference
 tuple.
 \item Eventual \Lean{PeriodClassSumWeyl}.
 \item Eventual \Lean{TwistedPeriodAdmissible}.
 \item There exist \(C\ge0\) and \(y_5\in\R\) such that every
 member with \(x\ge y_5\) and \(x>\mathfrak F.x_1\) satisfies
 \Lean{Serv} with this \(C\).
 \item Eventually the reference sharp bound and \Lean{TwistedColumnBoundAtT}
 hold for every \(u\).
 \item Eventually \(T_u\nmid\Delta_u\) for every \(u\).
 \item There is \(C_1\in\R\) such that the fixed-prime tuple bound and
 \Lean{SgAtG} each hold eventually with this same \(C_1\).
\end{enumerate}
Then there exists \(x_0\in\R\) such that both \(\Lean{adm2}\)
and \(\Lean{adm12}\) hold in every member with
\(x\ge x_0\), \(x>\mathfrak F.x_1\) and \(i\in\mathfrak F.I_x\).
If the eight inputs are closed on a class \(\mathcal C\), both
eventual conclusions are closed there. The closed inputs form
\(\Lean{Adm2InU}\,\mathcal C\).
\end{theorem}
\leanref{Formalism.\allowbreak{}L2.\allowbreak{}G.\allowbreak{}Adm2InU.\allowbreak{}famU}{Formalism.\allowbreak{}LocalToGlobal.\allowbreak{}GroupLattices.\allowbreak{}AdmissibleAtGeneratorPlaces}{263--301}
\leanref{Formalism.\allowbreak{}L2.\allowbreak{}G.\allowbreak{}Adm2InU.\allowbreak{}admissible\_\allowbreak{}family\_\allowbreak{}of\_\allowbreak{}inclusion}{Formalism.\allowbreak{}LocalToGlobal.\allowbreak{}GroupLattices.\allowbreak{}AdmissibleAtGeneratorPlaces}{315--339}

\begin{proof}
Let \(y_1,\ldots,y_4\) be thresholds for the point and period
hypotheses, \(y_5\) the service threshold, \(y_a,y_b\) the sharp-bound
and coprimality thresholds, and \(y_6,y_7\) the fixed-prime thresholds.
For \(e_z=\mathfrak F.\mathrm{ez}>0\), take
\[
 x_0=\max\{y_1,y_2,y_3,y_4,y_5,y_a,y_b,y_6,y_7,3^{5/e_z}\}.
\]
For \Lean{ServedHasModel} in such a member, suppose
\(p=p_w\ne T_u\) and \(W_p\ne\varnothing\). By \ref{B3.51}, \(p\)
is an inert-class candidate with \(\operatorname{ord}(h_p)\ge
x^{\varepsilon/4}\). An odd-order \(h_p\) would have zero image in
the cyclic \(2\)-group and order dividing \(o\). Since
\(o\le x^{\varepsilon/100}<x^{\varepsilon/4}\), the order is even.
The second component of arithmetic Frobenius at each prime above
\(p\) has this order.

The inert class excludes \(2\), and
\(\mathrm{splitBound}\ge\lfloor x^{e_z/5}\rfloor\ge3\) gives
\(\bigl(\frac{-d}{3}\bigr)=1\), which excludes \(3\).
Every host level prime is split, while \(p\) is inert. Hence \(p\)
lies outside the host level primes. The distinct primes \(p\) and
\(T_u\) also satisfy \(p\nmid T_u\). These are the local index
conditions of \ref{B4.07}.

The family supplies \Lean{NuInM}, \Lean{NuServedDef},
\Lean{NuAtTDef}, \Lean{ErrDefs} and \Lean{Ties9}.
Lemma~\ref{B4.20} gives
\Lean{LevelPrimeSizes}. An \(\mathrm{IsCs}\) prime satisfies
\Lean{NoGain} and has inert class, so it differs from every split
level prime and belongs to \(\mathrm{drop}\) by the family criterion.
With \Lean{ParametersFill} and the linearity of the period map,
all hypotheses of \ref{B4.29} hold. For closedness apply the
finite-conjunction rule of (4.46) to the eight inputs and use
the same threshold construction.
\end{proof}

\begin{theorem}[The selected family has divisible periods]\label{B4.31}
Let \(\mathcal A\) be a family class and
\(\mathrm{sel}:\Lean{TwistSelection}\,\mathcal A\) a twist selection.
Put
\[
 \mathcal D_3=\Lean{D3}\,3\,\mathcal A
 =\Lean{Des}\,(\Lean{WithChN}\,\mathcal A)\,3.
\]
Assume the eventual point, period and sharp-column hypotheses of
\ref{B4.30} are closed on \(\mathcal D_3\). Assume also that the
existence of \(C_1\in\R\) for which the fixed-prime tuple bound
and \Lean{SgAtG} hold eventually, uniformly in \(i\), is closed
on \(\mathcal D_3\). Suppose \(\Lean{SelectionServed}\) holds there
for the restriction of \(\mathrm{sel}\).

Then there is \(\varepsilon_1>0\) such that, for every
\(0<\varepsilon\le\varepsilon_1\), there is a size-threshold record
\(\mathcal K\) with the following property. For every
\(\mathfrak F:\Lean{FamP}\,\varepsilon\) in \(\mathcal D_3\) and
\(C_H\in\R\) with \(\mathfrak F.\mathrm{Large}(\mathcal K,C_H)\),
there exists \(x_0\in\R\) such that
\[
 \begin{gathered}
 \forall x\in\R,\quad
 (x\ge x_0\land x>\mathfrak F.x_1)\ \Longrightarrow\\
 \forall i\in\mathfrak F.I_x,\quad
 \forall u,w,\quad \forall\phi\in N_u,\quad
 \forall H\in\widetilde M_{\mu_w}\cap V_0(\mu_w),\\
 \exists y\in K,\qquad
 \iota(y)=P_u^{\mu_w}(\phi,H),\qquad
 y=0\ \text{or}\
 v_w(y)\ge
 \left\lfloor e_w
 \sum_{j=1}^{m}\bigl(\gainD_j(w)+F_j(w)-E_{u,j}(w)\bigr)
 \right\rfloor .
 \end{gathered}
\]
The field, lattices and periods in the display belong to
\(D_{x,i},L_{x,i}\).
\end{theorem}
\leanref{GlobalTheory.\allowbreak{}Top.\allowbreak{}closed\_\allowbreak{}family\_\allowbreak{}admissible\_\allowbreak{}twist\_\allowbreak{}link\_\allowbreak{}alt}{Development.\allowbreak{}AutomorphicArithmetic.\allowbreak{}ThetaCorrespondence.\allowbreak{}ThetaPairLifts.\allowbreak{}ThirdPlaceSeries.\allowbreak{}ClosedFamilyAdmissibleTwistLink}{12--31}

\begin{proof}
The definition of \(\mathcal D_3\) includes \(\Lean{DGamSize}\),
which combines with \(\Lean{SelectionServed}\) to give the closed
service \(\Lean{ServF}\) by \ref{B4.16}. The design condition gives
closed eventual coprimality by \ref{B4.21}. With the other closed
inputs this gives \(\Lean{Adm2InU}\,\mathcal D_3\). Apply
\ref{B4.30} and expand its closed eventual \(\Lean{adm2}\) conclusion.
\end{proof}

For the constructed family, the closed rows of \ref{B8.09} and the
fixed-prime row of \ref{B8.06} supply the arithmetic hypotheses of
\ref{B4.31}.


% ---- p4c4 ----
\chapter{Gluing row lattices across Hecke characters}\label{ch:hecke-row-gluing}

\section{Characters and the participating family}

\begin{definition}[Allowed finite-order characters]\label{B4.32}
Fix a number-field extension $E/F$ and a coefficient field $K$. Let $S\subseteq\N$. Put
\[
 \mathcal A=E\otimes_F\mathbb A_{F,f},\qquad
 \mathcal C=\operatorname{Hom}(\mathcal A^\times,K^\times).
\]
Let $R\subseteq\mathcal A$ be the subring generated by the images of $\mathcal O_E$ and $\prod_{v<\infty}\mathcal O_{F_v}$. For an ideal $J\subseteq R$, let $U(J)$ be the image in $\mathcal A^\times$ of $\ker(R^\times\to(R/J)^\times)$.

Let $v$ range over the nonzero prime ideals of $\mathcal O_F$. Write $Nv=\#(\mathcal O_F/v)$ and let $p(v)$ be the residue characteristic. Set
\[
 U_v=\{a\in\mathcal A^\times:a_{v'}=1\text{ for every }v'\ne v\},
 \qquad I_v=R^\times\cap U_v,\qquad
 W_v=U(\sqrt{p(v)R})\cap U_v.
\]
Here $a_{v'}$ denotes the projection to $(E\otimes_F F_{v'})^\times$. We identify $R^\times$ with its image in $\mathcal A^\times$.

Call $v$ good for $S$ if there are distinct primes $\mathfrak P,\mathfrak P'$ of $\mathcal O_E$ such that
\[
 v\mathcal O_E=\mathfrak P\mathfrak P',\qquad
 Nv\text{ is prime},\qquad Nv\notin v^2,\qquad p(v)\notin S.
\]
In the third condition, $Nv$ is viewed as an element of $\mathcal O_F$.

Let $\mathcal C_{\mathrm{fin}}\subseteq\mathcal C$ consist of the finite-order characters trivial on the image of $E^\times$ under $e\mapsto e\otimes1$ whose kernels contain $U(NR)$ for some positive integer $N$. Let $\mathcal T(E/F,K,S)\subseteq\mathcal C_{\mathrm{fin}}$ be the subgroup of characters trivial on the base ideles $1\otimes a$ and on every $W_v$, with trivial restriction to $I_v$ when $v$ fails the good-place condition. Its reference character is $1$.
\end{definition}
\leanref{Formalism.\allowbreak{}L2.\allowbreak{}DIV.\allowbreak{}finChar}{Formalism.\allowbreak{}LocalToGlobal.\allowbreak{}Divisibility.\allowbreak{}Tuples}{107--109}
\leanref{Formalism.\allowbreak{}L2.\allowbreak{}DIV.\allowbreak{}tupGrp}{Formalism.\allowbreak{}LocalToGlobal.\allowbreak{}Divisibility.\allowbreak{}Tuples}{112--115}
\leanref{Formalism.\allowbreak{}L2.\allowbreak{}DIV.\allowbreak{}Tup}{Formalism.\allowbreak{}LocalToGlobal.\allowbreak{}Divisibility.\allowbreak{}Tuples}{118--119}

\begin{definition}[The family used by the division construction]\label{B4.33}
For a number field $K$ and a finite subset $TU\subseteq\mathcal T(E/F,K,S)$, the row and column types form $\mathcal T_U=\{1\}\cup TU$. The division kernel datum gives a bijection $e:T\to\mathcal T_U$ from the abstract index set, with $e(t_0)=1$ at the reference index. Write $\tau_t=\operatorname{pr}_1(e(t))$ for the underlying character and $\sigma:K\hookrightarrow\C$ for the arithmetic embedding. For an idele $a$ of $E$, view its finite component $a_f$ in $\mathcal A^\times$. The Hecke character $\chi_{\theta(t)}$ of the arithmetic label $\theta(t)$ satisfies $\chi_{\theta(t)}(a)=\sigma(\tau_t(a_f))$.
\end{definition}
\leanref{Formalism.\allowbreak{}L2.\allowbreak{}DIV.\allowbreak{}TupF}{Formalism.\allowbreak{}LocalToGlobal.\allowbreak{}Divisibility.\allowbreak{}LiftLatInst}{58--59}
\leanref{Formalism.\allowbreak{}L2.\allowbreak{}DIV.\allowbreak{}colIn3}{Formalism.\allowbreak{}LocalToGlobal.\allowbreak{}Divisibility.\allowbreak{}HeckeCSDICommTrRep}{130--141}
\leanref{Formalism.\allowbreak{}L2.\allowbreak{}DIV.\allowbreak{}DivisionKernelData}{Development.\allowbreak{}AutomorphicArithmetic.\allowbreak{}ThetaCorrespondence.\allowbreak{}ThetaPairLifts.\allowbreak{}ThirdPlaceSeries.\allowbreak{}DIVRecBridgeSel}{32--57}
\leanref{GlobalTheory.\allowbreak{}Mem.\allowbreak{}Col2.\allowbreak{}Reads}{Development.\allowbreak{}AbelianVarietyTheory.\allowbreak{}Definitions.\allowbreak{}AutomorphicMembership.\allowbreak{}TrueColumnChoicesHeckeCharReferenceEqOne}{48--49}

\section{The common row lattice}

Put $\mathcal O=\mathcal O_K$. For a finite place $w$ of $K$, write $\mathfrak p_w$ for its prime ideal. The localization of an $\mathcal O$-submodule $L$ inside its ambient $K$-space is
\[
 L_{(w)}=\mathcal O_{\mathfrak p_w}L
 =\{x:\text{there is }a\in\mathcal O\setminus\mathfrak p_w
                       \text{ with }ax\in L\}.
\]
A full lattice is a finitely generated $\mathcal O$-submodule whose $K$-span is the ambient space.

\begin{definition}[Local choice and global intersection]\label{B4.34}
Fix $r\in\N$. For each $t\in T$, take $n_t,d_t\in\N$ and an $\mathcal O$-submodule $M_t\subseteq K^{n_t}$. Choose $K$-linear maps
\[
 \pi_t:K^{n_t}\longrightarrow K^{d_t},\qquad
 \iota_t:K^{d_t}\longrightarrow K^r.
\]
Put $\Lambda_t=\pi_t(M_t)$ and $L_t=\iota_t(\Lambda_t)$. Choose $t_w\in T$ at each finite place and put
\[
                      RR=\bigcap_{w<\infty}(L_{t_w})_{(w)}
                      \subseteq K^r.
\]
\end{definition}
\leanref{Formalism.\allowbreak{}L2.\allowbreak{}Indef.\allowbreak{}DefSide.\allowbreak{}abelian₂}{Formalism.\allowbreak{}LocalToGlobal.\allowbreak{}IndefiniteForms.\allowbreak{}IndefSide}{68--69}
\leanref{Formalism.\allowbreak{}L2.\allowbreak{}Indef.\allowbreak{}IndefData.\allowbreak{}RR}{Formalism.\allowbreak{}LocalToGlobal.\allowbreak{}IndefiniteForms.\allowbreak{}IndefSide}{281--283}
\leanref{Formalism.\allowbreak{}DivisionValues.\allowbreak{}glueLat}{Formalism.\allowbreak{}QuotientMetricDivisors.\allowbreak{}LatticeGlueUnique}{11--14}

\begin{theorem}[Gluing finitely many local changes]\label{B4.35}
Fix $r\in\N$ and a full lattice $L_0\subseteq K^r$. Choose a full lattice $L(w)\subseteq K^r$ at each finite place $w$. Suppose
\[
 \{w:(L(w))_{(w)}\ne(L_0)_{(w)}\}
\]
is finite. Then $G=\bigcap_w(L(w))_{(w)}$ is a full lattice and
\[
                         G_{(w)}=(L(w))_{(w)}
                         \qquad\text{for every }w.
\]
\end{theorem}
\leanref{Formalism.\allowbreak{}DivisionValues.\allowbreak{}lattice\_\allowbreak{}gluing}{Formalism.\allowbreak{}QuotientMetricDivisors.\allowbreak{}HermitianTwistDivisors.\allowbreak{}SecondStepGluing}{256--259}
\leanref{Formalism.\allowbreak{}DivisionValues.\allowbreak{}GluingInclusion.\allowbreak{}glue\_\allowbreak{}spec}{Formalism.\allowbreak{}QuotientMetricDivisors.\allowbreak{}HermitianTwistDivisors.\allowbreak{}SecondStepGluing}{74--183}

This is the existence and local recovery part of Voight's lattice correspondence \cite[Theorem 9.4.9]{voight-2021}, specialized to the Dedekind domain $\mathcal O$ and the ambient space $K^r$.
\begin{proof}
At each exceptional place $v$, clear denominators to choose nonzero $a_v,b_v\in\mathcal O$ with $a_vL_0\subseteq L(v)$ and $b_vL(v)\subseteq L_0$. Their products $a,b$ give $aL_0\subseteq G\subseteq b^{-1}L_0$, since $L_0=\bigcap_v(L_0)_{(v)}$. Thus $G$ spans $K^r$ and is finitely generated by Noetherianity.

Fix $w$ and $x\in L(w)$. Choose $0\ne c\in\mathcal O$ with $cx\in aL_0$. Factor $(c)=\mathfrak p_w^k\mathfrak d$ with $\mathfrak d$ prime to $\mathfrak p_w$, and take $s\in\mathfrak d\setminus\mathfrak p_w$. For $v\ne w$ we have $s/c\in\mathcal O_{\mathfrak p_v}$, so $sx=(s/c)cx\in(L(v))_{(v)}$. At $w$ we have $sx\in L(w)$. Hence $sx\in G$ and $x\in G_{(w)}$. Localizing gives $(L(w))_{(w)}\subseteq G_{(w)}$. The reverse inclusion follows from $G\subseteq (L(w))_{(w)}$.
\end{proof}

\begin{proposition}[Recovery of the chosen row lattice]\label{B4.36}
In Definition \ref{B4.34}, suppose each $\Lambda_t$ is full and each $\iota_t$ is surjective. Suppose $\{w:t_w\ne t_0\}$ is finite. Then $RR$ is a full lattice and
\[
                RR_{(w)}=(\iota_{t_w}(\Lambda_{t_w}))_{(w)}
                \qquad(w<\infty).
\]
\end{proposition}
\leanref{Formalism.\allowbreak{}L2.\allowbreak{}Indef.\allowbreak{}IndefRows.\allowbreak{}RR\_\allowbreak{}spec}{Formalism.\allowbreak{}LocalToGlobal.\allowbreak{}IndefiniteForms.\allowbreak{}IndefUse}{219--230}
\leanref{Formalism.\allowbreak{}L2.\allowbreak{}Indef.\allowbreak{}IndefColsR.\allowbreak{}RR\_\allowbreak{}spec}{Formalism.\allowbreak{}LocalToGlobal.\allowbreak{}IndefiniteForms.\allowbreak{}PairingConditionOfInverseMatrixIdentity}{332--343}
\begin{proof}
The image of a finite generating set of $\Lambda_t$ generates $L_t$, and
\[
 KL_t=\iota_t(K\Lambda_t)=\iota_t(K^{d_t})=K^r.
\]
Apply \ref{B4.35} with $L_0=L_{t_0}$ and $L(w)=L_{t_w}$. The localizations agree wherever $t_w=t_0$.
\end{proof}

\begin{lemma}[The type at conjugate places]\label{B4.36a}
In the family of \ref{B4.33}, let $\mathrm{Tw}$ be the image of a finite conjugation-stable set $W$ of arithmetic places among the finite places of $K$. Write $\bar w$ for the conjugate of $w$ under the arithmetic involution of $K$ and $p_w$ for its residue characteristic. Let $\delta,\delta'$ be the auxiliary Hecke characters, each of order $15$ on the local units at every place of $E$ above its respective rational prime $p_\delta,p_{\delta'}$. For $w\in\mathrm{Tw}$, take Hecke characters $\psi_w$ unramified above both primes with $\psi_{\bar w}=\psi_w$ under arithmetic conjugation. Set $\xi_w=\delta'\psi_w$ if $p_w=p_\delta$, and $\xi_w=\delta\psi_w$ otherwise. Assume each $\xi_w$ lies in the image of $t\mapsto\chi_{\theta(t)}$.

Both the row and column coordinate systems use the same selector $w\mapsto t_w$, with $\chi_{\theta(t_w)}=\xi_w$ on $\mathrm{Tw}$ and $t_w=t_0$ off $\mathrm{Tw}$. In each system suppose every $\Lambda_t$ is full and every $\iota_t$ is an isomorphism. Then
\[
 t_{\bar w}=t_w,\qquad w\in\mathrm{Tw}\ \Longleftrightarrow\ t_w\ne t_0.
\]
In either system, fix $w$ and put $t=t_w$. Then $RR_{(z)}=(L_t)_{(z)}$ for $z=w,\bar w$.
\end{lemma}
\leanref{GlobalTheory.\allowbreak{}Mem.\allowbreak{}Top.\allowbreak{}TwFin.\allowbreak{}tw\_\allowbreak{}fin}{Development.\allowbreak{}AbelianVarietyTheory.\allowbreak{}Proofs.\allowbreak{}IwasawaDecomposition.\allowbreak{}ResCharConj}{58--79}
\leanref{GlobalTheory.\allowbreak{}Mem.\allowbreak{}R3.\allowbreak{}tupL\_\allowbreak{}of\_\allowbreak{}twisted}{Development.\allowbreak{}AbelianVarietyTheory.\allowbreak{}Proofs.\allowbreak{}IwasawaDecomposition.\allowbreak{}MemHeightOneSpectrumIffAdditiveValuationPos}{82--189}
\leanref{Formalism.\allowbreak{}IndefR.\allowbreak{}Links.\allowbreak{}RealDat3.\allowbreak{}height\_\allowbreak{}change\_\allowbreak{}of\_\allowbreak{}D0}{Formalism.\allowbreak{}LocalToGlobal.\allowbreak{}IndefiniteReal.\allowbreak{}BaseLiftsPackageOfPrimitiveLinks}{105--109}
\leanref{Formalism.\allowbreak{}L2.\allowbreak{}Indef.\allowbreak{}IndefData.\allowbreak{}H0}{Formalism.\allowbreak{}LocalToGlobal.\allowbreak{}IndefiniteForms.\allowbreak{}IndefSide}{294--296}
\begin{proof}
The valuation of $a\in\mathcal O$ at the conjugate arithmetic place equals that of $\bar a$ at the original place. Thus the place correspondence respects conjugation, so $\mathrm{Tw}$ is stable and $t_{\bar w}=t_w=t_0$ off $\mathrm{Tw}$.

On $\mathrm{Tw}$, conjugation fixes $p_w$, so $\xi_{\bar w}=\xi_w$. At a place of $E$ above the selected auxiliary prime, $\psi_w$ is trivial on the local units and $\xi_w$ has order $15$ there. The surjection $a\mapsto a_f$ and the embedding $\sigma$ make the character map in \ref{B4.33} injective, with $\chi_{\theta(t_0)}=1$. Both places therefore select the same nonreference type.

By \ref{B4.36}, $RR_{(z)}=(L_{t_z})_{(z)}$ in each system. At $z=\bar w$, substitute $t_{\bar w}=t_w$.
\end{proof}

\section{Passing the period bound to coordinate rows}

Let $V_\mu$ be the complex row space for the arithmetic label $\mu$. Put $V_0(\mu)=\operatorname{im}(e_{\mathrm{op},\mu})$ for the idempotent row operator $e_{\mathrm{op},\mu}$. If $A\subseteq\mathcal O$ is the image of the arithmetic coefficient ring, the test set $\widetilde M_\mu$ is the $\sigma(A)$-span of the finite generator set $\mathcal G_\mu$.

\begin{lemma}[Integral decomposition of test vectors]\label{B4.37}
Fix $\mu$ and suppose
\[
 \sigma(a)y\in\widetilde M_\mu
 \qquad(a\in\mathcal O,\ y\in\widetilde M_\mu).
\]
For every $x\in V_0(\mu)$ admitting an expression
\[
                    x=\sum_{v\in\mathcal G_\mu}\sigma(c_v)v,
                    \qquad c_v\in\mathcal O,
\]
there are a finite set $S_x\subseteq\widetilde M_\mu\cap V_0(\mu)$ and coefficients $b_y\in\mathcal O$ such that
\[
                         x=\sum_{y\in S_x}\sigma(b_y)y.
\]
This decomposition property is denoted \Lean{CutUp}.
\end{lemma}
\leanref{Formalism.\allowbreak{}L2.\allowbreak{}G.\allowbreak{}ServJoint.\allowbreak{}CutUp}{Formalism.\allowbreak{}LocalToGlobal.\allowbreak{}GroupLattices.\allowbreak{}JointServices}{65--71}
\leanref{Formalism.\allowbreak{}L2.\allowbreak{}G.\allowbreak{}ServJoint.\allowbreak{}cutUp\_\allowbreak{}of\_\allowbreak{}stable}{Formalism.\allowbreak{}LocalToGlobal.\allowbreak{}GroupLattices.\allowbreak{}JointServices.\allowbreak{}Cut}{33--46}
\begin{proof}
Scalar stability and closure under sums give $x\in\widetilde M_\mu$. Take $S_x=\{x\}$ and $b_x=1$.
\end{proof}

For a finite place $w$ and $s\in\R$, put
\[
 \mathcal V_{w,s}=\{0\}\cup\{a\in K^\times:v_w(a)\ge s\},
\]
where $v_w(a)$ is the exponent of $\mathfrak p_w$ in $a\mathcal O$. For $a=b/c\in K^\times$ with $b,c\in\mathcal O$, both $v_w$ and the additive valuation in \ref{B4.29} give $\operatorname{ord}_{\mathfrak p_w}(b)-\operatorname{ord}_{\mathfrak p_w}(c)$. The valuation inequality for sums and $v_w(b)\ge0$ for $b\in\mathcal O\setminus\{0\}$ make $\mathcal V_{w,s}$ an $\mathcal O$-submodule.

\begin{proposition}[Period bounds on the coordinate row lattice]\label{B4.38}
Fix a column $u$. Write $V_u$ for the $K$-space containing $\mathrm{NuLat}_u$ and let $E_u:K^{a_u}\to V_u$ be the column bridge. For each $t\in T$, let $J_t:K^{n_t}\to V_{\theta(t)}$ be the row bridge.
Let $j_N$ embed $V_u$ into its complex realization and write $P_{u,t}(\phi,H)\in K$ for the coordinate bilinear pairing. Assume that for every $t\in T$ the identity
\[
 \sigma(P_{u,t}(\phi,H))
   =\operatorname{perCol}_u(\theta(t))(j_NE_u\phi,J_tH)
\]
holds for all $\phi\in K^{a_u}$ and $H\in K^{n_t}$. Assume also that for every $t\in T$ and $H\in M_t$,
\[
 J_tH\in V_0(\theta(t)),\qquad
 J_tH=\sum_{v\in\mathcal G_{\theta(t)}}\sigma(c_v)v
 \quad\text{for some }c_v\in\mathcal O.
\]
Fix a type $t$, a finite place $w$, a threshold $s\in\R$ and a subset $N\subseteq K^{a_u}$. Suppose \Lean{CutUp} holds for $\theta(t)$ and, for every $x\in N$ and $y\in\widetilde M_{\theta(t)}\cap V_0(\theta(t))$, there is $a\in\mathcal V_{w,s}$ with
\[
             \sigma(a)=
             \operatorname{perCol}_u(\theta(t))(j_NE_ux,y).
\]
Then
\[
 P_{u,t}(\phi,H)\in\mathcal V_{w,s}
 \qquad
 (\phi\in\operatorname{span}_{\mathcal O}N,\ H\in M_t).
\]

Let $\mu_0$ be the reference label and $\mu_\star,\mu'_\star$ the two auxiliary labels. Let $\mu_w$ be the arithmetic label used in the period bound of \ref{B4.29}. Assume the column bridge satisfies $\theta(t_w)=\mu_w$ at every finite place $w$. Suppose \Lean{CutUp} holds for every
$\mu\in\{\mu_0,\mu_\star,\mu'_\star\}\cup\{\mu_w:w<\infty\}$.
If $\mathcal N_u\subseteq K^{a_u}$ is an $\mathcal O$-submodule with $E_u(\mathcal N_u)\subseteq\mathrm{NuLat}_u$ and the period bound \Lean{adm2} of \ref{B4.29} holds, then
\[
 P_{u,t_w}(\phi,H)=0
 \quad\text{or}\quad
 v_w(P_{u,t_w}(\phi,H))\ge s_{u,w},
 \qquad
 s_{u,w}=\lfloor e_w\sg(u,w)\rfloor
\]
for every finite $w$ and for all $\phi\in\mathcal N_u$ and $H\in M_{t_w}$. Here $e_w$ is the ramification index of $w$ over its rational residue characteristic.
\end{proposition}
\leanref{Formalism.\allowbreak{}L2.\allowbreak{}G.\allowbreak{}ServJointF.\allowbreak{}ColBrF.\allowbreak{}Tied}{Formalism.\allowbreak{}LocalToGlobal.\allowbreak{}GroupLattices.\allowbreak{}ServJointF}{30--40}
\leanref{Formalism.\allowbreak{}L2.\allowbreak{}G.\allowbreak{}ServJointF.\allowbreak{}perToVal2F}{Formalism.\allowbreak{}LocalToGlobal.\allowbreak{}GroupLattices.\allowbreak{}ServJointF}{42--83}
\leanref{Formalism.\allowbreak{}L2.\allowbreak{}G.\allowbreak{}ServJointF.\allowbreak{}val2AllAtRF}{Formalism.\allowbreak{}LocalToGlobal.\allowbreak{}GroupLattices.\allowbreak{}ServJointF}{84--103}
\leanref{Formalism.\allowbreak{}L2.\allowbreak{}G.\allowbreak{}ServJoint.\allowbreak{}vAdd\_\allowbreak{}eq\_\allowbreak{}vw}{Formalism.\allowbreak{}LocalToGlobal.\allowbreak{}GroupLattices.\allowbreak{}JointServices.\allowbreak{}VAddEqVw}{37--52}
\leanref{Formalism.\allowbreak{}L2.\allowbreak{}G.\allowbreak{}ServJoint.\allowbreak{}perCol\_\allowbreak{}lin}{Formalism.\allowbreak{}LocalToGlobal.\allowbreak{}GroupLattices.\allowbreak{}JointServices.\allowbreak{}ValIntNonneg}{12--23}
\begin{proof}
Fix $H\in M_t$. Apply \Lean{CutUp} to obtain $J_tH=\sum_{y\in S}\sigma(b_y)y$ with $S\subseteq\widetilde M_{\theta(t)}\cap V_0(\theta(t))$ finite and $b_y\in\mathcal O$.
For $x\in N$, choose the period representative $a_y\in\mathcal V_{w,s}$ for each $y\in S$. The reference and twisted period formulas are complex-linear in the row vector. By the bridge identity,
\[
 \sigma(P_{u,t}(x,H))=\sum_{y\in S}\sigma(b_y)\sigma(a_y)
 =\sigma\left(\sum_{y\in S}b_ya_y\right).
\]
Injectivity of $\sigma$ gives $P_{u,t}(x,H)=\sum_{y\in S}b_ya_y\in\mathcal V_{w,s}$. Since $\phi\mapsto P_{u,t}(\phi,H)$ is $\mathcal O$-linear, the bound extends to $\operatorname{span}_{\mathcal O}N$.

For the selected type $t_w$, use $N=\mathcal N_u$ and $s=s_{u,w}$. The inclusion $E_u(\mathcal N_u)\subseteq\mathrm{NuLat}_u$ and $\theta(t_w)=\mu_w$ give the period representatives by \ref{B4.29}. The first part applies since $\mathcal N_u$ is an $\mathcal O$-submodule.
\end{proof}


% ---- p4c5 ----
\chapter{The pairing after division}\label{ch:pairing-after-division}

Let $K$ be the coefficient number field and put $\mathcal O=\mathcal O_K$.
For a finite place $w$, write $\mathfrak p_w$ for its prime ideal and
$v_w:K^\times\to\Z$ for the normalized additive valuation. For an
$\mathcal O$-submodule $M$ of a $K$-vector space, use the localization
\[
 M_{(w)}=\{x:\text{there is }t\in\mathcal O\setminus\mathfrak p_w
                    \text{ with }tx\in M\}.
\]
A lattice means a finitely generated $\mathcal O$-submodule spanning its
ambient $K$-vector space. For $R\subset W$ put
$R^\vee=\{\lambda\in W^*:\lambda(R)\subset\mathcal O\}$.

\section{The divided column module}

Fix the arithmetic data $D$ and the size data $L$.
For each word $u$ and finite place $w$, let
\[
 S_{u,w}=\sg(u,w),\qquad
 s_{u,w}=\lfloor e_w S_{u,w}\rfloor ,
 \qquad e_w=\text{the ramification index of }w\text{ over }\Q.
\]
Definition~\ref{B4.05} gives the budget $S_{u,w}$, and
Lemma~\ref{B4.28} passes from the period bound with threshold $e_wS_{u,w}$
to the bound with integer threshold $s_{u,w}$.
Both quantities vanish outside the finite set $L.\mathrm{supp}$.
Write $V_u$ for the column space and $N_u=\mathrm{NuLat}_u\subset V_u$
for its input submodule.

\begin{definition}[Division ideal and divided submodule]\label{B4.39}
For each word $u$, define
\[
 I_u=\prod_{w\in L.\mathrm{supp}}\mathfrak p_w^{-s_{u,w}},
 \qquad
 I_uN_u=\operatorname{span}_{\mathcal O}
       \{a\phi:a\in I_u,\ \phi\in N_u\}.
\]
The division condition for $(D,L)$ is the equality
$N'_u=I_uN_u$ for every $u$, where $N'_u=(D.\mathrm{div.Nu}\,u).L$
is the recorded column submodule.
\end{definition}
\leanref{Formalism.\allowbreak{}BasicNotions.\allowbreak{}DV.\allowbreak{}Iu}{Formalism.\allowbreak{}LocalRingCardinality.\allowbreak{}DataDIVDefs}{17--20}
\leanref{Formalism.\allowbreak{}BasicNotions.\allowbreak{}DV.\allowbreak{}NprimeDef}{Formalism.\allowbreak{}LocalRingCardinality.\allowbreak{}DataDIVDefs}{22--24}

At a place with $s_{u,w}>0$, this operation permits additional
denominators of order $s_{u,w}$. A negative $s_{u,w}$ imposes additional
divisibility.

\begin{proposition}[Realization of local shifts]\label{B4.40}
Let $N\subset K^r$ be a lattice, with $r\in\N$.
Fix a finitely supported integer function $w\mapsto\tau_w$. Put
$I=\prod_w\mathfrak p_w^{-\tau_w}$. Then $IN$ is a lattice. For every
$w$ there exists $a_w\in K^\times$ with $v_w(a_w)=-\tau_w$, and for
every such choice
\[
                  (IN)_{(w)}=(a_wN)_{(w)}.
\]
Moreover, any $\mathcal O$-submodule $N'\subset K^r$ with these
localizations equals $IN$.

More generally, if $J$ is any nonzero fractional ideal and
$a\in K^\times$ satisfies $v_w(a)=v_w(J)$, where $v_w(J)$ is the
exponent of $\mathfrak p_w$ in $J$, then $(JN)_{(w)}=(aN)_{(w)}$
for every $\mathcal O$-submodule $N\subset K^r$.
\end{proposition}
\leanref{Formalism.\allowbreak{}DivisionValues.\allowbreak{}localizationAt\_\allowbreak{}fractionalIdeal\_\allowbreak{}smul\_\allowbreak{}eq\_\allowbreak{}smul}{Formalism.\allowbreak{}QuotientMetricDivisors.\allowbreak{}HermitianTwistDivisors.\allowbreak{}TwistedLatticeNaturality}{547--572}
\leanref{Formalism.\allowbreak{}DivisionValues.\allowbreak{}shift\_\allowbreak{}existence}{Formalism.\allowbreak{}QuotientMetricDivisors.\allowbreak{}HermitianTwistDivisors.\allowbreak{}TwistedLatticeNaturality}{649--662}

Local principality follows from the discrete valuation ring criterion
\cite[Lemma~10.120.17, Tag~034X]{stacks-2026}.
Voight gives the local criterion and reconstruction for lattices over a
Dedekind domain \cite[Corollary~9.4.7 and Theorem~9.4.9]{voight-2021}.
The archive specializes reconstruction to prescribed shifts over
$\mathcal O_K$ and proves uniqueness among all submodules. For uniqueness,
apply \cite[Lemma~10.23.1, Tag~00HN]{stacks-2026} to $(N'+IN)/N'$ and
$(N'+IN)/IN$.

For arbitrary $\mathcal O$-submodules $N,N'$ and any integer function
$\tau$, call $N'$ a \emph{shift of $N$ by $\tau$} if for every $w$
there is $a_w\in K^\times$ of valuation $-\tau_w$ such that
$N'_{(w)}=(a_wN)_{(w)}$.
\leanref{Formalism.\allowbreak{}DivisionValues.\allowbreak{}IsShiftAt}{Formalism.\allowbreak{}QuotientMetricDivisors.\allowbreak{}DivIntoDualShift}{18--21}
\leanref{Formalism.\allowbreak{}DivisionValues.\allowbreak{}IsShift}{Formalism.\allowbreak{}QuotientMetricDivisors.\allowbreak{}DivIntoDualShift}{23--24}

\begin{lemma}[The arithmetic and coordinate shifts agree]\label{B4.41}
Fix $u$ and a $K$-linear isomorphism $e:K^r\to V_u$.
Let $N=e^{-1}(N_u)$ and let $N^c\subset K^r$ be an $\mathcal O$-submodule
with $e(N^c)=N'_u$.
Assume the division condition of Definition~\ref{B4.39}.
At the prime corresponding to $w$, put
$\mathrm{sgRec}_u(w)=s_{u,w}$ viewed in $\R$.
For a real-valued function $g$ on the finite places, put
$d_w=\mathrm{shiftZ}(g)(w)=\lfloor g(w)\rfloor$. When $d$ has finite support,
$\mathrm{shiftLat}(N,g)=(\prod_{d_w\ne0}\mathfrak p_w^{-d_w})N$.
Then
\[
 N^c=\mathrm{shiftLat}(N,\mathrm{sgRec}_u)
     =I_uN,\qquad
 \mathrm{shiftZ}(\mathrm{sgRec}_u)(w)=s_{u,w}.
\]
In particular $N^c$ is the shift of $N$ by $w\mapsto s_{u,w}$.
\end{lemma}
\leanref{Formalism.\allowbreak{}L2.\allowbreak{}DIV.\allowbreak{}sgRec}{Formalism.\allowbreak{}LocalToGlobal.\allowbreak{}Divisibility.\allowbreak{}RecBridge}{387--388}
\leanref{Formalism.\allowbreak{}L2.\allowbreak{}DIV.\allowbreak{}shiftZ}{Formalism.\allowbreak{}LocalToGlobal.\allowbreak{}Divisibility.\allowbreak{}Inputs}{381--381}
\leanref{Formalism.\allowbreak{}L2.\allowbreak{}DIV.\allowbreak{}shiftLat}{Formalism.\allowbreak{}LocalToGlobal.\allowbreak{}Divisibility.\allowbreak{}Inputs}{383--385}
\leanref{Formalism.\allowbreak{}L2.\allowbreak{}DIV.\allowbreak{}Cons.\allowbreak{}shift\_\allowbreak{}nuOf}{Formalism.\allowbreak{}LocalToGlobal.\allowbreak{}Divisibility.\allowbreak{}DivisibilityConstruction.\allowbreak{}ShiftOfColTie}{53--79}

\begin{proof}
Since $s_{u,w}$ is an integer, the $w$-factor of the shift ideal is
$\mathfrak p_w^{-s_{u,w}}$. Their product is $I_u$ by Definition~\ref{B4.39}.
By $K$-linearity, $e(I_uN)=I_ue(N)=I_uN_u=N'_u=e(N^c)$.
Since $e$ is injective, $I_uN=N^c$. At $w$, a power of a local uniformizer gives
$a_w\in K^\times$ with $v_w(a_w)=-s_{u,w}=v_w(I_u)$.
Proposition~\ref{B4.40} then gives $N^c_{(w)}=(a_wN)_{(w)}$.
\end{proof}

\section{Cancellation and local integrality}

\begin{lemma}[Cancellation on a shifted column]\label{B4.42}
Fix a finite place $w$ and $s\in\Z$. Let $N,N'\subset K^r$, $r\in\N$,
be $\mathcal O$-submodules with
$N'_{(w)}=(aN)_{(w)}$ for some $a\in K^\times$ of valuation $-s$.
Let $f:K^r\to K$ be $K$-linear. If
\[
 \forall\psi\in N,\qquad f(\psi)=0
                     \quad\text{or}\quad v_w(f(\psi))\ge s,
\]
then for every $\phi\in N'$ one has
$f(\phi)=0$ or $v_w(f(\phi))\ge0$.
\end{lemma}
\leanref{Formalism.\allowbreak{}DivisionValues.\allowbreak{}shift\_\allowbreak{}pairing\_\allowbreak{}formula}{Formalism.\allowbreak{}QuotientMetricDivisors.\allowbreak{}HermitianTwistDivisors.\allowbreak{}TwistedLatticeNaturality}{684--718}

\begin{proof}
For $\phi\in N'$, the equality of localizations gives
$t\in\mathcal O\setminus\mathfrak p_w$ and $\psi\in N$ such that
$t\phi=a\psi$. Hence
\[
                    t f(\phi)=a f(\psi),\qquad v_w(t)=0.
\]
If $f(\psi)=0$, then $f(\phi)=0$. If $f(\phi)\ne0$, both sides are
nonzero and
\[
 v_w(f(\phi))=v_w(a)+v_w(f(\psi))-v_w(t)
             =-s+v_w(f(\psi))\ge0.
\]
\end{proof}

\begin{lemma}[From a local row pairing to global integers]\label{B4.43}
Fix $w$, integers $n,h\ge0$, an $\mathcal O$-submodule $M\subset K^n$,
and $K$-linear maps $\pi:K^n\to K^h$, $\iota:K^h\to K^h$.
Let $b:K^h\to K$ and $P:K^n\to K$ be $K$-linear.
Fix $\zeta\in K^\times$ with $v_w(\zeta)=0$. Suppose
$P(F)=\zeta\,b(\iota\pi F)$ for every $F\in K^n$.
For any $c\in\Z$, if $P(F)=0$ or $v_w(P(F))\ge c$ for every $F\in M$,
then
\[
 x\in(\iota\pi M)_{(w)}
 \quad\Longrightarrow\quad
 b(x)=0\ \text{or}\ v_w(b(x))\ge c.
\]

For the global assertion, fix $h$, a $K$-linear functional $b:K^h\to K$,
and an $\mathcal O$-submodule $R\subset K^h$. Suppose the local hypotheses
hold at every $w$ with this $b$, source dimension $n_w$, and $c=0$.
If $R_{(w)}=(\iota_w\pi_wM_w)_{(w)}$ for every $w$, then
$b(R)\subset\mathcal O$.
\end{lemma}
\leanref{Formalism.\allowbreak{}DivisionValues.\allowbreak{}reduction\_\allowbreak{}first\_\allowbreak{}step}{Formalism.\allowbreak{}QuotientMetricDivisors.\allowbreak{}HermitianTwistDivisors.\allowbreak{}VwEqZeroOfNotMem}{34--60}
\leanref{Formalism.\allowbreak{}DivisionValues.\allowbreak{}divisor\_\allowbreak{}into\_\allowbreak{}dual\_\allowbreak{}lattice}{Formalism.\allowbreak{}QuotientMetricDivisors.\allowbreak{}HermitianTwistDivisors.\allowbreak{}DivIntoDual}{49--61}

\begin{proof}
For $x\in(\iota\pi M)_{(w)}$, choose
$t\in\mathcal O\setminus\mathfrak p_w$ and $F\in M$ with
$tx=\iota\pi F$. Then $P(F)=\zeta\,t\,b(x)$. If $b(x)\ne0$, this is
nonzero, and $v_w(\zeta)=v_w(t)=0$ gives
$v_w(b(x))=v_w(P(F))\ge c$.

Take $x\in R$. The first assertion gives
$v_w(b(x))\ge0$ at every finite place when $b(x)\ne0$.
The principal fractional ideal $(b(x))$ therefore has nonnegative
exponent at every prime, so it is an integral ideal and
$b(x)\in\mathcal O$. The value zero also belongs to $\mathcal O$.
\end{proof}

\section{The pairing into the common dual}

\begin{theorem}[Integral pairing after division]\label{B4.44}
Fix $r,h\in\N$ and a $K$-linear map $B:K^r\to(K^h)^*$.
Take $\mathcal O$-submodules $N,N'\subset K^r$ and $R\subset K^h$.
For each finite place $w$, fix $n_w\in\N$ and an $\mathcal O$-submodule
$M_w\subset K^{n_w}$. Let
\[
 \pi_w:K^{n_w}\to K^h,\quad \iota_w:K^h\to K^h,\quad
 P_w:K^r\times K^{n_w}\to K
\]
be $K$-linear maps and a $K$-bilinear pairing, respectively.
Suppose $\zeta_w\in K^\times$ has $v_w(\zeta_w)=0$ and
\[
 P_w(\phi,F)=\zeta_w B(\phi)(\iota_w\pi_wF)
                   \quad(\phi\in K^r,\ F\in K^{n_w}),
 \qquad
 R_{(w)}=(\iota_w\pi_wM_w)_{(w)}.
\]
Let $\sigma_w\in\R$ and $\tau_w\in\Z$ be arbitrary functions of $w$.
Suppose $N'$ is a shift of $N$ by $\tau$, with
$\tau_w\le\lceil\sigma_w\rceil$ at every $w$. Assume at each $w$ that
\[
 \phi\in N,\ F\in M_w
 \quad\Longrightarrow\quad
 P_w(\phi,F)=0\ \text{or}\ v_w(P_w(\phi,F))\ge\sigma_w.
\]
Then $B(N')\subset R^\vee$.
\end{theorem}
\leanref{Formalism.\allowbreak{}DivisionValues.\allowbreak{}divisor\_\allowbreak{}into\_\allowbreak{}dual\_\allowbreak{}shift}{Formalism.\allowbreak{}QuotientMetricDivisors.\allowbreak{}HermitianTwistDivisors.\allowbreak{}TwistedLatticeNaturality}{720--730}

\begin{proof}
Fix $w$ and $F\in M_w$. Each nonzero value of
$f_F(\phi)=P_w(\phi,F)$ on $N$ has integer valuation, hence
\[
 v_w(f_F(\phi))\ge\lceil\sigma_w\rceil\ge\tau_w.
\]
Apply Lemma~\ref{B4.42} to $f_F$. Fix any $\phi\in N'$.
The lemma gives $P_w(\phi,F)=0$ or $v_w(P_w(\phi,F))\ge0$
for every $w$ and $F\in M_w$.
Lemma~\ref{B4.43} with $b=B(\phi)$ gives $B(\phi)(R)\subset\mathcal O$.
\end{proof}

\begin{theorem}[The selected characters give one dual lattice]\label{B4.45}
Use the column record and the row data of Proposition~\ref{B4.36},
with selected index $t_w$ and common coordinate row lattice $R^c\subset K^h$.
Put $A_w=\iota_{t_w}\pi_{t_w}$, where $\pi_t$ is the row map of
Definition~\ref{B4.34}. Proposition~\ref{B4.36} gives
$R^c_{(w)}=(A_wM_{t_w})_{(w)}$.
For a column index $i$, let $B^c_i:K^{a_i}\to(K^h)^*$ be its coordinate
column map. The column record gives
\[
 P_{i,t_w}(\phi,F)=\zeta_{i,t_w}B^c_i(\phi)(A_wF),
 \qquad \zeta_{i,t_w}\in K^\times,\quad v_w(\zeta_{i,t_w})=0.
\]
Let $N,N'\subset K^{a_i}$ be $\mathcal O$-submodules and let
$\sigma_w\in\R$, $\tau_w\in\Z$. Suppose $N'$ is a shift of $N$ by
$\tau$, with $\tau_w\le\lceil\sigma_w\rceil$ at every $w$. Assume at each $w$ that
\[
 \phi\in N,\ F\in M_{t_w}
 \quad\Longrightarrow\quad
 P_{i,t_w}(\phi,F)=0\ \text{or}\ v_w(P_{i,t_w}(\phi,F))\ge\sigma_w.
\]
Then $B^c_i(N')\subset (R^c)^\vee$.

For the arithmetic application, assume $\theta(t_w)=\mu_w$ at every
finite place $w$, with the arithmetic labels of Definition~\ref{B4.33}.
For an arithmetic word $u$, write $B^c_u$ for the coordinate map at
$\mathrm{wU}(u)$ and $a_u$ for its domain dimension. Let $B_u:V_u\to W^*$
be the period map and $RR\subset W$ the common row lattice in the
arithmetic row space $W$. Fix $K$-linear isomorphisms $E_u:K^{a_u}\to V_u$
and $E_R:K^h\to W$, and an $\mathcal O$-submodule $N^c_u\subset K^{a_u}$.
Assume
\[
 E_u(N^c_u)=N'_u,\qquad E_R(R^c)=RR,\qquad
 B_u(E_u x)(E_R y)=B^c_u(x)(y)\quad(x\in K^{a_u},\ y\in K^h).
\]
Under the division condition of Definition~\ref{B4.39} and the hypotheses
of Proposition~\ref{B4.38} for $E_u$ and the column at $\mathrm{wU}(u)$,
one has $B_u(N'_u)\subset RR^\vee$.
\end{theorem}
\leanref{Formalism.\allowbreak{}L2.\allowbreak{}Indef.\allowbreak{}IndefColsR.\allowbreak{}intoDual}{Formalism.\allowbreak{}LocalToGlobal.\allowbreak{}IndefiniteForms.\allowbreak{}PairingConditionOfInverseMatrixIdentity}{393--408}
\leanref{Formalism.\allowbreak{}L2.\allowbreak{}DIV.\allowbreak{}oneMap\_\allowbreak{}of\_\allowbreak{}coord}{Formalism.\allowbreak{}LocalToGlobal.\allowbreak{}Divisibility.\allowbreak{}RecBridge}{33--50}

\begin{proof}
Apply Theorem~\ref{B4.44} with $R=R^c$, $B=B^c_i$, $M_w=M_{t_w}$,
$\pi_w=A_w$, $\iota_w=\mathrm{id}$ and $P_w=P_{i,t_w}$.

Put $N=E_u^{-1}(N_u)$.
Proposition~\ref{B4.38} gives the bound on $N\times M_{t_w}$ with
$\sigma_w=s_{u,w}$ viewed in $\R$. By Lemma~\ref{B4.41}, $N^c_u$
is the shift of $N$ by $\tau_w=s_{u,w}=\lceil\sigma_w\rceil$.
The first assertion gives $B^c_u(N^c_u)\subset(R^c)^\vee$.
For $\phi\in N'_u$ and $r\in RR$, take $x\in N^c_u$ and $y\in R^c$
mapping to $\phi$ and $r$. Then $B_u(\phi)(r)=B^c_u(x)(y)\in\mathcal O$.
\end{proof}


% ---- p4c6 ----
\chapter{Slope change and what is handed to step 5}\label{ch:p4c6}

Write \(n_K=[K:\Q]\) and, at a finite place \(w\), put
\[
 \ell_w=\frac{\log N\mathfrak p_w}{n_K},\qquad
 e_w=e(\mathfrak p_w/p_w),\qquad p_w=\operatorname{resChar}(w).
\]
Normalize \(v_w\) to take the value \(1\) on a uniformizer.
Use \(S_{u,w},s_{u,w},N_u,I_u\) from \ref{B4.39} and give \(N_u\)
the ambient basis and frames of \(N'_u\), as in \(\Lean{NuH}\).

\section{The exact change of slope}

Let \(V\) be a finite-dimensional space over a totally complex number
field \(K\), with basis \(b\) and invertible frames \(X_v\) at infinity.
For \(y\in V\), let \(\rho_{b,X}(y)\) collect the real and imaginary
parts of \(X_v^{-1}v([y]_b)\) over those places, where \([y]_b\) is the
coordinate vector of \(y\) and \(v\) acts by its chosen embedding.
For a nonzero finitely generated \(\mathcal O_K\)-submodule \(M\subset V\),
put
\[
 \mu_{b,X}(M)
 =2\,\operatorname{slope}(\rho_{b,X}M)-\log2+
 \frac{\log|\operatorname{disc}K|}{n_K}.
\]
Here \(\operatorname{slope}\) is minus the log covolume in the real
span divided by the \(\Z\)-rank.

\begin{lemma}[Multiplication by an integral ideal]\label{B4.46}
Let \(M\subset V\) be a nonzero finitely generated
\(\mathcal O_K\)-module, and let \(\mathfrak b\subset\mathcal O_K\)
be a nonzero ideal. Then
\[
 \mu_{b,X}(\mathfrak bM)
 =\mu_{b,X}(M)-\frac{2}{n_K}\log N\mathfrak b.
\]
\end{lemma}
\leanref{Formalism.\allowbreak{}MainNotions.\allowbreak{}J.\allowbreak{}mu\_\allowbreak{}smul\_\allowbreak{}ideal}{Formalism.\allowbreak{}SharpRestriction.\allowbreak{}AdmissibleParameterJoins.\allowbreak{}SumLogPrimeFactorsLe}{30--72}

\begin{lemma}[Multiplication by a fractional ideal]\label{B4.47}
Let \(I\) be a nonzero fractional ideal of \(K\), and let
\(M,N\subset V\) be \(\mathcal O_K\)-submodules with
\[
 N=IM=\operatorname{span}_{\mathcal O_K}
 \{ay:a\in I,\ y\in M\}.
\]
Assume that \(N\) is finitely generated and nonzero. Then
\[
 \mu_{b,X}(N)=\mu_{b,X}(M)-\frac{2}{n_K}\log N(I),
\]
where \(N(I)\in\Q_{>0}\) is the absolute fractional-ideal norm.
\end{lemma}
\leanref{Formalism.\allowbreak{}MainNotions.\allowbreak{}J.\allowbreak{}mu\_\allowbreak{}frac}{Formalism.\allowbreak{}SharpRestriction.\allowbreak{}AdmissibleParameterJoins.\allowbreak{}SumLogPrimeFactorsLe}{73--171}

In \ref{B4.47}, \(M=I^{-1}N\) is nonzero and finitely generated.
Give \(I\) the absolute-value metric at each infinite place. Its normalized
Arakelov degree is \(-\log N(I)/n_K\). Both lemmas are the ideal case of
Bost's tensor formula \cite[Proposition 5.5(3)]{fonseca-2019}, which gives
the slope change \(-2\log N(I)/n_K\) in this normalization.

\begin{lemma}[Norm of the division ideal]\label{B4.48}
For every word \(u\), the fractional ideal \(I_u\) is nonzero and
\[
 \log N(I_u)
 =-\sum_{w\in L.\mathrm{supp}}s_{u,w}\log N\mathfrak p_w.
\]
\end{lemma}
\leanref{Formalism.\allowbreak{}MainNotions.\allowbreak{}J.\allowbreak{}Iu\_\allowbreak{}ne\_\allowbreak{}zero}{Formalism.\allowbreak{}SharpRestriction.\allowbreak{}AdmissibleParameterJoins.\allowbreak{}SumLogPrimeFactorsLe}{352--357}
\leanref{Formalism.\allowbreak{}MainNotions.\allowbreak{}J.\allowbreak{}log\_\allowbreak{}absNorm\_\allowbreak{}Iu}{Formalism.\allowbreak{}SharpRestriction.\allowbreak{}AdmissibleParameterJoins.\allowbreak{}SumLogPrimeFactorsLe}{358--379}

\begin{proof}
Each factor \(\mathfrak p_w^{-s_{u,w}}\) is nonzero and has norm
\((N\mathfrak p_w)^{-s_{u,w}}\). Multiply the norms and take logarithms.
\end{proof}

\begin{theorem}[Slope difference in the arithmetic data]\label{B4.49}
Assume \(\Lean{NprimeDef}\), so that \(N'_u=I_uN_u\) for every \(u\).
Fix a word \(u\). Suppose that the frames of \(N'_u\) are invertible
at every infinite place and that \(N'_u\) is finitely generated and
nonzero. Then
\[
 \Delta\mu_u:=\mu(N'_u)-\mu(N_u)
 =2\sum_{w\in L.\mathrm{supp}}\ell_ws_{u,w}.
\]
\end{theorem}
\leanref{Formalism.\allowbreak{}MainNotions.\allowbreak{}J.\allowbreak{}mu\_\allowbreak{}diff}{Formalism.\allowbreak{}SharpRestriction.\allowbreak{}AdmissibleParameterJoins.\allowbreak{}SumLogPrimeFactorsLe}{380--398}
\leanref{Formalism.\allowbreak{}BasicNotions.\allowbreak{}SizeData₂.\allowbreak{}NuH}{Formalism.\allowbreak{}LocalRingCardinality.\allowbreak{}DataParams₂}{236--237}

\begin{proof}
By \ref{B4.47} and \ref{B4.48},
\[
 \Delta\mu_u=-\frac{2}{n_K}\log N(I_u)
 =2\sum_{w\in L.\mathrm{supp}}\ell_ws_{u,w}.
\]
\end{proof}

For a full lattice \(N\subset K^d\) over a totally complex field,
let \(\operatorname{covol}_X(N)\) be its Euclidean covolume in the
frames \(X\). Its normalized logarithmic covolume is
\[
 \operatorname{lcv}_X(N)=\frac{2}{n_K}
 \log\!\left(
 2^{r_2d}|\operatorname{disc}K|^{-d/2}
 \operatorname{covol}_X(N)\right),
 \qquad r_2=\frac{n_K}{2}.
\]
Write \(\operatorname{idx}_w(N,N')\) for the valuation of the determinant
of a local basis of \(N'\) relative to one of \(N\).
A shift by \(\tau_w\in\Z\) at \(w\) means
\[
 N'_{(w)}=a_wN_{(w)},\qquad
 a_w\in K^\times,\qquad v_w(a_w)=-\tau_w.
\]

\begin{theorem}[Slope difference in coordinates]\label{B4.50}
Let \(K\) be a number field with the given involution, \(d\ge0\),
and \(N\subset K^d\) a full \(\mathcal O_K\)-lattice. If an
\(\mathcal O_K\)-submodule \(N'\subset K^d\) is a local shift of \(N\)
by \(\tau_w\in\Z\) at \(w\), then
\[
 \operatorname{idx}_w(N,N')=-d\tau_w.
\]
Take \(K\) totally complex and fix invertible frames \(X_v\).
For \(d>0\), one has \(\mu_X(N)=-\operatorname{lcv}_X(N)/d\).
If \(N'\) is full and the shift holds at every finite place,
then for \(d>0\) the shifts have finite support and
\[
 \begin{aligned}
 \operatorname{lcv}_X(N')-\operatorname{lcv}_X(N)
 &=-2d\sum_w\ell_w\tau_w,\\
 \mu_X(N')-\mu_X(N)&=2\sum_w\ell_w\tau_w.
 \end{aligned}
\]
For \(d=0\), both normalized logarithmic covolumes are zero.
\end{theorem}
\leanref{Formalism.\allowbreak{}DivisionValues.\allowbreak{}shift\_\allowbreak{}index\_\allowbreak{}formula}{Formalism.\allowbreak{}QuotientMetricDivisors.\allowbreak{}HermitianTwistDivisors.\allowbreak{}TwistedLatticeNaturality}{90--98}
\leanref{Formalism.\allowbreak{}DivisionValues.\allowbreak{}mu\_\allowbreak{}log\_\allowbreak{}conductor}{Formalism.\allowbreak{}QuotientMetricDivisors.\allowbreak{}HermitianTwistDivisors.\allowbreak{}TwistedLatticeNaturality}{101--129}
\leanref{Formalism.\allowbreak{}DivisionValues.\allowbreak{}shift\_\allowbreak{}slope\_\allowbreak{}mu}{Formalism.\allowbreak{}QuotientMetricDivisors.\allowbreak{}HermitianTwistDivisors.\allowbreak{}TwistedLatticeNaturality}{145--183}

\begin{proof}
The local change of basis is \(a_w\) times the identity, hence its
determinant has valuation \(-d\tau_w\). Full lattices agree locally away
from the divisors of a common denominator, so the relative indices have
finite support, as do the shifts if \(d>0\).
Factoring the indices relative to a common sublattice gives
\[
 \operatorname{lcv}_X(N')-\operatorname{lcv}_X(N)
 =2\sum_w\ell_w\operatorname{idx}_w(N,N').
\]
The real realization has \(\Z\)-rank \(n_Kd\), which gives
\(\mu_X=-\operatorname{lcv}_X/d\) for \(d>0\). In dimension zero
the covolume and its normalizing factor are \(1\).
\end{proof}

\section{Rounding and the rational primes}

Define the unrounded budget by
\[
 A_u=\sum_{w\in L.\mathrm{supp}}2\ell_we_wS_{u,w}.
\]

\begin{lemma}[Rounding error]\label{B4.51}
Under the hypotheses of \ref{B4.49},
\[
 \left|\Delta\mu_u-A_u\right|
 \le 2\sum_{\substack{w\in L.\mathrm{supp}\\S_{u,w}\ne0}}\ell_w.
\]
\end{lemma}
\leanref{Formalism.\allowbreak{}MainNotions.\allowbreak{}J.\allowbreak{}round\_\allowbreak{}le}{Formalism.\allowbreak{}SharpRestriction.\allowbreak{}AdmissibleParameterJoins.\allowbreak{}SumLogPrimeFactorsLe}{401--430}

\begin{proof}
By \ref{B4.49}, the difference is the sum of the rounding terms
\(2\ell_w(\lfloor e_wS_{u,w}\rfloor-e_wS_{u,w})\). Each vanishes
when \(S_{u,w}=0\) and has absolute value at most \(2\ell_w\).
\end{proof}

\begin{lemma}[Total weight above a rational prime]\label{B4.52}
Let \(K\) be any number field and \(p\) any rational prime.
Let \(W_p\) be the finite set of all finite places of \(K\) above \(p\),
and let \(f_w=[\mathcal O_K/\mathfrak p_w:\F_p]\). Then
\[
 \sum_{w\in W_p}e_wf_w=n_K,\qquad
 \sum_{w\in W_p}e_w\ell_w=\log p.
\]
\end{lemma}
\leanref{Formalism.\allowbreak{}MainNotions.\allowbreak{}J.\allowbreak{}sum\_\allowbreak{}e\_\allowbreak{}f}{Formalism.\allowbreak{}SharpRestriction.\allowbreak{}AdmissibleParameterJoins.\allowbreak{}SumLogPrimeFactorsLe}{202--217}
\leanref{Formalism.\allowbreak{}MainNotions.\allowbreak{}J.\allowbreak{}sum\_\allowbreak{}lw\_\allowbreak{}e}{Formalism.\allowbreak{}SharpRestriction.\allowbreak{}AdmissibleParameterJoins.\allowbreak{}SumLogPrimeFactorsLe}{218--236}

The first identity uses the ramification degree formula for \((p)\) in
\(\Z\subseteq\mathcal O_K\) and \(\operatorname{rank}_{\Z}\mathcal O_K=n_K\).
This is the case of \(\Z_{(p)}\) and \(K/\Q\) in
\cite[Remark 15.113.6, Tag 09E8]{stacks-2026b}. The second identity uses
\(N\mathfrak p_w=p^{f_w}\).

\begin{lemma}[Support of the gain and loss in a parameter family]\label{B4.52a}
Let \(\mathfrak F\in\mathrm{FamP}(\varepsilon)\),
\(x>\mathfrak F.x_1\), and \(i\in\mathfrak F.I_x\).
At every finite place \(w\) of this member with \(p=p_w>X''\), the raw
gain vanishes for \(1\le j\le m\).
For each word \(u\), let \(M_u\) be its stabilizer denominator as in
\ref{B4.03}. If also \(p\notin\{T_u,g_0,g_1,d\}\) and \(p\nmid M_u\),
then \(E_{u,j}(w)=0\) for \(1\le j\le m\).
Thus \(\Lean{ZeroAbove}\) holds for every member at \(x\ge\mathfrak F.x_1+1\).
\end{lemma}
\begin{proof}
The gain vanishes off \(L.\mathrm{supp}\). On that support, at primes of
class \(\Lean{IsResidueClassZero}\) or \(\Lean{IsSg}\), the filling condition
identifies it with the prime gain supported on \(p\le X''\).
The non-inert clause gives zero in the other classes.
In \(\Lean{EDef}\) of \ref{B4.03}, the exceptional-prime branches require
\(p\in\{T_u,g_0,g_1,d\}\), and the two inert-prime branches require
\(p\le X''\). The remaining term is \(v_p(M_u)=0\).
\end{proof}
\leanref{Formalism.\allowbreak{}MainNotions.\allowbreak{}WeightDegreeSizes.\allowbreak{}gain\_\allowbreak{}above}{Formalism.\allowbreak{}SharpRestriction.\allowbreak{}Divisibility.\allowbreak{}WeightDegreeSizes}{294--304}
\leanref{Formalism.\allowbreak{}BasicNotions.\allowbreak{}Eval\_\allowbreak{}above}{Formalism.\allowbreak{}LocalRingCardinality.\allowbreak{}DataErr}{174--188}
\leanref{Formalism.\allowbreak{}MainNotions.\allowbreak{}WeightDegreeSizes.\allowbreak{}zero\_\allowbreak{}gain\_\allowbreak{}above\_\allowbreak{}family}{Formalism.\allowbreak{}SharpRestriction.\allowbreak{}Divisibility.\allowbreak{}WeightDegreeSizes}{399--401}
\leanref{Formalism.\allowbreak{}L2.\allowbreak{}DIV.\allowbreak{}Wire.\allowbreak{}za\_\allowbreak{}of}{Formalism.\allowbreak{}LocalToGlobal.\allowbreak{}Divisibility.\allowbreak{}RecBridgeWire}{19--21}

\begin{theorem}[Rational-prime bound for the error]\label{B4.53}
Retain the hypotheses of \ref{B4.49} and assume \(X''\ge0\).
Assume \(\Lean{ZeroAbove}\). For a parameter family \(\mathfrak F\),
\ref{B4.52a} supplies this hypothesis for every member with \(x>\mathfrak F.x_1\).
For each word \(u\) and layer \(1\le j\le m\), it gives vanishing of the
raw gain and \(E_{u,j}(w)\) at places of residue characteristic \(p\)
satisfying
\[
 p>X'',\qquad p\notin\{T_u,g_0,g_1,d\},\qquad p\nmid M_u.
\]
Under these conditions \(S_{u,w}=0\), and
\[
 \left|\Delta\mu_u-A_u\right|
 \le2\left(
 X''\log4+\log T_u+\log g_0+\log g_1+\log d+\log M_u
 \right).
\]
\end{theorem}
\leanref{Formalism.\allowbreak{}MainNotions.\allowbreak{}J.\allowbreak{}sg\_\allowbreak{}eq\_\allowbreak{}zero}{Formalism.\allowbreak{}SharpRestriction.\allowbreak{}AdmissibleParameterJoins.\allowbreak{}SumLogPrimeFactorsLe}{431--447}
\leanref{Formalism.\allowbreak{}MainNotions.\allowbreak{}J.\allowbreak{}nprimeMu\_\allowbreak{}at}{Formalism.\allowbreak{}SharpRestriction.\allowbreak{}AdmissibleParameterJoins.\allowbreak{}SumLogPrimeFactorsLe}{448--493}

\begin{proof}
Form \(P_u\) from the primes \(p\le X''\) and the four exceptional
primes. Adjoin the prime divisors of \(M_u\).
At a place above \(p\notin P_u\), \(\Lean{ZeroAbove}\) gives
\(\gainD_j(w)=E_{u,j}(w)=0\), and the support of \(F_j\) gives
\(F_j(w)=0\). Thus \(S_{u,w}=0\).
By \ref{B4.52} and \(e_w\ge1\), the sum of \(\ell_w\) over any subset of
\(W_p\) is at most \(\log p\). Hence \ref{B4.51} bounds the error by
\(2\sum_{p\in P_u}\log p\).

Put \(P(n)=\prod_{p\le n}p\). For \(k\ge1\), the primes in
\((k+1,2k+1]\) divide \(\binom{2k+1}{k}\le4^k\), since the two equal
central coefficients belong to a sum of \(2^{2k+1}\).
Thus \(P(2k+1)\le P(k+1)4^k\), while \(P(2k)=P(2k-1)\) for \(k\ge2\).
Induction from \(n=0,1,2\) gives \(P(n)\le4^n\). Use \(n=\lfloor X''\rfloor\)
and add the four exceptional logs and \(\log M_u\), which bounds the
sum over its distinct prime divisors.
\end{proof}

\begin{theorem}[Uniform error in the family parameters]\label{B4.54}
Assume the arithmetic size conditions \(\Lean{SizeQD}\) and
\(\Lean{SizeG}\), together with \(\Lean{NprimeDef}\) and
\(\Lean{ZeroAbove}\). Suppose that each \(N'_u\) is finitely generated
and has positive \(K\)-rank. Let \(C_T,C_\Gamma\in\R\) satisfy,
for every word \(u\),
\[
 \log T_u\le C_TX'',\qquad
 \log M_u\le C_\Gamma X''.
\]
Then, for every \(u\),
\[
 \left|\Delta\mu_u-A_u\right|
 \le2(\log4+3+C_T+C_\Gamma)X''.
\]
\end{theorem}
\leanref{Formalism.\allowbreak{}MainNotions.\allowbreak{}J.\allowbreak{}nprimeMuR\_\allowbreak{}at}{Formalism.\allowbreak{}SharpRestriction.\allowbreak{}AdmissibleParameterJoins.\allowbreak{}SumLogPrimeFactorsLe}{694--710}

\begin{proof}
The goodness of \(D\) gives invertible frames, and positive rank gives
\(N'_u\ne0\). Here \(X''=X'xn\ge x>1\) since \(X',n\ge1\).
Put \(\theta=\varepsilon/100\).
The size bounds and \(0<\varepsilon<1/52\) give
\[
 \log g_i\le x^\theta x^\varepsilon
 =x^{\theta+\varepsilon}\le x\le X''
 \quad(i=0,1),\qquad
 \log d\le x^{\varepsilon/5+\theta}\le x\le X''.
\]
Apply \ref{B4.53}.
\end{proof}

\section{The data passed to the height argument}

Take a reference row system and a selected column system forming an
\(\Lean{IndefCoreR}\), with both sides well formed and with the pairing
data of \ref{B4.45}.
They share the layer eigenvalues and metric factors \(\kappa_v>0\).
Their common quotient type is \(t_w\), with \(t_{\bar w}=t_w\) as in \ref{B4.36a}.
Write \(R_0,R_{\mathrm c}\) for the reference and selected glued row lattices.
Fix a column \(u\), and write \(B\) for its map to the selected row dual, with
frames \(X_{\mathrm c}\) on the source and \(X_R^\vee\) on the target. Put
\[
 n_\infty=\#\{v\mid\infty\},\qquad
 \langle z_v\rangle_v=\frac1{n_\infty}\sum_{v\mid\infty}z_v,\qquad
 \bar\kappa=\langle\log\kappa_v\rangle_v,\qquad
 \mu_0=\mu(R_0^\vee)-\bar\kappa,\qquad
 \mu_C=\mu(R_{\mathrm c}^\vee).
\]
Outside the twisting set \(\mathrm{Tw}\), assume that the reference
layer divisor contributions and \(F_{0,j,w}\) vanish. Let
\(\delta,h_{\mathrm r},\lambda_{\mathrm r}\) be the passage parameters
of the row data. For \(0\le\delta\le1/3\), the row data give the two
inequalities of \ref{B4.55a} for every \(U\ge\mu_{\max}(R_0)\), with
\(R_1=R_{\mathrm c}\) and \(\lambda=\lambda_{\mathrm r}\).
For the constructed family these follow from the character-level
embeddings of \ref{B4.55a}.
Assume \(\mu_{\max}(R_0)\le h_{\mathrm r}-\bar\kappa\).

In the chain of \ref{B5.07}, let \(i(j)\) be the step assigned to
layer \(j\) and \(X_{\mathrm{last},v}\) the last frame.
At the base type \(t_0\), choose an invertible frame \(X_{0,v}\) with
\[
 (X_{0,v}^{-1})^*X_{0,v}^{-1}=v(H_{V,t_0}),
\]
where \(H_{V,t_0}\) is the base Hermitian matrix on the selected side
and \(v\) acts entrywise. Use the numerical data of \ref{B5.08}.
Set \(H=H^*\) and \(\mathrm{pq}_j=\beta_j\), and let \(R_*\) denote
the right-hand side of that lemma's conclusion.
Use the height data of \ref{B8.35}, with \(h_P=A\) and \(k_{j,v}=s_{j,v}\).
For the terminal map \(P_u^{\rm term}\) of \ref{B5.07}, require
\[
 \|P_u^{\rm term}\|_{X_{\mathrm{last},v}\to X_{0,v}^{\vee}}^2\le e^{h_P}
 \qquad(v\mid\infty).
\]
For the coupled kernel data \(q_{iv}=\mathrm{CBData.sq}(i,v)\), require
\(\log q_{iv}\le\sum_{j:i(j)=i}k_{j,v}\) for each \(i,v\), with \(\log0=0\).
At the base type, require \(l_{6,j}\le\log\lambda_{j,\alpha,v}\) at every
infinite place \(v\) for each layer \(j\) and spectral index \(\alpha\).
Require \(\sum_j H_{X,j}\le\sum_j l_{6,j}\).
The construction in \ref{B8.36} supplies the terminal and coefficient bounds
of \ref{B8.35}. Lemma~\ref{B5.07a} gives the resulting kernel norm and
height bounds, and \ref{B5.07b} gives the spectral bounds.

For each quotient type \(t\) and each finite place \(w\), choose a basis
\(e_{t,w}\) of the reference quotient lattice localized at \(w\).
Write \(\bar w\) for the conjugate place.
Let \(G_w\) be the Gram determinant in the basis \(e_{t_w,w}\).
Take \(S_y\supseteq\mathrm{Tw}\) finite with \(v_w(G_w)=0\) off \(S_y\).
Write \(D_R>0\) for the reference row dimension. For the layer lattice
\(T=T_{j,t_w}\) with Hermitian form \(H_T=H_{T,j,t_w}\), put
\(J_{j,w}=\operatorname{idx}_w(T,T^{\vee_{H_T}})\).
In a basis \(f\) valid at both \(w\) and \(\bar w\), this is
\(-v_w(\det\operatorname{Gram}_{H_T}(f))\).
Require real bounds \(L_G,E_3,E_{2,j}\) and real-valued functions
\(g_w,\eta_{w,j}\) with
\[
 \begin{gathered}
 v_w(G_w)\le D_Rg_w\quad(w\in S_y),\qquad
 \sum_{w\in S_y}\ell_wg_w\le L_G,\\
 J_{j,w}\le e_wJ_{0,j,w}D_R+\eta_{w,j}
 \quad(w\in\mathrm{Tw}),\qquad
 \sum_{w\in\mathrm{Tw}}\frac{\ell_w\eta_{w,j}}{D_R}\le E_{2,j},\\
 \sum_{w\in S_y}\ell_w\frac{2e_w}{p_w^2}\le E_3.
 \end{gathered}
\]
Here \(J_{0,j,w}\) is the prescribed main term for the pairing index.
For each reference layer \(j\), require
\(\operatorname{lcv}(M_j)/D_R\le l_{3,j}\) in its reference frame.
For \(w\in\mathrm{Tw}\), let \(A_{j,t_w}\) be the image of \(T_{j,t_w}\)
in the reference layer space. Put
\[
 i_{jw}=\operatorname{idx}_w(M_j,A_{j,t_w})
 +\operatorname{idx}_{\bar w}(M_j,A_{j,t_w}),\qquad
 e'_{jw}=e_w\,\mathrm{err}_0(w,j),
\]
where \(\mathrm{err}_0(w,j)\ge0\) is the reference row error. The row data give
\[
 D_R^{-1}\sum_j i_{jw}\le\sum_j(2e_wF_{0,j,w}+e'_{jw})+2e_w/p_w^2,\qquad
 \sum_{w\in\mathrm{Tw}}\ell_w(e'_{jw}+e_wJ_{0,j,w})\le E_{1,j}.
\]
Put
\[
 F_*=\sum_j\sum_{w\in\mathrm{Tw}}\ell_we_wF_{0,j,w},
 \qquad
 B_*=\sum_j l_{3,j}+2F_*+L_G+\sum_j(E_{1,j}+E_{2,j})+E_3.
\]
Require real numbers \(E_{l3},Q_{\rm pass}\) such that
\[
 \begin{gathered}
 \sum_j l_{3,j}\le
 \sum_{j,\rho}\mathrm{pq}_j w_{j,\rho}\log\Omega+E_{l3},\\
 E_b:=E_{l3}+L_G+\sum_j(E_{1,j}+E_{2,j})+E_3
 \le c_emx^\varepsilon(B_0+N_c),\\
 0\le\delta,\qquad Q_{\rm pass}\ge1,\qquad 3\delta Q_{\rm pass}\le1,\\
 B_*+h_{\mathrm r}+28\lambda_{\mathrm r}\le Q_{\rm pass},\qquad
 \lambda_{\mathrm r}\le\lambda_0,\qquad 28\le C_0.
 \end{gathered}
\]
For sufficiently large \(x\) in the selected family, take
\(Q_{\rm pass}=Q_0\) from \ref{B5.09a}.

\begin{theorem}[Four conclusions in coordinates]\label{B4.55}
Let \(K\) be totally complex and take the systems and numerical inputs
above. Let \(N\subset K^{d_u}\) be a full lattice with \(d_u>0\) and
invertible column frames \(X_{\mathrm c,v}\). Let \(\sigma_w\in\R\)
have floors \(\tau_w=\lfloor\sigma_w\rfloor\) of finite support.
At each finite place \(w\), suppose that every \(\varphi\in N\) and every
\(F\) in the selected coordinate row lattice \(M_{t_w}\) satisfy
\[
 P_{u,t_w}(\varphi,F)=0
 \quad\text{or}\quad
 v_w(P_{u,t_w}(\varphi,F))\ge\sigma_w.
\]
Let \(N^\sigma\) be the shift of \(N\) by \(\tau\).
Then \(B(N^\sigma)\subseteq R_{\mathrm c}^{\vee}\). There are positive
numbers \(c_v\) with \(\|B\|_v^2\le c_v\) at every infinite place
and \(\langle\log c_v\rangle_v=H-\bar\kappa\). Moreover,
\[
 \begin{gathered}
 \mu_C-\bar\kappa+H-2F_*\le R_*,\\
 \mu_{X_{\mathrm c}}(N^\sigma)-\mu_{X_{\mathrm c}}(N)
 =2\sum_w\ell_w\tau_w.
 \end{gathered}
\]
\end{theorem}
\leanref{Formalism.\allowbreak{}L2.\allowbreak{}DIV.\allowbreak{}divClausesR\_\allowbreak{}of}{Formalism.\allowbreak{}LocalToGlobal.\allowbreak{}Divisibility.\allowbreak{}ConstantGeTwentyEight}{184--248}
\leanref{Formalism.\allowbreak{}L2.\allowbreak{}DIV.\allowbreak{}DivClausesR}{Formalism.\allowbreak{}LocalToGlobal.\allowbreak{}Divisibility.\allowbreak{}ConstantGeTwentyEight}{111--122}
\leanref{Formalism.\allowbreak{}L2.\allowbreak{}DIV.\allowbreak{}HtNums.\allowbreak{}le\_\allowbreak{}hu}{Formalism.\allowbreak{}LocalToGlobal.\allowbreak{}Divisibility.\allowbreak{}DivisionHeightBoundStatement}{113--123}

\begin{proof}
By \ref{B4.40}, \(N^\sigma\) exists and is full. Apply \ref{B4.45}
with \(\tau_w\le\lceil\sigma_w\rceil\) to obtain the inclusion.

Since \(\|K_i\|_v^2\le q_{iv}\), the log bound gives
\(\|K_i\|_v^2\le\exp(\sum_{j:i(j)=i}k_{j,v})\), also when \(q_{iv}=0\).
For \(B\ne0\), \ref{B5.07} gives \(\log\|B\|_v^2\le t_v\), where
\[
 t_v=h_P+\sum_j k_{j,v}-\log\kappa_v
 -\min_\alpha\sum_j\log\lambda_{j,\alpha,v}.
\]
The layer bounds hold simultaneously for each \(\alpha\). With
\(\bar t=\langle t_v\rangle_v\), they give
\[
 \bar t\le h_P+\\sum_j\langle k_{j,v}\rangle_v-\sum_j l_{6,j}
 -\bar\kappa\le H-\bar\kappa.
\]
The last inequality is \(\Lean{HtNums.\allowbreak{}le\_\allowbreak{}hu}\) with
\(m_\lambda=\sum_j l_{6,j}\ge\sum_j H_{X,j}\).
Take \(c_v=\exp(t_v+H-\bar\kappa-\bar t)\), or
\(c_v=\exp(H-\bar\kappa)\) if \(B=0\).
Both choices give the norm bound and the stated mean.

The reference-row divisor estimate \ref{B5.08x} gives \(\mu_0\le B_*\).
Since \(Q_{\rm pass}\ge1\) and \(3\delta Q_{\rm pass}\le1\), we have
\(0\le\delta\le1/3\).
Use the second passage inequality with \(U=h_{\mathrm r}-\bar\kappa\) to get
\[
 \mu_C-\bar\kappa\le\mu_0+L_{\rm pass},\qquad
 L_{\rm pass}=3\delta(\mu_0+h_{\mathrm r}+28\lambda_{\mathrm r})
 +28\lambda_{\mathrm r}.
\]
The bounds on \(B_*\) and \(\lambda_{\mathrm r}\) give
\(L_{\rm pass}\le1+C_0m\lambda_0\).
The layer budgets give
\[
 \mu_0\le\sum_{j,\rho}\mathrm{pq}_jw_{j,\rho}\log\Omega+2F_*+E_b.
\]
Apply \ref{B5.08} with these bounds to obtain
\(\mu_C-\bar\kappa+H-2F_*\le R_*\). The slope difference is \ref{B4.50}.
\end{proof}

\begin{lemma}[Passage between character levels]\label{B4.55a}
Let \(K\) be totally complex. Let a finite group of order \(M\) act
unitarily at infinity on a metrized \(K\)-space and preserve a full lattice.
Suppose the retained irreducible types have distinct characters and scalar
endomorphisms. Assume their central projectors sum to the identity.
Take linear \(K\)-valued characters \(\chi_i\) of subgroups of order
\(h_i\) at levels \(i=0,1\). On type \(s\), \(\chi_0\) has multiplicity
one and \(\chi_1\) has multiplicity \(n_s\), with
\[
 (1-\delta)\bar n\le n_s\le(1+\delta)\bar n,\qquad
 \bar n>0,\quad 0\le\delta\le1/3.
\]
Let \(P_i\) be the ambient lattice intersected with the \(\chi_i\)-space.
Suppose the row embedding identifies \(R_i\) with a nonzero lattice
between \(h_iP_i\) and \(P_i\) in the ambient frame multiplied by
\(\sqrt{h_i}\). For \(U\ge\mu_{\max}(R_0)\) and \(\lambda\ge\log M\),
\[
 \mu_{\max}(R_1)\le U+14\lambda,\qquad
 \mu(R_1^\vee)\le\mu(R_0^\vee)+
 3\delta\bigl(\mu(R_0^\vee)+U+28\lambda\bigr)+28\lambda.
\]
\end{lemma}
\begin{proof}
Put \(l=\log M\). The lattice inclusions change either slope by an amount
in \([-2\log h_i,0]\). Scaling the frame adds \(\log h_i\), so the mean
and maximum slopes of \(R_i\) differ from those of \(P_i\) by at most \(l\).

Omit types with zero ambient isotypic space. They contribute zero to the
rank and degree sums. Write \(r_s,d_s\) for the rank and degree of the
base cut in type \(s\), with degree equal to rank times slope. The total
ranks are \(R=\sum_s r_s>0\) and \(W=\sum_s n_sr_s\).
Central projectors have denominator dividing \(M\), so at either level
the orthogonal sum of the type cuts lies between \(MP_i\) and \(P_i\),
at a cost of at most \(2l\) per rank.

In an irreducible type, translates of the base character line span the
type. Project onto the selected character space and choose \(n_s\)
translates whose projections form a basis. On the isotypic space the
projected group operators give an isomorphism from \(n_s\) copies of the
base space to the selected space. The inverse translates followed by
base projections give the adjoint map, also an isomorphism.
Each map uses at most a central projector and the two subgroup averages,
so multiplication by \(M^3\) clears denominators. The resulting maps have
norms at most \(M^3\sqrt{n_s}\) by unitarity and Cauchy--Schwarz.
Since \(n_s^2\le(\dim s)^2\le M\), their log squared norms are at most
\(13l/2\). Comparing degrees gives
\[
 \mu(P_0)R\le\sum_s d_s+2lR,\qquad
 \mu(P_1)W\ge\sum_s n_sd_s-\tfrac{17}{2}lW.
\]
The reverse injection bounds the maximum slope of each selected cut by
\(\mu_{\max}(P_0)+13l/2\). The orthogonal sum and lattice inclusions give
\(\mu_{\max}(P_1)\le\mu_{\max}(P_0)+17l/2\).
The two frame comparisons give
\(\mu_{\max}(R_1)\le U+21l/2\le U+14\lambda\).

Set \(A=U+l\) and \(z_s=Ar_s-d_s\ge0\), since the base cut lies in \(P_0\).
Then \(\sum_s z_s/R\le A-\mu(R_0)+3l\), and
\[
 \frac{\sum_s n_sz_s}{W}
 \le\frac{1+\delta}{1-\delta}\frac{\sum_s z_s}{R}
 \le(1+3\delta)\frac{\sum_s z_s}{R}.
\]
The lower degree bound at level \(1\) and the frame comparison give
\[
 -\mu(R_1)\le-\mu(R_0)+\tfrac{25}{2}l+
 3\delta\bigl(U-\mu(R_0)+4l\bigr).
\]
Use \(\mu(R_i^\vee)=-\mu(R_i)\) and \(l\le\lambda\).
\end{proof}
\leanref{Formalism.\allowbreak{}IndefR.\allowbreak{}LatticeTransfer.\allowbreak{}levelAmb\_\allowbreak{}transfer}{Formalism.\allowbreak{}LocalToGlobal.\allowbreak{}IndefiniteReal.\allowbreak{}LevelAmbTransfer}{13--55}
\leanref{Formalism.\allowbreak{}BranchingDeterminantCap.\allowbreak{}meanTransferBound}{Formalism.\allowbreak{}BranchSandwichDivisors.\allowbreak{}BranchIsometryProofs.\allowbreak{}TransferMean}{19--65}
\leanref{Formalism.\allowbreak{}IndefR.\allowbreak{}CxPass.\allowbreak{}In3.\allowbreak{}passage\_\allowbreak{}of\_\allowbreak{}steps}{Formalism.\allowbreak{}LocalToGlobal.\allowbreak{}IndefiniteReal.\allowbreak{}CutPassage3SOfSteps}{49--109}

\begin{theorem}[Return to the original period data]\label{B4.56}
Fix \(D,L\) and \(C,C_{\mathrm{abs}}\in\R\). Take a compatible division
package. For each word \(u\) it supplies the inputs of
\(\Lean{DivInputs0R}\) above for the coordinate lattice \(\mathcal N\)
and the shift \(\sigma_w=s_{u,w}\), as well as the period-bridge
identities of \ref{B4.38}. Its \(K\)-linear isometries identify the
shifted coordinate column with \(N'_u\) and the glued row lattice
with \(RR\). The isometries preserve the pairings. Assume \(\Lean{adm2}\) of \ref{B4.29} and \(\Lean{NprimeDef}\).
Assume \(\Lean{CutUp}\) of \ref{B4.37} for every used character.

For the arithmetic data, put \(a=xe^{x^\varepsilon}\) and
\(L_a=\log(a\,\mathrm{rd}(E))\), where \(\mathrm{rd}(E)\) is the root
discriminant of the host field \(E\).
With \(H_u,R_{*,u}\) the coordinate bounds, assume
\[
 H_u-\bar\kappa=h_u,\qquad
 2F_*=\Fo,\qquad
 R_{*,u}\le\mathrm{MAIN}(u)+Cn_TaL_a
 +Cm(B_0+N_c)x^{C_{\mathrm{abs}}\varepsilon},
\]
where \(n_T=\#TT\) and, with \(\phi(z)=z\log z-z\) and \(\phi(0)=0\),
\[
 \mathrm{MAIN}(u)=\sum_{t\in TT}
 \bigl[m(\phi(a_t)-a_t\log\Omega)-a_t\log T_u\bigr]
\]
uses the arithmetic weights and period constant. Then for every word \(u\),
\[
 \begin{gathered}
 B_u(N'_u)\subseteq RR^\vee,\\
 B_u\ne0\ \Longrightarrow\
 h(B_u):=\langle\log\|B_u\|_v^2\rangle_v\le h_u,\\
 \mu(RR^\vee)+h_u-\Fo
 \le\mathrm{MAIN}(u)+Cn_TaL_a
 +Cm(B_0+N_c)x^{C_{\mathrm{abs}}\varepsilon}.
 \end{gathered}
\]
\end{theorem}
\leanref{Formalism.\allowbreak{}L2.\allowbreak{}DIV.\allowbreak{}DivisionPackageFace.\allowbreak{}wordFacts}{Formalism.\allowbreak{}LocalToGlobal.\allowbreak{}Divisibility.\allowbreak{}RecBridge}{633--642}
\leanref{Formalism.\allowbreak{}L2.\allowbreak{}DIV.\allowbreak{}WordFacts}{Formalism.\allowbreak{}LocalToGlobal.\allowbreak{}Divisibility.\allowbreak{}RecBridge}{128--134}

\begin{proof}
Fix \(u\) with column isomorphism \(e\). The shifts are \(s_{u,w}\),
zero off \(L.\mathrm{supp}\), so the shifted lattice is \(I_u\mathcal N\).
The shift identification and \(\Lean{NprimeDef}\) give
\(I_ue(\mathcal N)=I_uN_u\).
Multiply by \(I_u^{-1}\), which exists by \ref{B4.48}, to get
\(e(\mathcal N)=N_u\).
Since \(\theta(t_w)=\mu_w\), \ref{B4.38} uses \(\Lean{adm2}\) and
\(\Lean{CutUp}\) for \(\mu_w\) to give the valuation bound on \(M_{t_w}\).

By \ref{B4.55}, \(B_u(\varphi)(r)=B(x)(y)\in\mathcal O_K\) for
\(\varphi\in N'_u\), \(r\in RR\) with coordinates \(x,y\).
If \(B_u\ne0\), the isometries identify its norms with those of \(B\), so
\(h(B_u)\le\langle\log c_v\rangle_v=H_u-\bar\kappa=h_u\).
Since \(\mu_C=\mu(RR^\vee)\), the third conclusion of \ref{B4.55}
gives the budget.
\end{proof}


\part{Two bounds on one slope}

% ---- p5c1 ----
\chapter{The lattice after division}\label{ch:lattice-after-division}

Fix a window point \(x>1\) and a parameter \(0<\varepsilon<1/52\). Let \(u_0\) be a translation. The words form a finite set \(W(u_0)\). The host CM field is \(E\), and \(f=[E:\Q]/2\) is the number of column positions \(\rho\).

Let \(K\) be the totally complex coefficient field. At a finite place \(w\), let \(\mathfrak p_w,p_w,e_w\) be its prime ideal, residue characteristic and ramification index. Put
\[
 \ell_w=\frac{\log N\mathfrak p_w}{[K:\Q]},\qquad c_w=2e_w\ell_w.
\]
Write \(S\) for the finite support of the local data. At infinity use the metrics of Part~IV. For \(B\ne0\) set \(h(B)=\langle\log\|B\|_v^2\rangle_\infty\), where \(\langle\cdot\rangle_\infty\) is the mean over infinite places. Slope differences and \(h(B)\) are twice the corresponding normalized Arakelov slope differences and archimedean height of \cite[\S\S2.4 and 2.6]{bost-chen-2013}.

For the arithmetic chain data \(D\), write \(N_u^0=\mathrm{NuLat}_u\) for the input lattice and \(N'_u\) for the divided lattice. Set \(M_u=(D.\mathrm{fm}\,u).\mathrm{lat}(0)\). The ramified comparison lattice \(M_{H,u}\) is the inverse image of \(M_u\) under \(\nu\mathrm{In}_u\), with the metric of \(N_u^0\). Let \(V_u,W\) be the ambient \(K\)-spaces of \(N'_u,RR\). The period map and target lattice are
\[
 B_u:V_u\longrightarrow W^\vee,\qquad
 C^\vee=RR^\vee=\{\ell\in W^\vee:\ell(RR)\subseteq\mathcal O_K\},
 \qquad \mu_C=\mu(C^\vee).
\]
Put
\[
 L=x^\varepsilon,\quad b=e^L,\quad a=xb,\quad
 L_a=\log(a\,\mathrm{rd}E),\quad n_T=\#TT,\quad S_a=\sum_{t\in TT}a_t .
\]
The integer top weights satisfy \(a\le a_t\le2a\). The chain has \(m\) lifts and starting rank \(k\). The auxiliary prime of \(u\) is \(T_u\). The window data supply the budgets \(B_0,N_c\) and scales \(b',X''\).

\section{Prime powers and finite-place budgets}

\begin{definition}[The required prime-power sum]\label{B5.01}
For the finite abelian group \(A\) of the host, put
\[
 G=(\Z/2\Z)^2\times A,\qquad \Phi=\{((e_1,e_2),t)\in G:e_1=0\},
 \qquad r_t(g)=((1,1),0)+g+((0,0),t).
\]
Let \(F_p\in G\) be the chosen Frobenius element and \(I_p\le G\) the inertia subgroup. Write \(q_1,d\) for the exceptional primes of the host. For \(N\in\N\), \(a_t\in\R\) and \(w_{u_0j}:G\to\R\), set
\[
\begin{aligned}
 A_{tj}^{(N)}
 ={}&
 \sum_{\substack{p\le N,\ p\ {\rm prime}\\
 p\notin\{2,q_1,d\}\\ F_p\notin\Phi}}
 2\log p
 \sum_{\substack{1\le i\le N\\iF_p\notin\Phi}}
 \left[\frac{a_t}{p^i}-w_{u_0j}(r_t(iF_p))\right]_+
 \\ &+
 \log d\sum_{i=1}^{N}
 \sum_{\substack{g\in iF_d+I_d\\g\notin\Phi}}
 \left[\frac{a_t}{d^i}-w_{u_0j}(r_t(g))\right]_+ ,
 \qquad [z]_+=\max(z,0).
\end{aligned}
\]
Then \(\ASK(u_0)=\sum_{t\in TT}\sum_{j=0}^{m-1} A_{tj}^{(N)}\), with index \(j\) denoting geometric layer \(j+1\).
\end{definition}
\leanref{Formalism.\allowbreak{}LFunctionChainParameters.\allowbreak{}characterSumData.\allowbreak{}ASK}{Formalism.\allowbreak{}SizeWindowBrackets.\allowbreak{}FirstPrimeNeSecondPrime}{84--88}
\leanref{Formalism.\allowbreak{}LFunctionChainParameters.\allowbreak{}characterSumData.\allowbreak{}askP}{Formalism.\allowbreak{}SizeWindowBrackets.\allowbreak{}FirstPrimeNeSecondPrime}{75--77}
\leanref{Formalism.\allowbreak{}LFunctionChainParameters.\allowbreak{}characterSumData.\allowbreak{}askD}{Formalism.\allowbreak{}SizeWindowBrackets.\allowbreak{}FirstPrimeNeSecondPrime}{79--82}
\leanref{Formalism.\allowbreak{}MainNotions.\allowbreak{}Lnk.\allowbreak{}Mod.\allowbreak{}data}{Formalism.\allowbreak{}SharpRestriction.\allowbreak{}GlobalSchurBounds.\allowbreak{}QuadraticLogDerivCongr}{464--505}

The first binary coordinate of \(F_p\) is \(1\) in the ordinary-prime sum, so \(i\) is odd. If each admissible weight \(w>0\) satisfies \(a_t\le Nw\), a positive summand requires \(p^i<a_t/w\le N\), which bounds \(p\) and \(i\). The same argument bounds the exponents at \(d\). In the application \(w\ge1\), so \(N=\lceil2a\rceil\) suffices.

\begin{definition}[Local multiplicities and their global totals]\label{B5.02}
Let \(g_{jw},F_{jw},E_{ujw}\) be the gain, base term and word-dependent loss of layer \(1\le j\le m\) from Parts~III and~IV. For the selected set \(\mathrm{drop}\) of rational primes, set \(\gainD_j(w)=0\) on \(p_w\in\mathrm{drop}\) and \(\gainD_j(w)=g_{jw}\) elsewhere. Put
\[
 \sg(u,w)=\sum_{j=1}^{m}\bigl(\gainD_j(w)+F_{jw}-E_{ujw}\bigr).
\]
Let \(e_w\tau_w\) be the containment exponent of \ref{B5.06} and \(\mathrm{io}_w\) its extra length allowance. Define
\[
\begin{aligned}
 \Ga&=\sum_{w\in S}c_w\left(\sum_{j=1}^{m}\gainD_j(w)-\tau_w\right),
 &\mathcal T&=\sum_{w\in S}c_w\tau_w,\\
 \Fo&=\sum_{j,w}c_wF_{jw},
 &\mathrm{IO}&=\sum_{w\in S}e_w\ell_w\,\mathrm{io}_w,\\
 \Ee_u&=\sum_{j,w}c_wE_{ujw},
\end{aligned}
\]
where the double sums run over the layers and \(S\). The quantities \(\Ga,\Fo\) are common to the words, and \(\sum_{w\in S}c_w\sg(u,w)=\Ga+\Fo-\Ee_u+\mathcal T\).
\end{definition}
\leanref{Formalism.\allowbreak{}BasicNotions.\allowbreak{}SizeData₂.\allowbreak{}sg}{Formalism.\allowbreak{}LocalRingCardinality.\allowbreak{}DataParams₂}{215--217}
\leanref{Formalism.\allowbreak{}BasicNotions.\allowbreak{}SizeData₂.\allowbreak{}sz}{Formalism.\allowbreak{}LocalRingCardinality.\allowbreak{}DataParams₂}{205--206}
\leanref{Formalism.\allowbreak{}BasicNotions.\allowbreak{}SizeData₂.\allowbreak{}gainD}{Formalism.\allowbreak{}LocalRingCardinality.\allowbreak{}DataParams₂}{210--213}
\leanref{Formalism.\allowbreak{}BasicNotions.\allowbreak{}SizeData₂.\allowbreak{}Fo}{Formalism.\allowbreak{}LocalRingCardinality.\allowbreak{}DataParams₂}{218--219}
\leanref{Formalism.\allowbreak{}BasicNotions.\allowbreak{}SizeData₂.\allowbreak{}Ga}{Formalism.\allowbreak{}LocalRingCardinality.\allowbreak{}DataParams₂}{220--221}
\leanref{Formalism.\allowbreak{}BasicNotions.\allowbreak{}SizeData₂.\allowbreak{}Ee}{Formalism.\allowbreak{}LocalRingCardinality.\allowbreak{}DataParams₂}{222--223}
\leanref{Formalism.\allowbreak{}MainNotions.\allowbreak{}Lnk.\allowbreak{}IOsum}{Formalism.\allowbreak{}SharpRestriction.\allowbreak{}GlobalSchurBounds.\allowbreak{}QuadraticLogDerivCongr}{163--164}

\begin{proof}
Interchange the layer and place sums and add back \(\mathcal T\).
\end{proof}

\begin{definition}[The ideal prescribed by period valuations]\label{B5.03}
An unrounded period bound means that for every \(u,w\), \(\phi\in N_u^0\) and \(F\in\widetilde M_{\mu_w}\cap V_0(\mu_w)\), there is \(y\in K\) with specified complex image \(\operatorname{perCol}_u(\mu_w)(j_N\phi,F)\) such that \(y=0\) or \(v_w(y)\ge e_w\sg(u,w)\). Here \(\mu_w\) is the arithmetic character and \(j_N\) the column embedding. The integer-valued additive valuation is \(v_w\). Set
\[
 s_{uw}=\lfloor e_w\sg(u,w)\rfloor,\qquad
 I_u=\prod_{w\in S}\mathfrak p_w^{-s_{uw}}.
\]
The condition \Lean{NprimeDef} is \(N'_u=I_uN_u^0=\operatorname{span}_{\mathcal O_K}\{a\phi:a\in I_u,\ \phi\in N_u^0\}\).
\end{definition}
\leanref{Formalism.\allowbreak{}MainNotions.\allowbreak{}LocalQuadraticParts.\allowbreak{}UnroundedValuationBound}{Formalism.\allowbreak{}SharpRestriction.\allowbreak{}Divisibility.\allowbreak{}LocalQuadraticParts}{212--217}
\leanref{Formalism.\allowbreak{}BasicNotions.\allowbreak{}DV.\allowbreak{}Iu}{Formalism.\allowbreak{}LocalRingCardinality.\allowbreak{}DataDIVDefs}{17--21}
\leanref{Formalism.\allowbreak{}BasicNotions.\allowbreak{}DV.\allowbreak{}NprimeDef}{Formalism.\allowbreak{}LocalRingCardinality.\allowbreak{}DataDIVDefs}{22--25}

The same \(y\) satisfies the threshold \(s_{uw}\). In the arithmetic construction, \ref{B4.29} gives the rounded bound on test vectors and \ref{B4.38} extends it to coordinate rows. The shift of \ref{B4.40} is identified with \(N'_u\) in \ref{B4.41}. The pairing of \ref{B4.44} gives \(B_u(N'_u)\subseteq RR^\vee\) by \ref{B4.45}.

\section{The two slopes of the divided lattice}

\begin{lemma}[The exact change of mean slope]\label{B5.04}
Let \(N\subseteq M\) be nonzero \(\mathcal O_K\)-lattices of rank \(d\) in the same Hermitian space, with \(M\) finitely generated. For a finite set \(S_0\) containing the support of \(M/N\), let \(\lambda_w\) be the length of \(M_{\mathfrak p_w}/N_{\mathfrak p_w}\) over \(\mathcal O_{K,\mathfrak p_w}\). Then
\[
 \mu(N)=\mu(M)-\frac2d\sum_{w\in S_0}\ell_w\lambda_w.
\]
Suppose \(N_u^0\subseteq M_{H,u}\) are commensurable, with quotient supported on \(S\) and local lengths \(e_w\,\mathrm{lenMN}_{uw}\). For \(d_u=\operatorname{rk}_K N'_u>0\),
\[
 \mu(N'_u)=\mu(M_{H,u})+2\sum_{w\in S}\ell_w s_{uw}
 -\frac2{d_u}\sum_{w\in S}\ell_we_w\,\mathrm{lenMN}_{uw}.
\]
\end{lemma}
\leanref{Formalism.\allowbreak{}L2.\allowbreak{}GlobalTheoryTheory.\allowbreak{}Places.\allowbreak{}mu\_\allowbreak{}index\_\allowbreak{}sum}{Formalism.\allowbreak{}LocalToGlobal.\allowbreak{}GlobalTheory.\allowbreak{}PlacesP}{61--103}
\leanref{Formalism.\allowbreak{}L2.\allowbreak{}GlobalTheoryTheory.\allowbreak{}Places.\allowbreak{}mu\_\allowbreak{}eq}{Formalism.\allowbreak{}LocalToGlobal.\allowbreak{}GlobalTheory.\allowbreak{}PlacesP}{129--170}

\begin{proof}
The local composition factors have cardinality \(N\mathfrak p_w\), so the index is \([M:N]=\prod_{w\in S_0}(N\mathfrak p_w)^{\lambda_w}\). It multiplies the covolume and changes the slope by \(-2\log[M:N]/([K:\Q]d)\). Commensurability and twisting preserve the span, so all ranks in the second assertion are \(d_u\). Apply the first formula and use \ref{B4.47} and \ref{B4.48} for the twist, which adds \(-2\log N(I_u)/[K:\Q]=2\sum_{w\in S}\ell_ws_{uw}\).
\end{proof}

\begin{lemma}[A local comparison for the maximum slope]\label{B5.05}
Suppose \(N_u^0\) and \(M_{H,u}\) have the same \(K\)-span and \(d_u>0\). For \(k_w\in\N\) supported on \(S\), assume at every finite place that
\[
 (N_u^0)_{\mathfrak p_w}\subseteq
 \mathfrak p_w^{k_w}(M_{H,u})_{\mathfrak p_w}.
\]
Then
\[
 \mu_{\max}(N'_u)\le\mu_{\max}(M_{H,u})+2\sum_{w\in S}\ell_w(s_{uw}-k_w).
\]
\end{lemma}
\leanref{Formalism.\allowbreak{}L2.\allowbreak{}GlobalTheoryTheory.\allowbreak{}Places.\allowbreak{}muMax\_\allowbreak{}le}{Formalism.\allowbreak{}LocalToGlobal.\allowbreak{}GlobalTheory.\allowbreak{}PlacesP}{171--238}

\begin{proof}
For \(J=\prod_{w\in S}\mathfrak p_w^{k_w}\), localization gives \(N_u^0\subseteq JM_{H,u}\), hence \(N'_u\subseteq I_uJM_{H,u}\). Use monotonicity of maximum slope with the induced metric and the bound \(\mu_{\max}(AM_{H,u})\le\mu_{\max}(M_{H,u})-2\log N(A)/[K:\Q]\) for \(A=I_uJ\).
\end{proof}

\begin{lemma}[Cost at the level prime]\label{B5.05a}
Use the eventual level models of \ref{B8.37}, supplied by real-place second-layer compatibility and \Lean{TuSize}. Put \(T=T_u\), \(H=T-1\), \(d=(D.\mathrm{mem}\,u).A.\mathrm{du}\) and \(d_M=\dim_\C M^0_{u,\C}\) for the zeroth column space \(M^0_{u,\C}\). Then
\[
 d=d_M(1-1/H)^f,\qquad
 \mu(M_{H,u})=\mu(M_u)-\frac{f\log H}{H-1},\qquad
 \mu_{\max}(M_{H,u})\le\mu_{\max}(M_u).
\]
The containment of \ref{B5.05} holds at \(w\mid T\) with exponent zero. For every finite \(F\subseteq\{w:w\mid T\}\) and \(d>0\),
\[
 \sum_{w\in F}\frac{2\ell_we_w\,\mathrm{lenMN}_{uw}}d
 \le36n_Tm(a/k+b')\log T .
\]
\end{lemma}
\leanref{Formalism.\allowbreak{}L2.\allowbreak{}GlobalTheoryTheory.\allowbreak{}LocalLatticeBases.\allowbreak{}congruence\_\allowbreak{}no\_\allowbreak{}kernel\_\allowbreak{}minimal}{Formalism.\allowbreak{}LocalToGlobal.\allowbreak{}GlobalTheory.\allowbreak{}GL6.\allowbreak{}CongruenceNoKernelOfHeckeRamified}{75--159}
\leanref{Formalism.\allowbreak{}MainNotions.\allowbreak{}J.\allowbreak{}division\_\allowbreak{}rank\_\allowbreak{}slopes\_\allowbreak{}of\_\allowbreak{}lines}{Formalism.\allowbreak{}SharpRestriction.\allowbreak{}GlobalSchurBounds.\allowbreak{}LinkMin}{497--583}
\leanref{Formalism.\allowbreak{}MainNotions.\allowbreak{}J.\allowbreak{}lemma\_\allowbreak{}five\_\allowbreak{}one\_\allowbreak{}of\_\allowbreak{}theorem\_\allowbreak{}two}{Formalism.\allowbreak{}SharpRestriction.\allowbreak{}GlobalSchurBounds.\allowbreak{}LinkMin}{587--613}

\begin{proof}
The split diamond group \(((\Z/T)^\times)^f\) acts with free orbits. Its translated slices are integrally isomorphic. They are orthogonal by disjoint support and isometric by the finite-character metric condition. In each coordinate the ramified lattice is the kernel of summation. The basis \(e_i-e_H\), \(1\le i<H\), has Gram determinant \(\det(I+\mathbf1\mathbf1^{\mathsf T})=H\). Tensoring in the \(f\) coordinates gives the rank fraction and increases the log Gram determinant per rank by \(f\log H/(H-1)\). Inclusion gives the maximum-slope bound.

At \(w\mid T\) write \(S_0=(\Lambda_{u,0})_{(w)}\) and \(R=S_0\cap V^{\rm ram}\). For an active position \(v\), put \(b_v=b_{\rm top}(v)\). The intermediate lattice \(B_v\) contains \(T^{b_v}S_0\) with colength \(e_wd_M(\sum_i\lambda_{v,i}/k+b_v)\) by the weight determinant. The operators \(D_v,U_v\) preserve the ramified space and satisfy \(D_vS_0\subseteq B_v\), \(U_vB_v\subseteq S_0\), and \(U_vD_v=T^{k-1}\) on that space. Thus \(U_v(B_v\cap V^{\rm ram})\) contains \(T^{k-1}R\) and has colength at most \(e_w(k-1)d_M\) in \(R\). Its quotient by \(U_vT^{b_v}R\) is an image of a submodule of \(B_v/T^{b_v}S_0\).

The congruence lattice \((\prod_v T^{b_v}U_v)^mR\) is contained in \(R\), with equality when there are no active positions. Lengths add under successive finite-index inclusions. The weight window bounds the cost of a position by
\[
 e_wd_M(a_{\max}/k+2\mathrm{Wt}_v+k+b_v)\le9e_wd_M(a/k+b'),
\]
where \(a_{\max}\) is the maximum top weight and \(\mathrm{Wt}_v\) the weight allowance. There are at most \(n_T\) positions per lift, and the local length is \(e_w\,\mathrm{lenMN}_{uw}\). Since \(T>2f+1\), Bernoulli's inequality gives \(d/d_M\ge1/2\). Sum with \(\sum_{w\in F}\ell_we_w\le\log T\) and multiply by \(2/d\).
\end{proof}

\begin{theorem}[Mean and maximum slopes after division]\label{B5.06}
Take the local comparison data of Part~IV. For each word put \(d=d_u\), \(d_M=\dim_\C M^0_{u,\C}\), \(\delta=1-d/d_M\), and \(S_T=\{w\in S:p_w=T_u\}\). Define
\[
 \sigma'_w=s_{uw}/e_w,\qquad
 R_u=\sum_{w\in S}c_w\bigl(\sg(u,w)-\sigma'_w\bigr),\qquad
 P_u=\sum_{w\in S_T}c_w\,\mathrm{lenMN}_{uw}/d.
\]
Suppose \(d_M>0\), \(T_u>2\), and the rank and slope comparisons of \ref{B5.05a} hold. Use the lengths and metrics of \ref{B5.04} and the containments of \ref{B5.05} with \(k_w=e_w\tau_w\). Assume \(\tau_w\in\N\), \(\mathrm{io}_w\ge0\), \(\tau_w=0\) on \(S_T\), and \(\mathrm{lenMN}_{uw}\le d_M(\tau_w+\mathrm{io}_w)\) off \(S_T\). Assume also
\[
\begin{gathered}
 \mathcal T\le b^2,\qquad
 \frac{2f}{T_u}\le\frac12,\qquad
 \frac{4f}{T_u}b^2+\frac{f\log(T_u-1)}{T_u-2}\le1,\\
 P_u\le c_1n_Tm(a/k+b')\log T_u,\qquad R_u\le c_1X'',
 \qquad c_1\ge4,\quad B_0\ge0,\quad\mathrm{IO}\ge0.
\end{gathered}
\]
With \(J_u=\Ga+\Fo-\Ee_u\) and \(\mathrm{Lo}_u=c_1\mathrm{IO}+c_1n_Tm(a/k+b')\log T_u+c_1X''+1\), one has
\[
 \mu(M_u)+J_u-\mathrm{Lo}_u\le\mu(N'_u),\qquad
 \mu_{\max}(N'_u)\le\mu_{\max}(M_u)+J_u+c_1mB_0.
\]
\end{theorem}
\leanref{Formalism.\allowbreak{}MainNotions.\allowbreak{}Lnk.\allowbreak{}slopesN\_\allowbreak{}at\_\allowbreak{}of}{Formalism.\allowbreak{}SharpRestriction.\allowbreak{}GlobalSchurBounds.\allowbreak{}TopDensityNGeSixteen}{1529--1540}
\leanref{Formalism.\allowbreak{}LFunctionChainParameters.\allowbreak{}WeightedSlopeBounds.\allowbreak{}QuotientSlopeBoundsWith}{Formalism.\allowbreak{}SizeWindowBrackets.\allowbreak{}QuotientSlopeBoundsWith}{19--23}
\leanref{Formalism.\allowbreak{}L2.\allowbreak{}GlobalTheoryTheory.\allowbreak{}Places.\allowbreak{}places\_\allowbreak{}at\_\allowbreak{}of}{Formalism.\allowbreak{}LocalToGlobal.\allowbreak{}GlobalTheory.\allowbreak{}PlacesP}{297--436}

\begin{proof}
Bernoulli's inequality gives \(0\le\delta\le f/(T_u-1)\le2f/T_u\le1/2\). By \ref{B5.02} and \ref{B5.04},
\[
 \mu(N'_u)=\mu(M_{H,u})+J_u-R_u-\mathcal B_u,\qquad
 \mathcal B_u=\sum_{w\in S}c_w\left(\frac{\mathrm{lenMN}_{uw}}d-\tau_w\right).
\]
Off \(S_T\), the parenthesis is at most \((\mathrm{io}_w+\delta\tau_w)/(1-\delta)\), and the contribution on \(S_T\) is \(P_u\). Hence
\[
 \mathcal B_u\le P_u+\frac{2\mathrm{IO}+\delta\mathcal T}{1-\delta}
 \le P_u+4\mathrm{IO}+\frac{4f}{T_u}b^2.
\]
The mean-slope comparison and the bound on the two errors give \(\mu(N'_u)\ge\mu(M_u)+J_u-R_u-P_u-4\mathrm{IO}-1\). For the maximum slope, \ref{B5.05} adds at most \(\sum_{w\in S}c_w(\sigma'_w-\tau_w)\le J_u\) to \(\mu_{\max}(M_{H,u})\), since \(\sigma'_w\le\sg(u,w)\). Use \(\mu_{\max}(M_{H,u})\le\mu_{\max}(M_u)\) and \(c_1mB_0\ge0\).

In the arithmetic family, the weight-sum bound gives
\[
 \mathcal T\le Cmn_Ta\log X'+b^{\delta(x)}\le b^2.
\]
Indeed \(\delta(x)\le1\) eventually, and \(m,n_T\le2x^3\), \(\log X'\le x\) give \(Cmn_Ta\log X'\le4|C|x^8 b\). Take the threshold so that \(4|C|x^8+1\le b\). The size bounds also give \(16f\le b\), \(4\le b\), and \(b^3\le T_u\). They imply \(2f/T_u\le1/2\) and \(4fb^2/T_u\le1/4\). Since \(b\le\sqrt{T_u}\) and \(\log(T_u-1)\le2\sqrt{T_u}\), the other error is at most \(T_u/(8(T_u-2))\le1/4\).

The divided lattice spans \(V_u\), and the column identification gives \(d_u=(D.\mathrm{mem}\,u).A.\mathrm{du}\). Apply \ref{B5.05a} with \(F=S_T\). Since \(X'\le T_u\), the selection at \(T_u\) is empty and \(\tau_w=0\) for \(w\mid T_u\). Away from \(T_u\), \ref{B5.20a} supplies the length and containment bounds, with \(\lfloor e_w\tau_w\rfloor_+=e_w\tau_w\) by \ref{B3.52}. The remaining error
\[
 R_u=2\sum_{w\in S}\ell_w\bigl(e_w\sg(u,w)-\lfloor e_w\sg(u,w)\rfloor\bigr)
 \le2(\log4+3+C_T+C_\Gamma)X''
\]
is the ideal-twist rounding error of \ref{B4.50}, bounded in \ref{B4.54}. Take \(c_1\ge\max\{36,2(\log4+3+C_T+C_\Gamma)\}\).
\end{proof}

\section{The height of a chain}

Let \(q\) map the test space onto its Hermitian quotient \(Q\) and \(\iota:Q\to W\) be the base comparison isomorphism. Let \((Q_p)_{p\in I}\) be a spanning family of spectral subspaces with \(I\) finite and nonempty. For \(z\in Q_p\) and \(y\in Q\), the metric identity is
\[
 \langle\iota z,\iota y\rangle_v
 =\kappa_v\prod_{j=1}^{m}\lambda_{jpv}\,\langle z,y\rangle_v,
 \qquad \kappa_v,\lambda_{jpv}>0 .
\]

\begin{theorem}[The archimedean chain estimate]\label{B5.07}
Let \(B_u\ne0\). Suppose its pullback to the test space factors as a chain of kernel maps \(K_i\) followed by the terminal map \(P_u^{\rm term}\). Let \(i(j)\) assign layer \(j\) to a kernel step. At every infinite place assume
\[
 \|P_u^{\rm term}\|_v^2\le e^{h_P},\qquad
 \|K_i\|_v^2\le\exp\left(\sum_{j:i(j)=i}h_{K,jv}\right).
\]
Then
\[
 \log\|B_u\|_v^2\le z_v
 :=h_P+\sum_jh_{K,jv}-\log\kappa_v-\min_{p\in I}\sum_j\log\lambda_{jpv}.
\]
If \(\sum_jL_{6,j}\le\langle\min_{p\in I}\sum_j\log\lambda_{jpv}\rangle_\infty\), then
\[
 h(B_u)\le\langle z_v\rangle_\infty
 \le h_P+\sum_j\langle h_{K,jv}\rangle_\infty
 -\sum_jL_{6,j}-\langle\log\kappa_v\rangle_\infty .
\]

For the chain and height table of \ref{B8.36} on the realization of \ref{B8.38}, use the common size constant \(C_{\rm sz}\). For the archimedean division data \(Z\), assume the height and metric identifications
\[
 D.\mathrm{div}.h(u)=\mathrm{hExpl6}(D,L,Z,C_{\rm sz},u),
 \qquad Z.\mathrm{kapR}(\sigma)=\kappa_{[\sigma]}\quad(\sigma:K\hookrightarrow\C).
\]
Here \(L\) in \(\mathrm{hExpl6}\) denotes the size datum and \([\sigma]\) is the infinite place of \(\sigma\). For a member of \(F_m\) at \(x>F_m.x_1\), \Lean{Fm.\allowbreak{}h\_\allowbreak{}eq} gives the first identity. Then \(h(B_u)\le h_u=D.\mathrm{div}.h(u)\), where
\[
 h_u=H_{P,u}+\sum_jH_{K,uj}-\sum_jL_{6,j}-\langle\log\kappa_v\rangle_\infty .
\]
With \(S_c,\eta,t_{\rm ch}\) as in \ref{B4.03}, put \(N_j=k+j(t_{\rm ch}+1)\), \(n_V=k+m-1+t_{\rm ch}\), and \(v_E=\operatorname{rd}(E)/2\). Write \(\nu_{uj\rho i}\) for the row weights, \(\varphi(t)=t\log t-t\) for \(t>0\), and \(\varphi(0)=0\). For the quadratic character \(\chi_d\) modulo \(d\), put \(D_1=\Re(L'(1,\chi_d)/L(1,\chi_d))\) and \(\log\Omega=\gamma+\log(\pi/v_E)-D_1\) as in \ref{B5.07x}. The budgets are
\[
\begin{aligned}
 H_{P,u}
 &=-S_a\log T_u+C_{\rm sz}\bigl(n_Tb'N_m+f(m\eta+S_c)N_m+fN_m^2\bigr)L_a,\\
 H_{K,uj}
 &=\sum_\rho\mathrm{kerPlace}_{uj\rho}+C_{\rm sz}fN_m^2(L_a+m^3L),\\
 \mathrm{kerPlace}_{uj\rho}
 &=\sum_{i=0}^{N_{j-1}-1}
 \bigl(\varphi(\nu_{uj\rho i})-\nu_{uj\rho i}\log\Omega\bigr),\\
 L_{6,j}
 &=\sum_\rho\log\operatorname{Hook}(a^+_{j\rho})-C_{\rm sz}fn_V^2L',
 \qquad L'=\log(b\,\mathrm{rd}E\,g_0g_1 n_V).
\end{aligned}
\]
The plus weight \(a^+_{j\rho}\) is the tuple \((\kappa^{\rm big}_{j,\rho,i})_{0\le i<k+j-1}\) of \ref{B4.02}. For a plus weight \(\alpha\) of length \(p\),
\[
 \log\operatorname{Hook}(\alpha)
 =\sum_{i=0}^{p-1}\log\Gamma(\alpha_i+p-i)
 -\sum_{0\le i<l<p}\log(\alpha_i-\alpha_l+l-i).
\]
\end{theorem}
\leanref{Formalism.\allowbreak{}L2.\allowbreak{}DIV.\allowbreak{}height\_\allowbreak{}chain\_\allowbreak{}atR}{Formalism.\allowbreak{}LocalToGlobal.\allowbreak{}Divisibility.\allowbreak{}ConstantGeTwentyEight}{123--183}
\leanref{Formalism.\allowbreak{}BasicNotions.\allowbreak{}DV.\allowbreak{}hExpl6}{Formalism.\allowbreak{}LocalRingCardinality.\allowbreak{}KernelPlaceMeanBound}{25--28}
\leanref{Formalism.\allowbreak{}BasicNotions.\allowbreak{}DV.\allowbreak{}logHook}{Formalism.\allowbreak{}LocalRingCardinality.\allowbreak{}KernelPlaceMeanBound}{9--11}
\leanref{Formalism.\allowbreak{}BasicNotions.\allowbreak{}DV.\allowbreak{}Lprime}{Formalism.\allowbreak{}LocalRingCardinality.\allowbreak{}KernelPlaceMeanBound}{13--14}
\leanref{Formalism.\allowbreak{}BasicNotions.\allowbreak{}DV.\allowbreak{}L6lo}{Formalism.\allowbreak{}LocalRingCardinality.\allowbreak{}KernelPlaceMeanBound}{16--18}
\leanref{Formalism.\allowbreak{}BasicNotions.\allowbreak{}DV.\allowbreak{}kernelPlaceMeanBound}{Formalism.\allowbreak{}LocalRingCardinality.\allowbreak{}KernelPlaceMeanBound}{20--23}

\begin{proof}
The \(Q_p\) span \(Q\) and lie in eigenspaces of \(\iota^*\iota\), so \(\iota^*\iota\ge a_v\,\mathrm{id}_Q\), where \(a_v=\kappa_v\min_p\prod_j\lambda_{jpv}\). The same lower bound on squared norms holds for the dual map \(\iota^\vee\), and \(q^\vee\) is isometric. Thus
\[
 a_v\|B_u\|_v^2\le \|(\iota q)^\vee B_u\|_v^2
 \le\|P_u^{\rm term}\|_v^2\prod_i\|K_i\|_v^2 .
\]
Take logs and average. For the actual chain use \(h_P=A\) and \(h_{K,jv}=s_{jv}\) from \ref{B5.07a} and the spectral subtraction of \ref{B5.07b}. Each infinite place has two embeddings. The metric identification therefore turns the embedding mean of \(\log\mathrm{kapR}\) into \(\langle\log\kappa_v\rangle_\infty\), and the height identity gives the bound.
\end{proof}

\begin{lemma}[Terminal and kernel budgets]\label{B5.07a}
Use the height table of \ref{B8.35} with \(N=N_m\ge k+m\), \(k,m\ge1\), \(\#\mathrm{Pl}\le f\), \(1\le L\le L_a\), and the lower bounds there on \(C_h=C_{\rm sz}\). Its terminal entry \(A\) and layer entries \(s_{jv}\) satisfy \(A\le H_{P,u}\) and \(\langle s_{jv}\rangle_\infty\le H_{K,uj}\). In the frames of \ref{B8.36},
\[
 \|P_u^{\rm term}\|_v^2\le e^A,\qquad
 \|K_i\|_v^2\le\exp\left(\sum_{j:i(j)=i}s_{jv}\right).
\]
\end{lemma}
\leanref{Formalism.\allowbreak{}DivisorInequality.\allowbreak{}height\_\allowbreak{}bound\_\allowbreak{}at\_\allowbreak{}set}{Formalism.\allowbreak{}DivisorInequality.\allowbreak{}DivisorInequalityProofs.\allowbreak{}EvenDegreeWeightBounds}{154--288}
\leanref{Formalism.\allowbreak{}DivisorInequality.\allowbreak{}UniformCovolumeTail.\allowbreak{}stepDiSharp}{Formalism.\allowbreak{}DivisorInequality.\allowbreak{}DivisorInequalityProofs.\allowbreak{}EvenDegreeWeightBounds}{27--149}
\leanref{Formalism.\allowbreak{}L2.\allowbreak{}DIV.\allowbreak{}nK\_\allowbreak{}of\_\allowbreak{}cb}{Formalism.\allowbreak{}LocalToGlobal.\allowbreak{}Divisibility.\allowbreak{}DivisionHeightBoundStatement}{565--584}

\begin{proof}
Put \(U=S_c+m\eta\) and \(I_\rho={\bf1}_{TT}(\rho)\). The interlacing bounds of \ref{B8.35} give
\[
 D_{0,\rho}\le U+2I_\rho b',\quad
 D_{1,\rho}\le4U+6I_\rho b',\quad
 -n_\rho\le-I_\rho a_\rho+(N-1)/2+U+2I_\rho b'.
\]
Multiply the last error by \(0\le\log T_u\le C_TNL_a\) and use \(\log(2a)\le2L_a\) in the determinant terms. With the evaluation error this gives
\[
 A+S_a\log T_u\le L_a\bigl((2C_T+24)n_Tb'N+(C_T+16)fUN
 +(C_T/2+1+C_{\rm ev})fN^2\bigr).
\]
Each coefficient is at most \(C_h\). In the mean of \eqref{eq:B8-height-size}, the period term contributes \(-n_{j\rho}\log\Omega\). Since \(g_j^0\le N^2\), the logs with coefficients \(c_5,c_6\) cost at most \((c_5+2c_6)fN^2L_a\). The other bounds there give
\[
\begin{aligned}
 \langle s_{jv}\rangle_\infty-\sum_\rho\mathrm{kerPlace}_{uj\rho}
 \le fN^2\bigl(& (C_\theta+c_5+2c_6+1+2C_4+2C_eC_8)L_a\\
 &+2(C_2+C_4)(C_T+C_8)m^3L\bigr).
\end{aligned}
\]
Both coefficients are at most \(C_h\). Point evaluation gives the terminal norm bound. Cauchy--Schwarz in the weighted coefficient identities gives \(\|K_i\|_v^2\le\mathrm{CB.sq}(i,v)\). A zero coefficient bound forces \(K_i=0\); a positive bound gives the result by exponentiating its log estimate.
\end{proof}

\begin{lemma}[Subtraction on the base spectrum]\label{B5.07b}
On the realization of \ref{B8.38}, use its layer size bounds and hook scale \(C_{\rm sz}\ge8\). At every infinite place \(v\),
\[
 \sum_jL_{6,j}\le\min_{p\in I}\sum_j\log\lambda_{jpv}.
\]
The left side is \(\sum_jH_{X,j}\) in the record of \ref{B5.08}.
\end{lemma}
\leanref{GlobalTheory.\allowbreak{}LowRow.\allowbreak{}scalar\_\allowbreak{}bound\_\allowbreak{}lowC}{Development.\allowbreak{}AbelianVarietyTheory.\allowbreak{}Proofs.\allowbreak{}DivisorProofs.\allowbreak{}AdelCOfInnerProductFormula}{52--87}
\leanref{UpperTheory.\allowbreak{}Mu.\allowbreak{}numeric\_\allowbreak{}bound}{Development.\allowbreak{}AutomorphicArithmetic.\allowbreak{}CyclotomicBounds.\allowbreak{}RootsOfUnity.\allowbreak{}MassNumeric}{6--81}
\leanref{GlobalTheory.\allowbreak{}LowRow.\allowbreak{}lamLow\_\allowbreak{}of\_\allowbreak{}weak}{Development.\allowbreak{}AbelianVarietyTheory.\allowbreak{}Proofs.\allowbreak{}DivisorProofs.\allowbreak{}LiftEigenvalueLowerBoundOfWeakCharacterisation}{20--41}
\leanref{Formalism.\allowbreak{}L2.\allowbreak{}DIV.\allowbreak{}levelSix\_\allowbreak{}fixed\_\allowbreak{}eq}{Formalism.\allowbreak{}LocalToGlobal.\allowbreak{}Divisibility.\allowbreak{}EigenvalueSumLeClosedOf}{50--56}
\leanref{Formalism.\allowbreak{}L2.\allowbreak{}Indef.\allowbreak{}IndefColsR.\allowbreak{}log\_\allowbreak{}eigenvalue\_\allowbreak{}sum\_\allowbreak{}le\_\allowbreak{}min}{Formalism.\allowbreak{}LocalToGlobal.\allowbreak{}IndefiniteForms.\allowbreak{}PairingConditionOfInverseMatrixIdentity}{328--331}

\begin{proof}
Put \(H_j^+=\sum_\rho\log\operatorname{Hook}(a^+_{j\rho})\), \(D_j=H_j^+-8fn_V^2L'\) and \(r=k+m-1\). At the distinguished embedding the lift times its adjoint acts by a positive scalar on each occurring Hecke component. In \ref{B8.13}, the index times the squared norm of the test vector is \(d_S^{-1}\), so the finite factor is \(E/(d_Od_S)\), with \(E\ge e^{-2rf}\) and \(0<d_O,d_S\le e^{4rf}\). The mass and formal degree, including the roots of unity, cost at most \(2fn_V^2L'+4n_Vf\) in the log. Both hook products enter with positive sign, and the minus-weight hooks are positive integers. Thus the log eigenvalue is at least
\[
 H_j^+-10rf-2fn_V^2L'-4n_Vf\ge D_j
\]
by \(2r\le n_V,\ n_V\ge8,\ L'\ge5.5\). The Hecke components span the compressed space, so \(e^{D_j}\) bounds the quadratic form of the lift times its adjoint from below there.

At another place the conjugate-lift estimate on a represented block gives \(\|y\|^2=\lambda_{jpv}a\) with \(a\ne0\) and \(\|y\|^2\ge e^{D_j}\Re a\). Positivity of \(\lambda_{jpv}\) gives \(a>0\) and \(\log\lambda_{jpv}\ge D_j\). On an absent block the prescribed scalar is \(\exp(L_{6,j})\). In either case \(L_{6,j}\le\log\lambda_{jpv}\). Sum at fixed \(p\) and take the minimum. Reindex the lifts to obtain the sum in the record.
\end{proof}

\begin{lemma}[Metric of a CM generator]\label{B5.07x}
Let \(K_d=\Q(\sqrt{-d})\) have discriminant \(-d<-4\). Take a nonzero integral ideal \(\mathfrak c\) with norm prime to \(2d\) and put \(\Lambda=\delta\mathfrak c\), where \(\delta^{12}=\operatorname{disc}(\mathfrak c)\) is its Weierstrass discriminant. Let \(M\subset\C\) be a number field containing \(g_2(\Lambda),g_3(\Lambda)\) and the prescribed images of \(E\). If \(s_\tau\in M^\times\) generates the host differential ideal at \(\tau\) over \(\mathcal O_M\), then
\[
 \frac1{[M:\Q]}\sum_{\sigma:M\hookrightarrow\C}
 \log\frac{\operatorname{area}(\Lambda_\sigma)}{\pi|\sigma s_\tau|^2}
 =\gamma+\log(\pi/v_E)-D_1 ,
\]
where \(\Lambda_\sigma\) has invariants \(\sigma g_2(\Lambda),\sigma g_3(\Lambda)\).
\end{lemma}
\leanref{GlobalTheory.\allowbreak{}MeanPeriodClosed.\allowbreak{}meanPeriodAt}{Development.\allowbreak{}AbelianVarietyTheory.\allowbreak{}Proofs.\allowbreak{}DivisorProofs.\allowbreak{}MeanLogPeriodChowlaSelberg}{58--82}
\leanref{GlobalTheory.\allowbreak{}MeanPeriodGenerator.\allowbreak{}mean\_\allowbreak{}logMet\_\allowbreak{}eq}{Development.\allowbreak{}AbelianVarietyTheory.\allowbreak{}Proofs.\allowbreak{}DivisorProofs.\allowbreak{}MeanLogMetricGeneratorEqMeanPeriod}{45--57}

\begin{proof}
The logarithmic form of the Chowla--Selberg product formula gives the mean of \(\log(4\pi\,\Im z_C\,|\eta(z_C)|^4)\) over ideal-class points \(z_C\) as \(\gamma+\log(2\pi/\sqrt d)-D_1\), with Dedekind eta and the Kronecker character \(\chi_d\) \cite{selberg-1967}. In the normalization \(g_2^3-27g_3^2=1\), the modular discriminant formula gives \(\operatorname{area}/\pi=4\pi\,\Im z\,|\eta(z)|^4\). CM conjugacy gives equal class multiplicities in a common extension. Scaling discriminant \(u\in\mathcal O_M^\times\) to one adds \(\tfrac16\log|\sigma u|\) to the log area, with mean zero by the product formula. Restricting embeddings preserves the mean.

The host ideal satisfies \(\mathfrak b\bar{\mathfrak b}=\mathfrak D_{E/K_d}\). The differential ideal extends \(\mathfrak b^{-1}\mathfrak D_{E/K_d}\) under \(\tau\). If \(B\) is the absolute norm of this ideal over \(E\), then \(\log|\mathrm N_{M/\Q}s_\tau|/[M:\Q]=\log B/[E:\Q]\). We have \(B^2=\mathrm N\mathfrak D_{E/K_d}\) and \(|\operatorname{disc}E|=d^{[E:K_d]}\mathrm N\mathfrak D_{E/K_d}\). Thus the squared generator term replaces \(\sqrt d\) by \(\operatorname{rd}(E)=2v_E\) in the denominator. Multiplication by \(-\sum_i\nu_i\) gives the kernel period term in \ref{B5.07}. The corresponding flat line contribution for rows is in \ref{B8.24a}.
\end{proof}

\section{Combining the divisor and height budgets}

Let \(J,\mathcal P\) be finite with \(\#J=m,\ \#\mathcal P\le f\). Fix \(\mathcal U\subseteq\mathcal P\) of size \(n_T\). The real size parameters of the numerical record satisfy
\[
\begin{gathered}
 x>1,\quad\varepsilon>0,\quad L=x^\varepsilon,\quad b=e^L,\quad a=xb,\quad
 r_d\ge1,\quad \log r_d,\log x\le L,\quad L_a=\log(ar_d),\\
 B_0=fmbL,\quad N_c\ge0,\quad m,f,k\ge1,\quad k+m\le n_*,\\
 S_c,t_c,\eta,b',w'',\lambda_0,n_V,L_p\ge0,\qquad n_V\le n_* .
\end{gathered}
\]
Here \(L_p\) is the hook-error parameter. Fix \(r\in\N\), real \(\log T,\log\Omega\) and weights \(a_\rho\) for \(\rho\in\mathcal P\). For \(j\in J,\rho\in\mathcal P\), use real row weights \(\nu_{j\rho i}\), \(0\le i<r\), a log hook budget \(\log H_{j\rho}\), and real \(\beta_j,w_{j\rho},e_{j\rho},\Delta_{j\rho}\). The coefficient \(\beta_j\) is \(\Lean{pq j}\) in the record. For nonzero real \(t\), set \(\varphi(t)=t\log|t|-t\). Assume
\[
\begin{gathered}
 0\le w_{j\rho}\le2b,\qquad |\beta_j-k|\le m,\qquad|\log\Omega|\le L_a,\\
 m\lambda_0\le b,\quad n_*^2\le b,\quad
 m^3n_*^2L\le b,\quad m\eta n_*L\le b,\quad kn_*\le b,\\
 kS_c\le w'',\quad m(t_c+1)\eta\le w'',\quad
 n_*w''\le b,\quad b'n_*\le a,\quad mkb'\le a,\\
 1\le C_0\le c_e,\qquad 0\le C_h\le C_0,\qquad C_hL_p\le C_0L_a .
\end{gathered}
\]
For each \(j,\rho\), write \(I_\rho={\bf1}_{\mathcal U}(\rho)\) and assume
\[
\begin{aligned}
 \sum_i\varphi(\nu_{j\rho i})-\log H_{j\rho}&\le I_\rho\varphi(a_\rho)+e_{j\rho},\\
 \sum_i\nu_{j\rho i}&=I_\rho a_\rho+kw_{j\rho}+\Delta_{j\rho},\\
 e_{j\rho}&\le C_0\bigl(w_{j\rho}+S_c+n_*^2+m(t_c+1)\eta\bigr)L_a,\\
 |\Delta_{j\rho}|&\le C_0\bigl(w_{j\rho}+k(n_*+S_c)+m(t_c+1)\eta+I_\rho kb'\bigr).
\end{aligned}
\]
Define
\[
\begin{aligned}
 H_P&=-\sum_{\rho\in\mathcal U}a_\rho\log T
 +C_h\bigl(n_Tb'n_*+f(m\eta+S_c)n_*+fn_*^2\bigr)L_a,\\
 H_{B,j}&=\sum_{\rho,i}\bigl(\varphi(\nu_{j\rho i})-\nu_{j\rho i}\log\Omega\bigr)
 +C_hfn_*^2(L_a+m^3L),\\
 H_{X,j}&=\sum_\rho\log H_{j\rho}-C_hfn_V^2L_p,\qquad
 H^*=H_P+\sum_j(H_{B,j}-H_{X,j}),\\
 \mathcal D(T)&=\sum_{\rho\in\mathcal U}
 \bigl(m(\varphi(a_\rho)-a_\rho\log\Omega)-a_\rho\log T\bigr).
\end{aligned}
\]

\begin{lemma}[The divisor slope and the height budget]\label{B5.08}
For the preceding data suppose \(\widetilde\mu_C,\mu_0,F_*,L_{\rm pass},E_b\) satisfy
\[
\begin{gathered}
 \widetilde\mu_C\le\mu_0+L_{\rm pass},\qquad
 \mu_0\le\sum_{j,\rho}\beta_jw_{j\rho}\log\Omega+2F_*+E_b,\\
 L_{\rm pass}\le1+C_0m\lambda_0,\qquad E_b\le c_emL(B_0+N_c).
\end{gathered}
\]
Then
\[
 \widetilde\mu_C+H^*-2F_*
 \le \mathcal D(T)+(50C_0+10)n_TaL_a+60c_emL(B_0+N_c).
\]
\end{lemma}
\leanref{Formalism.\allowbreak{}L2.\allowbreak{}DIV.\allowbreak{}RecNums.\allowbreak{}divIneq}{Formalism.\allowbreak{}LocalToGlobal.\allowbreak{}Divisibility.\allowbreak{}DivisionHeightBoundStatement}{131--181}
\leanref{Formalism.\allowbreak{}DivisorInequality.\allowbreak{}slope\_\allowbreak{}bounds}{Formalism.\allowbreak{}DivisorInequality.\allowbreak{}DivisorInequalityProofs.\allowbreak{}BaseLevelBoundsWithCovolumeHolds}{47--70}
\leanref{Formalism.\allowbreak{}DivisionValues.\allowbreak{}lemmaRb\_\allowbreak{}of}{Formalism.\allowbreak{}QuotientMetricDivisors.\allowbreak{}HermitianTwistDivisors.\allowbreak{}ReductionLemmaSecond}{163--287}
\leanref{Formalism.\allowbreak{}L2.\allowbreak{}Indef.\allowbreak{}IndefSide.\allowbreak{}boundsMu}{Formalism.\allowbreak{}LocalToGlobal.\allowbreak{}IndefiniteForms.\allowbreak{}IndefUse}{271--335}

\begin{proof}
The weight relations give
\[
\begin{aligned}
 &\sum_i\varphi(\nu_{j\rho i})-\log H_{j\rho}
 -\sum_i\nu_{j\rho i}\log\Omega+\beta_jw_{j\rho}\log\Omega\\
 &\qquad\le I_\rho\bigl(\varphi(a_\rho)-a_\rho\log\Omega\bigr)
 +e_{j\rho}+\bigl((\beta_j-k)w_{j\rho}-\Delta_{j\rho}\bigr)\log\Omega .
\end{aligned}
\]
Summing the indicator terms and adding the terminal main term gives \(\mathcal D(T)\). Put \(Z=n_TaL_a\) and \(V=mL(B_0+N_c)\), so \(L\le L_a\le3L\) and \(V\ge fm^2bL^2\ge b\).
\[
\begin{array}{rcl}
 L_{\rm pass}&\le&2C_0V,\\
 E_b&\le&c_eV,\\
 C_h\bigl(n_Tb'n_*+f(m\eta+S_c)n_*+fn_*^2\bigr)L_a&\le&C_0Z+9C_0V,\\
 \displaystyle\sum_{j,\rho}e_{j\rho}&\le&15C_0V,\\
 \displaystyle\sum_{j,\rho}|(\beta_j-k)w_{j\rho}|L_a&\le&6V,\\
 \displaystyle\sum_{j,\rho}|\Delta_{j\rho}|L_a&\le&C_0Z+15C_0V,\\
 mC_hfn_*^2(L_a+m^3L)&\le&4C_0V,\\
 mC_hfn_V^2L_p&\le&3C_0V .
\end{array}
\]
For the terminal error use \(b'n_*\le a\) and \(m\eta n_*,S_cn_*,n_*^2\le b\). In the \(\Delta\)-sum the indicator term costs at most \(C_0mn_Tkb'L_a\le C_0Z\), and the rest costs at most \(5C_0bL_a\) per pair. The kernel and hook terms use \(m^3n_*^2L\le b,\ n_V\le n_*\), and \(C_hL_p\le C_0L_a\). The total is at most \(2C_0Z+55c_eV\).

For the row application set \(\mu_0=\mu(R_0^\vee)-\langle\log\kappa_{\rm base}\rangle_\infty\). Combine \ref{B5.08x} with \(\sum_j l_{3,j}\le\sum_{j,\rho}\beta_jw_{j\rho}\log\Omega+E_{l3}\) from \ref{B4.55}. This gives the second slope inequality with \(E_b=E_{l3}+L_G+\sum_j(E_{1,j}+E_{2,j})+E_3\).
\end{proof}

\begin{lemma}[The two local indices of a change of row type]\label{B5.08a}
Take the reference row system of \ref{B4.55}, with base type \(t_0\) and glued lattice \(R_0\). Let \(\Lambda_t\) be the quotient lattice of type \(t\) and \(L_t=\iota_t\Lambda_t\) its image. Let \(A_{jt}=\theta_{jt}T_{jt}\) be the layer image in the space of the reference lattice \(M_j\). For commensurable lattices put
\[
 [A:B]_w=\operatorname{length}_w(A/(A\cap B))-\operatorname{length}_w(B/(A\cap B)),
\]
with lengths after localization at \(\mathfrak p_w\). Write \(\bar w\) for the conjugate place and \(\sharp\) for the Hermitian dual for the layer pairing. Let \(G_{t,w}\in K\) represent at the distinguished embedding \(v_0\) the quotient Gram determinant in a basis of \(\Lambda_t\) at \(w,\bar w\). Put
\[
 C_w(t)=v_w(G_{t,w})+
 \sum_j\bigl([M_j:A_{jt}]_w+[M_j:A_{jt}]_{\bar w}
 +[T_{jt}:T_{jt}^{\sharp}]_w\bigr).
\]
For the selected type \(t_w\) one has
\[
 [L_{t_0}:R_0]_w+[L_{t_0}:R_0]_{\bar w}=C_w(t_w)-C_w(t_0).
\]
\end{lemma}
\leanref{Formalism.\allowbreak{}DivisionValues.\allowbreak{}StepBAux.\allowbreak{}idx\_\allowbreak{}change\_\allowbreak{}img}{Formalism.\allowbreak{}QuotientMetricDivisors.\allowbreak{}HermitianTwistDivisors.\allowbreak{}ReductionSecondStep}{379--399}
\leanref{Formalism.\allowbreak{}DivisionValues.\allowbreak{}ConjugateWeightGram.\allowbreak{}local\_\allowbreak{}idW}{Formalism.\allowbreak{}QuotientMetricDivisors.\allowbreak{}HermitianTwistDivisors.\allowbreak{}ConjugateWeightLatticeOfGeneral}{174--318}
\leanref{Formalism.\allowbreak{}DivisionValues.\allowbreak{}stepBW\_\allowbreak{}of}{Formalism.\allowbreak{}QuotientMetricDivisors.\allowbreak{}HermitianTwistDivisors.\allowbreak{}ConjugateWeightLatticeOfGeneral}{324--499}

\begin{proof}
Fix \(w\) and put \(t=t_w\). Localize away from \(\mathfrak p_w\cup\mathfrak p_{\bar w}\), a semilocal principal ideal domain, and choose bases of the quotient lattices there. Let \(a\) express by its determinant the image basis of \(L_t\) in that of \(L_{t_0}\). By \ref{B4.36a}, \(t_{\bar w}=t_w=t\), so the image basis of \(L_t\) generates \(R_0\) at both places. The elementary divisors give \([L_{t_0}:R_0]_z=v_z(a)\) for \(z=w,\bar w\). For layer \(j\) choose bases of \(T_{j,t_0},T_{jt}\) with Gram determinants \(h_{j0},h_{jt}\). The change determinant \(b_j\) between their image bases has
\[
 v_z(b_j)=[M_j:A_{jt}]_z-[M_j:A_{j,t_0}]_z\qquad(z=w,\bar w).
\]
The semilocal ring is stable under conjugation. Inverse Gram matrices give \([T_{jt}:T_{jt}^\sharp]_w=-v_w(h_{jt})\) and the corresponding formula at \(t_0\).

At the base type the squared determinant ratio at \(v\) is \(\kappa_v^{D_C}\) times the product of the layer ratios, where \(D_C\) is the row rank. At \(v_0\) the identity holds for every type. A spectral summand of dimension \(m_\alpha\) contributes \((\kappa_{v_0}\prod_j\lambda_{j,\alpha,v_0})^{m_\alpha}\) to the row ratio and \(\lambda_{j,\alpha,v_0}^{m_\alpha}\) to layer \(j\). Multiply over the summands to obtain the identity with type scalar one.
\leanref{Formalism.\allowbreak{}L2.\allowbreak{}Indef.\allowbreak{}IndefRows.\allowbreak{}detId}{Formalism.\allowbreak{}LocalToGlobal.\allowbreak{}IndefiniteForms.\allowbreak{}IndefUse}{196--206}

A basis change multiplies the Gram determinant by the image of \(\bar a a\) at \(v_0\). Dividing the two determinant-ratio identities cancels \(\kappa_{v_0}^{D_C}\) and gives
\[
 \bar a a\,G_{t_0,w}\prod_j h_{jt}=G_{t,w}\prod_j(\bar b_j b_j h_{j0})
\]
in \(K\), since the embedding is injective and respects involution. Take \(v_w\) and use \(v_w(\bar b)=v_{\bar w}(b)\). The row index occurs twice at a self-conjugate place.
\end{proof}

\begin{lemma}[The reference row divisor]\label{B5.08x}
Take the well-formed reference side of \ref{B4.55} with the selector of \ref{B4.36a}. Use the untwisted identities and local bounds. Let \(D_C=\operatorname{rk}R_0>0\) and \(\bar\kappa_0=\langle\log\kappa_{\rm base}\rangle_\infty\). Write \(l_{3,j}\) for the upper bound on \(\operatorname{lcv}(M_j)/D_C\), where \(\operatorname{lcv}\) is the log covolume in the slope normalization. Then
\[
 \mu(R_0^\vee)-\bar\kappa_0
 \le\sum_j l_{3,j}+2F_*+L_G+\sum_j(E_{1,j}+E_{2,j})+E_3 .
\]
\end{lemma}
\leanref{Formalism.\allowbreak{}DivisionValues.\allowbreak{}ConjugateWeightGram.\allowbreak{}local\_\allowbreak{}idW}{Formalism.\allowbreak{}QuotientMetricDivisors.\allowbreak{}HermitianTwistDivisors.\allowbreak{}ConjugateWeightLatticeOfGeneral}{174--318}
\leanref{Formalism.\allowbreak{}DivisionValues.\allowbreak{}second\_\allowbreak{}step\_\allowbreak{}refined\_\allowbreak{}weak}{Formalism.\allowbreak{}QuotientMetricDivisors.\allowbreak{}HermitianTwistDivisors.\allowbreak{}ConjugateWeightLatticeOfGeneral}{506--517}
\leanref{Formalism.\allowbreak{}L2.\allowbreak{}Indef.\allowbreak{}IndefSide.\allowbreak{}boundsMu}{Formalism.\allowbreak{}LocalToGlobal.\allowbreak{}IndefiniteForms.\allowbreak{}IndefUse}{271--335}
\leanref{Formalism.\allowbreak{}L2.\allowbreak{}DIV.\allowbreak{}DivInputsDef3.\allowbreak{}toDivInputs0R}{Formalism.\allowbreak{}LocalToGlobal.\allowbreak{}Divisibility.\allowbreak{}TwistedPlaceEtaSumLeExtended}{247--287}

\begin{proof}
Since \(\ell_{\bar w}=\ell_w\), \ref{B5.08a} gives
\[
 \operatorname{lcv}(R_0)-\operatorname{lcv}(L_{t_0})
 =2\sum_w\ell_w[L_{t_0}:R_0]_w
 =\sum_w\ell_w\bigl(C_w(t_w)-C_w(t_0)\bigr).
\]
For the base quotient map \(P\) and test Gram matrix \(H\), put \(g=\det(PH^{-1}P^*)^{-1}\) in the standard quotient basis. The star here is over \(K\). The determinant is nonzero because \(P\) has full row rank and the distinguished embedding respects involution. The entries of \(H\) lie in the subfield where all embeddings respect involution, and \(H\) is positive definite at infinity.

For \(\sigma:K\hookrightarrow\C\), put \(\sigma'(a)=\overline{\sigma(\bar a)}\). Write \(d_\sigma=\det(P_\sigma H_\sigma^{-1}P_\sigma^*)^{-1}\), now with complex adjoint. Since \(H_{\sigma'}=H_\sigma\), the mixed Gram determinant of the rows of \(P_\sigma,P_{\sigma'}\) in the metric \(H_\sigma^{-1}\) is \(\sigma(g)^{-1}\). Cauchy--Schwarz on the exterior products gives \(d_\sigma d_{\sigma'}\le|\sigma(g)|^2\). The permutation \(\sigma\mapsto\sigma'\) and the product formula bound the mean of \(\log d_\sigma\) by \(\sum_w\ell_wv_w(g)\). A change determinant \(c\) to the semilocal quotient basis changes the valuation by \(v_w(c)+v_{\bar w}(c)\). The weighted sum gives the change in covolume, so \(\operatorname{lcv}(\Lambda_{t_0})\le\sum_w\ell_w v_w(G_{t_0,w})\).

Each base layer has its algebraic Hermitian metric at all embeddings. The product formula gives \(\operatorname{lcv}(T_{j,t_0})=-\sum_w\ell_w[T_{j,t_0}:T_{j,t_0}^\sharp]_w\). For the squared determinant ratio \(D_{jv}\) of \(\theta_{j,t_0}\), the image and index formulas give
\[
 \langle\log D_{jv}\rangle_\infty
 =\operatorname{lcv}(M_j)+
 \sum_w\ell_w\bigl(2[M_j:A_{j,t_0}]_w+[T_{j,t_0}:T_{j,t_0}^\sharp]_w\bigr).
\]
The covolumes satisfy \(\operatorname{lcv}(L_{t_0})=\operatorname{lcv}(\Lambda_{t_0})+D_C\bar\kappa_0+\sum_j\langle\log D_{jv}\rangle_\infty\). Add the index formula for \(R_0\) to cancel the base-type finite terms. With \(S_y\) containing the support of the selected terms,
\[
 \mu(R_0^\vee)=\frac{\operatorname{lcv}(R_0)}{D_C}
 \le \bar\kappa_0+\sum_j\frac{\operatorname{lcv}(M_j)}{D_C}
 +\frac1{D_C}\sum_{w\in S_y}\ell_w C_w(t_w).
\]

Put \(i_{jw}=[M_j:A_{j,t_w}]_w+[M_j:A_{j,t_w}]_{\bar w}\), \(q_{jw}=[T_{j,t_w}:T_{j,t_w}^\sharp]_w\), and \(G_w=G_{t_w,w}\). The Gram sum divided by \(D_C\) is at most \(L_G\). The layer terms vanish outside \(\mathrm{Tw}\subseteq S_y\), and there the bounds of \ref{B4.55} give
\[
\begin{aligned}
 D_C^{-1}\sum_j i_{jw}
 &\le \sum_j(2e_wF^0_{jw}+\mathrm{err}_{jw})+2e_w/p_w^2,\\
 D_C^{-1}q_{jw}&\le e_wJ^0_{jw}+\eta_{jw}/D_C,\\
 \sum_{w\in\mathrm{Tw}}\ell_w(\mathrm{err}_{jw}+e_wJ^0_{jw})&\le E_{1,j},
 &\sum_{w\in\mathrm{Tw}}\ell_w\eta_{jw}/D_C&\le E_{2,j}.
\end{aligned}
\]
Thus the layer terms cost at most \(2F_*+\sum_j(E_{1,j}+E_{2,j})+E_3\), where \(F_*=\sum_{j,w\in\mathrm{Tw}}\ell_we_wF^0_{jw}\) and \(E_3\) bounds \(\sum_{w\in S_y}2\ell_we_w/p_w^2\).

In the arithmetic application \(M\) is the level group order and \(s>0\) the integral denominator. Multiplication by \(Ms\) maps the test dual into the test lattice. On quotient determinants this gives \(v_w(G_{t,w})\le D_Cv_w(Ms)\), so \(L_G=\lambda_0+\log s\) suffices for \(\log M\le\lambda_0\). The twisted pairing index is at most
\[
 D_Ce_wJ^0_{jw}+a_{jw}+D_C\bigl(v_w(d_{j,t_w})+v_w(s)\bigr),
\]
where the first term bounds the cut exponent times the rank, \(a_{jw}\) is the length of the lift lattice modulo its cut, and \(d_{jt}\) is the layer denominator. If \(d_{jt}\mid d_g\) and \(\sum_{w\in\mathrm{Tw}}\ell_wa_{jw}\le D_Cb_j\), the product formula gives \(E_{2,j}=b_j+\log d_g+\log s\). Grouping by rational primes gives \(E_3=1\) by \(\sum_p2\log p/p^2\le1\).
\end{proof}

\section{The four conclusions of division}

\begin{theorem}[The division package]\label{B5.09}
Use the coordinate lattices and valuation hypotheses of \ref{B4.55} for \(N_u^0\) and the shift \(N'_u\) with exponents \(s_{uw}=\lfloor\sigma_w\rfloor\) of finite support. Let \(H^*\) and \(\mathrm{rhs}\) be the height and right-hand side of \ref{B5.08}. Assume that the averaged terminal and kernel bounds after spectral subtraction are at most \(H^*\). Put \(\bar\kappa=\langle\log\kappa_v\rangle_\infty\), and let \(\mu_0=\mu(R_0^\vee)-\langle\log\kappa_{\rm base}\rangle_\infty\) and \(B_*\) be its upper bound in \ref{B5.08x}. Assume the numerical data of \ref{B5.08} and the passage bounds
\[
\begin{gathered}
 \delta_0\ge0,\quad Q_0\ge1,\quad 3\delta_0Q_0\le1,\quad
 B_*+H_0+28\lambda\le Q_0,\quad
 \lambda\le\lambda_0,\quad C_0\ge28,\\
 \mu(RR^\vee)-\bar\kappa
 \le\mu_0+3\delta_0(\mu_0+H_0+28\lambda)+28\lambda .
\end{gathered}
\]
Then the four conclusions of \ref{B4.55} hold for \(B_u,N'_u,RR\) with \(H=H^*\) and \(R_*=\mathrm{rhs}\).
\end{theorem}
\leanref{Formalism.\allowbreak{}L2.\allowbreak{}DIV.\allowbreak{}DivClausesR}{Formalism.\allowbreak{}LocalToGlobal.\allowbreak{}Divisibility.\allowbreak{}ConstantGeTwentyEight}{111--122}
\leanref{Formalism.\allowbreak{}L2.\allowbreak{}DIV.\allowbreak{}divClausesR\_\allowbreak{}of}{Formalism.\allowbreak{}LocalToGlobal.\allowbreak{}Divisibility.\allowbreak{}ConstantGeTwentyEight}{184--248}

\begin{proof}
Since \(\mu_0\le B_*\), the passage allowance is at most \(1+28\lambda\le1+C_0m\lambda_0\). Thus \ref{B5.08} supplies the numerical slope bound in \ref{B4.55}, with passage parameters \((\delta,Q_{\rm pass},h_{\rm r},\lambda_{\rm r})=(\delta_0,Q_0,H_0,\lambda)\). The height and shift clauses apply with the stated budgets. In the arithmetic application, the coordinate period bound of \ref{B4.38} allows \(\sigma_w=s_{uw}\) from \ref{B5.03}.
\end{proof}

\begin{lemma}[The unit budget for row passage]\label{B5.09a}
Use the local error bounds of \ref{B4.55} for a word on the selected arithmetic family. Put \(g=\min(g_0,g_1)\), \(r=k+m-1\), and let \(M\) be the order of the finite row group. Set
\[
 A_3=\sum_{j,\rho}\beta_jw_{j\rho}\log\Omega,\quad
 B=A_3+\Fo+c_emL(B_0+N_c),\quad Q_0=\max\{1,B+x^{C_H}b+28\log M\}.
\]
For fixed family constants and sufficiently large \(x\), the relative multiplicity bound holds with \(\delta_0=8fr/g\), and
\[
 B_*+H_0+28\lambda\le Q_0,\qquad 1\le Q_0\le3b^2,\qquad 3\delta_0Q_0\le1 .
\]
\end{lemma}
\leanref{GlobalTheory.\allowbreak{}ColumnMultSum.\allowbreak{}columnMult}{Development.\allowbreak{}AbelianVarietyTheory.\allowbreak{}Proofs.\allowbreak{}DivisorProofs.\allowbreak{}ReordEqConj}{35--78}
\leanref{Formalism.\allowbreak{}L2.\allowbreak{}DIV.\allowbreak{}Arith4S.\allowbreak{}q\_\allowbreak{}rec}{Formalism.\allowbreak{}LocalToGlobal.\allowbreak{}Divisibility.\allowbreak{}EigenvalueSumLeClosedOf}{1885--1954}
\leanref{Formalism.\allowbreak{}L2.\allowbreak{}DIV.\allowbreak{}hg\_\allowbreak{}of\_\allowbreak{}sizeG}{Formalism.\allowbreak{}LocalToGlobal.\allowbreak{}Divisibility.\allowbreak{}EigenvalueSumLeClosedOf}{295--325}
\leanref{Formalism.\allowbreak{}L2.\allowbreak{}DIV.\allowbreak{}delta\_\allowbreak{}mul\_\allowbreak{}Q\_\allowbreak{}le\_\allowbreak{}one}{Formalism.\allowbreak{}LocalToGlobal.\allowbreak{}Divisibility.\allowbreak{}EigenvalueSumLeClosedOf}{126--139}

\begin{proof}
At each of the \(2f\) selected places let \(q_w=|\mathbf k_w|\) and let \(n_w\) be the local multiplicity. The local character sum divided by the subgroup order is the trace of the character projector. The local estimate for the cuspidal factors gives
\[
 1-q_w^{-1}\le n_w/q_w^k\le1+(2r-1)/q_w .
\]
The retained characters and the projector sum factor by place. Thus \(n/\bar n=\prod_w(n_w/q_w^k)\), with \(\bar n=\prod_wq_w^k\). Since \(q_w\ge g\) and \(4fr/g\le1\), the product lies between \(1-2f/g\) and \(e^{4fr/g}\le1+8fr/g\). This gives \(\delta_0\). The layer bounds and \(2F_*=\Fo\) give \(B_*\le B\). Take \(H_0=x^{C_H}b\) and \(\lambda=\log M\).

Put \(P=1+f+k+m+t_{\rm ch}+N_m+L\). The bounds \(|\beta_j|\le k+m,\ w_{j\rho}\le2b,\ |\log\Omega|\le L_a\le4P\) give \(A_3\le16P^4b\). The place bound \(16P^7b/p\) for the base terms gives at most \(64P^7b\log(X''+1)\) per layer after summing over eligible primes. Since \(m\le P\) and \(\log(X''+1)\le4P\), \(\Fo\le256P^9b\). Also \(L(B_0+N_c)\le2P^4b\) and \(\log M\le2mL(B_0+N_c)\). Hence \(B\le(272+2c_e)P^9b\) and \(\log M\le4P^5b\).

For fixed family constants \(P\) is polynomially bounded in \(x\). Eventually \(B\le b^2,\ x^{C_H}\le b,\ 28\log M\le b^2\), so \(Q_0\le3b^2\). Also \(72fr\le b^6\), and \(g_i\ge(kb)^8\). Hence \(24frQ_0\le72frb^2\le b^8\le g\).
\end{proof}

\begin{theorem}[Division data for the actual period family]\label{B5.10}
Suppose isomorphisms identify the coordinate and actual word families and preserve their lattices and metrics. Assume they preserve the pairings as well. For each word use the clauses of \ref{B5.09} with \(H^*-\bar\kappa=h_u,\ 2F_*=\Fo\), and \(\mathrm{rhs}\le\mathcal D(T_u)+E_{\rm div}\), where
\[
 \mathcal D(T_u)=m\sum_{t\in TT}\bigl(\varphi(a_t)-a_t\log\Omega\bigr)-S_a\log T_u,\qquad
 E_{\rm div}=c_{\rm Div}\bigl(n_TaL_a+m(B_0+N_c)x^{e_{\rm Div}\varepsilon}\bigr).
\]
Assume \(RR\) is finitely generated over \(\mathcal O_K\) and spans \(W\). At each finite place, assume for every word and layer that \(g_{jw}=E_{ujw}=0\) when \(p_w>X''\), \(p_w\notin\{T_u,g_0,g_1,d\}\) and \(p_w\nmid D_{\Gamma,u}\), where \(D_{\Gamma,u}\) is the positive integer divisible by the stabilizer orders.

On \(K\)-vectors, assume the complex period is the image of the \(K\)-valued period under the specified embedding, as in Part~IV. Suppose a common complex target isometry \(V\), column isometries \(U_u\), and scalars \(a_u\in\C^\times\) satisfy \(B_u^\C(\phi,F)=a_u\,\mathrm{bes}_u(U_u\phi,VF)\) for all complex vectors \(\phi,F\). Assume the frame bound
\[
 s_{\rm fr}+\mu_{\max}(RR)\le x^{C_H}b,\qquad
 s_{\rm fr}=\frac{\sum_ud_u(\mu(N'_u)-h_u)}{\sum_ud_u},
\]
and the layerwise and adjusted-height comparisons of \ref{B5.13} and \ref{B5.14}, both with constant \(c_{\rm Div}\). Then the family has the full division data \Lean{LinkDIV\_\allowbreak{}at}, with wordwise clauses
\[
 B_u(N'_u)\subseteq RR^\vee,\qquad
 B_u\ne0\Longrightarrow h(B_u)\le h_u,\qquad
 \mu_C+h_u-\Fo\le\mathcal D(T_u)+E_{\rm div}.
\]
\end{theorem}
\leanref{Formalism.\allowbreak{}MainNotions.\allowbreak{}LinkDIV\_\allowbreak{}at}{Formalism.\allowbreak{}SharpRestriction.\allowbreak{}GlobalTheoremInputs}{77--90}
\leanref{Formalism.\allowbreak{}L2.\allowbreak{}DIV.\allowbreak{}wordFacts\_\allowbreak{}of\_\allowbreak{}face}{Formalism.\allowbreak{}LocalToGlobal.\allowbreak{}Divisibility.\allowbreak{}RecBridge}{204--226}
\leanref{Formalism.\allowbreak{}L2.\allowbreak{}DIV.\allowbreak{}linkDIV\_\allowbreak{}at\_\allowbreak{}words}{Formalism.\allowbreak{}LocalToGlobal.\allowbreak{}Divisibility.\allowbreak{}RecBridge}{327--340}
\leanref{Formalism.\allowbreak{}L2.\allowbreak{}DIV.\allowbreak{}DivisionPackageFace.\allowbreak{}wordFacts}{Formalism.\allowbreak{}LocalToGlobal.\allowbreak{}Divisibility.\allowbreak{}RecBridge}{633--641}

\begin{proof}
The isometries of \ref{B4.56} transport the integrality and norm bounds of \ref{B5.09}. They also transport the target slope. Averaging the log norm bounds gives \(h_u\), and the numerical slope bound becomes \(\mu_C+h_u-\Fo\le\mathrm{rhs}\) with \(\mathrm{MAIN}(u)=\mathcal D(T_u)\). The remaining clauses are the target, Bessel, vanishing, frame and comparison conditions in the hypotheses.

For the selected family use the height bounds of \ref{B5.07a} and \ref{B5.07b} and the reference slope of \ref{B5.08x}. Apply the passage estimate of \ref{B4.55a} with the budget of \ref{B5.09a}. The frame bound is \ref{B5.10a}. Take the larger of the divisor constant and \(c_E\). Its maximum with \(20+6/\varepsilon\) gives a common \(c_{\rm Div}\).
\end{proof}

\begin{lemma}[Choosing the height exponent]\label{B5.10a}
There is \(\varepsilon_1>0\) such that for every \(0<\varepsilon\le\varepsilon_1\) in the chapter's range there are family thresholds for the row estimate. Use the selected class with the functional and row hypotheses of \ref{B8.22}, \ref{B8.23}, \ref{B8.38} and the level and stabilizer bounds. Use the estimates of \ref{B5.06}, \ref{B5.07}, \ref{B5.12}, \ref{B5.13}. The divided lattices are full of positive rank. Fix the outer constants \(c_{\rm out}\), including the chain-parameter bound \(C_t\), and let \(M=\max(C_t,5)+3\). Fix the row exponent \(F_0(\varepsilon,c_{\rm out})\) for these data. Then choose
\[
 C_H\ge1+\max\{4M+1,9M+2,F_0,9M+1\}.
\]
Under the family size conditions \Lean{Large} for the chosen thresholds and \(C_H\), there is \(x_0>F_m.x_1\) such that \(s_{\rm fr}+\mu_{\max}(RR)\le x^{C_H}b\) for every member at \(x\ge x_0\).
\end{lemma}
\leanref{Formalism.\allowbreak{}L2.\allowbreak{}DIV.\allowbreak{}FunSharpTop.\allowbreak{}closed\_\allowbreak{}funSharp\_\allowbreak{}alt}{Development.\allowbreak{}AutomorphicArithmetic.\allowbreak{}ThetaCorrespondence.\allowbreak{}ThetaPairLifts.\allowbreak{}ThirdPlaceSeries.\allowbreak{}HefunFunTop}{88--116}
\leanref{Formalism.\allowbreak{}MainNotions.\allowbreak{}DIV.\allowbreak{}Proofs.\allowbreak{}height\_\allowbreak{}bound\_\allowbreak{}of\_\allowbreak{}component\_\allowbreak{}bounds}{Formalism.\allowbreak{}SharpRestriction.\allowbreak{}Divisibility.\allowbreak{}SlopeExcess.\allowbreak{}HeightBoundOfComponentBounds}{136--151}

\begin{proof}
Put \(q_\kappa=\langle\log\kappa_v\rangle_\infty\). The functional estimate gives \(\mu_{\max}(RR)+q_\kappa\le x^{a_{\rm row}}b+14\lambda_0\) with \(a_{\rm row}=a_{\rm row}(\varepsilon,c_{\rm out})\) and \(\lambda_0\le4b\). For \(x\ge56\) the bound is at most \(x^{F_0}b\) with \(F_0=\max(a_{\rm row},1)+1\).

Put \(P=1+f+k+m+t_{\rm ch}+N_m+L\le x^M\) by the rank windows. By \ref{B5.06},
\[
 \mu(N'_u)\le\mu_{\max}(M_u)+\Ga+\Fo-\Ee_u+C mB_0 .
\]
The original slope is at most \(x^{4M+1}b\) eventually by \ref{B5.12}. The gain bound is \(\Ga\le64P^4b\), and \ref{B5.09a} gives \(\Fo\le256P^9b\). By \ref{B5.13}, \(-\Ee_u\le c_E mB_0\). Since \(mB_0\le P^4b\), these three terms sum to at most \(x^{9M+2}b\) eventually.

In the height formula of \ref{B5.07}, the size bounds give \(S_a\log T_u\le2|C_{Tu}|P^4b\) and \(n_Tb'N_m+f(m\eta+S_c)N_m+fN_m^2\le4P^3b\). With \(L_a\le4P\) this bounds \(-H_{P,u}\) by \((2|C_{Tu}|+16|C_{\rm sz}|)P^9b\). The row and hook estimates bound \(-\mathrm{kerPlace}_{uj\rho}\) and \(\log\operatorname{Hook}(a^+_{j\rho})\) by \(300P^7b\) and \(160P^7b\), respectively. Since \(f,m\le P\), these contribute at most \(300P^9b\) and \(160P^9b\) to the sums. The kernel corrections cost at most \(5|C_{\rm sz}|P^9b\), by \(fN_m^2L_a\le4P^4\) and \(fm^3N_m^2L\le P^7\) per layer. The hook correction costs at most \(6|C_{\rm sz}|P^9b\), by \(fn_V^2L'\le6P^5\) per layer. Consequently
\[
 -h_u-q_\kappa\le(500+2|C_{Tu}|+30|C_{\rm sz}|)P^9b
 \le x^{9M+1}b
\]
for sufficiently large \(x\). Adding the row bound cancels \(q_\kappa\). Each of the five bounds, including that for \(C mB_0\), is at most \(x^{C_H-1}b\). Their sum is at most \(x^{C_H}b\) for \(x\ge5\). Average with weights \(d_u/\sum_{u'}d_{u'}\).
\end{proof}

\section{The auxiliary prime and the original lattice}

For each word put
\[
\begin{aligned}
 d_u^\star&=\operatorname{mass}_k\,
 \operatorname{idx}_{\Pi}(g_0,g_1)\prod_\rho\dim V(u,\rho),\\
 D^\star&=\operatorname{vol}(Y_0)\,
 \operatorname{idx}_{\Sigma}(g_0,g_1)\prod_\rho d_{\Sigma,\rho}.
\end{aligned}
\]
The two indices are finite-level subgroup indices and \(d_{\Sigma,\rho}\) are spectral degrees. The adjacent-rank mass quotient for the standard unimodular lattices is given by \ref{B8.48}. The volume induction \ref{B8.51} uses \ref{B8.49} and the rank-two value \ref{B8.50}. With \(\Gamma_u=\Lean{GamP}\), put \(X_u=(D^\star/d_u^\star)^{1/(fk)}\) and \(T_u^\star=X_u\Gamma_u^{-2/(fk)}\).

\begin{theorem}[Three counting bounds for the auxiliary primes]\label{B5.11}
Fix \(C_0,C_Y\). Above the common threshold for the row bounds and spectral gap, take a member satisfying the separated-weight and row-size conditions \Lean{DataA} of \ref{B8.51} and the mass and volume identities there. Assume the spectral estimate of \ref{B8.50a}. Identify its recorded \(d_u^\star,D^\star\) with the two products above. Assume \(L\le L_a,\ L_a\ge1,\ k\ge3,\ f\ge1\), and \(\Gamma_u=\Gamma_{u'}>1\) with \(\log\Gamma_u/(fk)\le1\) for all words. Let \(T_u^\star\le T_u\le(1+1/(fk))T_u^\star\). Put
\[
 \mathcal W=\sum_{j=1}^{m}\frac1f\sum_\rho\log(w_{j\rho}v_E/\pi),\qquad
 \ell_g=\log(g_0g_1).
\]
Let \(C\) be a common constant for the torus bounds and the variation bound \(|\log X_u-\log X_{u'}|\le C\log k\) below. Set
\[
 c_{\rm Count}=\max(C,0)+2,\qquad
 \lambda_5=\mathcal W+m\ell_g-c_{\rm Count}(L_a+m^2L_a/k).
\]
Then
\[
 \lambda_5\le\log T_u\le c_{\rm Count}mL,\qquad
 |\log T_u-\log T_{u'}|\le c_{\rm Count}\log k .
\]
For an admissible family with the volume comparison \ref{B8.49} and rank-two identification \ref{B8.50}, the mass and spectral estimates give uniform \(c_{\rm Count}\) and \(x_0\) for these conclusions at every member with \(x\ge x_0\).
\end{theorem}
\leanref{Formalism.\allowbreak{}MainNotions.\allowbreak{}GlobalTheoryTheory.\allowbreak{}count\_\allowbreak{}at\_\allowbreak{}of}{Formalism.\allowbreak{}SharpRestriction.\allowbreak{}GlobalSchurBounds.\allowbreak{}JointsL}{20--143}
\leanref{Formalism.\allowbreak{}MainNotions.\allowbreak{}StandardLattice.\allowbreak{}division\_\allowbreak{}level\_\allowbreak{}prime\_\allowbreak{}family}{Formalism.\allowbreak{}SharpRestriction.\allowbreak{}GlobalSchurBounds.\allowbreak{}LinkMin}{425--454}
\leanref{Formalism.\allowbreak{}MainNotions.\allowbreak{}GlobalTheoryTheory.\allowbreak{}countF\_\allowbreak{}fam}{Formalism.\allowbreak{}SharpRestriction.\allowbreak{}GlobalSchurBounds.\allowbreak{}JointsL}{1152--1155}

\begin{proof}
At \(\rho\) let \(z_a=z_{\rho,a}(u)\), \(1\le a\le k\), be the shifted column entries and let \(\Delta_\rho(u)=\prod_{a<b}(z_a-z_b)\). With spacing parameter \(s>0\), gaps with \(a=1\) change by at most \(s\) between words and are at least \(4ks\). The other gaps change by at most \(2s\) and are at least \(3(b-a)s\). Apply \(|\log y-\log y'|\le|y-y'|/\min(y,y')\) and sum by \(b-a\). This gives \(|\log\Delta_\rho(u)-\log\Delta_\rho(u')|\le\tfrac14+\tfrac23k(1+\log k)\). Since \(d_u^\star\) is proportional to \(\prod_\rho\Delta_\rho(u)\) with a common factor, and \(D^\star\) is common,
\[
 |\log X_u-\log X_{u'}|\le\frac1{4k}+\frac23(1+\log k)\le2\log k .
\]

In the selected family one can take \(C_0=C_Y=\max(5,C_{YL}+1)\), where \(C_{YL}\) bounds the type-prime sizes. The separated weights give the opposite-sign gap of \ref{B8.50a} by \Lean{gapB\_\allowbreak{}ev}. Put \(r=k+m-1\), \(F_n=\prod_{i=0}^{n-1}i!\), and let \(\lambda_\rho\) denote the row parameter. Since \(e^{-r}\le1/3\), that estimate gives \(d_{\Sigma,\rho}\ge(2/3)\dim V(\lambda_\rho)\).

Let \(v_1,\ldots,v_{k-1}\) be the nonnegative row entries interlacing the column, and \(-b_{j\rho}\) the negative auxiliary entries. Write \(\Delta(v)\) for the row entries' Vandermonde. The \(2k-3\) column differences involving \(z_1\) or \(z_k\) are at most \(8a\). For \(2\le i<j\le k-1\), interlacing gives \(z_i-z_j\le v_{i-1}-v_j\), and the unused adjacent row gaps are at least one. In the row Vandermonde, each cross difference is \(v_i+b_{j\rho}\ge b_{j\rho}\), and the negative entries have gaps at least one. Thus
\[
 \Delta_\rho(u)\le(8a)^{2k-3}\Delta(v),\qquad
 \Delta(\lambda_\rho)\ge\Delta(v)\prod_{j=1}^m b_{j\rho}^{\,k-1}.
\]
By Weyl's dimension formula and \(2(\log8+1+\log x)\le L\) at the common threshold, the log loss from \((8a)^{2k-3}\) is at most \(5kL\). Therefore
\[
 \log\frac{d_{\Sigma,\rho}F_r}{\Delta_\rho(u)}
 \ge \log(2/3)+(k-1)\sum_j\log b_{j\rho}-5kL.
\]
The weight comparison gives \(k\log w_{j\rho}-(k-1)\log b_{j\rho}\le k/x+2L\). Since \(m\le x\) and \(m\le k\), we obtain
\[
 S_{\rm dim}-k\sum_{\rho,j}\log w_{j\rho}\ge-f(1+k+7kL),
 \qquad S_{\rm dim}=\sum_\rho\log\frac{d_{\Sigma,\rho}F_r}{\Delta_\rho(u)} .
\]

Put \(\ell_E=\log(v_E/\pi)\). The factorials in the product of adjacent-rank mass quotients cancel the quotient of factorial products in the dimension formula. The volume and local indices give
\[
\begin{aligned}
 fk\log X_u={}&
 f\left((m-1)k+\frac{m(m-1)}2\right)\ell_E
 -\frac{f(m-1)}2\log(2v_E)+S_L\\
 &+f\left(km+\frac{m(m-1)}2\right)\ell_g+fS_g+S_{\rm dim},
\end{aligned}
\]
where \(S_L=\sum_{n=k+1}^{r}\log L_\eta(E,n)\) and \(S_g=\sum_{i=0}^1\sum_{j=k+1}^{r}\log(1-g_i^{-j})\). The Euler factors in \ref{B8.48} give \(|\log L_\eta(E,n)|\le f\log\zeta(n)\le f\), so \(-f(m-1)\le S_L\le fm\). For each \(i\) the local product is between \(1-\sum_{j\ge k+1}g_i^{-j}\ge1/2\) and \(1\), hence \(-2\le S_g\le0\). Subtract \(fk(\mathcal W+m\ell_g)\) in the product identity. Since \(0\le\ell_E\le L,\ \ell_g\ge0,\ \log(2v_E)\le L\), the difference is at least
\[
 -f\left(8kL+\frac{m-1}{2}L+m+2+k\right)\ge-11fkL
 \qquad(44\le m\le k,\ L\ge1).
\]
Thus \(\log X_u\ge\mathcal W+m\ell_g-11L\).

For the upper bound, interlacing gives \(v_i-v_j\le z_i-z_{j+1}\) in one-based indexing. The unused adjacent column gaps are at least one, so \(\Delta(v)\le\Delta_\rho(u)\). In the compact Weyl numerator of \ref{B8.50a}, at most \(2km\) factors remain after those of \(\Delta(v)\), each at most \(H+2k\) for the maximum row difference \(H\). The compact denominator \(F_pF_q\) is at least \(F_k\), where \(p,q\) is the row signature. With \(H\le4e^L,\ pq\le km,\ r\le2k,\ \log(2k)\le L\), the spectral upper bound gives
\[
 \log\frac{d_{\Sigma,\rho}}{\dim V(u,\rho)}
 \le\log4+2km\log(H+2k)+km\log(3pqrH)\le km(12+10L).
\]
Also \(\log(F_r/F_k)\le(r-k)r\log r\le2mkL\), so \(S_{\rm dim}\le fkm(12+12L)\). Use \(S_L\le fm,\ S_g\le0,\ \ell_g\le2C_YL\) in the product identity. After division by \(fk\), the \(\ell_E\) term contributes at most \(\tfrac32mL\) and the \(\ell_g\) term at most \(3C_YmL\). Since \(m/k+m(12+12L)+\tfrac32mL\le15mL\) for \(L\ge9,\ k\ge5\), this gives \(\log X_u\le(15+3C_Y)mL\).
\leanref{Formalism.\allowbreak{}MainNotions.\allowbreak{}StandardLattice.\allowbreak{}part\_\allowbreak{}five\_\allowbreak{}family\_\allowbreak{}of\_\allowbreak{}mass\_\allowbreak{}volume}{Formalism.\allowbreak{}SharpRestriction.\allowbreak{}GlobalSchurBounds.\allowbreak{}LinkMin}{318--423}

Since \(g_i\ge(kb)^8\), \(\mathcal W\ge0\) eventually and \(\log\Gamma_u/(fk)\le1\), the lower bound gives \(\log T_u^\star\ge mL\) for sufficiently large \(x\). Apply \ref{B8.51a} with \(V=fk\) and modulus \(Q_Sqd(g_0g_1)^2\), where \(Q_Sqd\) is the host modulus factor. This congruence ensures splitting by \ref{B8.54}.

The interval for \(T_u\) gives
\[
 \log T_u=\log X_u-\frac{2\log\Gamma_u}{fk}+\epsilon_u,\qquad
 0\le\epsilon_u\le\log(1+1/(fk))\le1.
\]
Take \(C\ge\max\{11,16+3C_Y,2\}\). The upper bound is \(CmL\), and the lower error below \(\mathcal W+m\ell_g\) is at most \(C(L+m^2L/k)+2\log\Gamma_u/(fk)\le c_{\rm Count}(L_a+m^2L_a/k)\). The common \(\Gamma_u\) cancels in the difference between words. Thus \(|\log T_u-\log T_{u'}|\le C\log k+1\le c_{\rm Count}\log k\). Increase the common threshold to exceed the lower endpoint \(x_1\).
\end{proof}

\begin{lemma}[Slopes before summing the column weights]\label{B5.11a}
For a word with a column model, write the integral highest weight at \(\rho\) as \(\lambda_i=\nu_i+t_\rho\), \(1\le i\le k\), with \(\nu_1\ge\cdots\ge\nu_k=0\). Put \(d_\rho=\sum_i\nu_i,\ s_\rho=\nu_1,\ l=\log T_u\). Then
\[
\begin{aligned}
 \mu(M_u)&\ge-\sum_\rho\bigl[(d_\rho/k+t_\rho)l+d_\rho\log k+k^2\log(s_\rho+k)\bigr],\\
 \mu_{\max}(M_u)&\le\sum_\rho\bigl[4d_\rho\log(kv_E)+k^2\log(s_\rho+k)+|t_\rho|l\bigr].
\end{aligned}
\]
\end{lemma}
\leanref{Formalism.\allowbreak{}L2.\allowbreak{}GlobalTheoryTheory.\allowbreak{}TraceIntegrality.\allowbreak{}bc\_\allowbreak{}mu}{Formalism.\allowbreak{}LocalToGlobal.\allowbreak{}GlobalTheory.\allowbreak{}GL4.\allowbreak{}ColumnLatticeExactSlope}{98--177}
\leanref{Formalism.\allowbreak{}L2.\allowbreak{}GlobalTheoryTheory.\allowbreak{}TraceIntegrality.\allowbreak{}eLat\_\allowbreak{}mu}{Formalism.\allowbreak{}LocalToGlobal.\allowbreak{}GlobalTheory.\allowbreak{}GL4.\allowbreak{}MemDualCrdLat}{71--128}
\leanref{Formalism.\allowbreak{}L2.\allowbreak{}GlobalTheoryTheory.\allowbreak{}TraceIntegrality.\allowbreak{}lower\_\allowbreak{}slope\_\allowbreak{}inequality\_\allowbreak{}family}{Formalism.\allowbreak{}LocalToGlobal.\allowbreak{}GlobalTheory.\allowbreak{}GL4.\allowbreak{}LowerSlopeInequalityFamily}{14--15}
\leanref{Formalism.\allowbreak{}L2.\allowbreak{}GlobalTheoryTheory.\allowbreak{}TraceIntegrality.\allowbreak{}slope\_\allowbreak{}inequality\_\allowbreak{}constant\_\allowbreak{}four}{Formalism.\allowbreak{}LocalToGlobal.\allowbreak{}GlobalTheory.\allowbreak{}GL4.\allowbreak{}PStmtFree}{14--15}

\begin{proof}
The Hermitian covolume formula gives the degree \(-\log[A^\vee:A]/[E:\Q]\) for a class lattice \(A\subseteq A^\vee\) over \(E\). The local composition factors give the index. Under isometric scalar extension the finite log norm of the determinant ideal and archimedean log determinant sum both multiply by the extension degree. Thus the normalized degree is preserved. Here \([A^\vee:A]=T_u^{2f}\) and \([E:\Q]=2f\). The column rank is \(k\), so its slope is \(-l/k\).

Put \(h_\nu=\|\mathrm{hw}_\nu\|^2\), the Fock norm squared of the highest-weight vector, and \(c_\nu=h_\nu/\prod_i\nu_i!\). Write \(\mu_\nu\) for minus the log Gram determinant per rank of the standard Weyl lattice. Normalized log determinants add under tensor product. Symmetry of weights gives determinant character \(\det^{r'\sum_i\lambda_i/k}\) in dimension \(r'\). Applying this at the column frames and using the equal-rank orthogonal classes gives
\[
 \mu(M_u)=\sum_\rho\bigl[\mu_\nu-(d_\rho/k+t_\rho)l\bigr].
\]
In integral monomial coordinates \(V\), the dual polynomial lattice has Gram matrix \(V^{\mathsf T}\operatorname{diag}(m!/h_\nu)V\), with \(m!=\prod_{a,i}m_{ai}!\), \(\sum_a m_{ai}=\nu_i\). The multinomial theorem gives \(m!/h_\nu\ge k^{-d_\rho}/c_\nu\). A maximal minor of \(V\) is a nonzero integer, so Cauchy--Binet bounds the dual determinant below by \((k^{-d_\rho}/c_\nu)^{\operatorname{rk}V}\). Inverting gives \(\mu_\nu\ge-d_\rho\log k-\log c_\nu\). In the hook norm
\[
 h_\nu=\frac{\prod_i(\nu_i+k-i)!}{\prod_{i<j}(\nu_i-\nu_j+j-i)}
\]
the denominator is a positive integer and each extra factorial factor is at most \(s_\rho+k\). Thus \(c_\nu\le(s_\rho+k)^{k(k-1)/2}\), which proves the mean bound.

For the maximum slope, the trace lattice on \(A^\vee\) has Euclidean dimension \(ek\), \(e=[E:\Q]\). The trace--norm inequality for the inverse different gives squared first minimum at least \(e/\operatorname{rd}(E)^2\) for its Euclidean dual. Banaszczyk's transference inequality is \(\lambda_n(\Lambda)\lambda_1(\Lambda^*)\le n\) for a Euclidean lattice of rank \(n\) \cite{banaszczyk-1993}. It gives \(ek\) rationally independent vectors with squared trace length at most \(ek^2(2v_E)^2\). Choose \(k\) independent \(E\)-vectors \(y_a\) from this spanning list. Pairing with them embeds each class in monomial coordinates with its determinant twist. The tensor norm bound gives squared multiplier \((\sum_a h_\sigma(y_a,y_a))^{d_\rho}\) on the polynomial part. Concavity of \(\log\) and the trace bounds give mean log multiplier at most \(d_\rho\log(k^3(2v_E)^2)\).

The determinant twist contributes \(kt_\rho\mu(Q)=-t_\rho l\), where \(Q\) is the column lattice. A nonzero minor on the determinant line of any sublattice gives the monomial allowance \(d_\rho\log k+\log c_\nu\) per rank by the Gram bound above. Integral orthogonal class projections give the same upper bound on \(M_u\). Use \(4k^4v_E^2\le(kv_E)^4\) for \(v_E\ge2\) and \(-t_\rho l\le|t_\rho|l\). Sum over \(\rho\).
\end{proof}

\begin{theorem}[Mean and maximum slopes of the original lattice]\label{B5.12}
For every word \(u\) and column model, the zeroth-layer lattice \(Q_{u,c,\rho}\) of class \(c\) and column embedding \(\rho\) has slope \(-\log T_u/k\). Let an admissible family satisfy \ref{B5.11} and have column models for every word beyond one family threshold. Assume \(|D_1|\le x^\varepsilon\) eventually. Then there are \(c_{\rm Slope}\ge0\) and \(x_0\) such that, for every member at \(x\ge x_0\) and every word,
\[
 -\mu(M_u)\le c_{\rm Slope}n_TaL_a(1+m/k),\qquad
 \mu_{\max}(M_u)\le c_{\rm Slope}n_TaL_a.
\]
\end{theorem}
\leanref{Formalism.\allowbreak{}L2.\allowbreak{}GlobalTheoryTheory.\allowbreak{}TraceIntegrality.\allowbreak{}q\_\allowbreak{}lattice\_\allowbreak{}slope}{Formalism.\allowbreak{}LocalToGlobal.\allowbreak{}GlobalTheory.\allowbreak{}GL4.\allowbreak{}ColumnLatticeExactSlope}{183--203}
\leanref{Formalism.\allowbreak{}MainNotions.\allowbreak{}flipTwoThree.\allowbreak{}slopesMF\_\allowbreak{}of\_\allowbreak{}count\_\allowbreak{}fam}{Formalism.\allowbreak{}SharpRestriction.\allowbreak{}GlobalSchurBounds.\allowbreak{}GlobalRankLatticeProofs.\allowbreak{}SlopesMFOfCountFam}{25--33}
\leanref{Formalism.\allowbreak{}MainNotions.\allowbreak{}GlobalTheoryTheory.\allowbreak{}JSL.\allowbreak{}slopesM\_\allowbreak{}at\_\allowbreak{}of}{Formalism.\allowbreak{}SharpRestriction.\allowbreak{}GlobalSchurBounds.\allowbreak{}JointsL}{1518--1589}

\begin{proof}
Use the column calculation in \ref{B5.11a}. Put \(C_n=\max(3,c_{\rm Count},C_{YL}+1)\) and \(c_L=3\). Increase the threshold so \(x^{14}\le X'=\exp(x^{\varepsilon/3})\), \(7\log x\le L\), and \(m\ge c_L+c_{\rm Count}+6C_n+9\). The parameter bounds give \(\log(xbv_E)\le3L,\ N_mm^4\le b'\), and \ref{B5.11} gives \(l_T=\log T_u\le c_{\rm Count}mL\). With \(|D_1|\le L\) these are the numerical bounds with log argument \(c_LL\).

At a top column the first weight is \(\alpha_{u\rho}\in[a,2a]\) and the negative of the last weight lies in \([b',2b')\). The interior entries lie in \([0,S_c]\). Other columns have entries in \([0,C_nS_c]\). With \(B=b'+S_c\), the decomposition of \ref{B5.11a} gives
\[
 \sum_\rho d_\rho\le2n_Ta+O_{C_n}(fkB),\quad
 \sum_\rho|t_\rho|\ll_{C_n} fB,\quad
 \sum_{\rho,i}\lambda_i\le\sum_{\rho\ {\rm top}}\alpha_{u\rho}+O_{C_n}(fkS_c).
\]
Here \(\log(s_\rho+k)\ll_{C_n}c_LL\) by \(S_c\le b',\ b'\le b,\ k\le x^5\). At a top position retain \(d_\rho\le\alpha_{u\rho}+O_{C_n}(kB)\) in the \(\log k\) term. Substituting, with \(S_1=\sum_{\rho\ {\rm top}}\alpha_{u\rho}l_T/k\), gives constants \(C_5,C_6\) for fixed \(C_n\) such that
\[
\begin{aligned}
 \mu(M_u)&\ge
 -S_1-C_5\bigl(n_Ta\log k+fkB\log k+fS_cl_T+fk^2c_LL\bigr),\\
 \mu_{\max}(M_u)&\le
 C_6\bigl(n_Ta\log(kv_E)+fB(k\log(kv_E)+l_T)+fk^2c_LL\bigr).
\end{aligned}
\]
\leanref{Formalism.\allowbreak{}MainNotions.\allowbreak{}flipTwoThree.\allowbreak{}nums\_\allowbreak{}wts\_\allowbreak{}of\_\allowbreak{}fam}{Formalism.\allowbreak{}SharpRestriction.\allowbreak{}GlobalSchurBounds.\allowbreak{}GlobalRankLatticeProofs.\allowbreak{}SevenLogLeLOfPowerBound}{698--747}
\leanref{Formalism.\allowbreak{}MainNotions.\allowbreak{}flipTwoThree.\allowbreak{}numsL\_\allowbreak{}of\_\allowbreak{}fam}{Formalism.\allowbreak{}SharpRestriction.\allowbreak{}GlobalSchurBounds.\allowbreak{}GlobalRankLatticeProofs.\allowbreak{}SlopesMFOfCountFam}{17--23}

Beyond a common threshold, \(m\le k\), \(L\le L_a\), \(\log k\le5L_a\), \(\log(kv_E)\le6L_a\), \(fkB\le n_Ta\) and \(fk^2\le n_Ta\). In units of \(n_TaL_a\), \(S_1\) is at most \(2c_{\rm Count}m/k\). Since \(S_c\le B\), the four errors multiplying \(C_5\) are bounded by \(5,5,c_{\rm Count},c_L\). The errors multiplying \(C_6\) are bounded by \(6,6+c_{\rm Count},c_L\). Thus \(c_{\rm Slope}=2c_{\rm Count}+\max(C_5,C_6,0)(12+c_{\rm Count}+c_L)\) suffices.
\end{proof}

\section{Comparing words}

\begin{lemma}[The layerwise \(E\)-comparison]\label{B5.13}
Use the local \(E\)-data of \ref{B4.03}, including the norm decomposition and degree count at \(g_0,g_1\). Assume \(T_u\ne g_i\), \(i=0,1\), and take \(c_g\ge0\) for the type-prime size bounds. Let \(C_T,C_\Gamma,C_d,C_I\) be the level-prime, stabilizer, discriminant and inert-prime loss constants, and put \(c_E=2c_g+C_T+C_\Gamma+C_d+C_I\). For every word and \(1\le j\le m\),
\[
 \left|\sum_{w\in S}c_wE_{ujw}-S_a\ell_g\right|\le c_EB_0,\qquad
 \sum_{\substack{w\in S\\p_w\notin\{g_0,g_1\}}}|c_wE_{ujw}|\le c_EB_0 .
\]
Consequently \(|\Ee_u-mS_a\ell_g|\le mc_EB_0\).
\end{lemma}
\leanref{Formalism.\allowbreak{}MainNotions.\allowbreak{}LocalQuadraticParts.\allowbreak{}adm11}{Formalism.\allowbreak{}SharpRestriction.\allowbreak{}Divisibility.\allowbreak{}LocalQuadraticParts}{316--322}
\leanref{Formalism.\allowbreak{}MainNotions.\allowbreak{}WeightDegreeSizes.\allowbreak{}secondComplement\_\allowbreak{}of\_\allowbreak{}inputs}{Formalism.\allowbreak{}SharpRestriction.\allowbreak{}Divisibility.\allowbreak{}WeightDegreeSizes}{677--681}
\leanref{Formalism.\allowbreak{}MainNotions.\allowbreak{}Lnk.\allowbreak{}Mod.\allowbreak{}data\_\allowbreak{}EComp5}{Formalism.\allowbreak{}SharpRestriction.\allowbreak{}GlobalSchurBounds.\allowbreak{}QuadraticLogDerivCongr}{594--622}

\begin{proof}
Put \(v_{i,u}=v_{g_i}(D_{\Gamma,u})\), \(r_g=c_g'f(N_m^2+m^2\eta)\) for the local coefficient \(c_g'\), and \(\nu^{\rm abs}_{uj\rho}=\sum_{h=0}^{N_{j-1}-1}\nu_{uj\rho h}\). The type-prime formula and \(\sum_{w\mid g_i}e_w\ell_w=\log g_i\) from \ref{B4.52} give
\[
 \sum_{\substack{w\in S\\p_w=g_i}}c_wE_{ujw}
 =\left(2v_{i,u}-\mathrm{fockDeg}_j+2r_g+\sum_\rho\nu^{\rm abs}_{uj\rho}\right)\log g_i ,
\]
using \(2F_{{\rm type},j}=-\mathrm{fockDeg}_j\). Subtract \(S_a\log g_i\) and apply the degree count and type-prime size bounds to get absolute error at most \(c_gB_0\).

The fiber of \(T_u\) costs at most \(C_TB_0\). Grouping the other fibers by rational prime uses \(\sum_{w\in S,\ p_w=p}e_w\ell_w\le\log p\). The stabilizer terms then cost at most \(2\log D_{\Gamma,u}\le C_\Gamma B_0\) and the discriminant term costs \(2\log d\,|\mathrm{EatD}_j|\le C_dB_0\). For a finite set \(P\) of primes \(p\le X''\) in the \(\Lean{IsResidueClassZero}\) or \(\Lean{IsSg}\) cases, the remaining budget is \(\sum_{p\in P}2\log p\,|\mathrm{EatCI}_j(p)|\le C_I B_0\). Here \(\mathrm{EatCI}_j=\mathrm{EatC}_j\) in the inert case with \(h_P=0\) and \(\mathrm{EatI}_j\) otherwise. Apply this to the eligible primes outside the level and type fibers. The off-type sum is bounded by \((C_T+C_\Gamma+C_d+C_I)B_0\). Add the two type-prime errors and sum over the layers.
\end{proof}

\begin{lemma}[The difference of adjusted heights]\label{B5.14}
For the heights of \ref{B5.07} put \(H_u=h_u+S_a\log T_u\). For every pair of words,
\[
 H_u-H_{u'}=\sum_{j=1}^{m}\sum_\rho
 \bigl(\mathrm{kerPlace}_{uj\rho}-\mathrm{kerPlace}_{u'j\rho}\bigr).
\]
Under the parameter and size hypotheses \Lean{ParameterChoice'} and \Lean{SizesAt} of the actual construction,
\[
 |H_u-H_{u'}|\le\left(20+\frac6\varepsilon\right)(n_TaL_a+mB_0).
\]
\end{lemma}
\leanref{Formalism.\allowbreak{}MainNotions.\allowbreak{}DIV.\allowbreak{}HookProducts.\allowbreak{}RBCompl3.\allowbreak{}compl3\_\allowbreak{}diff}{Formalism.\allowbreak{}SharpRestriction.\allowbreak{}Divisibility.\allowbreak{}HookProducts.\allowbreak{}Compl3Record}{15--22}
\leanref{Formalism.\allowbreak{}MainNotions.\allowbreak{}DIV.\allowbreak{}HookProducts.\allowbreak{}RBCompl3.\allowbreak{}compl3\_\allowbreak{}record}{Formalism.\allowbreak{}SharpRestriction.\allowbreak{}Divisibility.\allowbreak{}HookProducts.\allowbreak{}Compl3Record}{217--281}
\leanref{Formalism.\allowbreak{}MainNotions.\allowbreak{}Lnk.\allowbreak{}Mod.\allowbreak{}data\_\allowbreak{}HComp}{Formalism.\allowbreak{}SharpRestriction.\allowbreak{}GlobalSchurBounds.\allowbreak{}QuadraticLogDerivCongr}{624--639}

\begin{proof}
The adjustment cancels the terminal term \(-S_a\log T_u\) in \ref{B5.07}. All remaining terms outside the row sums cancel in \(H_u-H_{u'}\).

The first and last source coordinates and the common tail are fixed. At a middle index \(1\le r,\ r+1<k\), the source coordinates belong to \([0,S_c]\) and the shift is common. Thus \(|\nu_{uj\rho r}-\nu_{u'j\rho r}|\le S_c\). Both entries are positive integers bounded by \(S_c+w_{j\rho}\), since the layer term \((N_{j-1}+N_j-1-2r)/2\) is nonnegative and subtracted from the shifted source coordinate. If \(M\ge1\) bounds \(S_c+w_{j\rho}\), integration of \(\varphi'(t)=\log t\) on \([1,M]\) gives
\[
 \left|(\varphi(y)-y\log\Omega)-(\varphi(y')-y'\log\Omega)\right|
 \le |y-y'|(\log M+|\log\Omega|).
\]
There are at most \(k\) middle indices per position and layer. The sum is at most \(mfkS_c(\log M+|\log\Omega|)\). The size bounds give \(M=6x^6b\) as an admissible bound and \(kS_c\le b\). The log derivative and discriminant bounds give \(|\log\Omega|\le6+2L\). By \(\log x\le L/\varepsilon\) and \(L\ge1\), the sum of the two log bounds is at most \((14+6/\varepsilon)L\). Hence \(|H_u-H_{u'}|\le(14+6/\varepsilon)B_0\), which implies the stated estimate. The translation preserves \(h_u,S_a,T_u\).
\end{proof}


% ---- p5c2 ----
\chapter{Prime sums and net gain}\label{ch:prime-sums-net-gain}

At a window point \(x\) and translation \(u_0\) of Alp\"oge's construction
\cite{alpoge-2026}, let \(\ASK,\Ga\) be the sums in \ref{B5.01} and
\ref{B5.02}. Let \(n_T=|TT|\) be the number of top positions and \(m\) the number of layers. Put \(Q=n_Tmb\) with \(b=e^{x^\varepsilon}\), \(0<\varepsilon<1/52\).
The weights in \ref{B5.22} give \(\ASK-\Ga,\mathrm{IO}\le c_{\rm Loss}Q\)
uniformly for large \(x\).

\section{The deficit at a rational prime}

Let \(J\) be a set of \(m\) layers and \(P_s\) a finite set of
rational primes. In the model, arithmetic layer \(j\in\mathrm{Fin}(m)\)
represents lift \(j+1\). For \(j\in J\), \(t\in TT\), \(p\in P_s\),
write \(A_{jtp}\) for the ordinary contribution before multiplication
by \(2\log p\) and \(D_{jt}\) for the contribution at the ramified
prime \(d\) with \(\log d\) included. Set
\[
 \ASK=\sum_{j\in J}\sum_{t\in TT}
 \left(2\sum_{p\in P_s}\log p\,A_{jtp}+D_{jt}\right),
 \qquad \ASK_d=\sum_{j,t}D_{jt}.
\]

\begin{lemma}[Decomposition over rational primes]\label{B5.15}
Give each \(w\) in a finite set \(W\) a real weight \(\lambda_w\).
Write \(p(w)\) for the rational prime associated to \(w\). Suppose
\[
 \sum_{\substack{w\in W\\p(w)=p}}\lambda_w=\log p
 \quad(p\in P_s).
\]
Suppose the local gains and deductions are constant above each
\(p\in P_s\), with values \(g_{jp}\) and \(\tau_p\), and vanish above
primes outside \(P_s\). Put
\[
 \Ga=2\sum_{w\in W}\lambda_w
 \left(\sum_{j\in J}g_{j,p(w)}-\tau_{p(w)}\right),
 \qquad
 \Delta_p=\sum_{j\in J}\sum_{t\in TT}A_{jtp}
 -\left(\sum_{j\in J}g_{jp}-\tau_p\right).
\]
Then
\[
 \ASK-\Ga
 =2\sum_{p\in P_s}\log p\,\Delta_p+\ASK_d.
 \tag{5.15.1}
\]
For a model in an admissible \(F_m\) at \(x\ge\max(F_m.x_1+1,3)\),
let \(P_s\) be the signal and inert-class primes at most
\[
 B=\max\!\left\{
 \left\lfloor a+\frac{a}{kx^{13}}\right\rfloor,\lfloor X''\rfloor
 \right\},
\]
with scales as in (5.20.1) and starting rank \(k\).
Then \(\Ga\) is the geometric gain of \ref{B5.02} and \(\ASK\)
is the model sum. In a consistent presentation as in \ref{B5.21},
this sum equals that of \ref{B5.01} at cutoff \(N\in\N\) if
\(a_t/w_j(\rho)\le N\) for all \(j\in\mathrm{Fin}(m)\),
\(t\in TT\), and \(\rho\in\Phi\).
\end{lemma}
\leanref{Formalism.\allowbreak{}LFunctionChainParameters.\allowbreak{}DeficitSplitLemmas.\allowbreak{}splitEq\_\allowbreak{}valid}{Formalism.\allowbreak{}SizeWindowBrackets.\allowbreak{}ChainSlopeWindows.\allowbreak{}AdmissibleSchurLemma}{63--68}
\leanref{Formalism.\allowbreak{}L2.\allowbreak{}GlobalTheoryTheory.\allowbreak{}Loss.\allowbreak{}Ga\_\allowbreak{}at}{Formalism.\allowbreak{}LocalToGlobal.\allowbreak{}GlobalTheory.\allowbreak{}PowFourLtCeilThresholdPrime}{210--245}
\leanref{Formalism.\allowbreak{}L2.\allowbreak{}GlobalTheoryTheory.\allowbreak{}Loss.\allowbreak{}ASK\_\allowbreak{}at}{Formalism.\allowbreak{}LocalToGlobal.\allowbreak{}GlobalTheory.\allowbreak{}PowFourLtCeilThresholdPrime}{386--395}

\begin{proof}
Grouping by rational prime gives (5.15.1). For the finite-sum
model take \(W=P_s\), \(p(w)=w\), and \(\lambda_w=\log w\).
The geometric support contains the full fibers over \(p\in P_s\)
with \(p\le X''\). For the coefficient field \(K\),
\(N\mathfrak p_w=p^{f_w}\) and \(\sum_{w\mid p}e_wf_w=[K:\Q]\)
\cite[Chap.~I, Prop.~8.2]{neukirch-1999}. Thus
\[
 \sum_{w\mid p}\frac{e_w\log N\mathfrak p_w}{[K:\Q]}=\log p.
\]
Here \(e_w,f_w\) are the ramification index and residue degree.
The adjusted gains and deductions are constant on each full fiber
by \ref{B3.52}. All other geometric summands vanish. For \(p\in P_s\) above \(X''\), \(g_{jp}=\tau_p=0\). For the ordinary sum,
\(a_t\le B\) by the top window and integrality. Lemma~\ref{B5.17a}
and \(w''\ge\sqrt b\ge1\) give zero contributions for \(p>a_t\).
The cutoff and model sums agree by \ref{B5.21}.
\leanref{Formalism.\allowbreak{}MainNotions.\allowbreak{}Lnk.\allowbreak{}askIsF}{Formalism.\allowbreak{}SharpRestriction.\allowbreak{}GlobalSchurBounds.\allowbreak{}FrameCodeOfAdditive}{546--556}
\end{proof}

Partition \(P_s\) into served primes \(S\) and discarded classes \(P_1,P_2,P_h\).

\begin{lemma}[Covered positions and discarded primes]\label{B5.16}
For each \(p\in S\), let \(TT_p\subseteq TT\), and let
\(\ell_p,r_p\in\R\). Suppose
\[
 \sum_{j\in J}\sum_{t\in TT_p}(A_{jtp}-\ell_p)-mr_p
 \le \sum_{j\in J}g_{jp}-\tau_p. \tag{5.16.1}
\]
For \(p\in P_s\setminus S\), suppose \(g_{jp}=0\) for every \(j\)
and \(\tau_p=0\). Then
\[
 \Delta_p\le
 \sum_{j\in J}\sum_{t\in TT\setminus TT_p}A_{jtp}
 +m|TT_p|\ell_p+mr_p \quad(p\in S),
 \tag{5.16.2}
\]
and \(\Delta_p=\sum_{j,t}A_{jtp}\) for \(p\in P_s\setminus S\).
Under the fiber-weight and gain hypotheses of \ref{B5.15},
\[
\begin{split}
 \ASK-\Ga\le{}&
 2\sum_{p\in S}\log p
 \left(\sum_{j\in J}\sum_{t\in TT\setminus TT_p}A_{jtp}
 +m|TT_p|\ell_p+mr_p\right)\\
 &+2\sum_{p\in P_1\cup P_2\cup P_h}\log p\sum_{j,t}A_{jtp}
 +\ASK_d.
\end{split}\tag{5.16.3}
\]
\end{lemma}
\leanref{Formalism.\allowbreak{}LFunctionChainParameters.\allowbreak{}DeficitSplitLemmas.\allowbreak{}split\_\allowbreak{}classes\_\allowbreak{}count}{Formalism.\allowbreak{}SizeWindowBrackets.\allowbreak{}ChainSlopeWindows.\allowbreak{}AdmissibleSchurLemma}{110--139}

\begin{proof}
Subtract (5.16.1) from the full sum at \(p\in S\). At a discarded
prime the gain and deduction vanish. Multiply by \(2\log p\) and
sum using \ref{B5.15}.
\end{proof}

In the model, discarded primes above \(X''\) form \(P_h\). Among
the others, candidates of \ref{B3.51} with
\(\Sigma_p>\theta=x^{-\varepsilon}\) form \(P_2\), and the rest form \(P_1\).
\leanref{Formalism.\allowbreak{}L2.\allowbreak{}GlobalTheoryTheory.\allowbreak{}Loss.\allowbreak{}clAt}{Formalism.\allowbreak{}LocalToGlobal.\allowbreak{}GlobalTheory.\allowbreak{}SuppSat}{43--47}

\section{The finite-sum loss bound}

Write \([y]_+=\max(y,0)\) and \(\log_p y=\log y/\log p\).
Let \(C,f,S_k\ge0\) and let \(N_0\in\R\) be the rounding-error size parameter. Put
\[
\begin{aligned}
 F_p&=Cf\{N_0^2\log_p a+(\log_p a)^2\},&
 C_p&=CS_k\log_p(2a),\\
 R_p&=\begin{cases}Q_1,&p<p_0,\\Q_2,&p\ge p_0,\end{cases}&
 O_p&=C\ell_x\left\{\frac{\mathcal L a}{pk}
 +\begin{cases}ap^{-L_1(p)},&p<p_0,\\a/p^2,&p\ge p_0.\end{cases}\right\}.
\end{aligned} \tag{5.17.1}
\]
Here \(\ell_x,\mathcal L,Q_1,Q_2\) are nonnegative real parameters,
\(k,K_0>0\), and \(L_1\) is a real-valued function of \(p\).

\begin{theorem}[The losses from all prime classes]\label{B5.17}
Assume the hypotheses of \ref{B5.15} and \ref{B5.16}. Let
\(p_0,d\in\N\), \(p_0\ge3\), \(d\ge4\), and suppose
\[
 a\ge1,\quad X'\ge1,\quad X''\ge2,\quad w''>0,\quad
 X''w''=a,\quad 0\le\vartheta\le\tfrac12.
\]
For every \(t\in TT\), suppose \(0\le a_t\le a(1+\vartheta)\).
Suppose
\[
 A_{jtp}=\sum_{i\in I_{jtp}}[a_t/p^i-w_{jtpi}]_+,\qquad
 D_{jt}=\log d\sum_{i\in I^d_{jt}}[a_t/d^i-w^d_{jti}]_+,
 \tag{5.17.2}
\]
where \(I_{jtp}\) and \(I^d_{jt}\) are finite subsets of \(\N\).
Every \(i\in I_{jtp}\) is odd and has \(w_{jtpi}\ge w''\).
Every \(i\in I^d_{jt}\) is positive and has \(w^d_{jti}\ge0\).
Suppose further that
\[
\begin{gathered}
 p\le X''\ (p\in S),\qquad p>X''\ (p\in P_h),\\
 r_p\le F_p,\qquad
 \ell_p\le C_p+\mathbf1_{\{p<X'\}}(R_p+O_p)\quad(p\in S),\\
 K_0\le p^{L_1(p)-1}\quad(p\in P_s,\ p<p_0).
\end{gathered} \tag{5.17.3}
\]
Then
\[
 \ASK-\Ga\le
 E_{\rm fl}+E_{\rm cp}+E_{\rm pr}+E_{\rm or}
 +E_{\rm out}+E_1+E_2+E_h+E_d, \tag{5.17.4}
\]
where
\[
\begin{aligned}
 E_{\rm fl}&=2mCfX''\{N_0^2\log a+2(\log a)^2\},\\
 E_{\rm cp}&=2CS_k\log(2a)\,mn_TX'',\\
 E_{\rm pr}&=2\log4\,mn_T(p_0Q_1+X'Q_2),\\
 E_{\rm or}&=2C\ell_x a\,mn_T
 \left\{\left(\frac{\mathcal L}{k}+\frac1{K_0}\right)
 (\log X'+\log4)
 +\frac{2(\log(p_0-1)+1)}{p_0-1}\right\},\\
 E_{\rm out}&=4ma\sum_{p\in S}\frac{\log p}{p}|TT\setminus TT_p|,\\
 E_i&=4mn_Ta\sum_{p\in P_i}\frac{\log p}{p}\quad(i=1,2),\\
 E_h&=2mn_T\log(2X'')\{a\vartheta^2+w''\vartheta\},\\
 E_d&=2mn_Ta\,\frac{\log d}{d}.
\end{aligned} \tag{5.17.5}
\]
\end{theorem}
\leanref{Formalism.\allowbreak{}LFunctionChainParameters.\allowbreak{}DeficitSplitLemmas.\allowbreak{}loss\_\allowbreak{}bound}{Formalism.\allowbreak{}SizeWindowBrackets.\allowbreak{}ChainSlopeWindows.\allowbreak{}AdmissibleSchurLemma}{762--918}
\leanref{Formalism.\allowbreak{}LFunctionChainParameters.\allowbreak{}DeficitSplitLemmas.\allowbreak{}SecondKindLoss}{Formalism.\allowbreak{}SizeWindowBrackets.\allowbreak{}SecondKindBucket}{261--293}

\begin{proof}
Sum the geometric series over odd exponents at \(p\) and positive
exponents at \(d\). Since \(a_t\le3a/2\),
\[
 A_{jtp}\le a_t\frac{p}{p^2-1}\le\frac{2a}{p},\qquad
 D_{jt}\le\frac{a_t\log d}{d-1}\le\frac{2a\log d}{d}.
\]
These give \(E_{\rm out},E_1,E_2,E_d\) in (5.16.3).
For \(p>X''\), all \(i\ge2\) have \(p^i>2X''\) and
\(a_t/p^i<w''\). The first power contributes at most
\(w''\vartheta\) and vanishes for \(p\ge X''(1+\vartheta)\).
The interval contains at most \(X''\vartheta+1\) integers and
\(\log p\le\log(2X'')\) there. Its contribution per \(j,t\)
is at most \(2\log(2X'')w''\vartheta(X''\vartheta+1)\).
Summing over \(j,t\) gives \(E_h\).

For \(p\in S\), use \(1/\log p\le2\) to get
\[
 2\log p\,F_p\le2Cf\{N_0^2\log a+2(\log a)^2\},\qquad
 2\log p\,C_p=2CS_k\log(2a).
\]
Sum over at most \(X''\) primes, with factors \(m\) and \(mn_T\),
to obtain \(E_{\rm fl},E_{\rm cp}\). The upper bound on \(\ell_p\)
in (5.17.3) is nonnegative, so \(n_T\) may replace \(|TT_p|\).

For \(\Theta(y)=\sum_{p\le y}\log p\) with the sum over primes,
\[
 \Theta(Y)\le Y\log4,\qquad
 \sum_{p\le Y}\frac{\log p}{p}\le\log Y+\log4
 \quad(Y\ge1). \tag{5.17.6}
\]
Induction gives the first bound from
\[
 \Theta(2r)-\Theta(r)\le\log\binom{2r}{r}\le2r\log2,\qquad
 \Theta(2r+1)-\Theta(r+1)\le\log\binom{2r+1}{r}\le2r\log2.
\]
The two central coefficients in the odd case are equal.
For \(n=\lfloor Y\rfloor\), the second bound follows from
\[
 n\sum_{p\le n}\frac{\log p}{p}
 \le\sum_{p\le n}\lfloor n/p\rfloor\log p+\Theta(n)
 \le\log(n!)+\Theta(n)\le n(\log n+\log4).
\]
Apply the first bound separately to the two branches of \(R_p\)
with \(p<X'\) to obtain \(E_{\rm pr}\). For \(O_p\), use
\(p^{-L_1(p)}\le1/(K_0p)\) below \(p_0\) and
\[
 \sum_{n\ge p_0}\frac{\log n}{n^2}
 \le\frac{2(\log(p_0-1)+1)}{p_0-1}. \tag{5.17.7}
\]
For \(H(y)=(\log y+1)/y\) and \(n\ge3\),
\(\log n/n^2\le2(H(n-1)-H(n))\). The second bound in
(5.17.6) and (5.17.7) give \(E_{\rm or}\).
\end{proof}

\begin{lemma}[Lower bound for an ordinary partner]\label{B5.17a}
Let \(D,L,X\) be the model, size data and orbit instance of
\ref{B3.46}--\ref{B3.47}, with the model's integer weights \(b_s\).
For every \(\ell\in\N\) and \(z\in\{0,1\}\times A\),
\[
 \operatorname{fill}_{\ell}(z)\ge w'',\qquad w''=\frac{b'}{xn}.
\]
In the model sum of \ref{B5.15}, for \(1\le\ell\le m\), \(t\in TT\),
\(p\in P_s\) and \(i\in I_{\ell tp}\), the exponent \(i\) is odd and
\(w_{\ell tpi}=\operatorname{fill}_{\ell}(z_i(t))\ge w''\).
\end{lemma}
\leanref{Formalism.\allowbreak{}L2.\allowbreak{}GlobalTheoryTheory.\allowbreak{}Loss.\allowbreak{}pd\_\allowbreak{}fill}{Formalism.\allowbreak{}LocalToGlobal.\allowbreak{}GlobalTheory.\allowbreak{}PowFourLtCeilThresholdPrime}{1025--1037}
\leanref{Formalism.\allowbreak{}L2.\allowbreak{}GlobalTheoryTheory.\allowbreak{}Loss.\allowbreak{}pd\_\allowbreak{}wz}{Formalism.\allowbreak{}LocalToGlobal.\allowbreak{}GlobalTheory.\allowbreak{}PowFourLtCeilThresholdPrime}{1038--1045}

\begin{proof}
Since \(n\ge1\), \(a=xb\) and \(b'=a/X'\) with \(X'\ge1\),
we have \(0<w''\le b'/x=b/X'\le b\). At a marked position
\((0,s)\) the fill is \(b_s+\operatorname{off}_{\ell}(0,s)\),
with \(b_s\ge b'\) and offset at least \(b'/x\).
At a thick position the fill is its offset, at least \(b\).
At every remaining position the offset is at least \(w''\) by
\ref{B3.47}. The ordinary sum retains the odd indices and uses
\(\operatorname{fill}_\ell(z_i(t))\) by \ref{B3.46} and \ref{B3.52}.
\end{proof}

\section{Exceptional primes and exceptional positions}

For \(e\ge1\) and odd \(p\), put
\[
 \mathcal H_e=(\Z/2^{e+2}\Z)^\times/\langle-1\rangle,
 \qquad h_e(p)=[p]\in\mathcal H_e .
\]

\begin{lemma}[Small order and narrow prime-power bands]\label{B5.18}
There is an absolute constant \(C_{\rm ord}\ge0\) with the following
property. Let \(e,J_0\ge1\) be integers,
\(Y_0\ge1\), \(X'\ge2\), \(p_0\ge8\), and
\(0\le\rho\le1/2\). Let \(T\) be a finite set of odd primes. For each
\(p\in T\), assume
\[
 \begin{gathered}
 \bigl(p\le Y_0\ \land\ \operatorname{ord}(h_e(p))\le J_0\bigr)
 \quad\text{or}\\
 \bigl(p\ge p_0\ \land\ \exists i\in\N_{\ge1},\
 X'(1-\rho)\le p^i\le X'(1+\rho)\bigr).
 \end{gathered}
\]
Then
\[
 \sum_{p\in T}\frac{\log p}{p}
 \le C_{\rm ord}J_0\frac{\log^2(2^eY_0)}{2^e}
 +\lfloor\log X'\rfloor
 (4\rho+p_0^{-1})\log(2X'). \tag{5.18.1}
\]
One may take \(C_{\rm ord}=12\).

For a first class satisfying the small-order or band condition with
\[
 Y_0=\max(Y,X''),\quad
 J_0=\lfloor x^{\varepsilon/4}\rfloor,\quad
 I_b=\lfloor x^{\varepsilon/3}\rfloor=\lfloor\log X'\rfloor,
\]
and the strict order bound
\(\operatorname{ord}(h_e(p))<x^{\varepsilon/4}\) in the first
alternative, its loss obeys
\[
 E_1\le E_{\rm ord}+E_{\rm band}, \tag{5.18.2}
\]
where
\[
\begin{aligned}
 E_{\rm ord}
 &=4mn_Ta\,C_{\rm ord}J_0
 \frac{\log^2(2^e\max(Y,X''))}{2^e},\\
 E_{\rm band}
 &=4mn_Ta\,I_b(4\rho+p_0^{-1})\log(2X').
\end{aligned}
\]
\end{lemma}
\leanref{Formalism.\allowbreak{}LFunctionChainParameters.\allowbreak{}PrimeWindowSizeBounds.\allowbreak{}first\_\allowbreak{}kind\_\allowbreak{}bound}{Formalism.\allowbreak{}SizeWindowBrackets.\allowbreak{}ChainSlopeWindows.\allowbreak{}OneLeThresholdMulPowers}{955--986}
\leanref{Formalism.\allowbreak{}LFunctionChainParameters.\allowbreak{}PrimeWindowSizeBounds.\allowbreak{}first\_\allowbreak{}le}{Formalism.\allowbreak{}SizeWindowBrackets.\allowbreak{}ChainSlopeWindows.\allowbreak{}OneLeThresholdMulPowers}{1036--1100}

\begin{proof}
Put \(r=\min(e,\lfloor\log_2J_0\rfloor)\) and
\(M=2^{e+2-r}\). For an odd integer \(s>1\), the congruence
\[
 h_e(s)^{2^j}=1
 \quad\Longleftrightarrow\quad
 s\equiv1\text{ or }-1\pmod {2^{e+2-j}}
 \quad(0\le j\le e) \tag{5.18.3}
\]
follows by factoring successive differences of squares. For \(j\ge1\),
\[
 v_2(s^{2^j}-1)=v_2(s-1)+v_2(s+1)+j-1.
\]
One of \(s-1,s+1\) has valuation \(1\), and an odd square is
\(1\) modulo \(8\). This proves (5.18.3) for \(j\ge1\), while
\(j=0\) is the definition of the quotient.

The order in \(\mathcal H_e\) divides \(2^e\). If it is at most
\(J_0\), it divides \(2^r\), so \(p\equiv\pm1\pmod M\).
Also \(2^e\le J_0M\). In either progression write \(p=qM\pm1\),
\(q\ge1\). Since \(p\ge qM/2\),
\[
 \sum_{\substack{p\le Y_0\\p\equiv\pm1\ (M)}}\frac{\log p}{p}
 \le\frac{4\log Y_0}{M}
 \sum_{q\le\lfloor Y_0\rfloor}\frac1q
 \le\frac{4\log Y_0(1+\log Y_0)}{M}.
\]
Use \(\log Y_0(1+\log Y_0)\le3\log^2(2^eY_0)\) and
\(2^e\le J_0M\) for the first term of (5.18.1).
\leanref{Formalism.\allowbreak{}RankGapBounds.\allowbreak{}LocalQuadraticParts.\allowbreak{}ordersSumY\_\allowbreak{}of}{Formalism.\allowbreak{}RankGapBounds.\allowbreak{}RankGapProofs.\allowbreak{}PmModulusFacts}{252--287}
\leanref{Formalism.\allowbreak{}RankGapBounds.\allowbreak{}LocalQuadraticParts.\allowbreak{}powTwoPowEqOneIff\_\allowbreak{}holds}{Formalism.\allowbreak{}RankGapBounds.\allowbreak{}RankGapProofs.\allowbreak{}FourNotDvdOddSqAddOne}{132--138}

For the band condition,
\(2i\le i\log p\le\log(2X')\le2\log X'\), so
\(1\le i\le\lfloor\log X'\rfloor\).
Taking \(i\)-th roots puts \(p\) in
\([y_i(1-\rho),y_i(1+\rho)]\), where \(y_i=(X')^{1/i}\).
The interval contains at most \(2\rho y_i+1\) integers and
\(1/p\le\min(2/y_i,1/p_0)\). Hence in each band
\[
 \sum\frac{\log p}{p}
 \le(2\rho y_i+1)\min(2/y_i,1/p_0)\log(2y_i)
 \le(4\rho+p_0^{-1})\log(2X').
\]
Sum over \(i\) and add the small-order bound to get (5.18.1).
Multiplication by \(4mn_Ta\) gives (5.18.2).
\end{proof}

Let \(\mathcal P_c\) be a finite set of candidate primes with real
shares \(\Sigma_p\). Define the weighted share by
\[
 R_W=\sum_{p\in\mathcal P_c}\frac{\log p}{p}
 \bigl(2\mathbf1_{\{\Sigma_p>\theta\}}+4\Sigma_p\bigr).
 \tag{5.19.1}
\]

\begin{lemma}[Bad positions, second-class primes, and local indices]\label{B5.19}
Let \(m,a\ge0\) and let \(\mathcal P_R\) be a finite set of primes.
Suppose each missing position \(t\in TT\setminus TT_p\), \(p\in S\),
satisfies \(p\in\mathcal P_R\) and a condition \(\mathrm{BAD}(p,t)\).
If
\[
 \sum_{p\in\mathcal P_R}\frac{\log p}{p}
 |\{t\in TT:\mathrm{BAD}(p,t)\}|\le R_{\rm out},
\]
then \(E_{\rm out}\le4amR_{\rm out}\).

For the second-class estimate, suppose \(\Sigma_p\ge0\) on
\(\mathcal P_c\) and \(P_2\subseteq\{p\in\mathcal P_c:\Sigma_p>\theta\}\).
Then \(E_2\le2amn_TR_W\). For the index estimate, let the shares
be arbitrary reals and suppose
\[
 \mathrm{io}_p\le
 2\Sigma_p\,\frac{2mn_Ta}{p}
 \quad(p\in\mathcal P_c).
 \tag{5.19.2}
\]
Then
\[
 \sum_{p\in\mathcal P_c}\log p\,\mathrm{io}_p
 \le amn_TR_W.
 \tag{5.19.3}
\]
Both estimates allow any real \(T\ge0\) in place of \(n_T\)
and any upper bound \(R\ge R_W\) in place of \(R_W\).
\end{lemma}
\leanref{Formalism.\allowbreak{}LFunctionChainParameters.\allowbreak{}SmallOrderPrimeLosses.\allowbreak{}out\_\allowbreak{}discriminant\_\allowbreak{}bound}{Formalism.\allowbreak{}SizeWindowBrackets.\allowbreak{}ChainSlopeWindows.\allowbreak{}DiscriminantSettingSatisfied}{32--74}
\leanref{Formalism.\allowbreak{}LFunctionChainParameters.\allowbreak{}SmallOrderPrimeLosses.\allowbreak{}second\_\allowbreak{}discriminant\_\allowbreak{}bound}{Formalism.\allowbreak{}SizeWindowBrackets.\allowbreak{}ChainSlopeWindows.\allowbreak{}DiscriminantSettingSatisfied}{76--94}
\leanref{Formalism.\allowbreak{}LFunctionChainParameters.\allowbreak{}SmallOrderPrimeLosses.\allowbreak{}io\_\allowbreak{}discriminant\_\allowbreak{}bound}{Formalism.\allowbreak{}SizeWindowBrackets.\allowbreak{}ChainSlopeWindows.\allowbreak{}DiscriminantSettingSatisfied}{96--112}

\begin{proof}
Multiply the bad-position bound by \(4am\). On \(P_2\) the
indicator in (5.19.1) is one, so nonnegativity gives
\(E_2\le2amn_TR_W\). Multiply (5.19.2) by \(\log p\) and sum.
The indicators in \(R_W\) increase the bound.
\end{proof}

\section{Losses on the scale \(Q\)}

Let \(e_2\in\N\), and let \(n,k,N_0,t_0,r_G,S_c,\eta,\mathcal L\)
be real size parameters. Define
\[
\begin{gathered}
 z=x^\varepsilon,\quad u=x^{\varepsilon/3},\quad
 b=e^z,\quad a=xb,\quad X'=e^u,\quad b'=a/X',\\
 f=2n,\quad w''=\frac{b'}{xn},\quad X''=xnX',\quad
 \delta_1=x^{-\varepsilon/8},\quad\theta=x^{-\varepsilon},\\
 \ell_x=\log x,\quad S_k=S_c+k+\eta,\quad Q_1=Cr_Gb'\ell_x,\\
 Q_2=C\left\{\frac{b'}x+\eta+S_c+
 \frac{(w''+t_0)\log x}{\delta_1}\right\}+b_0.
\end{gathered} \tag{5.20.1}
\]
Here \(b_0\ge0\) is the additional local comparison error and
\(\mathcal L\) is \Lean{LaG} in the model.
Fix \(\nu\ge0\) and \(K_s\ge1\). The size window consists of
\[
\begin{gathered}
 x^2\le2^{e_2}\le n\le x^\nu,\qquad x^2\le k\le x^\nu,\\
 0\le N_0,t_0,r_G\le x^\nu,\qquad
 x^4+1\le p_0\le x^\nu,\qquad
 0\le\mathcal L\le x^{2\varepsilon},\\
 0\le S_c,\eta\le e^{x^{\varepsilon/4}},\qquad
 0\le\rho\le x^{-2},\qquad0\le\vartheta\le x^{-1},\\
 K_0\ge x^2,\qquad d\ge x^2,\qquad
 0\le C,C_{\rm ord}\le K_s,\qquad
 1\le Y\le2X',\qquad 0\le b_0\le e^{z/2}.
\end{gathered} \tag{5.20.2}
\]
Here \(p_0,d\) are integers. Assume
\[
\begin{gathered}
 x\ge1,\qquad100K_sx^{4\varepsilon}\le x,\\
 100K_sx^{4\nu+4}\le e^u,\qquad
 100K_sx^{4\nu+4}e^{u+x^{\varepsilon/4}}\le e^{z/2}.
\end{gathered} \tag{5.20.3}
\]
Put
\[
 \sigma_x=100K_sx^{4\varepsilon-1}.
 \tag{5.20.4}
\]

\begin{theorem}[Joint loss bounds in a size window]\label{B5.20}
An absolute choice of \(C_{\rm ord}\ge0\) has the following property.
Let \(0<\varepsilon\le1\) and assume (5.20.1)--(5.20.3).
Let \(TT\ne\varnothing\) and \(J\) be finite. Suppose every prime
in \(P_s\) is odd.
For every member \(\upsilon\) of an indexed family of finite-sum
data, assume the hypotheses of \ref{B5.17}.
Assume its first class satisfies the condition of \ref{B5.18} with
\(e=e_2\) and \(Y_0=\max(Y,X'')\), with
\(\operatorname{ord}(h_{e_2}(p))<x^{\varepsilon/4}\) in the
small-order alternative.
Suppose the hypotheses of \ref{B5.19} hold uniformly in \(\upsilon\),
with common shares and bad-position condition satisfying
\[
 R_{\rm out}\le n_Tx^{-2+\varepsilon},
 \qquad R_W\le x^{-9}. \tag{5.20.5}
\]
Then, with \(Q=n_T|J|b\), for every \(\upsilon\),
\[
\begin{aligned}
 \ASK_\upsilon-\Ga_\upsilon
 &\le(2C\log4+13\sigma_x)Q,\\
 \sum_{p\in\mathcal P_c}\log p\,\mathrm{io}_{\upsilon,p}
 &\le\sigma_xQ.
\end{aligned} \tag{5.20.6}
\]
In particular,
\[
 \ASK_\upsilon-\Ga_\upsilon\le(3C+13)Q,\qquad
 \sum_{p\in\mathcal P_c}\log p\,\mathrm{io}_{\upsilon,p}\le Q.
 \tag{5.20.7}
\]
For fixed \(\nu,K_s\) and \(0<\varepsilon<1/4\), the largeness
conditions hold for all sufficiently large \(x\).
\end{theorem}
\leanref{Formalism.\allowbreak{}LFunctionChainParameters.\allowbreak{}SmallOrderPrimeLosses.\allowbreak{}joint\_\allowbreak{}core}{Formalism.\allowbreak{}SizeWindowBrackets.\allowbreak{}ChainSlopeWindows.\allowbreak{}DiscriminantSettingSatisfied}{178--226}
\leanref{Formalism.\allowbreak{}LFunctionChainParameters.\allowbreak{}SmallOrderPrimeLosses.\allowbreak{}joint\_\allowbreak{}discriminant\_\allowbreak{}bound}{Formalism.\allowbreak{}SizeWindowBrackets.\allowbreak{}ChainSlopeWindows.\allowbreak{}DiscriminantSettingSatisfied}{228--232}
\leanref{Formalism.\allowbreak{}LFunctionChainParameters.\allowbreak{}SmallOrderPrimeLosses.\allowbreak{}joint\_\allowbreak{}discriminant\_\allowbreak{}eventually}{Formalism.\allowbreak{}SizeWindowBrackets.\allowbreak{}ChainSlopeWindows.\allowbreak{}DiscriminantSettingSatisfied}{234--241}

In (5.20.2) the archive parameters \(\beta,s_0\) are zero.

\begin{proof}
Take \(C_{\rm ord}=12\) by \ref{B5.18}. The window gives \(p_0\ge8\),
\(d\ge4,\ e_2\ge1,\ \vartheta\le1/2,\ X''w''=a\).
For \(J=\varnothing\), (5.16.1) gives \(\tau_p\le0\) on \(S\)
and \(\tau_p=0\) elsewhere. Apply (5.15.1) and (5.19.2).
Below, \(m>0\).

Put \(E=x^\nu\), \(S_x=e^{x^{\varepsilon/4}}\), \(B=e^{z/2}\) and
\(D_0=100K_sE^4x^4\). By (5.20.3),
\[
 X'\ge D_0,\quad B\ge D_0X'S_x,\quad b=B^2,\quad
 x\ge100,\quad z\ge1,\quad 0\le\sigma_x\le1. \tag{5.20.8}
\]
In particular \(\log x,\nu\log x\le u\le z\) and
\[
 \log a\le2z,\qquad \log(2a)\le\log2+2z\le x. \tag{5.20.9}
\]
Split \(E_{\rm pr}\) according to \(p_0Q_1\) and \(X'Q_2\). Substituting
\(n,N_0\le E,\ n_T\ge1,\ S_k\le3ES_x\) in (5.17.5) gives
\[
\begin{aligned}
 E_{\rm fl}/Q&\le12K_sE^4x^3X'/b,\\
 E_{\rm cp}/Q&\le6K_sE^2x^2X'S_x/b,\\
 E_{\rm pr,low}/Q&\le3K_sE^2x^2/X'.
\end{aligned} \tag{5.20.10}
\]
Each bound is at most \(\sigma_x\) by (5.20.8).
The high branch uses
\[
 \frac{X'Q_2}{b}
 = C+\frac{CX'(\eta+S_c)}b
 +\frac{C\log x\,x^{\varepsilon/8}}n
 +\frac{CX't_0\log x\,x^{\varepsilon/8}}b
 +\frac{X'b_0}b. \tag{5.20.11}
\]
Since \(x^{\varepsilon/8}\le z\) and \(z^2\le x\), the four terms
after \(C\) are bounded by \(2K_sX'S_x/b,\ K_sz^2/x^2,\
K_sX'Ex/b,\ X'/B\), respectively. Each is at most \(\sigma_x/4\).
Thus
\[
 E_{\rm pr,high}/Q\le2\log4(C+\sigma_x)
 \le2C\log4+3\sigma_x. \tag{5.20.12}
\]
In \(E_{\rm or}\), the bounds \(\mathcal L/k\le z^2/x^2\),
\(K_0^{-1}\le x^{-2}\), \(\log X'+\log4\le3z\) and
\(2(\log(p_0-1)+1)/(p_0-1)\le4z/x^2\) give
\(E_{\rm or}/Q\le20K_sz^4/x\le\sigma_x\).

The bounds \(x^2\le2^{e_2}\le n\le E\) and \(Y\le2X'\) give
\(\log(2^{e_2}\max(Y,X''))\le4z\). With \(J_0,I_b\le z\)
and \(\rho,p_0^{-1}\le x^{-2}\), \ref{B5.18} yields
\[
 E_{\rm ord}/Q\le64K_sz^3/x,\qquad
 E_{\rm band}/Q\le40z^2/x. \tag{5.20.13}
\]
By \ref{B5.19} and (5.17.5),
\[
\begin{aligned}
 E_{\rm out}/Q&\le4z/x,& E_2/Q&\le2x^{-8},\\
 E_h/Q&\le16z/x,& E_d/Q&\le4z/x.
\end{aligned} \tag{5.20.14}
\]
Here \(\log(2X'')\le4z\), \(a\vartheta^2,w''\vartheta\le b/x\),
and \(\log d/d\le2\log x/x^2\) by monotonicity for \(d\ge x^2\).

The ten terms outside \(E_{\rm pr,high}\) are each bounded by
\(\sigma_xQ\). Adding (5.20.12) gives the deficit bound.
The index bound is \(amn_Tx^{-9}=x^{-8}Q\le\sigma_xQ\) by
\ref{B5.19}. Since \(\sigma_x\le1\) and \(2\log4\le3\), this
gives (5.20.7). For fixed \(\nu,K_s\), \(4\varepsilon<1\) and
\(\log x=o(u)\), \(u+x^{\varepsilon/4}=o(z)\) give (5.20.3)
eventually.
\end{proof}

\begin{theorem}[Local comparison on the design class]\label{B5.20a}
Given the twist selection and family duties of \ref{B8.02}, the
following assertion is closed on \(\mathcal D_3(\mathcal A)\) in the
sense of \ref{B8.39}. There are \(C\ge0\) and \(x_0\), uniform in the
family, such that every model at \(x\ge x_0\) has the following properties.
\begin{enumerate}
\item At \(w\) of residue characteristic \(p\ne T_u\), the containment
of \ref{B5.05} holds at exponent \(\lfloor e_w\tau_w\rfloor_+\).
The normalized length satisfies
\(\mathrm{lenMN}_{uw}\le d_M(\tau_w+\mathrm{io}_w)\), with \(d_M\)
as in \ref{B4.07a}. If \(V_p=\varnothing\), then
\(\tau_w=\mathrm{io}_w=\mathrm{lenMN}_{uw}=0\).
\item For \(v\in V_p\), \(S_v\le2mN_va/p\). At every \(w\) over
a signal or inert-class prime \(p\le X''\), one has
\(0\le L.\mathrm{gain}(j,w)\le4n_Ta/p\) for \(1\le j\le m\).
The comparison (5.16.1) holds with the covered positions and gains
at \(w\). In (5.17.1) and (5.20.1) take \(N_0=N_m\), \(t_0=t\),
\(S_c=\mathrm{Sc}\), \(L_1(p)=1+\lceil\log_p(k^2x^{78})\rceil\)
and \(b_0=b^{x^{-\varepsilon/2}}\). Then \(r_p=F_p\) and
\(\ell_p=C_p+\mathbf1_{\{p\ {\rm inert\ class},\ p<X'\}}(R_p+O_p)\).
\end{enumerate}
On good candidate orbits the mean of \(\kappa_v\) over the local
integer weight box is at most \(N_vx^{-12}\). Below \(p_0\) the
bound holds pointwise.
\end{theorem}

\begin{proof}
The duties tie the lattices to the twist. Use \ref{B5.20b}
below \(p_0\) and the minimum in \ref{B3.52a} above the split.
Extend to the coefficient DVR. The bounds in \ref{B5.20b} and
\ref{B5.20c} give the claimed pointwise and mean estimates.

For \(V_p\ne\varnothing\), \(\Sigma_p\le x^{-\varepsilon}\le1/2\)
for large \(x\). Apply \ref{B4.07a} with the model sandwich and
saturated ramified part. The neighbours descend by rationality,
and flat extension preserves images and intersections in (4.6).
Thus the localized pair extends to \(R,N_P\).
Each composition factor at \(w\) acquires length
\(v_{\mathcal O}(p)/e_w\), so the extended quotient has length
\(v_{\mathcal O}(p)\mathrm{lenMN}_{uw}\).
Divide the bound in \ref{B4.07a} by \(v_{\mathcal O}(p)\) and descend
the inclusion. Off \(T_u\) the localized lattices agree for empty selection.

The price and digit loss are nonnegative. Since \(\bar\alpha_s\le6a/5\)
and \(p/(p^2-1)\le4/(3p)\), the unrounded bound for \(S_v\) is
\(8mN_va/(5p)\). For \(N_v>0\), \(a/p\ge b'\ge500\) absorbs rounding.
For \(N_v=0\), the ceiling is zero. The same series bounds
\(g_0,g_U\) of \ref{B3.48}--\ref{B3.49} by \(8N_va/(5p)\) per lift.
Sum the maxima of \ref{B3.51} over disjoint orbits.

For a fixed lift and top \(s\), put \(f_h=\operatorname{fill}_j(z_h(s))\).
The sums \(A_s=\sum_{h\ {\rm odd}}[a_s/p^h-f_h]_+\) and
\(M_s=\sum_{h\ {\rm odd}}\min(a_s/p^h,f_h)\) add to \(pa_s/(p^2-1)\).
Decreasing \(a_s\) by \(2c\) costs at most \(2c\) in the series.
Increasing the capacities by at most \(c\) costs at most
\(c\log_p(2a)\), since \(f_h\ge1\) and only \(p^h<2a\)
can contribute to the difference. Thus on a good candidate orbit
with \(N=N_v\),
\[
 \sum_s A_s+\sum_{s\ {\rm even}}M_s
 \le g_U+\bar\delta+3Nc\log_p(2a).
\]
For \(p_0\le p<X'\), let \(h\) be the least positive backward
distance from an odd mark \(y\) contributing to \(J\) to an even
mark. Pair \(y\) with \((y-h,h)\), with positions modulo \(n'\).
The even period makes \(h\) odd. The pair determines \(y\), and its
term in \(M_{y-h}\) pays
\(\min(b_y,a_{y-h}/p^h)\). The marked part of the price adds at
most \(N(b'/x+\eta)\).

Write \(G_1=C_{31}\log x/\delta_1\). Good orbits give \(n'\le NG_1\).
Below \(p_0\) the size bounds give
\[
 G_1(w''+\eta+c+N_m),\quad kq_p,\quad
 K_{\rm DII}(H^3+r_G^5+k^2)\log_p a\ \le b',
 \qquad c_0\le3b',
\]
with \(H=N_{\rm win}\). Substitution in \ref{B3.50} and \ref{B3.52a}
bounds \(\mathrm{Pr}+\lceil t_{II}\rceil_++D_{II}+1\) by
\(N(Q_1+O_p)\) with the internal comparison constant.
Put \(L_1^{\rm loc}(p)=\lceil\log_p((kx^{26})^4)\rceil\).
For \(p<p_0\), \(pk^2x^{78}\le(kx^{26})^4\), so
\(L_1(p)\le L_1^{\rm loc}(p)\) and
\(ap^{-L_1^{\rm loc}(p)}\le ap^{-L_1(p)}\). Above the split,
\(H=O_\varepsilon(\log x)\) bounds the \(H\)-term by a constant
times \(NO_p\), and \(n'w''\le NG_1w''\).
After pairing, the remaining costs per mark are at most
\(Y_0^{13}\le b_0\), where \(Y_0=\exp(x^{23\varepsilon/48})\).
The bounds \(G'=O_\varepsilon(x^{\varepsilon/8}\log x)\) and
\(\log p\le x^{\varepsilon/3}\) give \(q_p\le Y_0\).
The size windows bound \(c,\eta,N_m\) and the polynomial factors
by \(Y_0\). Expanding \(D_I,c_6\) gives degree at most twelve
in these bounds.

For each selected orbit and lift, these estimates give
\[
 \sum_{s\in TT\cap v}(A_s-\ell_p^{\rm int})
 \le g_U+\bar\delta-\mathrm{Pr}-\lceil t_v\rceil_+-D_p-1,
\]
where \(\ell_p^{\rm int}\) uses the internal comparison constant
\(C_{\rm int}\). Sum over the selected orbits and lifts. The retained
gain dominates \(g_U-\mathrm{Pr}\), and the presentation costs sum to
\(\tau_p-X.\mathrm{ep}(p)\). Since \(\sum_{\rm reps}n_p\le f\),
\(\sum_v N_v\le n_T\le f\) and \(k\le N_m\), the formulas in
\ref{B3.47} and \ref{B3.52a} give
\(\sum_v\bar\delta(v)+X.\mathrm{ep}(p)/m\le F_p\) with an enlarged constant.

For empty selection \(\tau_p=0\). An inert-class prime \(p<X'\)
then has no covered positions. Signal primes have no excluded odd
partners. For \(p\ge X'\), an excluded partner is marked and
\(a_s/p^h\le b'+b'/x\le f_h\). If \(p\ge X'\) is of inert class
and \(h_p=0\), every odd partner of \(s\) is \((0,s)\), so
\(a_s/p^h\le b'+b'/x\le\operatorname{fill}_j(0,s)\). Thus
\(A_s=0\) for every top, and (5.16.1) holds with zero gain and
deduction. For signal primes and inert-class primes \(p\ge X'\)
with \(h_p\ne0\), the ordinary gain suffices with the same allowances.
The finite arithmetic sum retains the terms of \(A_s\) in its
exponent range, with the same weights by \ref{B3.52}.
With \(c=4\mathrm{Sc}\), take \(C=4|C_{\rm int}|+1\) in (5.17.1).
At a served prime the filling identities give \(TT_p,g_{jp},\tau_p\)
from any \(w\mid p\). Substituting these identities gives (5.16.1).
\end{proof}
\leanref{Formalism.\allowbreak{}MainNotions.\allowbreak{}G.\allowbreak{}Sv}{Formalism.\allowbreak{}SharpRestriction.\allowbreak{}GrowthBounds.\allowbreak{}IsResidueClassZero}{336--337}
\leanref{Formalism.\allowbreak{}MainNotions.\allowbreak{}G.\allowbreak{}Joints.\allowbreak{}linear\_\allowbreak{}programming\_\allowbreak{}input\_\allowbreak{}of\_\allowbreak{}hypotheses}{Formalism.\allowbreak{}SharpRestriction.\allowbreak{}GrowthBounds.\allowbreak{}Joints}{2040--2128}
\leanref{Formalism.\allowbreak{}MainNotions.\allowbreak{}AuxBlock8.\allowbreak{}Pf.\allowbreak{}gain\_\allowbreak{}bound\_\allowbreak{}at\_\allowbreak{}of\_\allowbreak{}star}{Formalism.\allowbreak{}SharpRestriction.\allowbreak{}GrowthBounds.\allowbreak{}LogarithmicSums.\allowbreak{}NotSingularOfCentral}{619--729}
\leanref{Formalism.\allowbreak{}L2.\allowbreak{}GlobalTheoryTheory.\allowbreak{}Loss.\allowbreak{}loss\_\allowbreak{}bound\_\allowbreak{}at}{Formalism.\allowbreak{}LocalToGlobal.\allowbreak{}GlobalTheory.\allowbreak{}PowFourLtCeilThresholdPrime}{766--817}
\leanref{GlobalTheory.\allowbreak{}Top.\allowbreak{}closed\_\allowbreak{}linkG\_\allowbreak{}of\_\allowbreak{}duty}{Development.\allowbreak{}AutomorphicArithmetic.\allowbreak{}ThetaCorrespondence.\allowbreak{}ThetaPairLifts.\allowbreak{}ThirdPlaceSeries.\allowbreak{}LinkG}{13--24}

\begin{lemma}[Projection at rank \(r_G\)]\label{B5.20b}
Under the family hypotheses of \ref{B5.20a}, eventually every good
candidate orbit below \(p_0\) passes the test of \ref{B3.52a} at
rank \(r=r_G\), cost \(\lceil t_{II}\rceil_+\), denominator \(D_{II}\)
and fraction \(\min(1,C'Nr^2k/x^{26})\), where \(N=N_v,\ C'=K.\mathrm{kII}\).
\end{lemma}
\leanref{Formalism.\allowbreak{}MainNotions.\allowbreak{}G.\allowbreak{}Joints.\allowbreak{}type\_\allowbreak{}two\_\allowbreak{}hypothesis\_\allowbreak{}of\_\allowbreak{}counts}{Formalism.\allowbreak{}SharpRestriction.\allowbreak{}GrowthBounds.\allowbreak{}Joints}{425--582}

\begin{proof}
Put \(Y=x^{26}\) and \(c=4\mathrm{Sc}\).
The orbit bounds give \(2L_1^{\rm loc}\le L_2<u_1,\ 3G'\le n'\)
and \(N\le40n'\delta_1\). The bounds \(n'\le NC_{31}\log x/\delta_1\)
and \(n'\ge x^{\varepsilon/4}\) give \(N\ge x^{\varepsilon/8}/(C_{31}\log x)\).
For large \(x\), \(N_1=N_{\rm win}=O_\varepsilon(\log x)\), fixed \(r\)
and \(k\ge x^{9/2}\) give \(N_1^2\le rN,\ N_1+3r\le k\) and
\(K.\mathrm{kII}'N_1^2\log a\le k\). Also \(p\ge13\) by splitting,
\(2\le r\le Y\) and \(256k^2NY^2\le k^2x^{78}\le p^{L_1^{\rm loc}-1}\).
The box gives \(a-c+\lfloor c\rfloor\le a_s\le2(a-c)-\lfloor c\rfloor\)
and \(b_s+2\lfloor c\rfloor\le c_0\).
Since \(p\le x^{38}\), \(\log q_p=O_\varepsilon(x^{\varepsilon/8}\log^2 x)\).
The formula in \ref{B3.52a} gives \(D_{II}=O(x^{16}b')\) from
\(n'\le x^3,\ k\le x^6,\ c,c_0=O(b')\). Thus
\[
 K.\mathrm{szII}\, k^3r^2Nq_p^5[8(D_{II}+1)+rNc_0+2rn'c]Y=o(a)
\]
because \(\log(a/b')=x^{\varepsilon/3}\). The plane estimate at height \(a-c\) gives loss \(CNr^2k/Y\), with cost and denominator bounded
by those of \ref{B3.52a} since the family constants dominate its constants.
Enlarge \(C\) to \(C'\) and cap the loss at one by allowing retained rank zero.
\end{proof}

\begin{lemma}[Mean projection loss at rank one]\label{B5.20c}
In the large families of \ref{B5.22a}, eventually each good candidate
orbit at \(p\ge p_0\) has mean least passing fraction in \ref{B3.52a}
at most \(C_m NN_1/Y\) in the uniform local box, with
\(N=N_v,\ N_1=N_{\rm win},\ Y=x^{13}\).
\end{lemma}
\leanref{Formalism.\allowbreak{}MainNotions.\allowbreak{}flipOneThree.\allowbreak{}Conclusion2}{Formalism.\allowbreak{}SharpRestriction.\allowbreak{}GrowthBounds.\allowbreak{}ProjectionLemmaBaseLevel}{78--91}

\begin{proof}
The tiles give \(L_2<u_1,\ 3G'\le n'\) and at least two cyclic
segments between even marks, of length at most \(G'\) and with at
most \(N_1\) interior marks each. The candidate test gives \(k^2Y^2\le p<X'\).
Put \(Y_b=x^\varepsilon/20\). From \(N\le n'\le x^3,\ k\ge x^{9/2}\)
and \(\mathcal L=O_\varepsilon(x^\varepsilon\log^2x)\) we get
\(2N_1+2\le k-1,\ 2Y_bN\mathcal L<k,\ K.\mathrm{kI}\,N\mathcal L^2\le k\).
The \(b\)-width \(\Delta\) satisfies
\(Y_bN\mathcal L b'/k\le\Delta\le \rho_B b'/4\) by \(N\le40n'\delta_1\).
Take the projection threshold \(CH\le\min(K.\mathrm{kI},K_{\rm DI},K.\rho/2)\).
Then \(2CHN_1\mathcal L/k+4\Delta/b'\le\rho_B=o(1)\), so the
candidate exclusion implies the projection's band condition.
Eventually \(CH\le Y_b\),
\(8\log(NY)\le\mathcal L,\ c\le\mathcal L b'/k\) and
\(kY^2Y_bN\mathcal L\log a\le X'\). The formula for \(D_I\) supplies
the denominator bound and gives \(D_I\le x^{49}\). From
\(q_p\le\exp(x^{23\varepsilon/48})\) and \(c=O(\exp(x^{\varepsilon/4}))\),
\[
 K.\mathrm{kI}\,k^3Nq_p^4[8(D_I+1)+(N+n')c]Y=o(b').
\]

The rank-one column estimates give a loss majorant common to all
dominant base weights, with margin at most \(C_sNN_1/Y\).
Fix the \(b\)-weights. A segment uses the \(a\)-weights at its \(s\)
interior marks and next even mark, a uniform tuple with mean loss
at most \(C_A(s+1)N_1/Y\). These sets of marks are disjoint, so
\(\sum(s+1)\le N\). The cost fits \ref{B3.52a} using \(b'\le a/p\).
The least passing fraction is bounded by the majorant, whose mean
is at most \((C_A+C_s)NN_1/Y\).
\leanref{Formalism.\allowbreak{}MainNotions.\allowbreak{}flipOneThree.\allowbreak{}Pf.\allowbreak{}segment\_\allowbreak{}mean\_\allowbreak{}kappa\_\allowbreak{}bound}{Formalism.\allowbreak{}SharpRestriction.\allowbreak{}GrowthBounds.\allowbreak{}ColumnSpanDensity.\allowbreak{}LocalHypothesisOfGlobal}{143--177}
\end{proof}

\section{Arithmetic models and common weights}

Let \(F_m\) be an admissible family with indices \(i\in I(x)\) and
gains and lattices as in \ref{B5.01}--\ref{B5.14}. Its local index is
\[
 \mathrm{IO}_i
 =\sum_{w\in\operatorname{supp}_i}
 \frac{e_w\log N\mathfrak p_w}{[K:\Q]}\,\mathrm{io}_{i,w}.
 \tag{5.21.1}
\]
A consistent presentation uses the model's frame, layer count,
host arithmetic datum and top labels with their weights. At
the code of an embedding \(\tau\), \(w_j\) reads the weight at
lift \(j+1\) and \(\tau\).

\begin{theorem}[Loss bounds for a small-share model]\label{B5.21}
Suppose \(F_m\) has the following eventual properties.
The condition \(\Lean{ChosenTT}\) gives
\[
 \frac{\delta_1n}{2}\le n_T\le2\delta_1n,
 \qquad\delta_1=x^{-\varepsilon/8}.
\]
Let \(\mathrm{BAD}(p,t)\) mean that the orbit of \(t\) at \(p\)
fails the good-orbit condition. By \(\Lean{ChosenBadp}\), the sum of
\((\log p/p)|\{t\in TT:\mathrm{BAD}(p,t)\}|\) over inert-class
primes \(p<X'\) with \(\operatorname{ord}(h_p)\ge x^{\varepsilon/4}\)
is at most \(n_Tx^{-2+\varepsilon}\), where \(h_p\) is the Frobenius
component of \ref{B3.46}. Assume the local comparison of \ref{B5.20a}
with fixed \(C_G\) and \(b_0=b^{\delta(x)}\), where \(\delta(x)\to0\).

Then there exist \(c_{\rm Loss}\ge0\) and \(x_0\) such that, for every
\(x\ge x_0\) and every \(i\in I(x)\) with \(R_W(i)\le x^{-9}\),
\[
 \mathrm{IO}_i\le c_{\rm Loss}Q_i,
 \qquad Q_i=|TT_i|m_i b_i. \tag{5.21.2}
\]
For every consistent arithmetic presentation of this model and
every cutoff \(N_{\rm cut}\in\N\) satisfying
\[
 a_t/w_j(\rho)\le N_{\rm cut}
 \quad(j\in\mathrm{Fin}(m_i),\ t\in TT_i,\ \rho\in\Phi),
 \tag{5.21.3}
\]
its arithmetic sum satisfies
\[
 \ASK_{N_{\rm cut}}-\Ga_i
 \le c_{\rm Loss}Q_i.
 \tag{5.21.4}
\]
\end{theorem}
\leanref{Formalism.\allowbreak{}L2.\allowbreak{}GlobalTheoryTheory.\allowbreak{}Loss7.\allowbreak{}losses\_\allowbreak{}of\_\allowbreak{}small\_\allowbreak{}share\_\allowbreak{}weight}{Formalism.\allowbreak{}LocalToGlobal.\allowbreak{}GlobalTheory.\allowbreak{}LossRequirementsOfSmallShares}{86--161}

\begin{proof}
The family size bounds give (5.20.2) uniformly for large \(x\),
with \(C=\max(C_G,0)\) and \(b_0=b^{\delta(x)}\le b^{1/2}\).
Increase the threshold for (5.20.3) and \(k^2x^{26}\le X'\).
Then \(p<p_0=\lceil k^2x^{26}\rceil\) implies \(p<X'\), and
\(\Lean{ChosenTT}\) gives \(TT_i\ne\varnothing\).
\leanref{Formalism.\allowbreak{}L2.\allowbreak{}GlobalTheoryTheory.\allowbreak{}Loss.\allowbreak{}windows\_\allowbreak{}ev}{Formalism.\allowbreak{}LocalToGlobal.\allowbreak{}GlobalTheory.\allowbreak{}PowFourLtCeilThresholdPrime}{1555--1595}

For (5.17.3) take \(K_0=k^2x^{78}\) and
\(L_1(p)=1+\lceil\log_p K_0\rceil\), so \(p^{L_1(p)-1}\ge K_0\).
The top window gives \(0\le a_t\le a(1+\vartheta)\) with
\(\vartheta=(kx^{13})^{-1}\). Lemma~\ref{B5.17a} supplies the ordinary
exponent and weight conditions of (5.17.2).
Ramified exponents are positive and their weights are nonnegative.
\leanref{Formalism.\allowbreak{}L2.\allowbreak{}GlobalTheoryTheory.\allowbreak{}Loss.\allowbreak{}side\_\allowbreak{}at}{Formalism.\allowbreak{}LocalToGlobal.\allowbreak{}GlobalTheory.\allowbreak{}PowFourLtCeilThresholdPrime}{1199--1209}

Take \(\rho=\rho_B\) and \(Y=X'(1+C_r\rho_B)\) as in \ref{B4.04}.
Each word prime \(T_u\) is split, so on \(P_s\) a prime is discarded
exactly when \(\mathrm{NoGain}(p)\) holds. Thus served primes satisfy
\(p\le X''\). At a discarded prime, the selection rule of \ref{B3.51}
gives \(V_p=\varnothing\). A prime in the odd-order branch fails
the order test because \(\operatorname{ord}(h_p)\) divides \(o\) and
\(o\le x^{\varepsilon/100}<x^{\varepsilon/4}\).
The other branches violate the candidate conditions or give
\(\Sigma_p>\theta\). Hence \(\tau_p=0\), and the adjusted gains vanish.
\leanref{Formalism.\allowbreak{}L2.\allowbreak{}GlobalTheoryTheory.\allowbreak{}Loss.\allowbreak{}nonserved\_\allowbreak{}zero}{Formalism.\allowbreak{}LocalToGlobal.\allowbreak{}GlobalTheory.\allowbreak{}PowFourLtCeilThresholdPrime}{687--698}

At served primes \ref{B5.20a} gives (5.16.1) and the error bounds
in (5.17.3). In the small-order branch of \(P_1\), \(p\le Y\)
and \([p]\in\mathcal H_{e_2}\) has order dividing
\(\operatorname{ord}(h_p)<x^{\varepsilon/4}\). For inert-class
\(p\le X''\) with nonzero \(h_p\) of odd order, \([p]=1\)
in the \(2\)-group. Apply \ref{B5.18} with \(Y_0=\max(Y,X'')\).
The band branch satisfies its interval condition.

Missing positions at served primes lie on bad orbits at the primes
counted by \(\Lean{ChosenBadp}\), which bounds their weighted count.
The shares of \ref{B3.51} and \ref{B3.52a} are nonnegative, so
\(\Sigma_p\ge0\). If \(V_p\ne\varnothing\), multiply the bound
\(\max_{v\in V_p}S_v\le2mn_Ta/p\) from \ref{B5.20a} by \(2\Sigma_p\)
to get (5.19.2). If \(V_p=\varnothing\), its left side is zero
and its right side is nonnegative.
\leanref{Formalism.\allowbreak{}L2.\allowbreak{}GlobalTheoryTheory.\allowbreak{}Loss7.\allowbreak{}loss\_\allowbreak{}setting\_\allowbreak{}family}{Formalism.\allowbreak{}LocalToGlobal.\allowbreak{}GlobalTheory.\allowbreak{}LossRequirementsOfSmallShares}{41--80}

The index vanishes off the candidate fibers. These lie below
\(X'\) and have full support, so \ref{B5.15} gives
\(\mathrm{IO}_i=\sum_{p\in\mathcal P_c}\log p\,\mathrm{io}_{i,p}\).
\leanref{Formalism.\allowbreak{}L2.\allowbreak{}GlobalTheoryTheory.\allowbreak{}Loss.\allowbreak{}IO\_\allowbreak{}at}{Formalism.\allowbreak{}LocalToGlobal.\allowbreak{}GlobalTheory.\allowbreak{}PowFourLtCeilThresholdPrime}{253--268}

By \ref{B5.21a}, \(w_j(\rho)\ge1\) on \(\Phi\), so (5.21.3) holds
in particular for \(N_{\rm cut}=\lceil2a\rceil\). By \ref{B5.01}, any cutoff satisfying (5.21.3) includes all positive terms. The ordinary primes of \ref{B5.01}, characterized by
\(p\notin\{2,q_1,d\}\) and \(F_p\notin\Phi\), are precisely the
signal and inert-class primes. At each such prime, \(iF_p\notin\Phi\)
exactly when \(i\) is odd. At \(d\), the coset \(iF_d+I_d\) contains
exactly one element outside \(\Phi\), the one used by the model
at exponent \(i\).

In both cases \(r_t(g)=c+g+\mathrm{top}(t)\) has first bit
\(1+1=0\). By \ref{B5.21a} it has a unique matching place, where
consistency fixes the weight. For odd \(i\) at an ordinary prime,
\(z_i(t)\) and \(r_t(iF_p)\) have \(A\)-coordinate \(t+ih_p\)
and remaining bit equal to the signal-class indicator.
At a place above \(p\), the two filling identities of \ref{B3.52}
give \(\operatorname{fill}_{j+1}(z_i(t))\) for both the ordinary
partner weight and the weight at the matching archimedean index.
\leanref{Formalism.\allowbreak{}MainNotions.\allowbreak{}Lnk.\allowbreak{}wPart\_\allowbreak{}eq}{Formalism.\allowbreak{}SharpRestriction.\allowbreak{}GlobalSchurBounds.\allowbreak{}FrameCodeOfAdditive}{464--478}
\leanref{Formalism.\allowbreak{}MainNotions.\allowbreak{}Lnk.\allowbreak{}reads\_\allowbreak{}w}{Formalism.\allowbreak{}SharpRestriction.\allowbreak{}GlobalSchurBounds.\allowbreak{}FrameCodeOfAdditive}{223--235}
Thus the cutoff and model sums agree term by term.
\leanref{Formalism.\allowbreak{}MainNotions.\allowbreak{}ZetaZeroFreeRegion.\allowbreak{}Br.\allowbreak{}isInert\_\allowbreak{}iff}{Formalism.\allowbreak{}SharpRestriction.\allowbreak{}GlobalSchurBounds.\allowbreak{}RamificationTwoAndConjugationAtQAndD}{371--381}
\leanref{Formalism.\allowbreak{}MainNotions.\allowbreak{}Lnk.\allowbreak{}askIsF}{Formalism.\allowbreak{}SharpRestriction.\allowbreak{}GlobalSchurBounds.\allowbreak{}FrameCodeOfAdditive}{546--556}
Applying (5.20.7) gives (5.21.2) and (5.21.4) with
\(c_{\rm Loss}=3|K_s|+13\).
\end{proof}

\begin{lemma}[Places and node sizes]\label{B5.21a}
For a model \(D\) in its own frame, the archimedean indices map
bijectively onto \(\Phi\). In an admissible family take
\(x\ge\max(x_1+1,3)\) and a model at \(x\), with its translation
and thick labels. For \(0\le j<m\), give a position with a matching
archimedean index its weight at lift \(j+1\) and give every unmatched
position weight \(1\). Then \(w_j(\rho)\ge1\) for \(\rho\in\Phi\).
At a thick position \(b\le w_j(\rho)\le b+\eta_{\rm sz}+1\), and
at every other position of \(\Phi\), \(w_j(\rho)\le2b'\).
\end{lemma}
\leanref{Formalism.\allowbreak{}MainNotions.\allowbreak{}Lnk.\allowbreak{}image\_\allowbreak{}plMap}{Formalism.\allowbreak{}SharpRestriction.\allowbreak{}GlobalSchurBounds.\allowbreak{}QuadraticLogDerivCongr}{351--360}
\leanref{Formalism.\allowbreak{}MainNotions.\allowbreak{}Lnk.\allowbreak{}Mod.\allowbreak{}w\_\allowbreak{}off}{Formalism.\allowbreak{}SharpRestriction.\allowbreak{}GlobalSchurBounds.\allowbreak{}QuadraticLogDerivCongr}{416--422}
\leanref{Formalism.\allowbreak{}MainNotions.\allowbreak{}Lnk.\allowbreak{}Mod.\allowbreak{}w\_\allowbreak{}plMap}{Formalism.\allowbreak{}SharpRestriction.\allowbreak{}GlobalSchurBounds.\allowbreak{}QuadraticLogDerivCongr}{423--426}
\leanref{Formalism.\allowbreak{}MainNotions.\allowbreak{}Lnk.\allowbreak{}numN\_\allowbreak{}at}{Formalism.\allowbreak{}SharpRestriction.\allowbreak{}GlobalSchurBounds.\allowbreak{}QuadraticLogDerivCongr}{1261--1313}
\leanref{Formalism.\allowbreak{}MainNotions.\allowbreak{}Lnk.\allowbreak{}nodes\_\allowbreak{}at}{Formalism.\allowbreak{}SharpRestriction.\allowbreak{}GlobalSchurBounds.\allowbreak{}QuadraticLogDerivCongr}{1198--1249}
\leanref{Formalism.\allowbreak{}MainNotions.\allowbreak{}Lnk.\allowbreak{}nodesF}{Formalism.\allowbreak{}SharpRestriction.\allowbreak{}GlobalSchurBounds.\allowbreak{}QuadraticLogDerivCongr}{1314--1322}

\begin{proof}
Let \(e_\tau\) be the bijection with the host subgroup \(\Psi\).
The position \(T(e_\tau(\tau))\) of \ref{B2.02} has first sign bit
zero. Its signs and \(A\)-coordinate recover the host element.
The image in \(\Phi\) has size \(f=2n=|\Phi|\), so the unmatched
positions are \(G\setminus\Phi\).

Write \(\eta=\eta_{\rm sz}\). The size and parameter bounds give
\(\sqrt b\le w''\) and \(t_0\eta\le w''/x\), where the rank after
\(\ell\) lifts is \(k+\ell(t_0+1)\) and \(t_0\ge1\). Thus
\(1\le w''\le b'\) and \(\eta\le b'/3\). Put
\(v=N_\ast x^\varepsilon\mathcal L/k\), with
\(N_\ast=2nx^{-\varepsilon/8}\). At a marked label,
\(b'\le b_s\le(1+v)b'\) and \(4v\le x^{-2}\), so \(v\le1/36\).
Its weight lies in \([b_s+b'/x,b_s+b'/x+\eta]\), whose upper
endpoint is at most \((1+1/36+2/3)b'<2b'\).
At a thin position the interval is \([w'',w''+\eta]\subseteq[1,2b']\).
At a thick position it is \([b,b+\eta+1]\), with \(b\ge1\).
In the coordinates of \ref{B2.02}, the thick positions are
\(((0,1),s_2+u_0)\) for thick labels \(s_2\in A\).
\end{proof}

Write \(\mathcal T(x)\subset\N\) for the top labels. A pair
\(a_v,b_v:\N\to\Z\) is admissible for a base model \(i_0\) when
\[
 a\le a_v(t)\le a+\frac{a}{kx^{13}},\qquad
 b'\le b_v(t)\le b_{\rm top}(x)
 \quad(t\in\mathcal T(x)).
 \tag{5.22.1}
\]
For \(x>F_m.x_1\), the family endpoint evaluated at \(i_0\) is
\(b_{\rm top}(x)=b'+2n\delta_1x^\varepsilon\mathcal L b'/k\).
\leanref{Formalism.\allowbreak{}MainNotions.\allowbreak{}FamCore.\allowbreak{}bTop\_\allowbreak{}eq}{Formalism.\allowbreak{}SharpRestriction.\allowbreak{}FamCore}{124--125}
A model realizes \(a_v,b_v\) if its top weights agree with them.

\begin{lemma}[A single weight vector for the translations]\label{B5.22}
Assume the hypotheses of \ref{B5.21}. Assume also the common-weight
selection property \(\Lean{ChosenWS}\). Beyond a fixed threshold,
for every \(x\) in the window there exist a base model \(i_0\)
and integer functions \(a_v,b_v\) satisfying (5.22.1), such
that some model realizes them and every model realizing them has
\(R_W\le x^{-9}\).

Then there exist \(c_{\rm Loss}\ge0\) and \(x_1\) such that, for every
window point \(x\ge x_1\), one can choose the same admissible
functions \(a_v,b_v\), realized by some model, and every model
realizing them satisfies (5.21.2) and (5.21.4).
The common-weight exponent is \(e_{\rm Loss}=0\).
\end{lemma}
\leanref{Formalism.\allowbreak{}MainNotions.\allowbreak{}Lnk.\allowbreak{}lossF\_\allowbreak{}of\_\allowbreak{}7}{Formalism.\allowbreak{}SharpRestriction.\allowbreak{}GlobalSchurBounds.\allowbreak{}LossesSmallShares}{308--313}
\leanref{Formalism.\allowbreak{}MainNotions.\allowbreak{}ChosenWS}{Formalism.\allowbreak{}SharpRestriction.\allowbreak{}GlobalSchurBounds.\allowbreak{}WeightVectorSmallSharesSeven}{51--59}

\begin{proof}
Choose the weights by \(\Lean{ChosenWS}\) from \ref{B5.22b}.
Since each realization has \(R_W\le x^{-9}\), \ref{B5.21} gives
the bounds with \(e_{\rm Loss}=0\). Covering carries the weights
to every integer translation \(0\le u_0<n\).
\leanref{Formalism.\allowbreak{}MainNotions.\allowbreak{}FamCore.\allowbreak{}coverS\_\allowbreak{}carry}{Formalism.\allowbreak{}SharpRestriction.\allowbreak{}FamCore}{187--205}
\end{proof}

\begin{proposition}[Small share in the integer weight box]\label{B5.22a}
There is an \(\varepsilon_0>0\) such that for each
\(0<\varepsilon\le\varepsilon_0\) there is a size threshold \(\mathcal K\)
with the following property. For every \(F_m,CH\) with
\(\operatorname{Large}(F_m,\mathcal K,CH)\) there is an \(x_0\)
such that, for every \(x\ge x_0\) and \(i\in F_m.I_x\), one can choose
integer functions \(A,B\) on the model's group, zero off \(TT\), with
\[
 a\le A_t\le a+a/(kx^{13}),\qquad
 b'\le B_t\le b'+2n\delta_1x^\varepsilon\mathcal L b'/k
 \quad(t\in TT)
\]
and
\[
 R_W(A,B)=\sum_{p\in\mathcal P_c}\frac{\log p}{p}
 \left(2\mathbf1_{\{\Sigma_p(A,B)>x^{-\varepsilon}\}}
       +4\Sigma_p(A,B)\right)\le x^{-9}.
\]
Here \(\mathcal L\) and \(\Sigma_p\) are as in \ref{B3.51}, with
\(\Sigma_p\) evaluated at the trial weights \(A,B\).
\end{proposition}

\begin{proof}
Choose \(\varepsilon_0<1/52\) and the family thresholds of
\ref{B5.20b} and \ref{B5.20c}. Give the nonempty product of integer
intervals the uniform law. On good candidate orbits with \(p<p_0\),
\(\kappa_v\le N_vx^{-12}\) pointwise by \ref{B3.51} and
\(K.\mathrm{kII}\,r_G^2k\le x^{14}\) eventually. For \(p\ge p_0\),
\ref{B5.20c} gives \(\mathbb E\kappa_v\le C_mN_vN_1/x^{13}
\le N_vx^{-12}\) for large \(x\).
\leanref{GlobalTheory.\allowbreak{}EX.\allowbreak{}RangeOne.\allowbreak{}rangeOne\_\allowbreak{}abs}{Development.\allowbreak{}AbelianVarietyTheory.\allowbreak{}Proofs.\allowbreak{}GlobalInterpolation.\allowbreak{}TensorOfPairLift}{637--640}
\leanref{Formalism.\allowbreak{}MainNotions.\allowbreak{}G.\allowbreak{}Joints.\allowbreak{}type\_\allowbreak{}one\_\allowbreak{}hypothesis\_\allowbreak{}of\_\allowbreak{}segments}{Formalism.\allowbreak{}SharpRestriction.\allowbreak{}GrowthBounds.\allowbreak{}Joints}{1283--1325}
\leanref{Formalism.\allowbreak{}MainNotions.\allowbreak{}AuxBlock8.\allowbreak{}Pf.\allowbreak{}restHypI\_\allowbreak{}ev}{Formalism.\allowbreak{}SharpRestriction.\allowbreak{}GrowthBounds.\allowbreak{}LogarithmicSums.\allowbreak{}MeanNum}{416--660}

The map \(j\mapsto t_0+jh_p\) is injective on a period, so each
local pair has equally many extensions. Disjointness of good orbits gives
\[
 \mathbb E\Sigma_p=m\sum_{v\ {\rm good}}\mathbb E\kappa_v
 \le mn_Tx^{-12}.
\]
Put \(u=x^{\varepsilon/3}=\log X'\),
\(I_p=\mathbf1_{\{\Sigma_p>x^{-\varepsilon}\}}\) and
\(Z=\sum_{p\in\mathcal P_c}(\log p/p)(I_p+\Sigma_p)\).
Since \(I_p\le x^\varepsilon\Sigma_p\) and
\(\sum_{p\in\mathcal P_c}\log p/p\le u\sum_{1\le j<X'}1/j\le u(1+u)\),
\[
 4\mathbb E Z\le4(x^\varepsilon+1)mn_Tx^{-12}u(2+u)
 \le48x^{-19/2+(9+2/100+2/3)\varepsilon}\le x^{-9}.
\]
Here \(m\le x^{1/2+\varepsilon/100}\) and \(n_T\le2x^{2+8\varepsilon+\varepsilon/100}\). For \(\varepsilon<1/52\), the exponent is below \(-9\), so the last bound holds for large \(x\). Some pair has \(Z\le\mathbb EZ\), and
\(R_W\le4Z\) term by term.
\end{proof}
\leanref{Formalism.\allowbreak{}MainNotions.\allowbreak{}AuxBlock8.\allowbreak{}Pf.\allowbreak{}sum\_\allowbreak{}Sig\_\allowbreak{}le}{Formalism.\allowbreak{}SharpRestriction.\allowbreak{}GrowthBounds.\allowbreak{}LogarithmicSums.\allowbreak{}SlopeEstimateOfLemmaBSlopeLine}{329--359}
\leanref{Formalism.\allowbreak{}MainNotions.\allowbreak{}AuxBlock8.\allowbreak{}Pf.\allowbreak{}meanEx\_\allowbreak{}of}{Formalism.\allowbreak{}SharpRestriction.\allowbreak{}GrowthBounds.\allowbreak{}LogarithmicSums.\allowbreak{}SlopeEstimateOfLemmaBSlopeLine}{435--481}
\leanref{Formalism.\allowbreak{}MainNotions.\allowbreak{}AuxBlock8.\allowbreak{}Pf.\allowbreak{}weight\_\allowbreak{}condition\_\allowbreak{}of\_\allowbreak{}mean\_\allowbreak{}excess}{Formalism.\allowbreak{}SharpRestriction.\allowbreak{}GrowthBounds.\allowbreak{}LogarithmicSums.\allowbreak{}SlopeEstimateOfLemmaBSlopeLine}{483--501}
\leanref{GlobalTheory.\allowbreak{}EX.\allowbreak{}G4.\allowbreak{}wClosed}{Development.\allowbreak{}AbelianVarietyTheory.\allowbreak{}Proofs.\allowbreak{}GlobalInterpolation.\allowbreak{}WindowShareRequirementClosedFormProof}{40--42}

\begin{lemma}[The weights of a model in the constructed family]\label{B5.22b}
Suppose \(F_m\) satisfies \(\Lean{ChosenWall}\), so that for some
\(x_W\) one has \(R_W(i)\le x^{-9}\) for every \(x\ge x_W\)
and \(i\in F_m.I_x\). Then \(F_m\) satisfies \(\Lean{ChosenWS}\)
of \ref{B5.22}.
\end{lemma}

\begin{proof}
Take \(x\in F_m.\mathrm{Win}\) with \(x\ge\max(x_W,F_m.x_1+1)\).
Choose \(i_0\in F_m.I_x\) and transport its weights along the top
bijection to \(\mathcal T(x)\), with zero elsewhere. The pair
satisfies (5.22.1) and \(i_0\) realizes it. The hypothesis bounds
\(R_W\) at every realization. Proposition~\ref{B8.02a} supplies
\(\Lean{ChosenWall}\) on the constructed family.
\end{proof}
\leanref{Formalism.\allowbreak{}MainNotions.\allowbreak{}FamCore.\allowbreak{}coverU\_\allowbreak{}hyp\_\allowbreak{}self}{Formalism.\allowbreak{}SharpRestriction.\allowbreak{}FamCore}{207--212}
\leanref{Formalism.\allowbreak{}MainNotions.\allowbreak{}chosenWS\_\allowbreak{}of\_\allowbreak{}wall}{Formalism.\allowbreak{}SharpRestriction.\allowbreak{}GlobalSchurBounds.\allowbreak{}WeightVectorSmallSharesEight}{22--31}


% ---- p5c3 ----
\chapter{A slope bound for families of maps}\label{ch:p5c3}

For a fixed translation \(u_0\), the divided lattices \(N'_u\), \(u\in W(u_0)\), map into \(RR^\vee\). Let \(s\) be the rank-weighted mean of \(\mu(N'_u)-h_u\), where \(h_u\) is the period-map height bound from \ref{B5.10}. If \(H\ge0\) and every target slope is at least \(s-H\), we bound the number below \(s-\tau\) for \(0\le\tau\le H\) and integrate in \(\tau\).

\section{Slopes and height}

Let \(K\) be a totally complex number field with ring of integers \(\mathcal O_K\). Put
\[
 n_K=[K:\Q],\qquad
 c_K=-\log 2+\frac{\log|\operatorname{disc}K|}{n_K}.
\]
In a finite dimensional \(K\)-space \(V\), a lattice is a finitely generated \(\mathcal O_K\)-submodule. Write \(r(M)=\dim_K\operatorname{span}_K M\) for the rank of a submodule. At each complex place \(v\), choose one embedding \(K\hookrightarrow\C\) and equip the scalar extension \(V_v\) with a Hermitian norm. The diagonal image of a lattice \(L\) in \(\bigoplus_{v\mid\infty}V_v\), with squared norm \(\sum_v\|x_v\|_v^2\), has real rank \(\rho(L)=n_Kr(L)\).

For \(L\ne0\), put
\[
 \mu_{\mathrm{real}}(L)=-\frac{\log\operatorname{covol}(L)}{\rho(L)},
\]
with covolume measured in the real span of \(L\). The Harder--Narasimhan (HN) filtration uses induced metrics on submodules and orthogonal quotient metrics. The maximum real slope is taken over nonzero submodules and the minimum over nonzero quotients by saturated submodules. A factor of \(K\)-rank \(k\) contributes \(n_Kk\) entries to the slope list. Use the normalization
\[
 \mu(L)=2\mu_{\mathrm{real}}(L)+c_K,
\]
with the same affine change for the extreme slopes \(\mu_{\min}(L)\) and \(\mu_{\max}(L)\). At \(L=0\), the real slopes are zero by convention and the normalized slopes are \(c_K\).

For a nonzero \(K\)-linear map \(\phi:V\to W\), define
\[
 h(\phi)=\frac{1}{\#\{v\mid\infty\}}
 \sum_{v\mid\infty}\log\|\phi_v\|_{\mathrm{op}}^2.
\]
Here the operator norms use the specified Hermitian metrics. The height may be negative.

\begin{proposition}[The slope inequality for one map]\label{B5.23}
Let \(V,W\) be finite dimensional \(K\)-spaces with the preceding metrics. Let \(A\subset V\) be a finitely generated lattice, and let \(\phi:V\to W\) be \(K\)-linear with \(\phi(A)\ne0\). Then
\[
 \mu_{\min}(A)-h(\phi)\le\mu_{\min}(\phi(A)).
\]
For every finitely generated lattice \(B\subset W\) containing \(\phi(A)\),
\[
 \mu_{\min}(A)\le\mu_{\max}(B)+h(\phi).
\]
If the restriction of \(\phi\) to \(A\) is injective, then also
\[
 \mu_{\max}(A)\le\mu_{\max}(B)+h(\phi).
\]
\end{proposition}
\leanref{Formalism.\allowbreak{}HN.\allowbreak{}slope\_\allowbreak{}inequality\_\allowbreak{}mean}{Formalism.\allowbreak{}HarderNarasimhan.\allowbreak{}HarderNarasimhanProofs.\allowbreak{}SlopeIneqMean}{249--252}
\leanref{Formalism.\allowbreak{}HN.\allowbreak{}SlopeIneqMean}{Formalism.\allowbreak{}HarderNarasimhan.\allowbreak{}StatementPl}{148--160}

This is the totally complex integral case of Bost and Chen \cite[Propositions~1.2 and 2.10]{bost-chen-2013}, with doubled slopes and twice the archimedean height bound. Invariance under finite extension identifies the absolute slopes with those over \(K\). The finite-place operator norms of the integral maps are at most one.

\begin{proof}
Put \(a=h(\phi)/2\). Multiply the target frame at \(v\) by \(\|\phi_v\|_{\mathrm{op}}e^{-a}\), so the target norm is divided by this factor. The factors have mean logarithm zero and preserve the covolumes of target sublattices and their quotients. The new operator norms are all \(e^a\). A nonzero quotient of \(\phi(A)\) by a saturated sublattice is isomorphic to the corresponding quotient of \(A\), with operator norm at most \(e^a\) at each place. In complex rank \(j\) the absolute value of the complex determinant is at most \(e^{ja}\). The real covolume factor across the places is thus at most \(e^{n_Kja}\), so the real slope drops by at most \(a\). Taking the infimum over the quotients gives the image inequality after doubling. The image slope lies between its minimum slope and \(\mu_{\max}(B)\). In the injective case the same volume bound applies to every sublattice of \(A\), with image in \(B\).
\end{proof}

\begin{proposition}[Counting slopes with the HN filtration]\label{B5.24}
Let \(L\) be a lattice, and let
\(\lambda_1(L),\ldots,\lambda_{\rho(L)}(L)\) be the slopes of its HN factors in the real normalization. For \(\theta\in\R\), let \(L_{\ge\theta}\) be the saturated HN step containing the slopes at least \(\theta\). Then
\[
 \#\{j:\lambda_j(L)<\theta\}
 =\rho(L)-\rho(L_{\ge\theta}).
\]
If \(\rho(L)>0\), then for every \(a\in\R\),
\[
 \sum_{j=1}^{\rho(L)}\bigl(2\lambda_j(L)+a\bigr)
 =\rho(L)\bigl(2\mu_{\mathrm{real}}(L)+a\bigr).
\]
A nonzero HN step satisfies
\[
 \mu_{\min,\mathrm{real}}(L_{\ge\theta})\ge\theta.
\]
Moreover, every nonzero \(\mathcal O_K\)-sublattice \(M\subseteq L\) with
\(\mu_{\min,\mathrm{real}}(M)\ge\theta\) is contained in \(L_{\ge\theta}\).
\end{proposition}
\leanref{Formalism.\allowbreak{}HN.\allowbreak{}frameB\_\allowbreak{}count\_\allowbreak{}lt}{Formalism.\allowbreak{}HarderNarasimhan.\allowbreak{}HarderNarasimhanProofs.\allowbreak{}FrameHermSharpB}{62--76}
\leanref{Formalism.\allowbreak{}HN.\allowbreak{}frameB\_\allowbreak{}sum\_\allowbreak{}sl}{Formalism.\allowbreak{}HarderNarasimhan.\allowbreak{}HarderNarasimhanProofs.\allowbreak{}FrameHermSharpB}{115--125}
\leanref{Formalism.\allowbreak{}HN.\allowbreak{}frameB\_\allowbreak{}hnStep}{Formalism.\allowbreak{}HarderNarasimhan.\allowbreak{}HarderNarasimhanProofs.\allowbreak{}FrameHermSharpB}{77--79}
\leanref{Formalism.\allowbreak{}HN.\allowbreak{}polygon\_\allowbreak{}extreme\_\allowbreak{}slopes}{Formalism.\allowbreak{}HarderNarasimhan.\allowbreak{}HarderNarasimhanProofs.\allowbreak{}HNExtreme}{114--116}

This is the Hermitian lattice case of Gaudron's canonical filtration \cite[\S5.2, Proposition~5.13]{gaudron-2008}. For \(L\ne0\), the extremes of \(2\lambda_j(L)+c_K\) are \(\mu_{\min}(L)\) and \(\mu_{\max}(L)\).

\begin{lemma}[The share below a threshold]\label{B5.25}
For each \(i\) in a finite set \(I\), let \(x_{ij}\) be real, \(1\le j\le m_i\), with \(m_i\) a nonnegative integer. Put \(R=\sum_i m_i\). Suppose \(U,\tau\ge0\), \(s\in\R\), and
\[
 \sum_{i,j}x_{ij}\ge Rs,\qquad x_{ij}\le s+U.
\]
Define
\[
 z_U(\tau)=
 \begin{cases}
 \dfrac{U}{U+\tau},&U>0,\\[4pt]
 0,&U=0.
 \end{cases}
\]
Then
\[
 \sum_i\#\{j:x_{ij}<s-\tau\}\le Rz_U(\tau).
\]
\end{lemma}
\leanref{Formalism.\allowbreak{}HN.\allowbreak{}thmF\_\allowbreak{}share}{Formalism.\allowbreak{}HarderNarasimhan.\allowbreak{}HarderNarasimhanProofs.\allowbreak{}TheoremFSharp}{11--34}

\begin{proof}
Let \(q=\sum_i\#\{j:x_{ij}<s-\tau\}\). If \(U>0\), splitting the sum at the threshold gives
\[
 Rs\le\sum_{i,j}x_{ij}
 \le q(s-\tau)+(R-q)(s+U)
 =R(s+U)-q(U+\tau).
\]
Thus \(q(U+\tau)\le RU\).

If \(U=0\), the sum and entry bounds force every entry to equal \(s\).
\end{proof}

\section{The image of a family at one complex place}

Let \(V_i\), for \(i\) in a finite set \(I\), and \(Z\) be finite dimensional complex Hilbert spaces. For each \(i\), let \(\widetilde B_i:V_i\to Z\) be complex-linear and let \(c_i>0\). Put
\[
 d_i=\dim_\C V_i,\qquad R=\sum_i d_i,\qquad D_Z=\dim_\C Z,
\]
and define the positive operators
\[
 P_i=c_i\widetilde B_i^*\widetilde B_i,\qquad
 Q_Z=\sum_i c_i\widetilde B_i\widetilde B_i^*.
\]
For an operator \(T\), write
\(\|T\|_{\mathrm{HS}}^2=\operatorname{tr}(T^*T)\).
The two frame estimates are
\begin{align}
 \sum_i\|P_i-1\|_{\mathrm{HS}}^2&\le eR,
 &0\le e\le\tfrac14, \label{eq:p5c3-frame-source}\\
 D_Z\operatorname{tr}(Q_Z^2)&\le
 (1+\eta)\bigl(\operatorname{tr}Q_Z\bigr)^2,
 &\eta\ge0. \label{eq:p5c3-frame-target}
\end{align}

\begin{theorem}[Codimension of the images]\label{B5.26}
Suppose \(R>0\) and the frame estimates hold. Let \(V'_i\subseteq V_i\) be complex subspaces, and let \(z\in\R\) satisfy
\[
 \sum_i\bigl(d_i-\dim_\C V'_i\bigr)\le zR.
\]
Then
\[
 \operatorname{codim}_Z\sum_i\widetilde B_iV'_i
 \le(\eta+4z+4\sqrt e)D_Z.
\]
\end{theorem}
\leanref{Formalism.\allowbreak{}HN.\allowbreak{}frame\_\allowbreak{}linear\_\allowbreak{}algebra\_\allowbreak{}first}{Formalism.\allowbreak{}HarderNarasimhan.\allowbreak{}HarderNarasimhanProofs.\allowbreak{}FrameLinearAlgebraFirst}{77--222}

\begin{proof}
Put \(T=\operatorname{tr}Q_Z\) and \(a_i=\|P_i-1\|_{\mathrm{HS}}\). By \(|\operatorname{tr}(P_i-1)|\le\sqrt{d_i}\,a_i\) and Cauchy--Schwarz,
\[
 T=\sum_i\operatorname{tr}P_i
 \ge R-\sum_i\sqrt{d_i}\,a_i
 \ge R-\sqrt{R\sum_i a_i^2}
 \ge(1-\sqrt e)R\ge R/2>0.
\]
Thus \(D_Z>0\). Set \(W=(\sum_i\widetilde B_iV'_i)^\perp\) and \(w=\dim_\C W\). Let \(\Pi_W\) project orthogonally onto \(W\) and put \(T_W=\operatorname{tr}(\Pi_W Q_Z)\). The adjoint \(\widetilde B_i^*\) maps \(W\) into \((V'_i)^\perp\). For the orthogonal projection \(E_i\) onto \((V'_i)^\perp\), Bessel's inequality gives
\[
 \operatorname{tr}\bigl(\Pi_Wc_i\widetilde B_i\widetilde B_i^*\bigr)
 \le\operatorname{tr}(E_iP_iE_i).
\]
With \(k_i=\dim_\C(V'_i)^\perp\), cyclicity of trace and the Hilbert--Schmidt inequality give
\[
 \operatorname{tr}(E_iP_iE_i)
 =k_i+\operatorname{tr}\bigl(E_i(P_i-1)\bigr)
 \le k_i+\sqrt{k_i}\,a_i.
\]
Since \(\sum_i k_i\le zR\), Cauchy--Schwarz gives
\begin{equation}\label{eq:p5c3-trace-above}
 T_W\le \sum_i k_i+\sqrt{\sum_i k_i}\sqrt{\sum_i a_i^2}
 \le R\bigl(z+\sqrt{ze}\bigr).
\end{equation}

Set \(A=Q_Z-(T/D_Z)1\). By self-adjointness and \eqref{eq:p5c3-frame-target},
\[
 \|A\|_{\mathrm{HS}}^2
 =\operatorname{tr}(Q_Z^2)-\frac{T^2}{D_Z}
 \le\eta\frac{T^2}{D_Z}.
\]
The Hilbert--Schmidt pairing with \(\Pi_W\) gives
\[
 T_W\ge\frac{T}{D_Z}w-\sqrt w\,\|A\|_{\mathrm{HS}}.
\]
With \(q=w/D_Z\), divide by \(T\) and use \eqref{eq:p5c3-trace-above} to get
\[
 q-\sqrt{\eta q}\le\frac{T_W}{T}
 \le2\bigl(z+\sqrt{ze}\bigr),
\]
since \(R/T\le2\). For \(z\ge1/4\), use \(q\le1\). For \(z<1/4\), the bounds \(\sqrt{\eta q}\le(\eta+q)/2\) and \(\sqrt{ze}\le\sqrt e\) give \(q\le\eta+4z+4\sqrt e\).
\end{proof}

\begin{lemma}[From complex dimension to \(K\)-rank]\label{B5.27}
Let \(B_i:V_i\to W\) be \(K\)-linear maps between finite dimensional \(K\)-spaces, and choose an infinite place \(v\) of \(K\). Put
\[
 d_i=\dim_KV_i,\qquad D_C=\dim_KW,\qquad R=\sum_i d_i>0.
\]
Suppose that, after one common invertible change of target coordinates and an invertible change of coordinates in each source at \(v\), there are lifts
\[
 \widetilde B_i:\C^{d_i}\longrightarrow
 Z=\C^{D_C}\oplus\C^{D_0}
\]
whose projections onto the first summand represent \(B_{i,v}\).
Assume the frame estimates
\eqref{eq:p5c3-frame-source}--\eqref{eq:p5c3-frame-target},
with \(D_Z=D_C+D_0\).

For \(\mathcal O_K\)-submodules \(N_i^\tau\subseteq V_i\), suppose
\[
 \sum_i\bigl(d_i-r(N_i^\tau)\bigr)\le zR.
\]
Then
\[
 D_C-r\left(\sum_iB_iN_i^\tau\right)
 \le(\eta+4z+4\sqrt e)D_Z.
\]
\end{lemma}
\leanref{Formalism.\allowbreak{}HN.\allowbreak{}thmF\_\allowbreak{}place}{Formalism.\allowbreak{}HarderNarasimhan.\allowbreak{}HarderNarasimhanProofs.\allowbreak{}TheoremFSharp}{35--160}

\begin{proof}
Let \(V'_i\) be the complex span of \(N_i^\tau\) in the chosen coordinates, and set \(M=\sum_iB_iN_i^\tau\). Scalar extension gives \(\dim_\C V'_i=r(N_i^\tau)\) and \(\dim_\C M_v=r(M)\), where \(M_v\) denotes the complex span of \(M\) in the target coordinates. For \(S=\sum_i\widetilde B_iV'_i\), the projection \(\pi:Z\to\C^{D_C}\) has \(\pi(S)=M_v\) and kernel dimension \(D_0\). Thus \(\dim S\le r(M)+D_0\), whence \(D_C-r(M)\le D_Z-\dim S\). Apply \ref{B5.26}.
\end{proof}

\section{Integration of the slope count}

\begin{proposition}[A finite layer-cake inequality]\label{B5.28}
Let \(D\) be a positive integer, \(s\in\R\), \(H\ge0\), and
\(x_1,\ldots,x_D\in\R\). Suppose
\[
 x_j\ge s-H\qquad(1\le j\le D).
\]
Let \(w:\R\to\R\) satisfy
\[
 \#\{j:x_j<s-\tau\}\le Dw(\tau)
 \qquad(0\le\tau\le H),
\]
and suppose that \(\tau\mapsto\min(1,w(\tau))\) is integrable on \([0,H]\). Then
\[
 s-\int_0^H\min(1,w(\tau))\,d\tau
 \le\frac1D\sum_{j=1}^D x_j.
\]
\end{proposition}
\leanref{Formalism.\allowbreak{}CorankShareLayerCake.\allowbreak{}FR.\allowbreak{}layer\_\allowbreak{}cake\_\allowbreak{}formula}{Formalism.\allowbreak{}CorankShare.\allowbreak{}CorankShareProofs.\allowbreak{}LayerCakeFormula}{57--98}

The counting-measure case of Lieb and Loss \cite[Theorem~1.13]{lieb-loss-2001} gives
\[
 \sum_j(s-x_j)_+
 =\int_0^H\#\{j:x_j<s-\tau\}\,d\tau.
\]
Here \((s-x_j)_+=\max(s-x_j,0)\). Bound the count by \(D\min(1,w(\tau))\) and use \(s-x_j\le(s-x_j)_+\).

For \(U,H\ge0\), define
\[
 \varphi(U,H)=
 \begin{cases}
 U\log(1+H/U),&U>0,\\
 0,&U=0,
 \end{cases}
\]
and, for \(0\le\delta<1\), \(\eta,e\ge0\), put
\begin{equation}\label{eq:p5c3-frame-error}
 E_{\mathrm{fr}}(U,H;\delta,\eta,e)
 =\frac{4\varphi(U,H)+(\eta+4\sqrt e)H}{1-\delta}.
\end{equation}

\begin{theorem}[The strong family slope inequality]\label{B5.29}\label{B5.30}
Let \(I\) be finite and nonempty. For each \(i\in I\), let \(N_i\) be a full lattice in a finite dimensional \(K\)-space \(V_i\). Let \(C^\vee\) be a full lattice in a finite dimensional \(K\)-space \(W\). Give these spaces Hermitian norms at every infinite place.

Let \(B_i:V_i\to W\) be \(K\)-linear, and let \(h_i\in\R\). Suppose
\[
 B_iN_i\subseteq C^\vee,\qquad
 B_i\ne0\ \Longrightarrow\ h(B_i)\le h_i.
\]
Set \(d_i=\dim_KV_i\) and \(R=\sum_i d_i\), and assume \(R>0\).

At one infinite place, suppose the maps admit lifts to a common complex space \(Z\) as in \ref{B5.27}. Write \(D_C=\dim_KW\) and \(D_Z=\dim_\C Z\). Assume the two frame estimates with
\[
 0\le e\le\tfrac14,\qquad \eta\ge0,
\]
and suppose
\[
 0\le\delta<1,\qquad (1-\delta)D_Z\le D_C.
\]
Let \(s\in\R\), \(U,H\ge0\), and assume
\begin{align*}
 Rs&\le\sum_i d_i\bigl(\mu(N_i)-h_i\bigr),\\
 \mu_{\max}(N_i)-h_i&\le s+U\qquad(i\in I),\\
 s-\mu_{\min}(C^\vee)&\le H.
\end{align*}
Then
\[
 s\le\mu(C^\vee)+E_{\mathrm{fr}}(U,H;\delta,\eta,e).
\]
Under the same hypotheses, one also has
\[
 s\le\mu(C^\vee)+\log\#I+E_{\mathrm{fr}}(U,H;\delta,\eta,e).
\]
\end{theorem}
\leanref{Formalism.\allowbreak{}HN.\allowbreak{}theoremFSharp\_\allowbreak{}of}{Formalism.\allowbreak{}HarderNarasimhan.\allowbreak{}HarderNarasimhanProofs.\allowbreak{}TheoremFSharp}{168--387}
\leanref{Formalism.\allowbreak{}HN.\allowbreak{}TheoremFSharp}{Formalism.\allowbreak{}HarderNarasimhan.\allowbreak{}StatementFrameF}{64--100}
\leanref{Formalism.\allowbreak{}HN.\allowbreak{}frame\_\allowbreak{}theorem\_\allowbreak{}sharp}{Formalism.\allowbreak{}HarderNarasimhan.\allowbreak{}HarderNarasimhanProofs.\allowbreak{}TheoremFSharp}{388--389}
\leanref{Formalism.\allowbreak{}HN.\allowbreak{}theoremF\_\allowbreak{}of\_\allowbreak{}sharp}{Formalism.\allowbreak{}HarderNarasimhan.\allowbreak{}HarderNarasimhanProofs.\allowbreak{}TheoremFSharp}{390--399}
\leanref{Formalism.\allowbreak{}HN.\allowbreak{}frame\_\allowbreak{}theorem\_\allowbreak{}log\_\allowbreak{}ideal}{Formalism.\allowbreak{}HarderNarasimhan.\allowbreak{}HarderNarasimhanProofs.\allowbreak{}TheoremFSharp}{400--401}

\begin{proof}
The trace estimate in the proof of \ref{B5.26} gives
\(\operatorname{tr}Q_Z\ge R/2>0\). Thus \(D_Z>0\), and the dimension hypothesis implies \(D_C>0\).

For \(\tau\ge0\), put
\[
 \theta_i(\tau)=\frac{s-\tau+h_i-c_K}{2},
 \qquad N_i^\tau=(N_i)_{\ge\theta_i(\tau)}.
\]
Consider the adjusted slopes
\[
 x_{ij}=2\lambda_j(N_i)+c_K-h_i,
 \qquad 1\le j\le n_Kd_i.
\]
By \ref{B5.24},
\[
 \sum_{i,j}x_{ij}
 =n_K\sum_i d_i\bigl(\mu(N_i)-h_i\bigr)
 \ge n_KRs,
\]
where a rank-zero source has an empty list. Each \(x_{ij}\le\mu_{\max}(N_i)-h_i\le s+U\). By \ref{B5.25}, at most \(n_KRz_U(\tau)\) entries lie below \(s-\tau\). By \ref{B5.24}, their number is \(n_K\sum_i(d_i-r(N_i^\tau))\). Thus
\begin{equation}\label{eq:p5c3-source-corank}
 \sum_i\bigl(d_i-r(N_i^\tau)\bigr)\le Rz_U(\tau).
\end{equation}

Define the target step by
\[
 \theta_C(\tau)=\frac{s-\tau-c_K}{2},
 \qquad S^\tau=(C^\vee)_{\ge\theta_C(\tau)}.
\]
If \(B_iN_i^\tau\ne0\), then \ref{B5.24} gives \(\mu_{\min}(N_i^\tau)\ge s-\tau+h_i\). By \ref{B5.23},
\[
 \mu_{\min}(B_iN_i^\tau)
 \ge\mu_{\min}(N_i^\tau)-h(B_i)\ge s-\tau.
\]
Hence \ref{B5.24} gives \(B_iN_i^\tau\subseteq S^\tau\), an inclusion that also holds for zero images.

Put \(M^\tau=\sum_iB_iN_i^\tau\). By \ref{B5.27} and
\eqref{eq:p5c3-source-corank},
\[
 D_C-r(M^\tau)
 \le\bigl(\eta+4z_U(\tau)+4\sqrt e\bigr)D_Z.
\]
Since \(M^\tau\subseteq S^\tau\) and \(D_Z\le D_C/(1-\delta)\),
\begin{equation}\label{eq:p5c3-target-corank}
 D_C-r(S^\tau)\le D_Cw(\tau),
 \qquad
 w(\tau)=\frac{\eta+4z_U(\tau)+4\sqrt e}{1-\delta}.
\end{equation}

The adjusted target slopes \(y_j=2\lambda_j(C^\vee)+c_K\) for \(1\le j\le n_KD_C\) have mean \(\mu(C^\vee)\) by \ref{B5.24}. They are at least \(\mu_{\min}(C^\vee)\ge s-H\), and \eqref{eq:p5c3-target-corank} gives
\[
 \#\{j:y_j<s-\tau\}\le n_KD_Cw(\tau).
\]
The function \(w\) is continuous on \([0,H]\) for \(U>0\) and constant for \(U=0\). By \ref{B5.28},
\[
 s-\mu(C^\vee)\le\int_0^H\min(1,w(\tau))\,d\tau
 \le\int_0^H w(\tau)\,d\tau.
\]
The integral of \(z_U\) on this interval is \(U\log(1+H/U)\) for \(U>0\) and zero for \(U=0\), so the last integral is \(E_{\mathrm{fr}}(U,H;\delta,\eta,e)\). The second bound follows from \(\log\#I\ge0\).
\end{proof}

\section{The frame supplied by the period maps}

Fix \(C_0\ge1\) and \(C_H>0\). Take \(x>1\) with \(\sigma=x^{-C_H}\le1/20\).
For the fixed translation, write \(\mathcal W=W(u_0)\), and set
\[
 d_u=r(N'_u),\qquad R=\sum_{u\in\mathcal W}d_u.
\]
The target is \(C^\vee=RR^\vee\). For each word \(u\), let \(f_u=\Lean{duS u}\) be its unrestricted source dimension. At one complex place, suppose the lifted period matrix factors as
\[
 \widetilde B_u=A_uJ_u,\qquad
 A_u:\C^{f_u}\longrightarrow Z,\qquad
 J_u:\C^{d_u}\longrightarrow\C^{f_u},\qquad J_u^*J_u=1.
\]
For positive scalars \(c_u\), put
\[
 K_{uv}=\sqrt{c_uc_v}\,A_u^*A_v,\qquad
 \mathcal R=\sum_u f_u.
\]
Thus \(K_{uu}=c_uA_u^*A_u\).

The frame input consists of numbers \(D_*>0\), \(\eta_0,\eta_1\), \(e_0\ge0\), and \(0\le d_e\le1/4\), with
\begin{align}
 \|K_{uv}\|_{\mathrm{HS}}^2
 &\le(1+\eta_0)\frac{f_uf_v}{D_*}
 &&(u\ne v), \label{eq:p5c3-cross}\\
 \|K_{uu}-1\|_{\mathrm{HS}}^2
 &\le e_0f_u
 &&(u\in\mathcal W), \label{eq:p5c3-diagonal}\\
 D_Z&\le(1+\eta_1)D_*,
 &\sum_u(f_u-d_u)&\le d_e\mathcal R, \label{eq:p5c3-rank-loss}\\
 -1\le\eta_0&\le\sigma/2,
 &\eta_1&\le\sigma/2,\qquad \sqrt{e_0}\le\sigma/2.
 \label{eq:p5c3-small-input}
\end{align}
Fix \(e_B\ge2e_0\). In the period model, \(b=e^L\) with \(L=x^\varepsilon\). The chain length is \(m\) and the starting rank is \(k\). Let \(f\) count the complex places of the host CM field. Use the counting scale \(\Gamma=\exp(C_\Gamma m^5fL)\), where \(C_\Gamma>0\), and assume
\begin{equation}\label{eq:p5c3-code-input}
 C_0^{-1}\le\frac{f_u\Gamma^2}{D_*}\le C_0
 \quad(u\in\mathcal W),\qquad
 \#\mathcal W\ge C_0\Gamma^2x^{C_H}.
\end{equation}
For the selected family, \ref{B5.31a} gives the Gram and rank estimates, and \ref{B5.32a} gives the word count and the smallness bounds.

\begin{theorem}[Application to the divided period family]\label{B5.31}
Assume the frame and counting inputs
\eqref{eq:p5c3-cross}--\eqref{eq:p5c3-code-input}.
Let \(N'_u\) and \(RR^\vee\) be full lattices with \(d_u>0\) for every word. Assume the period maps satisfy
\[
 B_uN'_u\subseteq RR^\vee,\qquad
 B_u\ne0\ \Longrightarrow\ h(B_u)\le h_u.
\]
Assume that the projections of the matrices \(\widetilde B_u\) represent \(B_u\) in invertible coordinates at the chosen place as in \ref{B5.27}, and that
\[
 0\le\delta<1,\qquad (1-\delta)D_Z\le r(RR^\vee).
\]
In the period model, \(D_*=D^\star\) is the common main term defined before \ref{B5.11}. Of the \(D_Z=D_{\rm all}\) rows, \(r(RR^\vee)=D_{\rm rows}=(1-\delta)D_{\rm all}\) are retained. The estimates of \ref{B8.52}--\ref{B8.53} give the label and count hypotheses for the bound on this \(\delta\) in \ref{B7.52}--\ref{B7.53}.
Define
\[
 s=\frac{\sum_ud_u(\mu(N'_u)-h_u)}{R},\qquad
 U=\max_u\{\mu_{\max}(N'_u)-h_u\}-s.
\]
Then the frame estimates hold with
\[
 \eta=5\sigma+3\sqrt{d_e},\qquad
 e=2e_0\le e_B,\qquad 0\le e\le\tfrac14.
\]
For every \(H>0\) satisfying \(s-\mu_{\min}(RR^\vee)\le H\),
\[
 s\le\mu(RR^\vee)+\log\#\mathcal W+
 E_{\mathrm{fr}}(U,H;\delta,\eta,e).
\]
\end{theorem}
\leanref{Formalism.\allowbreak{}MainNotions.\allowbreak{}Lnk.\allowbreak{}frame\_\allowbreak{}at\_\allowbreak{}of}{Formalism.\allowbreak{}SharpRestriction.\allowbreak{}GlobalSchurBounds.\allowbreak{}TopDensityNGeSixteen}{538--601}
\leanref{Formalism.\allowbreak{}LFunctionChainParameters.\allowbreak{}MatrixLatticeEstimates.\allowbreak{}joint\_\allowbreak{}a\_\allowbreak{}ramified\_\allowbreak{}holds}{Formalism.\allowbreak{}SizeWindowBrackets.\allowbreak{}ChainSlopeWindows.\allowbreak{}JointDModelAppliesUnitLattice}{254--259}
\leanref{Formalism.\allowbreak{}LFunctionChainParameters.\allowbreak{}MatrixLatticeEstimates.\allowbreak{}ramification\_\allowbreak{}linear\_\allowbreak{}case\_\allowbreak{}holds}{Formalism.\allowbreak{}SizeWindowBrackets.\allowbreak{}ChainSlopeWindows.\allowbreak{}JointDModelAppliesUnitLattice}{164--186}
\leanref{Formalism.\allowbreak{}MainNotions.\allowbreak{}Lnk.\allowbreak{}FrameReal\_\allowbreak{}at}{Formalism.\allowbreak{}SharpRestriction.\allowbreak{}GlobalSchurBounds.\allowbreak{}TopDensityNGeSixteen}{519--532}
\leanref{Formalism.\allowbreak{}MainNotions.\allowbreak{}GamL}{Formalism.\allowbreak{}SharpRestriction.\allowbreak{}GlobalSchurBounds.\allowbreak{}Chosen}{14--15}

\begin{proof}
By the lower bounds in \eqref{eq:p5c3-code-input},
\begin{equation}\label{eq:p5c3-large-total}
 \mathcal R\ge
 \#\mathcal W\,\frac{D_*}{C_0\Gamma^2}
 \ge\frac{D_*}{\sigma}.
\end{equation}

Put \(a=\sqrt{e_0}\), \(t=\sqrt{d_e}\), and define
\[
 Q_f=\sum_uc_uA_uA_u^*,\qquad
 Q_Z=\sum_uc_uA_uJ_uJ_u^*A_u^*.
\]

For the source estimate, the isometry identity gives
\[
 c_u\widetilde B_u^*\widetilde B_u-1
 =J_u^*(K_{uu}-1)J_u.
\]
Compression by isometries decreases the Hilbert--Schmidt norm. Hence
\[
 \sum_u\|c_u\widetilde B_u^*\widetilde B_u-1\|_{\mathrm{HS}}^2
 \le e_0\mathcal R.
\]
The rank-loss bound gives \(R\ge(1-d_e)\mathcal R\), so \(e_0\mathcal R\le2e_0R\). This gives \eqref{eq:p5c3-frame-source} with \(e=2e_0\le\sigma^2/2\le1/800\).

For the target estimate, the diagonal input implies
\[
 \|K_{uu}\|_{\mathrm{HS}}\le(1+a)\sqrt{f_u}.
\]
The trace expansion
\[
 \operatorname{tr}(Q_f^2)=\sum_{u,v}\|K_{uv}\|_{\mathrm{HS}}^2
\]
and the cross bounds give
\[
 \operatorname{tr}(Q_f^2)
 \le(1+\eta_0)\frac{\mathcal R^2}{D_*}
 +(1+a)^2\mathcal R.
\]
Here \(1+\eta_0\ge0\). By \eqref{eq:p5c3-large-total},
\begin{equation}\label{eq:p5c3-compressed-square}
 \operatorname{tr}(Q_f^2)
 \le\frac{\mathcal R^2}{D_*}
 \bigl((1+\eta_0)+(1+a)^2\sigma\bigr).
\end{equation}
The restricted cross matrices are \(J_u^*K_{uv}J_v\). Applying Hilbert--Schmidt contraction to each pair gives
\[
 \operatorname{tr}(Q_Z^2)\le\operatorname{tr}(Q_f^2).
\]

For the trace lower bound, the diagonal input gives
\[
 \operatorname{tr}Q_f
 =\sum_u\operatorname{tr}K_{uu}\ge(1-a)\mathcal R.
\]
The matrix \(E_u=1-J_uJ_u^*\) is an orthogonal projection of rank \(f_u-d_u\). Therefore
\begin{align*}
 \operatorname{tr}Q_f-\operatorname{tr}Q_Z
 &=\sum_u\operatorname{tr}(K_{uu}E_u)\\
 &\le\sum_u(f_u-d_u)
 +a\sum_u\sqrt{f_u(f_u-d_u)}\\
 &\le\bigl(d_e+a\sqrt{d_e}\bigr)\mathcal R.
\end{align*}
The last step uses Cauchy--Schwarz and
\eqref{eq:p5c3-rank-loss}. Consequently,
\begin{equation}\label{eq:p5c3-compressed-trace}
 \operatorname{tr}Q_Z\ge\beta\mathcal R,\qquad
 \beta=1-a-t^2-ta.
\end{equation}
The ranges \(a\le\sigma/2\le1/40\) and \(0\le t\le1/2\) give
\(\beta\ge57/80>0\).

To bound the coefficient in \eqref{eq:p5c3-compressed-square} after multiplication by \(D_Z\), use
\[
 (1+\eta_1)\bigl((1+\eta_0)+(1+a)^2\sigma\bigr)
 \le(1+5\sigma)(1-a)^2.
\]
By \eqref{eq:p5c3-rank-loss}, \(1+\eta_1\ge0\). The left side is nondecreasing in \(\eta_0,\eta_1,a\), and the right side is nonincreasing in \(a\). It therefore suffices to set all three parameters equal to \(\sigma/2\). At these values the difference of the right and left sides is
\[
 \sigma\left(2-\frac{13}{2}\sigma+\frac12\sigma^2-\frac18\sigma^3\right)
 \ge0\qquad(0<\sigma\le1/20).
\]

Put \(P=1+5\sigma+3t\). Set \(B=1-a\) and \(v=t(t+a)\). Since \(P\le11/4\) and \(t+a\le21/40\),
\[
 2P(t+a)\le\frac{231}{80}<\frac{117}{40}\le3B.
\]
Hence
\[
 P(B-v)^2-(1+5\sigma)B^2
 =tB\bigl(3B-2P(t+a)\bigr)+Pv^2\ge0.
\]
Since \(D_Z\le(1+\eta_1)D_*\),
\[
 D_Z\operatorname{tr}(Q_Z^2)
 \le(1+5\sigma+3t)\beta^2\mathcal R^2
 \le(1+5\sigma+3t)(\operatorname{tr}Q_Z)^2.
\]
This is \eqref{eq:p5c3-frame-target} with
\(\eta=5\sigma+3t\).

The weighted mean hypothesis in \ref{B5.29} holds with equality. Since \(\mu(N'_u)\le\mu_{\max}(N'_u)\), we have \(U\ge0\), and its definition gives the upper bound for every adjusted maximum slope. Apply \ref{B5.30}.
\end{proof}

\begin{lemma}[Gram estimates in the arithmetic coordinates]\label{B5.31a}
For the selected family with the spectral rows of \ref{B8.02}, the per-word dimension bounds in \eqref{eq:p5c3-code-input}, the level conditions and the Bessel identification of \ref{B5.10}, one may take, uniformly for large \(x\),
\[
 \eta_0=b^{-6},\qquad e_0=2C_0/\Gamma+2e^{-kf\log k},\qquad
 \eta_1=e^{-c_7rf\log r},\qquad d_e=2f/b
\]
in \eqref{eq:p5c3-cross}--\eqref{eq:p5c3-rank-loss}. Here \(r=k+m-1\) and \(c_7>0\) is the total-count constant.
\end{lemma}

\begin{proof}
The full normalized Gram operator off the diagonal is bounded by \ref{B8.10b}. For a fixed word \(u\), its diagonal block is \(K_0+K''\), where \(K_0\) contains the identity cell \(I\) and \(K''\) sums the other terms with \(G(t)=0\), in the notation of \ref{B8.10b}.

For \(1\le j\le r_S=(m-2)/2\), let \(D_\delta>0\) and \(S_\delta\ge1\) be the determinant norm and minor height of a rank-\(j\) label \(\delta\). Multiplying the amplitude and multiplicity estimates in the type bound of \ref{B8.02} gives the label factor \(D_\delta^{(j+1)/2}S_\delta^{(C_p-C_s)k}\), with multiplicity exponent \(C_p\) and decay exponent \(C_s\). By \(D_\delta\le e^{C_1m^2fL}S_\delta^2\), the remaining height exponent is \((C_p-C_s)k+j+1\le-3r_S\), since \(C_s\ge C_p+4\) and \(k\ge r_S+3\). The coset bound \(\sum_\delta S_\delta^{-3r_S}\le e^{C_2m^2fLr_S}\) on finite sets of occurring labels thus gives
\[
 a_j=e^{Cm^6fL}(C/k)^{j(k-1-j)f}
\]
as a majorant for the sum over all rank-\(j\) labels, with \(C\ge1\) fixed. For ranks \((j,j')\) the pair bound is \(f_ua_ja_{j'}\). Sum over rank pairs and use the finite-type Gram factor of at most two to get
\[
 \|K''\|_{\mathrm{HS}}^2
 \le 2f_u\Bigl(\sum_{j=1}^{r_S}a_j\Bigr)^2
 \le f_u e^{-kf\log k}.
\]
For the last inequality, \(j(k-1-j)\ge k-2\) bounds the sum by \(k e^{Cm^6 fL}(C/k)^{(k-2)f}\). By \(C_\ell m^6L\le k\log k\) and \(C_\ell\ge8C\), the squared exponential is at most \(e^{kf\log k/4}\). The squared power contributes at most \(e^{-3kf\log k/2}\) for large \(k\), and \(2k^2\le e^{kf\log k/4}\).

For \(K_0-I\) use the pair partition in the proof of \ref{B8.10b}. Put \(t=m^5fL\) and \(\Delta_{uu}=f_u^2/D^\star\). The equal-word scalar and nongeneric corank-one coefficients are bounded by \(e^{C_Bt}\) and \(e^{C_4t}\), respectively. The loop coefficients from \ref{B8.10a} are at most one. The generic and nongeneric squared norms therefore sum to at most
\[
 P_g\bigl(e^{C_Bt}+3+e^{C_4t}\bigr)\Delta_{uu},
 \qquad P_g=\prod_{j=0}^1(1+3/g_j)^f .
\]
Use \(P_g\le4^f\) and \(10\cdot4^f\le e^t\) to bound the prefactor. Since \(C_\Gamma\ge1+\max(C_B,C_4)\), twice this sum is at most \(10P_g e^{\max(C_B,C_4)t}\Delta_{uu}\le\Gamma\Delta_{uu}\). Applying the squared triangle bound to the generic and nongeneric parts gives \(\|K_0-I\|_{\mathrm{HS}}^2\le\Gamma\Delta_{uu}\). Apply the same bound to \(K_0-I\) and \(K''\). Since \(\Gamma f_u/D^\star\le C_0/\Gamma\), the full diagonal defect is at most \(e_0 f_u\).

Extend the retained orthonormal target basis to the full target. The Bessel identification writes the divided map in invertible coordinates as the retained block of \(A_uJ_u\), where \(J_u^*J_u=1\). Write \(a_u\ge0\) for the unfolding constant. In the normalized kernel the factor \((a_ua_v)^{-1/2}\) takes the value zero when \(a_ua_v=0\). If \(a_u=0\), the squared diagonal defect is therefore \(\|I\|_{\mathrm{HS}}^2=f_u\). The lower dimension bound gives \(f_u>0\), so the defect bound \(e_0f_u<f_u\) for large \(x\) gives \(a_u>0\). Take \(c_u=a_u^{-1}\). The normalized blocks have matrices \(\sqrt{c_uc_v}A_u^*A_v\) up to entrywise conjugation, which preserves the squared Hilbert--Schmidt norm.

For the total level idempotent on the selected row realization, the count in \ref{B8.53} uses rank \(r=k+m-1\) and counted dimension \(D_{\rm all}=D_Z\). Its formal-degree main term is \(D^\star\) from \ref{B5.11}. The relative count gives \(D_Z\le(1+\eta_1)D^\star\).

In the level models of \ref{B8.37}, commensurability and the ideal twist preserve the ramified column rank by \ref{B5.04}. Thus \ref{B5.05a} gives \(d_u=f_u(1-1/(T_u-1))^f\). Since \(T_u\ge b^{m/2}\ge b\ge3\), Bernoulli's inequality gives \(f_u-d_u\le f_uf/(T_u-1)\le2ff_u/b\). Sum over \(u\).
\end{proof}
\leanref{Formalism.\allowbreak{}L2.\allowbreak{}SG.\allowbreak{}RowsPos.\allowbreak{}kernel\_\allowbreak{}diagonal\_\allowbreak{}positive}{Formalism.\allowbreak{}LocalToGlobal.\allowbreak{}SiegelLatticeTheory.\allowbreak{}FlagStabilizerTensors.\allowbreak{}RowsPos}{202--209}
\leanref{Formalism.\allowbreak{}MainNotions.\allowbreak{}S.\allowbreak{}linkS\_\allowbreak{}concl}{Formalism.\allowbreak{}SharpRestriction.\allowbreak{}CuspSpanNormBounds.\allowbreak{}StandingHypothesesHoldFamily}{186--205}
\leanref{Formalism.\allowbreak{}MainNotions.\allowbreak{}Lnk.\allowbreak{}matK\_\allowbreak{}at\_\allowbreak{}of}{Formalism.\allowbreak{}SharpRestriction.\allowbreak{}GlobalSchurBounds.\allowbreak{}TopDensityNGeSixteen}{1097--1108}
\leanref{Formalism.\allowbreak{}MainNotions.\allowbreak{}product\_\allowbreak{}first\_\allowbreak{}zeroth\_\allowbreak{}parts.\allowbreak{}rank\_\allowbreak{}nat}{Formalism.\allowbreak{}SharpRestriction.\allowbreak{}GlobalSchurBounds.\allowbreak{}FiniteModelOfAdmissibleTransports}{782--841}

\section{Parameters common to the family}

\begin{theorem}[Uniform frame error and the height interval]\label{B5.32}
Fix an admissible family and constants \(C_H>0\), \(C_0\ge1\). Suppose that, for all sufficiently large \(x\), every member has the full lattices and period-map data of \ref{B5.31}, with the same value \(\sigma=x^{-C_H}\). Assume that the counting inputs hold uniformly over the members.

Suppose, in addition, that for all sufficiently large \(x\) there are common numbers \(d_e=d_e(x)\) and \(e_B=e_B(x)\) such that
\[
 0\le d_e\le\tfrac14,\qquad 0\le e_B\le\tfrac14,\qquad
 \bigl(3\sqrt{d_e}+4\sqrt{e_B}\bigr)x^{C_H}\le1,
\]
and every member admits the frame input of \ref{B5.31} with these values of \(d_e\) and with \(2e_0\le e_B\).

Then there is a common threshold beyond which one may take
\[
 \bar\eta=5x^{-C_H}+3\sqrt{d_e},\qquad \bar e=e_B.
\]
These parameters satisfy
\[
 \bar\eta\ge0,\qquad 0\le\bar e\le\tfrac14,\qquad
 (\bar\eta+4\sqrt{\bar e})x^{C_H}\le6.
\]
For each member, the conclusion of \ref{B5.31} holds with parameters
\(0\le\eta\le\bar\eta\) and \(0\le e\le\bar e\).

In the period model with the division data of \ref{B5.10}, write \(RR_{u_0}\) for the full positive-rank lattice of translation \(u_0\). All translations have the same weight \(b=e^{x^\varepsilon}\). For each of them one may take \(H=x^{C_H}b\). In particular, for each \(u_0\),
\[
 s(u_0)\le\mu(RR_{u_0}^\vee)+\log\#W(u_0)
 +\frac{4\varphi(U(u_0),x^{C_H}b)+6b}{1-\delta(u_0)}.
\]
\end{theorem}
\leanref{Formalism.\allowbreak{}MainNotions.\allowbreak{}Lnk.\allowbreak{}frameF\_\allowbreak{}of}{Formalism.\allowbreak{}SharpRestriction.\allowbreak{}GlobalSchurBounds.\allowbreak{}TopDensityNGeSixteen}{606--647}
\leanref{Formalism.\allowbreak{}MainNotions.\allowbreak{}Lnk.\allowbreak{}Mod.\allowbreak{}data\_\allowbreak{}HHBound}{Formalism.\allowbreak{}SharpRestriction.\allowbreak{}GlobalSchurBounds.\allowbreak{}QuadraticLogDerivCongr}{714--729}
\leanref{Formalism.\allowbreak{}MainNotions.\allowbreak{}Lnk.\allowbreak{}lat\_\allowbreak{}at}{Formalism.\allowbreak{}SharpRestriction.\allowbreak{}GlobalSchurBounds.\allowbreak{}QuadraticLogDerivCongr}{1323--1333}
\leanref{Formalism.\allowbreak{}MainNotions.\allowbreak{}Lnk.\allowbreak{}Mod.\allowbreak{}data}{Formalism.\allowbreak{}SharpRestriction.\allowbreak{}GlobalSchurBounds.\allowbreak{}QuadraticLogDerivCongr}{464--505}

\begin{proof}
Take \(x\) above the common cutoffs and \(20^{1/C_H}\). By \ref{B5.31},
\[
 \eta=5x^{-C_H}+3\sqrt{d_e}=\bar\eta,\qquad
 e=2e_0\le e_B=\bar e.
\]
The common numerical hypothesis gives
\[
 (\bar\eta+4\sqrt{\bar e})x^{C_H}
 =5+\bigl(3\sqrt{d_e}+4\sqrt{e_B}\bigr)x^{C_H}\le6.
\]

For each translation in the period model, target duality and the height clause of \ref{B5.10} give \(s-\mu_{\min}(RR^\vee)=s+\mu_{\max}(RR)\le H=x^{C_H}b\). In \ref{B5.31} the term \((\eta+4\sqrt e)H\) is therefore at most \(6b\).
\end{proof}

\begin{corollary}[The code and the shared precision]\label{B5.32a}
Under the design hypotheses of \ref{B8.02}, for every sufficiently small \(\varepsilon>0\) and every prescribed threshold record, the construction supplies a family and an exponent \(C_H>0\) with \(C_0=3\) and \(\#W(u_0)=\lceil3\Gamma^2x^{C_H}\rceil\) at each selected point and translation. For the family of \ref{B5.31a}, the common upper bounds in \ref{B5.32} may be taken, for large \(x\), as
\[
 d_e(x)=E^2,\qquad e_B(x)=2E^2,\qquad E=\exp(-x^{\varepsilon/8}).
\]
\end{corollary}

\begin{proof}
For fixed \(\varepsilon\), take the supremum of the required threshold records. In the standard design, \(C_F\) is the word-count exponent. After fixing the outer list, set \(C_H=C_F=C_\Gamma\) above the middle thresholds. These include the height threshold of \ref{B5.10a}. Evaluate the last thresholds on these two lists.

At each point put \(Y=3\Gamma^2x^{C_F}\). The design inequalities are \(Y\le\#W(u_0)\le Y+1\). The rank and host bounds give \(m^5fL=O(x^{9/2+9\varepsilon})\) and \(k\ge x^{9/2+10\varepsilon}\). Hence \(\log Y=o(k)\), so \(M_W=\lceil Y\rceil\le 2^{k-3}\) eventually. Choose \(M_W\) distinct binary words of length \(k-2\). The word set at each translation is in bijection with \(\{0,\ldots,M_W-1\}\). The class-number formula of \ref{B8.54} and the rank identification used in \ref{B7.54} give \(1/3\le f_u\Gamma^2/D^\star\le3\). Multiplying the word count by the per-word lower bound gives \eqref{eq:p5c3-large-total}.

Put \(y=x^{\varepsilon/8}\). Here \(L=y^8\). Since \(k\ge8\log\Gamma\) eventually and \(r\ge k\), the errors \(\eta_0,\eta_1,\sqrt{e_0}\) in \ref{B5.31a} are eventually at most \(E\). In particular,
\[
 e_0\le(2C_0+2)e^{-C_\Gamma y^8}\le E^2.
\]
Also \(2f\le2x^5\le e^{L/2}\) and \(4y\le L\), so \(2f/b\le E^2\). For large \(x\ge18\) we have \(x^{C_H+1}\le e^y\), hence \(Ex^{C_H}\le1/18\). This gives \(E\le\sigma/2\) for \eqref{eq:p5c3-small-input}, and
\[
 (3\sqrt{d_e}+4\sqrt{e_B})x^{C_H}
 \le(3+4\sqrt2)/18<1.
\]
The bounds \(d_e,e_B\le1/4\) follow from \(E\le1/18\).
\end{proof}
\leanref{UpperTheory.\allowbreak{}Top.\allowbreak{}designS}{Development.\allowbreak{}AutomorphicArithmetic.\allowbreak{}CharacterBoxes.\allowbreak{}DesignS}{21--99}
\leanref{Formalism.\allowbreak{}MainNotions.\allowbreak{}Lnk.\allowbreak{}matRealF\_\allowbreak{}of}{Formalism.\allowbreak{}SharpRestriction.\allowbreak{}GlobalSchurBounds.\allowbreak{}TopDensityNGeSixteen}{750--806}


% ---- p5c4 ----
\chapter{Meeting the two sides; the eight error terms}\label{ch:p5c4}

At a window point the division estimates and period maps bound the same rank-weighted mean from opposite sides. The target slope cancels when the bounds are combined. The arithmetic loss estimate then bounds \(\ASK\) by \(\Main\) plus eight error terms.

\section{Common data and their realisation}

Let \(U_0\) be a finite set of translations. For each \(u_0\in U_0\), let \(W(u_0)\) be a finite nonempty set of words. Let \(TT\) be a finite set of top positions, \(n_T=|TT|\), and let \(m\geq0\) and \(f>0\) be integers. Layers are indexed by \(0\leq j<m\), and the \(f\) positions used in an archimedean average are indexed by \(\rho\). In the arithmetic realisation these positions are the elements of \(\Phi\), so \(f=|\Phi|\). Write \(w_{u_0j\rho}\) for the weight at position \(\rho\) in layer \(j\) and translation \(u_0\).

The common top weights are \(a_t\geq0\). The abstract numerical statements allow real \(\varepsilon\). The window family has \(0<\varepsilon<1/52\). Put
\[
 S_a=\sum_{t\in TT}a_t,\qquad Q=n_Tmb,\qquad
 \langle z_\rho\rangle_\rho=\frac1f\sum_{\rho=1}^f z_\rho .
\]
For a fixed translation, write \(M_u\) for the original lattice and \(N'_u\) for the divided lattice. Set
\[
 d_u=\operatorname{rk}N'_u,\qquad
 \mu_u=\mu(N'_u),\quad \mu_u^+=\mu_{\max}(N'_u),\qquad
 \mu_{M,u}=\mu(M_u),\quad \mu_{M,u}^+=\mu_{\max}(M_u).
\]
The target \(C^\vee=RR^\vee\) has mean slope \(\mu_C\) and minimum slope \(\mu_{C,\min}\). Let \(T_u>0\) be the numerical level of the word \(u\). In the arithmetic model it is the auxiliary prime of that word. Write \(h_u\) for the height allowance and \(\Ee_u\) for the total layer loss. These quantities may depend on \(u_0\). Define
\[
 \langle z_u\rangle_d=\frac{\sum_{u\in W(u_0)}d_uz_u}
 {\sum_{u\in W(u_0)}d_u},\qquad
 s=\langle\mu_u-h_u\rangle_d,\qquad
 U=\max_{u\in W(u_0)}(\mu_u^+-h_u)-s .
\]
The numerical arguments also apply when the slopes and \(d_u\) are real data satisfying the conditions below, with \(d_u>0\).
\leanref{Formalism.\allowbreak{}RankGapBounds.\allowbreak{}WeightedMeanData.\allowbreak{}Data.\allowbreak{}s}{Formalism.\allowbreak{}RankGapBounds.\allowbreak{}MeanBoundSecondForm}{92--95}
\leanref{Formalism.\allowbreak{}MainNotions.\allowbreak{}Lnk.\allowbreak{}Mod.\allowbreak{}data}{Formalism.\allowbreak{}SharpRestriction.\allowbreak{}GlobalSchurBounds.\allowbreak{}QuadraticLogDerivCongr}{464--505}

The real scales \(a,b,b',k,L,L_a,B_0,N_c\) are as in the preceding chapters. Let \(k>0\) throughout the numerical formulas. The real parameter \(X''\) is common to the translations. In the window model the group is \(G=E_0\times A\), where \(A\) is the translation group. Put \(n=|A|\). Then \(X''=a/w''\), where \(w''=b'/(xn)\), and \(\ell_g=\log(g_0g_1)\). Let
\[
 c_{\rm Div},e_{\rm Div},c_E,c_H,c_{\rm Count},
 c_{\rm Slope},c_1,c_{\rm Loss},e_{\rm Loss},c_{HH}
\]
be the constants in the estimates. The symbol \(c_1\) denotes the archive's \(c_{\rm Sl1}\). The abstract implications use real constants with the signs specified below. The concrete constants are nonnegative. Write \(\mathrm{IO}\) for the local-index total of \ref{B5.21} and set
\begin{align}
 \log\Omega&=\gamma+\log(\pi/v)-L_pL,\nonumber\\
 S_{\rm div}&=m\sum_{t\in TT}
 (a_t\log a_t-a_t-a_t\log\Omega),\label{eq:p5c4-Sdiv}\\
 E_{\rm div}&=c_{\rm Div}n_TaL_a+
 c_{\rm Div}m(B_0+N_c)x^{e_{\rm Div}\varepsilon},\nonumber\\
 P&=c_{\rm Loss}Qx^{e_{\rm Loss}\varepsilon},\label{eq:p5c4-P}\\
 \mathrm{Lo}_u&=c_1\mathrm{IO}
 +c_1n_Tm(a/k+b')\log T_u+c_1X''+1,\qquad
 J=\langle\mathrm{Lo}_u\rangle_d .\label{eq:p5c4-Lo}
\end{align}
Here \(v>0\) and \(\gamma,L_pL\in\R\). In the arithmetic model \(v=\operatorname{rd}(E)/2\), where \(\operatorname{rd}(E)\) is the root discriminant of the host field, and \(\gamma\) is Euler's constant. A top weight equal to zero contributes zero to \(S_{\rm div}\) and to \(\Main\).
\leanref{Formalism.\allowbreak{}RankGapBounds.\allowbreak{}WeightedMeanData.\allowbreak{}Data.\allowbreak{}SDiv}{Formalism.\allowbreak{}RankGapBounds.\allowbreak{}RankGapProofs.\allowbreak{}WeightedMeanData}{129--133}
\leanref{Formalism.\allowbreak{}RankGapBounds.\allowbreak{}WeightedMeanData.\allowbreak{}Data.\allowbreak{}LoP}{Formalism.\allowbreak{}RankGapBounds.\allowbreak{}MeanBoundSecondForm}{105--106}
\leanref{Formalism.\allowbreak{}LFunctionChainParameters.\allowbreak{}WeightedSlopeBounds.\allowbreak{}Lo5}{Formalism.\allowbreak{}SizeWindowBrackets.\allowbreak{}QuotientSlopeBoundsWith}{15--17}

For the uniform family inputs, the same constants work beyond one threshold for every \(x\) and every model \(i\in I_x\). The common frame bounds may be chosen after \(x\), before the model.

Fix a loss exponent \(e_L\geq0\). The common-weight loss hypothesis gives \(c_{\rm Loss}\geq0\) and a threshold \(x_L\). For every window point \(x\geq x_L\), it supplies a base model \(i_0\) and functions \(a_v,b_v:\N\to\Z\) admissible for \(i_0\) in the intervals of \ref{B5.22}. Some model \(i_1\) realises both functions on the top labels. Every model realising them satisfies the two loss bounds of \ref{B5.21} with \(c_{\rm Loss}Qx^{e_L\varepsilon}\) on the right.

Each of the following hypotheses involving words holds for every translation and every word at that translation. The counting threshold \(\lambda_5\) is defined in \ref{B5.34}.
\begin{enumerate}[label=\(\mathrm{H}\arabic*\),leftmargin=*]
\item\label{hyp:p5c4-setting}
The signs and order conditions are
\[
 x,b,k,v>0,\quad a,b'\geq0,\quad c_1\geq0,\quad
 a_t\geq0,\quad w_{u_0j\rho}>0,\quad d_u,T_u>0,\quad
 \mu_u\leq\mu_u^+.
\]
Moreover,
\[
 \bar\delta<1,\qquad \bar\eta\geq0,\qquad
 0\leq\bar e\leq\tfrac14 .
\]
\item\label{hyp:p5c4-div}
The division and height estimate is
\[
 \mu_C+h_u-\Fo\leq S_{\rm div}-S_a\log T_u+E_{\rm div}.
\]
\item\label{hyp:p5c4-E}
The total layer loss satisfies
\[
 |\Ee_u-mS_a\ell_g|\leq mc_EB_0 .
\]
\item\label{hyp:p5c4-H}
For any two words, the adjusted heights satisfy
\[
 |(h_u+S_a\log T_u)-(h_{u'}+S_a\log T_{u'})|
 \leq c_H(n_TaL_a+mB_0).
\]
\item\label{hyp:p5c4-count}
The numerical levels satisfy
\[
 \lambda_5(u_0)\leq\log T_u\leq c_{\rm Count}mL,\qquad
 |\log T_u-\log T_{u'}|\leq c_{\rm Count}\log k .
\]
\item\label{hyp:p5c4-M}
The original lattice slopes satisfy
\[
 -c_{\rm Slope}n_TaL_a(1+m/k)\leq\mu_{M,u},\qquad
 \mu_{M,u}^+\leq c_{\rm Slope}n_TaL_a .
\]
\item\label{hyp:p5c4-N}
The divided lattice satisfies
\begin{align*}
 \mu_u^+&\leq\mu_{M,u}^++\Ga+\Fo-\Ee_u+c_1mB_0,\\
 \mu_{M,u}+\Ga+\Fo-\Ee_u-\mathrm{Lo}_u&\leq\mu_u .
\end{align*}
\item\label{hyp:p5c4-loss}
The two arithmetic losses satisfy
\[
 \ASK(u_0)-\Ga(u_0)\leq P,\qquad \mathrm{IO}(u_0)\leq P .
\]
\item\label{hyp:p5c4-frame}
At each translation, let \(\delta,\eta,e,\mathcal H\) be the values in its data. If
\[
 0\leq\delta<1,\qquad \eta\geq0,\qquad
 0\leq e\leq\tfrac14,\qquad
 s-\mu_{C,\min}\leq\mathcal H,\quad \mathcal H>0,
\]
then the family slope estimate gives
\[
 s\leq\mu_C+\log|W(u_0)|
 +\frac{4\vartheta(U,\mathcal H)+(\eta+4\sqrt e)\mathcal H}
 {1-\delta}.
\]
Here, for \(H>0\),
\[
 \vartheta(u,H)=
 \begin{cases}
 u\log(1+H/u),&u>0,\\
 0,&u=0 .
 \end{cases}
\]
\item\label{hyp:p5c4-HH}
The height parameter obeys
\[
 s-\mu_{C,\min}\leq\mathcal H,\qquad
 0<\mathcal H\leq x^{c_{HH}}b .
\]
\item\label{hyp:p5c4-parameters}
At every translation the frame parameters satisfy
\[
 0\leq\delta\leq\bar\delta,\qquad
 0\leq\eta\leq\bar\eta,\qquad
 0\leq e\leq\bar e .
\]
\end{enumerate}
The quantities \(\Ga,\Fo,\mathrm{IO}\) are common to the words of a fixed translation. Hypotheses \(\mathrm{H}2\)--\(\mathrm{H}4\) come from \ref{B5.10}, \ref{B5.13} and \ref{B5.14}. Hypotheses \(\mathrm{H}5\)--\(\mathrm{H}8\) take the forms given in \ref{B5.11}, \ref{B5.12}, \ref{B5.06} and \ref{B5.21}.
\leanref{Formalism.\allowbreak{}RankGapBounds.\allowbreak{}WeightedMeanData.\allowbreak{}Data.\allowbreak{}HidaSetting}{Formalism.\allowbreak{}RankGapBounds.\allowbreak{}MeanBoundSecondForm}{111--117}
\leanref{Formalism.\allowbreak{}LFunctionChainParameters.\allowbreak{}WeightedSlopeBounds.\allowbreak{}PointwiseBoundVariant}{Formalism.\allowbreak{}SizeWindowBrackets.\allowbreak{}QuotientSlopeBoundsWith}{72--75}

\begin{theorem}[Common realisation at a window point]\label{B5.33}
Let \(\mathcal F\) be a window family as in \ref{B5.22}. Fix
\[
 C_{\rm abs},C_H,e_L\geq0,\qquad C_5\in\R,\qquad c_6>0.
\]
Assume the uniform lattice-slope and counting estimates of \ref{B5.06} and \ref{B5.10}--\ref{B5.14}. Assume also the uniform division estimate from these references with exponent \(C_{\rm abs}\). Use \(C_H\) in the height clause of \ref{B5.10}, as in Lemma~\ref{B5.10a} for the selected design. At each model assume \(d_u>0\) and \(\mu(N'_u)\leq\mu_{\max}(N'_u)\) for every word, and \(\mu_{\min}(RR^\vee)=-\mu_{\max}(RR)\) for its target.

Assume the family frame estimate with exponent \(C_H\) and the irreducibility bound. Assume also the common-weight loss estimate with exponent \(e_L\). In the selected design \ref{B5.32a} supplies the frame parameters. For the same target, \ref{B8.52}--\ref{B8.53} give the good-label injection and the relative cylinder counts on a row system satisfying \Lean{AllJoinA}. Green's local weight laws and the design bounds with size exponent \(C_s=6\) complete the hypotheses of \ref{B7.53}, which gives the required defect estimate.

For the arithmetic data assume the chosen top-set density, the size-window input, the chosen \(q,d\) bounds with constant \(C_5\), and the chosen thin-character estimate with constant \(c_6\). Lemmas~\ref{B5.21a} and~\ref{B5.33a} supply the node and character inputs for this family.

Then there are real constants \(K,c_{HH},e_K,C_L,x_0\). For each window point \(x\geq x_0\), one can choose a pair of integer weight functions \(a_v,b_v\) and, at each translation, a model realising that pair. One can take
\[
 c_{HH}=C_H,\qquad C_L=C_5,\qquad
 e_{\rm Div}=C_{\rm abs},\qquad e_{\rm Loss}=e_L .
\]
The frame, \(m,f,q,d\), and all global scales agree across the selected models. Their lattice and arithmetic quantities satisfy \(\mathrm{H}1\)--\(\mathrm{H}11\).

In particular, with \(b=e^{x^\varepsilon}\), \(a=xb\), the common top weights satisfy
\[
 a\leq a_t\leq2a,\qquad |a_t/b-x|\leq1,
\]
and the common summation limit may be chosen as
\[
 N=\lceil2a\rceil,\qquad
 \frac{a_t}{w_{u_0j\rho}}\leq N\qquad(t\in TT,\ 0\le j<m,\ \rho\in\Phi).
\]
The realisation gives the arithmetic identifications for \ref{B5.46}--\ref{B5.49} and the model-dependent windows of \ref{B5.45}, with \(1\leq k\leq x^{e_K}\). It gives the bounds \(x^{1/2-\varepsilon}\leq m\leq x^{1/2+\varepsilon/2}\) of \ref{B5.50}.

Let \(S_2\) be the common set of integer thick labels of the selected models, and let \(a_0\) be the frame's generator of \(A_0\). Put \(T_2=\{sa_0:s\in S_2\}\subseteq A_0\). The labels satisfy \(0\leq s<\operatorname{ord}(a_0)\), so \(n_2=|T_2|=|S_2|\). A position is thick at \(u_0\) if it is of the form \(((0,1),s_2+u_0)\) for \(s_2\in T_2\). With the model's common width \(\eta_{\rm sz}\geq0\), the weights satisfy
\[
 b\leq w_{u_0j\rho}\leq b+\eta_{\rm sz}+1
 \quad\text{at thick positions},\qquad
 0<w_{u_0j\rho}\leq2b'\quad\text{at the other positions of }\Phi .
\]
Put \(X'=\exp(x^{\varepsilon/3})\) and
\(\theta=\exp(-c_6x^{\varepsilon/8})\). With this width and \(\theta\), the replacement-cost bounds of Lemma~\ref{B5.49} are
\[
 2f(x+1)\theta\leq 1/x,\qquad
 6f(x+1)(\eta_{\rm sz}+1)\leq b/x .
\]
Also,
\[
 b'X'=xb,\qquad n_T\geq\tfrac12 nx^{-\varepsilon/8},\qquad
 n_2/n\geq x^{-\varepsilon},\qquad k\leq n^2x^{1/2-\varepsilon}.
\]
For the family of Lemma~\ref{B5.22} one may take \(e_L=0\).
\end{theorem}
\leanref{Formalism.\allowbreak{}MainNotions.\allowbreak{}Lnk.\allowbreak{}linkCore}{Formalism.\allowbreak{}SharpRestriction.\allowbreak{}GlobalSchurBounds.\allowbreak{}LossesSmallShares}{38--297}
\leanref{Formalism.\allowbreak{}MainNotions.\allowbreak{}linkGLOBAL7\_\allowbreak{}items}{Formalism.\allowbreak{}SharpRestriction.\allowbreak{}GlobalSchurBounds.\allowbreak{}LinkMin}{693--722}
\leanref{Formalism.\allowbreak{}MainNotions.\allowbreak{}ChainTied}{Formalism.\allowbreak{}SharpRestriction.\allowbreak{}GlobalSchurBounds.\allowbreak{}TieGlobal}{97--115}

\begin{proof}
Take \(e_K\geq0\) from the size-window input for \(c_6\). Let \(c_{\rm Count},c_{\rm Slope},c_1,c_{\rm Loss}\) be the nonnegative constants from the family estimates. The budgets multiplying the division constant \(C_d\) are nonnegative, so take
\[
 c_{\rm Div}=c_E=c_H=C_d^+=\max(C_d,0),\qquad
 K=C_d^++c_{\rm Count}+c_{\rm Slope}+c_1+c_{\rm Loss}+6 .
\]
If \(x_1\) is the family threshold, take \(x_0\geq\max(x_1+1,3)\) above the thresholds in the hypotheses and in \ref{B5.45}. Fix \(x\geq x_0\) in the window and choose \(i_0,a_v,b_v,i_1\) from the common-weight loss hypothesis. The covering property of \ref{B5.22} gives a model \(i(u_0)\) at each translation \(u_0\in A\) with the same pair of top-weight functions as \(i_1\). Take the frame and global scales from \(i_0\). Use its thick labels \(S_2\) as well.

At this \(x\), the assumed family frame estimate supplies common bounds \(\bar\eta,\bar e\) satisfying
\[
 \bar\eta\geq0,\qquad 0\leq\bar e\leq\tfrac14,\qquad
 (\bar\eta+4\sqrt{\bar e})x^{C_H}\leq6 .
\]
At each model \(0\leq\eta\leq\bar\eta\) and \(0\leq e\leq\bar e\). Irreducibility gives \(0\leq\delta\leq\bar\delta=\exp(-x^{\varepsilon_z/7})\), where \(\varepsilon_z=\varepsilon/J_{\mathcal F}\) and \(J_{\mathcal F}\geq1\) is the family's scale parameter. This gives \(\mathrm{H}11\), with \(\bar\delta\leq\exp(-1)<1/2\) since \(x\geq1\).

The height clause of \ref{B5.10} and target duality give
\[
 s-\mu_{\min}(RR^\vee)=s+\mu_{\max}(RR)\leq x^{C_H}b .
\]
Take \(\mathcal H=x^{C_H}b\) in \(\mathrm{H}10\). The slope inequality supplied by the frame estimate at this height gives \(\mathrm{H}9\). The signs and lattice hypotheses give \(\mathrm{H}1\). The division and counting estimates together with the lattice slope estimates give \(\mathrm{H}2\)--\(\mathrm{H}7\). For \(\mathrm{H}3\), sum \ref{B5.13} over the layers.

Admissibility and the size window give \(a\leq a_t\leq2a\) and \(|a_t/b-x|\leq1\). Lemma~\ref{B5.21a} gives \(w_{u_0j\rho}\geq1\) on \(\Phi\) and the two bounds with width \(\eta_{\rm sz}\). Thus
\[
 \frac{a_t}{w_{u_0j\rho}}\leq a_t\leq2a\leq\lceil2a\rceil.
\]
Read \(a_t\) from \(a_v\) at the natural-number label of \(t\), and use the layer \(j+1\) weights on \(\Phi\) and the extension by \(1\) on \(G\setminus\Phi\) from \ref{B5.21a}. The common frame, primes, modulus and Artin map identify the Frobenius and inertia data with the datum \(D_{\rm ar}\) of \(i_0\). With \(N=\lceil2a\rceil\) this gives
\[
 \ASK(u_0)=\sum_{t\in TT}\sum_{j=0}^{m-1}
 \ASK(D_{\rm ar},N,a_t,t,w_{u_0j}) .
\]
By \ref{B2.04} the partner \(\rho\) of each contributing \(g\notin\Phi\) lies in \(\Phi\), both for ordinary primes and for \(d\). A positive summand requires \(p^i<a_t/w_{u_0j\rho}\leq N\), so \(p,i\leq N\) for \(p\geq2\) and \(i\geq1\). The common-weight loss hypothesis applies to each \(i(u_0)\) with these weights and cutoff and gives \(\mathrm{H}8\).

The global identifications give
\[
 L=x^\varepsilon,\qquad a=xb,\qquad
 B_0=2nmbx^\varepsilon,\qquad
 N_c=\frac{n_2kbx^\varepsilon}{n(nx^{-\varepsilon})}.
\]
The chosen top-set density gives \(n_T\geq\tfrac12 nx^{-\varepsilon/8}>0\), hence \(n_T\geq1\). Since \(x\geq1\) and \(K\geq6\),
\[
 (\bar\eta+4\sqrt{\bar e})x^{C_H}b
 \leq6b\leq Kx^\varepsilon n_Ta .
\]
The definition \(b'=a/X'\) gives \(b'X'=xb\). The size-window input gives the replacement costs and the bounds on \(m,k,n_2/n\). It also gives the remaining rank and word-count bounds of \ref{B5.45}.

Apply Lemma~\ref{B5.33a} to the common arithmetic datum at \(i_0\) with the chosen thin-character and \(q,d\) estimates.
\end{proof}

\begin{lemma}[Character constants for a host]\label{B5.33a}
Let \(F_m\) be an admissible family of exponent \(\varepsilon\), and fix \(c_6,C_5\in\R\). For every \(x\ge0\) and every model at \(x\), take the arithmetic datum of its host in the same frame. There exists a function \(\ell:\widehat G\to\C\) with \(L_pL=\Re\ell(\chi_C)\), where \(\chi_C\) is the \(d\)-sign character. If \(\xi\ne\mathbf1\) and a primitive Dirichlet character \(\chi\) of positive modulus satisfies \(\chi(n)=\operatorname{primVal}(\xi,n)\) for every \(n\ge1\), then \(\ell(\xi)=L'(1,\chi)/L(1,\chi)\).

Put \(X'=\exp(x^{\varepsilon/3})\) and \(\theta=\exp(-c_6x^{\varepsilon/8})\). If the host thin-character bound holds, then \(|G_\xi(y)|\le y\theta\) for every \(\xi\) and \(y\ge X'/3\). If the host \(q,d\) bounds hold with constant \(C_5\), then \(|G_\xi(x)|\le x^{1/2+\varepsilon}\) and \(|\ell(\xi)|\le C_5\log x\) for the two sign characters and their product.
\end{lemma}
\leanref{Formalism.\allowbreak{}MainNotions.\allowbreak{}LnkBr.\allowbreak{}prim\_\allowbreak{}unique}{Formalism.\allowbreak{}SharpRestriction.\allowbreak{}GlobalSchurBounds.\allowbreak{}TopDensityNGeSixteen}{133--158}
\leanref{Formalism.\allowbreak{}MainNotions.\allowbreak{}Lnk.\allowbreak{}char\_\allowbreak{}at}{Formalism.\allowbreak{}SharpRestriction.\allowbreak{}GlobalSchurBounds.\allowbreak{}TopDensityNGeSixteen}{476--505}
\leanref{Formalism.\allowbreak{}MainNotions.\allowbreak{}Lnk.\allowbreak{}charF}{Formalism.\allowbreak{}SharpRestriction.\allowbreak{}GlobalSchurBounds.\allowbreak{}TopDensityNGeSixteen}{506--508}

\begin{proof}
Pull \(\xi\) back to the host by the change of coordinates \(T\) of \ref{B2.02}. Its model representative \(\chi_\xi\) and the primitive reduction of the Artin character take the same values at primes unramified in the host field \(E\), where both evaluate \(\xi\) on Frobenius. Let \(M\) be the product of their conductors and \(|\operatorname{disc} E|\). Each unit class modulo \(M>0\) contains a prime by Dirichlet's theorem \cite[Chap.~4]{davenport-1980}. This prime is unramified. Thus the two characters raised to level \(M\) are equal. Their common conductor is each original modulus by primitivity, so the representatives coincide. In particular \(\chi_\xi(n)=\operatorname{primVal}(\xi,n)\) for \(n\ge1\), by \ref{B2.15}.

Take \(\ell(\xi)\) to be the host constant \(\Lean{Lp}(\xi\circ T)\). For \(\xi\ne\mathbf1\) this is \(L'(1,\chi_\xi)/L(1,\chi_\xi)\). Any primitive \(\chi\) with the prescribed values agrees with \(\chi_\xi\) on every positive integer and hence coincides with it by the common-conductor argument. Lemma~\ref{B2.15} gives \(G_\xi=g_{\chi_\xi}\), with \(G_{\mathbf1}=G_z\). The \(d\)-sign cuts out \(\Q(\sqrt{-d})\), whose quadratic character gives the model constant \(L_pL=\Re\ell(\chi_C)\).

Each \(\chi_\xi\) is a primitive character of the host, so the thin-character hypothesis applies at every \(y\ge X'/3\). On positive integers coprime to the host modulus \(Q_Sqd\) of \ref{B2.01}, the Artin map identifies the two signs with the Kronecker characters \(({-q}/\cdot)\) and \(({-d}/\cdot)\). Their product gives the product sign on the same units. Each comparison character is defined on the units modulo \(8qd\), where passing to the primitive reduction preserves its values. This primitive reduction and the corresponding model representative agree on the units modulo \((Q_Sqd)(8qd)\), so primitive uniqueness identifies them. Thus the agreement holds for every positive integer coprime to \(8qd\). The chosen \(q,d\) bounds for these representatives give the two inequalities at \(x\).
\end{proof}

\section{The main term}

\begin{definition}[Counting threshold and geometric main term]\label{B5.34}
For each translation define
\begin{align}
 r&=c_{\rm Count}(L_a+m^2L_a/k),\label{eq:p5c4-r}\\
 \lambda_5(u_0)
 &=\sum_{j=0}^{m-1}
 \left\langle\log(w_{u_0j\rho}v/\pi)\right\rangle_\rho
 +m\ell_g-r,\label{eq:p5c4-lambda}\\
 \Main(u_0)
 &=\sum_{j=0}^{m-1}\sum_{t\in TT}
 a_t\left(\log a_t-1-\gamma+L_pL
 -\left\langle\log w_{u_0j\rho}\right\rangle_\rho\right).
 \label{eq:p5c4-main}
\end{align}
\end{definition}
\leanref{Formalism.\allowbreak{}LFunctionChainParameters.\allowbreak{}WeightedSlopeBounds.\allowbreak{}Lam5}{Formalism.\allowbreak{}SizeWindowBrackets.\allowbreak{}ChainSlopeWindows.\allowbreak{}LevelOfCongruence}{17--19}
\leanref{Formalism.\allowbreak{}RankGapBounds.\allowbreak{}WeightedMeanData.\allowbreak{}Data.\allowbreak{}Main}{Formalism.\allowbreak{}RankGapBounds.\allowbreak{}MeanBoundSecondForm}{107--110}

\begin{lemma}[Exact cancellation of the geometric main terms]\label{B5.35}
For \(v>0\), \(f>0\), and \(w_{u_0j\rho}>0\),
\[
 S_{\rm div}-S_a\lambda_5(u_0)+mS_a\ell_g
 =\Main(u_0)+S_ar .
\]
\end{lemma}
\leanref{Formalism.\allowbreak{}LFunctionChainParameters.\allowbreak{}WeightedSlopeBounds.\allowbreak{}key\_\allowbreak{}identity5}{Formalism.\allowbreak{}SizeWindowBrackets.\allowbreak{}ChainSlopeWindows.\allowbreak{}LevelOfCongruence}{105--111}

\begin{proof}
Fix \(u_0\) and put
\[
 Z=\sum_j\langle\log w_{u_0j\rho}\rangle_\rho,\qquad
 B=\sum_t a_t(\log a_t-1-\gamma+L_pL).
\]
Positivity of \(v,w_{u_0j\rho}\) gives
\[
 \sum_j\left\langle\log(w_{u_0j\rho}v/\pi)\right\rangle_\rho
 =Z+m(\log v-\log\pi).
\]
Since \(\log\Omega=\gamma+\log\pi-\log v-L_pL\),
\[
 S_{\rm div}=mB-mS_a(\log\pi-\log v),\qquad
 \Main(u_0)=mB-S_aZ.
\]
In \(S_{\rm div}-S_a\lambda_5+mS_a\ell_g\) the terms with \(\log\pi-\log v\) and \(\ell_g\) cancel by \eqref{eq:p5c4-lambda}. The remainder is \(mB-S_aZ+S_ar=\Main(u_0)+S_ar\). This holds with empty sums as well.
\end{proof}

For the earlier threshold, replace \(r\) in the identity by \(c_{\rm Count}(L+m^2L/k)+2\log\Gamma/(fk)\), where \(\log\Gamma\) is the auxiliary logarithm \Lean{logGam}.
\leanref{Formalism.\allowbreak{}RankGapBounds.\allowbreak{}WeightedMeanData.\allowbreak{}Data.\allowbreak{}key\_\allowbreak{}identity}{Formalism.\allowbreak{}RankGapBounds.\allowbreak{}RankGapProofs.\allowbreak{}WeightedMeanData}{170--217}

\section{The lower bound and the slope spread}

\begin{lemma}[Lower bound for the height-adjusted mean slope]\label{B5.36}
Assume \(\mathrm{H}1,\mathrm{H}2,\mathrm{H}3,\mathrm{H}5,\mathrm{H}6,\mathrm{H}7\). Put
\[
 K_5(u_0)=c_{\rm Slope}n_TaL_a(1+m/k)
 +m(S_a\ell_g+c_EB_0)
 +S_{\rm div}-S_a\lambda_5(u_0)+E_{\rm div}.
\]
Then, for every translation,
\[
 \Ga+\mu_C-K_5(u_0)-J\leq s .
\]
\end{lemma}
\leanref{Formalism.\allowbreak{}LFunctionChainParameters.\allowbreak{}WeightedSlopeBounds.\allowbreak{}s\_\allowbreak{}lower5}{Formalism.\allowbreak{}SizeWindowBrackets.\allowbreak{}ChainSlopeWindows.\allowbreak{}LevelOfCongruence}{50--66}

\begin{proof}
The second inequality of \(\mathrm{H}7\) gives
\[
 \mu_u-h_u\geq
 \mu_{M,u}+\Ga-\Ee_u+\Fo-h_u-\mathrm{Lo}_u .
\]
By \(\mathrm{H}2\),
\[
 \Fo-h_u\geq\mu_C-S_{\rm div}+S_a\log T_u-E_{\rm div}.
\]
As \(S_a\geq0\), the counting lower bound gives
\(S_a\log T_u\geq S_a\lambda_5\). The bounds in \(\mathrm{H}3\) and \(\mathrm{H}6\) give
\[
 \Ee_u\leq m(S_a\ell_g+c_EB_0),\qquad
 \mu_{M,u}\geq-c_{\rm Slope}n_TaL_a(1+m/k).
\]
Consequently
\[
 \mu_u-h_u\geq\Ga+\mu_C-K_5(u_0)-\mathrm{Lo}_u .
\]
Take the \(d\)-weighted mean.
\end{proof}

\begin{lemma}[Variation between two words]\label{B5.37}
Assume \(\mathrm{H}1,\mathrm{H}3,\mathrm{H}4,\mathrm{H}5\). Define
\[
 V=2mc_EB_0+c_H(n_TaL_a+mB_0)
 +S_ac_{\rm Count}\log k .
\]
For all \(u,u'\in W(u_0)\),
\[
 (-\Ee_u-h_u)-(-\Ee_{u'}-h_{u'})\leq V .
\]
\end{lemma}
\leanref{Formalism.\allowbreak{}LFunctionChainParameters.\allowbreak{}WeightedSlopeBounds.\allowbreak{}Vvar\_\allowbreak{}bound5}{Formalism.\allowbreak{}SizeWindowBrackets.\allowbreak{}ChainSlopeWindows.\allowbreak{}LevelOfCongruence}{68--79}

\begin{proof}
The two sides of the absolute-value bound in \(\mathrm{H}3\) give
\[
 \Ee_{u'}-\Ee_u\leq2mc_EB_0.
\]
Set \(H_u=h_u+S_a\log T_u\). Then
\[
 h_{u'}-h_u
 =H_{u'}-H_u+S_a(\log T_u-\log T_{u'}).
\]
Hypothesis \(\mathrm{H}4\) bounds the first difference by
\(c_H(n_TaL_a+mB_0)\). Since \(S_a\geq0\), the two-word estimate in \(\mathrm{H}5\) bounds the second term by \(S_ac_{\rm Count}\log k\). Adding the bounds gives \(V\).
\end{proof}

\begin{lemma}[Maximum slope above the weighted mean]\label{B5.38}
Assume \(\mathrm{H}1,\mathrm{H}3\)--\(\mathrm{H}7\). Put
\[
 K_3=c_{\rm Slope}n_TaL_a(2+m/k)+V+c_1mB_0 .
\]
Then
\[
 0\leq U\leq K_3+J .
\]
\end{lemma}
\leanref{Formalism.\allowbreak{}LFunctionChainParameters.\allowbreak{}WeightedSlopeBounds.\allowbreak{}u\_\allowbreak{}function\_\allowbreak{}le\_\allowbreak{}mean}{Formalism.\allowbreak{}SizeWindowBrackets.\allowbreak{}ChainSlopeWindows.\allowbreak{}LevelOfCongruence}{81--103}
\leanref{Formalism.\allowbreak{}RankGapBounds.\allowbreak{}WeightedMeanData.\allowbreak{}Data.\allowbreak{}uF\_\allowbreak{}nonneg}{Formalism.\allowbreak{}RankGapBounds.\allowbreak{}RankGapProofs.\allowbreak{}WeightedMeanData}{159--169}

\begin{proof}
Subtract the respective heights in the upper bound of \(\mathrm{H}7\) for \(u\) and the lower bound for \(u'\). The common \(\Ga+\Fo\) cancels in their difference. By \(\mathrm{H}6\) and Lemma~\ref{B5.37},
\begin{align*}
 (\mu_u^+-h_u)-(\mu_{u'}-h_{u'})
 &\leq \mu_{M,u}^+-\mu_{M,u'}\\
 &\quad+(-\Ee_u-h_u)-(-\Ee_{u'}-h_{u'})
 +c_1mB_0+\mathrm{Lo}_{u'}\\
 &\leq c_{\rm Slope}n_TaL_a(2+m/k)
 +V+c_1mB_0+\mathrm{Lo}_{u'}\\
 &=K_3+\mathrm{Lo}_{u'}.
\end{align*}
Average over \(u'\) with weights \(d_{u'}\) and take the maximum over \(u\) to obtain \(U\leq K_3+J\). By \(\mathrm{H}1\), every \(\mu_u-h_u\) is at most \(\max_{u'}(\mu_{u'}^+-h_{u'})\). The weighted mean is at most the same maximum, so \(U\geq0\).
\end{proof}

\section{The logarithmic correction and the upper bound}

\begin{lemma}[Uniform logarithmic correction]\label{B5.39}
Let \(u,H,X,b,U_*\) be real numbers satisfying
\[
 0\leq u\leq U_*,\qquad H>0,\qquad H\leq Xb,\qquad X,b>0 .
\]
With the zero convention in \(\vartheta\),
\[
 \vartheta(u,H)\leq (U_*+b)\log(1+X).
\]
\end{lemma}
\leanref{Formalism.\allowbreak{}RankGapBounds.\allowbreak{}WeightedMeanData.\allowbreak{}phi\_\allowbreak{}bound}{Formalism.\allowbreak{}RankGapBounds.\allowbreak{}RankGapProofs.\allowbreak{}WeightedMeanData}{81--118}

\begin{proof}
Put \(L_X=\log(1+X)>0\). The bound holds for \(u=0\). For \(u>0\), apply \(\log z\leq z-1\) with \(z=(1+H/u)/(1+X)\) and multiply by \(u\). This gives
\[
 \vartheta(u,H)
 \leq uL_X+\frac{H}{1+X}
 -u\left(1-\frac1{1+X}\right)
 \leq U_*L_X+\frac{Xb}{1+X}.
\]
The same inequality with \(z=(1+X)^{-1}\) gives \(X/(1+X)\leq L_X\).
\end{proof}

\begin{theorem}[Upper bound for the net gain]\label{B5.40}
Assume \(\mathrm{H}1\)--\(\mathrm{H}7\) and
\(\mathrm{H}9\)--\(\mathrm{H}11\). Write
\[
 X=x^{c_{HH}}>0,\qquad
 K_0=S_ar+E_{\rm div}+mc_EB_0
 +c_{\rm Slope}n_TaL_a(1+m/k).
\]
For every translation,
\begin{align*}
 \Ga(u_0)\leq{}&
 \Main(u_0)+\log|W(u_0)|+K_0+J\\
 &+\frac{4\log(1+X)(K_3+J+b)
 +(\bar\eta+4\sqrt{\bar e})Xb}
 {1-\bar\delta}.
\end{align*}
\end{theorem}
\leanref{Formalism.\allowbreak{}LFunctionChainParameters.\allowbreak{}WeightedSlopeBounds.\allowbreak{}Ga\_\allowbreak{}bound5}{Formalism.\allowbreak{}SizeWindowBrackets.\allowbreak{}ChainSlopeWindows.\allowbreak{}LevelOfCongruence}{113--161}

\begin{proof}
The assumptions on \(\delta,\eta,e,\mathcal H\) in \(\mathrm{H}9\) follow from \(\mathrm{H}1,\mathrm{H}10,\mathrm{H}11\). Thus Lemma~\ref{B5.36} and the family slope estimate give
\[
 \Ga+\mu_C-K_5-J
 \leq s
 \leq\mu_C+\log|W(u_0)|
 +\frac{4\vartheta(U,\mathcal H)
 +(\eta+4\sqrt e)\mathcal H}{1-\delta}.
\]
Subtract \(\mu_C\) and use \(K_5(u_0)=\Main(u_0)+K_0\) from Lemma~\ref{B5.35}. Lemma~\ref{B5.38} gives \(0\leq U\leq K_3+J\), and \(\mathrm{H}10\) gives \(0<\mathcal H\leq Xb\). Thus Lemma~\ref{B5.39} and the frame parameter bounds give
\[
 4\vartheta(U,\mathcal H)\leq4\log(1+X)(K_3+J+b),\qquad
 (\eta+4\sqrt e)\mathcal H\leq(\bar\eta+4\sqrt{\bar e})Xb.
\]
Both terms of the original numerator are nonnegative. Since \(0\leq\delta\leq\bar\delta<1\), the reciprocal \(1/(1-\delta)\) is at most \(1/(1-\bar\delta)\).
\end{proof}

\begin{lemma}[Uniform bound for the mean quotient loss]\label{B5.41}
Assume \(\mathrm{H}1,\mathrm{H}5\). Then
\[
 J\leq c_1\mathrm{IO}
 +c_1n_Tm(a/k+b')c_{\rm Count}mL+c_1X''+1 .
\]
If also \(\mathrm{H}8\) holds, define
\[
 M_{\rm Lo}=c_1P+
 c_1n_Tm(a/k+b')c_{\rm Count}mL+c_1X''+1 .
\]
Then \(J\leq M_{\rm Lo}\) at every translation.
\end{lemma}
\leanref{Formalism.\allowbreak{}LFunctionChainParameters.\allowbreak{}WeightedSlopeBounds.\allowbreak{}weighted\_\allowbreak{}mean\_\allowbreak{}lower\_\allowbreak{}bound\_\allowbreak{}le}{Formalism.\allowbreak{}SizeWindowBrackets.\allowbreak{}ChainSlopeWindows.\allowbreak{}LevelOfCongruence}{162--171}
\leanref{Formalism.\allowbreak{}LFunctionChainParameters.\allowbreak{}WeightedSlopeBounds.\allowbreak{}each\_\allowbreak{}unit\_\allowbreak{}bound\_\allowbreak{}statement}{Formalism.\allowbreak{}SizeWindowBrackets.\allowbreak{}ChainSlopeWindows.\allowbreak{}LevelOfCongruence}{188--191}

\begin{proof}
In \eqref{eq:p5c4-Lo} the coefficient of \(\log T_u\) is nonnegative by \(\mathrm{H}1\). Use \(\log T_u\leq c_{\rm Count}mL\) and average over \(u\). Since \(c_1\geq0\), substituting \(\mathrm{IO}\leq P\) gives the second inequality.
\end{proof}

\section{The eight error terms}

\begin{definition}[Error at one translation]\label{B5.42}
Put
\[
 \bar U=K_3+M_{\rm Lo},\qquad X=x^{c_{HH}}.
\]
The error at \(u_0\) is
\begin{align}
 E_8(u_0)={}&\!
 \underbrace{S_ac_{\rm Count}(L_a+m^2L_a/k)}_{\text{\rm 1}}
 \nonumber\\
 &+\underbrace{c_{\rm Div}n_TaL_a+
 c_{\rm Div}m(B_0+N_c)x^{e_{\rm Div}\varepsilon}}
 _{\text{\rm 2}}
 \nonumber\\
 &+\underbrace{mc_EB_0}_{\text{\rm 3}}
 +\underbrace{c_{\rm Slope}n_TaL_a(1+m/k)}_{\text{\rm 4}}
 \nonumber\\
 &+\underbrace{\log|W(u_0)|}_{\text{\rm 5}}
 +\underbrace{c_{\rm Loss}n_Tmbx^{e_{\rm Loss}\varepsilon}}
 _{\text{\rm 6}}
 \nonumber\\
 &+\underbrace{M_{\rm Lo}}_{\text{\rm 7}}
 +\underbrace{\frac{4\log(1+X)(\bar U+b)
 +(\bar\eta+4\sqrt{\bar e})Xb}
 {1-\bar\delta}}_{\text{\rm 8}} .
 \label{eq:p5c4-eight}
\end{align}
The expressions inside terms \(7\) and \(8\) are
\begin{align*}
 M_{\rm Lo}={}&c_1c_{\rm Loss}n_Tmbx^{e_{\rm Loss}\varepsilon}
 +c_1n_Tm(a/k+b')c_{\rm Count}mL+c_1X''+1,\\
 \bar U={}&c_{\rm Slope}n_TaL_a(2+m/k)
 +2mc_EB_0+c_H(n_TaL_a+mB_0)\\
 &+S_ac_{\rm Count}\log k+c_1mB_0+M_{\rm Lo}.
\end{align*}
The translation dependence is in \(\log|W(u_0)|\). In the concrete family each word set is in bijection with the same \(W(x)\), so this cardinality is common as well.
\end{definition}
\leanref{Formalism.\allowbreak{}LFunctionChainParameters.\allowbreak{}WeightedSlopeBounds.\allowbreak{}ErrX5at}{Formalism.\allowbreak{}SizeWindowBrackets.\allowbreak{}QuotientSlopeBoundsWith}{45--55}
\leanref{Formalism.\allowbreak{}LFunctionChainParameters.\allowbreak{}WeightedSlopeBounds.\allowbreak{}UBar5}{Formalism.\allowbreak{}SizeWindowBrackets.\allowbreak{}QuotientSlopeBoundsWith}{41--43}

Term \(1\) is the counting remainder in \ref{B5.35}. Terms \(2\) through \(4\) come from \(\mathrm{H}2,\mathrm{H}3,\mathrm{H}6\) in the lower bound \ref{B5.36}. Term \(5\) adds the nonnegative \(\log|W(u_0)|\) to the frame bound of \ref{B5.29}. Term \(6\) bounds \(\ASK-\Ga\), and term \(7\) bounds \(J\) by \ref{B5.41}.

The kernel in \ref{B5.29} has integral \(\int_0^{\mathcal H}U/(U+\tau)\,d\tau=\vartheta(U,\mathcal H)\), with value zero at \(U=0\). Terms \(7\) and \(8\) both contain \(M_{\rm Lo}\) because \(J\) enters the lower slope bound and the bound for \(U\) in \ref{B5.38}.

\begin{theorem}[The pointwise bound for \(\ASK\)]\label{B5.43}
For the finite sets specified above, assume \(\mathrm{H}1\)--\(\mathrm{H}11\). Then for every \(u_0\in U_0\),
\[
 \boxed{\ \ASK(u_0)\leq\Main(u_0)+E_8(u_0)\ } .
\]
In particular this holds for every translation of every sufficiently large window point in the realisation of Theorem~\ref{B5.33}.
\end{theorem}
\leanref{Formalism.\allowbreak{}LFunctionChainParameters.\allowbreak{}WeightedSlopeBounds.\allowbreak{}each\_\allowbreak{}unit\_\allowbreak{}bound\_\allowbreak{}statement}{Formalism.\allowbreak{}SizeWindowBrackets.\allowbreak{}ChainSlopeWindows.\allowbreak{}LevelOfCongruence}{175--197}
\leanref{Formalism.\allowbreak{}LFunctionChainParameters.\allowbreak{}WeightedSlopeBounds.\allowbreak{}PointwiseBoundVariant}{Formalism.\allowbreak{}SizeWindowBrackets.\allowbreak{}QuotientSlopeBoundsWith}{72--75}

\begin{proof}
The coefficient of \(J\) in \ref{B5.40} is \(1+4\log(1+X)/(1-\bar\delta)>0\). Substitute \(J\leq M_{\rm Lo}\) from \ref{B5.41} and use \(\mathrm{H}8\) to obtain
\begin{align*}
 \ASK-\Main
 &\leq P+\Ga-\Main\\
 &\leq P+\log|W(u_0)|+K_0+M_{\rm Lo}\\
 &\quad+\frac{4\log(1+X)(K_3+M_{\rm Lo}+b)
 +(\bar\eta+4\sqrt{\bar e})Xb}
 {1-\bar\delta}.
\end{align*}
Expanding \(K_0\) and using \(\bar U=K_3+M_{\rm Lo}\) gives \eqref{eq:p5c4-eight}. Theorem~\ref{B5.33} gives the window case.
\end{proof}

Under the identifications of Lemma~\ref{B5.46}, \(\Main(u_0)\) is the sum over top positions and layers of the real parts of the arithmetic main terms. Under the window hypotheses of Theorem~\ref{B5.45},
\[
 E_8(u_0)\leq x^{(e_{\rm Div}+e_{\rm Loss}+3)\varepsilon}
 \bigl[n_Ta(1+m^2/k)+m(B_0+N_c)\bigr].
\]


% ---- p5c5 ----
\chapter{Four terms and the square-root bound}\label{ch:p5c5}

For the chains of Alp\"oge's archive \cite{alpoge-2026},
Theorem~\ref{B5.43} gives \(\ASK(u_0)-\Main(u_0)\le E_8(u_0)\),
with \(E_8\) as in Definition~\ref{B5.42}. When every element of the
averaging subgroup has a geometric lift, Theorem~\ref{B5.47} transfers
a uniform bound on \(E_8\) to the mean arithmetic difference.

Write \(n_T\) for the number of top positions, \(m\) for the chain
length, \(k\) for the initial rank, and \(a\) for the common scale of the
top weights. Set
\begin{equation}\label{eq:four-bracket}
 \mathcal E=n_Ta\left(1+\frac{m^2}{k}\right)+m(B_0+N_c),
 \qquad C_{\rm err}=e_{\rm Div}+e_{\rm Loss}+3.
\end{equation}

\section{Combining the eight terms}

\begin{theorem}[A real-variable error bound]\label{B5.44}
Let all the parameters below be real. Suppose
\[
\begin{gathered}
 y,y_D,y_L\ge1,\qquad n_T,a\ge1,\qquad b,m\ge0,\qquad k>0,\\
 b\le a,\quad m\le k,\quad mb\le a,\quad
 0\le L,L_a\le2y,\\
 b'\ge0,\quad kb'\le a,\quad 0\le X''\le a,\quad
 0\le S_a\le2n_Ta,\quad B_0,N_c\ge0,\\
 \ell_W\le a,\quad 0\le\ell_k\le gy,\quad
 0\le\ell_H\le hy,\quad 0\le d_I\le2,\quad
 \tau\le Ky\,n_Ta.
\end{gathered}
\]
Assume that each of
\[
 c_{\rm Div},c_E,c_H,c_{\rm Count},c_{\rm Slope},c_1,c_{\rm Loss}
\]
lies in \([0,K]\), that \(h,g\ge1\), and that
\[
 200(1+K)^2hg\le y.
\]
Define
\[
\begin{aligned}
 P_0&=c_{\rm Loss}n_Tmb\,y_L,\\
 J_0&=c_1P_0+
 c_1n_Tm(a/k+b')\,c_{\rm Count}mL+c_1X''+1,\\
 V_0&=2mc_EB_0+c_H(n_TaL_a+mB_0)+S_ac_{\rm Count}\ell_k,\\
 U_0&=c_{\rm Slope}n_TaL_a(2+m/k)+V_0+c_1mB_0+J_0.
\end{aligned}
\]
Let
\[
\begin{aligned}
 R_1&=S_ac_{\rm Count}(L_a+m^2L_a/k),\\
 R_2&=c_{\rm Div}n_TaL_a+c_{\rm Div}m(B_0+N_c)y_D,\\
 R_3&=mc_EB_0,\qquad
 R_4=c_{\rm Slope}n_TaL_a(1+m/k),\\
 R_5&=\ell_W,\qquad R_6=P_0,\qquad R_7=J_0,\\
 R_8&=d_I\{4\ell_H(U_0+b)+\tau\}.
\end{aligned}
\]
Then
\[
 \sum_{i=1}^8R_i\le y_Dy_Ly^3
 \left\{n_Ta(1+m^2/k)+m(B_0+N_c)\right\}.
\]
\end{theorem}
\leanref{Formalism.\allowbreak{}LFunctionChainParameters.\allowbreak{}LErr.\allowbreak{}core\_\allowbreak{}ineq}{Formalism.\allowbreak{}SizeWindowBrackets.\allowbreak{}ChainSlopeWindows.\allowbreak{}SubThreeMulLogTwoLe}{13--208}

\begin{proof}
Put \(r=m^2/k\), \(A_1=n_Ta\), \(A_3=mB_0\),
\(A_4=mN_c\), and
\[
 E=A_1(1+r)+A_3+A_4,\qquad
 Y_0=y_Dy_Ly,\qquad P=Y_0E.
\]
The summands of \(E\) are nonnegative, with \(E\ge A_1\ge a\ge1\).
Since \(Y_0\ge y,y_D,y_L,1\), we have \(P\ge E\).

Since \(S_a(1+r)\le2E\) and \(L_a\le2Y_0\),
\[
 R_1\le4KP,\qquad R_2\le3KP,\qquad R_3\le KP.
\]
Using \(m/k\le1\) and \(mb\le a\) gives
\[
 R_4\le4KP,\qquad R_5\le a\le P,\qquad R_6\le KP.
\]

For the middle term of \(J_0\), use \(b'\le a/k\). Then
\[
 n_Tm^2(a/k+b')\le2n_Ta\,m^2/k=2A_1r\le2E,
\]
and consequently
\[
 c_1n_Tm(a/k+b')\,c_{\rm Count}mL\le4K^2P.
\]
The remainder \(c_1P_0+c_1X''+1\) is at most \((K^2+K+1)P\).
Thus
\begin{equation}\label{eq:J0-bound}
 J_0\le(5K^2+K+1)P.
\end{equation}
The slope term in \(U_0\) is at most \(6KP\). The first two
summands of \(V_0\) contribute at most \(5KP\) together, and
\(S_ac_{\rm Count}\ell_k\le2KgP\).
Also \(c_1mB_0\le KP\) and \(b\le P\).
Together with \eqref{eq:J0-bound}, these give
\[
 U_0+b\le uP,\qquad u=5K^2+13K+2Kg+2.
\]
Since \(\tau\le KyA_1\le KP\),
\[
 R_8\le d_I\{4hy\,uP+KP\}\le2\{4hy\,uP+KP\}.
\]
Using \(y\ge1\) gives
\[
 \sum_i R_i\le (5K^2+14K+2)P+2\{4hy\,uP+KP\}
 \le (5K^2+16K+2+8hu)yP.
\]
Since \(h,g\ge1\),
\[
\begin{aligned}
 5K^2+16K+2+8hu
 &\le hg(45K^2+136K+18)\\
 &\le200(1+K)^2hg\le y.
\end{aligned}
\]
We obtain \(\sum_iR_i\le y^2P=y_Dy_Ly^3E\).
\end{proof}

\begin{theorem}[The error bound in a chain window]\label{B5.45}
For the geometric data of Definition~\ref{B5.42}, take finite sets
\(U_0,\mathrm{Top}\). For each \(u_0\in U_0\), let \(W(u_0)\) be
finite and nonempty. Let \(m,f\in\N\), \(n_T=\#\mathrm{Top}\), and
\(S_a=\sum_{t\in\mathrm{Top}}a_t\).
Assume that the window data satisfy
\[
\begin{gathered}
 x\ge1,\quad 0<\varepsilon\le1,\quad b\ge1,\quad xb\le a,\quad n_T\ge1,\\
 m\le x^{1/2+\varepsilon/2},\quad xm\le k,\quad
 e_K\ge0,\quad 1\le k\le x^{e_K},\quad k\le a,\\
 0\le L,L_a\le2x^\varepsilon,\quad b'\ge0,\quad kb'\le a,\quad
 0\le X''\le a,\\
 0\le a_t\le2a\quad(t\in\mathrm{Top}),\qquad B_0,N_c\ge0,\\
 \log\#W(u_0)\le(k-3)\log2\quad(u_0\in U_0),\\
 \bar\delta\le\tfrac12,\qquad \bar\eta,\bar e\ge0.
\end{gathered}
\]
Suppose the seven coefficients of Theorem~\ref{B5.44} lie in \([0,K]\),
\(e_{\rm Div},e_{\rm Loss},c_{HH}\ge0\), and
\[
 (\bar\eta+4\sqrt{\bar e})x^{c_{HH}}b
 \le Kx^\varepsilon n_Ta.
\]
Assume also
\[
 \log x\le x^\varepsilon,\qquad
 200(1+K)^2(1+c_{HH})(1+e_K)\le x^\varepsilon.
\]
Then, for every \(u_0\in U_0\),
\[
 E_8(u_0)\le x^{C_{\rm err}\varepsilon}\mathcal E.
\]
\end{theorem}
\leanref{Formalism.\allowbreak{}LFunctionChainParameters.\allowbreak{}LChain.\allowbreak{}errX5at\_\allowbreak{}le}{Formalism.\allowbreak{}SizeWindowBrackets.\allowbreak{}ChainSlopeWindows.\allowbreak{}PChain}{613--673}

\begin{proof}
Apply Theorem~\ref{B5.44} with
\[
\begin{gathered}
 y=x^\varepsilon,\quad y_D=x^{e_{\rm Div}\varepsilon},\quad
 y_L=x^{e_{\rm Loss}\varepsilon},\quad
 g=1+e_K,\quad h=1+c_{HH},\\
 \ell_W=\log\#W(u_0),\quad \ell_k=\log k,\quad
 \ell_H=\log(1+x^{c_{HH}}),\\
 d_I=(1-\bar\delta)^{-1},\qquad
 \tau=(\bar\eta+4\sqrt{\bar e})x^{c_{HH}}b.
\end{gathered}
\]
The powers \(y,y_D,y_L\) are at least \(1\). Since
\(\varepsilon\le1\), we have \(m\le x\), so \(mb\le xb\le a\).
Also \(m\le xm\le k\) and \(1\le b\le xb\le a\).
The bounds on \(a_t\) give \(0\le S_a\le2n_Ta\). The logarithms satisfy
\[
\begin{aligned}
 \log\#W(u_0)&\le(k-3)\log2\le k\le a,\\
 0\le\log k&\le e_K\log x\le(1+e_K)y,\\
 0\le\log(1+x^{c_{HH}})&\le\log2+c_{HH}\log x\le(1+c_{HH})y.
\end{aligned}
\]
Here \(0\le\log2\le1\), \(x^{c_{HH}}\ge1\), and \(\log x\le y\).
The bound on \(\bar\delta\) gives \(0<d_I\le2\).
The terms \(R_1,\ldots,R_8\) sum to \(E_8(u_0)\) by
Definition~\ref{B5.42}, and \(y_Dy_Ly^3=x^{C_{\rm err}\varepsilon}\).
\end{proof}

\section{The arithmetic difference and its average}

Let \(H\) be a finite abelian group, written additively, and put
\[
 n=\#H,\qquad \Gamma=(\Z/2\Z)^2\times H,\qquad
 \Phi=\{((0,e),h):e\in\Z/2\Z,\ h\in H\}.
\]
Thus \(f=\#\Phi=2n\). Write
\[
 r_t(g)=((1,1),0)+g+((0,0),t),\quad
 e_1(t)=((1,0),t),\quad
 \chi_C((e_0,e_1),h)=(-1)^{e_0}.
\]
Characters of \(\Gamma\) take values in \(\C^\times\), and
\(\mathbf1\) denotes the trivial character.
For arithmetic data \(D\) and a function
\(\ell:\widehat\Gamma\to\C\), use the centred functions \(G_\psi\)
of Definition~\ref{B2.04}, with \(G_{\mathbf1}=G_z\).
Fix a cutoff \(N\in\N\), top weights \(a_t\in\R\) for \(t\in H\),
and local weights \(w_{s,j}:\Gamma\to\R\) for \(s\in H\) and \(0\le j<m\).
Write \(\ASK(D,N,A,t,w)\) for the sum \(\ASK_N(A,t,w)\) in
Definition~\ref{B2.04}, with \(D\) explicit.
For \(A\in\R\) and \(w:\Gamma\to\R\), the one-layer main term is
\begin{equation}\label{eq:arith-main}
 \mathrm{MAIN}(\ell,A,w)
 =A\left(\log A-1-\gamma+\ell(\chi_C)
 -\frac1f\sum_{\rho\in\Phi}\log w(\rho)\right).
\end{equation}
The geometric main term \(\Main(u_0)\) is real.
For the algebraic identities with arbitrary real weights, use the
archive's convention \(\log r=\log|r|\) for \(r\ne0\) and \(\log0=0\).

For the two comparisons, use the geometric data of
Definition~\ref{B5.34}, a map \(\mathrm{tr}:U_0\to H\), a set
\(TT\subseteq H\), and bijections
\[
 e_T:\mathrm{Top}\longrightarrow TT,\qquad
 e_P:\{0,\ldots,f-1\}\longrightarrow\Phi.
\]
For the models of Theorem~\ref{B5.33}, choose \(e_P\) using \(f=|\Phi|\).
The geometric array indexes the weights \(w\) via \(e_P\).
By Lemma~\ref{B5.21a}, summing their logarithms over \(\Phi\)
is the same as summing over the model's archimedean places.
Choose \(\ell\) by Lemma~\ref{B5.33a} at the base model \(i_0\) of
Theorem~\ref{B5.33}. The same \(\ell\) serves every translation, and
\(L_pL=\Re\ell(\chi_C)\).
\leanref{Formalism.\allowbreak{}MainNotions.\allowbreak{}Lnk.\allowbreak{}Mod.\allowbreak{}sum\_\allowbreak{}places}{Formalism.\allowbreak{}SharpRestriction.\allowbreak{}GlobalSchurBounds.\allowbreak{}QuadraticLogDerivCongr}{458--463}
\leanref{Formalism.\allowbreak{}MainNotions.\allowbreak{}Lnk.\allowbreak{}Mod.\allowbreak{}ident}{Formalism.\allowbreak{}SharpRestriction.\allowbreak{}GlobalSchurBounds.\allowbreak{}QuadraticLogDerivCongr}{792--798}
For general data, assume the weight and constant identities
\[
\begin{aligned}
 a^{\rm geom}_\tau&=a_{e_T(\tau)},&
 w^{\rm geom}_{u_0j}(\rho)&=w_{\mathrm{tr}(u_0),j}(e_P\rho),\\
 \gamma_{\rm geom}&=\gamma,& L_pL&=\Re\ell(\chi_C)
\end{aligned}
\]
for all geometric indices in their respective ranges. For every
\(u_0\in U_0\), assume also
\[
 \ASK_{\rm geom}(u_0)=
 \sum_{t\in TT}\sum_{j=0}^{m-1}
 \ASK(D,N,a_t,t,w_{\mathrm{tr}(u_0),j}).
\]

\begin{lemma}[Identification of the main term]\label{B5.46}
Under these identifications, every \(u_0\in U_0\) satisfies
\[
 \Main(u_0)=
 \sum_{t\in TT}\sum_{j=0}^{m-1}
 \Re\mathrm{MAIN}(\ell,a_t,w_{\mathrm{tr}(u_0),j}).
\]
\end{lemma}
\leanref{Deduction.\allowbreak{}ch\_\allowbreak{}top\_\allowbreak{}main\_\allowbreak{}eq}{Development.\allowbreak{}AutomorphicArithmetic.\allowbreak{}DedekindLFunctions.\allowbreak{}LandauTheorems.\allowbreak{}ChTop}{11--36}

\begin{proof}
By Definition~\ref{B5.34}, the geometric main term is
\[
 \sum_j\sum_{\tau\in\mathrm{Top}}a^{\rm geom}_\tau
 \left(\log a^{\rm geom}_\tau-1-\gamma_{\rm geom}+L_pL
 -\frac1f\sum_{\rho=0}^{f-1}
 \log w^{\rm geom}_{u_0j}(\rho)\right).
\]
Reindex the logarithm sum by \(e_P\) and the top sum by \(e_T\).
With \(\gamma_{\rm geom}=\gamma\) and \(L_pL=\Re\ell(\chi_C)\),
each summand is the real part of \eqref{eq:arith-main}.
Exchange the two outer sums.
\end{proof}

Let \(T_2\subseteq H\), and put \(n_T=|TT|\) and \(n_2=|T_2|\).
For a real top-weight scale \(a_E\), define
\[
\begin{aligned}
 B_*(x)&=2nmbx^\varepsilon,&
 N_*(x)&=\frac{n_2kbx^\varepsilon}{n(nx^{-\varepsilon})},\\
 \mathcal B_x(a_E)&=
 n_Ta_E(1+m^2/k)+m\{B_*(x)+N_*(x)\}.
\end{aligned}
\]
For \(u_0\in H\), set
\[
 \mathrm{gap}(u_0)=
 \sum_{t\in TT}\sum_{j=0}^{m-1}
 \{\ASK(D,N,a_t,t,w_{u_0j})
 -\Re\mathrm{MAIN}(\ell,a_t,w_{u_0j})\}.
\]

\begin{theorem}[The geometric difference is the arithmetic difference]
\label{B5.47}
Assume the identifications stated before Lemma~\ref{B5.46}.
Suppose the eleven hypotheses of Theorem~\ref{B5.43}, the window
hypotheses of Theorem~\ref{B5.45}, and
\[
 B_0\le B_*(x),\qquad N_c\le N_*(x)
\]
hold for the same geometric data. Then every \(u_0\in U_0\) satisfies
\[
 \mathrm{gap}(\mathrm{tr}(u_0))
 =\ASK_{\rm geom}(u_0)-\Main(u_0)
 \le x^{C_{\rm err}\varepsilon}\mathcal B_x(a).
\]
\end{theorem}
\leanref{Deduction.\allowbreak{}ch\_\allowbreak{}top}{Development.\allowbreak{}AutomorphicArithmetic.\allowbreak{}DedekindLFunctions.\allowbreak{}LandauTheorems.\allowbreak{}ChTop}{37--73}

\begin{proof}
The assumed identity for \(\ASK_{\rm geom}\) and
Lemma~\ref{B5.46} give the equality. By
Theorems~\ref{B5.43} and~\ref{B5.45}, the geometric difference is at most
\(E_8(u_0)\le x^{C_{\rm err}\varepsilon}\mathcal E\).
The top-position bijection gives \(\#\mathrm{Top}=n_T\).
Since \(m\ge0\), the bounds on \(B_0,N_c\) give
\(\mathcal E\le\mathcal B_x(a)\).
\end{proof}

\begin{theorem}[The one-layer Fourier identity]\label{B5.48}
Let \(D\) have Frobenius elements \(F_p\in\Gamma\) and inertia subgroups
\(I_p\le\Gamma\). Its distinguished integers are \(q_1,d\).
Assume the local conditions of Lemma~\ref{B2.03} for a finite set
\(S\subseteq\N\). Thus \(d\) is prime, \(2,q_1\in S\), and
\(d\notin S\). The inertia and Frobenius conditions are
\[
\begin{gathered}
 I_p=0\quad\text{for primes }p\notin S\cup\{d\},\\
 I_p\subseteq\Phi,\quad F_p\in\Phi\quad(p\in S),\qquad
 I_d=\{0,((1,0),0)\}.
\end{gathered}
\]
Let \(\ell:\widehat\Gamma\to\C\), \(A>0\), \(t\in H\),
\(N\in\N\), and \(w:\Gamma\to\R\).
If
\[
 w(\rho)>0,\qquad A/w(\rho)\le N\qquad(\rho\in\Phi),
\]
then
\[
 \ASK(D,N,A,t,w)-\mathrm{MAIN}(\ell,A,w)
 =\frac1f\sum_{g\notin\Phi}\sum_{\psi\in\widehat\Gamma}
 \overline{\psi(g)}\,w(r_t(g))
 G_\psi\!\left(\frac A{w(r_t(g))}\right)
\]
as an equality in \(\C\).
\end{theorem}
\leanref{Deduction.\allowbreak{}ch\_\allowbreak{}balance}{Development.\allowbreak{}AutomorphicArithmetic.\allowbreak{}DedekindLFunctions.\allowbreak{}LandauTheorems.\allowbreak{}ChBalance}{200--261}

\begin{proof}
Apply Theorem~\ref{B2.07} with \(a=A\) and expand \(\mathrm{De}\)
by Definition~\ref{B2.04}.
\end{proof}

For \(A_0\le H\), write
\[
 \Psi_A=\{\psi\in\widehat\Gamma:
 \psi((0,0),s)=1\ \text{for every }s\in A_0\}.
\]
A position \(\rho\) is thick at \(u_0\) if
\(\rho=((0,1),s+u_0)\) for some \(s\in T_2\).
Put \(Q=n_Tmb\) and \(\delta_2=n_2/n\).

\begin{lemma}[Averaging over a subgroup]\label{B5.49}
Let the host conditions of Theorem~\ref{B5.48} hold, and let
\(TT,T_2\subseteq H\), \(T_2\subseteq A_0\le H\), \(m,N\in\N\).
Suppose \(x\ge3\), \(b,b',X'>0\), \(\eta,\theta\ge0\), and
\(b'X'\le xb\). Assume that for every \(t\in TT\),
\[
 |a_t/b-x|\le1.
\]
For every \(u_0\in A_0\) and \(0\le j<m\), assume
\[
\begin{gathered}
 w_{u_0j}(\rho)>0\quad(\rho\in\Phi),\\
 a_t/w_{u_0j}(\rho)\le N\quad((t,\rho)\in TT\times\Phi),\\
 b\le w_{u_0j}(\rho)\le b+\eta+1\quad\text{at thick positions},\\
 w_{u_0j}(\rho)\le2b'\quad\text{at the other positions of }\Phi.
\end{gathered}
\]
Suppose the thin-scale estimate and the replacement costs satisfy
\[
\begin{gathered}
 |G_\psi(y)|\le y\theta
 \quad(\psi\in\widehat\Gamma,\ y\ge X'/3),\\
 2f(x+1)\theta\le1/x,\qquad
 6f(x+1)(\eta+1)\le b/x.
\end{gathered}
\]
Let \(E,Z\in\R\). If
\[
 \mathrm{gap}(u_0)\le E\quad(u_0\in A_0)
\]
and
\[
 n_TZ\le
 \sum_{t\in TT}\Re\sum_{\psi\in\Psi_A}
 \psi(e_1(t))G_\psi(a_t/b),
\]
then
\[
 Q\frac{\delta_2}{2}Z\le E+\frac{2Q}{x}.
\]
\end{lemma}
\leanref{Deduction.\allowbreak{}CharacterBalanceMoves.\allowbreak{}core}{Development.\allowbreak{}AutomorphicArithmetic.\allowbreak{}DedekindLFunctions.\allowbreak{}LandauTheorems.\allowbreak{}ChCore}{283--339}
\leanref{Deduction.\allowbreak{}CharacterBalanceMoves.\allowbreak{}th\_\allowbreak{}le}{Development.\allowbreak{}AutomorphicArithmetic.\allowbreak{}DedekindLFunctions.\allowbreak{}LandauTheorems.\allowbreak{}ChCore}{75--153}

\begin{proof}
For \(t\in TT\), we have \(a_t\ge(x-1)b>0\).
The proof of Theorem~\ref{B2.10} supplies the weight Lipschitz estimate
for every \(D,\ell\). With \(a=a_t\), the replacement error in
Lemma~\ref{B2.09} is at most \(2b/x\), since
\[
 2f\{3(x+1)(\eta+1)+(x+1)b\theta\}\le b/x+b/x.
\]
Apply Theorem~\ref{B2.10} with \(\mathscr E=E\) and \(Y_*=Z\).
\end{proof}

\section{The four normalized terms}

\begin{lemma}[Both ends of the chain-length window]\label{B5.50}
Let \(m\in\N\) be the length of a chain at a scale \(x>1\),
\(0<\varepsilon<1/52\). Suppose its size condition is
\[
 x^{-\varepsilon/100}\sqrt{x\,x^{-\varepsilon/8}}
 \le m\le
 x^{\varepsilon/100}\sqrt{x\,x^{-\varepsilon/8}}.
\]
Then
\[
 x^{1/2-\varepsilon}\le m\le x^{1/2+\varepsilon/2}.
\]
\end{lemma}
\leanref{Formalism.\allowbreak{}MainNotions.\allowbreak{}Lnk.\allowbreak{}winA}{Formalism.\allowbreak{}SharpRestriction.\allowbreak{}GlobalSchurBounds.\allowbreak{}QuadraticLogDerivCongr}{1414--1416}
\leanref{Formalism.\allowbreak{}MainNotions.\allowbreak{}Lnk.\allowbreak{}winL}{Formalism.\allowbreak{}SharpRestriction.\allowbreak{}GlobalSchurBounds.\allowbreak{}QuadraticLogDerivCongr}{1417--1419}

\begin{proof}
The endpoints are \(x^{1/2-\varepsilon/16-\varepsilon/100}\) and
\(x^{1/2-\varepsilon/16+\varepsilon/100}\).
Their exponents lie in \([1/2-\varepsilon,1/2+\varepsilon/2]\),
so \(x>1\) gives the claimed bounds.
\end{proof}

\begin{lemma}[The top-weight term]\label{B5.51}
Let \(x,X,Y,a_E,n_T,m,b\in\R\), with \(n_T,m,b\ge0\).
If \(x\le Xm\), \(X\le Y\), and \(a_E\le2xb\), then
\[
 n_Ta_E\le2Y(n_Tmb).
\]
In particular, for \(x\ge1\), \(\varepsilon>0\), and
\(m\ge x^{1/2-\varepsilon}\), one may take
\(X=x^{1/2+\varepsilon}\) and
\(Y=x^{1/2+13\varepsilon/8}\).
\end{lemma}
\leanref{Formalism.\allowbreak{}LFunctionChainParameters.\allowbreak{}ThickSetMeanBounds.\allowbreak{}SZ.\allowbreak{}T1}{Formalism.\allowbreak{}SizeWindowBrackets.\allowbreak{}ChainSlopeWindows.\allowbreak{}LipschitzWeightAndMinimum}{31--40}

\begin{proof}
The assumptions give \(x\le Xm\le Ym\), so
\[
 n_Ta_E\le2n_Txb\le2n_TYmb.
\]
For the specified powers,
\(x=x^{1/2+\varepsilon}x^{1/2-\varepsilon}\le Xm\)
and \(X\le Y\).
When \(Q>0\), the intermediate bound reads
\(n_Ta_E/Q\le2x/m\).
\end{proof}

\begin{lemma}[The inverse-rank term]\label{B5.52}
Let \(x>0\), \(Y\ge1\), \(n_T,m,b\ge0\), \(k>0\), and
\(a_E\in\R\). If \(xm\le k\) and \(a_E\le2xb\), then
\[
 n_Ta_E\frac{m^2}{k}\le2Y(n_Tmb).
\]
\end{lemma}
\leanref{Formalism.\allowbreak{}LFunctionChainParameters.\allowbreak{}ThickSetMeanBounds.\allowbreak{}SZ.\allowbreak{}T2}{Formalism.\allowbreak{}SizeWindowBrackets.\allowbreak{}ChainSlopeWindows.\allowbreak{}LipschitzWeightAndMinimum}{41--52}

\begin{proof}
As \(m\ge0\) and \(xm\le k\), one has \(m^2/k\le m/x\).
Therefore
\[
 n_Ta_E\frac{m^2}{k}
 \le n_T(2xb)\frac m x=2Q\le2YQ.
\]
When \(Q>0\), the sharper normalized estimate is
\(n_Ta_Em^2/(kQ)\le2xm/k\le2\).
\end{proof}

\begin{lemma}[The layer-position term]\label{B5.53}
Let \(x>0\), \(\varepsilon\in\R\), \(m,b,n\ge0\), and
\(n_T\in\R\). Set \(Y=x^{1/2+13\varepsilon/8}\).
If
\[
 m\le x^{1/2+\varepsilon/2},\qquad
 \tfrac12 nx^{-\varepsilon/8}\le n_T,
\]
then
\[
 m(2nmbx^\varepsilon)\le4Y(n_Tmb).
\]
\end{lemma}
\leanref{Formalism.\allowbreak{}LFunctionChainParameters.\allowbreak{}ThickSetMeanBounds.\allowbreak{}SZ.\allowbreak{}T3}{Formalism.\allowbreak{}SizeWindowBrackets.\allowbreak{}ChainSlopeWindows.\allowbreak{}LipschitzWeightAndMinimum}{53--76}

\begin{proof}
The density assumption gives \(n_T\ge0\) and
\(n\le2n_Tx^{\varepsilon/8}\). Thus
\[
\begin{aligned}
 2nm x^\varepsilon
 &\le4n_Tm x^{\varepsilon+\varepsilon/8}\\
 &\le4n_Tx^{1/2+\varepsilon/2+\varepsilon+\varepsilon/8}
 =4n_TY.
\end{aligned}
\]
Multiply by \(mb\ge0\).
\end{proof}

\begin{lemma}[The thick-set rank term]\label{B5.54}
Let \(x\ge1\), \(\varepsilon>0\), \(m,b\ge0\), \(n>0\), \(k\ge0\),
and \(n_T,n_2\in\R\).
Set \(Y=x^{1/2+13\varepsilon/8}\). If
\[
 \tfrac12 nx^{-\varepsilon/8}\le n_T,\qquad
 n_2\le n,\qquad k\le n^2x^{1/2-\varepsilon},
\]
then
\[
 m\,\frac{n_2kbx^\varepsilon}{n(nx^{-\varepsilon})}
 \le2Y(n_Tmb).
\]
\end{lemma}
\leanref{Formalism.\allowbreak{}LFunctionChainParameters.\allowbreak{}ThickSetMeanBounds.\allowbreak{}SZ.\allowbreak{}T4}{Formalism.\allowbreak{}SizeWindowBrackets.\allowbreak{}ChainSlopeWindows.\allowbreak{}LipschitzWeightAndMinimum}{77--114}

\begin{proof}
Since \(x^\varepsilon x^{-\varepsilon}=1\), the left side equals
\[
 (mb)\frac{n_2k}{n^2}x^{2\varepsilon}.
\]
Since \(k\ge0\), we have \(n_2k\le nk\). Apply the upper bound on \(k\)
and then the density bound to obtain
\[
 \frac{n_2k}{n^2}x^{2\varepsilon}
 \le n x^{1/2+\varepsilon}
 \le2n_Tx^{1/2+9\varepsilon/8}\le2n_TY.
\]
Multiply by \(mb\). For positive \(Q\), the exact normalized
expression is
\[
 \frac{mN_*}{Q}=\frac{n_2}{n_T}\frac{k}{n^2}x^{2\varepsilon}.
\]
\end{proof}

\section{The square-root estimate}

\begin{theorem}[The size-window bound]\label{B5.55}
Let \(C\ge0\), \(\varepsilon>0\), \(x\ge3\), and
\(n,n_T,n_2,m\in\N\), with \(n>0\). Let \(b>0\),
\(a_E\ge0\), and \(k,E,G_x\in\R\).
Suppose
\[
\begin{gathered}
 a_E\le2xb,\qquad
 \tfrac12 nx^{-\varepsilon/8}\le n_T,\qquad
 x^{-\varepsilon}\le\delta_2=\frac{n_2}{n},\qquad n_2\le n,\\
 x^{1/2-\varepsilon}\le m\le x^{1/2+\varepsilon/2},\qquad
 xm\le k\le n^2x^{1/2-\varepsilon},\\
 E\le x^{C\varepsilon}\mathcal B_x(a_E),\qquad
 26\le x^{3\varepsilon/8}.
\end{gathered}
\]
If
\[
 Q\frac{\delta_2}{2}
 \left(G_x-4x^{1/2+\varepsilon}\right)
 \le E+\frac{2Q}{x},
 \qquad Q=n_Tmb,
\]
then
\[
 G_x\le x^{1/2+(C+3)\varepsilon}.
\]
\end{theorem}
\leanref{Formalism.\allowbreak{}LFunctionChainParameters.\allowbreak{}ThickSetMeanBounds.\allowbreak{}proposition\_\allowbreak{}size\_\allowbreak{}bound}{Formalism.\allowbreak{}SizeWindowBrackets.\allowbreak{}ChainSlopeWindows.\allowbreak{}LipschitzWeightAndMinimum}{507--593}

\begin{proof}
The lower bounds imply \(n_T,m,k>0\), hence \(Q>0\).
Put \(P=x^{C\varepsilon}\), \(X=x^{1/2+\varepsilon}\), and
\(Y=x^{1/2+13\varepsilon/8}\).
Then \(P,Y\ge1\) and \(X\le Y\).
Lemmas~\ref{B5.51}--\ref{B5.54} give
\[
\begin{aligned}
 \mathcal B_x(a_E)
 &=n_Ta_E+n_Ta_E\frac{m^2}{k}+mB_*+mN_*\\
 &\le(2+2+4+2)YQ=10YQ.
\end{aligned}
\]
Thus \(E/Q\le10PY\). Multiplying the assumed inequality by \(2/Q\)
gives
\[
 \delta_2(G_x-4X)\le20PY+\frac4x\le22PY.
\]
The last step uses \(x\ge3\) and \(PY\ge1\).
The density lower bound gives \(0<\delta_2\) and
\(\delta_2^{-1}\le x^\varepsilon\). Hence
\[
 G_x\le4X+22PYx^\varepsilon.
\]
Substituting the powers gives
\[
 PYx^\varepsilon=x^{1/2+(C+21/8)\varepsilon},
 \qquad
 X\le x^{1/2+(C+21/8)\varepsilon}.
\]
Consequently
\[
 G_x\le26x^{1/2+(C+21/8)\varepsilon}
 \le x^{3\varepsilon/8}
 x^{1/2+(C+21/8)\varepsilon}
 =x^{1/2+(C+3)\varepsilon}.
\]
\end{proof}

For a finite set \(K_p\subseteq\Psi_A\) of retained characters and
\(c\in\R\), define
\[
 \omega_c(\psi)=
 \begin{cases}
 1,&\psi\text{ is trivial on }H,\\
 c,&\text{otherwise},
 \end{cases}
 \qquad
 G_m^A(x)=\sum_{\psi\in K_p}\omega_c(\psi)\Re G_\psi(x),
\]
and
\[
 \mathrm{Sig}_A(x)=
 \sum_{\psi\in\Psi_A\setminus K_p}|G_\psi(x)|.
\]
Theorem~\ref{B2.14} denotes this remainder by \(\Sigma_A\).

\begin{theorem}[Return to the retained character sum]\label{B5.56}
Fix real numbers \(C,\varepsilon,C_E,K_o,c,M\), with
\(C\ge0\), \(\varepsilon>0\), and \(|c|\le1\).
For some \(x_1\) depending on these constants, the following holds
for every \(x\ge x_1\).

Let \(H\) be a finite abelian group, \(A_0\le H\), \(TT\subseteq H\),
and \(T_2\subseteq A_0\). Impose all the assumptions of
Lemma~\ref{B5.49} preceding the introduction of \(E\) and \(Z\).
Let \(a_E\ge0\), \(k>0\), and suppose
\[
\begin{gathered}
 a_E\le2xb,\qquad
 \tfrac12 nx^{-\varepsilon/8}\le n_T,\qquad
 x^{-\varepsilon}\le n_2/n,\\
 x^{1/2-\varepsilon}\le m\le x^{1/2+\varepsilon/2},\qquad
 xm\le k\le n^2x^{1/2-\varepsilon}.
\end{gathered}
\]
Let \(K_p\subseteq\Psi_A\) contain \(\mathbf1\). Suppose
\[
 \#\Psi_A\le M,\qquad
 \psi(((1,0),0))=1\quad(\psi\in K_p),
\]
and assume the character bounds
\[
\begin{aligned}
 \left|\sum_{t\in TT}\psi((0,0),t)-cn_T\right|
 &\le Mx^{-1+\varepsilon}n_T
 &&(\psi\in K_p,\ \psi|_H\ne\mathbf1),\\
 |\ell(\psi)|&\le K_o
 &&(\psi\in K_p\setminus\{\mathbf1\}),\\
 |\ell(\psi)|&\le C_E\log x
 &&(\psi\in\Psi_A\setminus K_p).
\end{aligned}
\]
If every \(u_0\in A_0\) satisfies
\[
 \mathrm{gap}(u_0)\le x^{C\varepsilon}\mathcal B_x(a_E),
\]
then
\[
 G_m^A(x)\le
 x^{1/2+(C+3)\varepsilon}+\mathrm{Sig}_A(x).
\]
\end{theorem}
\leanref{Deduction.\allowbreak{}ch\_\allowbreak{}core}{Development.\allowbreak{}AutomorphicArithmetic.\allowbreak{}DedekindLFunctions.\allowbreak{}LandauTheorems.\allowbreak{}ChCore}{343--443}
\leanref{Deduction.\allowbreak{}CharacterBalanceMoves.\allowbreak{}one\_\allowbreak{}char}{Development.\allowbreak{}AutomorphicArithmetic.\allowbreak{}DedekindLFunctions.\allowbreak{}LandauTheorems.\allowbreak{}ChCore}{202--282}

\begin{proof}
Take \(x_1\) from the proof of Theorem~\ref{B2.14}. For \(x\ge x_1\),
equation~\eqref{eq:p2c2-top-lower} gives the lower bound in
Lemma~\ref{B5.49} with
\[
 Z=G_m^A(x)-\mathrm{Sig}_A(x)-4x^{1/2+\varepsilon}.
\]
Apply Lemma~\ref{B5.49} with \(E=x^{C\varepsilon}\mathcal B_x(a_E)\).
Since \(T_2\subseteq A_0\subseteq H\), we have \(n_2\le n\).
Theorem~\ref{B5.55} applies with
\(G_x=G_m^A(x)-\mathrm{Sig}_A(x)\).
\end{proof}

Fix a family as in Theorem~\ref{B5.33}, with division and loss exponents
bounded by \(e_D,e_L\) and height exponent bounded by \(c_{HH}\).
Fix \(\Lambda,C_E\), put \(C=e_D+e_L+3\), and take \(x\) large enough
that
\[
 \log x\le x^\varepsilon,\qquad
 200(1+K)^2(1+c_{HH})(1+e_K)\le x^\varepsilon.
\]
Assume \(\mathrm{CoreAtH}(H_{\rm host},TT,e_D,e_L,K,c_{HH},e_K,
\varepsilon,x)\) as in Theorem~\ref{B2.20}. Use its chain datum with the accompanying \(N,w,\ell\). These satisfy the weight and cutoff bounds of Lemma~\ref{B5.49}.
The lemma's thin-character estimate also holds.
In the family of Theorem~\ref{B5.33}, the construction takes
\(N=\lceil2a\rceil\) and uses the \(\ell\) chosen at \(i_0\).

The height bound \(c_{HH}\) gives both large-\(x\) inequalities for the chain. By the covering property,
every \(s\in A_0\) has the form \(\mathrm{tr}(u)\). Apply
Theorem~\ref{B5.47} at \(u\). Since \(C_{\rm err}\le C\) and
\(\mathcal B_x(a)\ge0\), the result is the gap bound in
Theorem~\ref{B5.56} with \(a_E=a\).

For the same host and top set, assume
\(\mathrm{EAtH}(\Lambda,H_{\rm host},TT,x,\varepsilon,C_E)\)
as in Theorem~\ref{B2.20}. Lemma~\ref{B2.17} gives the character
bounds for these data and the same \(\ell\). Theorem~\ref{B5.56} gives
\[
 G_m^A(x)-\mathrm{Sig}_A(x)
 \le x^{1/2+(e_D+e_L+6)\varepsilon}
\]
for sufficiently large window points. Theorem~\ref{B2.18} identifies
\(G_m^A(x)\) with the target sum. Theorem~\ref{B2.19} bounds the
remainder by \(4o\,x^{1/2+\varepsilon}\), where \(o\) is the order
parameter of the target level. The proof of Theorem~\ref{B2.20}
uses \(1+4o\le x^\varepsilon\) to obtain the target upper bound
\(x^{1/2+(e_D+e_L+7)\varepsilon}\).


\part{The window bound VWin}

% ---- p6c1 ----
\chapter{The window bound and its inputs}\label{ch:window-inputs}

\begin{definition}[Local kernels and the total bound]\label{B6.01}
For a number field $E$ and prime $\ell$, set $S_\ell=\{v:v\mid\ell\}$.
Let $G_{E,S_\ell}$ be the Galois group of the maximal extension of $E$
unramified at finite places outside $S_\ell$. For $k\in\Z$ set
\[
 \Sha_{\ell,k}(E)=\ker\left(H^1(G_{E,S_\ell},\Q_\ell/\Z_\ell(-k))
 \to\prod_{v\in S_\ell}H^1(G_{E_v},\Q_\ell/\Z_\ell(-k))\right).
\]
Use continuous cohomology with discrete twists and let $s_{\ell,k}(E)$
be the kernel's $\Z_\ell$-length. For $t\in\N$, $B\in\R$,
$\mathrm{VVle}(E,t,B)$ requires that for every algebraic closure $\Omega/E$
and choices of places above $S_\ell$ and descended cyclotomic characters
at each $\ell$, the lengths for $1\le k<t$ are finite and
\[
 \sum_{\ell\in S}\log\ell\sum_{1\le k<t}s_{\ell,k}(E)\le B
\]
for every finite set $S$ of rational primes.
\end{definition}
\leanref{Formalism.\allowbreak{}GaloisCohomology.\allowbreak{}Sha}{Formalism.\allowbreak{}GaloisCohomology.\allowbreak{}Statement}{50--79}
\leanref{Formalism.\allowbreak{}GaloisCohomology.\allowbreak{}shaLen}{Formalism.\allowbreak{}GaloisCohomology.\allowbreak{}Statement}{78}
\leanref{UpperTheory.\allowbreak{}Top.\allowbreak{}VVle}{Development.\allowbreak{}AutomorphicArithmetic.\allowbreak{}CharacterBoxes.\allowbreak{}ExistsPageN}{12--18}

\begin{definition}[Fields, prime pairs and page functions]\label{B6.02}
Fix a primitive Dirichlet character $\psi$ modulo $M$, $a\in\N$ and
primes $q,d$. Then $\mathrm{IsChiField}(\psi,a,q,d,E)$ means
$E\simeq_\Q\Q(\zeta_N)^H$ for a positive common multiple $N$ of
$2^{a+2},M,q,d$, with $\zeta_N$ primitive and
\[
 H=\{u\in(\Z/N\Z)^\times:u\equiv\pm1\pmod{2^{a+2}},\
 \psi(u)=1,\ (u/q)=(u/d)=1\}.
\]
The symbols are Jacobi symbols. For $0<\varepsilon<1/52$ and $U>0$,
set $A=U^{\varepsilon/5}$. A window pair $\mathrm{WinPair}(\varepsilon,U,q,d)$
consists of primes with
\[
\begin{gathered}
 q\equiv3\pmod8,\qquad d\equiv7\pmod{16},\\
 A/2<\log q\le A,\qquad A(\log U)^2/2<\log d\le A(\log U)^2,\\
 q\mid d+1,\qquad p\mid d+1\quad\text{for every odd prime }p\le A.
\end{gathered}
\]
A page function $a:\R\to\N$ satisfies
$x^{2+8\varepsilon}\le2^{a(x)}\le2x^{2+8\varepsilon}$ for $x\ge1$.
\end{definition}
\leanref{UpperTheory.\allowbreak{}Top.\allowbreak{}IsChiField}{Development.\allowbreak{}AutomorphicArithmetic.\allowbreak{}CharacterBoxes.\allowbreak{}ExistsPageN}{19--25}
\leanref{UpperTheory.\allowbreak{}Top.\allowbreak{}WinPair}{Development.\allowbreak{}AutomorphicArithmetic.\allowbreak{}CharacterBoxes.\allowbreak{}ExistsPageN}{27--37}
\leanref{UpperTheory.\allowbreak{}Top.\allowbreak{}PageN}{Development.\allowbreak{}AutomorphicArithmetic.\allowbreak{}CharacterBoxes.\allowbreak{}ExistsPageN}{39}

\begin{definition}[The window assertion]\label{B6.03}
Put $t(x)=\lfloor x^{C_t}\rfloor_\N$ for $C_t\in\R$.
In $\mathrm{VWinAt}(\psi,C_t)$ each $0<\varepsilon<1/52$ and page
function $a$ admit reals $J\ge1,X_0$ such that for every $X\ge X_0$
and window pair $(q,d)$ at $X^{1/J}$, some $Y\ge X^{1/J}$ satisfies $Y^2\le X$
and $\mathrm{VVle}(E,t(x),e^{x^\varepsilon/2})$ for all $x\in[Y,Y^2]$
and $E$ with $\mathrm{IsChiField}(\psi,a(x),q,d,E)$.
Write $\mathrm{VWin}$ for $\mathrm{VWinAt}(\psi,401)$ for every primitive $\psi$.
\end{definition}
\leanref{UpperTheory.\allowbreak{}Top.\allowbreak{}VAt}{Development.\allowbreak{}AutomorphicArithmetic.\allowbreak{}CharacterBoxes.\allowbreak{}ExistsPageN}{41--42}
\leanref{UpperTheory.\allowbreak{}Top.\allowbreak{}VWinAt}{Development.\allowbreak{}AutomorphicArithmetic.\allowbreak{}CharacterBoxes.\allowbreak{}ExistsPageN}{44--47}
\leanref{UpperTheory.\allowbreak{}Top.\allowbreak{}VWin}{Development.\allowbreak{}AutomorphicArithmetic.\allowbreak{}CharacterBoxes.\allowbreak{}ExistsPageN}{49}

\begin{theorem}[The representation input gives the window]\label{B6.04}
The representation assertion $\mathcal R$ of \ref{B9.02} with conditions
\ref{B9.04}--\ref{B9.06} implies $\mathrm{VWin}$.
\end{theorem}
\leanref{UpperTheory.\allowbreak{}One.\allowbreak{}vwin\_\allowbreak{}of\_\allowbreak{}row}{Development.\allowbreak{}AutomorphicArithmetic.\allowbreak{}AutomorphicVarieties.\allowbreak{}ModularityArgument.\allowbreak{}ThresholdArguments.\allowbreak{}DenseThresholdCases.\allowbreak{}VwinOfRow}{4--5}
\leanref{UpperTheory.\allowbreak{}FinalTheorems.\allowbreak{}everyCharacter}{Development.\allowbreak{}AutomorphicArithmetic.\allowbreak{}CharacterBoxes.\allowbreak{}MainConclusion}{7--11}
\leanref{UpperTheory.\allowbreak{}Gal.\allowbreak{}Int.\allowbreak{}galoisLocalGlobalStrong\_\allowbreak{}recParameter}{Development.\allowbreak{}AutomorphicArithmetic.\allowbreak{}Foundations.\allowbreak{}GaloisRepresentations.\allowbreak{}IntegralStructures.\allowbreak{}RowSrc}{6}
\begin{proof}
For $h:\mathcal R$, apply \ref{B6.05} with $O=\langle h\rangle$ and $h$.
For the base change input, use \ref{B6.04a}.
\end{proof}
\leanref{UpperTheory.\allowbreak{}LiteratureTheorems.\allowbreak{}OpenImageCondition}{Development.\allowbreak{}AutomorphicArithmetic.\allowbreak{}AutomorphicVarieties.\allowbreak{}ModularityArgument.\allowbreak{}ThresholdArguments.\allowbreak{}DenseThresholdCases.\allowbreak{}TemperedProductOfSquareIntegrable}{207--208}

\begin{lemma}[One base change for all places]\label{B6.04a}
Let $F/F^+$ be CM and unramified at finite places, and $\pi$ automorphic
on the standard definite unitary group of rank $n\ge1$. There are
discrete $\GL_{n_i}(\mathbf A_F)$ components, $\sum_i n_i=n$, with
$n_i=a_i d_i$ and $a_i,d_i\ge1$. In the archive's induction convention,
component $i$ embeds in the induction of
$\rho_i|\det|^{(a_i-1)((d_i-1)/2-j)}$ for $0\le j<d_i$, with $\rho_i$
normalized, unitary, conjugate self-dual and cuspidal on $\GL_{a_i}(\mathbf A_F)$.
The local induced product at finite $w$ has an irreducible subquotient
$P_w$, spherical if $\pi$ is unramified there. For every unramified
Satake datum $(\alpha_i)_{1\le i\le\lfloor n/2\rfloor}$ of $\pi$ at an
inert $w$, $P_w$ has multiset $\{\alpha_i^{\pm1}\}_i$ with an extra $1$
for odd $n$. At split $w$, $P_w$ occurs in the local restriction of
$\pi$. For each weight $\lambda$ of $\pi$ at $\tau:F\hookrightarrow\C$,
the multiset union of component infinitesimal characters is
$\{\lambda_j+(n-1)/2-j:0\le j<n\}$.
\end{lemma}
\leanref{PositivityTheory.\allowbreak{}Rg.\allowbreak{}BcN.\allowbreak{}IsBaseChangeN}{Development.\allowbreak{}WeilBound.\allowbreak{}RepresentationTheoryGL.\allowbreak{}ArchimedeanBaseChange.\allowbreak{}ArchimedeanBaseChangeProof.\allowbreak{}BaseChangeIsobaricNOf}{21--37}
\leanref{PositivityTheory.\allowbreak{}Rg.\allowbreak{}BcN.\allowbreak{}BaseChangeIsobaricN}{Development.\allowbreak{}WeilBound.\allowbreak{}RepresentationTheoryGL.\allowbreak{}ArchimedeanBaseChange.\allowbreak{}ArchimedeanBaseChangeProof.\allowbreak{}BaseChangeIsobaricNOf}{39--46}
\leanref{UpperTheory.\allowbreak{}BaseChangeIdentities.\allowbreak{}baseChangeN}{Development.\allowbreak{}AutomorphicArithmetic.\allowbreak{}Foundations.\allowbreak{}BaseChangeIdentities.\allowbreak{}DefiniteNHolds}{5--6}
\begin{proof}
For a set $S$ of at most one split place, fix a level $K$ and
finite-dimensional archimedean type spaces $W$. The Hecke comparison
tests $S\cup T$ for each finite set $T$ of nonsplit places where $\pi$
is unramified. If no matching spectral character is nonzero on its
block's squared weight operator, a separating test gives a target
matrix $A$ with eigenvalue $1$ and weighted comparison norms at
most $1/2$. The trace-power bound gives
$|\operatorname{tr}(A^k)|\le C(1+k)^N2^{j-k}$ for fixed $C,N,j$ and
large $k$. Hence $\sum_{k\ge1}\operatorname{tr}(A^k)z^k$ converges on
$|z|<2$. Near zero it equals $\sum_\lambda m_\lambda\lambda z/(1-\lambda z)$
with eigenvalue multiplicities $m_\lambda$. The term for $\lambda=1$
gives a pole at $z=1$, a contradiction. The matching character gives a
family $D_T$ of the required weight at $K,W$ that matches at $S\cup T$
and almost all split places.

Fixed $K,W$ and weight give a finite list of families up to local
equivalence and permutation. If each fails at an unramified nonsplit
place, take one such place per family for $T$. Since $D_T$ matches
throughout $T$, one family matches at every unramified nonsplit place.
Choose it with $S=\varnothing$ and repeat with $S=\{w\}$ for each split $w$.
Since $\lambda_j\ge\lambda_{j+1}$, consecutive entries of
$\lambda_j+(n-1)/2-j$ differ by at least one. Cofinite Satake equality
gives local rigidity up to permutation, so the second family's $P_w$
works for the first. At nonsplit $w$ use the matching spherical
constituent if $\pi$ is unramified and any irreducible subquotient
otherwise. Satake comparison gives conjugate-dual stability. The
discrete-spectrum string decomposition gives the cuspidal data.
\end{proof}
\leanref{PositivityTheory.\allowbreak{}Rg.\allowbreak{}BcB.\allowbreak{}One.\allowbreak{}fieldN\_\allowbreak{}of}{Development.\allowbreak{}WeilBound.\allowbreak{}RepresentationTheoryGL.\allowbreak{}ArchimedeanBaseChange.\allowbreak{}ArchimedeanBaseChangeProof.\allowbreak{}FieldNOf}{52--82}
\leanref{PositivityTheory.\allowbreak{}Rg.\allowbreak{}BcB.\allowbreak{}Plug2.\allowbreak{}bci\_\allowbreak{}of}{Development.\allowbreak{}WeilBound.\allowbreak{}RepresentationTheoryGL.\allowbreak{}ArchimedeanBaseChange.\allowbreak{}ArchimedeanBaseChangeProof.\allowbreak{}SatakeExponents.\allowbreak{}BciOf}{22--26}

\begin{proposition}[Assembling the four groups of inputs]\label{B6.05}
Given the record $O$ with sole field $\mathcal R$, the base change of
\ref{B6.04a} and a proof of $\mathcal R$, one has $\mathrm{VWin}$.
\end{proposition}
\leanref{UpperTheory.\allowbreak{}One.\allowbreak{}Join.\allowbreak{}automorphic\_\allowbreak{}window\_\allowbreak{}of\_\allowbreak{}inputs}{Development.\allowbreak{}AutomorphicArithmetic.\allowbreak{}AutomorphicVarieties.\allowbreak{}ModularityArgument.\allowbreak{}ThresholdArguments.\allowbreak{}DenseThresholdCases.\allowbreak{}TopMainZj}{71--74}
\leanref{GlobalTheory.\allowbreak{}Mem.\allowbreak{}LiteralAutomorphy.\allowbreak{}LiteratureRankNHypothesis}{Development.\allowbreak{}AutomorphicArithmetic.\allowbreak{}AutomorphicVarieties.\allowbreak{}ModularityArgument.\allowbreak{}ThresholdArguments.\allowbreak{}DenseThresholdCases.\allowbreak{}LitAut}{12--63}
\leanref{GlobalTheory.\allowbreak{}Mem.\allowbreak{}LiteralAutomorphy.\allowbreak{}LitCoh}{Development.\allowbreak{}AbelianVarietyTheory.\allowbreak{}Definitions.\allowbreak{}AutomorphicMembership.\allowbreak{}Brauer}{105--117}
\begin{proof}
Let $h_N=(h_{\rm loc},h_{\rm aut},h_{\rm coh},h_{\rm CM})$ collect the
local representation theory, automorphic theory, Galois cohomology
and CM inputs. Use $O$ for the local and cohomology fields. The
automorphic field uses \ref{B6.04a} and $\mathcal R$. For the CM field,
abelian-scheme duality and smooth lifting give crystalline--de Rham
duality. Use Gauss--Manin and Serre--Tate in addition to period comparison.
For weak descent, ramification is allowed at a finite set $S$ of
split places of $F^+$ containing all places above a rational prime.
Thus the cuspidal representation is spherical at every nonsplit
place. Apply \ref{B6.06} with $h_N$ and the descent of \ref{B9.36a}.
\end{proof}

\begin{proposition}[From automorphic inputs to counting]\label{B6.06}
Given $h_N$ of \ref{B6.05} and weak descent in ranks two through four,
there is a record $h_H$ of type \Lean{HypothesesH'}. Every such $h_H$
implies $\mathrm{VWin}$.
\end{proposition}
\leanref{UpperTheory.\allowbreak{}V.\allowbreak{}HV.\allowbreak{}DescentVanishingReductions.\allowbreak{}published\_\allowbreak{}hypotheses\_\allowbreak{}from\_\allowbreak{}patching\_\allowbreak{}alt}{Development.\allowbreak{}AutomorphicArithmetic.\allowbreak{}AutomorphicVarieties.\allowbreak{}ModularityArgument.\allowbreak{}ThresholdArguments.\allowbreak{}DenseThresholdCases.\allowbreak{}PublishedHypothesesFromPatching}{98--103}
\leanref{UpperTheory.\allowbreak{}V.\allowbreak{}HV.\allowbreak{}DescentVanishingReductions.\allowbreak{}vwin\_\allowbreak{}of\_\allowbreak{}hypotheses\_\allowbreak{}H\_\allowbreak{}E\_\allowbreak{}alt}{Development.\allowbreak{}AutomorphicArithmetic.\allowbreak{}AutomorphicVarieties.\allowbreak{}ModularityArgument.\allowbreak{}ThresholdArguments.\allowbreak{}DenseThresholdCases.\allowbreak{}WindowBoundOfLiteratureLevel}{25--34}
\begin{proof}
Adjoin cohomological identities and Borel injectivity. Quadratic base
change and admissibility give the totally real rank-two Galois input.
Strong multiplicity one gives duality. For a newform $f$ of weight
$k\ge2$ and $j\in\Z$, cuspidal semisimplicity and admissibility give a
cuspidal $\GL_2(\mathbf A_\Q)$ representation with infinitesimal
character $\{j+\tfrac12,j+\tfrac32-k\}$. At $p$ prime to the level its
Satake multiset $B_p$ satisfies
$\prod_{\alpha\in B_p}(X-p^{j+1/2}\alpha)=P_{f,p}(X)$,
the Hecke polynomial of $f$ at $p$.

For Galois $L/K$ of prime degree $p$ with $K$ totally complex,
\Lean{cyc} gives a cuspidal lift of regular algebraic cuspidal $\GL_n$
data over $K$, $n\ge1$, unless they have a self-twist by a nontrivial
character satisfying the norm conditions below. Regular algebraic
means distinct infinitesimal weights in $\Z+(n-1)/2$ at each embedding.
Such data over $L$ descend if, for almost every $w$, their Satake
multiset agrees with that at every $w'$ above the same place of $K$.
Infinitesimal characters match by restriction of embeddings. At almost
every $w\mid v$, the lift raises Satake parameters to the residue
degree $f(w/v)$.

Two descents admit one Hecke character $\eta$ of $K$, possibly trivial,
with $\eta^p=1$ and $\eta(v)^{f(w/v)}=1$ for almost every $w\mid v$.
At almost every $v$, $\eta$ is unramified and the Satake multisets obey
$B_{1,v}=\eta(v)B_v$. Arthur and Clozel's weak matching gives the Satake comparison.
Their strong form matches at every place
\cite[Ch.~III, Thms.~4.2 and~5.1]{arthur-clozel-1989}.

For \Lean{cyc2}, take regular algebraic cuspidal rank-two data over an
imaginary quadratic field with infinitesimal weight gap an integer
at least two at each embedding. Under the same Satake invariance,
a cuspidal descent to $\Q$ matches at almost every split prime
and at infinity.
For \Lean{quat}, use the quaternion algebra ramified exactly at the
real places of an even-degree totally real field, with a finite adelic splitting.

Set $h_N=h_H.\mathrm{litN}$ and take $h_K$ from \ref{B6.15}.
For coefficients of absolute ramification index one and fixed
$\iota:\overline{\Q}_\ell\simeq\C$, \ref{B6.52}(1)--(3) and local
cohomology from $h_N$ give one marking for finiteness and the Type~I
and Type~II attachments at the deformation points.
With $h_H,h_K$ and the Kato vanishings \ref{B6.43}, this gives I
(\ref{B6.30}). Use \ref{B6.49} for II and III (\ref{B6.31}).
For V (\ref{B6.34}), use the Gaussian lift \ref{B6.51} from
\ref{B6.52}(4). Apply \ref{B6.07}.
\end{proof}
\leanref{UpperTheory.\allowbreak{}V.\allowbreak{}Z.\allowbreak{}Tf.\allowbreak{}SR.\allowbreak{}descWeakAll}{Development.\allowbreak{}AutomorphicArithmetic.\allowbreak{}AutomorphicVarieties.\allowbreak{}ModularityArgument.\allowbreak{}TraceFormula.\allowbreak{}SplitRigidity.\allowbreak{}HasOfSplit}{26--27}
\leanref{UpperTheory.\allowbreak{}V.\allowbreak{}Z.\allowbreak{}Bc.\allowbreak{}cycC\_\allowbreak{}holds}{Development.\allowbreak{}AutomorphicArithmetic.\allowbreak{}AutomorphicVarieties.\allowbreak{}ModularityArgument.\allowbreak{}BaseChangeUnits.\allowbreak{}CycFree}{87--88}

\begin{proposition}[Four counts suffice]\label{B6.07}
Given $h_N$, the small-twist bound $h_K$ (\Lean{LemmaK} for some
$C_K\in\R$), and counts I, II, III and V, one has $\mathrm{VWin}$.
\end{proposition}
\leanref{UpperTheory.\allowbreak{}V.\allowbreak{}HV.\allowbreak{}Core.\allowbreak{}of\_\allowbreak{}four\_\allowbreak{}alt}{Development.\allowbreak{}AutomorphicArithmetic.\allowbreak{}AutomorphicVarieties.\allowbreak{}ModularityArgument.\allowbreak{}ThresholdArguments.\allowbreak{}DenseThresholdCases.\allowbreak{}HVCore}{18--22}
\begin{proof}
At $C_t=401$, use \ref{B6.27} with $h_K$ and with duality and local finiteness
from $h_N$. For the rank bound, use \ref{B6.17} \cite{ferrero-washington-1979}.
On this window \ref{B6.39} gives the plus bound from I, II, III and V
with \ref{B6.32}, \ref{B6.35} and \ref{B6.36}. Use Poitou--Tate duality
and unramified orthogonality in \ref{B6.13} for the minus bound.
Apply \ref{B6.14}.
\end{proof}


% ---- p6c2 ----
\chapter{Three length bounds}\label{ch:vwin-three-length-bounds}
Use $s_\ell(E,k)=s_{\ell,k}(E)$ from \ref{B6.01} and $n=[E:\Q]$.

\section{Complex conjugation and arithmetic duality}
\begin{proposition}[Splitting and addition of the bounds]\label{B6.08}
Let $E$ be CM and $k\in\N$. Write $s_\ell^\pm$ for the lengths of
$\Sha_{\ell,k}^\pm(E)$ in conjugation signs $(-1)^k,(-1)^{k+1}$.
For odd $\ell$, $s_\ell=s_\ell^++s_\ell^-$. Fix $t\in\N$.
For each choice of data in \ref{B6.01}, assume $s_2(E,k)$ and both
$s_\ell^\pm(E,k)$ for every odd $\ell$ are finite for $1\le k<t$.
Assume also that, for every finite set $T$ of odd primes,
\[
\sum_{1\le k<t}s_2(E,k)\le B_2,\qquad
\sum_{\ell\in T}\log\ell\sum_{1\le k<t}s_\ell^\pm(E,k)\le B_\pm.
\]
Then $\mathrm{VVle}(E,t,B_2\log2+B_-+B_+)$ holds.
\end{proposition}
\leanref{UpperTheory.\allowbreak{}V.\allowbreak{}shaLen\_\allowbreak{}eq\_\allowbreak{}add}{Development.\allowbreak{}AutomorphicArithmetic.\allowbreak{}AutomorphicVarieties.\allowbreak{}LengthEqAddEigen}{88--104}
\leanref{UpperTheory.\allowbreak{}V.\allowbreak{}vvle\_\allowbreak{}of\_\allowbreak{}parts}{Development.\allowbreak{}AutomorphicArithmetic.\allowbreak{}AutomorphicVarieties.\allowbreak{}LengthEqAddEigen}{127--190}
\begin{proof}
Conjugation $c$ permutes the places above $\ell$. At odd $\ell$,
$(1\pm(-1)^kc)/2$ split the kernel. Add lengths with weights $\log\ell$.
\end{proof}

\begin{theorem}[Finite-coefficient Poitou--Tate duality]\label{B6.09}
Let $K$ be a number field, $\ell$ prime and $r\ge1$, with $\ell\ne2$
or $K$ totally complex. Set $R=\Z/\ell^r\Z$. For primitive $\zeta\in\mu_{\ell^r}$,
take conjugation-compatible invariants $\mathrm{inv}_v:H^2(K_v,\mu_{\ell^r})\simeq R$
at finite $v$. Fix a finite set $S$ of finite places containing those
above $\ell$ and choose lifts. For a finite continuous $R[G_{K,S}]$-module $A$ set
$A^\dagger=\mathrm{Hom}_R(A,\mu_{\ell^r})$,
$P^i(A)=\bigoplus_{v\in S}H^i(K_v,A)$,
$D_R=\mathrm{Hom}_R(-,R)$ and $H_G^i=H^i(G_{K,S},-)$.
\[
\begin{aligned}
0&\to H_G^0(A)\to P^0(A)\to D_RH_G^2(A^\dagger)\\
 &\to H_G^1(A)\to P^1(A)\to D_RH_G^1(A^\dagger)\\
 &\to H_G^2(A)\to P^2(A)\to D_RH_G^0(A^\dagger)\to0
\end{aligned}
\]
The sequence is exact. The duality maps sum the local cup invariants,
with the restriction of the dual class first for $i=1$ and second for $i=0,2$.
The cup of a uniformiser's Kummer class with the unramified character
sending arithmetic Frobenius to $1$ has invariant $1$.
Localization kernels $\Sha_S^i(A)$ are finite for $i=1,2$, with
$\Sha_S^2(A)\simeq D_R\Sha_S^1(A^\dagger)$, and $H_G^i(A)=0$ for $i\ge3$.
\end{theorem}
\leanref{GlobalTheory.\allowbreak{}Published.\allowbreak{}PoitouTateDuality}{Development.\allowbreak{}AbelianVarietyTheory.\allowbreak{}PublishedTheorems.\allowbreak{}GaloisCohomology₂.\allowbreak{}ExistsPrimitiveRootPow}{13--28}
\leanref{UpperTheory.\allowbreak{}V.\allowbreak{}HV.\allowbreak{}U.\allowbreak{}pt\_\allowbreak{}s}{Development.\allowbreak{}AutomorphicArithmetic.\allowbreak{}AutomorphicVarieties.\allowbreak{}ModularityArgument.\allowbreak{}ThresholdArguments.\allowbreak{}DenseThresholdCases.\allowbreak{}HVFinalU}{11--11}

This is Milne's sequence with vanishing archimedean Tate terms
\cite[Theorem~I.4.10]{milne-2006}.

\section{The negative component}
For a Dirichlet character $\theta$ of conductor $f$ and Bernoulli polynomials $B_j$, set
\[
B_{j,\theta^{\rm prim}}=f^{j-1}\sum_{a=1}^f\theta^{\rm prim}(a)B_j(a/f),\qquad
A_{N,j}(\theta)=N\mathrm{primorial}(j+1)B_{j,\theta^{\rm prim}},
\]
where $\mathrm{primorial}(h)=\prod_{p\le h,\ p\ {\rm prime}}p$.
For $j\ge2$ and $\theta(-1)=(-1)^j$, the functional equation gives
nonvanishing from $L(j,\bar\theta^{\rm prim})\ne0$
\cite[Proposition~4.1 and Theorem~4.2]{washington-1997}.
For $f\mid N$, $N>0$, the primorial clears denominators in $f^jB_j(a/f)$,
so $A_{N,j}$ is integral. For each complex embedding $\sigma$,
\[
|\sigma A_{N,j}(\theta)|\le b(N,j)=16N^{j+1}j!
\tag{6.1}
\]
by $|B_{j,\theta^{\rm prim}}|\le4f^jj!/(2\pi)^j$ and $\mathrm{primorial}(j+1)\le4^{j+1}$.
\leanref{UpperTheory.\allowbreak{}V.\allowbreak{}Bern.\allowbreak{}b\_\allowbreak{}primmem}{Development.\allowbreak{}AutomorphicArithmetic.\allowbreak{}AutomorphicVarieties.\allowbreak{}GaloisLibrary.\allowbreak{}BernoulliNumbers.\allowbreak{}Primmem}{17--55}
\leanref{UpperTheory.\allowbreak{}V.\allowbreak{}Bern.\allowbreak{}b\_\allowbreak{}nzC}{Development.\allowbreak{}AutomorphicArithmetic.\allowbreak{}AutomorphicVarieties.\allowbreak{}GaloisLibrary.\allowbreak{}BernoulliNumbers.\allowbreak{}BNzC}{10--70}
\leanref{UpperTheory.\allowbreak{}V.\allowbreak{}Bern.\allowbreak{}b\_\allowbreak{}denupper}{Development.\allowbreak{}AutomorphicArithmetic.\allowbreak{}AutomorphicVarieties.\allowbreak{}GaloisLibrary.\allowbreak{}BernoulliNumbers.\allowbreak{}Denupper}{28--39}

\begin{theorem}[A Bernoulli bound for one character]\label{B6.10}
Let $E\subseteq\Q(\mu_N)$ be CM, $\ell$ odd prime, $\ell\nmid n$,
and $O\subseteq\overline{\Q}_\ell$ a DVR finite over $\Z_\ell$ with
fraction field $K_O$. For $k\ge1$, let $\theta:\Gal(E/\Q)\to O^\times$
have conjugation sign $(-1)^{k+1}$ and level-$N$ pullback $\widetilde\theta$.
In $H^1(\Q,(K_O/O)(\theta^{-1}\varepsilon_\ell^{-k}))$ let $C_\theta$
be the kernel of restriction to finite inertia and decomposition at $\ell$,
with $\varepsilon_\ell$ cyclotomic. Restriction identifies $C_\theta$ with
the $\theta$-part of $\Sha_{\ell,k}(E)\otimes_{\Z_\ell}O$, and
\[
\mathrm{length}_O C_\theta\le v_O(A_{N,k+1}(\widetilde\theta))
\tag{6.2}
\]
for the normalized valuation $v_O$.
\end{theorem}
\leanref{UpperTheory.\allowbreak{}V.\allowbreak{}ML.\allowbreak{}iv}{Development.\allowbreak{}AutomorphicArithmetic.\allowbreak{}AutomorphicVarieties.\allowbreak{}PrimitiveFrobeniusFixed.\allowbreak{}Iv}{4--7}
\begin{proof}
Untwist and average over $\Gal(E/\Q)$ to identify the character part.
Inflate $\theta$ to $G_\Q$. Since $A_{N,k+1}=(N/f)A_{f,k+1}$, work
at $f=\mathrm{cond}(\widetilde\theta)$.
Set $W=K_O/O$, $\alpha=\theta^{-1}\varepsilon_\ell^{-k}$.
Conjugation $\delta$ acts by $-1$, so each class has a unique cocycle
with $c(\delta)=0$. Write $c=(\alpha-1)m_\ell$ on decomposition $D_\ell$.
Choose $a_0=\pi^h$ killing $c(G_\Q)$ and $m_\ell$, with $\pi$ a uniformiser.
Let $T(a_0)$ be normalized $W[a_0]$-valued cocycles unramified outside
$f\ell$ and coboundaries in $W[a_0]$ on decomposition groups $D_p$ for $p\mid f\ell$.
At $p\mid f$, $p\ne\ell$, primitivity and prime-to-$\ell$ order give
$g\in I_p$ with $\alpha(g)-1$ a unit. Subtract the inertia coboundary
in $W[a_0]$ to leave $D_p$-values in $W^{I_p}=0$.

Suppose $\alpha(g_0)-\varepsilon_\ell(g_0)$ is a unit for some $g_0$.
Set $M_0=\ell^m$ with $m\ge\max(1,v_\ell(f))$ and $H=f^{(\ell')}M_0$,
where $f^{(\ell')}$ is the prime-to-$\ell$ part of $f$.
Write $[t]_M$ for least nonnegative residues. On $\mu_H$, let $\sigma_t$
lift the action of $t$ and let $1\le b<H$ represent $g_0$.
\[
h_{M,b}(t)=\lfloor b[t]_M/M\rfloor-\lfloor b[-t]_M/M\rfloor,\qquad
z=\sum_{t\in(\Z/H\Z)^\times}h_{H,b}(t)\alpha(\sigma_t)^{-1}.
\tag{6.3}
\]
Put $\chi=\widetilde\theta^{\rm prim}$, $j=k+1$, $C_{M,j}(u)=M^{j-1}B_j(u/M)$
and $d_j=2j\mathrm{primorial}(j+1)$. For $u=[t]_H$, $a=\lfloor bu/H\rfloor$,
$v=bu-Ha$, use $h_{H,b}(t)=2a-(b-1)$ and expand the first two
Bernoulli terms $v^j/H-jv^{j-1}/2$. Clearing denominators gives
\[
d_j\bigl(b^kh_{H,b}(t)u^k+2(C_{H,j}(v)-b^jC_{H,j}(u))/j\bigr)\in H\Z.
\]
Sum with weights $\chi_H(bt)$ for $\chi$ induced to level $H$.
Reindex the Bernoulli sums using $\alpha(\sigma_t)^{-1}\equiv u^k\chi(t)\bmod M_0O$ to get
\[
d_j\bigl(b^k\chi(b)z-2(\chi(b)b^j-1)B_{j,\chi_H}/j\bigr)\in M_0O.
\tag{6.4}
\]
The choice of $g_0$ makes $\chi(b)b^j-1$ a unit. The Euler factor
for $\ell\nmid f$ is a unit. Take $|M_0|_\ell<|a_0d_jB_{j,\chi}/j|_\ell$. Then
\[
|z|_\ell=|B_{j,\chi}/j|_\ell,\qquad a_0z\mid M_0,\qquad
z\mid A_{f,j}(\widetilde\theta).
\tag{6.5}
\]

Fix a place $A_0$ above a prime $r\equiv1\bmod H$. Let $d$ be squarefree
with prime factors $q\equiv1\bmod H$, $q\ne r$, $\ell\nmid\mathrm{ord}_q(r)$.
Put $M=Hd$ and let $P$ lie below $A_0$ at level $M$. Write $\sigma_{M,t}$
for the $t$-power action on $\mu_M$. By Chebotarev choose degree-one
$P'\nmid M$ with $PP'=(y)$. For $p=\mathrm NP'$ and $D=(p-1)/M$,
lift the $D$-th residue power to $\eta$ with values in $\mu_M$.
For $1\le i\le M-2$, $J_i=J(\eta,\eta^i)$ reduces at $\sigma_{M,t}^{-1}P'$ to
$\sum_{x\in\F_p}x^{D[t]_M}(1-x)^{D[it]_M}$. The sole possible nonzero power
sum in the expansion has degree $p-1$ and coefficient prime to $p$.
It occurs for $[t]_M+[it]_M\ge M$. Since $J_i\bar J_i=p$ and $p$
is unramified, the valuation is $1$ minus the carry. The carries telescope
for $J=\prod_{i=1}^{b-1}J_i$, whose valuation is $b-1-\lfloor b[t]_M/M\rfloor$.
Set $X_d=\prod_t(\sigma_{M,t}^{-1}y)^{h_{M,b}(t)}/(\bar J/J)$.
The valuations of $\bar J/J$ at the conjugates of $P'$ are $h_{M,b}(t)$, so
\[
(X_d)=\prod_t(\sigma_{M,t}^{-1}P)^{h_{M,b}(t)},\qquad X_d\bar X_d=1
\]
by oddness of $h_{M,b}$.

Lift arithmetic Frobenius at a new $q$ to $\phi\in D_q$, acting by
$q$-th powers on $\mu_M$. The relative residue degree $s=s(d,q)$
at $r$ divides $\mathrm{ord}_q(r)$. Since $b<q$, summing $h_{qM,b}$
over lifts of $t$ gives $h(t)-h(q^{-1}t)$ for $h=h_{M,b}$. Thus
$\mathrm{Norm}(X_{qd})/(X_d/\phi^{-1}X_d)^s$ is a unit of modulus $1$
at all embeddings, a $2M$-th root of unity.

Choose $P'$ and $y$ at level $qM$. At level $M$
use the relative norms of $P'$ and $y$. Both characters are defined
on $\F_p^\times$ and satisfy $\eta_{qM}^{\,q}=\eta_M$. The upper prime
over $r$ has norm $P^s$, so the resulting lower element agrees with
$X_d^s$ up to a $2M$-th root of unity.

The residue $F_t$ of $\sigma_{qM,\tilde t}^{-1}y$ is independent
of the lift by total ramification. Frobenius gives $F_t=F_{qt}^q$.
The upper principal factor reduces to
$\prod_t F_t^{h(t)-h(q^{-1}t)}=\prod_t F_{qt}^{(q-1)h(t)}$,
the lower factor for $\sigma_{M,q}^{-1}\mathrm{Norm}(y)$.
Above $q$, the relation between the characters gives the inverse
Frobenius action on the Jacobi quotient $\bar J/J$. Divide to get
$X_{qd}\equiv\zeta'(\phi^{-1}X_d)^s$ above $q$, with $(\zeta')^{2qM}=1$.

Put $P_0(v)=g_0(v)/v^{u_0}$ with $u_0$ the exponent on $\mu_{M_0}$.
For $x_d=P_0(X_d)^{2f^{(\ell')}d}$ the root factors vanish and the norm
and residue exponents become $S=qs(d,q)$. Let $L$ list the primes of $d$,
and let $q::L$ prepend the new prime $q$.
For a layer generator $\tau_p$, put $\mathcal D_p=\sum_{i=0}^{p-2}i\tau_p^i$.
Apply $\prod_{p\in L}\mathcal D_p$ multiplicatively to $x_d,x_{qd}$
to define $y_L,w_L$. With $N_q$ the norm, $(\tau_q-1)\mathcal D_q=q-1-N_q$ gives
\[
\tau_q(y_{q::L})/y_{q::L}=w_L^{q-1}(\phi^{-1}y_L/y_L)^S.
\]
Since $M_0\mid q-1$ and $\phi\in U=G_{\Q(\mu_H)}$, induction gives
$M_0$-th roots of $g(y_L)/y_L$ at level $Hd$ for every $g\in U$.

Choose $Y_L^{M_0}=y_L$ and $z_g^{M_0}=g(y_L)/y_L$ with $z_g$ at level $Hd$.
Take primitive $\zeta\in\mu_{M_0}$ and an $O$-module generator $w$ of $W[M_0]$. Set
\[
\zeta^{\kappa_L(g)}=\frac{gP_0(Y_L)}{P_0(Y_L)P_0(z_g)},\quad
\nu_L(g)=\sum_t\beta(\sigma_t)\kappa_L(\sigma_t^{-1}g\sigma_t)w,
\quad \beta=\varepsilon_\ell\alpha^{-1}.
\]
The root $z_gg(z_h)$ gives additivity since $P_0$ commutes with
cyclotomic actions and kills $\mu_{M_0}$. Reindexing gives equivariance.
Kummer units give unramifiedness outside $Hrd$.

At $A_0$ adjoining $\rho^{M_0}=r$ is totally tamely ramified, so
take $\tau_r\in I_{A_0}$ with $\tau_r\rho=\zeta\rho$. Tame conjugation
gives weights $\beta(\sigma_t)/\varepsilon_\ell(\sigma_t)=\alpha(\sigma_t)^{-1}$.
Each $P_0$ multiplies the weighted valuation sum by $\alpha(g_0)-u_0$
modulo $M_0O$. At $q$ choose inertia $\tau_q$ inducing the layer generator,
and $z_{\tau_q}=w_L^{(q-1)/M_0}z_{\phi^{-1}}^S$. For $B=P_0(Y_L)$, the residue
$w_L\equiv(\phi^{-1}y_L)^S$ gives
$P_0(w_L)^{(q-1)/M_0}\equiv(\phi^{-1}B)^{S(q-1)}$. Inertia fixes residues
and $B\equiv(\phi^{-1}B)^q$, so the quotient reduces to
$(\phi^{-1}B/(BP_0(z_{\phi^{-1}})))^S$. Reduction embeds $\mu_{M_0}$.
Summing at conjugate places gives
\[
\nu_\varnothing(\tau_r)=2f^{(\ell')}(\alpha(g_0)-u_0)^2zw,\qquad
\nu_{q::L}(\tau_q)=-q\,s(d,q)\nu_L(\Frob_q).
\tag{6.6}
\]
The multipliers and $\beta(g_0)-1$ are units. Equivariance extends $\nu$ by
$\widetilde\nu(g)=-(\beta(g_0)-1)^{-1}\nu(gg_0g^{-1}g_0^{-1})$.

Set $(A_L)=\mathrm{Ann}(I_L)$, $I_L=\langle\nu_L(U)\rangle_O$.
Impose zero on decomposition at $r$ and $L$ at every lift to obtain
$S_L\subset T(a_0)$. The actions on $U$ are trivial modulo $M_0$.
A cocycle zero there descends to an abelian quotient where $c(\delta)=0$
forces $c=0$. Let $V\subset U$ fix $\rho$. If $c|_V=0$, conjugation on
$c(U)$ acts by $\alpha$ and by $\varepsilon_\ell$ via the radical
quotient. Their unit difference at $g_0$ forces $c(U)=0$.
For $\nu_L\bmod\pi I_L$ use $\beta(\delta)-\varepsilon_\ell(\delta)=2$.
For $0\ne c\in S_L$, the two kernels in $V$ are thus proper, as is
the subgroup fixing $\mu_{H\ell}$ by ramification at $r$. A cover
by three proper subgroups requires index two, while $\ker(c|_V)$
has odd index. By Chebotarev take a new $q$ with Frobenius in $V$
outside all three. Then $c(\Frob_q)\ne0$ and $\nu_L(\Frob_q)$ generates $I_L$.
As $r$ is an $M_0$-th power modulo $q$ and $v_\ell(q-1)=m$,
$\ell\nmid\mathrm{ord}_q(r)$.

Put $a'=A_{q::L}=A_Le$, with $a_0\mid A_L\mid a'\mid M_0$.
In $B'=O/(a')\simeq W[a']$, $y=\nu_{q::L}(\tau_q)$ has annihilator
$(A_L)$ by (6.6), so generates $(e)$. Cup $\widetilde\nu_{q::L}$
with $c'\in S_L$ by multiplication and separating $R$-linear forms
$\lambda:B'\to R$, $R=\Z/M_0\Z$. The twists give $R(1)$.
Local triviality kills cups at $r$, $L$ and $p\mid f\ell$.
At $p\ne\ell$, unramified subgroups for finite unramified
$\ell$-primary coefficients and their Tate dual are exact annihilators
\cite[Theorem~I.2.6]{milne-2006}. Thus cups away from $q$ vanish,
and by \ref{B6.09} so does the cup at $q$.

At $q$ the actions are trivial. For $\phi=\Frob_q$ and unramified
$R$-valued $u$, $u(h)c'(\phi)=u(\phi)c'(h)$ by agreement on inertia
and on $\phi$. Hence $g\mapsto\lambda(\widetilde\nu_{q::L}(g)c'(\phi))$
cups trivially with each $u$ and is unramified by orthogonality.
Thus $\lambda(yc'(\phi))=0$ for all $\lambda$, so $yc'(\phi)=0$.
Evaluation at $\phi$ maps to $W[e]$ with kernel $S_{q::L}$, since
$c'$ is unramified and vanishing on $D_q$ passes to conjugate lifts.
\[
\mathrm{length}_O S_L\le\mathrm{length}_O S_{q::L}+v_O(e).
\tag{6.7}
\]
At $r$, write $A_\varnothing=(M_0/z)e_0$, with $a_0\mid M_0/z$ by (6.5).
The same pairing using (6.6) gives quotient length at most $v_O(e_0)$.
Since $c(\phi)\ne0$, $v_O(e)>0$. Thus $A_L\mid M_0$ forces termination
at $S_{L_*}=0$. Hence
\[
\mathrm{length}_O T(a_0)\le v_O(e_0)+v_O(A_{L_*})-v_O(A_\varnothing)
=v_O(A_{L_*})-v_O(M_0/z)\le v_O(z).
\]
The images of $T(\pi^h)$ increase with $h$ and have union $C_\theta$.
By (6.5) their lengths are bounded by $v_O(A_{f,j})$, so the sequence
stabilizes and the same bound holds for $C_\theta$.

If $\alpha\equiv\varepsilon_\ell\bmod\pi$, prime-to-$\ell$ order makes
$\theta$ unramified off $\ell$. Normalize at $t\in I_\ell$ acting by
$-1$ on $\mu_\ell$, so the local vector vanishes. Each $\F_\ell$-linear
form on $W[\pi]$ gives a Kummer class of $\pm\ell^s$ for a $\pi$-torsion
cocycle. Local triviality gives $\ell\mid s$. Since $\ell$ is odd,
the class is zero. The projected cocycle vanishes since $c(t)=0$
and $\alpha(t)-1$ is a unit. Separation and induction on the
power of $\pi$ killing $c$ give $C_\theta=0$.
\end{proof}

\begin{lemma}[Cyclotomic annihilation on unramified homomorphisms]\label{B6.10a}
Let $N\ge1$, $K=\Q(\mu_N)$ and $G=(\Z/N\Z)^\times$, with $\tau_b\in\Gal(K/\Q)$
acting by $b$-th powers on $\mu_N$. For an odd positive integer $a$ prime to $N$, set
\[
j_b=\frac{a[a^{-1}b]_N-[b]_N}{N},\qquad E_b=\frac{a-1}{2}-j_b .
\]
For every ideal class $c$ of $K$, written additively,
$\sum_{b\in G}E_b\tau_b^{-1}c=0$.
If $N>2$ and $h:G_K\to X$ is a locally constant homomorphism to an additive
abelian group, zero on finite inertia, then for every $u\in G_K$ and every
choice of lifts $\sigma_b\in G_\Q$ of $\tau_b$,
\[
\sum_{b\in G}E_bh(\sigma_b^{-1}u\sigma_b)=0 .
\]
\end{lemma}
\leanref{UpperTheory.\allowbreak{}V.\allowbreak{}Stick.\allowbreak{}coef}{Development.\allowbreak{}AutomorphicArithmetic.\allowbreak{}AutomorphicVarieties.\allowbreak{}GaloisLibrary.\allowbreak{}StickelbergerTheorem.\allowbreak{}ImportsAlt}{9--17}
\leanref{UpperTheory.\allowbreak{}V.\allowbreak{}Odd.\allowbreak{}clKill}{Development.\allowbreak{}AutomorphicArithmetic.\allowbreak{}AutomorphicVarieties.\allowbreak{}OddResidueBounds.\allowbreak{}ClKill}{17--56}
\leanref{UpperTheory.\allowbreak{}V.\allowbreak{}Odd.\allowbreak{}homKill}{Development.\allowbreak{}AutomorphicArithmetic.\allowbreak{}AutomorphicVarieties.\allowbreak{}OddResidueBounds.\allowbreak{}HomKill}{9--50}
\begin{proof}
For $N\le2$ the class group is trivial. In the Stickelberger element
$\Theta_N=\sum_b([b]_N/N)\tau_b^{-1}$, multiplication by $a\tau_a^{-1}-1$
gives coefficients $j_b$. When $N$ is the conductor of $K$,
$(\beta\Theta_N)\operatorname{Cl}(K)=0$ for $\beta,\beta\Theta_N\in\Z[G]$
by Stickelberger's theorem \cite[Theorem~6.10]{washington-1997}.

For $N>2$, Chebotarev in the Hilbert class field gives a positive density
of primes in $c$. Primes of higher residue degree have density zero,
so choose a degree-one prime $P\nmid N$ in $c$
\cite[Ch.~VII, Theorem~13.4]{neukirch-1999}.
Then $p=\mathrm NP$ splits completely in $K$. Put $P_t=\tau_t^{-1}P$ and
$C_i(t)=\lfloor([it]_N+[t]_N)/N\rfloor$. For $i\not\equiv0,-1\bmod N$,
the Jacobi sum from the proof of \ref{B6.10} for $i\bmod N$ divides $p$ and has
valuation $1-C_i(t)$ at $P_t$. Thus $\prod_tP_t^{C_i(t)}$ is principal.
For $i\equiv0,-1$ it is $(1),(p)$, respectively. The identity
$NC_i(t)=[it]_N+[t]_N-[(i+1)t]_N$ gives
$\sum_{0\le i<a}C_i(t)=\lfloor a[t]_N/N\rfloor$.
Apply $\tau_a^{-1}$ to the product over $0\le i<a$. At $P_b$ the exponent is
$\lfloor a[a^{-1}b]_N/N\rfloor=j_b$. Since $\prod_bP_b=(p)$,
the fractional ideal with exponents $E_b$ is principal.

The open kernel of $h$ fixes an unramified abelian extension of the
totally complex field $K$. Thus $h=\lambda\circ A$, where
$A:G_K\to\operatorname{Cl}(K)$ is restriction to the Hilbert class
field followed by inverse Artin reciprocity. Equivariance gives
$A(\sigma_b^{-1}u\sigma_b)=\tau_b^{-1}A(u)$
\cite[Appendix, \S3]{washington-1997}.
Apply $\lambda$ to the class identity for $A(u)$.
\end{proof}

\begin{lemma}[Summation away from the degree]\label{B6.11}
For CM $E\subseteq\Q(\mu_N)$, $N,k\ge1$, $s_\ell^-(E,k)$ is finite
at odd primes $\ell\nmid n$. For finite sets $T$ of these primes,
\[
\sum_{\ell\in T}s_\ell^-(E,k)\log\ell\le n\log b(N,k+1).
\tag{6.8}
\]
\end{lemma}
\leanref{UpperTheory.\allowbreak{}V.\allowbreak{}Odd.\allowbreak{}primeBound\_\allowbreak{}of}{Development.\allowbreak{}AutomorphicArithmetic.\allowbreak{}AutomorphicVarieties.\allowbreak{}OddResidueBounds.\allowbreak{}ThmBcOfPsiBound}{293--294}
\begin{proof}
Let $D$ be the product of $A_{N,k+1}(\widetilde\theta)$ over characters
of sign $(-1)^{k+1}$. Each factor is nonzero by the parity criterion
preceding (6.1). Their product is a Galois-invariant algebraic integer,
hence lies in $\Z$.
By (6.1), $\log|D|\le n\log b(N,k+1)$. Sum (6.2) in a splitting $O$
and divide by $v_O(\ell)$ to get $s_\ell^-\le v_\ell(D)$. Sum over $T$.
\end{proof}

\begin{proposition}[Primes dividing the degree]\label{B6.12}
Let $E\subseteq\Q(\mu_N)$ be CM, $N,k\ge1$, and $\ell$ odd prime.
For some $r\in\N$,
\[
\ell^r\Sha^-_{\ell,k}(E)=0,\qquad r\log\ell\le30n(k+1)\log(k\ell N).
\tag{6.9}
\]
The length is finite and
\[
s_\ell^-(E,k)\log\ell\le240\ell n^2(k+1)\log^2(k\ell N).
\tag{6.10}
\]
For $M=(m,H_M)$, $m>0$, $H_M\le(\Z/m\Z)^\times$, put
$K_M=\Q(\mu_m)^{H_M}$ and $o=[K_M:\Q]$. Take primes
$q\equiv3\bmod8$, $d\equiv7\bmod16$, $\gcd(qd,m)=1$.
For $a\in\N$, $\mathrm{IsCompField}(M,a,q,d,E)$ means
$E\simeq_\Q K_M\Q(\mu_{2^{a+2}})^+\Q(\sqrt{-q},\sqrt{-d})$.
If $\mathrm{IsCompField}(M,a,q,d,E)$ holds, every odd prime dividing
$n$ divides $o$.
\end{proposition}
\leanref{UpperTheory.\allowbreak{}V.\allowbreak{}Odd.\allowbreak{}killPow}{Development.\allowbreak{}AutomorphicArithmetic.\allowbreak{}AutomorphicVarieties.\allowbreak{}OddResidueBounds.\allowbreak{}KillPow}{31--105}
\leanref{UpperTheory.\allowbreak{}V.\allowbreak{}Odd.\allowbreak{}stickBound}{Development.\allowbreak{}AutomorphicArithmetic.\allowbreak{}AutomorphicVarieties.\allowbreak{}OddResidueBounds.\allowbreak{}ThmBcOfPsiBound}{291--292}
\leanref{UpperTheory.\allowbreak{}V.\allowbreak{}Odd.\allowbreak{}deg}{Development.\allowbreak{}AutomorphicArithmetic.\allowbreak{}AutomorphicVarieties.\allowbreak{}OddResidueBounds.\allowbreak{}DegExistsLift}{33--100}
\begin{proof}
Put $H=N\ell$, $b=2H+1$, $j=k+1$, $\Gamma=\Gal(E/\Q)$.
With $\sigma_t$ the cyclotomic action on $E$, set
\[
\Theta_j=-\sum_{t\in(\Z/H\Z)^\times}C_{H,j}([t]_H)[\sigma_t]/j,
\qquad Y=(1-b^j[\sigma_b])\Theta_j.
\]
Bernoulli distribution projects the level-$L$ element to $Y$ for $L=H\ell^m$.
For $v=[b^{-1}t]_L$, $a_t=(bv-[t]_L)/L$, the first two Bernoulli
terms give the coefficient congruence before projection
\[
-\frac{C_{L,j}([t]_L)-b^jC_{L,j}(v)}j
\equiv -[t]_L^k((b-1)/2-a_t)\pmod{(L/d_j)\Z}.
\]
Sum over fibres to $\Gamma$. For $v_\ell(L)\ge u+v_\ell(d_j)$
the integral sum represents $z=Y\in\Z_\ell[\Gamma]$ modulo $\ell^u$.

For $x\in\Sha_{\ell,k}(E)[\ell^u]$, choose a continuous cocycle $\xi$
representing $x$. Enlarge $L$ to trivialize the twist on its values and
local cobounding vectors. Then $h=\xi|_{G_{\Q(\mu_L)}}$ is an unramified
homomorphism. Let $U_\infty=\ker(\varepsilon_\ell|_{G_E})$ and
$f=\operatorname{res}_{U_\infty}(zx)$, an additive homomorphism.
For $g\in U_\infty\cap G_{\Q(\mu_H)}$, both $g$ and its conjugates fix
$\mu_L$. The weight $[t]_L^k$ cancels the twist on the restricted values.
With $\sigma_t$ lifting the action of $t$ on $\mu_L$, we obtain
\[
 f(g)=-\sum_t\bigl((b-1)/2-a_t\bigr)h(\sigma_t^{-1}g\sigma_t)=0
\]
by \ref{B6.10a}. Since $g^{\varphi(H)}\in U_\infty\cap G_{\Q(\mu_H)}$
for all $g\in U_\infty$, we have $\varphi(H)f=0$.
Restriction is injective by inflation--restriction, as a generator of
the procyclic image minus $1$ acts surjectively on the nonzero twist.
Hence $z\varphi(H)x=0$.

In $\Z_\ell[\Gamma]/(c-(-1)^j)$, the adjugate gives
$\Delta\varphi(H)\Sha^-=0$ for the determinant $\Delta$ of multiplication
by $z$. With $d_Y=Hj(j+1)!$ and $u_Y=2b^jH^jj!$, clearing denominators
gives $D=d_Y^{n/2}\Delta\in\Z$. The level-$H$ eigenvalues
$(1-b^j\theta(b))(-B_{j,\theta}/j)$ have nonzero Euler factors
$1-\theta^{\rm prim}(p)p^{j-1}$. By parity and the Bernoulli bound,
$0<|D|\le(d_Yu_Y)^{n/2}$. Use $r=v_\ell(D\varphi(H))$ and
$\log(d_Yu_YH)\le30j\log(k\ell N)$.

Lift $\ell$-torsion classes to residual cocycles and restrict along the
prime-to-$\ell$ extension $E'=E(\mu_\ell)$ to trivialize the twist.
For $m'=[E':\Q]$ and $S$ the places above $\ell$, the $S$-unit and
$S$-class sequence bounds the Kummer classes by $\ell^{2m'}h(E')$.
Conductor-discriminant gives $\mathfrak D=|\operatorname{disc}E'|\le H^{m'}$
and $m'\le H$. Ideal counting gives
\[
\log h(E')\le\tfrac12\log\mathfrak D+
m'\log(1+\tfrac12\log\mathfrak D)\le3m'\log H.
\]
As $m'\le\ell n$, the $\ell$-torsion dimension is at most $8\ell n\log H$.
Multiplication by $\ell^i$ embeds each successive torsion quotient
in the $\ell$-torsion. Sum using (6.9).

With $H_M$ in place of $\ker\psi$, \ref{B6.02} gives the compositum.
Its Galois group embeds in a $2$-group times $\Gal(K_M/\Q)$.
\end{proof}

\section{The negative budget and the common window}
\begin{theorem}[Uniform negative budget]\label{B6.13}
Use $M,o,q,d$ from \ref{B6.12}, $a,t\ge2$ and CM $E$ with
$\mathrm{IsCompField}(M,a,q,d,E)$. Let $F$ be the least positive
integer with $E\hookrightarrow\Q(\mu_F)$. The lengths $s_\ell^-(E,k)$
are finite for odd $\ell$ and $1\le k<t$. For a finite odd prime set $T$,
\[
\begin{split}
\sum_{\ell\in T}\log\ell\sum_{1\le k<t}s_\ell^-(E,k)
 &\le B_-(240,o,n,F,t),\\
B_-(C,o,n,F,t)
 &=3nt^2\log(tF)+C\sum_{\substack{\ell\mid o\\\ell\ {\rm odd\ prime}}}
\ell n^2t^2\log^2(t\ell F).
\end{split}\tag{6.11}
\]
For real $C$ and $o\in\N$, put $K=3+|C|\sum_{\ell\mid o,\ \ell\ {\rm odd\ prime}}\ell^3$
with empty sum at $o=0$. If $F,t\ge1$, $n\le F$, and $\log(Ft)\le2W$,
then $B_-\le Ke^{8W}$.
\end{theorem}
\leanref{UpperTheory.\allowbreak{}V.\allowbreak{}Odd.\allowbreak{}thmBcFree\_\allowbreak{}of}{Development.\allowbreak{}AutomorphicArithmetic.\allowbreak{}AutomorphicVarieties.\allowbreak{}OddResidueBounds.\allowbreak{}ThmBcOfPsiBound}{295--296}
\leanref{UpperTheory.\allowbreak{}V.\allowbreak{}minusBound\_\allowbreak{}le\_\allowbreak{}exp\_\allowbreak{}of\_\allowbreak{}log\_\allowbreak{}le}{Development.\allowbreak{}AutomorphicArithmetic.\allowbreak{}AutomorphicVarieties.\allowbreak{}CMFieldExponentialBounds.\allowbreak{}minusle}{6--76}
\begin{proof}
Sum (6.8) at $N=F$ using $\log b(F,k+1)\le(2t+5)\log(tF)$ and
$(t-1)(2t+5)\le3t^2$. Sum (6.10) for odd $\ell\mid n$, using $\ell\mid o$.
For $u=Ft$, use $nt\le u$, $t\log u\le u^2$, $\log(\ell u)\le\ell u$.
The first term is at most $3u^4$ and each summand before $C$ at most
$\ell^3u^4$. $u^4\le e^{8W}$.
\end{proof}

\begin{theorem}[Absorption into the window bound]\label{B6.14}
For $C_t\ge1$, $\mathrm{JointWin}(C_t)$ from \ref{B6.27} and \ref{B6.39}
implies $\mathrm{VWinAt}(\psi,C_t)$ of \ref{B6.03} for every primitive $\psi$.
\end{theorem}
\leanref{UpperTheory.\allowbreak{}V.\allowbreak{}vwinAt\_\allowbreak{}of\_\allowbreak{}joint}{Development.\allowbreak{}AutomorphicArithmetic.\allowbreak{}AutomorphicVarieties.\allowbreak{}CMFieldExponentialBounds.\allowbreak{}AdmissibleOfPageN}{41--130}
\leanref{UpperTheory.\allowbreak{}V.\allowbreak{}log\_\allowbreak{}conductor\_\allowbreak{}mul\_\allowbreak{}height\_\allowbreak{}le\_\allowbreak{}of\_\allowbreak{}pageN}{Development.\allowbreak{}AutomorphicArithmetic.\allowbreak{}AutomorphicVarieties.\allowbreak{}CMFieldExponentialBounds.\allowbreak{}LogConductorMulHeightLeOfPageN}{6--57}
\leanref{UpperTheory.\allowbreak{}V.\allowbreak{}error\_\allowbreak{}terms\_\allowbreak{}le\_\allowbreak{}sqrt\_\allowbreak{}exp\_\allowbreak{}rpow}{Development.\allowbreak{}AutomorphicArithmetic.\allowbreak{}AutomorphicVarieties.\allowbreak{}CMFieldExponentialBounds.\allowbreak{}ErrorTermsLeSqrtExpRpow}{6--89}
\begin{proof}
Use the data from \ref{B6.03} and set $m=\mathrm{cond}\psi$.
Take $A_1=4$ and $t(x)=\lfloor x^{C_t}\rfloor$. For every $x\ge2$,
\[
 4\le x^2\le2^{a(x)}\le2x^{2+8\varepsilon}\le2x^3\le x^4,
 \qquad 2\le t(x)\le x^{C_t}.
\]
Thus $a,t$ are admissible.
For the window constants $J_0>0,C,P,X_1$, set $J=\max(J_0,1)$
and $U=X^{1/J}$. For large $X\ge X_1$, take a pair satisfying
$\mathrm{WinPair}(\varepsilon,U,q,d)$.
By \ref{B6.02}, $q,d>\max(2X^{C_t}+1,P,m)$ and
$\log(qd)\le U^{\varepsilon/5}(1+\log^2 U)\le U^{\varepsilon/4}
\le X^{\varepsilon/(4J_0)}$. Thus $\gcd(qd,m)=1$. Apply
$\mathrm{JointWin}(C_t)$ to obtain $U\le X^{1/J_0}\le Y$ and $Y^2\le X$.

On $[Y,Y^2]$, use the sharper upper bound $2^{a(x)}\le2x^{2+8\varepsilon}$.
By \ref{B6.02}, $E$ embeds at level
$\operatorname{lcm}(2^{a(x)+2},m,q,d)$ and is CM since it contains
$\Q(\sqrt{-q})$. Thus $n\le F\le2^{a(x)+2}mqd$ and
\[
\log(Ft)\le3\log2+\log m+(2+8\varepsilon+C_t)\log x+\log(qd)
\le2x^{\varepsilon/4}
\]
for large $U$, since $\log(qd)\le U^{\varepsilon/4}$. With $W=x^{\varepsilon/4}$
and $b=x^{4\varepsilon/5}$, combine \ref{B6.08}, \ref{B6.13} and the window bounds
\[
e^b\log2+Ke^{8W}+(Ft)^C+Ce^b
\le(\log2+K+1+|C|)e^{(9+2|C|)b}\le e^{x^\varepsilon/2},
\]
uniformly as $U\to\infty$ since $x^\varepsilon/b\to\infty$.
\end{proof}


% ---- p6c3 ----
\chapter{The cyclotomic tower and a common window}\label{ch:B6-cyclotomic-window}

For a number field $L$, a prime $\ell$, and $k\in\Z$, let
$S_\ell=\{v\text{ finite in }L:v\mid\ell\}$ and put
\[
 s_\ell(L,k)=\operatorname{length}_{\Z_\ell}\Sha_{\ell,k}(L),
 \qquad
 V_\ell(L,k)=\operatorname{length}_{\Z_\ell}
 H^2(G_{L,S_\ell},\Z_\ell(k+1)).
\]
The group $\Sha_{\ell,k}(L)$ is defined in \ref{B6.01}. Both lengths
may be infinite. Put $D_x=\exp(x^{4\varepsilon/5})$ for $x,\varepsilon>0$ and
$j_0(\ell,R)=\lfloor\log_\ell R\rfloor+3$, with $\lfloor\log_\ell0\rfloor=0$.
The cyclotomic layer $L_m/L_0$ of index $m$ has degree $\ell^m$.
Fix the record $h_N$ constructed under the hypotheses of Proposition~\ref{B6.05}.

\section{A bound for small twists}

\begin{theorem}[The small-twist $H^2$ bound]\label{B6.15}
There is $C_K\ge1$ such that whenever $L$ is totally complex and
contains $\mu_\ell$ for odd $\ell$ or $i$ for $\ell=2$, for each
integer $k\ge1$ and choice of local data the group
$H^2(G_{L,S_\ell},\Z_\ell(k+1))$ is finite and, with $n=[L:\Q]$
and $D_L=\operatorname{disc}(L)$,
\begin{equation}\label{eq:B6-small-H2}
 \log\#H^2(G_{L,S_\ell},\Z_\ell(k+1))
 \le\exp\!\left((2kn)^{C_K}\log(2|D_L|)\right).
\end{equation}
\end{theorem}
\leanref{UpperTheory.\allowbreak{}V.\allowbreak{}ZK.\allowbreak{}lemmaK\_\allowbreak{}litN\_\allowbreak{}alt}{Development.\allowbreak{}AutomorphicArithmetic.\allowbreak{}AutomorphicVarieties.\allowbreak{}ModularityArgument.\allowbreak{}ThresholdArguments.\allowbreak{}DenseThresholdCases.\allowbreak{}KFinal}{11--11}
\leanref{UpperTheory.\allowbreak{}V.\allowbreak{}LemmaK}{Development.\allowbreak{}AutomorphicArithmetic.\allowbreak{}AutomorphicVarieties.\allowbreak{}GeneratorCountTuples.\allowbreak{}LemmaK}{13--18}

\begin{proof}
By \ref{B6.16}, $s_\ell(L,k)\log\ell\le\exp((2kn)^C\log(2|D_L|))$
for an absolute $C$. Put $A=\Q_\ell/\Z_\ell(-k)$. By local finiteness, apply the Poitou--Tate
and local Tate dualities supplied by $h_N$ to $A[\ell^r]$ and its dual
$\Z/\ell^r\Z(k+1)$. Pass to the direct limit
for $A[\ell^r]$ and the inverse limit for the dual coefficients.
Degree-two localization has kernel dual to $\Sha_{\ell,k}(L)$ and
cokernel dual to $H^0(L,A)$. Its local targets are dual to $H^0(L_v,A)$.
\begin{equation}\label{eq:B6-H2-local-correction}
 s_\ell(L,k)\le V_\ell(L,k)\le s_\ell(L,k)+
 \sum_{v\mid\ell}\operatorname{length}H^0(L_v,A).
\end{equation}

The local cyclotomic image has index $d_v\le[L_v:\Q_\ell]\le n$
and contains $(1+\ell^2)^{d_v}$. Lifting the exponent gives
invariant length at most $2+v_\ell(d_vk)\le2+\lfloor\log_\ell n\rfloor+k$.
Thus
\begin{equation}\label{eq:B6-local-H0}
 \sum_{v\mid\ell}\operatorname{length}H^0(L_v,A)
 \le n(2+2\lfloor\log_\ell n\rfloor+k-\lfloor\log_\ell n\rfloor).
\end{equation}
By $\ell\le n+1$, this correction times $\log\ell$ is at most
$n(3+k)\log(n+1)$. Increase $C$ to absorb it.
\end{proof}
\leanref{UpperTheory.\allowbreak{}V.\allowbreak{}ZK.\allowbreak{}sFinal}{Development.\allowbreak{}AutomorphicArithmetic.\allowbreak{}AutomorphicVarieties.\allowbreak{}ModularityArgument.\allowbreak{}ClassCounting.\allowbreak{}SFinal}{100--163}
\leanref{UpperTheory.\allowbreak{}V.\allowbreak{}LDataCardLeFinrank.\allowbreak{}vFormula\_\allowbreak{}of\_\allowbreak{}classicalTheorems}{Development.\allowbreak{}AutomorphicArithmetic.\allowbreak{}AutomorphicVarieties.\allowbreak{}GaloisLibrary.\allowbreak{}GaloisCohomologyLength.\allowbreak{}Formula}{15--24}

\begin{theorem}[Class-group counting for a small twist]\label{B6.16}
There is $C$ such that for every prime $\ell$, every totally complex
$L$ with the roots of unity of \ref{B6.15}, and every integer $k\ge1$,
\[
 s_\ell(L,k)<\infty,\qquad
 s_\ell(L,k)\log\ell\le\exp\!\left((2k[L:\Q])^C\log(2|D_L|)\right).
\]

Let $F/L$ be cyclic over a totally complex $L$, with generator
$\sigma$ of order $\ell^{\nu-c}$, where $\ell$ is prime,
$1\le c\le\nu$ and $c\ge2$ if $\ell=2$. Suppose $F$ contains a
primitive $\ell^\nu$th root $\zeta$ with $\sigma(\zeta)=\zeta^{1+\ell^c}$.
Fix an integer $k\ge1$ and $b,B\in\N$. Put
\[
 h=\operatorname{ord}(\sigma),\qquad a=(1+\ell^c)^k,\qquad
 \mathcal N_a(\alpha)=\prod_{j=0}^{h-1}\sigma^j(\alpha)^{a^j}.
\]
Let $\mathcal E_{\nu,k}$ consist of functions $\varphi$ on the
finite primes of $F$, valued in $\Z/\ell^\nu\Z$, with
\[
 \sum_P\operatorname{ord}_P(\alpha)\varphi(P)=0
 \quad(\alpha\in F^\times),\qquad \varphi(\sigma^{-1}P)=a\varphi(P).
\]
Suppose every finite set $S\subset F^\times$ has size at most $B$
if its elements admit divisors $\mathfrak d_\alpha$ with
\[
 \operatorname{ord}_P(\alpha)
 =\mathfrak d_\alpha(\sigma P)-a\mathfrak d_\alpha(P)\quad(\text{all }P)
\]
and their $\mathcal N_a(\alpha)^{\ell^b}$ are distinct modulo
$F^{\times\ell^\nu}$. Then $\mathcal E_{\nu,k}$ is finite and
\begin{equation}\label{eq:B6-finite-class-count}
 \#\mathcal E_{\nu,k}\,(\ell^\nu)^{r_2(L)}
 \le B\,\ell^{(b+v_\ell(k))g(L)+c+v_\ell(k)},\qquad
 g(L)=2[L:\Q]+\lfloor\log_2 h_L\rfloor+2.
\end{equation}
Here $h_L$ is the class number and $r_2(L)$ counts complex places.
\end{theorem}
\leanref{UpperTheory.\allowbreak{}V.\allowbreak{}ZK.\allowbreak{}sCount}{Development.\allowbreak{}AutomorphicArithmetic.\allowbreak{}AutomorphicVarieties.\allowbreak{}ModularityArgument.\allowbreak{}ClassCounting.\allowbreak{}SelmerOrderCounts}{16--111}
\leanref{UpperTheory.\allowbreak{}V.\allowbreak{}ZK.\allowbreak{}sNum}{Development.\allowbreak{}AutomorphicArithmetic.\allowbreak{}AutomorphicVarieties.\allowbreak{}ModularityArgument.\allowbreak{}ClassCounting.\allowbreak{}NumThree}{211--378}

\begin{proof}
Write $Y=F^\times$ additively and $\partial:Y\to I=\operatorname{Div}(F)$
for the divisor map. Put $g=\sigma^{-1}$ and let it act on $I$ by
$(g\mathfrak d)(P)=\mathfrak d(\sigma P)$. Set $\delta=g-a$.
\[
 U=\partial^{-1}(\delta I),\qquad W=U/(\delta Y+Y_{\mathrm{tors}}),
 \qquad C=I/(\partial Y+\delta I),
 \qquad E=\frac{\ker\partial}{\delta(\ker\partial)+Y_{\mathrm{tors}}}.
\]
Since $g^h=1$ and $a>1$, $\delta$ is injective on $I$.
The snake lemma gives $0\to E\to W\to\ker(\delta|\mathrm{Cl}(F))\to0$.
The last group has order $\#C$. On complex places $g$ has $r_2(L)$
cycles of length $h$. The logarithmic unit lattice omits the constant
direction of eigenvalue $1-a$ for $\delta$. Thus
\begin{equation}\label{eq:B6-W-factor}
 \#W=\#E\,\#C,
\end{equation}
and
\begin{equation}\label{eq:B6-unit-det}
 \#E\,(a-1)=(a^h-1)^{r_2(L)}.
\end{equation}
We have $\mathcal E_{\nu,k}=\operatorname{Hom}(C,\Z/\ell^\nu\Z)$, so its
order is at most $\ell^{v_\ell(\#C)}$. Lifting the exponent gives
$v_\ell(a-1)=c+v_\ell(k)$ and $v_\ell(a^h-1)=\nu+v_\ell(k)$.

In additive notation $\mathcal N_a=\sum_{j=0}^{h-1}a^jg^{-j}$ and
$\delta\mathcal N_a=\mathcal N_a\delta=(1-a^h)g$. Thus
\[
 f:W\longrightarrow F^\times/(F^{\times\ell^\nu}\mu(F)),\qquad
 [\alpha]\longmapsto[\mathcal N_a(\alpha)^{\ell^b}]
\]
is well defined, where $\mu(F)$ is the group of roots of unity.
For $[u]\in\ker f$, write $\ell^b\mathcal N_a u=\ell^\nu y$ modulo
torsion. Apply $\delta$ and cancel $\ell^\nu$ modulo torsion. As $g$
is invertible on $W$, $\ker f$ is killed by $\ell^b(a^h-1)/\ell^\nu$,
of valuation $b+v_\ell(k)$. Distinct image classes give $\#\operatorname{im} f\le B$.

For $[u]\in W[\ell]$, choose $y$ with $\ell u=\delta y$ modulo torsion.
The classes of $y$ modulo $\ell Y+Y_{\mathrm{tors}}$ distinguish
the $[u]$'s, since $y-y'=\ell z$ modulo torsion implies
$u-u'=\delta z$ modulo torsion. Injectivity of $\delta$ on $I$ gives
$\partial y\in\ell I$, and $a\equiv1\bmod\ell$ gives invariance.
For a representative $\alpha$ and a generator $\rho$ of $\mu(F)$, write
\[
 \sigma(\alpha)=\alpha v^\ell\rho^r,\qquad 0\le r<\ell,\qquad
 (N_{F/L}v)^\ell=N_{F/L}(\rho)^{-r}.
\]
There are at most $\ell^2$ pairs $(r,N_{F/L}v)$. For equal pairs,
Hilbert 90 on the ratio of the $v$'s makes the ratio of representatives
descend to $L^\times$ modulo $\ell$th powers.
The order of $1+\ell^c$ modulo $\ell^\nu$ is $h$, so $F=L(\zeta)$
and the descended classes have valuations divisible by $\ell$
away from $\ell$. The valuations above $\ell$ contribute at most
$[L:\Q]$ dimensions. The classes with all valuations divisible by
$\ell$ map to $\mathrm{Cl}(L)[\ell]$, with kernel the units modulo
$\ell$th powers, of dimension at most $[L:\Q]$. Hence
$\#W[\ell]\le\ell^{g(L)}$ and
\[
 \ell^{v_\ell(\#W)}
 \le\#\operatorname{im} f\,\ell^{(b+v_\ell(k))g(L)}
 \le B\,\ell^{(b+v_\ell(k))g(L)}.
\]
Equations \eqref{eq:B6-W-factor} and \eqref{eq:B6-unit-det} give
\eqref{eq:B6-finite-class-count}.

For the weighted norms assume
\begin{equation}\label{eq:B6-room}
 2k+1<\ell^c,\qquad 2k+1<r_2(L).
\end{equation}
Put $G_r=\GL_r(\mathcal O_L)$ and let $T_1,T_2$ be the orders of
$H_i(G_{4k+6},\Z)_{\mathrm{tors}}$ for $i=2k+1,2k$, respectively.
Lemma~\ref{B6.16c} gives their finiteness.
The counting hypothesis holds with
\begin{equation}\label{eq:B6-few}
 b=2k+(2k+1)(c+2k+1),\qquad B=(\ell^\nu)^{r_2(L)}T_1T_2.
\end{equation}
Put $q_0=\ell^\nu$, $A_\nu=\Z/q_0\Z$, $\kappa=1+\ell^c$, and $M_*=2k+1$.
Define
\[
 P(T)=\prod_{\substack{0\le j<M_*\\j\ne k}}(T-\kappa^j),\qquad
 Q(T)=(T-a)P(T),\qquad p_{k,c}=P(a).
\]
Lifting the exponent gives $v_\ell(\kappa^d-1)=c+v_\ell(d)<c+M_*$
for $0<d<M_*$. The factors of $P(a)$ thus give, in $A_\nu$,
\begin{equation}\label{eq:B6-pe}
 p_{k,c}\mid\ell^{M_*(c+M_*)}.
\end{equation}

Let $C_{q_0}=\langle\zeta\rangle$ and $D=C_{q_0}\times C_{q_0}$.
For $0\le t\le M_*$ let $d_t(s)=(s,s^t)$. With coefficients in $A_\nu$,
the cyclic bar cycle $z$ is the sum of
$[\zeta^{i_1}\mid\zeta\mid\cdots\mid\zeta^{i_k}\mid\zeta]$
over $i_1,\ldots,i_k\bmod q_0$. Translation of the indices cancels
the boundary modulo $q_0$. The tensor product of the two periodic
resolutions gives $2k+1$ basis vectors in degree $2k$ with zero
differentials over $A_\nu$. A primitive integral relation on lifts
of the $2k+2$ classes $(d_t)_*[z]$ therefore gives
\[
 \sum_t n_t(d_t)_*z=\partial\xi,\qquad n_{t_0}\in A_\nu^\times
\]
for some $t_0$.

Identify $D$ with $A_\nu^2$ in exponent coordinates and write
$\pi_t(i,j)=i+tj$. The required weight $\eta\in\Z[A_\nu^2]$ has
\[
 (\pi_t)_*\eta=0\ (t\ne t_0),\qquad
 (\pi_{t_0})_*\eta=\sum_i[T^i]Q(T)[\kappa^i].
\]
Let $J=(X^{\ell^c}-1)$ in $\Z[X]$. Substitution
$S(f)(X)=f(X^\kappa)$ acts by $\kappa^j$ on $J^j/J^{j+1}$.
$S(f)\equiv f\bmod J$ and
\[
 S(X^{\ell^c}-1)=(X^{\ell^c}-1)
 \sum_{i=0}^{\kappa-1}X^{i\ell^c}.
\]
Thus $Q(S)X\in J^{M_*}$. The map sending $X$ to $[1]$ in $\Z[A_\nu]$ gives a target
pushforward divisible by $([\ell^c]-[0])^{M_*}$.
For $t\ne t_0$, $t_0-t$ divides $\ell^c$ in $A_\nu$ by
\eqref{eq:B6-room}, so geometric summation gives
$[t_0]-[t]\mid[\ell^c]-[0]$. Divide the target by the product of
these factors and lift the quotient by $i\mapsto(i,0)$. Multiply
the lift by $\prod_{t\ne t_0}([(0,1)]-[(t,0)])$. Its pushforward
vanishes at every $t\ne t_0$ and gives the target at $t_0$.
\leanref{UpperTheory.\allowbreak{}V.\allowbreak{}ZK.\allowbreak{}sEta}{Development.\allowbreak{}AutomorphicArithmetic.\allowbreak{}AutomorphicVarieties.\allowbreak{}ModularityArgument.\allowbreak{}ClassCounting.\allowbreak{}ExpStep}{116--182}

For $\partial\alpha=(g-a)\mathfrak d$, put
$\beta=\prod_i\sigma^{-i}(\alpha)^{[T^i]P(T)}$. Then
$\partial\beta=P(g)(g-a)\mathfrak d=Q(g)\mathfrak d$.
On $F^D$ let $\rho_0$ act diagonally, with character
$\zeta^{is_1+js_2}$ at coordinate $(i,j)$ for $(s_1,s_2)\in A_\nu^2$.
It preserves the standard lattice $\mathcal O_F^D$ and the lattice
whose coordinate $(i,j)$ is the ideal of divisor $\eta(i,j)\mathfrak d$.

Construct an $L$-linear map $G_t$ between the two lattices using
scalar blocks $\pi_t=x$. If $t\ne t_0$, every block has total
target divisor zero. At $t_0$, first apply $\sigma^{-i}$ to a
coordinate with $\pi_{t_0}=\kappa^i$, with $i$ taken modulo $h$.
The scalar becomes $\zeta$ and its divisor is transformed by $g^i$.
Leave the other coordinates fixed. The exponent-$1$ block now has
total divisor $Q(g)\mathfrak d=\partial\beta$, while the other
blocks have divisor zero. Steinitz classification gives a lattice
isomorphism in each block. A unit on an ideal summand adjusts the
determinant to $\beta$ on the distinguished block and $1$ elsewhere.
Undo the coordinate conjugations to obtain $G_t$ commuting with
$\rho_0d_t$.
\leanref{UpperTheory.\allowbreak{}V.\allowbreak{}ZK.\allowbreak{}Cl.\allowbreak{}cBetaDiv}{Development.\allowbreak{}AutomorphicArithmetic.\allowbreak{}AutomorphicVarieties.\allowbreak{}ModularityArgument.\allowbreak{}ClassCounting.\allowbreak{}LatticeBlockClasses.\allowbreak{}CBetaDiv}{33--62}
\leanref{UpperTheory.\allowbreak{}V.\allowbreak{}ZK.\allowbreak{}Cl.\allowbreak{}cGtEq}{Development.\allowbreak{}AutomorphicArithmetic.\allowbreak{}AutomorphicVarieties.\allowbreak{}ModularityArgument.\allowbreak{}ClassCounting.\allowbreak{}LatticeBlockClasses.\allowbreak{}GtEqZeta}{54--106}

Let $\iota$ be the inclusion of the standard lattice stabilizer $H$.
Put $G_*=G_{t_0}$, $\rho_1=G_*^{-1}\rho_0G_*$ and $b_t=G_*^{-1}G_t$.
Both representations take values in $H$, and $b_t\in H$ with
$\operatorname{conj}_{b_t^{-1}}\rho_1d_t=\rho_0d_t$, where
$\operatorname{conj}_g(v)=gvg^{-1}$. The bar prism $p_g$ inserts $g$
with alternating signs and conjugates the following entries by $g^{-1}$.
It satisfies $\partial p_g+p_g\partial=(\operatorname{conj}_{g^{-1}})_*-1$.
Thus
\[
 y=(\rho_1)_*\xi-(\rho_0)_*\xi+
 \sum_t n_t p_{b_t}((\rho_1d_t)_*z)
\]
is a cycle. Its last sum has boundary $(\rho_0)_*\partial\xi-(\rho_1)_*\partial\xi$.

Choose $d\in\mathcal O_L\setminus\{0\}$ that clears denominators of
the maps and $\alpha$ as well as those of their inverses. Work in the group preserving
$\mathcal O_F[1/d]^D$. Put $Z_t=(\rho_0d_t)_*z$ and
$C_*=(\operatorname{conj}_{G_*^{-1}})_*$. Prism multiplication gives
\[
 p_{G_t}Z_t\equiv p_{G_*}Z_t+p_{b_t}C_*Z_t
 \pmod{\operatorname{im}\partial}.
\]
Apply the prism boundary identity to $(\rho_0)_*\xi$. It gives
$(\rho_1)_*\xi-(\rho_0)_*\xi\equiv
p_{G_*}(\rho_0)_*(\partial\xi)$ modulo boundaries. Summing the
multiplication formula gives
\[
 \iota_* y\equiv\sum_t n_t p_{G_t}Z_t
 \pmod{\operatorname{im}\partial}.
\]
Over $\mathcal O_L$ the standard lattice is projective. Extend the
$H$-action trivially to a complementary summand in a finite free
module and push $n_{t_0}^{-1}y$ into its general linear group to get
a cycle $y_\alpha$ of degree $2k+1$.
\leanref{UpperTheory.\allowbreak{}V.\allowbreak{}ZK.\allowbreak{}Cl.\allowbreak{}cChain}{Development.\allowbreak{}AutomorphicArithmetic.\allowbreak{}AutomorphicVarieties.\allowbreak{}ModularityArgument.\allowbreak{}ClassCounting.\allowbreak{}LatticeBlockClasses.\allowbreak{}CChain}{8--63}
\leanref{UpperTheory.\allowbreak{}V.\allowbreak{}ZK.\allowbreak{}Cl.\allowbreak{}cJoin}{Development.\allowbreak{}AutomorphicArithmetic.\allowbreak{}AutomorphicVarieties.\allowbreak{}ModularityArgument.\allowbreak{}ClassCounting.\allowbreak{}LatticeBlockClasses.\allowbreak{}CJoin}{5--187}

Take an unramified prime $P$ of $F/L$ of residue degree one,
away from $\ell$ and $d$. Put $K=k(P)$ and $d_K=\#K-1$. The reduction
of $\zeta$ has order $q_0$. Choose a generator of $K^\times$ with
$(d_K/q_0)$th power $\bar\zeta$. Write $L_K$ for its least nonnegative
logarithm and put $\operatorname{ind}=L_K\bmod q_0$ and
\[
 U_K(v,w)=\frac{L_K(v)+L_K(w)-L_K(vw)}{d_K}\pmod{q_0}.
\]
Cancellation of logarithms makes $U_K$ a symmetric two-cocycle.
In the prism evaluation of $U_K^{\smile k}\smile\operatorname{ind}$,
adjacent insertions in each $U_K$ pair cancel by symmetry. The
last contributes $\operatorname{ind}(g)$. For scalar exponent $x$,
each pair contributes
\[
 \sum_{i\bmod q_0}U_K(\bar\zeta^{xi},\bar\zeta^x)
 =q_0 L_K(\bar\zeta^x)/d_K=x\quad\text{in }A_\nu.
\]

For a finite-dimensional $K$-space $V$, transfer the cocycle over
the points of $\mathbf P(V)$, the one-dimensional subspaces.
Choose a nonzero vector on each line. For $(g_1,\ldots,g_d)$ and
initial line $\lambda_0$, follow $\lambda_i=g_i^{-1}\lambda_{i-1}$.
Record the scalars of $g_i$ between the chosen vectors and sum the
cocycle on these scalar tuples over $\lambda_0$ to define $\theta_{P,V}$.
Scalar multiplication agrees with the bar faces. Inserting
change-of-vector scalars gives a cochain homotopy. The cocycle
kills boundaries and is invariant under isomorphism on cycles.
For $V\oplus W$, lines in $W$ give its contribution. Each line of
$V$ has $(\#K)^{\dim W}\equiv1\bmod q_0$ lifts with vectors chosen
by projection. This gives additivity, and a trivial summand contributes zero.

Prism multiplication makes evaluation for a scalar action of
exponent $x$ a homomorphism in the inserted matrix $g$. It factors
through the determinant since $\mathrm{SL}$ is generated by
elementary transvections, which are commutators with diagonal matrices.
The scalar calculation gives $x^k\operatorname{ind}(\det g)$.
\leanref{UpperTheory.\allowbreak{}V.\allowbreak{}ZK.\allowbreak{}sCoc}{Development.\allowbreak{}AutomorphicArithmetic.\allowbreak{}AutomorphicVarieties.\allowbreak{}ModularityArgument.\allowbreak{}ClassCounting.\allowbreak{}FProd}{101--150}
\leanref{UpperTheory.\allowbreak{}V.\allowbreak{}ZK.\allowbreak{}sTrans}{Development.\allowbreak{}AutomorphicArithmetic.\allowbreak{}AutomorphicVarieties.\allowbreak{}ModularityArgument.\allowbreak{}ClassCounting.\allowbreak{}STrans}{411--413}
\leanref{UpperTheory.\allowbreak{}V.\allowbreak{}ZK.\allowbreak{}sThetaOf}{Development.\allowbreak{}AutomorphicArithmetic.\allowbreak{}AutomorphicVarieties.\allowbreak{}ModularityArgument.\allowbreak{}ClassCounting.\allowbreak{}SThetaOf}{192--273}

Let $r:\mathcal O_F\to K$ be reduction at $P$. The maps
$r\circ\sigma^j$ identify $\mathcal O_F\otimes_{\mathcal O_L}K$
with $\prod_{j=0}^{h-1}K$ by the Chinese remainder theorem for the
split prime below $P$. For each $t$ reduce the commuting pair
$(\rho_0d_t,G_t)$ in its chosen block coordinates. At $t=t_0$ the
distinguished block in component $j$ has exponent $\kappa^j$ and
determinant $\overline{\sigma^j(\beta)}$. The other determinants
are $1$. The free representation adds a trivial summand. Write
$\theta_P$ for evaluation in the resulting representation. The chain
congruence and normalization give
\[
 \theta_P(y_\alpha)=\sum_{j=0}^{h-1}a^j\operatorname{ind}(\sigma^j(\beta))
 =P(a)\sum_{j=0}^{h-1}a^j\operatorname{ind}(\sigma^j(\alpha))
 =p_{k,c}\operatorname{ind}(\mathcal N_a(\alpha)).
\]
For the middle equality expand $\beta$ and shift $j$ by $i$
in the term $[T^i]P(T)$. The shift uses $a^h\equiv1\bmod q_0$.
\leanref{UpperTheory.\allowbreak{}V.\allowbreak{}ZK.\allowbreak{}Cl.\allowbreak{}cBetaVal}{Development.\allowbreak{}AutomorphicArithmetic.\allowbreak{}AutomorphicVarieties.\allowbreak{}ModularityArgument.\allowbreak{}ClassCounting.\allowbreak{}LatticeBlockClasses.\allowbreak{}CBetaVal}{31--122}

On a finite set of $\alpha$'s, let $\Phi=(\theta_P)_P$ over the
common good primes. Equality of $\ell^{2k}\Phi(y_\alpha)$ and
$\ell^{2k}\Phi(y_{\alpha'})$ gives by \eqref{eq:B6-pe} that
\[
 \mathcal N_a(\alpha)^{\ell^b}\mathcal N_a(\alpha')^{-\ell^b}
\]
is a $q_0$th power at every good prime. All residue roots then lie
in $k(P)$, and reduction is injective on $\mu_{q_0}$. Frobenius thus
fixes a $q_0$th root in the Kummer extension. Chebotarev at degree-one
primes gives the global $q_0$th power relation.
\leanref{LateralTheory.\allowbreak{}Q.\allowbreak{}SDetectCut.\allowbreak{}goal}{Development.\allowbreak{}EtaleCohomology.\allowbreak{}CyclotomicDiscriminants.\allowbreak{}SDetectCut}{101--142}

Stabilize the cycles to a common rank $r\ge4k+6$. Additivity and
vanishing on trivial summands preserve their $\Phi$ values. The sequence
$0\to\Z\xrightarrow{q_0}\Z\to A_\nu\to0$ gives a Bockstein with image
in $H_{2k}(G_r,\Z)[q_0]$. By Lemma~\ref{B6.16a}, since
$2(2k)+4\le4k+6\le r$, this subgroup has at most $T_2$ elements.
Differences within a fibre lift to integral homology in degree $2k+1$.
Denote a lift by $x$. Let $D_\ell$ repeat a matrix $\ell$ times in
block diagonal form, and let $S$ stabilize it to the same rank by
adding an identity. In integral homology set
\[
 z_0(x)=x,\qquad
 z_{t+1}(x)=(D_\ell)_*z_t(x)-\ell^{t+2}S_*z_t(x).
\]
Each step multiplies the $\Phi$ value by $\ell(1-\ell^{t+1})$.
The second factor is a unit in $A_\nu$, so $\ell^{2k}\Phi(x\bmod q_0)$
comes from an integer multiple of $z_{2k}(x)$.

A stable real cohomology functional of weight $j$ is multiplied
by $\ell^j$ under repetition by $\ell$. The recurrence acts on
it by $\prod_{t=0}^{2k-1}(\ell^j-\ell^{t+2})$, which vanishes for
$2\le j\le2k+1$. Under \eqref{eq:B6-room}, the stable comparison in
degree $2k+1$ expresses each cohomology evaluation as a sum of
these higher weights and a linear combination of one weight-one
functional per complex place. A class in the recurrence image
annihilated by these $r_2(L)$ functionals therefore pairs to zero with
every stable cohomology class. The pairing over $\R$ and
Lemma~\ref{B6.16a} give vanishing over $\R$ at the final matrix rank,
so the integral image has rank at most $r_2(L)$. This upper bound
equals $\dim_{\R}(K_{2k+1}(\mathcal O_L)\otimes\R)$ for totally complex
$L$ by Borel's formula \cite[Prop.~12.2]{borel-1974}.
\leanref{PositivityTheory.\allowbreak{}KB.\allowbreak{}sBStarOdd\_\allowbreak{}of}{Development.\allowbreak{}AutomorphicArithmetic.\allowbreak{}AutomorphicVarieties.\allowbreak{}LiteratureResults.\allowbreak{}KTheoryStability.\allowbreak{}ReductionBoundsTorsion.\allowbreak{}OddSpanRank.\allowbreak{}SBStarOddOf}{13--78}

At the final rank $r'=\ell^{2k}r$, Lemma~\ref{B6.16a} gives
$H_{2k+1}(G_{r'},\Z)$ torsion of order $T_1$, since
$2(2k+1)+4=4k+6\le r'$. In each Bockstein fibre subtract a
representative and take the subgroup generated by the integral
recurrence images of the lifted differences. It has rank at most
$r_2(L)$ and torsion of order at most $T_1$. For the chosen representative
$y_0$, the values $\ell^{2k}(\Phi(y_\alpha)-\Phi(y_0))$ belong to the image
under $\Phi$ of this subgroup modulo $q_0$. That image has at most
$q_0^{r_2(L)}T_1$ elements. There are at most $T_2$ fibres, so the
separation argument gives \eqref{eq:B6-few}.
\leanref{UpperTheory.\allowbreak{}V.\allowbreak{}ZK.\allowbreak{}zSeq}{Development.\allowbreak{}AutomorphicArithmetic.\allowbreak{}AutomorphicVarieties.\allowbreak{}ModularityArgument.\allowbreak{}ClassCounting.\allowbreak{}HasVal}{56--60}
\leanref{UpperTheory.\allowbreak{}V.\allowbreak{}ZK.\allowbreak{}BOdd}{Development.\allowbreak{}AutomorphicArithmetic.\allowbreak{}AutomorphicVarieties.\allowbreak{}ModularityArgument.\allowbreak{}ClassCounting.\allowbreak{}LemmaKOfSteps}{11--15}
\leanref{UpperTheory.\allowbreak{}V.\allowbreak{}ZK.\allowbreak{}selmer\_\allowbreak{}upper\_\allowbreak{}bound\_\allowbreak{}of\_\allowbreak{}inputs}{Development.\allowbreak{}AutomorphicArithmetic.\allowbreak{}AutomorphicVarieties.\allowbreak{}ModularityArgument.\allowbreak{}ClassCounting.\allowbreak{}SUpperOf}{21--62}
\leanref{UpperTheory.\allowbreak{}V.\allowbreak{}ZK.\allowbreak{}sUpperOf}{Development.\allowbreak{}AutomorphicArithmetic.\allowbreak{}AutomorphicVarieties.\allowbreak{}ModularityArgument.\allowbreak{}ClassCounting.\allowbreak{}SUpperOf}{64--66}
\leanref{UpperTheory.\allowbreak{}V.\allowbreak{}ZK.\allowbreak{}bOdd}{Development.\allowbreak{}AutomorphicArithmetic.\allowbreak{}AutomorphicVarieties.\allowbreak{}ModularityArgument.\allowbreak{}ClassCounting.\allowbreak{}BOdd}{3--3}
\leanref{LateralTheory.\allowbreak{}Q.\allowbreak{}SFewOfCut.\allowbreak{}sFewOf}{Development.\allowbreak{}EtaleCohomology.\allowbreak{}CyclotomicDiscriminants.\allowbreak{}SFewOfCut}{153--211}

Substitute \eqref{eq:B6-few} into \eqref{eq:B6-finite-class-count} to get
\begin{equation}\label{eq:B6-E-bound}
 \#\mathcal E_{\nu,k}
 \le T_1T_2\,\ell^{(b+v_\ell(k))g(L)+c+v_\ell(k)}.
\end{equation}
Take $c$ maximal with $\mu_{\ell^c}\subset L$. Put
$H_\nu=G_{L(\mu_{\ell^\nu})}$ and $H_\infty=\bigcap_\nu H_\nu$ inside
$G_L$. By \ref{B6.20}, restriction injects Sha into equivariant
characters of $H_\infty$ with the stated local conditions.

Let $\phi$ be the restriction of a Sha class to $H_\infty$.
Choose an open normal $N\subset G_L$ such that $H_\infty\cap N$
lies in its kernel. Extend it to $\widetilde\phi$ on $H_\infty N$
via $H_\infty N/N=H_\infty/(H_\infty\cap N)$. Compactness puts
$H_\nu$ inside $H_\infty N$ for large $\nu$. Inertia away from
$\ell$ already lies in $H_\infty$. In $H_\infty N$, the elements
$u$ belonging to a decomposition group above $\ell$ with
$\widetilde\phi(u)\ne0$ form a compact set. There are finitely many
places over $\ell$ in $L$, and the conjugates of each decomposition
group are parametrized by the compact group $G_L$. A second compact
set consists of pairs $(\gamma,u)\in G_L\times H_\infty N$ where
$\widetilde\phi(\gamma u\gamma^{-1})\ne
\chi_\ell(\gamma)^{-k}\widetilde\phi(u)$, with $\chi_\ell$ the
cyclotomic character. The defining conditions are locally constant.
The two sets miss $H_\infty$ and $G_L\times H_\infty$, respectively,
so both miss the corresponding finite-layer sets for large $\nu$.
The image is finite and is killed by $\ell^\nu$ after increasing $\nu$.

For a finite set of Sha classes take a common such $\nu\ge\max(c,2)$.
Class field theory gives $\varphi(P)=\widetilde\phi(\Frob_P)$ over
$F=L(\mu_{\ell^\nu})$ and the principal divisor relation.
Since $\Frob_{\sigma^{-1}P}$ is conjugate to $\sigma^{-1}\Frob_P\sigma$,
equivariance gives $\varphi(\sigma^{-1}P)=a\varphi(P)$. Restriction
recovers the original characters. Thus \eqref{eq:B6-E-bound} bounds
every finite subset. The Sha length $s$ satisfies
\begin{equation}\label{eq:B6-Sha-count}
 \ell^s\le T_1T_2\,\ell^{(b+v_\ell(k))g(L)+c+v_\ell(k)}.
\end{equation}
\leanref{UpperTheory.\allowbreak{}V.\allowbreak{}Kum.\allowbreak{}xdualExt}{Development.\allowbreak{}AutomorphicArithmetic.\allowbreak{}AutomorphicVarieties.\allowbreak{}GaloisLibrary.\allowbreak{}KummerTheory.\allowbreak{}XdualExt}{82--117}
\leanref{UpperTheory.\allowbreak{}V.\allowbreak{}Kum.\allowbreak{}eqExt}{Development.\allowbreak{}AutomorphicArithmetic.\allowbreak{}AutomorphicVarieties.\allowbreak{}GaloisLibrary.\allowbreak{}KummerTheory.\allowbreak{}CycZConjGal}{132--177}
\leanref{PositivityTheory.\allowbreak{}QuaternionShaBounds.\allowbreak{}ShaAssumptions.\allowbreak{}good\_\allowbreak{}of}{Development.\allowbreak{}AutomorphicArithmetic.\allowbreak{}AutomorphicVarieties.\allowbreak{}ModularityArgument.\allowbreak{}ClassCounting.\allowbreak{}RingEquivalences.\allowbreak{}OrderOfOneAddPrimePowZmodPrimePow}{390--410}
\leanref{PositivityTheory.\allowbreak{}QuaternionShaBounds.\allowbreak{}ShaAssumptions.\allowbreak{}card\_\allowbreak{}le}{Development.\allowbreak{}AutomorphicArithmetic.\allowbreak{}AutomorphicVarieties.\allowbreak{}ModularityArgument.\allowbreak{}ClassCounting.\allowbreak{}RingEquivalences.\allowbreak{}OrderOfOneAddPrimePowZmodPrimePow}{412--448}
\leanref{PositivityTheory.\allowbreak{}QuaternionShaBounds.\allowbreak{}ShaAsm\_\allowbreak{}holds}{Development.\allowbreak{}AutomorphicArithmetic.\allowbreak{}AutomorphicVarieties.\allowbreak{}ModularityArgument.\allowbreak{}ClassCounting.\allowbreak{}RingEquivalences.\allowbreak{}OrderOfOneAddPrimePowZmodPrimePow}{454--486}

Lemma~\ref{B6.16c} at $N=4k+6$ bounds $\log T_1$ and $\log T_2$
separately by the right side of \eqref{eq:B6-homology-torsion}.

To ensure \eqref{eq:B6-room}, pass to the cyclotomic layer $L'/L$
of index $m=\lceil\log_\ell(4k+4)\rceil$. Put $\delta=[L':L]$,
$n=[L:\Q]$, $n'=\delta n$, and take $c$ maximal with $\mu_{\ell^c}\subset L'$. Then
\[
 2k+1<\min(\ell^c,r_2(L')),\qquad
 \delta\le(n+1)(4k+4),\qquad \ell^c\le2n',\qquad
 |D_{L'}|\le |D_L|^\delta(\ell^c)^{n'}.
\]
For the last inequality use the derivative of $X^{\ell^c}-1$
to bound the relative different. Use $X=X_\ell(L)$ and the operator
$T=u^k\tau$ of Theorem~\ref{B6.20}. Since $T^{\ell^m}-1$ is divisible
by $T-1$, the quotient $X/(T^{\ell^m}-1)X$ surjects onto $X/(T-1)X$.
Thus $s_\ell(L,k)\le s_\ell(L',k)$.
\leanref{LateralTheory.\allowbreak{}Q.\allowbreak{}SLayerCut.\allowbreak{}goal\_\allowbreak{}of}{Development.\allowbreak{}EtaleCohomology.\allowbreak{}CyclotomicDiscriminants.\allowbreak{}SLayerCut}{55--107}

Put $v=2kn'$ and $H'=\log(2|D_{L'}|)$. The class-number bound gives
$g(L')\le2n'+4n'H'+2$. With $c,\log\ell,v_\ell(k)\le v$ and
$b\le v+2v^2$, \eqref{eq:B6-Sha-count} and \eqref{eq:B6-homology-torsion} give
\[
 s_\ell(L',k)\log\ell\le2\exp(v^{24}H')+v^9H'\le\exp(v^{26}H').
\]
Since $v\le(2kn)^5$ and $H'\le(2kn)^{11}\log(2|D_L|)$,
this gives the first assertion with $C=141$.
\end{proof}
\leanref{UpperTheory.\allowbreak{}V.\allowbreak{}ZK.\allowbreak{}lemmaK\_\allowbreak{}of\_\allowbreak{}steps}{Development.\allowbreak{}AutomorphicArithmetic.\allowbreak{}AutomorphicVarieties.\allowbreak{}ModularityArgument.\allowbreak{}ClassCounting.\allowbreak{}LemmaKOfSteps}{170--175}
\leanref{UpperTheory.\allowbreak{}V.\allowbreak{}ZK.\allowbreak{}lemmaK}{Development.\allowbreak{}AutomorphicArithmetic.\allowbreak{}AutomorphicVarieties.\allowbreak{}ModularityArgument.\allowbreak{}ClassCounting.\allowbreak{}LemmaKOfLitN}{28--31}

\begin{lemma}[Integral stability over the ring of integers]\label{B6.16a}
Let $L$ be a totally complex number field and $G_r=\GL_r(\mathcal O_L)$.
For $i,r,r'\in\N$ with $2i+4\le r\le r'$, adjoining an identity block
induces an isomorphism $H_i(G_r,\Z)\simeq H_i(G_{r'},\Z)$. The torsion
at rank $r'$ is finite and has order $\#H_i(G_r,\Z)_{\mathrm{tors}}$.
\end{lemma}
\begin{proof}
To see that $\mathcal O_L$ has Bass stable rank at most two, take a
unimodular row $(a_1,\ldots,a_t,b)$ with $t\ge2$. If all $a_j$ vanish,
adding the unit $b$ to the first entry suffices. Otherwise choose
$a_j\ne0$ and $h\ne j$. At each prime dividing $a_j$, prescribe $u=1$
modulo the prime if it divides $a_h$, and $u=0$ otherwise. The Chinese
remainder theorem gives such a $u$. Replace $a_h$ by $a_h+ub$.
A prime containing the first $t$ altered entries would divide $a_j$.
If it divided $a_h$, the relation $u\equiv1$ would force it to contain
$b$ and hence the whole original row. If it did not divide $a_h$, then
$u\equiv0$ would make the altered entry nonzero modulo that prime.
Thus the first $t$ entries are unimodular.

For $i\ge1$ and stable dimension at most one, van der Kallen's
Theorem~4.11(i) gives an isomorphism on integral homology from rank
$r$ to $r+1$ for $r\ge2i+1$ \cite{kallen-1980}. In degree zero,
stabilization induces the identity on $\Z$.
The archive's stability theorem applies for $r\ge2i+s+2$
over a commutative ring of Bass stable rank at most $s$. Take $s=2$
and compose the consecutive isomorphisms. They restrict to
isomorphisms on torsion, which is finite by Lemma~\ref{B6.16c}.
\end{proof}
\leanref{UpperTheory.\allowbreak{}V.\allowbreak{}KL.\allowbreak{}dedekindSR}{Development.\allowbreak{}AutomorphicArithmetic.\allowbreak{}AutomorphicVarieties.\allowbreak{}TensorHomCubes.\allowbreak{}DedekindSR}{50--111}
\leanref{GlobalTheory.\allowbreak{}HomStab.\allowbreak{}kStability\_\allowbreak{}main}{Development.\allowbreak{}AutomorphicArithmetic.\allowbreak{}AutomorphicVarieties.\allowbreak{}LiteratureResults.\allowbreak{}KTheoryStability.\allowbreak{}HomologicalStability.\allowbreak{}LinkExactnessStability.\allowbreak{}HomologicalStabilityTheorem}{10--27}
\leanref{UpperTheory.\allowbreak{}V.\allowbreak{}ZK.\allowbreak{}Up.\allowbreak{}torsStab}{Development.\allowbreak{}AutomorphicArithmetic.\allowbreak{}AutomorphicVarieties.\allowbreak{}ModularityArgument.\allowbreak{}ClassCounting.\allowbreak{}TorsionRankCounts.\allowbreak{}TorsStab}{16--39}

\begin{lemma}[An exact complex with bounded orbits]\label{B6.16b}
Let $L$ be totally complex, $n=[L:\Q]$, $D=|D_L|$, $N\in\N$ and
$G=\GL_N(\mathcal O_L)$. Put $d=nN$, $I=D^Nd^d$, $R=N^NI^2$ and
\[
 U_p=I^{1+dN}\bigl((2R+1)^d+1\bigr)^{dN(p+1)}.
\]
There are $G$-sets $X_p$ of ordered vertex tuples with deletion faces whose
augmented complex $\Z[X_\bullet]\to\Z$ is exact. In degree $p\ge0$ there
are at most $U_p$ orbits. Each stabilizer has order at most $3^{nN^2}$.
For $N\ge1$,
\[
 U_p\le (2Dd)^{12(p+1)d^3N}.
\]
\end{lemma}
\begin{proof}
For $N=0$ take tuples of the unique point, with $0^0=1$. Suppose $N\ge1$.
For a tuple $h=(h_w)$ of Hermitian forms at the complex places put
$h[v]=\sum_w v_w^*h_wv_w$. Let $P$ be cut out by $h[v]\ge1$ for
$0\ne v\in\mathcal O_L^N$. A large multiple of the standard form lies
in $P$, since the embedded lattice is discrete.

In an integral basis let $B$ be the real Gram matrix of $h\in P$.
Rational directions are dense, so $B$ is positive semidefinite.
For $\delta>0$ write $B+\delta\,1=S^tS$. Minkowski's box theorem gives
$0\ne z\in\Z^d$ with $\|Sz\|_\infty\le2|\det S|^{1/d}$.
Hence $1\le4d|\det S|^{2/d}$. Passing to $\delta=0$ gives
$(4d)^d\det B\ge1$. In particular all $h_w$ are positive definite.
The embedding matrix of an integral basis of $\mathcal O_L$ has
squared absolute determinant $D$ when both embeddings in each
conjugate pair are used. Passing to real and imaginary parts gives
real covolume $2^{-n/2}\sqrt D$. In $N$ copies, each complex Hermitian
determinant occurs twice in the real determinant, whence
\[
 2^d\det B=D^N\Bigl(\prod_w\det h_w\Bigr)^2,\qquad
 \prod_w\det h_w\ge I^{-1}.
\]
The inequality uses $D\ge1$ and $d\ge2$.

Let $X_p$ consist of the ordered $(p+1)$-tuples of vertices of $P$
at which a common real linear functional attains its minimum on $P$.
Require the functional to be positive on every nonzero positive
semidefinite tuple. Repetitions are allowed. Such a functional
dominates a positive multiple of the norm on the semidefinite cone
by compactness of its unit section. Its sublevels in $P$ are compact,
so it attains a minimum at a vertex. Near any $h\in P$ the forms
have a uniform positive lower eigenvalue, so finitely many lattice
inequalities can become active. At a vertex the active evaluations
span the dual space, since one could otherwise perturb in both
directions while preserving all inequalities.

For a finite cycle choose a witnessing functional for each tuple and
let $K$ be their convex hull. Fix $h_0\in P$. A vertex minimizing a
member of $K$ lies in at least one of the sublevels $f(h)\le f(h_0)$
for the chosen witnesses $f$, by taking their convex combination.
The compact union of these sublevels contains finitely many vertices,
since finitely many vectors can be minimal there. The minimizing cones
of these vertices intersect $K$ in a finite closed convex cover of
$K$. Its nerve has the reduced homology of $K$, which is zero. The
ordered-tuple chains give the same homology by the acyclic carrier
theorem. Over each simplex the tuple complex contracts by prepending
one of its vertices. Thus the cycle bounds. A minimizing vertex
gives surjectivity of the augmentation. Change of basis preserves
$P$ and the tuple condition.

At the average of a tuple the active vectors are exactly its common
minimal vectors. If these failed to span $L^N$, a nonzero $L$-linear
functional $\lambda$ vanishing on them would give a nonzero
semidefinite tuple $u=(|\lambda_w|^2)_w$. Subtracting a small multiple
of $u$ from the average stays in $P$ by local finiteness and decreases
the common minimizing functional. Choose $N$ independent common
minima as the columns of $M$. For any vertex $h$ of the tuple put
$H_w=M_w^*h_wM_w$ and $i=|N_{L/\Q}(\det M)|$. Each $H_w$ has diagonal
entries at most one. Hadamard's inequality and the determinant formula give
\[
 i/I\le\prod_w\det H_w=i\prod_w\det h_w\le1.
\]
Thus $1\le i\le I$ and each $\det H_w\ge I^{-1}$. The matrix
$C=iM^{-1}$ is integral, since $\mathcal O_L^N/M\mathcal O_L^N$
has order $i$. For fixed $i$, equal residues of $C,C'$ modulo $i$
give integral $M'M^{-1}=1+M'(C-C')/i$ and its inverse. Hence the
matrices of common minima have at most $\sum_{i=1}^I i^{nN^2}\le I^{1+dN}$
left $G$-classes.

Fix a representative $M$. Each $H_w$ has trace at most $N$, so
its least eigenvalue is at least $I^{-1}N^{1-N}$. For any minimal
$v$ at a vertex with these columns, $h[v]=1$ gives
$\|M_w^{-1}v_w\|_2^2\le IN^{N-1}$. Each conjugate coordinate of
$Cv$ therefore has absolute value at most $R$. Distinct algebraic
integers are at distance at least one in the maximum embedding norm
by the field norm. Packing open balls of radius $1/2$ in the real
space of dimension $n$ bounds the number per coordinate by $(2R+1)^n$.
Since $v\mapsto Cv$ is injective, at most $(2R+1)^d$ vectors can be
minimal at a vertex with these columns. A vertex is determined by at most
$dN$ independent active equations. Listing them with an empty slot for
padding gives at most $\bigl((2R+1)^d+1\bigr)^{dN}$ vertices per representative.
Counting the $p+1$ entries gives $U_p$. With $B_0=2Dd$, we have
$I\le B_0^{2d}$ and $R\le B_0^{5d}$. The exponent of $B_0$ is at most
\[
 2d(1+dN)+\bigl((5d+2)d+1\bigr)dN(p+1)\le12(p+1)d^3N.
\]

The stabilizer fixes the first form. Its columns are lattice points
of bounded length, so it is finite. A nontrivial kernel of reduction
modulo $3$ would contain an element of prime order $q$. Write it as
$1+3^aY$ with $a\ge1$ maximal. For $q\ne3$ its $q$th power differs
from $1$ by $q3^aY$ modulo $3^{a+1}$. For $q=3$ its cube differs
from $1$ by $3^{a+1}Y$ modulo $3^{a+2}$. Both contradict its order.
Thus reduction embeds the stabilizer in the reduction of
$M_N(\mathcal O_L)$, which has $3^{nN^2}$ elements.
\end{proof}
\leanref{UpperTheory.\allowbreak{}V.\allowbreak{}KT.\allowbreak{}sx\_\allowbreak{}exact}{Development.\allowbreak{}AutomorphicArithmetic.\allowbreak{}AutomorphicVarieties.\allowbreak{}HermiteNormalForms.\allowbreak{}Assembly}{14--17}
\leanref{UpperTheory.\allowbreak{}V.\allowbreak{}KT.\allowbreak{}gramdet}{Development.\allowbreak{}AutomorphicArithmetic.\allowbreak{}AutomorphicVarieties.\allowbreak{}HermiteNormalForms.\allowbreak{}Gramdet}{49--101}
\leanref{UpperTheory.\allowbreak{}V.\allowbreak{}KT.\allowbreak{}orb}{Development.\allowbreak{}AutomorphicArithmetic.\allowbreak{}AutomorphicVarieties.\allowbreak{}HermiteNormalForms.\allowbreak{}ColsMinSmul}{34--76}
\leanref{UpperTheory.\allowbreak{}V.\allowbreak{}KT.\allowbreak{}orbnum}{Development.\allowbreak{}AutomorphicArithmetic.\allowbreak{}AutomorphicVarieties.\allowbreak{}HermiteNormalForms.\allowbreak{}Orbnum}{80--92}
\leanref{UpperTheory.\allowbreak{}V.\allowbreak{}KT.\allowbreak{}stab}{Development.\allowbreak{}AutomorphicArithmetic.\allowbreak{}AutomorphicVarieties.\allowbreak{}HermiteNormalForms.\allowbreak{}ImportsAlt}{8--26}

\begin{lemma}[Torsion at a fixed rank]\label{B6.16c}
With the notation of Lemma~\ref{B6.16b}, let $m\ge0$, $A=U_m$ and
$s=3^{nN^2}$. The torsion subgroup of $H_m(G,\Z)$ is finite of order
$T\le(m+2)^A s^{mAs^m}$. If $k\ge1$, $1\le N\le4k+6$ and
$m\le2k+1$, then
\begin{equation}\label{eq:B6-homology-torsion}
 \log T\le\exp\!\bigl((2kn)^{24}\log(2D)\bigr).
\end{equation}
\end{lemma}
\begin{proof}
Write $C_p=\Z[X_p]$, $K_0=\Z$, and $K_{p+1}=\ker(C_p\twoheadrightarrow K_p)$.
Exactness gives $0\to K_{p+1}\to C_p\to K_p\to0$. For $p<m$, put
$j=m-p$. By Shapiro's lemma $H_j(G,C_p)$ is the sum of $H_j(S,\Z)$
over orbit stabilizers $S$. Transfer kills each summand by $|S|$,
since $j>0$. The bar cycles lie in a free abelian group of rank
$|S|^j$, so the summand needs at most $|S|^j$ generators and has
order at most $|S|^{|S|^j}$. There are at most $A$ orbits and so
$\#H_j(G,C_p)\le s^{As^j}$. In the exact sequence
\[
 H_j(G,C_p)\longrightarrow H_j(G,K_p)\longrightarrow H_{j-1}(G,K_{p+1})
\]
the torsion in the middle has at most $\#H_j(G,C_p)$ elements over
each torsion value on the right.

The group $H_0(G,K_m)=(K_m)_G$ is the cokernel of
$(C_{m+1})_G\to(C_m)_G$. The target is free on at most $A$ orbits.
Each column is the alternating sum of $m+2$ faces, so its
$\ell^1$-norm is at most $m+2$. If the columns span rank $t$, the
order of the torsion divides any nonzero $t$-rowed minor by Smith
normal form. Hadamard bounds the minor by $(m+2)^t$. Multiplying
the bounds along the $m$ exact sequences gives
$(m+2)^A s^{A\sum_{j=1}^m s^j}\le(m+2)^A s^{mAs^m}$.

Put $u=kn$, $H=\log(2D)$ and $a=12(m+1)d^3N\log(2Dd)$.
By Lemma~\ref{B6.16b}, $A\le e^a$. Taking logarithms of the order
bound and using $\log(m+2)\le m+1$, $\log s\le s$ gives
\[
 \log T\le A(2m+1)s^{m+1}
 \le\exp\!\bigl(a+2(m+1)nN^2+2m\bigr).
\]
Here $d,N\le10u$, $m+1\le4u$ and $\log(2Dd)=H+\log d\le16uH$
since $H\ge\log2$. Substitution gives
$12d^3N\log(2Dd)+2nN^2\le1\,920\,300u^5H$ and $2m\le12u^6H$.
The exponent is at most $7\,681\,212u^6H\le(2u)^{24}H$.
\end{proof}
\leanref{UpperTheory.\allowbreak{}V.\allowbreak{}KH.\allowbreak{}shift}{Development.\allowbreak{}AutomorphicArithmetic.\allowbreak{}AutomorphicVarieties.\allowbreak{}TorsionBoundsChains.\allowbreak{}LogLeExp}{14--47}
\leanref{UpperTheory.\allowbreak{}V.\allowbreak{}KT.\allowbreak{}tors\_\allowbreak{}le}{Development.\allowbreak{}AutomorphicArithmetic.\allowbreak{}AutomorphicVarieties.\allowbreak{}HermiteNormalForms.\allowbreak{}Assembly}{25--30}
\leanref{UpperTheory.\allowbreak{}V.\allowbreak{}ZK.\allowbreak{}torsAll}{Development.\allowbreak{}AutomorphicArithmetic.\allowbreak{}AutomorphicVarieties.\allowbreak{}ModularityArgument.\allowbreak{}ClassCounting.\allowbreak{}ExponentLe}{70--118}

\section{Generator bounds and quadratic extensions}

An abelian datum $M$ specifies $N_M\ge1$ and a subgroup of
$(\Z/N_M\Z)^\times$, with fixed field $K_M$ of degree $n_M$.
For primes $q,d$ with
\begin{equation}\label{eq:B6-prime-pair}
 q\equiv3\pmod8,\qquad d\equiv7\pmod{16},\qquad \gcd(qd,N_M)=1,
\end{equation}
put
\[
 K_{M,q,d}=K_M(i,\sqrt{-q},\sqrt{-d}),\qquad
 E_{M,a,q,d}=K_M\,\Q(\zeta_{2^{a+2}})^+
 \Q(\sqrt{-q})\,\Q(\sqrt{-d}).
\]
For $M=\operatorname{AbF.ofChar}(\psi)$ the second field is the field of
\ref{B6.02} at page index $a$.

Put $L_\infty=L_0(\mu_{\ell^\infty})$ for a base $L_0$ containing
the roots of unity of \ref{B6.15}. Let $\mathcal D_\ell(L_0)$ be the
continuous $\Q_\ell/\Z_\ell$-valued characters of $G_{L_\infty}$
unramified everywhere and split above $\ell$, with compact dual
$X_\ell(L_0)$. Write $\mathrm{XX}$ for the compact dual with splitting
omitted. Imposing splitting gives $\mathrm{XX}\twoheadrightarrow
\mathrm{XP}=X_\ell(L_0)$.

\begin{theorem}[Ferrero--Washington in finite layers]\label{B6.17}
Let $\ell$ be prime and $L/\Q$ finite abelian. Assume $\mu_\ell\subset L$
for odd $\ell$ and $i\in L$ for $\ell=2$. There is an integer $r\ge0$
such that for every $L'=L(\mu_{\ell^n})$, $n\ge0$,
\[
 \#\bigl(\mathrm{Cl}(L')/\mathrm{Cl}(L')^\ell\bigr)\le\ell^r.
\]
\end{theorem}
\leanref{UpperTheory.\allowbreak{}V.\allowbreak{}Published.\allowbreak{}FerreroWashington}{Development.\allowbreak{}AutomorphicArithmetic.\allowbreak{}AutomorphicVarieties.\allowbreak{}PublishedTheorems.\allowbreak{}FerreroWashington}{7--12}
\leanref{UpperTheory.\allowbreak{}V.\allowbreak{}FW.\allowbreak{}ferreroWashington}{Development.\allowbreak{}AutomorphicArithmetic.\allowbreak{}AutomorphicVarieties.\allowbreak{}GaloisLibrary.\allowbreak{}FerreroWashingtonIndex.\allowbreak{}FerreroWashington}{104--119}

Ferrero and Washington prove the vanishing of $\mu_\ell$ for cyclotomic
$\Z_\ell$-extensions of abelian fields \cite{ferrero-washington-1979}.
For the base fields in Theorem~\ref{B6.17}, this vanishing is
equivalent to a uniform bound on the $\ell$-ranks of the finite-layer
class groups.

\begin{theorem}[Uniform numbers of tower generators]\label{B6.18}
For odd prime $\ell$ and abelian datum $A$ with $\mu_\ell\subset K_A$,
some $R_\ell(A)\ge1$ bounds the $\Z_\ell$-generator numbers of both
tower modules for every realization of $A$.

For each $M$, some $R_0(M)\ge1$ bounds both generator numbers
at $2$ for $K_{M,q,d}$, uniformly in \eqref{eq:B6-prime-pair}.
If $g$ bounds the unramified mod-$2$ character dimension in the
infinite tower of $K_M(i)$, one may take $R_0(M)=4(g+6n_M)$.
\end{theorem}
\leanref{UpperTheory.\allowbreak{}V.\allowbreak{}Two.\allowbreak{}towerFixed}{Development.\allowbreak{}AutomorphicArithmetic.\allowbreak{}AutomorphicVarieties.\allowbreak{}QuadraticExtensions.\allowbreak{}IsAlgClosureRat}{98--111}
\leanref{UpperTheory.\allowbreak{}V.\allowbreak{}Two.\allowbreak{}towerTwo}{Development.\allowbreak{}AutomorphicArithmetic.\allowbreak{}AutomorphicVarieties.\allowbreak{}QuadraticExtensions.\allowbreak{}IsAlgClosureRat}{87--97}
\leanref{UpperTheory.\allowbreak{}V.\allowbreak{}Two.\allowbreak{}baseRk}{Development.\allowbreak{}AutomorphicArithmetic.\allowbreak{}AutomorphicVarieties.\allowbreak{}QuadraticExtensions.\allowbreak{}IsAlgClosureRat}{69--86}

\begin{proof}
By \ref{B6.17} and class field theory the finite-layer unramified
mod-$\ell$ character groups have size at most $\ell^r$. Their
restriction images increase and exhaust the infinite-layer group by
the finite-quotient extension and inertia compactness of \ref{B6.16}.
Compact Nakayama gives at most $r$ generators of $\mathrm{XX}$ and
its quotient $\mathrm{XP}$, so take $R_\ell(A)=\max(1,r)$.

Allow ramification above $q,d$ in the tower over $K_M(i)$.
Each place adds at most one mod-$2$ inertia dimension by tameness.
The congruences give decomposition indices at most $2,4$ over
$\Q(i)$, hence at most $6n_M$ places after base change.

At a quadratic adjunction, write $U\subset V$ for the index-two
Galois subgroups and choose $\sigma\in V\setminus U$. An invariant
character $\phi$ of $U$ extends precisely if $\phi(\sigma^2)=0$.
The obstruction has at most two values and a successful extension
has two choices for the value on $\sigma$. Both preserve the
ramification conditions, since the excluded inertia lies in $U$.
Thus the invariants on $U$ have dimension at most that on $V$.
Conjugation satisfies $(\iota-1)^2=0$, so the whole space on $U$
has at most twice that dimension. The two adjunctions give
$4(g+6n_M)$ for the unramified subspace and hence for both duals.
\end{proof}

\begin{lemma}[The cost of a quadratic extension]\label{B6.19}
If $E\subset E'$ are number fields with $[E':E]\mid2$, then for every
integer $k$,
\begin{equation}\label{eq:B6-quadratic-loss}
 s_2(E,k)\le1+s_2(E',k).
\end{equation}
For $a\ge2$ and a prime pair \eqref{eq:B6-prime-pair}, every realization
$E$ of $E_{M,a,q,d}$ admits a layer $L_m/K_{M,q,d}$ with $m\le a+2$
and $s_2(E,k)\le1+s_2(L_m,k)$ for every $k\ge1$. The base field has
\[
 [K_{M,q,d}:\Q]\le8n_M,\qquad
 K_{M,q,d}\hookrightarrow\Q(\zeta_{4N_Mqd}).
\]
\end{lemma}
\leanref{UpperTheory.\allowbreak{}V.\allowbreak{}Two.\allowbreak{}shaTwo}{Development.\allowbreak{}AutomorphicArithmetic.\allowbreak{}AutomorphicVarieties.\allowbreak{}QuadraticExtensions.\allowbreak{}ShaTwo}{11--28}
\leanref{UpperTheory.\allowbreak{}V.\allowbreak{}Two.\allowbreak{}twoBase}{Development.\allowbreak{}AutomorphicArithmetic.\allowbreak{}AutomorphicVarieties.\allowbreak{}QuadraticExtensions.\allowbreak{}IsAlgClosureRat}{112--129}

\begin{proof}
Restriction preserves the Sha conditions. For a quadratic extension
its kernel on $H^1$ is $\ker(1+\sigma:A\to A)/(\sigma-1)A$ by
inflation--restriction, with $A=(\Q_2/\Z_2(-k))^{G_{E'}}$ and
$\sigma$ generating the quotient. This subquotient of a locally
cyclic group is killed by $2$, so has length at most one.

Adjoin $\mu_{2^{a+2}}$ to $K_{M,q,d}$ to get $L_m$ with $m\le a+2$.
Under a common embedding, $E\subset L_m$ and $[L_m:E]\mid2$ by
the index of the real cyclotomic subfield. The three quadratic
adjunctions give the base degree and the displayed embedding.
\end{proof}

\section{Twisted cokernels and their centres}

For a topological generator $\sigma$ of $\Gal(L_\infty/L_0)$,
put $u=\chi_\ell(\sigma)$ and write $\tau$ for its action on $X_\ell(L_0)$.
Then
\begin{equation}\label{eq:B6-unit-level}
 v_\ell(u-1)=c,\qquad c\ge1,\qquad c\ge2\ \text{if }\ell=2,
\end{equation}
with $\mu_{\ell^c}\subset L_0$. For an endomorphism $\tau$ of a
$\Z_\ell$-module $X$, put
\[
 w_{\tau,u}(k,m)=\operatorname{length}_{\Z_\ell}
 X/((u^k\tau)^{\ell^m}-1)X.
\]

\begin{theorem}[Descent to a twisted cokernel]\label{B6.20}
Let $L_0$ contain the roots of unity in \ref{B6.15}. For the tower
module and generator above, every layer $L_m/L_0$ and every integer
$k\ge1$ satisfy
\begin{equation}\label{eq:B6-cokernel-descent}
 s_\ell(L_m,k)=w_{\tau,u}(k,m).
\end{equation}
\end{theorem}
\leanref{UpperTheory.\allowbreak{}V.\allowbreak{}PropA.\allowbreak{}descends}{Development.\allowbreak{}AutomorphicArithmetic.\allowbreak{}AutomorphicVarieties.\allowbreak{}AdditiveValuationNorms.\allowbreak{}Descends}{8--15}
\leanref{UpperTheory.\allowbreak{}V.\allowbreak{}PropA.\allowbreak{}Descends}{Development.\allowbreak{}AutomorphicArithmetic.\allowbreak{}AutomorphicVarieties.\allowbreak{}AdditiveValuationNorms.\allowbreak{}ABoundOfDescends}{16--19}
\leanref{UpperTheory.\allowbreak{}V.\allowbreak{}PropA.\allowbreak{}twLen}{Development.\allowbreak{}AutomorphicArithmetic.\allowbreak{}AutomorphicVarieties.\allowbreak{}AdditiveValuationNorms.\allowbreak{}CokLenLeOfRangeLe}{26--27}

\begin{proof}
Write $\Gamma_m=\Gal(L_\infty/L_m)$ and $A=\Q_\ell/\Z_\ell(-k)$.
On $A$, the generator of $\Gamma_m$ minus the identity acts by
the surjective scalar $u^{-k\ell^m}-1$. Hence $H^1(\Gamma_m,A)=0$,
and $H^2(\Gamma_m,A)=0$ by cohomological dimension one. Restriction
therefore identifies global classes with equivariant characters on
$G_{L_\infty}$, where the coefficients are trivial. Away from $\ell$,
inertia lies in this subgroup. Above $\ell$, a class whose
character vanishes on the kernel in a decomposition group
descends locally to an open subgroup of $\Gamma_m$. Its $H^1$
vanishes by the same nonzero-scalar argument. Restriction therefore
identifies $\Sha_{\ell,k}(L_m)$ with $(\mathcal D_\ell(L_0)(-k))^{\Gamma_m}$.
Duality preserves length and turns invariants into coinvariants for
$(u^k\tau)^{\ell^m}$.
\end{proof}
\leanref{UpperTheory.\allowbreak{}V.\allowbreak{}Kum.\allowbreak{}resInfOnto}{Development.\allowbreak{}AutomorphicArithmetic.\allowbreak{}AutomorphicVarieties.\allowbreak{}GaloisLibrary.\allowbreak{}KummerTheory.\allowbreak{}ResInfOnto}{102--166}

\begin{theorem}[Centres of twisted cokernels]\label{B6.21}
Let $X$ have at most $R$ generators over $\Z_\ell$, let
$\tau$ be a $\Z_\ell$-linear endomorphism, and let $u,c$ satisfy
\eqref{eq:B6-unit-level}. For each $j\ge0$ there are $s_i\in\Z_\ell$
for $1\le i\le R$ such that
\begin{equation}\label{eq:B6-centres-fixed}
 w_{\tau,u}(k,j)\le\sum_{i=1}^R(j+c+v_\ell(s_i-k))
 \qquad(k\in\N),
\end{equation}
where $v_\ell(0)=\infty$. If $\tau$ is unipotent modulo $\ell$,
some $s_1,\ldots,s_R\in\Z_\ell$ satisfy
\begin{equation}\label{eq:B6-centres-all}
 w_{\tau,u}(k,m)\le\sum_{i=1}^R(\max(m,j_0)+c+v_\ell(s_i-k))
 \qquad(k,m\in\N),\qquad j_0=j_0(\ell,R).
\end{equation}
If $|s_i-k|_\ell\ge\ell^{-b}$ for every $i$, with $b\in\N$, then
\begin{equation}\label{eq:B6-centres-far}
 w_{\tau,u}(k,m)\le R(\max(m,j_0)+c+b).
\end{equation}
\end{theorem}
\leanref{UpperTheory.\allowbreak{}V.\allowbreak{}PropA.\allowbreak{}exists\_\allowbreak{}centres\_\allowbreak{}level}{Development.\allowbreak{}AutomorphicArithmetic.\allowbreak{}AutomorphicVarieties.\allowbreak{}AdditiveValuationNorms.\allowbreak{}ExistsCentresLevel}{23--33}
\leanref{UpperTheory.\allowbreak{}V.\allowbreak{}PropA.\allowbreak{}exists\_\allowbreak{}centres}{Development.\allowbreak{}AutomorphicArithmetic.\allowbreak{}AutomorphicVarieties.\allowbreak{}AdditiveValuationNorms.\allowbreak{}ExistsCentresLevel}{34--50}
\leanref{UpperTheory.\allowbreak{}V.\allowbreak{}PropA.\allowbreak{}exists\_\allowbreak{}centres\_\allowbreak{}norm}{Development.\allowbreak{}AutomorphicArithmetic.\allowbreak{}AutomorphicVarieties.\allowbreak{}AdditiveValuationNorms.\allowbreak{}ExistsCentresLevel}{57--70}

\begin{proof}
Lift $\tau^{\ell^j}$ to a matrix $S$ on $\Z_\ell^R$ along a
surjection to $X$. For $t=u^{k\ell^j}$ the cokernel of $tS-1$
surjects onto the one counted by $w_{\tau,u}(k,j)$. Its length
is $\sum_i v_\ell(tz_i-1)$ by Smith normal form, where the
$z_i$ are the eigenvalues of $S$ with multiplicity.

For units $A,B$ and algebraic $z$,
$|A-B|_\ell\le\max(|Az-1|_\ell,|Bz-1|_\ell)$ by the ultrametric
inequality for $|z|_\ell\ge1$, and both right-hand norms are $1$
for $|z|_\ell<1$. Lifting the exponent gives
\[
 |(u^{\ell^j})^a-(u^{\ell^j})^b|_\ell
 =\ell^{-(j+c)}|a-b|_\ell\qquad(a,b\in\N).
\]
For $z=z_i$ the balls centred at $a\in\N$ with radii
$r_a=\ell^{j+c}|u^{a\ell^j}z-1|_\ell$ are nested by
$|a-b|_\ell\le\max(r_a,r_b)$. Compactness in $\Z_\ell$ gives
$s_i$ in their intersection. Taking $a=k$ gives
\eqref{eq:B6-centres-fixed}.

Put $g=u^k\tau$ and $h=\lfloor\log_\ell R\rfloor+1$.
Modulo $\ell$, $g-1$ is nilpotent of index at most $R$ if $R>0$,
and $X=0$ if $R=0$. The binomial theorem then gives
$(g^{\ell^h}-1)X\subseteq\ell X$ and
$(g^{\ell^{h+1}}-1)X\subseteq\ell^2 X$. If $(w-1)X\subseteq\ell^2 X$,
then $(w^\ell-1)=(w-1)(1+\cdots+w^{\ell-1})$. The second factor
has image $\ell X$ by Nakayama, since it agrees with $\ell$
modulo $\ell^2 X$. Its cokernel has length at most $R$.
Iterate from $j_0=h+2$ and use divisibility for $m\le j_0$ to get
\begin{equation}\label{eq:B6-growth}
 w_{\tau,u}(k,m)\le w_{\tau,u}(k,j_0)+R(m-j_0)_+.
\end{equation}
Use \eqref{eq:B6-centres-fixed} at $j_0$ and then bound each valuation
by $b$ under the norm condition.
\end{proof}

\begin{theorem}[The two tower bounds]\label{B6.22}
Let $L_0$ be totally complex with the roots of unity of \ref{B6.15},
and let $X_\ell(L_0)$ have at most $R$ generators, with $R\ge1$.
Then \eqref{eq:B6-centres-far} holds for the Sha lengths by
\ref{B6.20}. For $k\ge1$ and any layers of indices $m$ and
$j_0=j_0(\ell,R)$,
\begin{equation}\label{eq:B6-second-bound}
 s_\ell(L_m,k)\le V_\ell(L_{j_0},k)+R(m-j_0)_+.
\end{equation}
\end{theorem}
\leanref{UpperTheory.\allowbreak{}V.\allowbreak{}PropA.\allowbreak{}aBound\_\allowbreak{}of\_\allowbreak{}isTower}{Development.\allowbreak{}AutomorphicArithmetic.\allowbreak{}AutomorphicVarieties.\allowbreak{}AdditiveValuationNorms.\allowbreak{}ABoundOfIsTower}{7--11}
\leanref{UpperTheory.\allowbreak{}V.\allowbreak{}PropA.\allowbreak{}aBound\_\allowbreak{}of\_\allowbreak{}descends}{Development.\allowbreak{}AutomorphicArithmetic.\allowbreak{}AutomorphicVarieties.\allowbreak{}AdditiveValuationNorms.\allowbreak{}ABoundOfDescends}{36--56}
\leanref{UpperTheory.\allowbreak{}V.\allowbreak{}Win.\allowbreak{}SecondBound}{Development.\allowbreak{}AutomorphicArithmetic.\allowbreak{}AutomorphicVarieties.\allowbreak{}LayeredIsomorphisms.\allowbreak{}SecondBound}{13--18}

\begin{proof}
Put $X=X_\ell(L_0)$ and $\Gamma=\Gal(L_\infty/L_0)\simeq\Z_\ell$.
The continuous action of $\Gamma$ on the finite space $X/\ell X$
factors through a cyclic group of order $\ell^a$ for some $a\ge0$.
Hence $(\tau-1)^{\ell^a}=\tau^{\ell^a}-1=0$ on $X/\ell X$.
Apply \ref{B6.20} and \ref{B6.21}, then use
$s_\ell\le V_\ell$ from \eqref{eq:B6-H2-local-correction}.
Layers of equal index are isomorphic over $L_0$.
\end{proof}
\leanref{UpperTheory.\allowbreak{}V.\allowbreak{}PropA.\allowbreak{}shaLeV}{Development.\allowbreak{}AutomorphicArithmetic.\allowbreak{}AutomorphicVarieties.\allowbreak{}AdditiveValuationNorms.\allowbreak{}ShaLeV}{7--16}

\section{Exceptional twists}

\begin{lemma}[At most $R$ exceptions below each cutoff]\label{B6.23}
Fix the centres of \ref{B6.22} and the corresponding level $c$.
For each $Z\in\N$ some $\mathcal E_Z\subseteq\{1,\ldots,Z\}$ has
size at most $R$. For $1\le k\le Z$ outside it and any layer
$L_m/L_0$,
\begin{equation}\label{eq:B6-first-bound}
 s_\ell(L_m,k)\le R\bigl(\max(m,j_0(\ell,R))+c+\lfloor\log_\ell Z\rfloor\bigr).
\end{equation}
\end{lemma}
\leanref{UpperTheory.\allowbreak{}V.\allowbreak{}Win.\allowbreak{}exc}{Development.\allowbreak{}AutomorphicArithmetic.\allowbreak{}AutomorphicVarieties.\allowbreak{}LayeredIsomorphisms.\allowbreak{}Exc}{7--42}
\leanref{UpperTheory.\allowbreak{}V.\allowbreak{}Win.\allowbreak{}FirstBound}{Development.\allowbreak{}AutomorphicArithmetic.\allowbreak{}AutomorphicVarieties.\allowbreak{}LayeredIsomorphisms.\allowbreak{}SecondBound}{27--31}

\begin{proof}
For $Z>0$, put $b=\lfloor\log_\ell Z\rfloor$ and exclude those
$1\le k\le Z$ with $|s_i-k|_\ell<\ell^{-b}$ for any $i$. Each
centre excludes at most one integer, since $\ell^{b+1}>Z$.
Apply \eqref{eq:B6-centres-far}. For $Z=0$ take the empty set.
\end{proof}

\begin{lemma}[A common interval for finitely many exceptions]\label{B6.24}
Let $R\ge0$ be an integer, $0<\eta\le C_t$, $X>1$, and
$\mathcal E\subset\N$ with $\#\mathcal E\le R$. Put
$J=2^{R+2}(C_t/\eta)^R$. There is $Y$ with $X^{1/J}\le Y$ and $Y^2\le X$
such that, for every $Y\le x\le Y^2$,
\begin{equation}\label{eq:B6-good}
 k\in\mathcal E,\quad k\le x^{C_t}\quad\Longrightarrow\quad k\le x^\eta.
\end{equation}
\end{lemma}
\leanref{UpperTheory.\allowbreak{}V.\allowbreak{}Win.\allowbreak{}window1}{Development.\allowbreak{}AutomorphicArithmetic.\allowbreak{}AutomorphicVarieties.\allowbreak{}LayeredIsomorphisms.\allowbreak{}Window1}{6--58}
\leanref{UpperTheory.\allowbreak{}V.\allowbreak{}Win.\allowbreak{}window}{Development.\allowbreak{}AutomorphicArithmetic.\allowbreak{}AutomorphicVarieties.\allowbreak{}LayeredIsomorphisms.\allowbreak{}Window}{7--43}

\begin{proof}
Put $U=\log X$, $B=C_t/\eta$ and
$T=\{\log k/C_t:k\in\mathcal E,\ k\ge1\}$. Seek $y$ with
$U/(2^{R+2}B^R)\le y\le U/2$ such that
\begin{equation}\label{eq:B6-log-window}
 a\in T,\quad y\le u\le2y,\quad a\le u\quad\Longrightarrow\quad Ba\le u.
\end{equation}
If $T$ misses $(U/(2B),U]$, take $y=U/2$. Otherwise remove a
point of $T$ in that interval and apply induction on $R$ with
upper endpoint $U/(2B)$. The new upper bound for $2y$ excludes
the removed point, and the lower bound for $y$ is as required.
Take $Y=e^y$ and put $u=\log x$ in \eqref{eq:B6-log-window}.
For $k=0$, $0\le x^\eta$.
\end{proof}

\section{Small exceptions and the pointwise budget}

Write $P_{C,\varepsilon}(x)=C(1+\log x)+2\exp(Cx^{3\varepsilon/4})$.
Take $C_K>0$ as in \ref{B6.15} and write $\eta=\varepsilon/(2C_K)$.
Use $x\in[X^{1/J},X]$ and cutoff $Z=\lfloor X^{C_t}\rfloor$.

\begin{lemma}[A pointwise bound and its sum]\label{B6.25}
Fix a prime $\ell$, $R,c_1,n_0\in\N$, $c_0,C_t,A_0\ge0$, and
$J,\varepsilon>0$. There is $C\ge0$ depending on these and $C_K$
with the following property. Assume the roots of unity of \ref{B6.15}
and \eqref{eq:B6-second-bound} for $L_0$, and
\eqref{eq:B6-first-bound} for $\mathcal E$ at $Z$ with $c\le c_1$.
Let $[L_{j_0}:\Q]\le n_0$. For $x\ge1$ in the stated range with
\eqref{eq:B6-good} and $\log(2|D_{L_{j_0}}|)\le c_0x^{\varepsilon/4}$,
each $1\le k\le x^{C_t}$ and $m\le A_0(1+\log x)$ satisfy
\begin{equation}\label{eq:B6-point-bound}
 s_\ell(L_m,k)<\infty,\qquad s_\ell(L_m,k)+1\le P_{C,\varepsilon}(x).
\end{equation}
For every $C_t,C\in\R$ and $\varepsilon>0$, some $x_0\ge2$ satisfies
\begin{equation}\label{eq:B6-final-absorption}
 x^{C_t}P_{C,\varepsilon}(x)\le D_x\qquad(x\ge x_0).
\end{equation}
\end{lemma}
\leanref{UpperTheory.\allowbreak{}V.\allowbreak{}Win.\allowbreak{}point2}{Development.\allowbreak{}AutomorphicArithmetic.\allowbreak{}AutomorphicVarieties.\allowbreak{}LayeredIsomorphisms.\allowbreak{}Point2}{8--49}
\leanref{UpperTheory.\allowbreak{}V.\allowbreak{}Win.\allowbreak{}final}{Development.\allowbreak{}AutomorphicArithmetic.\allowbreak{}AutomorphicVarieties.\allowbreak{}LayeredIsomorphisms.\allowbreak{}Final}{6--68}

\begin{proof}
We have $k\le Z$. Outside $\mathcal E$ use \eqref{eq:B6-first-bound}.
Inside $\mathcal E$, $k\le x^\eta$ gives
\[
 (2k[L_{j_0}:\Q])^{C_K}\log(2|D_{L_{j_0}}|)
 \le (2n_0)^{C_K}c_0x^{3\varepsilon/4}.
\]
Apply \ref{B6.15}, divide by $\log\ell>1/2$, and use
\eqref{eq:B6-second-bound}. Both cases give
\begin{equation}\label{eq:B6-point-before-logs}
 s_\ell(L_m,k)\le R(m+j_0+c+\lfloor\log_\ell Z\rfloor)
 +2\exp\bigl((2n_0)^{C_K}c_0x^{3\varepsilon/4}\bigr).
\end{equation}
Use $\lfloor\log_\ell Z\rfloor\le2C_tJ\log x$ and take
$C=R(A_0+j_0+c_1+2C_tJ)+1+(2n_0)^{C_K}c_0$.
The product is bounded above for $x\ge2$ by
\[
 3\exp\!\left((|C_t|+|C|)\log x+|C|+|C|x^{3\varepsilon/4}\right).
\]
The logarithm of this bound is $o(x^{4\varepsilon/5})$.
\end{proof}

\begin{theorem}[Bounds at $2$ and at a fixed odd prime]\label{B6.26}
Fix $J,\varepsilon>0$, $C_t\ge0$ and $A_1\in\R$.
Call $a,t:\R\to\N$ admissible if for every $x\ge2$
\begin{equation}\label{eq:B6-admitted}
 a(x)\ge2,\qquad 2^{a(x)}\le x^{A_1},\qquad 2\le t(x)\le x^{C_t}.
\end{equation}
Each $x_0$ is uniform in admissible $a,t$ and $X\ge2$.
The bounds hold for $x\in[X^{1/J},X]$, $x\ge x_0$, satisfying
\eqref{eq:B6-good} for the corresponding exceptional set.

At $2$, fix $M$ as in \ref{B6.18}. There is $x_0\ge2$ such that
for every pair \eqref{eq:B6-prime-pair} with
$\log(qd)\le X^{\varepsilon/(4J)}$, a set $\mathcal E_2$ of size at
most $R_0(M)$ gives, for every realization $E$ of $E_{M,a(x),q,d}$,
\begin{equation}\label{eq:B6-two-sum}
 s_2(E,k)<\infty\quad(1\le k<t(x)),\qquad
 \sum_{1\le k<t(x)}s_2(E,k)\le D_x.
\end{equation}

At an odd prime $\ell$, fix $A$ with $\mu_\ell\subset K_A$.
Put $R_\ell=R_\ell(A)$ and $j_\ell=j_0(\ell,R_\ell)$.
There is $x_0\ge2$ such that for each $a,t,X$ some $\mathcal E_\ell$
of size at most $R_\ell$ gives, for every realization $L_0$ of $A$
and layer $L_{j_\ell}/L_0$,
\begin{equation}\label{eq:B6-odd-layer-sum}
 s_\ell(L_{j_\ell},k)<\infty\quad(1\le k<t(x)),\qquad
 \sum_{1\le k<t(x)}s_\ell(L_{j_\ell},k)\le D_x.
\end{equation}
\end{theorem}
\leanref{UpperTheory.\allowbreak{}V.\allowbreak{}Win.\allowbreak{}partA}{Development.\allowbreak{}AutomorphicArithmetic.\allowbreak{}AutomorphicVarieties.\allowbreak{}LayeredIsomorphisms.\allowbreak{}PartA}{24--78}
\leanref{UpperTheory.\allowbreak{}V.\allowbreak{}Win.\allowbreak{}partB}{Development.\allowbreak{}AutomorphicArithmetic.\allowbreak{}AutomorphicVarieties.\allowbreak{}LayeredIsomorphisms.\allowbreak{}PartB}{50--95}

\begin{proof}
At $2$, take $L_0=K_{M,q,d}$ and $j=j_0(2,R_0(M))$.
For constants $n_0,c_0$ depending on $M,j$,
\begin{equation}\label{eq:B6-layer-disc}
 [L_j:\Q]\le n_0,\qquad
 \log(2|D_{L_j}|)\le c_0w\quad(w\ge1,\ \log(qd)\le w).
\end{equation}
Take $n_0=8n_M2^j$. If $L_j=L_0(\mu_{2^r})$ with $r$
minimal, then $2^r\le2n_0$ and $L_j$ embeds in a cyclotomic field
of conductor dividing $2^r(4N_Mqd)$. The discriminant formula
in towers gives $\log(2|D_{L_j}|)\le\log2+n_0\log(8n_0N_Mqd)$.

Take $\mathcal E_2$ by \ref{B6.18}, \ref{B6.22} and \ref{B6.23}.
For the stated range of $x$, $\log(qd)\le x^{\varepsilon/4}$.
Since $2^{c-1}\le[L_0:\Q]$, $c\le16n_M$. By
\ref{B6.19} and admissibility, each $E$ has a layer with
$m\le a(x)+2\le(2|A_1|+2)(1+\log x)$ and
$s_2(E,k)\le s_2(L_m,k)+1$. Apply \ref{B6.25} with these
uniform data.

At an odd prime, choose one realization and take $\mathcal E_\ell$
at the same cutoff by \ref{B6.23}. Use \ref{B6.25} with
$n_0=[L_{j_\ell}:\Q]$, $c_0=\log(2|D_{L_{j_\ell}}|)$ and $A_0=j_\ell$.
Transport by layer isomorphism gives the same set and bound for
every realization. In each case sum over at most $x^{C_t}$ terms
and use \eqref{eq:B6-final-absorption}.
\end{proof}
\leanref{UpperTheory.\allowbreak{}V.\allowbreak{}Win.\allowbreak{}layeremb}{Development.\allowbreak{}AutomorphicArithmetic.\allowbreak{}AutomorphicVarieties.\allowbreak{}LayeredIsomorphisms.\allowbreak{}Layeremb}{7--34}
\leanref{UpperTheory.\allowbreak{}V.\allowbreak{}Win.\allowbreak{}disc}{Development.\allowbreak{}AutomorphicArithmetic.\allowbreak{}AutomorphicVarieties.\allowbreak{}LayeredIsomorphisms.\allowbreak{}Imports}{6--44}

\section{Selection of the common window}

Use $C_K$ from \ref{B6.15} and $R_0(M),R_\ell(A)$ from \ref{B6.18},
with $R_\ell(A)=0$ outside the odd-prime case.

\begin{theorem}[The common Iwasawa window]\label{B6.27}
Fix $C_t\ge1$, an abelian datum $M$, a finite set $B'$ of odd
primes and, for each $\ell\in B'$, a datum $A_\ell$ with
$\mu_\ell\subset K_{A_\ell}$. Fix $0<\varepsilon<1/52$ and
$A_1\in\R$. Put
\[
 R=R_0(M)+\sum_{\ell\in B'}R_\ell(A_\ell),\qquad
 J=2^{R+2}(C_t/\eta)^R.
\]
There is $X_0$ such that for every admissible $a,t$, $X\ge X_0$,
and pair \eqref{eq:B6-prime-pair} with
\begin{equation}\label{eq:B6-qd-window-size}
 \log(qd)\le X^{\varepsilon/(4J)},
\end{equation}
there is $Y$ with $X^{1/J}\le Y$ and $Y^2\le X$. For every
$x\in[Y,Y^2]$, \eqref{eq:B6-two-sum} holds for every realization
$E$ of $E_{M,a(x),q,d}$. For each $\ell\in B'$,
\eqref{eq:B6-odd-layer-sum} holds for every realization of $A_\ell$
and its layers of index $j_0(\ell,R_\ell(A_\ell))$.
\end{theorem}
\leanref{UpperTheory.\allowbreak{}V.\allowbreak{}Win.\allowbreak{}thmBab\_\allowbreak{}of\_\allowbreak{}lemmaK}{Development.\allowbreak{}AutomorphicArithmetic.\allowbreak{}AutomorphicVarieties.\allowbreak{}LayeredIsomorphisms.\allowbreak{}ThmBabOfLemmaK}{9--13}
\leanref{UpperTheory.\allowbreak{}V.\allowbreak{}Win.\allowbreak{}thmBab\_\allowbreak{}of}{Development.\allowbreak{}AutomorphicArithmetic.\allowbreak{}AutomorphicVarieties.\allowbreak{}LayeredIsomorphisms.\allowbreak{}ThmBabOf}{11--103}

\begin{proof}
Let $x_1$ be the maximum of $2$ and the thresholds of \ref{B6.26}
and set $X_0=\max(2,x_1^J)$. The union
$\mathcal E_2\cup\bigcup_{\ell\in B'}\mathcal E_\ell$ has size at
most $R$. Since $0<\eta\le C_t$, \ref{B6.24} supplies $Y$ with
\eqref{eq:B6-good} for this union throughout $[Y,Y^2]$.
Here $x\ge X^{1/J}\ge x_1$, so \ref{B6.26} gives the two bounds.
\end{proof}

At $C_t=401$, Theorem~\ref{B6.27} supplies the cyclotomic window estimate
used in Proposition~\ref{B6.07}. That proposition combines it with
the automorphic counts to obtain $\mathrm{VWin}$.
\leanref{UpperTheory.\allowbreak{}FinalTheorems.\allowbreak{}everyCharacter\_\allowbreak{}of\_\allowbreak{}row}{Development.\allowbreak{}AutomorphicArithmetic.\allowbreak{}CharacterBoxes.\allowbreak{}MainConclusion}{7--8}
\leanref{UpperTheory.\allowbreak{}V.\allowbreak{}HV.\allowbreak{}DescentVanishingReductions.\allowbreak{}vwin\_\allowbreak{}of\_\allowbreak{}hypotheses\_\allowbreak{}H\_\allowbreak{}E\_\allowbreak{}alt}{Development.\allowbreak{}AutomorphicArithmetic.\allowbreak{}AutomorphicVarieties.\allowbreak{}ModularityArgument.\allowbreak{}ThresholdArguments.\allowbreak{}DenseThresholdCases.\allowbreak{}WindowBoundOfLiteratureLevel}{25--29}
\leanref{UpperTheory.\allowbreak{}V.\allowbreak{}HV.\allowbreak{}Core.\allowbreak{}of\_\allowbreak{}four\_\allowbreak{}alt}{Development.\allowbreak{}AutomorphicArithmetic.\allowbreak{}AutomorphicVarieties.\allowbreak{}ModularityArgument.\allowbreak{}ThresholdArguments.\allowbreak{}DenseThresholdCases.\allowbreak{}HVCore}{18--22}

For a primitive Dirichlet character $\psi$, use
$M=\operatorname{AbF.ofChar}(\psi)$, $B'=\mathcal B_\psi$ from
\ref{B6.28}, and $A_\ell=\operatorname{AbF.cyc}(\ell\,\operatorname{cond}\psi)$.
The bounds supply the term at $2$ in \ref{B6.14} and the fixed-layer
estimates for \ref{B6.32}, where the odd-prime lengths acquire factors
$\log\ell$.


% ---- p6c4 ----
\chapter{Summing the plus component by classes}\label{ch:b6-classwise}

Write $\chi$ for the fixed primitive character $\psi$, of conductor $f$ and order $o$. For a CM field $E$, an odd prime $\ell$, and a finite set $S$ of odd primes, put
\[
 s^+_{\ell,k}(E)=\operatorname{length}_{\Z_\ell}\Sha^+_{\ell,k}(E),
 \qquad
 P^+(E,t;S)=\sum_{\ell\in S}\log\ell
 \sum_{1\le k<t}s^+_{\ell,k}(E).
\]
The plus component has complex conjugation eigenvalue $(-1)^k$. All bounds use the algebraic closures and local data of \ref{B6.01} and \ref{B6.08}.

\section{Character components of a host}

\subsection{The character classes}\label{B6.28}

Let $a\ge2$ be an integer and let $q,d$ be primes with
$q\equiv3\pmod8$, $d\equiv7\pmod{16}$, and $\gcd(qd,f)=1$.
The host is given by
\[
 N=\operatorname{lcm}(2^{a+2},f)\,qd,\qquad
 H=\left\{u\in(\Z/N\Z)^\times:
 \begin{array}{l}
 u\bmod 2^{a+2}\in\{1,-1\},\quad\chi(u\bmod f)=1,\\
 u\bmod q\text{ and }u\bmod d\text{ are squares}
 \end{array}\right\}.
\]
Its field is the fixed field of $H$ in $\Q(\mu_N)$, up to $\Q$-isomorphism. For $t\in\N$, the host is large for $t$ if
\[
 q,d>2t+1,\qquad q,d>\max\{f,o,p_K,h_*,p_1\}.
\]
Here $p_K$ is the least prime at least $7$ congruent to $3\bmod4$ and prime to $f$. Put $h_*=\#\operatorname{Cl}(\Q(\sqrt{-p_K}))$. Let $p_1$ be the least prime congruent to $1\bmod o_{\mathrm{odd}}$ and prime to $2fp_K$, where $o_{\mathrm{odd}}$ is the odd part of $o$. The fixed set of primes is
\[
 \mathcal B_\chi=\{3\}\cup
 \{r:r\text{ is an odd prime and }r\mid ofp_Kh_*p_1\}.
\]

A coefficient ring $O$ at $\ell$ is a DVR in $\overline{\Q}_\ell$ finite over $\Z_\ell$. Host characters $(\Z/N\Z)^\times\to O^\times$ are trivial on $H$. Put
\[
 X_{\ell,O}(k)=
 \{\theta\in\operatorname{Hom}((\Z/N\Z)^\times,O^\times):
 H\subseteq\ker\theta,\ \theta(-1)=(-1)^k\}.
\]
Write $\theta=\theta_0\kappa$, where $\kappa$ has $\ell$-power order and $\theta_0$ has order prime to $\ell$. For a prime $r$, write $N=r^uN'$ with $\gcd(r,N')=1$. Under the Chinese remainder decomposition of the units, put $\theta_r(x_r,x')=\theta(x_r,1)$, with $\theta_r=1$ if $u=0$. Set
\[
 \theta_P=\theta_0(\theta_0)_\ell^{-1},\qquad
 \theta_1=\theta_0(\theta_0)_2^{-1},\qquad
 c_2(\theta)=v_2(\operatorname{cond}(\theta_2)).
\]
The modulus $\operatorname{lcm}(4,\operatorname{ord}\theta_1)$ is prime to $\ell$. Define
\[
 f_0=\operatorname{ord}_{(\Z/\operatorname{lcm}(4,\operatorname{ord}\theta_1)\Z)^\times}(\ell),
 \qquad n_0=v_2(\ell^{f_0}-1),\qquad
 R_1(\ell,\theta)=16o(\ell f)^2\,4^{n_0}.
\]
For a host large for $t$, the table classifies the triples $(\ell,k,\theta)$ with $1\le k<t$ and $\theta\in X_{\ell,O}(k)$.
\[
\begin{array}{c|l}
 \mathrm{I}&\ell\notin\mathcal B_\chi,\quad \ell-1\nmid2k\\
 \mathrm{II}&k\ge2,\quad \operatorname{ord}\theta_P\text{ even},
 \quad(\ell\in\mathcal B_\chi\text{ or }\ell-1\mid2k)\\
 \mathrm{III}&\ell\notin\mathcal B_\chi,\quad \ell-1\mid2k,\quad
 \theta\ne1,\quad\operatorname{ord}\theta\text{ odd}\\
 \mathrm{IV}&\ell\in\mathcal B_\chi,\quad k\ge2,\quad
 \operatorname{ord}\theta_P\text{ odd}\\
 \mathrm{V}&\ell\in\mathcal B_\chi,\quad k=1,\quad
 R_1(\ell,\theta)\le2^{c_2(\theta)}\\
 \mathrm{Vs}&\ell\in\mathcal B_\chi,\quad k=1,\quad
 2^{c_2(\theta)}<R_1(\ell,\theta)\\
 \mathrm{Z}&\ell\notin\mathcal B_\chi,\quad \ell-1\mid2k,\quad\theta=1.
\end{array}
\]
\leanref{UpperTheory.\allowbreak{}V.\allowbreak{}HV.\allowbreak{}Host.\allowbreak{}X}{Development.\allowbreak{}AutomorphicArithmetic.\allowbreak{}AutomorphicVarieties.\allowbreak{}HeightAndVanishing.\allowbreak{}DirichletPrimeToLPartMulLPowerPart}{188--189}
\leanref{UpperTheory.\allowbreak{}V.\allowbreak{}HV.\allowbreak{}chiP}{Development.\allowbreak{}AutomorphicArithmetic.\allowbreak{}AutomorphicVarieties.\allowbreak{}HeightAndVanishing.\allowbreak{}DirichletPrimeToLPartMulLPowerPart}{86}
\leanref{UpperTheory.\allowbreak{}V.\allowbreak{}Bchi}{Development.\allowbreak{}AutomorphicArithmetic.\allowbreak{}AutomorphicVarieties.\allowbreak{}StmtFields}{66--67}
\leanref{UpperTheory.\allowbreak{}V.\allowbreak{}HV.\allowbreak{}cases}{Development.\allowbreak{}AutomorphicArithmetic.\allowbreak{}AutomorphicVarieties.\allowbreak{}HeightAndVanishing.\allowbreak{}Cases}{8--48}
\leanref{UpperTheory.\allowbreak{}V.\allowbreak{}HV.\allowbreak{}CountV}{Development.\allowbreak{}AutomorphicArithmetic.\allowbreak{}AutomorphicVarieties.\allowbreak{}HeightAndVanishing.\allowbreak{}Counts}{69}
\leanref{UpperTheory.\allowbreak{}V.\allowbreak{}HV.\allowbreak{}CountVs}{Development.\allowbreak{}AutomorphicArithmetic.\allowbreak{}AutomorphicVarieties.\allowbreak{}HeightAndVanishing.\allowbreak{}Counts}{71}

The quotient by $H$ embeds in the product of the cyclic quotient $(\Z/2^{a+2}\Z)^\times/\{\pm1\}$ of order $2^a$, the image of $\chi$, and the square-class groups at $q,d$. Its order divides $2^{a+2}o$ and its exponent divides $2^{a+1}o$. At odd $r$ the local image lies in the image of $\chi$ for $r\mid f$ and in a square-class group otherwise. Thus $\operatorname{ord}\theta_r\mid o$ or $\operatorname{ord}\theta_r\mid2$, respectively. Also $16\mid N$.

Outside $\mathcal B_\chi$, $\ell\nmid\operatorname{ord}\theta$ and $\theta_0=\theta$. If $\ell-1\mid2k$, then $5\le\ell\le2k+1<q,d$, so $k\ge2$, $\ell\nmid N$ and $\theta_P=\theta$.

\subsection{Lengths and coefficient rings}\label{B6.29}

Let $\varepsilon_\ell:G_\Q\to\Z_\ell^\times$ be the cyclotomic character. Inflate $\theta$ via the action on $\mu_N$ and put
\[
 A_O=O\otimes_{\Z_\ell}(\Q_\ell/\Z_\ell),\qquad
 M_{\ell,k,\theta}=A_O(\theta^{-1}\varepsilon_\ell^{-k}).
\]
Let $\mathcal S_{\ell,k,\theta}$ consist of classes in $H^1(G_\Q,M_{\ell,k,\theta})$ vanishing on finite inertia and on decomposition groups at $p\mid\ell N$. Define
\[
 L_{\ell,k,\theta}=\operatorname{length}_O\mathcal S_{\ell,k,\theta}
 \in\N\cup\{\infty\},\qquad
 e_O=\operatorname{length}_O(O/\ell O).
\]
For finite $L_{\ell,k,\theta}$ put $a_{\ell,k,\theta}=L_{\ell,k,\theta}/e_O$. A character bound $B$ means finite length and $a_{\ell,k,\theta}\log\ell\le B$.
\leanref{UpperTheory.\allowbreak{}V.\allowbreak{}Chr.\allowbreak{}rho}{Development.\allowbreak{}AutomorphicArithmetic.\allowbreak{}AutomorphicVarieties.\allowbreak{}CharacterLevelBounds.\allowbreak{}ShaFamMono}{17--18}
\leanref{UpperTheory.\allowbreak{}V.\allowbreak{}Chr.\allowbreak{}VS}{Development.\allowbreak{}AutomorphicArithmetic.\allowbreak{}AutomorphicVarieties.\allowbreak{}CharacterLevelBounds.\allowbreak{}ShaFamMono}{30}
\leanref{UpperTheory.\allowbreak{}V.\allowbreak{}Chr.\allowbreak{}V'}{Development.\allowbreak{}AutomorphicArithmetic.\allowbreak{}AutomorphicVarieties.\allowbreak{}CharacterLevelBounds.\allowbreak{}ShaFamMono}{34}
\leanref{UpperTheory.\allowbreak{}V.\allowbreak{}Chr.\allowbreak{}nrm}{Development.\allowbreak{}AutomorphicArithmetic.\allowbreak{}AutomorphicVarieties.\allowbreak{}CharacterLevelBounds.\allowbreak{}ShaFamMono}{71}
\leanref{UpperTheory.\allowbreak{}V.\allowbreak{}HV.\allowbreak{}VLe}{Development.\allowbreak{}AutomorphicArithmetic.\allowbreak{}AutomorphicVarieties.\allowbreak{}HeightAndVanishing.\allowbreak{}DirichletPrimeToLPartMulLPowerPart}{120}
\leanref{UpperTheory.\allowbreak{}V.\allowbreak{}HV.\allowbreak{}coefIndep}{Development.\allowbreak{}AutomorphicArithmetic.\allowbreak{}AutomorphicVarieties.\allowbreak{}HeightAndVanishing.\allowbreak{}CoefIndep}{68--77}

Finite free base change $O\subseteq O'$ commutes with continuous cochains and restriction kernels. It preserves finiteness and scales finite lengths and $e_O$ by $e(O'/O)$. Compare rings in a common finite extension.

\section{The automorphic counting bounds}

Use the hypotheses $\mathscr H$ of \ref{B6.05} and \ref{B6.06} and the small-twist bound of \ref{B6.15}.

\begin{theorem}[Class I, summed over primes]\label{B6.30}
For every $\chi$ there are $C_{\mathrm I}\in\N$ and distinct primes $r_1,r_2$. Let a host for $\chi$ be large for $t$ with $q,d>r_1,r_2$. Fix $1\le k<t$ and a finite set $S$ of odd primes with $\ell\notin\mathcal B_\chi$ and $\ell-1\nmid2k$ for each $\ell\in S$. For any $O_\ell$ all lengths for $\theta\in X_{\ell,O_\ell}(k)$ are finite, and
\[
 \sum_{\ell\in S}\log\ell
 \sum_{\theta\in X_{\ell,O_\ell}(k)}a_{\ell,k,\theta}
 \le(Nt)^{C_{\mathrm I}}.
\]
\end{theorem}
\leanref{UpperTheory.\allowbreak{}V.\allowbreak{}U3.\allowbreak{}countI}{Development.\allowbreak{}AutomorphicArithmetic.\allowbreak{}AutomorphicVarieties.\allowbreak{}UnitaryThreeGroups.\allowbreak{}CountI}{22--136}
\leanref{UpperTheory.\allowbreak{}V.\allowbreak{}Z.\allowbreak{}J.\allowbreak{}count\_\allowbreak{}i\_\allowbreak{}of\_\allowbreak{}printed\_\allowbreak{}bound\_\allowbreak{}k\_\allowbreak{}kato\_\allowbreak{}automorphy}{Development.\allowbreak{}AutomorphicArithmetic.\allowbreak{}AutomorphicVarieties.\allowbreak{}ModularityArgument.\allowbreak{}ThresholdArguments.\allowbreak{}DenseThresholdCases.\allowbreak{}LinkedUnlE}{49--54}

\begin{proof}
Use the constants $c,A$ of \ref{B6.42} for $K=\Q(\sqrt{-p_K})$, of discriminant $-p_K$. Distinct primes $r_1,r_2>30fp_K+1$ satisfy $r_i\nmid30N|D_K|$ when $q,d>r_1,r_2$. For each character kernel $J$ in $\Gamma=(\Z/N\Z)^\times$, set
\[
 S_J=\{\ell\in S:\ker\theta=J\text{ for some }\theta\in X_{\ell,O_\ell}(k)\}.
\]
Apply \ref{B6.42} at $r_1,r_2$ for each $J$ with $S_J\ne\varnothing$ to obtain nonzero integers $D_{J,i}$ with
\[
 \log|D_{J,i}|\le c\,\varphi([\Gamma:J])[N(k+1)]^A\log r_i.
\]
For a witnessing $\theta$ at $\ell\in S_J$, \ref{B6.28} gives $\ell\nmid\operatorname{ord}\theta$ and $\theta_\ell^2=1$. Inertia at $p_K$ fixes $\bar\varepsilon_\ell$ and $\theta_P$. It acts nontrivially on $K$. The remaining conditions of \ref{B6.42} follow from $\ell\nmid p_Kh_*$, $p_K\nmid N$, $2\mid N$ and the parity of $\theta$.

Use an $r_i\ne\ell$ to obtain finite lengths in the orbit and $\sum_{\ker\theta=J}a_{\ell,k,\theta}\le 2v_\ell(D_{J,1})+2v_\ell(D_{J,2})$. Equations (6.42.3) and (6.42.4) bound the full sum by $2c(\log r_1+\log r_2)N[N(k+1)]^A$. Since $N\ge16$ and $k+1\le t$, a power of $Nt$ absorbs the fixed factor.
\end{proof}

\begin{theorem}[Classes II and III]\label{B6.31}
For every $\chi$ there are $C_{\mathrm{II}},C_{\mathrm{III}}\in\N$ such that for every host large for $t$, every odd prime $\ell$, every coefficient ring $O$, and every $1\le k<t$, a character $\theta\in X_{\ell,O}(k)$ of class II or III has finite length and, respectively,
\[
 a_{\ell,k,\theta}\log\ell\le(\ell Nk)^{C_{\mathrm{II}}},
 \qquad
 a_{\ell,k,\theta}\log\ell\le(\ell Nk)^{C_{\mathrm{III}}}.
\]
\end{theorem}
\leanref{UpperTheory.\allowbreak{}V.\allowbreak{}Th.\allowbreak{}sp4Main\_\allowbreak{}alt}{Development.\allowbreak{}AutomorphicArithmetic.\allowbreak{}AutomorphicVarieties.\allowbreak{}ModularityArgument.\allowbreak{}ThresholdArguments.\allowbreak{}DenseThresholdCases.\allowbreak{}ConductorAtMapInjectiveCoeff}{131--171}
\leanref{UpperTheory.\allowbreak{}V.\allowbreak{}HV.\allowbreak{}DescentVanishingReductions.\allowbreak{}countII\_\allowbreak{}and\_\allowbreak{}countIII\_\allowbreak{}of\_\allowbreak{}printsH\_\allowbreak{}of\_\allowbreak{}sp4In}{Development.\allowbreak{}AutomorphicArithmetic.\allowbreak{}AutomorphicVarieties.\allowbreak{}ModularityArgument.\allowbreak{}ThresholdArguments.\allowbreak{}DenseThresholdCases.\allowbreak{}AdjointLiftingOfCohomology}{33--39}

\begin{proof}
Set
\[
\begin{split}
 n_{\mathrm{II}}&=
 \left\lceil\log_2\max\{\ell N4^{n_0},
 2^{c_2(\theta)+6}o\ell^3f^2\}\right\rceil,\\
 n_{\mathrm{III}}&=
 \left\lceil\log_2\max\{\ell fp_1 4^{n_0},8\ell^2\}\right\rceil .
\end{split}
\]
Choose $(n',p')=(n_{\mathrm{II}},1)$ in II and $(n_{\mathrm{III}},p_1)$ in III. After base change take a primitive odd $\xi_2$ modulo $2^{n'}$. Take $\xi_0=1$ in II. In III, $\operatorname{ord}\theta\mid o_{\mathrm{odd}}\mid p_1-1$, so take $\xi_0$ modulo $p_1$ of order $\operatorname{ord}\theta$. Put $\xi=\xi_0\xi_2$ and let $\omega$ be the Teichmüller lift of $\bar\varepsilon_\ell$.

The level formulas give $c_2(\theta)<n'$. Let $T$ consist of $\ell$ and the primes $v\ne\ell$ where $\kappa_v\ne1$. In II, $v\le\ell f$ for $v\in T$, and the $2$-parts of the orders of $\theta,\xi_0$ divide $2^s$ with $s=c_2(\theta)+1+v_2(o)$. In III, $T=\{\ell\}$ and $s=0$ suffices. In both cases $\gcd(p',v)=1$ and $v^2 2^s\le2^{n'}$ for $v\in T$.

Put $\Psi=\theta\varepsilon_\ell^k$, with bars for reduction. For $v\in T$ choose an arithmetic Frobenius lift $\sigma_v$ with $\varepsilon_\ell(\sigma_\ell)=1$ at $\ell$. At other $v\in T$, $\kappa_v\ne1$ forces $v\equiv1\bmod\ell$, so $\bar\varepsilon_\ell(\sigma_v)=1$ on $T$. Take $u$ to be an odd common multiple of the odd parts of the orders of $\theta,\xi_0$. If $\overline{\xi\Psi^i}(\sigma_v)=1$ for some $i\in\Z$, taking $2^su$-th powers kills the $\theta,\xi_0$ factors. Reduction is injective on $2$-power roots of unity, so $\xi_2(v)^{2^su}=1$. Squaring and using faithfulness on the units congruent to $1\bmod4$ gives $2^{n'}\mid v^{2^{s+1}u}-1$. The valuation is $s+v_2(v^2-1)$, contrary to $2^s(v^2-1)<2^{n'}$.

To apply \ref{B6.59}, use $f_0$ as the exponent in (6.59.3). At a common level, the lifts of $\bar\xi$, $\bar\Psi\bar\xi^{-1}$ and $\bar\Psi\bar\xi$ have odd-modulus and dyadic factors
\[
 (\xi_0,\xi_2),\qquad
 (\omega^k\theta_1\xi_0^{-1},\theta_2\xi_2^{-1}),\qquad
 (\omega^k\theta_1\xi_0,\theta_2\xi_2).
\]
Here $\kappa$ reduces to $1$. Each odd-modulus factor is induced by a primitive character of conductor $f'\mid p'\ell\operatorname{cond}(\theta_1)$. Its order divides $\ell^{f_0}-1$, which is divisible by $\ell-1$ and $\operatorname{ord}\theta_1$. In III, $\theta_1=\theta$ and $\xi_0$ has that same order. Each dyadic factor is primitive because it is nontrivial at $1+2^{n'-1}$, where $\theta_2$ is trivial. The products are odd since $\Psi$ is even and $\xi$ is odd. Since $\operatorname{cond}(\theta_1)$ divides $N$ in II and $f$ in III, the level formulas and the definition of $f_0$ give
\[
 4\mid\ell^{f_0}-1,\qquad
 f'4^{n_0}\le p'\ell\operatorname{cond}(\theta_1)4^{n_0}\le2^{n'}.
\]

For (6.59.0), at places in $T$ a local coboundary vanishes on inertia in the character kernel. Away from $T$, the local condition kills inertia in $\ker\bar\rho$, where $\bar\rho$ is the residual diagonal representation of the datum. This kernel has index prime to $\ell$ because the image lies in a product of multiplicative groups of the residue field. Hence \ref{B6.59} gives vanishing for all three reductions and their inverses.

In II, $\theta_P$ has even order. In III it equals the nontrivial $\theta$. Its conductor and order are prime to $\ell$, so the Chinese remainder theorem gives $g\in G_\Q$ fixing $\mu_\ell$ with $\bar\theta_P(g)\ne1$. Then $\bar\Psi(g)=\bar\theta_P(g)$ since $(\theta_0)_\ell$ factors through the action on $\mu_\ell$. Since $\Psi$ is even and $\xi$ is odd, either $g$ or its product with complex conjugation gives a marking $g_0$ with $\bar\xi(g_0)\ne1$ and $\bar\Psi(g_0)\ne1$. In the good datum $\mathcal D$, take $S_{\mathcal D}=\{p:p\mid2\ell Np'\}$.
\leanref{UpperTheory.\allowbreak{}V.\allowbreak{}X.\allowbreak{}xvan}{Development.\allowbreak{}AutomorphicArithmetic.\allowbreak{}AutomorphicVarieties.\allowbreak{}LevelCohomology.\allowbreak{}HcXvan}{19--41}
\leanref{UpperTheory.\allowbreak{}V.\allowbreak{}GradeZeroInterpolatingElement.\allowbreak{}dataEx}{Development.\allowbreak{}AutomorphicArithmetic.\allowbreak{}AutomorphicVarieties.\allowbreak{}SymplecticGSp4.\allowbreak{}DataEx}{366--450}

Let $L(\mathcal D)$ be the length after imposing decomposition conditions at $S_{\mathcal D}$ on the group defining $L_{\ell,k,\theta}$. At added primes the Frobenius scalar is a root of unity times $p^{-k}$. Its difference from $1$ acts surjectively on the divisible module, so the unramified $H^1$ vanishes. Thus $L_{\ell,k,\theta}\le L(\mathcal D)$.

Put $\mathcal N=\ell2^{n'}Np'$ and $h=(\ell-1)2^{n'}(2o)$. The finite parts of the distinguished characters are $\theta\xi\omega^{-1}$, $\theta$, $\xi\omega^{-1}$ and $1$. They factor through $\Gal(\Q(\mu_{\mathcal N})/\Q)$ and have orders dividing $h$. Take arithmetic Frobenius $\sigma$ at the least prime $p_0\ge5$ prime to $\ell Np'$. The trace of the distinguished point $\rho_0$ at $\sigma$ is
\[
 p_0^{k+1}z_1+p_0^kz_2+p_0z_3+1,\qquad z_i^h=1 .
\]

By \ref{B6.49}, the zero-weight algebra is finite over $O$ and the irreducible points have the required automorphic attachments. Each representation has level $(\mathcal N;(k+1,k,1,0))$ and multiplier $\varepsilon_\ell^{k+1}\theta\xi\omega^{-1}$. At each place above $\ell$, choose an ordinary flag whose graded characters $\kappa_i$ satisfy $\kappa_i\ne\kappa_j\varepsilon_\ell$ for all $i,j$. By \ref{B6.50} for singletons, each trace difference at $\sigma$ is nonzero.

Use \ref{B6.49} for proper span at block points and \ref{B6.48} for rigidity at irreducible points. By \ref{B6.09} and \ref{B6.15}, with local finiteness and Tate duality, strict Selmer groups with divisible coefficients of character $\eta^{-1}\varepsilon_\ell^{-e}$ for finite-order $\eta$ and $e\ge1$ have finite length. The corresponding strict classes over $\overline{\Q}_\ell$ vanish.

The unit $\Psi^{-1}(g_0)-1$ gives unique normalized cocycles vanishing at $g_0$, which add. Take a nonzero $b_0\in O$ annihilating the finite-length Selmer module and choose $s_v$ in inertia at each $v\in T$ with $\Psi(s_v)\ne1$. Use the infinite cyclotomic image at $\ell$ and $\kappa_v\ne1$ elsewhere in $T$. Evaluation at $s_v$ shows that $b_0(\Psi^{-1}(s_v)-1)$ kills the local coboundary witness. Hence $b=b_0\prod_{v\in T}(\Psi^{-1}(s_v)-1)$ kills the cocycles and the witnesses, which can be taken in $A_O[b]\simeq O/(b)$.

The associated derivations lift integrally in weight $-1$ and kill the defining equations modulo $b$. Their reductions add on the universal matrix entries because the values there are fixed multiples of the cocycle values. Those entries generate the deformation algebra, and evaluation on the $(3,1)$ entries recovers the cocycles up to units. The resulting additive injection gives $L(\mathcal D)\le\operatorname{length}_O C_{-1}(F)$ by (6.58.2) of \ref{B6.58}. Apply (6.58.3) to $a=\operatorname{tr}\rho_{\mathrm{univ}}(\sigma)-\operatorname{tr}\rho_0(\sigma)$. The representatives of the two-sided orbits form a pairwise nonisomorphic family $\Xi$ because conjugate irreducible representations belong to the same orbit. Thus
\[
 \frac{L(\mathcal D)}{e_O}\log\ell
 \le 2\sum_{\rho\in\Xi}
 -\log|\operatorname{tr}\rho(\sigma)-\operatorname{tr}\rho_0(\sigma)|_\ell,
 \qquad |\ell|_\ell=\ell^{-1}.
\]
Let $C_0,A_0$ denote the constants $c,A$ in (6.50.1). By \ref{B6.50} for $\Xi$, the nonzero integer $D$ gives the upper bound $2C_0h[\mathcal N(k+1)]^{A_0}\log p_0$, since $v_\ell(D)\log\ell\le\log|D|$. Since $\operatorname{ord}\theta_1\mid2o$, one has $f_0\le F_0=\varphi(4o)$ and $2^{n_0}<\ell^{F_0}$. Substitute the level formulas and the bounds $2^{c_2(\theta)}\mid N$, $p'\le1+p_1$ and $p_0\le2\ell Np'+10$ to get $2C_0h[\mathcal N(k+1)]^{A_0}\log p_0\le(\ell Nk)^{C_\chi}$ for a constant depending on $\chi$. Coefficient independence gives the result over the original ring.
\end{proof}

\section{Class IV and its cyclotomic layer}

\begin{theorem}[Transfer of the class IV sum]\label{B6.32}
Fix $\chi$, $\ell\in\mathcal B_\chi$, and $j\in\N$. Let $t\ge2$, $B\in\R$, and write $s_{\ell,k}(L)=\operatorname{length}_{\Z_\ell}\Sha_{\ell,k}(L)$ for number fields $L$. For every realization $L_0$ of $\Q(\mu_{\ell f})$ and cyclotomic layer $L_1/L_0$ of degree $\ell^j$, assume finite lengths for $1\le k<t$ and
\[
 \sum_{1\le k<t}s_{\ell,k}(L_1)\le B
\]
for every choice of local data. Then for every host for $\chi$ large for $t$ and every $O$-valued host character $\theta$ with $\operatorname{ord}\theta_P$ odd, all $L_{\ell,k,\theta}$ with $1\le k<t$ are finite, and
\[
 \sum_{1\le k<t}a_{\ell,k,\theta}
 \le B+(Nt)^{C_4},\qquad C_4=\ell f\ell^j+3.
\]
\end{theorem}
\leanref{UpperTheory.\allowbreak{}V.\allowbreak{}HV.\allowbreak{}LayerLe}{Development.\allowbreak{}AutomorphicArithmetic.\allowbreak{}AutomorphicVarieties.\allowbreak{}HeightAndVanishing.\allowbreak{}LayerLe}{12--14}
\leanref{UpperTheory.\allowbreak{}V.\allowbreak{}HV.\allowbreak{}CountIV}{Development.\allowbreak{}AutomorphicArithmetic.\allowbreak{}AutomorphicVarieties.\allowbreak{}HeightAndVanishing.\allowbreak{}LayerLe}{16--20}
\leanref{UpperTheory.\allowbreak{}V.\allowbreak{}HV.\allowbreak{}countIV}{Development.\allowbreak{}AutomorphicArithmetic.\allowbreak{}AutomorphicVarieties.\allowbreak{}HeightAndVanishing.\allowbreak{}CountIV}{39--89}

\begin{proof}
Since $\ell\ne q,d$, the local factors at $2,q,d$ retain their $2$-power orders in $\theta_P$ and hence are trivial. Thus $\operatorname{cond}\theta\mid f$, and \ref{B6.33} gives $L_1$ of degree $n\le\ell f\ell^j$ trivializing $\theta$.

For $k\ge1$, restriction to any abelian $L_1/\Q$ of degree $n$ that trivializes $\theta$ gives the extended-length inequality
\[
 L_{\ell,k,\theta}\le e_O\bigl(s_{\ell,k}(L_1)+n\nu_\ell(N,k)\bigr),
 \tag{6.IV}\label{eq:b6-transfer}
\]
where
\[
 \nu_\ell(N,k)=
 \begin{cases}
 \max\{v_\ell(N),1\}+v_\ell(k),&\ell\mid N\text{ or }\ell-1\mid k,\\
 0,&\text{otherwise}.
 \end{cases}
\]
Write $G=G_\Q$, $U=G_{L_1}$, and $M=M_{\ell,k,\theta}$, with scalar action $\lambda$. In the inflation--restriction sequence \cite[II, Proposition 1.34]{milne-2020}, the kernel is represented by the module $Z$ of cocycles vanishing on $U$. Evaluation on $G/U$ embeds $Z$ into $M^n$. By abelianness and the cocycle identity,
\[
 (\lambda(g)-1)z(h)=(\lambda(h)-1)z(g)\qquad(z\in Z).
\]
Thus $(\lambda(g)-1)z$ is a coboundary. If $a$ generates $(\lambda(g)-1:g\in G)\subset O$, the restriction kernel is a quotient of $Z/aZ$.

Lift the generators of any finitely generated submodule of $Z/aZ$ to a finite-length submodule $W\subseteq Z$. Its image has length at most $\operatorname{length}_O(W/aW)=\operatorname{length}_O W[a]$. Thus $\operatorname{length}_O(Z/aZ)\le\operatorname{length}_O Z[a]$, and evaluation embeds $Z[a]$ into $(M^G)^n$.
\leanref{UpperTheory.\allowbreak{}V.\allowbreak{}Base.\allowbreak{}kerResAb}{Development.\allowbreak{}AutomorphicArithmetic.\allowbreak{}AutomorphicVarieties.\allowbreak{}GaloisLibrary.\allowbreak{}SelmerTransport.\allowbreak{}KerResAb}{46--151}
\leanref{UpperTheory.\allowbreak{}V.\allowbreak{}Chr.\allowbreak{}upV}{Development.\allowbreak{}AutomorphicArithmetic.\allowbreak{}AutomorphicVarieties.\allowbreak{}CharacterLevelBounds.\allowbreak{}LiftGens}{24--57}

To bound $M^G$, let $s=\max\{v_\ell(N),1\}$. Inertia at $\ell$ contains a lift of $1+\ell^s$ under $\varepsilon_\ell$, and $\theta$ is trivial on this lift. The scalar minus $1$ has valuation $s+v_\ell(k)$. When $\ell\nmid N$ and $\ell-1\nmid k$, a lift of a residual cyclotomic generator has scalar minus $1$ a unit. Thus $\operatorname{length}_O M^G\le e_O\nu_\ell(N,k)$. On $U$ the twist is $\varepsilon_\ell^{-k}$. The inertia conditions and decomposition conditions above $\ell$ persist under restriction, so the image lies in $O\otimes_{\Z_\ell}\Sha_{\ell,k}(L_1)$ and has length at most $e_Os_{\ell,k}(L_1)$. This proves \eqref{eq:b6-transfer}.

By the layer hypothesis and \eqref{eq:b6-transfer}, the lengths are finite. Sum after dividing by $e_O$ and use $\nu_\ell(N,k)\le2Nk$. The bound is $B+2nNt^2\le B+(Nt)^{n+3}\le B+(Nt)^{C_4}$, since $2n\le2^{n+1}$.
\end{proof}

\begin{lemma}[A field that trivializes the transferred characters]\label{B6.33}
For each odd prime $\ell$ and $j\in\N$, every algebraic closure of $\Q$ contains abelian fields $L_0\subseteq L_1$ with
\[
 L_0\simeq\Q(\mu_{\ell f}),\qquad
 L_1/L_0\text{ a cyclotomic layer},\qquad
 [L_1:L_0]=\ell^j,\qquad [L_1:\Q]\le\ell f\ell^j.
\]
For every integral domain $A$, every $N\ge1$, and every $A$-valued Dirichlet character of level $N$ whose conductor divides $f$, the associated Galois character is trivial on $G_{L_1}$.
\end{lemma}
\leanref{UpperTheory.\allowbreak{}V.\allowbreak{}HV.\allowbreak{}ivField}{Development.\allowbreak{}AutomorphicArithmetic.\allowbreak{}AutomorphicVarieties.\allowbreak{}HeightAndVanishing.\allowbreak{}IvField}{32--82}

\begin{proof}
Take $L_0=\Q(\mu_{\ell f})$ and $L_1=\Q(\mu_{\ell^{j+1}f})$. The degree ratio is $\varphi(\ell^{j+1}f)/\varphi(\ell f)=\ell^j$ and $[L_1:\Q]=\ell^j\varphi(\ell f)\le\ell^j\ell f$. Characters of conductor dividing $f$ factor through the action on $\mu_f\subset L_0$.
\end{proof}

\section{The two classes at twist one}

\begin{theorem}[Class V]\label{B6.34}
For every $\chi$ there is $C_{\mathrm V}\in\N$ such that for every host large for $t\ge2$, every $\ell\in\mathcal B_\chi$, and every $\theta\in X_{\ell,O}(1)$ with $R_1(\ell,\theta)\le2^{c_2(\theta)}$, one has
\[
 L_{\ell,1,\theta}<\infty,\qquad
 a_{\ell,1,\theta}\log\ell\le(\ell N)^{C_{\mathrm V}}.
\]
\end{theorem}
\leanref{UpperTheory.\allowbreak{}V.\allowbreak{}Z.\allowbreak{}J.\allowbreak{}count\_\allowbreak{}v\_\allowbreak{}of\_\allowbreak{}printed\_\allowbreak{}bound\_\allowbreak{}k\_\allowbreak{}reduced\_\allowbreak{}unitary\_\allowbreak{}lift}{Development.\allowbreak{}AutomorphicArithmetic.\allowbreak{}AutomorphicVarieties.\allowbreak{}ModularityArgument.\allowbreak{}ThresholdArguments.\allowbreak{}DenseThresholdCases.\allowbreak{}LinkedUnlE}{55--61}
\leanref{UpperTheory.\allowbreak{}V.\allowbreak{}UnitaryThreeInertia.\allowbreak{}countV\_\allowbreak{}of\_\allowbreak{}finiteness\_\allowbreak{}duality\_\allowbreak{}euler\_\allowbreak{}host}{Development.\allowbreak{}AutomorphicArithmetic.\allowbreak{}AutomorphicVarieties.\allowbreak{}ModularityArgument.\allowbreak{}ThresholdArguments.\allowbreak{}DenseThresholdCases.\allowbreak{}CountVOf}{33--60}

\begin{proof}
The threshold on $R_1$ gives the Gaussian host datum over $K=\Q(i)$. By \ref{B6.51} and \ref{B6.52}, the weight-zero submodule of the marked universal lift $u$ is finite over $O$ and each Type I point with matrix restriction $r$ to $G_K$ has an automorphic attachment for $r\varepsilon_\ell^{-1}$. Local comparison gives level $(N;(1,0,-1))$ and the separated flags of \ref{B6.41}.

By \ref{B6.51}, $r$ is irreducible on every open subgroup of $G_K$. Reducibility on $G_K\cap\ker\varepsilon_\ell$ would give a nontrivial self-twist of order dividing three by Clifford theory. Schur's lemma on its open kernel forces a scalar intertwiner and hence a trivial twist. First-order deformations give extended adjoint cocycles that are inertia coboundaries away from $\ell$ and lie in $H^1_f$ at $\ell$ by ordinary comparison. Apply \ref{B6.53}.

Since $12\ell N-1\equiv3\pmod4$, it has a prime divisor $p\equiv3\pmod4$. Then $p\nmid6\ell N$ and $7\le p\le12\ell N$. Take $\sigma$ to be inverse Frobenius at the Gaussian prime of norm $p^2$. The distinguished trace is $(p^2)^{-1}z+1+p^2z^{-1}$, with $z^{\varphi(N)}=1$ under the coefficient identification. Apply \ref{B6.41} to singletons with the indexed family of $\varphi(N)$-th roots of unity. Each trace difference at $\sigma$ is nonzero.

Put $\Psi=\theta\varepsilon_\ell$ and $c=c_2(\theta)$. Let $\eta_4$ cut out $K$ and let $\omega$ be the Teichmüller lift of $\bar\varepsilon_\ell$. At odd primes outside $f$ the local factors have order dividing two, so $\operatorname{cond}(\theta_1^2)\mid f$. The threshold gives $c\ge4$ and $2\ell f4^{n_0}\le2^c$. The primitive character inducing $(\omega\theta_1)^2$ has odd conductor $f'\mid\ell f$ and order dividing $\ell^{f_0}-1$. The character $\theta_2^2\eta_4$ has conductor $2^{c-1}$ because squaring halves the order on the cyclic units congruent to $1\bmod4$, where $\eta_4$ is trivial. The product of the two primitive factors is odd and reduces to $\eta_4\bar\Psi^2$. Thus $4\mid\ell^{f_0}-1$ and $f'4^{n_0}\le2^{c-1}$.

Away from $\ell$, the local condition kills inertia in $\ker(\eta_4,\bar\Psi)$, whose index divides $2|\mathbf k^\times|$ for the residue field $\mathbf k$. At $\ell$ the cocycle is locally a coboundary, so it vanishes on inertia in the character kernel. These are the conditions (6.59.0). The inverse-character conclusion of \ref{B6.59} gives residual Selmer vanishing for $\eta_4\bar\Psi^{-2}$ in weight $-2$.
\leanref{UpperTheory.\allowbreak{}V.\allowbreak{}X.\allowbreak{}dirData}{Development.\allowbreak{}AutomorphicArithmetic.\allowbreak{}AutomorphicVarieties.\allowbreak{}LevelCohomology.\allowbreak{}HcDirData}{29--110}
\leanref{UpperTheory.\allowbreak{}V.\allowbreak{}X.\allowbreak{}xpp}{Development.\allowbreak{}AutomorphicArithmetic.\allowbreak{}AutomorphicVarieties.\allowbreak{}LevelCohomology.\allowbreak{}Xpp}{11--85}

At $v\in T_b$ from (6.51.0), the frame supplies an element with trivial sign and trivial residual cyclotomic value. On this element $\bar\Psi^2\ne1$. Its action minus the identity is invertible in nonzero weights. Local duality gives $H^2=0$ there, and the local Euler formula gives $H^1=0$ there for $v\ne\ell$.

In the normalized Poitou--Tate Euler formula the terms away from $\ell$ cancel. At $\ell$ the local condition takes coboundaries in negative weights and all cocycles in positive weights. Write $d_j$ for the residual dimension and $d_j^+$ for the dimension fixed by complex conjugation in weight $j$. The difference of Selmer and obstruction dimensions is $-d_j^+$ for $j<0$ and $d_j-d_j^+$ for $j>0$. Here $(d_j,d_j^+)=(2,1)$ for $|j|=1$ and $(1,0)$ for $|j|=2$. The tame Euler formula gives zero at weight zero. Thus the difference is $-1$ at $-1$, $1$ at $1$ and $2$, and zero elsewhere. The Gaussian class number one gives weight-zero cocycle vanishing. For integral cocycles in negative weights use \ref{B6.09}, \ref{B6.15} and local duality. By \ref{B6.58}, an integer $V$ bounds the frame's Selmer lengths over $O/(b)$ for every $b\ne0$. The Type I representatives $\rho_i$ have pairwise nonisomorphic restrictions to $G_K$ and satisfy
\[
 V\le2e_O\sum_i v_\ell\bigl(
 \operatorname{tr}\rho_i(\sigma)-\operatorname{tr}\rho_0(\sigma)\bigr).
\]
For $b=\ell^m$, represent a class killed by $b$ by $z$ with $bz$ the coboundary of $m_0$. Subtract the coboundary of a $b$-division point of $m_0$ to make $z$ take values in $A_O[b]\simeq O/(b)$. Where the frame requires a local coboundary, evaluation at its local element $g$ with $\Psi(g)-1\in O^\times$ shows that the coboundary comes from an element killed by $b$. On the other specified inertia subgroups $\Psi=1$, so $z$ vanishes there. Hence $\mathcal S_{\ell,1,\theta}[b]$ comes from the finite-coefficient group and has length at most $V$. Finite image of continuous cocycles gives exhaustion by these torsion subgroups, so $L_{\ell,1,\theta}\le V$.

Group the representatives by the two embeddings of $K$. By \ref{B6.41} with the same roots of unity in each group, the sum of valuations times $\log\ell$ is at most $2a_0\varphi(N)(2N)^{b_0}\log p$ for fixed $a_0,b_0$. Divide by $e_O$ and use $\varphi(N)\le N$ and $p\le12\ell N$.
\end{proof}

\begin{theorem}[Class Vs]\label{B6.35}
For every $\chi$ there is $C_{\mathrm{Vs}}\in\N$ such that for every host large for $t\ge2$, every $\ell\in\mathcal B_\chi$, and every $\theta\in X_{\ell,O}(1)$ with $2^{c_2(\theta)}<R_1(\ell,\theta)$, one has
\[
 L_{\ell,1,\theta}<\infty,\qquad
 a_{\ell,1,\theta}\log\ell\le(\ell N)^{C_{\mathrm{Vs}}}.
\]
\end{theorem}
\leanref{UpperTheory.\allowbreak{}V.\allowbreak{}HV.\allowbreak{}countVs}{Development.\allowbreak{}AutomorphicArithmetic.\allowbreak{}AutomorphicVarieties.\allowbreak{}HeightAndVanishing.\allowbreak{}CountVs}{86--171}

\begin{proof}
Put $F_0=\varphi(4o)$ and $M_\chi=\max\mathcal B_\chi$. By the order bounds of \ref{B6.28} and \ref{B6.31},
\[
 \operatorname{ord}\theta\le2o\,2^{c_2(\theta)}
 < 2oR_1(\ell,\theta)\le32o^2(\ell f)^2\ell^{2F_0}.
\]
Let $L$ be the compositum of the field cut out by $\theta$ with $\Q(\mu_\ell)$. Then $L/\Q$ is abelian and totally imaginary. The field $L$ trivializes $\theta$. We have
\[
 n=[L:\Q]\le\ell\operatorname{ord}\theta
 \le n_\chi=32o^2(M_\chi f)^2 M_\chi^{2F_0+1},\qquad
 L\subseteq F=\Q(\mu_{\ell N}).
\]
Put $D=|D_L|$ and $X=\ell N$. Then $D^{[F:L]}\le |D_F|\le X^{[F:\Q]}$, so $\log D\le n\log X$.

The dual of the strict kernel for $L$ embeds in $H^2(G_{L,S_\ell},\Z_\ell(2))$ by \ref{B6.09}, local finiteness and local Tate duality. By \ref{B6.15}, this $H^2$ is finite and the logarithm of its order is at most $X^A$, where $A=(2n_\chi)^{\lceil C_K\rceil}(n_\chi+1)$. Apply \eqref{eq:b6-transfer} with $k=1$ to obtain finiteness and
\[
 a_{\ell,1,\theta}\log\ell
 \le X^A+n\nu_\ell(N,1)\log\ell\le X^A+2n_\chi X.
\]
\end{proof}

\section{The trivial character}

\begin{theorem}[Vanishing in class Z]\label{B6.36}
Let $\ell\ge5$ be prime and let $k\ge2$ be even with $\ell-1\mid2k$.
For every coefficient ring $O$, the subgroup of $H^1(G_\Q,A_O(-k))$
whose classes restrict to zero on the inertia group at every finite prime and on a decomposition group at $\ell$ has length zero. Consequently, for every host large for $t$ and every class Z triple with $1\le k<t$,
\[
 L_{\ell,k,1}=0,\qquad a_{\ell,k,1}=0.
\]
\end{theorem}
\leanref{UpperTheory.\allowbreak{}V.\allowbreak{}X.\allowbreak{}propZ}{Development.\allowbreak{}AutomorphicArithmetic.\allowbreak{}AutomorphicVarieties.\allowbreak{}LevelCohomology.\allowbreak{}PropZ}{11--16}
\leanref{UpperTheory.\allowbreak{}V.\allowbreak{}X.\allowbreak{}zOne}{Development.\allowbreak{}AutomorphicArithmetic.\allowbreak{}AutomorphicVarieties.\allowbreak{}LevelCohomology.\allowbreak{}TorsClass}{97--205}
\leanref{UpperTheory.\allowbreak{}V.\allowbreak{}X.\allowbreak{}zQuad}{Development.\allowbreak{}AutomorphicArithmetic.\allowbreak{}AutomorphicVarieties.\allowbreak{}LevelCohomology.\allowbreak{}ZqField}{228--233}
\leanref{UpperTheory.\allowbreak{}V.\allowbreak{}HV.\allowbreak{}countZ}{Development.\allowbreak{}AutomorphicArithmetic.\allowbreak{}AutomorphicVarieties.\allowbreak{}HeightAndVanishing.\allowbreak{}CountZ}{5--32}

\begin{proof}
Put $A=A_O$ and $\alpha=\varepsilon_\ell^{-k}$. Finite image of continuous cocycles in $A(\alpha)$ reduces the proof to $\ell$-torsion.

If $\ell-1\mid k$, the action on $A[\ell]$ is trivial. For a class killed by $\ell$, write $\ell z$ as the coboundary of $m$. Subtract the coboundary of an $\ell$-division point of $m$ from $z$. The resulting homomorphism to $A[\ell]$ vanishes on inertia away from $\ell$ by the local condition and trivial inertia action. Kronecker--Weber \cite[I, Theorem 4.16]{milne-2020} puts its field in a cyclotomic field. Trivial inertia at $p\ne\ell$ removes each $p$-power conductor factor. Both $z$ and $\alpha$ factor through $\Gal(\Q(\mu_{\ell^\infty})/\Q)$. The decomposition group at $\ell$ surjects onto this group, so the local coboundary formula holds globally.

If $\ell-1\nmid k$, then $k=(\ell-1)u/2$ with $u$ odd, and evenness of $k$ gives $\ell\equiv1\pmod4$. The residual character cuts out $K=\Q(\sqrt\ell)$ of discriminant $\ell$. Choose $\delta$ in inertia at $\ell$ with $\alpha(\delta)-1\in O^\times$ and normalize $z$ by a coboundary so that $z(\delta)=0$. For a class killed by $\ell$, evaluation at $\delta$ makes the coboundary $\ell z$ zero. Thus $z$ takes values in $A[\ell]$.

The restriction to $G_K$ is a homomorphism, unramified away from $\ell$. On each decomposition group over $\Q_\ell$, write $z(g)=(\alpha(g)-1)m$ by local vanishing. Its inertia contains $t$ with $\alpha(t)-1$ a unit because the residual cyclotomic character is surjective on inertia. Then $\ell z(t)=0$ gives $\ell m=0$. Since $G_K$ acts trivially on $A[\ell]$, the restriction vanishes above $\ell$. At real places the odd-primary homomorphism vanishes. Thus $z|_{G_K}$ factors through $\operatorname{Cl}(K)$ by unramified reciprocity \cite[V, Example 3.9]{milne-2020}.

Minkowski gives an integral ideal of norm at most $\sqrt\ell/2$ in each class \cite[Theorem 4.3]{milne-2020a}. Prime-power factorization bounds the number of ideals of norm $n$ by $d(n)$, the number of positive divisors of $n$. With $T=\sqrt\ell$,
\[
 \#\operatorname{Cl}(K)\le\sum_{n\le T}d(n)
 =\sum_{a\le T}\left\lfloor T/a\right\rfloor
 \le T(1+\log T)<T^2=\ell.
\]
Hence $z|_{G_K}=0$. Since $\delta$ represents the nontrivial coset, $z=0$.

In class Z, $\ell\ge5$ and $\theta=1$ forces $k$ even, so the subgroup with the conditions at $p\mid N$ also vanishes.
\end{proof}

\section{Summing back to the field}

\begin{theorem}[The primes in the fixed set]\label{B6.37}
Fix $\chi$, $\ell\in\mathcal B_\chi$, and $j\in\N$. There is $C_{\ell,j}\in\N$ depending on these data. Let $t\ge2$ and let the host be large for $t$. Assume $B\ge0$ satisfies the layer hypothesis of \ref{B6.32}. For each corresponding CM field $E$, the lengths $s^+_{\ell,k}(E)$ are finite for $1\le k<t$, and
\[
 \log\ell\sum_{1\le k<t}s^+_{\ell,k}(E)
 \le(Nt)^{C_{\ell,j}}+C_{\ell,j}B.
\]
\end{theorem}
\leanref{UpperTheory.\allowbreak{}V.\allowbreak{}HV.\allowbreak{}sumIn}{Development.\allowbreak{}AutomorphicArithmetic.\allowbreak{}AutomorphicVarieties.\allowbreak{}HeightAndVanishing.\allowbreak{}FinalLe}{86--263}
\leanref{UpperTheory.\allowbreak{}V.\allowbreak{}HV.\allowbreak{}redAll}{Development.\allowbreak{}AutomorphicArithmetic.\allowbreak{}AutomorphicVarieties.\allowbreak{}HeightAndVanishing.\allowbreak{}RedAll}{6--21}

\begin{proof}
Choose $O$ containing the $\varphi(N)$-th roots of unity. Write
$X(k)=X_{\ell,O}(k)$, $M=Nt$, and $n=[E:\Q]=[(\Z/N\Z)^\times:H]\le N$.
On the sign eigenspace, a cyclic factor generator $T$ of order $m$ satisfies
$T^m-1=\prod_{\zeta^m=1}(T-\zeta)=0$.
Since $\ker(AB)/\ker B$ embeds in $\ker A$, the length is at most the sum of the lengths of the kernels of $T-\zeta$. Iterate on the kernels over the cyclic factors. As $2\in O^\times$, the surviving characters are those in $X(k)$.

For each $\theta$, put $W=M_{\ell,k,\theta}$ and $\Delta=\Gal(E/\Q)$. The cokernel of $H^1(G_\Q,W)\to H^1(G_E,W)^\Delta$ injects by transgression into $H^2(\Delta,W^{G_E})$. The scalar-action bound for the abelian group $\Delta$ gives length at most $n^2\operatorname{length}_O W^{G_\Q}\le e_On^2\nu_\ell(N,k)$. The restriction kernel for decomposition at $\ell$ costs at most $e_On\nu_\ell(N,k)$, and each remaining inertia group at $p\mid N$ costs at most $e_Onv_\ell(n)$. The additional decomposition conditions at $p\mid N$ cost at most $e_O\omega(N)v_\ell(n)$ in total, where $\omega(N)$ counts distinct prime factors. Thus
\[
 e_Os^+_{\ell,k}(E)
 \le\sum_{\theta\in X(k)}
 \bigl(L_{\ell,k,\theta}+e_O\delta_\ell(N,n,k)\bigr),
 \tag{6.R}\label{eq:b6-redall}
\]
where
\[
 \delta_\ell(N,n,k)
 =(n+n^2)\nu_\ell(N,k)+\omega(N)(n+1)v_\ell(n).
\]
For $U=Nk\ge16$, use $n,v_\ell(n),\omega(N)\le U$ and $\nu_\ell(N,k)\le2U$ to get
\[
 \delta_\ell(N,n,k)\le3U^2+3U^3\le U^6 .
 \tag{6.L}\label{eq:b6-loss}
\]

Let $\mathcal O$ be the union over $1\le k<t$ of the characters in $X(k)$ with $\operatorname{ord}\theta_P$ odd. By \ref{B6.32} their conductors divide $f$, so change of level gives $|\mathcal O|\le\varphi(f)\le f$. Each such $\theta$ contributes at most $B+M^{C_4}$ over all twists by \ref{B6.32}. For the other characters, \ref{B6.31}, \ref{B6.34}, and \ref{B6.35} give
\[
 a_{\ell,k,\theta}\log\ell\le(\ell Nk)^b\le M^{2b},
 \qquad b=\max\{C_{\mathrm{II}},C_{\mathrm V},C_{\mathrm{Vs}}\}.
\]
Here $\ell<q\le N$, by largeness of the host and $\ell\in\mathcal B_\chi$. Every character length is finite, so \eqref{eq:b6-redall} gives finite field lengths.

Divide by $e_O$ and sum with weight $\log\ell$. Extend each $\mathcal O$ character sum to all twists by nonnegativity. For the remaining terms use $|X(k)|\le N$ and \eqref{eq:b6-loss} to obtain
\[
\begin{split}
 \log\ell\sum_{1\le k<t}s^+_{\ell,k}(E)
 &\le(\log\ell)f(B+M^{C_4})+tN M^{2b}+(\log\ell)NtM^6\\
 &\le \ell fB+M^{C_4+2}+M^{2b+1}+M^8.
\end{split}
\]
Since $M\ge3$, take $C_{\ell,j}=C_4+2b+9+\ell f$.
\end{proof}

\begin{theorem}[The primes outside the fixed set]\label{B6.38}
For each $\chi$ there are $C_{\mathrm{out}}\in\N$ and primes $r_1,r_2$. Let the host be large for $t\ge2$ with $q,d>r_1,r_2$, let $E$ be a corresponding CM field, and let $S$ be a finite set of odd primes disjoint from $\mathcal B_\chi$. Then $s^+_{\ell,k}(E)<\infty$ for $\ell\in S$ and $1\le k<t$, and
\[
 P^+(E,t;S)\le(Nt)^{C_{\mathrm{out}}}.
\]
\end{theorem}
\leanref{UpperTheory.\allowbreak{}V.\allowbreak{}HV.\allowbreak{}sumOut}{Development.\allowbreak{}AutomorphicArithmetic.\allowbreak{}AutomorphicVarieties.\allowbreak{}HeightAndVanishing.\allowbreak{}SumOut}{50--231}
\leanref{UpperTheory.\allowbreak{}V.\allowbreak{}HV.\allowbreak{}redPrime}{Development.\allowbreak{}AutomorphicArithmetic.\allowbreak{}AutomorphicVarieties.\allowbreak{}HeightAndVanishing.\allowbreak{}RedPrime}{11--57}

\begin{proof}
Here $n=[E:\Q]\mid2^{a+2}o$ and $\ell\nmid n$. Over $O_\ell$ containing the $\varphi(N)$-th roots of unity the character projectors are integral. Restriction identifies the eigenspaces with the twisted strict Selmer groups with the decomposition condition at $\ell$. The additional conditions at $p\mid N$ cost zero by $v_\ell(n)=0$. Thus \eqref{eq:b6-redall} holds with loss zero.

Take $r_1,r_2$ from \ref{B6.30}. Put $M=Nt$ and let $b=\max\{C_{\mathrm{II}},C_{\mathrm{III}}\}$. For each $k$, class I contributes at most $M^{C_{\mathrm I}}$. The other primes have $\ell\le2t-1$. At each of them class Z contributes zero by \ref{B6.36}, and \ref{B6.31} bounds each remaining character by $M^{3b}$ since $\ell Nk\le M^3$. There are at most $2t$ such primes and $N$ characters per prime. The finite sum over twists is at most $t(M^{C_{\mathrm I}}+2M^{3b+1})\le M^{C_{\mathrm I}+3b+4}$.
\end{proof}

\begin{theorem}[The pointwise bound and its common window]\label{B6.39}
For fixed $\chi$ and $j:\N\to\N$, there are $C\in\N$ and primes $r_1,r_2$. Let a host be large for $t\ge2$ with $q,d>r_1,r_2$. Assume $B\ge0$ satisfies the layer hypothesis of \ref{B6.32} with index $j(\ell)$ at each $\ell\in\mathcal B_\chi$. Then for every corresponding CM field $E$, odd prime $\ell$, and $1\le k<t$, the plus length is finite. For each finite set $S$ of odd primes,
\[
 P^+(E,t;S)\le(Nt)^C+CB.
 \tag{6.P}\label{eq:b6-point}
\]

Fix tower data $K=(R_0,R_\ell,C_K)$ with $C_K>0$. Use $M$ cut out by $\chi$ and $A_\ell=\Q(\mu_{\ell f})$ on $\mathcal B_\chi$. Set $j(\ell)=\lfloor\log_\ell R_\ell(A_\ell)\rfloor+3$ there, with the convention of \ref{B6.27}, and extend $j$ to $\N$. Take the corresponding $C,r_1,r_2$. Fix $C_t>0$, $0<\varepsilon<1/52$ and $A_1\in\R$, with $J$ as in \ref{B6.27}. Put $D_x=\exp(x^{4\varepsilon/5})$.

Let $a,t$ be admissible as in \ref{B6.27}. At every sufficiently large scale $X$, uniformly in $a,t$, take a prime pair as in \ref{B6.28} with
\[
 q,d>2X^{C_t}+1,\qquad
 q,d>\max\{f,o,p_K,h_*,p_1,r_1,r_2\},\qquad
 \log(qd)\le X^{\varepsilon/(4J)}.
\]
On any $K$-window $[Y,Y^2]$ satisfying the window conditions of \ref{B6.27}, for all $x$ in the window, all corresponding CM fields $E$, and all finite sets $S$ of odd primes, the plus lengths for $1\le k<t(x)$ are finite and
\[
 P^+(E,t(x);S)\le
 \bigl(\operatorname{cond}(E)t(x)\bigr)^C+CD_x.
 \tag{6.W}\label{eq:b6-pluswindow}
\]
For the tower data of \ref{B6.27} and $C_t\ge1$, such a window exists. At $C_t=401$ one may take $t(x)=\lfloor x^{401}\rfloor$, with the $2$-primary bound on the same interval.
\end{theorem}
\leanref{UpperTheory.\allowbreak{}V.\allowbreak{}HV.\allowbreak{}pointHV\_\allowbreak{}of\_\allowbreak{}counts}{Development.\allowbreak{}AutomorphicArithmetic.\allowbreak{}AutomorphicVarieties.\allowbreak{}HeightAndVanishing.\allowbreak{}JointWinOfCounts}{16--18}
\leanref{UpperTheory.\allowbreak{}V.\allowbreak{}HV.\allowbreak{}point}{Development.\allowbreak{}AutomorphicArithmetic.\allowbreak{}AutomorphicVarieties.\allowbreak{}HeightAndVanishing.\allowbreak{}PointArith}{36--92}
\leanref{UpperTheory.\allowbreak{}V.\allowbreak{}HV.\allowbreak{}thmHV}{Development.\allowbreak{}AutomorphicArithmetic.\allowbreak{}AutomorphicVarieties.\allowbreak{}HeightAndVanishing.\allowbreak{}ThmHV}{6--40}
\leanref{UpperTheory.\allowbreak{}V.\allowbreak{}HV.\allowbreak{}jointWin\_\allowbreak{}of\_\allowbreak{}counts}{Development.\allowbreak{}AutomorphicArithmetic.\allowbreak{}AutomorphicVarieties.\allowbreak{}HeightAndVanishing.\allowbreak{}JointWinOfCounts}{25--28}

\begin{proof}
Put $r=|\mathcal B_\chi|$, $D=\sum_{\ell\in\mathcal B_\chi}C_{\ell,j(\ell)}$ and $C_0=C_{\mathrm{out}}$. Split $S$ at $\mathcal B_\chi$ and apply \ref{B6.37} and \ref{B6.38}. The sum is at most $r(Nt)^D+DB+(Nt)^{C_0}$, which gives \eqref{eq:b6-point} with $C=C_0+D+r$ since $r+1\le(Nt)^r$. Singletons give finiteness at every odd prime.

On the window, $X^{1/J}\le Y$ and $Y^2\le X$, so $2\le x\le X$ for $X\ge2^J$. The prime inequalities make the host large for $t(x)$ and ensure $q,d>r_1,r_2$. For each $\ell\in\mathcal B_\chi$, the odd-layer bound on the given $K$-window supplies the hypothesis of \eqref{eq:b6-point} with $B=D_x$.

The character group of $E$ contains $\chi$ and the quadratic characters of conductors $q,d$. It contains a faithful character of $(\Z/2^{a+2}\Z)^\times/\{\pm1\}$, of conductor $2^{a+2}$ by its nontrivial value at $1+2^{a+1}$. These conductors have least common multiple $N$. Since $E$ is realizable in $\Q(\mu_N)$, $\operatorname{cond}(E)=N$ and \eqref{eq:b6-pluswindow} follows.
\end{proof}
\leanref{UpperTheory.\allowbreak{}V.\allowbreak{}HV.\allowbreak{}hostCond}{Development.\allowbreak{}AutomorphicArithmetic.\allowbreak{}AutomorphicVarieties.\allowbreak{}HeightAndVanishing.\allowbreak{}HostCond}{256--323}


% ---- p6c5 ----
\chapter{Automorphic counts and deformation rigidity}\label{ch:vwin-automorphic}

Write \(\mathscr H\) for the inputs of \ref{B6.05} and \ref{B6.06}, including the normalised Poitou--Tate theorem of \ref{B6.09}. The automorphic-to-Galois attachment in \ref{B9.02} starts with a cuspidal representation, as in Chenevier and Harris \cite[Theorem~3.2.3]{chenevier-harris-2013}. Lemma~\ref{B6.48a} obtains cuspidal attachments for irreducible arithmetic points of an ordinary Hecke quotient.

Put \(G_F=\Gal(\overline F/F)\) and \(C_\ell=\overline{\Q}_\ell\). Write \(\varepsilon_\ell\) for the cyclotomic character, with \(v_\ell(\ell)=1\) and \(|\ell|_\ell=\ell^{-1}\). Over a coefficient DVR \(O\subset C_\ell\) finite over \(\Z_\ell\), use the Selmer group \(\mathcal S_{\ell,k,\vartheta}\) of length \(L_{\ell,k,\vartheta}\) and the quantity \(a_{\ell,k,\vartheta}\) from \ref{B6.29}.

A continuous representation has \emph{level \((M;\mathbf m)\)} if away from \(\ell\) its inertia restrictions are diagonal with characters factoring through the action on \(\mu_M\). At every place above \(\ell\) it must have an ordinary flag with characters \(\kappa_i\) satisfying
\[
 \kappa_i|_I=\varepsilon_\ell^{m_i}\chi_i|_I,\qquad
 \chi_i|_{G_{F(\mu_M)}}=1.
\]
At such a place either \(\kappa_i\ne\kappa_j\varepsilon_\ell\) for all \(i,j\), or the representation is potentially crystalline. It is \emph{separated ordinary} if a flag satisfying these inequalities can be chosen at every place above \(\ell\).

\leanref{UpperTheory.\allowbreak{}V.\allowbreak{}Af.\allowbreak{}Level}{Development.\allowbreak{}AutomorphicArithmetic.\allowbreak{}AutomorphicVarieties.\allowbreak{}AutomorphicForms.\allowbreak{}BlochKatoSelmerVanishes}{111--115}
\leanref{UpperTheory.\allowbreak{}V.\allowbreak{}Th.\allowbreak{}Unl}{Development.\allowbreak{}AutomorphicArithmetic.\allowbreak{}AutomorphicVarieties.\allowbreak{}TheoremsMuLinked.\allowbreak{}AutP}{10--14}

Fix \(\iota:C_\ell\simeq\C\). Over a CM field the witness \Lean{WitA} requires potential unlinking at every unlinked place above \(\ell\). At a potentially unlinked place each prolongation admits a flag with \(\kappa_i^m\ne(\kappa_j\varepsilon_\ell)^m\) for all \(i,j\) and \(m\ge1\). Unlinking requires the condition for \(m=1\). For each regular algebraic cuspidal attachment through \(\iota\), the witness requires a common \(t\in\R\) at all \(v\nmid\ell\) and all potentially unlinked \(v\mid\ell\). There the local component has toral support and
\[
 q_v^{t-1/2}<|\xi(\varpi_v)|<q_v^{t+1/2}
\]
for every support character \(\xi\), where \(q_v\) is the residue cardinality and \(\varpi_v\) a uniformiser. Over \(\Q\), \Lean{WitAQ} requires this witness on each irreducible quadratic CM restriction. The condition \Lean{IsAutQA} for a continuous representation requires an attachment with witness on each such restriction.

\section{Trace differences and the first count}

Fix an imaginary quadratic field \(K\). A member \(\rho\) of \(\mathcal A_3(K,\iota;k,N)\) is irreducible and separated ordinary of level \((N;(k,0,-k))\). Its twist \(\rho\varepsilon_\ell^{-1}\) is conjugate self-dual automorphic with witness. In the block family \(\mathcal A_{2+1}(K,\iota;k,N|D_K|)\), a member is conjugate to
\[
 \begin{pmatrix} W_{00}&0&W_{01}\\0&1&0\\ W_{10}&0&W_{11}\end{pmatrix}\Big|_{G_K}.
\]
Here \(W\) is irreducible on \(G_\Q\), of level \((N|D_K|;(k,-k))\) and determinant of cyclotomic weight zero. The twist \(W\varepsilon_\ell^{-k}\) has a newform attachment of weight \(2k+1\), and \(W\) satisfies \Lean{WitAQ}.

\leanref{UpperTheory.\allowbreak{}V.\allowbreak{}Z.\allowbreak{}J.\allowbreak{}Adm3U\_\allowbreak{}s}{Development.\allowbreak{}AutomorphicArithmetic.\allowbreak{}AutomorphicVarieties.\allowbreak{}ModularityArgument.\allowbreak{}ThresholdArguments.\allowbreak{}DenseThresholdCases.\allowbreak{}NumberU3US}{27--29}
\leanref{UpperTheory.\allowbreak{}V.\allowbreak{}Z.\allowbreak{}Win.\allowbreak{}Dn.\allowbreak{}Adm2A}{Development.\allowbreak{}AutomorphicArithmetic.\allowbreak{}AutomorphicVarieties.\allowbreak{}ModularityArgument.\allowbreak{}AutomorphyLifting.\allowbreak{}AutomorphyDescent.\allowbreak{}AutNewformW}{14--20}

\begin{theorem}[A single-character trace bound]\label{B6.40}
Assume the cohomological inputs for the three-dimensional deformation problem, including \ref{B6.15}. Assume also the algebraic assertions of \ref{B6.58} and the vanishing of \ref{B6.43}. Assume rigidity for Types I and II and admissibility for Type I. Assume the endpoint automorphy assertions of \ref{B6.52} in the form used for the marked problem.

Let \(\ell\) be an odd prime, \(k\ge1\) and \(N\ge1\) even. Suppose \(\vartheta:(\Z/N)^\times\to O^\times\) satisfies
\[
 \ell\nmid\operatorname{ord}(\vartheta),\quad
 \operatorname{comp}_\ell(\vartheta)^2=1,\quad
 \vartheta(-1)=(-1)^k,\quad
 \ell\nmid |D_K|h_K,\quad (N,D_K)=1,\quad \ell-1\nmid2k.
\]
Here \(h_K\) is the class number and \(\operatorname{comp}_\ell\) is the modulus component at \(\ell\). Let \(\vartheta_P\) be the prime-to-\(\ell\) order part with its \(\ell\)-local component removed, and suppose
\[
 \ker\overline\varepsilon_\ell\cap\ker\vartheta_P\not\subseteq G_K.
\]
Put \(\Psi=\varepsilon_\ell^k\vartheta\). Fix \(\iota:C_\ell\simeq\C\) and \(\sigma\in G_K\), and put \(T=\Psi(\sigma)+1+\Psi(\sigma)^{-1}\). If \(\operatorname{tr}\rho(\sigma)\ne T\) throughout \(\mathcal A_3\cup\mathcal A_{2+1}\), then \(L_{\ell,k,\vartheta}<\infty\), and some finite pairwise nonisomorphic set \(X\) in this union satisfies
\[
 a_{\ell,k,\vartheta}
 \le2\sum_{\rho\in X}v_\ell\bigl(\operatorname{tr}\rho(\sigma)-T\bigr).
 \tag{6.40.1}
\]
\end{theorem}

\leanref{UpperTheory.\allowbreak{}V.\allowbreak{}Z.\allowbreak{}J.\allowbreak{}oneCharU\_\allowbreak{}of\_\allowbreak{}katoTwoW\_\allowbreak{}rig\_\allowbreak{}admI\_\allowbreak{}automorphy}{Development.\allowbreak{}AutomorphicArithmetic.\allowbreak{}AutomorphicVarieties.\allowbreak{}ModularityArgument.\allowbreak{}ThresholdArguments.\allowbreak{}DenseThresholdCases.\allowbreak{}FOneCharU}{75--95}
\leanref{UpperTheory.\allowbreak{}V.\allowbreak{}Z.\allowbreak{}J.\allowbreak{}OneCharU\_\allowbreak{}s}{Development.\allowbreak{}AutomorphicArithmetic.\allowbreak{}AutomorphicVarieties.\allowbreak{}ModularityArgument.\allowbreak{}ThresholdArguments.\allowbreak{}DenseThresholdCases.\allowbreak{}NumberU3US}{49--59}

\begin{proof}
The values of \(\vartheta\) have order prime to \(\ell\) and generate an unramified coefficient ring. Work over it by \ref{B6.29}. The auxiliary-field conditions give the marked problem with \(r_0|_{G_K}=\operatorname{diag}(\Psi,1,\Psi^{-1})\). By \ref{B6.40a} and (6.58.2), every nonzero \(b\in O\) gives \(\operatorname{length}_O\mathcal S_{\ell,k,\vartheta}[b]\le\operatorname{length}_O C_{-1}(F)\). A continuous cocycle from the compact Galois group to the discrete divisible coefficient module has finite image killed by a power of the uniformiser. The Selmer group is the increasing union of its uniformiser-power torsion submodules. Taking the supremum gives
\[
 L_{\ell,k,\vartheta}\le\operatorname{length}_O C_{-1}(F).
 \tag{6.40.2}
\]
Apply the orbit estimate of \ref{B6.58} to
\[
 a=\operatorname{tr}r(\sigma)-\bigl(\Psi(\sigma)+1+\Psi(\sigma)^{-1}\bigr).
 \tag{6.40.3}
\]
The marked residue-characteristic endpoint of \ref{B6.52} supplies weight-zero finiteness. At Type I points use \ref{B6.53}. At Type II points second-case rigidity treats the block adjoint and \ref{B6.43} treats the off-diagonal \(W\) and \(W^\vee\simeq W\otimes\det W\). The cohomological input gives the scalar and one-sided conditions, so the tangent space at each two-sided point has dimension at most one.

Type I admissibility and the Type II extension of \ref{B6.52} place the two-sided points in the two families. Thus \(a\) is nonzero there, and \ref{B6.58} bounds \(\operatorname{length}_O C_{-1}(F)\) by \(2\sum_\xi v_O(\xi(a))\) over a finite set of orbit representatives. Isomorphic \(G_K\)-representations give points in the same marked torus orbit. The representatives therefore give a pairwise nonisomorphic set \(X\). Since \(v_O=v_\ell\) and \(O/\ell O\) has length one, (6.40.2) gives (6.40.1) with the coefficient comparison of \ref{B6.29}.
\end{proof}

\begin{lemma}[Integral lift on the ambient algebra]\label{B6.40a}
In the marked problem of \ref{B6.40}, let \(S\) be the normalised global ring, \(x_0:S\to O\) the point of \(r_0\), and \(q:A\to S\) the graded model map. Put \(x=x_0q\). For \(0\ne b\in O\), \(\mathcal S_{\ell,k,\vartheta}[b]\) injects additively into the \(O/(b)\)-valued \(x\)-derivations of \(A\) killing the equations \(F_i\). Each derivation in the image lifts to an \(O\)-valued \(x\)-derivation on \(A\) supported in weight \(-1\).
\end{lemma}
\begin{proof}
Put \(\alpha=\Psi^{-1}\) and let \(\eta\) be the quadratic character of \(K\). Choose the marking \(g_1\in\ker\eta\) with \(\Psi(g_1)^2-1\in O^\times\), and choose \(h_*\) in the distinguished local group with \(\Psi(h_*)^2-1\in O^\times\). Since \(\alpha(g_1)-1\) is a unit, subtracting the coboundary of \(z(g_1)/(\alpha(g_1)-1)\) gives the unique representative vanishing at \(g_1\). This choice is additive, so a class killed by \(b\) has representative killed by \(b\). Choose a generator identifying the divisible module's \(b\)-torsion with \(O/(b)\), and write \(t\) for the resulting cocycle. On the distinguished local group \(t(h)=s-\alpha(h)s\) for \(s\in O/(b)\). The Selmer condition supplies a cobounding element in the divisible module, and evaluation at \(h_*\) shows that \(b\) kills it. On the prescribed kernels \(\alpha=1\) and the local coboundaries vanish. From \(c^2=1\) and \(\Psi(c)=1\) we get \(2t(c)=0\), hence \(t(c)=t(cg_1^{-1}c)=0\).

For matrix units numbered from one, put \(V=e_{21}-e_{32}\). The polarised adjoint action on \(V\) is \(\alpha\), since negative reflected transpose fixes \(V\). Thus \(\gamma=tV\) is a normalised adjoint cocycle. Lift \(s\) to \(\widetilde s\in O\). Conjugating \(r_0\) by \(1+\epsilon\widetilde sV\), with \(\epsilon^2=0\), gives the integral local cocycle \(\widetilde sV-\operatorname{Ad}(r_0(h))\widetilde sV\) reducing to \(\gamma(h)\). The local condition contains \(r_0\) and is stable under the conjugation. The coefficient of \(\epsilon\) gives an integral derivation \(D_B\) on the local condition ring \(B\).

Over \(O\oplus\epsilon O/(b)\), \((1+\epsilon\gamma)r_0\) satisfies the local condition by reduction and the marking and kernel equations by normalisation. Its universal map is \((x_0,D)\) with
\[
 D(r(g)_{ij})=\Psi(g)^{w_j}\gamma(g)_{ij},\qquad w=(1,0,-1).
\]
Write \(B=O[[X]]/H\) in centred local entry coordinates. The ambient ring \(P=O[[X,Y]]/HO[[X,Y]]\) has new global entry coordinates \(Y\) and a surjection \(\widehat q:P\to S\) restricting to \(q\) on the sum \(A\) of its homogeneous parts. On the variables \(Z_i=(X,Y)_i\), take values \(d_i=D_B(X_i)\) for old variables, arbitrary lifts of \(D(\widehat qY_i)\) for new variables of weight \(-1\), and zero for the other new variables. These reduce to the prescribed derivatives. For the dual block use \(d(U^{-1})=-U_0^{-1}(dU)U_0^{-1}\) at the base matrix, followed by reflected transpose.

The formula \(E(f)=\sum_i[Z_i]f\,d_i\), with \([Z_i]f\) the linear coefficient of \(Z_i\), defines a derivation at the origin. On the old series it agrees with the pullback of \(D_B\), so \(E(h)=h(0)=0\) for \(h\in H\) and \(E(uh)=0\) for every \(u\). It descends to \(P\). Modulo \(b\) it agrees with \(D\widehat q\) on the variables and hence on all series by the linear coefficient formula. On \(A\), \(E\) has weight \(-1\), and its reduction \(Dq\) kills the \(F_i\in\ker q\). The universal entries have preimages in \(A\) and determine the square-zero point. Their derivatives give additivity and recover \(t\).
\end{proof}
\leanref{UpperTheory.\allowbreak{}V.\allowbreak{}U3.\allowbreak{}sha\_\allowbreak{}inj\_\allowbreak{}gen}{Development.\allowbreak{}AutomorphicArithmetic.\allowbreak{}AutomorphicVarieties.\allowbreak{}UnitaryThreeGroups.\allowbreak{}ShaInjGen}{338--362}
\leanref{UpperTheory.\allowbreak{}exists\_\allowbreak{}base\_\allowbreak{}pointDerivation\_\allowbreak{}of\_\allowbreak{}weight\_\allowbreak{}homogeneous}{Development.\allowbreak{}AutomorphicArithmetic.\allowbreak{}AutomorphicVarieties.\allowbreak{}UnitaryThreeGroups.\allowbreak{}ExistsBasePointDerivationOfWeightHomogeneous}{48--135}
\leanref{UpperTheory.\allowbreak{}V.\allowbreak{}U3.\allowbreak{}model\_\allowbreak{}lower}{Development.\allowbreak{}AutomorphicArithmetic.\allowbreak{}AutomorphicVarieties.\allowbreak{}UnitaryThreeGroups.\allowbreak{}ModelLower}{373--402}

\begin{theorem}[An integer for three-dimensional trace differences]\label{B6.41}
Assume the automorphic inputs of \(\mathscr H\). For every imaginary quadratic \(K\) there are \(c,A\in\N\) with the following property. Let \(k,N\ge1\), let \(p\ge7\) be a prime with \(p\nmid N|D_K|\), and take \(w\mid p\) in \(K\). Put \(q_w=\mathrm Nw\). Let \(z:\mathcal O\to\C\) be a finite indexed family of roots of unity permuted by every automorphism of \(\C\). There is \(D\in\Z\setminus\{0\}\), chosen before \(\ell\), with
\[
 \log|D|\le c\,\#\mathcal O\,[N(k+1)]^A\log p.
 \tag{6.41.1}
\]
For every prime \(\ell\ne2,p\), every \(\iota:C_\ell\simeq\C\), every Frobenius \(\sigma\) at \(w\) and all finite sets \(X_u\subset\mathcal A_3\cup\mathcal A_{2+1}\) of pairwise nonisomorphic representations, indexed by \(u\in\mathcal O\), the differences
\[
 \Delta_{u,\rho}=\operatorname{tr}\rho(\sigma^{-1})
 -\iota^{-1}\!\left(q_w^{-k}z_u+1+q_w^kz_u^{-1}\right)
\]
are nonzero and satisfy
\[
 \sum_{u\in\mathcal O}\sum_{\rho\in X_u}v_\ell(\Delta_{u,\rho})\le v_\ell(D).
 \tag{6.41.2}
\]
\end{theorem}

\leanref{UpperTheory.\allowbreak{}GalN.\allowbreak{}V.\allowbreak{}Win.\allowbreak{}N.\allowbreak{}numberU3U\_\allowbreak{}of\_\allowbreak{}prints}{Development.\allowbreak{}AutomorphicArithmetic.\allowbreak{}AutomorphicVarieties.\allowbreak{}ModularityArgument.\allowbreak{}ThresholdArguments.\allowbreak{}DenseThresholdCases.\allowbreak{}KeepNumU}{79--80}
\leanref{UpperTheory.\allowbreak{}V.\allowbreak{}Z.\allowbreak{}J.\allowbreak{}NumberU3U\_\allowbreak{}s}{Development.\allowbreak{}AutomorphicArithmetic.\allowbreak{}AutomorphicVarieties.\allowbreak{}ModularityArgument.\allowbreak{}ThresholdArguments.\allowbreak{}DenseThresholdCases.\allowbreak{}NumberU3US}{30--45}

\begin{proof}
The automorphic families give conjugation-stable label sets \(S_1,S_2\). Each label \(s\in S=S_1\sqcup S_2\) determines the isomorphism class and records \(t(s)=\iota(\operatorname{tr}\rho(\sigma^{-1}))\). In the block case the trace includes the trivial summand. The bounds are
\[
 |t(s)|\le\sqrt{q_w}+2,\qquad q_w^kt(s)\text{ is an algebraic integer},
 \tag{6.41.3}
\]
\[
 \#S_1\le c_1[N(k+1)]^{A_1},\qquad
 \#S_2\le c_2[N|D_K|(k+1)]^{A_2}.
 \tag{6.41.4}
\]
Put \(q=q_w\) and \(R=q^k\). The integral factor at \((u,s)\) is
\[
 F_{u,t(s)}=R\bigl(t(s)-R^{-1}z_u-1-Rz_u^{-1}\bigr)
 =Rt(s)-z_u-R-R^2z_u^{-1}.
 \tag{6.41.5}
\]
As \(R\ge q\ge7\), \(\sqrt q\le q/2\) and \(3+1/R<q/2\). Thus \(|Rt(s)-z_u-R|\le R(\sqrt q+3)+1<R^2\). The last term has absolute value \(R^2\), so the factor is nonzero, with
\[
 \log|F_{u,t(s)}|\le2k\log q_w+1.
 \tag{6.41.6}
\]
The product \(D=\prod_{u,s}F_{u,t(s)}\) is a nonzero algebraic integer fixed by every complex field automorphism. It belongs to \(\Z\) and
\[
 \log|D|\le\#\mathcal O\,\#S\,(2k\log q_w+1).
 \tag{6.41.7}
\]
Since \(q_w=p^f\), \(1\le f\le2\), (6.41.4) gives (6.41.1) with \(c=5(c_1+c_2|D_K|^{A_2})\), \(A=A_1+A_2+1\). For each \(u\) the elements of \(X_u\) have distinct labels, with corresponding factors \(q_w^k\Delta_{u,\rho}\) under \(\iota^{-1}\). The multiplier is an \(\ell\)-adic unit and all other factors have nonnegative valuation. This gives (6.41.2).

\leanref{UpperTheory.\allowbreak{}GalN.\allowbreak{}V.\allowbreak{}Win.\allowbreak{}frobenius\_\allowbreak{}trace\_\allowbreak{}norm\_\allowbreak{}nonzero\_\allowbreak{}log\_\allowbreak{}bound}{Development.\allowbreak{}AutomorphicArithmetic.\allowbreak{}Foundations.\allowbreak{}GaloisRepresentations.\allowbreak{}AdmissibleResidues.\allowbreak{}FrobeniusTraceCountBoundOfRatEigensystemCount}{329--400}
\end{proof}

\begin{theorem}[One integer for a character orbit]\label{B6.42}
Assume \ref{B6.40} and \ref{B6.41}. For every number field \(K\) there are \(c_K,A_K\in\N\) as follows. Let \(k,N\ge1\), put \(G=(\Z/N)^\times\), and take \(H\le G\). For every prime \(p\nmid30N|D_K|\) there is \(D\in\Z\setminus\{0\}\) with
\[
 \log|D|\le c_K\varphi([G:H])[N(k+1)]^{A_K}\log p.
 \tag{6.42.1}
\]
For every prime \(\ell\ne p\) and coefficient ring \(O\), if the \(O\)-valued orbit of kernel \(H\) contains a character satisfying the auxiliary-field conditions of \ref{B6.40}, all lengths in the orbit are finite and
\[
 \sum_{\ker\vartheta=H}a_{\ell,k,\vartheta}\le2v_\ell(D).
 \tag{6.42.2}
\]
\end{theorem}

\leanref{UpperTheory.\allowbreak{}V.\allowbreak{}Z.\allowbreak{}J.\allowbreak{}propGK\_\allowbreak{}of\_\allowbreak{}numberU3U\_\allowbreak{}of\_\allowbreak{}oneCharU}{Development.\allowbreak{}AutomorphicArithmetic.\allowbreak{}AutomorphicVarieties.\allowbreak{}ModularityArgument.\allowbreak{}ThresholdArguments.\allowbreak{}DenseThresholdCases.\allowbreak{}FPropGKU}{110--172}
\leanref{UpperTheory.\allowbreak{}V.\allowbreak{}U3.\allowbreak{}PropGK}{Development.\allowbreak{}AutomorphicArithmetic.\allowbreak{}AutomorphicVarieties.\allowbreak{}UnitaryThreeGroups.\allowbreak{}KatoUse}{36--43}

\begin{proof}
For \(K\) outside the imaginary quadratic case the good data are empty and \(D=1\) works. Otherwise choose \(w\mid p\). The complex characters \(\mathcal O\) of kernel \(H\) are permuted by complex field automorphisms. There are \(\varphi([G:H])\) if \(G/H\) is cyclic and zero otherwise. Use \ref{B6.41} with \(z_\vartheta=\vartheta(\mathrm Nw)^{-1}\).

A good character of kernel \(H\) supplies the roots of unity for the full \(O\)-valued orbit, identified by \(\iota\) with \(\mathcal O\). Goodness passes to the orbit. At inverse Frobenius the reducible trace is \(\iota^{-1}(q_w^{-k}z_\vartheta+1+q_w^kz_\vartheta^{-1})\). Apply \ref{B6.40} to each character using nonvanishing in \ref{B6.41}, and sum by (6.41.2).

Take two distinct auxiliary primes with integers \(D,D'\). At least one bound applies for each \(\ell\) in the set \(S_H\) of primes admitting a good character of kernel \(H\). Thus
\[
 \sum_{\ell\in S_H}\log\ell\sum_{\ker\vartheta=H}a_{\ell,k,\vartheta}
 \le2\log|D|+2\log|D'|.
 \tag{6.42.3}
\]
The subsets of characters counted by the occurring kernels are disjoint, so
\[
 \sum_{H\ \mathrm{occurring}}\varphi([G:H])\le\#G\le N.
 \tag{6.42.4}
\]

\leanref{UpperTheory.\allowbreak{}V.\allowbreak{}U3.\allowbreak{}corGK}{Development.\allowbreak{}AutomorphicArithmetic.\allowbreak{}AutomorphicVarieties.\allowbreak{}UnitaryThreeGroups.\allowbreak{}CorGSumLog}{37--74}
\leanref{UpperTheory.\allowbreak{}V.\allowbreak{}U3.\allowbreak{}countI}{Development.\allowbreak{}AutomorphicArithmetic.\allowbreak{}AutomorphicVarieties.\allowbreak{}UnitaryThreeGroups.\allowbreak{}CountI}{22--137}
\end{proof}

\section{The two-dimensional vanishing input}

A character of \(G_\Q\) factors through \(\mu_M\) if its kernel contains \(G_{\Q(\mu_M)}\).

\begin{theorem}[Kato vanishing for two quadratic companions]\label{B6.43}
Assume \(\mathscr H\). Let \(\ell\) be odd, \(k\ge1\), and \(\ell-1\nmid2k\). Let \(K\) be imaginary quadratic with quadratic character \(\theta\), and suppose
\[
 \ell\nmid |D_K|h_K,\qquad N\ge1,\quad 2\mid N,\quad (N,D_K)=1.
\]
Let \(\vartheta:G_\Q\to C_\ell^\times\) have order prime to \(\ell\), factor through \(\mu_N\), and have square trivial on each inertia group above \(\ell\). Put \(\Psi=\varepsilon_\ell^k\vartheta\) and require \(\Psi(c)=1\) for a complex conjugation \(c\). Let \(W:G_\Q\to\GL_2(C_\ell)\) have level \((N|D_K|;(k,-k))\), determinant \(\theta\), and traces
\[
 \operatorname{tr}W(s)=\Psi(s)+\Psi(s)^{-1}
 \qquad(s\in I_A,\ A\mid\ell),
 \tag{6.43.1}
\]
\[
 \left|\operatorname{tr}W(g)-\bigl(\theta(g)\Psi(g)+\Psi(g)^{-1}\bigr)\right|_\ell<1
 \qquad(g\in G_\Q).
 \tag{6.43.2}
\]
Require the following.
\begin{enumerate}
\item For some \(A\mid\ell\) and \(M\ge1\), the characteristic polynomial of \(W(s)\) is \((X-\varepsilon_\ell(s)^k)(X-\varepsilon_\ell(s)^{-k})\) whenever \(s\in I_A\) fixes \(\mu_M\).
\item At some \(A\mid\ell\), \(\theta|_{I_A}=1\) and \(W|_{D_A}\) has an ordinary flag with ordered inertia characters \(\varepsilon_\ell^k\chi,\varepsilon_\ell^{-k}\chi\), where \(\chi\) factors through \(\mu_\ell\).
\item The restriction \(W|_{G_K\cap\ker\varepsilon_\ell}\) is irreducible.
\end{enumerate}
Fix \(\iota:C_\ell\simeq\C\). If \(W\varepsilon_\ell^{-k}\) has a newform attachment of weight \(2k+1\) and \(W\) satisfies \Lean{WitAQ}, then
\[
 H_f^1(\Q,W)=H_f^1(\Q,W\otimes\theta)=0.
 \tag{6.43.3}
\]
\end{theorem}

\leanref{UpperTheory.\allowbreak{}V.\allowbreak{}Z.\allowbreak{}Kato.\allowbreak{}WeilSphericalRecords.\allowbreak{}KatoTwoW}{Development.\allowbreak{}AutomorphicArithmetic.\allowbreak{}AutomorphicVarieties.\allowbreak{}ModularityArgument.\allowbreak{}KatoQuadraticCharacters.\allowbreak{}SphericalDecomposition.\allowbreak{}KatoTwoW}{10--15}
\leanref{UpperTheory.\allowbreak{}V.\allowbreak{}Z.\allowbreak{}Kato.\allowbreak{}Used}{Development.\allowbreak{}AutomorphicArithmetic.\allowbreak{}AutomorphicVarieties.\allowbreak{}ModularityArgument.\allowbreak{}KatoQuadraticCharacters.\allowbreak{}KatoTwoOfUseN}{11--42}
\leanref{UpperTheory.\allowbreak{}V.\allowbreak{}Z.\allowbreak{}Kato.\allowbreak{}WeilSphericalRecords.\allowbreak{}katoTwoW\_\allowbreak{}of\_\allowbreak{}printsH}{Development.\allowbreak{}AutomorphicArithmetic.\allowbreak{}AutomorphicVarieties.\allowbreak{}ModularityArgument.\allowbreak{}ThresholdArguments.\allowbreak{}DenseThresholdCases.\allowbreak{}KatoTwoWOfPublishedH}{48--49}
\leanref{UpperTheory.\allowbreak{}V.\allowbreak{}Z.\allowbreak{}Kato.\allowbreak{}WeilSphericalRecords.\allowbreak{}katoTwoW\_\allowbreak{}of\_\allowbreak{}printsH\_\allowbreak{}of\_\allowbreak{}rT\_\allowbreak{}of\_\allowbreak{}quadBC2}{Development.\allowbreak{}AutomorphicArithmetic.\allowbreak{}AutomorphicVarieties.\allowbreak{}ModularityArgument.\allowbreak{}ThresholdArguments.\allowbreak{}DenseThresholdCases.\allowbreak{}KatoTwoWOfPublishedH}{38--47}

\begin{proof}
The newform attachment and \Lean{Used} give ordinary admissibility. For each finite cyclotomic twist passing the flag test, local compatibility and the witness give a spherical component at \(\ell\). The spherical-source theorem supplies the source of \ref{B6.44}, which applies with \(\det W=\theta\).
\end{proof}

\begin{lemma}[Applying the ordinary theorem twice]\label{B6.44}
Assume the cohomological inputs of \(\mathscr H\) and the unramified-source theorem for admissible spherical sources. Take \Lean{Used} and the newform attachment of \ref{B6.43}. Suppose \(W\) satisfies \Lean{SphAt}. Namely, for every \(M\ge1\) and every \(\phi\) factoring through \(\mu_M\), at a place \(A\mid\ell\) where \(W\otimes\phi\) has a flag with ordered inertia characters \(\varepsilon_\ell^k,\varepsilon_\ell^{-k}\), every regular algebraic cuspidal attachment of \(W\otimes\phi\) is spherical at the prime below \(A\). Then \(H_f^1(\Q,W)=H_f^1(\Q,W\otimes\det W)=0\).
\end{lemma}

\leanref{UpperTheory.\allowbreak{}V.\allowbreak{}Z.\allowbreak{}Kato.\allowbreak{}WeilSphericalRecords.\allowbreak{}twoS}{Development.\allowbreak{}AutomorphicArithmetic.\allowbreak{}AutomorphicVarieties.\allowbreak{}ModularityArgument.\allowbreak{}KatoQuadraticCharacters.\allowbreak{}SphericalDecomposition.\allowbreak{}SphAtTw}{16--35}
\leanref{GlobalTheory.\allowbreak{}Ord2.\allowbreak{}vanishOrdL\_\allowbreak{}of\_\allowbreak{}zetaUr}{Development.\allowbreak{}AbelianVarietyTheory.\allowbreak{}Proofs.\allowbreak{}FoundationProofs.\allowbreak{}VanishOrdLOfZetaUr}{20--26}

\begin{proof}
The unramified source puts \(W\) in the normal form for \ref{B6.46} and \ref{B6.47}. Put \(\delta=\det W\) and choose \(M_\delta\) with \(\delta\) factoring through \(\mu_{M_\delta}\). The test for \(W\otimes\delta\) at \(\phi\) of modulus \(M\) is the test for \(W\) at \(\phi\delta\) of modulus \(MM_\delta\). Determinant-twist admissibility supplies the other hypotheses, since \(\det(W\otimes\delta)=\delta\) by \(\delta^2=1\). The same two theorems apply to its unramified source.
\end{proof}

At \(A\mid\ell\) write \(W|_{D_A}\) in its ordinary flag basis with diagonal characters \(\alpha,\beta\), in that order. The normal form \(\operatorname{Nf}(k,S,A,W)\) means that \(\ell\) is odd, \(k\ge1\), and \(W\) is continuous over an integral coefficient ring and unramified outside the finite set \(S\ni\ell\). Its determinant is quadratic and odd, and it acts irreducibly on \(\ker\varepsilon_\ell\cap\ker\det W\). For some \(h\ge1\),
\[
 \alpha(s)^h=\varepsilon_\ell(s)^{kh}\quad(s\in I_A),
 \qquad \alpha\ne\varepsilon_\ell\quad\text{on }D_A .
 \tag{6.44.1}
\]
A source \(\operatorname{SrcUr}(\iota,k,A,W,g)\) is a newform \(g\) of weight \(2k+1\) and level \(M_g\) prime to \(\ell\), with \(|\iota^{-1}(a_\ell(g))|_\ell=1\). For some \(\chi\) factoring through \(\mu_\ell\), it is attached to \(W\otimes\varepsilon_\ell^{-k}(\det W)^{-1}\chi\), with \(\beta\varepsilon_\ell^k(\det W)^{-1}\chi=1\) on \(I_A\).

Number coordinates \(0,1\) and put
\[
 \operatorname{Gr}_A(y)\iff
 \exists b\in C_\ell,\quad y(g)_1=(\beta(g)-1)b\quad(g\in D_A).
 \tag{6.44.2}
\]
Write \(\operatorname{Unr}(y)\) for local triviality on inertia at each prime different from \(\ell\).

\leanref{GlobalTheory.\allowbreak{}Ord2.\allowbreak{}Nf}{Development.\allowbreak{}AbelianVarietyTheory.\allowbreak{}Definitions.\allowbreak{}Foundations.\allowbreak{}FinalTheorems}{53--67}
\leanref{GlobalTheory.\allowbreak{}Ord2.\allowbreak{}Gr}{Development.\allowbreak{}AbelianVarietyTheory.\allowbreak{}Definitions.\allowbreak{}Foundations.\allowbreak{}FinalTheorems}{72--73}
\leanref{GlobalTheory.\allowbreak{}Ord2.\allowbreak{}Unr}{Development.\allowbreak{}AbelianVarietyTheory.\allowbreak{}Definitions.\allowbreak{}Foundations.\allowbreak{}FinalTheorems}{68--70}

\begin{theorem}[Integral cyclotomic Euler classes]\label{B6.45}
Fix \(\operatorname{Nf}(k,S,A,W)\) and \(\operatorname{SrcUr}(\iota,k,A,W,g)\). Choose an analytic modular-function tower \(X\) with holomorphic \(j\)-function on the upper half-plane, a Hecke projection \(\operatorname{pr}:X.\operatorname{Fn}_\ell\to C_\ell^2\), and \(N\ge1\) prime to \(\ell\). Take \(W\) over \(O\), a nonzero weight function \(c_w:\Z/N\to O\) and \(c,d>1\) prime to \(6\ell N\). For each \(r_0\in(\Z/\ell)^\times\) take \(g_{r_0}\in\Gamma_1(N)\) with \((0,1)g_{r_0}\equiv t(r_0,-1)\bmod\ell\) for some \(t\in(\Z/\ell)^\times\). Take a character function \(\psi_0\) on \(\Z/\ell\) with \(\beta\varepsilon_\ell^k(\det W)^{-1}(\psi_0\circ\bar\varepsilon_\ell)=1\) on \(I_A\). Require \(\operatorname{pr}\) to be \(\Z_\ell\)-linear, integral on the weight-\(2k-1\) module and equivariant there for \(\rho_k\) below. For primes \(q\nmid\ell N\) require
\[
 \operatorname{pr}(T_qb)=\iota^{-1}(a_q(g))\operatorname{pr}(b),\qquad
 \operatorname{pr}(\langle q\rangle b)=\delta(q)\operatorname{pr}(b)
\]
for scalars \(\delta(q)\).

There are \(c_0,m\ge0\), a coefficient ring \(O'\supseteq O\), a finite set \(S'\), and functions \(Y_r:G_\Q\times(\Z/\ell^r)^\times\to C_\ell^2\) for \(r\ge0\) and \(y:G_\Q\to C_\ell^2\) with the following properties.
\begin{enumerate}
\item The coordinates of \(Y_r\) are continuous \(O\)-valued functions of \(s\). With \(\bar\varepsilon_{\ell^r}\) the cyclotomic character modulo \(\ell^r\),
\[
 \begin{split}
 \rho_k(s)&=\bigl(\varepsilon_\ell(s)^{-k}\psi_0(\bar\varepsilon_\ell(s))\bigr)^{-1}W(s),\\
 Y_r(st,u)&=Y_r(s,u)+\rho_k(s)Y_r(t,u\bar\varepsilon_{\ell^r}(s)^{-1}).
 \end{split}
\]
Write \(Y_r=\sum_x c_w(x)Y_{x,r}\). Lifts of the smoothed modular-unit classes at exponent \(c_0\) give projected representatives \(Y^{\rm proj}_{x,r,n}\) with integral zero-cochains \(b_{x,r,n}\) such that
\[
 Y_{x,r}(s,u)\equiv Y^{\rm proj}_{x,r,n}(s,u)
 +\rho_k(s)b_{x,r,n}(u\bar\varepsilon_{\ell^r}(s)^{-1})-b_{x,r,n}(u)
 \pmod{\ell^n}.
\]
\item For reduction \(\operatorname{red}_r:(\Z/\ell^{r+1})^\times\to(\Z/\ell^r)^\times\),
\[
 Y_r(s,u)=\sum_{\operatorname{red}_r(u')=u}Y_{r+1}(s,u').
 \tag{6.45.1}
\]
\item At every \(r\) the second coordinates are uniformly locally constant to each prescribed precision on a finite Galois quotient. For the least nonnegative representative \(\widetilde u\) of \(u\), put \(\chi_r(u)=\psi_0(\widetilde u\bmod\ell)\). Then \(\sum_u(\widetilde u^{\,k})^{-1}\chi_r(u)Y_r(s,u)\equiv y(s)\bmod\ell^r\), with inversion at zero taken to be zero at \(r=0\).
\item The class \(\ell^my\) satisfies the global tame Euler relations outside \(S'\) over \(O'\). The functions \(\ell^mY_r(s,u)_1\) form the local Iwasawa system whose weighted specialisations converge to \(\ell^my(s)_1\) uniformly for \(s\in I_A\cap\ker\varepsilon_\ell\).
\end{enumerate}
\end{theorem}

\leanref{UpperTheory.\allowbreak{}Kd.\allowbreak{}K4.\allowbreak{}kato\_\allowbreak{}eulerSystem\_\allowbreak{}integrality}{Development.\allowbreak{}AutomorphicArithmetic.\allowbreak{}Foundations.\allowbreak{}IwahoriReduction.\allowbreak{}NormProducts.\allowbreak{}AdmissibleWeightedDataMonoCoef}{208--209}
\leanref{UpperTheory.\allowbreak{}V.\allowbreak{}Z.\allowbreak{}Kato.\allowbreak{}Joint.\allowbreak{}KatoEulerSystemIwasawaClassStatement}{Development.\allowbreak{}AutomorphicArithmetic.\allowbreak{}AutomorphicVarieties.\allowbreak{}ModularityArgument.\allowbreak{}KatoQuadraticCharacters.\allowbreak{}FrobeniusHodgeTateLaws.\allowbreak{}SmoothedZetaClassHodgeTateLawStatement}{11--30}

\begin{proof}
For fixed \(x\) take the complete weight lattice \(B=X.B(N,\ell,2k-1)\), \(A_r=(\Z/\ell^r)^\times\), and \(\kappa_r=\bar\varepsilon_{\ell^r}^{-1}\). Kato's newform norm relation is in \(H^1\), with Euler factors for newly omitted primes \cite[\S8.11 and Prop.~8.12]{kato-2004}. Here the ordinary projection is applied after the integral representatives have been chosen. The modular-unit distribution identity gives \(P_r F_{r+1,n}\sim F_{r,n}\) at precision \(\ell^n\), where \(P_r\) sums over fibres and \(\sim\) denotes equality up to coboundary. The top-level reductions commute, so the successive sums count each lift once. Reducing the lifted symbol from precision \(n+1\) to \(n\) gives \(F_{r,n+1}\sim F_{r,n}\).

The lifts have a common denominator exponent \(c_L\). Evaluating the weight polynomial at \((i,1)\), \(0\le i\le2k-1\), and multiplying by the adjugate of \(V=(i^j)\) recovers its coefficients times \(\det V\). Increasing the exponent by \(v_\ell(\det V)\) suffices uniformly to obtain both correcting coboundaries over \(B\). Lemma~\ref{B6.45a} gives compatible representatives. Project and sum over \(x\) with \(c_w(x)\) at a common exponent to obtain (6.45.1). Apply the projected tame relation to these representatives. Its common exponent \(a\) replaces \(Y\) by \(\ell^aY\) and raises the representative exponent from \(c_0\) to \(c_0+a\). Apply \Lean{middle\_\allowbreak{}glue2} to the inverse-weight specialisations of this rescaled tower. It gives \(y\) with the congruences of assertion 3 and an exponent \(m\) for which \(\ell^my\) satisfies the tame relations and \(\ell^mY_r(s,u)_1\) has the local specialisation of assertion 4.
\end{proof}
\leanref{UpperTheory.\allowbreak{}V.\allowbreak{}Z.\allowbreak{}Kato.\allowbreak{}Nb.\allowbreak{}kato\_\allowbreak{}norm\_\allowbreak{}compatible\_\allowbreak{}cyclotomic\_\allowbreak{}tower}{Development.\allowbreak{}AutomorphicArithmetic.\allowbreak{}AutomorphicVarieties.\allowbreak{}ModularityArgument.\allowbreak{}KatoQuadraticCharacters.\allowbreak{}EulerFactorBounds.\allowbreak{}OneLeLevel}{175--222}
\leanref{UpperTheory.\allowbreak{}V.\allowbreak{}Z.\allowbreak{}Kato.\allowbreak{}Nb.\allowbreak{}katoZetaElement\_\allowbreak{}normCompatible\_\allowbreak{}in\_\allowbreak{}lTower}{Development.\allowbreak{}AutomorphicArithmetic.\allowbreak{}AutomorphicVarieties.\allowbreak{}ModularityArgument.\allowbreak{}KatoQuadraticCharacters.\allowbreak{}EulerFactorBounds.\allowbreak{}ZetaClassCohomologyCompat}{210--315}
\leanref{UpperTheory.\allowbreak{}Kd.\allowbreak{}K4.\allowbreak{}mlOf2}{Development.\allowbreak{}AutomorphicArithmetic.\allowbreak{}Foundations.\allowbreak{}IwahoriReduction.\allowbreak{}NormProducts.\allowbreak{}IsCompactIntegralCoefficients}{86--212}
\leanref{UpperTheory.\allowbreak{}Kd.\allowbreak{}K4.\allowbreak{}g3}{Development.\allowbreak{}AutomorphicArithmetic.\allowbreak{}Foundations.\allowbreak{}IwahoriReduction.\allowbreak{}NormProducts.\allowbreak{}AdmissibleWeightedDataMonoCoef}{15--37}

\begin{lemma}[Compatible integral representatives]\label{B6.45a}
Let \(\Gamma\) act linearly on a separated complete \(\ell\)-adic \(\Z_\ell\)-module \(B\) with locally constant orbit maps modulo every \(\ell^nB\). Let \(A_r\) be finite groups with surjections \(\pi_r:A_{r+1}\to A_r\) and compatible locally constant homomorphisms \(\kappa_r:\Gamma\to A_r\). Write \(P_r\) for fibre summation on zero- or one-cochains and set \((\partial b)(s,u)=s b(u\kappa_r(s))-b(u)\).

For \(r,n\ge0\) let \(F_{r,n}:\Gamma\times A_r\to B\) be locally constant in \(s\) modulo \(\ell^nB\) and satisfy the twisted cocycle identity there. If \(F_{r,n+1}-F_{r,n}\) and \(P_rF_{r+1,n}-F_{r,n}\) are integral coboundaries modulo \(\ell^nB\), there are continuous twisted cocycles \(Y_r=P_rY_{r+1}\) with \(Y_r\equiv F_{r,n}+\partial b_{r,n}\pmod{\ell^nB}\) for integral zero-cochains \(b_{r,n}\).
\end{lemma}
\begin{proof}
A section of \(\pi_r\) lifts a zero-cochain by supporting its value at the selected point of each fibre, so \(P_r\) is onto in degree zero. Multiplication by \(\kappa_{r+1}(s)\) identifies the fibres over \(u\) and \(u\kappa_r(s)\). Hence \(P_r\partial=\partial P_r\).

Construct \(D_j=F_{j,j}+\partial b_j\) with \(b_0=0\). Modulo \(\ell^jB\) write \(F_{j+1,j+1}-F_{j+1,j}=\partial a\) and \(P_jF_{j+1,j}-F_{j,j}=\partial c\). Lift \(b_j-c-P_ja\) along \(P_j\) to \(b_{j+1}\). Then \(P_jD_{j+1}\equiv D_j\bmod\ell^jB\). For fixed \(r\), the pushforwards to \(A_r\) for \(j\ge r\) stabilise successively modulo \(\ell^jB\) and converge to \(Y_r\). Finite summation commutes with the limit. At any \(j\ge\max(r,n)\), lower the precision to \(n\) and push down to \(r\). The two coboundary hypotheses give \(F_{r,n}+\partial b_{r,n}\) at that precision. These residues are locally constant and satisfy the twisted cocycle identity, so the limit is a continuous cocycle.
\end{proof}
\leanref{UpperTheory.\allowbreak{}Kd.\allowbreak{}K4.\allowbreak{}mittagLeffler\_\allowbreak{}algebraic}{Development.\allowbreak{}AutomorphicArithmetic.\allowbreak{}Foundations.\allowbreak{}IwahoriReduction.\allowbreak{}NormProducts.\allowbreak{}AdicEqZeroOfMemAll}{6--78}
\leanref{UpperTheory.\allowbreak{}Kd.\allowbreak{}K4.\allowbreak{}ml3}{Development.\allowbreak{}AutomorphicArithmetic.\allowbreak{}Foundations.\allowbreak{}IwahoriReduction.\allowbreak{}NormProducts.\allowbreak{}AdicEqZeroOfMemAll}{118--223}

For an integral series \(F(T)=\sum_nF_nT^n\) and a primitive \(\ell^r\)-th root \(\xi\), define
\[
 \operatorname{mass}(F,r,\xi,x)=
 \sum_{t=0}^{\ell^r-1}\xi^{-tx}\sum_{n=0}^{\ell^{3r}-1}F_n(\xi^t-1)^n,
 \qquad
 M_r(F,\xi)=\sum_{\substack{0\le x<\ell^r\\\ell\nmid x}}
 \phi(x\bmod\ell)x^k\operatorname{mass}(F,r,\xi,x).
\]
Take \(\phi\) to be the inverse of \(\psi_0\) on units.

\begin{theorem}[Local nonvanishing of the Euler class]\label{B6.46}
For each normal form \(W\) with source \(\operatorname{SrcUr}(\iota,k,A,W,g)\), there are a coefficient ring \(O\) containing the entries of \(W\), an integral cocycle \(y:G_\Q\to W\), and a finite set \(S'\) such that
\[
 \operatorname{Unr}(y),\qquad \neg\operatorname{Gr}_A(y).
 \tag{6.46.1}
\]
Let \(n\ge1\), \(w\notin S'\) with \(\ell^n\mid\mathrm Nw-1\), and let \(\chi:G_\Q\to\Z/\ell^n\) be surjective and factor through \(\mu_{\mathrm Nw}\). There are a continuous \(O\)-valued twisted cocycle \(Y:G_\Q\times\Z/\ell^n\to W\), scalars \(P\in O\), \(u\) with \(|u|_\ell=1\), and \(b\in O^2\) such that
\[
 Y(st,x)=Y(s,x)+W(s)Y(t,x-\chi(s)),
 \tag{6.46.2}
\]
\[
 \sum_xY(s,x)=Py(s)+(W(s)-1)b,
 \tag{6.46.3}
\]
\[
 P\equiv u\det(W(\sigma_w)-1)\pmod{\ell^n}
 \tag{6.46.4}
\]
for each Frobenius \(\sigma_w\) at \(w\). For each prime different from \(w,\ell\) there are a place \(B\) above it and \(t:\Z/\ell^n\to O^2\) with \(Y(s,x)=(W(s)-1)t(x)\) for \(s\in I_B\).
\end{theorem}

\leanref{UpperTheory.\allowbreak{}V.\allowbreak{}Z.\allowbreak{}Kato.\allowbreak{}Final.\allowbreak{}zetaUr}{Development.\allowbreak{}AutomorphicArithmetic.\allowbreak{}AutomorphicVarieties.\allowbreak{}ModularityArgument.\allowbreak{}KatoQuadraticCharacters.\allowbreak{}UnramifiedZeta.\allowbreak{}ZetaUr}{9--13}
\leanref{GlobalTheory.\allowbreak{}Ord2.\allowbreak{}ES}{Development.\allowbreak{}AbelianVarietyTheory.\allowbreak{}Definitions.\allowbreak{}Foundations.\allowbreak{}FinalTheorems}{98--103}
\leanref{GlobalTheory.\allowbreak{}Ord2.\allowbreak{}Tame}{Development.\allowbreak{}AbelianVarietyTheory.\allowbreak{}Definitions.\allowbreak{}Foundations.\allowbreak{}FinalTheorems}{88--96}

\begin{proof}
Use \ref{B6.45} with \(N\ge5\), \((g_{r_0})_{10}\equiv-r_0(g_{r_0})_{00}\bmod\ell\), and unit \(U_\ell\)-eigenvalue for \(\operatorname{pr}\) on the weight-\(2k-1\) module. Fix \(j:\overline\Q\to C_\ell\) inducing \(A\). Choose a diamond character \(\chi_d\) with \(\operatorname{pr}(\langle d\rangle b)=\chi_d(d)\operatorname{pr}(b)\) for \((d,N)=1\). At depth \(n\ge c_L\) a finitely supported linear functional \(\lambda\) of the embedded local expansion recovers the second projection coordinate modulo \(\ell^{n-c_L}\) on this eigenspace, up to a coherent unit.

Apply \ref{B6.46a} with \(\psi=\iota\circ\psi_0\) on units. Take \(c_w(x)=\iota^{-1}(\widehat\theta(x)/2)\), where \(\widehat\theta(x)=\sum_{t\bmod N}\theta(t)e^{2\pi itx/N}\). Choose \(\lambda\) jointly with the admissible projection. On the coefficient sequences of forms of weight \(2k+1\), level \(N\ell\) and character \(\omega\), \(\lambda\) is a fixed nonzero scalar times the ordinary functional of \ref{B6.46a} transported by \(\iota^{-1}\).

\leanref{UpperTheory.\allowbreak{}V.\allowbreak{}Z.\allowbreak{}Kato.\allowbreak{}Joint.\allowbreak{}kato\_\allowbreak{}global\_\allowbreak{}admissibility\_\allowbreak{}of\_\allowbreak{}admissible}{Development.\allowbreak{}AutomorphicArithmetic.\allowbreak{}AutomorphicVarieties.\allowbreak{}ModularityArgument.\allowbreak{}KatoQuadraticCharacters.\allowbreak{}FrobeniusHodgeTateLaws.\allowbreak{}KatoGlobalAdmissibilityOfAdmissible}{35--39}

Put \(d_0=D=\ell+1\). For \(Y_r(s,u)_1-d_0Y_r(s,ud_0)_1\), the Hodge--Tate law supplies compatible primitive roots \(\zeta_m\in\overline\Q\) of order \(\ell^m\) and integers \(m_0,c_1\). For each \(m\ge m_0\) take its polynomial presentation \((P_i,o_i)\) with values
\[
 \operatorname{htv}_m(b)=\sum_i o_i\,j\!\left(
 \frac{\zeta_m^bP_i'(\zeta_m^b)}{P_i(\zeta_m^b)}\right).
\]
For every primitive \(\chi\) of conductor \(\ell^m\), the same presentation gives
\[
 \begin{split}
 \sum_{b\in(\Z/\ell^m)^\times}\chi(b)\operatorname{htv}_m(b)
 &\equiv(1-d_0\chi^{-1}(d_0))v_m^{(0)}\kappa_0
 \left(\sum_{a\bmod\ell^m}\chi(a)j(\zeta_m)^a\right)
 \chi^{-1}(N)\lambda(\mathcal E_{m,\chi})
 \pmod{\ell^{\max(m-c_1,0)}} .
 \end{split}
 \tag{6.46.5}
\]
Here \(\kappa_0\ne0\), the \(v_m^{(0)}\) are units with \(v_{m+1}^{(0)}\equiv v_m^{(0)}\bmod\ell^m\), and \(\mathcal E_{m,\chi}\) is the weighted Eisenstein coefficient vector.

The eight translates \((b_i,t_i)\) of \ref{B6.46a} give the multiplier \(\sum_i\iota^{-1}(b_i)\chi^{-1}(t_i)\). In (6.46.5) the weight function \(c_w\) yields this multiplier times the Gauss sum and the character value of the Eisenstein distribution, up to a common nonzero scalar and a power of the ordinary eigenvalue. The parity projector is \((1+\iota^{-1}(s)\chi(-1))/2\). Thus the Eisenstein side uses the series \(L\) of \ref{B6.46a}. If \(h\) is the scaling exponent of \ref{B6.45}, the local system and corresponding global class are
\[
 Z_r(\sigma,u)=\ell^h\bigl(Y_r(\sigma,u)_1-DY_r(\sigma,Du)_1\bigr),
 \qquad y_*=(1-D^{k+1})\ell^hy.
\]
In the inverse-weight specialisation the change \(u\mapsto Du\) contributes \(D^k\bmod\ell^r\) and preserves \(\chi_r\) since \(D\equiv1\bmod\ell\). Hence \(Z\) specialises to \((y_*)_1\) on \(I_A\cap\ker\varepsilon_\ell\). Scaling the tame data with the global class gives (6.46.2)--(6.46.4) and preserves the local trivialisations away from \(\ell,w\) and \(\operatorname{Unr}(y_*)\).

\leanref{UpperTheory.\allowbreak{}V.\allowbreak{}Z.\allowbreak{}Kato.\allowbreak{}Joint.\allowbreak{}glue2}{Development.\allowbreak{}AutomorphicArithmetic.\allowbreak{}AutomorphicVarieties.\allowbreak{}ModularityArgument.\allowbreak{}KatoQuadraticCharacters.\allowbreak{}FrobeniusHodgeTateLaws.\allowbreak{}AdmissibleHeckeAlgebraOfGenerated}{17--231}

Let \(G\) be the series for the local Iwasawa class and put
\[
 S_m(F,\chi,z)=\sum_{b\in(\Z/\ell^m)^\times}\chi(b)
 \sum_{i<\ell^{3m}}F_i(z^b-1)^i.
\]
Apply the local Kummer calculation to \(Z\) with the additional smoothing integer \(2\). The \(P_i\) in (6.46.5) have integral inertia-fixed coefficients and unit values at \(\zeta_m\). The weighted Kummer characters of these values represent the smoothed coordinate at \(u=1\) modulo \(\ell^{\max(m-c_1,0)}\). Translation of the logarithmic derivatives supplies the multiplier \(1-2\chi^{-1}(2)\). For all large \(m\) there is a primitive \(\ell^m\)-th root \(z\) with
\[
 (1-2\chi^{-1}(2))v'_m\kappa'(a')^m S_m(L,\chi,z)
 \equiv v_m a^m S_m(G,\chi,z)\pmod{\ell^{m-C}}
\]
for every primitive \(\chi\). Here \(a,a'\) are units, \(\kappa'\ne0\), \(C\) is fixed and each of the unit sequences \(v_m,v'_m\) agrees with its successor modulo \(\ell^m\).

Write \(A_{F,d}\) for the Mellin series on the \(\phi\)-component with argument multiplied by \(d\). With \(T_*=(1+\ell)^k-1\), Fourier inversion gives
\[
 d^k\phi(d)\ell^{-r}M_r(F,\xi)\equiv A_{F,d}(T_*)\pmod{\ell^r}.
\]
By \ref{B6.46a}, \(A_{L,1}(T_*)\ne0\). The series \(B=A_{L,1}-2A_{L,2}\) has nonzero value at \(T_*\), since its multiplier \(1-2^{k+1}\phi(2)\) is nonzero by raising a putative equality to the power \(\ell-1\).

For \(j\ge1\) and a primitive \(\ell^j\)-th root \(\eta\), choose \(\chi\) of conductor \(\ell^{j+1}\) with \(\chi^{-1}(1+\ell)=\eta\) and \(\chi^{-1}(b)=b^k\phi(b\bmod\ell)\) on the Teichmüller units \(b\). These prescriptions determine the character on the tame and wild factors, and the order of \(\eta\) gives the conductor. For \(m=j+1\), the direct and inverse Gauss sums satisfy \(\tau\tau'=\ell^m\) and \(v_\ell(\tau')\ge\lfloor m/2\rfloor\). Hence \(v_\ell(\tau)\le\lceil m/2\rceil\). After division by \(\tau\), the comparison at \(T=\eta-1\) is between \(A_{G,1}\) and \(w_j(a'/a)^j\kappa'B\), with coherent units \(w_j\) and precision at least \(\lfloor j/2\rfloor-C'\) after clearing a fixed denominator.

Put \(A_1=A_{G,1}\), \(u=a'/a\) and \(H(T)=(1+T)^\ell-1\). Successive comparisons show that \(A_1(B\circ H)-u(A_1\circ H)B\) tends to zero at these primitive roots. For a nonzero bounded series over a finite extension of \(\Q_\ell\), the first coefficient of maximal norm gives the unique dominant term as \(|\eta-1|_\ell\) tends to one. The difference therefore vanishes identically. The same dominant-term test with \(B\ne0\) gives \(A_1\ne0\). A leading Laurent power \(T^j\) of \(A_1/B\) gives \(1=u\ell^j\), so \(j=0\) and \(u=1\). Subtract the constant term and apply this identity to a possible next power to conclude that \(A_1/B\) is constant. Evaluation at \(T_*\) gives
\[
 M_r(G,\xi)\not\equiv0\pmod{\ell^{r+n_1}}
\]
for some fixed \(n_1\), every large \(r\) and every primitive \(\xi\).

\leanref{GlobalTheory.\allowbreak{}KZ.\allowbreak{}Cmp.\allowbreak{}nonvanishing\_\allowbreak{}transfer}{Development.\allowbreak{}AbelianVarietyTheory.\allowbreak{}Proofs.\allowbreak{}FoundationProofs.\allowbreak{}NonvanishingTransferCharacterSumsPowerSeries}{281--360}

If \((y_*)_1\) vanished on \(I_A\cap\ker\varepsilon_\ell\), the inverse-weight specialisations of \(Z\) would tend uniformly to zero there. The local division formula then gives \(g\) prime to \(\ell\), a fixed \(b\ge0\) and an integral system \(Z'\) with
\[
 \ell^bZ_r(\sigma,x)=Z'_r(\sigma,gx)-\phi(g)g^k Z'_r(\sigma,x)
 \qquad(\sigma\in I_A,\ \sigma|_{\mu_{\ell^r}}=1).
\]
The Kummer series respects translation, whose weighted mass is \(\phi(g)g^k\) times the original modulo \(\ell^{2r}\). Hence \(\ell^bM_r(G,\xi)\equiv0\bmod\ell^{2r}\), contrary to the bound for large \(r\ge n_1+b+1\).

Take \(\sigma\in I_A\cap\ker\varepsilon_\ell\) with \(y_*(\sigma)_1\ne0\). It fixes \(\mu_\ell\) and every quadratic field. Inertia fixes square roots of \(\ell\)-adic units, and the quadratic Gauss sum in \(\Q(\mu_\ell)\) supplies a square root of \((-1)^{(\ell-1)/2}\ell\). Thus \(\det W(\sigma)=1\) and the inertia identity of \(\operatorname{SrcUr}\) gives \(\beta(\sigma)=1\). By (6.44.2), the same class \(y_*\) that satisfies the tame relations fulfils \(\neg\operatorname{Gr}_A(y_*)\).

\leanref{UpperTheory.\allowbreak{}Kd.\allowbreak{}C7.\allowbreak{}pCore}{Development.\allowbreak{}AutomorphicArithmetic.\allowbreak{}Foundations.\allowbreak{}IwahoriReduction.\allowbreak{}LangPolynomials.\allowbreak{}MemOfMulMem}{253--353}
\leanref{GlobalTheory.\allowbreak{}quot\_\allowbreak{}eq\_\allowbreak{}one\_\allowbreak{}of\_\allowbreak{}cycBar\_\allowbreak{}eq\_\allowbreak{}one}{Development.\allowbreak{}AbelianVarietyTheory.\allowbreak{}Proofs.\allowbreak{}FoundationProofs.\allowbreak{}QuotientCharTrivialInertiaCyclotomicTower}{231--242}
\leanref{UpperTheory.\allowbreak{}V.\allowbreak{}Z.\allowbreak{}Kato.\allowbreak{}K5.\allowbreak{}smoothedZetaClass\_\allowbreak{}hodgeTateLaw}{Development.\allowbreak{}AutomorphicArithmetic.\allowbreak{}AutomorphicVarieties.\allowbreak{}ModularityArgument.\allowbreak{}KatoQuadraticCharacters.\allowbreak{}SmoothedZetaClasses.\allowbreak{}SmoothedZetaClassHodgeTateLaw}{4--5}
\leanref{UpperTheory.\allowbreak{}Kd.\allowbreak{}C7.\allowbreak{}pLocu}{Development.\allowbreak{}AutomorphicArithmetic.\allowbreak{}Foundations.\allowbreak{}IwahoriReduction.\allowbreak{}LangPolynomials.\allowbreak{}NormFactorLeOne}{162--165}
\end{proof}

\begin{lemma}[The Eisenstein moment at the prescribed weight]\label{B6.46a}
Let \(g\) be a normalised cuspidal newform of weight \(2k+1\) and level \(M_g\), where \(k\ge1\), \(\ell\) is odd, \(\ell\nmid M_g\) and \(|\iota^{-1}(a_\ell(g))|_\ell=1\). Fix a character \(\psi\bmod\ell\) and put \(\phi=\iota^{-1}\circ\psi^{-1}\). One can choose \(N\ge5\) divisible by \(M_g\) and prime to \(\ell\), a character \(\theta\bmod N\) of sign \(s=(-1)^k\psi(-1)\) and an ordinary Eisenstein functional at this level. With \(D=\ell+1\), the resulting smoothed distribution has a series \(L\in\mathcal O_E[[T]]\) for a finite extension \(E/\Q_\ell\), and there are \(n_0,R\) with \(M_r(L,\xi)\not\equiv0\pmod{\ell^{r+n_0}}\) for every \(r\ge\max(1,R)\) and every primitive \(\ell^r\)-th root \(\xi\).
\end{lemma}
\begin{proof}
Choose a prime \(p\ge\ell+7\), take \(N=M_g(M_gp)^2\), and use the depleted form \(g'\) with coefficients \(a_n(g)\) for \((n,M_gp)=1\) and zero otherwise. Let \(\omega\) be its character at level \(N\). There are at least three characters modulo \(p\) of sign \(s\). Inflate them to \(N\) and choose \(\theta\ne\omega^{\pm1}\), so \((\theta\omega)\psi\) is nontrivial at level \(N\ell\).

Unfold the pairing of the cusp form with coefficients \(\overline{g'(n)}\) against the product of the weight-\(k\) Eisenstein series with divisor weights \(\psi,\theta\) and the weight-\(k+1\) lattice Eisenstein series of character \((\theta\omega)\psi\). It gives
\[
 P=\frac{\Gamma(2k)}{(4\pi)^{2k}}L(g'\otimes\psi,k+1)L(g'\otimes\theta,2k).
\]
The twisted values are right limits. Both arguments lie in the nonvanishing range \(w\ge(K+1)/2\) for a depleted eigenform of odd weight \(K=2k+1\), including \(k=1\). Dilation of the cusp form by \(\ell\) gives \(\theta(\ell)P/\ell^{2k}\).

Let \(a,b\) have sum \(\omega^{-1}(\ell)a_\ell(g)\) and product \(\omega^{-1}(\ell)\ell^{2k}\), with \(a\) the unit root. The stabilised functional is a nonzero scalar times \(P(\ell\theta(\ell)-b)/\ell^{2k+1}\), which is nonzero since \(v_\ell(b)=2k\). Depletion of the Eisenstein factors at \(\ell\) also requires \(a\ne\ell^k\omega^{-1}(\ell)\) and \(a\ne\theta(\ell)\). The first follows from ordinarity. Equality in the second would give complex absolute values \(|a|=1\), \(|b|=\ell^{2k}\). The recurrence
\[
 \omega(\ell)^{-j}g'(\ell^j)=\frac{b^{j+1}-a^{j+1}}{b-a}
\]
would contradict the Rankin bound \(O(\ell^{j(k+1/2)})\). Thus the depleted product has nonzero ordinary value. The Atkin--Lehner transformation at \(N\) interchanges the two presentations of the Eisenstein product up to nonzero scalars. The transformed functional has eigenvalue \(\omega(\ell)a\), and its ordinary value times \(2/\varphi(N)\) gives the interpolation value \(v\ne0\).

Write \(\nu_r^0\) for this Eisenstein distribution after inversion on units and projection to sign \(s\). For \(c,d>1\) prime to \(6N\ell\), put \(A_c=\theta(c)\), \(B_d=(\theta\omega)^{-1}(d)\) and \([t]h(y)=h(ty)\). Expand
\[
 [N](c^2-c^{2-2k}A_c[c])(d^2-dB_d[d]).
\]
The four translates act on \(t\mapsto\psi^{-1}(t)t^k\) by the scalar
\[
 Q_4=N^k\psi^{-1}(N)(c^2-c^{2-k}A_c\psi^{-1}(c))
 (d^2-d^{k+1}B_d\psi^{-1}(d)).
\]
The first difference is nonzero since \(c^k>1\). For \(k>1\) the terms of the second have different absolute values. For \(k=1\) use the Chinese remainder theorem to choose \(d\equiv1\bmod\ell\) prime to \(6N\ell\) with \((\theta\omega)(d)\ne1\). Append to each pair \((b_i,t_i)\) of coefficient and translate the pair \((-Db_i,Dt_i)\). The eight terms act by \(Q_8=(1-D^{k+1})Q_4\ne0\).

The unscaled point masses are bounded over a finite extension. A fixed power \(\kappa\) of \(\ell\) makes \(\nu_r(h)=\kappa\sum_i\iota^{-1}(b_i)\nu_r^0(h(t_i\,\cdot))\) integral. Write \(h_r(y)=\phi(\widetilde y)\widetilde y^k\) on units and zero elsewhere. Interpolation gives
\[
 |\nu_r(h_r)-A_*|_\ell\le C|\ell|_\ell^r,\qquad A_*=\kappa\iota^{-1}(vQ_8)\ne0.
\]
For units \(a',d'\) with \(a'y=d'\bmod\ell^r\), the residue of \(t_i y\) and \(t_i d'/a'\) have congruent \(k\)-th powers modulo \(\ell^r\) and the same character value. The contribution at \(-y\) has sign \(s\).

Take \(L(T)=\int(1+T)^y\,d\nu(y)\). At \(T=\xi^t-1\), the truncation in \(\operatorname{mass}\) gives \(\nu_r(y\mapsto\xi^{ty})\bmod\ell^{2r}\). For a root of order \(\ell^j>1\) the tail has valuation at least \(\ell^{3r}/\varphi(\ell^j)\ge2r\), and at \(T=0\) it vanishes. The sum of \(\xi^{t(y-x)}\) over \(t\) is \(\ell^r\) for \(y=x\) and zero elsewhere. Thus \(\ell^{-r}M_r(L,\xi)\equiv\nu_r(h_r)\bmod\ell^r\) uniformly in \(\xi\). For large \(r\) this normalised sum has the absolute value of \(A_*\). Take \(n_0>v_\ell(A_*)\).
\end{proof}
\leanref{GlobalTheory.\allowbreak{}EisM.\allowbreak{}twisted\_\allowbreak{}lvalue\_\allowbreak{}ne\_\allowbreak{}zero\_\allowbreak{}edge}{Development.\allowbreak{}AbelianVarietyTheory.\allowbreak{}Proofs.\allowbreak{}FoundationProofs.\allowbreak{}TwistedLvalueNeZeroEdge}{167--220}
\leanref{GlobalTheory.\allowbreak{}EisM.\allowbreak{}rankin\_\allowbreak{}values\_\allowbreak{}at\_\allowbreak{}zero}{Development.\allowbreak{}AbelianVarietyTheory.\allowbreak{}Proofs.\allowbreak{}FoundationProofs.\allowbreak{}RankinValuesAtZero}{58--211}
\leanref{GlobalTheory.\allowbreak{}EisM.\allowbreak{}eigF\_\allowbreak{}value\_\allowbreak{}ne\_\allowbreak{}zero}{Development.\allowbreak{}AbelianVarietyTheory.\allowbreak{}Proofs.\allowbreak{}FoundationProofs.\allowbreak{}ExistsHeckeEigenfunctionalDepletedEisensteinNeZe}{72--204}
\leanref{GlobalTheory.\allowbreak{}EisM.\allowbreak{}chapterM\_\allowbreak{}of\_\allowbreak{}level}{Development.\allowbreak{}AbelianVarietyTheory.\allowbreak{}Proofs.\allowbreak{}FoundationProofs.\allowbreak{}OrdinaryEisensteinEigenfunctionalExists}{39--108}
\leanref{UpperTheory.\allowbreak{}V.\allowbreak{}Z.\allowbreak{}Kato.\allowbreak{}Joint.\allowbreak{}smNz}{Development.\allowbreak{}AutomorphicArithmetic.\allowbreak{}AutomorphicVarieties.\allowbreak{}ModularityArgument.\allowbreak{}KatoQuadraticCharacters.\allowbreak{}FrobeniusHodgeTateLaws.\allowbreak{}SmNzNorm}{15--71}
\leanref{UpperTheory.\allowbreak{}V.\allowbreak{}Z.\allowbreak{}Kato.\allowbreak{}Joint.\allowbreak{}meas3}{Development.\allowbreak{}AutomorphicArithmetic.\allowbreak{}AutomorphicVarieties.\allowbreak{}ModularityArgument.\allowbreak{}KatoQuadraticCharacters.\allowbreak{}FrobeniusHodgeTateLaws.\allowbreak{}DistributionApply}{153--211}

\begin{theorem}[Euler data imply Bloch--Kato vanishing]\label{B6.47}
Assume the cohomological inputs of \(\mathscr H\). For every \(\operatorname{Nf}(k,S,A,W)\), Euler data as in \ref{B6.46} imply \(H_f^1(\Q,W)=0\).
\end{theorem}

\leanref{GlobalTheory.\allowbreak{}Ord2.\allowbreak{}es\_\allowbreak{}vanish}{Development.\allowbreak{}AbelianVarietyTheory.\allowbreak{}Proofs.\allowbreak{}FoundationProofs.\allowbreak{}LocPublishedOfLitCoh}{16--24}

\begin{proof}
Put \(W^D=W^\vee(1)\simeq W\otimes\varepsilon_\ell(\det W)^{-1}\). Suppose \(z\) represents a nonzero strict class for \(W^D\), unramified away from \(\ell\) and trivial on \(D_A\). Scale \(y,z\) and enlarge coefficients to make their coordinates and local trivialisations integral. Form \(y\cup z\) using the alternating pairing \(\langle\, ,\,\rangle:W\times W^D\to C_\ell(1)\). By strictness at \(\ell\) and unramified orthogonality elsewhere, it is locally a coboundary modulo every \(\ell^n\). Poitou--Tate gives an integral trivialising cochain \(u_n\) factoring modulo \(\ell^n\) through \(G_{\Q,S}\), the group unramified outside \(S\).

Put \(N_0=\ker\varepsilon_\ell\cap\ker W\) and \(H_0=\ker\varepsilon_\ell\cap\ker\det W\). If \(y\) vanished on \(N_0\) it would be Greenberg-ordinary by the restriction criterion. The dual criterion gives a coboundary for \(z\) if it vanishes there. Thus both restrictions are nonzero. Irreducibility and the pairing give \(g,h\in N_0\) with
\[
 b=\langle y(h),z(g)\rangle-\langle y(g),z(h)\rangle\ne0 .
 \tag{6.47.1}
\]
Irreducibility also gives \(\phi\in[H_0,H_0]\) with \(d=\det(W(\phi)-1)\ne0\), of norm at most one by integrality. Define
\[
 h_n(s)=\det(W(s)-1)u_n(s)-
 \left\langle\operatorname{adj}(W(s)-1)y(s),z(s)\right\rangle .
\]
For \(\sigma\in[H_0,H_0]\), choose by Chebotarev \(w\notin S\cup S'\) with \(\ell^n\mid\mathrm Nw-1\) and every \(\mathrm Nv\bmod\mathrm Nw\) for \(v\in S\) an \(\ell^n\)-th power. Take a Frobenius \(\tau\) of determinant one with
\[
 \det(W(\tau)-1)h_n(\tau)\equiv\det(W(\sigma)-1)h_n(\sigma)\pmod{\ell^n}.
\]
A surjective \(\chi\) of conductor \(\mathrm Nw\) and order \(\ell^n\) is trivial on each old \(D_v\). The tame relation gives an auxiliary cochain \(d_w\). On \(G_{\Q,S\cup\{w\}}\) use the two-cocycle modulo \(\ell^n\)
\[
 \Xi(s,t)=(d_w\cup z)(s,t)-\widetilde{\chi(s)}P\varepsilon_\ell(s)u_n(t),
\]
with the least nonnegative representative \(\widetilde{\chi(s)}\). At the old places the second term vanishes and the first is a coboundary by strictness or unramified orthogonality. Reciprocity gives a coboundary at \(w\). The local height identity yields \(Ph_n(\tau)\equiv0\bmod\ell^n\), so (6.46.4) gives \(\det(W(\sigma)-1)h_n(\sigma)\equiv0\bmod\ell^n\).

Apply this at \(\phi,\phi[g,h]\), which both have determinant factor \(d\) since \(W(g)=W(h)=1\). Put \(t_0=h_n(\phi)\), \(t_1=h_n(\phi[g,h])\). Then \(|dt_i|_\ell\le\ell^{-n}\) and the height commutator identity gives \(|t_1-t_0-db|_\ell\le\ell^{-n}\). From \(d^2b=d(db-(t_1-t_0))+dt_1-dt_0\) we get
\[
 |d^2b|_\ell\le\ell^{-n}\qquad(n\ge1).
 \tag{6.47.2}
\]
This contradicts \(db\ne0\), so the strict dual group vanishes.

For an ordinary dual class, take \(z\) with \(z_1=0\) on \(D_A\) and scale the cocycles and local corrections to be integral. The cup product is zero modulo every \(\ell^n\) away from \(\ell\), hence at \(\ell\) by reciprocity. There it pairs \(y_1\) with \(z_0\) up to sign. Their nontrivial characters \(\beta,\varepsilon_\ell\beta^{-1}\) give one-dimensional cohomology spaces in perfect Tate duality. By (6.46.1) the class of \(y_1\) is nonzero, so \(z_0\) is a coboundary. Subtract it in the first line to make \(z\) strict.

By Greenberg--Wiles duality, a nonzero ordinary class for \(W\) would give dual finite-level classes of unbounded order, contrary to the dual vanishing. The local Bloch--Kato comparison places a Bloch--Kato class in this ordinary Selmer group.

\leanref{GlobalTheory.\allowbreak{}Ord2.\allowbreak{}StrictKill.\allowbreak{}main}{Development.\allowbreak{}AbelianVarietyTheory.\allowbreak{}Proofs.\allowbreak{}FoundationProofs.\allowbreak{}StrictVanishingDualOfLiterature}{162--282}
\leanref{GlobalTheory.\allowbreak{}Ord2.\allowbreak{}Upgrade.\allowbreak{}upgrade}{Development.\allowbreak{}AbelianVarietyTheory.\allowbreak{}Proofs.\allowbreak{}FoundationProofs.\allowbreak{}OrdinarySelmerDualVanishesOfStrict}{161--277}
\leanref{GlobalTheory.\allowbreak{}Ord2.\allowbreak{}grZero\_\allowbreak{}of\_\allowbreak{}grZero\_\allowbreak{}dual1}{Development.\allowbreak{}AbelianVarietyTheory.\allowbreak{}Proofs.\allowbreak{}FoundationProofs.\allowbreak{}OrdinarySelmerVanishesOfDualVanishes}{10--17}
\end{proof}

\section{Four-dimensional symplectic points}

For a datum \(D\) of case II or III of \ref{B6.31}, an allowed lift is a specialisation of a lift over a complete Noetherian local coefficient algebra satisfying the fixed similitude, normalisation and local conditions of \(D\). Write \Lean{D.\allowbreak{}PtIn} for this condition and \Lean{D.\allowbreak{}EPt} for its characteristic-zero instance. The zero-weight algebra is generated by the matrix-coefficient monomials of total weight zero for the diagonal action. A block point preserves the coordinate subspaces \(\{0,2\}\) and \(\{1,3\}\).

\begin{theorem}[Infinitesimal rigidity of an automorphic symplectic point]\label{B6.48}
Assume the unramified inputs derived from \(\mathscr H\) and \ref{B6.53}. Let \(D\) be of case II or III with finite characters \(\vartheta,\xi,\omega\) factoring through \(\mu_N\), \(N\ge1\). Fix \(\iota:C_\ell\simeq\C\) and an irreducible characteristic-zero point \(r\) of \(D\). On every quadratic CM field \(K\) where \(r\) remains irreducible, assume \(r|_{G_K}\) has a regular algebraic cuspidal attachment with witness. Then every allowed lift \(r':G_\Q\to M_4(C_\ell[\epsilon]/\epsilon^2)\) reducing to \(r\) has the form
\[
 r'(g)=r(g)+\epsilon(Xr(g)-r(g)X)\qquad(g\in G_\Q)
 \tag{6.48.1}
\]
for \(X\in\mathfrak{sp}_4(C_\ell)\).
\end{theorem}

\leanref{UpperTheory.\allowbreak{}V.\allowbreak{}Th.\allowbreak{}rigSp4\_\allowbreak{}alt}{Development.\allowbreak{}AutomorphicArithmetic.\allowbreak{}AutomorphicVarieties.\allowbreak{}ModularityArgument.\allowbreak{}ThresholdArguments.\allowbreak{}DenseThresholdCases.\allowbreak{}RigSp4}{181--194}
\leanref{UpperTheory.\allowbreak{}V.\allowbreak{}GradeZeroInterpolatingElement.\allowbreak{}Raw.\allowbreak{}RigidAt}{Development.\allowbreak{}AutomorphicArithmetic.\allowbreak{}AutomorphicVarieties.\allowbreak{}SymplecticGSp4.\allowbreak{}DecompSigmaOrd}{182--184}

\begin{proof}
Write \(r'(g)=(1+\epsilon z(g))r(g)\), so
\[
 z(gh)=z(g)+r(g)z(h)r(g)^{-1}.
 \tag{6.48.2}
\]
Similitude puts \(z\) in \(\mathfrak{sp}_4\). It vanishes on a finite-index inertia subgroup away from \(\ell\), so averaging gives an inertia coboundary. At \(\ell\), \Lean{tb\_\allowbreak{}atL} gives the finite local condition. Choose an imaginary quadratic \(K\) where \(r\) remains irreducible and an ordinary model over a solvable CM extension \(L/K\). Solvable base change transfers the attachment and witness. Apply \ref{B6.53} over \(L\) and \Lean{adDown} to the extended adjoint over \(\Q\).

Set
\[
 J=\begin{pmatrix}0&0&0&1\\0&0&1&0\\0&-1&0&0\\-1&0&0&0\end{pmatrix},
 \qquad Y^\dagger=J^{-1}Y^tJ,\qquad P(Y)=\frac{Y-Y^\dagger}{2}.
\]
The image of \(P\) is \(\mathfrak{sp}_4\), and the similitude identity gives
\[
 P(r(g)Yr(g)^{-1})=r(g)P(Y)r(g)^{-1}.
 \tag{6.48.3}
\]
For the quadratic character \(\eta\) of \(K\), the extended action is
\[
 R(g)Y=\operatorname{Ad}(r(g))P(Y)+
 \eta(g)\operatorname{Ad}(r(g))(Y-P(Y)).
 \tag{6.48.4}
\]
It is conjugation on \(G_K\) and the prescribed transpose-dual action on the other coset. The inclusion of \(\mathfrak{sp}_4\) preserves local conditions, so adjoint vanishing gives \(z(g)=R(g)Y-Y\). Apply \(P\) and multiply by \(r(g)\) on the right to get (6.48.1) with \(X=-P(Y)\).

\leanref{UpperTheory.\allowbreak{}V.\allowbreak{}Z.\allowbreak{}Win.\allowbreak{}Dn.\allowbreak{}exists\_\allowbreak{}cm\_\allowbreak{}field\_\allowbreak{}twist\_\allowbreak{}h1}{Development.\allowbreak{}AutomorphicArithmetic.\allowbreak{}AutomorphicVarieties.\allowbreak{}ModularityArgument.\allowbreak{}ThresholdArguments.\allowbreak{}DenseThresholdCases.\allowbreak{}ExistsCmFieldTwistH1}{28--103}
\end{proof}

\begin{lemma}[Irreducible arithmetic points of the Hecke quotient]\label{B6.48a}
Assume \(\mathscr H\). In a definite unitary ordinary Hecke model of rank \(n\ge1\) over a CM field \(L\), take level with the prescribed ramification \(S'=S\). Let \(J\) be its Hecke ideal in the universal ordinary determinant ring \(B\). Fix \(\iota:C_\ell\simeq\C\). Let \(x:B\to C_\ell\) kill \(J\) and have locally algebraic ordinary characters with strictly increasing weights at each embedding. If the determinant law of an irreducible \(r\) on \(G_L\) equals the pullback of \(x.D\), then \(r\) has a regular algebraic conjugate self-dual cuspidal attachment. Suppose also that \(r\) is continuous. Then every unlinked place above \(\ell\) is potentially unlinked. For each regular algebraic cuspidal attachment, one exponent \(t\) gives the local support bound at every place \(v\) with \(v|_{L^+}\ne v_0\) and with either \(v\nmid\ell\) or potential unlinking at \(v\). Here \(v_0\) is the distinguished place of \(L^+\).
\end{lemma}
\begin{proof}
The existence of \(x\) makes \(J\) proper, hence the kernel of the surjection onto the ordinary Hecke algebra. At its algebraic weight the induced eigensystem specialises to classical automorphic data with semisimple Galois determinant \(x.D\). Equality of determinant laws identifies semisimple representations, so irreducibility gives the entire rank and the Frobenius polynomials of the cuspidal attachment. The ordinary comparison gives potential unlinking and the required support at the corresponding places above \(\ell\). Away from \(\ell,v_0\) the local comparison gives toral support. The bound for unlinked toral support gives one exponent for each attachment.
\end{proof}
\leanref{UpperTheory.\allowbreak{}Gal.\allowbreak{}projectedFormsSystemPart.\allowbreak{}irrArithPointsAutomorphicAway\_\allowbreak{}of\_\allowbreak{}isJ}{Development.\allowbreak{}AutomorphicArithmetic.\allowbreak{}AutomorphicVarieties.\allowbreak{}ModularityArgument.\allowbreak{}ThresholdArguments.\allowbreak{}DenseThresholdCases.\allowbreak{}EndoscopyRank2Ordinary}{48--72}

\begin{lemma}[Minimal components through a residual point]\label{B6.48b}
Assume the cohomological input for the four-dimensional frame. Let \(T\) be a fitted tower of case II or III with \(\vartheta|_{G_{T.F}}\) of order prime to \(\ell\), \(B\to A\) its ordinary coefficient and matrix rings over the extended DVR \(\Lambda\), and \(J\subset B\) a potential modularity ideal for \(T,\iota\). Put \(d=[T.F:\Q]\) and let \(s\) count the local flags. Suppose \(s\), the strict root cocycle dimensions and the global cocycle dimensions of weights \(0,\pm\psi\) are at most \(C\). For each odd root \(w\), suppose that the global root cocycle space and the product \(\prod_v Z_{\mathrm{loc},v}(w)\) of the local root cocycle spaces each have dimension at most \(d+C\). Assume \(d\ge30(C+1)^2\).

Let \(x:A\to k_\Lambda\) be a residual point whose induced representation and flags satisfy the frame conditions \Lean{IsRes} of pairing, diagonal reduction and continuity. Take coordinate weights \((1,1),(0,1),(1,0),(0,0)\) and entry weights their differences. Write \(E(w)\) if some representation entry of weight \(w\) is nonzero. Put \(\alpha=(1,0)\), \(\beta=(-1,1)\), \(\psi=(0,1)\), \(\theta=(1,1)\). Supports refer to off-diagonal entries. Write \(V\) for support in \(\{\beta,-\alpha,-\theta\}\) with local flag matrices \(n_v\) supported in \(\{-\alpha,-\psi,-\theta\}\) and \((n_v)_{20}\ne0\) for every \(v\). Write \(H\) for support in \(\{-\beta,-\psi,-\theta\}\) with every \((n_v)_{10}\ne0\). Assume one of
\[
 \begin{gathered}
 E(-\psi)\land(E(\beta)\lor E(-\alpha)),\qquad
 E(-\alpha)\land E(-\beta),\qquad V\land E(-\alpha),\\
 V\land E(\beta)\land E(-\theta)\quad
 \text{with }\bar\chi_0/\bar\chi_1\ne1\text{ on every }G_v,\\
 H\land\bigl(E(-\psi)\lor(E(-\beta)\land E(-\theta))\bigr),
 \end{gathered}
\]
where \(\bar\chi_i\) are the residual diagonal characters. If one minimal prime of \(A\) below \(\ker x\) contains \(JA\), every such prime does.
\end{lemma}
\begin{proof}
Lift the two primes to the flat completed local deformation ring \(R\) at \(x\). Put \(b=2d+1-2s\) and \(d_v=d/s\). The power-series presentation over \(\Lambda\) has at least \(b+1\) more variables than relations. Connectedness joins the minimal primes by a chain with \(\dim R/(P+Q)\ge b+1\) for adjacent primes.

Start an adjacent step with \(JR\subset P\). The reducible support loci for the first four alternatives have dimension at most \(d+20C+1\). The same bound holds for \(H\) except for support \(Y=\{\pm\beta,\pm\theta\}\). In the tangent estimate at the diagonal limit, one odd entry and two diagonal entries contribute at most \(d+3C\). In the kernel the ten symplectic coordinates contribute at most \(10C\) and the six lower flag coordinates at most \(6s\). Evaluation on the matrix and flag entries is injective on derivations. Passing to ring dimension adds one.

On the Hecke special fibre, coincidence of two characters on \(U_v=I_v\cap\ker\varepsilon_\ell\cap\ker\bar r\) leaves at most \(2d+1-d_v\) parameters in \(B\). Use \(d_v\) inertia generators at each place for one character, at the other places for a second, and one character difference at \(v\). Similitude determines their partners. Off the reducible loci the fibre in \(R\) adds at most \(2+4s\) dimensions. This bound also holds on \(Y\) away from the other reducible supports. The diagonal projectors at the marking isolate the entries, and traces of products recover opposite entries from one nonzero entry in each pair. On \(Y\), the flag with \((n_v)_{10}\ne0\) gives \((\chi_1-\chi_0)(\chi_1-\chi_3)=0\) on \(G_v\). Two proper subgroups cannot cover the group, so one pair coincides on \(G_v\).

Both \(d+20C+1\) and \(2d+1-d_v+2+4s\) are less than \(b\) by \(s\le C\) and \(d_v\ge30(C+1)\). Cutting \(V(P+Q)\) by the uniformiser leaves dimension at least \(b\). Prime avoidance and successive parameter cuts give a one-dimensional quotient outside the excluded loci at its generic point. Normalisation gives a curve \(y:A\to k'[[t]]\) with \(k'\) finite, centred at \(x\), so the constant term of \(y|_B\) kills \(\mathfrak m_B\). The generic representation is irreducible and its characters are separated on the \(U_v\). The curve kills \(J\) because it factors through \(R/(P+Q)\). Apply the ordinary deformation quotient comparison to each of the two Hecke ideals defining the tower ideal. Their images in \(B\) lie in every prime \(\mathfrak q\subseteq\ker(y|_B)\). Since \(Q\cap B\subseteq\ker(y|_B)\), their sum \(J\) lies in \(Q\cap B\). Induction along the chain and contraction to \(A\) give the claim.
\end{proof}
\leanref{UpperTheory.\allowbreak{}V.\allowbreak{}Cu.\allowbreak{}chain\_\allowbreak{}rank\_\allowbreak{}four\_\allowbreak{}of\_\allowbreak{}red\_\allowbreak{}a\_\allowbreak{}y\_\allowbreak{}coin\_\allowbreak{}fib\_\allowbreak{}y}{Development.\allowbreak{}AutomorphicArithmetic.\allowbreak{}AutomorphicVarieties.\allowbreak{}ModularityArgument.\allowbreak{}ThresholdArguments.\allowbreak{}DenseThresholdCases.\allowbreak{}ReductionTypeChainOfRedAYCoinFibY}{213--268}
\leanref{UpperTheory.\allowbreak{}V.\allowbreak{}Cu.\allowbreak{}RankFourReductionTypeChain.\allowbreak{}ch\_\allowbreak{}core4}{Development.\allowbreak{}AutomorphicArithmetic.\allowbreak{}AutomorphicVarieties.\allowbreak{}ModularityArgument.\allowbreak{}ThresholdArguments.\allowbreak{}DenseThresholdCases.\allowbreak{}ReductionTypeChainOfRedAYCoinFibY}{78--211}
\leanref{UpperTheory.\allowbreak{}V.\allowbreak{}Ws.\allowbreak{}residual\_\allowbreak{}charpoly\_\allowbreak{}rank\_\allowbreak{}four}{Development.\allowbreak{}AutomorphicArithmetic.\allowbreak{}AutomorphicVarieties.\allowbreak{}ModularityArgument.\allowbreak{}ThresholdArguments.\allowbreak{}DenseThresholdCases.\allowbreak{}ClassicalFixedByIwahori}{90--126}

\begin{theorem}[Finiteness and automorphy of symplectic points]\label{B6.49}
Assume \(\mathscr H\). Let \(h\) satisfy \(h.\operatorname{Big}(t)\), let \(\ell\) be odd and take \(0<k<t\), \(\vartheta\in X_\ell(k)\). For every good four-dimensional datum \(D\) of case II or III:
\begin{enumerate}
\item For every allowed lift over a complete local coefficient algebra its zero-weight algebra is finite over \(O\).
\item For every characteristic-zero block point \(r\),
\[
 \operatorname{span}_{C_\ell}
 \{r(g)_{\{0,2\},\{0,2\}}:g\in G_\Q\}\ne M_2(C_\ell).
 \tag{6.49.1}
\]
\item There is one \(\iota:C_\ell\simeq\C\) for which every irreducible characteristic-zero point has an attachment with witness on every irreducible quadratic CM restriction.
\end{enumerate}
\end{theorem}

\leanref{UpperTheory.\allowbreak{}V.\allowbreak{}HV.\allowbreak{}DescentVanishingReductions.\allowbreak{}sp4In0\_\allowbreak{}alt}{Development.\allowbreak{}AutomorphicArithmetic.\allowbreak{}AutomorphicVarieties.\allowbreak{}ModularityArgument.\allowbreak{}ThresholdArguments.\allowbreak{}DenseThresholdCases.\allowbreak{}Sp4In0}{18--19}
\leanref{UpperTheory.\allowbreak{}V.\allowbreak{}Aut.\allowbreak{}SymplecticFourImageDense}{Development.\allowbreak{}AutomorphicArithmetic.\allowbreak{}AutomorphicVarieties.\allowbreak{}ModularityArgument.\allowbreak{}ThresholdArguments.\allowbreak{}DenseThresholdCases.\allowbreak{}IsAut4}{14--17}

\begin{proof}
Fix \(\iota\) for the Hecke models and their restrictions. Use the ordinary models and seeds of ranks two and four with the class-group bound of \ref{B6.17}. In case II the deep tower joins the residual model to the reduction of the rank-four symplectic seed for case II by a chain satisfying the first four alternatives of \ref{B6.48b}. In case III take a tensor-factor tower adapted to the residual labels and enlarge it to remove the tame obstructions. For the constant \(C\) of the label and cohomology bounds take depth at least
\[
 30(C+1)^2.
 \tag{6.49.2}
\]
This supplies the off-diagonal nilpotence and dimension bounds for
\[
 q_0\xrightarrow{P}q_3\xrightarrow{M}q_4
 \xrightarrow{P}q_2\xrightarrow{M}q_{\mathrm{seed}}.
 \tag{6.49.3}
\]
The \(P\)-links preserve \((2,0)\) entries of the representation and flags and the \(M\)-links preserve \((1,0)\). The points \(q_0,q_3,q_4\) satisfy alternatives among the first four of \ref{B6.48b}. The point \(q_2\) has \(H\) with \(E(-\psi)\), and the seed has \(H\) with \(E(-\beta),E(-\theta)\).

Each link is a polynomial family \(A\to k[X]\) over the residue field. A minimal prime below its kernel lies below both evaluation kernels. Thus containment of \(JA\) in every component at one end implies containment in some component at the other. In the notation of the frame, these are \Lean{AllIn} and \Lean{SomeIn}. Start with the seed component and alternate \ref{B6.48b} with this link argument back to \(q_0\). For a domain-valued lift \(y\), choose a minimal prime below \(\ker y\subseteq\ker q_0\). It contains \(JA\), so \(y\) kills \(JA\).

\leanref{UpperTheory.\allowbreak{}V.\allowbreak{}Cu.\allowbreak{}chainIV\_\allowbreak{}alt}{Development.\allowbreak{}AutomorphicArithmetic.\allowbreak{}AutomorphicVarieties.\allowbreak{}ModularityArgument.\allowbreak{}ThresholdArguments.\allowbreak{}DenseThresholdCases.\allowbreak{}CuPChainIV}{16--82}
\leanref{UpperTheory.\allowbreak{}V.\allowbreak{}Ws.\allowbreak{}conn}{Development.\allowbreak{}AutomorphicArithmetic.\allowbreak{}AutomorphicVarieties.\allowbreak{}WeakShimuraTheory.\allowbreak{}CompactIntegralFields.\allowbreak{}AlgHomResUnique}{56--68}

In characteristic \(\ell\), the Hecke finiteness relation kills a power of the maximal coefficient ideal after inertial specialisation. The characteristic polynomials on \(G_{T.F}\) are therefore the residual products. Every \(g\) has a power in \(G_{T.F}\), so its eigenvalues are algebraic over \(\F_\ell\). The difference between its trace and the sum of residual characters lies in the maximal valuation ideal and is algebraic over \(\F_\ell\), hence zero. Therefore
\[
 \operatorname{tr}\bigl(r(x)[r(g),r(h)]\bigr)=0\qquad(x,g,h\in G_\Q),
\]
with the matrix commutator in brackets. Full matrix span and the nondegenerate trace pairing would make every pair of image matrices commute. At a distinguished block endpoint write \(A_g,B_g\) for the blocks, with \(B_g\) spanning \(M_2\). Their polynomials at the marking \(g_0\) are coprime. The polynomial \(f=\det(X-A_{g_0})\) kills \(A_{g_0}\) and \(f(B_{g_0})\) is invertible. By linearity in products of group elements, insert \(f(r(g_0))\) as an extra factor in the trace identity. The trace pairing on \(M_2\) gives the same contradiction on the \(B\) block. If the zero-weight algebra were infinite over \(O\), the valuation criterion would give a residue-characteristic point with a nonconstant zero-weight value. Such values are constant at triangular points, and the other possibilities give the two excluded cases. 

In characteristic zero the promodular model gives the purity comparison. For automorphy take \(J=\phi_1(J_1)B+\phi_2(J_2)B\) from two ordinary determinant rings with the same twist and different auxiliary primes. The lift kills both ideals. On inertia in the common open kernel of the finite characters its ordinary characters satisfy \(\psi_i^2=\varepsilon_\ell^{2w_i}\) with \((w_i)=(k+1,k,1,0)\), and the model has the prescribed similitude. The residual characters remain pairwise distinct on the tower CM field, so the abelian restriction criterion preserves irreducibility there. The two arithmetic evaluations represent the same twisted restriction. Apply \ref{B6.48a} to each and use the other comparison at its exceptional place. Remove the algebraic twist and use solvable descent with the similitude pairing to reach the tower's quadratic CM subfield. Base change and descent across quadratic CM fields give the attachment on every irreducible quadratic restriction. The local flag and inertia comparisons carry the witness.

\leanref{UpperTheory.\allowbreak{}V.\allowbreak{}Ws.\allowbreak{}contradiction\_\allowbreak{}star\_\allowbreak{}promodular}{Development.\allowbreak{}AutomorphicArithmetic.\allowbreak{}AutomorphicVarieties.\allowbreak{}ModularityArgument.\allowbreak{}ThresholdArguments.\allowbreak{}DenseThresholdCases.\allowbreak{}ContradictionStarPromodular}{183--222}
\leanref{UpperTheory.\allowbreak{}Gal.\allowbreak{}projectedFormsSystemPart.\allowbreak{}promodular\_\allowbreak{}existence\_\allowbreak{}of\_\allowbreak{}local\_\allowbreak{}zero}{Development.\allowbreak{}AutomorphicArithmetic.\allowbreak{}AutomorphicVarieties.\allowbreak{}ModularityArgument.\allowbreak{}ThresholdArguments.\allowbreak{}DenseThresholdCases.\allowbreak{}PmTwo}{21--51}
\leanref{UpperTheory.\allowbreak{}V.\allowbreak{}Ws.\allowbreak{}aut\_\allowbreak{}n\_\allowbreak{}alt}{Development.\allowbreak{}AutomorphicArithmetic.\allowbreak{}AutomorphicVarieties.\allowbreak{}ModularityArgument.\allowbreak{}ThresholdArguments.\allowbreak{}DenseThresholdCases.\allowbreak{}MapRepComp}{201--243}
\leanref{UpperTheory.\allowbreak{}V.\allowbreak{}Z.\allowbreak{}Win.\allowbreak{}Dn.\allowbreak{}pmDown\_\allowbreak{}alt}{Development.\allowbreak{}AutomorphicArithmetic.\allowbreak{}AutomorphicVarieties.\allowbreak{}ModularityArgument.\allowbreak{}ThresholdArguments.\allowbreak{}DenseThresholdCases.\allowbreak{}DnPmDown}{32--87}

For (6.49.1), full block span permits passage to finite coefficients and a companion point with the same characteristic polynomials for the purity comparison. On respective open subgroups the block determinants are \(\varepsilon_\ell^{k+2}\) and \(\varepsilon_\ell^k\). Take the Frobenius of \Lean{nePure\_\allowbreak{}w} in their intersection with the purity subgroup. For \(Q=p^f\) the purity inequality becomes \(Q^{k+2}<Q^2Q^k=Q^{k+2}\), a contradiction.

\leanref{UpperTheory.\allowbreak{}V.\allowbreak{}Cu.\allowbreak{}modularity\_\allowbreak{}rank\_\allowbreak{}four\_\allowbreak{}of\_\allowbreak{}promodularity\_\allowbreak{}case\_\allowbreak{}three}{Development.\allowbreak{}AutomorphicArithmetic.\allowbreak{}AutomorphicVarieties.\allowbreak{}ModularityArgument.\allowbreak{}ThresholdArguments.\allowbreak{}DenseThresholdCases.\allowbreak{}ModularityRankFourOfPromodularity}{20--38}
\leanref{UpperTheory.\allowbreak{}V.\allowbreak{}Cu.\allowbreak{}promodularity\_\allowbreak{}case\_\allowbreak{}three\_\allowbreak{}of\_\allowbreak{}chain\_\allowbreak{}reduction}{Development.\allowbreak{}AutomorphicArithmetic.\allowbreak{}AutomorphicVarieties.\allowbreak{}ModularityArgument.\allowbreak{}ThresholdArguments.\allowbreak{}DenseThresholdCases.\allowbreak{}PromodularityCaseThreeOfLemmas}{131--153}
\leanref{UpperTheory.\allowbreak{}V.\allowbreak{}Ws.\allowbreak{}fin0\_\allowbreak{}no\_\allowbreak{}end\_\allowbreak{}all\_\allowbreak{}irr\_\allowbreak{}of\_\allowbreak{}good\_\allowbreak{}con\_\allowbreak{}purity\_\allowbreak{}aut}{Development.\allowbreak{}AutomorphicArithmetic.\allowbreak{}AutomorphicVarieties.\allowbreak{}ModularityArgument.\allowbreak{}ThresholdArguments.\allowbreak{}DenseThresholdCases.\allowbreak{}VWinPm}{75--106}
\end{proof}

\begin{theorem}[An integer for four-dimensional trace differences]\label{B6.50}
Assume \(\mathscr H\). There are \(c,A\in\N\) as follows. Let \(k\ge2\), \(M\ge1\) and \(p\ge5\) be prime with \(p\nmid M\). Let \(h\ge1\) and \(z_i\in\C\) with \(z_i^h=1\) for \(1\le i\le3\). There is \(D\in\Z\setminus\{0\}\) with
\[
 \log|D|\le c\varphi(h)[M(k+1)]^A\log p.
 \tag{6.50.1}
\]
Let \(\ell\ne2,p\) be prime with \(\ell\mid M\), and fix \(\iota:C_\ell\simeq\C\) and a Frobenius \(\sigma_p\). Let \(X\) be a finite pairwise nonisomorphic set of irreducible representations of \(G_\Q\) satisfying \Lean{IsAutQA}, of level \((M;(k+1,k,1,0))\) with an unlinked flag at each place above \(\ell\). For each \(r\in X\) require a finite cyclotomic character \(\chi\) and \(P\in\GL_4(C_\ell)\) with \(r(g)^tPr(g)=\varepsilon_\ell(g)^{k+1}\chi(g)P\) for all \(g\in G_\Q\). Then
\[
 \Delta_r=\operatorname{tr}r(\sigma_p)-\iota^{-1}(p^{k+1}z_1+p^kz_2+pz_3+1)
\]
is nonzero for all \(r\in X\) and
\[
 \sum_{r\in X}v_\ell(\Delta_r)\le v_\ell(D).
 \tag{6.50.2}
\]
\end{theorem}

\leanref{UpperTheory.\allowbreak{}GalN.\allowbreak{}V.\allowbreak{}numberSp4\_\allowbreak{}n\_\allowbreak{}alt}{Development.\allowbreak{}AutomorphicArithmetic.\allowbreak{}AutomorphicVarieties.\allowbreak{}ModularityArgument.\allowbreak{}ThresholdArguments.\allowbreak{}DenseThresholdCases.\allowbreak{}VWLevel}{18--20}
\leanref{UpperTheory.\allowbreak{}V.\allowbreak{}Af.\allowbreak{}NumberSp4Alt}{Development.\allowbreak{}AutomorphicArithmetic.\allowbreak{}AutomorphicVarieties.\allowbreak{}ModularityArgument.\allowbreak{}ThresholdArguments.\allowbreak{}DenseThresholdCases.\allowbreak{}AfStmt}{34--43}

\begin{proof}
Level and unlinking give the local safety condition for the weights \(1,0\) differing by one. The counting input uses the given attachments to assign distinct isomorphism classes distinct labels in a conjugation-stable set \(S\), of size at most \(c_0[M(k+1)]^{A_0}\), with eigenvalue multisets \(D_{\rm eig}.\beta(s)\). Set
\[
 t_s=\sum_{\beta\in D_{\rm eig}.\beta(s)}\beta,\qquad
 F_{s,a}=t_s-(p^{k+1}z_1^a+p^kz_2^a+pz_3^a+1)
 \quad(a\in(\Z/h)^\times).
\]
The factors are integral. Automorphisms permute the labels and change the three roots by a common unit exponent modulo \(h\), so permute the factors.

For \(x=p^{(k+1)/2}\), \(y=p^k\), purity gives \(|t_s|\le3\sqrt p\,x+x\le\tfrac72 y\) since \(\sqrt p\,x\le y\) and \(2x\le y\). Since \(p+1<y/2\), we have \(|t_s-(p^kz_2^a+pz_3^a+1)|\le\tfrac92y+p+1<py\). The leading term has norm \(py\), so
\[
 F_{s,a}\ne0,\qquad |F_{s,a}|\le2py\le p^{2(k+1)}.
 \tag{6.50.3}
\]
The product \(D=\prod_{s,a} F_{s,a}\) is a nonzero integer with log size at most \(2(k+1)\#S\,\varphi(h)\log p\). Inverting the eigenvalues of the attachment's polynomial at inverse Frobenius identifies \(\iota(\operatorname{tr}r(\sigma_p))=t_s\) at the label \(s\) of \(r\). Hence \(\Delta_r=\iota^{-1}(F_{s,1})\). All other factors have nonnegative valuation.

\leanref{UpperTheory.\allowbreak{}GalN.\allowbreak{}V.\allowbreak{}Win.\allowbreak{}N.\allowbreak{}number\_\allowbreak{}sp4\_\allowbreak{}of\_\allowbreak{}admissible\_\allowbreak{}to\_\allowbreak{}admissible\_\allowbreak{}n}{Development.\allowbreak{}AutomorphicArithmetic.\allowbreak{}AutomorphicVarieties.\allowbreak{}ModularityArgument.\allowbreak{}ThresholdArguments.\allowbreak{}DenseThresholdCases.\allowbreak{}VWinNumN}{94--155}
\end{proof}

\section{Three-dimensional endpoints over the Gaussian field}

The three-dimensional deformation group is \(\GL_3(R)\rtimes\{1,-1\}\) with the polarising involution. Its sign on \(G_\Q\) is the quadratic character \(\eta_4\) of \(\Q(i)\). Its matrix part is a representation on the sign kernel.

Write \(\vartheta_{\ell'}\) for the prime-to-\(\ell\) order part and \(c_2(\vartheta)\) for the conductor exponent at \(2\). With \(\kappa_\ell\) as in \ref{B6.28}, put
\[
 S_b=\{v\mid N:v\ne\ell,\ \operatorname{comp}_v(\kappa_\ell(\vartheta))\ne1\},
 \qquad T_b=\{\ell\}\cup S_b .
\]
A host datum has \(\ell\ne2\), \(k\ge1\), \(2\mid N\), \(\vartheta(-1)=(-1)^k\), \(o_1\ge1\), and
\[
 \operatorname{comp}_r(\vartheta_{\ell'})^{2o_1}=1\quad(r\text{ odd prime}),\qquad
 8o_1v^2<2^{c_2(\vartheta)}\quad(v\in T_b),\qquad c_2(\vartheta)\ge4.
 \tag{6.51.0}
\]
The frame uses the places dividing \(\ell N\), sign \(\eta_4\) and \(\Psi=\varepsilon_\ell^k\vartheta\). At \(\ell\) it prescribes an ordinary upper-triangular flag. At \(v\in S_b\) a single basis congruent to the identity makes \(D_v\cap\ker\eta_4\) diagonal with prescribed diagonal on \(I_v\cap\ker\eta_4\). At the remaining places it kills inertia on the common kernel of the sign and residual \(\Psi\).

\begin{theorem}[The Gaussian three-dimensional lifting assertion]\label{B6.51}
Assume the simultaneous marking, residue-characteristic polynomial identity and characteristic-zero attachment of \ref{B6.52}(4). Let \(K\simeq\Q(i)\), fix a host datum and its frame \(F\) on \(\Gamma_{\Q,S}\), and fix \(\iota:C_\ell\simeq\C\). There is \(g_1\in\Gamma_{\Q,S}\) with
\[
 \eta_4(g_1)=1,\qquad\Psi(g_1)^2-1\in O^\times
 \tag{6.51.1}
\]
such that for every universal lift \(u\) of the marked normalised frame over \(R\), \(R_0=\operatorname{wpart}(O,u,0)\subset R\) is finite over \(O\). Let \(\xi\) be an \(O\)-algebra map to \(C_\ell\) from the subalgebra of \(R\) generated by the weight parts, and \(r\) an associated characteristic-zero point. If the matrices on \(\ker\eta_4\) act irreducibly, then \((r|_{G_K})_{\mathrm{mat}}\otimes\varepsilon_\ell^{-1}\) is conjugate self-dual automorphic with witness.
\end{theorem}

\leanref{UpperTheory.\allowbreak{}V.\allowbreak{}Aut.\allowbreak{}ThmS6Core\_\allowbreak{}s}{Development.\allowbreak{}AutomorphicArithmetic.\allowbreak{}AutomorphicVarieties.\allowbreak{}ModularityArgument.\allowbreak{}ThresholdArguments.\allowbreak{}DenseThresholdCases.\allowbreak{}ThmS6CoreS}{27--36}
\leanref{UpperTheory.\allowbreak{}V.\allowbreak{}Pr.\allowbreak{}symmetric\_\allowbreak{}cube\_\allowbreak{}theorem\_\allowbreak{}of\_\allowbreak{}endomorphisms\_\allowbreak{}alt}{Development.\allowbreak{}AutomorphicArithmetic.\allowbreak{}AutomorphicVarieties.\allowbreak{}ModularityArgument.\allowbreak{}ThresholdArguments.\allowbreak{}DenseThresholdCases.\allowbreak{}SymmetricCubeTheoremOfEndomorphisms}{76--131}

\begin{proof}
The endpoint supplies \(g_1\) and the reducible characteristic polynomials on the sign kernel of a finite-index \(U\). Pullback preserves the marking and local conditions. Over complete DVRs in residue characteristic, these polynomials and the marking give triangularity on the sign kernel, so the strict generators of one sign vanish. The condition \Lean{frame\_\allowbreak{}class\_\allowbreak{}one} fixes the reducible diagonals. The graded valuation criterion gives finiteness of \(R_0\).

In characteristic zero, weight-zero finiteness gives a finite coefficient DVR model by diagonal conjugation. For the normal core \(U^{\rm core}\) of any open subgroup in \(\ker\eta_4\), choose \(\sigma\in I_\ell\) with cyclotomic value \(1+\ell\). Then \(u=\sigma^{[G_\Q:U^{\rm core}]\varphi(N)}\) lies in \(I_\ell\cap U^{\rm core}\) with \(\vartheta(u)=1\), \(\Psi(u)^2\ne1\). By (6.51.0) take \(x\in D_\ell\cap\ker\eta_4\) with \(\Psi(x)^2-1\in O^\times\). The local criterion with \(x,u\) gives irreducibility on the core and hence on the open subgroup. Apply the characteristic-zero endpoint and transport the conclusion through diagonal conjugation on the image of \(G_K\).

\leanref{UpperTheory.\allowbreak{}V.\allowbreak{}UnitaryThreeInertia.\allowbreak{}universal\_\allowbreak{}diagonal\_\allowbreak{}of\_\allowbreak{}triangular\_\allowbreak{}points}{Development.\allowbreak{}AutomorphicArithmetic.\allowbreak{}AutomorphicVarieties.\allowbreak{}UnitaryThreeDeformations.\allowbreak{}UniversalDiagonalOfTriangularPoints}{22--73}
\end{proof}

\begin{lemma}[Points on the linked CM frame]\label{B6.51a}
Assume \(\mathscr H\). Put \(\Psi=\varepsilon_\ell^e\vartheta\) with \(e\ge1\) and \(N\) even. Let \(\ell N\mid M\) and suppose every prime ramified in the sign \(\eta\) divides \(M\). Fix \(B\mid\ell\), \(\iota:C_\ell\simeq\C\) and a Galois link for \((\eta,\Psi,c,g_1,B,M)\) to a frame \(\Phi\) over an abelian CM field \(L\). Take the residual frame and seed data \Lean{Ok} and \Lean{SeedClause}, with \(d=\sum_v d_v>3C_0+7\) and every \(d_v>5\). Put \(w_{3,i}=1-i\) and \(w_{2,i}=1-2i\) for \(i\) starting at zero.

For \(n=2,3\), take an adically continuous framed lift over a coefficient DVR in \(C_\ell\) finite over \(O\), with prescribed diagonal reduction and inertia characters the pullbacks of \(\Psi^{w_{n,i}}\). If the induced representation \(r\) on \(G_L\) is irreducible, \(r\kappa_n\) is conjugate self-dual automorphic with witness, where \(\kappa_3=\varepsilon_\ell^{-1}\) and one algebraic \(\kappa_2\) works for all rank-two points. In rank three, over a complete coefficient DVR \(A\) with finite residue field of characteristic \(\ell\) and \(\mathfrak m_OA=0\), impose the same residual condition and require the inertia characters to kill \(\ker\bar\Psi\) in inertia. The polynomial on the sign kernel of the frame is then \(\prod_{i=0}^2(X-\bar\Psi(g)^{1-i})\).
\end{lemma}
\begin{proof}
The seed data give a quadratic field \(K_0\) and characters \(\chi'\ne\eta'\) with \(\eta'=\eta\), \(\ell\nmid D_{K_0}\operatorname{ord}\chi'\) and \(\varepsilon_\ell^j\chi'\equiv\Psi\) modulo the maximal ideal. Use the ordinary rank-two seed. In rank three take the trace-zero adjoint of the two-dimensional Eisenstein-congruence lift of weight \(m+1\), with \(m=j+(\ell-1)t\) above the seed threshold. On \(G_L\), \(\eta=1\) and \(\bar\varepsilon_\ell^{m-j}=1\). The adjoint has residual characters \(\bar\Psi,1,\bar\Psi^{-1}\) and weights \(m,0,-m\). The adjoint lift gives its cuspidal attachment.

In rank \(n\) the seed gives a prime over the Hecke ideal of quotient dimension \(1+nd\) in the coefficient ring. Among the finitely many matrix-ring components meeting the irreducible locus, take the union of those contained in the Hecke locus and the union of the others. Each component meets the connected residual locus cut out by the maximal coefficient ideal with the upper triangular locus removed. The two unions cover it, so if both occur they share a closed point there. In its completion the ordinary curve comparison propagates containment between irreducible branches. The block intersections in rank three use rank-two containment and ordinary splitting. The seed gives containment at all irreducible points, first in rank two and then in rank three.

The resulting evaluation has determinant that of \(r\kappa_n\) and weights \(a_\tau-e w_{n,i}\) for \(a_\tau\) the local weight of \(\kappa_n\). These increase strictly in \(i\). Apply \ref{B6.48a} with two distinct auxiliary rational primes and the same twist. Together the comparisons cover all places. Unlinked toral support supplies one exponent for each attachment.

In residue characteristic the rank-two block gives an ideal \(I_3\) in rank three. Off the upper and lower triangular loci, evaluation kills \(I_3\cap J_3\), where \(J_3\) is the propagated ideal. Imposing the inertia parameters makes the radical contain the maximal coefficient ideal, so weight-zero coefficients take their residual values. In the determinant expansion the off-diagonal two-cycles and three-cycles vanish and the diagonals give the stated polynomial. At triangular points the ratios of diagonal characters to their residual lifts have finite exponent by the inertia generators and abelian class-group bound. Principal units of an equal-characteristic DVR are torsion-free, so these ratios are one.
\end{proof}
\leanref{UpperTheory.\allowbreak{}Gal.\allowbreak{}projectedFormsSystemPart.\allowbreak{}linked\_\allowbreak{}endomorphism\_\allowbreak{}dense}{Development.\allowbreak{}AutomorphicArithmetic.\allowbreak{}AutomorphicVarieties.\allowbreak{}ModularityArgument.\allowbreak{}ThresholdArguments.\allowbreak{}DenseThresholdCases.\allowbreak{}LinkedEndE}{18--67}
\leanref{UpperTheory.\allowbreak{}V.\allowbreak{}Pr.\allowbreak{}absMain}{Development.\allowbreak{}AutomorphicArithmetic.\allowbreak{}AutomorphicVarieties.\allowbreak{}CurveDVRSpecialization.\allowbreak{}AbsMain}{12--30}
\leanref{UpperTheory.\allowbreak{}Gal.\allowbreak{}projectedFormsSystemPart.\allowbreak{}twoRuns\_\allowbreak{}alt}{Development.\allowbreak{}AutomorphicArithmetic.\allowbreak{}AutomorphicVarieties.\allowbreak{}ModularityArgument.\allowbreak{}ThresholdArguments.\allowbreak{}DenseThresholdCases.\allowbreak{}TwoRuns}{80--109}

\begin{theorem}[The endpoints used in cases I and V]\label{B6.52}
Assume \(\mathscr H\). For each good imaginary-quadratic datum of the first count, each \(\iota:C_\ell\simeq\C\) and each \(B\mid\ell\), there is one marking \(g_1\) for assertions 1--3. Write \(\eta\) for the sign and \(c\) for complex conjugation. An ordinary point has sign \(\eta\), is adically continuous and has matrix reduction equal to that of \(r_0\). Its matrix at \(c\) is \(1\), and the \(01,10\) entries of \((r(g_1))_{\rm mat}-(r(cg_1^{-1}c))_{\rm mat}\) vanish. Away from \(\ell\) it kills inertia elements with \(\eta(s)=\Psi(s)=1\). At \(B\) a basis congruent to the identity makes its matrices upper triangular on \(D_B\cap\ker\eta\), with diagonal \((\Psi,1,\Psi^{-1})\) on \(I_B\cap\ker\eta\).
\begin{enumerate}
\item There is a finite-index \(U\le G_\Q\) on which every ordinary point over a complete DVR \(A\) with finite residue field and \(\mathfrak m_OA=0\) has the reducible characteristic polynomials on \(U\cap\ker\eta\).
\item Let \(r\) be ordinary over a DVR \(O'\subset C_\ell\) finite over \(O\). If its matrix part is irreducible on every open \(U\le G_\Q\) in \(\ker\eta\), its restriction to the specified imaginary quadratic field twisted by \(\varepsilon_\ell^{-1}\) is conjugate self-dual automorphic with witness.
\item Let \(r\) be Type II ordinary over such \(O'\), with trivial middle character on \(\ker\eta\). For every \(W:G_\Q\to\GL_2(C_\ell)\) agreeing with the corner of \(r\) on \(\ker\eta\) and with \(\operatorname{tr}W(c)=0\), \(W\varepsilon_\ell^{-k}\) has a newform attachment of weight \(2k+1\) and \(W\) satisfies \Lean{WitAQ}.
\item For each Gaussian host satisfying (6.51.0), each \(\iota\), \(c\) and \(B\mid\ell\), a single \(g_1\) satisfying (6.51.1) makes assertions 1 and 2 hold for \((\eta_4,\Psi,c,g_1,B)\). The lifting conclusion follows by \ref{B6.51}.
\end{enumerate}
\end{theorem}

\leanref{UpperTheory.\allowbreak{}V.\allowbreak{}Z.\allowbreak{}Win.\allowbreak{}Dn.\allowbreak{}ends0\_\allowbreak{}alt}{Development.\allowbreak{}AutomorphicArithmetic.\allowbreak{}AutomorphicVarieties.\allowbreak{}ModularityArgument.\allowbreak{}ThresholdArguments.\allowbreak{}DenseThresholdCases.\allowbreak{}DnEndsTop}{15--16}
\leanref{UpperTheory.\allowbreak{}V.\allowbreak{}Z.\allowbreak{}Win.\allowbreak{}Dn.\allowbreak{}ends\_\allowbreak{}of\_\allowbreak{}alt}{Development.\allowbreak{}AutomorphicArithmetic.\allowbreak{}AutomorphicVarieties.\allowbreak{}ModularityArgument.\allowbreak{}ThresholdArguments.\allowbreak{}DenseThresholdCases.\allowbreak{}DnEnds}{35--39}
\leanref{UpperTheory.\allowbreak{}V.\allowbreak{}Pr.\allowbreak{}GdEndAlt}{Development.\allowbreak{}AutomorphicArithmetic.\allowbreak{}AutomorphicVarieties.\allowbreak{}ModularityArgument.\allowbreak{}ThresholdArguments.\allowbreak{}DenseThresholdCases.\allowbreak{}AutEnd3AltAlt}{35--40}

\begin{proof}
Take \(M=\ell N|D_K|\) for the first count and \(M=\ell N\) for the Gaussian host. The sign ramifies at primes dividing \(D_K\) in the first case and at \(2\mid N\) in the second. Parity gives \(\Psi(c)=1\), and \(c^2=1\), \(\eta(c)=-1\). The frame construction chooses the seed data and a class-rank bound \(C_0\). It also chooses a layer with \(\ell^{3C_0+7}>3C_0+7\). It gives the link of \ref{B6.51a} and \(g_1\in G_L\) with \(\bar\Psi(g_1)^2\ne1\).

For (1) take \(U=G_L\). The ordinary inertia characters satisfy \ref{B6.51a} after reduction. For (2), \(G_L\) is open in \(G_\Q\) and lies in \(G_K=\ker\eta\), so the restriction is irreducible and \ref{B6.51a} applies. The frame gives the twisted pairing with multiplier \(\varepsilon_\ell^{-2}\), and local comparison gives abelian inertial parameters and ordinary flags with unlinking preserved upon restriction. Solvable descent over the abelian extension \(L/K\) gives the attachment and witness over \(K\).

For (3), the corner \(W|_{G_K}\) is irreducible by Type II. If its restriction to \(G_L\) were reducible, the rank-two self-twist criterion for the abelian quotient \(G_K/G_L\) would give a nontrivial character \(\lambda\) of that quotient with \(W|_{G_K}\simeq W|_{G_K}\otimes\lambda\). Choose \(g\in G_K\) with \(\lambda(g)\ne1\). Since \(g_1\in G_L\), both \(W(g)\) and \(W(gg_1)\) have trace zero. The residual corner has diagonal characters \(\bar\Psi,\bar\Psi^{-1}\), so \(\bar\Psi(g)^2=\bar\Psi(gg_1)^2=-1\). This forces \(\bar\Psi(g_1)^2=1\), contrary to the marking. Apply the rank-two part of \ref{B6.51a} over \(L\). For the extension \(W\) with \(\operatorname{tr}W(c)=0\), odd rank-two descent gives the newform and the quadratic-restriction witness. In the Gaussian case, (1) and (2) with the chosen marking follow from the same construction and give the premise of \ref{B6.51}.

\leanref{UpperTheory.\allowbreak{}V.\allowbreak{}Fr.\allowbreak{}host\_\allowbreak{}ex}{Development.\allowbreak{}AutomorphicArithmetic.\allowbreak{}AutomorphicVarieties.\allowbreak{}FrobeniusSmoothVectors.\allowbreak{}HostEx}{8--26}
\leanref{UpperTheory.\allowbreak{}V.\allowbreak{}Z.\allowbreak{}Win.\allowbreak{}Dn.\allowbreak{}autEnd3\_\allowbreak{}alt}{Development.\allowbreak{}AutomorphicArithmetic.\allowbreak{}AutomorphicVarieties.\allowbreak{}ModularityArgument.\allowbreak{}ThresholdArguments.\allowbreak{}DenseThresholdCases.\allowbreak{}AutEnd3Alt}{13--106}
\leanref{UpperTheory.\allowbreak{}V.\allowbreak{}Z.\allowbreak{}Win.\allowbreak{}Dn.\allowbreak{}hostEnd\_\allowbreak{}of\_\allowbreak{}loc3\_\allowbreak{}frame\_\allowbreak{}linkedEnd\_\allowbreak{}downCsdK}{Development.\allowbreak{}AutomorphicArithmetic.\allowbreak{}AutomorphicVarieties.\allowbreak{}ModularityArgument.\allowbreak{}ThresholdArguments.\allowbreak{}DenseThresholdCases.\allowbreak{}DnHost}{17--67}
\leanref{UpperTheory.\allowbreak{}V.\allowbreak{}Z.\allowbreak{}Win.\allowbreak{}Dn.\allowbreak{}AutEnd2Ao}{Development.\allowbreak{}AutomorphicArithmetic.\allowbreak{}AutomorphicVarieties.\allowbreak{}ModularityArgument.\allowbreak{}AutomorphyLifting.\allowbreak{}AutomorphyDescent.\allowbreak{}AutEnd2Ao}{18--26}
\end{proof}

\section{Adjoint vanishing and ordinary derivations}

Let \(L\) be CM with totally real subfield \(L^+\). A polarised ordinary frame gives a quotient \(\Gamma\) of \(G_{L^+}\), its index-two CM subgroup \(\Gamma_0\), a complex conjugation \(c\), the cyclotomic character and the local maps and subgroups. Write \(I_0,D_0\) for the distinguished inertia and decomposition subgroups. In an admissible residual ordinary datum the diagonal characters kill \(I_0\) and coincide on \(D_0\), where the residual cyclotomic character is nontrivial.

For the universal ordinary determinant ring \(B\), a characteristic-zero point \(x\) has determinant law \(x.D\) and ordered characters \(x.\psi_{v,i}\). For a representation \(\rho\) of \(\Gamma_0\) and polarisation \(P\), the extended adjoint on \(M_n(C_\ell)\) is
\[
 R(g)X=\rho(g)X\rho(g)^{-1}\quad(g\in\Gamma_0),
 \qquad R(c)X=-P^{-1}X^tP.
 \tag{6.53.0}
\]

\begin{theorem}[Bloch--Kato vanishing for the extended adjoint]\label{B6.53}
Assume \(\mathscr H\). Take a CM Hida setting \(s\) and frame \(V\) with \(V.S'=s.S\). Suppose \(L\) contains no primitive \(\ell\)-th root of unity. Let \(O\subset C_\ell\) contain the images of \(\mathcal O_L\) under every embedding \(L\to C_\ell\), and take \(2\le n\le4\). Fix an admissible residual ordinary datum with universal determinant ring \(B\), an isomorphism \(\iota:C_\ell\simeq\C\), an arithmetic point \(x\), and \(\rho:\Gamma_0\to\GL_n(C_\ell)\) of determinant law \(x.D\). Put \(\Gamma'=\Gamma_0\cap\ker\varepsilon_\ell\) and require the following.
\begin{enumerate}
\item The continuous \(\rho\) is trivial on \(I_0\), and its pullback to \(G_L\) is conjugate self-dual automorphic with witness.
\item The ordinary flags of \(\rho\) have ordered characters \(\kappa_{v,i}=x.\psi_{v,i}\).
\item The restriction \(\rho|_{\Gamma'}\) is irreducible.
\item Every open \(U\subseteq\Gamma_0\) meeting \(\Gamma'\) contains \(\sigma\in U\cap\Gamma'\) of separable characteristic polynomial.
\item At each \(v\mid\ell\) some local element mapping to \(\Gamma'\) gives infinite order for every ratio of distinctly indexed ordinary characters.
\item At each \(v\mid\ell\) there are strictly decreasing integers \(m_i\) and \(h\ge1\) with
\[
 \left(\frac{\kappa_{v,i}(s)}{\kappa_{v,j}(s)}
 \varepsilon_\ell(s)^{m_j-m_i}\right)^h=1
 \qquad(s\in I_v,\ i,j).
 \tag{6.53.1}
\]
\end{enumerate}
For every \(P\in\GL_n(C_\ell)\) with \(\rho(g)^tP\rho(cgc^{-1})=\varepsilon_\ell(g)^{1-n}P\) on \(\Gamma_0\), and every extended action (6.53.0), let \(z:\Gamma\to M_n(C_\ell)\) be a continuous cocycle satisfying the Bloch--Kato local conditions on pullback to \(G_{L^+}\). Then some \(X\in M_n(C_\ell)\) satisfies
\[
 z(g)=R(g)X-X\qquad(g\in\Gamma).
 \tag{6.53.2}
\]
\end{theorem}

\leanref{UpperTheory.\allowbreak{}V.\allowbreak{}Th.\allowbreak{}AdZero}{Development.\allowbreak{}AutomorphicArithmetic.\allowbreak{}AutomorphicVarieties.\allowbreak{}TheoremsMuLinked.\allowbreak{}OrdZero}{68--73}
\leanref{UpperTheory.\allowbreak{}V.\allowbreak{}Th.\allowbreak{}AdjointLiftingTheorem}{Development.\allowbreak{}AutomorphicArithmetic.\allowbreak{}AutomorphicVarieties.\allowbreak{}ModularityArgument.\allowbreak{}ThresholdArguments.\allowbreak{}DenseThresholdCases.\allowbreak{}AutomorphicPointGalois}{35--39}
\leanref{UpperTheory.\allowbreak{}V.\allowbreak{}HV.\allowbreak{}DescentVanishingReductions.\allowbreak{}adjoint\_\allowbreak{}lifting\_\allowbreak{}of\_\allowbreak{}cohomology}{Development.\allowbreak{}AutomorphicArithmetic.\allowbreak{}AutomorphicVarieties.\allowbreak{}ModularityArgument.\allowbreak{}ThresholdArguments.\allowbreak{}DenseThresholdCases.\allowbreak{}AdjointLiftingOfCohomology}{31--31}

\begin{proof}
Away from \(\ell\) the inertia coboundary vanishes on residual inertia where \(\rho\) is trivial. At \(v\mid\ell\), (6.53.1) and the infinite-order ratios give \(Y_v\) such that \(z-dY_v\) is upper triangular on \(D_v\) with zero diagonal on \(I_v\) in the flag basis. Write \(dX(g)=R(g)X-X\) and \(X_0=-z(c)/2\). Since \(R(c)z(c)=-z(c)\), \(z'=z-dX_0\) vanishes at \(c\) and satisfies \(z'(cgc^{-1})=R(c)z'(g)\). Its local corrections are \(Y_v-X_0\).

Cyclotomic irreducibility gives full matrix span. The infinite-order ratios on \(\Gamma'\) exclude \(\kappa_{v,i}=\kappa_{v,j}\varepsilon_\ell\) for \(i\ne j\). For equal indices use local cyclotomic nontriviality. The finite-dimensional space of cyclotomic-twisted adjoint cocycles admits finitely many trace tests on separable elements of \(\Gamma'\) with the same common kernel as all such tests, by \Lean{finDet}. Hypothesis 4 and continuity then give \(\operatorname{tr}(u(g)\rho(g))=0\) on \(\Gamma'\) for a cocycle in the common kernel. By \Lean{detG} it is a coboundary. This is determinant regularity for \ref{B6.54}.

Subtract the coboundary given by \ref{B6.54} on \(\Gamma_0\), extended to \(\Gamma\) with the same matrix. The resulting cocycle factors through the order-two quotient with values in the \(\Gamma_0\)-invariants, where averaging gives a coboundary.

\leanref{UpperTheory.\allowbreak{}V.\allowbreak{}Th.\allowbreak{}adjointLiftingTheorem\_\allowbreak{}of\_\allowbreak{}locPrints\_\allowbreak{}of\_\allowbreak{}core}{Development.\allowbreak{}AutomorphicArithmetic.\allowbreak{}AutomorphicVarieties.\allowbreak{}ModularityArgument.\allowbreak{}ThresholdArguments.\allowbreak{}DenseThresholdCases.\allowbreak{}AdLOf}{104--150}
\end{proof}

An ordinary cocycle on \(\Gamma_0\) is continuous and satisfies
\[
 z(gh)=z(g)+\rho(g)z(h)\rho(g)^{-1},\qquad
 Pz(cgc^{-1})=-z(g)^tP .
 \tag{6.54.0}
\]
It vanishes on prescribed residual inertia away from \(\ell\). At each \(v\mid\ell\), in an ordinary flag basis for \(\rho\) with characters \(x.\psi_{v,i}\), subtract a local coboundary to make it upper triangular on \(D_v\) with zero diagonal on \(I_v\).

Determinant regularity requires finitely many separable \(s_j\in\Gamma'\) such that a continuous \(u\) satisfying
\[
 u(gh)=u(g)+\varepsilon_\ell(g)\rho(g)u(h)\rho(g)^{-1},\qquad
 Pu(cgc^{-1})=u(g)^tP
\]
is a twisted coboundary whenever \(\operatorname{tr}(u(s_j)\rho(s_j)^a)=0\) for every \(j,a\ge0\).

\begin{theorem}[Ordinary cocycles are coboundaries]\label{B6.54}
Assume \(\mathscr H\). For \(2\le n\le4\), take an admissible residual ordinary datum in a Hida level frame \(V\) with \(V.S'=s.S\), and fix \(\iota:C_\ell\simeq\C\) and an arithmetic weight \(h\). Let \(x\) be an ordinary determinant point of weight \(h\) and \(\rho\) continuous with determinant law \(x.D\). Suppose \(\rho\) is conjugate self-dual automorphic with witness and polarisation \(P\), trivial on \(I_0\), and
\[
 \operatorname{span}_{C_\ell}\rho(\Gamma_0)=M_n(C_\ell),\qquad
 x.\psi_{v,i}\ne x.\psi_{v,j}\varepsilon_\ell\quad(v\mid\ell,\ i,j).
 \tag{6.54.1}
\]
Assume determinant regularity. Then every ordinary cocycle \(z\) has the form
\[
 z(g)=\rho(g)Y\rho(g)^{-1}-Y\qquad(g\in\Gamma_0)
 \tag{6.54.2}
\]
for one \(Y\in M_n(C_\ell)\).
\end{theorem}

\leanref{UpperTheory.\allowbreak{}V.\allowbreak{}Th.\allowbreak{}OrdVan2Alt}{Development.\allowbreak{}AutomorphicArithmetic.\allowbreak{}AutomorphicVarieties.\allowbreak{}ModularityArgument.\allowbreak{}ThresholdArguments.\allowbreak{}DenseThresholdCases.\allowbreak{}OrdVan2}{13--20}
\leanref{UpperTheory.\allowbreak{}V.\allowbreak{}Th.\allowbreak{}ordinary\_\allowbreak{}vanishing\_\allowbreak{}two\_\allowbreak{}of\_\allowbreak{}printed\_\allowbreak{}hida}{Development.\allowbreak{}AutomorphicArithmetic.\allowbreak{}AutomorphicVarieties.\allowbreak{}ModularityArgument.\allowbreak{}ThresholdArguments.\allowbreak{}DenseThresholdCases.\allowbreak{}TOrdVan}{159--192}

\begin{proof}
The determinant law of \(\rho_\epsilon=(1+\epsilon z)\rho\) and its ordered characters give a derivation at \(x\). Equation (6.54.0) imposes polarisation and the local condition kills inertia derivatives. Extend coefficients and normalise the flags using (6.54.1) to get an integral point \(x'\) and \(x'\)-derivation \(B'\to O'\). Arithmetic specialisation of the Hida family with these flags gives \(J\subseteq\ker x'\) for the level ideal \(J\). With the separated roots and trace detection of determinant regularity, \ref{B6.55} gives
\[
 \operatorname{tr}(z(g)\rho(g))=0\qquad(g\in\Gamma_0).
 \tag{6.54.3}
\]

Set \(\partial(\rho(g))=z(g)\rho(g)\). For a finite relation \(\sum_g a_g\rho(g)=0\), differentiate \(\operatorname{tr}(\sum_g a_g\rho_\epsilon(g)\rho_\epsilon(h))\). Its derivative is zero by (6.54.3) at \(gh\), and the term with \(\partial(\rho(h))\) vanishes by the relation. Hence
\[
 \operatorname{tr}\!\left(\left(\sum_g a_gz(g)\rho(g)\right)\rho(h)\right)=0
 \qquad(h\in\Gamma_0).
\]
The full span and trace pairing force the sum inside the pairing to be zero, so \(\partial\) is well-defined on the matrix algebra. The cocycle identity gives the Leibniz rule. Write the inner derivation as \(\partial(A)=AY-YA\) and multiply by \(\rho(g)^{-1}\) to obtain (6.54.2).
\end{proof}

\begin{theorem}[Inertia derivatives determine trace derivatives]\label{B6.55}
Assume the rank-\(n\) automorphic and cohomological inputs in \(\mathscr H\). Let \(n\ge1\) and take a Hida level frame \(V\) with \(V.S'=s.S\). Let \(B\) be a universal ordinary determinant ring with universal pair \((D^{\rm univ},\psi^{\rm univ})\), \(x:B\to O\) an integral point of arithmetic weight \(h\), and \(J\subseteq\ker x\). Require the ordinary Hecke level-ideal condition, with \(D=n[L^+:\Q]\) weight parameters represented by inertia elements, classical auxiliary modules and control at every integral arithmetic specialisation.

Take an integral point \(y\) with \(x_*(D^{\rm univ},\psi^{\rm univ})=(D_{y.\rho},y.\psi)\). Its representation \(y.\rho:\Gamma_0\to\GL_n(O)\) is adically continuous, with integral symmetric or alternating polarisation satisfying the identity of \ref{B6.53}. It kills the designated inertia groups and its image spans an \(O\)-module containing \(cM_n(O)\) for some \(c\ne0\). At each \(v\mid\ell\) take an integral flag basis with ordered characters \(y.\psi_{v,i}\). Their values at some local element are distinct, and \(y.\psi_{v,i}\ne y.\psi_{v,j}\varepsilon_\ell\) for \(i<j\). Require for \(Z\in M_n(O)\)
\[
 \varepsilon_\ell(g)y.\rho(g)Z=Zy.\rho(g)\quad(g\in D_0)
 \quad\Longrightarrow\quad Z=0.
 \tag{6.55.1}
\]
Choose \(s_1,\ldots,s_{q_0}\in\Gamma_0\cap\ker\varepsilon_\ell\) with \(\det(X-y.\rho(s_j))=\prod_i(X-\alpha_{j,i})\) having distinct \(\alpha_{j,i}\in O\) for each \(j\), and assume determinant regularity on these elements. Then every \(O\)-linear \(x\)-derivation \(\delta:B\to O\) killing the universal ordinary characters on inertia satisfies
\[
 \delta\bigl(D^{\rm univ}.\operatorname{tr}(g)\bigr)=0\qquad(g\in\Gamma_0).
 \tag{6.55.2}
\]
\end{theorem}

\leanref{UpperTheory.\allowbreak{}V.\allowbreak{}Th.\allowbreak{}core\_\allowbreak{}o\_\allowbreak{}of\_\allowbreak{}literal\_\allowbreak{}automorphy\_\allowbreak{}n}{Development.\allowbreak{}AutomorphicArithmetic.\allowbreak{}AutomorphicVarieties.\allowbreak{}ModularityArgument.\allowbreak{}ThresholdArguments.\allowbreak{}DenseThresholdCases.\allowbreak{}TCoreO}{129--159}
\leanref{UpperTheory.\allowbreak{}V.\allowbreak{}Th.\allowbreak{}OrdinaryMainCase}{Development.\allowbreak{}AutomorphicArithmetic.\allowbreak{}AutomorphicVarieties.\allowbreak{}ModularityArgument.\allowbreak{}ThresholdArguments.\allowbreak{}DenseThresholdCases.\allowbreak{}ThDefs}{53--64}

\begin{proof}
Use \ref{B6.56} at the \(s_j\), with its adapted residual and local kernels. The determinant coefficients and characters generate \(B\) adically. The classical auxiliary modules give the operator identities in \ref{B6.57}, and arithmetic control gives its annihilation condition.

On \(H\) the residual image is unipotent, so the lifted image is pro-\(\ell\) and kills \(\mathrm{PL}_\ell(H)\) on the CM subgroup. Write \(\bar\rho_{r,m}\) for the pullback at level \(r\) reduced modulo \(\mathfrak m_O^m\) and \(\mathcal S_{r,m}^{\operatorname{tr}}\) for the functions \(g\mapsto\operatorname{tr}(z(g)\bar\rho_{r,m}(g))\) from ordinary Selmer cocycles. The single-layer cohomology count under (6.55.1) and determinant regularity gives at each \(m\), for large \(r\),
\[
 \#\mathcal S_{r,m}^{\operatorname{tr}}Q^{m((n+b)s+bd)}
 \le Q^{C+m((n+b)(s+d)+nq_0)},\qquad Q=\#(O/\mathfrak m_O),
 \tag{6.55.3}
\]
with \(b=\binom n2\), \(s=\#\{v\mid\ell\}\) in \(L^+\), \(d=\sum_{v\mid\ell}d_v=[L^+:\Q]\), \(d_v=[L_v^+:\Q_\ell]\). A place above \(\ell\) contributes at most \(m(n+b)(1+d_v)+C_v\) to the exponent and an auxiliary place at most \(mn+C_j\). The constants and bounded global terms are uniform in \(m\). Cancel the local flag factor to obtain the exponent \(m(D+nq_0)+C\). Since \(\delta\) kills the selected inertia weights, \ref{B6.57} applies.
\end{proof}

\begin{theorem}[Taylor--Wiles auxiliary places]\label{B6.56}
Assume the local and global cohomological inputs. Fix a CM Hida setting with level frame \(V_0.S'=s.S\) and residual datum. Suppose every open subgroup of \(\Gamma\) has finitely many continuous additive \(\F_\ell\)-characters. For \(s_1,\ldots,s_{q_0}\in\Gamma_0\cap\ker\varepsilon_\ell\), there is an auxiliary system of distinct places \(v_{r,j}\) of \(L^+\) outside \(s.S\) splitting in \(L\), with chosen \(w_{r,j}\mid v_{r,j}\) and
\[
 \mathrm Nw_{r,j}\equiv1\pmod{\ell^r}.
 \tag{6.56.1}
\]
For open normal \(H\le\Gamma\), put
\[
 \mathrm{PL}_\ell(H)=
 \bigcap_{\substack{U\triangleleft_{\rm open}\Gamma\\ H/(H\cap U)\text{ an }\ell\text{-group}}}U .
\]
Writing \(\bar\chi_i\) for the residual diagonal characters, the system uses
\[
 H=\{g\in\Gamma_0:\bar\varepsilon_\ell(g)=1,\ \bar\chi_i(g)=1\text{ for all }i\}
\]
and gives enlarged groups \(h_r:\Gamma_r\twoheadrightarrow\Gamma\). Its inverse-Frobenius and tame inertia elements \(\phi_{r,j},\tau_{r,j}\in(\Gamma_r)_0\) have \(h_r(\phi_{r,j})s_j^{-1}\in H\), in every open normal \(U\supseteq\mathrm{PL}_\ell(H)\) for all large \(r\). Moreover \(h_r(\tau_{r,j})=1\) and
\[
 \phi_{r,j}^{-1}\tau_{r,j}\phi_{r,j}
 \equiv\tau_{r,j}^{\mathrm Nw_{r,j}}\pmod{\mathrm{PL}_\ell(h_r^{-1}H)}.
 \tag{6.56.2}
\]
The kernel of \(h_r\) lies in the closed normal subgroup generated by the \(\tau_{r,j}\) and \(\mathrm{PL}_\ell(h_r^{-1}H)\). The continuous \(\F_\ell\)-character spaces of \(h_r^{-1}H\) have uniformly bounded cardinality.
\end{theorem}

\leanref{UpperTheory.\allowbreak{}V.\allowbreak{}Pa.\allowbreak{}arPl}{Development.\allowbreak{}AutomorphicArithmetic.\allowbreak{}AutomorphicVarieties.\allowbreak{}PatchingModules.\allowbreak{}AutomorphicDetection.\allowbreak{}ArPl}{13--180}
\leanref{UpperTheory.\allowbreak{}V.\allowbreak{}Pa.\allowbreak{}ArPl}{Development.\allowbreak{}AutomorphicArithmetic.\allowbreak{}AutomorphicVarieties.\allowbreak{}PatchingModules.\allowbreak{}TaylorWilesObstructionBound}{119--125}

\begin{proof}
The residual kernel lemma \Lean{hEx} makes \(H\) open normal. The Frattini construction gives open normal \(W_r\subseteq H\) eventually contained in each open normal \(U\supseteq\mathrm{PL}_\ell(H)\). Lift \(s_j\) to \(g_j\in G_L\) and write \(\rho:G_L\to\Gamma_0\) for the quotient. Chebotarev modulo \(\rho^{-1}(W_r)\cap\ker(\varepsilon_\ell\bmod\ell^r)\) gives the split places with Frobenius lifts \(\sigma_{r,j}\) such that
\[
 \sigma_{r,j}g_j\in\rho^{-1}(W_r)\cap\ker(\varepsilon_\ell\bmod\ell^r).
 \tag{6.56.3}
\]
Since \(\varepsilon_\ell(g_j)=1\), this gives (6.56.1), and the image of \(\sigma_{r,j}^{-1}\) gives \(\phi_{r,j}\) with \(h_r(\phi_{r,j})s_j^{-1}\in W_r\). Tame \(\ell\)-inertia gives the generators and (6.56.2). The kernel of the enlarged level is the closed normal closure of the new inertia groups, generated modulo \(\mathrm{PL}_\ell(h_r^{-1}H)\) by these tame generators.

A continuous \(\F_\ell\)-character of \(h_r^{-1}H\) is conjugation invariant there. At most \(q_0[\Gamma:H]\) values on conjugate tame generators determine its restriction to the kernel. Characters vanishing on that kernel descend to \(H\). The bound is
\[
 \#\operatorname{Hom}_{\mathrm{cont}}(H,\F_\ell)\,\ell^{q_0[\Gamma:H]}.
\]
\end{proof}

\begin{lemma}[Counting derivations in finite quotients]\label{B6.57}
Let \(\ell\ne2\), \(n\ge1\), and let \(O\) be a complete coefficient DVR with finite residue field \(\mathbf k\). Fix an adapted Taylor--Wiles system \(T\), an ordinary determinant pair \(u\) over a complete Noetherian local \(O\)-algebra \(B\) with local structure map surjective on residue fields, an ideal \(J\) and an integral point \(x:B\to O\). Take \(y\) as in \ref{B6.55}, with \(u\) specialising at \(x\) to its determinant law and ordered characters. Suppose:
\begin{enumerate}
\item The determinant coefficients and characters generate \(B\) adically, and \(J\subseteq\ker x\).
\item The Hecke datum has \(D\) weight parameters. At level \(r\), write \(Q_j=\mathrm Nw_{r,j}\) and \(u_{ji}\) for the splitting-algebra roots. Put
\[
 f=\prod_j\prod_{i\ne a}(u_{ji}-u_{ja})(u_{ji}-Q_j u_{ja}).
\]
Between the splitting algebra tensored with the base module and the level operator module there are additive maps in both directions intertwining root-polynomial actions, with both composites the action of \(f\). There is one \(E\ge0\) such that at every level \(f^E\) times each trace
\[
 \operatorname{tr}\!\left[
 \left(\prod_{a\ne i}(\phi_{r,j}-u_{ja})\right)^{E+1}(\tau_{r,j}-1-y_{ji})\right]
\]
acts by zero. The \(y_{ji}\) are the tame variables on the level module.
\item For the weight map \(\lambda_B:\Lambda\to B\) and Hecke module \(M\), some \(t\in B\) with \(x(t)\ne0\) satisfies \(taM\subseteq\ker(x\circ\lambda_B)M\) for every \(a\in\ker x\).
\item The representation \(y.\rho\) kills the CM part of \(\mathrm{PL}_\ell(H)\) and has distinct roots \(\alpha_{j,i}\in O\) at each \(s_j\).
\item For some \(C\), at each \(m\) and every sufficiently large auxiliary level the ordinary Selmer trace functions modulo \(\mathfrak m_O^m\) have cardinality at most
\[
 |\mathbf k|^{m(D+nq_0)+C}.
 \tag{6.57.1}
\]
\end{enumerate}
Then every \(O\)-linear \(x\)-derivation \(\delta:B\to O\) killing the selected \(D\) weight characters kills every universal trace.
\end{lemma}

\leanref{UpperTheory.\allowbreak{}V.\allowbreak{}Th.\allowbreak{}count}{Development.\allowbreak{}AutomorphicArithmetic.\allowbreak{}AutomorphicVarieties.\allowbreak{}TheoremsMuLinked.\allowbreak{}GaloisRepresentationTypes.\allowbreak{}MapRingHomSingleMulPartialCharpoly}{131--249}
\leanref{UpperTheory.\allowbreak{}V.\allowbreak{}Th.\allowbreak{}Count}{Development.\allowbreak{}AutomorphicArithmetic.\allowbreak{}AutomorphicVarieties.\allowbreak{}TheoremsMuLinked.\allowbreak{}OrdinaryHeckeAlgebraIntegral}{134--143}

\begin{proof}
Suppose \(d=\delta(\operatorname{tr}u(g_0))\ne0\), and set \(D_0=D+nq_0\). A fixed \(c_1\ne0\) takes trace derivatives modulo \(\mathfrak m_O^m\) to ordinary Selmer trace functions. When a derivation kills traces, a fixed \(c_2\ne0\) annihilates its ordinary-character derivatives.

The operator and annihilation hypotheses give \Lean{Split} with fixed nonzero \(c_3\in O\). For each \(m\) at large \(r\) it gives roots \(\alpha^{(r,m)}_{j,a}\in O\) and derivations \(e_b\) with values \(c_3\delta_{ab}\bmod\mathfrak m_O^m\) on the \(D_0\) weight and auxiliary tests indexed by \(a\). These diagonal values may be nonunits. The auxiliary test sets \(y_{ji}=0\) and is the determinant-law trace of
\[
 \left(\prod_{a\ne i}(\phi_{r,j}-\alpha^{(r,m)}_{j,a})\right)^{E+1}(\tau_{r,j}-1).
 \tag{6.57.2}
\]
The pullback \(e_\ast\) of \(\delta\) kills this test because \(\tau_{r,j}\) maps to \(1\), and kills the weight-coordinate tests by the weight hypothesis.

For fixed \(m\) choose \(r\) large enough for the tests and (6.57.1). With \(I\) the test indices together with \(\ast\), the map
\[
 \Phi(\mathbf t)(g)=c_1\sum_{a\in I}t_a e_a(\operatorname{tr}u_r(g)),\qquad \mathbf t\in O^I
 \tag{6.57.3}
\]
has image in the finite Selmer trace set. If \(\Phi(\mathbf t)=0\), the combined derivation after multiplication by \(c_1\) kills traces, so multiplication by \(c_2\) kills the weight derivatives. The auxiliary derivatives vanish by trace linearity on the group algebra. The diagonal tests give \(c_1c_2c_3t_a\in\mathfrak m_O^m\) for \(a\ne\ast\). Evaluating at a lift of \(g_0\) and multiplying by \(c_2c_3\) gives the bound also for \(d t_\ast\). Therefore
\[
 ct_a\in\mathfrak m_O^m\quad(a\in I),\qquad c=c_1c_2c_3d\ne0.
 \tag{6.57.4}
\]
Put \(a_0=v_O(c)\), \(b=a_0D_0+C+1\), \(m=a_0+b\). By (6.57.4) representatives modulo \((\mathfrak m_O^b)^I\) have distinct images. Thus \(|\mathbf k|^{b(D_0+1)}\le\#\operatorname{im}\Phi\le|\mathbf k|^{mD_0+C}\), contrary to \(b(D_0+1)-mD_0-C=1\).
\end{proof}

\section{The algebraic length comparison}

Let \(A=\bigoplus_j A_j\) be graded over a domain \(O\) and \(x:A\to O\) an \(O\)-point. Put
\[
 \operatorname{sat}((\ker x)^2)=\{a\in A:ca\in(\ker x)^2\text{ for some }0\ne c\in O\}.
\]
Write \(C_j\) for the image of \(A_j\cap\ker x\) in the quotient by this saturated square. For equations \(F_i\) put
\[
 C_j(F)=C_j/(C_j\cap\langle[F_i]\rangle_O).
 \tag{6.58.0}
\]
The saturated quotient embeds the cotangent lattice into its scalar extension to \(\operatorname{Frac}O\). A two-sided point is nonzero on some element of each strict sign of weight. For \(K\)-points a scalar \(t\in K^\times\) multiplies a weight-\(j\) value by \(t^j\).

\begin{lemma}[From Selmer length to trace differences]\label{B6.58}
\emph{Framed presentation.}
Let \(O\) be complete Noetherian local with finite residue field \(\mathbf k\) of odd characteristic \(\ell\). Fix a compact totally disconnected group whose open subgroups have finitely many continuous additive \(\F_\ell\)-characters. Give each basis vector a weight in a finite free weight lattice. Fix a polarised residual representation with normalised lift \(r_0\) and framed local data. Each local condition holds at \(r_0\), is given by a closed ideal on the allowed lifts, is invariant under allowed basis changes and torus conjugation, and lifts along small extensions. For two allowed lifts with the same image under a small extension, if one satisfies the local condition, the other does exactly when they agree on the prescribed subgroup \(N_q\) of trivial residual action.

Let \(\operatorname{Ob}_j\) be the finite-dimensional quotient of locally constant weight-\(j\) residual adjoint \(2\)-cocycles with local trivialising \(1\)-cochains by framed \(2\)-coboundaries. A framed boundary is the differential of a global \(1\)-cochain with restrictions corrected by local coboundaries at the fixed frames and by local cocycles vanishing on \(N_q\) at the local conditions. Let \(B\) be a complete Noetherian local \(O\)-algebra with local structure map surjective on residue fields, and a presentation \(O[[T_i]]\twoheadrightarrow B\) of homogeneous kernel with finitely many weighted variables in \(\mathfrak m_B\). Fix allowed local lifts over \(B\) with entries of the prescribed weights. The universal framed algebra \(R\) has a homogeneous relative presentation over \(B\). Its new variables are entries of \(r-r_0\) or \(\alpha-1\) for the universal representation and frame matrices. In weight \(j\) it has \(\dim_{\mathbf k}H^1_{\mathrm{fr},j}\) variables and, allowing zero relations,
\[
 \dim_{\mathbf k}\operatorname{Ob}_j\text{ relations}.
 \tag{6.58.1}
\]
Here \(H^1_{\mathrm{fr},j}\) consists of locally constant weight-\(j\) residual adjoint cocycles with local framing vectors expressing the framed restrictions as coboundaries. The distinguished vector is zero and the cocycles kill the prescribed local kernels.

\emph{Cotangent comparison.}
Let \(O\) be a DVR with finite residue field, \(x(A_j)=0\) and \(F_i\in A_j\). Suppose an \(O\)-module \(M\) injects additively into \(O/(c)\)-valued \(x\)-derivations for some \(c\ne0\), each killing the \(F_i\) and lifting to an \(O\)-valued derivation supported in weight \(j\). Then
\[
 \operatorname{length}_O M\le\operatorname{length}_O C_j(F).
 \tag{6.58.2}
\]

\emph{Orbit estimate.}
Let \(O\) be a complete characteristic-zero DVR, \(K\) an algebraically closed extension of \(\operatorname{Frac}O\), and \(v\) a valuation extending the normalised valuation of \(O\). Suppose \(A_0\) is complete Noetherian local over \(O\) with local map surjective on residue fields, and \(A\) is generated over \(A_0\) by finitely many elements of weights \(-2,-1,1,2\). Let \(x\) kill the nonzero weights. Let \(I\) be generated by
\[
 F_1,\ldots,F_{r+1}\in A_{-1},\quad G_1,\ldots,G_s\in A_{-2},\quad
 E_1,\ldots,E_{r'}\in A_1,\quad E'_1,\ldots,E'_{s'}\in A_2,\quad
 Z_1,\ldots,Z_m\in A_0\cap\ker x .
\]
Impose the following hypotheses.
\begin{enumerate}
\item \(r\ge1\) and \(x\) has a graded formal chart over \(\operatorname{Frac}O\) with constant term \(x\). Its linear coefficients identify the scalar-extended cotangent space with coordinates
\[
 \begin{array}{c|ccccc}
 j&-2&-1&0&1&2\\\hline
 \dim T_j&s&r&m&r'+1&s'+1 .
 \end{array}
\]
There are no other weights. Each cotangent-lattice element of weight \(-1,-2,0\), respectively, has a nonzero \(O\)-multiple in the span of the classes of \(F_i,G_i,Z_i\), respectively.
\item The image of \(A_0\) in \(A/I\) is finite over \(O\).
\item If a point of \(A/I\) kills the nonzero weights and differs from \(x_K\), the fibre over its weight-zero restriction has points of at most one strict sign. A strict-sign point has a nonzero value of that sign and kills all weights of the opposite sign.
\item At every two-sided \(K\)-point of \(A/I\), any two tangent derivations are linearly dependent.
\item \(a\in A_0\) has \(x(a)=0\) and \(\xi(a)\ne0\) at each two-sided point of \(A/I\).
\end{enumerate}
There exists a finite set \(T\) of two-sided \(K\)-points of \(A/I\) meeting every two-sided \(K^\times\)-orbit. The length of \(C_{-1}(F)\) is finite, and every such \(T\) satisfies
\[
 \operatorname{length}_O C_{-1}(F)\le2\sum_{\xi\in T}v(\xi(a)).
 \tag{6.58.3}
\]
\end{lemma}

\leanref{UpperTheory.\allowbreak{}V.\allowbreak{}Alg.\allowbreak{}lemmaP}{Development.\allowbreak{}AutomorphicArithmetic.\allowbreak{}AutomorphicVarieties.\allowbreak{}CommutativeAlgebra.\allowbreak{}LemmaP}{256--425}
\leanref{UpperTheory.\allowbreak{}V.\allowbreak{}Alg.\allowbreak{}M.\allowbreak{}mprime\_\allowbreak{}holds}{Development.\allowbreak{}AutomorphicArithmetic.\allowbreak{}AutomorphicVarieties.\allowbreak{}CommutativeAlgebra.\allowbreak{}GradedModulesLength.\allowbreak{}Mprime}{38--81}
\leanref{UpperTheory.\allowbreak{}V.\allowbreak{}Alg.\allowbreak{}M.\allowbreak{}main\_\allowbreak{}holds}{Development.\allowbreak{}AutomorphicArithmetic.\allowbreak{}AutomorphicVarieties.\allowbreak{}CommutativeAlgebra.\allowbreak{}GradedModulesLength.\allowbreak{}MainCotSpan}{32--71}

\begin{proof}
The relative tangent space consists of the residual cocycles and framing vectors of the statement. Entry evaluations separate them. Choose a dual basis in each weight from these evaluations and lift it to entries of \(r-r_0,\alpha-1\) to obtain the relative power-series variables. Each entry has weight the difference of the two basis weights.

Let \(P_R\) be the power-series ring in the old and new variables, \(Q_R=\ker(P_R\to R)\) and \(Q_B^e\) the extended ideal of relations of \(B\). The relative minimal relation space is \(Q_R/(Q_B^e+\mathfrak m_{P_R}Q_R)\). Push the corresponding square-zero extension out along a homogeneous functional on this space to a one-dimensional kernel. Its obstruction records global multiplication defects and the discrepancies at the local lifts. The local small-extension hypothesis and the conditions on \(N_q\) allow local corrections. Vanishing of the obstruction gives a framed boundary and hence a section over \(B\). Relative relations have zero linear part in the new variables. Evaluation on the section forces the functional to vanish. Thus the dual relation space injects into the obstructions in each weight. Padding by zero relations gives (6.58.1).

An integral derivation at \(x\) kills \((\ker x)^2\), hence its saturation by torsion-freeness of \(O\). Its reduction modulo \(c\) kills the \(F_i\) and factors through \(C_j(F)\) in weight \(j\). This gives an additive injection into \(\operatorname{Hom}_O(C_j(F),O/(c))\), since the derivation is supported in weight \(j\). If \(C_j(F)=\bigoplus_iO/(\pi^{b_i})\) has finite length, its dual here has length \(\sum_i\min(b_i,v_O(c))\le\sum_i b_i\). Cardinalities over the finite residue field give (6.58.2). Infinite length on the right gives the extended-length inequality.

For a two-sided point some coefficients of weights \(-2\) and \(2\) are nonzero, by squaring a weight-\(\pm1\) coefficient if needed. Choose from finite generating lists for the two weight modules. Rescale a negative coefficient to one using a square root in \(K\) and invert the positive coefficient. For the finitely many restrictions to \(A_0\) these slices are of finite type. Fixing the negative value cuts the orbit tangent, so hypothesis 4 makes the slices zero-dimensional. They give finite orbit covers.

The chart and finite generation give a free \(O\)-lattice \(C_{-1}\) of rank \(r\). By elementary divisors, change the equation generators within their \(O\)-span and choose \(y_i\in A_{-1}\), with classes a basis, and \(b_i\ge0\) such that
\[
 F_i\equiv\pi^{b_i}y_i\pmod{\operatorname{sat}((\ker x)^2)}\quad(1\le i\le r),
 \qquad F_{r+1}\in\operatorname{sat}((\ker x)^2).
 \tag{6.58.4}
\]
Then
\[
 \operatorname{length}_O C_{-1}(F)=\sum_i b_i .
 \tag{6.58.5}
\]
Cut the chart by the \(Z_i\), which have full linear rank in weight zero, and call the resulting map \(\phi\). The section construction gives a weight-zero representative \(a'\equiv a\bmod I\) with \(x(a')=0\) and \(\phi(a')\ne0\). Use it for \(a\) in the calculation.

Let \(R\) be the triply graded Rees algebra over \(A_0\) with generating modules \(A_{-2}, A_2, A_0F_{r+1}^2+A_{-2}^2A_2\) in the three unit multidegrees. The sections in these degrees are respectively \((F_1^2,\ldots,F_r^2,G_1,\ldots,G_s)\), \((E_1^2,\ldots,E_{r'}^2,E'_1,\ldots,E'_{s'})\) and \(F_{r+1}^2\). Let \(\mathbf t\) be their list, \(J\) the coefficient kernel of \(\phi\) in \(R\), and \(M=R/J\). Write \(\chi_L(N)_d\) for the Koszul Euler characteristic of a list \(L\) on the degree-\(d\) complex. For a section \(s\) of degree \(e\) it obeys
\[
 \chi_{(s,L)}(N)_d=\chi_L(N/sN)_d-\chi_L(0:_N s)_{d-e}.
\]
For the empty list use degree-piece length. Admissibility as in \ref{B6.58a} gives finite terminal lengths in large multidegrees. By \ref{B6.58b}, \(\mathbf t\) is admissible on \(M/aM\), so \ref{B6.58a} gives
\[
 2^{r+r'}\sum_i b_i\le\chi_{\mathbf t}(M/aM)_d .
 \tag{6.58.6}
\]
The comparison of orders and the valuation calculation in \ref{B6.58b} give
\[
 \chi_{\mathbf t}(M/aM)_d
 \le\chi_{(a)}(M/\mathbf t M)_d
 \le2^{r+1+r'}\sum_{\xi\in T}v(\xi(a)).
 \tag{6.58.7}
\]
Divide by \(2^{r+r'}\) and use (6.58.5).

\leanref{UpperTheory.\allowbreak{}V.\allowbreak{}Alg.\allowbreak{}M.\allowbreak{}core\_\allowbreak{}holds}{Development.\allowbreak{}AutomorphicArithmetic.\allowbreak{}AutomorphicVarieties.\allowbreak{}CommutativeAlgebra.\allowbreak{}GradedModulesLength.\allowbreak{}AlgMemCoef}{92--144}
\end{proof}

\begin{lemma}[The lower intersection number]\label{B6.58a}
Use \(\phi\), the elementary divisors and \(R\) of \ref{B6.58}, with \(M=R/J\). A list \(L\) is admissible on \(N\) if \((\mathfrak m_{A_0}R\,R_{++})^k\subset\operatorname{Ann}N+(L)\) for some \(k\), where \(R_{++}=\prod_{i=0}^2(R_{e_i})\) is the product of the ideals generated in the three unit multidegrees. If \(\mathbf t\) is admissible on \(M/aM\), then (6.58.6) holds for all coordinatewise sufficiently large \(d\).
\end{lemma}
\begin{proof}
For a chart monomial, count the absolute weights of the negative variables with multiplicity, and likewise the positive ones. In degree \(d\) put \(u=2d_0+4d_2,\ v=2d_1+2d_2\). Negative degree is at least \(u\) in each chart coefficient of \(R_d\). The first two generating modules contribute bounds \(2,0\). In the third, \(F_{r+1}\) has no linear part by (6.58.4), so weight \(-1\) forces negative degree at least two. Both summands of the third module therefore contribute at least four. At negative degree \(u\) the weight forces positive degree \(v\).

Quotient \(R\) by the homogeneous ideal of coefficients with negative degree strictly above \(u\) in each multidegree. The lower bounds add, so this is an ideal. It contains \(J+(a)\), since a weight-zero series vanishing at the origin has negative degree at least one. The resulting \(P_d\) is the lattice of leading coefficients of bidegree \((u,v)\) with \(A_0\) acting by \(x\). For an absolute-weight-one coordinate the chart provides a homogeneous lift with that coordinate at its least bidegree up to scalar. For absolute weight two, subtract lifts of the products of two same-sign weight-one coordinates at the least bidegree and clear denominators. Every monomial of bidegree \((u,v)\) is a product of \(d_0+2d_2\) blocks of weight \(-2\) and \(d_1+d_2\) of weight \(2\), by pairing the weight-one variables. The summand \(A_{-2}^2A_2\) in direction three supplies the products in \(R_d\). Thus \(P_d\) is a full, finite free lattice of rank \([X^uY^v]H\), with
\[
 H=\frac1{(1-X)^r(1-X^2)^s(1-Y)^{r'+1}(1-Y^2)^{s'+1}}.
\]

In \(P\), the first \(r\) sections equal \(\pi^{2b_i}y_i^2\) with the same degree tags, since the difference of squares in (6.58.4) has negative degree at least three. Replacing \(z\) by \(\pi^kz\) in a list \((z,L)\) adds \(k\chi_L(P/\pi P)_d\). This follows from the exact triangle for composition of the two multiplications and regularity of \(\pi\) on \(P\). For \(P/\pi P\) the alternating sum of Koszul terms applies the backward differences for the degrees of \(L\) to its Hilbert function. Removing the uniformiser powers one at a time gives the same differences for every \(i\), with generating series
\[
 H(1-X^2)^{r-1+s}(1-Y^2)^{r'+s'}(1-X^4Y^2)
 =(1+X)^{r-1}(1+Y)^{r'}\left(
 \frac1{(1-X)(1-Y)}+
 \frac{Y^2(1+X+X^2+X^3)}{(1-Y)(1-Y^2)}\right).
\]
At large \(u,v\) the second term has zero coefficient and the first gives \(2^{r-1+r'}\). The final list after removing the powers is admissible because its ideal contains that of the original list on \(P\). By eventual nonnegativity and additivity as in \ref{B6.58b}, first for this list and then for \(M/aM\twoheadrightarrow P\), the lower bound is \(\sum_i2b_i\,2^{r-1+r'}\).
\end{proof}
\leanref{UpperTheory.\allowbreak{}V.\allowbreak{}Alg.\allowbreak{}M.\allowbreak{}low\_\allowbreak{}holds}{Development.\allowbreak{}AutomorphicArithmetic.\allowbreak{}AutomorphicVarieties.\allowbreak{}CommutativeAlgebra.\allowbreak{}GradedModulesLength.\allowbreak{}AdmissibleMonoIdeal}{191--323}
\leanref{UpperTheory.\allowbreak{}V.\allowbreak{}Alg.\allowbreak{}M.\allowbreak{}hf\_\allowbreak{}holds}{Development.\allowbreak{}AutomorphicArithmetic.\allowbreak{}AutomorphicVarieties.\allowbreak{}CommutativeAlgebra.\allowbreak{}GradedModulesLength.\allowbreak{}ReesGradedPieceSurjection}{199--271}
\leanref{UpperTheory.\allowbreak{}V.\allowbreak{}Alg.\allowbreak{}M.\allowbreak{}num\_\allowbreak{}holds}{Development.\allowbreak{}AutomorphicArithmetic.\allowbreak{}AutomorphicVarieties.\allowbreak{}CommutativeAlgebra.\allowbreak{}GradedModulesLength.\allowbreak{}FrobeniusInclusionXEqMulQ}{206--284}

\begin{lemma}[The upper intersection number]\label{B6.58b}
Under the orbit hypotheses of \ref{B6.58}, \((\mathbf t,a)\) is admissible on \(M\), and (6.58.7) holds for every finite orbit cover \(T\) in sufficiently large multidegrees.
\end{lemma}
\begin{proof}
Put \(B=A_0,\ L=I\cap B,\ D=A_{-2}A_2\subset B\), and \(W_d=(M/\mathbf t M)_d\). Write \(T_i\) for the degree tags. For a prime \(P\supset J+(\mathbf t)\) avoiding \(R_{++}\), choose \(fT_0,gT_1,hT_2\notin P\) in the three generating modules. The identity
\[
 (F_{r+1}^2T_0)(hT_2)=(hT_0)(F_{r+1}^2T_2)
\]
puts \(F_{r+1}^2T_0\) in \(P\) along with all the other \(F_i^2T_0\). To see \(L\subset P\cap B\), test a weight-zero multiple \(cz\) of a homogeneous equation. For \(z=G_i\), multiply by \((fT_0)(gT_1)\) to get \((G_iT_0)(cfgT_1)\). For \(z=F_i\) use \((cz)^2\) instead to get \((F_i^2T_0)(c^2fgT_1)\). Interchange the tags for the positive equations. The \(Z_i\) already belong to \(J\).

Suppose \(P\supset D\). Since \(B/L\) is finite over \(O\), \(P\cap B\) is either the closed point or the kernel of an \(O\)-map \(\sigma:B\to K\). In the latter case send \(u\in A_{-2n}\) to \(uT_0^n\bmod P\) for \(n\ge0\) in the even-weight subring, and kill positive weights. Grouping monomials in the generators gives \(A_{\pm2n}=A_{\pm2}^n\). For \(v\in A_{2k}\), \(0<k\le n\), the proposed image of \(vu\) times \((fT_0)^k\) equals \((vf^k)(uT_0^n)=0\bmod P\), since \(vf^k\in D^k\). This proves multiplicativity. The kernel contains the even equations and the squares of the odd ones. It misses \(f\). The full ring is integral over the even subring. Lying over and the Nullstellensatz on the fibre over \(\sigma\) with \(f\) inverted give a strictly negative point of \(A/I\). The construction with \(gT_1\) gives a strictly positive point.

Weight-zero projection followed by \(\sigma\) gives a fixed point in the same fibre. Every product of opposite nonzero weights of total weight zero has square in \(D\). Away from \(x_K\), the two strict signs contradict hypothesis 3. Over \(x_K\), a negative point kills the weight-\(-1\) part of the augmentation square because each product has a nonnegative-weight factor that vanishes. It kills the saturated square there as well. The rational spanning by the \(F_i\) kills \(A_{-1}\). In weight \(-2\) the extra splitting \((-1,-1)\) now also vanishes, and spanning by the \(G_i\) kills \(f\), a contradiction. Thus \(P\supset D\) forces \(P\cap B=\mathfrak m_B\).

Radicals and powers give \(L^nW_d=0\) and \(\mathfrak m_B^nW_d\subseteq DW_d\) for large \(d\). The corresponding powers of \(LR_{++}\) and \(\mathfrak m_BR_{++}\) lie in \(J+(\mathbf t)\) and \(J+(\mathbf t)+DR\), and generation in unit multidegrees gives the stated eventual containments. Therefore \(B/\operatorname{Ann}W_d\) is finite over \(O\). At \(\ker\sigma\supset D\) the uniformiser is invertible and Nakayama gives zero local length. Suppose \(\sigma(L)=0\) and \(D\not\subseteq\ker\sigma\). The weight-zero projection splits \(B/L\hookrightarrow A/I\) as modules, so the fibre over \(\sigma\) is nonempty. Since \(\sigma(fg)\ne0\) for some \(f\in A_{-2}\) and \(g\in A_2\), every point of the fibre is two-sided. Since \(a\) is nonzero at that restriction, every prime over \(J+(\mathbf t,a)\) contains \(R_{++}\) or \(\mathfrak m_B\). This proves admissibility.

At a two-sided restriction take \(f,g\) of weights \(-2,2\) with \(\sigma(fg)\ne0\) and invert \(fg\). At a point of the geometric generic fibre with \(f=1\), the orbit derivation sends \(f\) to \(-2\). By hypothesis 4 it spans the tangent space. Hence the slice has zero tangent space there and is reduced. Adjoin a square root of the unit \(f\). Rescaling identifies this cover with the slice times \(\mathbb G_m\), so the localised generic \(A/I\) is reduced and its finite weight-zero algebra at \(\ker\sigma\) is a field. For any weight \(w\) take a unit \(e\) of weight \(-2w\) using \(f,g\). For a nonzero weight-\(w\) element \(h_0\), \(h_0^2e\) is a nonzero scalar and every \(h\) of that weight satisfies
\[
 h=(hh_0e)(h_0^2e)^{-1}h_0 .
\]
Each weight piece has local length at most one.

Let \(I_2\) replace each \(F_i,E_i\) in \(I\) by its square. On replacing \(z\) by \(z^2\), the kernel of the map to the old quotient is an image of that quotient shifted by \(\deg z\). Each weight of \(A/I_2\) consequently has local length at most \(2^{r+1+r'}\). After inverting \(fg\), the three coefficient modules are generated by \(f,g,f\). In the third module use \(f^2g\). For \(d_i\ge1\) each section imposes its coefficient times the full complementary even-weight piece. Therefore localised \(W_d\) is a quotient of the weight \(2(d_1-d_0-d_2)\) piece of \(A/I_2\).

For \(m_\sigma=\operatorname{length}_{B_{\ker\sigma}}(W_d)_{\ker\sigma}\), a prime filtration gives
\[
 \chi_{(a)}(W_d)=\sum_{\sigma\in T|_B}m_\sigma v(\sigma(a)),
\]
counting each restriction once. On a vertical factor kernel and cokernel have equal length. A horizontal domain factor is finite free over \(O\), and its cokernel length is the valuation of the determinant of multiplication by \(a\), the product of its values over all fraction-field embeddings. Lengths over \(B\) and \(O\) agree since their residue fields agree. The valuations are nonnegative by integrality, so the upper bound of (6.58.7) follows.

Multiplication by \(a\) on \(M\) is injective because \(J\) is the coefficient kernel and \(\phi(a)\ne0\) in a domain. With \(a\) first the Euler characteristic is \(\chi_{\mathbf t}(M/aM)_d\). With \(a\) last the Euler recursion subtracts at each step the Euler characteristic on the shifted annihilator of the removed section. These terms are eventually nonnegative. To see this, additivity reduces an admissible list on a finite graded module to homogeneous prime factors. Factors supported on the irrelevant locus vanish in high degrees. On a domain, a nonzero section is regular. For a zero section the recursion gives the difference for a degree shift. Multiplying homogeneous generators outside the prime gives a regular section of the same degree, and its exact sequence identifies that difference with the Euler characteristic on its quotient. Induction on the list ends in ordinary length. Dropping the subtracted terms gives the first inequality of (6.58.7), since permutation preserves the Koszul complex.
\end{proof}
\leanref{UpperTheory.\allowbreak{}V.\allowbreak{}Alg.\allowbreak{}M.\allowbreak{}zii\_\allowbreak{}holds}{Development.\allowbreak{}AutomorphicArithmetic.\allowbreak{}AutomorphicVarieties.\allowbreak{}CommutativeAlgebra.\allowbreak{}GradedModulesLength.\allowbreak{}SameDegreeZeroOfRestriction}{15--43}
\leanref{UpperTheory.\allowbreak{}V.\allowbreak{}Alg.\allowbreak{}M.\allowbreak{}up\_\allowbreak{}holds}{Development.\allowbreak{}AutomorphicArithmetic.\allowbreak{}AutomorphicVarieties.\allowbreak{}CommutativeAlgebra.\allowbreak{}GradedModulesLength.\allowbreak{}TwoSigned}{17--96}

\section{Residual vanishing from a Bernoulli unit}

For \(\gamma:G_\Q\to\mathbf k^\times\) and a cocycle \(z\) with coefficients in \(\mathbf k(\gamma)\), write \(\operatorname{UnrEv}_\ell(z)\) if for each place \(A\) there is a subgroup \(H\le G_\Q\) such that, with \(J_A=I_A\cap\ker\gamma\),
\[
 [J_A:J_A\cap H]\text{ is finite and prime to }\ell,\qquad z|_{J_A\cap H}=0.
 \tag{6.59.0}
\]

\begin{lemma}[Residual Bernoulli-unit criterion]\label{B6.59}
Let \(\ell\) be an odd prime and \(O\subset\overline\Q_\ell\) a DVR finite over \(\Z_\ell\). Fix a ring homomorphism \(\operatorname{red}:O\to\mathbf k\) to a discrete field of characteristic \(\ell\). Let \(N\ge1\), and let \(\vartheta:(\Z/N)^\times\to O^\times\) be nontrivial of order prime to \(\ell\) with reduction \(\bar\vartheta\) on \(G_\Q\). For a nontrivial character \(\lambda\) displayed at level \(L\), put \(B_{1,\lambda}=L^{-1}\sum_{t\in(\Z/L)^\times}[t]_L\lambda(t)\), with least nonnegative residues. Assume
\[
 \bar\vartheta\ne\bar\varepsilon_\ell,\qquad |B_{1,\vartheta^{-1}}|_\ell=1 .
 \tag{6.59.1}
\]
Every continuous \(1\)-cocycle \(z:G_\Q\to\mathbf k(\bar\vartheta)\) satisfying \(\operatorname{UnrEv}_\ell(z)\) has the form
\[
 z(g)=(\bar\vartheta(g)-1)b\quad(g\in G_\Q)
 \tag{6.59.2}
\]
for some \(b\in\mathbf k\).

For the product assertion let \(d\ge1\), \(f\) odd and \(\phi\bmod f,\xi\bmod 2^n\) be primitive \(O\)-valued characters. Suppose
\[
 4\mid\ell^d-1,\qquad \operatorname{ord}(\phi)\mid\ell^d-1,\qquad
 \phi(-1)\xi(-1)=-1,\qquad f4^{v_2(\ell^d-1)}\le2^n.
 \tag{6.59.3}
\]
Then \(|B_{1,\phi\xi}|_\ell=1\) by \ref{B6.59a}. Every continuous cocycle satisfying (6.59.0) for either \(\gamma=\overline{\phi\xi}\) or \(\gamma=(\overline{\phi\xi})^{-1}\) has the form \(z(g)=(\gamma(g)-1)b\).
\end{lemma}

\leanref{UpperTheory.\allowbreak{}V.\allowbreak{}X.\allowbreak{}residual\_\allowbreak{}bernoulli\_\allowbreak{}nonvanishing}{Development.\allowbreak{}AutomorphicArithmetic.\allowbreak{}AutomorphicVarieties.\allowbreak{}LevelCohomology.\allowbreak{}ResidualBernoulliNonvanishing}{174--264}
\leanref{UpperTheory.\allowbreak{}V.\allowbreak{}X.\allowbreak{}corX}{Development.\allowbreak{}AutomorphicArithmetic.\allowbreak{}AutomorphicVarieties.\allowbreak{}LevelCohomology.\allowbreak{}ChrInvNeg}{197--211}

\begin{proof}
Put \(\gamma=\bar\vartheta\) and \(F=\Q(\mu_N)\). On \(J_A\) the cocycle is additive, with image killed by \(\ell\) and order dividing the prime-to-\(\ell\) index in (6.59.0). Its restriction to \(G_F\subseteq\ker\gamma\) is therefore an unramified homomorphism.

Choose lifts \(\sigma_b\in G_\Q\) of \(b\in(\Z/N)^\times\). Prime-to-\(\ell\) roots of unity reduce injectively, so \(N>2\) and there is \(b_0\) with
\[
 \gamma(\sigma_{b_0})\ne1 .
 \tag{6.59.4}
\]
Choose odd positive \(a\) prime to \(N\) with
\[
 \operatorname{red}(\vartheta(a))\ne\operatorname{red}(a).
 \tag{6.59.5}
\]
For \(\ell\nmid N\) take \(a=\ell\). For \(\ell\mid N\), (6.59.1) supplies the residue class modulo \(N\), which has an odd representative. Write \(y=[a^{-1}b]_N\). The integral smoothing coefficients satisfy
\[
 E_1(N,a,b)=\frac{a-1}{2}-\frac{ay-[b]_N}{N},\qquad
 \sum_b E_1(N,a,b)\vartheta^{-1}(b)=(1-a\vartheta^{-1}(a))B_{1,\vartheta^{-1}}.
 \tag{6.59.6}
\]
Substitute \(b\equiv ay\bmod N\) in the terms with \(y\) and use that \(\vartheta^{-1}\) sums to zero. Both factors on the right are units by (6.59.5) and (6.59.1).

For \(u\in G_F\) the cocycle identity gives
\[
 z(\sigma_b^{-1}u\sigma_b)=\bar\vartheta(b)^{-1}z(u).
 \tag{6.59.7}
\]
Apply the cyclotomic annihilation of \ref{B6.10a} with \(K=F\), \(X=(\mathbf k,+)\) and \(h=z|_{G_F}\). This is the integral Stickelberger relation on unramified homomorphisms \cite[Theorem~6.10]{washington-1997}. It gives \(\sum_b E_1(N,a,b)z(\sigma_b^{-1}u\sigma_b)=0\), which by (6.59.7) is the reduction of the sum in (6.59.6) times \(z(u)\). Thus \(z|_{G_F}=0\).

\leanref{UpperTheory.\allowbreak{}V.\allowbreak{}Odd.\allowbreak{}homKill}{Development.\allowbreak{}AutomorphicArithmetic.\allowbreak{}AutomorphicVarieties.\allowbreak{}OddResidueBounds.\allowbreak{}HomKill}{9--50}

The commutator of \(g\) and \(\sigma_{b_0}\) lies in \(G_F\). Comparing \(z(g\sigma_{b_0})\) and \(z(\sigma_{b_0}g)\) gives \((\gamma(\sigma_{b_0})-1)z(g)=(\gamma(g)-1)z(\sigma_{b_0})\). Divide using (6.59.4).

Inversion preserves the conductors and the product sign. It also preserves the order of \(\phi\). Use \ref{B6.59a} on \((\phi^{-1},\xi^{-1})\) and on \((\phi,\xi)\) to obtain the units in (6.59.1) for \(\vartheta=\phi\xi\) and \(\vartheta=(\phi\xi)^{-1}\), respectively. Both products have order prime to \(\ell\). Choose a Galois element acting by \(1+2^{n-1}\) on \(\mu_{2^n}\) and trivially on \(\mu_{\ell f}\). Primitivity gives value \(-1\) for both products there, whereas \(\bar\varepsilon_\ell\) is \(1\). Hence (6.59.2) applies to both reductions.
\end{proof}

\begin{lemma}[An effective dyadic Bernoulli unit]\label{B6.59a}
Let \(\ell\) be prime, \(d\ge1\), \(f\ge1\) odd and \(n\ge0\). Let \(\phi,\xi\) be primitive Dirichlet characters over \(\overline\Q_\ell\) of levels \(f,2^n\). Put \(Q=\ell^d\) and \(e=v_2(Q-1)\). If \(4\mid Q-1\), \(\operatorname{ord}\phi\mid Q-1\), \(\phi(-1)\xi(-1)=-1\), and \(f4^e\le2^n\), then \(|B_{1,\phi\xi}|_\ell=1\) at level \(f2^n\).
\end{lemma}
\begin{proof}
Put \(M=2^n\), \(m=n-e\), \(s=f2^e\), so \(m\ge e\ge2\) and \(1<s\le2^m\). Extend characters by zero off the units. For \(u\bmod M\) define
\[
 A_u=\sum_{\substack{0\le j<fM\\j\equiv u\bmod M}}j\phi(j),\qquad
 g(u)=\frac{A_u-A_0}{fM},\qquad
 c=\frac1M\sum_{j\bmod f}\phi(j).
\]
In \(A_{u-ft}\) translate \(j\) by \(ft\). In each full period of length \(fM\), a class modulo \(M\) visits every class modulo \(f\) once. Translation contributes \(fM\phi(j)\) from each \(0\le j<ft\) in class \(u\), with repeated periods counted with multiplicity. Thus for every integer \(t\ge0\),
\[
 g(u)-g(u-ft)=tc-\sum_{\substack{0\le j<ft\\j\equiv u\bmod M}}\phi(j).
\]
The values of \(\phi\) and \(M^{-1}\) are integral and their reductions are fixed by the \(Q\)-power map. For each \(u\) choose \(t\ge0\) with \(ft\equiv u\bmod M\). Since \(g(0)=0\), all \(g(u)\) are integral with reductions fixed by that map. This also holds when \(\ell\mid f\). Oddness makes \(\phi\xi\) nontrivial and \(\sum_u\xi(u)=0\) by primitivity at \(M\ge16\). Therefore
\[
 B_{1,\phi\xi}=\frac1{fM}\sum_{0\le j<fM}j\phi(j)\xi(j)=\sum_u\xi(u)g(u).
\]

Reduce in the residue field and suppress the bars. If the sum vanishes, reindexing gives \(\sum_u\xi(u)g(uy)=0\) for every unit \(y\). Successive squaring gives \(v_2(Q^{2^j}-1)=e+j\). Thus the powers of \(Q\) modulo \(M\) are \(1+2^e h\) for \(0\le h<2^m\). Average the corresponding Frobenius conjugates. The geometric sum in \(\xi(u)^{2^e}\) is \(2^m\) when \(\xi(u)^{2^e}=1\) and zero otherwise. Since \(\ell\) is odd, it follows that
\[
 \sum_{u\in U}\xi(u)g(uy)=0,\qquad
 U=\{u\in(\Z/M)^\times:\xi(u)^{2^e}=1\}.
\]
Set \(h_0=1+M/2\). It generates the kernel of reduction to level \(M/2\), so \(\xi(h_0)=-1\) by primitivity, also after reduction. Hence \(\xi\) is faithful on the cyclic subgroup of units \(1\bmod4\), whose unique involution is \(h_0\). For \(u\in U\), this gives \(u^{2^e}=1\bmod M\). In \((u-1)(u+1)\prod_{i=1}^{e-1}(u^{2^i}+1)\), the last \(e-1\) factors and one of \(u-1,u+1\) each have valuation one. Thus \(u\equiv\delta_u\bmod 2^m\) for a unique \(\delta_u=\pm1\).

Subtract the restricted sums at \(y=1\) and \(y=1-s\). Since \(2^e\mid s\), \(u(1-s)\equiv u-\delta_us\bmod M\). For \(\delta_u=1\), the shift formula with \(t=2^e\) has a boundary term exactly when \([u]_M<s\), hence at \(u=1\) because \(s\le2^m\). For \(\delta_u=-1\) use the formula at \(u+s\) with opposite sign. The condition \([u+s]_M<s\) forces \([u+s]_M=s-1\), hence \(u=-1\). Both \(1,-1\) are in \(U\), and the difference gives
\[
 0=2^ec\sum_{u\in U}\delta_u\xi(u)-\phi(1)+\xi(-1)\phi(s-1).
\]
Multiplication by \(h_0\) preserves \(U\) and \(\delta_u\) and negates \(\xi(u)\). Thus the sum cancels. Since \(s\equiv0\bmod f\) and the product character is odd, the remaining terms add to \(-2\). This contradiction in odd characteristic proves that the integral Bernoulli value is a unit.
\end{proof}
\leanref{UpperTheory.\allowbreak{}V.\allowbreak{}X.\allowbreak{}lx\_\allowbreak{}partial}{Development.\allowbreak{}AutomorphicArithmetic.\allowbreak{}AutomorphicVarieties.\allowbreak{}LevelCohomology.\allowbreak{}LxPartial}{15--115}
\leanref{UpperTheory.\allowbreak{}V.\allowbreak{}X.\allowbreak{}lx\_\allowbreak{}core}{Development.\allowbreak{}AutomorphicArithmetic.\allowbreak{}AutomorphicVarieties.\allowbreak{}LevelCohomology.\allowbreak{}LxCore}{65--221}
\leanref{UpperTheory.\allowbreak{}V.\allowbreak{}X.\allowbreak{}lemmaX}{Development.\allowbreak{}AutomorphicArithmetic.\allowbreak{}AutomorphicVarieties.\allowbreak{}LevelCohomology.\allowbreak{}LemmaX}{78--303}

Count I combines \ref{B6.42}, \ref{B6.43} and the first endpoint of \ref{B6.52}. Counts II and III use \ref{B6.49} and \ref{B6.59} for the six residual characters of \ref{B6.31}. Count V uses \ref{B6.51} and the product-character criterion of \ref{B6.59} in \ref{B6.34}. There the constructed product reduces to \(\eta_4\bar\Psi^2\), and its inverse gives \(\eta_4\bar\Psi^{-2}\).

\leanref{UpperTheory.\allowbreak{}V.\allowbreak{}Z.\allowbreak{}J.\allowbreak{}count\_\allowbreak{}i\_\allowbreak{}of\_\allowbreak{}printed\_\allowbreak{}bound\_\allowbreak{}k\_\allowbreak{}kato\_\allowbreak{}automorphy}{Development.\allowbreak{}AutomorphicArithmetic.\allowbreak{}AutomorphicVarieties.\allowbreak{}ModularityArgument.\allowbreak{}ThresholdArguments.\allowbreak{}DenseThresholdCases.\allowbreak{}LinkedUnlE}{49--54}
\leanref{UpperTheory.\allowbreak{}V.\allowbreak{}HV.\allowbreak{}DescentVanishingReductions.\allowbreak{}countII\_\allowbreak{}and\_\allowbreak{}countIII\_\allowbreak{}of\_\allowbreak{}printsH\_\allowbreak{}of\_\allowbreak{}sp4In}{Development.\allowbreak{}AutomorphicArithmetic.\allowbreak{}AutomorphicVarieties.\allowbreak{}ModularityArgument.\allowbreak{}ThresholdArguments.\allowbreak{}DenseThresholdCases.\allowbreak{}AdjointLiftingOfCohomology}{33--39}
\leanref{UpperTheory.\allowbreak{}V.\allowbreak{}Z.\allowbreak{}J.\allowbreak{}count\_\allowbreak{}v\_\allowbreak{}of\_\allowbreak{}printed\_\allowbreak{}bound\_\allowbreak{}k\_\allowbreak{}reduced\_\allowbreak{}unitary\_\allowbreak{}lift}{Development.\allowbreak{}AutomorphicArithmetic.\allowbreak{}AutomorphicVarieties.\allowbreak{}ModularityArgument.\allowbreak{}ThresholdArguments.\allowbreak{}DenseThresholdCases.\allowbreak{}LinkedUnlE}{55--61}


% ---- p6c6 ----
\chapter{The exponent 401 and the downstream window}\label{ch:B6-6}

\begin{lemma}[The rank cutoff]\label{B6.60}\label{B6.61}
For a family satisfying the page normalization, $e_P=12$, so
$c_{TO}=28e_P+41=377$ and $C_t=c_{TO}+24=401$.

Let $0<\varepsilon<1/52$ and $x\ge2^{100}$. At scale $x$, put
$T_O=\lceil x^{377}\rceil$. Let $R$ be a rank datum with rank parameters
$m,k,t\in\N$. Write $n=|A|$, where $A$ is the finite abelian group of
its frame.
Put $f=2n$, $r=k+m-1$, and $L=x^\varepsilon$. Thus $m\ge2$ and
$m+3\le k$.
Define
\[
 n_W=8\left\lceil\frac{10r}{8}
 \left(f+\frac{f^2(k-1)^3L^2}{m}\right)\right\rceil .
\]
Assume
\[
\begin{aligned}
 x^{2+8\varepsilon}&\le n\le2x^{2+8\varepsilon},\\
 x^{1/2-\varepsilon/16}&\le m\le2x^{1/2-\varepsilon/16},\\
 x^{9/2+10\varepsilon}&\le k\le2x^{9/2+10\varepsilon},
 \qquad r+t=8fr+4+T_On_W.
\end{aligned}
\]
Then $t\le x^{400}$.
\end{lemma}
\leanref{Formalism.\allowbreak{}MainNotions.\allowbreak{}PageFix}{Formalism.\allowbreak{}SharpRestriction.\allowbreak{}PageFix}{9--18}
\leanref{GlobalTheory.\allowbreak{}EX.\allowbreak{}TOv}{Development.\allowbreak{}AbelianVarietyTheory.\allowbreak{}Definitions.\allowbreak{}GlobalTwisting.\allowbreak{}PairRowIffQdRows}{17--17}
\leanref{GlobalTheory.\allowbreak{}EX.\allowbreak{}nWv}{Development.\allowbreak{}AbelianVarietyTheory.\allowbreak{}Definitions.\allowbreak{}GlobalTwisting.\allowbreak{}PairRowIffQdRows}{19--20}
\leanref{GlobalTheory.\allowbreak{}EX.\allowbreak{}RankVal}{Development.\allowbreak{}AbelianVarietyTheory.\allowbreak{}Definitions.\allowbreak{}GlobalTwisting.\allowbreak{}PairRowIffQdRows}{28--32}
\leanref{UpperTheory.\allowbreak{}Top.\allowbreak{}eventually\_\allowbreak{}rank\_\allowbreak{}le\_\allowbreak{}rpow\_\allowbreak{}of\_\allowbreak{}lt\_\allowbreak{}one\_\allowbreak{}div\_\allowbreak{}fiftytwo}{Development.\allowbreak{}AutomorphicArithmetic.\allowbreak{}CharacterBoxes.\allowbreak{}EventuallyRankLeRpowOfLtOneDivFiftytwo}{5--126}
\begin{proof}
Put $u=x^{1/100}\ge2$. The bounds on $\varepsilon$ give
\[
 f\le4u^{216},\quad k\le2u^{470},\quad r\le2k\le4u^{470},
 \quad m\ge u^{49},\quad L^2\le u^4.
\]
Indeed, $2+8\varepsilon\le216/100$,
$9/2+10\varepsilon\le470/100$,
$1/2-\varepsilon/16\ge49/100$, and $2\varepsilon\le4/100$.
Hence
\[
 \frac{f^2(k-1)^3L^2}{m}
 \le128u^{2\cdot216+3\cdot470+4-49}=128u^{1797},
 \qquad
 f+\frac{f^2(k-1)^3L^2}{m}\le132u^{1797}.
\]
Using $\lceil v\rceil<v+1$ in the definition of $n_W$ gives
\[
 n_W\le5280u^{2267}+8\le5288u^{2267},
 \qquad T_O\le2u^{37700},
 \qquad T_On_W\le10576u^{39967}.
\]
The rank identity now yields
\[
 t\le128u^{686}+4+10576u^{39967}
 \le10708u^{39967}\le u^{40000}=x^{400},
\]
where the last inequality uses $10708\le2^{33}\le u^{33}$.
\end{proof}

\begin{theorem}[Windows at arbitrarily large scales]\label{B6.62}
Fix a primitive character $\psi$, $0<\varepsilon<1/52$, and $c\ge1$.
Assume $\mathrm{VWinAt}(\psi,401)$ of \ref{B6.03} and the
prime-pair existence assertion of \ref{B2.43}.
Then there is $J\ge1$ such that, for every $h\in\R$, there is $Z\ge h$
with the following property. For every $x\in[Z,Z^{3/2}]$ there exist
$U\in\R$, primes $q,d$, and $a\in\N$ such that
\[
 x^{1/J}\le U\le x,\qquad
 \operatorname{PairGood}(\operatorname{cond}\psi,\varepsilon,U,x,q,d),
 \qquad
 x^{2+8\varepsilon}\le c2^a\le2x^{2+8\varepsilon}.
\]
Here $\operatorname{PairGood}$ is the condition of \ref{B2.42}.
For these $q,d,a$, if a number field $E$ satisfies
$\mathrm{IsChiField}(\psi,a,q,d,E)$ as in Definition~\ref{B6.02}, then
\[
 \mathrm{VVle}\bigl(E,\lfloor x^{400}\rfloor,\exp(x^\varepsilon/2)\bigr)
\]
holds as in Definition~\ref{B6.01}.
\end{theorem}
\leanref{UpperTheory.\allowbreak{}Top.\allowbreak{}WinDatum}{Development.\allowbreak{}AutomorphicArithmetic.\allowbreak{}CharacterBoxes.\allowbreak{}TwistForm}{19--26}
\leanref{UpperTheory.\allowbreak{}Top.\allowbreak{}winPairs}{Development.\allowbreak{}AutomorphicArithmetic.\allowbreak{}CharacterBoxes.\allowbreak{}WinPairs}{5--81}
\begin{proof}
Apply Theorem~\ref{B2.48} with $M=\operatorname{cond}\psi$. Its witnesses
$Y,q,d,e$ give $U,q,d,a$ here. Put $\lambda=c^{1/(2+8\varepsilon)}$.
In that proof the upstream window starts at $W\ge\lambda^{400}$, and
$Z=\lambda W$. For each $x\in[Z,Z^{3/2}]$, the corresponding point is
$v=x/\lambda\ge W$. Thus
\[
 x^{400}=\lambda^{400}v^{400}\le v^{401}.
\]
\end{proof}

In Theorem~\ref{B2.59}, the rank $t$ is a natural number, so
Lemma~\ref{B6.61} gives $t\le\lfloor x^{400}\rfloor$ at the scales used
there. The window estimate therefore covers every required index
$1\le j<t$.

\part{Input one: Green}

% ---- p7c1 ----
\chapter{Green's classification and its use}\label{ch:green}
\begingroup
\footnotesize
\setlength{\topsep}{1pt plus 1pt}
\setlength{\abovedisplayskip}{4pt plus 1pt}
\setlength{\belowdisplayskip}{4pt plus 1pt}
\setlength{\abovedisplayshortskip}{2pt plus 1pt}
\setlength{\belowdisplayshortskip}{3pt plus 1pt}

\section{Green's theorem}
For a prime $p$ put $G_d=\GL_d(\F_p)$, $N_d=p^d-1$, and $\Delta_d=\prod_{1\le j<d}(p^j-1)$. Let $\mu_N$ denote the $N$th roots of unity.

\begin{definition}[Cuspidality]\label{B7.01}
A cuspidal of $G_d$ ($d>0$) is an irreducible character with mean zero on the standard unipotent radicals of types $(a,d-a)$, $0<a<d$.
\end{definition}
\leanref{GlobalTheory.\allowbreak{}GLq.\allowbreak{}IsCusp}{Development.\allowbreak{}AbelianVarietyTheory.\allowbreak{}Definitions.\allowbreak{}NonGaloisDescent.\allowbreak{}LabelDatum}{64--66}

\begin{definition}[Parameters]\label{B7.02}
Let $\mathcal P_d$ consist of monic irreducibles $P$ of degree $d$ over $\F_p$ with $P(0)\ne0$. For $[K:\F_p]=d$, a character $\theta\in\widehat{K^\times}$ is regular if the $\theta^{p^i}$, $0\le i<d$, are distinct.
\end{definition}
\leanref{GlobalTheory.\allowbreak{}GLq.\allowbreak{}IsPrimeLab}{Development.\allowbreak{}AbelianVarietyTheory.\allowbreak{}Definitions.\allowbreak{}NonGaloisDescent.\allowbreak{}LabelDatum}{91--92}
\leanref{GlobalTheory.\allowbreak{}Published.\allowbreak{}IsRegularChar}{Development.\allowbreak{}AbelianVarietyTheory.\allowbreak{}PublishedTheorems.\allowbreak{}GaloisCohomology₂.\allowbreak{}GreenRegularCharacters}{10--12}

\begin{definition}[Label datum]\label{B7.03}
A label datum $L$ consists of a unital multiplicative $t:\overline{\F}_p\to\C$ injective off zero and generators $u_d$ of $\mu_{N_d}(\overline{\F}_p)$ with $u_{de}^{N_{de}/N_d}=u_d$ for $d,e>0$. Set $\Theta_d(x,u_d^a)=t(x)^a$ for $x\in\mu_{N_d}$, with value zero off $\langle u_d\rangle$ in the second argument.
\end{definition}
\leanref{GlobalTheory.\allowbreak{}GLq.\allowbreak{}LabelDatum}{Development.\allowbreak{}AbelianVarietyTheory.\allowbreak{}Definitions.\allowbreak{}NonGaloisDescent.\allowbreak{}LabelDatum}{81--87}
\leanref{GlobalTheory.\allowbreak{}GLq.\allowbreak{}thetaOf}{Development.\allowbreak{}AbelianVarietyTheory.\allowbreak{}Definitions.\allowbreak{}NonGaloisDescent.\allowbreak{}LabelDatum}{88--90}

\begin{definition}[Prime label]\label{B7.04}
An elliptic $y\in G_d$ has irreducible characteristic polynomial. A cuspidal $\chi$ has label $P\in\mathcal P_d$ if for every root $x$ of $P$, every elliptic $y$, and every characteristic root $z$ of $y$,
\begin{equation}\label{eq:green-label}
 \chi(y)=(-1)^{d-1}\sum_{i<d}\Theta_d(x,z^{p^i}).
\end{equation}
\end{definition}
\leanref{GlobalTheory.\allowbreak{}GLq.\allowbreak{}IsCuspOf}{Development.\allowbreak{}AbelianVarietyTheory.\allowbreak{}Definitions.\allowbreak{}NonGaloisDescent.\allowbreak{}LabelDatum}{93--97}

\begin{definition}[Regular form]\label{B7.05}
For every prime $p$ and $[K:\F_p]=d>0$, a family $\tau_\theta$ on $\widehat{K^\times}$ induces a bijection from regular Frobenius orbits to cuspidals. For regular $\theta$, $\tau_\theta(1)=\Delta_d$ and $\tau_\theta(y)=(-1)^{d-1}\sum_{i<d}\theta(z^{p^i})$ for elliptic $y$ and each characteristic root $z\in K^\times$.
\end{definition}
\leanref{GlobalTheory.\allowbreak{}Published.\allowbreak{}GreenRegularCharacters}{Development.\allowbreak{}AbelianVarietyTheory.\allowbreak{}PublishedTheorems.\allowbreak{}GaloisCohomology₂.\allowbreak{}GreenRegularCharacters}{13--21}

\begin{definition}[Polynomial form]\label{B7.06}
Green asserts that for every prime $p$, datum $L$ and $d>0$, a family $\tau_P$ on $\F_p[X]$ restricts to a bijection from $\mathcal P_d$ to cuspidals. For $P\in\mathcal P_d$, $\tau_P$ has label $P$ and degree $\Delta_d$.
\end{definition}
\leanref{GlobalTheory.\allowbreak{}GLq.\allowbreak{}Green}{Development.\allowbreak{}AbelianVarietyTheory.\allowbreak{}Definitions.\allowbreak{}NonGaloisDescent.\allowbreak{}LabelDatum}{108--113}

\begingroup
\linespread{0.88}\selectfont
\setlength{\topsep}{0pt}
\subsection{Construction of the regular family}
Let $V$ have dimension $n$ over a finite field $F$ of size $v$. Put $H=\GL(V)$, $c_g=\det(XI-g)$, and $I(\phi,\psi)=\sum_{g\in H}\phi(g)\overline{\psi(g)}$.

\begin{lemma}[Frobenius]\label{B7.07}
For irreducible $P\in\F_p[X]$ with root $x$, the roots are $(x^{p^i})_{0\le i<\deg P}$, and $x^{p^j}=x$ iff $\deg P\mid j$ for $j\ge0$. This is the prime-field orbit theorem \cite[Ch.~2]{lidl-niederreiter-1997}.
\end{lemma}
\leanref{GlobalTheory.\allowbreak{}GreenExists.\allowbreak{}pow\_\allowbreak{}eq\_\allowbreak{}self\_\allowbreak{}iff}{Development.\allowbreak{}AbelianVarietyTheory.\allowbreak{}Proofs.\allowbreak{}InertiaEigenvalues.\allowbreak{}CuspidalCharExistsOfGreen}{60--67}
\leanref{GlobalTheory.\allowbreak{}GreenExists.\allowbreak{}rootsBar\_\allowbreak{}eq}{Development.\allowbreak{}AbelianVarietyTheory.\allowbreak{}Proofs.\allowbreak{}InertiaEigenvalues.\allowbreak{}CuspidalCharExistsOfGreen}{68--95}

\begin{lemma}[Lift]\label{B7.08}
There is a unital multiplicative $t:\overline F\to\C$ injective off zero.
\end{lemma}
\leanref{GlobalTheory.\allowbreak{}CuspGL.\allowbreak{}TeichEmb.\allowbreak{}exists\_\allowbreak{}teichmuller}{Development.\allowbreak{}AbelianVarietyTheory.\allowbreak{}Proofs.\allowbreak{}CharacterTheory.\allowbreak{}ExistsTeichmullerMapAlgebraicClosureFiniteField}{24--77}
Reduction distinguishes the Teichmüller lifts. Embed the algebraic part of the Witt fraction field, which contains the nonzero lifts as roots of unity, into $\C$.

\begin{definition}[Root sums]\label{B7.09}
For $h\ge0$, set $s_h(P)=\sum_{P(\alpha)=0}t(\alpha)^h$, counting multiplicities and putting $s_h(0)=0$. Then $s_h$ is additive on monic products and $s_h((X-1)^j)=j$.
\end{definition}
\leanref{GlobalTheory.\allowbreak{}CuspGL.\allowbreak{}rootSum}{Development.\allowbreak{}AbelianVarietyTheory.\allowbreak{}Definitions.\allowbreak{}CommutativeAlgebra.\allowbreak{}InvariantSubspaceMobiusInverseRecEqOfFinrankLt}{65--67}
\leanref{GlobalTheory.\allowbreak{}CuspGL.\allowbreak{}rootSum\_\allowbreak{}isAdditive\_\allowbreak{}and\_\allowbreak{}pow\_\allowbreak{}X\_\allowbreak{}sub\_\allowbreak{}one}{Development.\allowbreak{}AbelianVarietyTheory.\allowbreak{}Proofs.\allowbreak{}CharacterTheory.\allowbreak{}RootPowerSumCharacterAdditiveAndXSubOnePow}{9--21}

\begin{definition}[Subspace inversion]\label{B7.10}
For $s:F[X]\to\C$, set $\Psi_s(V,g)=s(c_g)-\sum_{W\subsetneq V,\ gW\subseteq W}\Psi_s(W,g|_W)$. Write $\Psi_h=\Psi_{s_h}$ and $R_h=s_h\circ c$.
\end{definition}
\leanref{GlobalTheory.\allowbreak{}CuspGL.\allowbreak{}psi}{Development.\allowbreak{}AbelianVarietyTheory.\allowbreak{}Definitions.\allowbreak{}CommutativeAlgebra.\allowbreak{}InvariantSubspaceMobiusInverseRecEqOfFinrankLt}{35--38}
\leanref{GlobalTheory.\allowbreak{}CuspGL.\allowbreak{}psi\_\allowbreak{}eq}{Development.\allowbreak{}AbelianVarietyTheory.\allowbreak{}Definitions.\allowbreak{}CommutativeAlgebra.\allowbreak{}InvariantSubspaceMobiusInverseRecEqOfFinrankLt}{55--64}

\begin{lemma}[Primitivity]\label{B7.11}
The recursion is conjugacy invariant. For additive $s$, it satisfies the primitivity identity $\sum_{h:V/A\to A}\Psi_s(g+h)=0$ for $gA\subseteq A$ and $0<A<V$. View the linear maps $h$ as endomorphisms via quotient and inclusion.
\end{lemma}
\leanref{GlobalTheory.\allowbreak{}CuspGL.\allowbreak{}psi\_\allowbreak{}conj}{Development.\allowbreak{}AbelianVarietyTheory.\allowbreak{}Proofs.\allowbreak{}CharacterTheory.\allowbreak{}InvariantSubspaceMobiusInverseCharpolyLinearEqui}{43--52}
\leanref{GlobalTheory.\allowbreak{}CuspGL.\allowbreak{}isPrimitive\_\allowbreak{}psi}{Development.\allowbreak{}AbelianVarietyTheory.\allowbreak{}Proofs.\allowbreak{}CharacterTheory.\allowbreak{}IsPrimitiveEndClassFunctionInvariantSubspaceMobi}{188--200}

\begin{lemma}[Values]\label{B7.12}
If $n>0$ and $s((X-1)^j)=j$ for all $j\ge0$, then $\Psi_s(1)=\prod_{1\le j<n}(1-v^j)$. For additive $s$, irreducible $c_g$ gives $\Psi_s(g)=s(c_g)$.
\end{lemma}
\leanref{GlobalTheory.\allowbreak{}CuspGL.\allowbreak{}psi\_\allowbreak{}one}{Development.\allowbreak{}AbelianVarietyTheory.\allowbreak{}Proofs.\allowbreak{}CharacterTheory.\allowbreak{}InvariantSubspaceMobiusInverseOneProdIco}{74--78}
\leanref{GlobalTheory.\allowbreak{}CuspGL.\allowbreak{}psi\_\allowbreak{}of\_\allowbreak{}irreducible\_\allowbreak{}charpoly}{Development.\allowbreak{}AbelianVarietyTheory.\allowbreak{}Proofs.\allowbreak{}CharacterTheory.\allowbreak{}InvariantSubspaceMobiusInverseCharpolyOfIrreduci}{52--57}

\begin{lemma}[Virtual characters]\label{B7.13}
On $H$, $R_h$ is virtual by Brauer induction, and $\Psi_h$ by grouping the summands of \ref{B7.10} into parabolic inductions.
\end{lemma}
\leanref{GlobalTheory.\allowbreak{}CuspGL.\allowbreak{}rootSum\_\allowbreak{}charpoly\_\allowbreak{}mem\_\allowbreak{}virt}{Development.\allowbreak{}AbelianVarietyTheory.\allowbreak{}Proofs.\allowbreak{}CharacterTheory.\allowbreak{}RootPowerSumCharpolyMemVirtualCharactersGL}{213--234}
\leanref{GlobalTheory.\allowbreak{}CuspGL.\allowbreak{}psi\_\allowbreak{}rootSum\_\allowbreak{}charpoly\_\allowbreak{}mem\_\allowbreak{}virt}{Development.\allowbreak{}AbelianVarietyTheory.\allowbreak{}Proofs.\allowbreak{}CharacterTheory.\allowbreak{}IsCharacterCuspidalRootSumRegular}{55--89}

\begin{lemma}[Mass]\label{B7.14}
Write $M_\phi(P)=\sum_{c_g=P}\phi(g)$ over endomorphisms. For additive $s$ and monic $P$ of degree $n>0$, this mass for $\Psi_s$ is $|H|s(Q)/(v^n-1)$ when $P=Q^a$ for monic irreducible $Q$ and $a>0$, and zero otherwise.
\end{lemma}
\leanref{GlobalTheory.\allowbreak{}CuspGL.\allowbreak{}mass\_\allowbreak{}psi\_\allowbreak{}eq\_\allowbreak{}glCard\_\allowbreak{}mul\_\allowbreak{}mangoldt\_\allowbreak{}div}{Development.\allowbreak{}AbelianVarietyTheory.\allowbreak{}Proofs.\allowbreak{}CharacterTheory.\allowbreak{}EndomorphismMassMobiusInverseMonicEq}{27--53}

\begin{lemma}[Torus pairing]\label{B7.15}
For primitive class $\phi$ and $n>0$, put $M=v^n-1$, $P_\alpha=\min_F(\alpha)$ and $a_\alpha=n/\deg P_\alpha$. Summing over $\alpha\in\mu_M(\overline F)$ gives $I(\phi,R_l)=\sum_\alpha a_\alpha M_\phi(P_\alpha^{a_\alpha})\overline{t(\alpha)^l}$ and $\sum_\alpha t(\alpha)^h\overline{t(\alpha)^l}=M\mathbf1_{h\equiv l\bmod M}$.
\end{lemma}
\leanref{GlobalTheory.\allowbreak{}CuspGL.\allowbreak{}ipGL\_\allowbreak{}rootSum\_\allowbreak{}eq\_\allowbreak{}sum\_\allowbreak{}torus}{Development.\allowbreak{}AbelianVarietyTheory.\allowbreak{}Proofs.\allowbreak{}CharacterTheory.\allowbreak{}InvertibleEndInnerProductRootPowerSumEqTorusSum}{12--25}
\leanref{GlobalTheory.\allowbreak{}CuspGL.\allowbreak{}torus\_\allowbreak{}orthogonality}{Development.\allowbreak{}AbelianVarietyTheory.\allowbreak{}Proofs.\allowbreak{}CharacterTheory.\allowbreak{}FrobeniusTorusCharacterOrthogonality}{66--93}

\begin{lemma}[Mixed pairing]\label{B7.16}
For $n>0$, $I(\Psi_h,R_l)=|H|\#\{i<n:hv^i\equiv l\bmod M\}$.
\end{lemma}
\leanref{GlobalTheory.\allowbreak{}CuspGL.\allowbreak{}ipGL\_\allowbreak{}psi\_\allowbreak{}rootSum}{Development.\allowbreak{}AbelianVarietyTheory.\allowbreak{}Proofs.\allowbreak{}CharacterTheory.\allowbreak{}InvertibleEndInnerProductRootPowerSumMobius}{30--97}
Insert \ref{B7.14} and $a_\alpha s_h(P_\alpha)=\sum_{i<n}t(\alpha)^{hv^i}$ in \ref{B7.15}, then apply torus orthogonality.

\begin{proposition}[Expansion]\label{B7.17}
For $n>0$ and primitive class $\phi$, $n|H|\phi=\sum_{0\le l<M}I(\phi,R_l)\Psi_l$ on $H$.
\end{proposition}
\leanref{GlobalTheory.\allowbreak{}CuspGL.\allowbreak{}expand\_\allowbreak{}primitive}{Development.\allowbreak{}AbelianVarietyTheory.\allowbreak{}Proofs.\allowbreak{}CharacterTheory.\allowbreak{}PrimitiveClassFunctionExpansionRootSums}{156--221}
The difference is primitive by \ref{B7.11}. For each $i<n$ there is a unique $k\bmod M$ with $kv^i\equiv l\bmod M$. In \ref{B7.15} the permutation $\alpha\mapsto\alpha^{v^i}$ preserves $P_\alpha$ and gives $I(\phi,R_k)=I(\phi,R_l)$. Thus \ref{B7.16} gives zero pairing of the difference with every $R_l$. Fourier inversion gives zero mass on each invertible primary fibre. A primitive class function on $H$ is supported on these fibres and determined by their masses.

\begin{lemma}[Inverted pairing]\label{B7.18}
For primitive $\phi$ and additive $s$, $I(\phi,\Psi_s)=I(\phi,s\circ c)$.
\end{lemma}
\leanref{GlobalTheory.\allowbreak{}CuspGL.\allowbreak{}PrimOrth.\allowbreak{}ipGL\_\allowbreak{}psi\_\allowbreak{}eq}{Development.\allowbreak{}AbelianVarietyTheory.\allowbreak{}Proofs.\allowbreak{}CharacterTheory.\allowbreak{}InvertibleEndomorphismInnerProductPrimitiveMobiu}{92--162}
In \ref{B7.10}, the $W=0$ term is $s(1)=0$. For $0<W<V$, the term depends on $g|_W$, so primitivity cancels its pairing as the off-diagonal block varies with the diagonal blocks fixed.

\begin{lemma}[Regular exponent]\label{B7.19}
If $n>0$ and $hv^i\not\equiv h\bmod M$ for $0<i<n$, then $I(\Psi_h,\Psi_h)=|H|$ by \ref{B7.16} and \ref{B7.18}. By \ref{B7.13}, $\Psi_h|_H$ is a signed irreducible. Its value at $1$ fixes the sign $(-1)^{n-1}$ by \ref{B7.12}.
\end{lemma}
\leanref{GlobalTheory.\allowbreak{}CuspGL.\allowbreak{}isChar\_\allowbreak{}psi\_\allowbreak{}rootSum\_\allowbreak{}of\_\allowbreak{}regular}{Development.\allowbreak{}AbelianVarietyTheory.\allowbreak{}Proofs.\allowbreak{}CharacterTheory.\allowbreak{}IsCharacterCuspidalRootSumRegular}{92--185}

\begin{lemma}[Cuspidality]\label{B7.20}
Over $\F_p$, the character of \ref{B7.19} is cuspidal by \ref{B7.11} at $g=1$.
\end{lemma}
\leanref{GlobalTheory.\allowbreak{}ParamLabel.\allowbreak{}isCusp\_\allowbreak{}psi}{Development.\allowbreak{}AbelianVarietyTheory.\allowbreak{}Proofs.\allowbreak{}DivisorProofs.\allowbreak{}GreenCuspidalCharIsCuspidalOfMinpolyOfParam}{18--49}

\begin{theorem}[Regular family]\label{B7.21}
For $[K:\F_p]=d>0$, a family on regular parameters has the cuspidality and values of \ref{B7.05}, with equality exactly within Frobenius orbits.
\end{theorem}
\leanref{GlobalTheory.\allowbreak{}GreenRegular.\allowbreak{}exists\_\allowbreak{}cusp\_\allowbreak{}of\_\allowbreak{}regular}{Development.\allowbreak{}AbelianVarietyTheory.\allowbreak{}Proofs.\allowbreak{}InertiaEigenvalues.\allowbreak{}ExistsCuspidalCharacterOfRegularCharacterGLn}{65--226}
With $K\subset\overline{\F}_p$ and $\theta=t^h$ on $K^\times$, take $(-1)^{d-1}\Psi_h|_{G_d}$ by \ref{B7.19} and \ref{B7.20}. Its values follow from \ref{B7.07} and \ref{B7.12}. The pairings detect the orbit by \ref{B7.16} and determine the character by \ref{B7.17}.

\begin{lemma}[Orbit injection]\label{B7.22}
For $[K:\F_p]=d>0$, an isomorphism $K^\times\simeq\widehat{K^\times}$ sends chosen roots of $P\in\mathcal P_d$ to characters $c(P)$ in distinct regular orbits.
\end{lemma}
\leanref{GlobalTheory.\allowbreak{}PrimeLabOrbit.\allowbreak{}primeLab\_\allowbreak{}to\_\allowbreak{}regularChar}{Development.\allowbreak{}AbelianVarietyTheory.\allowbreak{}Proofs.\allowbreak{}InertiaEigenvalues.\allowbreak{}MonicIrreducibleInjectRegularCharacterOrbits}{125--152}

\begin{lemma}[Character injection]\label{B7.23}
By \ref{B7.21}, $L_d(P)=\tau_{c(P)}$ injects $\mathcal P_d$ into the cuspidals of $G_d$.
\end{lemma}
\leanref{GlobalTheory.\allowbreak{}GreenRegularCharactersInertiaEigenvalues.\allowbreak{}family}{Development.\allowbreak{}AbelianVarietyTheory.\allowbreak{}Proofs.\allowbreak{}InertiaEigenvalues.\allowbreak{}GreenCuspidalCharsByRegularCharsProof}{23--34}

Let $I(C)$ be the Harish--Chandra product of a cuspidal list $C$ of total degree $n$. For nontrivial additive $\eta$ of $\F_p$, set $\xi_\eta(u)=\eta(\sum_i u_{i,i+1})$ on upper unitriangular $U_n$. Write $\operatorname{wh}_\eta$ for the linearly extended $\xi_\eta$-multiplicity and $\langle\ ,\ \rangle$ for the normalized pairing.

\begin{theorem}[Generic constituent]\label{B7.24}
Every $I(C)$ has an irreducible constituent of positive Whittaker multiplicity, by cuspidal genericity \cite[\S1]{deriziotis-gotsis-1998} and multiplicativity under Harish--Chandra induction.
\end{theorem}
\leanref{GlobalTheory.\allowbreak{}GLq.\allowbreak{}GenericConstituent.\allowbreak{}exists\_\allowbreak{}generic\_\allowbreak{}constituent}{Development.\allowbreak{}AbelianVarietyTheory.\allowbreak{}Proofs.\allowbreak{}InertiaEigenvalues.\allowbreak{}ExistsGenericConstituentHarishChandraCuspidal}{98--126}

\begin{theorem}[Harish--Chandra pairing]\label{B7.25}
For equal total degrees, $\langle I(C),I(C')\rangle=0$ unless $C'$ permutes $C$. Distinct entries in $C$ give $\langle I(C),I(C)\rangle=1$. These are prime-field cases of the Harish--Chandra pairing \cite{zelevinsky-1981}.
\end{theorem}
\leanref{GlobalTheory.\allowbreak{}GLq.\allowbreak{}IpCuspProd.\allowbreak{}ip\_\allowbreak{}hcL\_\allowbreak{}cusp}{Development.\allowbreak{}AbelianVarietyTheory.\allowbreak{}Proofs.\allowbreak{}InertiaEigenvalues.\allowbreak{}CharacterPairingHarishChandraCuspidalMackey}{259--264}

\begin{theorem}[Gelfand--Graev norm]\label{B7.26}
For $n>0$, $\Gamma_\eta=\operatorname{Ind}_{U_n}^{G_n}\xi_\eta$ has $\operatorname{wh}_\eta(\Gamma_\eta)=p^n-p^{n-1}$ and $\langle\chi,\Gamma_\eta\rangle=\operatorname{wh}_\eta(\chi)$ for class functions $\chi$. These are the prime-field norm and reciprocity formulas \cite[\S1]{deriziotis-gotsis-1998}.
\end{theorem}
\leanref{GlobalTheory.\allowbreak{}GLq.\allowbreak{}GGNorm.\allowbreak{}gg\_\allowbreak{}norm}{Development.\allowbreak{}AbelianVarietyTheory.\allowbreak{}Proofs.\allowbreak{}InertiaEigenvalues.\allowbreak{}GelfandGraevCharacterGLnNormAndReciprocity}{88--92}

\begin{corollary}[Generic count]\label{B7.27}
For $n>0$, each generic irreducible contributes at least one to the norm $p^n-p^{n-1}$ in \ref{B7.26}, which bounds their number.
\end{corollary}
\leanref{GlobalTheory.\allowbreak{}GLq.\allowbreak{}card\_\allowbreak{}generic\_\allowbreak{}irr\_\allowbreak{}le}{Development.\allowbreak{}AbelianVarietyTheory.\allowbreak{}Proofs.\allowbreak{}InertiaEigenvalues.\allowbreak{}NcardGenericIrreducibleCharactersGLnLe}{49--89}

\begin{lemma}[Support count]\label{B7.28}
For $n>0$ there are at most $p^n-p^{n-1}$ cuspidal multisets of total degree $n$. Choose generic constituents by \ref{B7.24}. They are distinct by \ref{B7.25} and counted in \ref{B7.27}.
\end{lemma}
\leanref{GlobalTheory.\allowbreak{}GLq.\allowbreak{}CuspMultisetCount.\allowbreak{}card\_\allowbreak{}le}{Development.\allowbreak{}AbelianVarietyTheory.\allowbreak{}Proofs.\allowbreak{}InertiaEigenvalues.\allowbreak{}CardCuspidalMultisetsTotalDegreeLe}{30--62}

\begin{lemma}[Polynomial count]\label{B7.29}
For $n>0$, the lower coefficients give $(p-1)p^{n-1}$ monic degree-$n$ polynomials with nonzero constant term.
\end{lemma}
\leanref{GlobalTheory.\allowbreak{}GrPolys.\allowbreak{}card\_\allowbreak{}monic\_\allowbreak{}coeff\_\allowbreak{}zero\_\allowbreak{}ne\_\allowbreak{}zero}{Development.\allowbreak{}AbelianVarietyTheory.\allowbreak{}Proofs.\allowbreak{}InertiaEigenvalues.\allowbreak{}NcardMonicNatDegreeCoeffZeroNeZero}{38--63}

\begin{lemma}[Exhaustion]\label{B7.30}
Injective cuspidal families $L_e$ on $\mathcal P_e$ for all $e>0$ exhaust the cuspidals.
\end{lemma}
\leanref{GlobalTheory.\allowbreak{}GreenRegularCharactersInertiaEigenvalues.\allowbreak{}count}{Development.\allowbreak{}AbelianVarietyTheory.\allowbreak{}Proofs.\allowbreak{}InertiaEigenvalues.\allowbreak{}GreenCuspidalCharsByRegularCharsProof}{56--132}
The substitution $P\mapsto(\deg P,L_{\deg P}(P))$ sends the polynomials of \ref{B7.29} injectively to the multisets of \ref{B7.28} and attains the bound. A missing degree-$n$ cuspidal would add a singleton.\par
\endgroup

\begin{theorem}[Green's classification]\label{B7.31}
The regular form \ref{B7.05} holds. It takes the prime-field cuspidals and their values at the identity and at elliptic elements from Green's full classification \cite[\S\S7--8, Thm.~14]{green-1955}.
\end{theorem}
\leanref{GlobalTheory.\allowbreak{}greenRegularCharacters\_\allowbreak{}holds}{Development.\allowbreak{}AbelianVarietyTheory.\allowbreak{}Proofs.\allowbreak{}InertiaEigenvalues.\allowbreak{}GreenCuspidalCharsByRegularCharsProof}{137--169}
With the prescribed $K$ in degree $d$, apply \ref{B7.30} to \ref{B7.23}.

\section{Polynomial labels}
Fix $L$ and $\iota:K=\F_{p^d}\hookrightarrow\overline{\F}_p$, $d>0$.

\begin{lemma}[Torus]\label{B7.32}
The image of $\iota$ is the root set of $X^{p^d}-X$. The preimage $v$ of $u_d$ generates $K^\times$ by its order.
\end{lemma}
\leanref{GlobalTheory.\allowbreak{}GreenBridge.\allowbreak{}exists\_\allowbreak{}emb\_\allowbreak{}eq}{Development.\allowbreak{}AbelianVarietyTheory.\allowbreak{}Proofs.\allowbreak{}InertiaEigenvalues.\allowbreak{}SumRangeEqRoots}{32--40}
\leanref{GlobalTheory.\allowbreak{}GreenBridge.\allowbreak{}exists\_\allowbreak{}unit\_\allowbreak{}gen}{Development.\allowbreak{}AbelianVarietyTheory.\allowbreak{}Proofs.\allowbreak{}InertiaEigenvalues.\allowbreak{}SumRangeEqRoots}{53--69}

\begin{lemma}[Duality]\label{B7.33}
The map $x\mapsto\theta_x$, defined by $\theta_x(v^a)=t(x)^a$, is a bijection from $\mu_{N_d}(\overline{\F}_p)$ onto $\widehat{K^\times}$. It satisfies $\theta_x(k)=\Theta_d(x,\iota(k))$.
\end{lemma}
\leanref{GlobalTheory.\allowbreak{}GreenBridge.\allowbreak{}th\_\allowbreak{}apply}{Development.\allowbreak{}AbelianVarietyTheory.\allowbreak{}Proofs.\allowbreak{}InertiaEigenvalues.\allowbreak{}SumRangeEqRoots}{122--125}
\leanref{GlobalTheory.\allowbreak{}GreenBridge.\allowbreak{}th\_\allowbreak{}inj}{Development.\allowbreak{}AbelianVarietyTheory.\allowbreak{}Proofs.\allowbreak{}InertiaEigenvalues.\allowbreak{}SumRangeEqRoots}{142--147}
\leanref{GlobalTheory.\allowbreak{}GreenBridge.\allowbreak{}th\_\allowbreak{}surj}{Development.\allowbreak{}AbelianVarietyTheory.\allowbreak{}Proofs.\allowbreak{}InertiaEigenvalues.\allowbreak{}SumRangeEqRoots}{148--161}
Evaluation at $v$ identifies the character group with $\mu_{N_d}(\C)$. The $N_d$ distinct values $t(x)$ exhaust this group.

\begin{lemma}[Regularity]\label{B7.34}
Evaluation at $v$ gives $\theta_{x^a}=\theta_x^a$ for $x\in\mu_{N_d}$, $a\ge0$. Thus $\theta_x$ is regular iff the $x^{p^i}$, $i<d$, are distinct.
\end{lemma}
\leanref{GlobalTheory.\allowbreak{}GreenBridge.\allowbreak{}th\_\allowbreak{}pow}{Development.\allowbreak{}AbelianVarietyTheory.\allowbreak{}Proofs.\allowbreak{}InertiaEigenvalues.\allowbreak{}SumRangeEqRoots}{137--141}
\leanref{GlobalTheory.\allowbreak{}GreenBridge.\allowbreak{}regular\_\allowbreak{}iff}{Development.\allowbreak{}AbelianVarietyTheory.\allowbreak{}Proofs.\allowbreak{}InertiaEigenvalues.\allowbreak{}SumRangeEqRoots}{207--212}

\begin{lemma}[Label roots]\label{B7.35}
Each $P\in\mathcal P_d$ has $d$ distinct roots in $\mu_{N_d}$ by \ref{B7.07} and $P(0)\ne0$.
\end{lemma}
\leanref{GlobalTheory.\allowbreak{}GreenBridge.\allowbreak{}rootOf\_\allowbreak{}spec}{Development.\allowbreak{}AbelianVarietyTheory.\allowbreak{}Proofs.\allowbreak{}InertiaEigenvalues.\allowbreak{}SumRangeEqRoots}{189--196}
\leanref{GlobalTheory.\allowbreak{}GreenBridge.\allowbreak{}root\_\allowbreak{}unit}{Development.\allowbreak{}AbelianVarietyTheory.\allowbreak{}Proofs.\allowbreak{}InertiaEigenvalues.\allowbreak{}SumRangeEqRoots}{197--206}
\leanref{GlobalTheory.\allowbreak{}GreenBridge.\allowbreak{}pow\_\allowbreak{}inj}{Development.\allowbreak{}AbelianVarietyTheory.\allowbreak{}Proofs.\allowbreak{}InertiaEigenvalues.\allowbreak{}SumRangeEqRoots}{180--186}

\begin{lemma}[Root sums]\label{B7.36}
For irreducible $f$ with root $z$, summing a function along $(z^{p^i})_{i<\deg f}$ sums it over the roots. For $x,w\in\mu_{N_d}$ and $a\ge0$, $\Theta_d(x^a,w)=\Theta_d(x,w^a)$ by \ref{B7.03}.
\end{lemma}
\leanref{GlobalTheory.\allowbreak{}GreenBridge.\allowbreak{}sum\_\allowbreak{}range\_\allowbreak{}eq\_\allowbreak{}roots}{Development.\allowbreak{}AbelianVarietyTheory.\allowbreak{}Proofs.\allowbreak{}InertiaEigenvalues.\allowbreak{}SumRangeEqRoots}{213--217}
\leanref{GlobalTheory.\allowbreak{}GreenBridge.\allowbreak{}thetaOf\_\allowbreak{}pow}{Development.\allowbreak{}AbelianVarietyTheory.\allowbreak{}Proofs.\allowbreak{}InertiaEigenvalues.\allowbreak{}SumRangeEqRoots}{218--224}

\begin{theorem}[Relabelling]\label{B7.37}
The regular form \ref{B7.05} implies \ref{B7.06}.
\end{theorem}
\leanref{GlobalTheory.\allowbreak{}GreenBridge.\allowbreak{}green\_\allowbreak{}of\_\allowbreak{}regular}{Development.\allowbreak{}AbelianVarietyTheory.\allowbreak{}Proofs.\allowbreak{}InertiaEigenvalues.\allowbreak{}SumRangeEqRoots}{225--305}
\begin{proof}
Choose a root $x_P$ of each $P\in\mathcal P_d$. By \ref{B7.34} and \ref{B7.35}, $\tau_P=\tau_{\theta_{x_P}}$ is cuspidal of degree $\Delta_d$. Every characteristic root $z$ of elliptic $y$ lies in $\iota(K^\times)$ by \ref{B7.32} and \ref{B7.35}. The regular formula and \ref{B7.33} give \eqref{eq:green-label} at $x_P$. For another root $x=x_P^{p^a}$, \ref{B7.36} and cycling the roots give
\[
 \sum_{i<d}\Theta_d(x,z^{p^i})
 =\sum_{i<d}\Theta_d(x_P,z^{p^{i+a}})
 =\sum_{i<d}\Theta_d(x_P,z^{p^i}).
\]
Equal characters give conjugate parameters and hence conjugate roots by \ref{B7.33} and \ref{B7.34}. Their minimal polynomials agree. Conversely a cuspidal is $\tau_{\theta_x}$ with $\theta_x$ regular. The degree $e$ of $x$ divides $d$. If $e<d$, then $x^{p^e}=x$ repeats the first conjugate, contrary to \ref{B7.34}. Thus $e=d$. The minimal polynomial of $x$ belongs to $\mathcal P_d$, and its chosen root represents the same orbit. Extend the family arbitrarily off $\mathcal P_d$.
\end{proof}
The unitary-group hypotheses use the regular classification of \ref{B7.31}. Relabelling supplies the polynomial classification used to prove the conclusion for every character.
\leanref{UpperTheory.\allowbreak{}LiteratureTheorems.\allowbreak{}green\_\allowbreak{}regular\_\allowbreak{}characters\_\allowbreak{}of\_\allowbreak{}literature}{Development.\allowbreak{}AutomorphicArithmetic.\allowbreak{}ThetaCorrespondence.\allowbreak{}ThetaPairLifts.\allowbreak{}ThirdPlaceSeries.\allowbreak{}EveryCharacterOfLiteratureOfVWin}{399--400}
\leanref{UpperTheory.\allowbreak{}LiteratureTheorems.\allowbreak{}published\_\allowbreak{}results\_\allowbreak{}eigen}{Development.\allowbreak{}AutomorphicArithmetic.\allowbreak{}ThetaCorrespondence.\allowbreak{}ThetaPairLifts.\allowbreak{}ThirdPlaceSeries.\allowbreak{}EveryCharacterOfLiteratureOfVWin}{666--694}
\leanref{UpperTheory.\allowbreak{}Lit.\allowbreak{}every\_\allowbreak{}character\_\allowbreak{}of\_\allowbreak{}theta\_\allowbreak{}integrals\_\allowbreak{}local\_\allowbreak{}matching}{Development.\allowbreak{}AutomorphicArithmetic.\allowbreak{}ThetaCorrespondence.\allowbreak{}ThetaPairLifts.\allowbreak{}ThirdPlaceSeries.\allowbreak{}SiegelEisensteinSecondTermHypotheses}{69--73}

\section{Series weights}
\begin{lemma}[Proper subfields]\label{B7.38}
For $p\ge5,d>0$, the elements $B$ of degree $<d$ in $\F_{p^d}^\times$ satisfy $2d|B|<N_d$.
\end{lemma}
\leanref{GlobalTheory.\allowbreak{}GLm.\allowbreak{}card\_\allowbreak{}proper\_\allowbreak{}subfield\_\allowbreak{}elements\_\allowbreak{}lt}{Development.\allowbreak{}AbelianVarietyTheory.\allowbreak{}Proofs.\allowbreak{}FlagMeasures.\allowbreak{}CardProperSubfieldElementsLt}{67--103}
Put $h=\lfloor d/2\rfloor$. Then $|B|\le S_h=\sum_{e=1}^h(p^e-1)$. For $h\ge1$, induction gives $4S_h+4h+5\le5p^h$ and $5h\le p^h$, hence $4hS_h<p^{2h}-1$. For odd $d>1$, $2dS_h<8hS_h<2p^{2h}<N_d$. For $d=1$, $B=\varnothing$.

\begin{lemma}[Power sums]\label{B7.39}
Let two lists of length $d>0$ in $\mu_M(\C)$, $M>1$, have equal power sums outside $B\subset\Z/M$, with $2d|B|<M$. Then the first entry of the second list occurs in the first list.
\end{lemma}
\leanref{GlobalTheory.\allowbreak{}GLq.\allowbreak{}LabelUnique.\allowbreak{}exists\_\allowbreak{}eq\_\allowbreak{}of\_\allowbreak{}sum\_\allowbreak{}eq}{Development.\allowbreak{}AbelianVarietyTheory.\allowbreak{}Proofs.\allowbreak{}InertiaEigenvalues.\allowbreak{}CuspidalCharacterLabelUniqueGLn}{28--110}
The multiplicity differences $c_\zeta$ have support of size at most $2d$ and Fourier support in $B$. Parseval and Cauchy--Schwarz give $M\sum_\zeta|c_\zeta|^2\le2d|B|\sum_\zeta|c_\zeta|^2$, so all differences vanish.

\begin{theorem}[Unique labels]\label{B7.40}
For $p\ge5$, two prime labels of one cuspidal agree. Under Green, two cuspidals with the same prime label agree.
\end{theorem}
\leanref{GlobalTheory.\allowbreak{}GLq.\allowbreak{}LabelUnique.\allowbreak{}label\_\allowbreak{}eq}{Development.\allowbreak{}AbelianVarietyTheory.\allowbreak{}Proofs.\allowbreak{}InertiaEigenvalues.\allowbreak{}CuspidalCharacterLabelUniqueGLn}{174--258}
\leanref{GlobalTheory.\allowbreak{}GLq.\allowbreak{}isCuspOf\_\allowbreak{}label\_\allowbreak{}unique}{Development.\allowbreak{}AbelianVarietyTheory.\allowbreak{}Proofs.\allowbreak{}InertiaEigenvalues.\allowbreak{}CuspidalCharacterLabelUniqueGLn}{263--274}
Evaluate \eqref{eq:green-label} on multiplication by $u_d^a$ of degree $d$. The lifted root lists have equal power sums outside the exponents corresponding to $B$ in \ref{B7.38}. By \ref{B7.39} and injectivity of $t$, the two minimal polynomials share a root. Under Green the unique label fixes the member of its family.

\begin{definition}[Series]\label{B7.41}
An irreducible $\rho$ has label $f$ if it occurs in $I(C)$ for a cuspidal list whose prime labels multiply to $f$. Sum the corresponding central idempotents $z_\rho$ to get $E_f$. For $a\in\C[G_n]$ put $w_a(f)=\Re[aE_f]_1/\Re[a]_1$, with identity coefficient $[\,]_1$ and $0^{-1}=0$.
\end{definition}
\leanref{GlobalTheory.\allowbreak{}GLq.\allowbreak{}HasLabel}{Development.\allowbreak{}AbelianVarietyTheory.\allowbreak{}Definitions.\allowbreak{}NonGaloisDescent.\allowbreak{}LabelDatum}{98--101}
\leanref{GlobalTheory.\allowbreak{}GLq.\allowbreak{}eSerC}{Development.\allowbreak{}AbelianVarietyTheory.\allowbreak{}Definitions.\allowbreak{}NonGaloisDescent.\allowbreak{}LabelDatum}{102--104}
\leanref{GlobalTheory.\allowbreak{}GLq.\allowbreak{}wtC}{Development.\allowbreak{}AbelianVarietyTheory.\allowbreak{}Definitions.\allowbreak{}NonGaloisDescent.\allowbreak{}LabelDatum}{105--107}

\begin{theorem}[Squarefree series]\label{B7.42}
For $p\ge5$ under Green, each irreducible of $G_n$ has a unique monic degree-$n$ label with nonzero constant term. For squarefree such $f$, its series consists of $\rho_f=I((\deg P,\tau_P)_{P\mid f})$, over monic irreducible factors.
\end{theorem}
\leanref{GlobalTheory.\allowbreak{}GLq.\allowbreak{}Label.\allowbreak{}label\_\allowbreak{}main}{Development.\allowbreak{}AbelianVarietyTheory.\allowbreak{}Proofs.\allowbreak{}InertiaEigenvalues.\allowbreak{}HarishChandraSeriesLabelExistsUniqueGLn}{271--279}
A minimal Levi with nonzero Harish--Chandra restriction supplies cuspidal factors by transitivity and support by reciprocity. By \ref{B7.25} and \ref{B7.40} the support fixes the label. For squarefree $f$, every such list is a permutation of the same distinct cuspidals. Induction respects permutation and has norm one by \ref{B7.25}.
\leanref{GlobalTheory.\allowbreak{}gl\_\allowbreak{}list\_\allowbreak{}product\_\allowbreak{}perm\_\allowbreak{}char}{Development.\allowbreak{}AbelianVarietyTheory.\allowbreak{}Proofs.\allowbreak{}InertiaEigenvalues.\allowbreak{}HarishChandraProductListIsCharacterAndPermInvari}{74--77}

\begin{lemma}[Arithmetic relabelling]\label{B7.43}
Assume $p\ge5$ and Green. Let $O$ label simple cuspidal representations by $\mathcal P_d$ in each degree $d>0$, with equality equivalent to isomorphism. Let $E^O_{r,f}$ be the rank-$r$ series projector for the product $f$ of $O$-labels. There are inverse degreewise permutations $\sigma,\pi$ of the prime labels. On monic polynomials $f$ with $f(0)\ne0$, their multiplicative extensions preserve degree and squarefreeness, with $\pi_*\sigma_*f=f$ and $E^O_{r,f}=E_{\sigma_*f}$.
\end{lemma}
\leanref{GlobalTheory.\allowbreak{}Top.\allowbreak{}Lq.\allowbreak{}exists\_\allowbreak{}perm}{Development.\allowbreak{}AbelianVarietyTheory.\allowbreak{}Proofs.\allowbreak{}GlobalInterpolation.\allowbreak{}AdmissibleLabelWeightsOfGreenCuspidalParametrisa}{89--101}
\leanref{GlobalTheory.\allowbreak{}Top.\allowbreak{}Lq.\allowbreak{}ext\_\allowbreak{}back}{Development.\allowbreak{}AbelianVarietyTheory.\allowbreak{}Proofs.\allowbreak{}GlobalInterpolation.\allowbreak{}AdmissibleLabelWeightsOfGreenCuspidalParametrisa}{129--134}
\leanref{GlobalTheory.\allowbreak{}Top.\allowbreak{}Lq.\allowbreak{}eSer\_\allowbreak{}eq}{Development.\allowbreak{}AbelianVarietyTheory.\allowbreak{}Proofs.\allowbreak{}GlobalInterpolation.\allowbreak{}AdmissibleLabelWeightsOfGreenCuspidalParametrisa}{233--250}
For a simple cuspidal $V_P$ of character $\tau_P$, put $\pi(P)=O(V_P)$. The isomorphism criterion for $O$ and Green make $\pi$ well-defined and injective, hence bijective by finiteness. Its inverse $\sigma$ gives $\sigma_*f=\prod_P\sigma(P)^{v_P(f)}$. The substitution preserves degree and multiplicities. The projectors coincide because it preserves cuspidal support.

For a subgroup $H$ and a degree-one character $\xi$, put $e(H,\xi)=|H|^{-1}\sum_{u\in H}\xi(u)^{-1}u$. Let $e_U=e(U_n,\xi_\eta)$. For $n=k+2s+1$, $k\ge1$, $s\ge0$, use zero-based indices. Let $H_J\subset U_n$ impose $u_{a,s+k}=0$ for $s\le a<s+k$. Set $e_J=e(H_J,\eta\circ\phi_J)$, where $\phi_J$ sums consecutive entries along $(0,\dots,s-1,s+k,\dots,n-1)$ and $(s,\dots,s+k-1)$. Write $w_U=w_{e_U}$ and $w_J=w_{e_J}$.

\begin{lemma}[Idempotents]\label{B7.44}
The subgroup averages $e_U,e_J$ are idempotent with identity coefficients $p^{-b},p^{k-b}$, respectively, for $b=n(n-1)/2$. For idempotent $a$ with $\Re[a]_1\ge0$, $w_a(f)\ge0$. The coefficient $[az_\rho]_1$ is $\rho(1)/|G_n|$ times the rank of $a$ in the representation of $\rho$.
\end{lemma}
\leanref{GlobalTheory.\allowbreak{}T.\allowbreak{}whittakerJacquetIdempotents\_\allowbreak{}properties\_\allowbreak{}primeField}{Development.\allowbreak{}AbelianVarietyTheory.\allowbreak{}Proofs.\allowbreak{}InertiaEigenvalues.\allowbreak{}WhittakerJacquetIdempotentsPropertiesPrimeFieldStatement}{8--9}
\leanref{GlobalTheory.\allowbreak{}T.\allowbreak{}harishChandraSeriesWeight\_\allowbreak{}nonneg\_\allowbreak{}of\_\allowbreak{}idempotent\_\allowbreak{}glnPrime}{Development.\allowbreak{}AbelianVarietyTheory.\allowbreak{}Proofs.\allowbreak{}InertiaEigenvalues.\allowbreak{}HarishChandraSeriesWeightNonnegOfIdempotentGlnPrStatement}{8--9}

\begin{lemma}[Multiplicity one]\label{B7.45}
A cuspidal $\chi$ of degree $\Delta_d$ has $\operatorname{wh}_\eta(\chi)=1$. This is the prime-field cuspidal multiplicity-one result with the degree specified \cite[\S1]{deriziotis-gotsis-1998}.
\end{lemma}
\leanref{GlobalTheory.\allowbreak{}GLq.\allowbreak{}wh\_\allowbreak{}cusp\_\allowbreak{}eq\_\allowbreak{}one}{Development.\allowbreak{}AbelianVarietyTheory.\allowbreak{}Proofs.\allowbreak{}InertiaEigenvalues.\allowbreak{}WhittakerDimensionCuspidalGreenDegreeEqOne}{66--80}

\begin{theorem}[Squarefree weight]\label{B7.46}
Under \ref{B7.42}, for monic squarefree $f$ of degree $n$ with $f(0)\ne0$, $w_U(f)=\prod_{P\mid f}(p^{\deg P}-1)^{-1}$ over monic irreducibles.
\end{theorem}
\leanref{ProductIntervalSplit.\allowbreak{}CompanionTheory.\allowbreak{}Member.\allowbreak{}harishChandraSeriesWeight\_\allowbreak{}squarefree\_\allowbreak{}label\_\allowbreak{}unipotent}{Development.\allowbreak{}AbelianVarietyTheory.\allowbreak{}Proofs.\allowbreak{}InertiaEigenvalues.\allowbreak{}HarishChandraSeriesWeightSquarefreeLabelGelfandG}{238--294}
The parabolic index times the cuspidal degrees gives $\rho_f(1)=|G_n|p^{-b}\prod_{P\mid f}(p^{\deg P}-1)^{-1}$. By \ref{B7.45} and multiplicativity, $\rho_f(e_U)=1$. Apply \ref{B7.44} to the sole constituent $\rho_f$.

\begin{theorem}[Lower comparison]\label{B7.47}
For any prime $p$, datum $L$, $n=k+2s+1$ with $k\ge1$, and every $f$, $(1-p^{-1})w_U(f)\le w_J(f)$.
\end{theorem}
\leanref{GlobalTheory.\allowbreak{}GLq.\allowbreak{}lowerWtComparison}{Development.\allowbreak{}AbelianVarietyTheory.\allowbreak{}Proofs.\allowbreak{}InertiaEigenvalues.\allowbreak{}HarishChandraSeriesWeightLowerComparison}{49--98}
\leanref{GlobalTheory.\allowbreak{}GLq.\allowbreak{}tr\_\allowbreak{}eJ\_\allowbreak{}eq\_\allowbreak{}tr\_\allowbreak{}eJ'}{Development.\allowbreak{}AbelianVarietyTheory.\allowbreak{}Proofs.\allowbreak{}InertiaEigenvalues.\allowbreak{}GroupAlgebraTracePatternIdempotentExchange}{13--28}
\leanref{GlobalTheory.\allowbreak{}GLq.\allowbreak{}ExactMult.\allowbreak{}tr\_\allowbreak{}eJ'\_\allowbreak{}eq\_\allowbreak{}sum\_\allowbreak{}wh\_\allowbreak{}der}{Development.\allowbreak{}AbelianVarietyTheory.\allowbreak{}Proofs.\allowbreak{}InertiaEigenvalues.\allowbreak{}TraceUnitriangularColumnVanishingIdempotentEqSum}{69--93}
Let $D_i^\eta$ take unipotent restriction to blocks $(i,n-i)$ and then the Whittaker multiplicity space in the second block. Column exchanges give the $\xi_\eta$-average on $u_{a,k}=0$ for $a<k$ with the same trace. Mirabolic decomposition and the Gelfand--Graev norms in rank $k-i$ give
\begin{equation}\label{eq:mirabolic-trace}
 \chi(e_J)=\sum_{i<k}(p^{k-i}-p^{k-i-1})\operatorname{wh}_\eta(D_i^\eta\chi)
 +\operatorname{wh}_\eta(D_k^\eta\chi).
\end{equation}
All terms are nonnegative for irreducible $\chi$. The $i=0$ term is $(p^k-p^{k-1})\chi(e_U)$. Sum by \ref{B7.44} and divide by the coefficient ratio $p^k$.

\begin{lemma}[Derivative counts]\label{B7.48}
Let $C$ have total degree $n$ and cuspidal factors of Whittaker multiplicity one. Then $\operatorname{wh}_\eta(I(C))=1$, and $\operatorname{wh}_\eta(D_i^\eta I(C))$ counts subsets of positions of total degree $i$, for $0\le i\le n$. In Mackey restriction, cuspidality eliminates the terms that split a factor. Each remaining term contributes one by Whittaker multiplicativity.
\end{lemma}
\leanref{GlobalTheory.\allowbreak{}GLq.\allowbreak{}Leibniz.\allowbreak{}main}{Development.\allowbreak{}AbelianVarietyTheory.\allowbreak{}Proofs.\allowbreak{}InertiaEigenvalues.\allowbreak{}BernsteinZelevinskyDerivativeHarishChandraProduct}{102--143}

\begin{theorem}[Upper comparison]\label{B7.49}
Under \ref{B7.46}, if $n=k+2s+1$, $k\ge1$, then $w_J(f)\le(1+p^{-1})^nw_U(f)$.
\end{theorem}
\leanref{GlobalTheory.\allowbreak{}GLq.\allowbreak{}UpperWeight.\allowbreak{}comb\_\allowbreak{}bound}{Development.\allowbreak{}AbelianVarietyTheory.\allowbreak{}Proofs.\allowbreak{}InertiaEigenvalues.\allowbreak{}HarishChandraSeriesWeightLeSquarefree}{30--81}
\leanref{GlobalTheory.\allowbreak{}GLq.\allowbreak{}UpperWeight.\allowbreak{}wtC\_\allowbreak{}eJ\_\allowbreak{}le}{Development.\allowbreak{}AbelianVarietyTheory.\allowbreak{}Proofs.\allowbreak{}InertiaEigenvalues.\allowbreak{}HarishChandraSeriesWeightLeSquarefree}{88--140}
For the subset counts $c_i$ of \ref{B7.48}, \eqref{eq:mirabolic-trace} gives
\[
 \frac{w_J(f)}{w_U(f)}
 =p^{-k}\rho_f(e_J)\le\sum_{i\le k}c_ip^{-i}
 \le\prod_{P\mid f}(1+p^{-\deg P})\le(1+p^{-1})^n.
\]

\begin{theorem}[Arithmetic weights]\label{B7.50}
Assume Green. Let $D$ be of type $D7$ with compatible normalized types $T$ for an embedding $j_E:D.\mathrm{host}.E\hookrightarrow\overline{\Q}$. At each type prime $p_j$ and for nontrivial $\eta$, the arithmetic $U,J$ weights are nonnegative and supported on monic degree-$r$ labels with nonzero constant term. They satisfy \ref{B7.47} on these labels and \ref{B7.46}, \ref{B7.49} on squarefree ones.
\end{theorem}
\leanref{GlobalTheory.\allowbreak{}Top.\allowbreak{}Lq.\allowbreak{}admissibleLabelWeights\_\allowbreak{}of\_\allowbreak{}greenCuspidalParametrisation}{Development.\allowbreak{}AbelianVarietyTheory.\allowbreak{}Proofs.\allowbreak{}GlobalInterpolation.\allowbreak{}AdmissibleLabelWeightsOfGreenCuspidalParametrisa}{255--300}
\leanref{Formalism.\allowbreak{}MainNotions.\allowbreak{}ArithmeticCylinders.\allowbreak{}admLq}{Formalism.\allowbreak{}SharpRestriction.\allowbreak{}LevelLetterCocycles.\allowbreak{}ArithmeticCylinders}{340--352}
Apply \ref{B7.43} at $p_j\ge5$ with lift $D.\mathrm{ntypes}.\mathrm{teichC}(j)$ and generators $T.\zeta(j)$. Here $r=k+2s+1$, $k\ge1$, and $s=\lfloor m/2\rfloor-1$. The arithmetic projector is $E_{\sigma_*f}$ on admissible labels and zero elsewhere. The substitution preserves the degrees and squarefreeness used in the weight formulas.

\section{Bad rows}
Write $\nu$ for the number of places over each type prime and $m$ for the chain length. For arithmetic row weights $w_{A,j,v}$ set $P_j(S)=\sum_{t\in S}\prod_v w_{A,j,v}(t_v)$ and
\[
 W(\nu,m,r,p,q,d)
 =\frac{3\nu}{p}+\frac{23\nu\log p}{\sqrt p}
 +2\nu m\left(\frac1{4\min(q,d)}+\frac{1+\log r}{p-1}\right)
 +\frac r2(7/8)^\nu+e^{-\nu/(5\log r)}.
\]
For a prime $\ell$ and a finite-order element $x$, write $x=x_\ell x_{\ell'}$ in $\langle x\rangle$, with $x_\ell$ of $\ell$-power order and $x_{\ell'}$ of order prime to $\ell$.

At $p=p_j$, a good tuple has squarefree admissible polynomials $t_v$. Within each $t_v$ the prime-to-$\ell$ parts of its roots are distinct for primes $\ell$ outside
$\{p\}\cup\{\ell:\exists\,1\le b\le r,\ p^b<\ell^{4v_\ell(p^b-1)}\}$.
For every prime $\ell\ne p$ these parts avoid those of the prescribed $\chi_{v,a}^{\pm1}\in\F_p^\times$, $1\le a\le m$. For each $1\le b\le r/2$ some $t_v$ has no degree-$b$ divisor. Some $t_v$ has an irreducible factor of prime degree $b$ with $r<2b$ and $b+3\le r$.
\leanref{Formalism.\allowbreak{}CyclotomicNorms.\allowbreak{}IsGood}{Formalism.\allowbreak{}CyclotomicNorms.\allowbreak{}StatementT}{19--30}
\leanref{Formalism.\allowbreak{}MainNotions.\allowbreak{}ArithmeticCylinders.\allowbreak{}goodAt}{Formalism.\allowbreak{}SharpRestriction.\allowbreak{}LevelLetterCocycles.\allowbreak{}ArithmeticCylinders}{355--359}

\begin{lemma}[Bad-label weight]\label{B7.51}
Let $D$ satisfy \ref{B7.50}, \Lean{TieAdmG}, $r\ge100$, and $p_j\ge4096$ for $j=0,1$. If the row idempotents above $p_j$ are $e_U(\eta_j)$ or $e_J(\eta_j)$ for one nontrivial $\eta_j$, their weights are nonnegative and $0\le P_j(S)\le W(\nu,m,r,p_j,q,d)$ for each finite set $S$ of bad tuples.
\end{lemma}
\leanref{Formalism.\allowbreak{}L2.\allowbreak{}GlobalTheoryTheory.\allowbreak{}ConcNWeight.\allowbreak{}weightAllRows\_\allowbreak{}of}{Formalism.\allowbreak{}LocalToGlobal.\allowbreak{}GlobalTheory.\allowbreak{}ConcNWeight}{345--356}
\leanref{Formalism.\allowbreak{}L2.\allowbreak{}GlobalTheoryTheory.\allowbreak{}AllJoin.\allowbreak{}weight\_\allowbreak{}of\_\allowbreak{}join}{Formalism.\allowbreak{}LocalToGlobal.\allowbreak{}HilbertRowConditions.\allowbreak{}AllJoinP}{22--25}
\leanref{Formalism.\allowbreak{}IrreducibleCountingLemmas.\allowbreak{}GaloisIrreducibility.\allowbreak{}property\_\allowbreak{}W\_\allowbreak{}of\_\allowbreak{}estimates}{Formalism.\allowbreak{}PadicSubsetSums.\allowbreak{}StableGradedCounts.\allowbreak{}PropertyWOfEstimates}{7--8}
\begin{proof}
The host gives distinct odd primes $q,d$ with $\max(q,d)\ge16$ and $16,q,d\mid p-1$. The projectors sum to one, so the $U$-law $P$ and row laws $Q_v$ are probabilities. Let $S'$ consist of the tuples in $S$ with all entries admissible. Every tuple in $S\setminus S'$ has weight zero, so $P_j(S)=P_j(S')$. On $S'$ the lower $r$ coefficients determine each polynomial. By \Lean{good\_\allowbreak{}of}, the coefficient image of $S'$ lies in the exceptional event of the polynomial estimate.

By the arithmetic squarefree law of \ref{B7.50}, the union bounds for non-squarefree or inadmissible entries and for root and twist collisions cost the first three terms of $W$ with $2\nu/p$ in the first term. For each $1\le b\le\lfloor r/2\rfloor$, the squarefree admissible polynomials with a degree-$b$ divisor have $P$-mass at most $7/8$. The product across places and the union over these degrees cost at most $\lfloor r/2\rfloor(7/8)^\nu$. Squarefree admissible polynomials missing the required prime-degree factor have $P$-mass at most $1-1/(5\log r)$, so missing it at all places costs at most $e^{-\nu/(5\log r)}$.

By \ref{B7.50}, $Q_v\ge(1-p^{-1})P$. Apply the product lower bound on the complement of the exceptional event. The exceptional probability grows by at most $1-(1-p^{-1})^\nu\le\nu/p$.
\end{proof}

\begin{theorem}[Defect bound]\label{B7.52}
Fix $0<\varepsilon\le1$, $c>0$, $C_s\in\R$, and $F_m$ with $J\ge1$, $\alpha=\varepsilon/J$. For all sufficiently large $x$ and each $i\in F_m.\mathrm{Ix}(x)$, suppose one row system $A$ satisfies \Lean{AllJoinA}, the laws of \ref{B7.50}, and
\[
 p_j\ge e^{8x^\alpha},\qquad x^2\le\nu\le x^{C_s},\qquad
 m\le x^{C_s},\qquad100\le r\le x^{C_s},\qquad
 \min(q,d)\ge e^{x^{\alpha/6}}.
\]
Assume the map from retained to all-row indices preserves labels and is injective with image exactly the indices good at both primes. For each set of label pairs, assume its dimension equals the number of all-row indices whose pair of labels belongs to that set.

For each cylinder $\mathrm{cyl}_j S$ let $N_j(S)$ count its indices. Put $D_j(S)=\mathrm{DstarE}(A,A.e_T(\mathrm{cyl}_j S))$, where $\mathrm{DstarE}(A,e)=d_\infty\operatorname{vol}(\mathrm{fund})\Re e(1)$ and $d_\infty$ is the degree factor. Assume $D_*=\mathrm{DstarE}(A,A.B.e)>0$ and, for every cylinder,
\[
 |N_j(S)-D_j(S)|\le\epsilon_C D_j(S),\qquad
 \epsilon_C=e^{-cr\nu\log r}.
\]
Eventually, for every $i$, the relation $D_{\rm rows}=(1-\delta_{x,i})D_{\rm all}$ gives $0\le\delta_{x,i}\le e^{-x^{\alpha/7}}$.
\end{theorem}
\leanref{GlobalTheory.\allowbreak{}Top.\allowbreak{}Irr.\allowbreak{}linkIRR\_\allowbreak{}of\_\allowbreak{}join\_\allowbreak{}ex}{Development.\allowbreak{}AbelianVarietyTheory.\allowbreak{}Proofs.\allowbreak{}GlobalInterpolation.\allowbreak{}LinkIRROfJoinEx}{17--33}
\leanref{Formalism.\allowbreak{}L2.\allowbreak{}GlobalTheoryTheory.\allowbreak{}ConcN.\allowbreak{}CountAllRows}{Formalism.\allowbreak{}LocalToGlobal.\allowbreak{}GlobalTheory.\allowbreak{}IrreducibleOfCountSizes}{183--187}
\begin{proof}
The full cylinder gives $D_{\rm all}\ge(1-\epsilon_C)D_*$. Let $S_j$ be the finite set of $j$-labels of indices bad at $j$. The retained-index condition gives $0\le D_{\rm all}-D_{\rm rows}\le\sum_jN_j(S_j)$. For finite $S$, \Lean{eT\_\allowbreak{}one} gives the real part of the cylinder kernel at the identity as $\Re((A.B.e)(1))P_j(S)$. Hence $D_j(S)=D_*P_j(S)$.

The level condition in \Lean{AllJoinA} gives one nontrivial $\eta_j$ with local idempotent $e_J(\eta_j)$ at every place above $p_j$. For each $j$, the $D7$ data give primality of $p_j$ and complete splitting in $E$ and $M_{\rm ray}$. They also give $p_j\equiv1\bmod Q_Sqd$. These are the conditions of \Lean{TieAdmG}. For $x^\alpha\ge3/2$, apply \ref{B7.51} and the relative count to obtain
\[
 0\le\delta_{x,i}\le
 \frac{1+\epsilon_C}{1-\epsilon_C}\sum_{j=0}^1 W(\nu,m,r,p_j,q,d).
\]
Eventually the prefactor is at most $3$. The first three terms of $W$ total $O(x^{3C_s+3}e^{-x^{\alpha/6}})$ by the size bounds and the eventual decrease of $(\log p)/\sqrt p$. The last two are $O(x^{C_s}\exp(-c'x^2/\log x))$ for some $c'>0$, since $\nu/\log r\ge x^2/(C_s\log x)$. As $0<\alpha\le1$, both bounds are $o(e^{-x^{\alpha/7}})$, uniformly in the index $i$.
\end{proof}

\begin{theorem}[Closed defect bound]\label{B7.53}
Assume Green. Fix $c>0$, a family class $\mathcal A$ and a twist selection. Assume the following assertion is closed on $\mathcal D_3(\mathcal A)$. There is an $x_c$ such that for every $x\ge x_c$ and every $i\in F_m.\mathrm{Ix}(x)$, one row system satisfies the index, dimension, join and relative-count conditions of \ref{B7.52}. Then the defect bound of \ref{B7.52} is closed on this class in the sense of \ref{B8.39}.
\end{theorem}
\leanref{GlobalTheory.\allowbreak{}Top.\allowbreak{}Irr.\allowbreak{}closed\_\allowbreak{}linkIRR\_\allowbreak{}of\_\allowbreak{}join\_\allowbreak{}count}{Development.\allowbreak{}AutomorphicArithmetic.\allowbreak{}ThetaCorrespondence.\allowbreak{}ThetaPairLifts.\allowbreak{}ThirdPlaceSeries.\allowbreak{}LinkIrr}{22--43}
\begin{proof}
For each sufficiently small $\varepsilon>0$, use the hypothesis's size threshold $K$. For eligible $F_m,C_H$ with $\mathrm{Large}(F_m,K,C_H)$, take the maximum of $x_c$, the selection threshold, $F_m.x_1+1$, $4$, $(3/2)^{1/\alpha}$ and the numerical threshold in \ref{B7.52}. At each subsequent point the selected normalized types give \ref{B7.50} for the same rows, and the design sizes give $C_s=6$ in \ref{B7.52}. 
\end{proof}

\begin{theorem}[Connected families]\label{B7.54}
Green gives $\mathrm{TopShapeAlt}(\mathrm{True},\mathrm{True},\mathcal C\mapsto\mathrm{Closed}(\mathcal C,\mathrm{LinkN}),\mathrm{TopG})$. Under the common hypotheses of \ref{B8.02}, assume \Lean{LinkN} is closed on $\mathcal D_3(\mathcal A)$. Fix $\lambda>1$ and a request $G$ with $\operatorname{HostExW}(\lambda,G.\mathrm{req})$. If $\operatorname{HypV}(F_m.\mathrm{toFamCore})$ holds for every $F_m\in G.\mathrm{cls}(\lambda)$, then $\operatorname{FamilyClassConnected}(G.\mathrm{cls}(\lambda))$.
\end{theorem}
\begin{proof}
Put $\mathcal E=G.\mathrm{cls}(\lambda)$ and $\mathcal A_G=\mathcal A\cap\mathcal E$. The design field of \ref{B8.02}, with the rational Minkowski bound, gives choices with $\Gamma_u\mid M(2fN(m))$ eventually. The prime factors of this Minkowski multiple are at most $2fN(m)+1$, so the design is smooth in the sense of \ref{B8.02}. The divisor bound also gives $\log\Gamma_u\le12 fN(m)\log(fN(m))$, the estimate \Lean{DGamSize}. On $\mathcal W_1(\mathcal A_G)$, the $T_u$ field gives \Lean{TuSize} under $C_{Tu}\ge c_T(\varepsilon,C_{YL})$. The class-number formula \ref{B8.54} and the selected rank ties give \Lean{ChosenProviso} at $3$, and \ref{B8.47} gives \Lean{ChosenThin}.

Apply the design field to the componentwise maximum of the requested threshold $\mathcal K$ and those for these closed conditions. This gives $F_m\in\mathcal W_1(\mathcal A_G)$ and $C_H\ge0$ with $\operatorname{Large}(F_m,\mathcal K,C_H)$, $\operatorname{ChosenE}(F_m,3,C_H)$ and both size bounds.

On $\mathcal D_3(\mathcal A)$, the row conditions of \ref{B8.02} and Green give closed \Lean{LinkIRR} by \ref{B7.53}. The top-set estimates from \Lean{ChosenE} and the compatible local groups from the selected second layer give closed \Lean{LinkG}. Since $F_m\in\mathcal A$, the fifth duty gives \Lean{ChosenWall}. Proposition~\ref{B8.02a} proves this duty for the constructed class. By \ref{B5.22b}, \Lean{ChosenWS} supplies a common weight pair with $R_W\le x^{-9}$ in every model realizing it. The top-set and bad-orbit bounds in \Lean{ChosenE} and the local comparison of \ref{B5.20a} give the modelwise estimate of \ref{B5.21}. Applying \ref{B5.22} to the common weights gives $e_L=0$. The other choices in \Lean{ChosenE} supply the arithmetic inputs to the loss and window estimates of \ref{B5.33}. The real-place second-layer conditions and the volume and mass fields of \ref{B8.02} supply the closed rank and global implications.

The selected twist data give the eventual local covering required by \Lean{LinkEta}. The spectral and counting estimates of \ref{B8.02} give the closed implication $\operatorname{TuSize}(F_m)\Rightarrow\operatorname{LinkS}(F_m,3)$.

Combine the closed estimates on $\mathcal D_3(\mathcal A)$ by \ref{B8.39}. Take the minimum of their $\varepsilon$ thresholds and the design threshold. For each smaller positive $\varepsilon$, take the componentwise maximum $\mathcal K$ of the closed estimates' size thresholds before choosing the design. The family furnished above lies in both $\mathcal D_3(\mathcal A)$ and $\mathcal E$.

Combine \Lean{HypV} and \Lean{LinkN} with the small-level rank equality from \Lean{zL} to obtain \Lean{DIVinN}. With \Lean{LinkEta} and the size bounds, the conditional estimate for the selected rows gives division with $C_{\rm abs}=2$. Division, full rank, the slope-rank condition and the size bounds give the slope comparison. The global implication for gain, frame, defect, division and slope yields $K,c_{HH},e_K,C_L,x_0$ such that $\operatorname{ChainTied}(F_m,2,0,K,c_{HH},e_K,C_L,x)$ holds for every window point $x\ge x_0$. Thus $\operatorname{FamilyClassConnected}(\mathcal E)$ holds with $(e_D,e_L)=(2,0)$, as required for \Lean{TopG} in \ref{B1.18}.
\end{proof}
\leanref{UpperTheory.\allowbreak{}Top.\allowbreak{}shape\_\allowbreak{}closed\_\allowbreak{}link\_\allowbreak{}n\_\allowbreak{}of\_\allowbreak{}green}{Development.\allowbreak{}AutomorphicArithmetic.\allowbreak{}ThetaCorrespondence.\allowbreak{}ThetaPairLifts.\allowbreak{}ThirdPlaceSeries.\allowbreak{}TopShape}{33--55}
\leanref{Formalism.\allowbreak{}MainNotions.\allowbreak{}designExists\_\allowbreak{}of\_\allowbreak{}choices}{Formalism.\allowbreak{}SharpRestriction.\allowbreak{}DesignExistsOfChoices}{54--68}
\leanref{GlobalTheory.\allowbreak{}Top.\allowbreak{}designExists\_\allowbreak{}of\_\allowbreak{}tuSize\_\allowbreak{}aux4n\_\allowbreak{}litThin}{Development.\allowbreak{}AutomorphicArithmetic.\allowbreak{}ThetaCorrespondence.\allowbreak{}ThetaPairLifts.\allowbreak{}ThirdPlaceSeries.\allowbreak{}LinkEx}{15--37}
\leanref{GlobalTheory.\allowbreak{}Top.\allowbreak{}closed\_\allowbreak{}linkG\_\allowbreak{}of\_\allowbreak{}duty}{Development.\allowbreak{}AutomorphicArithmetic.\allowbreak{}ThetaCorrespondence.\allowbreak{}ThetaPairLifts.\allowbreak{}ThirdPlaceSeries.\allowbreak{}LinkG}{13--24}
\leanref{Formalism.\allowbreak{}L2.\allowbreak{}GlobalTheoryTheory.\allowbreak{}Loss7.\allowbreak{}loss\_\allowbreak{}requirements\_\allowbreak{}of\_\allowbreak{}small\_\allowbreak{}shares}{Formalism.\allowbreak{}LocalToGlobal.\allowbreak{}GlobalTheory.\allowbreak{}LossRequirementsOfSmallShares}{337--340}
\leanref{GlobalTheory.\allowbreak{}Top.\allowbreak{}link\_\allowbreak{}gl}{Development.\allowbreak{}AbelianVarietyTheory.\allowbreak{}Proofs.\allowbreak{}GlobalInterpolation.\allowbreak{}LinkGl}{14--31}
\leanref{UpperTheory.\allowbreak{}Top.\allowbreak{}topDesWin}{Development.\allowbreak{}AutomorphicArithmetic.\allowbreak{}CharacterBoxes.\allowbreak{}DesWin}{11--31}
\endgroup


\part{Input three: the witness for HypsAlt}

% ---- p8c1 ----
\chapter{A common witness for the 46 conditions}\label{ch:p8c1}

At scale \(x>F_m.x_1\), a family \(F_m\) of parameter \(\varepsilon\)
has points \(i\in I_x=F_m.\mathrm{Ix}(x)\), arithmetic and size data
\(D=D_{x,i}\), \(L=L_{x,i}\), object \(O=O_{x,i}\) and weight
\(b=\exp(x^\varepsilon)\). The coefficient field has embedding
\(\iota_D:K_D\hookrightarrow\C\). The bundle \(\beta\) ties the
arithmetic and spectral realizations to the rows and columns.
Write \(P\) for the unitary-group third-place inputs and \(F\) for
the local theta inputs. See \ref{B8.11}--\ref{B8.18} for the theta
integrals and \ref{B8.12}--\ref{B8.13} for their normalizations.

\section{Existence and the common data}

\begin{theorem}[Existence of the third input]\label{B8.01}
Assume \(P,F\), the theta-integral comparison and local matching
of \ref{B8.11}, and the covolume bound of \ref{B8.24}.
Fix a witness \(h_L\) that a local correspondence satisfying
\(\Lean{IsRecT}\) exists for every number field \(E\), algebraic
closure \(\Omega/E\), finite place \(w\) and valuation subring of
\(\Omega\) above \(w\).

For every \(c_{\mathrm{Dep}}\in\R\), assume the linked row on
\(p_0(h_L,c_{\mathrm{Dep}})\) with conclusion
\[
 \operatorname{Closed}_{\mathcal D_3(\mathcal A)}
 \bigl[(F_m,CH)\mapsto\operatorname{BoundAtTwoThreeFamilyFonAlt}(F_m)\bigr].
\]
\ref{B8.41} gives this row from \(P.\mathrm{nLink}\).
Then
\(\operatorname{HypsAlt}(\mathrm{NL},\operatorname{DefiniteSideCondition})\)
is inhabited.
\end{theorem}
\leanref{UpperTheory.\allowbreak{}Lit.\allowbreak{}nonempty\_\allowbreak{}hyps\_\allowbreak{}definite\_\allowbreak{}side\_\allowbreak{}of\_\allowbreak{}closed\_\allowbreak{}bound}{Development.\allowbreak{}AutomorphicArithmetic.\allowbreak{}ThetaCorrespondence.\allowbreak{}ThetaPairLifts.\allowbreak{}ThirdPlaceSeries.\allowbreak{}SiegelEisensteinSecondTermHypotheses}{44--68}

\begin{proof}
Unramified vacuum generation and convergent Siegel--Weil determine every
starting column by \ref{B8.15}. Apply \ref{B8.18} with \ref{B8.11} at
all admissible data over these columns to get \(\operatorname{RowRat}\).

Let \(C_F\) be the maximum of the grid threshold in \ref{B8.23}
from \ref{B8.22} and the lower-weight input, and the existence
threshold in \ref{B8.21}. At \(c_{\mathrm{Dep}}=C_F\),
\(\mathrm{GridA}(C_F)\) is eventually available. Transfer the open
rows from the list extension via the origin by \ref{B8.08}.
The projections to the origin and then to the law give the
functional row of \ref{B8.23} and the covolume row of \ref{B8.34}.
Apply \ref{B8.03} with the linked \(2,3\)-bound.
\end{proof}

Use the closed quantifiers of \ref{B8.39}, with point cutoffs
depending on \(F_m,CH\). Put
\[
 \mathcal W_1(\mathcal A)=
 \operatorname{WithTw}\bigl(\operatorname{WithSm}(\operatorname{WithChN}\mathcal A)\bigr),
 \qquad
 \mathcal D_3(\mathcal A)=\operatorname{D3}(3,\mathcal A).
\]
The added conditions of \(\operatorname{WithChN}\) are
\(\Lean{DGamDvdF}\), tame distinctness, \(\Lean{BoxesF}\),
\(\Lean{QDRowsF}\) and \(C_{YL}\le9\).
For character order \(\mathrm{ord}\) and rank increment \(t\), put
\(f=2\,\mathrm{ord}\) and \(N(j)=k+j(t+1)\). The Minkowski multiple is
\[
 M(n)=\prod_{\substack{p\le n+1\\p\ {\rm prime}}}
 p^{\sum_{j\ge0}\lfloor n/(p^j(p-1))\rfloor}.
\]
Eventually at every point and for every word \(u\),
\(\Lean{DGamDvdF}\) requires \(\Gamma_u=(D.\mathrm{mem}\,u).C.\mathrm{DGam}\)
to divide \(M(2fN(m))\), and \(\operatorname{WithSm}\) requires its prime
divisors to be at most \(2fN(m)+1\).
The operation \(\operatorname{WithTw}\) adds good twisting.
The class \(\mathcal D_3(\mathcal A)\) adds
\(\operatorname{ChosenE}(F_m,3,CH)\) for some \(CH\),
\(\operatorname{TuSize}(F_m)\) and \(\operatorname{DGamSize}(F_m)\).
The construction classes \(W_1,W_3\) coincide with
\(\mathcal W_1,\mathcal D_3\), respectively.

\begin{definition}[The common data]\label{B8.02}
Let \(N\) be a family-class predicate and \(DR(D,O,c,\beta,C,CH)\)
a predicate with \(c,C,CH\in\R\).
The \(46\)-field record \(\operatorname{HypsAlt}(N,DR)\) uses a
family class \(\mathcal A\) with raw selection
\[
 \sigma(F_m,h,x,h_x,i)\in
 \operatorname{BundleCondition}(D_{x,i},O_{x,i},F_m.C_{\mathrm{sz}}),
 \qquad h:\mathcal A(F_m),\quad h_x:F_m.x_1<x,\quad i\in I_x. \tag{8.2}
\]
Project the family selection from \(\sigma\) with cutoff \(F_m.x_1\).
Put \(\mathcal D_3=\mathcal D_3(\mathcal A)\) in the table.

\begin{longtable}{@{}p{0.21\textwidth}p{0.74\textwidth}@{}}
\toprule
Fields & Requirement\\
\midrule
\endhead
\(\mathcal A,\sigma,c_T\)
 & Family class, raw selection (8.2), and \(c_T:\R^2\to\R\).\\
\(\Lean{hEx}\)
 & Assume Minkowski division. For \(\lambda>1\) and an arithmetic
 requirement \(G\) supplied on its windows, put
 \(\mathcal A_G(F_m)=\mathcal A(F_m)\land G.\mathrm{cls}(\lambda,F_m)\).
 Require \(\operatorname{DesignChoices}(\operatorname{WithTw}(\mathcal A_G),3,c_T)\).
 For all sufficiently small \(\varepsilon>0\) and every threshold
 record \(\mathcal K\), choose \(F_m\) in the design class and
 \(CH\ge0,\ C_5,C_b>0\), with \(F_m.\operatorname{Large}(\mathcal K,CH)\)
 and \(C_{Tu}\ge c_T(\varepsilon,C_{YL})\).
 Require \(\Lean{ThickSetEventually},\Lean{ChosenTT},
 \Lean{ChosenBad}(C_b),\Lean{ChosenQD}(C_5)\),
 \(\Lean{ChosenCode}\) and \(\Lean{CodeUpF}\) at \(3,CH\), together
 with all the added conditions of \(\operatorname{WithChN}\) and \(F_m.C_F=CH\).\\
\(\Lean{hduty}\)
 & Seven duties for every \(F_m\in\mathcal A\). The fifth gives
 \(\Lean{ChosenWall}\) by \ref{B8.02a}.\\
\(\Lean{hRB},\Lean{hTw}\)
 & Bounded indefinite rows and twisting tied to the size datum.\\
\(\Lean{hPA},\Lean{hLamP}\), \(\Lean{hDivG},\Lean{hDcPtG}\)
 & The pointwise conditions in \ref{B8.08}.\\
\(\Lean{hTu}\)
 & Closed on \(\mathcal W_1(\mathcal A)\),
 \(C_{Tu}\ge c_T(\varepsilon,C_{YL})\Rightarrow\operatorname{TuSize}(F_m)\).\\
\(\Lean{hn4}\)
 & Closed on \(\mathcal W_1(\mathcal A)\), at every \(x>F_m.x_1\),
 \(\Lean{GBig}\) and splitting at \(g_0,g_1\) give \ref{B8.54}
 for every word \(u\) with \(T_u\) split.\\
\(\Lean{hLit},\Lean{hWu},\Lean{hRF}\)
 & \(\operatorname{LitThin}\) of \ref{B8.44}, Wu decay under the
 specified gap, and finite-dimensional row spaces.\\
\(\Lean{hTF},\Lean{sLp},\Lean{sO}\)
 & The level-prime row of \ref{B8.10} and its loop and distinct-word
 rows from \ref{B8.10a} and \ref{B8.10b}, respectively.\\
\(\Lean{c30},\Lean{hc30},\Lean{h30}\)
 & Constants \(c_{30}\) from the outer list with \(C_{30}\ge0\),
 \(A_{30}\ge2(C_{YL}+1)\) for every family. Closed on \(\mathcal D_3\),
 they give the idempotent counts of \ref{B8.52}.\\
\(\Lean{sTy},\Lean{sNG},\Lean{sM}\)
 & Type estimate, nongeneric-row estimate, diagonal identity.\\
\(\Lean{hInstR}\)
 & At each selected point, equal-rank spaces of the prescribed
 signatures and measures have equal fundamental-domain volumes.
 At rank \(\operatorname{rank}Y_0\), require genus and class mass
 and class finiteness.\\
\(\Lean{selJ}\)
 & Choose the compatible joins of \ref{B8.08}.\\
\(\Lean{jOn}\)
 & Closed on \(\mathcal D_3\), eventually
 \(\iota_{\mu_0}(V_0(\mu_0))\) is the whole target space.\\
\(\Lean{hMt}\)
 & For every \(F_m\in\mathcal A\), eventually all
 \(\widetilde M_\mu\) at all points are stable under \(\iota_D(\mathcal O_{K_D})\).\\
\(\Lean{rRows}\)
 & The simultaneous bounds (8.4).\\
\(\Lean{aPt0},\Lean{aPtTw}\), \(\Lean{aPer0},\Lean{aPerTw}\),
\(\Lean{tie},\Lean{aP10}\)
 & The six closed, eventual assertions of \ref{B8.09} in order.\\
\(\Lean{aAtG}\)
 & Closed on \(\mathcal D_3\), one \(C_1\) bounds the eventual
 used-tuple data and gives the family-size bound at type primes.\\
\(\Lean{aC2}\)
 & Closed on \(\mathcal D_3\), some \(C_{\deg},C_g,C_T,C_G,C_d,C_I,x_0\in\R\)
 with \(C_g\ge0\) give \(\Lean{SecondComplementInputs}\) at every
 \(x\ge x_0\) with \(x>F_m.x_1\).\\
\(\Lean{hServ}\)
 & \(\operatorname{SelectionServed}\) of \ref{B8.56} on \(\mathcal D_3\).\\
\(\Lean{cIc},\Lean{hcI},\Lean{irrC}\)
 & A constant \(c_{\rm Ic}>0\). Closed on \(\mathcal D_3\), eventually
 each point has arithmetic rows with injective good labels,
 compatible joins, and the cylinder-count error of \ref{B8.52}.\\
\(\Lean{zL}\)
 & Closed on \(\mathcal D_3\), eventually \(\mathrm{lPrj}_j\)
 acts surjectively on vectors at \(\mu_0\) for every \(0\le j<m\)
 at each selected point.\\
\(\Lean{hU},\Lean{hG}\)
 & \ref{B8.49} for each \(D\) satisfying \(\Lean{VolTie}\), and
 \ref{B8.50} for every \(D\).\\
\(\Lean{nRest}\)
 & \(N(\mathcal D_3)\).\\
\bottomrule
\end{longtable}

For \(\Lean{rRows}\) set \(c=F_m.C_{\mathrm{sz}}\) and
\(\beta=\sigma(F_m,h|_{\mathcal A},x,h_x,i)\), with \(h|_{\mathcal A}\)
the membership induced by \(h\). Require
\[
 \begin{split}
 \operatorname{Closed}_{\mathcal D_3}
 \bigl[(F_m,CH)\mapsto
 \forall h:\mathcal D_3(F_m)\ \exists C\ge0\quad
 &\text{eventually along the selection,}\\[-2pt]
 \operatorname{IndefBounds}\bigl(\beta.B.Q,C,CH;
 \operatorname{DIVinN}(F_m),\operatorname{LeafUniqueMap}(\beta)\bigr)
 &\land DR(D,O,c,\beta,C,CH)\bigr].
 \end{split} \tag{8.4}
\]
\(\operatorname{IndefBounds}\) denotes the nine conditions of \ref{B8.38}.
\end{definition}
\leanref{GlobalTheory.\allowbreak{}Fc.\allowbreak{}HypsAlt}{Development.\allowbreak{}AutomorphicArithmetic.\allowbreak{}ThetaCorrespondence.\allowbreak{}ThetaPairLifts.\allowbreak{}ThirdPlaceSeries.\allowbreak{}Hyps}{22--90}
\leanref{GlobalTheory.\allowbreak{}Fc.\allowbreak{}R1.\allowbreak{}SelectionDefiniteRows}{Development.\allowbreak{}AutomorphicArithmetic.\allowbreak{}ThetaCorrespondence.\allowbreak{}ThetaPairLifts.\allowbreak{}ThirdPlaceSeries.\allowbreak{}CountsIrreducibilityConstants}{91--93}

\begin{proposition}[The fifth construction duty]\label{B8.02a}
Let \(\mathcal A\) be the class constructed in \ref{B8.05}.
For \(F_m\in\mathcal A\), \(x>F_m.x_1\) and \(i\in F_m.I_x\),
the model's own weights satisfy \(R_W(i)\le x^{-9}\).
In particular \(x_W=F_m.x_1+1\) works in \ref{B5.22b}.
\end{proposition}
\begin{proof}
Fix a request with \(\lambda>1\). On its supplied windows take
an arithmetic host \(H\) with \(TT\) as in \ref{B2.54}.
Choose integer-valued functions \(a,b\) in the weight boxes and set
both to zero off \(TT\). Fix an injective word code. Construct a valid
canonical point for each \(0\le u_0<n\). These points have common host
data and word labels and cover all translations. Choose the parameters
of this preliminary family above the size threshold of \ref{B5.22a}.
At each window scale beyond its cutoff, apply \ref{B5.22a} at any index
to get a trial pair \(A,B\).

With window datum \(s=(x,\varepsilon,J)\) and fixed size parameters
\(c_D,c_h\), the host share uses the tuple
\[
 (R,s,q,d,\mathrm{cls},TT,S_c,\eta),\qquad
 S_c=4k\operatorname{Delv}(s,c_D),\quad
 \eta=m\operatorname{h0v}(s,c_h).
\]
The point has these same data and \(h_p\), so its candidate tests
agree with the host's. Both use the least integer label in each
coset of \(\langle h_p\rangle\). Fix the common constants
\(K=F_m.\mathrm{GK}\). For each good representative \(t_0\), the
marked positions and restricted arrays agree. Formula \ref{B3.51}
gives the same \(\kappa_v(t_0)\) for \(p<p_0\). For \(p\ge p_0\)
both values are infima of the same admissible set. Hence at every
candidate prime the two sums \(\Sigma_p(A,B)=m\sum_{t_0}\kappa_v(t_0)\)
over the good representatives agree.
The two box memberships transfer as well, and the chosen pair
gives the host bound \(\Lean{WOK}\).
\leanref{UpperTheory.\allowbreak{}Top.\allowbreak{}hostWW}{Development.\allowbreak{}AutomorphicArithmetic.\allowbreak{}CharacterBoxes.\allowbreak{}HostWW}{27--87}
\leanref{GlobalTheory.\allowbreak{}EX.\allowbreak{}weightWindowRequirement\_\allowbreak{}host\_\allowbreak{}of\_\allowbreak{}point}{Development.\allowbreak{}AbelianVarietyTheory.\allowbreak{}Proofs.\allowbreak{}GlobalInterpolation.\allowbreak{}WeightWindowRequirementHostOfPoint}{17--54}

Construct new points over the same host with \(A,B\) and the
common word code at every translation. Their own weights, extended
by zero, are \(A,B\). The equality of shares gives \(R_W(i)\le x^{-9}\)
with the summand in (5.19.1). The box memberships complete
\(\Lean{WallPt}\). The class condition in \ref{B8.05} retains
\(\Lean{WallPt}\) at every point above \(F_m.x_1\).
\end{proof}
\leanref{GlobalTheory.\allowbreak{}EX.\allowbreak{}wreq\_\allowbreak{}of\_\allowbreak{}wok}{Development.\allowbreak{}AbelianVarietyTheory.\allowbreak{}Proofs.\allowbreak{}GlobalInterpolation.\allowbreak{}WeightRequirementOfValid}{32--69}
\leanref{UpperTheory.\allowbreak{}Top.\allowbreak{}kitS}{Development.\allowbreak{}AutomorphicArithmetic.\allowbreak{}CharacterBoxes.\allowbreak{}KitS}{14--62}
\leanref{UpperTheory.\allowbreak{}Top.\allowbreak{}designS}{Development.\allowbreak{}AutomorphicArithmetic.\allowbreak{}CharacterBoxes.\allowbreak{}DesignS}{21--98}
\leanref{GlobalTheory.\allowbreak{}EX.\allowbreak{}DQ.\allowbreak{}duty\_\allowbreak{}seven\_\allowbreak{}library}{Development.\allowbreak{}AbelianVarietyTheory.\allowbreak{}Proofs.\allowbreak{}GlobalInterpolation.\allowbreak{}DutySevenLibrary}{15--23}

Use \(DR\) from \ref{B8.35}. Under its hypotheses, let
\(\mathrm{NuC}\) be any system of integral face submodules and
\(I\) any three-sided division datum tied to it that agrees with
\(\beta.P.I\) under transport. Put \(e_u=I.eNu(u)\). For every word
\(u\), \(\Lean{RowsU}\) holds at \(I.w_U(u)\). It gives a tied
column chain and makes \(e_u^{-1}(\mathrm{NuLat}_u)\) a lattice.
Put \(\mathrm{NL}(\mathcal C)=
\operatorname{Closed}_{\mathcal C}[(F_m,CH)\mapsto\operatorname{LinkN}(F_m)]\).
The finiteness premise in \ref{B8.40} ranges over all primes.

\section{From member parameters to the common data}

For \(h_L,c_{\mathrm{Dep}}\) and a point condition \(A\), put
\[
 \begin{split}
 \mathcal N_{\rm orig}&=\operatorname{extensionWithOrigin}(h_L,c_{\mathrm{Dep}}),\\
 \mathcal N_A&=\operatorname{extensionWithAtProperty}(h_L,\mathcal N_{\rm orig},A),\\
 p_A&=\operatorname{parametersOfExtension}
 (\mathcal N_A,\mathrm{none},\mathrm{geoCl}(\mathrm{anyGeometry})).
 \end{split} \tag{8.6}
\]
Write \(Q_A,Q'_A\) for the constructed and good member conditions
on the slots \(p_A.s\), and \(L_A\) for their linkage.
For \(p_0(h_L,c_{\mathrm{Dep}})\), replace \(\mathcal N_A\) in the
last line of (8.6) by \(\operatorname{extensionWithList}(h_L,c_{\mathrm{Dep}})\).
The condition \(\operatorname{EvLgG}(A)\) makes \(A\) eventual on
the base, frame and twist.

\begin{theorem}[Assembly on an augmented extension]\label{B8.03}
Assume \(P\) and the rational row tests \(\operatorname{RowRat}\)
of \ref{B8.18}. Fix \(h_L\) as in \ref{B8.01}, \(c_{\mathrm{Dep}}\in\R\),
and \(A\) satisfying \(\operatorname{EvLgG}(A)\). Assume the linked
\(2,3\)-bound of \ref{B8.01} at \(c_{\mathrm{Dep}}\).
Along \(p_A\), assume two linked rows closed on \(W_3(\mathcal A)\).
For each eligible \(F_m,CH\), the covolume row gives \(C_0\ge0\)
such that for every \(C\ge C_0\) the condition \(\Lean{L3}\) at
\(C,CH\) holds eventually at the selected points. Under
\(\operatorname{DIVinN}(F_m)\), the functional row gives the same
quantifiers for \(\Lean{FUN}\). Each point cutoff can depend on
\(C\). Then \(\operatorname{HypsAlt}(\mathrm{NL},\operatorname{DefiniteSideCondition})\)
is inhabited.
\end{theorem}
\leanref{GlobalTheory.\allowbreak{}Fc.\allowbreak{}S0.\allowbreak{}StageU.\allowbreak{}hypotheses\_\allowbreak{}nonempty\_\allowbreak{}of\_\allowbreak{}unramified\_\allowbreak{}exponents\_\allowbreak{}alt}{Development.\allowbreak{}AutomorphicArithmetic.\allowbreak{}ThetaCorrespondence.\allowbreak{}ThetaPairLifts.\allowbreak{}ThirdPlaceSeries.\allowbreak{}HypothesesNonemptyOfUnramifiedExponents}{41--61}

\begin{proof}
\ref{B8.07} supplies (8.14). The literature inputs in \(P\) and its
geometric Galois hypotheses give the third-layer and series hypotheses
of \ref{B8.43}. Combine them with the norm input of \(P\), depth-zero
support, local finiteness, Tate duality, the local Euler characteristic
and the linked \(2,3\)-bound. Then \ref{B8.43} gives the linked row
\[
 \mathcal A\longmapsto\mathrm{NL}(\mathcal D_3(\mathcal A)). \tag{8.7}
\]
Transfer this and the open rows via the origin by \ref{B8.08}.
Apply \ref{B8.06} with \(NN=\mathrm{NL}\), the functional and
covolume rows, and the surjectivity and lower-weight inputs of \(P\).
For \ref{B8.04}, take \(\operatorname{LitThin}\) from \ref{B8.44},
Wu decay, finite dimensionality and the type estimate. Use the
second convergent unitary Siegel--Weil input for the nongeneric
and diagonal rows and for \ref{B8.49}. The remaining mass and
supply inputs are \ref{B8.50} and (8.14).
\end{proof}

\begin{theorem}[Filling the common record from all rows]\label{B8.04}
Let \(KP\) be a condition on a base and its frame, \(Q,Q'\) member
conditions, and \(L_k\) a bundle linkage. Fix \(N,DR\) as in
\ref{B8.02}. Assume
\begin{enumerate}
 \item the adjacent-rank mass formula of \ref{B8.48};
 \item \(\operatorname{LitThin}\), Wu decay, finite dimensionality of
 the row spaces, the type and nongeneric-row estimates, and the
 diagonal main-term identity;
 \item the volume ratio of \ref{B8.49} for each \(D\) satisfying
 \(\Lean{VolTie}\), and the positive-definite rank-two genus
 mass of \ref{B8.50} for every \(D\);
 \item that Minkowski division implies \(\operatorname{BaseCond}(KP)\),
 \(\operatorname{FirstG}(KP,Q)\), \(\operatorname{ExportF}(KP,Q)\),
 and \(\operatorname{SecondG}(KP,Q,Q')\).
 \(\Lean{BaseCond}\) supplies \(KP\)-bases with specified finite
 sets of algebraic elements and preserves \(KP\) under
 frame-preserving enlargement. On untwisted \(KP\)-bases,
 \(\Lean{FirstG}\) gives a \(Q\)-member if the required primes
 are present, and \(\Lean{ExportF}\) gives export data for every
 \(Q\)-member. From a member with export data \(\Lean{SecondG}\)
 enlarges the base with the same frame and prime set to meet
 new avoidance and coefficient conditions. It gives one member
 satisfying \(Q\), twisting compatibility, and \(Q'\).
 \item \(\operatorname{ExBL}(KP,Q,Q',L_k)\) and
 \(\operatorname{AllRowsAlt}(KP,Q,Q',L_k,N,DR)\), including the
 definite-side row for \(DR\).
\end{enumerate}
Then \(\operatorname{HypsAlt}(N,DR)\) is inhabited with
\[
 c_T(\varepsilon,y)=\max\bigl(11,\ 16+3\max(5,y+1)\bigr). \tag{8.8}
\]
\end{theorem}
\leanref{GlobalTheory.\allowbreak{}Fc.\allowbreak{}nonempty\_\allowbreak{}hyps\_\allowbreak{}of\_\allowbreak{}allRows\_\allowbreak{}alt}{Development.\allowbreak{}AutomorphicArithmetic.\allowbreak{}ThetaCorrespondence.\allowbreak{}ThetaPairLifts.\allowbreak{}ThirdPlaceSeries.\allowbreak{}NonemptyHypsOfAllRows}{19--101}

\begin{proof}
Fix \(c_{30},c_{\rm Ic}\) and strengthen the point condition for
their counts. Item~3 gives the \(T_u\)-row of \ref{B8.51} with
(8.8) and point condition \(\top\). Apply \ref{B8.05} to the
conjunction \(H\) of these and the other member rows to obtain
\[
 \mathcal A,\qquad \sigma,\qquad H(\mathcal A,\sigma).
 \tag{8.9}
\]
Choose the joins in \(H\) and transfer the closed rows by the
working-class equivalences. The design and duties give
\(\Lean{hEx},\Lean{hduty}\). Lattice stability and the analytic
and mass hypotheses give the remaining fields.
\end{proof}

\section{Choosing the family after imposing the rows}

Let \(L_k\) link constructed points to bundles of type
\(\mathrm{Bd}(D,O,c)\). \(\operatorname{LargeBaseCondition}(KP,Q,R)\)
requires \(R\) at every \(KP,Q\)-point beyond the construction thresholds.
The class \(\mathrm{cls}_7(U)\) adds \(U(D_{x,i},O_{x,i})\) for
all \(x>F_m.x_1\) and \(i\in I_x\) to \(\mathrm{cls}_7(\top)\).

A row \(\operatorname{Row}(Q,L_k,H)\) gives a bundle condition
\(E(\kappa,D,O,c,\beta)\), with construction parameters \(\kappa\).
The condition \(\operatorname{RowE}(E)\) gives \(E\) on every
linked bundle at \(Q\)-points beyond the parameter and scale thresholds.
If each \(F_m\in\mathcal A\) belongs to \(\mathrm{cls}_7(\top)\)
and has one linked \(\kappa\) for which every selected bundle
comes from a linked \(Q\)-point and satisfies \(E\), the row
gives \(H(\mathcal A,\sigma)\). Put
\[
 \operatorname{RowL}(KP,Q,Q',L_k,H)
 \ \Longleftrightarrow\
 \exists R\ \bigl(\operatorname{LargeBaseCondition}(KP,Q,R)
 \land\operatorname{Row}(Q'\land R,L_k,H)\bigr). \tag{8.10}
\]
In \(\operatorname{ExBL}(KP,Q,Q',L_k)\), the last conjunct of
(8.10) is \(\operatorname{ExB}(Q'\land R,L_k)\).
It supplies a linked bundle for every \(c\in\R\) at every
\(Q'\land R\)-point beyond the parameter and scale thresholds.

\begin{theorem}[A common family and selection]\label{B8.05}
Assume the mass formula of item~1 of \ref{B8.04} and the volume
ratio of \ref{B8.49} for each \(D\) satisfying \(\Lean{VolTie}\).
Fix \(KP,Q,Q',\mathrm{Bd},L_k\), a property \(H(\mathcal A,\sigma)\)
of a family class and raw selection, and \(c_T:\R^2\to\R\).
Assume the implication from Minkowski division to the four member
conditions of \ref{B8.04}, and suppose \(\operatorname{ExBL}\) and
\(\operatorname{RowL}(H)\) hold for these data. There exist
\(\mathcal A,\sigma\) such that
\begin{enumerate}
 \item Minkowski division implies \(\operatorname{ExD}(\mathcal A,c_T)\),
 the construction-class form of the design supply in \(\Lean{hEx}\).
 \item Each \(F_m\in\mathcal A\) satisfies the seven construction duties
 and has stable coefficient lattices.
 \item Each \(F_m\in\mathcal A\) satisfies \(\mathrm{cls}_7(\top)\)
 and has one linked parameter \(\kappa\). For all \(x>F_m.x_1\)
 and \(i\in I_x\), the selected bundle is built from a linked
 \(Q'\)-point with this same \(\kappa\).
 \item \(H(\mathcal A,\sigma)\) holds.
\end{enumerate}
\end{theorem}
\leanref{UpperTheory.\allowbreak{}Top.\allowbreak{}classSel}{Development.\allowbreak{}AutomorphicArithmetic.\allowbreak{}CharacterBoxes.\allowbreak{}ClassSel}{14--29}

\begin{proof}
For two rows intersect the additional point and bundle conditions.
Take the minimum \(\varepsilon\)-cutoff and componentwise maximum
size threshold. Take the maximum for \(c_D\), then for \(c_h\)
at this \(c_D\). Take the scale maximum last. Unite the finite sets of required
algebraic elements and required or avoided primes. Induction gives
finite conjunctions.

Intersect the strengthenings from \(\operatorname{ExBL}\) and
(8.10). For the resulting \(R\) put \(Q^*=Q'\land R\). Then
\[
 \operatorname{ExB}(Q^*,L_k),\qquad
 \operatorname{Row}(Q^*,L_k,H). \tag{8.11}
\]
Let \(E\) be the bundle condition of this row. Write \(\kappa.\mathrm{Pos}\)
for \(C_5,c_6,C_b>0\) in the parameter record \(\kappa\). Define
\(\mathcal Q_E(\kappa,D,O)\) by
\[
 \kappa.\mathrm{Pos}\ \Longrightarrow\
 \forall c\in\R\ \exists\beta\in\mathrm{Bd}(D,O,c)\quad
 \begin{cases}
 \beta\text{ is built from a linked }Q^*\text{-point with }\kappa,\\
 E(\kappa,D,O,c,\beta).
 \end{cases} \tag{8.12}
\]
Beyond both cutoffs, \(\operatorname{ExB}\) supplies the bundle
and \(\operatorname{RowE}\) gives \(E\). Thus \(\mathcal Q_E\)
holds eventually on \(Q^*\)-points.

Set
\[
 \mathcal A(F_m)\ \Longleftrightarrow\
 \exists\kappa\,\bigl(\operatorname{Link}(\kappa,F_m)\land
 \mathrm{cls}_7(\mathcal Q_E(\kappa))(F_m)\bigr).
\]
Choose \(\kappa\) once per family. The link gives
\(\kappa.\mathrm{Pos}\), so choose \(\sigma\) by (8.12) at
\(c=F_m.C_{\mathrm{sz}}\). Then (8.11) gives \(H\) and
\(Q^*\Rightarrow Q'\) gives item~3. The class gives the duties,
with the fifth as in \ref{B8.02a}, and lattice stability.

Under Minkowski division the eventual \(R\) strengthens the second
member condition of \ref{B8.04} to \(Q^*\). With the two volume
inputs, the eventual point property gives \(\operatorname{RecExP}\)
for \(\mathcal Q_E(\kappa)\land QD\) by (8.12), where \(QD\) requires
tuple use, twisting compatibility and coefficient-lattice stability.
Apply \(\Lean{recS}\) for \(\operatorname{ExD}(\mathcal A,c_T)\).
\end{proof}
\leanref{GlobalTheory.\allowbreak{}Fc.\allowbreak{}RowL.\allowbreak{}and}{Development.\allowbreak{}AbelianVarietyTheory.\allowbreak{}Definitions.\allowbreak{}GlobalTwisting.\allowbreak{}EventuallyAllScaleContextsAndOfBoth}{76--81}
\leanref{GlobalTheory.\allowbreak{}Fc.\allowbreak{}Row.\allowbreak{}elim}{Development.\allowbreak{}AbelianVarietyTheory.\allowbreak{}Definitions.\allowbreak{}GlobalTwisting.\allowbreak{}ValidForConstructedPointClassesEventuallyWorking}{364--367}

\section{The member rows}

For a property \(r\) of a member condition, put
\[
 \operatorname{Lg1}(KP,Q,Q';r)
 \quad\Longleftrightarrow\quad
 \exists R\ \bigl(\operatorname{LargeBaseCondition}(KP,Q,R)
 \land r(Q'\land R)\bigr). \tag{8.13}
\]
Along \(p\), take \(Q=\operatorname{IsBuiltE}(p.s)\),
\(Q'=\operatorname{IsGoodE}(p.s)\) and \(L_k=\operatorname{BundleLinkage}(p)\).
The condition \(KP_x(B,f)\) requires \(B.\mathrm{KsExact}(\delta,f)\)
for some field data \(\delta\). The standard embedding to level \(m\)
must carry each starting frame vector to its counterpart.

\begin{theorem}[The twenty-six rows]\label{B8.06}
Assume global Artin reciprocity, the second convergent unitary
Siegel--Weil input, discrete-series cover and classification,
Harish--Chandra modules, cusp decay, the Casselman--Wallach input,
surjectivity onto the point space and the lower-weight bound.

Fix \(h_L\) as in \ref{B8.01}, a norm-choice extension \(\mathcal N\),
an extra extension \(S\) and a geometry condition \(\Gamma\).
Along \(p=\operatorname{parametersOfExtension}(\mathcal N,S,\Gamma)\)
assume the open conditions of \ref{B8.08}, the two rows of \ref{B8.03},
and the linked row \(\mathcal A\mapsto NN(\mathcal D_3(\mathcal A))\)
for a predicate \(NN\) on family classes. Then
\(\operatorname{AllRowsAlt}(KP_x,Q,Q',L_k,NN,\operatorname{DefiniteSideCondition})\)
holds for the member data of \(p\). Its \(26\) rows have form (8.13).
Where applicable they close on \(W_1,W_3\). The join type is nonempty,
and the counting row supplies both estimates with their constants.
\end{theorem}
\leanref{GlobalTheory.\allowbreak{}Fc.\allowbreak{}AllAtMemberStage0.\allowbreak{}allRows\_\allowbreak{}of\_\allowbreak{}open0\_\allowbreak{}alt}{Development.\allowbreak{}AutomorphicArithmetic.\allowbreak{}ThetaCorrespondence.\allowbreak{}ThetaPairLifts.\allowbreak{}ThirdPlaceSeries.\allowbreak{}LinkROfIsLinked}{50--86}
\leanref{GlobalTheory.\allowbreak{}Fc.\allowbreak{}AllRowsAlt}{Development.\allowbreak{}AutomorphicArithmetic.\allowbreak{}ThetaCorrespondence.\allowbreak{}ThetaPairLifts.\allowbreak{}ThirdPlaceSeries.\allowbreak{}AllRows}{12--39}

\begin{proof}
\ref{B8.09} and \ref{B8.10} give the arithmetic and spectral rows.
Use the class-number formula \ref{B8.54} for \(\Lean{hn4}\)
and the hyperplane volume ratio for \(\Lean{hInstR}\).
The second complement gives \(\Lean{aC2}\). For \(\Lean{aAtG}\),
\(p.s.R_x=RX\) identifies the row forms with the period map
through the columns. Take \(R=\top\) in (8.13).

Keep the five pointwise rows of \ref{B8.08} named in \ref{B8.02}.
Primitivity gives \(\Lean{jOn},\Lean{zL}\), and the comparison
gives \(\Lean{hTw}\). Use \ref{B8.55} for \(\Lean{hRB}\) and
\ref{B8.56} with reciprocity and linkage for \(\Lean{hServ}\).
Apply \ref{B8.52} with \(\Lean{lab}\) and Siegel--Weil for
\(\Lean{c30}\). The norm row gives \(\Lean{nRest}\).

For \(\Lean{rRows}\), conjoin \(\Lean{hPA},\Lean{hDivG},
\Lean{good},\Lean{geo}\), the covolume and functional rows, and
\(\Lean{mi},\Lean{kap},\Lean{df}\). Apply \ref{B8.38} on the
common strengthening with the compatible summand choices.
The \(\Lean{ON}\) and \(\Lean{FUN}\) conditions use
\(\Lean{DIVinN}\), and \(\Lean{L6}\) uses \(\Lean{LeafUniqueMap}\).
Combine the constant and scale thresholds to obtain (8.4).
\end{proof}

\begin{theorem}[Member and bundle supply on the same extension]\label{B8.07}
Assume \(P\) and \(\operatorname{RowRat}\). For every \(h_L\) as in
\ref{B8.01}, \(c_{\mathrm{Dep}}\in\R\), and \(A\) satisfying
\(\operatorname{EvLgG}(A)\), the parameters in (8.6) satisfy
\[
 \operatorname{SiegelMassRatioCM},\qquad
 \operatorname{MemberE}(p_A.s),\qquad
 \operatorname{ExBL}(KP_x,Q_A,Q'_A,L_A). \tag{8.14}
\]
The member assertion consists of the four member conditions of
\ref{B8.04}.
\end{theorem}
\leanref{GlobalTheory.\allowbreak{}T.\allowbreak{}large\_\allowbreak{}unramified\_\allowbreak{}exponent\_\allowbreak{}bundles\_\allowbreak{}exist\_\allowbreak{}alt}{Development.\allowbreak{}AutomorphicArithmetic.\allowbreak{}ThetaCorrespondence.\allowbreak{}ThetaPairLifts.\allowbreak{}ThirdPlaceSeries.\allowbreak{}MemberBundlesLargeUnramifiedExponents}{69--70}
\leanref{GlobalTheory.\allowbreak{}Fc.\allowbreak{}StageOneU.\allowbreak{}member\_\allowbreak{}bundles\_\allowbreak{}large\_\allowbreak{}unramified\_\allowbreak{}exponents\_\allowbreak{}alt}{Development.\allowbreak{}AutomorphicArithmetic.\allowbreak{}ThetaCorrespondence.\allowbreak{}ThetaPairLifts.\allowbreak{}ThirdPlaceSeries.\allowbreak{}MemberBundlesLargeUnramifiedExponents}{34--63}

\begin{proof}
Artin reciprocity, CM curves and ray fields, real representation
theory and Siegel--Weil give the compatible-system hypotheses \(h_A\).
Its mass ratio is \(\Lean{hA.\allowbreak{}a18}\). The theta and doubling inputs
with the rational row tests give the multiplicity cover.

After saturation each eligible starting column eventually has a
state \(P_w\) and extension \(s_w\) preserving the starting data,
with empty obstruction type and all \(\eta_{j,v}(s_w)=0\).
Untwisted \(P_w.Wt\) is closed under lying over the same weight-one place.
Multiplicative closure gives multiplicity on the origin extension.
Use the elliptic and local conditions from geometric Galois, with
its local correspondence witness identified with \(h_L\) by proof
irrelevance. The base-lattice construction gives origin member supply.

Write \(Q_{\rm orig}\) for the origin member condition. For a point
\(X\) with base \(X.B\), frame \(f\), twist \(T\) and auxiliary
datum \(c\) in scale context \(K\),
\[
 Q_A(K,X,f,T,c)\ \Longleftrightarrow\
 Q_{\rm orig}(K,X,f,T,c)\land A(K,X.B,f,T). \tag{8.15}
\]
This also holds for good members. Inherit the base-field clause
and export, and intersect the initial-member thresholds with those
for \(A\). Retwisting gives \(A\) on the new data and (8.15) gives
the new members.

The origin slots satisfy \(\Lean{PositionMv}\) and the origin
records project through the list to its law. With \(A\), the
compatible system, the multiplicity cover and the analytic inputs
of \(P\), the Weil block construction gives bundle supply on
\(\mathcal N_A\).
\end{proof}
\leanref{GlobalTheory.\allowbreak{}Mem.\allowbreak{}LiteralAutomorphy.\allowbreak{}memberE\_\allowbreak{}extensionWithAtProperty\_\allowbreak{}of\_\allowbreak{}memberE}{Development.\allowbreak{}AutomorphicArithmetic.\allowbreak{}ThetaCorrespondence.\allowbreak{}ThetaPairLifts.\allowbreak{}ThirdPlaceSeries.\allowbreak{}IsBuiltEPrmAIff}{35--53}

\begin{theorem}[The twelve open conditions]\label{B8.08}
Assume \(P\). For every \(h_L\) as in \ref{B8.01} and
\(c_{\mathrm{Dep}}\in\R\), the parameters \(p_0(h_L,c_{\mathrm{Dep}})\)
satisfy the linked rows
\[
\begin{array}{c|p{0.76\textwidth}}
 \Lean{hPA} & Primitive rows and separated types.\\
 \Lean{hLamP} & An assignment of finite types to row indices. For each
 layer and row index the assigned isotypic spaces form a direct sum
 equal to the transported row space. At used tuples the nonzero
 components have the prescribed \(\lambda\). For each row index the
 sum of its assigned row counts equals the base row dimension.\\
 \Lean{hDivG} & Compatibility of both \(RR\) lattices with their
 indefinite realizations, rational-lift spanning at the used tuples,
 and agreement of the theta maps, numerical data and column pairings.\\
 \Lean{hDcPtG} & At every used tuple the embedded definite-column
 evaluation equals the arithmetic column period on transported inputs.\\
 \Lean{selJ} & Existence of a compatible join at every point.\\
 \Lean{tie} & Existence of a compatible indefinite recursion at every point.\\
 \Lean{lab} & A Hida realization \(X\), good operator \(G\), and
 isometry \(\iota\) from the spectral space to the model space,
 satisfying \(\Lean{AutTie},\Lean{RowDictTie},\Lean{LabTie}\).\\
 \Lean{good} & Eventually \(\Lean{GoodFTAll}\) at the selected points.\\
 \Lean{geo} & Compatible point summands with class-sum pinning,
 recursion ties and an evaluation grid.\\
 \Lean{mi} & Eventual multiplicity identity.\\
 \Lean{kap} & Eventually \(Y.\mathrm{kapR}(e)=P.S_1.\kappa(v_e)\)
 for every embedding \(e:K_D\hookrightarrow\C\) with infinite
 place \(v_e\). Here \(Y\) is the division datum and \(P=\beta.P\)
 the row-side family.\\
 \Lean{df} & Eventual definite-side condition for each \(C\ge C_0\).
\end{array} \tag{8.16}
\]
The rows \(\Lean{good},\Lean{mi},\Lean{kap},\Lean{df}\) close on
\(W_3(\mathcal A)\). For \(\Lean{df}\) choose \(C_0\ge0\) per
eligible \(F_m,CH\), then a point cutoff per \(C\ge C_0\).

For extensions with the same \(h_L,S,\Gamma\), a pointwise map
\(N'\to\operatorname{Nonempty}(N)\) at the same indices transfers
every linked row with conclusion \(H(\mathcal A,\sigma)\) from
\(N\) to \(N'\).
\end{theorem}
\leanref{GlobalTheory.\allowbreak{}T.\allowbreak{}open0\_\allowbreak{}of\_\allowbreak{}unitary\_\allowbreak{}group\_\allowbreak{}hypotheses}{Development.\allowbreak{}AutomorphicArithmetic.\allowbreak{}ThetaCorrespondence.\allowbreak{}ThetaPairLifts.\allowbreak{}ThirdPlaceSeries.\allowbreak{}OpenOrbitOfUnramifiedExponents}{75--76}
\leanref{GlobalTheory.\allowbreak{}Fc.\allowbreak{}S0.\allowbreak{}Open0ProofU.\allowbreak{}open\_\allowbreak{}orbit\_\allowbreak{}of\_\allowbreak{}unramified\_\allowbreak{}exponents\_\allowbreak{}alt}{Development.\allowbreak{}AutomorphicArithmetic.\allowbreak{}ThetaCorrespondence.\allowbreak{}ThetaPairLifts.\allowbreak{}ThirdPlaceSeries.\allowbreak{}OpenOrbitOfUnramifiedExponents}{39--69}

\begin{proof}
The parameters give \(\Lean{hPA},\Lean{hDcPtG}\). Weil
representation and discrete series over the law give \(\Lean{tie}\).
These rows give the joins. The theta main theorem with
vacuum generation, doubling and strong approximation gives the row
identity and \(\Lean{mi}\). The identity and divisor law give
\(\Lean{hDivG}\). The real-group and Harish--Chandra constructions
give \(\Lean{hLamP}\) and then \(\Lean{kap}\).

Green's prime-field formula gives the polynomial parametrization
\cite{green-1955} by \ref{B7.37}. The Zelevinsky input and projection
to the law give the finite types and labels. The comparison and
divisor rows with the ray-class-field input give \(\Lean{geo}\).
Use \ref{B8.37} for \(\Lean{df}\).

Project the constructed and good points with their linkage.
Pull back the additional point condition at the same thresholds.
A selection over \(N'\) then satisfies the selection premise over \(N\).
\end{proof}
\leanref{GlobalTheory.\allowbreak{}Glob.\allowbreak{}open\_\allowbreak{}and\_\allowbreak{}linked\_\allowbreak{}of\_\allowbreak{}stronger\_\allowbreak{}extension\_\allowbreak{}alt}{Development.\allowbreak{}AutomorphicArithmetic.\allowbreak{}ThetaCorrespondence.\allowbreak{}ThetaPairLifts.\allowbreak{}ThirdPlaceSeries.\allowbreak{}OpenAndLinkedOfStrongerExtension}{28--33}

\section{Evaluation, periods, and the selected spectrum}

\begin{theorem}[Six arithmetic rows]\label{B8.09}
For any slots \(\mathcal S\) and linkage \(L_k\), take \(KP=KP_x\),
\(Q=\operatorname{IsBuiltE}(\mathcal S)\) and
\(Q'=\operatorname{IsGoodE}(\mathcal S)\). The six
\(\operatorname{Lg1}\)-rows are closed on \(W_3(\mathcal A)\)
and eventual at selected points.
\begin{enumerate}
 \item For each word \(u\), reference evaluation is injective on the
 specified column labels and induces bijections on their
 stabilizers. Its complex point map \(\mathrm{Pt}_u\) takes
 \(F_i(m)\) into \(\mathrm{ptC}_u\). There is an integral
 rational realization \(P_u\) such that for every infinite
 place \(v\) and rational level-\(m\) vector \(z\),
 \[
 q_{u,\mathrm{pt},v}(P_u z)
 \le\mathrm{ptB}(D,u)\,q_{u,m,v}(z),\qquad
 \mathrm{ptB}(D,u)=\prod_\tau T_u^{-\deg_\tau}\mathrm{beOn}_\tau^2.
 \tag{8.17}
 \]
 Here \(q_v\) is the real diagonal of the majorant, and
 \(\deg_\tau,\mathrm{beOn}_\tau\) are the chain degree and norm
 factor at archimedean type \(\tau\). The same constant gives
 \(\|\mathrm{Pt}_u\phi\|^2\le\mathrm{ptB}(D,u)\|\phi\|^2\)
 on the complex test space \(F_i(m)\).
 \item Each used nonreference tuple has a rational integral twisted
 evaluation for every word.
 \item Assume \(\Lean{DGamDef}\). For every \(u,\mu\), the twisted
 integer equals \(\Gamma_{u_0}\) for the distinguished word \(u_0\).
 Each such integer is a multiple of all class stabilizer orders
 at levels \(0\le j\le m\) for all words and tuples.
 For each \(u\), each \(z\) in its reference
 point lattice, and \(H\in\widetilde M_{\mu_0}\cap V_0(\mu_0)\),
 there is one \(y\in K_D\) representing the period through
 \(\iota_D\). For each finite place \(w\) put
 \(p=\operatorname{resChar}(w)\) and \(vDG_{u,w}=v_p(\Gamma_u)\).
 With ramification index \(e_w\), require
 \[
 y=0\quad\text{or}\quad v_w(y)\ge-\mathrm{ur}_w(vDG_{u,w}),
 \qquad \mathrm{ur}_w(a)=e_w a. \tag{8.18}
 \]
 \item Under the same \(\Lean{DGamDef}\), take \(z\) in the point
 lattice of any used tuple \(\mu\ne\mu_0\) and
 \(H\in\widetilde M_\mu\cap V_0(\mu)\). One representative
 \(y\in K_D\) of the twisted period satisfies (8.18) at every
 \(w\) with \(\operatorname{resChar}(w)\notin Q_p(\mu)\).
 \item Each point satisfies \(\Lean{RealTied}\) and admits a compatible
 symmetric column system. For every word its level-zero
 character has finite exponent and takes values in
 \(\iota_D(K_D)\).
 \item Let \(\operatorname{resChar}(w)=T_u\) and
 \(E^T_{u,w}=E^{\rm atT}_{u,w}\), the loss in \ref{B4.22}.
 For each input in \(N_{\lambda,u,w}\) and \(c\in K_D\) of
 finite integral valuation \(v_w(c)\ge e_wmE^T_{u,w}\),
 scaling the reference column by \(\iota_D(c)\) puts it
 in the image of the local lattice \(S_m(w)\).
 For each used \(\mu\ne\mu_0\) with
 \(\operatorname{resChar}(w)\notin Q_p(\mu)\), the twisted
 column of each such input is represented under its level-\(m\)
 embedding by \(az\). Here \(z\) lies in the twisted level-\(m\)
 lattice localized at \(w\), and \(a\in K_D\) is zero or
 satisfies \(v_w(a)\ge-e_wmE^T_{u,w}\).
\end{enumerate}
\end{theorem}
\leanref{GlobalTheory.\allowbreak{}T.\allowbreak{}period\_\allowbreak{}entries\_\allowbreak{}tie\_\allowbreak{}and\_\allowbreak{}level\_\allowbreak{}valuations\_\allowbreak{}at\_\allowbreak{}built}{Development.\allowbreak{}AutomorphicArithmetic.\allowbreak{}ThetaCorrespondence.\allowbreak{}ThetaPairLifts.\allowbreak{}ThirdPlaceSeries.\allowbreak{}SixRowsAtMemberStatement}{12--13}
\leanref{GlobalTheory.\allowbreak{}Fc.\allowbreak{}six\_\allowbreak{}rows\_\allowbreak{}at\_\allowbreak{}member\_\allowbreak{}alt}{Development.\allowbreak{}AutomorphicArithmetic.\allowbreak{}ThetaCorrespondence.\allowbreak{}ThetaPairLifts.\allowbreak{}ThirdPlaceSeries.\allowbreak{}SixRowsAtMember}{11--24}

\begin{proof}
Take \(R=\top\) in (8.13). Good members carry the columns for
the reference forms map and twisted evaluation. The period
theorems at the two kinds of tuple give (8.18) under
\(\Lean{DGamDef}\). The real comparison gives item~5 and the sharp
column bounds give item~6. These point conclusions apply to
every linkage and to the selection in \ref{B8.04}.
\end{proof}

For words \(u,u'\) and a set \(S\) of term pairs, put
\[
 \operatorname{PairAbs}_{u,u'}(S)
 =\Gamma_1\sum_{(t,t')\in S}|\langle v_t,v_{t'}\rangle|, \tag{8.19}
\]
with term vector \(v_t\) and normalization \(\Gamma_1\).
Among same-row pairs with nonzero scalars \(G(t),G(t')\), let
\(S_{\mathrm{big}}\) consist of pairs of corank at least \(2\) and
\(S_{\mathrm{nsc}}\) of nonscalar generic pairs of corank \(1\). Put
\[
 \epsilon_L(D)=\exp\!\left(-\frac13 kf\log k\right),\qquad
 \Delta_{u,u'}=\frac{d_u^C d_{u'}^C}{D_C^\star}, \tag{8.20}
\]
where \(d_u^C,d_{u'}^C\) are the selected model dimensions and
\(D_C^\star\) the model dimension main term. The arithmetic
counterparts are \(d_u^S,d_{u'}^S,D^\star\).

\begin{theorem}[Three estimates for the selected spectral rows]\label{B8.10}
Under the second convergent unitary Siegel--Weil input,
(8.21)--(8.23) give \(\operatorname{Lg1}\)-rows for every
\(KP,Q,Q',L_k\), closed on \(\mathcal D_3(\mathcal A)\) and
eventual at the points. Put
\[
 \begin{gathered}
 C'=sC'(F_m)=F_m.C_{\rm kLow},\qquad
 C_\ell=sCl(F_m)=\min(C_{\rm fL},C_w,C_{\rm Del},C_{\rm Count},
 C_{\rm h0r},C_{\rm gap}),\\
 p=(C_{YL}+1,\tfrac14,3,\tfrac12,C_{Tu},3).
 \end{gathered}
\]
Fix a real-valued function \(a_F\) of the outer list. Under the
closed size thresholds, for each eligible \(F_m,CH\) there is a
cutoff \(x_0\) for (8.21) and (8.22). Both estimates hold under
their respective hypotheses for every \(x\ge x_0\) with \(x>F_m.x_1\)
at or above the selection cutoff, and for every \(i\in I_x\).
For (8.21), assume for the pair \(u,u'\) that
\(\mathrm{Cnt.Hyp}(D,L,u,u';C',C_\ell,p,a_F(F_m.c_{\mathrm{out}}))\)
and \(\Lean{ConcTie}(X,Y)\) hold for the selected models \(X,Y\).
The latter requires \(d_u^C=d_u^S\) for all \(u\) and \(D_C^\star=D^\star\),
level-subgroup normalization, neatness for each word, Wu decay and
phase integrability. It also requires compatible measures, a standard
additive character, the local conditions and \(\Lean{TFrameAt}\)
of (8.23). For the term operators \(\operatorname{op}(t)\), assume
\[
 \sum_t\|\operatorname{op}(t)\|<\infty,\qquad
 \sum_t\|v_t\|<\infty,\qquad \Gamma_1>0.
\]
Then
\[
 \begin{split}
 \operatorname{PairAbs}_{u,u'}(S_{\mathrm{big}})
 &\le\epsilon_L(D)\Delta_{u,u'},\\
 \operatorname{PairAbs}_{u,u'}(S_{\mathrm{nsc}})
 &\le\epsilon_L(D)\Delta_{u,u'}.
 \end{split} \tag{8.21}
\]
For each word pair put
\[
 K_0=\operatorname{part}(\mathrm{idSet})+
 \operatorname{part}(\mathrm{genSet})+
 \operatorname{part}(\mathrm{ngSet}),\qquad
 \operatorname{part}(S)=\sum_{t\in S}\operatorname{op}(t).
\]
The sets are the identity terms and the generic and nongeneric
terms with \(G(t)\ne0\), respectively.
Under \(\Lean{AdmissibleBoundsFromSource.\allowbreak{}hyp}\) at
\(3,C',C_\ell,C_{YL}+1,\tfrac14,3,\tfrac12,C_{Tu}\), distinct \(u,u'\) satisfy
\[
 \|K_0(u,u')\|_{\mathrm{HS}}^2\le(1+b^{-6})\frac{d_u^S d_{u'}^S}{D^\star}.
 \tag{8.22}
\]
Under \(\mathrm{Cnt.Hyp0}(D,L;C',C_\ell,p)\), every level prime \(T_w\) satisfies
\[
 1\le T_w\left(\frac{D.\mathrm{host.ka}}{g_0g_1}\right)^{2r_S}.
 \tag{8.23}
\]
Here \(r_S=m/2-1\) for even \(m\ge2\).
\end{theorem}
\leanref{GlobalTheory.\allowbreak{}LocalRows.\allowbreak{}lg1\_\allowbreak{}sLp\_\allowbreak{}sO\_\allowbreak{}hTF}{Development.\allowbreak{}AutomorphicArithmetic.\allowbreak{}ThetaCorrespondence.\allowbreak{}ThetaPairLifts.\allowbreak{}ThirdPlaceSeries.\allowbreak{}AtChoice}{27--28}
\leanref{Formalism.\allowbreak{}L2.\allowbreak{}SG.\allowbreak{}ThmSOwn.\allowbreak{}LoopRowS}{Formalism.\allowbreak{}LocalToGlobal.\allowbreak{}SiegelLatticeTheory.\allowbreak{}FlagStabilizerTensors.\allowbreak{}ThmSOwn}{1015--1025}

\begin{proof}
On \(\mathcal D_3(\mathcal A)\) the twisting selection has
\(\operatorname{TuSize}\). The convergent Siegel--Weil formula gives
the hyperplane volume ratio, whence the closed rank bounds and (8.23).
Use the large-gap Bessel multiplicity and Whittaker--Bessel
results with the discrete-series coefficient bound. Wu decay and
finite dimensionality supply the finite spectral spaces.

Put \(c_{\rm out}=F_m.c_{\rm out}\) and let \(L_*\) be its
concatenation with \(F_m.c_{\rm mid}(CH)\). Write \(C_{\rm row}(a_F,L_*)\)
for the row-bound threshold. Evaluate the two loop thresholds at
\((C_{YL},C_{Tu},a_F(c_{\rm out}),C_{\rm row}(a_F,L_*))\). Require
each entry of the last list to be at least their maximum and
combine this with the row and volume thresholds. If \(x_L\) is the
loop cutoff and \(y_R\) the raised selection cutoff, take
\(x\ge\max(x_L,y_R+1)\). Lemma~\ref{B8.10a} gives (8.21).

For distinct words take \(a_F=3e^{7/4}\). The rank bounds give
the size hypotheses. Entry \(17\), counted from zero, of the outer
and middle lists concatenated is \(C_\Gamma>0\). This list
determines the choices of \(k_0,C_n,C'_0,C_{\ell,0}\). Impose
\[
 sC'(F_m),sCl(F_m)\ge\max(C'_0,C_{\ell,0},29). \tag{8.24}
\]
Beyond the resulting point cutoff, the source hypotheses give the
gap and counting bounds, and the selection gives compatibility
and absolute convergence. Lemma~\ref{B8.10b} uses \ref{B8.10a} to
give (8.22). Take \(x_0\) to be the maximum of the cutoffs for
(8.21) and (8.22), and \(R=\top\) in (8.13) for the three rows.
\end{proof}
\leanref{GlobalTheory.\allowbreak{}LocalRows.\allowbreak{}at\_\allowbreak{}choice\_\allowbreak{}alt}{Development.\allowbreak{}AutomorphicArithmetic.\allowbreak{}ThetaCorrespondence.\allowbreak{}ThetaPairLifts.\allowbreak{}ThirdPlaceSeries.\allowbreak{}AtChoice}{16--23}
\leanref{Formalism.\allowbreak{}L2.\allowbreak{}SG.\allowbreak{}ThmSOwn.\allowbreak{}offDiagonalK0\_\allowbreak{}of\_\allowbreak{}rows}{Formalism.\allowbreak{}LocalToGlobal.\allowbreak{}SiegelLatticeTheory.\allowbreak{}FlagStabilizerTensors.\allowbreak{}LinkSFinalP}{58--115}

\begin{lemma}[Summing the loop pairs]\label{B8.10a}
Under the hypotheses and thresholds for (8.21) in \ref{B8.10},
\(\operatorname{PairAbs}_{u,u'}(S_{\rm big}\cup S_{\rm nsc})
\le e^{-P/3}\Delta\), where \(P=kf\log k\) and
\(\Delta=\Delta_{u,u'}\).
\end{lemma}
\begin{proof}
Put \(r=r_S\), \(m=2r+2\), \(g=g_0g_1\), \(L_b=\log b\) and
\(f_w=d_w^S\) for each word \(w\). Bound each absolute inner product
by \(\|v_t\|\|v_{t'}\|\). Within each row, group finite subsums by reduced Gram
matrices and corank \(s\). The sum of reciprocal stabilizer orders
over rows, times the column index \(\dim\pi_u\) and \(\Gamma_1\),
equals \(f_u\). The nondegenerate row bound has the factor
\(R'=f_{u'}/D^\star\). For the degenerate row use
\(1\le C_0\Gamma^2R'\), where \(C_0=3\) in \ref{B8.10}.
Thus both bounds use \(\Delta=f_uR'\).

Put \(a=a_F(F_m.c_{\rm out})\) and choose \(C\ge\max(1,4\sqrt a)\)
above the row and mass thresholds. Set
\(H=b^{Cm^2f}\), \(N=T_u^f\) and \(v=D.\mathrm{host}.v\).
For \(2\le s\le k-2\) let \(\Xi_s=M_sM_{k-s}/M_k\), with \(M_j\)
the rank-\(j\) unitary mass. For \(s=k-1,k\) take
\(\Xi_s=(Ck)^{kf}(H^{Cm}N)^{Ck}\). The reduced-pair count and
the amplitude bounds give the factors
\[
 \begin{split}
 F_s&=(kv)^{kf}\Xi_s C^{kf}H^{C(s+m)m^2}(Csf)^{6sf}\Gamma,\\
 U_s&=(2/g)^{(s-1)kf}F_s,\qquad
 V_s=(2/g)^{(s-2)kf}F_s e^{Cm^6fL_b}(C/k)^{kf}.
 \end{split}
\]
The level congruences supply \(s-1\) powers of \((2/g)^{kf}\)
for nondegenerate pairs and \(s-2\) for degenerate pairs. In the
latter case the amplitude supplies \((C/k)^{kf}\).
At corank one count partners with each term fixed in turn.
Injectivity of the term label and uniqueness of the lattice
parameter bound the fibres. The counting hypothesis gives
\(T_{u'}/T_u\le a\), so the weight ratio in the small-amplitude
bound satisfies
\[
 4\sqrt{T_{u'}/T_u}\,g^{1-m}k^{-r}\le Cg^{-(2r+1)}k^{-r}.
\]
The two fibre bounds and Cauchy--Schwarz give the \(kf/2\) power
of the ratio in the bound
\[
 J=C^fH^{Cm^3}\Gamma e^{Cm^6 fL_b}
 (Cg^{-(2r+1)}k^{-r})^{kf/2}.
\]
Each finite subsum is bounded by
\(\max(1,C_0)\Delta[\sum_{s=2}^k(U_s+\Gamma^2V_s)+J]\).
The same bound holds for the full sum by monotone convergence.

For \(\Gamma=e^{C_\Gamma m^5fL_b}\) put
\(\gamma=e^{\min(C,C_\Gamma)m^5fL_b}\). Replace \(\Gamma\) by
\(\gamma\) in the bracket and multiply the outside factor by
\((\Gamma/\gamma)^3\), since the highest power is three and
\(\Gamma\ge\gamma\). Write \(A=\max(1,C_0)(\Gamma/\gamma)^3\).
The bounds give \(\log(\Gamma/\gamma)\le k/8\) and the last
threshold gives \(6\log(3A)\le P\).
For \(s\le k-2\) use \((kv)^{kf}\Xi_s\le H^C C^{skf}\).
For the endpoints use \(N\le b^{Cmf}\), \(v\le b\) and
\(s-2\ge k-3\). With \(4\le m\le f\le k\),
\(g\ge(kb)^8\) and \(k\ge16C^{32}m^5fL_b\), the logarithms at
\(\gamma\) satisfy
\[
 \log U_s\le-\tfrac12kf\log g-\log k,\qquad
 \log V_s\le-P/2-2\log\gamma-\log k,\qquad
 \log J\le-P
\]
for positive terms. At \(s=2\), the saving in \(V_s\) comes
from \(-kf\log k\) in the amplitude. Summing over \(s\) gives
\(\sum_s U_s\le g^{-kf/2}\) and \(\gamma^2\sum_s V_s\le e^{-P/2}\).
The remaining term satisfies \(J\le e^{-P}\). Since \(g^{-kf/2}\le e^{-4P}\),
their sum gives \(3A e^{-P/2}\Delta\le e^{-P/3}\Delta\).
\end{proof}
\leanref{GlobalTheory.\allowbreak{}LoopRowBoundsGap.\allowbreak{}rows\_\allowbreak{}along\_\allowbreak{}family}{Development.\allowbreak{}AbelianVarietyTheory.\allowbreak{}Proofs.\allowbreak{}FlagMeasures.\allowbreak{}LoopRowBoundsAlongFamilyLargeGaps}{25--44}
\leanref{Formalism.\allowbreak{}L2.\allowbreak{}SG.\allowbreak{}LoopsTies.\allowbreak{}loop\_\allowbreak{}evaluation\_\allowbreak{}threshold\_\allowbreak{}zero}{Formalism.\allowbreak{}LocalToGlobal.\allowbreak{}SiegelLatticeTheory.\allowbreak{}FlagStabilizerTensors.\allowbreak{}RootLatticeThresholds.\allowbreak{}LoopEvaluationThresholdZero}{44--120}
\leanref{Formalism.\allowbreak{}FourierDecayAlt.\allowbreak{}P.\allowbreak{}GrowthPolynomialOperator.\allowbreak{}t\_\allowbreak{}LoopsJoint}{Formalism.\allowbreak{}HermitianCurvature.\allowbreak{}HermitianCurvatureProofs.\allowbreak{}GrowthOfBasic}{1191--1229}

\begin{lemma}[The scalar row and the full kernel]\label{B8.10b}
At a selected point under the source and size hypotheses for
(8.22) of \ref{B8.10}, distinct words \(u,v\) satisfy the bound
\((1+b^{-6})\Delta_{u,v}\) for the squared Hilbert--Schmidt norms
of both \(K_0\) and the full normalized Gram operator.
\end{lemma}
\begin{proof}
The kernel expansion represents the normalized Gram operator in
isometric column coordinates by \(\sum_t\operatorname{op}(t)\).
Absolute convergence allows partition into the four cells.
For distinct words the identity cell vanishes. In the nonidentity
cell with \(G(t)=0\), unfolding gives zero for irrelevant terms.
For each relevant term the orbital integral intertwines unequal
word weights and vanishes by Schur's lemma. Hence the sum equals \(K_0\).

Put \(d=d_u^S,\ d'=d_v^S,\ r=r_S\) and \(L_b=\log b\).
Let \(c_{\rho p}\) be the row coefficients on terminal flags \(p\)
and \(V_p\) their archimedean vectors. Put \(s_d=\sup\mathrm{Dens}/\alpha\),
where \(\mathrm{Dens}\) is the set of global densities and \(\alpha\)
is the product of the local normalization factors. Fibre counting
bounds each eligible fibre's normalized density by \(s_d\). The
change to flag coordinates acts on both kernels by the same
diagonal congruence. Thus for every finite terminal flag set \(Y\)
and finite row set \(S\), the matrix
\[
 \left[s_d K_{\rm top}(p,q)-\sum_{\rho\in S}
 c_{\rho p}\overline{c_{\rho q}}\right]_{p,q\in Y}
\]
is positive semidefinite, with \(K_{\rm top}\) the top flag kernel.
Rescale the top kernel by \(e_1=\max(1,s_d)\).
The lower density multipliers equal one.

At a position where the words differ, let \(\Theta^{-f}\) be the
largest product of local flag thresholds. The size estimate at
\((C_{\rm Del},c_T,C_b)=(5,\tfrac12,3)\) gives \(\Theta\ge k^{1/16}\).
Here \(C_f\ge256\) and \(C_f\log b\le f\log k\).
In the recursion for the square root of the flag energy, divide
out the square roots of the density multipliers. A jump of \(j\ge2\)
levels then has coefficient at most \(h^{j-1}\), with \(h=(C_1/\Theta)^f\).
The one-level coefficient is at most \(1+h\). The running maximum
grows by a factor at most \(1+h+h/(1-h)\) per level. Its square
over \(r+1\) levels is \(1+O((r+1)h)\) for small \((r+1)h\).
Since \(r+1\le f\), enlarge \(k_0\) so that \(C^{f+1}(r+1)\le k^{f/32}\).
Then
\[
 \delta_s=C(r+1)(C/\Theta)^f\le k^{-f/32}\le b^{-8}\le b^{-7}.
\]
Thus for every finite row set \(S\),
\[
 \sum_{\rho\in S}\Big\|\sum_p c_{\rho p}V_p\Big\|^2\le A I_N,
 \qquad A=(1+\delta_s)e_1,
\]
where \(I_N\) is the normalized archimedean integral. Each rational
prime has at most \(f\) density factors with logarithm at most
\(4q^{-j}\) for residue cardinality \(q\) and \(j=k-2r-2\).
Since \(\sum_{n\ge2}n^{-j}\le2^{-j}+\int_2^\infty t^{-j}\,dt
\le3\,2^{-j}\), it follows that \(\log e_1\le12f\,2^{-j}\).

The signed generic scalar sum of corank one is \(\Gamma_1\rho_0K_V\),
where \(K_V\) is the orbital quadratic form. The row bound gives
\(cc'K_V\le n_0 A I_N\) with unfolding constants \(c,c'\).
The Jacobian and formal degree give
\(I_N\le V_{ZZ}/(V_{\rm loc} C_\infty\deg)\).
The row mass gives \(n_0\rho_0=V_B C_\infty M/(V_{ZZ}\mu_K)\).
The local mass and volume factors obey
\[
 \frac{M}{V_{\rm loc}\mu_K}
 =\frac{\mu_{K''}}{\Gamma_1\mu_{S''}\mu_{S''}'},
 \qquad
 \frac1c=\frac{d\mu_{S''}}{V_{H,\infty}},\qquad
 \frac1{c'}=\frac{d'\mu_{S''}'}{V'_{H,\infty}}.
\]
Since \(V_{H,\infty}V'_{H,\infty}=V_B V_G\), the bound becomes
\[
 \Gamma_1\rho_0K_V\le\frac{A dd'\mu_{K''}}{\deg V_G}
 =A\Delta_{u,v}.
\]

Put \(P_g=\prod_{j=0}^1(1+3/g_j)^f\), \(\ell=e^{-kf\log k/3}\)
and \(n=2^{-(k-2r-5)}e^{C_nm^5fL_b}\).
The local factors \(\prod_{i=1}^k(1-g_j^{-i})^{-1}\le1+3/g_j\)
bound the finite-type Gram matrix by \(P_g\Gamma_1 I\).
Pair this matrix inequality with the positive Gram matrix of
the term vectors. Row orthogonality removes cross-row terms.
For a same-row pair with nonzero \(G\), the \(2r+2\) vectors in
the two frames are linearly independent. The pair subspace lies
in the orthogonal complement of their span in dimension \(k+2r+1\).
It has dimension at most \(k-1\), so the pair has corank at least
one. For generic terms use the scalar bound and (8.21) on the
two loop classes. For nongeneric terms the absolute sum over
\(S_{\rm big}\) is at most \(\ell\Delta_{u,v}\), and the signed
corank-one sum is at most \(n\Delta_{u,v}\) by \(\Lean{sNG}\)
of \ref{B8.02}. The triangle inequality gives
\[
 \|K_0\|_{\rm HS}^2
 \le\big( \sqrt{P_g(A+2\ell)}+\sqrt{P_g(\ell+n)}\big)^2\Delta_{u,v}.
\]
Put \(z=b^{-7}\). Use \(k\ge\max(32,4r+12,70\log b)\),
\((kb)^8\le g_j\), \(f\le k\) and \(8C_nm^5fL_b\le k\).
Then \(P_g\le e^{6z}\) and \(\ell,n,\log e_1\le e^{-k/5}\le z^2\).
Thus \(A\le(1+z)e^{z^2}\). For \(b\ge16\) we have \(z\le10^{-3}\).
The radicands are at most \(1+11z\) and \(3z^2\), so the sum of
their square roots is at most \(1+(11/2+7/4)z\).
Its square is at most \(1+16z\le1+b^{-6}\).
\end{proof}
\leanref{Formalism.\allowbreak{}FourierDecayAlt.\allowbreak{}P.\allowbreak{}GrowthPolynomialOperator.\allowbreak{}t\_\allowbreak{}MainRow}{Formalism.\allowbreak{}HermitianCurvature.\allowbreak{}HermitianCurvatureProofs.\allowbreak{}GrowthOfBasic}{1070--1184}
\leanref{Formalism.\allowbreak{}FourierDecayAlt.\allowbreak{}P.\allowbreak{}GrowthPolynomialOperator.\allowbreak{}t\_\allowbreak{}PairBound}{Formalism.\allowbreak{}HermitianCurvature.\allowbreak{}HermitianCurvatureProofs.\allowbreak{}GrowthOfBasic}{525--551}
\leanref{Formalism.\allowbreak{}L2.\allowbreak{}SG.\allowbreak{}TheoremSConclusion.\allowbreak{}theoremS\_\allowbreak{}of\_\allowbreak{}conj}{Formalism.\allowbreak{}LocalToGlobal.\allowbreak{}SiegelLatticeTheory.\allowbreak{}FlagStabilizerTensors.\allowbreak{}TheoremSConclusion}{68--83}
\leanref{Formalism.\allowbreak{}L2.\allowbreak{}SG.\allowbreak{}FixedDecomposition.\allowbreak{}kppOffC\_\allowbreak{}of\_\allowbreak{}steps0}{Formalism.\allowbreak{}LocalToGlobal.\allowbreak{}SiegelLatticeTheory.\allowbreak{}CokernelGLCounts.\allowbreak{}RankPointCounts.\allowbreak{}FixedDecomposition}{755--775}


% ---- p8c2 ----
\chapter{Theta correspondence, test vectors and evaluation}\label{ch:p8c2}

Let \(E/F\) be a CM extension with conjugation \(\sigma\) and adele rings \(\mathbb A_F,\mathbb A_E\). Write \([H]=H(F)\backslash H(\mathbb A_F)\) and \(\operatorname{avg}_{[H]}=\operatorname{vol}([H])^{-1}\int_{[H]}\). Let \(\mathscr P=\)\Lean{ThetaLocalInputs} from Theorem~\ref{B8.01} and \(\mathscr J=\)\Lean{JoinOut} collect the local theta inputs and the toric, denominator, lattice-norm, transport and reflection inputs, respectively. At scale \(x>1\) put \(L=x^\varepsilon=\log b\), with \(0<\varepsilon<1/52\).

\section{Matching theta integrals}

Let \(V,V'\) be Hermitian and \(W\) skew-Hermitian, of dimensions \(a,a',a+a'\) and signatures \((a-p,p),(q,a'-q),(a+a'-p-q,p+q)\) at each infinite place. Require
\begin{equation}\label{eq:p8c2-triple-range}
 a'<a,\qquad a+a'\equiv0\pmod2,\qquad
 2(p+q)+1<\frac{a-a'}2,\qquad 2(p+q)+1<a'.
\end{equation}
Use the prescribed Cartan involutions. Take the prescribed self-dual lattices with their split frames. Assume \(E/\Q\) Galois and \(E/F\) unramified at finite places. The additive character \(\psi\) of \(\mathbb A_F/F\) has infinite part \(y\mapsto\exp(-2\pi i\operatorname{Tr}_{F_\infty/\R}y)\). Put \(T=S_0\sqcup Q\subseteq P\), where \(P\) is the set of prescribed split places with chosen embeddings. At inert \(v\) require odd residue characteristic and \(Nv>3\). At inert places and on \(P\), \(\psi_v\) has conductor \(\mathcal O_v\). Fix \(v_1\notin T\) above a rational prime split completely in \(E\).

The unitary Hecke characters \(\alpha,\beta\) restrict to \(\eta^{a'},\eta^a\), with \(\eta\) quadratic for \(E/F\), and their infinite exponents satisfy \(m_\alpha-m_\beta>a+a'\). Both are unramified off \(T\), and \(\nu=\alpha/\beta\) is ramified above \(T\). The Weil models \(\omega,\omega'\) on \(V\otimes_E W,V'\otimes_E W\) use \((\beta,\beta/\alpha),(\alpha,\alpha/\beta)\), respectively.

Let \(\mathrm{SoLoc}(t)\) be the primed Schwartz vectors with full-rank support in an allowed barred model for \(t:E\hookrightarrow F_v\). For the chosen \(t_1\) over \(v_1\), let \(\mathrm{So}(t_1)\) be the span of restricted pure tensors with \(v_1\)-factor in \(\mathrm{SoLoc}(t_1)\). Put \(I(\Phi,h)=\int_{[\U(V)]}\theta_\omega(g,h;\Phi)\,dg\), and define \(I'\) using \(V',\omega'\). Real polynomial vectors are linear combinations of Heisenberg derivatives of Gaussians.

The maps \(p_\infty,p'_\infty\) on real polynomial vectors and \(p_v,p'_v\) on finite Schwartz vectors surject onto common nonzero irreducible quotients \(\Pi_\infty,\Pi_v\). They intertwine \(\U(W)\), at infinity through its Lie algebra and maximal compact. They kill the source coinvariant relations at finite places. At infinity they kill the source Lie-algebra action and are invariant under the source maximal compact. Off \(T\) the vacua have the same nonzero image. At \(v\in T\) choose split and barred model maps \(\iota,\iota'\) sending vacua to integral-matrix indicators. Then \(p_v\phi=p'_v\phi'\) is equivalent to \(\mathrm{bigP}(\iota\phi)=\mathrm{smallP}(\iota'\phi')\) on \(\GL_{a+a'}(F_v)\).

A polynomial match has equal images under \(p_\infty\otimes\bigotimes_{v\in S}p_v\) and its primed counterpart, where \(S\) contains \(T\) and the non-vacuum places. A general match is a limit in the real Schwartz topology of polynomial matches with fixed finite factors and primed vectors in \(\mathrm{So}(t_1)\).

\begin{theorem}[Matching and proportional theta integrals]\label{B8.11}
Assume Theorem~\ref{B8.12}. Matching maps exist for every triple above. Fix such maps. The integral \(I\) converges on Schwartz vectors. For \(\Phi'\in\mathrm{So}(t_1)\), put \(\operatorname{qht}(g)=\inf_{\gamma\in\U(V')(F)}\operatorname{ht}(\gamma g)\), with \(\operatorname{ht}\) the adelic height. For every \(N\in\N\), some \(C_N\in\R,\ k_N\in\N\) satisfy
\[
 \|\theta_{\omega'}(g,h;\Phi')\|\operatorname{qht}(g)^N
 \le C_N\operatorname{ht}(h)^{k_N}\qquad(g,h\text{ arbitrary}).
\]
The integral \(I'\) converges with moderate growth on this space and is nonzero for some test there. A unique \(c\in\C^\times\) satisfies
\begin{equation}\label{eq:p8c2-matched-integrals}
 I(\Phi,h)=cI'(\Phi',h)
\end{equation}
for every match and every \(h\).

At \(v\notin T\) over a prime split completely in \(E\), the kernels are the local coinvariant relations, and \(p'_v(\mathrm{SoLoc}(t))=\Pi_v\) for each \(t:E\hookrightarrow F_v\). A Schwartz functional killing the primed relations factors through \(p'_v\). The primed Weil action preserves \(\mathrm{So}(t_1)\). The sign triples in the step formulas below are uniform in the base fields and belong to \(\{\pm1\}^3\).
\end{theorem}
\leanref{UpperTheory.\allowbreak{}LiteratureTheorems.\allowbreak{}theta\_\allowbreak{}integrals\_\allowbreak{}agree\_\allowbreak{}and\_\allowbreak{}local\_\allowbreak{}matching\_\allowbreak{}of\_\allowbreak{}literature}{Development.\allowbreak{}AutomorphicArithmetic.\allowbreak{}ThetaCorrespondence.\allowbreak{}ThetaPairLifts.\allowbreak{}ThirdPlaceSeries.\allowbreak{}EveryCharacterOfLiteratureOfVWin}{565--566}
\leanref{PositivityTheory.\allowbreak{}ThKR.\allowbreak{}theoremK\_\allowbreak{}head2}{Development.\allowbreak{}WeilBound.\allowbreak{}ThetaLifts.\allowbreak{}Reductions.\allowbreak{}FiniteClosure.\allowbreak{}TheoremKOfMod}{26--27}
\leanref{GlobalTheory.\allowbreak{}IndefiniteRowSideFullHecke.\allowbreak{}WeilTensorModels.\allowbreak{}ThetaIntegralsAgree}{Development.\allowbreak{}AbelianVarietyTheory.\allowbreak{}BlockTheory.\allowbreak{}DivisorBase.\allowbreak{}WeilTensorModels}{585--592}
\leanref{GlobalTheory.\allowbreak{}IndefiniteRowSideFullHecke.\allowbreak{}WeilTensorModels.\allowbreak{}LocalMatching}{Development.\allowbreak{}AbelianVarietyTheory.\allowbreak{}BlockTheory.\allowbreak{}DivisorBase.\allowbreak{}WeilTensorModels}{581--583}

For fractional ideals \(\mathfrak b,\mathfrak c\), signs \(e_1,e_2,s\) and \(z\in\Z\), set
\begin{equation}\label{eq:p8c2-step-factor}
 \mathfrak f(e_1,e_2,s;\mathfrak b,\mathfrak c;z)
 =\nu(\mathfrak b)^{e_1}\nu(\mathfrak c)^{e_2}
   \prod_{v\in T}\nu_v(s)\epsilon_v(\psi,\nu_v,z).
\end{equation}
Use the prescribed ideal character and the component \(\nu_v\) at the chosen embedding. For uniformizer \(\varpi_v\) and conductor exponent \(c_v\), put
\[
 \epsilon_v(\psi,\nu_v,z)
 =(Nv)^{-c_vz}\nu_v(\varpi_v)^{c_v}
 \sum_{u\in(\mathcal O_v/\varpi_v^{c_v})^\times}
    \nu_v(u)^{-1}\psi_v(u/\varpi_v^{c_v}).
\]

In the \(\beta\)-step let \(a'=b_0+2\), \(\ell=(a-b_0-2)/2\). Keep \(V\) and remove adapted planes from \(V',W\) at frame positions \((0,b_0+1),(a,a+b_0+1)\). The lower indices are \((p,q-1)\). With lower constant \(c_1\),
\begin{equation}\label{eq:p8c2-beta-step}
 cL^{S_0}(\nu,\ell+1)=
 \frac{c_1}{C\,\mathfrak f(e_{\beta,1},e_{\beta,2},s_\beta;\mathfrak b,\mathfrak c;-\ell)} .
\end{equation}
In the \(\alpha\)-step let \(q=0\), \(a=b_0+2\), \(\ell=(b_0+2-a')/2\). Keep \(V'\) and remove planes from \(V,W\) at \((0,b_0+1)\) on each side. The lower indices are \((p-1,0)\), and
\begin{equation}\label{eq:p8c2-alpha-step}
 c=C\,\mathfrak f(e_{\alpha,1},e_{\alpha,2},s_\alpha;\mathfrak b,\mathfrak c;1-\ell)
 L^{S_0}(\nu,\ell)c_1 .
\end{equation}
For source and target plane vectors \(x,x'\) as in Lemma~\ref{B8.11b}, set \(\mathfrak b=\{s\in E:sx\in\Lambda_{\rm source}\}\) and \(\mathfrak c=\{s\in E:sx'\in\Lambda_W\}\). The \(L\)-function omits places over \(S_0\). With models, measures, domains, planes and real maps fixed, choose \(C\ne0\) before the characters, Weil actions and finite matching maps.
\leanref{GlobalTheory.\allowbreak{}IndefiniteRowSideFullHecke.\allowbreak{}WeilTensorModels.\allowbreak{}StepBeta}{Development.\allowbreak{}AbelianVarietyTheory.\allowbreak{}BlockTheory.\allowbreak{}DivisorBase.\allowbreak{}WeilTensorModels}{559--567}
\leanref{GlobalTheory.\allowbreak{}IndefiniteRowSideFullHecke.\allowbreak{}WeilTensorModels.\allowbreak{}StepAlpha}{Development.\allowbreak{}AbelianVarietyTheory.\allowbreak{}BlockTheory.\allowbreak{}DivisorBase.\allowbreak{}WeilTensorModels}{571--579}

\begin{proof}
The local and real quotients give the matching maps. They also give the split kernels and full-rank surjectivity. Siegel-set bounds give convergence and decay. At fixed finite factors they also give continuity in the real Schwartz topology.

Induct on \(p+q\) for polynomial comparison and nonvanishing. At zero, the sources and target are definite. Theorem~\ref{B8.12} and local doubling give nonzero tests on both sides. For a polynomial match \(y\), write \(I_y,I'_y\) for its integrals. The doubled inner-product identities give constants \(c,\lambda\), uniform in \(y\), such that
\[
 D_y=I_y-cI'_y,\qquad
 \langle D_y,I'_y\rangle=0,\qquad
 \|D_y\|^2=\lambda\|I'_y\|^2 .
\]
Choose a match with \(I'_y\ne0\). The reverse lift on the primed source with the primed test of \(y\) is constant for \(D_y\) and for \(I'_y\). Fubini gives the respective constants \(\langle D_y,I'_y\rangle/\operatorname{vol}[\U(V')]\) and \(\|I'_y\|^2/\operatorname{vol}[\U(V')]\). The doubling coefficient identity equates the first constant to \(\lambda\) times the second. Thus \(\lambda=0\), and every \(D_y\) vanishes by continuity. The unprimed witness gives \(c\ne0\).

For positive index sum choose planes by Lemma~\ref{B8.11a}. In Lemma~\ref{B8.11b} take \(c=c_1/(CS)\) for \(q>0\) and \(c=CS c_1\) for \(q=0<p\). The primitive Gauss sums and the Euler product at \(z>1\) give \(S\ne0\). Apply the Levi identities to compact translates of a polynomial match with \(\Phi'\in\mathrm{So}(t_1)\). The Levi laws and Iwasawa decomposition give \(B=cB'\) for \(q>0\) and \(A=cA'\) for \(q=0<p\). The difference \(f=I(\Phi)-cI'(\Phi')\) belongs to \(\Sigma\) of Lemma~\ref{B8.11c}.

Put \(d=a+a'-2\). If \(q=0<p\), definiteness gives \(B'=0\). The unprimed \(B\) is a zero-evaluation integral with source \(V\) and target \(W_1\). The real vanishing criterion applies because the smaller target signature index is \(p-1<p\), \(|m_\beta-m_\alpha|>d\) and \(d+2p<2a\). This gives vanishing on all Schwartz tests, so \(B=0\).

For \(q>0\), the line coefficient of \(f\) has pure \(\alpha\) type. Induct on \(p\) for differences in \(\Sigma\) with this property. At \(p=0\), definiteness gives \(A=0\). On the primed side replace the test by one with the same integral whose zero evaluation still belongs to \(\mathrm{So}(t_1)\) for \(W_1\). The zero-evaluation formula for \(A'\) now has smaller target index \(q-1<q\) and \(|m_\alpha-m_\beta|>d\). The real vanishing criterion on these full-rank tests gives \(A'=0\).

At \(p>0\), Witt transitivity takes the line to \(Ex'\subset W\), the target line in the step that removes a plane from \(V\). At line coordinate one the coefficient of each compact translate is a lower difference at \((p-1,q)\). Take the next line term and split it into its two Levi types. Let \(u\in\mathbb A_F^\times\) have components \(2\) at infinity and \(1\) at finite places, viewed in \(\mathbb A_E^\times\). Opposite scalings \(u,u^{-1}\) on the lines define \(a_u\). They fix the first type and multiply the second by \(\lambda\), with \(|\lambda|=|u|_E^{(a'-(a-2))/2}<1\). For the two-step term \(c_2\), the function \(c_2(a_u g)-c_2(g)\) thus scales by \(\lambda\). Reversing the order of integration gives \(c_2\) as the unipotent average over \(E\backslash\mathbb A_E\) of the two-plane term on \(\mathrm{SL}_2(\mathbb A_E)\). Its forward values \(c_2(a_u^k g)\) are bounded for \(k\ge0\) uniformly on compact sets of \(g\). Apply the smoothing estimate of Lemma~\ref{B8.11c} on this ray. For the bounded smoothed function, backward iteration bounds the difference by \(O(|\lambda|^k)\) for every \(k\). Shrinking the smoothing support gives \(c_2(a_u g)=c_2(g)\). Since \(1-\lambda\ne0\), the second type vanishes. Induction kills the lower differences, and Iwasawa decomposition kills the upper line term. Lemma~\ref{B8.11c} now gives \(f=0\) in both cases. Transport the identity to the original frames.

For \(q>0\), match the nonzero \(B\) test with a primed test in \(\mathrm{So}(t_1)\). For \(q=0\), start with the nonzero definite primed test there and take an unprimed match. These witnesses give uniqueness, and continuity gives the general comparison. For any prescribed adapted planes, cancel the nonzero \(B'(m)\) at the \(\beta\) witness. At the \(\alpha\) witness, replace the primed \(v_1\)-factor by a full-rank preimage of its quotient image. The definite integral factors through those coinvariants, so \(A=cA'\) at that witness. Cancel \(A'(m)\ne0\) to obtain \eqref{eq:p8c2-alpha-step}; the \(\beta\) cancellation gives \eqref{eq:p8c2-beta-step}.
\end{proof}
\leanref{PositivityTheory.\allowbreak{}ThKG.\allowbreak{}Leaves.\allowbreak{}const\_\allowbreak{}of}{Development.\allowbreak{}WeilBound.\allowbreak{}ThetaLifts.\allowbreak{}ConeHypotheses.\allowbreak{}CtSmallCone}{88--136}
\leanref{PositivityTheory.\allowbreak{}ThKG.\allowbreak{}Leaves.\allowbreak{}isConst\_\allowbreak{}of}{Development.\allowbreak{}WeilBound.\allowbreak{}ThetaLifts.\allowbreak{}ConeHypotheses.\allowbreak{}CtSmallCone}{179--189}
\leanref{PositivityTheory.\allowbreak{}ThKG.\allowbreak{}Leaves.\allowbreak{}kConst}{Development.\allowbreak{}WeilBound.\allowbreak{}ThetaLifts.\allowbreak{}ConeHypotheses.\allowbreak{}CtSmallCone}{191--198}
\leanref{PositivityTheory.\allowbreak{}ThKG.\allowbreak{}Leaves.\allowbreak{}kBeta}{Development.\allowbreak{}WeilBound.\allowbreak{}ThetaLifts.\allowbreak{}ConeHypotheses.\allowbreak{}CtSmallCone}{200--218}
\leanref{PositivityTheory.\allowbreak{}ThKG.\allowbreak{}Leaves.\allowbreak{}kAlpha}{Development.\allowbreak{}WeilBound.\allowbreak{}ThetaLifts.\allowbreak{}ConeHypotheses.\allowbreak{}CtSmallCone}{220--231}
\leanref{PositivityTheory.\allowbreak{}ThKEnd.\allowbreak{}Two.\allowbreak{}twoStep\_\allowbreak{}of'}{Development.\allowbreak{}WeilBound.\allowbreak{}ThetaLifts.\allowbreak{}EigenvalueEndgame.\allowbreak{}SignatureSplitting.\allowbreak{}TwoStepOf}{22--102}

\begin{lemma}[Removing a hyperbolic plane]\label{B8.11a}
For \(q>0\) one can choose Cartan involutions and split frames on the spaces of Theorem~\ref{B8.11} so the \(\beta\)-planes are adapted and the complements give \((p,q-1)\). For \(q=0<p\) the \(\alpha\)-step gives \((p-1,0)\). Removal on the \(V\)-side is possible for every \(p>0\) with \(q\) fixed.
\end{lemma}
\begin{proof}
Each space losing a plane has dimension \(n\ge3\) and both signs at each infinite place. For isotropy, multiply a skew form by a nonzero trace-zero scalar. After diagonalization the quadratic form over \(F\) is \(\sum_{i=1}^n d_iN_{E/F}(z_i)\), with \(d_i\in F^\times\) of both signs at each real place. Its dimension \(2n\ge6\) gives isotropy at finite places. Hasse--Minkowski gives an isotropic line \(Ex\). \cite[Ch.~6]{scharlau-1985}

Write \(h\) for the original form, linear in the first argument, and \(\Lambda\) for its self-dual lattice. The dual is taken with values in \(\mathcal O_E\) for a Hermitian source and in the absolute codifferent \(\mathfrak D_{E/\Q}^{-1}\) for the skew-Hermitian target. Scale \(x\) so the ideal \(J=h(\Lambda,x)\) has valuation zero at the specified split primes and their conjugates. Since \(JJ^{-1}=\mathcal O_E\), choose \(y_0\in J^{-1}\Lambda\) with \(h(x,y_0)=1\). For \(z\in\Lambda\), the pairing of \(h(z,y_0)x\) with \(w\in\Lambda\) is \(h(z,h(w,x)y_0)\) up to the Hermitian sign. By self-duality, \(h(z,y_0)x\in\Lambda\). Unramifiedness gives integral \(\theta\) with \(\theta+\bar\theta=1\). Put \(y=y_0-\theta h(y_0,y_0)x\). Then \(h(y,y)=0\) and \(h(x,y)=1\). Applying the projection formula to \(Jy_0\subseteq\Lambda\) gives \(Jy\subseteq\Lambda\). The same duality calculation with \(y\) gives \(h(z,y)x\in\Lambda\). Thus both projections preserve the lattice and the orthogonal complement inherits self-duality. The line ideals are units at the prescribed primes, so the two lines extend to split frames in the required positions. Take Cartan involutions adapted to the planes.

Removal subtracts one from each sign. In \eqref{eq:p8c2-triple-range} the common left side of the strict size bounds falls by two and their right sides fall by at most two. The character restrictions persist by \(\eta^2=1\).
\end{proof}
\leanref{PositivityTheory.\allowbreak{}ThKG.\allowbreak{}Tow.\allowbreak{}Pln.\allowbreak{}exists\_\allowbreak{}plane}{Development.\allowbreak{}WeilBound.\allowbreak{}ThetaLifts.\allowbreak{}TowerModels.\allowbreak{}PlaneMain}{26--136}
\leanref{PositivityTheory.\allowbreak{}ThKG.\allowbreak{}Tow.\allowbreak{}towerB\_\allowbreak{}of}{Development.\allowbreak{}WeilBound.\allowbreak{}ThetaLifts.\allowbreak{}TowerModels.\allowbreak{}TowerCodimensionTwoOf}{32--73}
\leanref{PositivityTheory.\allowbreak{}ThKG.\allowbreak{}Tow.\allowbreak{}towerA\_\allowbreak{}of}{Development.\allowbreak{}WeilBound.\allowbreak{}ThetaLifts.\allowbreak{}TowerModels.\allowbreak{}TowerCodimensionTwoOf}{75--126}
\leanref{PositivityTheory.\allowbreak{}ThKG.\allowbreak{}Tow.\allowbreak{}towerCodimensionTwo\_\allowbreak{}of}{Development.\allowbreak{}WeilBound.\allowbreak{}ThetaLifts.\allowbreak{}TowerModels.\allowbreak{}TowerCodimensionTwoOf}{128--172}

\begin{lemma}[The two Levi identities]\label{B8.11b}
Fix an adapted step and its upper and lower real maps. Assume nonzero polynomial comparison and nonzero polynomial witnesses on both sides for smaller index sums, with the primed witness in \(\mathrm{So}(t_1)\). Let \(c_1\) be the lower comparison scalar. Write \(A+B,A'+B'\) for the line terms in order of characters \(\alpha|\cdot|_E^{a'/2},\beta|\cdot|_E^{a/2}\). Put \(z=\ell+1\) for \(\beta\) and \(z=\ell\) for \(\alpha\), and \(S=\mathfrak f(e_1,e_2,s;\mathfrak b,\mathfrak c;1-z)L^{S_0}(\nu,z)\) with uniform signs. There is \(C\ne0\), fixed before the characters, Weil actions and finite maps, such that at every Levi lift \(m\) of \((1,h_1)\),
\[
 \begin{array}{ll}
 \beta:& CS B(m)=c_1B'(m)\quad\text{for polynomial matches with }\Phi'\in\mathrm{So}(t_1),\\
 \alpha:& B'=0,\quad A(m)=CS c_1A'(m)\quad\text{for every polynomial match}.
 \end{array}
\]
Some unprimed polynomial test has \(B(m)\ne0\) in the \(\beta\)-step. Some polynomial match has \(A'(m)\ne0\) in the \(\alpha\)-step.
\end{lemma}
\begin{proof}
Let \(e_0\) be evaluation at zero in the removed target variable. On the side losing a source plane \(H=(x,y)\), use mixed-polarization evaluation \(e\) for \(H,H'=(x',y')\). With source Cartan and lattice stabilizer \(K\) and Heisenberg translation \(\rho\), put
\[
 \mathcal T\phi=\int_K\int_{\mathbb A_E^\times}
 e\bigl(\rho(u(t))\omega(k,1)\phi\bigr)\nu(t)|t|_E^z\,d^\times t\,dk,
 \qquad u_\alpha(t)=txy'^{\,t},\quad u_\beta(t)=\bar t\,yx'^{\,t}.
\]
Character orthogonality in the centre of the line unipotent retains the isotropic vectors. The zero orbit gives \(e_0\). Witt transitivity unfolds the nonzero orbit at a vector stabilizer and leaves the line coordinate and compact integral in \(\mathcal T\). For the primed nonzero orbit, the two partial Fourier expansions on each rational line, with their zero terms, satisfy
\[
 \Sigma_A^{\rm all}=\kappa\Sigma_B^{\rm all},\qquad
 \kappa=C_A^{-1}C_B\ne0,
\]
where \(C_A,C_B\) are the Fourier normalization constants. Reindex the second nonzero sum by \(t\mapsto\bar t^{-1}\). Full-rank support gives \(\Theta_B(0)=0\), hence
\[
 \Sigma_A^\times=\kappa\Sigma_B^\times-\Theta_A(0).
\]
The subtracted term obeys the \(A'\) law. A definite source contributes the zero orbit. The two coefficients are separated by an idele of norm different from one, since their eigenvalue moduli differ. At \(m\) the coefficients are given by the lower integrals at \(h_1\)
\[
 \begin{array}{c|cc}
 &\text{unprimed}&\text{primed}\\ \hline
 \beta&B=C_0I_1(e_0\Phi)&B'=C_1I'_1(\mathcal T\Phi')\\
 \alpha&A=C_1I_1(\mathcal T\Phi)&A'=C_0I'_1(e_0\Phi'),\quad B'=0 .
 \end{array}
\]
Here \(C_0,C_1\ne0\) incorporate measures and evaluations, with \(\kappa\) included in \(C_1\) in the \(\beta\)-row. At \(mh\) use the translated tests.

Factor \(\mathcal T\) on pure tensors. Summing the valuation shells of the vacua gives the unramified Euler factors. Conjugation inverts \(\nu\), and the codifferent has value one. The starting shells from the line ideals therefore give powers \((1,1)\) for \(u_\alpha\) and \((-1,-1)\) for \(u_\beta\). At \(v\in T\) let \(d=\operatorname{condExp}(\nu_v)>0\) and \(Q_v=Nv\). With additive volume one on \(\mathcal O_v\),
\[
 \int_{\mathcal O_v^\times}\nu_v(u)^{-1}\psi_v(u/\varpi_v^d)\,du
 =Q_v^{-d}\sum_{u\bmod\varpi_v^d}^{\times}\nu_v(u)^{-1}\psi_v(u/\varpi_v^d).
\]
Multiplying by \(\nu_v(\varpi_v)^dQ_v^{dz}\) gives \(\epsilon_v(\psi,\nu_v,1-z)\). The standard matched vectors realize this ratio with \(s=1\) on a nonzero lower quotient image. Local uniqueness extends it to the quotient. The common vacuum fixes the normalization off \(T\), and the real intertwiner contributes a nonzero scalar fixed by the real maps. The lower matches are
\[
 (C_2 S e_0\Phi,\mathcal T\Phi')\quad(\beta),\qquad
 (\mathcal T\Phi,C_2 S e_0\Phi')\quad(\alpha)
\]
for \(C_2\ne0\). In the first case the lower primed test is in \(\mathrm{So}(t_1)\). In the second take a full-rank preimage at \(v_1\), which preserves the definite integral by coinvariance. The lower comparison gives both identities with \(C=C_2C_1/C_0\). Zero evaluation surjects onto lower polynomial tests. Lift a nonzero lower unprimed test for \(\beta\), or a nonzero lower primed test for \(\alpha\) and choose a match. The entry with \(C_0\) gives the witness.
\end{proof}
\leanref{PositivityTheory.\allowbreak{}ThKCT.\allowbreak{}genSmall\_\allowbreak{}of}{Development.\allowbreak{}WeilBound.\allowbreak{}ThetaLifts.\allowbreak{}SmallCones.\allowbreak{}GenSmallOf}{18--128}
\leanref{PositivityTheory.\allowbreak{}ThKG.\allowbreak{}Loc.\allowbreak{}Cone.\allowbreak{}Id.\allowbreak{}idealId'}{Development.\allowbreak{}WeilBound.\allowbreak{}ThetaLifts.\allowbreak{}LocalTheory.\allowbreak{}ConeShifts.\allowbreak{}IdealId0}{49--73}
\leanref{PositivityTheory.\allowbreak{}ThKG.\allowbreak{}Std.\allowbreak{}eps\_\allowbreak{}eq}{Development.\allowbreak{}WeilBound.\allowbreak{}ThetaLifts.\allowbreak{}LocalTheory.\allowbreak{}StandardValuations.\allowbreak{}ExistsVol}{85--99}
\leanref{PositivityTheory.\allowbreak{}ThKG.\allowbreak{}Loc.\allowbreak{}locB\_\allowbreak{}of}{Development.\allowbreak{}WeilBound.\allowbreak{}ThetaLifts.\allowbreak{}LocalTheory.\allowbreak{}LocBOf}{17--249}
\leanref{PositivityTheory.\allowbreak{}ThKG.\allowbreak{}Loc.\allowbreak{}locA\_\allowbreak{}of}{Development.\allowbreak{}WeilBound.\allowbreak{}ThetaLifts.\allowbreak{}LocalTheory.\allowbreak{}LocAOf}{17--173}
\leanref{PositivityTheory.\allowbreak{}ThKG.\allowbreak{}keyB\_\allowbreak{}of}{Development.\allowbreak{}WeilBound.\allowbreak{}ThetaLifts.\allowbreak{}ConeHypotheses.\allowbreak{}KeyAOf}{37--70}
\leanref{PositivityTheory.\allowbreak{}ThKG.\allowbreak{}keyA\_\allowbreak{}of}{Development.\allowbreak{}WeilBound.\allowbreak{}ThetaLifts.\allowbreak{}ConeHypotheses.\allowbreak{}KeyAOf}{88--126}

\begin{lemma}[A line test for theta-integral differences]\label{B8.11c}
For a triple as above, let
\[
 \Sigma=\{I(\Phi)-I'(\Phi'):\Phi,\Phi'\text{ polynomial},\ \Phi'\in\mathrm{So}(t_1)\}.
\]
For nonzero totally isotropic rational \(U\subset W\) let \(C_U\) be the average over the unipotent radical of its stabilizer. If \(f\in\Sigma\) and \(C_{Ex}f=0\) on the line of one rational hyperbolic plane, then \(C_Uf=0\) for every such \(U\). If also \(p+q>0\), then \(f=0\).
\end{lemma}
\begin{proof}
Witt transitivity and rational invariance give the line case. Induct on \(j=\dim_E U\). A dual isotropic frame gives the Levi embedding \(m\) of \(\GL_j(\mathbb A_E)\). For \(g\mapsto C_Uf(m(g)h)\), each proper standard block average factors through the constant term for the initial subspace by changing the integration order. These terms vanish by induction. A nonzero restriction has a nonzero Whittaker coefficient at a translate. Expand along successive last columns of the upper unipotent. The zero frequencies are block averages. Rational basis changes take each nonzero frequency to the last coordinate, and Fourier uniqueness on the compact additive quotients gives the assertion. Write \(W_f(h)\) for the coefficient.

Choose a split finite place outside the exceptional places of the tests and frame and outside those of \(h\). Fix the other factors. On each split matrix model the coefficient is a functional of the Schwartz test in \(Y\), with a left determinant character. Write \(Y_i\) for column \(i\). The shear \(Y\mapsto Y+Y_1\rho\), \(\rho_1=0\), acts by \(\xi(\rho_2)\), with \(\xi\) nontrivial additive. For \(j\ge3\) the shear with \(Y_2\) and \(\rho_1=\rho_2=0\) acts by \(\xi(\rho_3)\).

Fix \(t\) with \(\xi(t)\ne1\). For independent \(Y_1,Y_2\) take \(\lambda Y_1=0,\ \lambda Y_2=1\). The shear \(\rho=t\lambda Y\) followed by \(1-tY_1\lambda\) on the left fixes \(Y\) with character \(\xi(t)\). The functional is supported on dependent pairs. If \(j\ge3\) and \(Y_2=sY_1\), \(Y_1\ne0\), add \(tY_2\) to column three and subtract \(tsY_1\) there. If \(Y_1=0\), add \(tY_1\) to column two. Both stabilizers have character \(\xi(t)\). To pass to tests, the compensating row or ratio is continuous on a fixed nonzero minor. For commuting square-zero parts \(A,B\) of the operations, \(L+AL+BL+ABL\) gives a compact open lattice preserved by the operators and their nearby values. Subdivide a compact open test set into small cosets of dilates of such lattices. A coset meeting the indicated locus has a stabilizer with nontrivial character and hence gives zero on the remaining support. This proves \(W_f=0\) for \(j\ge3\).

For \(j=2\), put \(Q=Nv\), let \(d\) be the conductor exponent of \(\xi\), and put \(r=\nu_v(\varpi)Q^{(a-a')/2}\). Let \(D(P,M)\) be the unprimed functional on the box with first two columns bounded by \(\varpi^P\mathcal O_v,\varpi^M\mathcal O_v\) and other entries integral. Set \(\Delta(P,M)=D(P,M)-D(P+1,M)\). If \(M-P<d\), a shear of nontrivial character preserves the box and \(D(P,M)=0\). Partition the shell for \(\Delta\) by the first pivot in column one. Scaling that row by \(\varpi\) gives \(\Delta(P,M)=r\Delta(P+1,M)\) when \(M-P>d\). The residue-count factor multiplying the row character is \(Q^{n-1}/Q^{a+a'-1}\), with \(n=a\) unprimed and \(n=a'\) primed. On dependent pairs the second pivot-row entry determines column two. For \(M-P\ge d\), shears of trivial character exchange its \(Q\) residue classes, so \(\Delta(P,M)=Q\Delta(P,M+1)\). Combining the recurrences at \(P=M-d\) gives \(D(M-d,M)=(Qr)^{-M}D(-d,0)\) for \(M\ge0\). Summing the shells gives the vacuum values up to a scalar,
\[
 Y_i(r)=(Qr)^{-i}\sum_{k=0}^{2i-d}r^k,\qquad i\ge0,
\]
with empty sum zero. The translate \(b_i\) is \(\operatorname{diag}(\varpi^i,\varpi^{-i})\) at the chosen place of \(E\) and the identity elsewhere. The primed parameter is \(r^{-1}\), and reversing the sum gives \(Y_i(r^{-1})=r^dY_i(r)\). Thus \(W_f(m(b_i)h')=c(h')Y_i(r)\) near \(h\), where the local right translate stays integral invertible. By \eqref{eq:p8c2-triple-range}, \(|r|>Q\), so \(|Y_i(r)|\to\infty\).

Restrict \(C_Uf\) to \(\mathrm{SL}_2(\mathbb A_E)\) in the Levi at the right translate \(h\). Its unipotent average vanishes by the line case. More generally, suppose the unipotent average is bounded along the positive Siegel ray uniformly over compact right factors. Convolve on the right with a nonnegative bump of integral one. On the shell of height \(2^k\), rational translation brings each unipotent displacement back to the shell with compact right factor and changes the bump by \(O(2^{-k})\). If the original function is \(O(2^{kM})\) on compact right factors, the difference between its convolution and the convolved unipotent average is \(O(2^{k(M-1)})\), by integration on a fixed compact set. For \(M\ge1\) the bounded average gives the same bound for the convolution. Iterate with arbitrarily small total support. Reduction to the Siegel set makes the resulting convolution bounded on the group. Fubini preserves the unipotent averages and the relation
\[
 Y_0 W_f(m(b_i)h')=Y_i W_f(h')
\]
under convolution supported in the neighborhood above. The smoothed coefficient is bounded, so \(|Y_i|\to\infty\) forces its value at \(h\) to vanish. Shrinking the supports gives \(W_f(h)=0\). The Whittaker criterion completes the induction.

For \(p+q>0\), pair the resulting cusp form with its defining theta integrals. The reverse lift is continuous and rationally invariant. At each completely split good place outside \(T\cup\{v_1\}\), the local calculation gives invariance by the local source group. Its closed stabilizer contains the adelic special unitary group by strong approximation at a split place. The rational points and these split local groups generate a dense subgroup of the norm-one determinant torus, so the stabilizer is the entire source group. On an indefinite source, the rapid decay of \(f\) and the full-rank theta bound give a uniform bound on the constant reverse lift times the quotient height. The quotient height is unbounded, so the lift is zero. At \(v_1\), every coinvariant class has a full-rank representative, which extends the vanishing to other tests. Fubini gives orthogonality. On a definite source a nonzero pairing would, by definite doubling against \(f\), make every polynomial theta integral there cuspidal. Zero evaluation gives their line coefficients as nonzero multiples of lower definite integrals and reaches a nonzero witness. Thus the definite side is orthogonal as well. Subtracting the pairings gives \(\int|f|^2=0\), and continuity gives \(f=0\).
\end{proof}
\leanref{PositivityTheory.\allowbreak{}ThKEnd.\allowbreak{}Cusp.\allowbreak{}Key.\allowbreak{}keyThree\_\allowbreak{}of\_\allowbreak{}locThree}{Development.\allowbreak{}WeilBound.\allowbreak{}ThetaLifts.\allowbreak{}EigenvalueEndgame.\allowbreak{}CuspidalAverages.\allowbreak{}ActsOfReads}{495--505}
\leanref{PositivityTheory.\allowbreak{}ThKEnd.\allowbreak{}Cusp.\allowbreak{}Key.\allowbreak{}keyTwo\_\allowbreak{}of\_\allowbreak{}locTwo\_\allowbreak{}vacBar}{Development.\allowbreak{}WeilBound.\allowbreak{}ThetaLifts.\allowbreak{}EigenvalueEndgame.\allowbreak{}CuspidalAverages.\allowbreak{}ActsOfReads}{629--704}
\leanref{PositivityTheory.\allowbreak{}ThKEnd.\allowbreak{}Cusp.\allowbreak{}kill\_\allowbreak{}coef}{Development.\allowbreak{}WeilBound.\allowbreak{}ThetaLifts.\allowbreak{}EigenvalueEndgame.\allowbreak{}CuspidalAverages.\allowbreak{}KillCoef}{14--60}
\leanref{PositivityTheory.\allowbreak{}ThKEnd.\allowbreak{}Cusp.\allowbreak{}cuspOfLine\_\allowbreak{}of\_\allowbreak{}keys}{Development.\allowbreak{}WeilBound.\allowbreak{}ThetaLifts.\allowbreak{}EigenvalueEndgame.\allowbreak{}CuspidalAverages.\allowbreak{}CuspOfLineOfKeys}{169--198}
\leanref{PositivityTheory.\allowbreak{}ThKEnd.\allowbreak{}zero\_\allowbreak{}of\_\allowbreak{}line}{Development.\allowbreak{}WeilBound.\allowbreak{}ThetaLifts.\allowbreak{}EigenvalueEndgame.\allowbreak{}TowerCodimensionTwo}{73--85}

\section{The integral formulas used in the construction}

\begin{theorem}[Doubling continuation and the definite Siegel--Weil value]\label{B8.12}
Let \(Y\) be Hermitian or skew-Hermitian of dimension \(n\ge1\) over the good base data, with common signature, prescribed Cartan involutions, self-dual lattice and split frames. Fix a doubling frame. Fix the Iwasawa height \(a\) and standard maximal compact \(K\). Let the Hecke character \(\chi\) restrict to \(\eta^c\) for \(c\in\N\). For continuous right \(K\)-finite \(f\) with \(f(px)=\chi(\det t)|\det t|_E^{s_1+n/2}f(x)\) on the Siegel parabolic with Levi coordinate \(t\), the series
\[
 E(s,x)=\sum_\gamma a(\gamma x)^{s-s_1}f(\gamma x)
\]
over rational Siegel-parabolic cosets is holomorphic in a right half-plane. Near each \(s_0\), some \(k\in\N\) gives an extension of \((s-s_0)^kE(s,\cdot)\) to continuous moderate-growth functions. Its evaluation by any linear functional bounded on each class \(|h(x)|\le\operatorname{ht}(x)^N\) is holomorphic there.

Let \(A,Y\) have opposite Hermitian signs and dimensions \(m,n\ge1\) with the stated model conditions. Let \(A\) have common definite signature at infinity. Use \(\chi|_{\mathbb A_F^\times}=\eta^m\) and compatible doubling data with functional \(l\). For \(K\)-finite Schwartz \(\Phi\), the section defined by \(x\mapsto l(\omega(x)\Phi)\) at \(s_0=(m-n)/2\) has holomorphic Eisenstein value
\begin{equation}\label{eq:p8c2-definite-sw}
 E(s_0,x;\Phi)=
 \begin{cases}
 \operatorname{avg}_{[\U(A)]}\theta(x,g;\Phi),&n<m,\\
 2\operatorname{avg}_{[\U(A)]}\theta(x,g;\Phi),&n\ge m
 \end{cases}
\end{equation}
for any Haar measure and fundamental domain.
\end{theorem}
\leanref{GlobalTheory.\allowbreak{}ThK.\allowbreak{}Dbl.\allowbreak{}EisBL}{Development.\allowbreak{}WeilBound.\allowbreak{}ThetaLifts.\allowbreak{}LatticeBasics.\allowbreak{}HypothesesE}{46--53}
\leanref{GlobalTheory.\allowbreak{}ThK.\allowbreak{}Dbl.\allowbreak{}SiegelWeilStd}{Development.\allowbreak{}WeilBound.\allowbreak{}ThetaLifts.\allowbreak{}LatticeBasics.\allowbreak{}HypothesesS}{11--21}
\leanref{UpperTheory.\allowbreak{}LiteratureTheorems.\allowbreak{}eis\_\allowbreak{}bl\_\allowbreak{}of\_\allowbreak{}literature}{Development.\allowbreak{}AutomorphicArithmetic.\allowbreak{}ThetaCorrespondence.\allowbreak{}ThetaPairLifts.\allowbreak{}ThirdPlaceSeries.\allowbreak{}EveryCharacterOfLiteratureOfVWin}{561--562}
\leanref{UpperTheory.\allowbreak{}LiteratureTheorems.\allowbreak{}siegel\_\allowbreak{}weil\_\allowbreak{}std\_\allowbreak{}of\_\allowbreak{}literature}{Development.\allowbreak{}AutomorphicArithmetic.\allowbreak{}ThetaCorrespondence.\allowbreak{}ThetaPairLifts.\allowbreak{}ThirdPlaceSeries.\allowbreak{}EveryCharacterOfLiteratureOfVWin}{563--564}

The continuation specializes Langlands's theorem to unitary doubling sections. \cite{langlands-1976}
Definiteness gives Witt index zero, so Gan--Qiu--Takeda's anisotropic formula applies with the quotient-average normalization. \cite[Theorem~7.1 and \S3.1]{gan-qiu-takeda-2014}
\leanref{PositivityTheory.\allowbreak{}SiegelWeilFormula.\allowbreak{}Fin.\allowbreak{}siegelWeilStd}{Development.\allowbreak{}WeilBound.\allowbreak{}CoreStatements.\allowbreak{}AutomorphicKernels.\allowbreak{}SiegelWeilFormula.\allowbreak{}SiegelWeilStandard.\allowbreak{}SiegelWeilStd}{8--8}

\begin{theorem}[Convergent Siegel--Weil and Rallis formulas]\label{B8.13}
For the first and third formulas let \(A,B\) be Hermitian of dimensions \(a,b\) on the same datum \(D\) of type \Lean{D7}, and let \(r_B\) be the Witt index of \(B\). For the first formula assume
\begin{equation}\label{eq:p8c2-convergent-range}
 a<b\le2a,\qquad a+r_B<b.
\end{equation}
Use compatible Iwasawa and Cartan data, doubled Weil characters \(1,\chi_B\), and rational additive covolume one. For every Haar measure on \(\U(B)(\mathbb A_F)\) and allowed Fock vector \(\Phi\), the descended kernel is integrable at every doubled group point \(g\). Its quotient average is \(E(g,f_\Phi,s_0)\) for the flat section \(f_\Phi\), with holomorphy at \(s_0=(b-a)/2\).

Over a CM extension let \(J,T\) be nonsingular Hermitian matrices of sizes \(M,n\) with \(2n<M\). Then \([\U(J)]\) has finite volume. Take \(\Phi=\Phi_\infty\prod_{v<\infty}\Phi_v\) with smooth rapidly decreasing compactly supported infinite factor and finite Schwartz--Bruhat factors. Fix an \(F\)-basis of \(E\). At almost every finite \(v\), let \(\Phi_v\) be the indicator of tuples whose coordinates in this basis lie in \(\mathcal O_v\). Admissible Haar measures and fundamental domains exist, as do nontrivial characters of \(\mathbb A_F/F\) with self-dual measures on tuples and Hermitian matrices. For every such choice the Gram pushforwards have continuous densities \(g_\infty,g_v\) on the nonsingular locus, and
\begin{equation}\label{eq:p8c2-density}
 \operatorname{avg}_{[\U(J)]}\sum_{\operatorname{Gram}(X)=T}\Phi(g^{-1}X)
 =g_\infty(T)\prod_{v<\infty}g_v(T)
\end{equation}
with convergent product.

For the inner product require \(a<b\le2a\) and either \(r_B=0\) or \(a+r_B<b\). Use compatible Weil models \(\omega,\omega_d\), Iwasawa and covolume-one data, and Haar measures, right invariant on \(\U(A)(\mathbb A_F)\). Fix the infinite-place realization and a Cartan frame of \(A\) with prescribed signatures. Let \(\pi\) be irreducible cuspidal on \(\U(A)\). Its convergent lifts to spaces of the prescribed signatures obtained by removing \(i\ge1\) hyperbolic planes from \(B\) must vanish for compatible models and the specified smooth forms. For allowed vectors \(\Phi_1,\Phi_2\) require the doubled vector
\[
 \delta(\Phi_1,\Phi_2)(x)
 =\Phi_1(\operatorname{res}_L x)\overline{\Phi_2(-\operatorname{res}_R x)}
\]
to be \(K\)-finite, where \(\operatorname{res}_L,\operatorname{res}_R\) are the two restrictions. For \(f_1,f_2\) in the smooth realization the integrals \(\theta_{\Phi_i} f_i(h)=\int_{[\U(A)]}\theta_\omega(g,h;\Phi_i)\overline{f_i(g)}\,dg\) converge and belong to \(L^2([\U(B)])\). The doubled zeta integral of the flat section satisfies
\begin{equation}\label{eq:p8c2-rallis}
 \operatorname{CT}_{s=s_0}Z_{\chi_B^{-1}}(s,\delta(\Phi_1,\Phi_2),f_1,f_2)
 =\operatorname{vol}[\U(B)]^{-1}
 \left\langle\theta_{\Phi_1}f_1,\theta_{\Phi_2}f_2\right\rangle_{[\U(B)]}.
\end{equation}
Here \(\operatorname{CT}\) denotes the constant Laurent coefficient and \(Z_\mu\) uses \(\mu(\det g_2)^{-1}\) at its second group argument.
\end{theorem}
\leanref{GlobalTheory.\allowbreak{}Published.\allowbreak{}SiegelWeilUnitaryConvergent}{Development.\allowbreak{}AbelianVarietyTheory.\allowbreak{}PublishedTheorems.\allowbreak{}SiegelWeil₂.\allowbreak{}SiegelWeilUnitaryConvergent}{18--32}
\leanref{HeckeCharacterZeroLocus.\allowbreak{}SiegelWeilU.\allowbreak{}SiegelWeilUnitaryConvergent}{Formalism.\allowbreak{}UnipotentAdelicQuotients.\allowbreak{}SiegelWeilUnitaryConvergent}{46--49}
\leanref{HeckeCharacterZeroLocus.\allowbreak{}SiegelWeilU.\allowbreak{}SiegelDensities}{Formalism.\allowbreak{}UnipotentAdelicQuotients.\allowbreak{}SiegelWeilUnitaryConvergent}{18--40}
\leanref{GlobalTheory.\allowbreak{}Published.\allowbreak{}RallisInnerProductConvergent}{Development.\allowbreak{}AbelianVarietyTheory.\allowbreak{}PublishedTheorems.\allowbreak{}TwistedTraceFormula.\allowbreak{}LowerLiftsZeroOfWittIdxEqZero}{41--59}
\leanref{UpperTheory.\allowbreak{}LiteratureTheorems.\allowbreak{}siegel\_\allowbreak{}weil\_\allowbreak{}unitary\_\allowbreak{}convergent\_\allowbreak{}of\_\allowbreak{}literature}{Development.\allowbreak{}AutomorphicArithmetic.\allowbreak{}ThetaCorrespondence.\allowbreak{}ThetaPairLifts.\allowbreak{}ThirdPlaceSeries.\allowbreak{}EveryCharacterOfLiteratureOfVWin}{481--482}
\leanref{UpperTheory.\allowbreak{}LiteratureTheorems.\allowbreak{}siegelWeilU\_\allowbreak{}unitary\_\allowbreak{}convergent\_\allowbreak{}of\_\allowbreak{}literature}{Development.\allowbreak{}AutomorphicArithmetic.\allowbreak{}ThetaCorrespondence.\allowbreak{}ThetaPairLifts.\allowbreak{}ThirdPlaceSeries.\allowbreak{}EveryCharacterOfLiteratureOfVWin}{581--582}
\leanref{UpperTheory.\allowbreak{}LiteratureTheorems.\allowbreak{}rallis\_\allowbreak{}inner\_\allowbreak{}product\_\allowbreak{}convergent\_\allowbreak{}of\_\allowbreak{}literature}{Development.\allowbreak{}AutomorphicArithmetic.\allowbreak{}ThetaCorrespondence.\allowbreak{}ThetaPairLifts.\allowbreak{}ThirdPlaceSeries.\allowbreak{}EveryCharacterOfLiteratureOfVWin}{485--486}

The first formula specializes Yamana's Theorem~2 to \(\delta=1,\ m=b,\ n=a,\ r=r_B\) in the range \(b-r_B>a\). \cite[p.~237]{yamana-2013}
The Gram identity specializes Weil's orbital formula to nonsingular unitary fibers with the measure of \S44, Theorem~2. \cite[\S52, Theorem~5, equation~(41), p.~77]{weil-1965}
The third formula uses the Rallis identity in the convergent range with lower-lift vanishing. \cite[\S3.1 and Propositions~11.1(ii) and~11.2]{gan-qiu-takeda-2014}

At fixed finite level approximate the doubled vector by Fock vectors in the Sobolev seminorms and average over \(K\) with a finite matrix-coefficient kernel reproducing its \(K\)-span. The projected sections restrict to a fixed matrix-coefficient space on \(K\). Iwasawa determines the sections from those restrictions, so they lie in a common finite-dimensional space. Sobolev bounds give pointwise convergence of the sections and theta averages. The graph of the Eisenstein value on this finite-dimensional section space is closed in the pointwise topology, by linearity and uniqueness of the value. The first formula therefore extends to the doubled vector. The product of the two theta kernels is the doubled kernel times \(\chi_B(\det g_2)\). Integrating against the cusp forms gives the displayed zeta value with the inverse quotient volume.
\leanref{UpperTheory.\allowbreak{}Lit.\allowbreak{}Rallis.\allowbreak{}RallisIPProof.\allowbreak{}rallis\_\allowbreak{}of\_\allowbreak{}dm}{Development.\allowbreak{}AutomorphicArithmetic.\allowbreak{}ThetaCorrespondence.\allowbreak{}FockModels.\allowbreak{}RallisDM}{24--110}

\section{Vacuum vectors and separation}

\begin{lemma}[Unramified vacuum generation]\label{B8.14}
Let \(L_{\rm loc}\) be a complete discretely valued field of characteristic zero, with integers \(\mathcal O\), uniformizer \(\varpi\) and residue cardinality \(q^2\) for odd \(q\ge2\). Its involution \(\sigma\) preserves the valuation, fixes \(\varpi\) and has \(\sigma u-u\in\mathcal O^\times\) for some \(u\in\mathcal O\). Let \(J_S\in\GL_n(\mathcal O)\), \(J_T\in\GL_{n'}(\mathcal O)\) be Hermitian and let \(\delta\in\mathcal O^\times\) satisfy \(\sigma\delta=-\delta\). Let \(\psi:L_{\rm loc}\to\C^\times\) be additive, trivial on \(\mathcal O^\sigma\) and nontrivial on \(\varpi^{-1}\mathcal O^\sigma\).

On a complex space \(M\) take operators \(\rho(x)\) for \(x\in\operatorname{Mat}_{n'\times n}(L_{\rm loc})\) and a representation \(\omega\) of \(\U(J_S)\times\U(J_T)\). With \(H(x,y)=\delta\operatorname{tr}(J_S^{-1}\,{}^t\!\sigma(y)J_Tx)\), assume
\[
 \rho(0)=1,\qquad
 \rho(x)\rho(y)=\psi((H(x,y)+\sigma H(x,y))/2)\rho(x+y),\qquad
 \omega(g,h)\rho(x)=\rho(hxg^{-1})\omega(g,h).
\]
Translates of integral-translation-fixed vectors span \(M\), and the integral unitary groups \(K_S,K_T\) fix each such vector. Suppose one such \(s\ne0\) has translates spanning \(M\). Then
\[
 M^{K_S}\subseteq\operatorname{span}_\C\{\omega(1,h)s:h\in\U(J_T)\},\qquad
 M^{K_T}\subseteq\operatorname{span}_\C\{\omega(g,1)s:g\in\U(J_S)\}.
\]
\end{lemma}
\leanref{GlobalTheory.\allowbreak{}Hk3.\allowbreak{}howeUnram\_\allowbreak{}holds}{Development.\allowbreak{}AbelianVarietyTheory.\allowbreak{}Proofs.\allowbreak{}SatakeParameters.\allowbreak{}RhoMemOfMkEq}{250--256}
\leanref{GlobalTheory.\allowbreak{}Hk3.\allowbreak{}HoweUnram}{Development.\allowbreak{}AbelianVarietyTheory.\allowbreak{}PublishedTheorems.\allowbreak{}GaloisGeometry₂.\allowbreak{}HoweUnram}{8--19}

\begin{proof}
Put \(\Lambda=\operatorname{Mat}_{n'\times n}(\mathcal O)\). On \(V_{\bar x}=\C\rho(x)s\) the translation by \(u\in\Lambda\) acts by \(\psi(H(u,x)+\sigma H(u,x))\). Unimodularity and the conductor distinguish classes modulo \(\Lambda\). Hence \(M=\bigoplus_{\bar x}V_{\bar x}\) with trivial-character line \(\C s\).

The group \(K_T\) permutes the lines. Put \(S=\operatorname{span}_\C\{\omega(g,1)s:g\in\U(J_S)\}\subseteq M^{K_T}\). On an orbit supporting an invariant, every stabilizer \(h\) of \(\bar x\) acts trivially on \(V_{\bar x}\). For \(u=hx-x\in\Lambda\), covariance gives the phase \(\psi(-(H(x,u)+\sigma H(x,u))/2)=1\). Write \(D_x=xJ_S^{-1}\,{}^t\!\sigma(x)\), \(C^*={}^t\!\sigma(C)\). The Cayley transform yields \(\psi((z+\sigma z)/2)=1\) for \(z=\delta\operatorname{tr}(CD_x)\) whenever \(C^*=-C,\ C_{ij}\in\varpi\mathcal O\) and \(Cx\) is integral. The integral-Gram lemma then gives an integral \(D_x\) after replacing \(x\) modulo \(\Lambda\).

The finite Gram-class construction gives representatives \(Y\) containing \(x\) and characters \(b\mapsto\chi_b(y)\) of a finite phase group \(B\). Its shifts give the self-dual lattice vacua \(w_b=\sum_{y\in Y}\chi_b(y)\rho(y)s\), whose lattice translations permute the summands with the Heisenberg phases. Write the lattice frame as \(gk\) with \(g\in\U(J_S)\) and \(k\) integral invertible. Then \(\omega(g^{-1},1)w_b\in\C s\), hence \(w_b\in S\). Averaging gives
\[
 |B|^{-1}\sum_b\chi_b(x)^{-1}w_b
 =\sum_{\chi_\bullet(y)=\chi_\bullet(x)}\rho(y)s\in S.
\]
For retained \(y\), \Lean{HomA} gives \(\ker_{\mathcal O}(x)\subseteq\ker_{\mathcal O}(y)\) modulo the target lattice, and \(yx^\ast\) is integral for the \(J_S,J_T\)-adjoint. By \Lean{GramEq}, if \(m,m'\in\mathcal O^n,\ c\in\mathcal O\) and \(x(cm),x(cm')\) are integral, then
\[
 c\bigl(\langle ym,ym'\rangle_{J_T}-\langle xm,xm'\rangle_{J_T}\bigr)\in\mathcal O.
\]
For \(\mu(\bar x)=\#\operatorname{im}(x:\mathcal O^n\to L_{\rm loc}^{n'}/\mathcal O^{n'})\) this gives \(\mu(\bar y)\le\mu(\bar x)\). At equality the Gram-class isometry gives \(y-hx\in\Lambda\) for some \(h\in K_T\). The lower-height part of the average is invariant and belongs to \(S\) by induction. The remaining part is nonzero on \(\bar x\) and spans the orbit invariants, which are determined on that line. Thus \(M^{K_T}\subseteq S\). The Hermitian-adjoint model exchanges the spaces.
\end{proof}
\leanref{GlobalTheory.\allowbreak{}Hk3.\allowbreak{}locVac\_\allowbreak{}span\_\allowbreak{}classes\_\allowbreak{}indep}{Development.\allowbreak{}AbelianVarietyTheory.\allowbreak{}Proofs.\allowbreak{}SatakeParameters.\allowbreak{}LatticeModelVacuumTranslateClassesIndependent}{8--98}
\leanref{GlobalTheory.\allowbreak{}Hk3.\allowbreak{}HoweG.\allowbreak{}gramClass\_\allowbreak{}orbit\_\allowbreak{}mod\_\allowbreak{}lower}{Development.\allowbreak{}AbelianVarietyTheory.\allowbreak{}Proofs.\allowbreak{}SatakeParameters.\allowbreak{}ExistsConjugateModelExtremalDualParameter}{16--129}
\leanref{GlobalTheory.\allowbreak{}Hk3.\allowbreak{}HoweDuality.\allowbreak{}fixed\_\allowbreak{}mem\_\allowbreak{}span}{Development.\allowbreak{}AbelianVarietyTheory.\allowbreak{}Proofs.\allowbreak{}SatakeParameters.\allowbreak{}RhoMemOfMkEq}{145--240}

For a real realization \(R\) put \(\mathcal V=\operatorname{SmT}(R.MY_0)\), \(\mathcal E_0=R.EEat(R.triv)\). Write \(t(w)\) for the tensor-product linearization at \(w\in\mathcal V\). Let \(H\) be the top definite group with embedding \(\iota\) into the last row group, and \(N\) the Bessel unipotent with character \(\psi_N\). For real test \(\rho\) and grid datum \(\Phi\) let \(F_{\rho,\Phi,u}\) be the successive theta lift of \(u\). For coefficient dual \(\xi\), top column \(\Psi\) and coefficient action \(\rho_{\rm coef}\), put \(\phi_{\xi,\Psi}(h)=\xi(\rho_{\rm coef}(h^{-1})\Psi(h_f))\), where \(h_f\) is the finite component. With the given quotient measures define
\[
 \mathcal B_{\rho,\xi,\Psi,\Phi}(u)=
 \int_{[H]}\int_{[N]}\phi_{\xi,\Psi}(h)
 F_{\rho,\Phi,u}(n\iota(h))\overline{\psi_N(n)}\,dn\,dh .
\]

The conditions \Lean{SizeRows} use weight-gap constant \(10\) and the prescribed rank and weight data \(N(m),d,S_c,\operatorname{gap}(D,1)\), with
\begin{gather}
 m^4\le f,\quad 44\le m,\quad
 x^2m^3L\log\operatorname{rd}(E)\le k,\quad fmN(m)^2d\le b,
 \label{eq:p8c2-size-rows-a}\\
 fm\le(r+1)^3,\quad \log g_i\le9L+1,\quad \log q_0\le L,
 \label{eq:p8c2-size-rows-b}\\
 \log\operatorname{rd}(E)\le\log\operatorname{cond}(E)\le L/100,\quad
 S_c+(m+1)\operatorname{gap}(D,1)+\tfrac12\le b .
 \label{eq:p8c2-size-rows-c}
\end{gather}
Require \(\mathfrak c_K=\mathcal O_K\) and \(b_{\tau,l}>0,\ e_{\tau,l}\le4b\) at every infinite place \(\tau\) and weight position \(l\).

\begin{theorem}[Real functionals separate rows]\label{B8.15}
Assume \(\mathscr P\), Lemma~\ref{B8.14} and the first formula of Theorem~\ref{B8.13}. For every \(\varepsilon\in\R,\ C_F\ge6510\) there is \(x_0\) such that for data of exponent \(\varepsilon\) at \(x\ge x_0\) with \Lean{SizeRows} and finite-group series compatibility, every real realization \(R\) with the level and group-record conditions, grid at \((1,C_F,b)\), Weil system and compatible row measures satisfies
\[
 \forall w\in\mathcal V\setminus\{0\}\quad
 \exists\rho,\xi\quad
 \forall t\in\mathcal E_0\setminus\{0\}\quad
 \exists\Psi,\Phi\quad \mathcal B_{\rho,\xi,\Psi,\Phi}(t(w))\ne0
\]
with convergence of the lifts and both integrations.
\end{theorem}
\leanref{UpperTheory.\allowbreak{}Lit.\allowbreak{}functionalsSeparate\_\allowbreak{}of\_\allowbreak{}howeUnram\_\allowbreak{}of\_\allowbreak{}siegelWeil}{Development.\allowbreak{}AutomorphicArithmetic.\allowbreak{}ThetaCorrespondence.\allowbreak{}ThetaPairLifts.\allowbreak{}LocalTypeDictionary.\allowbreak{}FunctionalsSeparateOfHoweUnramOfSiegelWeil}{17--19}
\leanref{UpperTheory.\allowbreak{}Lit.\allowbreak{}real\_\allowbreak{}functionals\_\allowbreak{}separate\_\allowbreak{}of\_\allowbreak{}hypothesis}{Development.\allowbreak{}AutomorphicArithmetic.\allowbreak{}ThetaCorrespondence.\allowbreak{}ThetaPairLifts.\allowbreak{}LocalTypeDictionary.\allowbreak{}RealFunctionalsSeparateOfHypothesis}{15--30}
\leanref{GlobalTheory.\allowbreak{}IndefiniteRowSideFullHecke.\allowbreak{}LevelPlusSubgroup.\allowbreak{}FunctionalsSeparate}{Development.\allowbreak{}AbelianVarietyTheory.\allowbreak{}BlockTheory.\allowbreak{}DivisorBase.\allowbreak{}LevelPlusSubgroup}{83--91}

\begin{proof}
For \(t,w\ne0\), branching and \Lean{exists\_\allowbreak{}kf} give a core form \(F\) with \(\operatorname{Per}_0(F,t(w))\ne0\). Doubling gives surjectivity on the core by the first formula of Theorem~\ref{B8.13} and Lemma~\ref{B8.14}. Hence \(F\) lies in the span of theta-lifted column chains. A chain with nonzero period gives by seesaw ascent
\begin{equation}\label{eq:p8c2-weak-separation}
 \forall t,w\ne0\quad\exists\rho,\xi,\Psi,\Phi\quad
 \mathcal B_{\rho,\xi,\Psi,\Phi}(t(w))\ne0 .
\end{equation}
Fix \(w\ne0\). The lift evaluates a multiplicity map depending on \(t,\Phi\) at the common top tensor \(v_\rho(w)\). At each real place opposite-sign Harish--Chandra parameters differ by at least \(120(\operatorname{rank}V_{\mathrm{top\ row}})^3\). The Bessel-form space with definite coefficients therefore has dimension at most one. Separate equivariance on pure tensors gives a scalar multiple of the product \(\ell\) of local spanning forms, with zero if a local space is zero. For \(A_w(\rho,\xi)=\ell(\xi,v_\rho(w))\),
\begin{equation}\label{eq:p8c2-common-real-factor}
 \mathcal B_{\rho,\xi,\Psi,\Phi}(t(w))=A_w(\rho,\xi)c_{t,\Psi,\Phi},
\end{equation}
where \(c_{t,\Psi,\Phi}\) is independent of \(\rho,\xi\). If \(\mathcal E_0=0\), any pair works. Otherwise \eqref{eq:p8c2-weak-separation} at one \(t_0\ne0\) fixes a pair with \(A_w\ne0\). At each \(t\ne0\) the same equation supplies a nonzero \(c_{t,\Psi,\Phi}\). Take the larger period and gap threshold, and use the tower bound for convergence.
\end{proof}
\leanref{GlobalTheory.\allowbreak{}SepRows.\allowbreak{}weakSep\_\allowbreak{}of\_\allowbreak{}steps}{Development.\allowbreak{}AbelianVarietyTheory.\allowbreak{}Proofs.\allowbreak{}TwistedOperators.\allowbreak{}RowSeparationByRealFunctionalsOfChainSteps}{71--85}
\leanref{UpperTheory.\allowbreak{}Lit.\allowbreak{}Bur.\allowbreak{}separate\_\allowbreak{}of\_\allowbreak{}burGapPos}{Development.\allowbreak{}AutomorphicArithmetic.\allowbreak{}ThetaCorrespondence.\allowbreak{}BesselModels.\allowbreak{}SepGap}{222--257}

\section{Twisted square models}

Fix \(D:\Lean{D7}\) with CM extension \(E/F\) and infinite-place realization \(\theta\). Let \(W,W_1,X\) have ranks \(u,v,d=u+v\) with \(d\) even and signatures \((u-p,p),(v-q,q),(d-p-q,p+q)\). Require \(v+2(2(p+q)+1)<u\) and \(2(p+q)+1<v\). The unitary characters \(\alpha,\beta\) restrict to \(\eta^v,\eta^u\) and have rational infinity exponents \(2(a_\tau-b_\tau)<-d\). At each finite component both are unramified or their quotient is ramified at a split component prime to the different of \(F/\Q\). Let \(\mathcal P(W,X)\) be the polynomial Fock vectors. Algebraic Fock functions are finite sums of their products with permitted finite factors.

\begin{theorem}[Convergence and matching for twisted squares]\label{B8.16}
Assume \(\mathscr P\), Theorem~\ref{B8.11} and the unitary Siegel-set theorem. For either Hermitian sign and compatible Weil data for \(M,X\), suppose at some infinite place \(\operatorname{sig}(M)=(\operatorname{rank}M-p,p)\) and \(\operatorname{rank}X<2(\operatorname{rank}M-p)\). For every permitted theta vector \(\Phi\), Haar measure and fundamental domain,
\begin{equation}\label{eq:p8c2-square-moderate}
 \int_{[\U(M)]}|\theta_\omega(g,h;\Phi)|\,dg\le C H_X(h)^N
 \qquad\text{for all }h
\end{equation}
for some \(C\in\R,\ N\in\N\) and the prescribed height \(H_X\).

Fix the standard additive character of sign \(-1\), frames and exponents \(a_\tau,b_\tau\). There are surjections \(q,q'\) from \(\mathcal P(W,X),\mathcal P(W_1,X)\) to a common \(P_{\mathrm{real}}\). Matching data exist for each twisted-square triple with these exponents, with signs \(+,-\) and splitting characters \((\beta,\beta/\alpha),(\alpha,\alpha/\beta)\). The common quotient \(\Pi\ne0\) is irreducible for the joint finite adelic and real Lie-algebra actions of \(\U(X)\). The quotient maps \(p,p'\) on stable algebraic Fock spaces intertwine these actions and kill the source finite differences and infinitesimal action. For each match there are an isomorphism \(e:P_{\mathrm{real}}\otimes_\C P_{\mathrm{fin}}\to\Pi\) and linear maps \(r,r'\) onto \(P_{\mathrm{fin}}\) on permitted finite functions, with
\[
 p(\operatorname{Fock}(f,P))=e(q(P)\otimes r(f)),\qquad
 p'(\operatorname{Fock}'(f',P'))=e(q'(P')\otimes r'(f')) .
\]

For any additive character with a twisted-square triple, matching data and compatible measures and domains, there is \(c\ne0\) such that every primed algebraic Fock function \(\Phi'\) has an algebraic representative \(\Phi_0\) modulo the finite coinvariant relations \(\mathcal R_{\rm fin}\) satisfying \eqref{eq:p8c2-square-moderate}. For every unprimed algebraic Fock function \(\Phi\),
\[
 p(\Phi)=p'(\Phi_0)\quad\Longrightarrow\quad
 \theta\operatorname{Int}_\omega(\Phi)=c\,\theta\operatorname{Int}_{\omega'}(\Phi_0).
\]
\end{theorem}
\leanref{GlobalTheory.\allowbreak{}WeilConv.\allowbreak{}twisted\_\allowbreak{}weil\_\allowbreak{}representation\_\allowbreak{}convergence}{Development.\allowbreak{}AutomorphicArithmetic.\allowbreak{}ThetaCorrespondence.\allowbreak{}ThetaPairLifts.\allowbreak{}LocalTypeDictionary.\allowbreak{}TwistedWeilRepresentationConvergence}{16--17}
\leanref{GlobalTheory.\allowbreak{}T.\allowbreak{}twisted\_\allowbreak{}local\_\allowbreak{}factors\_\allowbreak{}match}{Development.\allowbreak{}AutomorphicArithmetic.\allowbreak{}ThetaCorrespondence.\allowbreak{}ThetaPairLifts.\allowbreak{}LocalTypeDictionary.\allowbreak{}TwistedSquareLocalMatching}{88--89}
\leanref{GlobalTheory.\allowbreak{}T.\allowbreak{}twisted\_\allowbreak{}square\_\allowbreak{}theta\_\allowbreak{}integrals\_\allowbreak{}eq}{Development.\allowbreak{}AutomorphicArithmetic.\allowbreak{}ThetaCorrespondence.\allowbreak{}ThetaPairLifts.\allowbreak{}LocalTypeDictionary.\allowbreak{}TwistedSquareThetaIntegralsEq}{48--49}

\begin{proof}
Write \(n=\operatorname{rank}M,\ d_X=\operatorname{rank}X,\ f_F=[F:\Q]\) and let \(r\le p\) be the Witt index. On a Siegel set the absolute kernel is bounded by \(C H_X(h)^N\prod_{i<r}t_i^{f_Fd_X}\), and the measure denominator is \(\prod_{i<r}t_i^{1+2f_F(n-2i-1)}\). On \(t_i\ge c t_{i+1},\ t_{r-1}\ge c>0\), the resulting exponents \(A_i-1\) give
\[
 \sum_{i=0}^j A_i=f_F(j+1)(d_X-2n+2(j+1))<0\quad(0\le j<r).
\]
Successively integrating from \(t_0\), using \(\int_b^\infty t^{A-1}dt=b^A/(-A)\), gives a finite constant. For \(r=0\) the cone factor is one.

If the exponents admit no match, take \(P_{\mathrm{real}}=0\). Otherwise real Howe duality and equality of kernels give a common simple quotient \(P_{\mathrm{real}}\). It has at most countable dimension, so Dixmier's lemma gives scalar endomorphisms. For each permitted \(f\), \(P\mapsto p(\operatorname{Fock}(f,P))\) descends to a real intertwiner, as does its primed counterpart. On each side their images span \(\Pi\), which is then a direct sum of simple copies of \(P_{\mathrm{real}}\). An intertwiner has finite support in that sum, by evaluation on a nonzero generator, and has scalar coordinates. Thus the intertwiners from either side span the same space \(P_{\mathrm{fin}}\). Evaluation is surjective. A kernel on a finite independent family of intertwiners would contain a simple copy whose scalar coordinates give a linear relation. Hence evaluation is an isomorphism and the factor maps extend linearly.

Choose \(\Phi_0\) modulo \(\mathcal R_{\rm fin}\) with seesaw representative in \(\mathrm{So}(t_1)\). Theorem~\ref{B8.11} gives the bound and scalar \(c_1\ne0\). Equality in \(\Pi\) matches the first seesaw representative with \(s\) times the second for the twisted-kernel scalar \(s\ne0\). Hence \(c=c_1s\).
\end{proof}
\leanref{GlobalTheory.\allowbreak{}WeilConv.\allowbreak{}weilConvergence\_\allowbreak{}of\_\allowbreak{}siegelSet}{Development.\allowbreak{}AbelianVarietyTheory.\allowbreak{}Proofs.\allowbreak{}TwistedOperators.\allowbreak{}WeilConvergenceOfSiegelSet}{24--82}
\leanref{GlobalTheory.\allowbreak{}SiegelSet.\allowbreak{}lintegral\_\allowbreak{}cone\_\allowbreak{}modulus\_\allowbreak{}lt\_\allowbreak{}top}{Development.\allowbreak{}AbelianVarietyTheory.\allowbreak{}Proofs.\allowbreak{}TwistedOperators.\allowbreak{}LintegralSiegelConeInvModulusLtTopWeil}{160--207}
\leanref{GlobalTheory.\allowbreak{}SquareLM.\allowbreak{}twisted\_\allowbreak{}square\_\allowbreak{}local\_\allowbreak{}matching}{Development.\allowbreak{}AutomorphicArithmetic.\allowbreak{}ThetaCorrespondence.\allowbreak{}ThetaPairLifts.\allowbreak{}LocalTypeDictionary.\allowbreak{}TwistedSquareLocalMatching}{18--81}
\leanref{GlobalTheory.\allowbreak{}SquareAgree.\allowbreak{}twisted\_\allowbreak{}theta\_\allowbreak{}integrals\_\allowbreak{}agree}{Development.\allowbreak{}AutomorphicArithmetic.\allowbreak{}ThetaCorrespondence.\allowbreak{}ThetaPairLifts.\allowbreak{}LocalTypeDictionary.\allowbreak{}TwistedSquareThetaIntegralsEq}{17--41}

\section{Fields of definition and finite tests}

A complete test choice fixes a grid and Weil system, row measures, a nonzero smooth tensor and real tests, square constants and a finite test lattice. Let \(P,f\) be the preliminary datum and frame, and \(\iota_X:K_X\hookrightarrow\C\) the coefficient embedding of a completion \(X\).

\begin{theorem}[A finite field requirement for a finite character set]\label{B8.17}
Assume \(\mathscr P,\mathscr J\), the Hecke--Tate functional equation and Theorems~\ref{B8.11}, \ref{B8.15}, \ref{B8.16}. There is \(C_*\in\R\) such that for every \(\varepsilon\in\R,\ C_F\ge C_*\) there is a threshold \(x_0\). Given a scale \(s\) of exponent \(\varepsilon\) with \(x(s)\ge x_0\), datum \(P\), frame \(f\) and finite set \(\Xi\) of Hecke characters of \(E\), there is a finite \(S\subset\overline{\Q}\cap\C\). For each completion \(X\) of \(P,f\) with \(S\subseteq\iota_X(K_X)\), size datum over \(X\) at \(s\), and compatible Hida realization \(R\) with frame \(f\), twisting places away from \(q_0\) and characters in \(\Xi\), a complete test choice at \((1,C_F,b(s))\) satisfies the field condition.
\end{theorem}
\leanref{GlobalTheory.\allowbreak{}T.\allowbreak{}theta\_\allowbreak{}lift\_\allowbreak{}extension\_\allowbreak{}field\_\allowbreak{}exists}{Development.\allowbreak{}AutomorphicArithmetic.\allowbreak{}ThetaCorrespondence.\allowbreak{}ThetaPairLifts.\allowbreak{}LocalTypeDictionary.\allowbreak{}ThetaLiftExtensionFieldExists}{50--51}
\leanref{GlobalTheory.\allowbreak{}RowEx.\allowbreak{}FieldFor}{Development.\allowbreak{}AbelianVarietyTheory.\allowbreak{}Definitions.\allowbreak{}TwistedInnerProducts.\allowbreak{}TwistPlacesOffPrimeQZero}{231--236}
\leanref{GlobalTheory.\allowbreak{}RowEx.\allowbreak{}ExField}{Development.\allowbreak{}AbelianVarietyTheory.\allowbreak{}Definitions.\allowbreak{}TwistedInnerProducts.\allowbreak{}TwistPlacesOffPrimeQZero}{241--242}

\begin{proof}
Take the largest parameter and scale thresholds. The level key is \(\kappa(R)=\{(\xi_t,Q_t):t\in T_R\}\), with \(Q_t\subseteq\operatorname{Ram}(\xi_t)\subseteq Q_t\cup\overline{Q_t}\). Finiteness of \(\Xi\) gives finitely many pairs and keys. For each realized key take a representative with grid \(G\) and Weil system. Complete the choice by separation and convergent square matching. The root field contains eighth roots of unity and roots of order \((Nw)^{\mathrm{depth}(w)}\) on \(\bigcup_t\mathrm{locPl}(G,t)\). It also contains \(\sqrt{Nw}\) there and on the grid-choice set \(PA\). Row sums and the dictionary identities give the phase and jet fields. Toric transport and the Hecke inputs give the normalized \(L\)-ratio field. Form their compositum \(N_\kappa\) with the character-value and bottom norm-sum fields.

In another realization of that key, choose an index for each pair. Injectivity identifies the base indices through the trivial character. Reindex and identify frames to transport the grid, Weil system and underlying choice. Completion preserves the transported evaluations. The union of rational bases of the \(N_\kappa\) supplies \(S\), since a field containing a basis contains the corresponding number field.
\end{proof}
\leanref{GlobalTheory.\allowbreak{}RowEx.\allowbreak{}exKey\_\allowbreak{}of}{Development.\allowbreak{}AbelianVarietyTheory.\allowbreak{}Proofs.\allowbreak{}TwistedOperators.\allowbreak{}ExKeyOf}{23--89}
\leanref{GlobalTheory.\allowbreak{}RowEx.\allowbreak{}exField\_\allowbreak{}of\_\allowbreak{}keyFin\_\allowbreak{}exKey}{Development.\allowbreak{}AbelianVarietyTheory.\allowbreak{}Proofs.\allowbreak{}TwistedOperators.\allowbreak{}ExFieldOfKeyFinExKey}{22--36}

\begin{theorem}[Finite rational test functionals]\label{B8.18}
Assume the inputs of Theorem~\ref{B8.17}. Also assume the unitary Siegel-set theorem and the finite-group cuspidal character input. In the eventual large-field construction take a base and frame satisfying \Lean{KPx} and twisting data satisfying \Lean{Forced}. Take a pinned column witness \(w\). Fix \(Z\in w.Lt\), admissible fields, a set of places and a stage \(z\) with \Lean{GrpTrue Z.\allowbreak{}N f}, \Lean{GrpAt}, \Lean{ExAt} for its realization and \Lean{InGeo} for its geometry.

Let \(EE\) be the row space of \(z.R.toCx3(z.F)\), \(FR_t\) the prescribed space over the coefficient field \(K\), \(j_t\) its complex realization and \(\Theta_t\) the map to \(EE\). For the stage embedding \(v_0:K\hookrightarrow\C\), there are \(N\in\N\) and \(\psi_i\in EE^\vee\) with
\begin{equation}\label{eq:p8c2-row-tests}
 \bigcap_{i\in\operatorname{Fin}(N)}\ker\psi_i=0,\qquad
 \psi_i(\Theta_t(j_t x))\in v_0(K)
 \quad(t\text{ arbitrary},\ i\in\operatorname{Fin}(N),\ x\in FR_t).
\end{equation}
\end{theorem}
\leanref{GlobalTheory.\allowbreak{}T.\allowbreak{}test\_\allowbreak{}vector\_\allowbreak{}assembly\_\allowbreak{}holds}{Development.\allowbreak{}AutomorphicArithmetic.\allowbreak{}ThetaCorrespondence.\allowbreak{}ThetaPairLifts.\allowbreak{}LocalTypeDictionary.\allowbreak{}TestVectorAtDatumOfJoin}{39--40}
\leanref{GlobalTheory.\allowbreak{}RowTestAssembly.\allowbreak{}test\_\allowbreak{}vector\_\allowbreak{}at\_\allowbreak{}datum\_\allowbreak{}of\_\allowbreak{}join}{Development.\allowbreak{}AutomorphicArithmetic.\allowbreak{}ThetaCorrespondence.\allowbreak{}ThetaPairLifts.\allowbreak{}LocalTypeDictionary.\allowbreak{}TestVectorAtDatumOfJoin}{20--31}
\leanref{Formalism.\allowbreak{}IndefR.\allowbreak{}Rat.\allowbreak{}RowTest}{Formalism.\allowbreak{}LocalToGlobal.\allowbreak{}IndefiniteReal.\allowbreak{}RationalAdjoints.\allowbreak{}RatZOfTests}{18--20}

\begin{proof}
At the larger test and field threshold the finite-group input gives a bare realization \(R\) with the same group. The test lattice and coefficient basis give convergent tests
\[
 \ell_i(u)=\mathcal B_{\rho,\xi,\Psi_i,\Phi_i}((R.HE.e(u))(w))
\]
on \(R.HE.H\) at \(w\ne0\). By \Lean{RealOK}, simultaneous vanishing persists for every real pair. For \(u\ne0\), Theorem~\ref{B8.15} supplies a real pair, datum and column with nonzero pairing. Since the chosen columns span at that datum, the common kernel is zero.

Let \(\iota_H:K\hookrightarrow\C\) be the Hilbert-model embedding. Projection of \(x\in FR_t\) to the good summand \(x'\) preserves its theta lift. Theorems~\ref{B8.11} and \ref{B8.16} give the square evaluation
\begin{equation}\label{eq:p8c2-test-formula}
 \iota_H(z_t)c_t\,\ell_i(R.\mathrm{Th}_t(R.jX_t x))=\kappa v_i(t,x'),
 \qquad c_t=\prod_{i',j'}c_{i'j't}\ne0 .
\end{equation}
Here \(z_t\in K\) is a root of unity and \(\kappa\) is common. By the field and local sum formulas, \(v_i(t,x'),c_t/c_{\mu_0}\in\iota_H(K)\). For \(\kappa\ne0\) set \(L_i=(c_{\mu_0}/\kappa)\ell_i\), with theta values \(v_i(t,x')/(\iota_H(z_t)(c_t/c_{\mu_0}))\). For \(\kappa=0\), \(L_i=\ell_i\) has zero theta values. In both cases the common kernel is zero.

The complex row space of \(R\) is \(\mathrm{Tw}_s(R.HE.H)\), with \(v_0=\mathrm{cj}_s\circ\iota_H\), where \(\mathrm{cj}_s\) is identity or conjugation. The linear tests \(\psi_i(\mathrm{Tw}_su)=\mathrm{cj}_s(L_i(u))\) separate and take theta values in \(v_0(K)\). Transfer along the common group gives the tests on \(EE\).
\end{proof}
\leanref{GlobalTheory.\allowbreak{}RowTestFromClaims.\allowbreak{}rowTest\_\allowbreak{}of\_\allowbreak{}choice}{Development.\allowbreak{}AbelianVarietyTheory.\allowbreak{}Proofs.\allowbreak{}TwistedOperators.\allowbreak{}RowTestHilbertModelsOfChoiceClaims}{104--133}
\leanref{GlobalTheory.\allowbreak{}Mem.\allowbreak{}R3.\allowbreak{}Rw.\allowbreak{}RowAt.\allowbreak{}atDat\_\allowbreak{}rowTest}{Development.\allowbreak{}AbelianVarietyTheory.\allowbreak{}Proofs.\allowbreak{}IwasawaDecomposition.\allowbreak{}AtDatRowTest}{44--64}
\leanref{GlobalTheory.\allowbreak{}Mem.\allowbreak{}LiteralAutomorphy.\allowbreak{}rowRatW\_\allowbreak{}of\_\allowbreak{}rowTest}{Development.\allowbreak{}AbelianVarietyTheory.\allowbreak{}Proofs.\allowbreak{}IwasawaDecomposition.\allowbreak{}RowRatWOfRowTest}{9--12}

\section{Preserving a prescribed grid}

\begin{lemma}[Finite field requirements for full grids]\label{B8.20}\label{B8.19}
Put \(\gamma(G)=(v_0,w(v_0),(\chi_{a+1})_{a\in\operatorname{Fin}(m)})\). A grid is full over \(R\) if it admits a complete compatible choice with the evaluation and square-field conditions. The latter asks, at \(\mu_0\), for bottom pairs with norms in \(\iota(K)\) for all \(i,j<m\), with common \((\delta,t_0)\) if bottom witnesses exist. For each \(i,j<m\), \(1\le s'\le\operatorname{halfRank}(D,i,j)\) and either twist, each corresponding untwisted and twisted toric pair \(x,y\) must have \(y/x\in\iota(K)\).

In \(\operatorname{GridAt}(s,C_F,X,f,\Xi)\) restrict the realizations of Theorem~\ref{B8.17} to square-recursive size data, with \Lean{SizeRows}, the host condition, \(F=E^+\) and \(\lvert\operatorname{logDerivativeAtOne}\rvert\le x^\varepsilon\). For every \(G_0\) at \((1,C_F,b(s))\) require a full \(G\) with \(\gamma(G)=\gamma(G_0)\). Requiring this at each completion of \(C,W,V,S,N,G\) gives \(\operatorname{GridA}(C_F)\). The finite-set assertion of Theorem~\ref{B8.17} with conclusion \(\operatorname{GridAt}\) is \(\operatorname{ExGrid}\).

Suppose that at the requisite thresholds every permitted grid admits a Weil system and an initial choice extending to a complete choice over a number field \(N_1\subset\C\) with the same underlying grid choice and Weil system. Suppose the square-field condition holds over a number field \(N_2\subset\C\) for each square-recursive datum, admissible level character system with base index and grid at \((1,C_F,b)\). Assume level agreement transports the grid and Weil system together with the evaluation data and both field conditions. If level keys for finite character sets and grid keys at fixed gap are finite, then \(\operatorname{ExGrid}\) holds.
\end{lemma}
\leanref{GlobalTheory.\allowbreak{}Mem.\allowbreak{}LiteralAutomorphy.\allowbreak{}Full}{Development.\allowbreak{}AbelianVarietyTheory.\allowbreak{}Definitions.\allowbreak{}AutomorphicMembership.\allowbreak{}GKeyFin}{43--47}
\leanref{GlobalTheory.\allowbreak{}Mem.\allowbreak{}LiteralAutomorphy.\allowbreak{}GridAt}{Development.\allowbreak{}AbelianVarietyTheory.\allowbreak{}Definitions.\allowbreak{}AutomorphicMembership.\allowbreak{}GKeyFin}{49--54}
\leanref{GlobalTheory.\allowbreak{}Mem.\allowbreak{}LiteralAutomorphy.\allowbreak{}ExGrid}{Development.\allowbreak{}AbelianVarietyTheory.\allowbreak{}Definitions.\allowbreak{}AutomorphicMembership.\allowbreak{}GKeyFin}{56--59}
\leanref{GlobalTheory.\allowbreak{}Mem.\allowbreak{}LiteralAutomorphy.\allowbreak{}GridA}{Development.\allowbreak{}AbelianVarietyTheory.\allowbreak{}Definitions.\allowbreak{}AutomorphicMembership.\allowbreak{}GKeyFin}{61--62}
\leanref{GlobalTheory.\allowbreak{}Mem.\allowbreak{}LiteralAutomorphy.\allowbreak{}exGrid\_\allowbreak{}of}{Development.\allowbreak{}AbelianVarietyTheory.\allowbreak{}Proofs.\allowbreak{}IwasawaDecomposition.\allowbreak{}GridExistenceOfPieces}{34--143}

\begin{proof}
Use the two fields \(N_1,N_2\) on each representative. Apply the basis-union and transport argument of Theorem~\ref{B8.17} to the realized pairs of level and grid keys at gap \(1\). The transported initial choice completes on a grid with the prescribed key and satisfies both field conditions.
\end{proof}

\begin{theorem}[Eventual existence of evaluation grids]\label{B8.21}
Under the inputs of Theorem~\ref{B8.17} there is \(C_{F1}\in\R\) such that
\[
 \forall C_F\ge C_{F1},\qquad \operatorname{EvLgG}\bigl(\operatorname{GridA}(C_F)\bigr).
\]
\end{theorem}
\leanref{GlobalTheory.\allowbreak{}Mem.\allowbreak{}LiteralAutomorphy.\allowbreak{}evaluation\_\allowbreak{}grid\_\allowbreak{}of\_\allowbreak{}join}{Development.\allowbreak{}AutomorphicArithmetic.\allowbreak{}ThetaCorrespondence.\allowbreak{}ThetaPairLifts.\allowbreak{}LocalTypeDictionary.\allowbreak{}EvaluationGridOfJoin}{23--44}
\leanref{GlobalTheory.\allowbreak{}Mem.\allowbreak{}LiteralAutomorphy.\allowbreak{}evLgG\_\allowbreak{}gridA\_\allowbreak{}of\_\allowbreak{}exGrid}{Development.\allowbreak{}AbelianVarietyTheory.\allowbreak{}Proofs.\allowbreak{}IwasawaDecomposition.\allowbreak{}EventuallyLargeGrid}{14--41}

\begin{proof}
On a representative grid \(G_{00}\), complete the initial choice by convergent square matching and separation with its Weil system. Apply the compositum construction in the proof of Theorem~\ref{B8.17} to this completion of \(G_{00}\) to obtain \(N_1\). For \(N_2\), bottom integrality and smallness give algebraic norms with common \((\delta,t_0)\) when witnesses exist. By \Lean{fieldDat} and toric value existence there is at most one ratio \(y/x\) for each \(i,j,s'\) and twist. Use the same \Lean{IsIntLattice} period lattice \(\Lambda=\mathcal O_{\Q(i\sqrt d)}\). The normalization gives \((y/x)^{12}=\beta(y)/\beta(x)\) with \(\beta(x)\ne0\) and algebraic right side. Adjoin the ratios and norms to obtain \(N_2\). Square transport preserves its field condition.

There are finitely many pairs \(w(v_0)\mid v_0\mid q_0\). At fixed preliminary data and gap, infinity types are fixed and conductor exponents are at most two on a fixed finite set and zero elsewhere. Ray class finiteness gives finite grid keys. Lemma~\ref{B8.20} gives \(\operatorname{ExGrid}\) at threshold \(C_{F1}\).

For \(C_F\ge C_{F1}\) take \(\varepsilon_1=1/100\). Use zero auxiliary thresholds and the scale threshold of \(\operatorname{ExGrid}\). Take empty first three field requests and fourth request \(S(P,f,\Xi)\) for finite \(\Xi\), empty otherwise. The twisting set at each completed stage is finite, so its fourth request gives \(\operatorname{GridAt}\).
\end{proof}

\section{The functional bound}

Let \(RR\) be the glued row lattice of positive rank with infinite frames \(X_{E,v}\) and normalizations \(\kappa_v>0\). In squared norms the condition is
\begin{equation}\label{eq:p8c2-fun}
 \operatorname{FUN}\qquad
 \mu_{\max,RR}\le N.hh-\operatorname{meanLog}_v\kappa_v,\qquad
 N(C,C_H).hh=x^{C_H}b ,
\end{equation}
where \(C_H\) is the family height exponent.

\begin{theorem}[The constant-function Fourier construction satisfies the functional bound]\label{B8.22}
Assume \(\mathscr P,\mathscr J\), the Hecke--Tate functional equation and Theorems~\ref{B8.11} and \ref{B8.15}. Assume the closed eventual condition \Lean{L6}, \(N.l6_j\le\log\lambda^{\rm true}_{j,\mu_0,p,v}\) for every \(C\), layer \(j\), block \(p\) and infinite place \(v\), with \(N=N(C,C_H)\).

There is \(C_{F0}'\in\R\) such that for each \(C_F\ge C_{F0}'\) there is \(\varepsilon_1>0\) and for each \(0<\varepsilon\le\varepsilon_1\) there are size thresholds \(K\). If \(F_m,C_H\) satisfy \(\operatorname{FamHyp}(F_m)\), \(F_m.\operatorname{Large}(K,C_H)\) and \(\operatorname{DIVinN}(F_m)\), there is \(x_0\) after which every admissible \(p\) with \(\operatorname{HasGrid}(p,C_F)\) satisfies
\[
 \forall C\in\R,\qquad (p.S_{\rm false}).\operatorname{FUN}(p.N(C,C_H)).
\]
Here \(\operatorname{HasGrid}\) requires a fitted grid with Weil system and a complete choice satisfying \Lean{ESerTie} and the two field conditions.
\end{theorem}
\leanref{GlobalTheory.\allowbreak{}BigTop.\allowbreak{}Fun.\allowbreak{}Close.\allowbreak{}fourier\_\allowbreak{}coefficient\_\allowbreak{}fit\_\allowbreak{}of\_\allowbreak{}join\_\allowbreak{}alt}{Development.\allowbreak{}AutomorphicArithmetic.\allowbreak{}ThetaCorrespondence.\allowbreak{}ThetaPairLifts.\allowbreak{}ThirdPlaceSeries.\allowbreak{}FourierCoefficientFitOfJoin}{12--17}
\leanref{GlobalTheory.\allowbreak{}BigTop.\allowbreak{}Fun.\allowbreak{}ConstantFunctionFits}{Development.\allowbreak{}AutomorphicArithmetic.\allowbreak{}ThetaCorrespondence.\allowbreak{}ThetaPairLifts.\allowbreak{}ThirdPlaceSeries.\allowbreak{}LeafInfinitePlaceBound}{52--53}
\leanref{Formalism.\allowbreak{}L2.\allowbreak{}Indef.\allowbreak{}IndefData.\allowbreak{}FUN}{Formalism.\allowbreak{}LocalToGlobal.\allowbreak{}IndefiniteForms.\allowbreak{}IndefSide}{331--332}

\begin{proof}
Let \(\iota_t\) be the row realization and \(P_t\) its rational-row quotient. Local sums give a Hermitian lattice \(M\) and evaluation matrices \(V_t\). They also give a map \(\Psi\) from the glued row space. The square product \(r_t=\prod_{i,j}r_{ij,t}\) is one at \(\mu_0\) and nonzero at the selected finite indices \(\mu(w)\). Absorb the root of unity into \(V_t\), a finite-place unit equal to one at the base. Then
\begin{equation}\label{eq:p8c2-psi-comparison}
 \Psi\iota_tP_t=r_t^{-1}V_t\qquad(t=\mu_0\text{ or }t=\mu(w)).
\end{equation}
Separation gives \(\ker V_{\mu_0}=\ker P_{\mu_0}\), so \(\Psi\) is injective by surjectivity of \(\iota_{\mu_0}P_{\mu_0}\). If \(d_F(w)\in\Z\) is the evaluation exponent supported on \(T_F\), put \(a_w=v_w(r_{\mu(w)})+d_F(w)\). Equation~\eqref{eq:p8c2-psi-comparison} gives the local inclusion in \(\mathfrak p_w^{-a_w}M_w\), with \(a_w=0\) off \(T_{\rm twist}\cup T_F\). Hence
\begin{equation}\label{eq:p8c2-denominator-inclusion}
 \Psi(RR)\subseteq\left(\prod_w\mathfrak p_w^{-a_w}\right)M.
\end{equation}

At infinity use \(\sqrt{\kappa_v}X_{E,v}\). Let \(N_*=k+m(t+k+m+6)\) with auxiliary rank \(t\), and \(D'_0=3C_1fN_*^2\log b\) with estimate constant \(C_1\ge1\). The base \(\lambda\)'s agree by \Lean{LevelLink}. The term \(N.l6_j\) is a nonnegative sum of hook-length logarithms minus \(C_{\rm sz}fn_V^2L_p\). Since \(L_p\le3c\log b,\ n_V\le N_*,\ C_1\ge c\max(0,C_{\rm sz})\), with family constants \(c,C_{\rm sz}\), we have \(N.l6_j\ge-D'_0\). On each nonzero base block,
\[
 \sum_{j\in\mathcal J}\log\lambda^{\rm false}_{j,\mu_0,p,v}
 \ge-|\mathcal J|D'_0\ge-mD'_0 .
\]
The quotient-norm comparison multiplies the bound by \(\exp(mD'_0/2)\). In squared norms the height therefore satisfies \(h(\Psi)\le H_I=H_{I,0}+mD'_0\), where \(H_{I,0}\) is the direct evaluation bound.

Put \(\ell_w=\log N\mathfrak p_w/[K:\Q]\) for the coefficient field \(K\). The estimates give \(2\sum_w\ell_w d_F(w)\le B_F\) and \(\mu_{\max}(M)\le E_M\). Let \(A_0,D_0\) bound bottom sizes and denominators and let \(B_0\) bound higher heights and prime counts. Write \(\Sigma_0=\sum_q e_q\log q\) for the modulus contribution. The product formula gives a signed bottom valuation sum at most \(A_0\) with negative part at most \(3D_0\). Its absolute sum is at most \(A_0+6D_0\). By \Lean{LeafTHB}, \(\sum_w\ell_w|v_w(z)+g_T(w)|\le B_0\) for \(z=z_T,z_{T,d}\). Set \(B=\max(A_0+6D_0,B_0)\). Each pair \(i,j\) contributes two bottom and at most \(8N_*\) shifted higher valuations per twist class, counted with multiplicity. The prime count has coefficient five on at most \(N_*\) indices per class. The weighted reciprocal-square sum is at most one after grouping by rational primes, so its degree bound costs \(2N_*\) per pair. The allowance \(d_\ell(w)\ge0\) and the modulus cost at most \(6D_0\) and \(4m\Sigma_0\). Thus
\begin{equation}\label{eq:p8c2-grid-budget}
 \sum_w\ell_w v_w(r_{\mu(w)})\le
 G_*:=m^2\bigl(2(2+8N_*)B+10N_*B_0+2N_*+6D_0+4m\Sigma_0\bigr).
\end{equation}
The change of frames adds \(\operatorname{meanLog}\kappa\) to the source slope, and the ideal in \eqref{eq:p8c2-denominator-inclusion} adds \(2\sum_w\ell_w a_w\) to the target slope. Injectivity gives
\begin{equation}\label{eq:p8c2-functional-budget}
 \mu_{\max,RR}+\operatorname{meanLog}\kappa\le E_M+H_I+2G_*+B_F.
\end{equation}

Let \(R_*=k+m-1\) and \(P_{\rm pf}\) be the prime-power modulus budget. Use \(1\le f\le k,\ 1\le m\le N_*,\ R_*\le N_*,\ mfN_*^2\le b,\ L\ge1,\ L+\log x\le3L\), and \(\Sigma_0\le2P_{\rm pf},\ c_S\le45fL,\ 1\le B_w\le C_1b,\ \log B_w\le C_1L\). Put \(l_1=\log(2fN_*),\ l_2=\log(2N_*B_w),\ l_3=\log(2fN_*B_w)\), so \(l_1\le1+L,\ l_2,l_3\le C_1L+1+L\). With \(U=C_1^3fm^2R_*N_*^3bL\) and \(V=C_1^3fm^2N_*^2P_{\rm pf}b\), substitution gives
\[
\begin{array}{l@{\quad}l}
\text{Quantity or bounding expression}&\text{Upper bound}\\
2C_1fN_*(l_1+B_wL+2N_*B_wc_S)\quad(D_0)&186C_1^2f^2N_*^2bL\\
C_1fN_*[2B_w(L+l_2)+N_*(L+l_3)]+C_1N_*^2c_S\quad(A_0)&57C_1^3fN_*^2bL\\
B_0=C_1fb(L+\log x)&3C_1fbL\\
E_M=4fn_pR_p\log(n_pR_p)&32U\\
H_I=C_1fmN_*bL+C_1m^2N_*^2\Sigma_0+mD'_0&4U+2V\\
B_F=C_1fbm^2(N_*+N_*^2)(L+\Sigma_0)&2U+4V\\
G_*&(23520+30+2+1116)U+8V .
\end{array}
\]
Use \(n_p\le2mR_*,\ R_p\le C_1mN_*^3b,\ \log(n_pR_p)\le4L\). The last row uses \(B\le1176C_1^3f^2N_*^2bL,\ 2(2+8N_*)\le20N_*\) and \(f\le R_*\). The polynomial-size condition gives
\begin{equation}\label{eq:p8c2-functional-polynomial}
 E_M+H_I+2G_*+B_F\le100000C_1^3fm^2N_*^2(R_*N_*L+P_{\rm pf})b
 \le x^{C_H}b .
\end{equation}
Take \(C_1\) above the joint family estimate threshold and impose the point conditions there. The fitted choice supplies the square terms and estimates on the false realization. Substituting in \eqref{eq:p8c2-functional-budget} gives \(\operatorname{FUN}\) for every \(C\).
\end{proof}
\leanref{GlobalTheory.\allowbreak{}BigTop.\allowbreak{}Fun.\allowbreak{}Close.\allowbreak{}fittedFunction\_\allowbreak{}at\_\allowbreak{}zero\_\allowbreak{}close\_\allowbreak{}alt}{Development.\allowbreak{}AutomorphicArithmetic.\allowbreak{}ThetaCorrespondence.\allowbreak{}ThetaPairLifts.\allowbreak{}ThirdPlaceSeries.\allowbreak{}FittedFunctionAtZeroClose}{42--67}
\leanref{GlobalTheory.\allowbreak{}BigTop.\allowbreak{}Fun.\allowbreak{}fun0Fit\_\allowbreak{}of\_\allowbreak{}leaves\_\allowbreak{}alt}{Development.\allowbreak{}AutomorphicArithmetic.\allowbreak{}ThetaCorrespondence.\allowbreak{}ThetaPairLifts.\allowbreak{}ThirdPlaceSeries.\allowbreak{}LeafInfinitePlaceBound}{83--120}
\leanref{Formalism.\allowbreak{}L2.\allowbreak{}Indef.\allowbreak{}Tree.\allowbreak{}functional\_\allowbreak{}lower\_\allowbreak{}bound\_\allowbreak{}of\_\allowbreak{}conditions}{Formalism.\allowbreak{}LocalToGlobal.\allowbreak{}IndefiniteForms.\allowbreak{}IndefiniteTree.\allowbreak{}FunctionalLowerBoundOfConditionsNumeric}{940--1033}
\leanref{Formalism.\allowbreak{}L2.\allowbreak{}Indef.\allowbreak{}Tree.\allowbreak{}size\_\allowbreak{}arith}{Formalism.\allowbreak{}LocalToGlobal.\allowbreak{}IndefiniteForms.\allowbreak{}IndefiniteTree.\allowbreak{}FunctionalLowerBoundOfConditionsNumeric}{243--397}

\section{The functional row on the chosen members}

\begin{theorem}[The functional condition along members]\label{B8.23}
Assume discrete-series coverage, Harish--Chandra realization, rapid decay of cusp forms, Casselman--Wallach globalization, the ray-class-field input and Theorem~\ref{B8.22}. There is \(C_{F0}\in\R\). For \(C_F\ge C_{F0}\), take an allowed local Langlands system \(h_L\), norm-choice extension \(N\), auxiliary extension \(S\) and geometry condition \(\Gamma\). Each value of \(N\) admits the extension with list at \(C_F\). Strengthen \(N\) by \(\operatorname{GridA}(C_F)\) and assume the open conditions of Theorem~\ref{B8.08}.

A condition \(R\) eventually holds for every constructed member over the permitted base frames. On good members satisfying \(R\), a row condition \(E\) eventually holds for every linked bundle. For a compatible selection on \(\mathcal A\) for this row, assume the closed hypotheses on \(W_3(\mathcal A)\) and \(\operatorname{DIVinN}\). There is \(C_0\ge0\) such that for each \(C\ge C_0\) the selected bundles eventually satisfy \(\operatorname{FUN}\) with \(N(C,C_H)\), with point threshold allowed to depend on \(C\). The class \(W_3(\mathcal A)\) requires \(W_1(\mathcal A)\) and a level-\(3\) height choice. It also requires eventual bounds on \(T_u,D_{\Gamma,u}\). For the latter some \(C',x_0\) give \(\log D_{\Gamma,u}\le C'fN(m)\log(fN(m))\) on every member at \(x\ge x_0\) and every word \(u\).
\end{theorem}
\leanref{GlobalTheory.\allowbreak{}Fc.\allowbreak{}StageOne.\allowbreak{}rowFun\_\allowbreak{}at\_\allowbreak{}gridA\_\allowbreak{}alt}{Development.\allowbreak{}AutomorphicArithmetic.\allowbreak{}ThetaCorrespondence.\allowbreak{}ThetaPairLifts.\allowbreak{}ThirdPlaceSeries.\allowbreak{}RowFunAtGridA}{32--51}
\leanref{GlobalTheory.\allowbreak{}Fc.\allowbreak{}S0.\allowbreak{}FunctionalCondition}{Development.\allowbreak{}AutomorphicArithmetic.\allowbreak{}ThetaCorrespondence.\allowbreak{}ThetaPairLifts.\allowbreak{}ThirdPlaceSeries.\allowbreak{}GeometricWeakTies}{34--35}

\begin{proof}
Use the threshold of Theorem~\ref{B8.22}. Ray-class-field existence and real-place compatibility give the geometry row from the divisor row. Intersect with the goodness row and with the primitive and divisor rows for \(R,E\). Compatible geometries give fitted grids on the summands, and discrete series and globalization give the eventual point-summand conditions from the real realization, quotient and division data. The selection hypotheses give \Lean{Duty7}. Apply Theorem~\ref{B8.22} on the fitted points under \(\operatorname{DIVinN}\). Take \(C_0=0\) in the guarded-point transfer and use compatibility to pass to bundles.
\end{proof}
\leanref{GlobalTheory.\allowbreak{}Fc.\allowbreak{}Ind.\allowbreak{}evWC\_\allowbreak{}fun\_\allowbreak{}of\_\allowbreak{}fun0Fit\_\allowbreak{}alt}{Development.\allowbreak{}AutomorphicArithmetic.\allowbreak{}ThetaCorrespondence.\allowbreak{}ThetaPairLifts.\allowbreak{}ThirdPlaceSeries.\allowbreak{}EvWCFunOfFun0Fit}{14--21}

Use Theorem~\ref{B8.18} in the member supply of Theorem~\ref{B8.07}. Apply Theorem~\ref{B8.03} at \(c_{\rm Dep}=C_F=\max(C_{F0},C_{F1})\) with the grid and functional rows, the linked \(2,3\)-bound of Theorem~\ref{B8.01} and the covolume row of Theorem~\ref{B8.34}.
\leanref{UpperTheory.\allowbreak{}Lit.\allowbreak{}nonempty\_\allowbreak{}hyps\_\allowbreak{}definite\_\allowbreak{}side\_\allowbreak{}of\_\allowbreak{}closed\_\allowbreak{}bound}{Development.\allowbreak{}AutomorphicArithmetic.\allowbreak{}ThetaCorrespondence.\allowbreak{}ThetaPairLifts.\allowbreak{}ThirdPlaceSeries.\allowbreak{}SiegelEisensteinSecondTermHypotheses}{44--68}


% ---- p8c3 ----
\chapter{Covolume bounds}\label{ch:B8-covolume}

A nonzero map of operator norm at most \(e^a\) at every infinite place has height at most \(2a\) in the squared-norm convention. For a rank-\(r\) lattice \(\Lambda\) over a number field \(K\) with \(s_K\) infinite places, let \(\operatorname{ht}(X,\Lambda)\) be its \(\Z\)-covolume in the frame \(X\) times \(2^{s_Kr}|\operatorname{disc}K|^{-r/2}\). Set
\[
 \operatorname{lcv}(X,\Lambda)=\frac{2}{[K:\Q]}\log\operatorname{ht}(X,\Lambda).
\]
\leanref{Formalism.\allowbreak{}HN.\allowbreak{}ht}{Formalism.\allowbreak{}HarderNarasimhan.\allowbreak{}StatementPl}{29--31}
\leanref{Formalism.\allowbreak{}HN.\allowbreak{}htMap}{Formalism.\allowbreak{}HarderNarasimhan.\allowbreak{}StatementPl}{60--62}
\leanref{GlobalTheory.\allowbreak{}DescentFamily.\allowbreak{}htMap\_\allowbreak{}le\_\allowbreak{}of\_\allowbreak{}points}{Development.\allowbreak{}AbelianVarietyTheory.\allowbreak{}Proofs.\allowbreak{}DivisorProofs.\allowbreak{}HtMapLeOfPoints}{34--38}
\leanref{Formalism.\allowbreak{}DivisionValues.\allowbreak{}lcv}{Formalism.\allowbreak{}QuotientMetricDivisors.\allowbreak{}Statement}{34--37}

At a point \(p\) and layer \(0\le i<m\), put \(E_p=\operatorname{unitE}(p)\) and \(j=i+1\). For the layer weights \(w_{j,\rho}\), set
\[
 M_{p,i}=(k+j-1-(m-j))\log\Omega\sum_\rho w_{j,\rho},
 \qquad \ell_{3,i}(C,CH)=M_{p,i}+CE_p.
 \tag{8.24a}
\]
The recorded root discriminant and logarithmic derivative give the factor of \ref{B5.07},
\[
 \log\Omega=\gamma+\log(2\pi/\operatorname{rd}(E))-L_E,\qquad
 L_E=\Re\frac{L'(1,\chi_d)}{L(1,\chi_d)},
\]
where \(E\) is the host field and \(\chi_d\) is the Kronecker character of \(\Q(\sqrt{-d})\). At fixed \(p,i,C\) the bound is independent of \(CH\).
\leanref{Formalism.\allowbreak{}L2.\allowbreak{}DIV.\allowbreak{}ZscOfNum}{Formalism.\allowbreak{}LocalToGlobal.\allowbreak{}Divisibility.\allowbreak{}ZscOfNum}{59--72}
\leanref{Formalism.\allowbreak{}L2.\allowbreak{}DIV.\allowbreak{}ofRecord\_\allowbreak{}logOm}{Formalism.\allowbreak{}LocalToGlobal.\allowbreak{}Divisibility.\allowbreak{}NumbersLemmas}{786--789}
\leanref{GlobalTheory.\allowbreak{}CovTop.\allowbreak{}unitE}{Development.\allowbreak{}AbelianVarietyTheory.\allowbreak{}Definitions.\allowbreak{}Divisors.\allowbreak{}CovolXOfCovol}{22--23}
\leanref{GlobalTheory.\allowbreak{}CovTop.\allowbreak{}mainT}{Development.\allowbreak{}AbelianVarietyTheory.\allowbreak{}Definitions.\allowbreak{}Divisors.\allowbreak{}CovolXOfCovol}{24--25}
\leanref{Formalism.\allowbreak{}L2.\allowbreak{}Indef.\allowbreak{}IndefScale.\allowbreak{}num}{Formalism.\allowbreak{}LocalToGlobal.\allowbreak{}IndefiniteForms.\allowbreak{}IndefSide}{224--232}

\section{The estimate and its geometric inputs}

Let \(\mathcal U\) be the class of all families. A family in \(W_3(\mathcal U)\) belongs to \(\mathcal F_{\rm cov}=\Lean{FamHyp}\) precisely when it satisfies the fixed-page and window conditions as well as the chosen-wall condition, with \(q_0<X'\) above its threshold.
\leanref{GlobalTheory.\allowbreak{}BigTop.\allowbreak{}FamHyp}{Development.\allowbreak{}AbelianVarietyTheory.\allowbreak{}Definitions.\allowbreak{}Divisors.\allowbreak{}TupleReadBaseTupleEqFamily}{503--507}

At \(p\), with datum \(D\) and size data \(L\), admissibility means
\[
 \Lean{RowsTAlt}(p)\quad\Longleftrightarrow\quad
 p.\Lean{PointConditions'}\land
 \bigl(\Lean{FieldSmall}(D,L)\land\Lean{ESerTie}(D)\bigr).
\]
\leanref{GlobalTheory.\allowbreak{}CovTop.\allowbreak{}RowsTAlt}{Development.\allowbreak{}AutomorphicArithmetic.\allowbreak{}ThetaCorrespondence.\allowbreak{}ThetaPairLifts.\allowbreak{}ThirdPlaceSeries.\allowbreak{}LiftLambdaBound}{17--17}
\leanref{GlobalTheory.\allowbreak{}CovTop.\allowbreak{}RecTied}{Development.\allowbreak{}AbelianVarietyTheory.\allowbreak{}Definitions.\allowbreak{}Divisors.\allowbreak{}CovolXOfCovol}{161--161}

In the closed range of \ref{B8.39} on \(\mathcal F_{\rm cov}\), fix \(F_m,CH\). The assertion \(\Lean{CovolumeBoundOfX}\) requires \(C_0\ge0\) and \(x_0\in\R\) such that
\begin{equation}
 \frac{\operatorname{lcv}(X_{W,i},\Lambda_i)}{D_{\rm norm}}\le M_{p,i}+CE_p
 \tag{8.24b}\label{B8.24}
\end{equation}
for \(x>F_m.x_1\), \(x\ge x_0\), \(\iota\in F_m.Ix(x)\), admissible \(p\) over \((x,\iota)\), \(C\ge C_0\), and \(0\le i<m\). In \(S=p.S(\mathrm{false})\), \(D_{\rm norm}=S.D\) is the dimension of \(X_E\), with the quotient set to zero when \(S.D=0\). For fixed \(S,N\), \(L3\) uses \(N.l3_i\) as the right-hand side at layer \(i\).
\leanref{GlobalTheory.\allowbreak{}CovTop.\allowbreak{}CovolumeBoundOfX}{Development.\allowbreak{}AutomorphicArithmetic.\allowbreak{}ThetaCorrespondence.\allowbreak{}ThetaPairLifts.\allowbreak{}ThirdPlaceSeries.\allowbreak{}LiftLambdaBound}{60--61}
\leanref{GlobalTheory.\allowbreak{}CovTop.\allowbreak{}EventuallyAtAllPoints}{Development.\allowbreak{}AutomorphicArithmetic.\allowbreak{}ThetaCorrespondence.\allowbreak{}ThetaPairLifts.\allowbreak{}ThirdPlaceSeries.\allowbreak{}LiftLambdaBound}{18--19}
\leanref{Formalism.\allowbreak{}L2.\allowbreak{}Indef.\allowbreak{}IndefData.\allowbreak{}L3}{Formalism.\allowbreak{}LocalToGlobal.\allowbreak{}IndefiniteForms.\allowbreak{}IndefSide}{315--317}
\leanref{Formalism.\allowbreak{}MainNotions.\allowbreak{}FamP.\allowbreak{}Large}{Formalism.\allowbreak{}SharpRestriction.\allowbreak{}ParamsHeightRankConstant}{138--140}

\begin{lemma}[The flat line in the row metric]\label{B8.24a}
For row data satisfying the flat-degree law on a nonempty component system, let \((P,Q)\) be the signature and \(\tau\) a CM embedding. In the host normalisation the flat first Chern value is
\[
 c_{1,\tau}=-\tfrac{P-Q}{2}\log\Omega-Q(\log d-\log\operatorname{rd}(E)).
\]
\end{lemma}
\begin{proof}
At a good-reduction CM point the archimedean frame has degree \((P-Q)\Re\log\Omega(E,\chi_d)/2\), the difference of the conjugate Lie eigenlattice degrees. Subtract the local different defects from the frame valuations to obtain a flat section. Its degree has the additional term \(Q(\log d-\log\operatorname{rd}(E))\), and the first Chern value is minus this degree. For the matching CM host, the period of \ref{B5.07x} identifies \(\Re\log\Omega(E,\chi_d)\) with the recorded \(\log\Omega\), since \(2\pi/\operatorname{rd}(E)=\pi/v_E\). The finite term vanishes when \(Q=0\).

On a constructed cover \(|\log d-\log\operatorname{rd}(E)|\le3L_x\) with \(L_x=L.\mathrm{sc}.L\). Let \(a_\tau\) be the zero-based \(k-2\) entry of the unshifted layer weight and \(\rho(\tau)\) its matching row index. Put \(\beta=P-Q\). Rounding gives
\[
 |a_\tau-w_{j,\rho(\tau)}|\le L.\mathrm{Sc}+n_V+m+t+2,
 \qquad n_V=P+Q+t,\quad t=D.R.t.
\]
Use \((P+Q)L.\mathrm{Sc}\le16b\) and \((P+Q)n_V\le b\) to get \((P+Q)|a_\tau-w_{j,\rho(\tau)}|\le20b\). Since \(|a_\tau|\le4b\), \(|\log\Omega|\le3L_x\) and \(Q\le m\), the errors below are at most \(60fbL_x\) and \(4fmbL_x\), with \(f=D.R.f\). \(E_p=fmbL_x\log\log x\) and \(\log\log x\ge1\) give
\[
 |M_{p,i}-\beta\log\Omega\textstyle\sum_\tau a_\tau|\le100E_p,\qquad
 QL_x\sum_\tau|a_\tau|\le100E_p .
\]
In the squared norm the first Chern contribution has a factor of two. Thus
\[
 \left|2\sum_\tau a_\tau c_{1,\tau}+\beta\log\Omega\sum_\tau a_\tau\right|
 \le6QL_x\sum_\tau|a_\tau|,\qquad
 2\sum_\tau a_\tau c_{1,\tau}\ge-M_{p,i}-700E_p .
\]
The different term in the twist costs at most \(4QL_x\sum_\tau|a_\tau|\). The determinant and bad-prime terms cost at most \(3E_p\) and \(80E_p\), respectively. The flatness loss in (8.29k) is therefore \((11\cdot100+83)E_p\). The slope main term is \(-M_{p,i}\), which gives \(+M_{p,i}\) in the covolume bound.
\end{proof}
\leanref{GlobalTheory.\allowbreak{}FlatLine.\allowbreak{}FlatDegree.\allowbreak{}c1\_\allowbreak{}eq}{Development.\allowbreak{}AbelianVarietyTheory.\allowbreak{}Definitions.\allowbreak{}Foundations.\allowbreak{}ExistsFlatSecRankZero}{788--790}
\leanref{GlobalTheory.\allowbreak{}FlatCore.\allowbreak{}flatDegreeSt\_\allowbreak{}of\_\allowbreak{}prismaticDieudonne}{Development.\allowbreak{}AbelianVarietyTheory.\allowbreak{}Proofs.\allowbreak{}DivisorProofs.\allowbreak{}FlatDegreeFromPrismatic}{23--35}
\leanref{GlobalTheory.\allowbreak{}FlatSplit.\allowbreak{}twice\_\allowbreak{}twist\_\allowbreak{}ge}{Development.\allowbreak{}AbelianVarietyTheory.\allowbreak{}Proofs.\allowbreak{}DivisorProofs.\allowbreak{}FlatBoundOfAnchor}{76--102}
\leanref{GlobalTheory.\allowbreak{}CovTop.\allowbreak{}Period.\allowbreak{}period\_\allowbreak{}eq\_\allowbreak{}kronecker}{Development.\allowbreak{}AbelianVarietyTheory.\allowbreak{}Proofs.\allowbreak{}DivisorProofs.\allowbreak{}LogPeriodEqKroneckerCharacterLogPeriodOfHost}{13--17}
\leanref{GlobalTheory.\allowbreak{}CovTop.\allowbreak{}AnchorRd.\allowbreak{}anchorAt\_\allowbreak{}of}{Development.\allowbreak{}AbelianVarietyTheory.\allowbreak{}Proofs.\allowbreak{}DivisorProofs.\allowbreak{}AnchorBoundStatement}{77--209}
\leanref{GlobalTheory.\allowbreak{}FlatSplit.\allowbreak{}flatAt\_\allowbreak{}of}{Development.\allowbreak{}AbelianVarietyTheory.\allowbreak{}Proofs.\allowbreak{}DivisorProofs.\allowbreak{}FlatBoundOfAnchor}{125--187}

For each eligible \(F_m,CH\), one choice of \(C\ge0,x_0\) works at every point and layer for every good cover \(Y\in\operatorname{Cov}(p,i)\), \(i+1<m\). The condition \(Y.\Lean{Good}(\lfloor3L_x/\log3\rfloor_{\N},391(frL_x)^2)\) includes the unit-minor and model conditions, where \(f=D.R.f\) and \(r=D.R.r\). Let \(P_*\) denote the family parameter \(PP\) at \(p\).
\leanref{GlobalTheory.\allowbreak{}CovTop.\allowbreak{}LiftsCondition}{Development.\allowbreak{}AutomorphicArithmetic.\allowbreak{}ThetaCorrespondence.\allowbreak{}ThetaPairLifts.\allowbreak{}ThirdPlaceSeries.\allowbreak{}LiftLambdaBound}{26--28}
\leanref{GlobalTheory.\allowbreak{}CovTop.\allowbreak{}CoversCondition}{Development.\allowbreak{}AutomorphicArithmetic.\allowbreak{}ThetaCorrespondence.\allowbreak{}ThetaPairLifts.\allowbreak{}ThirdPlaceSeries.\allowbreak{}LiftLambdaBound}{29--32}
\leanref{Formalism.\allowbreak{}MainNotions.\allowbreak{}Closed}{Formalism.\allowbreak{}SharpRestriction.\allowbreak{}ParamsHeightRankConstant}{159--162}
\leanref{GlobalTheory.\allowbreak{}CovTop.\allowbreak{}EventuallyAtAllPoints.\allowbreak{}and}{Development.\allowbreak{}AutomorphicArithmetic.\allowbreak{}ThetaCorrespondence.\allowbreak{}ThetaPairLifts.\allowbreak{}ThirdPlaceSeries.\allowbreak{}LiftLambdaBound}{20--25}

\begin{theorem}[Geometric inputs for the covolume estimate]\label{B8.25}
Assume characteristic-zero resolution and principalisation, invariant Poincar\'e--Lelong, \(HSUse\), the inequality of \ref{B8.27}, Fargues exactness, Zarhin's statement, and nef Riemann--Roch. The input \(HSUse\) combines arithmetic Kleiman with the Hilbert--Samuel upper bound \(G7\) for data generated by small sections and the nef arithmetic-volume bound.

Assume also integral Siegel and prescribed-fan toroidal compactification, semiabelian Lie-bundle seminegativity, point separation, the arithmetic degree formula, \(\Lean{PublishedHodgeBundleTheorems}\), and Grothendieck deformation theory. Then \(\Lean{CovolumeBoundOfX}\) holds.
\end{theorem}
\leanref{SelfIntersectionVolumeEstimates.\allowbreak{}topCovolSlim2\_\allowbreak{}use\_\allowbreak{}alt\_\allowbreak{}alt}{Development.\allowbreak{}AutomorphicArithmetic.\allowbreak{}ThetaCorrespondence.\allowbreak{}ThetaPairLifts.\allowbreak{}ThirdPlaceSeries.\allowbreak{}TopCovolSlim2Use}{16--22}
\leanref{PositivityTheory.\allowbreak{}MbHsUse.\allowbreak{}HSUse}{Development.\allowbreak{}WeilBound.\allowbreak{}CoreStatements.\allowbreak{}AutomorphicKernels.\allowbreak{}MoriwakiBigness.\allowbreak{}HilbertSamuelApplications.\allowbreak{}SmallGeneratedCountUpperOfUse}{12--13}

\begin{proof}
Let \(\mathcal H_P\) be the Hodge-bundle package. Fan compactification, Lie-bundle seminegativity, point separation, the Fubini--Study Chow-height formula and Green's parametrisation (\ref{B7.31}) give \(\mathcal H_X\). Combine integral compactification with the fan charts and the degree formula to obtain \(\mathcal H_H\). Chow-form comparison bounds self-intersections and toroidal nonvanishing gives the small-section injection. Section comparison uses Kawamata--Viehweg vanishing, log-canonical \(L^2\)-extension, GAGA and nef numerics, with integral Chow forms and intersection lengths. It also uses hyperplane compatibility.

The Hodge-bundle results give Gram and measure comparisons and deep-cover existence. Use \(\Lean{cov\_\allowbreak{}size\_\allowbreak{}bound}\) for (8.28b), (8.29l) and (8.32b), and \ref{B8.24a} for the flat twist. Deformation theory and flat degree give the cover degrees. Comparison with component ranks gives the shares. Apply \ref{B8.26} and \ref{B8.28} for heights and \ref{B8.29} for slopes, with the terminal defect from the Gram and anchor estimates. The Gram and measure comparisons give (8.32a) and the maximum-slope bound, and rationality gives the first inequality of (8.32c). Use \ref{B8.30} and \ref{B8.31} for truncation, then apply \ref{B8.33}.
\end{proof}
\leanref{SelfIntersectionVolumeEstimates.\allowbreak{}topCovolG\_\allowbreak{}alt\_\allowbreak{}alt}{Development.\allowbreak{}AutomorphicArithmetic.\allowbreak{}ThetaCorrespondence.\allowbreak{}ThetaPairLifts.\allowbreak{}ThirdPlaceSeries.\allowbreak{}TopCovolG}{39--58}

\section{Arithmetic volume and component heights}

Let \(v\) have signature \((p_{\rm sig},q_{\rm sig})\) and embedding set \(\Psi\). Put \(n=p_{\rm sig}+q_{\rm sig}\) and set
\[
 \|v\|_1=\sum_{\tau,a}|v_{\tau,a}|,\qquad
 \operatorname{spr}(v)=\sum_\tau
 \bigl(\max_a v_{\tau,a}-\min_a v_{\tau,a}\bigr).
\]
A gap \(g>0\) requires every \(v_\tau\) to be strictly decreasing and \(g\le |v_{\tau,a}-v_{\tau,b}|\) for \(p_{\rm sig}\le a<n\), \(0\le b<n\), \(a\ne b\). Write \(\|v\|_\infty=\max_{\tau,a}|v_{\tau,a}|\) and set
\[
 \begin{split}
 B(v,g)&=4|\Psi|n^2(1+\|v\|_\infty/g),\\
 \mathcal E_s(v,g)&=(\log B(v,g)+s)\|v\|_1+q_{\rm sig}\operatorname{spr}(v),\\
 \ell(D)&=\log\operatorname{rd}(E)+|L_E|.
 \end{split}
 \tag{8.26a}
\]
Let \(\operatorname{vol}_Z(t)\) be the mixed volume on a component \(Z\) of \(Y\) in degrees divisible by \(8\), with height valuations above \(S_B\) below and record valuations elsewhere. The archimedean scale is \(\xi_\R(D)=|\phi_0(\xi_0)|\) at the distinguished host embedding. For \(\deg_Z(v)>0\), define the extended-real slope by
\[
 \nu_Z(v;B)=\inf_{t\in\R}\left(\operatorname{vol}_Z(t)/\deg_Z(v)-t\right).
\]
\leanref{GlobalTheory.\allowbreak{}CovHt.\allowbreak{}WtGap}{Development.\allowbreak{}AbelianVarietyTheory.\allowbreak{}Definitions.\allowbreak{}Divisors.\allowbreak{}SemiposAtOfPositive}{88--89}
\leanref{GlobalTheory.\allowbreak{}CovHt.\allowbreak{}bndB}{Development.\allowbreak{}AbelianVarietyTheory.\allowbreak{}Definitions.\allowbreak{}Divisors.\allowbreak{}SemiposAtOfPositive}{91--92}
\leanref{GlobalTheory.\allowbreak{}CovHt.\allowbreak{}errX}{Development.\allowbreak{}AbelianVarietyTheory.\allowbreak{}Definitions.\allowbreak{}Divisors.\allowbreak{}SemiposAtOfPositive}{94--95}
\leanref{GlobalTheory.\allowbreak{}CompHt.\allowbreak{}compVolSlope}{Development.\allowbreak{}AbelianVarietyTheory.\allowbreak{}Definitions.\allowbreak{}Divisors.\allowbreak{}IsHomogeneousCompIdeal}{280--282}

\begin{proposition}[The deep-cover mean estimate]\label{B8.26}
Assume the first eight hypotheses of \ref{B8.25}, through nef Riemann--Roch. There is \(C_{\rm deep}\ge0\) uniform in \(D\) and core theta chains \(T\) with the chain laws and \(m^2\le k\). Take \(1\le j<m\) and a deep cover \(Y\) at level \(j\) with \(Y.\Lean{Good}(e,o)\), \(e\in\N,\ o\in\R\). Suppose the level parameter \(M'\) of \(Y\) has at most \(15\) distinct prime divisors. For every integral weight \(v\) at level \(j\) with gap \(g>0\), there are real \(s_Z\) for the components of \(Y\) with
\[
 s_Z\le\nu_Z(v;B(v,g)),\qquad
 \frac{1}{\#\operatorname{Comp}(Y)}\sum_Zs_Z
 \ge-C_{\rm deep}\mathcal E_{\ell(D)}(v,g).
 \tag{8.26b}
\]
\end{proposition}
\leanref{PositivityTheory.\allowbreak{}HtPt.\allowbreak{}Ah.\allowbreak{}deep\_\allowbreak{}slim\_\allowbreak{}use}{Development.\allowbreak{}WeilBound.\allowbreak{}CoreStatements.\allowbreak{}AutomorphicKernels.\allowbreak{}MoriwakiBigness.\allowbreak{}HilbertSamuelApplications.\allowbreak{}AhSlimU}{9--15}
\leanref{GlobalTheory.\allowbreak{}CompanionTheory.\allowbreak{}Ar.\allowbreak{}Plug.\allowbreak{}DeepStR}{Development.\allowbreak{}AutomorphicSpectrum.\allowbreak{}Definitions.\allowbreak{}GlobalGeometry.\allowbreak{}ModelSmall}{63--67}
\leanref{GlobalTheory.\allowbreak{}CovHt.\allowbreak{}MeanGe}{Development.\allowbreak{}AbelianVarietyTheory.\allowbreak{}Definitions.\allowbreak{}Divisors.\allowbreak{}SemiposAtOfPositive}{97--99}
\leanref{GlobalTheory.\allowbreak{}CovHt.\allowbreak{}MeanAll}{Development.\allowbreak{}AbelianVarietyTheory.\allowbreak{}Definitions.\allowbreak{}Divisors.\allowbreak{}SemiposAtOfPositive}{101--101}

\begin{proof}
Combine resolution and principalisation with invariant Poincar\'e--Lelong to obtain the intersection system with restriction. For the large-degree line, \(HSUse\) gives the lower successive-minima inequality. The fixed-variety inequality \ref{B8.27} gives the upper successive-minima and lower volume bounds. Semiabelian extension, GAGA for sections, Fargues's Hodge--Tate comparison, finite-flat quotients and Zarhin's construction give the good root tower \(\mathscr T\). Goodness gives primewise good level and level parameter \(A=1\). The function-field condition splits the components of the global model. Choose a good origin with the chart conditions of \ref{B8.26c} and identify the components of \(Y\) with those of \(\mathscr T\).

Set \(B=B(v,g)\), \(N=\lfloor B\rfloor\). Here \(q_{\rm sig}=m-j\ge1\), \(n=k+m-1\), and \(|\Psi|=2n_0\) with \(n_0=D.R.n\ge1\). The bound \(m^2\le k\) gives \(n\ge3q_{\rm sig}+1\ge4\), so \(B\ge128\). Among primes dividing \(2|\operatorname{disc}E|M'\), at most \(15+3\) exceed \(n_0\) by the host bound, and at most \(n_0+1\) are at most \(n_0\). Hence \(\#\mathscr T.\mathrm{bad}\le4|\Psi|n^2\le N\). Proposition~\ref{B8.26c} applies to the model slopes \(a_Z(v;N)\).

For the mixed valuations set
\[
 \begin{gathered}
 S_B=\{p\text{ prime}:p\le B\text{ or }p\mid2|\operatorname{disc}E|M'\},\\
 h_E(p)=\frac1{p-1}+\frac{40v_p(|\operatorname{disc}E|)}{[E:\Q]},\qquad
 c_2(B)=\sum_{p\in S_B}(h_E(p)+\mathbf1_{p=2})\log p .
 \end{gathered}
\]
Let \(m_\tau(t),n_\tau(t)\) count model and mixed small sections at norm bound \(e^{\tau t}\). Given \(\eta>0\), set \(t'=t-2c_2(B)\|v\|_1-\eta\). Multiply model sections by the \(d_0\) generic frame sections of degree \(k_0\), where \(d_0\) is the frame extension degree and \(8\mid k_0\). Trace pairing gives joint injectivity. After clearing denominators the log norms are at most \(\tau t'+2(\tau+k_0)c_2(B)\|v\|_1+O(1)\le(\tau+k_0)t\) for large \(\tau\). Hence \(m_\tau(t')\le n_{\tau+k_0}(t)^{d_0}\). The degree shift preserves the subsequence \(8\mid\tau\) and its asymptotic normalisation. The extra model field degree cancels \(d_0\). Nef Riemann--Roch and birational pullback and pushforward give \(\operatorname{vol}_Z(t)\ge\deg_Z(v)(a_Z(v;N)+t')\). Let \(\eta\downarrow0\) and take the infimum to obtain
\[
 a_Z(v;N)-2c_2(B)\|v\|_1\le\nu_Z(v;B).
 \tag{8.26d}
\]
\leanref{GlobalTheory.\allowbreak{}CompanionTheory.\allowbreak{}Ar.\allowbreak{}Plug.\allowbreak{}SlopeVolM.\allowbreak{}slope\_\allowbreak{}sub\_\allowbreak{}cost\_\allowbreak{}le\_\allowbreak{}compVolSlope}{Development.\allowbreak{}AutomorphicSpectrum.\allowbreak{}Proofs.\allowbreak{}GeometryProofs.\allowbreak{}UniformAllNormHeightBound}{300--372}

The prime sum of \(\log p/(p-1)\) up to \(N\) is at most \(\log N+2.2\) by \(\Lean{primeLogSum\_\allowbreak{}nat}\). At most \(N\) further terms occur, each at most \(\log(N+1)/N\). Add \(40\log\operatorname{rd}(E)\) from the discriminant and \(\log 2\) from \(p=2\). Since \(B\ge128\),
\[
 0\le c_2(B)\le60\bigl(\log B+\log\operatorname{rd}(E)\bigr).
 \tag{8.26e}
\]
Take \(s_Z=a_Z(v;N)-2c_2(B)\|v\|_1\). By \ref{B8.26c}, (8.26e) and \(N\le B\), (8.26b) holds with \(C_{\rm deep}=C_T+120\).
\end{proof}
\leanref{GlobalTheory.\allowbreak{}CompanionTheory.\allowbreak{}Ar.\allowbreak{}Plug.\allowbreak{}deepStR\_\allowbreak{}of\_\allowbreak{}use}{Development.\allowbreak{}WeilBound.\allowbreak{}CoreStatements.\allowbreak{}AutomorphicKernels.\allowbreak{}MoriwakiBigness.\allowbreak{}HilbertSamuelApplications.\allowbreak{}DeepStROfUse}{12--24}
\leanref{GlobalTheory.\allowbreak{}CompanionTheory.\allowbreak{}Ar.\allowbreak{}Link.\allowbreak{}good\_\allowbreak{}tower\_\allowbreak{}volume\_\allowbreak{}theorem\_\allowbreak{}bis}{Development.\allowbreak{}WeilBound.\allowbreak{}CoreStatements.\allowbreak{}AutomorphicKernels.\allowbreak{}MoriwakiBigness.\allowbreak{}HilbertSamuelApplications.\allowbreak{}GoodTowerVolumeTheoremBis}{9--20}
\leanref{GlobalTheory.\allowbreak{}CompanionTheory.\allowbreak{}Ar.\allowbreak{}Plug.\allowbreak{}top\_\allowbreak{}meanAll}{Development.\allowbreak{}AutomorphicSpectrum.\allowbreak{}Proofs.\allowbreak{}GeometryProofs.\allowbreak{}AbelianCmCompLawsBound}{54--138}
\leanref{GlobalTheory.\allowbreak{}CompanionTheory.\allowbreak{}Ar.\allowbreak{}Plug.\allowbreak{}deepStR\_\allowbreak{}of\_\allowbreak{}volume}{Development.\allowbreak{}AutomorphicSpectrum.\allowbreak{}Proofs.\allowbreak{}GeometryProofs.\allowbreak{}DeepStROfVolume}{21--59}
\leanref{GlobalTheory.\allowbreak{}CompanionTheory.\allowbreak{}Ar.\allowbreak{}hostLval}{Development.\allowbreak{}AutomorphicSpectrum.\allowbreak{}Definitions.\allowbreak{}GlobalGeometry.\allowbreak{}TowerOrigin}{123--126}
\leanref{Formalism.\allowbreak{}BasicNotions.\allowbreak{}Host.\allowbreak{}card\_\allowbreak{}bad\_\allowbreak{}gt\_\allowbreak{}le}{Development.\allowbreak{}AutomorphicArithmetic.\allowbreak{}PrimeDiscriminants.\allowbreak{}BadPrimes}{74--78}
\leanref{GlobalTheory.\allowbreak{}CompanionTheory.\allowbreak{}Ar.\allowbreak{}Plug.\allowbreak{}costTwo}{Development.\allowbreak{}AutomorphicSpectrum.\allowbreak{}Definitions.\allowbreak{}GlobalGeometry.\allowbreak{}ModelSmall}{54--54}
\leanref{ReciprocityTheory.\allowbreak{}ChartRestriction.\allowbreak{}G.\allowbreak{}H1.\allowbreak{}primeLogSum\_\allowbreak{}nat}{ReciprocityTheory.\allowbreak{}ReductiveGroups.\allowbreak{}ArithmeticGeometry.\allowbreak{}FirstCohomology.\allowbreak{}PrimeLogSum}{270--271}

Retain the first eight hypotheses of \ref{B8.25}. Write \(r=p_{\rm sig}\), \(s=q_{\rm sig}\), \(n=r+s\), \(f=|\Psi|\) and \(\lambda_N=\log(N\operatorname{rd}(E))+|L_E|\). The constants are uniform for a global integral model over an abelian CM host with a root tower \(\mathscr T\) of good origin and a full intersection system. Require local holomorphic lifts of the projective point maps with injective projective differential and model parameter \(4\) at the selected primes. Assume \(s\ge1\), \(3s+1\le n\) and \(\#\mathscr T.\mathrm{bad}\le4fn^2\).
\leanref{GlobalTheory.\allowbreak{}CompanionTheory.\allowbreak{}Ar.\allowbreak{}Top.\allowbreak{}GoodOrigin.\allowbreak{}Extra}{Development.\allowbreak{}AutomorphicSpectrum.\allowbreak{}Statements.\allowbreak{}TowerExistence.\allowbreak{}GoodTowerExistsX}{11--20}

Put \(d=f\binom n2\), \(e=frs\), and \(L_v=-\sum_{\tau,i}v_{\tau i}X_{\tau i}\). Let \(D\) be integration on the product of full flags in the Chern roots \(X_{\tau i}\). Obtain \(P_{N,Z}\) by dividing the limiting arithmetic intersection at cutoff \(N\) on component \(Z\) by \([F:\Q]\deg(\mathrm{Fl})\), where \(F\) is the coefficient field and \(\deg(\mathrm{Fl})\) is the degree of the flag cover. Its geometric integral is \(V_ZD\), with \(V_Z>0\). Put \(P_N=\sum_ZP_{N,Z}\) and \(V=\sum_ZV_Z\).

\begin{lemma}[A uniform lower bound for the Hodge test class]\label{B8.26a}
Let \(\mu=\sum_{\tau,i\ge r}X_{\tau i}\), and let \(J\) be the product of \(X_{\tau b}-X_{\tau a}\) over \(a<b\) in the same signature block at each \(\tau\). Put \(F_*=\mu^{e+1}J\) and \(H_*=(e+1)D(\mu^eJ)>0\). There are absolute constants \(C_h\ge0\) and \(C_u=4\) such that for every integer \(N\ge\max(128,4fn^2,\#\mathscr T.\mathrm{bad})\),
\[
 P_N(F_*)\ge -VH_*(C_h+C_u/2)fs\lambda_N .
\]
\end{lemma}
\begin{proof}
Regularise the Hodge metric corrected by half the cusp weight \(U\), with an ample perturbation of coefficient \(a/b\), \(1\le a\le b\). At a closed point, choose \(\mathcal O_E\)-linearly independent \(v_1,\ldots,v_t\in I\) spanning a finite-index \(\mathcal O_E\)-submodule of the period lattice \(I\). Complete their projections to bases of the smaller eigenspaces with vectors from the integral eigenlattices. Index and determinant comparison bounds the labelled fibre degree by
\[
 fsX-ft\log\operatorname{rd}(E)+x+ft\log(4\pi)/2,\qquad
 X=2\sum_{\substack{\ell\ {\rm prime}\\\ell\le N\ {\rm or}\ \ell\in\mathscr T.\mathrm{bad}}}
             \frac{\log\ell}{\ell-1}+2\log\operatorname{rd}(E).
\]
Here \(x\) is the sum of the smaller abelian eigenspace degrees. By the prime sum estimate in \ref{B8.26}, the sum is at most \(C_1\log N\) for an absolute \(C_1\).

Let \(y\) sum the conjugate eigenspace degrees and \(h\) be minus the total abelian degree. The determinant inclusion gives \(x+y\le-h\). Flat degree comparison and the maximum-slope bound in rank at most \(s\) give
\[
 |x-y|\le cfn\lambda_N,\qquad x\le fsc\log(2+fn+\max(h,0)),
\]
where \(c=C_s+(1+\gamma+|\log(2\pi)|)(2+C_1)\) and \(C_s\) is the absolute maximum-slope constant. For \(x>0\), these bounds give \(h\le cfn\lambda_N\). Since \(fn\le e^{\lambda_N}\) and \(\lambda_N\ge1\), it follows that \(x\le C_afs\lambda_N\) with \(C_a\ge0\) depending only on \(c\). The bound also holds for \(x\le0\). Using \(t\le s\) and \(\log(4\pi)\le\lambda_N\) bounds the labelled degree by \(C_hfs\lambda_N\), where \(C_h=2C_1+3+C_a\). The regularised metric dominates the adapted determinant metric. The point height is therefore at least \(-C_hfs\lambda_N-(a/b)K_A\), including on the boundary, where \(-K_A\) is a lower bound for the auxiliary height.
\leanref{LevelTowerAlt.\allowbreak{}height\_\allowbreak{}comparison\_\allowbreak{}global\_\allowbreak{}model}{Development.\allowbreak{}WeilBound.\allowbreak{}CoreStatements.\allowbreak{}PlurisubharmonicCones.\allowbreak{}HeightComparisonGlobalModel}{80--245}
\leanref{GlobalTheory.\allowbreak{}CompanionTheory.\allowbreak{}Ar.\allowbreak{}Cusp.\allowbreak{}TopHt.\allowbreak{}htU\_\allowbreak{}of\_\allowbreak{}given}{Development.\allowbreak{}AutomorphicSpectrum.\allowbreak{}Proofs.\allowbreak{}GeometryProofs.\allowbreak{}CuspHeightUpper}{120--223}

Normalise the ball density on the arithmetic group's fundamental set to a probability measure. Let \(\lambda_1\) be the first lattice minimum. There are \(2f\) independent integral elements with conjugates bounded by \(v_0=2f\operatorname{rd}(E)\). Short-vector scaling and packing give
\[
 \Pr(Y_1>t)\le\min\{1,2\cdot36^{a_0} e^{-b_0t}\}\qquad(t\ge0),
\]
where \(Y_1=\log^+(v_0^{-1/2}/\lambda_1)\), \(a_0=2fs\) and \(b_0=2f(n-1)\). Since \(b_0\ge3a_0\) and \(a_0\ge1\), integration gives \(\mathbb E Y_1\le(\log(2\cdot36^{a_0})+1)/b_0\le3\). The pointwise estimate \(U\le fs(2\log^+(1/\lambda_1)+\log(2fs\operatorname{rd}(E)))\) gives
\[
 \operatorname{mean}(U)\le fs(6+\log v_0+\log(2fs\operatorname{rd}(E)))\le4fs\lambda_N,
\]
using \(2fs\le N\) and \(6\le2\log(N\operatorname{rd}(E))\). Sum the fundamental-set inequalities with their sheet multiplicities as weights, then divide by the positive weighted volume.
\leanref{GlobalTheory.\allowbreak{}CompanionTheory.\allowbreak{}Ar.\allowbreak{}Cusp.\allowbreak{}fund\_\allowbreak{}set\_\allowbreak{}mean\_\allowbreak{}cuspW}{Development.\allowbreak{}AutomorphicSpectrum.\allowbreak{}Proofs.\allowbreak{}MetricProofs.\allowbreak{}CuspWeightIntegrableOnFundamentalSet}{230--328}

For sufficiently small smoothing parameter the perturbed line is smooth and conjugation-invariant, with semipositive curvature. It is nef on the fibres and ample on the generic fibre. By the point-height inequality supplied by \(HSUse\), a common lower bound for the point heights also bounds the self-intersection divided by \((e+1)\) times the geometric degree from below. Zhang gives this lower bound as part of a two-sided inequality for a relatively semiample Hermitian line bundle with smooth metric and ample generic fibre \cite[Th.~5.2]{zhang-1995}. Let \(a/b\to0\) in the Stokes formula and subtract half the mean of \(U\) to remove the metric correction. Flag pushforward multiplies the resulting normalised height by \(V_ZH_*\). Sum over \(Z\).
\end{proof}
\leanref{GlobalTheory.\allowbreak{}CompanionTheory.\allowbreak{}Ar.\allowbreak{}Top.\allowbreak{}slope\_\allowbreak{}lower\_\allowbreak{}of\_\allowbreak{}base\_\allowbreak{}settings}{Development.\allowbreak{}AutomorphicSpectrum.\allowbreak{}Proofs.\allowbreak{}DivisibilityProofs.\allowbreak{}SlopeLowerBaseSettings}{110--118}

\begin{lemma}[The two coefficients at a cutoff]\label{B8.26b}
Take \(N\) in the range of \ref{B8.26a} and put \(C_0=\max(C_h+C_u/2,0)\). The block and total translation coefficients \(A,b\) in the cutoff identity satisfy
\[
 A\ge-K\lambda_N,\qquad |rA+nb|\le4n\lambda_N,\qquad K=2(C_0+4)+3.
\]
\end{lemma}
\begin{proof}
For homogeneous \(F\) of degree \(d+1\), define \(S(F)=D(\sum_{\tau,i<r}\partial_{\tau i}F)\) and \(T(F)=D(\sum_{\tau,i}\partial_{\tau i}F)\). With \(h_{\rm fl}\) the arithmetic flag height, the cutoff identity is
\[
 P_N(F)/V=h_{\rm fl}(F)-AS(F)-bT(F)+R_N(F).
\]
Decompose \(F=\sum_{\tau,k}p_{\tau k}Q_{\tau k}\) with \(p_{\tau k}=\sum_i X_{\tau i}^k\), \(1\le k\le n\), since the flag coinvariant algebra has top degree \(d\). On \(L_v^{d+1}\), \(P_M-Vh_{\rm fl}\) alternates within blocks. At a wall, shift the first block of the specified embedding by \(2\theta\), with \(\theta=5f/M\), to obtain \(v'\). Use the weight \(v'-z\nu\), where \(z=\theta^2\) and \(\nu\) is the fixed tower weight, with smoothing parameter \(\exp(-1/z^2)\). The perturbed geometric degree is \(O(\theta)\). The height bound \(O(\log M)\) and displacement bound \(2\theta+z\|\nu\|_\infty\) give an \(O(\theta\log M)\) bound on the wall.

Fix the tower and restrict block-alternating polynomials of degree \(d+1\) to compact cross-block wall patches with relative interior. Permutations within blocks carry the adjacent wall bound to every patch. Vanishing on a patch gives its wall factor. Together with the within-block factors, these give the full Vandermonde, with block-symmetric linear quotient. On a complement of this kernel the wall suprema control the coefficient norm. The flag height vanishes on the walls, so the coefficient error is \(O(\log M/M)\). Differentiation of \(L_v^{d+1}\) identifies the Vandermonde times block-symmetric linear forms with the translation functionals. Averaging over the transitive embedding permutations gives scalar coefficients. The fixed-tower interpolation constant disappears in the cutoff limit. On changing the cutoff from \(M\) to \(N\), absorb the degree-one prime terms into these coefficients. For \(k\ge2\) the tail on \(p_{\tau k}Q\) is
\[
 \sum_{\substack{\ell>N\ {\rm prime}\\\ell\notin\mathscr T.\mathrm{bad}}}
 k\log\ell\sum_{h\ge1}\ell^{-hk}D(P_{\ell,\tau,h,k-1}Q).
\]
If \(\operatorname{frob}(\ell,\tau,h-1)=\operatorname{inl}(\sigma)\), then \(P_{\ell,\tau,h,k-1}=\sum_{i<r}X_{\sigma i}^{k-1}\). For \(\operatorname{inr}(\sigma)\), use \(P_{\ell,\tau,h,k-1}=\sum_{i\ge r}(-X_{\sigma i})^{k-1}\). The series converges absolutely. At fixed tower, the tests \(p_{\tau1}Q_1,p_{\tau2}Q_2\) with shift pairs \((r,n),(S_2,0)\), \(S_2\ne0\), force convergence of both coefficients as \(M\to\infty\). Taking limits gives the identity.
\leanref{GlobalTheory.\allowbreak{}CompanionTheory.\allowbreak{}Ar.\allowbreak{}Top.\allowbreak{}exists\_\allowbreak{}identAt}{Development.\allowbreak{}AutomorphicSpectrum.\allowbreak{}Proofs.\allowbreak{}DivisibilityProofs.\allowbreak{}TowerIdentityAtCutoff}{155--190}

On the degree-one test \(h_{\rm fl}=R_N=0\), so \(rA+nb\) is minus the \(V_Z\)-weighted mean of the flat heights \(h_{\rm rec}+\sum d_\ell\). With \(u=\log\operatorname{rd}(E)\), the local bounds are
\[
 |h_{\rm rec}|\le n(3u+|L_E|)/2,\qquad
 |d_\ell|\le n((\ell-1)^{-1}+v_\ell)\log\ell,\qquad
 \sum v_\ell\log\ell\le u .
\]
For these primes \(\sum\log\ell/(\ell-1)\le3\log N+1\), so the sum of the bounds is at most \(4n\lambda_N\).

Block translations fix \(J\), so \(S(F_*)=0\) and \(T(F_*)=fsH_*\). The dense Grassmannian points with root-of-unity chart matrices \(Z\) have squared Pl\"ucker norm \(\det(1+ZZ^*)\le(1+r)^s\) per embedding, hence point heights at most \(fs\log(n+1)/2\). Upper successive minima and flag pushforward give \(h_{\rm fl}(F_*)\le H_*fs\log(n+1)/2\). For \(X=2(rs+n)\le1.5n^2\), the test-polynomial power-sum bound at \(\ell>N\) is \(1.5fXH_*\log\ell/\ell^2\). The prime sum of \(\log\ell/\ell^2\) is at most \(2\log4/N\), so \(N\ge4fn^2\) gives \(|R_N(F_*)|\le1.6H_*\). By \ref{B8.26a} and \(fs\ge1\),
\[
 b\le C_0\lambda_N+\tfrac12\log(n+1)+1.6 .
\]
Since \(n/r\le2\), \(\log(n+1)\le\lambda_N\) and \(1.6\le\lambda_N\), the bound \(rA+nb\ge-4n\lambda_N\) yields \(A\ge-K\lambda_N\).
\end{proof}
\leanref{GlobalTheory.\allowbreak{}CompanionTheory.\allowbreak{}Ar.\allowbreak{}Top.\allowbreak{}test\_\allowbreak{}height\_\allowbreak{}bound\_\allowbreak{}slope\_\allowbreak{}flat}{Development.\allowbreak{}AutomorphicSpectrum.\allowbreak{}Proofs.\allowbreak{}DivisibilityProofs.\allowbreak{}TestHeightBound}{32--142}

\begin{proposition}[The mean of the component slopes]\label{B8.26c}
For every real weight \(v\) with the strict decrease and gap \(g>0\) of (8.26a), and every \(N\in\N\) with \(N\ge\max(128,\#\mathscr T.\mathrm{bad})\) and \(B(v,g)<N+1\), the model slopes satisfy
\[
 \frac1{\#\operatorname{Comp}}\sum_Za_Z(v;N)
 \ge-C_T\bigl[\lambda_N\|v\|_1+s\operatorname{spr}(v)\bigr],
 \qquad C_T=2C_0+17.
 \tag{8.26c}
\]
\end{proposition}
\begin{proof}
Put \(H_v=(d+1)D(L_v^d)\). The flag degree formula
\[
 D(L_v^d)=d!\prod_\tau\prod_{a<b}\frac{v_{\tau a}-v_{\tau b}}{b-a}
\]
gives \(H_v>0\). For \(F=L_v^{d+1}\), set \(s_\tau=\sum_{i<r}v_{\tau i}\), \(t_\tau=\sum_i v_{\tau i}\), \(a_\tau=\sum_i|v_{\tau i}|\). Differentiation gives
\[
 -AS(F)-bT(F)=H_v\sum_\tau(As_\tau+bt_\tau).
\]
Dominance gives \(0\le s_\tau-rt_\tau/n\le a_\tau\) and \(|t_\tau|\le a_\tau\). Split each summand as \(A(s_\tau-rt_\tau/n)+(rA+nb)t_\tau/n\). Since \(4fn^2\le B(v,g)<N+1\), \ref{B8.26b} bounds the sum below by \(-(K+4)\lambda_N\|v\|_1\).

Write \(\Delta_\tau=\max_i v_{\tau i}-\min_i v_{\tau i}\), \(B_*=B(v,g)\), \(D_*=n(1+2\|v\|_\infty/g)\). For \(2\le k\le n\), the flag insertion estimate gives \(k|D(P_{\ell,\tau,h,k-1}Q_{\tau k})|\le H_vs\Delta_\tau(2D_*)^{k-1}/k!\) for every output embedding \(\sigma\) and either block. The spread \(\Delta_\tau\) comes from \(Q_{\tau k}\). The geometric series and prime tail give a factor at most \(5.56/B_*^{k-1}\). Since \(2D_*/B_*\le1/2\) and \(\sum_{k\ge2}2^{1-k}/k!\le0.3\), one gets \(|R_N(F)|\le2H_vs\operatorname{spr}(v)\). Dominance gives \(h_{\rm fl}(F)\ge0\). The cutoff identity gives
\[
 \frac{P_N(F)}{VH_v}\ge-(K+4)\lambda_N\|v\|_1-2s\operatorname{spr}(v).
\]
Compute the geometric Hodge degree by integration over sheets and by degree pushforward. Cancelling the covering degree times the coefficient field degree gives \(c_0V_Z=\kappa I\), where \(I\) is the common ball volume of the base fundamental sets and \(c_0\ne0,\kappa\) depend only on \(f,r,s\). Thus \(V_Z=V/\#\operatorname{Comp}\). The unweighted mean of \(a_Z(v;N)=P_{N,Z}(F)/(V_ZH_v)\) is therefore \(P_N(F)/(VH_v)\). Both terms of the error bracket are nonnegative, so \(C_T=K+4+2\) suffices.
\end{proof}
\leanref{GlobalTheory.\allowbreak{}CompanionTheory.\allowbreak{}Ar.\allowbreak{}Cusp.\allowbreak{}volumeTheoremX\_\allowbreak{}of\_\allowbreak{}inputs}{Development.\allowbreak{}AutomorphicSpectrum.\allowbreak{}Proofs.\allowbreak{}DivisibilityProofs.\allowbreak{}VolumeTheoremOfInputs}{72--121}
\leanref{GlobalTheory.\allowbreak{}CompanionTheory.\allowbreak{}Ar.\allowbreak{}Top.\allowbreak{}tower\_\allowbreak{}T2\_\allowbreak{}of\_\allowbreak{}inputs}{Development.\allowbreak{}AutomorphicSpectrum.\allowbreak{}Proofs.\allowbreak{}DivisibilityProofs.\allowbreak{}TowerT2}{69--164}
\leanref{RootFieldNumerics.\allowbreak{}rest\_\allowbreak{}sum\_\allowbreak{}total\_\allowbreak{}bound}{Development.\allowbreak{}WeilBound.\allowbreak{}CoreStatements.\allowbreak{}FlagImmersions.\allowbreak{}TwoMulThreePowLeFactorial}{112--163}
\leanref{SheetDecomposition.\allowbreak{}volEq\_\allowbreak{}of}{Development.\allowbreak{}WeilBound.\allowbreak{}CoreStatements.\allowbreak{}VolumeEquality.\allowbreak{}VolEqOf}{57--97}
\leanref{RootFieldNumerics.\allowbreak{}mean\_\allowbreak{}slope\_\allowbreak{}of\_\allowbreak{}proportional\_\allowbreak{}volume}{Development.\allowbreak{}WeilBound.\allowbreak{}CoreStatements.\allowbreak{}FlagImmersions.\allowbreak{}MeanSlopeOfProportionalVolume}{13--39}

\begin{theorem}[Arithmetic volume in the semipositive case]\label{B8.27}
Let \(y\) be birational multiprojective coordinate data in a field \(\Omega\), with arithmetic variety of dimension \(d+1\) defined by \(P=\operatorname{mCoordIdeal}(y)\). Take \(P_0\subseteq P\) and an intersection system \(II\) satisfying the intersection axioms. Let \(A\) be a smooth conjugation-invariant Hermitian line datum of multidegree \(A.a\), nef on the geometric fibres and on the generic fibre. Assume the Levi form of \(\log A.h\) is nonnegative at the complex points of \(P\) in the cone-tangent directions of \(P_0\). Then
\[
 II(P;A^{d+1})\le
 \limsup_{r\to\infty}
 \frac{\log\#\{s\in\operatorname{chartLat}(y,A.a,r):
 \operatorname{IsSmallH}(P,y,A,r,s)\}}
 {r^{d+1}/(d+1)!}.
 \tag{8.27a}
\]
The chart lattice tests integrality in each admissible chart. For multihomogeneous \(g,h\) with \(g(y)=sh(y)\) and multidegree difference \(rA.a\), smallness requires \(|g(z)|\le |h(z)|A.h(z)^r\) at each complex cone point \(z\) of \(P\). Moriwaki's Corollary~6.4 gives this inequality for a smooth Hermitian line on a projective arithmetic variety, nef on every geometric fibre and with semipositive complex curvature \cite{moriwaki-2009}. Since \(P_0\subseteq P\), the curvature hypothesis tests every cone-tangent direction of \(P\).
\end{theorem}
\leanref{ReciprocityTheory.\allowbreak{}ChartRestriction.\allowbreak{}G.\allowbreak{}ArakelovGeometry.\allowbreak{}MoriwakiBigStatement}{ReciprocityTheory.\allowbreak{}ReductiveGroups.\allowbreak{}ArithmeticGeometry.\allowbreak{}ArakelovGeometry.\allowbreak{}MoriwakiSt}{29--37}
\leanref{PositivityTheory.\allowbreak{}Pf1.\allowbreak{}HilbertSamuelMoriwaki.\allowbreak{}moriwakiBig\_\allowbreak{}outright}{Development.\allowbreak{}WeilBound.\allowbreak{}CoreStatements.\allowbreak{}AutomorphicKernels.\allowbreak{}MoriwakiBigness.\allowbreak{}MoriwakiBigTheorem.\allowbreak{}Outright}{9--11}

\begin{proof}
In arithmetic dimension one the log count is at least \(r\,II(P;A)-O(1)\). In higher dimension, positive generic degree gives the lower bound along multiples by ample approximation. For generic degree zero, put \(A_q=qA+A_{\rm FS}\) with \(A_{\rm FS}\) the Fubini--Study line. The positive-degree case and the shifted small-section and rank bounds give
\[
 II(P;A_q^{d+1})\le\widehat{\operatorname{vol}}(A_q)
 \le C(q+1)\deg(A_{q,\Q}^d)\qquad(q\in\N,\ q\ge1)
\]
with \(C\) independent of \(q\). The generic-degree polynomial has degree at most \(d-1\). Divide by \(q^{d+1}\) and let \(q\to\infty\) to obtain \(II(P;A^{d+1})\le0\). The zero section gives nonnegative volume.
\end{proof}
\leanref{GroundTheory.\allowbreak{}Heights₂.\allowbreak{}Heights.\allowbreak{}HilbertSamuel.\allowbreak{}ArithmeticVolumeDegreeZero.\allowbreak{}proof}{GroundTheory.\allowbreak{}HeightsAndTheta.\allowbreak{}Heights.\allowbreak{}HilbertSamuel.\allowbreak{}CoeffLeOfEvalLeMul}{79--94}
\leanref{GroundTheory.\allowbreak{}Heights₂.\allowbreak{}Heights.\allowbreak{}Targets.\allowbreak{}MoriwakiBig.\allowbreak{}of\_\allowbreak{}found}{GroundTheory.\allowbreak{}HeightsAndTheta.\allowbreak{}Heights.\allowbreak{}OrbitLimits.\allowbreak{}MoriwakiBig}{11--16}
\leanref{UpperTheory.\allowbreak{}LiteratureTheorems.\allowbreak{}every\_\allowbreak{}character\_\allowbreak{}of\_\allowbreak{}literature\_\allowbreak{}of\_\allowbreak{}vWin}{Development.\allowbreak{}AutomorphicArithmetic.\allowbreak{}ThetaCorrespondence.\allowbreak{}ThetaPairLifts.\allowbreak{}ThirdPlaceSeries.\allowbreak{}EveryCharacterOfLiteratureOfVWin}{810--813}

At a nonterminal layer, fix a good cover \(Y\in\operatorname{Cov}(p,i)\) and put \(j=i+1\). Set
\[
 d_{\rm fl}=\dim\operatorname{Flag}(\operatorname{dat}(p,i)),\qquad
 T_*=\left\lceil P_*^2d_{\rm fl}^{\,2}\frac{k}{m}\right\rceil_{\N},
 \qquad \mathcal B=\binom{T_*-1}{d_{\rm fl}}.
\]
Let \(\lambda\) be the unshifted layer weight and \(c_0=c_0(D,j)\) the central shift. With zero-based indices, put \(u_{W,\tau,a}=\lambda_{\tau,a}-\lambda_{\tau,k-2}\), \(u_M=u_W-c_0\operatorname{hodgeWt}\), and \(\operatorname{kap}(p,i)=\sum_{\tau,a}|\lambda_{\tau,a}|\). The condition \(\Lean{Quo}(p,i,P_*,u_1)\) requires \(u_{1,\tau,k-2}=0\) and, for \(b=a+1\),
\[
 u_{1,\tau,a}-u_{1,\tau,b}
 =\left\lfloor\frac{u_{W,\tau,a}-u_{W,\tau,b}
 +\begin{cases}-c_0-(2n-2),&b=p_{\rm sig},\\2,&b\ne p_{\rm sig}\end{cases}}
 {T_*+d_{\rm fl}+1}\right\rfloor .
\]
Put \(u'=u_1-3n\operatorname{hodgeWt}\), where \(\operatorname{hodgeWt}_{\tau,a}\) equals \(0\) for \(a<p_{\rm sig}\) and \(-1\) otherwise. For \(V(u)=\prod_{\tau,a<b}(u_{\tau,a}-u_{\tau,b})\), set
\[
 \rho=1-\frac{V(u')}{V(u_1)},\qquad
 z=7\left(\frac{d_{\rm fl}^{\,2}}{T_*}+d_{\rm fl}\rho\right).
 \tag{8.28a}
\]
The size condition gives quotient weights and, for each such weight,
\[
 T_*\ge3d_{\rm fl}+5,\qquad 0\le\rho<1,\qquad z<1.
 \tag{8.28b}
\]
Let \(h_Z=Y.ht(u_1,Z)\) and \(\Delta_Z=Y.dg(u_1,Z)\).
\leanref{GlobalTheory.\allowbreak{}CovTop.\allowbreak{}lenT}{Development.\allowbreak{}AbelianVarietyTheory.\allowbreak{}Definitions.\allowbreak{}Divisors.\allowbreak{}CovolXOfCovol}{42--42}
\leanref{GlobalTheory.\allowbreak{}CovTop.\allowbreak{}zed}{Development.\allowbreak{}AbelianVarietyTheory.\allowbreak{}Definitions.\allowbreak{}Divisors.\allowbreak{}CovolXOfCovol}{44--45}
\leanref{GlobalTheory.\allowbreak{}CovTop.\allowbreak{}Quo}{Development.\allowbreak{}AbelianVarietyTheory.\allowbreak{}Definitions.\allowbreak{}Divisors.\allowbreak{}CovolXOfCovol}{49--49}
\leanref{GlobalTheory.\allowbreak{}CovTop.\allowbreak{}SizeAt}{Development.\allowbreak{}AbelianVarietyTheory.\allowbreak{}Definitions.\allowbreak{}Divisors.\allowbreak{}CovolXOfCovol}{75--79}

\begin{proposition}[Heights of the covering components]\label{B8.28}
Assume the self-intersection and small-section injection bounds, the packages \(\mathcal H_X,\mathcal H_P,\mathcal H_H\) and section comparison of \ref{B8.25}, and Grothendieck deformation theory. Assume \ref{B8.26} and uniform deep goodness, size, cover degrees and flat degrees.

There is a uniform \(C_{\rm ht}\ge0\) such that for every admissible \(u_1\) there are \(c_s,c_c,X\ge0\) and real \(s_Z\) with
\[
 \begin{gathered}
 h_Z+c_s\ge0,\qquad
 (1-\rho)\max(0,s_Z+c_s)\le h_Z+c_s+c_c,\\
 \frac1{\#\operatorname{Comp}(Y)}\sum_Zs_Z\ge-X,
 \end{gathered}
 \tag{8.28c}
\]
and
\[
 \begin{gathered}
 T_*X\le C_{\rm ht}E_p,\qquad T_*c_c\le C_{\rm ht}E_p,\\
 c_s\bigl(6d_{\rm fl}+10+T_*(z+2\rho)\bigr)\le C_{\rm ht}E_p.
 \end{gathered}
 \tag{8.28d}
\]
\end{proposition}
\leanref{SelfIntersectionVolumeEstimates.\allowbreak{}height\_\allowbreak{}lower\_\allowbreak{}bound\_\allowbreak{}general\_\allowbreak{}alt}{Development.\allowbreak{}AutomorphicArithmetic.\allowbreak{}ThetaCorrespondence.\allowbreak{}ThetaPairLifts.\allowbreak{}ThirdPlaceSeries.\allowbreak{}HeightLowerBoundGeneral}{21--22}
\leanref{GlobalTheory.\allowbreak{}CovTop.\allowbreak{}HeightAt}{Development.\allowbreak{}AbelianVarietyTheory.\allowbreak{}Definitions.\allowbreak{}Divisors.\allowbreak{}CovolXOfCovol}{146--152}

\begin{proof}
Deep goodness gives a finite nonempty component set, with \(\Delta_Z>0\) by the degree condition. Toroidal and semiabelian extension give goodness and birationality on \(Z\), with \(\Delta'_Z\ge(1-\rho)\Delta_Z\) for the degree of \(u'\). The common short-generator log bound \(c_s\) gives \(h_Z+c_s\ge0\).

The carrying sections \(\theta_Z\) of \(\Lean{Mult}\) cost at most \(c_{\rm mult}\ge0\). The chain laws and parameter choice give \(m^2\le k\). The good cover's budget gives \(\#\{\ell\text{ prime}:\ell\mid M'\}\le15\) and a gap \(g>0\) for \(u'\). Apply Proposition~\ref{B8.26}. For \(B=B(u',g)\), let \(H_{\rm loc}=\|u'\|_1\sum_{p\in S_B}h_E(p)\log p\) be the mixed valuation cost. Put
\[
 c_c=c_{\rm mult}+H_{\rm loc},\qquad X=C_{\rm deep}\mathcal E_{\ell(D)}(u',g).
 \tag{8.28e}
\]
At \(t=c_s\) the mixed volume is at least \(\max(0,\Delta'_Z(s_Z+c_s))\). By \(\Lean{HTCmp}\), multiplying the balanced scalar in degree \(8r\) by \(\prod_{p\in S_B}p^{\lceil8rh_E(p)\|u'\|_1\rceil}\) clears denominators at log cost at most \(8rH_{\rm loc}+\sum_{p\in S_B}\log p\). For the resulting scalar \(\alpha_r\), the map \(f\mapsto\alpha_r\theta_Z^rf\) adds at most \(8rc_c+O(1)\) to log norms and is injective by generic nonvanishing. For every \(\eta>0\) it lands in sections small at \(c_s+c_c+\eta\) for large \(r\). Comparison with the self-intersection upper bound \(\Lean{NoBase}\) gives
\[
 \Delta'_Z(s_Z+c_s)\le\Delta_Z(h_Z+c_s+c_c+\eta).
 \tag{8.28f}
\]
Let \(\eta\downarrow0\). For \(s_Z+c_s\ge0\) use the degree ratio, and for \(s_Z+c_s<0\) use \(h_Z+c_s\ge0\).

The generator, carrying-section and budget estimates give
\[
 \begin{gathered}
 c_s\bigl(6d_{\rm fl}+10+T_*(z+2\rho)\bigr)\le C_sE_p,\qquad
 T_*c_{\rm mult}\le C_mE_p,\\
 T_*H_{\rm loc},\,T_*\mathcal E_{\ell(D)}(u',g)\in[0,C_bE_p].
 \end{gathered}
\]
Take \(C_{\rm ht}=C_s+C_m+C_b+C_{\rm deep}C_b\).
\end{proof}
\leanref{GlobalTheory.\allowbreak{}CovHt.\allowbreak{}HtTopG.\allowbreak{}slope\_\allowbreak{}le\_\allowbreak{}height}{Development.\allowbreak{}AutomorphicArithmetic.\allowbreak{}ThetaCorrespondence.\allowbreak{}ThetaPairLifts.\allowbreak{}HodgeBundleWeights.\allowbreak{}HeightAddLogNonneg}{29--65}
\leanref{GlobalTheory.\allowbreak{}CovHt.\allowbreak{}HtTopG.\allowbreak{}heightAt\_\allowbreak{}eq\_\allowbreak{}of\_\allowbreak{}cover}{Development.\allowbreak{}AutomorphicArithmetic.\allowbreak{}ThetaCorrespondence.\allowbreak{}ThetaPairLifts.\allowbreak{}HodgeBundleWeights.\allowbreak{}HeightAddLogNonneg}{66--105}
\leanref{GlobalTheory.\allowbreak{}CovHt.\allowbreak{}HtTopG.\allowbreak{}height\_\allowbreak{}of\_\allowbreak{}meanAll\_\allowbreak{}bound}{Development.\allowbreak{}AutomorphicArithmetic.\allowbreak{}ThetaCorrespondence.\allowbreak{}ThetaPairLifts.\allowbreak{}ThirdPlaceSeries.\allowbreak{}HeightOfSteps}{19--38}
\leanref{GlobalTheory.\allowbreak{}CompHt.\allowbreak{}le\_\allowbreak{}compVolSlope\_\allowbreak{}iff}{Development.\allowbreak{}AbelianVarietyTheory.\allowbreak{}Definitions.\allowbreak{}Divisors.\allowbreak{}IsHomogeneousCompIdeal}{283--286}

\section{From component averages to the core lattice}

Put \(r_Z=Y.rkC(u_M,Z)\), \(d_Z=Y.degC(u_M,Z)\), and
\(\mu_Y(u_M)=Y.\mu(\operatorname{wtOf}(\operatorname{dat}(p,i),u_M,0))\).
The cover-degree and share estimates give, for some \(\Delta_*>0\),
\[
 \Delta_*\le\Delta_Z\le\frac{\Delta_*}{1-\rho},\qquad
 (1-z)r_Z\le\mathcal B\Delta_Z\le r_Z.
 \tag{8.29a}
\]
The height sum bound gives \(A>0,\ R\ge0\) with
\[
 \begin{split}
 d_Z&\ge\mathcal B\Delta_Z
 \bigl(T_*h_Z-3(d_{\rm fl}+1)|h_Z|-A\log(2+|h_Z|)-R\bigr),\\
 A\log\max(A,1)+R&\le C_{\rm HS}E_p,
 \end{split}
 \tag{8.29b}
\]
with uniform \(C_{\rm HS}\). The squared-norm degree inequality is
\[
 2\frac{\sum_Zd_Z}{\sum_Zr_Z}\le\mu_Y(u_M).
 \tag{8.29c}
\]
\leanref{GlobalTheory.\allowbreak{}CovTop.\allowbreak{}DegAt}{Development.\allowbreak{}AbelianVarietyTheory.\allowbreak{}Definitions.\allowbreak{}Divisors.\allowbreak{}CovolXOfCovol}{140--141}
\leanref{GlobalTheory.\allowbreak{}CovTop.\allowbreak{}ShareAt}{Development.\allowbreak{}AbelianVarietyTheory.\allowbreak{}Definitions.\allowbreak{}Divisors.\allowbreak{}CovolXOfCovol}{143--144}
\leanref{GlobalTheory.\allowbreak{}CovTop.\allowbreak{}HSAt}{Development.\allowbreak{}AbelianVarietyTheory.\allowbreak{}Definitions.\allowbreak{}Divisors.\allowbreak{}CovolXOfCovol}{153--156}
\leanref{GlobalTheory.\allowbreak{}CovTop.\allowbreak{}SumAt}{Development.\allowbreak{}AbelianVarietyTheory.\allowbreak{}Definitions.\allowbreak{}Divisors.\allowbreak{}CovolXOfCovol}{158--159}

\begin{proposition}[Lower bounds for the cover and core slopes]\label{B8.29}
Assume uniform deep goodness, size, cover degrees and shares, \ref{B8.28}, and (8.29b)--(8.29c). Then
\[
 \mu_Y(u_M)\ge-C_UE_p,\qquad
 C_U=2(3C_{\rm ht}+C_{\rm HS}+4).
 \tag{8.29d}
\]
Assume further the terminal defect bound, good-cover existence at each nonterminal layer, and the ring, integral map, cover-rank, short-generator and flatness comparisons. The core type lattice in its core norm then satisfies
\[
 \mu_{\rm core}(p,i)\ge-M_{p,i}-C_{\rm core}E_p.
 \tag{8.29e}
\]
For terminal constant \(C_{\rm term}\) and flatness constant \(C_F\ge0\) one may take
\[
 C_{\rm core}=\max\{C_{\rm term},\,2(C_F+C_U)+9\}.
 \tag{8.29f}
\]
\end{proposition}
\leanref{GlobalTheory.\allowbreak{}T.\allowbreak{}uniform\_\allowbreak{}lower\_\allowbreak{}bound\_\allowbreak{}of\_\allowbreak{}sums'}{Development.\allowbreak{}AutomorphicArithmetic.\allowbreak{}ThetaCorrespondence.\allowbreak{}ThetaPairLifts.\allowbreak{}ThirdPlaceSeries.\allowbreak{}ULowOfStatement}{11--12}
\leanref{GlobalTheory.\allowbreak{}T.\allowbreak{}liftLowBound\_\allowbreak{}of\_\allowbreak{}cover\_\allowbreak{}estimates}{Development.\allowbreak{}AutomorphicArithmetic.\allowbreak{}ThetaCorrespondence.\allowbreak{}ThetaPairLifts.\allowbreak{}ThirdPlaceSeries.\allowbreak{}LowOfCoverHypothesesStatement}{11--12}

\begin{proof}
Write \(T=T_*\), \(d=d_{\rm fl}\), \(w_Z=\mathcal B\Delta_Z\), \(a=T-3d-4\ge0\), and \(L_A=A\log\max(A,1)\). For \(t\ge0\), use \(\log y\le y-1\) at \(y=(2+t)/\max(A,1)\) to obtain
\[
 A\log(2+t)\le t+L_A+2.
 \tag{8.29g}
\]
Since \(|h_Z|\le h_Z+2c_s\), (8.29b) gives
\[
 d_Z\ge w_Z\bigl(a(h_Z+c_s)-N_0\bigr),\qquad
 N_0=ac_s+2c_s(3d+4)+L_A+2+R.
 \tag{8.29h}
\]
Put \(n_Y=\#\operatorname{Comp}(Y)\) and \(m_0=\max\{0,(1-\rho)c_s-X-c_c\}\). By (8.28c), \(\sum_Z(h_Z+c_s)\ge n_Ym_0\). Since \(h_Z+c_s\ge0\), the bounds \(w_Z\ge\mathcal B\Delta_*\) and \((1-\rho)\sum_Zw_Z\le n_Y\mathcal B\Delta_*\) give
\[
 \sum_Zd_Z\ge\left(\sum_Zw_Z\right)
 \bigl(a(1-\rho)m_0-N_0\bigr).
 \tag{8.29i}
\]
Using \((1-\rho)m_0\ge(1-2\rho)c_s-X-c_c\) and \(a\le T\) gives
\[
 \begin{split}
 a(1-\rho)m_0-N_0&\ge-H,\\
 H&=T(X+c_c)+c_s\bigl(6d+10+T(z+2\rho)\bigr)+L_A+R+4.
 \end{split}
 \tag{8.29j}
\]
By \(0<\sum_Zw_Z\le\sum_Zr_Z\) the mean degree is at least \(-H\). Equations (8.28d) and (8.29b) with \(E_p\ge1\) give \(H\le(3C_{\rm ht}+C_{\rm HS}+4)E_p\). Apply (8.29c).

For a good nonterminal cover and quotient weight write \(G_s(c)\) for shift by \(c\). Let \(A_-\ne0\) and \(B\) be the cover and deep type lattices at shifts \(-c_0\) and \(0\). By flatness and (8.29d),
\[
 \mu(A_-)\ge-M_{p,i}-(C_F+C_U)E_p,
 \tag{8.29k}
\]
where \ref{B8.24a} gives the twist estimate on the constructed covers. For \(\mathcal K=P_*\operatorname{kap}(p,i)\), size gives
\[
 0<b\le E_p,\qquad b\le\mathcal K,\qquad |M_{p,i}|\le\mathcal K.
 \tag{8.29l}
\]
The integral maps \(\phi_h:A_-\to B\), \(1\le h\le s\), have joint kernel zero and satisfy
\[
 \begin{gathered}
 \sum_h\|\phi_hf\|_\sigma^2\le e^{2b}\|f\|_\sigma^2,\\
 s\operatorname{rk}B\le Y.\operatorname{rk}(G_s(c_0(D,j)))
 \le(1+b/\mathcal K)\operatorname{rk}A_-.
 \end{gathered}
 \tag{8.29m}
\]
The injection into \(B^s\) has height at most \(2b\). Short generators span \(B\) over its coefficient field with norms at most \(e^{\mathcal K}\), hence \(\mu_{\min}(B^s)\ge-2\mathcal K\). For \(r=\operatorname{rk}A_-\) and \(N=s\operatorname{rk}B\), the quotient degree inequality gives
\[
 r\bigl(\mu(A_-)-2b\bigr)\le N\mu(B)+2(N-r)\mathcal K .
\]
Since \(r\le N\le(1+b/\mathcal K)r\), any \(t\le\mu(A_-)-2b\) yields
\[
 \mu(B)\ge t-\frac b{\mathcal K}(|t|+2\mathcal K).
 \tag{8.29n}
\]
Take \(t=-M_{p,i}-(C_F+C_U)E_p-2b\). By (8.29l),
\[
 \frac b{\mathcal K}|t|\le b+(C_F+C_U)E_p+2b.
\]
Thus \(\mu(B)\ge-M_{p,i}-2(C_F+C_U)E_p-7b\). The ring comparison gives \(\mu_{\rm core}\ge\mu(B)-2b\). Use \(b\le E_p\) and combine with the terminal bound to obtain (8.29e)--(8.29f).
\end{proof}
\leanref{GlobalTheory.\allowbreak{}CovTop.\allowbreak{}FromComp.\allowbreak{}core}{Development.\allowbreak{}AbelianVarietyTheory.\allowbreak{}Proofs.\allowbreak{}DivisorProofs.\allowbreak{}UniformLowerBoundOfSums}{15--33}
\leanref{Formalism.\allowbreak{}ExtensionFiltration.\allowbreak{}sumComponents'}{Formalism.\allowbreak{}MeanHeightDescent.\allowbreak{}AverageDegreeBounds.\allowbreak{}SumComponents}{8--119}
\leanref{GlobalTheory.\allowbreak{}DescentFamily.\allowbreak{}descent}{Development.\allowbreak{}AbelianVarietyTheory.\allowbreak{}Proofs.\allowbreak{}DivisorProofs.\allowbreak{}HtMapLeOfPoints}{69--144}
\leanref{GlobalTheory.\allowbreak{}CovTop.\allowbreak{}LowCover.\allowbreak{}lowAt\_\allowbreak{}of\_\allowbreak{}cover}{Development.\allowbreak{}AbelianVarietyTheory.\allowbreak{}Proofs.\allowbreak{}DivisorProofs.\allowbreak{}LowerBoundOfAllEstimates}{65--106}
\leanref{GlobalTheory.\allowbreak{}CovTop.\allowbreak{}LowCover.\allowbreak{}low\_\allowbreak{}of\_\allowbreak{}cover\_\allowbreak{}hypotheses}{Development.\allowbreak{}AutomorphicArithmetic.\allowbreak{}ThetaCorrespondence.\allowbreak{}ThetaPairLifts.\allowbreak{}ThirdPlaceSeries.\allowbreak{}LowOfCoverHypotheses}{16--32}

\section{Removing the bad part}

Let \(\Lambda_{\rm type}\subseteq\Lambda_{\rm all}\) be the type and ambient lattices in the rational realisation at layer \(i\), and let \(e=e_{\rm good}\) be the good projector. Set
\[
 \Lambda_{\rm good}=\Lambda_{\rm type}\cap\operatorname{im}e,\qquad
 \Lambda_{\rm bad}=\Lambda_{\rm type}\cap\ker e,\qquad
 \delta=\exp(-x^{e_z/7}).
 \tag{8.30a}
\]
These lattices are integral, and \(e_z=F_m.\mathrm{ez}=\varepsilon/J\) with \(J\ge1\). Here \(\delta\) controls discarded rank. In the arithmetic application of Lemma~\ref{B4.55a}, take \(\delta_0=8f(k+m-1)/\min(g_0,g_1)\). Lemma~\ref{B5.09a} gives \(3\delta_0Q_0\le1\).
\leanref{GlobalTheory.\allowbreak{}CovTop.\allowbreak{}latGood}{Development.\allowbreak{}AbelianVarietyTheory.\allowbreak{}Definitions.\allowbreak{}Divisors.\allowbreak{}NormNegCmConstantEqAtAllEmbeddings}{44--46}
\leanref{GlobalTheory.\allowbreak{}CovTop.\allowbreak{}latBad}{Development.\allowbreak{}AbelianVarietyTheory.\allowbreak{}Definitions.\allowbreak{}Divisors.\allowbreak{}NormNegCmConstantEqAtAllEmbeddings}{47--49}

\begin{proposition}[Rank of the bad lattice]\label{B8.30}
Assume \(\mathcal H_X,\mathcal H_P\) and section comparison from \ref{B8.25}, Grothendieck deformation theory, both compactifications, Lie-bundle seminegativity and Green's cuspidal parametrisation. Assume the uniform Gram, deep-cover existence and goodness, size, flat-degree, cover-degree and share conditions. For each eligible \(F_m,CH\) there is \(x_0\) such that at every admissible point with \(x>F_m.x_1\) and \(x\ge x_0\), every layer satisfies
\[
 \operatorname{rk}\Lambda_{\rm bad}\le3\delta\operatorname{rk}\Lambda_{\rm type}.
 \tag{8.30b}
\]
\end{proposition}
\leanref{GlobalTheory.\allowbreak{}T.\allowbreak{}liftBadSet\_\allowbreak{}of\_\allowbreak{}deep\_\allowbreak{}size\_\allowbreak{}flatDeg\_\allowbreak{}share}{Development.\allowbreak{}AutomorphicArithmetic.\allowbreak{}ThetaCorrespondence.\allowbreak{}ThetaPairLifts.\allowbreak{}ThirdPlaceSeries.\allowbreak{}CoveringBadLocusOfDescentSteps}{54--55}
\leanref{GlobalTheory.\allowbreak{}CovTop.\allowbreak{}BadAt}{Development.\allowbreak{}AbelianVarietyTheory.\allowbreak{}Definitions.\allowbreak{}Divisors.\allowbreak{}CovolXOfCovol}{69--71}

\begin{proof}
For large \(x\), \(\Lean{SizesRec}\) with exponent \(6\) and \(x^{e_z}\ge3/2\) give \(D.R.r\ge100\) and \(g_0,g_1\ge4096\). For standard nontrivial additive characters, \ref{B7.50} and label compatibility give \ref{B7.51} for both \(e_U\)-rows. For bad-label sets \(S_a\) this gives
\[
 u_a=\operatorname{locShare}_a(S_a)\le W(D.R.f,m,D.R.r,g_a,q,d)
 \le x^{21}e^{-x^{e_z/6}}\le\delta/20
\]
eventually, since \(21\log x-x^{e_z/6}+x^{e_z/7}\to-\infty\). In particular,
\[
 u_0+u_1\le\delta.
 \tag{8.30c}
\]

For \(i+1<m\), take a good deep cover with level group \(\Gamma\). Write \(G_0=\prod_{a=0,1}\prod_{v\mid g_a}\GL_{D.R.r}(\F_{g_a})\), with \(v\) in the prescribed place sets. The map \(\pi:\Gamma\to G_0\) has a homomorphic section \(s\) whose image commutes with \(\ker\pi\). Let \(P\) be upper unitriangular in each factor with a nondegenerate Whittaker character, and \(H_a\) upper unitriangular off \(g_a\) and unrestricted above \(g_a\). Put \(\widetilde P=\pi^{-1}(P)\), \(\widetilde H_a=\pi^{-1}(H_a)\). For the cover representation \(\rho\) on \(V\), the two ranks in (8.30b) are the dimensions of \(V_\lambda\cap\ker\rho(s_*e_0)\) and \(V_\lambda\), where \(V_\lambda\) is the pulled-back character space and \(e_0\) projects onto good types.

Write \(c_0=c_0(D,i+1)\). Rank growth gives
\[
 Y.\operatorname{rk}(G_s(2c_0))\le\tfrac32Y.\operatorname{rk}(G_s(-c_0)).
\]
Separation by regular cusp sections gives \(V|_{\widetilde H_a}\hookrightarrow\C[\widetilde H_a]^{\,n_a}\) with \(n_a|\widetilde H_a|\le\tfrac32|\widetilde P|\dim V_\lambda\). If \(b_a\) projects onto \(S_a\) at \(g_a\), then \(1-e_0=b_0+b_1-b_0b_1\) bounds the bad dimension by the sum of the ranks of the lifted \(b_a\) on \(V_\lambda\). In \(G_0\), \(b_a\) times the character projector is an idempotent supported on \(H_a\) with identity coefficient \(u_a/|P|\), by character independence of the \(e_U\) weights. Lift by \(s_*\) and average over \(\ker\pi\). The resulting projection has support in \(\widetilde H_a\) and identity coefficient \(u_a/|\widetilde P|\). Trace in the regular copies bounds its rank by
\[
 n_a|\widetilde H_a|u_a/|\widetilde P|\le\tfrac32u_a\dim V_\lambda .
\]
Sum over \(a\) and use (8.30c). At \(i+1=m\), the deformation-class summands are freely permuted, so \(n_a|\widetilde H_a|=|\widetilde P|\dim V_\lambda\). The same trace bound gives \(\operatorname{rk}\Lambda_{\rm bad}\le\delta\operatorname{rk}\Lambda_{\rm type}\).
\end{proof}
\leanref{GlobalTheory.\allowbreak{}CovTop.\allowbreak{}LocShareNum.\allowbreak{}local\_\allowbreak{}share\_\allowbreak{}count\_\allowbreak{}le\_\allowbreak{}exp\_\allowbreak{}alt}{Development.\allowbreak{}AutomorphicArithmetic.\allowbreak{}ThetaCorrespondence.\allowbreak{}ThetaPairLifts.\allowbreak{}ThirdPlaceSeries.\allowbreak{}LocalShareCountLeExp}{31--56}
\leanref{GlobalTheory.\allowbreak{}CovTop.\allowbreak{}LocShareNum.\allowbreak{}wg\_\allowbreak{}sum\_\allowbreak{}le}{Development.\allowbreak{}AbelianVarietyTheory.\allowbreak{}Proofs.\allowbreak{}DivisorProofs.\allowbreak{}LocalShareSumBound}{271--283}
\leanref{GlobalTheory.\allowbreak{}BadEmb.\allowbreak{}regular\_\allowbreak{}embeddings\_\allowbreak{}cover\_\allowbreak{}bad\_\allowbreak{}embedding\_\allowbreak{}locus\_\allowbreak{}alt}{Development.\allowbreak{}AutomorphicArithmetic.\allowbreak{}ThetaCorrespondence.\allowbreak{}ThetaPairLifts.\allowbreak{}ThirdPlaceSeries.\allowbreak{}RegularEmbeddingsCoverBadEmbeddingLocus}{38--86}
\leanref{GlobalTheory.\allowbreak{}BadIndef.\allowbreak{}bad\_\allowbreak{}components\_\allowbreak{}indefinite\_\allowbreak{}alt}{Development.\allowbreak{}AutomorphicArithmetic.\allowbreak{}ThetaCorrespondence.\allowbreak{}ThetaPairLifts.\allowbreak{}ThirdPlaceSeries.\allowbreak{}BadComponentsIndefinite}{24--54}
\leanref{GlobalTheory.\allowbreak{}CovTop.\allowbreak{}definite\_\allowbreak{}bad\_\allowbreak{}place\_\allowbreak{}lift\_\allowbreak{}alt}{Development.\allowbreak{}AutomorphicArithmetic.\allowbreak{}ThetaCorrespondence.\allowbreak{}ThetaPairLifts.\allowbreak{}ThirdPlaceSeries.\allowbreak{}ExistsLiftAtDefiniteBadPlaces}{21--53}
\leanref{GlobalTheory.\allowbreak{}CovTop.\allowbreak{}covering\_\allowbreak{}bad\_\allowbreak{}locus\_\allowbreak{}of\_\allowbreak{}descent\_\allowbreak{}steps\_\allowbreak{}alt}{Development.\allowbreak{}AutomorphicArithmetic.\allowbreak{}ThetaCorrespondence.\allowbreak{}ThetaPairLifts.\allowbreak{}ThirdPlaceSeries.\allowbreak{}CoveringBadLocusOfDescentSteps}{21--40}

\begin{lemma}[Slope loss under orthogonal truncation]\label{B8.31}
Over a totally complex number field \(K\), let \(V\) have Hermitian norms \(N_\sigma\) admitting a Gram frame. Take a finitely generated \(\mathcal O_K\)-module \(M\subset V\), a submodule \(0\ne L\subseteq M\), and a \(K\)-linear idempotent \(e\) with
\[
 N_\sigma(ex+y-ey)^2=N_\sigma(ex)^2+N_\sigma(y-ey)^2
 \quad\text{for all }\sigma,x,y.
 \tag{8.31a}
\]
Suppose \(g\in\N\), \(g\ge1\), and \(ge(L)\subseteq L\). Let \(0\le\delta\le1/12\) and \(\mathcal Q\in\R\). If
\[
 \begin{gathered}
 \operatorname{rk}(L\cap\ker e)\le3\delta\operatorname{rk}L,\\
 \mu_{\max}(M)\le\mathcal Q,\qquad -\mathcal Q\le\mu(L),\qquad
 2\log g\le\mathcal Q,
 \end{gathered}
 \tag{8.31b}
\]
then
\[
 \mu(L\cap\operatorname{im}e)\ge\mu(L)-12\delta\mathcal Q.
 \tag{8.31c}
\]
\end{lemma}
\leanref{GlobalTheory.\allowbreak{}CutNum.\allowbreak{}slope\_\allowbreak{}good\_\allowbreak{}ge}{Development.\allowbreak{}AbelianVarietyTheory.\allowbreak{}Proofs.\allowbreak{}DivisorProofs.\allowbreak{}MajAtAddRight}{111--143}

\begin{proof}
Put \(L_g=L\cap\operatorname{im}e\) and \(L_b=L\cap\ker e\), with ranks \(d_g,d_b\). Since \(gL\subseteq L_g\oplus L_b\subseteq L\),
\[
 \operatorname{rk}L=d_g+d_b.
 \tag{8.31d}
\]
Polarise (8.31a) and use multiplication by an element of \(K\) of nonreal image to obtain Hermitian orthogonality. The map \(1-e\) embeds \(L/(L_g\oplus L_b)\) into \(g^{-1}L_b/L_b\). Its order is at most \(g^{[K:\Q]d_b}\), so the normalised degree correction is at most \(2d_b\log g\). Bounding the degree of \(L_b\) by \(d_b\mu_{\max}(M)\) gives
\[
 (d_g+d_b)\mu(L)\le d_g\mu(L_g)+d_b\bigl(\mu_{\max}(M)+2\log g\bigr).
 \tag{8.31e}
\]
The rank bound gives \((1-3\delta)d_b\le3\delta d_g\), hence \(d_g>0\) and
\[
 d_b\le4\delta d_g.
 \tag{8.31f}
\]
Since \(\mathcal Q\ge2\log g\ge0\), (8.31e) and (8.31b) give \(d_g(\mu(L)-\mu(L_g))\le3d_b\mathcal Q\le12\delta d_g\mathcal Q\).
\end{proof}
\leanref{GlobalTheory.\allowbreak{}CutNum.\allowbreak{}numeric}{Development.\allowbreak{}AbelianVarietyTheory.\allowbreak{}Proofs.\allowbreak{}DivisorProofs.\allowbreak{}MajAtAddRight}{93--109}
\leanref{Formalism.\allowbreak{}ExtensionFiltration.\allowbreak{}goodBadCut}{Formalism.\allowbreak{}MeanHeightDescent.\allowbreak{}DivisionInequalityWithCovolume}{23--23}

\section{The layer inequality and the common constant}

Give \(\Lambda_{\rm all}\) the rational-realisation norm and restrict it to its sublattices. Gram and recursion give
\[
 \Lambda_{\rm all}\text{ finitely generated},\qquad
 \Lambda_{\rm good}\ne0,\qquad
 \frac{\operatorname{lcv}(X_{W,i},\Lambda_i)}{D_{\rm norm}}
 \le-\mu(\Lambda_{\rm good}).
 \tag{8.32a}
\]
The truncation condition makes \(e\) orthogonal with denominator \(g=|p.P.Gp|\), so \(ge(\Lambda_{\rm type})\subseteq\Lambda_{\rm type}\) and \(2\log g\le b\). The maximum-slope and size conditions supply family exponents \(A,B\) with
\[
 \mu_{\max}(\Lambda_{\rm all})\le x^A b,\qquad
 1\le b\le E_p\le x^B b,\qquad |M_{p,i}|\le x^B b.
 \tag{8.32b}
\]
Full size includes (8.29l). It also requires quotient weights at nonterminal layers and (8.28b) for each such weight. The core comparisons are
\[
 \mu(\Lambda_{\rm type})\ge\mu_{\rm core}-C_1E_p,\qquad
 \mu_{\rm core}\ge-M_{p,i}-C_2E_p.
 \tag{8.32c}
\]
\leanref{GlobalTheory.\allowbreak{}CovTop.\allowbreak{}GramAt}{Development.\allowbreak{}AbelianVarietyTheory.\allowbreak{}Definitions.\allowbreak{}Divisors.\allowbreak{}CovolXOfCovol}{53--55}
\leanref{GlobalTheory.\allowbreak{}CovTop.\allowbreak{}RecAt}{Development.\allowbreak{}AbelianVarietyTheory.\allowbreak{}Definitions.\allowbreak{}Divisors.\allowbreak{}CovolXOfCovol}{57--59}
\leanref{GlobalTheory.\allowbreak{}CovTop.\allowbreak{}CutAt}{Development.\allowbreak{}AbelianVarietyTheory.\allowbreak{}Definitions.\allowbreak{}Divisors.\allowbreak{}CovolXOfCovol}{61--67}
\leanref{GlobalTheory.\allowbreak{}CovTop.\allowbreak{}CoreAt}{Development.\allowbreak{}AbelianVarietyTheory.\allowbreak{}Definitions.\allowbreak{}Divisors.\allowbreak{}CovolXOfCovol}{92--94}
\leanref{GlobalTheory.\allowbreak{}CovTop.\allowbreak{}LowAt}{Development.\allowbreak{}AbelianVarietyTheory.\allowbreak{}Definitions.\allowbreak{}Divisors.\allowbreak{}CovolXOfCovol}{96--96}

\begin{lemma}[The covolume inequality at one layer]\label{B8.32}
Fix a point and layer with \(x>0\). Let \(A,B,P_*,C_1,C_2\in\R\) with \(C_1,C_2\ge0\). Assume the Gram, recursive, truncation and bad-rank conditions, the maximum-slope and full size conditions, and (8.32c). If
\[
 12e^{-x^{e_z/7}}(x^A+x^B+C_1+C_2+1)\le1,
 \tag{8.32d}
\]
then
\[
 \frac{\operatorname{lcv}(X_{W,i},\Lambda_i)}{D_{\rm norm}}
 \le M_{p,i}+(C_1+C_2+1)E_p.
 \tag{8.32e}
\]
\end{lemma}
\leanref{GlobalTheory.\allowbreak{}CovTop.\allowbreak{}Assembly.\allowbreak{}lcv\_\allowbreak{}le}{Development.\allowbreak{}AbelianVarietyTheory.\allowbreak{}Proofs.\allowbreak{}DivisorProofs.\allowbreak{}CovolumeFromGramAndBounds}{39--85}

\begin{proof}
Put \(H=x^A+x^B+C_1+C_2+1\) and \(\mathcal Q=HE_p\). By (8.32d), \(0<\delta\le1/12\). The rational realisation has totally complex coefficient field. Apply \ref{B8.31} with \(L=\Lambda_{\rm type}\), \(M=\Lambda_{\rm all}\). Gram and truncation give the lattice and projector hypotheses, and \ref{B8.30} gives the bad-rank bound. By (8.32b), \(\mu_{\max}(\Lambda_{\rm all})\le\mathcal Q\) and \(2\log g\le b\le\mathcal Q\). By (8.32c),
\[
 -\mu(\Lambda_{\rm type})\le M_{p,i}+(C_1+C_2)E_p
 \le x^B b+(C_1+C_2)E_p\le\mathcal Q.
\]
The loss \(12\delta HE_p\) is at most \(E_p\). Combine (8.32a) and (8.32c).
\end{proof}

\begin{proposition}[One constant for all points and layers]\label{B8.33}
Assume the family estimates
\[
 \begin{gathered}
 \Lean{GramCondition},\quad\Lean{LiftRecursive},\quad
 \Lean{LiftTruncation},\quad\Lean{LiftBadSet},\\
 \Lean{LiftLambdaBound},\quad\Lean{SizeCondition},\quad
 \Lean{CoreBound'},\quad\Lean{LiftLowBound}.
 \end{gathered}
\]
Then \(\Lean{CovolumeBoundOfX}\) holds. For each eligible family and height parameter one may take
\[
 C_0=C_1+C_2+1
 \tag{8.33a}
\]
with \(C_1,C_2\ge0\) from (8.32c).
\end{proposition}
\leanref{GlobalTheory.\allowbreak{}CovTop.\allowbreak{}Assembly.\allowbreak{}covolX\_\allowbreak{}of\_\allowbreak{}lift\_\allowbreak{}bounds\_\allowbreak{}and\_\allowbreak{}coreLe}{Development.\allowbreak{}AutomorphicArithmetic.\allowbreak{}ThetaCorrespondence.\allowbreak{}ThetaPairLifts.\allowbreak{}ThirdPlaceSeries.\allowbreak{}CovolXOfSteps}{16--40}

\begin{proof}
Choose \(A,B,C_1,C_2\) for \(F_m,CH\). Since \(e_z=F_m.ez>0\), \(e^{-x^{e_z/7}}x^a\to0\) for \(a=A,B,0\). Take \(x_0\) beyond the threshold in (8.32d) and the eight point thresholds. Apply \ref{B8.32} at \(x>F_m.x_1>1\) and \(x\ge x_0\). Since \(E_p\ge1\), the bound holds for every \(C\ge C_0\).
\end{proof}
\leanref{GlobalTheory.\allowbreak{}CovTop.\allowbreak{}Assembly.\allowbreak{}exists\_\allowbreak{}loss\_\allowbreak{}le}{Development.\allowbreak{}AbelianVarietyTheory.\allowbreak{}Proofs.\allowbreak{}DivisorProofs.\allowbreak{}CovolumeFromGramAndBounds}{25--35}

\section{Transfer to the selected indefinite realisation}

\begin{theorem}[Covolume bounds on the linked bundles]\label{B8.34}
Assume the discrete-series covering statement, Harish--Chandra module theory, rapid decay of cusp forms, the Casselman--Wallach statement, and the ray-class-field statement in the formulations for the linked bundles. Assume \(\Lean{CovolumeBoundOfX}\).

Let \(h_L\) give local Langlands parameters, \(N\) be a norm-choice extension, and \(S,\Gamma\) be as in \ref{B8.08}. Assume a height-counting extension with its law for each datum admitted by \(N\). If the parameters satisfy \(\Lean{OpenParameterStatements}\), then
\[
 \Lean{LinkedAt}\bigl(
 \operatorname{parametersOfExtension}(N,S,\Gamma),
 \Lean{CovolumeBoundThirdPlaceSeries}\bigr).
 \tag{8.34a}
\]
An allowed strengthening \(R\) has a common row witness. For \(\mathcal A,\sigma\) satisfying it, the estimate is closed on \(W_3(\mathcal A)\) as in \ref{B8.39}. For eligible \(F_m,CH\) there is \(C_0\ge0\) such that for every \(C\ge C_0\) some \(x_0\) gives
\[
 \frac{\operatorname{lcv}(X_{W,j},\Lambda_j)}{D_{\rm norm}}
 \le M_{p,j}+CE_p
 \tag{8.34b}
\]
at every selected index and layer for \(x>F_m.x_1,\ x\ge x_0\), in the selected indefinite realisation at \(p\).
\end{theorem}
\leanref{GlobalTheory.\allowbreak{}Glob.\allowbreak{}CovolRow.\allowbreak{}linked\_\allowbreak{}bundles\_\allowbreak{}l3\_\allowbreak{}of\_\allowbreak{}open\_\allowbreak{}alt}{Development.\allowbreak{}AutomorphicArithmetic.\allowbreak{}ThetaCorrespondence.\allowbreak{}ThetaPairLifts.\allowbreak{}ThirdPlaceSeries.\allowbreak{}LinkedBundlesL3OfOpen}{31--46}
\leanref{GlobalTheory.\allowbreak{}Fc.\allowbreak{}S0.\allowbreak{}CovolumeBoundThirdPlaceSeries}{Development.\allowbreak{}AutomorphicArithmetic.\allowbreak{}ThetaCorrespondence.\allowbreak{}ThetaPairLifts.\allowbreak{}ThirdPlaceSeries.\allowbreak{}GeometricWeakTies}{32--33}
\leanref{GlobalTheory.\allowbreak{}Fc.\allowbreak{}Row}{Development.\allowbreak{}AbelianVarietyTheory.\allowbreak{}Definitions.\allowbreak{}GlobalTwisting.\allowbreak{}ValidForConstructedPointClassesEventuallyWorking}{268--269}
\leanref{GlobalTheory.\allowbreak{}Fc.\allowbreak{}W3}{Development.\allowbreak{}AbelianVarietyTheory.\allowbreak{}Definitions.\allowbreak{}GlobalTwisting.\allowbreak{}ValidForConstructedPointClassesEventuallyWorking}{371--372}
\leanref{GlobalTheory.\allowbreak{}Fc.\allowbreak{}Ind.\allowbreak{}EvWC}{Development.\allowbreak{}AbelianVarietyTheory.\allowbreak{}Definitions.\allowbreak{}Divisors.\allowbreak{}BaseRankDimension}{90--91}

\begin{proof}
Combine the ray-class-field input and height-counting extensions with the open divisor row to obtain the class-sum divisor row. Conjoin it with the primitive, divisor, goodness and geometric rows to choose \(R\) and its witness. For \(\mathcal A,\sigma\) satisfying the witness, the geometric row gives compatible component data \(\tau\). The representation-theoretic inputs give the eventual point conditions for \(\sigma_\tau\), and the class-sum row gives \(\Lean{RecTied}\). The seven construction duties and working-class hypotheses give \(\Lean{FamHyp}\) on \(W_3(\mathcal A)\). Beyond the common threshold apply \ref{B8.24} and evaluate \(L3\) on the selected summand with the indefinite data at \(\mathrm{false}\) and the record for \(C,CH\).
\end{proof}
\leanref{GlobalTheory.\allowbreak{}Fc.\allowbreak{}Ind.\allowbreak{}eventual\_\allowbreak{}L3\_\allowbreak{}of\_\allowbreak{}covolume\_\allowbreak{}alt}{Development.\allowbreak{}AutomorphicArithmetic.\allowbreak{}ThetaCorrespondence.\allowbreak{}ThetaPairLifts.\allowbreak{}ThirdPlaceSeries.\allowbreak{}EventualL3OfCovolume}{14--20}
\leanref{GlobalTheory.\allowbreak{}Fc.\allowbreak{}Ind.\allowbreak{}evWC\_\allowbreak{}onPt\_\allowbreak{}of\_\allowbreak{}guarded\_\allowbreak{}alt}{Development.\allowbreak{}AutomorphicArithmetic.\allowbreak{}ThetaCorrespondence.\allowbreak{}ThetaPairLifts.\allowbreak{}ThirdPlaceSeries.\allowbreak{}EvWCOnPtOfGuarded}{12--32}
\leanref{GlobalTheory.\allowbreak{}Fc.\allowbreak{}Ind.\allowbreak{}summandL3Bound}{Development.\allowbreak{}AutomorphicArithmetic.\allowbreak{}ThetaCorrespondence.\allowbreak{}ThetaPairLifts.\allowbreak{}ThirdPlaceSeries.\allowbreak{}SummandFunctionalBound}{54--54}

In \ref{B8.01}, this linked row supplies the covolume hypothesis of \ref{B8.03}.
\leanref{UpperTheory.\allowbreak{}Lit.\allowbreak{}nonempty\_\allowbreak{}hyps\_\allowbreak{}definite\_\allowbreak{}side\_\allowbreak{}of\_\allowbreak{}closed\_\allowbreak{}bound}{Development.\allowbreak{}AutomorphicArithmetic.\allowbreak{}ThetaCorrespondence.\allowbreak{}ThetaPairLifts.\allowbreak{}ThirdPlaceSeries.\allowbreak{}SiegelEisensteinSecondTermHypotheses}{65--68}


% ---- p8c4 ----
\chapter{The definite side and local errors}\label{ch:B8-definite-local}

Use \ref{B8.02}. Write \(\mathcal L\) for the size datum of \(D\), \(K=D.\mathrm{what}_2.\mathrm{LL}\), \(L=x^\varepsilon\), \(b=e^L\) and \(a=xb\).
\section{The definite column condition}

For a word \(u\), let \(R\) be its numerical record on layers \(J\) and types \(\mathrm{Pl}\), with weights \(\nu_{\ell,\tau,i}\). Put \(N=R.Nn\), \(T=T_u\), \(L_a=\log(aR.\mathrm{rd})\), \(n_{\ell,\tau}=\sum_i\nu_{\ell,\tau,i}\) and \(p_\ell=|\mathrm{Pl}|g_\ell^0\). Bars denote averages over infinite places \(v\).
\begin{align*}
\bar\omega_\tau&=\log\Omega,\qquad
\bar\vartheta\le C_\theta L,\qquad g_\ast\le C_8L,\\
M_{\ell,v}&\le C_2fm^2N^2(g_\ast+\log T)+C_efN^2g_\ast,\\
h_\ell,\ h'_\ell&\le C_4fN^2(L+\log T+g_\ast),\qquad
0\le g_\ell^0\le N^2,\\
\log(8+32p_\ell V)&\le c_5L_a,\qquad
\log(2n_{\ell,\tau}+2)\le c_6L_a .
\end{align*}
The terminal log bound \(A\) satisfies
\begin{align*}
A&\le-\sum_\tau n_\tau\log T+\sum_\tau(2c_\tau^*+\xi_\tau^*)+e_{\mathrm{ev}},\\
c_\tau^*&\le (N-1)D_{1,\tau}\log(2a),\qquad
\xi_\tau^*\le\tfrac12N(N-1)\log(2a),\\
n_\tau&=\lambda_\tau-D_{0,\tau},\qquad
D_{1,\tau}\le3D_{0,\tau}+S+\eta,\\
D_{0,\tau}&\le S+m\eta+
\begin{cases}b_\tau,&\tau\in TT,\\0,&\tau\notin TT,\end{cases}\\
\lambda_\tau&=a_\tau-\tfrac{N-k}{2},\qquad b_\tau\le2b'
\quad(\tau\in TT),\\
\lambda_\tau&\ge S-\tfrac{N-1}{2}\quad(\tau\notin TT),\\
0&\le\log T\le C_TmL,\qquad e_{\mathrm{ev}}\le C_{\mathrm{ev}}fN^2L_a .
\end{align*}
The \(C\)-constants and \(c_5,c_6\) are nonnegative, with
\[
R.C_h\ge\max\{2(16+C_T+C_{\mathrm{ev}}),\
C_\theta+c_5+2c_6+1+2C_4+2C_eC_8,\
2(C_2+C_4)(C_T+C_8)\}.
\]
Put
\begin{align}
s_{\ell,v}
 &=\sum_\tau\left[\sum_i(\nu_{\ell,\tau,i}\log\nu_{\ell,\tau,i}-\nu_{\ell,\tau,i})
       -n_{\ell,\tau}\omega_{\tau,v}+1+2g_\ell^0\log(2n_{\ell,\tau}+2)\right]\notag\\
 &\quad+p_\ell\bigl(\vartheta_v+\log(8+32p_\ell V)\bigr)+2M_{\ell,v}+h_\ell+h'_\ell.
\label{eq:B8-height-size}
\end{align}
Use \(\log z=\log|z|\) for \(z\ne0\) and \(\log0=0\).
\leanref{Formalism.\allowbreak{}L2.\allowbreak{}DIV.\allowbreak{}HtNums}{Formalism.\allowbreak{}LocalToGlobal.\allowbreak{}Divisibility.\allowbreak{}Inputs}{210--311}
\leanref{Formalism.\allowbreak{}DimensionVariation.\allowbreak{}liftedMapLogSize}{Formalism.\allowbreak{}ThetaSupremumBounds.\allowbreak{}DivIneqExplicit6}{15--23}
\begin{definition}[Definite rows for the selected bundle]\label{B8.35}
Fix \(D,O,\mathbf b\) with size constant \(c_{\mathrm{sz}}\), and realisation constants \(C,C_H\). Each choice of primitive linkage at level one, separation, covolume, surjectivity, isometry and passage witnesses gives a package \(\mathscr P\) with definite face \(D_c\).

For a face member \(w\), \(\mathrm{RowsU}(R,w,\mathcal N)\) requires \(D_c.a(w)>0\) and an integral lattice \(\mathcal N\subset K^{D_c.a(w)}\). Its chain is
\[
K^{d_0}\xrightarrow{K_0}K^{d_1}\longrightarrow\cdots
\xrightarrow{K_{r-1}}K^{d_r},\qquad d_0=D_c.a(w).
\]
Take \(\operatorname{lf}:J\to\{0,\ldots,r-1\}\) and frames \(X_{\alpha,v}\in\GL_{d_\alpha}(\C)\) with \(X_{0,v}=D_c.X_c(w,v)\). In isometric coordinates starting with \(\mathrm{colsIn}\), the chain represents \(n\) theta kernel steps and the complementary map.

Require \(P_{\mathrm{end}}:K^{d_r}\to(K^{n_0})^\vee\) linear and frames \(Y_v\) for the form \(H_{\mu_0}\) of base dimension \(n_0\), with
\begin{align}
D_c.Pt(w,\mu_0)(\phi,F)&=P_{\mathrm{end}}(K_{r-1}\cdots K_0\phi)(F),\label{eq:B8-column-factor}\\
(Y_v^{-1})^\ast Y_v^{-1}&=v(H_{\mu_0}).\label{eq:B8-definite-frame}
\end{align}
Choose a height table \(\mathfrak h\) for \(R\) as above. Let \(p,q\) index coefficients in normed complex spaces of adjacent layers. With \(K_{\alpha,v}=v(K_\alpha)\) and \(T_{\alpha pq,v}\) continuous linear, require
\begin{align*}
0<\gamma_{\alpha p}&\le1,\qquad
\sum_p\gamma_{\alpha p}\|\Phi_{\alpha p,v}(z)\|^2=\|X_{\alpha,v}^{-1}z\|_2^2,\\
\Phi_{\alpha+1,q,v}(K_{\alpha,v}z)
 &=\sum_p\gamma_{\alpha p}\epsilon_{\alpha pq,v}T_{\alpha pq,v}\Phi_{\alpha p,v}(z),\\
|\epsilon_{\alpha pq,v}|&=1,\qquad \|T_{\alpha pq,v}\|^2\le t_{\max}(\alpha,v),\qquad t_{\max}(\alpha,v)\ge0 .
\end{align*}
Set \(\mathrm{CB.sq}(\alpha,v)=a_\alpha a_{\alpha+1}h_\alpha^{\mathrm{cl}}h_{\alpha+1}^{\mathrm{cl}}t_{\max}(\alpha,v)\), where \(a_\alpha\) counts groups and \(h_\alpha^{\mathrm{cl}}\) bounds their sizes. Require
\begin{align}
\|P_{\mathrm{end}}\|_{X_{r,v}\to Y_v^\vee}^{\,2}&\le e^A,\label{eq:B8-terminal-bound}\\
\log\mathrm{CB.sq}(\alpha,v)&\le
\sum_{\ell:\operatorname{lf}(\ell)=\alpha}s_{\ell,v}.\label{eq:B8-layer-bound}
\end{align}

Assume wordwise \(\Lean{NumFacts}\) at \(\lambda_0,c_0,c_e,S_w,C_1\in\R\) with \(C_h=c_{\mathrm{sz}}\). For each three-sided \(I\) tied to integral face submodules \(\mathrm{NuC}\) and agreeing with \(\mathbf b.P.I\) under transport, require \(\mathrm{RowsU}(R_u,I.w_U(u),e_u^{-1}(\mathrm{NuLat}_u))\) for every \(u\). Here \(e_u\) gives coordinates, and \(R_u=\Lean{RecNums.\allowbreak{}ofNum}\) uses the letter and level \(T_u\) together with these constants.
\end{definition}
\leanref{GlobalTheory.\allowbreak{}Fc.\allowbreak{}DefiniteSideCondition}{Development.\allowbreak{}AutomorphicArithmetic.\allowbreak{}ThetaCorrespondence.\allowbreak{}ThetaPairLifts.\allowbreak{}ThirdPlaceSeries.\allowbreak{}DefiniteSideCondition}{25--35}
\leanref{Formalism.\allowbreak{}L2.\allowbreak{}RowDataTransport.\allowbreak{}RowsU}{Formalism.\allowbreak{}LocalToGlobal.\allowbreak{}PtKOfReal}{141--159}
\leanref{Formalism.\allowbreak{}L2.\allowbreak{}DIV.\allowbreak{}ColChain}{Formalism.\allowbreak{}LocalToGlobal.\allowbreak{}Divisibility.\allowbreak{}Inputs}{319--341}
\leanref{Formalism.\allowbreak{}L2.\allowbreak{}DIV.\allowbreak{}FirstCols}{Formalism.\allowbreak{}LocalToGlobal.\allowbreak{}Divisibility.\allowbreak{}DivisionHeightBoundStatement}{504--507}
\begin{theorem}[Columns from point evaluation]\label{B8.36}
For \(C_Y,C_T\in\R\), \(C_T\ge0\), there is \(C_0\) uniform in \(D,O,c_{\mathrm{sz}},\mathbf b,C,C_H\). Assume \(\Lean{RealTied}\), wordwise symmetric columns satisfying \(\Lean{ChiFin0}\), and \(\Lean{ColumnCharactersCompatible}\) for the selected characters and evaluations. Let \(W\) agree with the bundle weights and satisfy \(\Lean{KerTies}\). Assume for all \(C_Y',C_T'\in\R\) that
\[
\mathrm{HtRowsS}(\mathcal L,C_Y'),\quad X'=\exp(x^{\varepsilon/3})\ge8,\quad x\ge2,\quad
(\forall u,\ 4a<T_u),\quad(\forall u,\ \log T_u\le C_T'mL)
\]
imply \(\Lean{KerSizes}\) for \(W\) at \(C_Y',C_T'\).

Assume the displayed hypotheses at \(C_Y,C_T\) and, for each \(u\),
\begin{equation}\label{eq:B8-divided-lattice}
N'_u=I_u\mathrm{NuLat}_u
\end{equation}
as integral submodules, with \(N'_u\) a lattice and \(\nu_u\) of positive rank. For \(c_{\mathrm{sz}}\ge C_0\) one has \(\mathrm{DefiniteSideCondition}(D,O,c_{\mathrm{sz}},\mathbf b,C,C_H)\).
\end{theorem}
\leanref{GlobalTheory.\allowbreak{}Fc.\allowbreak{}RL.\allowbreak{}DefSide.\allowbreak{}definiteSideCondition\_\allowbreak{}of\_\allowbreak{}nprimeDef\_\allowbreak{}of\_\allowbreak{}nuFull}{Development.\allowbreak{}AutomorphicArithmetic.\allowbreak{}ThetaCorrespondence.\allowbreak{}ThetaPairLifts.\allowbreak{}ThirdPlaceSeries.\allowbreak{}PointEvaluationSide}{49--64}
\leanref{GlobalTheory.\allowbreak{}ColRows.\allowbreak{}rowsU\_\allowbreak{}of\_\allowbreak{}colPins\_\allowbreak{}alt}{Development.\allowbreak{}AutomorphicArithmetic.\allowbreak{}ThetaCorrespondence.\allowbreak{}ThetaPairLifts.\allowbreak{}ThirdPlaceSeries.\allowbreak{}RowsUOfColPins}{17--84}
\begin{proof}
The theta steps \(q_j\) join \(\mathcal F_0,\ldots,\mathcal F_n\), with complement \(q_{\mathrm{cpl}}:\mathcal F_n\to\mathcal S\). Point evaluation gives the isometry \(f_0=\mathrm{colsIn}\) and \(g:\mathcal S\to(K^{n_0})^\vee\) with
\begin{equation}\label{eq:B8-period-kernel}
D_c.Pt(u,\mu_0)(\phi,F)=g(q_{\mathrm{cpl}}q_{n-1}\cdots q_0f_0\phi)(F).
\end{equation}
The kernel estimates give \(C_2=C_\theta=C_4=c_5=c_6=C_8=C_e=\Gamma(C_Y)\ge0\). The last map gives \(C_{\mathrm{ev}}=e_{\mathrm{ev}}=0\). Thus the threshold for \(R.C_h=c_{\mathrm{sz}}\) is
\begin{equation}\label{eq:B8-column-threshold}
C_0=\max\{2(16+C_T),\ 6\Gamma+1+2\Gamma^2,\ 4\Gamma(C_T+\Gamma)\}.
\end{equation}
\leanref{GlobalTheory.\allowbreak{}ColRowsTopS.\allowbreak{}rowsBody\_\allowbreak{}colFace}{Development.\allowbreak{}AbelianVarietyTheory.\allowbreak{}Proofs.\allowbreak{}DivisorProofs.\allowbreak{}ColumnsMemberRowsOfLogLevelLe}{25--58}
\leanref{GlobalTheory.\allowbreak{}ColRows.\allowbreak{}exists\_\allowbreak{}htNums\_\allowbreak{}of\_\allowbreak{}kerSide\_\allowbreak{}lastSide}{Development.\allowbreak{}AbelianVarietyTheory.\allowbreak{}Proofs.\allowbreak{}DivisorProofs.\allowbreak{}ExistsHtNumsOfKerSideLastSide}{11--27}

Choose framed coordinates \(e_j\) for \(1\le j\le n\) and \(e_{\mathcal S}\) on \(\mathcal S\). With \(e_0=f_0\), take \(K_j=e_{j+1}^{-1}q_je_j\) for \(j<n\), \(K_n=e_{\mathcal S}^{-1}q_{\mathrm{cpl}}e_n\) and \(P_{\mathrm{end}}=ge_{\mathcal S}\). Equation \eqref{eq:B8-period-kernel} gives \eqref{eq:B8-column-factor} and the chain ties.

Together with point evaluation, interlacing and same shift give the last-map bound in frames \(Y_v\) of \(H_{\mu_0}\),
\[
|v(g(\psi)(F))|^2\le e^A\|\psi\|_{\mathcal S,v}^2\|F\|_{Y_v,v}^2.
\]
Assign layer \(\ell\) to step \(\min(\mathrm{no}(\ell),n)\). Put \(t_{\alpha,v}=\sum_{\operatorname{lf}(\ell)=\alpha}s_{\ell,v}\) and \(\|z\|_{X_\alpha,v}=\|X_{\alpha,v}^{-1}v(z)\|_2\). Then
\[
\|K_\alpha z\|_{X_{\alpha+1},v}^2\le e^{t_{\alpha,v}}\|z\|_{X_\alpha,v}^2 .
\]
Since \(K\) is totally complex, \(v(K)\) is dense in \(\C\). Continuity gives \eqref{eq:B8-terminal-bound}.

Use one coefficient per layer on \(\C^{d_\alpha}\), with \(\gamma=\epsilon=1\) and
\[
\Phi_{\alpha,v}(z)=X_{\alpha,v}^{-1}z,\qquad
T_{\alpha,v}=X_{\alpha+1,v}^{-1}v(K_\alpha)X_{\alpha,v}.
\]
Use one group of size one and \(\mathrm{CB.sq}(\alpha,v)=t_{\max}=e^{t_{\alpha,v}}\).
\leanref{GlobalTheory.\allowbreak{}DefRowsChain.\allowbreak{}rowsBody\_\allowbreak{}of\_\allowbreak{}framed}{Development.\allowbreak{}AbelianVarietyTheory.\allowbreak{}Proofs.\allowbreak{}DivisorProofs.\allowbreak{}KerChainMatricesWithNormBounds}{227--269}
\leanref{GlobalTheory.\allowbreak{}DefRowsCB.\allowbreak{}exists\_\allowbreak{}cbData}{Development.\allowbreak{}AbelianVarietyTheory.\allowbreak{}Proofs.\allowbreak{}DivisorProofs.\allowbreak{}ExistsCbData}{13--56}

For \(0\ne a_u\in I_u\), \(a_u\mathrm{NuLat}_u\subset N'_u\) is finitely generated. Its \(K\)-span contains \(N'_u\) by \eqref{eq:B8-divided-lattice}. Thus \(\mathrm{NuLat}_u\) pulls back to a lattice, and \(D_c.a(u)=\dim_K(\nu_u).V>0\).
\leanref{GlobalTheory.\allowbreak{}ColRows.\allowbreak{}nuLat\_\allowbreak{}comap\_\allowbreak{}isLattice}{Development.\allowbreak{}AbelianVarietyTheory.\allowbreak{}Proofs.\allowbreak{}DivisorProofs.\allowbreak{}IsLatticeOfEqSpanFracSmul}{52--60}

Use the package ties with the division and word ties to transport the rows to \(I.w_U(u)\).
\end{proof}
\section{A common bound along the selected family}
\begin{theorem}[The definite condition along members]\label{B8.37}
For each local Langlands witness \(h_L\) and member parameter datum \(p\), \(\Lean{LinkedAt}\,p\,\Lean{RowDefectCondition}\) holds. For each linked selection \(\sigma\), the following is closed on \(W_3\mathcal A\).
\[
\begin{gathered}
\exists C_0\ge0,\ \forall C\ge C_0,\ \exists x_0,\
\forall x>F_m.x_1,\ x\ge x_0,\ \forall i\in F_m.Ix(x),\\
\mathrm{DefiniteSideCondition}
(D_{x,i},O_{x,i},F_m.C_{\mathrm{sz}},\sigma_{x,i},C,C_H).
\end{gathered}
\]
Under real second-layer ties, \(\mathrm{TuSize}\) and \(\mathrm{ChosenE}(F_m,3,C_H')\) give \(\mathrm{NuFullF}(F_m)\) and \(\mathrm{MuRank}(F_m)\), hence a lattice \(N'_u\) and positive rank of \(\nu_u\) eventually at each point and word.
\end{theorem}
\leanref{GlobalTheory.\allowbreak{}Fc.\allowbreak{}RL.\allowbreak{}DefSide.\allowbreak{}linkedAt\_\allowbreak{}df\_\allowbreak{}of\_\allowbreak{}parameters}{Development.\allowbreak{}AutomorphicArithmetic.\allowbreak{}ThetaCorrespondence.\allowbreak{}ThetaPairLifts.\allowbreak{}ThirdPlaceSeries.\allowbreak{}PointEvaluationSide}{65--83}
\leanref{GlobalTheory.\allowbreak{}Fc.\allowbreak{}RL.\allowbreak{}DefSide.\allowbreak{}full\_\allowbreak{}rank}{Development.\allowbreak{}AbelianVarietyTheory.\allowbreak{}Proofs.\allowbreak{}GlobalInterpolation.\allowbreak{}WeilUnitaryEntries}{25--32}
\begin{proof}
Transport point evaluation by linkage. Under second-layer ties, \(\Lean{family\_\allowbreak{}theorems\_\allowbreak{}two\_\allowbreak{}and\_\allowbreak{}five}\) uses integral pairing and ramified Hecke vanishing. With \(\mathrm{TuSize}\) this controls levels eventually. Commensurability and the chosen proviso give lattice and rank thresholds \(a_1,a_2\).

To set the middle threshold in Theorem~\ref{B8.36}, use outer entries \(6,13\), indexed from zero, as \(C_Y,C_T\), with \(C_T\) replaced by its nonnegative part. The selected bundle supplies \(n\in\N\) with \(D.s.Del=n\). With \(\mathrm{TuSize}(F_m)\), \(\Lean{famRowsSizes}\) gives the five size hypotheses for \(x\ge a_3\). Apply Theorem~\ref{B8.36} at \(x\ge\max(a_1,a_2,a_3)\) with \eqref{eq:B8-divided-lattice} and take \(C_0=0\).
\end{proof}
\begin{theorem}[One constant for both realisations]\label{B8.38}
Assume the propositions \(\Lean{DiscreteSeriesCover}\), \(\Lean{HarishChandraModules}\), \(\Lean{CuspRapidDecay}\), \(\Lean{CasselmanWallach}\) and \(\Lean{DiscreteSeriesClassification}\) of \ref{B8.01}, and the surjectivity and lower-weight inputs. At a slot \(Q\) with linkage \(L_k\), assume the primitive, divisor, goodness, geometric compatibility, third-layer covolume, functional, multiplicity, \(\kappa\) and defect rows.

For each class \(\mathcal A\) and raw selection \(\sigma\) satisfying the slot member condition, the following is closed on \(\mathrm D3\,3\,\mathcal A\). For each \(F_m,C_H\) satisfying the class and largeness conditions of \ref{B8.39}, some \(C\ge0\) gives
\[
\begin{gathered}
\mathrm{IndefiniteRealisationConstantBounds}
(\mathbf b.B.Q,C,C_H,\mathrm{DIVinN}(F_m),\mathrm{LeafUniqueMap}(\mathbf b))\\
{}\land \mathrm{DefiniteSideCondition}(D,O,c_{\mathrm{sz}},\mathbf b,C,C_H)
\end{gathered}
\]
for some \(x_0\), at every \(x>F_m.x_1\) at least \(x_0\) and the selection threshold, and every \(i\in F_m.Ix(x)\). Here \(\mathbf b=\sigma(F_m,h|_{\mathcal A},x,h_x,i)\) for \(h:\mathrm D3\,3\,\mathcal A(F_m)\), \(h_x:F_m.x_1<x\), and \(c_{\mathrm{sz}}=F_m.C_{\mathrm{sz}}\).
\end{theorem}
\leanref{GlobalTheory.\allowbreak{}Fc.\allowbreak{}StageZeroRow.\allowbreak{}row\_\allowbreak{}rRows\_\allowbreak{}stage0\_\allowbreak{}alt}{Development.\allowbreak{}AutomorphicArithmetic.\allowbreak{}ThetaCorrespondence.\allowbreak{}ThetaPairLifts.\allowbreak{}ThirdPlaceSeries.\allowbreak{}TupAllUsedAtPart}{46--67}
\leanref{GlobalTheory.\allowbreak{}Fc.\allowbreak{}row\_\allowbreak{}rRows\_\allowbreak{}of\_\allowbreak{}parts\_\allowbreak{}alt}{Development.\allowbreak{}AutomorphicArithmetic.\allowbreak{}ThetaCorrespondence.\allowbreak{}ThetaPairLifts.\allowbreak{}ThirdPlaceSeries.\allowbreak{}RowRRowsOfParts}{21--63}
\leanref{GlobalTheory.\allowbreak{}Fc.\allowbreak{}R1.\allowbreak{}SelectionDefiniteRows}{Development.\allowbreak{}AutomorphicArithmetic.\allowbreak{}ThetaCorrespondence.\allowbreak{}ThetaPairLifts.\allowbreak{}ThirdPlaceSeries.\allowbreak{}CountsIrreducibilityConstants}{91--93}

Miller and Schmid prove rapid decay on every Siegel set for smooth cuspidal automorphic forms of moderate growth \cite[Theorem~A]{miller-2012}. The unitary input assumes uniform moderate growth and gives \(\sup_g|\phi(g)|H_{\mathrm{quot}}(g)^N<\infty\) for every \(N\in\N\), with \(H_{\mathrm{quot}}\) the height infimum over rational left translates.

\begin{proof}
Fix an archimedean place \(\tau\) and \(K_{\tau'}\)-finite vectors in the other factors \(\tau'\ne\tau\). The source in the remaining variable is the irreducible unitary discrete-series model \(H_\tau\). Proposition~\ref{B8.38a} extends the map continuously and equivariantly to the smooth vectors of \(H_\tau\). If the map is nonzero on a \(K_\tau\)-finite vector, the lower seminorm \(q_S\) may depend on the map and on the list \(S\). Transport the primitive and divisor rows to compatible summands. Class sums and the real inputs give eventual point rows. Classification gives small scale. With the seven duties these rows give the norm bound via \(\Lean{More}\) and surjectivity, and the \(L6\) bound via lower weight. The slot supplies covolume and functional bounds.

Conjoin the closed rows on \(W_3\mathcal A\) and restrict to \(\mathrm D3\,3\,\mathcal A\). Take \(C=\max_{1\le r\le6}C_r\), where \(C_1\ge0\) comes from small scale. Take \(x_0\) to be the maximum of the six thresholds at \(C\) and those for multiplicity and \(\kappa\). The indefinite entries are
\[
\begin{array}{c|ccccccc}
\text{row}&1&2&3&4&5&7&8\\ \hline
\text{entries}&w1,\mathrm{seeIso},\mathrm{pass}&\mathrm{covol}&
\mathrm{onto}&\mathrm{lamLow}&\mathrm{fun0}&\mathrm{miStar}&\mathrm{kap}
\end{array}
\]
Rows \(3,5\) use \(\mathrm{DIVinN}\). Row \(4\) uses the leaf hypothesis and row \(6\) the definite condition. The unitary-group hypotheses supply surjectivity and lower weight for \ref{B8.06}.
\leanref{GlobalTheory.\allowbreak{}Fc.\allowbreak{}S0.\allowbreak{}ontoX\_\allowbreak{}and\_\allowbreak{}lamLow\_\allowbreak{}of\_\allowbreak{}unitaryGroupHypotheses'}{Development.\allowbreak{}AutomorphicArithmetic.\allowbreak{}ThetaCorrespondence.\allowbreak{}ThetaPairLifts.\allowbreak{}ThirdPlaceSeries.\allowbreak{}SurjectiveAndLowWeightOfUnramifiedExponents}{16--29}
\end{proof}
\begin{proposition}[Sobolev bounds for unitary intertwiners]\label{B8.38a}
Let \(r,s\) be nonnegative integers and put \(G=\U(r,s)\), with maximal compact subgroup \(K=\U(r)\times\U(s)\) and real Lie algebra \(\mathfrak g\). Let \((\sigma,H)\) be a strongly continuous unitary representation on a complex Hilbert space whose closed invariant subspaces are \(0\) and \(H\). Write \(V=H_{K\text{-fin}}\) and \(H^\infty\) for its \(K\)-finite and smooth vectors. Let \((\pi,E)\) be a smooth complex Fr\'echet representation with continuous derived operators. For each defining seminorm \(p_i\) assume there are a continuous seminorm \(b_i\) and \(N_i\in\N\) such that
\[
p_i(\pi(g)z)\le h(g)^{N_i}b_i(z),\qquad h(g)=1+\operatorname{tr}(g^*g)
\quad(g\in G,\ z\in E).
\]
Let \(T_0:H\to E\) be complex linear and intertwine \(K\) and \(\mathfrak g\) on \(V\). There is a linear \(T:H\to E\) agreeing with \(T_0\) on \(V\) whose restriction to \(H^\infty\) is continuous and \(G\)-equivariant. This restriction is unique and its image in \(E\) is closed.

If \(T_0v_0\ne0\) for some \(v_0\in V\), then for each finite list \(S\) of words in \(\mathfrak g\) there is a continuous seminorm \(q_S\) on \(E\) such that
\[
\sum_{D\in S}\|d\sigma(D)w\|\le q_S(Tw)\qquad(w\in H^\infty).
\]
Words act by ordered composition, with the identity for the empty word.
\end{proposition}
\leanref{GlobalTheory.\allowbreak{}Published.\allowbreak{}CasselmanWallach}{Development.\allowbreak{}AbelianVarietyTheory.\allowbreak{}PublishedTheorems.\allowbreak{}AutomorphicForms₂.\allowbreak{}HasDRepIffHasDerivAt}{98--111}
\leanref{PositivityTheory.\allowbreak{}CWB.\allowbreak{}casselmanWallach}{Development.\allowbreak{}WeilBound.\allowbreak{}CoreStatements.\allowbreak{}PublishedTheorems.\allowbreak{}CasselmanWallach.\allowbreak{}FinalB}{8--9}
\begin{proof}
For \(H=0\), take \(T=0\). Suppose \(H\ne0\). The Harish--Chandra module \(V\) has smooth globalization \(H^\infty\). Fix a Haar measure \(dg\). The Schwartz algebra \(\mathcal S(G)\) consists of smooth densities \(f\,dg\) for which \(h^nD_LD_R f\) lies in \(L^1(G,dg)\) for every \(n\in\N\) and all left- and right-invariant differential operators \(D_L,D_R\). The corresponding \(L^1\)-seminorms define its topology \cite[\S2.5]{bernstein-2014}. The growth bound gives its action on \(E\) by integration. Since \(T_0(V)\subset E_{K\text{-fin}}\) by \(K\)-equivariance, apply Bernstein and Kr\"otz \cite[Theorem~10.6 and Proposition~11.2]{bernstein-2014} to extend \(T_0|_V\) uniquely to a continuous \(G\)-map on \(H^\infty\), then extend linearly to \(H\).
\leanref{PositivityTheory.\allowbreak{}CW.\allowbreak{}lift\_\allowbreak{}of\_\allowbreak{}steps}{Development.\allowbreak{}WeilBound.\allowbreak{}CoreStatements.\allowbreak{}PublishedTheorems.\allowbreak{}CasselmanWallach.\allowbreak{}CasselmanWallachAdmissible}{38--59}

Suppose \(T_0v_0\ne0\). Choose \(p=p_i\) with \(p(Tv_0)>0\) and put \(N=N_i\). Set
\[
Q(z)=\sup_{g\in G} h(g)^{-N}p(\pi(g)z).
\]
This is continuous since \(Q\le b_i\). Submultiplicativity and right \(K\)-invariance of \(h\) give \(Q(\pi(g)z)\le h(g)^NQ(z)\) and \(Q(\pi(k)z)=Q(z)\) for \(k\in K\). To see that \(P=Q\circ T\) is a norm on \(V\), take \(0\ne u\in V\). Its translates span a dense subspace of \(H\) by irreducibility. For \(f\in C_c^\infty(G)\), convolution is continuous \(H\to H^\infty\). Approximating its integral on each translate shows that the smooth closure of their span contains \(\sigma(f)H\). An approximate identity gives density in \(H^\infty\). If \(P(u)=0\), then \(p\circ T\) vanishes on the span and hence at \(v_0\), a contradiction.
\leanref{PositivityTheory.\allowbreak{}CW.\allowbreak{}e1a}{Development.\allowbreak{}WeilBound.\allowbreak{}CoreStatements.\allowbreak{}PublishedTheorems.\allowbreak{}CasselmanWallach.\allowbreak{}ExistsFunctional}{27--80}
\leanref{PositivityTheory.\allowbreak{}CW.\allowbreak{}e1b}{Development.\allowbreak{}WeilBound.\allowbreak{}CoreStatements.\allowbreak{}PublishedTheorems.\allowbreak{}CasselmanWallach.\allowbreak{}ExponentialKSet}{38--68}

Let \(\Omega_K\) be the nonnegative compact Casimir. Its eigenspaces span \(V\). They are finite dimensional by admissibility and finiteness of the \(K\)-types per eigenvalue. For each \(0\ne a\in V\), Hahn--Banach gives \(\ell\) with \(|\ell|\le P\) and \(\ell(a)=P(a)\). Put
\[
B(x,y)=\int_K\overline{\ell(\sigma(k)x)}\,\ell(\sigma(k)y)\,dk
\]
for Haar probability \(dk\). Then \(|B(x,y)|\le P(x)P(y)\), and \(B(a,a)>0\) by positivity near the identity. Invariance makes \(\Omega_K\) symmetric for \(B\). On each eigenspace \(W\) the functionals \(B(x,\cdot)\), for \(x\in W\) and all these forms \(B\), separate points. They span \(W^\vee\) and vanish on the other eigenspaces. Each pairing \(\langle\eta,\cdot\rangle\), \(\eta\in V\), is therefore a finite linear combination of these functionals on \(V\). The bound for \(B\) and density of \(V\) in \(H^\infty\) give \(c_\eta\ge0\) with
\[
|\langle\eta,y\rangle|\le c_\eta Q(Ty)\qquad(y\in H^\infty).
\]
\leanref{PositivityTheory.\allowbreak{}CW.\allowbreak{}e2c}{Development.\allowbreak{}WeilBound.\allowbreak{}CoreStatements.\allowbreak{}PublishedTheorems.\allowbreak{}CasselmanWallach.\allowbreak{}SeminormApproximation}{46--119}
\leanref{PositivityTheory.\allowbreak{}CW.\allowbreak{}c5}{Development.\allowbreak{}WeilBound.\allowbreak{}CoreStatements.\allowbreak{}PublishedTheorems.\allowbreak{}CasselmanWallach.\allowbreak{}Imports}{11--52}

For a finite generating set \(I\) of \(V\), the map \((f_e)\mapsto\sum_{e\in I}\sigma(f_e)e\) from \(\mathcal S(G)^I\) onto \(H^\infty\) is a Fr\'echet quotient \cite[\S6.1 and Theorem~10.6]{bernstein-2014}. The infimum of
\(\sum_{e\in I}\int_G |f_e(g)|h(g)^N\,dg\) over representatives of \(v\) is therefore a continuous seminorm of \(v\). The full Casimir acts by a scalar on \(V\) and hence on \(H^\infty\), so the compact Sobolev norms define the smooth topology \cite[Proposition~3.9]{bernstein-2014}. The infimum is bounded by \(M\|Av\|\) for some \(M\ge0\) and \(A=(1+\Omega_K)^m\), \(m\in\N\). Pairing the quotient formula with \(w\in H\) gives the following matrix coefficient estimate for every \(B_0\ge0\).
\[
\bigl(|\langle w,\sigma(g)e\rangle|\le B_0 h(g)^N\quad(e\in I,\ g\in G)\bigr)
\ \Longrightarrow\
|\langle w,v\rangle|\le MB_0\|Av\|\quad(v\in V).
\]
\leanref{PositivityTheory.\allowbreak{}CWB.\allowbreak{}maxPr}{Development.\allowbreak{}WeilBound.\allowbreak{}CoreStatements.\allowbreak{}PublishedTheorems.\allowbreak{}CasselmanWallach.\allowbreak{}EtaleCohomologyNegativeVanishingDegreeZero}{91--92}

For \(y\in V\) take \(w=Ay\), \(v=A^{-1}y\) and \(c=\sum_{e\in I}c_e\). The inverse exists on \(V\) by nonnegativity of \(\Omega_K\). Unitarity and \(h(g^{-1})=h(g)\) give
\[
|\langle w,\sigma(g)e\rangle|\le c_e Q(\pi(g^{-1})Tw)
\le c\,h(g)^NQ(Tw).
\]
Let \(A_E\) be the corresponding polynomial in the derived operators on \(E\). Since \(T(Ay)=A_E Ty\) and \(A\) is symmetric, the estimate gives
\[
\|y\|^2=|\langle Ay,A^{-1}y\rangle|
\le Mc\,Q(A_E Ty)\,\|y\|.
\]
Hence \(q_0(z)=Mc\,Q(A_Ez)\) is continuous and \(\|y\|\le q_0(Ty)\) on \(V\).
\leanref{PositivityTheory.\allowbreak{}CW.\allowbreak{}e3c}{Development.\allowbreak{}WeilBound.\allowbreak{}CoreStatements.\allowbreak{}PublishedTheorems.\allowbreak{}CasselmanWallach.\allowbreak{}EigenspacesOrthogonal}{23--115}

The Lie algebra preserves \(V\), so the intertwining identity on \(V\) holds for each word \(D\) by composition. The seminorm \(q_S(z)=\sum_{D\in S}q_0(d\pi(D)z)\) is continuous. The bound extends from \(V\) to \(H^\infty\) by density. The empty word gives injectivity. If \(Tw_j\) converges in \(E\), the bounds make \((w_j)\) Cauchy in \(H^\infty\). By completeness \(w_j\to w\) in \(H^\infty\), and \(Tw=\lim_j Tw_j\). This proves closedness since \(E\) is metrizable. If \(T_0|_V=0\), density gives the closed image \(\{0\}\).
\end{proof}
\leanref{PositivityTheory.\allowbreak{}CW.\allowbreak{}e3d}{Development.\allowbreak{}WeilBound.\allowbreak{}CoreStatements.\allowbreak{}PublishedTheorems.\allowbreak{}CasselmanWallach.\allowbreak{}E3D}{10--56}
\leanref{GlobalTheory.\allowbreak{}CompanionTheory.\allowbreak{}SmoothClosed.\allowbreak{}isClosed\_\allowbreak{}image\_\allowbreak{}of\_\allowbreak{}sobolev\_\allowbreak{}lower}{Development.\allowbreak{}AutomorphicSpectrum.\allowbreak{}Proofs.\allowbreak{}BlockDecomposition.\allowbreak{}IsClosedImageOfSobolevLower}{26--60}
\leanref{PositivityTheory.\allowbreak{}CW.\allowbreak{}cw\_\allowbreak{}of\_\allowbreak{}steps}{Development.\allowbreak{}WeilBound.\allowbreak{}CoreStatements.\allowbreak{}PublishedTheorems.\allowbreak{}CasselmanWallach.\allowbreak{}CasselmanWallachAdmissible}{61--108}
\section{The quantifiers and the local error}
\begin{definition}[Closed assertions on family classes]\label{B8.39}
For a family predicate \(\mathcal C\) at each \(\varepsilon\), a threshold datum \(\mathcal K\) has a real \(\mathcal K_{\mathrm{pre}}\) and real-valued functions \(\mathcal K_t,\mathcal K_{\mathrm{mid}},\mathcal K_{\mathrm{last}}\) of real lists. Largeness \(\mathrm{Large}(F_m,\mathcal K,C_H)\) means
\begin{align*}
\mathcal K_{\mathrm{pre}}&\le c&&(c\in F_m.c_{\mathrm{pre}}),\\
\mathcal K_{\mathrm{mid}}(F_m.c_{\mathrm{out}})&\le c&&(c\in F_m.c_{\mathrm{mid}}(C_H)),\\
\mathcal K_{\mathrm{last}}(F_m.c_{\mathrm{out}}\mathbin{+\!\!+}F_m.c_{\mathrm{mid}}(C_H))&\le c&&(c\in F_m.c_{\mathrm{last}}).
\end{align*}
Define \(\mathrm{Closed}(\mathcal C,P)\) by
\[
\begin{gathered}
\exists\varepsilon_1>0,\ \forall\,0<\varepsilon\le\varepsilon_1,\
\exists\mathcal K,\ \forall F_m:\mathrm{FamP}(\varepsilon),\ \forall C_H\in\R,\\
\mathcal C(F_m)\land\mathrm{Large}(F_m,\mathcal K,C_H)
\ \Longrightarrow\ P(F_m,C_H).
\end{gathered}
\]
In particular, \(\mathrm{NL}(\mathcal C)=\mathrm{Closed}(\mathcal C,(F_m,C_H)\mapsto\mathrm{LinkN}(F_m))\).
\end{definition}
\leanref{Formalism.\allowbreak{}MainNotions.\allowbreak{}Closed}{Formalism.\allowbreak{}SharpRestriction.\allowbreak{}ParamsHeightRankConstant}{159--161}
\leanref{GlobalTheory.\allowbreak{}Fc.\allowbreak{}S0.\allowbreak{}NL}{Development.\allowbreak{}AbelianVarietyTheory.\allowbreak{}Definitions.\allowbreak{}GlobalTwisting.\allowbreak{}GeometricWeakTiesOfStrong}{36--36}

Closed assertions conjoin by taking the smaller cutoff and the componentwise maximum of the threshold data at each allowed \(\varepsilon\).
\begin{definition}[The local-error estimate]\label{B8.40}
At a local index \(i\), let \(O_i=D.nd.O(i)\) and \(\ell_i=D.nd.l(i)\). For finite-length \(O_i\)-modules put
\[
\lambda_i(M)=\frac{\operatorname{length}_{O_i}M}
{\operatorname{length}_{O_i}(O_i/\ell_i O_i)}.
\]
For \(\mathrm{Bd}\), use this rational quotient with infinite lengths and division by zero both set to zero. At twisted good points, \(\mathrm{KillAt}\) gives finite Selmer length.
Let \(M_i\) be the lattice image over the fraction field and \(V_Z,V_{Z^c}\) the eigenspace sums for \(Z=Z_{\mathrm{good}}(i)\) and its complement. Put
\[
U_i=V_Z\cap(M_i+V_{Z^c}),\qquad L_i=M_i\cap V_Z,\qquad
\eta_j(w)=\sum_{i\in\mathrm{locs}(j,w)}\lambda_i(U_i/L_i)
\quad(1\le j\le m).
\]
\(\mathrm{locs}(j,w)\) indexes layer \(j\) at \(w\).
The condition \(\mathrm{ShaFinOdd}(D)\) requires finite \(\operatorname{shaLen}(D.nd.ldat(\ell),e)\) at every prime \(\ell\), including \(2\), for \(1\le e<t\). Put \(v_\ell(D)=\log\ell\sum_{e=1}^{t-1}D.\mathrm{what}_2.\mathrm{shaLen}(\ell,e)\) with rational lengths at primes, and zero elsewhere. Define
\[
\eta_j(\mathcal L)=
\sum_{\substack{w\in\mathcal L.\mathrm{supp}\\\mu(w)\ne\mu_0}}
\operatorname{lw}(w)e(w)\eta_j(w),\qquad
\operatorname{lw}(w)=\frac{\log N\mathfrak p_w}{[K:\Q]}.
\]
Here \(\mathfrak p_w,e(w)\) are the prime and ramification index at \(w\), \(dZ0_j=\dim_{\C}Z_j(\mu_0)\), and \(N_c=\mathrm{budgetNc}(D.R,x,\varepsilon,|D.\mathrm{host}.T2|)\). The condition \(\mathrm{LinkN}_{\mathrm{at}}(\mathcal L,j,C)\) is
\begin{equation}\label{eq:B8-linkN}
\bigl(v_\ell(D)\bigr)_{\ell\ge0}\ \text{is summable}
\quad\Longrightarrow\quad
\eta_j(\mathcal L)\le dZ0_j
\left[C\{fmbL^2+(\log x)^2N_c\}+\sum_{\ell\ge0}v_\ell(D)\right].
\end{equation}
Write \(\mathrm{LinkN}(F_m)\) if some \(C,x_0\in\R\) give \eqref{eq:B8-linkN} for all \(x\ge x_0\), \(i\in F_m.Ix(x)\) with \(\mathrm{ShaFinOdd}(D_{x,i})\), and \(1\le j\le m\).
\end{definition}
\leanref{Formalism.\allowbreak{}MainNotions.\allowbreak{}LinkN}{Formalism.\allowbreak{}SharpRestriction.\allowbreak{}GlobalTheoremInputs}{45--47}
\leanref{Formalism.\allowbreak{}MainNotions.\allowbreak{}ShaFinOdd}{Formalism.\allowbreak{}SharpRestriction.\allowbreak{}GlobalTheoremInputs}{25--26}
\leanref{Formalism.\allowbreak{}MainNotions.\allowbreak{}vvTerm}{Formalism.\allowbreak{}SharpRestriction.\allowbreak{}GlobalTheoremInputs}{31--32}
\leanref{Formalism.\allowbreak{}MainNotions.\allowbreak{}etaSum}{Formalism.\allowbreak{}SharpRestriction.\allowbreak{}GlobalTheoremInputs}{36--38}
\leanref{Formalism.\allowbreak{}MainNotions.\allowbreak{}LinkN\_\allowbreak{}at}{Formalism.\allowbreak{}SharpRestriction.\allowbreak{}GlobalTheoremInputs}{40--43}
\section{The primes \(2\) and \(3\)}

Let \(\mathscr H_2=\Lean{Hypotheses₂}\) include Fontaine theory, global reciprocity, Bloch--Kato, integral Bhatt--Scholze and local cohomology.

For a representation family \(\rho\), put
\begin{equation}\label{eq:B8-Bd}
\mathrm{Bd}(\rho,i,\pi,B)\ \Longleftrightarrow\
v_{cc}(i,\pi)\le
s_\pi(\rho_i)+B+\mathrm{VV}_D(\ell_i)+v_{\ell_i}(h_E)+\ell_i^{-2}.
\end{equation}
Here \(v_{cc}(i,\pi)=\lambda_i(O_i/\mathrm{ccAt}_i(\pi))\), with \(\mathrm{ccAt}_i(\pi)\) the Hecke ideal vanishing off \(Z(i)\) evaluated at \(\pi\). The Selmer length \(s_\pi(\rho_i)=\lambda_i(\mathrm{Sel}_\pi)\) uses \(\rho_i(\pi)\otimes\chi_{i,1}^{-1}\) on \((\operatorname{Frac}(O_i)/O_i)^r\). Use Bloch--Kato above \(\ell_i\), and away from \(\ell_i\) allow all classes on \(S_i\cup Q_i\) and unramified classes elsewhere. The term \(\mathrm{VV}_D(\ell)\) sums rational Sha lengths at \(\ell\), and \(h_E\) is the class number of \(E\).

Along with the norm-choice facts, \(\mathrm{SetupAlt}\) requires a canonical hub with constants \(\mathrm{rowC}(C_2)\) and admissible parameters. It imposes \(\mathcal L.Sc+\mathcal L.\eta_{\mathrm{sz}}+t+2\le b\), \(L\ge8\) and \(\mathrm{splitBound}\ge60\), with \(\mathrm{KillAt}(X.\rho)\) and \(\mathrm{ShaFinOdd}(D)\).
\begin{theorem}[A family bound at \(2\) and \(3\)]\label{B8.41}
Assume \(\mathscr H_2\). For every \(\varepsilon\in\R\) and \(F_m:\mathrm{FamP}(\varepsilon)\) there are \(C,x_0\in\R\) with the following property. Take \(x\ge x_0\), \(h_x:F_m.x_1<x\), \(a\in F_m.Ix(x)\) and \(n_{\mathrm{ch}}=(F_m.ex(x,h_x,a)).nch\). For every primewise Wach-module system \(HW\), Fontaine datum \(F\) over \(HW\), comparison datum for \(F\), and hub \(X\) and \(C_2\in\R\) satisfying \(\mathrm{SetupAlt}\), one has
\[
\mathrm{Bd}(X.\rho,i,\pi,Cfmb\log x+Cx^CL^2)
\]
for every twisted index \(i\) with \(\ell_i=2\) or \((\ell_i=3\text{ and }3\notin D_{x,a}.\mathrm{host.small})\), and all \(\pi\in Z_{\mathrm{good}}(i)\).

This gives \(\mathrm{BoundAtTwoThreeFamilyFonAlt}(F_m)\) and the linked closed row on \(\mathrm D3\,3\,\mathcal A\) for each member parameter datum.
\end{theorem}
\leanref{UpperTheory.\allowbreak{}N.\allowbreak{}at23\_\allowbreak{}family\_\allowbreak{}alt}{Development.\allowbreak{}AutomorphicArithmetic.\allowbreak{}ThetaCorrespondence.\allowbreak{}ThetaPairLifts.\allowbreak{}ThirdPlaceSeries.\allowbreak{}At23OfAll}{19--20}
\leanref{UpperTheory.\allowbreak{}N.\allowbreak{}boundAtTwoThree\_\allowbreak{}linked\_\allowbreak{}of\_\allowbreak{}recT}{Development.\allowbreak{}AutomorphicArithmetic.\allowbreak{}ThetaCorrespondence.\allowbreak{}ThetaPairLifts.\allowbreak{}ThirdPlaceSeries.\allowbreak{}At23OfAll}{24--26}
\leanref{GlobalTheory.\allowbreak{}NLink.\allowbreak{}BoundAtTwoThree}{Development.\allowbreak{}AbelianVarietyTheory.\allowbreak{}Definitions.\allowbreak{}NonGaloisDescent.\allowbreak{}BoundAtTwoThreeFamily}{52--53}
\begin{proof}
Put \(H_B=b+\mathcal L.Sc+\mathcal L.\eta_{\mathrm{sz}}+t+2\), \(s_\ell=\max(1,v_\ell(\mathrm{condN}))\) and \(\lambda_\ell=s_\ell+\log_\ell(2H_B)+1\). Write
\[
\operatorname{cost}_{A,B}(\ell)=AfmH_Bs_\ell+Bf^2t^2n_V\lambda_\ell.
\]
Write \(s_0(i)=D.nd.s0(i)\) for the local level exponent. If \(\ell_i=2\), the tower gives \(4\mid QS\) and the local-level rule gives \(s_0(i)\ge2\). Congruence excludes the theta primes from the small-prime set. If \(\ell_i=3\) and \(3\notin D.\mathrm{host.small}\), then \(3\ne q,d\) gives \(3\nmid\mathrm{condN}\), \(s_3=1\) and ramification indices one above \(3\). The local-level rule gives \(s_0(i)=1\). Since \(3\le\mathrm{splitBound}\), the splitting condition gives \(\bigl(\frac{-d}{3}\bigr)=1\).

For \(v\in Q_i\) the split construction gives \(\zeta_{i,v}\ne1\). Killing and Sha-finiteness give finite lengths in the Sha inequality. The second clause of \(\mathrm{KillAt}\) gives
\[
t\sum_{v\in Q_i}\lambda_i\!\left(O_i/(\zeta_{i,v}-1)\right)\le\ell_i^{-2}.
\]
At both primes \(s_0(i)=s_{\ell_i}\) and \(h_0\le s_0(i)+\lfloor\log_{\ell_i} t\rfloor\) for \(t\ge1\). The heights \(\lfloor H_B\rfloor,\lfloor2b\rfloor+1,\lfloor2b\rfloor\) are each at most \(4H_B\). With \(c=1000\) the cyclotomic estimates give
\[
\begin{gathered}
e'\le cH_Bs_\ell,\quad
c_0(u)\le cf n_V\lambda_\ell,\quad h_{0,W}(u)\le crf n_V\lambda_\ell
\quad(u\text{ in the column set}),\\
c_B\le ctf n_V\lambda_\ell,\qquad h_0\le c\lambda_\ell .
\end{gathered}
\]
The split estimate gives \(\mathrm{Bd}(\rho,i,\pi,\operatorname{cost}_{A,B}(\ell_i))\) with \(A=6c=6000\), \(B=30c=30000\).
\leanref{UpperTheory.\allowbreak{}N.\allowbreak{}onAt\_\allowbreak{}eq\_\allowbreak{}two\_\allowbreak{}bounds\_\allowbreak{}of\_\allowbreak{}setup}{Development.\allowbreak{}AutomorphicArithmetic.\allowbreak{}ThetaCorrespondence.\allowbreak{}ThetaPairLifts.\allowbreak{}ThirdPlaceSeries.\allowbreak{}NRowTwo}{11--43}
\leanref{UpperTheory.\allowbreak{}N.\allowbreak{}rowThreeAt\_\allowbreak{}of\_\allowbreak{}setup\_\allowbreak{}alt}{Development.\allowbreak{}AutomorphicArithmetic.\allowbreak{}ThetaCorrespondence.\allowbreak{}ThetaPairLifts.\allowbreak{}ThirdPlaceSeries.\allowbreak{}NRowThree}{12--50}

Eventually \(f,n,t\le x^{C_{\mathrm{par}}}\) for some \(C_{\mathrm{par}}\ge1\). The conductor satisfies \(v_\ell(\mathrm{condN})=v_\ell(QS)\) at \(2\) because \(2\) is small, and at \(3\) since \(3\ne q,d\). Since \(2^{v_\ell(QS)}\le\ell^{v_\ell(QS)}\le QS\le n^2\le x^{2C_{\mathrm{par}}}\) and \(\log 2>2/3\), we have \(s_\ell\le3C_{\mathrm{par}}\log x\) for \(\log x\ge1\). The bound \(H_B\le6b\) gives \(\log_\ell(2H_B)\le2L+6\) and \(\lambda_\ell\le(3C_{\mathrm{par}}+9)xL\). With \(n_V\le2x^{C_{\mathrm{par}}}\), the costs are at most
\[
18AC_{\mathrm{par}}fmb\log x,\qquad
2B(3C_{\mathrm{par}}+9)x^{5C_{\mathrm{par}}+1}L .
\]
Since \(L\ge1\), their sum is bounded by \(Cfmb\log x+Cx^CL^2\) for \(C=18AC_{\mathrm{par}}+2B(3C_{\mathrm{par}}+9)+5C_{\mathrm{par}}+1\).
\leanref{UpperTheory.\allowbreak{}N.\allowbreak{}Num.\allowbreak{}cost\_\allowbreak{}le'}{Development.\allowbreak{}AutomorphicArithmetic.\allowbreak{}ResidueCharacters.\allowbreak{}NumCost}{55--129}

For \(x>1\) use the larger threshold and the larger nonnegative constant. In Definition~\ref{B8.39} take cutoff \(1\) and zero threshold data.
\leanref{UpperTheory.\allowbreak{}N.\allowbreak{}boundAtTwoThree\_\allowbreak{}of\_\allowbreak{}local\_\allowbreak{}bounds\_\allowbreak{}at\_\allowbreak{}two\_\allowbreak{}three}{Development.\allowbreak{}AutomorphicArithmetic.\allowbreak{}ThetaCorrespondence.\allowbreak{}ThetaPairLifts.\allowbreak{}ThirdPlaceSeries.\allowbreak{}FamilyBoundAtTwoThreeOfLocalBounds}{37--54}
\end{proof}
\section{Cutting to the good part}

For a finite discrete Galois module \(M\) over \(\Z/\ell^r\), with \(\ell\) prime and \(r>0\), cohomology at a nonarchimedean completion \(F_v\) is finite through degree \(2\) and vanishes above it. Cup products with \(M^\vee(1)\) are perfect, with \(H^2(\mu_{\ell^r})\simeq\Z/\ell^r\) over the coefficients. The local Euler characteristic gives
\[
|H^1(F_v,M)|=|H^0(F_v,M)|\,|H^2(F_v,M)|
\begin{cases}|M|^{[F_v:\Q_\ell]},&v\mid\ell,\\1,&v\nmid\ell.\end{cases}
\]
This is the decomposition-group form of \cite[Chapter VII, \S\S2--3]{neukirch-2008}.
\leanref{ReciprocityTheory.\allowbreak{}Gal.\allowbreak{}LocalFiniteness}{ReciprocityTheory.\allowbreak{}Hypotheses.\allowbreak{}GalCoh}{195--199}
\leanref{ReciprocityTheory.\allowbreak{}Gal.\allowbreak{}LocalTateDuality}{ReciprocityTheory.\allowbreak{}Hypotheses.\allowbreak{}GalCoh}{201--205}
\leanref{ReciprocityTheory.\allowbreak{}Gal.\allowbreak{}LocalEulerChar}{ReciprocityTheory.\allowbreak{}Hypotheses.\allowbreak{}GalCoh}{213--217}

Let \(F\) be a number field and \(G_F=\Gal(\overline F/F)\). Fix \(g\in G_F\) and an open normal subgroup \(U\subset G_F\). The primes of \(F\) unramified in the extension fixed by \(U\) whose Frobenius class is that of \(gU\) in \(G_F/U\) have positive Chebotarev density \cite[\S1A]{thorner-2019}. Primes of \(F\) of absolute residue degree greater than one have density zero, and the primes ramified over \(\Q\) form a finite set. Thus outside any prescribed finite set there is a prime of residue degree and ramification index one over \(\Q\) with a prolongation whose inertia lies in \(U\) and whose Frobenius maps to \(gU\).

A cut keeps selected indices with nonempty spectral sets \(X_i\). A good cut keeps exactly the indices with \(Z_{\mathrm{good}}(i)\ne\varnothing\) and requires
\[
Z_{\mathrm{good}}(i)\subset X_i,\qquad Z(i)\cap X_i\subset Z_{\mathrm{good}}(i).
\]
At retained \(i\), require an integral Hecke-commuting \(\alpha_i\) acting by \(M_0=D.\mathrm{what}_2.\mathrm{restRowsFiniteGroupCard}\) on good eigenspaces and by zero outside \(X_i\). Require \(\Lean{Nc}\) for automorphic lifts into \(X_i\) at each \(w\mid g_k\), \(k=0,1\). In dimension \(d_w\), each wild inertia element has characteristic polynomial \((X-1)^{d_w}\), and monodromy is zero. Each tame inertia generator has \(z^{Nw-1}=1\) for every characteristic root \(z\) of multiplicity at least \(\operatorname{thr}(w,d_w)\), with residue cardinality \(Nw\). The threshold is \(2\) if the prime \(p\) below \(w\) splits completely in \(E\) and \((d_w=1\text{ or }16d_w^2<p)\), and \(1\) otherwise.

The ties \(\mathrm{Tie4}(D)\) require paired-character compatibility and data with depth zero and eigenspan, polynomial compatibility and the local series condition. The Galois preconditions \(\mathrm{GP}(D,n_{\mathrm{ch}})\) include Galois rows, admissible decomposition, the first provision, \(Z\)-conclusion, residual irreducibility, tame eigenvalues, weight and block dictionaries, unramified and depth-zero conditions, base change, weight use, localisation uniqueness, coefficient largeness, norm-choice and twist ties. Each finite-type row must satisfy the good-label condition for the character and its inverse. Each parameter list on a good row must admit a triple satisfying the prescribed Bloch--Kato condition.
\leanref{UpperTheory.\allowbreak{}GalC.\allowbreak{}GaloisPreconditions}{Development.\allowbreak{}AutomorphicArithmetic.\allowbreak{}ThetaCorrespondence.\allowbreak{}ThetaPairLifts.\allowbreak{}ThirdPlaceSeries.\allowbreak{}GaloisPreconditions}{13--31}
\begin{theorem}[Good cuts and restoration of the family bound]\label{B8.42}
Assume the Chebotarev and decomposition-group cohomology inputs. Under \(\mathrm{GP}(D,n_{\mathrm{ch}})\) and \(\mathrm{Tie4}(D)\), some good cut \(c\) satisfies \(\mathrm{NormChoiceFacts}(D^c,n_{\mathrm{ch}}^c)\) and \(\mathrm{Tie4}(D^c)\) with transported norm choice \(n_{\mathrm{ch}}^c\).

Assume also \(\mathscr H_2\). For a twist-good family \(F_m\), suppose \(\mathrm{GP}(D_{x,i},n_{\mathrm{ch}})\land\mathrm{Tie4}(D_{x,i})\) holds for every \(h_x:F_m.x_1<x\) and \(i\in F_m.Ix(x)\), with \(n_{\mathrm{ch}}=(F_m.ex(x,h_x,i)).nch\). Then \(\mathrm{LinkN}(F_m)\) holds.
\end{theorem}
\leanref{UpperTheory.\allowbreak{}GalC.\allowbreak{}exists\_\allowbreak{}good\_\allowbreak{}cut\_\allowbreak{}facts\_\allowbreak{}tie4}{Development.\allowbreak{}AutomorphicArithmetic.\allowbreak{}ThetaCorrespondence.\allowbreak{}ThetaPairLifts.\allowbreak{}ThirdPlaceSeries.\allowbreak{}CFacts}{58--62}
\leanref{UpperTheory.\allowbreak{}GalC.\allowbreak{}linkN\_\allowbreak{}of\_\allowbreak{}twGood\_\allowbreak{}galoisPreconditions}{Development.\allowbreak{}AutomorphicArithmetic.\allowbreak{}ThetaCorrespondence.\allowbreak{}ThetaPairLifts.\allowbreak{}ThirdPlaceSeries.\allowbreak{}LinkNFam}{37--53}
\begin{proof}
With paired-character datum \(Q\) and \(e_0=D.\mathrm{ntypes}.e0\), take \(X_i\) to consist of points with a lift having a type vector \(F\) with \(e_0F\ne0\). The projector is the identity over good points and zero over \(Z(i)\setminus Z_{\mathrm{good}}(i)\).

Let \(G_{\mathrm{type}}\) be the product of the finite general linear groups at \(g_0,g_1\). For retained \(i\) and good tuple \(p\), choose central \(z_{i,p}\in O_i[G_{\mathrm{type}}]\) mapping to \(M_0\) times the tensor idempotent of \(p\). Then \(z_i=\sum_p z_{i,p}\) maps to \(M_0e_0\), and eigenspan gives the required action of the integral Hecke-commuting \(\alpha_i=Q.zO_i(z_i)\) over the fraction field.
\leanref{UpperTheory.\allowbreak{}GalC.\allowbreak{}exists\_\allowbreak{}unit\_\allowbreak{}one\_\allowbreak{}add\_\allowbreak{}generator}{Development.\allowbreak{}AutomorphicArithmetic.\allowbreak{}Foundations.\allowbreak{}GaloisRepresentations.\allowbreak{}HeckeGaloisCriteria.\allowbreak{}ExistsUnitOneAddGenerator}{17--56}
Fix \(x\in X_i\) and a lift with a type vector \(F\) such that \(e_0F\ne0\). This gives a good tuple and \(\Lean{Nc}\) at the places above both theta primes. The first provision identifies the parameter list of this lift with that of any other lift of \(x\), up to permutation. Hence \(\Lean{Nc}\) holds at those places for every lift of \(x\).
\leanref{UpperTheory.\allowbreak{}GalC.\allowbreak{}exists\_\allowbreak{}good\_\allowbreak{}alt}{Development.\allowbreak{}AutomorphicArithmetic.\allowbreak{}ThetaCorrespondence.\allowbreak{}ThetaPairLifts.\allowbreak{}ThirdPlaceSeries.\allowbreak{}ExistsGood}{15--54}

For \(w\mid\ell_i\), use the cut for \(\Lean{Nc}\) at \(g_0,g_1\) and unramified parameters otherwise. Local--global compatibility gives the level-zero weights as de Rham jumps. Under \(\mathrm{AdmPar}\) the weight and block dictionaries give \(\mathrm{AdmJAlt}\). The residual row gives \(\Lean{RbEx}\). Admissibility and \(\Lean{CutRec}\) give local--global compatibility. For \(w\in S_i\), write \(\gamma_{i,w}=D.nd.\mathrm{gam0}(i,w)\) and take a residual representation \(\bar\rho_i\) as in \(\Lean{RbEx}\). At a theta prime, admissibility excludes the reduced character value \(\bar\chi_{i,1}(\gamma_{i,w})\) from the spectrum of \(\bar\rho_i(\gamma_{i,w})\). At the auxiliary prime the exclusion follows from \(\Lean{EigPd}\), obtained under admissibility from compatibility and \(\Lean{RbEx}\). The exclusion gives the hypothesis of the local cohomology bound \(\Lean{HOneTwo}\). This proves the fifth provision. Admissibility and \(\Lean{RbEx}\) give the sixth. Both provisions with the \(Z\)-conclusion give the second. Reconstruct the first on the cut. Chebotarev gives the coefficient duties from largeness, the first provision, jumps and residual irreducibility. The four ties restrict.
\leanref{UpperTheory.\allowbreak{}GalC.\allowbreak{}facts\_\allowbreak{}of\_\allowbreak{}galois\_\allowbreak{}preconditions\_\allowbreak{}of\_\allowbreak{}cutRec}{Development.\allowbreak{}AutomorphicArithmetic.\allowbreak{}ThetaCorrespondence.\allowbreak{}ThetaPairLifts.\allowbreak{}ThirdPlaceSeries.\allowbreak{}CFacts}{45--57}
\leanref{UpperTheory.\allowbreak{}GalC.\allowbreak{}compatibilityTies\_\allowbreak{}eigenRestriction}{Development.\allowbreak{}AutomorphicArithmetic.\allowbreak{}Foundations.\allowbreak{}GaloisRepresentations.\allowbreak{}HeckeGaloisCriteria.\allowbreak{}CompatibilityTiesEigenRestriction}{86--88}

Use the good cut at each member with \(x>F_m.x_1\) and the full cut otherwise. The family \(F_m^c\) is twist-good. Its norm facts and ties give common paired-character constants. Apply Theorem~\ref{B8.41}. Past its thresholds, at each point with \(\mathrm{ShaFinOdd}\) choose the exceptional primes and a hub for them.

Combine the \(2,3\)-bound and the remaining small-prime bound. Congruence treats large primes away from the exceptional set and theta primes. Paired-character and ramified-place estimates give type and ramification, and the commuting-endomorphism condition gives passage. With the midrange, rank, multiplicity and numerical bounds this gives \(\mathrm{ConclF}(F_m^c)\). Pass to \(\mathrm{AdmF}(F_m^c)\) and sum with the place weights for \(\mathrm{LinkN}(F_m^c)\).
\leanref{GlobalTheory.\allowbreak{}NLinkClose.\allowbreak{}main\_\allowbreak{}alt\_\allowbreak{}alt}{Development.\allowbreak{}AutomorphicArithmetic.\allowbreak{}ThetaCorrespondence.\allowbreak{}ThetaPairLifts.\allowbreak{}ThirdPlaceSeries.\allowbreak{}ImportsFurtherTail}{22--109}
\leanref{GlobalTheory.\allowbreak{}LinkFam.\allowbreak{}main}{Development.\allowbreak{}AbelianVarietyTheory.\allowbreak{}Proofs.\allowbreak{}InertiaEigenvalues.\allowbreak{}LinkNAlongFamilyOfLocalBounds}{11--17}

Project \(U_i/L_i\) to the cut quotient. For a kernel representative \(u\) choose an integral lift \(m_0\) of its cut image. Since \(u-m_0\) lies in the discarded eigenspan, \(M_0u=\alpha_i m_0\) is integral and belongs to \(L_i\).

The kernel preimage \(P\subset U_i\) is finitely generated and torsion-free over the discrete valuation ring \(O_i\), hence free. Its rank is bounded by the good eigenspace dimension, at most \(\sum_{\pi\in Z_{\mathrm{good}}(i)}\mathrm{mult}(i,\pi)\). The basis maps to kernel generators killed by \(M_0\). Each contributes at most \(\lambda_i(O_i/M_0O_i)\le v_{\ell_i}(M_0)\) to the normalized length. The image length is at most the cut quotient length. Discarded indices have zero error, so
\begin{equation}\label{eq:B8-eta-cut}
\eta_j(w)\le\eta_j^c(w)+
\left(\sum_{i\in\mathrm{locs}(j,w)}
\sum_{\pi\in Z_{\mathrm{good}}(i)}\mathrm{mult}(i,\pi)\right)
v_{\ell(w)}(M_0).
\end{equation}
\(\Lean{goodPointMultiplicityLeRank\_\allowbreak{}proof}\) bounds the sum by the zero-space rank, equal to \(dZ0_j\) at twisted places by the rank row.
\leanref{UpperTheory.\allowbreak{}GalC.\allowbreak{}Alg.\allowbreak{}length\_\allowbreak{}le}{Development.\allowbreak{}AutomorphicArithmetic.\allowbreak{}Foundations.\allowbreak{}GaloisRepresentations.\allowbreak{}HeckeGaloisCriteria.\allowbreak{}LengthLeOfSpan}{186--292}
\leanref{UpperTheory.\allowbreak{}GalC.\allowbreak{}etaJ\_\allowbreak{}le}{Development.\allowbreak{}AutomorphicArithmetic.\allowbreak{}Foundations.\allowbreak{}GaloisRepresentations.\allowbreak{}HeckeGaloisCriteria.\allowbreak{}IsHeckeGaloisIff}{73--115}

Above \(\ell\), the twisted weights sum to at most \(\log\ell\). The weighted valuations thus sum to at most \(\sum_{\ell\mid M_0}v_\ell(M_0)\log\ell=\log M_0\), whence
\begin{equation}\label{eq:B8-cut-sum}
\eta_j(\mathcal L)\le\eta_j^c(\mathcal L)+dZ0_j\log M_0.
\end{equation}
\leanref{UpperTheory.\allowbreak{}GalC.\allowbreak{}sum\_\allowbreak{}val\_\allowbreak{}le}{Development.\allowbreak{}AutomorphicArithmetic.\allowbreak{}Foundations.\allowbreak{}GaloisRepresentations.\allowbreak{}HeckeGaloisCriteria.\allowbreak{}EvPowLeExp}{16--45}
\leanref{UpperTheory.\allowbreak{}GalC.\allowbreak{}etaSum\_\allowbreak{}le\_\allowbreak{}add}{Development.\allowbreak{}AutomorphicArithmetic.\allowbreak{}Foundations.\allowbreak{}GaloisRepresentations.\allowbreak{}HeckeGaloisCriteria.\allowbreak{}EvPowLeExp}{48--57}

Put \(B_x=fmbL^2+(\log x)^2N_c\). Admissibility gives \(\log M_0\le C_1x^{C_1}L\) eventually for some \(C_1\ge0\). At \((0,i_0)\) with \(i_0\in F_m.Ix(0)\) the scale has \(\varepsilon>0\). Eventually \(x^{C_1}\le b\) and \(L\ge1\), so \(x^{C_1}L\le bL^2\le B_x\) for \(f,m\ge1\). The cut preserves \(v_\ell(D),dZ0_j,B_x\). With \(C\) from \(\mathrm{LinkN}(F_m^c)\), finiteness and summability give
\[
\eta_j(\mathcal L)\le dZ0_j\left[(C+C_1)B_x+\sum_{\ell\ge0}v_\ell(D)\right].
\]
Take the maximum of the cut-family, admissibility and exponential-domination thresholds and \(F_m.x_1+1\).
\leanref{UpperTheory.\allowbreak{}GalC.\allowbreak{}linkN\_\allowbreak{}of\_\allowbreak{}gd}{Development.\allowbreak{}AutomorphicArithmetic.\allowbreak{}Foundations.\allowbreak{}GaloisRepresentations.\allowbreak{}HeckeGaloisCriteria.\allowbreak{}LinkNOfGd}{38--93}
\end{proof}
\section{The closed local condition in the witness}
\begin{theorem}[The closed \(\mathrm{LinkN}\) condition]\label{B8.43}
Assume \(\mathscr H_2\), \(\Lean{PublishedInputsThirdLayerAlt}\), \(\Lean{GeometricGaloisRepCompatibility}\), \Lean{GlobalTheory.\allowbreak{}HeightCounting.\allowbreak{}Up.\allowbreak{}SerN}, depth-zero support and the three decomposition-group cohomology inputs. For a local Langlands witness \(h_L\) and \(c_{\mathrm{Dep}}\in\R\), put \(p=\mathrm{parametersOfDependence}(h_L,c_{\mathrm{Dep}})\). The linked closed \(2,3\)-condition for \(p\) gives the linked closed \(\mathrm{LinkN}\) condition on \(\mathrm D3\,3\,\mathcal A\).
\end{theorem}
\leanref{GlobalTheory.\allowbreak{}HeightCounting.\allowbreak{}Up.\allowbreak{}SerN}{Development.\allowbreak{}AbelianVarietyTheory.\allowbreak{}Definitions.\allowbreak{}AutomorphicMembership.\allowbreak{}CenUsedOfCenReal}{56--62}
\leanref{UpperTheory.\allowbreak{}GalC.\allowbreak{}final\_\allowbreak{}alt\_\allowbreak{}alt}{Development.\allowbreak{}AutomorphicArithmetic.\allowbreak{}ThetaCorrespondence.\allowbreak{}ThetaPairLifts.\allowbreak{}ThirdPlaceSeries.\allowbreak{}CFinal}{19--22}
\leanref{UpperTheory.\allowbreak{}GalC.\allowbreak{}final\_\allowbreak{}of\_\allowbreak{}alt}{Development.\allowbreak{}AutomorphicArithmetic.\allowbreak{}ThetaCorrespondence.\allowbreak{}ThetaPairLifts.\allowbreak{}ThirdPlaceSeries.\allowbreak{}LinkNFam}{75--86}
\leanref{UpperTheory.\allowbreak{}GalC.\allowbreak{}linkN\_\allowbreak{}of\_\allowbreak{}boundAtTwoThree\_\allowbreak{}slot}{Development.\allowbreak{}AutomorphicArithmetic.\allowbreak{}ThetaCorrespondence.\allowbreak{}ThetaPairLifts.\allowbreak{}ThirdPlaceSeries.\allowbreak{}LinkNFam}{54--63}
\begin{proof}
At a member, the third-layer, geometric, series, depth-zero and local cohomology inputs give \(\mathrm{GP}(D,O.ex.nch)\). Height-cardinality with the third-layer and depth-zero inputs gives \(\mathrm{Tie4}\). Transport these rows along the extension with the prescribed list and geometric closure. Conjoin them with the closed \(2,3\)-bound and keep its cutoff \(\varepsilon_1\) and data \(\mathcal K_\varepsilon\). Apply Theorem~\ref{B8.42} on the common restriction to each large family in \(\mathrm D3\,3\,\mathcal A\).
\leanref{UpperTheory.\allowbreak{}GalS.\allowbreak{}galois\_\allowbreak{}preconditions\_\allowbreak{}at\_\allowbreak{}mem}{Development.\allowbreak{}AutomorphicArithmetic.\allowbreak{}ThetaCorrespondence.\allowbreak{}ThetaPairLifts.\allowbreak{}ThirdPlaceSeries.\allowbreak{}PreAtMem}{86--103}

In \ref{B8.03} and \ref{B8.04}, the published and geometric inputs of \(P\) supply the series hypotheses. The strengthened extension gives \(\mathrm{NL}(\mathrm D3\,3\,\mathcal A)\) in the common witness's \(\Lean{nRest}\) field.
\leanref{GlobalTheory.\allowbreak{}Fc.\allowbreak{}S0.\allowbreak{}StageU.\allowbreak{}hypotheses\_\allowbreak{}nonempty\_\allowbreak{}of\_\allowbreak{}unramified\_\allowbreak{}exponents\_\allowbreak{}alt}{Development.\allowbreak{}AutomorphicArithmetic.\allowbreak{}ThetaCorrespondence.\allowbreak{}ThetaPairLifts.\allowbreak{}ThirdPlaceSeries.\allowbreak{}HypothesesNonemptyOfUnramifiedExponents}{41--61}
\leanref{GlobalTheory.\allowbreak{}Fc.\allowbreak{}AllAtMemberStage0.\allowbreak{}allRows\_\allowbreak{}of\_\allowbreak{}open0\_\allowbreak{}alt}{Development.\allowbreak{}AutomorphicArithmetic.\allowbreak{}ThetaCorrespondence.\allowbreak{}ThetaPairLifts.\allowbreak{}ThirdPlaceSeries.\allowbreak{}LinkROfIsLinked}{50--86}
\leanref{GlobalTheory.\allowbreak{}Fc.\allowbreak{}nonempty\_\allowbreak{}hyps\_\allowbreak{}of\_\allowbreak{}allRows\_\allowbreak{}alt}{Development.\allowbreak{}AutomorphicArithmetic.\allowbreak{}ThetaCorrespondence.\allowbreak{}ThetaPairLifts.\allowbreak{}ThirdPlaceSeries.\allowbreak{}NonemptyHypsOfAllRows}{19--101}
\end{proof}


% ---- p8c5 ----
\chapter{Primes in short intervals, masses and counts}\label{ch:p8c5}

Let \(E=D.\mathrm{host}.E\), \(F=E^+\), \(f=[F:\Q]\). Write \(m\) for chain length, \(k\) for starting rank, \(r\) for spectral rank and \(T_u\) for the level prime at word \(u\). Set \(L=x^\varepsilon,\ b=e^L\) for \(0<\varepsilon<1/52\).

\section{Uniform estimates for Dirichlet \(L\)-functions}

For primitive \(\chi\bmod N\), \(N\ge1\), set
\[
 \psi_0(\chi,u)=\sum_{1\le n\le u}\Lambda(n)\chi(n)-{\bf1}_{\chi=1}u .
\]
Among distinct zeros in \(0<\Re\rho<1\), let \(Z_\chi(T)\) consist of those with \(|\Im\rho|<T\) and \(Z_\chi[t]\) of those with \(|\Im\rho-t|\le1\). Their multiplicities are \(m_\chi(\rho)\).

\begin{definition}[The uniform analytic input]\label{B8.44}
The assertion \(\Lean{LitThin}\) requires the positive constants of \ref{B8.45} and \ref{B8.46} to work for every datum \(D\).
\end{definition}
\leanref{Formalism.\allowbreak{}MainNotions.\allowbreak{}LitThin}{Formalism.\allowbreak{}SharpRestriction.\allowbreak{}GlobalSchurBounds.\allowbreak{}MainSup}{88--90}
\leanref{GlobalTheory.\allowbreak{}ThinScaleLiterature.\allowbreak{}litThin}{Development.\allowbreak{}AbelianVarietyTheory.\allowbreak{}Proofs.\allowbreak{}GlobalInterpolation.\allowbreak{}LitThin}{8--9}

\begin{theorem}[Density near \(1\) and logarithmic zero regions]\label{B8.45}
There are \(C_1,c>0\) such that for \(Q\ge3\) any finite set \(\mathcal T\) of triples \((N,\chi,\rho)\) with \(1\le N\le Q\), \(\chi\) primitive modulo \(N\), and
\[
 L(\rho,\chi)=0,\qquad \Re\rho\ge399/400,\qquad |\Im\rho|\le Q
\]
satisfies \(\sum_{(N,\chi,\rho)\in\mathcal T}m_\chi(\rho)\le C_1Q^{1/4}\).

If \(\chi\) is primitive modulo \(N\ge3\) and \(L(\rho,\chi)=0\) with
\[
 \Re\rho\ge1-\frac{c}{\log(\max(N,N|\Im\rho|))},
\]
then \(\chi\) is real and \(\rho\) is the unique zero there. The zero is real and simple. Every zero of \(\zeta\) in the open critical strip satisfies
\[
 \Re\rho\le1-\frac{c}{\log(|\Im\rho|+3)} .
\]
\end{theorem}
\leanref{GlobalTheory.\allowbreak{}ZerosNearOne.\allowbreak{}density\_\allowbreak{}and\_\allowbreak{}regions}{Development.\allowbreak{}AbelianVarietyTheory.\allowbreak{}Proofs.\allowbreak{}GlobalInterpolation.\allowbreak{}DensityAndRegions}{10--16}
\leanref{Formalism.\allowbreak{}MainNotions.\allowbreak{}ShortIntervalPrimes.\allowbreak{}adm6}{Formalism.\allowbreak{}SharpRestriction.\allowbreak{}GlobalSchurBounds.\allowbreak{}ShortIntervalPrimes}{128--139}

The fixed-strip estimate follows from Montgomery's density bound \cite[Ch.~12]{montgomery-1971}. There are \(O(Q^{1/8})\) occupied cells \((N,\chi,\lfloor\Im\rho\rfloor)\), each of multiplicity \(O(\log Q)\). The logarithmic regions use Davenport \cite[\S\S13--14]{davenport-1980} with a smaller absolute constant.

\begin{theorem}[Explicit formula and unit-window zero count]\label{B8.46}
There are \(C_7,C_{11}>0\) such that for every primitive \(\chi\) modulo \(N\ge1\) and \(2\le T\le u\), the set \(Z_\chi(T)\) is finite and
\[
 \left|\psi_0(\chi,u)+\sum_{\rho\in Z_\chi(T)}
        m_\chi(\rho)\frac{u^\rho}{\rho}\right|
 \le C_7\left(\frac{u\log^2(NuT)}{T}+\log(Nu)
       +\sum_{\rho\in Z_\chi(1)}\frac{m_\chi(\rho)}{|\rho|}\right).
 \tag{8.46.1}
\]
For every \(t\in\R\), the set \(Z_\chi[t]\) is finite and
\[
 \sum_{\rho\in Z_\chi[t]}m_\chi(\rho)
 \le C_{11}\log\bigl(N(|t|+2)\bigr).
 \tag{8.46.2}
\]
\end{theorem}
\leanref{GlobalTheory.\allowbreak{}ZeroBounds.\allowbreak{}explicit\_\allowbreak{}formula\_\allowbreak{}and\_\allowbreak{}window\_\allowbreak{}count}{Development.\allowbreak{}AbelianVarietyTheory.\allowbreak{}Proofs.\allowbreak{}GlobalInterpolation.\allowbreak{}ExplicitFormulaAndWindowCount}{9--13}
\leanref{Formalism.\allowbreak{}MainNotions.\allowbreak{}ShortIntervalPrimes.\allowbreak{}adm7}{Formalism.\allowbreak{}SharpRestriction.\allowbreak{}GlobalSchurBounds.\allowbreak{}ShortIntervalPrimes}{141--147}
\leanref{Formalism.\allowbreak{}MainNotions.\allowbreak{}ShortIntervalPrimes.\allowbreak{}adm11}{Formalism.\allowbreak{}SharpRestriction.\allowbreak{}GlobalSchurBounds.\allowbreak{}ShortIntervalPrimes}{158--161}

In Davenport's formula \cite[\S19, (9)]{davenport-1980}, bound the signed low-zero correction by the sum of its absolute terms. The conductor-one formula is in \S17 and the unit-window count in \S\S15--16.

Put \(N_E=D.\mathrm{host}.QS\,q\,d\) and \(\lambda_E=\max(1,\log N_E)\). A primitive character \(\chi\) belongs to \(D\), as in \(\Lean{IsCharED}\), if some character \(\omega:G\to\C^\times\) satisfies
\[
 \chi(a)=\omega(\operatorname{Art}_D(a))
 \quad(a\in\N,\ \gcd(a,D.\mathrm{host}.\mathrm{condN})=1).
 \tag{8.47.1}
\]
Set \(X_1(Y)=e^{Y^{\varepsilon/5}}\), \(X_2(Y)=e^{Y^{\varepsilon/5}(\log Y)^2}\), and let \(s_D\) be the odd part of \(D.\mathrm{host}.QS\). For a real primitive \(\chi\), write \(\mathrm{Good}(Q,\chi)\) for \(N\le Q\) and nonvanishing at \(\Re s>399/400,\ |\Im s|\le Q\). At each sufficiently large family point, \(\Lean{BoxesF}\) supplies \(Y\in[x^{1/J},x]\), with \(J\) the family exponent, for which every real \(\chi\) belonging to \(D\) satisfies
\[
 \begin{split}
 d\mid N&\ \Longrightarrow\ \mathrm{Good}(8s_D X_1(Y)X_2(Y),\chi),\\
 q\mid N,\ d\nmid N&\ \Longrightarrow\ \mathrm{Good}(8s_D X_1(Y),\chi).
 \end{split}
 \tag{8.47.2}
\]
At this \(Y\), the six \(\Lean{Boxes\_\allowbreak{}at}\) bounds use \(8X_1\) for the \(q\)-character and \(8X_1X_2\) for the \(d\)- and \(qd\)-characters. The same bounds apply to their \(\chi_8\)-twists. For some \(Y\) in the same interval, \(\Lean{QDRowsF}\) gives
\[
 \begin{gathered}
 q\equiv3\pmod8,\quad X_1(Y)^{1/2}<q\le X_1(Y),\\
 d\equiv7\pmod{16},\quad X_2(Y)^{1/2}<d\le X_2(Y),\\
 \left(\frac{-d}{q}\right)=1,\qquad
 \left(\frac{-d}{p}\right)=1\quad(p\le Y^{\varepsilon/5},\ p\text{ prime}),\\
 \log N_E\le2Y^{\varepsilon/5}(\log Y)^2 .
 \end{gathered}
 \tag{8.47.3}
\]
The symbols in (8.47.3) are Kronecker symbols.

\begin{proposition}[The character-sum estimate on a family]\label{B8.47}
Assume \(\Lean{LitThin}\). Let \(F_m\) satisfy \(\Lean{BoxesF}\) and \(\Lean{QDRowsF}\). There are \(c_6>0\) and \(x_0\in\R\) such that for every \(x\ge x_0\), \(i\in F_m.I_x\), and character \(\chi\) belonging to \(D=F_m.D(x,i)\),
\[
 |g_\chi(y)|\le y e^{-c_6x^{\varepsilon/8}}
 \qquad\left(y\ge e^{x^{\varepsilon/3}}/3\right).
 \tag{8.47.4}
\]
The closed form on any family class takes the additional premise
\(\exists C_5>0,\ \Lean{ChosenQD}(F_m,C_5)\).
\end{proposition}
\leanref{Formalism.\allowbreak{}MainNotions.\allowbreak{}chosenThin\_\allowbreak{}of\_\allowbreak{}boxesN}{Formalism.\allowbreak{}SharpRestriction.\allowbreak{}GlobalSchurBounds.\allowbreak{}MainSup}{105--114}
\leanref{Formalism.\allowbreak{}MainNotions.\allowbreak{}chosenThin\_\allowbreak{}closed}{Formalism.\allowbreak{}SharpRestriction.\allowbreak{}GlobalSchurBounds.\allowbreak{}MainSup}{124--127}

\begin{proof}
The conductor divides \(N_E=D.\mathrm{host}.\mathrm{condN}\). If \(\chi\) is real and neither \(q\) nor \(d\) divides \(N\), primitivity gives \(N\mid8s_D\). Since \(s_D\le F_m.Co^2\) for \(x>F_m.x_1\), the nonprincipal characters in this case have bounded conductors. Nonvanishing at \(1\) gives them a common zero-free neighborhood. For the other real characters (8.47.2) excludes the possible exceptional zero if the region constant is less than \(\log(3)/400\). By \ref{B8.45}, including its zeta region, there is a uniform \(c_0>0\) with
\[
 \Re\rho\le1-\frac{c_0}{\lambda_E+\log(|\Im\rho|+3)}
 \quad(0<\Re\rho<1).
 \tag{8.47.5}
\]
Apply (8.47.5) to the zero \(1-\rho\) of the inverse character, which also belongs to \(D\). This gives
\[
 |\Im\rho|\le2\ \Longrightarrow\
 \Re\rho\ge\kappa/\lambda_E,\qquad \kappa=c_0/3.
 \tag{8.47.6}
\]

For \(T\ge1\), group the zeros by \(j=\lfloor|\Im\rho|\rfloor\). By (8.46.2) each group has multiplicity at most \(2C_{11}(\lambda_E+\log(\lfloor T\rfloor+2))\). By (8.47.6) and \(|\rho|\ge|\Im\rho|\), the reciprocal of \(|\rho|\) is at most \({\bf1}_{j\le2}\lambda_E/\kappa+2/(j+1)\). The harmonic bound \(2+\log T\) yields
\[
 \sum_{\rho\in Z_\chi(T)}\frac{m_\chi(\rho)}{|\rho|}
 \le C_2(\lambda_E+\log T)^2,\qquad C_2=6C_{11}(3/\kappa+4).
 \tag{8.47.7}
\]
\leanref{Formalism.\allowbreak{}MainNotions.\allowbreak{}ShortIntervalPrimes.\allowbreak{}zeroSumE\_\allowbreak{}of}{Formalism.\allowbreak{}SharpRestriction.\allowbreak{}GlobalSchurBounds.\allowbreak{}FirstClaimLemmaTwoFourOfLemmaTwoThree}{2355--2377}

For \(v=\log u\ge1\), set \(T=e^{\sqrt v}\) and \(c_1=\min(\kappa,1)\). Put \(H=e^{-c_1v/(\lambda_E+\sqrt v)}\). We have \(T^{-1}\le H\) and \(uH\ge1\). On \(Z_\chi(T)\), \(\log(|\Im\rho|+3)\le\sqrt v+2\), so (8.47.5) gives \(u^{\Re\rho}\le uH\). Use (8.47.7) in (8.46.1) and \(\log(NuT)\le2(v+\lambda_E)\) to obtain
\[
 |\psi_0(\chi,u)|\le C_3u(v+\lambda_E)^2
 e^{-c_1v/(\lambda_E+\sqrt v)},\qquad C_3=C_7(5+C_2)+C_2 .
 \tag{8.47.8}
\]
\leanref{Formalism.\allowbreak{}MainNotions.\allowbreak{}ShortIntervalPrimes.\allowbreak{}psi\_\allowbreak{}thin\_\allowbreak{}core}{Formalism.\allowbreak{}SharpRestriction.\allowbreak{}GlobalSchurBounds.\allowbreak{}FirstClaimLemmaTwoFourOfLemmaTwoThree}{2028--2148}

The fixed-modulus estimate \(\psi_0(\chi,u)\ll_\chi u/(\log u)^2\) makes \(\psi_0/u^2\) integrable on \([1,\infty)\). Its integral is \(-L'(1,\chi)/L(1,\chi)\) for \(\chi\ne1\) and \(-1-\gamma\) for conductor one. Partial summation gives the uncentred sum \(y\int_1^y(\psi_0(\chi,u)+{\bf1}_{\chi=1}u)\,du/u^2\). Subtracting the centring term, which for \(G_z(y)\) is \(y(\log y-1-\gamma)\), gives
\[
 g_\chi(y)=-y\int_y^\infty\psi_0(\chi,u)\,\frac{du}{u^2}.
 \tag{8.47.9}
\]
\leanref{Formalism.\allowbreak{}MainNotions.\allowbreak{}ShortIntervalPrimes.\allowbreak{}l21a}{Formalism.\allowbreak{}SharpRestriction.\allowbreak{}GlobalSchurBounds.\allowbreak{}FirstClaimLemmaTwoFourOfLemmaTwoThree}{368--402}

For large \(x\), (8.47.3) gives
\[
 \lambda_E\le8x^{\varepsilon/5}(\log x)^2.
 \tag{8.47.10}
\]
With \(V=x^{\varepsilon/3}-\log3\le\log y\), integrate (8.47.8) after the substitution \(v=\log u\). For \(V\le v\le\lambda_E^2\) use \(e^{-c_1v/(2\lambda_E)}\). For \(v\ge\max(V,\lambda_E^2)\) use \(e^{-c_1\sqrt v/2}\) and integrate in \(\sqrt v\). The bounds are polynomial factors in \(V,\lambda_E\) times \(e^{-c_1V/(2\lambda_E)}\) and \(e^{-c_1\sqrt V/2}\), respectively. By (8.47.10), \(V/\lambda_E\gg x^{2\varepsilon/15}/(\log x)^2\) and \(\sqrt V\gg x^{\varepsilon/6}\), so their sum is at most \(e^{-c_6x^{\varepsilon/8}}\) beyond a common family threshold.
\leanref{Formalism.\allowbreak{}MainNotions.\allowbreak{}ShortIntervalPrimes.\allowbreak{}thin\_\allowbreak{}of\_\allowbreak{}zeroFree}{Formalism.\allowbreak{}SharpRestriction.\allowbreak{}GlobalSchurBounds.\allowbreak{}FirstClaimLemmaTwoFourOfLemmaTwoThree}{2379--2397}
\leanref{Formalism.\allowbreak{}MainNotions.\allowbreak{}ShortIntervalPrimes.\allowbreak{}thin\_\allowbreak{}of\_\allowbreak{}psiThin}{Formalism.\allowbreak{}SharpRestriction.\allowbreak{}GlobalSchurBounds.\allowbreak{}FirstClaimLemmaTwoFourOfLemmaTwoThree}{1949--1992}
\end{proof}

\section{Standard masses and hyperplane volumes}

Over a number field \(K\) with conjugation, put \(h_M(z,w)=\sum_{a=1}^M\overline{z_a}w_a\) and \(\Lambda_M=\mathcal O_K^M\). For a positive definite lattice set
\[
 \operatorname{mass}(h,\Lambda)=
 \sum_{[\Lambda']\in\operatorname{gen}(\Lambda)}
       |\operatorname{Aut}(h,\Lambda')|^{-1}.
 \tag{8.48.1}
\]
Write \(M_j=\operatorname{massI}(D,j)\) for the mass of the standard rank-\(j\) lattice over \(E\). For \(f_K\) the number of infinite places of \(K\), put \(v_K=|\operatorname{disc}K|^{1/(2f_K)}/2\), and for \(M\ge3\) set
\[
 L_\eta(K,M)=\prod_{\mathfrak P}
 \begin{cases}
 (1-(-1)^M(\sqrt{\mathrm N\mathfrak P})^{-M})^{-1},
       &\overline{\mathfrak P}=\mathfrak P,\\
 (1-(\mathrm N\mathfrak P)^{-M})^{-1/2},
       &\overline{\mathfrak P}\ne\mathfrak P .
 \end{cases}
 \tag{8.48.2}
\]
The product counts conjugate maximal ideals separately. For \(\Lean{IsCMUnram}\), \(K\) is totally complex with involution inducing complex conjugation at every embedding, and every \(\mathfrak P\) admits \(\theta\in\mathcal O_K\) with \(\theta-\bar\theta\notin\mathfrak P\).

\begin{theorem}[The adjacent-rank mass formula]\label{B8.48}
For every \(K\) satisfying \(\Lean{IsCMUnram}\) and integer \(M\ge3\),
\[
 \frac{\operatorname{mass}(h_M,\Lambda_M)}
      {\operatorname{mass}(h_{M-1},\Lambda_{M-1})}
 =
 \frac{\bigl((M-1)!(v_K/\pi)^M\bigr)^{f_K}}{(2v_K)^{f_K/2}}\,
 L_\eta(K,M).
 \tag{8.48.3}
\]
\end{theorem}
\leanref{GlobalTheory.\allowbreak{}Mass.\allowbreak{}massRatio\_\allowbreak{}all\_\allowbreak{}unramified\_\allowbreak{}cmFields}{Development.\allowbreak{}AbelianVarietyTheory.\allowbreak{}Proofs.\allowbreak{}GlobalInterpolation.\allowbreak{}MassRatioAllUnramifiedCmFields}{8--9}
\leanref{Formalism.\allowbreak{}HermitianLattices.\allowbreak{}SiegelMassBounds.\allowbreak{}siegel\_\allowbreak{}mass\_\allowbreak{}ratio}{Formalism.\allowbreak{}GL3MeanSquares.\allowbreak{}HookFormulaTorsion.\allowbreak{}MassUpperBoundClosed}{12--12}
\leanref{Formalism.\allowbreak{}MainNotions.\allowbreak{}StandardLattice.\allowbreak{}SiegelMassRatioCM}{Formalism.\allowbreak{}SharpRestriction.\allowbreak{}GlobalSchurBounds.\allowbreak{}StandardLattice}{281--284}

Propositions 4.4 and 4.5 of \cite{gan-hanke-yu-2001}, with local corrections equal to one, give the same adjacent-rank quotient. Write \(M_j(K)\) for the standard mass over \(K\). In each genus representative of rank \(M\), count the vectors of norm at most \(J^2\) at every archimedean embedding. Weight each count by \(1/|\operatorname{Aut}|\) and divide the sum by \(J^{2f_K M}\). The volume limit is \(M_M(K)(\pi^M/M!)^{f_K}/v_K^{f_KM}\). Grouping by represented norm gives \(M_{M-1}(K)M^{-f_K}L_\eta(K,M)/(2v_K)^{f_K/2}\) from the adjacent-rank mass constant and the local densities. Equating yields (8.48.3). For the binary representation input write \(K=K^+(\sqrt\delta)\) with \(\delta\) totally negative. For totally positive \(a,b,n\in K^+\), the form \(\langle a,-a\delta,b,-b\delta,-n\rangle\) is isotropic at infinity by indefiniteness and at finite places by dimension five. A global isotropic vector has nonzero last coordinate by positivity. Dividing by it gives \(as\bar s+bt\bar t=n\).
\leanref{Formalism.\allowbreak{}HermitianLattices.\allowbreak{}Ext.\allowbreak{}L3.\allowbreak{}massConst\_\allowbreak{}succ}{Formalism.\allowbreak{}GL3MeanSquares.\allowbreak{}GenusMeanAsymptotics.\allowbreak{}UnimodularClassRepresentatives.\allowbreak{}MassConstSucc}{99--131}
\leanref{Formalism.\allowbreak{}HermitianLattices.\allowbreak{}Ext.\allowbreak{}L3.\allowbreak{}siegelMassRatio\_\allowbreak{}of\_\allowbreak{}LI}{Formalism.\allowbreak{}GL3MeanSquares.\allowbreak{}GenusMeanAsymptotics.\allowbreak{}UnimodularClassRepresentatives.\allowbreak{}SiegelMassRatioOfLI}{35--55}
\leanref{GlobalTheory.\allowbreak{}Mass.\allowbreak{}binary\_\allowbreak{}hermitian\_\allowbreak{}represents}{Development.\allowbreak{}AbelianVarietyTheory.\allowbreak{}Proofs.\allowbreak{}FlagMeasures.\allowbreak{}BinaryHermitianRepresents}{12--66}

For compatible constant signature \((P,Q)\), normalize the finite-adelic measure by \(\mu_f(K_{\Lambda,f})=1\). An admissible Haar measure uses this finite factor and the prescribed compact-dual measures \(\mu^{\rm cd}_{P,Q}\) in archimedean signature frames. The condition \(\Lean{VolTie}\,D\) supplies such a measure for each \(W\), and its arithmetic fundamental domains have volume \(\operatorname{vol}_D(W)>0\) for every admissible measure.

\begin{theorem}[The volume ratio for a unit hyperplane]\label{B8.49}
Assume \ref{B8.13} and \(\Lean{VolTie}\,D\). Let \(W\) have rank \(M\ge3\) and compatible constant signature \((P,Q)\). For \(s=1\) with \(P\ge1\), or \(s=-1\) with \(Q\ge1\), there exist a vector \(z\in W\) of norm \(s\) and a Hermitian space \(W'\) isometric to \(z^\perp\). Its signature is \((P-1,Q)\) or \((P,Q-1)\), respectively. With \(v_E=D.\mathrm{host}.v\),
\[
 \frac{\operatorname{vol}_D(W')}{\operatorname{vol}_D(W)}
 =\left(\frac{(2\pi)^M}{(M-1)!}\right)^f
   \frac{(2v_E)^{f(1/2-M)}}{L_\eta(E,M)} .
 \tag{8.49.1}
\]
This is \(\Lean{MassRatioUnderVolume}\,D\).
\end{theorem}
\leanref{GlobalTheory.\allowbreak{}VolRatioFromSiegelWeil.\allowbreak{}volumeRatio}{Development.\allowbreak{}AbelianVarietyTheory.\allowbreak{}Proofs.\allowbreak{}FlagMeasures.\allowbreak{}UnitaryHyperplaneVolumeRatioSiegelWeil}{10--14}
\leanref{Formalism.\allowbreak{}MainNotions.\allowbreak{}StandardLattice.\allowbreak{}VolumeRatio₂}{Formalism.\allowbreak{}SharpRestriction.\allowbreak{}GlobalSchurBounds.\allowbreak{}ThickPlacesCorollaryOfAnyTranslate}{2330--2340}

\begin{proof}
The complement construction gives \(z,W'\) and a self-dual lattice \(\Lambda\) containing \(z\), with \(\Lambda\cap z^\perp\) the base lattice of \(W'\). In \ref{B8.13} take the finite lattice indicators and an archimedean bump positive at \(z\). Its integral \(a\) on the regular norm-\(s\) fibre is finite and positive. The mean is
\[
 a\prod_{v<\infty}\alpha_v(M).
 \tag{8.49.2}
\]
The stabilizer of \(z\) is the unitary group of \(W'\), with finite level given by the intersection lattice. The quotient-to-fibre measure comparison gives \(c_M^{-f}a\) for \(c_M=(2\pi)^M/(M-1)!\). The datum with the form of \(W\) and base lattice \(\Lambda\) has volume \(\operatorname{vol}_D(W)\) by isometry invariance. The unfolded mean is
\[
 \frac{\operatorname{vol}_D(W')}{\operatorname{vol}_D(W)}c_M^{-f}a .
 \tag{8.49.3}
\]

For the global additive character standard at infinity, let \(U=\mathcal O_v\) and \(U^\perp\) be its character dual at a finite place of \(F\). With \(I_v=U\cap U^\perp\), the ratio \([U^\perp:I_v]/[U:I_v]=i_v\) is the norm of the local different. The unit-fibre density is \(i_v^{1/2-M}(1-\eta_v^M q_v^{-M})\), where \(q_v\) is the residue cardinality and \(\eta_v=1,-1\) at split and inert places, respectively. Use \(\prod_v i_v=|\operatorname{disc}F|=(2v_E)^f\) and (8.48.2) to get
\[
 \prod_{v<\infty}\alpha_v(M)=\frac{(2v_E)^{f(1/2-M)}}{L_\eta(E,M)} .
 \tag{8.49.4}
\]
Equate (8.49.2) and (8.49.3).
\leanref{GlobalTheory.\allowbreak{}VolRatioUnfold.\allowbreak{}ratio\_\allowbreak{}of\_\allowbreak{}setting}{Development.\allowbreak{}AbelianVarietyTheory.\allowbreak{}Proofs.\allowbreak{}GlobalInterpolation.\allowbreak{}VolumeRatioOfUnfolding}{295--391}
\leanref{GlobalTheory.\allowbreak{}VolRatioUnfold.\allowbreak{}volumeRatio\_\allowbreak{}of\_\allowbreak{}unfolding}{Development.\allowbreak{}AbelianVarietyTheory.\allowbreak{}Proofs.\allowbreak{}GlobalInterpolation.\allowbreak{}VolumeRatioOfUnfolding}{393--421}
\end{proof}

\begin{proposition}[The rank-two genus mass]\label{B8.50}
For every \(D\) and every rank-two \(W\) of compatible signature \((2,0)\) at each embedding, the base lattice satisfies
\[
 \operatorname{mass}(W.h,W.L_0)=M_2 .
 \tag{8.50.1}
\]
\end{proposition}
\leanref{GlobalTheory.\allowbreak{}genusMassAt\_\allowbreak{}two}{Development.\allowbreak{}AbelianVarietyTheory.\allowbreak{}Proofs.\allowbreak{}GlobalInterpolation.\allowbreak{}GenusMassAtTwo}{7--8}

\begin{proof}
At each nonsplit finite place a self-dual Gram matrix has unit determinant, a norm from the unramified quadratic extension. The spaces agree at split places and have the same signatures at infinity, so Landherr's theorem \cite[Ch.~10]{scharlau-1985} gives an isometry \(e:W\to(E^2,h_2)\).
\leanref{GlobalTheory.\allowbreak{}DefiniteMass.\allowbreak{}exists\_\allowbreak{}isometry}{Development.\allowbreak{}AbelianVarietyTheory.\allowbreak{}Proofs.\allowbreak{}GlobalInterpolation.\allowbreak{}UnramifiedOfIsCMUnram}{45--94}

The isometry preserves duals. The local classification of unimodular lattices puts \(e(W.L_0)\) in the standard genus \cite{jacobowitz-1962}.
\leanref{GlobalTheory.\allowbreak{}DefiniteMass.\allowbreak{}sameGenus\_\allowbreak{}stdLat}{Development.\allowbreak{}AbelianVarietyTheory.\allowbreak{}Proofs.\allowbreak{}GlobalInterpolation.\allowbreak{}UnramifiedOfIsCMUnram}{126--145}

Transport preserves the automorphism orders in (8.48.1).
\leanref{GlobalTheory.\allowbreak{}DefiniteMass.\allowbreak{}mass\_\allowbreak{}eq\_\allowbreak{}std}{Development.\allowbreak{}AbelianVarietyTheory.\allowbreak{}Proofs.\allowbreak{}GlobalInterpolation.\allowbreak{}UnramifiedOfIsCMUnram}{147--166}
\end{proof}

\begin{lemma}[Degree of a separated row]\label{B8.50a}
In the family's compact-dual measure let the row signature be \((p,q)=(p_Y,q_Y)\) with \(p,q\ge1\), \(r=p+q\) and \(F_s=\prod_{h=0}^{s-1}h!\). At column \(\rho\) let \(\lambda\) be the decreasing Harish--Chandra parameter and \(s_a\) its signs. Put \(V(\lambda)=\prod_{a<b}(\lambda_a-\lambda_b)/F_r\), \(H=\max_{a,b}|\lambda_a-\lambda_b|\). For \(e^{-r}\le t<1/2\), suppose
\[
 |\lambda_a-\lambda_b|\ge8192\min(p,q)(pq)^3\log r/t^2
 \qquad(s_a\ne s_b).
\]
Then the formal degree \(d_{\Sigma,\rho}\) and lowest-type dimension \(d_\tau\) satisfy \((1-t)V(\lambda)\le d_{\Sigma,\rho}\le4d_\tau(3pqH)^{pq}\).
\end{lemma}
\begin{proof}
Here \(\lambda_a-(r-1)/2\in\Z\). The gap exceeds \(32rpq+2\), hence meets the large-gap discrete-series and multiplicity-one thresholds. In this measure the formal-degree formula is \(d_{\Sigma,\rho}=V(\lambda)\). Set \(\rho_N(a)\) equal to half the number of opposite signs after \(a\) minus half the number before it. The compact Weyl formula gives
\[
 d_\tau=\frac{\prod_{a<b,\ s_a=s_b}(\lambda_a-\lambda_b+\rho_N(a)-\rho_N(b))}{F_p F_q}.
\]
For equal signs the shift difference counts opposite signs strictly between \(a,b\). The \(pq\) opposite-sign factors of the full Vandermonde are at most \(H\), and \(F_pF_q\le F_r\). Thus \(V(\lambda)\le d_\tau H^{pq}\), which gives the bounds.
\end{proof}
\leanref{PI.\allowbreak{}UnitaryGroupMassFormula.\allowbreak{}HCFdegRGap.\allowbreak{}hCFdegRGap\_\allowbreak{}holds}{Formalism.\allowbreak{}PeriodIntegrals.\allowbreak{}UnitaryGroupMassFormula.\allowbreak{}FormalDegreeGap.\allowbreak{}HCFdegRGapHolds}{108--111}
\leanref{Formalism.\allowbreak{}MainNotions.\allowbreak{}GappedMultiplicityBounds.\allowbreak{}theorem\_\allowbreak{}three\_\allowbreak{}rows\_\allowbreak{}of\_\allowbreak{}real\_\allowbreak{}place}{Formalism.\allowbreak{}SharpRestriction.\allowbreak{}GlobalSchurBounds.\allowbreak{}TheoremThreeFamilyOfGappedMultiplicity}{131--193}
\leanref{Formalism.\allowbreak{}MainNotions.\allowbreak{}GappedMultiplicityBounds.\allowbreak{}explicit\_\allowbreak{}bound\_\allowbreak{}family}{Formalism.\allowbreak{}SharpRestriction.\allowbreak{}GlobalSchurBounds.\allowbreak{}TheoremThreeFamilyOfGappedMultiplicity}{200--201}

\begin{proposition}[Volume induction and the scale of \(T_u\)]\label{B8.51}
Assume \ref{B8.49} and \ref{B8.50} for every datum. For every member condition and bundle association they give \(\Lean{H\_\allowbreak{}hTu}\) with threshold function
\[
 c_{T,\min}(y)=\max(11,\ 16+3\max(5,y+1)).
 \tag{8.51.1}
\]
On \(\mathrm{WithTw}(\mathrm{WithSm}(\mathrm{WithChN}(\mathcal A)))\), the quantifiers of \ref{B8.39} give
\[
 C_{Tu}\ge c_{T,\min}(C_{YL})\ \Longrightarrow\
 \exists x_0\ \forall x\ge x_0\ \forall i\in F_m.I_x\ \forall u,\quad
 b^{m/2}\le T_u,\quad \log T_u\le C_{Tu}mL .
 \tag{8.51.2}
\]
At compatible points the base volume satisfies
\[
 \operatorname{vol}_D(Y_0)=M_r .
 \tag{8.51.3}
\]
\end{proposition}
\leanref{GlobalTheory.\allowbreak{}Fc.\allowbreak{}row\_\allowbreak{}hTu\_\allowbreak{}cTmin\_\allowbreak{}alt}{Development.\allowbreak{}AutomorphicArithmetic.\allowbreak{}ThetaCorrespondence.\allowbreak{}ThetaPairLifts.\allowbreak{}ThirdPlaceSeries.\allowbreak{}RowHTuCTmin}{11--21}
\leanref{Formalism.\allowbreak{}L2.\allowbreak{}GlobalTheoryTheory.\allowbreak{}GenusMassVolumes.\allowbreak{}volume\_\allowbreak{}eq\_\allowbreak{}mass\_\allowbreak{}rank\_\allowbreak{}two}{Formalism.\allowbreak{}LocalToGlobal.\allowbreak{}GlobalTheory.\allowbreak{}GL5.\allowbreak{}DiscreteSeriesOfFormalDegree}{104--110}

\begin{proof}
With \(\Lean{VolTie}\) from the realization, the definite class decomposition and \ref{B8.50} give volume \(M_2\) in rank two. By (8.48.3) and (8.49.1), with \(R_M=M_M/M_{M-1}\),
\[
 \operatorname{vol}_D(W')/\operatorname{vol}_D(W)=R_M^{-1},
 \tag{8.51.4}
\]
since \((2v_E)^{f/2}(\pi/v_E)^{Mf}=(2\pi)^{Mf}(2v_E)^{f(1/2-M)}\).
For \(P\ge2\), remove negative unit vectors and then positive unit vectors down to rank two. The reverse ratios multiply to \(M_{P+Q}\), in particular \(M_r\) for \(Y_0\).
\leanref{Formalism.\allowbreak{}L2.\allowbreak{}GlobalTheoryTheory.\allowbreak{}GenusMassVolumes.\allowbreak{}vol\_\allowbreak{}indef}{Formalism.\allowbreak{}LocalToGlobal.\allowbreak{}GlobalTheory.\allowbreak{}GL5.\allowbreak{}DiscreteSeriesOfFormalDegree}{57--84}

Take \(C_Y=C_0=\max(5,C_{YL}+1)\), \(\Delta=D.s.\mathrm{Del}\) and \(w_{j\rho}=\mathrm{wj}(D,j+1,\rho)\). Beyond a common threshold these weights satisfy the separated-weight and row-size conditions of \(\Lean{DataA}\). For the common row product \(D^\star\) and the column product \(d_u^\star\) of \ref{B5.11}, put \(X_u=(D^\star/d_u^\star)^{1/(fk)}\). The product identities give \(\operatorname{logXuD}(D,u)=\log X_u\). Apply \ref{B8.50a} and the product calculation of \ref{B5.11} to get
\[
 \log X_u\le(15+3C_Y)mL .
 \tag{8.51.5}
\]
The torus-count constant may be any \(C\) with
\[
 C\ge11,\qquad C\ge16+3C_Y .
 \tag{8.51.6}
\]
\leanref{Formalism.\allowbreak{}MainNotions.\allowbreak{}StandardLattice.\allowbreak{}HarishChandraReduction.\allowbreak{}explicitConstants\_\allowbreak{}family\_\allowbreak{}of\_\allowbreak{}volume\_\allowbreak{}eq\_\allowbreak{}mass}{Formalism.\allowbreak{}SharpRestriction.\allowbreak{}GlobalSchurBounds.\allowbreak{}TheoremThreeFamilyOfFormalDegreeInput}{83--84}
\leanref{Formalism.\allowbreak{}MainNotions.\allowbreak{}StandardLattice.\allowbreak{}part\_\allowbreak{}five\_\allowbreak{}all\_\allowbreak{}of\_\allowbreak{}explicit}{Formalism.\allowbreak{}SharpRestriction.\allowbreak{}GlobalSchurBounds.\allowbreak{}TheoremThreeOfExplicitAndCompact}{328--332}

With \(\Gamma_u>1\) the \(\Lean{GamP}\) parameter, put \(T_u^\star=X_u\Gamma_u^{-2/(fk)}\). The datum gives \(T_u\in[T_u^\star,(1+1/(fk))T_u^\star]\). With \(\mathcal W=\sum_{j<m}f^{-1}\sum_\rho\log(w_{j\rho}v_E/\pi)\ge0\), the lower torus estimate gives
\[
 \log T_u+\frac{2\log\Gamma_u}{fk}
 \ge \mathcal W+m\log(g_0g_1)-C(L+m^2L/k)
 \ge16mL-C(L+1).
 \tag{8.51.7}
\]
The last inequality uses \(\log g_\nu\ge8L\) and \(m^2L/k\le1\). Use \(\log\Gamma_u/(fk)\le1\), \(L\ge1\) and \(m\ge x^{1/4}\) to get \(\log T_u\ge mL/2\) beyond the common threshold. By (8.51.5) and the upper endpoint, \(\log T_u\le(15+3C_Y)mL+1\le CmL\). Take \(C=c_{T,\min}(C_{YL})\).
\leanref{Formalism.\allowbreak{}MainNotions.\allowbreak{}GlobalTheoryTheory.\allowbreak{}tuGe\_\allowbreak{}of}{Formalism.\allowbreak{}SharpRestriction.\allowbreak{}GlobalSchurBounds.\allowbreak{}JointsL}{322--347}
\leanref{Formalism.\allowbreak{}MainNotions.\allowbreak{}StandardLattice.\allowbreak{}HarishChandraReduction.\allowbreak{}tuSize\_\allowbreak{}min}{Formalism.\allowbreak{}SharpRestriction.\allowbreak{}GlobalSchurBounds.\allowbreak{}TheoremThreeFamilyOfFormalDegreeInput}{97--103}
\end{proof}

\begin{lemma}[Selecting a prime at the word scale]\label{B8.51a}
Fix \(0<\varepsilon<1/52\) and \(c_s,C_Y>0\). For sufficiently large \(x\) let \(d\ge100\) be prime, \(d\equiv3\bmod4\), and let \(N\) be a positive multiple of \(d\) with \(\log N\le3x^{\varepsilon/5}(\log x)^2\). Suppose the positive integers \(g_i,m\) satisfy \(\log g_i\le C_Yx^\varepsilon\), \(m\ge x^{1/3}\). Take \(1\le V\le x^{12}\) and \(S>0\) with \(\log S\ge c_s m x^\varepsilon\). Put \(M=N(g_0g_1)^2\). There is a rational prime \(S<T\le(1+1/V)S\) of the form \(T=\pi_T\bar\pi_T\), with \(\pi_T\equiv1\bmod M\mathcal O_{\Q(\sqrt{-d})}\).
\end{lemma}
\begin{proof}
Apply \ref{B2.50} with modulus \(M\) on the interval with left endpoint \(S\) and relative length \(1/V\). We have \(M\ge d\) and \(\sqrt{\log S}\ge\sqrt{c_s}\,x^{1/6+\varepsilon/2}\). Both \(12\log x\) and
\[
 \log(4dM^2)\le\log4+9x^{\varepsilon/5}(\log x)^2+8C_Yx^\varepsilon
\]
are \(o(x^{1/6+\varepsilon/2})\). Thus \(e^{-c\sqrt{\log S}}\le1/V\le1\) and the conductor bound holds for the same \(c\). The count is positive. The norm condition is \(A^2+dB^2=4T,\ M\mid B,\ 2M\mid A-2-B\). Take \(\pi_T=(A+B\sqrt{-d})/2\).

Before choosing any level prime, the data give a common row number \(D^\star\) and the column products \(d_u^\star\). The row bound is \(D^\star\ge(2/3)^fM_r\operatorname{idx}_{\Sigma}\prod_\rho V(\lambda_\rho)\). The product calculation in the proof of \ref{B5.11}, applied to these numbers, gives \(\log X_u\ge15mL\) from \(\mathcal W\ge0\), \(\log(g_0g_1)\ge16L\) and an error of size at most \(mL\). Since \(\log\Gamma_u=C_\Gamma m^5fL\) and \(2C_\Gamma m^4\le k\), we have \(\log T_u^\star\ge mL\). Take \(N=Q_Sqd,\ V=fk,\ S=T_u^\star\), with the stated modulus bound. The bounds \(m\ge x^{1/3},\ fk\le x^{12},\ g_i\le b^9\) allow \(c_s=1,C_Y=9\). The excluded primes lie below \(e^{x^{\varepsilon/3}}\le S\). Taking norms gives \(T_u\equiv1\bmod M\), hence complete splitting in \(E\subseteq\Q(\zeta_N)\) and \(F\).
\end{proof}
\leanref{UpperTheory.\allowbreak{}T3.\allowbreak{}levelPrimeN}{Development.\allowbreak{}AutomorphicArithmetic.\allowbreak{}ShortNorms.\allowbreak{}levelPrimeN}{36--161}
\leanref{GlobalTheory.\allowbreak{}EX.\allowbreak{}WinLow.\allowbreak{}window\_\allowbreak{}lower}{Development.\allowbreak{}AbelianVarietyTheory.\allowbreak{}Proofs.\allowbreak{}GlobalInterpolation.\allowbreak{}WindowLower}{177--295}
\leanref{GlobalTheory.\allowbreak{}EX.\allowbreak{}OB.\allowbreak{}exists\_\allowbreak{}levCh\_\allowbreak{}of\_\allowbreak{}levelPrime}{Development.\allowbreak{}AbelianVarietyTheory.\allowbreak{}Proofs.\allowbreak{}GlobalInterpolation.\allowbreak{}ExistsLevChOfLevelPrime}{20--112}

\section{Two counting statements on the same selected data}

Write \(\mathcal D_3=\mathrm{D3}(3,\mathcal A)\) for the class carrying the selected data and \(F_m.\mathrm{cOut}\) for the list of outer constants.

\begin{theorem}[Counting constants from the label condition]\label{B8.52}
Assume \ref{B8.13} and the label row of \ref{B8.08} for \(K_P,Q,Q',L_k\) with enlargement \(R\) and condition \(Q'\wedge R\). There are counting constants \(c_{30}(\ell)\) for finite real lists \(\ell\) and \(c_{\rm Ic}>0\) such that for every \(F_m\),
\[
 0\le(c_{30}(F_m.\mathrm{cOut})).C,\qquad
 2(C_{YL}+1)\le(c_{30}(F_m.\mathrm{cOut})).A .
 \tag{8.52.1}
\]
On \(\mathcal D_3\) with the quantifiers of \ref{B8.39}, the constants give \ref{B8.53} for the admissible level-character idempotents on the specified realization with \(\max_v\min(p_v,q_v)\ge1\), the condition \(\Lean{PLevT'}\).

At every sufficiently large point there is one arithmetic row system \(A\) with the selected base data and a common nontrivial additive character for the local idempotents at each distinguished prime. The label-preserving inclusion identifies the selected indices with the indices having good label pairs. For each subset \(S'\) of pairs,
\[
 \dim A(S')=\#\{\tau\in\mathrm{Fin}(D_{\rm all})\mid
                     (\mathrm{lab}_0(\tau),\mathrm{lab}_1(\tau))\in S'\}.
\]
For \(k_0\in\{0,1\}\) and any subset \(S\) of its type set, write \(\operatorname{cyl}_{k_0}S=\{t\mid t_{k_0}\in S\}\). The total main term \(D_A^\star(A.B.e)\) is positive, and each cylinder main term is nonnegative. Then
\[
 \left|\dim A(\operatorname{cyl}_{k_0}S)
       -D_A^\star(e_{\operatorname{cyl}_{k_0}S})\right|
 \le e^{-c_{\rm Ic}rf\log r}D_A^\star(e_{\operatorname{cyl}_{k_0}S}).
 \tag{8.52.2}
\]
\end{theorem}
\leanref{GlobalTheory.\allowbreak{}CountRowByName.\allowbreak{}row\_\allowbreak{}counts\_\allowbreak{}of\_\allowbreak{}labels\_\allowbreak{}alt}{Development.\allowbreak{}AutomorphicArithmetic.\allowbreak{}ThetaCorrespondence.\allowbreak{}ThetaPairLifts.\allowbreak{}ThirdPlaceSeries.\allowbreak{}RowCountsOfLabels}{13--14}
\leanref{GlobalTheory.\allowbreak{}CountRow.\allowbreak{}row\_\allowbreak{}counts\_\allowbreak{}of\_\allowbreak{}ties\_\allowbreak{}alt}{Development.\allowbreak{}AutomorphicArithmetic.\allowbreak{}ThetaCorrespondence.\allowbreak{}ThetaPairLifts.\allowbreak{}ThirdPlaceSeries.\allowbreak{}RowCountsOfTies}{12--21}

\begin{proof}
Put \(A_0(\ell)=2(\ell[6]+1)\) with zero-based indexing and default \(0\). For \(F_m.\mathrm{cOut}\) the entry is \(C_{YL}\). With \(n_0,C_0\) from \ref{B8.53}, take
\[
 c_{30}(\ell)=\left(0,n_0(A_0(\ell)),A_0(\ell),0,C_0(A_0(\ell)),1/16,1/16\right)
 \tag{8.52.3}
\]
in the order \((C_{\rm res},n_0,A,C_b,C,c,c_b)\).
\leanref{GlobalTheory.\allowbreak{}CountConstants.\allowbreak{}count\_\allowbreak{}constants\_\allowbreak{}of\_\allowbreak{}siegelWeil}{Development.\allowbreak{}AbelianVarietyTheory.\allowbreak{}Proofs.\allowbreak{}FlagMeasures.\allowbreak{}CountConstantsOfSiegelWeil}{11--24}

On \(\mathcal D_3\), use (8.51.2) and the count and gap estimates beyond the realization threshold and \(x>F_m.x_1\). The label row supplies \(X\) with its good operator and an isometry from the selected spectral space to \(HH_X\). Finite-dimensionality from \ref{B8.53e} and the gap estimate give compatible labels under this isometry. With the standard nontrivial characters \(\psi_\nu\) take \(A=\mathrm{allRowsOfL}(X,(\psi_\nu)_\nu)\). The label inclusion gives the subset dimension formula. The irreducibility row supplies \(c_{\rm Ic}\).
\leanref{GlobalTheory.\allowbreak{}IrrRowConcrete.\allowbreak{}irr\_\allowbreak{}row\_\allowbreak{}closed}{Development.\allowbreak{}AbelianVarietyTheory.\allowbreak{}Proofs.\allowbreak{}FlagMeasures.\allowbreak{}IrreducibilityRowOfConcreteAutomorphicData}{18--41}

For a cylinder at \(k_0\) use its local idempotent and take \(g_\circ=D.\mathrm{host}.g(1-k_0)\) for unipotent support. The dictionary identifies the image dimension with \(\dim A(\operatorname{cyl}_{k_0}S)\). Apply (8.53.4) at \(n=r\) with \(0<c_{\rm Ic}\le c\). The base data give positivity for the total main term, and positive type gives nonnegativity for each cylinder.
\leanref{Formalism.\allowbreak{}MainNotions.\allowbreak{}N.\allowbreak{}CylinderCountFamily.\allowbreak{}countCylAt\_\allowbreak{}of}{Formalism.\allowbreak{}SharpRestriction.\allowbreak{}LevelLetterCocycles.\allowbreak{}JointConditionsZpowers.\allowbreak{}JCountP}{16--22}
\leanref{GlobalTheory.\allowbreak{}AllRowsAtRecord.\allowbreak{}allRows\_\allowbreak{}at\_\allowbreak{}record}{Development.\allowbreak{}AbelianVarietyTheory.\allowbreak{}Proofs.\allowbreak{}FlagMeasures.\allowbreak{}AllRowsAtRecord}{67--95}
\end{proof}

Let \(I\) be a nondegenerate unitary datum of rank \(n\) over a CM extension \(E/F\) unramified at finite places. Set \(f=[F:\Q]\), \(\mathrm{rd}=|\operatorname{disc}E|^{1/(2f)}\). The stabilizer \(K_f\) of the completed self-dual lattice \(\widehat\Lambda\) has congruence subgroup \(K(g)\), the kernel of its action modulo \(g\). Take discrete series \(\sigma_v\) with Harish--Chandra parameter \(\lambda_v\) and unit lowest-type vector \(w_v\) at each real place. Their Hilbert tensor product is \(H_\sigma\). Use finite-adelic Haar measure \(\mu_f\) and a fundamental domain \(\mathscr F\) for product Haar measure. On automorphic \(L^2\), use the right regular action. For right convolution \(R_e\) by \(e\), set
\[
 D_{\rm count}=\dim_{\C}\{R_e\circ A\mid
      A\in\operatorname{Hom}^{\rm cont}_{G_\infty}(H_\sigma,L^2)\}.
\]
With \(d_\infty=\prod_{v\mid\infty}d(\sigma_v)\) the product of the formal degrees for the prescribed archimedean measures, put
\[
 D^\star=d_\infty\,\operatorname{vol}(\mathscr F)\,\Re e(1).
 \tag{8.53.1}
\]

\begin{lemma}[Counting idempotents]\label{B8.53}
Assume \ref{B8.13}. For each real \(A\) there are \(n_0\in\N\) and \(C_0\ge0\) with the following property for every additional level predicate \(P\), every \(C\ge C_0\) and \(c=c_b=1/16\). Take the data above with \(g,g_\circ\in\N\) and
\[
 \begin{gathered}
 n\ge n_0,\quad b\ge2,\quad \log\log\mathrm{rd}\le n,\quad
 |\lambda_{v,a}|\le b^2,\quad (p_v,q_v)=(p,q),\quad q^2\le n,\\
 1\le g\le b^A,\qquad g_\circ\text{ prime},\qquad 4n<g_\circ .
 \end{gathered}
 \tag{8.53.2}
\]
The prime \(g_\circ\) splits completely and unramified in \(I.E\). The admissible \(e\) is a locally constant self-adjoint convolution idempotent of positive type, supported in \(K_f\) and bi-invariant under \(K(g)\), with unipotent support modulo \(g_\circ\). Require \(P(I,\mu_f,g_\circ,e)\). At each \(v\) the coordinates of \(\lambda_v\) strictly decrease and lie in \((n-1)/2+\Z\). Require \(|\lambda_{v,a}-\lambda_{v,a'}|\ge h\) whenever \(a,a'\) have different signature signs, with
\[
 Cn^{19}f^5\log(b\,\mathrm{rd})\le h,\qquad
 \log f\le\sqrt n,\qquad
 \log\log b\le c_bnf\log n .
 \tag{8.53.3}
\]
Then
\[
 |D_{\rm count}-D^\star|\le e^{-cnf\log n}D^\star .
 \tag{8.53.4}
\]
\end{lemma}
\leanref{GlobalTheory.\allowbreak{}CountBridge.\allowbreak{}count\_\allowbreak{}unbundled}{Development.\allowbreak{}AbelianVarietyTheory.\allowbreak{}Proofs.\allowbreak{}FlagMeasures.\allowbreak{}UnitaryAutomorphicCountingUnbundled}{40--74}
\leanref{GlobalTheory.\allowbreak{}AllRowsAtRecord.\allowbreak{}countAllRows\_\allowbreak{}of}{Development.\allowbreak{}AbelianVarietyTheory.\allowbreak{}Proofs.\allowbreak{}FlagMeasures.\allowbreak{}AllRowsAtRecord}{37--65}
\leanref{Formalism.\allowbreak{}MainNotions.\allowbreak{}S.\allowbreak{}CountE.\allowbreak{}CountIdempotentsClass}{Formalism.\allowbreak{}SharpRestriction.\allowbreak{}CuspSpanNormBounds.\allowbreak{}BridgeCountE}{45--59}
\leanref{Formalism.\allowbreak{}L2.\allowbreak{}GlobalTheoryTheory.\allowbreak{}ConcN.\allowbreak{}CountAllRows}{Formalism.\allowbreak{}LocalToGlobal.\allowbreak{}GlobalTheory.\allowbreak{}IrreducibleOfCountSizes}{183--186}

\begin{proof}
By Schur orthogonality, evaluation at \(\bigotimes_v w_v\) identifies the counted intertwiners with the image of the projector given by right convolution by
\[
 \phi(y)=d_\infty\,
   \overline{\prod_{v\mid\infty}\langle w_v,\sigma_v(y_v)w_v\rangle}\ e(y_f).
\]
The kernel is
\[
 K(x,y)=\sum_{\gamma\in U(F)}\phi(x^{-1}\gamma y).
 \tag{8.53.5}
\]
The identity term on the diagonal is \(a=d_\infty\Re e(1)\ge0\). For definite signature the counting identity gives \(D_{\rm count}=D^\star\). Otherwise \(1\le q\le p\) since \(q^2\le n\) and \(n\ge8\).

Put \(H=nf\log n\), \(N=n^2f\), \(\alpha=e^{-H/4}\), \(\beta=e^{-N}\), and
\[
 X=nf\,\mathrm{rd}\,gb,\qquad \eta=(10^6n^8f^2)^{-1},\qquad
 \tau_*=X^{-11q}.
\]
Write \(\operatorname{least}(x)\) for the least Harder--Narasimhan size of a nonzero sublattice at \(x\), and \(\tau_1(x)=\min(1/2,\operatorname{least}(x))\). The moderate region has \(\operatorname{least}(x)\ge\tau_*\), and the deep region is its complement in \(\mathscr F\). Write \(\mathrm{nearAbs}\) and \(\mathrm{farSum}\) for the absolute diagonal sums over nonidentity elements with all \(d_v\le\eta\) and with some \(d_v>\eta\), respectively. Here \(d_v(y)=(\sum_a(\log s_a(y))^2)^{1/2}\) for the singular values \(s_a(y)\) in signature coordinates.

By \ref{B8.53e} the image has a finite orthonormal basis of continuous representatives \(F_\nu\), with
\[
 K(x,x)=\sum_\nu|F_\nu(x)|^2,\qquad
 D_{\rm count}=\int_{\mathscr F}K(x,x)\,dx.
 \tag{8.53.6}
\]
The regional estimates are
\[
 \begin{split}
 \int_{\rm moderate}\mathrm{nearAbs}&\le\alpha D^\star,\qquad
 \int_{\rm moderate}\mathrm{farSum}\le\beta D^\star,\\
 \operatorname{vol}(\mathrm{deep})&\le\alpha\operatorname{vol}(\mathscr F),
 \qquad \int_{\rm deep}K(x,x)\,dx\le\beta D_{\rm count}.
 \end{split}
 \tag{8.53.7}
\]
Lemma~\ref{B8.53a} and the positive-type inequality \(|e(y_f)|\le\Re e(1)\) bound each term \(y=x^{-1}\gamma x\) by
\[
 |\phi(y)|\le a\prod_{v\mid\infty}
 \exp\!\left[-\frac h8\min\!\left(d_v(y_v)^2,
                          \frac{d_v(y_v)}{\sqrt{2q}}\right)\right].
 \tag{8.53.8}
\]

For a near element in the support, the finite level makes \(\operatorname{tr}(\gamma^j)-n\) integral and unipotence modulo \(g_\circ\) makes it divisible by \(g_\circ\). Its absolute value is at most \(3n<g_\circ\) at each embedding for \(1\le j\le n\), so its norm forces it to vanish. Newton's identities give unipotence and the polar bound gives \(\|x^{-1}\gamma x-1\|_{\rm HS}\le\delta=3\sqrt{n\eta}<1/4\) at each real place. For \(\operatorname{least}(x)\ge1/2\), separation in \ref{B8.53a} leaves the identity. Otherwise unfold over flags with \(s\) blocks of sizes \(j_i\) and total \(1\le J\le q\). A root scale \(\lambda\in(0,1)\) of multiplicity \(m_\lambda\) contributes \((1+12\delta/\lambda)^{m_\lambda}\) to the count bound and \(\lambda^{m_\lambda}\) to the density. The product over the roots is at most \(e^{36\delta qnf}\le2\), since \(\lambda+12\delta\le e^{12\delta}\), \(\sum_\lambda m_\lambda\le3Jnf\) and \(\delta qnf\le1/100\). Integrate the remaining \(\prod_i 2fj_i\,dt_i/t_i\) over \([\tau_*,1/2]^s\). The integral for the type is at most
\[
 2\kappa D^\star\prod_i\left(2fj_i\log\frac1{2\tau_*}\right)
 \le 2e^{-JH/2}Y^sD^\star ,
\]
where \(\kappa\le e^{-JH/2}\) by \ref{B8.53b} and \(Y=22fq^2\log X\le e^{H/8}\) for large \(n\). Since \(s\le J\), summation over fewer than \(2^q\) types gives \(2^{q+1}e^{-3H/8}D^\star\le\alpha D^\star\).
\leanref{GlobalTheory.\allowbreak{}NearType.\allowbreak{}integral\_\allowbreak{}count\_\allowbreak{}dens}{Development.\allowbreak{}AbelianVarietyTheory.\allowbreak{}Proofs.\allowbreak{}FlagMeasures.\allowbreak{}NearTermsModeratePointsOneTypeBound}{40--73}
\leanref{GlobalTheory.\allowbreak{}NearMod.\allowbreak{}nearModerate}{Development.\allowbreak{}AbelianVarietyTheory.\allowbreak{}Proofs.\allowbreak{}FlagMeasures.\allowbreak{}OneTypeArith}{57--157}

Put \(A_x=(3\sqrt n/\tau_1(x)^2)^{2N}\). By \ref{B8.53a}, the arithmetic ball with radii \(R_v\ge0\) satisfies
\[
 \#\{\gamma\in U(F)\mid (x^{-1}\gamma x)_f\in K_f,\
                 d_v((x^{-1}\gamma x)_v)\le R_v\ \forall v\}
 \le A_x\exp\!\left(2n^2\sum_v R_v\right).
 \tag{8.53.9}
\]
Summing (8.53.8) over dyadic shells using (8.53.9) gives
\[
 \mathrm{farSum}(x)\le a A_x\,2f3^f e^{-h\eta^2/16}.
 \tag{8.53.10}
\]
For a place with \(d_v>\eta\), half its decay contributes \(e^{-h\eta^2/16}\) since \(\eta\le1/\sqrt{2q}\). The remaining decay bounds the shell counts by geometric tails. On the moderate region, \(A_x\le(3\sqrt n/\tau_*^2)^{2N}\). Under (8.53.2)--(8.53.3) the gap estimate gives
\[
 2N\log(3\sqrt n/\tau_*^2)+\log(2f)+f\log3+N\le h\eta^2/16 .
\]
Thus the pointwise bound is \(a\beta\), with integral at most \(\beta D^\star\).
\leanref{GlobalTheory.\allowbreak{}FarTerms.\allowbreak{}farBound\_\allowbreak{}of\_\allowbreak{}ptBound}{Development.\allowbreak{}AbelianVarietyTheory.\allowbreak{}Proofs.\allowbreak{}FlagMeasures.\allowbreak{}SummableOfSupport}{42--124}
\leanref{GlobalTheory.\allowbreak{}Deep.\allowbreak{}farModerate\_\allowbreak{}of\_\allowbreak{}farBound}{Development.\allowbreak{}AbelianVarietyTheory.\allowbreak{}Proofs.\allowbreak{}FlagMeasures.\allowbreak{}FarTermsModerateRegionOfFlagBound}{53--113}

By \ref{B8.53b}, \(\operatorname{vol}(\mathrm{deep})\le e^{-H/3}\operatorname{vol}(\mathscr F)\), since \(\tau_*<1/2\). The projector integral is bounded in \ref{B8.53e}.
\leanref{GlobalTheory.\allowbreak{}CountFinal.\allowbreak{}countingTheorem\_\allowbreak{}of\_\allowbreak{}volumeStep}{Development.\allowbreak{}AbelianVarietyTheory.\allowbreak{}Proofs.\allowbreak{}FlagMeasures.\allowbreak{}UnitaryAutomorphicCountingOfCompactDualVolumeRat}{37--49}

Take \(n_0\ge e^{15}\) large enough for (8.53.7). Integrating \(K(x,x)-a\) gives
\[
 |D_{\rm count}-D^\star|\le(2\alpha+\beta)D^\star+\beta D_{\rm count}.
 \tag{8.53.11}
\]
For \(\theta=e^{-H/16}\), one has \(\alpha,\beta\le\theta/8\) and \(\beta\le1/2\). Substitute \(D_{\rm count}\le D^\star+|D_{\rm count}-D^\star|\) on the right to get
\[
 (1-\beta)|D_{\rm count}-D^\star|\le(2\alpha+2\beta)D^\star
 \le \theta D^\star/2 .
\]
This proves (8.53.4). For the unbundled hypotheses take \(C_0=\max(0,8\cdot10^{14}(A+4))\).
\leanref{GlobalTheory.\allowbreak{}CountSum.\allowbreak{}countingTheorem\_\allowbreak{}of\_\allowbreak{}finite\_\allowbreak{}of\_\allowbreak{}bounds}{Development.\allowbreak{}AbelianVarietyTheory.\allowbreak{}Proofs.\allowbreak{}FlagMeasures.\allowbreak{}DiscreteSeriesCountingOfFiniteAndBounds}{34--142}
\leanref{GlobalTheory.\allowbreak{}abs\_\allowbreak{}sub\_\allowbreak{}le\_\allowbreak{}exp\_\allowbreak{}neg\_\allowbreak{}sixteenth\_\allowbreak{}mul}{Development.\allowbreak{}AbelianVarietyTheory.\allowbreak{}Proofs.\allowbreak{}FlagMeasures.\allowbreak{}AbsSubLeExpNegSixteenthMul}{7--43}
\end{proof}

\begin{lemma}[Local decay and arithmetic balls]\label{B8.53a}
At a real place of \ref{B8.53} put \(q=\min(p_v,q_v)\ge1\) and assume \(h\ge120nq^{3/2}\). On the lowest type the compression of \(\sigma_v(y)\) has norm at most \(\exp[-(h/8)\min(d^2,d/\sqrt{2q})]\), where \(d=d_v(y)\). For \(n\ge1\) and radii \(R_v\ge0\), the arithmetic ball is finite, and its cardinality satisfies (8.53.9).
\end{lemma}
\begin{proof}
Write \(y=k\exp Z\) with \(k\) compact and \(Z\) interchanging the signature spaces, \(\|Z\|_{\rm HS}=d\). For \(Y=l^{-1}Zl\), average over the maximal compact in \(l\) the rank-one projections onto translates of a highest-weight unit vector, weighted by the type dimension times
\[
 \prod_{i=1}^{n-1}\Delta_i(e^Y)^{-(\lambda_i-\lambda_{i+1}+1)},
\]
where \(\Delta_i\) is the leading minor of order \(i\). Schur orthogonality gives the value at the identity. The lowest-type differential equations identify the Taylor coefficients at \(t=0\) of the average and the compression along \(e^{tZ}\). Both functions of \(t\) are real analytic, so they agree along this curve.

By Cauchy--Binet, \(\Delta_i(e^{tY})=\mathbb E e^{tU}\) for sums \(U\) of \(i\) eigenvalues, weighted by squared eigenbasis minors. Differentiation at zero gives \(\mathbb E U=0,\ \mathbb E U^2=\sum_{a>i\ge b}|Y_{ab}|^2\). Since \(\operatorname{rank}Y\le2q\), \(|U|\le\sqrt{2q}\,d\). Take \(t=\min(1,1/(\sqrt{2q}\,d))\) for \(d>0\). From \(e^z\ge1+z+z^2/3\) on \([-1,1]\) and \(\mathbb E e^{tU}\le(\mathbb E e^U)^t\) one gets \(\log\Delta_i(e^Y)\ge t\mathbb E U^2/4\). Each nonzero \(Y_{ab}\) below the diagonal joins opposite signs and has total weight \(\sum_{b\le i<a}(\lambda_i-\lambda_{i+1}+1)\ge h\). Their squared moduli sum to \(d^2/2\), so the multiplier is at most \(e^{-htd^2/8}\). The dimension times the unweighted average of the projections is the identity, which gives the norm bound. At \(d=0\) use unitarity.

Put \(t_0=\tau_1(x)\). Two distinct points differ by a nonzero integral lattice map \(L\). The size inequality gives \((\prod_v\|L_v\|_{\rm op})^{1/f}\ge t_0^2\), since the largest source size is the reciprocal of the least target size. Scale the local matrices by \(e^{\bar R-R_v}\), with \(\bar R=f^{-1}\sum_v R_v\). The scales have product one, so the points are \(t_0^2\)-separated in a maximum-norm ball of radius \(\sqrt n\,e^{\bar R}\). Packing balls of radius \(t_0^2/2\) in real dimension \(2N\) gives \((1+2\sqrt n\,e^{\bar R}/t_0^2)^{2N}\), which implies (8.53.9).
\end{proof}
\leanref{GlobalTheory.\allowbreak{}DecGapBound.\allowbreak{}decGap\_\allowbreak{}of\_\allowbreak{}formulaGap}{Development.\allowbreak{}AbelianVarietyTheory.\allowbreak{}Proofs.\allowbreak{}FlagMeasures.\allowbreak{}LowestKTypeDecayGapOfCompressionFormulaGap}{11--32}
\leanref{GlobalTheory.\allowbreak{}Jet.\allowbreak{}FormGapFree.\allowbreak{}formulaGap}{Development.\allowbreak{}AbelianVarietyTheory.\allowbreak{}Proofs.\allowbreak{}FlagMeasures.\allowbreak{}LowestKTypeCompressionFormulaOfStrictLowestKType}{22--99}
\leanref{GlobalTheory.\allowbreak{}AutBall.\allowbreak{}autBallCount}{Development.\allowbreak{}AbelianVarietyTheory.\allowbreak{}Proofs.\allowbreak{}FlagMeasures.\allowbreak{}AutBallCount}{28--48}

\begin{lemma}[Two volume estimates]\label{B8.53b}
Let \(V=\operatorname{vol}(\mathscr F)\), \(r=\mathrm{rd}\), \(N=n^2f\), \(H=nf\log n\). Let \(c_v\mu_v\) have Hilbert--Schmidt density one at the identity. Set \(P_0=\prod_v c_v,\ k_R=\mu_f(K(g))\). Under the unit-vector recursion, for \(n=p+q\ge8,\ q^2\le n,\ Z\ge2,\ r\le Z,\ 2fr\le Z^2\), one has \(P_0V\le Z^N\mu_f(K_f)\le Z^Ng^{2N}k_R\) for \(g\ge1\).
If also \(\log n\ge15,\ \log r\le e^n,\ \log f\le\sqrt n\), the flag density constant for positive increments \(j_1,\ldots,j_a\) of sum \(1\le J\le q\) is at most \(e^{-JH/2}\), and \(\Pr(\operatorname{least}<1/2)\le e^{-H/3}\) for normalized volume.
\end{lemma}
\begin{proof}
Apply the one-vector Gram formula of \ref{B8.13} for norm \(\pm1\) at rank \(M\ge3\) \cite[\S52, Theorem~5, (41)]{weil-1965}. Removing the compact-dual factors in (8.49.1) gives gauge-volume multiplier \(r^{f(M-1/2)}L_\eta(E,M)\) on restoring the vector. Descend to signature \((2,0)\). Write \(Y=r^f,\ L=1+\log Y\). A short line gives at most \(YL^f\) possible ideals generated by real elements and \(2^fYL^{2f}\) choices for each of the two line classes with its unit squareclass. There are at most \(Y^2\) gluing choices. Multiply by \((8\pi^3)^f\). Since \(\prod_{M=3}^n L_\eta(E,M)\le(\prod_{M=3}^n\zeta(M))^f\le1.4^f\), the result is at most \(Z^{f(26+(n^2-4)/2)}\le Z^N\), using \(L\le Z^2\) and \(8\pi^3\cdot4\cdot1.4\le Z^{11}\). The level index is at most \(g^{2N}\) by reduction in \(\operatorname{End}_{\mathcal O_E}(\Lambda)/g\).

In the flag unfolding the residual \(\U(p-J,q-J)\) volume cancels against \(V\) by this recursion. Set \(S_m=\prod_{l=1}^m 2\pi^l/(l-1)!\) and \(d_U=4\sum_{i<k}j_ij_k+2J(n-2J)+\sum_i j_i^2\). The radical covolume and Levi volumes give
\[
 \kappa=(r/2)^{-fd_U/2}2^{-f(nJ-J^2+J/2)}
 \left({S_pS_q\over S_{p-J}S_{q-J}\prod_i S_{j_i}}\right)^f
 {(\operatorname{res}_1\zeta_E)^a\prod_i\prod_{l=2}^{j_i}\zeta_E(l)\over\prod_{M=n-2J+1}^n L_\eta(E,M)} .
\]
Put \(s=\sqrt n,\ l=\log n,\ c=\log(2\pi)\). The six log factors divided by \(fJ\) are bounded respectively by
\[
 n/3,\quad 0,\quad c-(n-2s)(l-c-1)+3.49+sl/4,\quad 2(3+n),\quad 2,\quad 1 .
\]
Use \(r/2\ge3/4,\ d_U\le2Jn\), the factorial products for \(S_m\) with \(J\le q\le s\), the residue bound \(\operatorname{res}_1\zeta_E\le(16\log r)^{2f}\), and the Euler products for the remaining factors. The sum is at most \(-nl/2\) since \(c\le2,\ l\ge15,\ 100s\le n\). On \(0<t_1<\cdots<t_a<1/2\) the radical factors are at most one and include \(t_i^{2fj_i(n-2J)}\) for each \(i\). Against \(\prod_i2fj_i\,dt_i/t_i\) the density integrates to at most \(\kappa\), by integrating \(\prod_i2fj_i t_i^{2fj_i-1}\) on the unit cube. A least size below \(1/2\) gives a nonempty flag. Summing over its possible compositions gives \(\sum_{J=1}^q2^{J-1}e^{-JH/2}\le e^{-H/3}\).
\end{proof}
\leanref{GlobalTheory.\allowbreak{}FlagMeas.\allowbreak{}AbsVol.\allowbreak{}absoluteVolume\_\allowbreak{}of\_\allowbreak{}volumeStep}{Development.\allowbreak{}AbelianVarietyTheory.\allowbreak{}Proofs.\allowbreak{}FlagMeasures.\allowbreak{}HermAbsoluteVolumeBoundOfVolumeRatioInduction}{37--69}
\leanref{GlobalTheory.\allowbreak{}FlagMeas.\allowbreak{}constantBound}{Development.\allowbreak{}AbelianVarietyTheory.\allowbreak{}Proofs.\allowbreak{}FlagMeasures.\allowbreak{}FlagExplicitConstantLeExpNegHalf}{82--228}
\leanref{GlobalTheory.\allowbreak{}NonzeroFlag.\allowbreak{}nonzeroFlagBound}{Development.\allowbreak{}AbelianVarietyTheory.\allowbreak{}Proofs.\allowbreak{}FlagMeasures.\allowbreak{}NonzeroFlagProbabilityBoundOfVolumeRatioInductio}{84--148}

\begin{lemma}[The deep-region exponent]\label{B8.53c}
Let \(n,f,q\) be integers with \(n\ge8,\ f,q\ge1,\ q^2\le n\). For \(r,g\ge1,\ b\ge2\) put \(X=nfrgb,\ N=n^2f,\ k=8qN\). If \(0<\tau\le X^{-11q}\), set
\[
 \rho={n^2\sqrt f\,gr\over4}(2\tau)^{1/q},\quad
 B_j=\sqrt n\,b^2+\tfrac32 n^{3/2}+3j,\quad B=B_{k+N},
\]
\[
 U=(2\rho\sqrt f B)^{2k}(6B)^Ne^f(17\sqrt n/\tau^2)^{2N}.
\]
Then \(\rho\le\log2\) and \(Ug^{2N}X^N\le e^{-N}\).
\end{lemma}
\begin{proof}
Write \(s=(2\tau)^{1/q},\,Y=X^q\). Since \(X\ge16,\ sX^{11}\le2,\ n^2\sqrt f\,gr\le X^2\), one has \(\rho\le1/2\). From \(q\le n\) one gets \(B\le30X^3\le X^6\) and \(n^2fgrB/2\le15X^5\le X^6\). Thus \((2\rho\sqrt fB)^{16q}\le Y^{96}(2\tau)^{16}\). Using \(f\le N,\ \tau Y^{11}\le1,\ Bng^2X\le Y^{10}\) gives
\[
 (Ug^{2N}X^N)^{1/N}\le 2^{16}\cdot6\cdot289\,e\,Y^{96}\tau^{12}Bng^2X
 \le 2^{16}\cdot6\cdot289\,e\,Y^{-26}\le e^{-1}.
\]
\end{proof}
\leanref{GlobalTheory.\allowbreak{}deep\_\allowbreak{}region\_\allowbreak{}exponents}{Development.\allowbreak{}AbelianVarietyTheory.\allowbreak{}Proofs.\allowbreak{}FlagMeasures.\allowbreak{}DeepRegionExponents}{54--204}

\begin{lemma}[A finite family of cuspidal isometries]\label{B8.53d}
Use the count data of \ref{B8.53} with \(n_0\ge8\), \(C\ge\max(0,8\cdot10^{14}(A+4))\), \(q\ge1\). Let \(T_i:H_\sigma\to L^2\), \(1\le i\le D\), be linear isometries with cuspidal images. Let \(V_j\) span derivatives through order \(j\) of lowest-type tensors, taken one place at a time. On these spaces assume \(F_i(u)\) is a continuous automorphic representative of \(T_i u\), is right \(K(g)\)-invariant and has right derivative \(F_i(d\sigma(Z)u)\) in direction \(Z\). Write \(F(u,y)=(F_i(u)(y))_i,\ w=\bigotimes_v w_v\). For \(\tau=\operatorname{least}(x)<X^{-11q}\), use \ref{B8.53c} with \(r=\mathrm{rd}\). Then
\[
 \|F(w,x)\|_2^2\le U D P_0/k_R .
\]
\end{lemma}
\begin{proof}
Differentiation maps \(V_j\) to \(V_{j+1}\). At one place the sum of squared derivative norms in an orthonormal Lie basis is the quadratic form of \(2\Omega_K-\Omega_G\). A compact highest weight \(\mu\) at order \(j\) satisfies \(\|\mu+\rho_K\|\le\|\lambda\|+\|\rho_{\rm nc}\|+\sqrt2j\), since one derivative adds a weight of the adjoint representation. Here the shifts are the compact and noncompact half-sums. With \(L=\|\lambda\|\) and \(a=\|\rho_{\rm nc}\|+\sqrt2j\), the Casimir difference is at most
\[
 2(L+a)^2-L^2+n^3/12\le(L+2a+0.3n^{3/2})^2\le B_j^2,
\]
using \(\|\rho_{\rm nc}\|\le0.58n^{3/2}\) and \(L\le\sqrt n\,b^2\). Cauchy--Schwarz over the places gives \(\|d\sigma(Z)u\|\le\sqrt f B_j\|Z\|_{\rm HS}\|u\|\) for the product norm.

Take the ordered exponential cube of side \(1/B\) in \(N\) coordinates. The gap satisfies \(h\le2b^2\), so the gap lower bound gives \(b^2\ge n^{19}\). The bounds \(12n^3/B,8n^5/B\le1\) give chart density at most \(e^fP_0\) relative to Haar measure. Thicken by \(K(g)\) to divide by \(k_R\). Overlaps differ by local operators of norm at most \(2\). The packing in \ref{B8.53a} bounds their multiplicity by \((1+4\sqrt n/\tau_1(y)^2)^{2N}\). In each coordinate use
\[
 |v(0)|^2\le(2/l)\int_0^l|v|^2+2l\int_0^l|v'|^2,\qquad l=1/B.
\]
Through order \(k+N\) each step costs at most \(2B+2B^{-1}B^2\le6B\), with \(\|T_i u'\|=\|u'\|\) for every differentiated vector \(u'\).

Append \(1/2\) to the flag sizes starting at \(\tau\). At most \(q\) consecutive ratios have product \(2\tau\), so some ratio \(t/t'\le s=(2\tau)^{1/q}\) and \(t^2\le s/4\). In the unipotent radical of the associated isotropic space, reduce the logarithm modulo \(2g\) times the integral first layer and the central remainder with its half-commutator modulo \(g\) times the centre lattice. The covering error is at most \(n^2\sqrt f\,grs(1/(4\sqrt2)+1/16)\le\rho\). Modulo rational translation and \(K(g)\), paths of length at most \(\rho\) to \(x\exp Z_1\exp Z_2\) cover the real orbit. Strong approximation in the radical gives the adelic cover. Along the paths \(\tau_1(y)\ge\tau/2\) since \(e^\rho\le2\). The cube estimate summed over \(i\) gives \(\|F(u,y)\|\le M\|u\|\) on \(V_k\), where
\[
 M^2=(6B)^Ne^f(17\sqrt n/\tau^2)^{2N}D P_0/k_R .
\]

Given \(M'\|u\|\) on \(V_{j+1}\) throughout the neighborhood, path integration bounds \(\|F(u,y)-F(u,x)\|\) on \(V_j\) by \(\rho\sqrt f B_jM'\|u\|\). The continuous constant term vanishes since it vanishes almost everywhere. Average the difference on the unipotent quotient with probability measure to get the same bound for \(\|F(u,x)\|\). The triangle inequality gives \(2\rho\sqrt f B_jM'\|u\|\) throughout the neighborhood. Iterate \(k\) times and square at \(w\).
\end{proof}
\leanref{GlobalTheory.\allowbreak{}SobolevBound.\allowbreak{}sobolev\_\allowbreak{}bound}{Development.\allowbreak{}AbelianVarietyTheory.\allowbreak{}Proofs.\allowbreak{}FlagMeasures.\allowbreak{}MeasurePreservingReindex}{239--285}
\leanref{GlobalTheory.\allowbreak{}DeepFamilyGap.\allowbreak{}deep\_\allowbreak{}bound\_\allowbreak{}family}{Development.\allowbreak{}AbelianVarietyTheory.\allowbreak{}Proofs.\allowbreak{}FlagMeasures.\allowbreak{}DeepPointBoundCuspidalIsometricIntertwinerFamily}{118--183}

\begin{corollary}[Mass of the counting projector]\label{B8.53e}
In the indefinite case of \ref{B8.53} the counted space is finite dimensional and \(\int_{\rm deep} K(x,x)\,dx\le e^{-N}D_{\rm count}\).
\end{corollary}
\begin{proof}
On a finite-dimensional subspace of counted intertwiners, \(T^*S=\langle T,S\rangle I\) by Schur's lemma. An orthonormal basis consists of isometries with orthogonal \(K(g)\)-fixed images. The translated discrete-series coefficient has zero real integral over the radical of every nonzero rational isotropic space. Its decay in log operator norm has rate \(h/(8\sqrt q)>12nq\), as required for the unipotent integral. Unfolding gives zero constant term for \(T_iw\). Its translates span a dense subspace of the image, and the cuspidal space is closed.

For \(u\in V_j\), express \(u\) as a smooth archimedean average of \(w\). Convolve \(T_iw\) with the averaging function times a normalized finite level bump. This gives the continuous representative \(F_i(u)\) of \(T_i u\), which is right \(K(g)\)-invariant by uniqueness of continuous representatives. Differentiating the convolution gives a continuous representative of \(T_i(d\sigma(Z)u)\), since \(T_i\) intertwines the action. Uniqueness gives \(R_ZF_i(u)=F_i(d\sigma(Z)u)\) pointwise for the same choice of representatives on all \(V_j\).
\leanref{GlobalTheory.\allowbreak{}SmoothRepsType.\allowbreak{}repsExist}{Development.\allowbreak{}AbelianVarietyTheory.\allowbreak{}Proofs.\allowbreak{}FlagMeasures.\allowbreak{}ExistsContinuousRowRepresentativesLowestKType}{87--137}

For this basis, \ref{B8.53c} and \ref{B8.53d} give \(\|F(w,x)\|^2\le e^{-N}DP_0/(g^{2N}X^Nk_R)\) at each deep point. By \ref{B8.53b} at \(Z=X\), this is at most \(e^{-N}D/V\). Elsewhere, evaluation on a lowest-type vector has norm at most \(C'\tau_1(x)^{-2N}\le C_1\) for the fixed datum, by smoothing with a compactly supported function fixing the type and counting overlaps. Applying this bound to \(\sum_i\overline{F_i(w)(x)}T_i\) gives \(\sum_i|F_i(w)(x)|^2\le C_1^2\). The integral of the sum over \(\mathscr F\) equals \(D\) by isometry. Hence \((1-e^{-N})D\le VC_1^2\) for every finite subspace. The counted space is therefore finite dimensional. For a full basis the orthogonal projector has kernel \(\sum_iF_i(w)(x)\overline{F_i(w)(y)}\). Its diagonal agrees with (8.53.5) by continuity. Integration over the deep region of volume at most \(V\) gives the bound.
\end{proof}
\leanref{GlobalTheory.\allowbreak{}RowsFinMult.\allowbreak{}rows\_\allowbreak{}finiteDimensional}{Development.\allowbreak{}AbelianVarietyTheory.\allowbreak{}Proofs.\allowbreak{}FlagMeasures.\allowbreak{}FiniteDimensionalRowsSpaceDiscreteSeriesCount}{238--344}
\leanref{GlobalTheory.\allowbreak{}DeepKernel.\allowbreak{}deep\_\allowbreak{}kernel\_\allowbreak{}bound\_\allowbreak{}of\_\allowbreak{}ptBound}{Development.\allowbreak{}AbelianVarietyTheory.\allowbreak{}Proofs.\allowbreak{}FlagMeasures.\allowbreak{}DeepPointKernelBoundOfPointwiseBound}{22--48}
\leanref{GlobalTheory.\allowbreak{}DeepRegion.\allowbreak{}deepBound}{Development.\allowbreak{}AbelianVarietyTheory.\allowbreak{}Proofs.\allowbreak{}FlagMeasures.\allowbreak{}DeepRegionKernelDiagonalBound}{28--150}

\section{Exact class counts}

Let \(U_k(\F_g)\) be the unitriangular subgroup of the column level. Set
\[
 \mathrm{idxColLoc}(g,k)=[\GL_k(\F_g):U_k(\F_g)],\qquad
 \mathrm{idxPi}(D,g_0,g_1)=
     \mathrm{idxColLoc}(g_0,k)^f\mathrm{idxColLoc}(g_1,k)^f .
 \tag{8.54.1}
\]
For a rational prime \(p\), \(\Lean{SplitsF}\,D\,p\) requires \(f\) places above \(p\), each of residue cardinality \(p\).

\begin{theorem}[The class number at constructed points]\label{B8.54}
For every member condition and bundle association, apply the family quantifiers of \ref{B8.39} on \(W_1(\mathcal A)\). At the constructed level for each \(x>F_m.x_1\) and \(i\in F_m.I_x\), assume \(\Lean{SplitsF}\,D\,g_\nu\) for both \(\nu=0,1\), and
\[
 \exists\nu\in\{0,1\},\qquad 2fk<g_\nu-1 .
 \tag{8.54.2}
\]
For every word \(u\), the datum records \(T_u\equiv1\bmod Q_Sqd\). Hence \(T_u\) splits completely in \(E\) and in \(F\). The class number is
\[
 \#\mathrm{cls}_u(0)=M_k\,\mathrm{idxPi}(D,g_0,g_1)(T_u^k-1)^f .
 \tag{8.54.3}
\]
\end{theorem}
\leanref{GlobalTheory.\allowbreak{}Fc.\allowbreak{}R0.\allowbreak{}classCountBound\_\allowbreak{}paramsLambda\_\allowbreak{}constructedPoints}{Development.\allowbreak{}AbelianVarietyTheory.\allowbreak{}Proofs.\allowbreak{}GlobalInterpolation.\allowbreak{}ClassCountBoundParamsLambdaConstructedPoints}{17--25}
\leanref{Formalism.\allowbreak{}MainNotions.\allowbreak{}StandardLattice.\allowbreak{}ClassNumberFormula₂}{Formalism.\allowbreak{}SharpRestriction.\allowbreak{}GlobalSchurBounds.\allowbreak{}StandardLattice}{453--456}
\leanref{Formalism.\allowbreak{}MainNotions.\allowbreak{}J.\allowbreak{}level\_\allowbreak{}prime\_\allowbreak{}splits\_\allowbreak{}totally\_\allowbreak{}real}{Formalism.\allowbreak{}SharpRestriction.\allowbreak{}AdmissibleParameterJoins.\allowbreak{}DIVfromGLz}{13--18}

\begin{proof}
Let \(K\) be the adelic stabilizer of the self-dual overlattice in the definite first space and \(U\subset K\) the column level. The class set is the corresponding double-coset space. The two unitriangular levels and the nonzero vector level at \(T_u\) give
\[
 [K:U]=\mathrm{idxPi}(D,g_0,g_1)(T_u^k-1)^f .
 \tag{8.54.4}
\]
For each genus representative \(\Lambda\), the stabilizers of \(\operatorname{Aut}(\Lambda)\) on the level cosets lie in conjugates of \(U\). At the \(g_\nu\) of (8.54.2) they embed into a pro-\(g_\nu\) group, by the unipotent reduction and the congruence kernel. A nontrivial stabilizer would contain an element of order \(g_\nu\), with rational characteristic polynomial divisible by \(\Phi_{g_\nu}\) of degree \(g_\nu-1>2fk\). Thus the action is free. By the self-dual mass identity the sum of \([K:U]/|\operatorname{Aut}(\Lambda)|\) is \([K:U]M_k\).
\leanref{GlobalTheory.\allowbreak{}ColClassNumber.\allowbreak{}classNumber}{Development.\allowbreak{}AbelianVarietyTheory.\allowbreak{}Proofs.\allowbreak{}FoundationProofs.\allowbreak{}MassSelfDualLatticeFirstSpaceEq}{102--130}
\leanref{GlobalTheory.\allowbreak{}ColClassNumber.\allowbreak{}classNumber\_\allowbreak{}closed}{Development.\allowbreak{}AbelianVarietyTheory.\allowbreak{}Proofs.\allowbreak{}GlobalInterpolation.\allowbreak{}ClassNumberClosed}{10--13}
\end{proof}

\section{Bounded real rows and reciprocity data}

\begin{proposition}[Bounded rows of the indefinite realization]\label{B8.55}
Assume the discrete-series covering, Harish--Chandra module, cusp rapid-decay and Casselman--Wallach inputs of the real theta construction. For every member condition \(Q\) and bundle association \(L_k\), the selected indefinite realization has the following properties at each selected point.

At each \(j<m\), the prescribed highest-weight module contains unit smooth extremal vectors \(u_\rho\) of the specified weights in the lowest \(K\)-types. A unit vector \(w\) in the paired space satisfies
\[
 \operatorname{ioP}_j\!\left(\bigotimes_\rho u_\rho,w\right)=\Delta_j .
\]
Fix the layer and its comparison datum \(O_j\) at \(t=\mu_0\). Its field is \(\iota(LL)\) for the fixed embedding \(\iota:LL\hookrightarrow\C\). The component coefficients \(c_j\) and phase \(\phi_j\) have modulus one. There are a central level element \(z\) and a root of unity \(\zeta_t\in\iota(LL)\) such that each allowed input \(x\) at \(t\) has a cuspidal theta class, and every admissible representative \(f\) satisfies
\[
 \phi_j\,\mathrm{thAd}_t(x)(q_0)
   =\zeta_t\,c_j(\mathrm{comp}(q_0))\,\mathrm{lieFac}_j\,
              \mathrm{jFn}_j(q_0,z\cdot f)
\]
at every \(q_0\in G_{\Q}(V)\backslash G_{\mathbb A}(V)\) for the comparison space \(V\). The dictionary gives \(\mathrm{comp}(q_0)\) from the chosen representative. The scalar \(\mathrm{lieFac}_j\) is the archimedean comparison factor and \(\mathrm{jFn}_j\) evaluates the section.

For each layer lift \((g,k_\ell)\), \(g\) preserves the prescribed cusp space and
\[
 \varphi_j(c_j,gf)(h)=\chi_{\det}(g)\,\varphi_j(c_j,f)(hk_\ell)
\]
for every form \(f\) there and adelic \(h\), with \(\chi_{\det}\) the product of the designated finite determinant characters. Each reference input \(x\) gives finitely many \(a_\nu\in LL\) and algebraic cusp forms \(f_\nu\) of the prescribed weight with
\[
 \phi_j\mathcal T_j(x)(h)
   =\sum_{\nu=1}^s\iota(a_\nu)\,\mathrm{lieFac}_j\,
                          \varphi_j(c_j,f_\nu)(h)
 \tag{8.55.1}
\]
for every \(h\). Use the selected reference vector for \(\mathcal T_j\). At every weight and degree the comparison-field cusp forms are spanned by base-field cusp forms. These scalar-extension identities hold over each common coefficient algebra and are preserved by the level action with the same coefficients.

The guard \(\Lean{QGuard}\,R\) requires the primes of each \(R.Q(t)\) above one rational prime to restrict alike to the quadratic subfield and to avoid \(QS\,q\,\mathrm N(\mathfrak c_K)\). Under it, one \(O'_j\) carries the pins for all tuples over a field containing \(\iota(LL)\). The theta identity holds for every tuple with common \(z\), its own \(\zeta_t\in\iota(LL)\), coefficient \(T.cE(j+1,O'_j.s0)\) and phase \(\phi_j\). The scalar-extension and level-action identities carry over. Formula (8.55.1) uses the reference \(\mu_0\) and the fixed \(\iota\).
\end{proposition}
\leanref{GlobalTheory.\allowbreak{}Fc.\allowbreak{}R1.\allowbreak{}row\_\allowbreak{}hRB\_\allowbreak{}univ\_\allowbreak{}alt}{Development.\allowbreak{}AutomorphicArithmetic.\allowbreak{}ThetaCorrespondence.\allowbreak{}ThetaPairLifts.\allowbreak{}ThirdPlaceSeries.\allowbreak{}RowHRBUniv}{19--24}
\leanref{Formalism.\allowbreak{}IndefR.\allowbreak{}IndefiniteRealisationBoundedRows}{Development.\allowbreak{}AutomorphicArithmetic.\allowbreak{}ThetaCorrespondence.\allowbreak{}ThetaPairLifts.\allowbreak{}ThirdPlaceSeries.\allowbreak{}IndefRRealTB}{39--41}

\begin{proof}
The member datum satisfies
\[
 n_V\ge8fr+4 .
 \tag{8.55.2}
\]
The real theta equality characterizes the realization. Wallach's smooth globalization \cite[Ch.~11]{wallach-1992} and cusp decay for smooth uniformly moderate forms \cite{miller-schmid-2012} apply to these unitary groups. By (8.55.2) an \(LL\)-basis maps to a complex basis at each layer and embedding. Evaluation is equivariant for the character twist, so the orthonormal-frame characterization gives the extremal vectors and delta pairing.
\leanref{GlobalTheory.\allowbreak{}RowsAtRealTBClose.\allowbreak{}rowsB\_\allowbreak{}of\_\allowbreak{}decay\_\allowbreak{}alt}{Development.\allowbreak{}AutomorphicArithmetic.\allowbreak{}ThetaCorrespondence.\allowbreak{}ThetaPairLifts.\allowbreak{}ThirdPlaceSeries.\allowbreak{}RowsBOfDecay}{16--20}

The \(\Lean{FWTie}\) witness gives chart independence and comparison over the two fields at the reference pins. The one-map comparison gives scalar extension, with level action equivariant for the determinant factor. Rationality gives (8.55.1). By (8.55.2) the reference class sums span over the second field, so under the guard the all-tuples comparison extends to this span. Phase compatibility gives the coefficient function and phase of \(O'_j\).
\leanref{GlobalTheory.\allowbreak{}RowsAtRealTB.\allowbreak{}rowsB\_\allowbreak{}of\_\allowbreak{}clauses\_\allowbreak{}alt}{Development.\allowbreak{}AutomorphicArithmetic.\allowbreak{}ThetaCorrespondence.\allowbreak{}ThetaPairLifts.\allowbreak{}ThirdPlaceSeries.\allowbreak{}RowsBOfClauses}{16--19}
\leanref{Formalism.\allowbreak{}IndefR.\allowbreak{}FWT.\allowbreak{}OneMapRowR}{Formalism.\allowbreak{}LocalToGlobal.\allowbreak{}IndefiniteReal.\allowbreak{}OneMapT}{151--155}
\end{proof}

\begin{proposition}[Reciprocity supplies the twisting and kernel conditions]\label{B8.56}
Assume global Artin reciprocity, admissible extension parameters \(p\) with local Langlands input, and a bundle association \(L_k\) compatible with the reciprocity record. The member condition \(\mathrm{IsGoodE}(p.s)\) supplies \(\Lean{SelectionServed}\).

Use the family quantifiers of \ref{B8.39} on \(\mathcal D_3\). Write \(D=F_m.D(x,i)\), \(D_x=F_m.X(x,i)\), \(X_G=F_m.Gx(x,i)\), \(\mathcal L=F_m.L(x,i)\). Beyond one family threshold, for every \(x>F_m.x_1\) and \(i\in F_m.I_x\), local data \(L_C\) satisfy the two conditions of \ref{B4.16},
\[
 \begin{gathered}
 \mathrm{CmpV}(\mathrm{cmp},D_x,\mathrm{tw},\mathrm{reW}(L_C)),\\
 \mathrm{SysTw}(D,\mathcal L,X_G,G_K,D_x,\mathrm{reW}(L_C),
                          \mathrm{lev},\mathrm{errX},c_X.\mathrm{err},1).
 \end{gathered}
 \tag{8.56.1}
\]
Here \(\mathrm{reW}(L_C)\) uses the local chains on the given filling. For \(x>F_m.x_1\) with \(x\ge(\mathrm{sel}\,F_m).1\) and \(i\in F_m.I_x\), the following hold. For each word and tuple, at every level-free \(w\) the adjacent class scalars have \(p\)-adic norm one for \(p=\operatorname{resChar}(w)\). For each used \(\mu\) and word \(u\), every adjacent class-pair map at each level-free \(w\) preserves the rational row spaces. The reference chain and every used twisted chain have \(\mathrm{Fi}_j\to\mathrm{Fi}_{j+1}\) for each lift and trivial initial stabilizers.

One collection \((e_w,\beta_w,\sigma_w)_w\) serves the following assertions at each level-free \(w\). If
\[
 (\mathrm{IsResidueClassZero}(D,p)\land\mathrm{HEv}(D,p))
       \lor\mathrm{IsSg}(D,p),
\]
then every admissible representative set \(R\) satisfies \(\mathrm{BeOK}(R)\) for these data. In this \(R\), at \(j<m\) with \(N_jN_{j+1}\) even the entries have the all-place theta expansion. At residue class zero with \(h_P=0\) they have the real-place expansion, and at \(\Lean{IsCs}\) they satisfy the all-real condition. For \(p\ne d\) these layers satisfy \(\mathrm{EntGe}(\mathcal E,u,\mu,w,j,0)\) for the selected entries \(\mathcal E\).

At level-free \(w\) of class \(\Lean{ramD}\), each entry at \(j<m\) has the form
\[
 \Theta=\exp\!\left(-S\sum_t u_t\sum_{a,b}
       X_{(0,t),a,b}X_{(1,t),a,b}\right)\Theta_0,\qquad
 |S|_p\le p^{1/2},\quad |u_t|_p\le1,\quad
 \left|\left(\prod_\alpha\nu_\alpha!\right)[X^\nu]\Theta_0\right|_p\le1
 \tag{8.56.2}
\]
for every multi-index \(\nu\), where \(p=\operatorname{resChar}(w)\).

Let \(p=\operatorname{resChar}(w)\) satisfy \(5\le p\le X''\), \(n_P(D,p)\) even and at least \(2\), residue class zero and \(h_P(D,p)\ne0\). For each \(u\), \(j<m\), \(t\in\mathrm{reps}(X_G,p)\), and each \(\Theta\) with the prescribed expansion at \((0,t)\), all coefficients of \(\operatorname{applyAt0}(\Theta,G_{xz}(I))\) have norm at most \(p^{-\max(g_0(t,j+1),0)}\) for each column choice \(I\). Here \(g_0\) is the ordinary gain of \ref{B3.48}.
\end{proposition}
\leanref{GlobalTheory.\allowbreak{}T.\allowbreak{}hServ\_\allowbreak{}of\_\allowbreak{}linkR}{Development.\allowbreak{}AutomorphicArithmetic.\allowbreak{}ThetaCorrespondence.\allowbreak{}ThetaPairLifts.\allowbreak{}ThirdPlaceSeries.\allowbreak{}RowHServOfLinkStatement}{14--15}
\leanref{GlobalTheory.\allowbreak{}Cp.\allowbreak{}SelectionServed}{Development.\allowbreak{}AutomorphicArithmetic.\allowbreak{}ThetaCorrespondence.\allowbreak{}ThetaPairLifts.\allowbreak{}ThirdPlaceSeries.\allowbreak{}CountsIrreducibilityConstants}{48--62}

\begin{proof}
The CM model supplied by reciprocity has algebraic torsion Taylor coefficients. The lattice and curve dictionary identifies its standard kernel with the selected kernel. At every level-free place the adjacent class scalars have norm one. For each used \(\mu\) and word \(u\), the class-pair comparison at each level-free \(w\) gives maps \(K_j(c,c')\) on the rational row spaces. The dictionary also gives the transfer and stabilizer identities. The inert model gives the parity expansions. The Taylor model gives integral jets and the quadratic factor (8.56.2). The resulting discriminant entry bound is in \ref{B4.16a}.
\leanref{GlobalTheory.\allowbreak{}Svd.\allowbreak{}servH\_\allowbreak{}row\_\allowbreak{}of\_\allowbreak{}artin}{Development.\allowbreak{}AbelianVarietyTheory.\allowbreak{}Proofs.\allowbreak{}SatakeParameters.\allowbreak{}BundleBothPartsOfStandardKernelOfGlobalArtinReci}{15--24}
\leanref{Formalism.\allowbreak{}BasicNotions.\allowbreak{}KerEnt.\allowbreak{}UnUnit}{Formalism.\allowbreak{}LocalRingCardinality.\allowbreak{}DataKerI}{129--130}
\leanref{Formalism.\allowbreak{}MainNotions.\allowbreak{}OrbitPlaces.\allowbreak{}KmapRat}{Formalism.\allowbreak{}SharpRestriction.\allowbreak{}Divisibility.\allowbreak{}OrbitPlaces}{1367--1369}
\leanref{Formalism.\allowbreak{}MainNotions.\allowbreak{}OrbitPlaces.\allowbreak{}JetIntegralityBis}{Formalism.\allowbreak{}SharpRestriction.\allowbreak{}Divisibility.\allowbreak{}OrbitPlaces}{339--341}
\leanref{Formalism.\allowbreak{}MainNotions.\allowbreak{}OrbitPlaces.\allowbreak{}ShapeD}{Formalism.\allowbreak{}SharpRestriction.\allowbreak{}Divisibility.\allowbreak{}OrbitPlaces}{343--349}

For the weight shape in \ref{B4.16}, write \(w_z=w_{\rm Pl}(D,j+1,z)\), with \(s_z,r_{z,\nu},c_z\) the size, row entries and cap at \(u,j,z\). Then \(w_z>0\). Off the top, \(s_z-(k-1)w_z\in[0,4w_z]\). At the top, \(r_{z,0}\ge2w_z\) and \(s_z-r_{z,0}-(k-3)w_z\in[0,4w_z]\). The bounds with size scales \(a,b\) include
\[
 \begin{gathered}
 s_z\le a^2,\qquad 2\log_5(a^2)+2\le k-1,\qquad X''\le a,\quad a>1,\\
 c_z\le b+\eta_{\rm Sz}+1+3\mathcal L.\mathrm{Sc},\qquad
 c_z\le\mathrm{fillG}(X_G,b_L,j+1,z)+c_G(\mathcal L),\\
 \mathrm{al}(\mathcal L,a_L,t)\le r_{(0,t),0}-r_{(0,t),1}\qquad(t\in TT).
 \end{gathered}
\]
At a selected point, \(x=\mathcal L.\mathrm{sc}.x\ge x_{\rm shape}(\varepsilon)\) ensures these bounds for each candidate entry datum. It also gives \(r_{(0,t),0}-r_{(0,t),1}\le2a\) for all \(u,j<m,t\in TT\) and \(\mathrm{splitBound}\ge3\). With \(d_{\ell1}\ge5\) and \(d_{\ell2}\ge4\), the orbit-letter estimates of \ref{B4.16} apply to the prescribed expansion.
\leanref{Formalism.\allowbreak{}MainNotions.\allowbreak{}OrbitPlaces.\allowbreak{}OrbitBoundWithoutOperator}{Formalism.\allowbreak{}SharpRestriction.\allowbreak{}Divisibility.\allowbreak{}OrbitPlaces}{333--337}
\leanref{GlobalTheory.\allowbreak{}Pdv.\allowbreak{}rowAt\_\allowbreak{}of\_\allowbreak{}stdTie}{Development.\allowbreak{}AbelianVarietyTheory.\allowbreak{}Proofs.\allowbreak{}SatakeParameters.\allowbreak{}RowAtOfStdTie}{13--36}

For (8.56.1) take the maximum of the shape threshold and the threshold for \(2r_G\le k\). For every used \(\mu\) and word \(u\), apply \ref{B8.56c} at each \(w\) with \(\Lean{GIdx}\) and \(p=\operatorname{resChar}(w)\notin Q'(\mu)\). The joint compares the same twist and chain on its given filling. It intertwines the lifts and transfers their terminal valuation. Each unrepresented position \(v\) belongs to a second-side orbit \((1,t_v)\). For \(p\le X''\) take \(f_{v,j}=F_{\rm Orb}(j,p,(1,t_v))-\mathrm{err}_F(j,p,(1,t_v))\). In the one-place letter coordinates the unit is \(1\) and the sign is \(-1\). The orbit estimate with constants \(11,1\) gives \(\operatorname{LowAt}(j-1,v,f_{v,j})\) for this system's pairing. For \(p>X''\) take \(f_{v,j}=0\) by \(\Lean{adm11}\). The summand is invariant under cyclic shift of \(t_v\). Summing once per complementary orbit and using \(\sum_{j<m}d_G(j)\le(m-1)v_p(D_{\Gamma,u})\) gives (4.11).
\leanref{GlobalTheory.\allowbreak{}TwU.\allowbreak{}SiteSys.\allowbreak{}sysTw\_\allowbreak{}sites}{Development.\allowbreak{}AbelianVarietyTheory.\allowbreak{}Proofs.\allowbreak{}SatakeParameters.\allowbreak{}LocalSystemsLowerBoundsSecondRealPlace}{48--73}
\leanref{GlobalTheory.\allowbreak{}Svd.\allowbreak{}rowTwU}{Development.\allowbreak{}AbelianVarietyTheory.\allowbreak{}Proofs.\allowbreak{}SatakeParameters.\allowbreak{}RowTwUHolds}{15--25}

The linkage and stabilizer-size hypothesis give \(\Lean{SelectionServed}\) in \ref{B4.16} from these two rows.
\leanref{GlobalTheory.\allowbreak{}row\_\allowbreak{}hServ\_\allowbreak{}of\_\allowbreak{}link\_\allowbreak{}alt}{Development.\allowbreak{}AutomorphicArithmetic.\allowbreak{}ThetaCorrespondence.\allowbreak{}ThetaPairLifts.\allowbreak{}ThirdPlaceSeries.\allowbreak{}RowHServOfLink}{17--23}
\leanref{GlobalTheory.\allowbreak{}Fc.\allowbreak{}R1.\allowbreak{}row\_\allowbreak{}hServ\_\allowbreak{}of\_\allowbreak{}kerServ\_\allowbreak{}alt}{Development.\allowbreak{}AutomorphicArithmetic.\allowbreak{}ThetaCorrespondence.\allowbreak{}ThetaPairLifts.\allowbreak{}ThirdPlaceSeries.\allowbreak{}RowHServOfKerServ}{13--22}
\end{proof}

\begin{lemma}[The corrected Taylor series in an orthonormal basis]\label{B8.56a}
Use the CM model of \ref{B8.56}. Let \(p\ge5\) be prime with \(\bigl(\frac{-d}{p}\bigr)=-1\). Take the theta datum for an ideal \(\mathfrak c\ne0\) of norm prime to \(2d\) and a twelfth root of the discriminant of its complex lattice. The datum uses a character \(\chi\) multiplicative on norm-one ideals with \(\chi((z))=\prod_\tau\tau(z)^{-1}\) for unitary \(z\). Fix \(\lambda:\overline{\Q}\hookrightarrow\C_p\). Let \(L,L'\) be self-dual lattices of even ranks. Let \(F\) be a finite \(\mathcal O_{\C_p}\)-linear combination of the \(\lambda\)-Taylor series of prime-to-\(p\) torsion translates of the theta sections for self-dual pairs locally equal to \(L,L'\) at \(p\).

For the period lattice \(\Lambda\), take any integral Weierstrass model \(W=C\cdot W_\Lambda\) of discriminant one. Write \(l_W\) for its normalized formal logarithm and put \(e=C.u^{-1}\). Set \(A(T)=\sigma_\Lambda^*(e\,l_W(T))/(eT)\), where \(\sigma_\Lambda^*\) is the normalized odd sigma series with coefficients embedded by \(\lambda\). For any basis \((B_i)_{i=1}^r\) at \(p\) of \(\mathfrak a_0\operatorname{Hom}(L,L')\) with its fixed ideal factor, orthonormal for \(\operatorname{Tr}_{E/K}\sum_{a,b}x_{ab}\overline{y_{ab}}\), one has
\[
 F\left(\left(e\sum_i\lambda(\tau(B_{i,ba}))l_W(T_i)\right)_{\tau,b,a}\right)\prod_i A(T_i)^{-1}\in\mathcal O_{\C_p}[[T_1,\ldots,T_r]].
\]
\end{lemma}
\leanref{Formalism.\allowbreak{}L2.\allowbreak{}KerInert.\allowbreak{}kernels\_\allowbreak{}integral\_\allowbreak{}inert}{Formalism.\allowbreak{}LocalToGlobal.\allowbreak{}InertKernelTheory.\allowbreak{}KernelsIntegralInert}{7--9}
\begin{proof}
For the elliptic factor, first translate the odd sigma function by \(t\) of exact order \(n\) in \(\C/\Lambda\), with \(n>1\) and \(p\nmid n\). Write the corrected translate as \(G_t=cH\), \(H(0)=1\). Iterate the sigma addition formula up to \(nt\in\Lambda\) to obtain, up to a unit scalar,
\[
 G_t^n\prod_{j=1}^{n-1}\{c^2(x(t)-x(T+jt))\}^{n-j}=T^n .
\]
Each abscissa difference is a power of \(T\) times an integral unit. The order is zero for distinct abscissas. For equal abscissas it is one away from two-torsion and two at a nonzero two-torsion point. These orders follow from the smooth Weierstrass equation and injectivity of reduction on prime-to-\(p\) torsion. Leading coefficients give \(|c|_p=1\) and \(H^n\in1+T\mathcal O_{\C_p}[[T]]\). The degree-\(k\) coefficient of \(H^n\) is \(n[T^k]H\) plus an integral polynomial in earlier coefficients. Since \(p\nmid n\), induction gives integrality of \(H\). At the origin the sigma series in the translation system is a unit times \(T\).

Choose a nonzero half-period \(h\) and put \(\theta_h=U_h\sigma^*/(U_h\sigma^*)(0)\), where \(U_t\) is the theta translation operator. The leading coefficient calculation gives \(|(U_h\sigma^*)(0)|_p=1\). If \(s\) has order \(N\) prime to \(p\), then \(2N(s+h)\in\Lambda\), and \(U_s\theta_h\) is a root of unity times \(U_{s+h}\sigma^*/(U_h\sigma^*)(0)\). Thus the corrected translates of \(\theta_h\) are integral. Passing to the height-two parameter substitutes an integral series of zero constant term. The exponential of the difference of the two phases is an integral unit.
\leanref{Formalism.\allowbreak{}L2.\allowbreak{}Th1.\allowbreak{}E.\allowbreak{}joint\_\allowbreak{}integrality}{Formalism.\allowbreak{}LocalToGlobal.\allowbreak{}TorsionIntegrality.\allowbreak{}TorsionTranslationPowers.\allowbreak{}NonsingularJetNorms.\allowbreak{}Joints}{128--129}
\leanref{Formalism.\allowbreak{}L2.\allowbreak{}Th1.\allowbreak{}E.\allowbreak{}Th0Analytic\_\allowbreak{}of}{Formalism.\allowbreak{}LocalToGlobal.\allowbreak{}TorsionIntegrality.\allowbreak{}TorsionTranslationPowers.\allowbreak{}NonsingularJetNorms.\allowbreak{}Taylor0ConstMul}{195--269}

In isometric coordinates the four-copy theta identity writes the fourth tensor power as a sum of products of elliptic translates. The denominator is the square root of a prime-to-\(p\) index, and character and torsion multipliers have norm one. The coordinate phases polarize against the identity Gram matrix, so the corrected bound is preserved under changes between orthonormal bases of the same localized lattice. For a coefficient \(a_\nu\) of one corrected copy the tensor coefficient at \((\nu,\nu,\nu,\nu)\) is \(a_\nu^4\). Hence \(|a_\nu|_p\le1\). The four entries give a standard \(2\)-by-\(2\) block. Products give the even-rank pairs. By weak approximation, unitary isometries move the two lattices to lattices standard at \(p\), whose sections expand in standard translates with prime-to-\(p\) denominators. Pullback transports the orthonormal basis and the torsion points. Sum with the given integral coefficients.
\leanref{Formalism.\allowbreak{}L2.\allowbreak{}KerInert.\allowbreak{}integral\_\allowbreak{}inert\_\allowbreak{}of\_\allowbreak{}weakApprox}{Formalism.\allowbreak{}LocalToGlobal.\allowbreak{}InertKernelTheory.\allowbreak{}IntegralInertOfWeakApprox}{8--13}
\end{proof}

\begin{lemma}[The change to paired variables]\label{B8.56b}
In the unramified CM model at \(p\ge5\), write \(P\) for the one-variable phase in height-two parameters and \(\ell\) for the formal logarithm. In the short model \(P(T)=-\log(\sigma^*(\epsilon\ell(T))/(\epsilon T))\), where \(\epsilon\) is the coordinate unit. Put \(x+_F y=\ell^{-1}(\ell(x)+\ell(y))\) and \(B_P(x,y)=P(x+_F y)-P(x)-P(y)\). Write \(A\sim B\) when their difference has zero constant term and integral exponential. The CM pairing satisfies \(b^\ell\sim B_P\), with the superscript denoting logarithmic substitution. Let \(f_e,f'_e\) be vectors of length \(r\), indexed by a finite set \(I\), with integral unramified entries fixed by the square of Frobenius \(\sigma\). For \(\langle u,v\rangle=\sum_i\sigma(u_i)v_i\), suppose
\[
 \langle f_e,f_h\rangle=\langle f'_e,f'_h\rangle=0,\qquad \langle f_e,f'_h\rangle=\delta_{eh}.
\]
Put \(L_i=\sum_e(f_{ei}X_e+f'_{ei}Y_e)\), \(S=\sum_e b(X_e,Y_e)\). If \(H(T)=\exp(\sum_iP(T_i))\Theta((\ell(T_i))_i)\) is integral, then \(U=\Theta(L)\exp S\) is integral after logarithmic substitution.
\end{lemma}
\leanref{Formalism.\allowbreak{}L2.\allowbreak{}KerInert.\allowbreak{}Rep.\allowbreak{}k1gen\_\allowbreak{}of\_\allowbreak{}orth}{Formalism.\allowbreak{}LocalToGlobal.\allowbreak{}InertKernelTheory.\allowbreak{}InertCoprimeLemmas.\allowbreak{}RepAsm}{17--24}
\begin{proof}
Put \(z_i=\ell^{-1}(L_i(\ell(X),\ell(Y)))\). These are formal sums of integral endomorphisms \([a]=\ell^{-1}(a\ell)\). Modulo \(\sim\), the CM norm identity gives \(P([a]x)=a\sigma(a)P(x)\), and its polarization is additive for \(+_F\) with \([a]\) adjoint to \([\sigma(a)]\). Expand \(\sum_iP(z_i)\) by the definition of \(B_P\). The diagonal terms from the \(f\)'s have norm coefficients \(\langle f_e,f_e\rangle=0\). The off-diagonal terms vanish by \(\langle f_e,f_h\rangle=0\). The same computation applies to \(f'\). The cross terms sum to \(\sum_{e,h} B_P(X_e,[\langle f_e,f'_h\rangle]Y_h)=\sum_eB_P(X_e,Y_e)\). Thus \(\sum_iP(z_i)\sim S^\ell\). In the full variable set,
\[
 U^\ell=H(z)\exp\!\left(S^\ell-\sum_iP(z_i)\right).
\]
Both factors have integral coefficients, and \(z_i(0)=0\). Write \(S^-(X,Y)=S(-X,-Y)\). The pairing laws give \((S^-)^\ell\sim S^\ell\). Thus \(\Phi=S\) and \(U=\Theta(L)\exp S\) satisfy the generator condition.
\leanref{Formalism.\allowbreak{}L2.\allowbreak{}KerInert.\allowbreak{}K1.\allowbreak{}hyperbolic\_\allowbreak{}of\_\allowbreak{}orthogonal}{Formalism.\allowbreak{}LocalToGlobal.\allowbreak{}InertKernelTheory.\allowbreak{}InertCoprimeLemmas.\allowbreak{}ConstantCoeffFormalSum}{402--439}
\leanref{Formalism.\allowbreak{}L2.\allowbreak{}KerInert.\allowbreak{}K1.\allowbreak{}sum\_\allowbreak{}beta\_\allowbreak{}even}{Formalism.\allowbreak{}LocalToGlobal.\allowbreak{}InertKernelTheory.\allowbreak{}InertCoprimeLemmas.\allowbreak{}ConstantCoeffFormalSum}{697--714}
\end{proof}

\begin{proposition}[The kernel on the selected chain]\label{B8.56c}
Let \(h_g\) witness \(\Lean{GIdx}\) for the word \(u\) at \(p=\operatorname{resChar}(w)\). Take the joint comparing the selected theta kernel and the local chain on its given filling. Assume \(\Lean{CurveEx}\) and \(\Lean{AdmAtEnt2}\), supplied by the CM model of \ref{B8.56}. For a used twist, trivial initial stabilizers and \(2r_G\le k\), there is an admissible system on the levels and chain of this joint, with \(d_G(0)=0\) and \(d_G(j)\le v_p(D_{\Gamma,u})\) for all \(j\).
\end{proposition}
\leanref{GlobalTheory.\allowbreak{}TwV.\allowbreak{}site\_\allowbreak{}cmpLite\_\allowbreak{}siteAdm}{Development.\allowbreak{}AbelianVarietyTheory.\allowbreak{}Proofs.\allowbreak{}SatakeParameters.\allowbreak{}JointSiteChainComparedCarriesAdmittedSystem}{77--89}
\begin{proof}
On the chosen joint, \(\Lean{CurveEx}\) gives the local curve and its height-two formal isomorphism. Apply \(\Lean{AdmAtEnt2}\) to this joint datum to obtain the admitted kernel on the same levels and chain.
\leanref{GlobalTheory.\allowbreak{}TwV.\allowbreak{}Lt.\allowbreak{}exists\_\allowbreak{}genusOne\_\allowbreak{}formal\_\allowbreak{}iso\_\allowbreak{}of\_\allowbreak{}unramified\_\allowbreak{}good\_\allowbreak{}reduction}{Development.\allowbreak{}AbelianVarietyTheory.\allowbreak{}Proofs.\allowbreak{}SatakeParameters.\allowbreak{}ExistsGenusOneFormalIsoOfUnramifiedGoodReduction}{17--24}
\leanref{GlobalTheory.\allowbreak{}TwV.\allowbreak{}Lt.\allowbreak{}inert\_\allowbreak{}theta\_\allowbreak{}admissibility\_\allowbreak{}of\_\allowbreak{}complex\_\allowbreak{}multiplication}{Development.\allowbreak{}AbelianVarietyTheory.\allowbreak{}Proofs.\allowbreak{}SatakeParameters.\allowbreak{}InertThetaAdmissibilityOfComplexMultiplication}{8--10}

At a source representative \(F=p^{e_j}\operatorname{Taylor}_\lambda(a_jG)\) before change of frame. The theta sum \(G\) uses the transported self-dual pair of even ranks, with level and character scalar \(a_j\). The finite transform has prime-to-\(p\) denominators and \(|p^{e_j}\lambda(a_j)|_p\le1\). Thus \ref{B8.56a} applies.

Before rescaling the good curve model, \(v=l_W^{-1}(e_{\rm LT}^{-1}\ell)\) is integral with \(v'(0)=e_{\rm LT}^{-1}\), where \(e_{\rm LT}\) is the unit of the height-two comparison. Substitute \(l_W(v)=e_{\rm LT}^{-1}\ell\) in the sigma quotient. The logarithmic phases differ up to sign by \(\log(e_{\rm LT}v(T)/T)\), whose exponential and inverse are integral. The coboundary gives \(b^\ell\sim B_P\). The parameter change of \ref{B8.56a} and one such unit per coordinate give \(H\) in \ref{B8.56b}, with frame scale \((g.C.u)^{-1}e_{\rm LT}^{-1}\).

Expand the framed matrix units in the orthonormal basis. The resulting coordinates \(f_e,f'_e\) are integral since the basis maps and their adjoints respect the two local lattices. Frobenius cycles the embeddings, so its square fixes the even and odd dual-basis sums. The trace pairing of units on the same side is zero. On opposite sides it is the matrix Kronecker delta times the Kronecker deltas for the place and the trace-dual basis. These are the Gram identities in \ref{B8.56b}. Reconstruction in the basis gives \(\Theta(L)\) for each frame. Transport to the source representative moves the frame and multiplies the series by a character value of norm one. This gives \(\Lean{adm2b}\) in every frame. Restriction to one place commutes with \(\ell\) and exponential. Since \(b(0,0)=0\), the same \(b\) gives \(\Lean{adm2}\).
\leanref{Formalism.\allowbreak{}L2.\allowbreak{}KerInert.\allowbreak{}RepLt.\allowbreak{}K1\_\allowbreak{}representatives\_\allowbreak{}LT\_\allowbreak{}of\_\allowbreak{}dictionary}{Formalism.\allowbreak{}LocalToGlobal.\allowbreak{}InertKernelTheory.\allowbreak{}InertCoprimeLemmas.\allowbreak{}K1RepresentativesLTOfDictionary}{253--284}
\leanref{Formalism.\allowbreak{}L2.\allowbreak{}KerInert.\allowbreak{}KGLt.\allowbreak{}first\_\allowbreak{}kernel\_\allowbreak{}law\_\allowbreak{}global\_\allowbreak{}small}{Formalism.\allowbreak{}LocalToGlobal.\allowbreak{}InertKernelTheory.\allowbreak{}InertCoprimeLemmas.\allowbreak{}FirstKernelSmallOfFrameTwist}{124--133}

The CM pairing laws give \(\Lean{adm2c}\). Their two additivity relations give factorial-integral jets with at least two derivatives in one variable. For the remaining mixed derivative use a Teichm\"uller scalar \(a\) with \(a-\sigma a\) a unit in the conjugation relation. The product rule gives \(\Lean{adm11}\). Expanding the integral multiplier as in (3.29.5) with all neighbour indices zero gives \(\Lean{adm3}\) for every output coefficient. Rational comparison gives \(\Lean{adm3b}\). The joint gives the two class sums of \ref{B3.30}, with the neighbour identity from the terminal class decomposition. The containment indicator gives the additive map in \(\Lean{adm5}\), and level equivariance gives the stabilizer identity. Counting oriented edges by their stabilizers gives \(\Lean{adm6}\). Bases adapted to each lower neighbour give the weight-lattice test of \(\Lean{adm7}\). The source Hecke operators commute and preserve their lattice. With \ref{B3.06} these give Definition~\ref{B3.05a}. Trivial initial stabilizers and the common stabilizer denominator give the stated losses.
\leanref{Formalism.\allowbreak{}L2.\allowbreak{}KerInert.\allowbreak{}AdmLoc2.\allowbreak{}admissible\_\allowbreak{}general\_\allowbreak{}of\_\allowbreak{}kernel\_\allowbreak{}law}{Formalism.\allowbreak{}LocalToGlobal.\allowbreak{}InertKernelTheory.\allowbreak{}InertCoprimeLemmas.\allowbreak{}AdmissibleEntriesOfKernelLaw}{62--73}
\end{proof}


% ---- p8c6 ----
\chapter{The chain length}\label{ch:p8c6}

Fix a scale $s=(x,\varepsilon,J)\in\R^3$ with $x>1$,
$0<\varepsilon<1/52$ and $J\ge1$. Put
\[
 L=x^\varepsilon,\qquad \delta_1=x^{-\varepsilon/8},\qquad
 \alpha=\sqrt{x\delta_1}=x^{1/2-\varepsilon/16},\qquad
 \beta=x^{9/2+10\varepsilon}.
\]

\phantomsection\label{B8.57}
For $y\in\R$ let $\lceil y\rceil_+=\max(0,\lceil y\rceil)$ and set
$e(y)=2\lceil y/2\rceil_+$. The member rank condition sets the
chain length, the number of lifts, to $m=m_*(s)=e(\alpha)$.
Theorem~\ref{B8.59} uses the same rounding formula to choose $m$
and the starting rank $k$ on a prescribed frame.
\leanref{GlobalTheory.\allowbreak{}Mem.\allowbreak{}Num.\allowbreak{}evenCeil}{Development.\allowbreak{}AbelianVarietyTheory.\allowbreak{}Definitions.\allowbreak{}AutomorphicMembership.\allowbreak{}EvenCeil}{11--11}
\leanref{GlobalTheory.\allowbreak{}Mem.\allowbreak{}Num.\allowbreak{}mAt}{Development.\allowbreak{}AbelianVarietyTheory.\allowbreak{}Definitions.\allowbreak{}AutomorphicMembership.\allowbreak{}EvenCeil}{21--21}
\leanref{Formalism.\allowbreak{}BasicNotions.\allowbreak{}Scale.\allowbreak{}delta1}{Formalism.\allowbreak{}LocalRingCardinality.\allowbreak{}DataParams₂}{45--45}

For $y>0$ we have $y/2\le\lceil y/2\rceil_+<y/2+1$,
so $y\le e(y)<y+2$. Every even natural number at least $y$
is $2a$ with $a\ge\lceil y/2\rceil_+$. Thus $m_*(s)$ is the
least even natural number at least $\alpha$, with
$\alpha\le m_*(s)<\alpha+2$.

\phantomsection\label{B8.58}
Let $C$ be the family constant package, and write
$c_W=\Lean{C.\allowbreak{}P.\allowbreak{}CnW}$, $c_T=\Lean{C.\allowbreak{}T.\allowbreak{}cTO}$ and
$c_o=\Lean{C.\allowbreak{}P.\allowbreak{}Cpde}$. For $n\in\N$ put $f=2n$ and define
\[
 \begin{aligned}
 k_*&=e(\beta),& r_*&=k_*+m_*-1,\\
 w_C&=8\left\lceil\frac{c_Wr_*}{8}
       \left(f+\frac{f^2(k_*-1)^3L^2}{m_*}\right)\right\rceil_+,
 &t_C&=\lceil x^{c_T}\rceil_+w_C+8fr_*+4-r_* .
 \end{aligned}
\]
Subtraction in $\N$ is truncated at zero.
For a rank record $R$ on a frame $\varphi$ with finite abelian group
$A$ and distinguished element $a_0$, put $n=|A|$ and
$o_\varphi=[A:\langle a_0\rangle]$. The condition
$\Lean{RChosen}(C,s,R)$ is
\[
 x^{2+8\varepsilon}\le n\le2x^{2+8\varepsilon},\qquad
 R.k=k_*,\quad R.m=m_*,\quad R.t=t_C,\quad o_\varphi\le c_o.
\]
The condition $\Lean{RanksOK}(C,s,n)$ requires $t_C$ to be odd and
\[
 \begin{gathered}
 2\le m_*,\quad m_*+3\le k_*,\quad r_*+5\le t_C,\quad
 3m_*<t_C,\\
 m_*\le f\le k_*,\qquad
 m_*^2\le k_*+m_*(t_C+1).
 \end{gathered}
\]
For every frame $\varphi$ of order $n$ satisfying
$\Lean{RanksOK}(C,s,n)$, $\Lean{ranksOf}$ extends $\varphi$ by
$(k_*,m_*,t_C)$. Even rounding gives the two required even ranks.
\leanref{GlobalTheory.\allowbreak{}Mem.\allowbreak{}Num.\allowbreak{}RChosen}{Development.\allowbreak{}AbelianVarietyTheory.\allowbreak{}Definitions.\allowbreak{}AutomorphicMembership.\allowbreak{}EventuallyThickCardLeCyclicOrder}{60--62}
\leanref{Formalism.\allowbreak{}HC.\allowbreak{}Frame.\allowbreak{}o}{Formalism.\allowbreak{}SharpRestriction.\allowbreak{}RamificationFrobenius.\allowbreak{}TotientGaloisClasses.\allowbreak{}ZsmulGeneratorEqZeroOfProjection}{70--70}
\leanref{GlobalTheory.\allowbreak{}Mem.\allowbreak{}Num.\allowbreak{}kAt}{Development.\allowbreak{}AbelianVarietyTheory.\allowbreak{}Definitions.\allowbreak{}AutomorphicMembership.\allowbreak{}EvenCeil}{23--23}
\leanref{GlobalTheory.\allowbreak{}Mem.\allowbreak{}Num.\allowbreak{}nWAt}{Development.\allowbreak{}AbelianVarietyTheory.\allowbreak{}Definitions.\allowbreak{}AutomorphicMembership.\allowbreak{}EvenCeil}{27--27}
\leanref{GlobalTheory.\allowbreak{}Mem.\allowbreak{}Num.\allowbreak{}tAt}{Development.\allowbreak{}AbelianVarietyTheory.\allowbreak{}Definitions.\allowbreak{}AutomorphicMembership.\allowbreak{}EvenCeil}{33--33}
\leanref{GlobalTheory.\allowbreak{}Mem.\allowbreak{}Num.\allowbreak{}RanksOK}{Development.\allowbreak{}AbelianVarietyTheory.\allowbreak{}Definitions.\allowbreak{}AutomorphicMembership.\allowbreak{}EvenCeil}{41--43}
\leanref{GlobalTheory.\allowbreak{}Mem.\allowbreak{}Num.\allowbreak{}ranksOf}{Development.\allowbreak{}AbelianVarietyTheory.\allowbreak{}Definitions.\allowbreak{}AutomorphicMembership.\allowbreak{}EvenCeil}{45--59}
\leanref{Formalism.\allowbreak{}BasicNotions.\allowbreak{}Ranks}{Formalism.\allowbreak{}LocalRingCardinality.\allowbreak{}Params}{6--23}

\begin{theorem}[Ranks at a window point]\label{B8.59}
For each $0<\varepsilon<1/52$ one may take $x_0=16$.
For every scale $s$ with this $\varepsilon$ and $x\ge x_0$, and every
arithmetic frame $\varphi$ of order $n$ satisfying
$x^{2+8\varepsilon}\le n\le2x^{2+8\varepsilon}$, there is a rank record
$R$ extending $\varphi$ such that
\[
 \alpha\le m\le2\alpha,\qquad
 \beta\le k\le2\beta,\qquad
 r+t=8fr+4+TW,
\]
where $m=R.m$, $k=R.k$, $t=R.t$, $f=2n$, $r=k+m-1$, and
\[
 T_O=T=\lceil x^{377}\rceil_+,\qquad
 n_W=W=8\left\lceil\frac{10r}{8}
       \left(f+\frac{f^2(k-1)^3L^2}{m}\right)\right\rceil_+.
\]
The bounds on $n,m,k$, together with the identity for $r+t$,
constitute $\Lean{RankVal}(s,R)$.
\end{theorem}
\leanref{GlobalTheory.\allowbreak{}EX.\allowbreak{}exists\_\allowbreak{}ranks\_\allowbreak{}rankVal\_\allowbreak{}frame}{Development.\allowbreak{}AbelianVarietyTheory.\allowbreak{}Proofs.\allowbreak{}GlobalInterpolation.\allowbreak{}ExistsRanksRankValFrame}{45--99}
\leanref{GlobalTheory.\allowbreak{}EX.\allowbreak{}TOv}{Development.\allowbreak{}AbelianVarietyTheory.\allowbreak{}Definitions.\allowbreak{}GlobalTwisting.\allowbreak{}PairRowIffQdRows}{17--17}
\leanref{GlobalTheory.\allowbreak{}EX.\allowbreak{}nWv}{Development.\allowbreak{}AbelianVarietyTheory.\allowbreak{}Definitions.\allowbreak{}GlobalTwisting.\allowbreak{}PairRowIffQdRows}{19--20}
\leanref{GlobalTheory.\allowbreak{}EX.\allowbreak{}RankVal}{Development.\allowbreak{}AbelianVarietyTheory.\allowbreak{}Definitions.\allowbreak{}GlobalTwisting.\allowbreak{}PairRowIffQdRows}{28--32}
\begin{proof}
By the ceiling bound above, $a\le e(a)<a+2\le2a$ for $a\ge2$.
Here $2\le\alpha\le x$ and $\beta\ge x^4$.
Indeed, $1/2\le1-\varepsilon/8\le2$ and $x\ge16$ give
$4\le x^{1-\varepsilon/8}\le x^2$.
Choose $m=e(\alpha)$ and $k=e(\beta)$.
Since $2\le2+8\varepsilon\le3$, the order bounds give
$x^2\le n\le2x^3$. Hence
\[
 2\le m\le2x\le2x^2\le2n\le4x^3\le x^4\le k.
\]
Also $m+3\le2x+3\le x^4\le k$.

Define $r=k+m-1$ and $t=8fr+4+TW-r$.
Since $f\ge m\ge2$ and $r\ge m$, the difference is at least
$15r+4$, so it has its ordinary integer value.
Since $k$ and $m$ are even, $r$ is odd. Since $8\mid W$,
the integer $t$ is odd. Moreover,
\[
 r+5\le t,\qquad 3m<t,\qquad
 m^2\le m(t+1)\le k+m(t+1).
\]
The bounds and parity conditions
give a rank record $R$ extending $\varphi$. The formula for $t$
gives the identity for $r+t$. The frame fixes $R.n=n$, so the
assumed bounds on $n$ carry over to $R$. The choices of $m$ and $k$
give their respective bounds.
\end{proof}
\leanref{GlobalTheory.\allowbreak{}EX.\allowbreak{}exists\_\allowbreak{}even\_\allowbreak{}between}{Development.\allowbreak{}AbelianVarietyTheory.\allowbreak{}Proofs.\allowbreak{}GlobalInterpolation.\allowbreak{}ExistsRanksRankValFrame}{10--16}
\leanref{GlobalTheory.\allowbreak{}EX.\allowbreak{}exists\_\allowbreak{}ranks\_\allowbreak{}of\_\allowbreak{}frame}{Development.\allowbreak{}AbelianVarietyTheory.\allowbreak{}Proofs.\allowbreak{}GlobalInterpolation.\allowbreak{}ExistsRanksRankValFrame}{24--43}


\part{Galois representations and local Langlands in the archive}

% ---- p9c1 ----
\chapter{Galois representations attached to polarized cuspidal representations}\label{ch:galois-polarized}
\section{The global theorem}

Let $E/F$ be CM with conjugation $c$. For an algebraic closure $\Omega$ put $G_E=\Gal(\Omega/E)$ and fix $\iota:\bar{\Q}_\ell\simeq\C$. A polarized input consists of a cuspidal $\pi$ on $\GL_a(\mathbf A_E)$, $a\ge1$, with distinct infinitesimal weights in $(a-1)/2+\Z$, and a Hecke character $\chi$ of $F$ with $\chi_v(-1)=(-1)^a$ at infinity and
\[
\pi^c\simeq(\pi\otimes(\chi\circ N_{E/F})^{-1})^\vee.
\]
For a Satake multiset $B$ at $w$, put $q_w=\#k(w)$ and $A_a(w,B)=q_w^{(a-1)/2}B$. Set $P_{w,B}(X)=\prod_{\alpha\in A_a(w,B)}(X-\alpha)$. We use inverse arithmetic Frobenius and transport coefficients by $\iota$. Galois representations are continuous and semisimple unless specified otherwise.

\begin{definition}[Literature package]\label{B9.01}
The two $\mathrm{Row}$ fields use \ref{B9.02}, and $\mathrm{RecTExists}$ uses \ref{B9.46}. The same row supplies the window.
\end{definition}
\leanref{UpperTheory.\allowbreak{}LiteratureTheorems.\allowbreak{}OpenLiteratureStatements}{Development.\allowbreak{}AutomorphicArithmetic.\allowbreak{}ThetaCorrespondence.\allowbreak{}ThetaPairLifts.\allowbreak{}ThirdPlaceSeries.\allowbreak{}EveryCharacterOfLiteratureOfVWin}{215--218}
\leanref{UpperTheory.\allowbreak{}FinalTheorems.\allowbreak{}everyCharacter\_\allowbreak{}of\_\allowbreak{}row}{Development.\allowbreak{}AutomorphicArithmetic.\allowbreak{}CharacterBoxes.\allowbreak{}MainConclusion}{7--9}
\leanref{UpperTheory.\allowbreak{}FinalTheorems.\allowbreak{}everyCharacter}{Development.\allowbreak{}AutomorphicArithmetic.\allowbreak{}CharacterBoxes.\allowbreak{}MainConclusion}{10--12}

\begin{theorem}[Global theorem]\label{B9.02}
Each polarized input admits $r:G_E\to\GL_a(\bar{\Q}_\ell)$, $w_{\mathrm{gt}}\in\Z$ and integer multisets $H_\tau$ as in \ref{B9.04}. These data satisfy \ref{B9.05} and \ref{B9.06} for every map of \ref{B9.47}.
\end{theorem}
\leanref{UpperTheory.\allowbreak{}Gal.\allowbreak{}Int.\allowbreak{}galoisLocalGlobalStrong\_\allowbreak{}recParameter}{Development.\allowbreak{}AutomorphicArithmetic.\allowbreak{}Foundations.\allowbreak{}GaloisRepresentations.\allowbreak{}IntegralStructures.\allowbreak{}RowSrc}{6--7}
\begin{proof}
Apply \ref{B9.03}.
\end{proof}

\begin{proposition}[Assembly]\label{B9.03}
Theorems \ref{B9.08} and \ref{B9.39} imply \ref{B9.02}.
\end{proposition}
\leanref{UpperTheory.\allowbreak{}Gal.\allowbreak{}Int.\allowbreak{}galoisLocalGlobalStrong\_\allowbreak{}of\_\allowbreak{}geometry\_\allowbreak{}lift}{Development.\allowbreak{}AutomorphicArithmetic.\allowbreak{}Foundations.\allowbreak{}GaloisRepresentations.\allowbreak{}IntegralStructures.\allowbreak{}Development}{30--32}
\leanref{UpperTheory.\allowbreak{}Gal.\allowbreak{}Int.\allowbreak{}galoisRepresentations\_\allowbreak{}of\_\allowbreak{}thetaLift}{Development.\allowbreak{}AutomorphicArithmetic.\allowbreak{}Foundations.\allowbreak{}GaloisRepresentations.\allowbreak{}IntegralStructures.\allowbreak{}GaloisRepresentationsOfThetaLift}{6--7}
\begin{proof}
Apply \ref{B9.28}, \ref{B9.31}, \ref{B9.35} and \ref{B9.38}, then \ref{B9.07} with the stated tame rank bound.
\end{proof}

\begin{definition}[Weights]\label{B9.04}
For $\tau:E\hookrightarrow\bar{\Q}_\ell$, $H_\tau=(a-1)/2-\operatorname{InfChar}_{\iota\circ\tau}(\pi)$ and $H_{\tau c}=w_{\mathrm{gt}}-H_\tau$.
\end{definition}
\leanref{UpperTheory.\allowbreak{}GalStar.\allowbreak{}PropsB}{Development.\allowbreak{}AutomorphicArithmetic.\allowbreak{}Foundations.\allowbreak{}GaloisRepresentations.\allowbreak{}TypeConditions.\allowbreak{}WeightsConjBis}{54--68}
\leanref{UpperTheory.\allowbreak{}GalStar.\allowbreak{}Row0b}{Development.\allowbreak{}AutomorphicArithmetic.\allowbreak{}Foundations.\allowbreak{}GaloisRepresentations.\allowbreak{}TypeConditions.\allowbreak{}WeightsConjBis}{69--74}

At $w\mid\ell$, Caraiani identifies the Frobenius-semisimplified Weil--Deligne representation, including monodromy, with the classical local parameter in the algebraic normalization for cohomological conjugate self-dual cuspidal input of rank $a\ge2$ \cite[Thm.~1.1]{caraiani-2014}. For the polarization above, compatibility holds for every weak parameter map under \ref{B9.05} and \ref{B9.06}.

\begin{definition}[Away from $\ell$]\label{B9.05}
For $w\nmid\ell$ and a prolongation put $x=\mathrm{WD}_\iota(r_w)$ and $y=rc(\pi_w\otimes|\det|_w^{(1-a)/2})$. Write $N_x,N_y$ for monodromy. The toral test requires polynomial $(X-1)^a$ on Weil commutators of $y$. The tame test requires it on wild inertia and requires $z^{q_w-1}=1$ for repeated roots at lifts of tame generators. The full tame test also requires $E_w\simeq\Q_p$ continuously and $a=1$ or $p>16a^2$. The relation $\operatorname{AwayB}(x,y)$ requires equal ranks. If $N_y=0$, it requires equality of classes under the toral test, and $N_x=0$ with equal inertia polynomials under the full tame test.
\end{definition}
\leanref{UpperTheory.\allowbreak{}GalStar.\allowbreak{}AwayB}{Development.\allowbreak{}AutomorphicArithmetic.\allowbreak{}Foundations.\allowbreak{}GaloisRepresentations.\allowbreak{}TypeConditions.\allowbreak{}WeightsConjBis}{26--32}
\leanref{UpperTheory.\allowbreak{}GalStar.\allowbreak{}Sqf}{Development.\allowbreak{}AutomorphicArithmetic.\allowbreak{}Foundations.\allowbreak{}GaloisRepresentations.\allowbreak{}TypeConditions.\allowbreak{}WeightsConjAlt}{29--34}

\begin{definition}[Above $\ell$]\label{B9.06}
At $w\mid\ell$ take any prolongation with a compatible period datum and bridge. For spherical $\pi_w$, $r_w$ is crystalline de Rham with jumps $H_\tau$ at induced embeddings $\tau$. If $N_y=0$ and either test of \ref{B9.05} holds in full, $r_w$ is potentially semistable de Rham with those jumps and its period parameter $x$ satisfies $\operatorname{AwayB}(x,y)$.
\end{definition}
\leanref{UpperTheory.\allowbreak{}GalStar.\allowbreak{}AboveB}{Development.\allowbreak{}AutomorphicArithmetic.\allowbreak{}Foundations.\allowbreak{}GaloisRepresentations.\allowbreak{}TypeConditions.\allowbreak{}WeightsConjBis}{39--53}

\begin{proposition}[Local clauses]\label{B9.07}
Statement \ref{B9.38} implies \ref{B9.04}--\ref{B9.06} with unrestricted rank in the tame clause.
\end{proposition}
\leanref{UpperTheory.\allowbreak{}Gal.\allowbreak{}Int.\allowbreak{}Row.\allowbreak{}galoisLocalGlobal\_\allowbreak{}of\_\allowbreak{}placeClassCompatibility}{Development.\allowbreak{}AutomorphicArithmetic.\allowbreak{}Foundations.\allowbreak{}GaloisRepresentations.\allowbreak{}IntegralStructures.\allowbreak{}InertCharpoly.\allowbreak{}GaloisLocalGlobalOfLocalGlobalCompatibility}{157--172}
\leanref{UpperTheory.\allowbreak{}Gal.\allowbreak{}Int.\allowbreak{}Row.\allowbreak{}away}{Development.\allowbreak{}AutomorphicArithmetic.\allowbreak{}Foundations.\allowbreak{}GaloisRepresentations.\allowbreak{}IntegralStructures.\allowbreak{}InertCharpoly.\allowbreak{}GaloisLocalGlobalOfLocalGlobalCompatibility}{96--106}
\leanref{UpperTheory.\allowbreak{}Gal.\allowbreak{}Int.\allowbreak{}Row.\allowbreak{}above}{Development.\allowbreak{}AutomorphicArithmetic.\allowbreak{}Foundations.\allowbreak{}GaloisRepresentations.\allowbreak{}IntegralStructures.\allowbreak{}InertCharpoly.\allowbreak{}GaloisLocalGlobalOfLocalGlobalCompatibility}{108--126}
\leanref{UpperTheory.\allowbreak{}Gal.\allowbreak{}Int.\allowbreak{}rest}{Development.\allowbreak{}AutomorphicArithmetic.\allowbreak{}Foundations.\allowbreak{}GaloisRepresentations.\allowbreak{}IntegralStructures.\allowbreak{}Development}{24--26}
\begin{proof}
Choose $r$ by \ref{B9.38} before fixing $rc$. Polarization gives weight symmetry. When $N_y=0$, apply \ref{B9.55a} to the generic normalized component. Toral comparison gives the sum of the unlinked characters by reciprocity. Tame depth-zero support gives first principal congruence invariants and inertial separation. Choose the auxiliary extension unramified at $w$ of residue degree $a!$, where the lifted support is toral. Each inertia element downstairs lifts to inertia upstairs. Evaluate the toral comparison on that lift and use compatibility of the support characters to recover its polynomial downstairs. Weil--Deligne restriction descends zero monodromy. Above $\ell$ the comparison gives potential semistability and the prescribed jumps. For \ref{B9.02}, use the spherical clause of \ref{B9.38} and impose the tame rank bound.
\end{proof}

\section{Geometric seeds in odd rank}

Let $a$ be odd, $x^2=-d$ in $E$ with $d\in\Z_{>0}$, $-d\in(\Q_\ell^\times)^2$, and $[F:\Q]$ even. Nonsplit finite places are unramified and away from $\ell$. Separated places have distinct restrictions to $F$. A split place over $p$ is usable if $-d\in(\Q_p^\times)^2$.

Take separated $w_0,v'\mid\ell$ agreeing on $\Q(x)$. Let $U$ be definite-unitary automorphic, unramified at nonsplit places off finite $T$ and spherical at $v'$ and at split places separated from $w_0$ and a finite set $S$. At $w_0$ it contains an elliptic type, a nontrivial character of a torsion-free compact open $J$ through $J\twoheadrightarrow\Z/\ell^m$ that forces cuspidality. Put $J(\lambda)=\{i-\lambda_i:0\le i<a\}$. A set is lacunary if disjoint nonempty subsets have different means.

\begin{theorem}[Geometric seed]\label{B9.08}
If $J(\lambda)$ is lacunary at an embedding over $v'$, there is a $\rho$ with a lacunary subset of size $a$ in an actual jump multiset. At usable $w\nmid\ell$ off a finite set, inertia is trivial and $P_{w,B}$ annihilates inverse Frobenius for every $B$ and prolongation. One can include any specified usable spherical $w\nmid\ell$. For spherical $w\mid\ell$ separated from $w_0,v'$ and agreeing with $v'$ on $\Q(x)$, two witnesses retain the common seed properties. The first is crystalline at $w$ and for each induced embedding its actual jumps contain every prescribed jump. The second gives the divisibility of \ref{B9.19} for each $R=A_a(w,B)$ and compatible period datum.
\end{theorem}
\leanref{UpperTheory.\allowbreak{}Gal.\allowbreak{}Int.\allowbreak{}Geo.\allowbreak{}GeoOne}{Development.\allowbreak{}AutomorphicArithmetic.\allowbreak{}Foundations.\allowbreak{}GaloisRepresentations.\allowbreak{}IntegralStructures.\allowbreak{}GeometricFlagBases.\allowbreak{}HodgeTrAt}{79--91}
\leanref{UpperTheory.\allowbreak{}Gal.\allowbreak{}Int.\allowbreak{}Geo.\allowbreak{}geoOne\_\allowbreak{}holds}{Development.\allowbreak{}AutomorphicArithmetic.\allowbreak{}Foundations.\allowbreak{}GaloisRepresentations.\allowbreak{}IntegralStructures.\allowbreak{}GeometricFlagBases.\allowbreak{}WeightOneGeometricRealization}{11--12}
\begin{proof}
For a plan of \ref{B9.09}, twist the piece of \ref{B9.15} by $\kappa_\ell$. Apply \ref{B9.17} away from $\ell$ and \ref{B9.18a} and \ref{B9.19} above it.
\end{proof}

\begin{proposition}[Geometric plans]\label{B9.09}
A plan gives a Hecke eigensystem $\theta$ and Hodge data at $v'$. It also gives an algebraic character $\kappa$. At usable $w\nmid\ell$ off a finite set, $\kappa$ is unramified and the correspondence roots are $\kappa(w)^{-1}A_a(w,B)$ for every $B$. A plan may include the local tests of \ref{B9.08}.
\end{proposition}
\leanref{UpperTheory.\allowbreak{}Gal.\allowbreak{}Int.\allowbreak{}Geo.\allowbreak{}PlanEx}{Development.\allowbreak{}AutomorphicArithmetic.\allowbreak{}Foundations.\allowbreak{}GaloisRepresentations.\allowbreak{}IntegralStructures.\allowbreak{}GeometricFlagBases.\allowbreak{}PackExOf}{82--92}
\leanref{UpperTheory.\allowbreak{}Gal.\allowbreak{}Int.\allowbreak{}Geo.\allowbreak{}Pl.\allowbreak{}planEx\_\allowbreak{}holds}{Development.\allowbreak{}AutomorphicArithmetic.\allowbreak{}Foundations.\allowbreak{}GaloisRepresentations.\allowbreak{}IntegralStructures.\allowbreak{}GeometricFlagBases.\allowbreak{}PrimeLevelLocalization.\allowbreak{}SpecTopHolds}{9--10}
\leanref{UpperTheory.\allowbreak{}Gal.\allowbreak{}Int.\allowbreak{}Geo.\allowbreak{}Plan.\allowbreak{}pack}{Development.\allowbreak{}AutomorphicArithmetic.\allowbreak{}Foundations.\allowbreak{}GaloisRepresentations.\allowbreak{}IntegralStructures.\allowbreak{}GeometricFlagBases.\allowbreak{}PackExOf}{101--126}
\begin{proof}
Use the spectral construction's auxiliary prime and realize the untwisted exponents on a fibre power.
\end{proof}

\begin{definition}[PEL objects]\label{B9.10}
A PEL object is an abelian scheme $A/S$ with prescribed action and Lie multiplicities. It also has a prescribed level. It has a symmetric line bundle $P$ on $A\times_S A$, trivializable on the axes and relatively ample on the diagonal, with $(b\times1)^*P\simeq(1\times b^*)^*P$ for order elements $b$. The universal object classifies by pullback up to the prescribed isomorphism.
\end{definition}
\leanref{GlobalTheory.\allowbreak{}UPel.\allowbreak{}HT.\allowbreak{}Gen.\allowbreak{}GObj}{Development.\allowbreak{}AbelianVarietyTheory.\allowbreak{}Definitions.\allowbreak{}Divisors.\allowbreak{}ExistsGDatumOneOneOne}{139--168}
\leanref{GlobalTheory.\allowbreak{}UPel.\allowbreak{}HT.\allowbreak{}Gen.\allowbreak{}GenModel}{Development.\allowbreak{}AbelianVarietyTheory.\allowbreak{}Definitions.\allowbreak{}Divisors.\allowbreak{}ExistsGDatumOneOneOne}{181--186}

\begin{theorem}[Representability]\label{B9.11}
Harris--Taylor data at $p$ over a Noetherian $\Q$-algebra with specified base datum and level $N\ge3$ prime to $p$ give representability of \ref{B9.10}.
\end{theorem}
\leanref{GlobalTheory.\allowbreak{}T.\allowbreak{}unitaryCharZeroGenericModel\_\allowbreak{}nonempty}{Development.\allowbreak{}AbelianVarietyTheory.\allowbreak{}Proofs.\allowbreak{}DivisorProofs.\allowbreak{}UnitaryCharZeroGenericModelNonemptyStatement}{8--10}
\leanref{GlobalTheory.\allowbreak{}UPel.\allowbreak{}HT.\allowbreak{}Gen.\allowbreak{}genModel\_\allowbreak{}nonempty}{Development.\allowbreak{}AbelianVarietyTheory.\allowbreak{}Proofs.\allowbreak{}DivisorProofs.\allowbreak{}UnitaryCharZeroGenericModelNonempty}{118--123}
This is the characteristic-zero case of the integral moduli theorem \cite[{\S}III.1]{harris-taylor-2001}.

\begin{proposition}[Fibre powers]\label{B9.12}
Over a characteristic-zero field, compatible compact data give smooth projective fibre powers of every nonnegative exponent with finite étale first Hecke projections of prescribed constant degree.
\end{proposition}
\leanref{UpperTheory.\allowbreak{}Gal.\allowbreak{}Int.\allowbreak{}Geo.\allowbreak{}KS}{Development.\allowbreak{}AutomorphicArithmetic.\allowbreak{}Foundations.\allowbreak{}GaloisRepresentations.\allowbreak{}IntegralStructures.\allowbreak{}GeometricFlagBases.\allowbreak{}IntegOfMem}{48--60}
\leanref{UpperTheory.\allowbreak{}Gal.\allowbreak{}G3.\allowbreak{}pKSEx}{Development.\allowbreak{}AutomorphicArithmetic.\allowbreak{}Foundations.\allowbreak{}GaloisRepresentations.\allowbreak{}PAdicWittLattices.\allowbreak{}ExistsDeeperNormalLevel}{100--173}
\leanref{UpperTheory.\allowbreak{}Gal.\allowbreak{}G3.\allowbreak{}kSNice\_\allowbreak{}holds}{Development.\allowbreak{}AutomorphicArithmetic.\allowbreak{}Foundations.\allowbreak{}GaloisRepresentations.\allowbreak{}PAdicWittLattices.\allowbreak{}FibrePowerTowerAdmissibleLevel}{152--153}
\leanref{UpperTheory.\allowbreak{}Gal.\allowbreak{}G3.\allowbreak{}kSEx\_\allowbreak{}holds}{Development.\allowbreak{}AutomorphicArithmetic.\allowbreak{}Foundations.\allowbreak{}GaloisRepresentations.\allowbreak{}PAdicWittLattices.\allowbreak{}FibrePowerTowerAdmissibleLevel}{172--173}
\begin{proof}
Base change the universal maps. Compactness gives projectivity and level indices give the étale degrees.
\end{proof}

\begin{definition}[Total cohomology]\label{B9.13}
Use geometric pullback on $\bigoplus_j\Q_\ell\otimes_{\Z_\ell}\varprojlim_mH^j_{\mathrm{\acute et}}(Z_\Omega,\Z/\ell^m)$ \cite{milne-1980}.
\end{definition}
\leanref{GlobalTheory.\allowbreak{}EtC.\allowbreak{}HlSub}{Development.\allowbreak{}AbelianVarietyTheory.\allowbreak{}Definitions.\allowbreak{}Divisors.\allowbreak{}PullHlComp}{216--224}
\leanref{GlobalTheory.\allowbreak{}EtC.\allowbreak{}HQ}{Development.\allowbreak{}AbelianVarietyTheory.\allowbreak{}Definitions.\allowbreak{}Divisors.\allowbreak{}PullHlComp}{239--240}
\leanref{GlobalTheory.\allowbreak{}GalV.\allowbreak{}Htot}{Development.\allowbreak{}AbelianVarietyTheory.\allowbreak{}Definitions.\allowbreak{}Divisors.\allowbreak{}IsCorrOpOfAvg}{11--12}
\leanref{GlobalTheory.\allowbreak{}GalV.\allowbreak{}galTotRep}{Development.\allowbreak{}AbelianVarietyTheory.\allowbreak{}Definitions.\allowbreak{}Divisors.\allowbreak{}IsCorrOpOfAvg}{58--63}

\begin{lemma}[Correspondences]\label{B9.14}
For a correspondence $(s,t)$ with $s$ finite étale of constant degree, set $T=\operatorname{Tr}_st^*$. A Galois cover $b$ splitting $s$ into branches $\sigma_j$ gives $b^*Tu=\sum_j(t\sigma_j)^*u$.
\end{lemma}
\leanref{UpperTheory.\allowbreak{}Gal.\allowbreak{}G3.\allowbreak{}pHopSpec}{Development.\allowbreak{}AutomorphicArithmetic.\allowbreak{}Foundations.\allowbreak{}GaloisRepresentations.\allowbreak{}PAdicWittLattices.\allowbreak{}PHopSpec}{17--51}
\leanref{UpperTheory.\allowbreak{}Gal.\allowbreak{}Int.\allowbreak{}Geo.\allowbreak{}hopSpec}{Development.\allowbreak{}AutomorphicArithmetic.\allowbreak{}Foundations.\allowbreak{}GaloisRepresentations.\allowbreak{}IntegralStructures.\allowbreak{}GeometricFlagBases.\allowbreak{}RepMain}{56--58}
This is finite étale trace base change over $\Q_\ell$ \cite{milne-1980}.

\begin{proposition}[Eigenspaces]\label{B9.15}
The common Hecke eigenspace over $\bar{\Q}_\ell$ is Galois-stable with finite-dimensional continuous semisimplification.
\end{proposition}
\leanref{UpperTheory.\allowbreak{}Gal.\allowbreak{}Int.\allowbreak{}Geo.\allowbreak{}repEx\_\allowbreak{}of}{Development.\allowbreak{}AutomorphicArithmetic.\allowbreak{}Foundations.\allowbreak{}GaloisRepresentations.\allowbreak{}IntegralStructures.\allowbreak{}GeometricFlagBases.\allowbreak{}RepAsm}{13--60}
\leanref{UpperTheory.\allowbreak{}Gal.\allowbreak{}Int.\allowbreak{}Geo.\allowbreak{}repEx\_\allowbreak{}core}{Development.\allowbreak{}AutomorphicArithmetic.\allowbreak{}Foundations.\allowbreak{}GaloisRepresentations.\allowbreak{}IntegralStructures.\allowbreak{}GeometricFlagBases.\allowbreak{}RepMain}{74--76}
\begin{proof}
Use \ref{B9.26} degreewise and equivariance of the correspondences. Take the diagonal blocks along a composition flag.
\end{proof}

\begin{lemma}[Ordinary points]\label{B9.16}
At an ordinary point let $T_j$ be the branch multiset indexed by $j$-planes of $H=k(w)^a$. For residue Frobenius $F_*$ there are multisets $A_j$ with $A_0=T_0$ and $A_a=0$. They satisfy $T_{j+1}=F_*A_j+q_w^{j+1}A_{j+1}$ for $0\le j<a$. Thus
\[
\sum_{j=0}^a(-1)^jq_w^{\binom j2}F_*^{a-j}T_j=0.
\]
Let $V$ be a valuation ring with algebraically closed fraction field and residue characteristic $p$ nonzero in $V$. If $y:\operatorname{Spec}V\to Y$ sends the closed point to a generic point of the special fibre, the reduced branches satisfy the same identity in the free abelian group on residue-field points.
\end{lemma}
\leanref{UpperTheory.\allowbreak{}Gal.\allowbreak{}crystallineComparison.\allowbreak{}ordDecomp\_\allowbreak{}of}{Development.\allowbreak{}AutomorphicArithmetic.\allowbreak{}Foundations.\allowbreak{}GaloisRepresentations.\allowbreak{}AnnihilatorModels.\allowbreak{}POrdDecomp}{98--170}
\leanref{UpperTheory.\allowbreak{}Gal.\allowbreak{}crystallineComparison.\allowbreak{}telescope\_\allowbreak{}of}{Development.\allowbreak{}AutomorphicArithmetic.\allowbreak{}Foundations.\allowbreak{}GaloisRepresentations.\allowbreak{}AnnihilatorModels.\allowbreak{}PTelescope}{47--77}
\leanref{UpperTheory.\allowbreak{}Gal.\allowbreak{}crystallineComparison.\allowbreak{}Mod.\allowbreak{}congKS\_\allowbreak{}link}{Development.\allowbreak{}AutomorphicArithmetic.\allowbreak{}Foundations.\allowbreak{}GaloisRepresentations.\allowbreak{}AnnihilatorModels.\allowbreak{}QuotObjLink}{13--13}
\leanref{UpperTheory.\allowbreak{}Gal.\allowbreak{}crystallineComparison.\allowbreak{}ordCount\_\allowbreak{}holds}{Development.\allowbreak{}AutomorphicArithmetic.\allowbreak{}Foundations.\allowbreak{}GaloisRepresentations.\allowbreak{}AnnihilatorModels.\allowbreak{}GenericMulOfProdHolds}{51--51}
\begin{proof}
The $(j+1)$-planes containing the distinguished line $C$ give the $j$-planes of $H/C$ with images $F_*A_j$. Above a quotient $(j+1)$-plane $\bar D$, the $\#\operatorname{Hom}(\bar D,C)=q_w^{j+1}$ complements give the same image. The terms in $A_j$ cancel by $\binom j2+j=\binom{j+1}2$. Integral quotients give the reductions.
\end{proof}

\begin{proposition}[Annihilator]\label{B9.17}
At $w\nmid\ell$ the plan's root identity gives trivial inertia and annihilator $P_{w,B}$ at inverse Frobenius on the twisted piece.
\end{proposition}
\leanref{UpperTheory.\allowbreak{}Gal.\allowbreak{}Mod.\allowbreak{}esStep}{Development.\allowbreak{}AutomorphicArithmetic.\allowbreak{}Foundations.\allowbreak{}GaloisRepresentations.\allowbreak{}BasicOpenSections.\allowbreak{}PEsStep}{49--95}
\leanref{UpperTheory.\allowbreak{}Gal.\allowbreak{}Mod.\allowbreak{}towES}{Development.\allowbreak{}AutomorphicArithmetic.\allowbreak{}Foundations.\allowbreak{}GaloisRepresentations.\allowbreak{}BasicOpenSections.\allowbreak{}PTowES}{147--221}
\leanref{UpperTheory.\allowbreak{}Gal.\allowbreak{}Int.\allowbreak{}Geo.\allowbreak{}CongFor}{Development.\allowbreak{}AutomorphicArithmetic.\allowbreak{}Foundations.\allowbreak{}GaloisRepresentations.\allowbreak{}IntegralStructures.\allowbreak{}GeometricFlagBases.\allowbreak{}HodgeTrAt}{54--57}
\leanref{UpperTheory.\allowbreak{}Gal.\allowbreak{}Int.\allowbreak{}Geo.\allowbreak{}Pack.\allowbreak{}rep}{Development.\allowbreak{}AutomorphicArithmetic.\allowbreak{}Foundations.\allowbreak{}GaloisRepresentations.\allowbreak{}IntegralStructures.\allowbreak{}GeometricFlagBases.\allowbreak{}MapInvMul}{43--87}
\begin{proof}
The deck transformations are Hecke maps with eigenvalue $1$, so averaging descends the eigenclasses to the smooth model. Apply \ref{B9.27} and \ref{B9.16}. Trace divides the lifted eigenvalue $m_j\vartheta_j$ by $m_j$ by \ref{B9.14}. Trivial inertia and the resulting annihilator $\sum_j(-1)^jq_w^{\binom j2}\vartheta_jX^{a-j}$ pass to the graded. Divide by $\vartheta_0\ne0$ and twist by $\kappa_\ell$.
\end{proof}

\begin{theorem}[Holomorphic Lefschetz]\label{B9.18}
Let $Z/\C$ be smooth projective of dimension $D$ and $Z\xleftarrow{c_1}Y\xrightarrow{c_2}Z$ a finite étale correspondence with finitely many transverse fixed points. Put $A_y=dc_2(dc_1)^{-1}$ at each fixed point. For $j\ge0$,
\[
\sum_{b=0}^D(-1)^b\operatorname{Tr}(\operatorname{Tr}_{c_1}c_2^*\mid H^b(Z,\Omega_Z^j))
=\sum_y\frac{[T^j]\det(1+TA_y)}{\det(1-A_y)}.
\]
\end{theorem}
\leanref{UpperTheory.\allowbreak{}Gal.\allowbreak{}LocalSmoothExtension.\allowbreak{}HolomorphicLefschetzFormula}{Development.\allowbreak{}AutomorphicArithmetic.\allowbreak{}Foundations.\allowbreak{}GaloisRepresentations.\allowbreak{}LefschetzSigmaSums.\allowbreak{}HodgeTraceFromLefschetzAndGroupSide}{102--112}
\leanref{UpperTheory.\allowbreak{}Gal.\allowbreak{}Int.\allowbreak{}Geo.\allowbreak{}weight\_\allowbreak{}one\_\allowbreak{}geometric\_\allowbreak{}realization}{Development.\allowbreak{}AutomorphicArithmetic.\allowbreak{}Foundations.\allowbreak{}GaloisRepresentations.\allowbreak{}IntegralStructures.\allowbreak{}GeometricFlagBases.\allowbreak{}WeightOneGeometricRealization}{7--9}
This is the Atiyah--Bott formula for the indicated bundles and correspondences \cite{atiyah-bott-1968}.

\begin{lemma}[Hodge degrees of a spectral test]\label{B9.18a}
For a dominant-weight plan of \ref{B9.09}, choose the fibre exponent $k$ from the weight, uniformly for its conjugates. For a prescribed jump $j$ at $\tau$, take the conjugate embedding and Hodge exponent $h_i=j+m_\tau$, where $m_\tau$ is the infinity exponent of $\kappa$. On the resulting fibre power $Z$, some $H^b(Z,\Omega^{h_i})$ contains a nonzero common eigenvector with eigenvalues $\theta$ transported by the corresponding complex embedding.
\end{lemma}
\leanref{UpperTheory.\allowbreak{}Gal.\allowbreak{}LocalSmoothExtension.\allowbreak{}groupTop\_\allowbreak{}of}{Development.\allowbreak{}AutomorphicArithmetic.\allowbreak{}Foundations.\allowbreak{}GaloisRepresentations.\allowbreak{}LefschetzSigmaSums.\allowbreak{}PGroupTop}{12--66}
\leanref{UpperTheory.\allowbreak{}Gal.\allowbreak{}LocalSmoothExtension.\allowbreak{}hodgeEig\_\allowbreak{}of}{Development.\allowbreak{}AutomorphicArithmetic.\allowbreak{}Foundations.\allowbreak{}GaloisRepresentations.\allowbreak{}LefschetzSigmaSums.\allowbreak{}PHodgeEig}{6--68}
\begin{proof}
Let $f$ be the spectral test and write $\theta$ for its transported eigenvalues. The Hecke operators commute on words applied to $f$. For a word polynomial $s$ with $s(\theta)\ne0$, the spectral condition gives a word polynomial $a$ for which $F=asf$ is regular with nonzero definite trace $D(F)$. It also gives two vanishing conditions. A geometric pair with a marked root and reducible polynomial contributes zero. On the definite unitary side the orbital terms outside the regular irreducible locus vanish.

Scale the fibres by nonzero even order elements $\alpha_u$. At a marked root $\mu$, the base tangent eigenvalues are $\mu/r$ for the other roots $r$. The fibre eigenvalues are twice algebraic integers since the isogeny composed with $\alpha_u/2$ preserves the lattice. Thus $1-A_{\rm fib}$ is an isogeny with kernel of order $|\det(1-A_{\rm fib})|^2$, the number of fibre fixed points over each fixed line with level. The finer-level multiplicity cancels the relative Hecke index in \ref{B9.18}. Expansion of the Schur alternants and interpolation in the Zariski-dense embedding coordinates of the $\alpha_u$ give a linear combination $C$ of multiplication pullbacks. Reindex the geometric sum by regular irreducible polynomials and their marked roots. For each polynomial the sum over its roots is $\prod_\tau s_{\lambda_\tau}(r_\tau)\sum_i y^{h_i}$, where $s_{\lambda_\tau}$ is the Schur character on the corresponding root list $r_\tau$. The definite unitary orbital trace formula sums these character weights to give
\[
\sum_{j,b}(-y)^j(-1)^b\operatorname{tr}(T_F C\mid H^b(Z,\Omega^j))
=D(F)\sum_i y^{h_i}.
\]
In Hodge degree $h_i$, choose $s(\theta)=1$ by interpolation to kill the other primary spaces in the Hecke span of $\operatorname{im}T_f$. The $y^{h_i}$ coefficient is nonzero since the $h_i$ are distinct. Thus $sT_f$ has nonzero image for some $b$, and the commuting nilpotents on the $\theta$-primary space have a nonzero common kernel.
\end{proof}

\begin{proposition}[Crystalline data]\label{B9.19}
At $w\mid\ell$ a surjective cover as in \ref{B9.17} gives crystallinity of the semisimplification. If the correspondences descend to one model with prescribed multiplicities, a common Hodge eigenclass of degree $h$ gives jump $h$ at the induced embedding. For each compatible period datum, the root identity for $R$ gives a crystalline piece whose Weil polynomial at some element of degree $-1$ divides $(\prod_{\alpha\in R}(X-\alpha))^m$ for some $m\in\N$.
\end{proposition}
\leanref{UpperTheory.\allowbreak{}Gal.\allowbreak{}Int.\allowbreak{}Geo.\allowbreak{}PhiK}{Development.\allowbreak{}AutomorphicArithmetic.\allowbreak{}Foundations.\allowbreak{}GaloisRepresentations.\allowbreak{}IntegralStructures.\allowbreak{}GeometricFlagBases.\allowbreak{}HodgeTrAt}{32--37}
\leanref{UpperTheory.\allowbreak{}Gal.\allowbreak{}Int.\allowbreak{}Geo.\allowbreak{}CrysStep}{Development.\allowbreak{}AutomorphicArithmetic.\allowbreak{}Foundations.\allowbreak{}GaloisRepresentations.\allowbreak{}IntegralStructures.\allowbreak{}GeometricFlagBases.\allowbreak{}HodgeTrAt}{112--115}
\leanref{UpperTheory.\allowbreak{}Gal.\allowbreak{}Int.\allowbreak{}Geo.\allowbreak{}JumpStep}{Development.\allowbreak{}AutomorphicArithmetic.\allowbreak{}Foundations.\allowbreak{}GaloisRepresentations.\allowbreak{}IntegralStructures.\allowbreak{}GeometricFlagBases.\allowbreak{}HodgeTrAt}{117--120}
\leanref{UpperTheory.\allowbreak{}Gal.\allowbreak{}Int.\allowbreak{}Geo.\allowbreak{}PhiStep}{Development.\allowbreak{}AutomorphicArithmetic.\allowbreak{}Foundations.\allowbreak{}GaloisRepresentations.\allowbreak{}IntegralStructures.\allowbreak{}GeometricFlagBases.\allowbreak{}HodgeTrAt}{122--126}
\leanref{UpperTheory.\allowbreak{}Gal.\allowbreak{}Int.\allowbreak{}Geo.\allowbreak{}jumpStep\_\allowbreak{}holds}{Development.\allowbreak{}AutomorphicArithmetic.\allowbreak{}Foundations.\allowbreak{}GaloisRepresentations.\allowbreak{}IntegralStructures.\allowbreak{}GeometricFlagBases.\allowbreak{}GeometricRealizationOfWeightOne}{18--19}
\leanref{UpperTheory.\allowbreak{}Ru.\allowbreak{}crystallineComparison.\allowbreak{}phiStep}{Development.\allowbreak{}AutomorphicArithmetic.\allowbreak{}CommutativeAlgebra.\allowbreak{}EtaleCohomology.\allowbreak{}SmoothModelCrystalline}{13--14}
\begin{proof}
Trace and pullback descend the eigenclasses with eigenvalues $\theta(g)/m_g$ for correspondence multiplicities $m_g$. Smooth proper comparison gives the crystalline periods and carries Hodge degree $h$ to jump $h$ \cite{bhatt-morrow-scholze-2018}. The descended eigenspace injects into the original piece. The jump survives semisimplification by additivity, and the annihilator on periods gives the divisibility. In \ref{B9.08} the twist subtracts $m_\tau$ from $h_i=j+m_\tau$.
\end{proof}

\section{Cohomological inputs}

Write $H^j(Z,i)=H^j_{\mathrm{\acute et}}(Z,\Z/b(i))$ for $b\ge1$ invertible on the base. The same notation applies to simplicial presheaves.

\begin{definition}[Core]\label{B9.20}
C2--C5 and D1 are \ref{B9.21}--\ref{B9.25}, respectively.
\end{definition}
\leanref{UpperTheory.\allowbreak{}Gal.\allowbreak{}KunnethLefschetz.\allowbreak{}core\_\allowbreak{}holds}{Development.\allowbreak{}AutomorphicArithmetic.\allowbreak{}Foundations.\allowbreak{}GaloisRepresentations.\allowbreak{}KunnethLefschetz.\allowbreak{}CoreHolds}{5--7}

\begin{theorem}[Affine curves, C2]\label{B9.21}
For a finite-type Dedekind $k$-algebra $A$ over an algebraically closed field, $H^j(\operatorname{Spec}A,i)=0$ for $j\ge2$.
\end{theorem}
\leanref{GlobalTheory.\allowbreak{}EtC.\allowbreak{}etC\_\allowbreak{}c2}{Development.\allowbreak{}AbelianVarietyTheory.\allowbreak{}Proofs.\allowbreak{}DivisorProofs.\allowbreak{}SmoothAffineCurveEtaleCohomologicalDimensionLeOn}{9--15}
This is affine-curve cohomological dimension for Tate coefficients \cite{milne-1980}.

\begin{proposition}[Generic restriction, C3]\label{B9.22}
For an absolutely integrally closed local domain $T$ with fraction field $K$ and étale $T[t]$-algebra $B$, restriction to $K\otimes_T B$ is an isomorphism in every degree and twist.
\end{proposition}
\leanref{UpperTheory.\allowbreak{}Cb.\allowbreak{}etale\_\allowbreak{}cohomology\_\allowbreak{}ind\_\allowbreak{}etale\_\allowbreak{}base\_\allowbreak{}change}{Development.\allowbreak{}AutomorphicArithmetic.\allowbreak{}Foundations.\allowbreak{}EtaleClopenDescent.\allowbreak{}EtaleCohomologyIndEtaleBaseChange}{272--309}
\leanref{UpperTheory.\allowbreak{}Cb.\allowbreak{}C3.\allowbreak{}local\_\allowbreak{}eqv}{Development.\allowbreak{}AutomorphicArithmetic.\allowbreak{}Foundations.\allowbreak{}EtaleClopenDescent.\allowbreak{}EtaleCohomologyIndEtaleBaseChange}{210--247}
\begin{proof}
On strictly henselian ind-étale localizations of $B$, Tate trivialization identifies $H^0$ with that of the generic fibre. Higher groups vanish by strict henselianity. On the generic fibre use the cover lemma in degree one and henselian-line vanishing above it. The pullback square gives the comparison on $B$.
\end{proof}

\begin{theorem}[Affine lines, C4]\label{B9.23}
Projection from the affine line over any ring induces isomorphisms in every degree and twist.
\end{theorem}
\leanref{GlobalTheory.\allowbreak{}EtC.\allowbreak{}HomotopyInv.\allowbreak{}homotopy\_\allowbreak{}invariance\_\allowbreak{}etale}{Development.\allowbreak{}AbelianVarietyTheory.\allowbreak{}Proofs.\allowbreak{}DivisorProofs.\allowbreak{}HomotopyInvarianceEtale}{36--46}
This is étale homotopy invariance \cite{artin-grothendieck-verdier-1973}.

\begin{theorem}[Laurent lines, C5]\label{B9.24}
For strictly henselian local $A$, the Laurent line has degree-one group $\Z/b$ additively and vanishing in degrees at least two, in every twist.
\end{theorem}
\leanref{GlobalTheory.\allowbreak{}EtC.\allowbreak{}laurentPolynomial\_\allowbreak{}cohomology\_\allowbreak{}strictlyHenselian\_\allowbreak{}vanishing}{Development.\allowbreak{}AbelianVarietyTheory.\allowbreak{}Proofs.\allowbreak{}DivisorProofs.\allowbreak{}LaurentPolynomialCohomologyVanishing}{8--9}
Cup product with the Kummer class of $t$ identifies degree one with $H^0(\operatorname{Spec}A,i-1)$ \cite{milne-1980}.

\begin{proposition}[Projective bundles, D1]\label{B9.25}
For the lines $\pi:\mathbf P(V)\to\mathcal E$ of a rank-$N$ bundle on a simplicial presheaf, let $L$ be tautological. If $b$ is invertible on schemes with zero-simplices, the map $\sum_j\pi^*u_jc_1(L)^j$ identifies $\bigoplus_{0\le j<N}H^{n-2j}(\mathcal E,i-j)$ with $H^n(\mathbf P(V),i)$.
\end{proposition}
\leanref{GlobalTheory.\allowbreak{}EtC.\allowbreak{}projectiveBundleFormula\_\allowbreak{}vectorBundle\_\allowbreak{}any\_\allowbreak{}space}{Development.\allowbreak{}AbelianVarietyTheory.\allowbreak{}Proofs.\allowbreak{}DivisorProofs.\allowbreak{}ProjectiveBundleFormulaVectorBundleAnySpace}{11--13}
\leanref{GlobalTheory.\allowbreak{}T.\allowbreak{}projectiveBundleFormula\_\allowbreak{}discrete\_\allowbreak{}space}{Development.\allowbreak{}AbelianVarietyTheory.\allowbreak{}Proofs.\allowbreak{}DivisorProofs.\allowbreak{}ProjectiveBundleFormulaDiscreteSpaceStatement}{5--6}
\leanref{GlobalTheory.\allowbreak{}PBFg.\allowbreak{}pbf\_\allowbreak{}of\_\allowbreak{}disc}{Development.\allowbreak{}AbelianVarietyTheory.\allowbreak{}Proofs.\allowbreak{}DivisorProofs.\allowbreak{}ProjectiveBundleFormulaOfDiscreteSpaces}{65--77}
\begin{proof}
Apply the scheme formula \cite{milne-1980} on representables and discrete presheaves degreewise, then compare bounded-below total complexes.
\end{proof}

\begin{theorem}[Duality]\label{B9.26}
For smooth projective integral $Z$ over algebraically closed $k$ with $\ell$ invertible, rational $\ell$-adic cohomology is finite-dimensional in degrees $0$ through $2\dim Z$ with perfect complementary trace pairings after Tate trivialization.
\end{theorem}
\leanref{UpperTheory.\allowbreak{}Gal.\allowbreak{}KunnethLefschetz.\allowbreak{}QDual}{Development.\allowbreak{}AutomorphicArithmetic.\allowbreak{}Foundations.\allowbreak{}GaloisRepresentations.\allowbreak{}KunnethLefschetz.\allowbreak{}OneMul}{191--193}
\leanref{UpperTheory.\allowbreak{}Gal.\allowbreak{}KunnethLefschetz.\allowbreak{}Core.\allowbreak{}qDual}{Development.\allowbreak{}AutomorphicArithmetic.\allowbreak{}Foundations.\allowbreak{}GaloisRepresentations.\allowbreak{}KunnethLefschetz.\allowbreak{}EtSideGood}{45--46}
This is rational smooth projective Poincaré duality \cite{milne-1980}.

\begin{theorem}[Specialization]\label{B9.27}
For smooth proper $Y/O$ and a valuation ring $V/O$ with algebraically closed fraction and residue fields, restriction from $Y_V$ to either fibre is an isomorphism for moduli invertible in the residue field. Rational specialization at such $\ell$ respects pullbacks and compatible automorphisms.
\end{theorem}
\leanref{UpperTheory.\allowbreak{}Gal.\allowbreak{}Mod.\allowbreak{}spSpecIf}{Development.\allowbreak{}AutomorphicArithmetic.\allowbreak{}Foundations.\allowbreak{}GaloisRepresentations.\allowbreak{}BasicOpenSections.\allowbreak{}PSpSpecIf}{18--108}
\leanref{UpperTheory.\allowbreak{}Gal.\allowbreak{}Mod.\allowbreak{}specGood}{Development.\allowbreak{}AutomorphicArithmetic.\allowbreak{}Foundations.\allowbreak{}GaloisRepresentations.\allowbreak{}BasicOpenSections.\allowbreak{}EsGenAll}{8--9}
This is proper base change with generic local-base comparison over $V$ \cite[Exp.~XVI, {\S}2]{artin-grothendieck-verdier-1973}.

\section{Rank extraction and congruence limits}

Ordered jumps have the stronger gaps if $j_{i+1}-j_0>a(j_i-j_0)$ for $0\le i<a-1$. Nonzero roots are in general position if, in every enumeration,
\[
\prod_{i\ne j}\prod_{h=1}^a(\alpha_i^h-\alpha_j^h)
\prod_{j\ne i,\ j\ne k}(\alpha_i\alpha_k-\alpha_j^2)
\prod_{\substack{\sum m_i=a,\ m_i\ge0\\(m_i)\ne(1,\ldots,1)}}
\left(\prod_i\alpha_i^{m_i}-\prod_i\alpha_i\right)\ne0,
\]
where $m_i$ are integers.

\begin{proposition}[Seed tests]\label{B9.28}
In the odd-rank field setting of \ref{B9.08}, fix separated split $w_0,v'\mid\ell$. Choose a reference $U$ and finite sets $S,Y$. There is a finite set $Z_0$ of usable places away from $\ell$ separated from $w_0,S$. For every $U'$ unramified at nonsplit places off $T$, with elliptic type at $w_0$, reference weights off the pair, stronger gaps at $v'$, the spherical conditions of \ref{B9.08}, and general position on $Z_0$, there is a rank-$a$ $r$ with the following data. At usable $w\nmid\ell$ separated from $S$, cofinitely $r$ is unramified with polynomial $P_{w,B}$ for $U'$ on inverses of Frobenius lifts in a dense open set. At such $w\in Y$ with general-position $B$, this holds at all lifts. At $w\mid\ell$ in $Y$ separated from $w_0,v',S$, $r$ is crystalline de Rham with the prescribed jumps and has unramified Weil polynomials from each general-position $A_a(w,B)$.
\end{proposition}
\leanref{GlobalTheory.\allowbreak{}CongL.\allowbreak{}SeedWU4}{Development.\allowbreak{}AbelianVarietyTheory.\allowbreak{}Definitions.\allowbreak{}GaloisGeometry.\allowbreak{}LimWU4OfLimWU3}{15--36}
\leanref{UpperTheory.\allowbreak{}Gal.\allowbreak{}Int.\allowbreak{}Div.\allowbreak{}seedFinOdd\_\allowbreak{}of\_\allowbreak{}geoOne}{Development.\allowbreak{}AutomorphicArithmetic.\allowbreak{}Foundations.\allowbreak{}GaloisRepresentations.\allowbreak{}IntegralStructures.\allowbreak{}ZetaCharacterDivisibility.\allowbreak{}MainFin}{13--16}
\leanref{UpperTheory.\allowbreak{}Gal.\allowbreak{}Int.\allowbreak{}Div.\allowbreak{}seedFinOdd\_\allowbreak{}of}{Development.\allowbreak{}AutomorphicArithmetic.\allowbreak{}Foundations.\allowbreak{}GaloisRepresentations.\allowbreak{}IntegralStructures.\allowbreak{}ZetaCharacterDivisibility.\allowbreak{}SeedFinOddOf}{44--79}
\leanref{UpperTheory.\allowbreak{}Gal.\allowbreak{}Int.\allowbreak{}Div.\allowbreak{}divOne\_\allowbreak{}of}{Development.\allowbreak{}AutomorphicArithmetic.\allowbreak{}Foundations.\allowbreak{}GaloisRepresentations.\allowbreak{}IntegralStructures.\allowbreak{}ZetaCharacterDivisibility.\allowbreak{}PDivOne}{35--73}
\begin{proof}
Take $Z_0=\varnothing$. Translate the least jump to zero. The maximum of two disjoint subsets exceeds $a$ times each smaller jump, so its subset has larger mean. Reflection preserves lacunarity. Conjugate the auxiliary places and local data to match the target on $\Q(x)$. Use \ref{B9.08}, \ref{B9.29} and \ref{B9.30}, and compare the extracted representations on their dense common separable Frobenius loci.
\end{proof}

\begin{proposition}[Rank extraction]\label{B9.29}
Suppose an actual jump multiset of $\rho$ above $\ell$ contains a lacunary subset of size $a$. Suppose that on a dense set $\rho(g)$ has a monic degree-$a$ annihilator $Q_g$ with $(-1)^aQ_g(0)=\eta(g)$ for a continuous character $\eta$. There is a rank-$a$ $r$ of determinant $\eta$ with dense open separable locus and $r\hookrightarrow\rho\hookrightarrow r^{\oplus m}$ for some $m$.
\end{proposition}
\leanref{UpperTheory.\allowbreak{}Gal.\allowbreak{}Int.\allowbreak{}Div.\allowbreak{}DivSet}{Development.\allowbreak{}AutomorphicArithmetic.\allowbreak{}Foundations.\allowbreak{}GaloisRepresentations.\allowbreak{}IntegralStructures.\allowbreak{}ZetaCharacterDivisibility.\allowbreak{}PhiOfDivOf}{121--130}
\leanref{LateralTheory.\allowbreak{}DivDb.\allowbreak{}Key}{Development.\allowbreak{}EtaleCohomology.\allowbreak{}DivisorCycles.\allowbreak{}CycleCount}{51--55}
\leanref{LateralTheory.\allowbreak{}DivDb.\allowbreak{}plus\_\allowbreak{}of}{Development.\allowbreak{}EtaleCohomology.\allowbreak{}DivisorCycles.\allowbreak{}DivPlus}{37--70}
\leanref{UpperTheory.\allowbreak{}Gal.\allowbreak{}Int.\allowbreak{}Div.\allowbreak{}divSet\_\allowbreak{}holds}{Development.\allowbreak{}AutomorphicArithmetic.\allowbreak{}Foundations.\allowbreak{}GaloisRepresentations.\allowbreak{}IntegralStructures.\allowbreak{}ZetaCharacterDivisibility.\allowbreak{}DivisionOneGeometricRepresentation}{52--52}
\leanref{LateralTheory.\allowbreak{}DivDb.\allowbreak{}KeyPf.\allowbreak{}Key\_\allowbreak{}of}{Development.\allowbreak{}EtaleCohomology.\allowbreak{}DivisorCycles.\allowbreak{}SumMulSum}{90--176}
\begin{proof}
Dependence of $1,M,\ldots,M^a$ is closed. Apply it to $(M-1)/t$ along arcs at the identity to obtain the Lie algebra degree bound. The Sen point $X$ and its lacunary roots $e_i$ solve the tangent equations and $\det(e_i-X)=0$ with finitely many lacunarity inequations. An algebraically closed extension point thus gives a $\bar{\Q}_\ell$-point. Take its semisimple spectral projectors $p_i$.

Commuting perturbations separate the block spectra, so the degree bound forces scalar diagonal Lie blocks. Distinct nonzero adjoint differences project the off-diagonal blocks into the Lie algebra. Products through distinct indices are brackets. Opposite products $XY=p_i[X,Y]p_i$ are scalar. The generated algebra $A$ therefore has scalar corners and is normalized by the group. The radical has nonzero invariant annihilator on each irreducible summand of $\rho$, so it vanishes.

For $h=\rho(g)$, match $p_i$ to $h^{-1}p_jh$ in simple blocks. Fix $i$ in the matching if $p_iAhp_i\ne0$. Factor the matching unit into transvections $u$ and a diagonal unit. The square-zero Lie blocks exponentiate, so $u$ lies in the Zariski closure of the image of every finite-index normal subgroup. Then $hu$ permutes the projectors. Multiply by a diagonal $t$ in the closure. On a cycle of length $b$ the spectrum consists of the $b$-th roots of the return spectrum scaled by the product of the cycle's diagonal factors. Differentiation along $\sum_i e_ip_i$ separates the cycle means. Choose $t$ to avoid collisions across cycles. Each cycle contributes at least its length in roots, so the degree bound gives exactly $a$ roots and scalar fixed blocks. Hence $p_iAhp_i\ne0$ implies $p_ih=c_ip_i(ut)^{-1}\in p_iA$.

The span of $A\rho(G_E)$ has scalar corners and is semisimple by the same radical argument. One standard module per factor gives rank $a$ and the two embeddings by matrix units. For an inclusion $I$ and left inverse $L$, $r(g)=L\rho(g)I$. The discriminant at $LhutI$ is nonzero, so Zariski density gives a separable point in every open coset. On the dense separable test set $Q_g$ is the characteristic polynomial. Compare constant terms and use continuity.
\end{proof}

\begin{lemma}[Multiplicities]\label{B9.30}
Let an $a$-by-$a$ matrix $M$ be annihilated by the root polynomial of a size-$a$ multiset $S$, and suppose $\det M=\prod_{\alpha\in S}\alpha$. If $S$ is the unique size-$a$ multiset supported in $\operatorname{supp}S$ with this product, then $M$ has spectrum $S$.
\end{lemma}
\leanref{UpperTheory.\allowbreak{}Gal.\allowbreak{}Int.\allowbreak{}Div.\allowbreak{}genPin\_\allowbreak{}of}{Development.\allowbreak{}AutomorphicArithmetic.\allowbreak{}Foundations.\allowbreak{}GaloisRepresentations.\allowbreak{}IntegralStructures.\allowbreak{}ZetaCharacterDivisibility.\allowbreak{}PGenPin}{28--48}
\leanref{UpperTheory.\allowbreak{}Gal.\allowbreak{}Int.\allowbreak{}Div.\allowbreak{}divAll\_\allowbreak{}holds}{Development.\allowbreak{}AutomorphicArithmetic.\allowbreak{}Foundations.\allowbreak{}GaloisRepresentations.\allowbreak{}IntegralStructures.\allowbreak{}ZetaCharacterDivisibility.\allowbreak{}DivisionOneGeometricRepresentation}{54--58}
\leanref{UpperTheory.\allowbreak{}Gal.\allowbreak{}Int.\allowbreak{}Div.\allowbreak{}detEx\_\allowbreak{}holds}{Development.\allowbreak{}AutomorphicArithmetic.\allowbreak{}Foundations.\allowbreak{}GaloisRepresentations.\allowbreak{}IntegralStructures.\allowbreak{}ZetaCharacterDivisibility.\allowbreak{}DivisionOneGeometricRepresentation}{49--49}
\leanref{UpperTheory.\allowbreak{}Gal.\allowbreak{}Int.\allowbreak{}Div.\allowbreak{}phiSeed\_\allowbreak{}of}{Development.\allowbreak{}AutomorphicArithmetic.\allowbreak{}Foundations.\allowbreak{}GaloisRepresentations.\allowbreak{}IntegralStructures.\allowbreak{}ZetaCharacterDivisibility.\allowbreak{}PPhiSeed}{20--62}
\begin{proof}
The annihilator gives support containment and the determinant gives the product. In the application the normalized central character gives $\prod A_a(w,B)$ at degree $-1$, and the last general-position factor gives uniqueness. Above $\ell$, the embeddings transfer crystallinity and jumps, and \ref{B9.19} supplies the support.
\end{proof}

\begin{proposition}[Hecke approximation]\label{B9.31}
For odd-rank $U$ unramified at nonsplit places off $T$, fix separated split $w_0,v'\mid\ell$. Choose levels $K_w$ supporting $U$ off the pair, hyperspecial at spherical split places there. For finite sets $Z,W$ of such spherical places with $W$ above $\ell$, there are nonzero $c_{w,n}$ for $w\in W$, $n\in\Z$, that work for every index set $J$. For each $j\in J$ take a finitely supported $\ell$-integral operator at a split place away from $\ell$. The operator must act by a scalar on each smooth irreducible's $K_w$-invariants. Augment the eigenvalues by
\[
\iota^{-1}\!\left(c_{w,n}\sum_{\beta\in A_a(w,B_w)}\beta^n\right).
\]
Systems of $U'$ approximate this system as in \ref{B9.32}. They retain levels and weights off the pair and nonsplit unramifiedness off $T$. They have elliptic type at $w_0$ and are spherical with stronger gaps at $v'$. They have general position on $Z\cup W$.
\end{proposition}
\leanref{GlobalTheory.\allowbreak{}CongL.\allowbreak{}CongUW4}{Development.\allowbreak{}AbelianVarietyTheory.\allowbreak{}Definitions.\allowbreak{}GaloisGeometry.\allowbreak{}LimWU4OfLimWU3}{37--56}
\leanref{GlobalTheory.\allowbreak{}CongL.\allowbreak{}PowSatU}{Development.\allowbreak{}AbelianVarietyTheory.\allowbreak{}Definitions.\allowbreak{}GaloisGeometry.\allowbreak{}HasLevelOffOf}{51--53}
\leanref{UpperTheory.\allowbreak{}Gal.\allowbreak{}Int.\allowbreak{}Cong.\allowbreak{}congFinOdd\_\allowbreak{}of\_\allowbreak{}ex}{Development.\allowbreak{}AutomorphicArithmetic.\allowbreak{}Foundations.\allowbreak{}GaloisRepresentations.\allowbreak{}IntegralStructures.\allowbreak{}CongruenceOddLevel.\allowbreak{}CongFinOddOfLvLac}{15--16}
\leanref{UpperTheory.\allowbreak{}Gal.\allowbreak{}Int.\allowbreak{}Cong.\allowbreak{}congUWG\_\allowbreak{}of}{Development.\allowbreak{}AutomorphicArithmetic.\allowbreak{}Foundations.\allowbreak{}GaloisRepresentations.\allowbreak{}IntegralStructures.\allowbreak{}CongruenceOddLevel.\allowbreak{}NotApartOfOverPair}{111--250}
\leanref{UpperTheory.\allowbreak{}Gal.\allowbreak{}Int.\allowbreak{}Cong.\allowbreak{}CongUWG.\allowbreak{}congUW4}{Development.\allowbreak{}AutomorphicArithmetic.\allowbreak{}Foundations.\allowbreak{}GaloisRepresentations.\allowbreak{}IntegralStructures.\allowbreak{}CongruenceOddLevel.\allowbreak{}NotApartOfOverPair}{252--264}
\begin{proof}
Adjoin the power sums and general-position tests. Make the test product nonzero. Impose type and gaps at precision preserving nonvanishing and then project. For rational operators, conjugate tuples give linearly independent forms in a finite-dimensional span. The orbit is finite, so all coordinates lie in a single finite extension of $\Q_\ell$ \cite{chenevier-harris-2013}.
\end{proof}

\begin{definition}[Polynomial approximation]\label{B9.32}
For $\vartheta\in k^J$ and $\mathcal S\subset k^J$ over an ultrametric field $k$, $\operatorname{PolyApprox}(\vartheta,\mathcal S)$ means that each $\delta>0$ admits finite $s\subset\mathcal S$ and $c_\eta\in k$ with
\[
\left|P(\vartheta)-\sum_{\eta\in s}c_\eta P(\eta)\right|\le\delta
\qquad(P\in\Z[X_j:j\in J]).
\]
\end{definition}
\leanref{GlobalTheory.\allowbreak{}Cong.\allowbreak{}PolyApprox}{Development.\allowbreak{}AbelianVarietyTheory.\allowbreak{}Definitions.\allowbreak{}GaloisGeometry.\allowbreak{}PolyApproxZeroPadic}{10--13}

\begin{proposition}[Trace limit]\label{B9.33}
Let $x:J\to\Gamma$ have dense range in a compact group and $\vartheta:J\to L$ with $[L:\Q_\ell]<\infty$. If sampled continuous rank-$a$ traces give \ref{B9.32}, then $\vartheta$ extends to a rank-$a$ trace with the same approximation.
\end{proposition}
\leanref{GlobalTheory.\allowbreak{}Cong.\allowbreak{}exists\_\allowbreak{}limit\_\allowbreak{}rep}{Development.\allowbreak{}AbelianVarietyTheory.\allowbreak{}Proofs.\allowbreak{}InertiaEigenvalues.\allowbreak{}ExistsLimitRep}{14--29}
\begin{proof}
Use $X_i-X_j$ and continuity of finite trace sums for the local Cauchy condition. Extend into $L$ and pass the inequalities to finite tuples. Send the error to zero in the pseudocharacter identities. Apply \ref{B9.34}.
\end{proof}

\begin{theorem}[Pseudocharacters]\label{B9.34}
Over an algebraically closed field of characteristic zero, every dimension-$a$ group pseudocharacter is the trace of a semisimple rank-$a$ representation. Over a topological coefficient field, continuity of its trace gives continuity of the representation.
\end{theorem}
\leanref{GlobalTheory.\allowbreak{}PsChar.\allowbreak{}exists\_\allowbreak{}ssRep}{Development.\allowbreak{}AbelianVarietyTheory.\allowbreak{}Proofs.\allowbreak{}InertiaEigenvalues.\allowbreak{}FiniteSemisimple}{64--146}
\leanref{GlobalTheory.\allowbreak{}SSCont.\allowbreak{}continuous\_\allowbreak{}of\_\allowbreak{}continuous\_\allowbreak{}trace}{Development.\allowbreak{}AbelianVarietyTheory.\allowbreak{}Proofs.\allowbreak{}InertiaEigenvalues.\allowbreak{}ContinuousOfContinuousTrace}{60--70}
This is Taylor--Rouquier for groups, with continuity from a trace-dual basis of the image algebra \cite{taylor-1991,rouquier-1996}.

\begin{proposition}[Removing gaps]\label{B9.35}
Under the odd-rank field hypotheses take separated split $w_0,v'\mid\ell$ and automorphic $U$ unramified at nonsplit places off $T$. There is a rank-$a$ $r$ with full spherical data at usable $w\nmid\ell$ and crystalline de Rham with prescribed jumps at spherical $w\mid\ell$ separated from the pair. For each such $w\mid\ell$, Satake multiset $B$, compatible period datum, prolongation, bridge and Weil element $\sigma$, there are seeds and common jumps giving \ref{B9.36} with $S_n=\iota^{-1}(\sum_{\alpha\in A_a(w,B)}\alpha^{-n\deg\sigma})$.
\end{proposition}
\leanref{UpperTheory.\allowbreak{}Gal.\allowbreak{}Lim.\allowbreak{}limFinOdd\_\allowbreak{}of}{Development.\allowbreak{}AutomorphicArithmetic.\allowbreak{}Foundations.\allowbreak{}GaloisRepresentations.\allowbreak{}TraceLimits.\allowbreak{}LimFinOddOf}{197--198}
\leanref{UpperTheory.\allowbreak{}Gal.\allowbreak{}Lim.\allowbreak{}exists\_\allowbreak{}trace\_\allowbreak{}limit}{Development.\allowbreak{}AutomorphicArithmetic.\allowbreak{}Foundations.\allowbreak{}GaloisRepresentations.\allowbreak{}TraceLimits.\allowbreak{}ExistsTraceLimit}{15--170}
\leanref{UpperTheory.\allowbreak{}Gal.\allowbreak{}Lim.\allowbreak{}limWU4\_\allowbreak{}of}{Development.\allowbreak{}AutomorphicArithmetic.\allowbreak{}Foundations.\allowbreak{}GaloisRepresentations.\allowbreak{}TraceLimits.\allowbreak{}LimFinOddOf}{16--196}
\leanref{GlobalTheory.\allowbreak{}CongL.\allowbreak{}crysDRU\_\allowbreak{}of\_\allowbreak{}linApprox}{Development.\allowbreak{}AbelianVarietyTheory.\allowbreak{}Proofs.\allowbreak{}SatakeParameters.\allowbreak{}BlowUpFibreSemilocal}{67--97}
\begin{proof}
Sample rational Hecke coordinates at inverse-Frobenius powers $x_j$. The finite split set without Satake data is conjugation-stable. For seed values $\eta_i$, the sets $\{j:\eta_i(j)=\operatorname{tr}\rho_i(x_j)\}$ have finite intersections with dense image by \ref{B9.28}. Restrict to samples near $g$ and use continuity of the finite trace sums. The samples are Cauchy in the finite extension $L/\Q_\ell$ of \ref{B9.31}. Completeness gives the limit and the product filters transfer the inequalities and local tests. Apply \ref{B9.33}. Away from $\ell$, $t(gs)=t(g)$ for $t=\operatorname{tr}r$ and inertia $s$ gives $r(s)=1$ by trace duality. Newton identities give the polynomial. Above $\ell$, retain the common jumps and use the coordinates of exponent $-n\deg\sigma$ with tolerance adjusted for the factors $c_{w,-n\deg\sigma}$. Apply \ref{B9.36}.
\end{proof}

\begin{lemma}[Uniform denominator for a Weil element]\label{B9.35a}
Under \ref{B9.36}, fix one Weil element. The limit and layer period operators $V,V_m$ have a common denominator $s\in A_{\rm cris}$ invertible in $B_{\rm cris}$. For the transition $\Pi_m$ to the limit and identity columns $e,e_m$, at every precision $L$ some $m\ge L$ satisfies $V\Pi_m\equiv\Pi_mV_m$ and $\Pi_me_m\equiv e\pmod{\ell^L}$.
\end{lemma}
\leanref{UpperTheory.\allowbreak{}Gal.\allowbreak{}ClosedPhiFromClosedTrace.\allowbreak{}coverJ\_\allowbreak{}holds}{Development.\allowbreak{}AutomorphicArithmetic.\allowbreak{}Foundations.\allowbreak{}GaloisRepresentations.\allowbreak{}ClosedPhiFunctoriality.\allowbreak{}CoverJ}{23--125}
\leanref{UpperTheory.\allowbreak{}Gal.\allowbreak{}ClosedPhiFromClosedTrace.\allowbreak{}frEx\_\allowbreak{}holds}{Development.\allowbreak{}AutomorphicArithmetic.\allowbreak{}Foundations.\allowbreak{}GaloisRepresentations.\allowbreak{}ClosedPhiFunctoriality.\allowbreak{}FrExHolds}{11--38}
\leanref{UpperTheory.\allowbreak{}Gal.\allowbreak{}ClosedPhiFromClosedTrace.\allowbreak{}opEx\_\allowbreak{}holds}{Development.\allowbreak{}AutomorphicArithmetic.\allowbreak{}Foundations.\allowbreak{}GaloisRepresentations.\allowbreak{}ClosedPhiFunctoriality.\allowbreak{}OpExHolds}{72--103}
\begin{proof}
Use the $\Z_\ell$-span of $r(g)$. Trace duality bounds the loss in recovering $x$ from $\operatorname{tr}(xr(g))$. Keep earlier seeds with weight zero when unused. The layer is the span of $(\rho_i(g))_i$ plus $\ell^b$ times a full stable product lattice. Test translates of tuple relations against the uniform trace estimate. Their evaluations at $r$ vanish modulo $\ell^m$ after absorbing the trace loss in the initial precision. For large $b$, elementary divisors put the intersection with the added lattice in any prescribed power of $\ell$ times the tuple span. Absorb the weight denominators in $b$ as well. Evaluation then defines $\Pi_m$ modulo $\ell^m$ with the identity congruence and the weighted basis-trace estimates. The layer's rational left action is a sum of scalar restrictions of the seeds.

Put $N=\max(\{0\}\cup\{|j|:j\in J_\tau\text{ for some }\tau\})$. The crystalline Breuil--Kisin frames have $AA'=E(u)^{2N}I$ for the uniformizer's Eisenstein polynomial $E(u)$. The limit frame is linked to the layers at arbitrarily high precision. For the Frobenius image $t$ of the integral period, the comparison matrices satisfy $YY'=t^{2N}I$ and $QQ'=\ell^{2k}I$. The divided-power bound at height $2N+1$ gives $k=(2N+1)n_0$ for the period datum, uniformly in rank. Set $Z=YQ$ and $Z'=Q'Y'$. With $b=t^{2N}\ell^{2k}$ we have $ZZ'=Z'Z=bI$. Let $n_1,n_2$ be the negative and positive parts of $[k(w):\F_\ell]\deg\sigma$ and let $\phi$ denote crystalline Frobenius. Then
\[
V=\phi^{n_2}(b)\phi^{n_1}(Z)\phi^{n_2}(Z'),
\qquad s=\phi^{n_2}(t)^{3N}\ell^{4k}\phi^{n_1}(t)^N.
\]
For the period basis $U=t^{-N}Z$ one has $V\phi^{n_2}(U)=s\phi^{n_1}(U)$. The frame links give $\Pi_m\phi^{n_i}(Z_m)\equiv\phi^{n_i}(Z)R\pmod{\ell^L}$ with common $R$. Multiply the link for $i=2$ on the left by $\phi^{n_2}(Z')$ and on the right by $\phi^{n_2}(Z'_m)$. Substitute in the link for $i=1$ to get $V\Pi_m\equiv\Pi_mV_m$. The formula for $s$ is independent of $L$.
\end{proof}

\begin{proposition}[Crystalline limits]\label{B9.36}
Let $r$ and the continuous seeds have the same positive rank, with $r$ semisimple. Fix compatible period data at $w\mid\ell$ for which every seed is crystalline with the same jumps $J_\tau$ at each induced embedding. Fix $\sigma\in W_{E_w}$. Suppose finite weights in $\bar{\Q}_\ell$ approximate $\operatorname{tr}r$ uniformly by seed traces and $S_n\in\bar{\Q}_\ell$ by transported seed Weil traces at $\sigma^n$, jointly to arbitrary precision for $0\le n\le N$ and each $N$. Then $r_w$ is crystalline and its Weil trace at $\sigma^n$ is $\iota(S_n)$ for every $n\ge0$.
\end{proposition}
\leanref{GlobalTheory.\allowbreak{}CongPhi.\allowbreak{}ClosedPhi}{Development.\allowbreak{}AbelianVarietyTheory.\allowbreak{}Definitions.\allowbreak{}NonGaloisDescent.\allowbreak{}JointApproxOfMem}{79--94}
\leanref{GlobalTheory.\allowbreak{}LimitTop.\allowbreak{}closedAtL}{Development.\allowbreak{}AbelianVarietyTheory.\allowbreak{}Proofs.\allowbreak{}SatakeParameters.\allowbreak{}ClosedAtL}{25--36}
\leanref{UpperTheory.\allowbreak{}Gal.\allowbreak{}ClosedPhiFromClosedTrace.\allowbreak{}trStep\_\allowbreak{}holds}{Development.\allowbreak{}AutomorphicArithmetic.\allowbreak{}Foundations.\allowbreak{}GaloisRepresentations.\allowbreak{}ClosedPhiFunctoriality.\allowbreak{}TrStep}{83--162}
\leanref{UpperTheory.\allowbreak{}Gal.\allowbreak{}ClosedPhiFromClosedTrace.\allowbreak{}trLim\_\allowbreak{}holds}{Development.\allowbreak{}AutomorphicArithmetic.\allowbreak{}Foundations.\allowbreak{}GaloisRepresentations.\allowbreak{}ClosedPhiFunctoriality.\allowbreak{}TrLimHolds}{8--58}
\leanref{UpperTheory.\allowbreak{}Gal.\allowbreak{}ClosedPhiFromClosedTrace.\allowbreak{}congPhi\_\allowbreak{}holds}{Development.\allowbreak{}AutomorphicArithmetic.\allowbreak{}Foundations.\allowbreak{}GaloisRepresentations.\allowbreak{}ClosedPhiFunctoriality.\allowbreak{}ClosedPhiHolds}{44--45}
\begin{proof}
Fixed-jump closedness for uniform weighted trace approximation gives crystallinity. Fix a power of $\sigma$. Write $p,\theta$ for its transported limit trace and target, and $p_k$ for its trace on seed $k$. Let $K^{\rm ur}$ be the maximal unramified algebraic extension of $\Q_\ell$ in $\bar{\Q}_\ell$, with period embedding $u:K^{\rm ur}\to B_{\rm cris}$. Write $t_i,t'_{jk}$ for the limit and seed components of the layer basis traces, respectively. The period comparison for $V$ gives, for each $K^{\rm ur}$-linear $\lambda:\bar{\Q}_\ell\to K^{\rm ur}$,
\[
\sum_i(Ve)_i u(\lambda(t_i))=s\,u(\lambda(p)).
\]
The corresponding layer identity has $V_me_m$, $t'_{jk}$ and $p_k$ in these positions, with the same $s$. Apply it for seed $k$ to $x\mapsto\lambda(c_kx)$, where $c_k$ is the layer weight. Thus the weights enter inside $\lambda$.

Suppress the subscript $m$ on $\Pi_m$. Put $\delta_j=\sum_i\Pi_{ij}t_i-\sum_k c_k t'_{jk}$ and $\eta=\theta-\sum_k c_kp_k$. Subtract the summed layer identities from the first identity and rearrange to get
\[
s\,u(\lambda(p-\theta))
=\sum_i(Ve-\Pi V_me_m)_i u(\lambda(t_i))
 +\sum_j(V_me_m)_j u(\lambda(\delta_j))-s\,u(\lambda(\eta)).
\]
Here $\Pi$ acts on period columns through the integral scalar embedding. The first column on the right is divisible by $\ell^L$ by \ref{B9.35a}. The basis and target errors satisfy $|\delta_j|,|\eta|\le\ell^{-L}$.

If $p\ne\theta$, choose $c_1,c_2\ge0$ so $\ell^{c_1}t_i$ and $z=\ell^{c_1}(p-\theta)$ are integral and $|z|\ge\ell^{-c_2}$. In a finite extension of $K^{\rm ur}$ containing these elements and the divided errors $\delta_j/\ell^L,\eta/\ell^L$, use an integral basis and project to a maximal coordinate of $z$. Extend the projection linearly to $\bar{\Q}_\ell$. Its values on the integral elements listed above are integral and $q_L=\lambda(z)$ satisfies $|q_L|\ge|z|$. Integral values in $K^{\rm ur}$ lift under $u$ to $A_{\rm cris}$. Multiply the displayed identity by $\ell^{c_1}$ and use the three error bounds to get $s\,u(q_L)\in\ell^L A_{\rm cris}$. The integral scalar $\ell^{c_2}/q_L$ also lifts. Hence $\ell^{c_2}s\in\ell^L A_{\rm cris}$ for every $L$. Separatedness and $\ell$-torsion freeness force $s=0$. This contradicts its invertibility in $B_{\rm cris}$. Thus $p=\theta$.
\end{proof}

\begin{lemma}[Weak descent]\label{B9.36a}
Suppose $E/F$ is CM and unramified at finite places. A regular algebraic conjugate self-dual cuspidal $\pi$ of positive rank, spherical at every nonsplit place, has a definite-unitary automorphic match at almost every split place, unramified at every nonsplit place.
\end{lemma}
\leanref{UpperTheory.\allowbreak{}V.\allowbreak{}Z.\allowbreak{}Tf.\allowbreak{}GaloisFromTraceFormula.\allowbreak{}asmTh\_\allowbreak{}of}{Development.\allowbreak{}AutomorphicArithmetic.\allowbreak{}AutomorphicVarieties.\allowbreak{}ModularityArgument.\allowbreak{}TraceFormula.\allowbreak{}GaloisFromTraceFormula.\allowbreak{}AsmThOf}{10--113}
\leanref{UpperTheory.\allowbreak{}Gal.\allowbreak{}Desc.\allowbreak{}weakAll}{Development.\allowbreak{}AutomorphicArithmetic.\allowbreak{}Foundations.\allowbreak{}GaloisRepresentations.\allowbreak{}BaseChangeDescent.\allowbreak{}BcBackOffT}{26--26}
\begin{proof}
Use the weight of $\pi$ and the nonsplit unramified test in the stabilized comparison. Suppose every unitary candidate and every regular cuspidal list omitting $\pi$ with inversion-invariant union at almost every nonsplit place differs from $\pi$ at a split place. A symmetric polynomial test at split places is one on $\pi$ and zero on the finite competing spectra. Its positive powers give zero unitary trace. Terms with two positive block sizes tend to zero on the stabilized side. The end terms all have sign $+1$ and tend to the same nonzero product of local traces. At least one end term occurs, a contradiction.

For a matching cuspidal list of length $j$, write each factor as a normalized cuspidal representation twisted by $|\det|^s$. The normalized exponents have total sum zero, so the reality of the matching exponents gives $s\in\R$. Comparing conjugate embeddings gives $2s\in\Z$. At a matching split place the normalized determinant has modulus one, while every matching root has modulus in $(q_w^{-1/2},q_w^{1/2})$ by the Jacquet--Shalika bounds for $\pi$. Take the geometric mean of the factor's root moduli to get $-1/2<s<1/2$. Thus $s=0$ and the original factors are unitary.

Take $M=16(j+1)^2$ twists of $\pi$ by powers of the unitary character $\mu^2$. At almost every nonsplit place $\mu^2$ has value one, whereas the sums of infinitesimal weights of the twists are distinct at infinity. The average of the twists has squared norm $1/M$ in the Rankin--Selberg pole pairing, given by prime-power sums divided by $-\log(s-1)$ as $s\downarrow1$. Since the twists coincide on the nonsplit part, the limit superior of the squared norm of $\pi$ there is at most $1/M$. If each factor differs from $\pi$, pairing $\pi$ with $\pi$ minus the sum of the factors gives one. Split power sums cancel, so Cauchy--Schwarz on the nonsplit part gives $1\le1/M+j/\sqrt M$, a contradiction.
\end{proof}

\begin{proposition}[Unitary descent]\label{B9.37}
For $\chi=\chi_{E/F}^a$, assume an imaginary quadratic subfield and even $[F:\Q]$. Assume the nonsplit finite places lie away from $\ell$ and are unramified and spherical. Given finite $Z$, there is a definite-unitary automorphic $U$ that matches $\pi$ at split places and has weights giving its infinitesimal characters. At nonsplit places it is unramified off a finite nonsplit $T$ disjoint from $Z$.
\end{proposition}
\leanref{GlobalTheory.\allowbreak{}Cong.\allowbreak{}DescOffT}{Development.\allowbreak{}AbelianVarietyTheory.\allowbreak{}Definitions.\allowbreak{}GaloisGeometry.\allowbreak{}MatchOfMatchOffTEmpty}{32--37}
\leanref{UpperTheory.\allowbreak{}Gal.\allowbreak{}Desc.\allowbreak{}descOffT}{Development.\allowbreak{}AutomorphicArithmetic.\allowbreak{}Foundations.\allowbreak{}GaloisRepresentations.\allowbreak{}BaseChangeDescent.\allowbreak{}BcBackOffT}{31--32}
\begin{proof}
The quadratic character $\eta_{E/F}=\chi_{E/F}$ is trivial on norms, so $\pi^c\simeq\pi^\vee$. Take $U$ from \ref{B9.36a}. A Speh factor of length at least two gives consecutive roots with modulus ratio $q_w$ at a matching split place. Both moduli lie strictly between $q_w^{-1/2}$ and $q_w^{1/2}$, so their ratio is strictly less than $q_w$. All lengths are one. Split-place rigidity identifies one factor with $\pi$. Its rank is $a$, the total rank. Multiplicity one identifies its realization and gives the split and archimedean data with $T=\varnothing$.
\end{proof}

\begin{lemma}[A ramified place for quadratic base change]\label{B9.37a}
Let $M/K$ be quadratic with $K$ totally complex. Let $\pi$ be cuspidal of rank $a\ge1$ over $K$. If $M/K$ ramifies over a spherical place of $\pi$, there is a cuspidal lift. At $w\mid v$ it has multiset $\{\alpha^{f(w/v)}:\alpha\in B_v\}$ whenever $B_v$ exists, with residue degree $f(w/v)$. Infinitesimal characters pull back by restriction of embeddings. Local components agree under continuous completion isomorphisms extending $K\to M$.
\end{lemma}
\leanref{GlobalTheory.\allowbreak{}QuadLiftBridge.\allowbreak{}quadLift}{Development.\allowbreak{}AbelianVarietyTheory.\allowbreak{}Proofs.\allowbreak{}SatakeParameters.\allowbreak{}QuadLift}{14--24}
\begin{proof}
The quadratic character is ramified at the designated spherical place. Twisting changes the unramified cuspidal support there, so strong multiplicity one excludes a quadratic self-twist. Prime-degree cyclic base change gives a cuspidal lift with the stated local data. The norm raises unramified parameters to the residue degree \cite[Ch.~III, Thms.~4.2 and 5.1]{arthur-clozel-1989}.
\end{proof}

\begin{proposition}[Spherical theorem]\label{B9.38}
Every polarized input has rank-$a$ $r$ unramified off a finite set, with inverse-Frobenius polynomial $P_{w,B_w}$ for some Satake multiset $B_w$ at every place outside that set. Every such $r$ has these data for all $B$ at spherical $w\nmid\ell$. At spherical $w\mid\ell$ it is crystalline de Rham with prescribed jumps and unramified Weil polynomials from $A_a(w,B)$.
\end{proposition}
\leanref{UpperTheory.\allowbreak{}Gal.\allowbreak{}Int.\allowbreak{}Core.\allowbreak{}toCoreK}{Development.\allowbreak{}AutomorphicArithmetic.\allowbreak{}Foundations.\allowbreak{}GaloisRepresentations.\allowbreak{}IntegralStructures.\allowbreak{}CrystallineDescent.\allowbreak{}ToCoreK}{7--10}
\leanref{UpperTheory.\allowbreak{}Gal.\allowbreak{}Int.\allowbreak{}Inp.\allowbreak{}coreAll\_\allowbreak{}of\_\allowbreak{}core}{Development.\allowbreak{}AutomorphicArithmetic.\allowbreak{}Foundations.\allowbreak{}GaloisRepresentations.\allowbreak{}IntegralStructures.\allowbreak{}CrystallineInputs.\allowbreak{}CoreAllOfRowStarCore}{4--6}
\leanref{UpperTheory.\allowbreak{}Gal.\allowbreak{}Int.\allowbreak{}Core.\allowbreak{}patchOne\_\allowbreak{}holds}{Development.\allowbreak{}AutomorphicArithmetic.\allowbreak{}Foundations.\allowbreak{}GaloisRepresentations.\allowbreak{}IntegralStructures.\allowbreak{}CrystallineDescent.\allowbreak{}PatchOne}{83--128}
\leanref{UpperTheory.\allowbreak{}Gal.\allowbreak{}Int.\allowbreak{}Core.\allowbreak{}stepDown\_\allowbreak{}of}{Development.\allowbreak{}AutomorphicArithmetic.\allowbreak{}Foundations.\allowbreak{}GaloisRepresentations.\allowbreak{}IntegralStructures.\allowbreak{}CrystallineDescent.\allowbreak{}JoinEndHolds}{92--135}
\leanref{UpperTheory.\allowbreak{}Gal.\allowbreak{}Int.\allowbreak{}Inp.\allowbreak{}coreAll\_\allowbreak{}of\_\allowbreak{}core\_\allowbreak{}inp}{Development.\allowbreak{}AutomorphicArithmetic.\allowbreak{}Foundations.\allowbreak{}GaloisRepresentations.\allowbreak{}IntegralStructures.\allowbreak{}CrystallineInputs.\allowbreak{}CoreAllOfRowStarCoreInp}{538--551}
\leanref{UpperTheory.\allowbreak{}Gal.\allowbreak{}Int.\allowbreak{}Core.\allowbreak{}joinU\_\allowbreak{}of}{Development.\allowbreak{}AutomorphicArithmetic.\allowbreak{}Foundations.\allowbreak{}GaloisRepresentations.\allowbreak{}IntegralStructures.\allowbreak{}CrystallineDescent.\allowbreak{}JoinEndHolds}{19--62}
\leanref{GlobalTheory.\allowbreak{}GoodCsdTwist.\allowbreak{}exists\_\allowbreak{}goodCsd\_\allowbreak{}twist}{Development.\allowbreak{}AbelianVarietyTheory.\allowbreak{}Proofs.\allowbreak{}SatakeParameters.\allowbreak{}ExistsGoodConjSelfDualTwistSplitUnramified}{65--111}
\begin{proof}
First work over an auxiliary field for quadratic polarization. Assume the hypotheses of \ref{B9.37}, the inert-place conditions of \ref{B9.41}, three pairwise separated places above $\ell$, and $x^2=-d$ in $E$ with $-d\in(\Q_\ell^\times)^2$. Take $U$ by \ref{B9.37}. Degree-one Chebotarev gives dense inverse Frobenius on the usable set $Q$ after finite removals. In odd rank apply \ref{B9.35} and \ref{B9.36} and use Newton identities. For a target above $\ell$, choose two of the three places separated from it and identify the resulting witness with the first by density. In even rank apply \ref{B9.45} on $Q$.

Descend the quadratic-polarization case by successive quadratic CM steps. Two preliminary steps split the places over $\ell$. The second extension and every further extension have three separated places. The final family of quadratic extensions provides the auxiliary inert-place and square conditions. In each descent choose infinitely many extensions split at the target and the places lacking Satake data, with ramification at a fresh spherical place. Requiring splitting at earlier ramification witnesses gives distinct extensions. Use \ref{B9.37a} for cuspidality. On pairwise composita, degree-one Frobenius gives equal traces of the restricted representations, and the same comparison gives conjugation invariance. Index-two patching gives a representation with the required restrictions outside finitely many choices. A splitting extension outside these exceptions has the same completion at the target, so the local data descend. Density and semisimplicity identify any two representations with the cofinite data.

For general polarization on a good field, assume an imaginary quadratic subfield and even $[F:\Q]$. Assume the nonsplit places away from $\ell$ are spherical and unramified. Let $b$ be the parallel infinity exponent of $\chi$ on positive real idèles. The character $\chi^{-1}\eta_{E/F}^{a+|b|}$ has the sign of exponent $-b$. Since the unramified norm is surjective on units, sphericality and polarization imply that it is unramified at nonsplit places. Extend it algebraically as $\kappa$ on $E$, unramified at the specified split place and at nonsplit places. Then $\kappa\kappa^c=(\chi\circ N_{E/F})^{-1}$ and $\pi\otimes\kappa$ has quadratic polarization. Twist its representation back by $(\kappa^{-1})_\ell$. For arbitrary fields use quadratic base change to good fields and the same splitting-place descent.
\end{proof}

\section{Two theta lifts and even rank}

\begin{proposition}[Stable range]\label{B9.39}
Let $E/F$ be CM unramified at finite places with an imaginary quadratic subfield and inert places of odd characteristic and norm $>3$. The datum $B$ selects complex embeddings and split embeddings at finite $B.P$. Its additive character is trivial on $F$ and equals $\exp(-2\pi i\operatorname{Tr}(x))$ at infinity. It is unramified at $B.P$ and inert places. For the definite Hermitian matrix $1_n$, $n\ge1$, there is a bound $d_0$ for definite skew-Hermitian targets $\delta1_d$. For $d\ge d_0$, use splitting characters $\chi_X,\chi_W$ restricting to quadratic powers $n,d$ with exponents $m_X,m_W$. An automorphic $U$ with $2\lambda_{v,k}\ge d+m_W(v)$ for $0\le k<n$ has a lift of weights
\[
2\lambda'_{v,k}=n+m_X(v)+
\begin{cases}
2\lambda_{v,k}-d-m_W(v)&k<n,\\
0&n\le k<d.
\end{cases}
\]
At finite split embeddings the components $\tau,\tau'$ admit a nonzero $\beta:\mathcal S(M_{n,d}(F_v))\times\tau\to\tau'$, linear in the first slot and conjugate-linear in the second, with $\beta(\omega(g,h)\Phi,\tau(g)u)=\tau'(h)\beta(\Phi,u)$ for the split Weil action.
\end{proposition}
\leanref{UpperTheory.\allowbreak{}TP.\allowbreak{}TLiftK}{Development.\allowbreak{}AutomorphicArithmetic.\allowbreak{}Foundations.\allowbreak{}GaloisRepresentations.\allowbreak{}ThetaParseval.\allowbreak{}HalfIndexLattices.\allowbreak{}TLiftK}{12--24}
\leanref{UpperTheory.\allowbreak{}TP.\allowbreak{}W.\allowbreak{}GoodK}{Development.\allowbreak{}AutomorphicArithmetic.\allowbreak{}Foundations.\allowbreak{}GaloisRepresentations.\allowbreak{}ThetaParseval.\allowbreak{}HalfIndexLattices.\allowbreak{}SplitStab}{18--25}
\leanref{UpperTheory.\allowbreak{}TP.\allowbreak{}LocPair}{Development.\allowbreak{}AutomorphicArithmetic.\allowbreak{}Foundations.\allowbreak{}GaloisRepresentations.\allowbreak{}ThetaParseval.\allowbreak{}SchurStep}{93--101}
\leanref{UpperTheory.\allowbreak{}TP.\allowbreak{}IsLiftWt}{Development.\allowbreak{}AutomorphicArithmetic.\allowbreak{}Foundations.\allowbreak{}GaloisRepresentations.\allowbreak{}ThetaParseval.\allowbreak{}SchurStep}{103--104}
\leanref{UpperTheory.\allowbreak{}TP.\allowbreak{}IsLift}{Development.\allowbreak{}AutomorphicArithmetic.\allowbreak{}Foundations.\allowbreak{}GaloisRepresentations.\allowbreak{}ThetaParseval.\allowbreak{}SchurStep}{106--112}
\leanref{UpperTheory.\allowbreak{}TP.\allowbreak{}K.\allowbreak{}tLiftK}{Development.\allowbreak{}AutomorphicArithmetic.\allowbreak{}Foundations.\allowbreak{}GaloisRepresentations.\allowbreak{}ThetaParseval.\allowbreak{}ThetaKernels.\allowbreak{}TLiftK}{4--6}
\begin{proof}
Choose $d_0$ for the stable range and height bound. The compact polynomial model gives the weights. Take a constituent of the nonzero lift of \ref{B9.40} and evaluate the pairing at split places.
\end{proof}

\begin{lemma}[Integral criterion]\label{B9.40}
In \ref{B9.39}, let $C_\Phi$ be the Weil coefficient of polynomial times finite Schwartz data and $M_\varphi$ the coefficient of a continuous left $\U(X)(F)$-invariant form. Suppose the doubling height satisfies $\int_{\U(X)(\mathbf A_F)}h^t<\infty$ for $t\ge t_0$, with $t_0<d/2$. Then $I=\int C_\Phi\overline{M_\varphi}$ converges absolutely and vanishes for a zero lift. If $d\ge2n$ and intertwiners from one finite-dimensional irreducible archimedean module send the same vector to nonzero polynomial $f_0$ and nonzero $\varphi$ in irreducible automorphic $U$, some finite Schwartz data give a nonzero lift.
\end{lemma}
\leanref{UpperTheory.\allowbreak{}TP.\allowbreak{}Rallis\_\allowbreak{}k}{Development.\allowbreak{}AutomorphicArithmetic.\allowbreak{}Foundations.\allowbreak{}GaloisRepresentations.\allowbreak{}ThetaParseval.\allowbreak{}ThetaKernels.\allowbreak{}GalTPDefs}{61--68}
\leanref{UpperTheory.\allowbreak{}TP.\allowbreak{}NonVan\_\allowbreak{}k}{Development.\allowbreak{}AutomorphicArithmetic.\allowbreak{}Foundations.\allowbreak{}GaloisRepresentations.\allowbreak{}ThetaParseval.\allowbreak{}ThetaKernels.\allowbreak{}GalTPDefs}{114--121}
\leanref{UpperTheory.\allowbreak{}TP.\allowbreak{}rallis\_\allowbreak{}k}{Development.\allowbreak{}AutomorphicArithmetic.\allowbreak{}Foundations.\allowbreak{}GaloisRepresentations.\allowbreak{}ThetaParseval.\allowbreak{}ThetaKernels.\allowbreak{}TPGRallis}{99--225}
\leanref{UpperTheory.\allowbreak{}TP.\allowbreak{}nonVan\_\allowbreak{}k}{Development.\allowbreak{}AutomorphicArithmetic.\allowbreak{}Foundations.\allowbreak{}GaloisRepresentations.\allowbreak{}ThetaParseval.\allowbreak{}ThetaKernels.\allowbreak{}TPGNonVan}{26--154}
\begin{proof}
Apply \ref{B8.12} to the doubled section of $\Phi\otimes\bar\Phi$ with definite target and splitting character of restriction $\eta_{E/F}^d$. At $s_0=(d-n)/2$ the Eisenstein value is the theta average with factor $1$ for $d>n$ and $2$ otherwise \cite[Thm.~7.1]{gan-qiu-takeda-2014}. With $\epsilon$ the doubling embedding and $[G]=\U(X)(F)\backslash\U(X)(\mathbf A_F)$, pair by
\[
\int_{[G]^2} E(s,\epsilon(g_1,g_2))\,\chi_W(\det g_2)^{-1}
       \varphi(g_2)\overline{\varphi(g_1)}\,dg_1\,dg_2 .
\]
The inverse character cancels the splitting character of the doubled Weil action. Unfolding far to the right gives a nonzero constant times $\int h^{s-s_0}C_\Phi\overline{M_\varphi}$. The bounds $|C_\Phi|\ll h^{d/2}$ and $|M_\varphi|\le M_\varphi(1)$ give holomorphy for $\Re(s-s_0)+d/2>t_0$. At $s_0$, Fubini on the compact quotients gives a multiple of the squared lift norm.

For $d\ge2n$, first choose a finite set $S$ containing the ramification of $\chi_W$ and the bad levels of $\varphi$, large enough that $\int_{K^c}h^{d/2}<\operatorname{vol}(K)$ for the integral subgroup $K$ in the adèles off $S$. Height integrability gives this choice. Take translated vacua at $S$ and vacua elsewhere. At $S$ the lattice construction gives coefficients supported in small congruence subgroups fixing $\varphi$. Outside such a subgroup, a lattice translation acts on the vector and its translate by distinct unitary characters, so their inner product vanishes. Archimedean Schur orthogonality gives a nonzero common factor. After dividing out the common factors from infinity and $S$, the integral over $K$ is $M_\varphi(1)\operatorname{vol}(K)>0$. The vacuum bound and Cauchy--Schwarz bound the remaining integral by $M_\varphi(1)\int_{K^c}h^{d/2}$. Thus $I\ne0$.
\end{proof}

\begin{proposition}[Two lifts]\label{B9.41}
Under the field conditions of \ref{B9.39}, with even $[F:\Q]$ and inert $v\nmid\mathfrak D_{F/\Q}$, take definite-unitary automorphic $U$ of positive even rank $n$ with weights $\lambda_\tau$ at every embedding and a split place $w_0$. For some odd $t$ there are two lifts $U_i$ of rank $n+t$ and algebraic characters $\kappa,\psi_{ij}$ with disjoint $\psi$-families and $\kappa,\psi_{1j}$ unramified at $w_0$. At split $w$ where $\kappa$ and the $i$-th family are unramified, each Satake $B$ admits $B_i$ with
\[
A_{n+t}(w,B_i)=\kappa(w)^{-1}A_n(w,B)\sqcup\{\psi_{ij}(w)\}_{0\le j<t}.
\]
For infinity exponents $m_\tau,p_{ij,\tau}$ the jumps are $(J(\lambda_\tau)+m_\tau)\sqcup\{-p_{ij,\tau}\}_j$.
\end{proposition}
\leanref{UpperTheory.\allowbreak{}Gal.\allowbreak{}TP.\allowbreak{}twoUpOK\_\allowbreak{}of\_\allowbreak{}steps}{Development.\allowbreak{}AutomorphicArithmetic.\allowbreak{}Foundations.\allowbreak{}GaloisRepresentations.\allowbreak{}ThetaParseval.\allowbreak{}TwoUpOKOfSteps}{13--32}
\leanref{UpperTheory.\allowbreak{}Gal.\allowbreak{}TP.\allowbreak{}TwoUpOK}{Development.\allowbreak{}AutomorphicArithmetic.\allowbreak{}Foundations.\allowbreak{}GaloisRepresentations.\allowbreak{}ThetaParseval.\allowbreak{}AllCoreOf}{18--22}
\leanref{UpperTheory.\allowbreak{}Gal.\allowbreak{}TP.\allowbreak{}InertOK}{Development.\allowbreak{}AutomorphicArithmetic.\allowbreak{}Foundations.\allowbreak{}GaloisRepresentations.\allowbreak{}ThetaParseval.\allowbreak{}AllCoreOf}{12--16}
\leanref{UpperTheory.\allowbreak{}Gal.\allowbreak{}TP.\allowbreak{}Sub.\allowbreak{}IsUp}{Development.\allowbreak{}AutomorphicArithmetic.\allowbreak{}Foundations.\allowbreak{}GaloisRepresentations.\allowbreak{}ThetaParseval.\allowbreak{}SplitLineSubtraction.\allowbreak{}RowDefOfOdd}{85--95}
\leanref{UpperTheory.\allowbreak{}Gal.\allowbreak{}TP.\allowbreak{}twoUpOK\_\allowbreak{}of\_\allowbreak{}k}{Development.\allowbreak{}AutomorphicArithmetic.\allowbreak{}Foundations.\allowbreak{}GaloisRepresentations.\allowbreak{}ThetaParseval.\allowbreak{}TwoUpOKOfSteps}{33--33}
\begin{proof}
Take $t=2(d_0+n)+1$. For a weight lower bound $b$, choose target exponents $\le2b-(n+t)$. Apply \ref{B9.39} and \ref{B9.40} to both character choices in dimension $d=n+t>2n$. Put $x=\chi_X(w)$ and $y=\chi_W(w)$. Write $q=Nw$. Split spherical lifting gives $(xy^{-1})B\sqcup\{xq^{(t-1)/2-j}:0\le j<t\}$. Scale by $q^{(n+t-1)/2}$ and use $\kappa(w)=yx^{-1}q^{-t/2}$. Subtract the weights of \ref{B9.39} from their indices. Use dual weights at conjugate embeddings.
\end{proof}

\begin{lemma}[Frobenius identity]\label{B9.42}
Let $Q$ be split places with dense inverse Frobenius away from $\ell$ after every finite removal. Suppose continuous equal-rank $r_i$ and continuous characters $\chi_{ij}$ are cofinitely unramified there, with inverse-Frobenius roots of $r_i$ a common multiset plus its own family's values. Then for all $g\in G_E$,
\[
\det(X-r_1(g))\prod_j(X-\chi_{2j}(g))
=\det(X-r_2(g))\prod_j(X-\chi_{1j}(g)).
\]
\end{lemma}
\leanref{UpperTheory.\allowbreak{}Gal.\allowbreak{}TP.\allowbreak{}frobIdQ}{Development.\allowbreak{}AutomorphicArithmetic.\allowbreak{}Foundations.\allowbreak{}GaloisRepresentations.\allowbreak{}ThetaParseval.\allowbreak{}CoreAtQIffRowQ}{52--87}
\begin{proof}
On the dense set both sides multiply the common polynomial by both line products. Use coefficientwise continuity.
\end{proof}

\begin{lemma}[Line subtraction]\label{B9.43}
Under \ref{B9.42} with semisimple $r_i$ and disjoint families, $r_1\simeq s\oplus\bigoplus_j\chi_{1j}$. For line roots $c_j$ and jumps $m_{j,\tau}$, the multisets $S\sqcup\{c_j\}_j$ and $J_\tau\sqcup\{m_{j,\tau}\}_j$ for $r_1$ become $S,J_\tau$ for $s$. At $w\mid\ell$ suppose $r_1$ and the lines are crystalline and their Weil polynomials at each $g$ use their respective roots raised to $-\deg g$. Then $s$ is crystalline with polynomial $\prod_{\alpha\in S}(X-\alpha^{-\deg g})$.
\end{lemma}
\leanref{UpperTheory.\allowbreak{}Gal.\allowbreak{}TP.\allowbreak{}Sub.\allowbreak{}peelAll\_\allowbreak{}holds}{Development.\allowbreak{}AutomorphicArithmetic.\allowbreak{}Foundations.\allowbreak{}GaloisRepresentations.\allowbreak{}ThetaParseval.\allowbreak{}SplitLineSubtraction.\allowbreak{}PPeelAll}{17--52}
\leanref{UpperTheory.\allowbreak{}Gal.\allowbreak{}TP.\allowbreak{}Sub.\allowbreak{}eigTransfer\_\allowbreak{}holds}{Development.\allowbreak{}AutomorphicArithmetic.\allowbreak{}Foundations.\allowbreak{}GaloisRepresentations.\allowbreak{}ThetaParseval.\allowbreak{}SplitLineSubtraction.\allowbreak{}PEigTransfer}{13--35}
\leanref{UpperTheory.\allowbreak{}Gal.\allowbreak{}TP.\allowbreak{}Sub.\allowbreak{}splitLine\_\allowbreak{}holds}{Development.\allowbreak{}AutomorphicArithmetic.\allowbreak{}Foundations.\allowbreak{}GaloisRepresentations.\allowbreak{}ThetaParseval.\allowbreak{}SplitLineSubtraction.\allowbreak{}PSplitLine}{100--133}
\begin{proof}
Brauer--Nesbitt identifies the enlarged sums. The first-family lines project trivially to the second-family lines and hence split off $r_1$. Cancel the line factors in the Frobenius and Weil polynomials and subtract their jumps by additivity.
\end{proof}

\begin{proposition}[One place]\label{B9.44}
Assume the density of \ref{B9.42} and odd-rank spherical compatibility on $Q$. Suppose there are two lifts of rank-$n$ $U$ at $w_0\in Q$. There is a rank-$n$ $r$ weakly compatible on $Q$. It is unramified with polynomial $P_{w,B}$ for every Satake $B$ cofinitely on $Q$. This holds at $w_0\nmid\ell$. At $w_0\mid\ell$ it is crystalline de Rham with jumps $J(\lambda_{\iota\circ\tau})$ at induced $\tau$ and the Weil polynomials of \ref{B9.43} for $S=A_n(w_0,B)$ for every $B$ and compatible period datum.
\end{proposition}
\leanref{UpperTheory.\allowbreak{}Gal.\allowbreak{}TP.\allowbreak{}atPlaceQ}{Development.\allowbreak{}AutomorphicArithmetic.\allowbreak{}Foundations.\allowbreak{}GaloisRepresentations.\allowbreak{}ThetaParseval.\allowbreak{}CoreAtQIffRowQ}{103--175}
\begin{proof}
Apply odd-rank compatibility to both lifts and realize the disjoint families as Galois characters. The cofinite Satake unions give \ref{B9.42}. Subtract by \ref{B9.43} and twist by $\kappa_\ell$ to cancel the root factor $\kappa(w)^{-1}$ and jump shift $m_{\iota\circ\tau}$.
\end{proof}

\begin{proposition}[All places]\label{B9.45}
With \ref{B9.44} available at every place of $Q$, weakly compatible representations exist and all have its local data throughout $Q$.
\end{proposition}
\leanref{UpperTheory.\allowbreak{}Gal.\allowbreak{}TP.\allowbreak{}subtractQ\_\allowbreak{}of}{Development.\allowbreak{}AutomorphicArithmetic.\allowbreak{}Foundations.\allowbreak{}GaloisRepresentations.\allowbreak{}ThetaParseval.\allowbreak{}CoreAtQIffRowQ}{200--210}
\leanref{UpperTheory.\allowbreak{}Gal.\allowbreak{}TP.\allowbreak{}weakIsoQ}{Development.\allowbreak{}AutomorphicArithmetic.\allowbreak{}Foundations.\allowbreak{}GaloisRepresentations.\allowbreak{}ThetaParseval.\allowbreak{}CoreAtQIffRowQ}{88--102}
\leanref{UpperTheory.\allowbreak{}Gal.\allowbreak{}Int.\allowbreak{}Core.\allowbreak{}subtractQ\_\allowbreak{}holds}{Development.\allowbreak{}AutomorphicArithmetic.\allowbreak{}Foundations.\allowbreak{}GaloisRepresentations.\allowbreak{}IntegralStructures.\allowbreak{}CrystallineDescent.\allowbreak{}LimOddOfLimFinOdd}{31--34}
\begin{proof}
Take the witness of \ref{B9.44} at the target. Remove both exceptional sets and the finitely many split places where $U$ has no Satake multiset. At each remaining place use the same multiset for the two representations. Density and Brauer--Nesbitt identify them.
\end{proof}


% ---- p9c2 ----
\chapter{A weak local Langlands correspondence}\label{ch:weak-local-langlands}

Let \(w\) be a finite place of a number field \(E\). Put \(K=E_w\), with residue cardinality \(q\), and fix a valuation ring \(A\) above \(w\) in an algebraic closure \(\Omega\) of \(E\). Let \(W_K\) be the Weil group, with inertia \(I_K\) and wild inertia \(P_K\subset I_K\). The degree \(\deg:W_K\to\Z\) sends arithmetic Frobenius to \(1\). Let \(\mathcal I(K)\) be the smooth irreducible complex classes \(x\) of \(\GL_{d(x)}(K)\), \(d(x)\ge0\). Cuspidal means supercuspidal of positive rank. In rank \(n\), put \(K_0=\GL_n(\mathcal O_K)\) and \(K_1=\ker(K_0\to\GL_n(\F_q))\).

Let \(\mathcal W(K)\) be the complex Weil--Deligne classes \(V=(\rho_V,N_V)\), with elementwise semisimple Weil action and open inertia kernel. The nilpotent \(N_V\) obeys \(\rho_V(\sigma)N_V\rho_V(\sigma)^{-1}=q^{\deg\sigma}N_V\). Irreducibility refers to \(\rho_V\), and tameness means \(P_K\) acts trivially.

Put \(\nu_K(z)=|z|_K=q^{-\operatorname{ord}_K z}\) and \(\nu_W(\sigma)=q^{\deg\sigma}\). Character twists of \(x\) use the determinant. For \(s\in\R\), write \(x\nu_K^s=x\otimes|\det|_K^s\) and \(V(s)=V\otimes\nu_W^s\). The self-dual \(\operatorname{Sp}_m\) has Weil characters \(\nu_W^{j-(m-1)/2}\), \(0\le j<m\), and one monodromy block.

\section{The existence theorem and its laws}

\begin{theorem}[Existence]\label{B9.46}
There is a map \(\operatorname{rec}_T:\mathcal I(K)\to\mathcal W(K)\) satisfying Definition~\ref{B9.47}.
\end{theorem}
\leanref{UpperTheory.\allowbreak{}Zero.\allowbreak{}LLC.\allowbreak{}recTExists}{Development.\allowbreak{}AutomorphicArithmetic.\allowbreak{}Foundations.\allowbreak{}LocalLanglands.\allowbreak{}LocalCentralUnits.\allowbreak{}RecTExists}{20--66}

\begin{proof}
Extend \(r\) of \ref{B9.53} to \(R\) by \ref{B9.54}. The block \(T_m(r(u))=r(u)((m-1)/2)\otimes\operatorname{Sp}_m\) has Weil constituents \(r(u\nu_K^j)=r(u)(j)\), \(0\le j<m\). For the cuspidal support \(s_x\), the multiset of segment terms, this gives
\[
 \det(X-\rho_{R(x)}(\sigma))
 =\prod_{v\in s_x}\det(X-\rho_{r(v)}(\sigma)).
 \tag{1}
\]
Degrees give the dimension law. For the central characters, \(\omega_x=\prod_{v\in s_x}\omega_v\) gives both determinant laws.

Twist the first terms for the twist laws. For good \(x\) the dual first term \((u\nu_K^{m-1})^\vee\) has parameter \(r(u)^\vee(1-m)\). Centering gives \(T_m(r(u))^\vee\) since \(\operatorname{Sp}_m\) is self-dual.

If \(R(x)=R(y)\) with \(y\) good, match the blocks and undo their shifts. By \ref{B9.53} the first terms agree, so \(x=y\).

If good \(x\) has \(K_1\)-fixed vectors, so does each support term by the depth-zero support theorem. By \ref{B9.53}, (1) equals \((X-1)^{d(x)}\) on \(P_K\), so semisimplicity gives trivial action. The rank-one and irreducibility laws follow from \(R=r\) on cuspidal classes. The inflated cuspidal type forces cuspidality, so the type law also follows from \(r\).
\end{proof}

\begin{definition}[The weak correspondence]\label{B9.47}
\(\Lean{IsRecT}\) requires an Artin map \(a\) as in Theorem~\ref{B9.50} and the following laws for \(R:\mathcal I(K)\to\mathcal W(K)\). For every \(x\) and \(s\in\R\),
\[
 \dim R(x)=d(x),\qquad R(x\nu_K^s)=R(x)(s).
\]
For smooth \(\chi:K^\times\to\C^\times\), \(R(\chi)=(\chi\circ a,0)\). For the Langlands multisegment \(\mathfrak m\) of \(x\),
\[
 R(x)=\bigoplus_{[u,m]\in\mathfrak m}
 R(u)((m-1)/2)\otimes\operatorname{Sp}_m .
 \tag{2}
\]
Cuspidal parameters are irreducible, and \(|\det\rho_{R(x)}|=|\omega_x\circ a|\) for every \(x\). For good \(x\) and smooth \(\chi\),
\[
 R(x\otimes\chi)=R(x)\otimes(\chi\circ a),\qquad
 \det\rho_{R(x)}=\omega_x\circ a,\qquad R(x^\vee)=R(x)^\vee .
\]
If \(R(x)=R(y)\) and \(y\) is good, then \(x=y\). A good \(x\) with \(K_1\)-fixed vectors has tame parameter. The map \(R\) satisfies the type law of \ref{B9.49}.
\end{definition}
\leanref{UpperTheory.\allowbreak{}Zero.\allowbreak{}LLC.\allowbreak{}LawsT}{Development.\allowbreak{}AutomorphicArithmetic.\allowbreak{}Foundations.\allowbreak{}LocalLanglands.\allowbreak{}IsRecTOfIsRec}{55--72}
\leanref{UpperTheory.\allowbreak{}Zero.\allowbreak{}LLC.\allowbreak{}IsRecT}{Development.\allowbreak{}AutomorphicArithmetic.\allowbreak{}Foundations.\allowbreak{}LocalLanglands.\allowbreak{}IsRecTOfIsRec}{76--76}

The full correspondence is bijective in positive rank and preserves all smooth twists and central characters \cite[Theorem~1.2(c--d)]{scholze-2013}\cite[\S4.4]{rodier-1982}.

\begin{definition}[Good cuspidal classes and good support]\label{B9.48}
A class \(u\) satisfies \(\Lean{GoodCusp}\) if \(d(u)=1\), or if \(K\simeq\Q_p\) continuously and \(u=u_0\otimes\chi\) with \(u_0^{K_1}\ne0\) and \(\chi\) smooth. Depth zero means \(u^{K_1}\ne0\). A class is \emph{good} if every term of its cuspidal support satisfies \(\Lean{GoodCusp}\).
\end{definition}
\leanref{UpperTheory.\allowbreak{}Zero.\allowbreak{}LLC.\allowbreak{}GoodCusp}{Development.\allowbreak{}AutomorphicArithmetic.\allowbreak{}Foundations.\allowbreak{}LocalLanglands.\allowbreak{}IsRecTOfIsRec}{20--22}
\leanref{UpperTheory.\allowbreak{}Zero.\allowbreak{}LLC.\allowbreak{}Good}{Development.\allowbreak{}AutomorphicArithmetic.\allowbreak{}Foundations.\allowbreak{}LocalLanglands.\allowbreak{}IsRecTOfIsRec}{24--24}

\begin{definition}[The depth-zero type law]\label{B9.49}
Fix a continuous isomorphism \(i:\Q_p\simeq K\) and \(n\ge1\). Let \(E'/E\) in \(\Omega\) have place \(w'\) induced by \(A\), with \(p\notin(w')^2\) and \(Nw'=p^n\). Choose local reciprocity \(a'\) at \(w'\) and a residue map \(\mathrm{red}:\mathcal O_{E'}\twoheadrightarrow k_n\) with kernel \(w'\) and \(|k_n|=p^n\). Take a character \(\theta:k_n^\times\to\C^\times\) of Frobenius orbit length \(n\) and an irreducible cuspidal \(\tau\) of \(\GL_n(\F_p)\) with
\[
 \operatorname{tr}\tau(g)=(-1)^{n-1}\sum_{j=0}^{n-1}\theta(z^{p^j})
 \tag{3}
\]
whenever \(g\) has irreducible characteristic polynomial with root \(z\in k_n^\times\). Suppose a smooth irreducible \(\pi\) of \(\GL_n(K)\) contains the inflation of \(\tau\) via \(i\) on \(\GL_n(\Z_p)\). Then \(N_{R(\pi)}=0\), and for \(t\in I_{E'_{w'}}\), \(b\in\mathcal O_{E'}\), and \(z=\mathrm{red}(b)\in k_n^\times\) with \(|a'(t)-b|_{w'}<1\),
\[
 \det(X-\rho_{R(\pi)}(t_E))=\prod_{j=0}^{n-1}(X-\theta(z^{p^j})).
 \tag{4}
\]
Here \(t_E\) denotes restriction of scalars to \(E\).
\end{definition}
\leanref{UpperTheory.\allowbreak{}Zero.\allowbreak{}LLC.\allowbreak{}TypeLaw}{Development.\allowbreak{}AutomorphicArithmetic.\allowbreak{}Foundations.\allowbreak{}LocalLanglands.\allowbreak{}IsRecTOfIsRec}{35--53}

Equation (3) is Green's formula \cite{green-1955}. The law uses the depth-zero inertia formula over \(\Q_p\) \cite[Addendum (A.1)]{silberger-zink-2008}.

\section{Cuspidal parameters}

\begin{theorem}[Local reciprocity]\label{B9.50}
There is a unique surjective homomorphism \(a:W_K\to K^\times\), continuous on \(I_K\), with
\[
 \operatorname{ord}_K(a(\sigma))=-\deg\sigma.
\]
Every smooth Weil character equals \(\chi\circ a\) for some smooth \(\chi\) of \(K^\times\). For each global Artin map \(\mathrm{Art}_E:\mathbf A_E^\times\to\Gal(\Omega/E)^{\mathrm{ab}}\),
\[
 \mathrm{Art}_E(\operatorname{loc}_w(a(\sigma)))=\overline{\sigma}.
\]
Here \(\sigma\in W_K\). The map \(\operatorname{loc}_w\) is the local idèle inclusion and \(\overline{\sigma}\) is the abelianized image.
\end{theorem}
\leanref{UpperTheory.\allowbreak{}Zero.\allowbreak{}LLC.\allowbreak{}artEx}{Development.\allowbreak{}AutomorphicArithmetic.\allowbreak{}Foundations.\allowbreak{}LocalLanglands.\allowbreak{}LocalCentralUnits.\allowbreak{}ArtEx}{7--15}
\leanref{GlobalTheory.\allowbreak{}LocLanglands.\allowbreak{}IsLocalArtin}{Development.\allowbreak{}AbelianVarietyTheory.\allowbreak{}PublishedTheorems.\allowbreak{}AutomorphicForms₂.\allowbreak{}HasExponentZeroIff}{145--152}

These are the local reciprocity laws in the stated Frobenius convention \cite{serre-1979,milne-2020}. The idèle identity determines \(a\) by injectivity of \(\mathrm{Art}_E\circ\operatorname{loc}_w\).

\begin{theorem}[Parameters on the good cuspidal classes]\label{B9.51}
Fix \(a\). Over \(K\simeq\Q_p\), every cuspidal \(u\) with \(u^{K_1}\ne0\) has a unique irreducible tame parameter \(V_u\) of dimension \(d(u)\) satisfying the type law for \(u\) and \(\det\rho_{V_u}=\omega_u\circ a\). The map \(u\mapsto V_u\) is injective.

The assignment \(r_0\) on good cuspidal classes commutes with smooth twists and is given by
\[
 r_0(\chi)=(\chi\circ a,0),\qquad
 r_0(u_0\otimes\chi)=V_{u_0}\otimes(\chi\circ a).
\]
Here \(\chi\) is smooth and \(u_0\) is depth zero.
\end{theorem}
\leanref{UpperTheory.\allowbreak{}Zero.\allowbreak{}LLC.\allowbreak{}parEx}{Development.\allowbreak{}AutomorphicArithmetic.\allowbreak{}Foundations.\allowbreak{}LocalLanglands.\allowbreak{}LocalCentralUnits.\allowbreak{}CentralUnit}{39--200}
\leanref{UpperTheory.\allowbreak{}Zero.\allowbreak{}LLC.\allowbreak{}parUniq}{Development.\allowbreak{}AutomorphicArithmetic.\allowbreak{}Foundations.\allowbreak{}LocalLanglands.\allowbreak{}LocalCentralUnits.\allowbreak{}ParUniq}{15--54}
\leanref{UpperTheory.\allowbreak{}Zero.\allowbreak{}LLC.\allowbreak{}goodMap}{Development.\allowbreak{}AutomorphicArithmetic.\allowbreak{}Foundations.\allowbreak{}LocalLanglands.\allowbreak{}LocalCentralUnits.\allowbreak{}GoodMap}{144--164}
\leanref{UpperTheory.\allowbreak{}Zero.\allowbreak{}LLC.\allowbreak{}twTame}{Development.\allowbreak{}AutomorphicArithmetic.\allowbreak{}Foundations.\allowbreak{}LocalLanglands.\allowbreak{}LocalCentralUnits.\allowbreak{}TwTame}{190--194}

\begin{proof}
Put \(n=d(u_0)\). Let \(\tau\) be the cuspidal \(\GL_n(\F_p)\)-representation on \(u_0^{K_1}\) with regular Green label \(\theta\) from \ref{B7.31}. Realize the unramified degree \(n\) extension of \(K\) as \(E'_{w'}\) and choose data \(D\) as in \ref{B9.49}. On common inertia \(R_D(t)=\overline{a'(t)}\in k_n^\times\) is onto. Set \(N=0\) and let degree-one Frobenius \(\phi\) cycle the lines of the distinct characters \(t\mapsto\theta(R_D(t)^{p^j})\), \(0\le j<n\). This gives irreducibility. Its matrix is determined up to diagonal conjugacy by the determinant \(\omega_{u_0}(a(\phi))\). On inertia,
\[
 \prod_{j=0}^{n-1}\theta(R_D(t)^{p^j})
 =\theta(N_{k_n/\F_p}R_D(t))=\omega_{u_0}(a(t)),
\]
by norm compatibility and the central character \(\theta|_{\F_p^\times}\) of \(\tau\).

Choose \(\beta\in\Omega\) with \(\beta^{p^n-1}=p\). The fields \(k_n\) identify with the same subfield of the residue field of \(A\), where tame reciprocity gives \(R_D(t)=\overline{t(\beta)/\beta}\). Density of \(\Q\) makes \(i\) unique, so the data use the same \(\tau\). For a generator \(\zeta\) of \(k_n^\times\), each \(\zeta^r\), \(1\le r\le n\), has degree \(n\). Indeed, each proper subfield \(\F_{p^d}\) has multiplicative index \((p^n-1)/(p^d-1)\ge 1+2^{n-d}>n\), since \(d\le n/2\). Apply (3) to the companion matrices of the minimal polynomials of these powers. The first \(n\) power sums determine \(\{\theta(\zeta^{p^j})\}_{j=0}^{n-1}\) by Newton's identities. Thus the transported labels lie in one Frobenius orbit and yield (4).

The inertia traces in (4) determine the inertia action. Together with the Frobenius determinant this gives uniqueness. Conversely, \(R_D\) is onto, so the inertia action gives the Green orbit and hence \(\tau\). The determinant recovers \(\omega_{u_0}\) since \(a\) is onto. Extend \(\tau\) by \(\widetilde\tau(zk)=\omega_{u_0}(z)\tau(\bar k)\) for \(z\in K^\times\), \(k\in K_0\). Then \(u_0\simeq\mathrm{c\text{-}Ind}_{K^\times K_0}^{\GL_n(K)}\widetilde\tau\) \cite[Proposition~6.6]{moy-prasad-1996}.

A tame twist \(\lambda\) multiplies the Green label by \(\bar\lambda\circ N_{k_n/\F_p}\), with \(\bar\lambda\) its reduction on units. This twists inertia by \(\lambda\circ a\). The central character becomes \(\omega_{u_0}\lambda^n\), so the Frobenius determinant gains \(\lambda(a(\phi))^n\). Hence \(V_{u_0\otimes\lambda}=V_{u_0}\otimes(\lambda\circ a)\).

The depth-zero twist criterion says that \(u_0\otimes\lambda\) has \(K_1\)-fixed vectors exactly when \(\lambda|_{1+\mathfrak p_K}=1\). If \(u_0\otimes\chi=u_1\otimes\chi_1\) for depth-zero bases, \(\chi\chi_1^{-1}\) is tame, so the formulas agree. In rank one \(V_{u_0}=(u_0\circ a,0)\).
\end{proof}

\begin{lemma}[Auxiliary irreducible parameters]\label{B9.52}
For \(n\ge2\) there is an irreducible \(J_n\) of dimension \(n\) with \(|\det\rho_{J_n}|=1\) on \(W_K\). If \(K\simeq\Q_p\), one can require \(J_n\ne V\otimes\xi\) for every tame parameter \(V\) and smooth Weil character \(\xi\).
\end{lemma}
\leanref{UpperTheory.\allowbreak{}Zero.\allowbreak{}LLC.\allowbreak{}junkEx}{Development.\allowbreak{}AutomorphicArithmetic.\allowbreak{}Foundations.\allowbreak{}LocalLanglands.\allowbreak{}LocalCentralUnits.\allowbreak{}JunkEx}{257--262}

\begin{proof}
Take \(L/K\) unramified of degree \(n\) with residue field \(k_n=\F_{q^n}\). For a uniformizer \(\varpi\) of \(K\), put \(M=q^n-1\) and define
\[
 b:\mathcal O_L^\times\longrightarrow k_n^+,\qquad
 b(u)=\overline{\frac{u^M-1}{M\varpi}}.
\]
Modulo \(\varpi^2\), multiplication gives additivity and \(b(1+\varpi x)=\bar x\). For the local Artin map \(a_L\), the map \(\pi=b\circ a_L|_{I_L}\) on \(I_L=I_K\) is onto with open kernel and satisfies
\[
 \pi(\phi t\phi^{-1})=f(\pi(t)),\qquad f(z)=z^q,
 \tag{5}
\]
for degree-one Frobenius \(\phi\), which fixes \(\varpi\) and acts by \(f\) on \(k_n\).

For the residue characteristic \(p\), take \(\lambda\in k_n\) of degree \(n\) over \(\F_q\). The character
\[
 \alpha(z)=\exp\!\left(\frac{2\pi i}{p}
 \operatorname{Tr}_{k_n/\F_p}(\lambda z)\right)
\]
has orbit length \(n\) under \(f\) by the trace pairing. Let inertia act on \(e_j\) by \(\eta_j=\alpha\circ f^j\circ\pi\), \(0\le j<n\). Set \(\rho(\phi)e_j=e_{j-1}\) for \(j>0\) and \(\rho(\phi)e_0=(-1)^{n-1}e_{n-1}\). Equation (5) gives the Weil relation. The Frobenius matrix has determinant \(1\). It has finite order and normalizes the finite inertia image, so \(\rho\) has finite image. Frobenius is transitive on the distinct inertia lines, so \(J_n=(\rho,0)\) is irreducible with unitary determinant.

The \(p\)-group \(k_n^+/\pi(P_K)\) is a quotient of tame inertia, hence trivial. Pick \(t\in P_K\) with \(\eta_0(t)\ne\eta_1(t)\). Then \(\rho_{J_n}(t)\) is nonscalar. Every character twist of a tame parameter has scalar action at \(t\).
\end{proof}
\leanref{UpperTheory.\allowbreak{}tameInertiaCharacter\_\allowbreak{}frobenius\_\allowbreak{}conj\_\allowbreak{}eq\_\allowbreak{}pow}{Development.\allowbreak{}AutomorphicArithmetic.\allowbreak{}WhittakerHecke.\allowbreak{}TameInertiaCharacterFrobeniusConjEqPow}{111--177}

\begin{proposition}[The assignment on all cuspidal classes]\label{B9.53}
Fix \(J_n\) as in Lemma~\ref{B9.52}. For cuspidal \(u\) of rank \(n\), define \(e(u)\in\R\) by \(|\omega_u|=\nu_K^{ne(u)}\). The assignment
\[
 r(u)=
 \begin{cases}
 r_0(u),&u\text{ satisfies }\Lean{GoodCusp},\\
 J_{d(u)}(e(u)),&\text{otherwise},
 \end{cases}
 \tag{6}
\]
satisfies the cuspidal laws of Definition~\ref{B9.47}. In particular, \(r(u)=r(v)\) with \(v\) good forces \(u=v\).
\end{proposition}
\leanref{UpperTheory.\allowbreak{}Zero.\allowbreak{}LLC.\allowbreak{}cuspEx}{Development.\allowbreak{}AutomorphicArithmetic.\allowbreak{}Foundations.\allowbreak{}LocalLanglands.\allowbreak{}LocalCentralUnits.\allowbreak{}CuspEx}{168--266}
\leanref{UpperTheory.\allowbreak{}Zero.\allowbreak{}LLC.\allowbreak{}expo}{Development.\allowbreak{}AutomorphicArithmetic.\allowbreak{}Foundations.\allowbreak{}LocalLanglands.\allowbreak{}LocalCentralUnits.\allowbreak{}ExponentMulAbsPow}{63--96}

\begin{proof}
Smooth central characters are unitary on units, so \(e(u)\) exists. In rank \(n=d(u)\) the central character of \(u\nu_K^s\) is \(\omega_u\nu_K^{ns}\). Thus
\[
 e(u\nu_K^s)=e(u)+s.
 \tag{7}
\]
Twists preserve goodness. On the second branch \(n\ge2\), (7) gives the real twist law. Twisting preserves dimension and irreducibility, and
\[
 |\det\rho_{J_n(e(u))}(\sigma)|
 =q^{ne(u)\deg\sigma}=|\omega_u(a(\sigma))|.
 \tag{8}
\]

Use \ref{B9.51} for twists and determinants on good classes. Duality inverts the rank-one character, or in higher rank the Green label and central character of a depth-zero base. Uniqueness and twisting give the duality law. For a good cuspidal class with \(K_1\)-fixed vectors, tameness follows from reciprocity in rank one and \ref{B9.51} in higher rank over \(K\simeq\Q_p\). The latter theorem gives the type law.

Suppose \(r(u)=r(v)\) with \(v\) good. The ranks agree, and in rank one \(u=v\) since \(a\) is onto. In higher rank, if \(u=u_0\otimes\chi\) and \(v=v_0\otimes\chi'\) are both good, put \(\lambda=\chi'\chi^{-1}\). Then \(V_{u_0}=V_{v_0}\otimes(\lambda\circ a)\). Tameness of both bases makes \(\lambda\circ a\) trivial on \(P_K\), hence \(\lambda\) trivial on \(1+\mathfrak p_K\). By \ref{B9.51}, \(u_0=v_0\otimes\lambda\), so \(u=v\).

For \(u\) in the second branch, a good \(v=v_0\otimes\chi\) of the same rank requires \(K\simeq\Q_p\). Equality would give
\[
 J_{d(u)}=V_{v_0}\otimes\bigl((\chi\circ a)\nu_W^{-e(u)}\bigr),
\]
contrary to Lemma~\ref{B9.52} since \(V_{v_0}\) is tame.
\end{proof}

\section{Multisegments and uniqueness}

\begin{theorem}[Langlands extension]\label{B9.54}
Each \(x\) has a unique Langlands multisegment. A segment \([u,m]\) consists of \(u\nu_K^j\), \(0\le j<m\), with \(u\) cuspidal and \(m\ge1\). Formula (2), using \(r(u)\) for each first term, extends any cuspidal assignment \(r\). The empty multisegment gives zero.
\end{theorem}
\leanref{UpperTheory.\allowbreak{}Zero.\allowbreak{}LLC.\allowbreak{}Pr.\allowbreak{}exh}{Development.\allowbreak{}AutomorphicArithmetic.\allowbreak{}Foundations.\allowbreak{}LocalLanglands.\allowbreak{}Multi}{11--11}
\leanref{UpperTheory.\allowbreak{}Zero.\allowbreak{}LLC.\allowbreak{}Pr.\allowbreak{}multi}{Development.\allowbreak{}AutomorphicArithmetic.\allowbreak{}Foundations.\allowbreak{}LocalLanglands.\allowbreak{}Multi}{12--12}
\leanref{GlobalTheory.\allowbreak{}LocLanglands.\allowbreak{}segSum}{Development.\allowbreak{}AbelianVarietyTheory.\allowbreak{}Definitions.\allowbreak{}NonGaloisDescent.\allowbreak{}Langlands}{17--18}
\leanref{GlobalTheory.\allowbreak{}LocLanglands.\allowbreak{}exists\_\allowbreak{}isLExt}{Development.\allowbreak{}AbelianVarietyTheory.\allowbreak{}Proofs.\allowbreak{}InertiaEigenvalues.\allowbreak{}LanglandsExtensionUniqueOnCuspidal}{34--56}

In the Langlands quotient parametrization \cite[Theorem~3]{rodier-1982}, each segment gives the unique irreducible quotient of normalized induction in increasing exponent order. On a common line of integer twists, \([a,b]\) precedes \([c,d]\) when \(a<c\le b+1\) and \(b<d\). Order the segments so that no earlier one precedes a later one. The normalized induction of their quotients has unique irreducible quotient \(x\). Zelevinsky gives the submodule parametrization \cite[Theorem~6.1]{zelevinsky-1980}. Formula (2) is permutation invariant and \([u,1]\) gives \(R(u)=r(u)\).

\begin{theorem}[Uniqueness on good representations]\label{B9.55}
Let \(R,R'\) satisfy \ref{B9.47}. Then \(R(x)=R'(x)\) for every good \(x\).
\end{theorem}
\leanref{UpperTheory.\allowbreak{}Zero.\allowbreak{}LLC.\allowbreak{}recTUnique}{Development.\allowbreak{}AutomorphicArithmetic.\allowbreak{}Foundations.\allowbreak{}LocalLanglands.\allowbreak{}LocalCentralUnits.\allowbreak{}RecTUnique}{28--53}

\begin{proof}
The common Artin map from \ref{B9.50} gives equality in rank one. Over \(K\simeq\Q_p\), use \ref{B9.51} on depth-zero cuspidal classes and the twist law on their twists. Then apply (2) to the first terms of good \(x\).
\end{proof}

\begin{lemma}[Support under the local tests]\label{B9.55a}
Let \(x\) be generic with \(N_{R(x)}=0\) for a map \(R\) as in \ref{B9.47}. If every Weil commutator has polynomial \((X-1)^{d(x)}\) on \(R(x)\), then \(x\) has unlinked rank-one support. Over \(K\simeq\Q_p\), trivial wild inertia on \(R(x)\) gives unlinked depth-zero support and \(x^{K_1}\ne0\). If also every repeated root \(z\) at a lift of a tame inertia generator satisfies \(z^q=z\), the support terms of rank at least two are pairwise inertially inequivalent. Both tests give good support.
\end{lemma}
\begin{proof}
Genericity gives an unlinked multisegment \cite[Theorem~9.7]{zelevinsky-1980}. In (2), zero monodromy forces singleton segments, so \(R(x)=\bigoplus_{u\in s_x}R(u)\). Under the toral test, semisimplicity and the eigenvalue \(1\) make every commutator act trivially. Each \(R(u)\) is then abelian and irreducible, hence one-dimensional.

In the tame case take a support term \(u\) of rank \(d\). Fix \(\phi\) of degree \(-1\) and \(D=\det\rho_{R(u)}(\phi)\). Tame irreducibles with these data inject into the Frobenius orbits of length \(d\) of characters of \(I_K/P_K\). Indeed, for a character of orbit length \(h\), an eigenline of \(\phi^h\) in its eigenspace gives an invariant span of dimension \(h\), so \(h=d\). The determinant fixes the last entry of the cyclic matrix after rescaling. Such characters factor through the quotient of order \(q^d-1\), identified by tame reciprocity with \(\F_{q^d}^{\times}\).

For each orbit take the Green type of \ref{B7.31}. Extend its central character by the value \(D\) at \(a(\phi)\) and compactly induce as in \ref{B9.51}. By \ref{B9.47}, the distinct good depth-zero classes so obtained have distinct tame images with the specified rank and determinant. The resulting injection of the finite orbit set into itself is onto. Thus \(R(u)\) equals the image of one of these good classes, which identifies \(u\) by \ref{B9.47}. The unlinked induction has \(K_1\)-fixed vectors by Iwasawa decomposition.

Inertially equivalent terms differ by an unramified twist, so their inertia polynomials agree by \ref{B9.47}. Their eigenvalues at a tame generator are repeated roots for \(R(x)\). The test fixes each under the \(q\)-power map, so the corresponding tame characters have orbit length one.
\end{proof}
\leanref{UpperTheory.\allowbreak{}Gal.\allowbreak{}Int.\allowbreak{}Row.\allowbreak{}plainTor\_\allowbreak{}of\_\allowbreak{}tor}{Development.\allowbreak{}AutomorphicArithmetic.\allowbreak{}Foundations.\allowbreak{}GaloisRepresentations.\allowbreak{}IntegralStructures.\allowbreak{}InertCharpoly.\allowbreak{}CharPolOfInert}{97--106}
\leanref{UpperTheory.\allowbreak{}Gal.\allowbreak{}Int.\allowbreak{}Row.\allowbreak{}dz\_\allowbreak{}of\_\allowbreak{}cusp}{Development.\allowbreak{}AutomorphicArithmetic.\allowbreak{}Foundations.\allowbreak{}GaloisRepresentations.\allowbreak{}IntegralStructures.\allowbreak{}InertCharpoly.\allowbreak{}CharPolOfInert}{114--139}
\leanref{UpperTheory.\allowbreak{}Gal.\allowbreak{}Int.\allowbreak{}Row.\allowbreak{}mixed\_\allowbreak{}of\_\allowbreak{}tame}{Development.\allowbreak{}AutomorphicArithmetic.\allowbreak{}Foundations.\allowbreak{}GaloisRepresentations.\allowbreak{}IntegralStructures.\allowbreak{}InertCharpoly.\allowbreak{}CharPolOfInert}{151--181}

In \ref{B9.07}, apply \ref{B9.55a} and \ref{B9.55} to \(x=\pi_w\nu_K^{(1-d(\pi_w))/2}\). For the integral type construction, \Lean{abConv} gives rank-one support when the inertia polynomials of \(R(x)\) equal those of a sum of Weil characters. Then \ref{B9.55} applies.


\begin{thebibliography}{999}\small

\bibitem{sutherland-2026} A.~V.~Sutherland, \emph{Verification files for the Lean proof of the Riemann Hypothesis}, \url{https://math.mit.edu/~drew/grh/}, 2026.

\bibitem{leanfro-2026} Lean and Mathlib communities, \emph{Lean verification of the Riemann Hypothesis}, joint statement endorsed by the Lean FRO and the Mathlib Initiative, \url{https://leanprover-community.github.io/blog/posts/riemann-hypothesis/}, 21 November 2026.

\bibitem{alpoge-2026} L.~Alpöge, remark on the proof and Lean archive, \url{https://alpo.ge/remark.pdf}, 20 November 2026.

\bibitem{apostol-1976} T.~M.~Apostol, \emph{Introduction to Analytic Number Theory}, Undergraduate Texts in Mathematics, Springer, New York, 1976, Chapter~12.

\bibitem{olver-2010} F.~W.~J.~Olver, D.~W.~Lozier, R.~F.~Boisvert and C.~W.~Clark (eds.), \emph{NIST Handbook of Mathematical Functions}, Cambridge University Press, 2010; online companion, DLMF, \S\S25.4, 25.15 and 5.2. % [Zeta reflection](https://dlmf.nist.gov/25.4), [Dirichlet L-functions](https://dlmf.nist.gov/25.15), [Gamma function](https://dlmf.nist.gov/5.2).

\bibitem{folland-1999} G.~B.~Folland, \emph{Real Analysis: Modern Techniques and Their Applications}, 2nd ed., John Wiley \& Sons (1999), Theorem~2.47, pp.~74--76. % [Text](https://upload.xuanyan.ws/Real%20Analysis%20Modern%20Techniques%20and%20Their%20Applications.pdf#page=87)

\bibitem{ahlfors-1979} L.~V.~Ahlfors, \emph{Complex Analysis}, 3rd ed., McGraw--Hill, 1979.

\bibitem{davenport-2000} H.~Davenport, \emph{Multiplicative Number Theory}, 3rd ed., revised by H.~L.~Montgomery, Springer, 2000.

\bibitem{widder-1941} D.~V.~Widder, \emph{The Laplace Transform}, Princeton Mathematical Series 6, Princeton University Press (1941), Chapter II, pp.~35--99.

\bibitem{rosser-schoenfeld-1962} J.~Barkley Rosser and L.~Schoenfeld, Approximate formulas for some functions of prime numbers, \emph{Illinois Journal of Mathematics} 6 (1962), 64--94.

\bibitem{montgomery-vaughan-2007} H.~L.~Montgomery and R.~C.~Vaughan, \emph{Multiplicative Number Theory I: Classical Theory}, Cambridge Studies in Advanced Mathematics 97, Cambridge University Press, 2007, \S2.2, Theorem~2.7(a). % [Publisher record](https://assets.cambridge.org/97805218/49036/frontmatter/9780521849036_frontmatter.pdf)

\bibitem{serre-1977} J.-P.~Serre, \emph{Linear Representations of Finite Groups}, Graduate Texts in Mathematics 42, Springer, New York, 1977, \S2.5, Theorem~7, and \S3.1, Theorem~9, pp.~10--26. % [Publisher records](https://link.springer.com/chapter/10.1007/978-1-4684-9458-7_2), (https://link.springer.com/chapter/10.1007/978-1-4684-9458-7_3)

\bibitem{benli-2026} K.~Benli, S.~Goel, H.~Twiss, and A.~Zaman, Explicit Deuring--Heilbronn phenomenon for Dirichlet \(L\)-functions, \emph{Proc. Amer. Math. Soc.} 154 (2026), no.~2, 509--525; revised text, arXiv:2410.06082v3.

\bibitem{thorner-2019} J.~Thorner and A.~Zaman, A unified and improved Chebotarev density theorem, \emph{Algebra Number Theory} 13 (2019), 1039--1068.

\bibitem{page-1935} A.~Page, On the number of primes in an arithmetic progression, \emph{Proc. London Math. Soc.} (2) 39 (1935), 116--141.

\bibitem{thorner-2024} J.~Thorner and A.~Zaman, Refinements to the prime number theorem for arithmetic progressions, \emph{Math. Z.} 306 (2024), no.~3, Paper No.~54.

\bibitem{montgomery-1971} H.~L.~Montgomery, \emph{Topics in Multiplicative Number Theory}, Lecture Notes in Mathematics 227, Springer, 1971.

\bibitem{hazewinkel-1979} M.~Hazewinkel, On formal groups. The functional equation lemma and some of its applications, \emph{Ast\'erisque} 63 (1979), 73--82.

\bibitem{krattenthaler-2015} C.~Krattenthaler and T.~W.~M\"uller, Truncated versions of Dwork's lemma for exponentials of power series and \(p\)-divisibility of arithmetic functions, \emph{Advances in Mathematics} 283 (2015), 489--529.

\bibitem{koblitz-1984} N.~Koblitz, \emph{\(p\)-adic Numbers, \(p\)-adic Analysis, and Zeta-Functions}, second edition, Graduate Texts in Mathematics 58, Springer-Verlag, New York (1984).

\bibitem{cohen-1946} I.~S.~Cohen and A.~Seidenberg, Prime ideals and integral dependence, \emph{Bull. Amer. Math. Soc.} 52 (1946), 252--261.

\bibitem{voight-2021} J.~Voight, Lattices and integral quadratic forms, in \emph{Quaternion Algebras}, Graduate Texts in Mathematics 288, Springer, Cham (2021), 139--154, Theorem 9.4.9. % [Publisher text](https://link.springer.com/chapter/10.1007/978-3-030-56694-4_9)

\bibitem{stacks-2026} The Stacks Project Authors, \emph{The Stacks Project} (2026).

\bibitem{fonseca-2019} T.~J.~Fonseca, Algebraic independence for values of integral curves, \emph{Algebra \& Number Theory} 13 (2019), 643--694, Proposition 5.5(3). \url{https://msp.org/ant/2019/13-3/ant-v13-n3-p03-s.pdf}.

\bibitem{stacks-2026b} The Stacks Project Authors, \emph{The Stacks Project}, More on Algebra, Remark 15.113.6, Tag 09E8 (2026). \url{https://stacks.math.columbia.edu/tag/09E8}.

\bibitem{bost-chen-2013} J.-B.~Bost and H.~Chen, Concerning the semistability of tensor products in Arakelov geometry, \emph{Journal de Mathématiques Pures et Appliquées} 99 (2013), 436--488. \url{https://webusers.imj-prg.fr/~huayi.chen/Recherche/bost_chen.pdf}.

\bibitem{gan-hanke-yu-2001} W.~T.~Gan, J.~P.~Hanke, and J.-K.~Yu, On an exact mass formula of Shimura, \emph{Duke Mathematical Journal} 107 (2001), 103--133. \url{https://web.math.princeton.edu/~jonhanke/Web-02/Mass-Formula-of-Shimura/shimura-mass.pdf}.

\bibitem{selberg-1967} A.~Selberg and S.~Chowla, On Epstein's zeta-function, \emph{J. Reine Angew. Math.} 227 (1967), 86--110. \url{https://doi.org/10.1515/crll.1967.227.86}.

\bibitem{banaszczyk-1993} W.~Banaszczyk, New bounds in some transference theorems in the geometry of numbers, \emph{Math. Ann.} 296 (1993), 625--635. \url{https://doi.org/10.1007/BF01445125}.

\bibitem{neukirch-1999} J.~Neukirch, \emph{Algebraic Number Theory}, Grundlehren der mathematischen Wissenschaften 322, Springer, 1999. \url{https://link.springer.com/book/10.1007/978-3-662-03983-0}.

\bibitem{gaudron-2008} \'E.~Gaudron, Pentes des fibr\'es vectoriels ad\'eliques sur un corps global, \emph{Rendiconti del Seminario Matematico della Universit\`a di Padova} 119 (2008), 21--95. \url{https://arxiv.org/abs/math/0605408}.

\bibitem{lieb-loss-2001} E.~H.~Lieb and M.~Loss, \emph{Analysis}, second edition, Graduate Studies in Mathematics 14, American Mathematical Society, Providence, RI, 2001, Theorem~1.13. \url{https://www.math.utoronto.ca/almut/MAT1000/LL-1.pdf}.

\bibitem{davenport-1980} H.~Davenport, \emph{Multiplicative Number Theory}, second edition, revised by H.~L.~Montgomery, Graduate Texts in Mathematics 74, Springer, New York, 1980. \url{https://doi.org/10.1007/978-1-4757-5927-3}.

\bibitem{ferrero-washington-1979} B.~Ferrero and L.~C.~Washington, The Iwasawa invariant $\mu_p$ vanishes for abelian number fields, \emph{Annals of Mathematics} (2) 109 (1979), 377--395. \url{https://doi.org/10.2307/1971116}

\bibitem{arthur-clozel-1989} J.~Arthur and L.~Clozel, \emph{Simple Algebras, Base Change, and the Advanced Theory of the Trace Formula}, Annals of Mathematics Studies 120, Princeton University Press, 1989. \url{https://www.claymath.org/library/cw/arthur/pdf/30.pdf}.

\bibitem{milne-2006} J.~S. Milne, \emph{Arithmetic Duality Theorems}, second edition, BookSurge, 2006, Chapter~I, \S\S2,4. \url{https://www.jmilne.org/math/Books/ADTnot.pdf}.

\bibitem{washington-1997} L.~C. Washington, \emph{Introduction to Cyclotomic Fields}, second edition, Graduate Texts in Mathematics 83, Springer, New York, 1997, Chapters~4--6 and Appendix, \S3. \url{https://doi.org/10.1007/978-1-4612-1934-7}.

\bibitem{borel-1974} A.~Borel, Stable real cohomology of arithmetic groups, \emph{Annales scientifiques de l'\'Ecole Normale Sup\'erieure} (4) 7 (1974), 235--272.

\bibitem{kallen-1980} W.~van der Kallen, Homology stability for linear groups, \emph{Inventiones Mathematicae} 60 (1980), 269--295. \url{https://doi.org/10.1007/BF01390018}

\bibitem{serre-1997} J.-P.~Serre, \emph{Galois Cohomology}, Springer Monographs in Mathematics, Springer, 1997. % [Publisher](https://link.springer.com/book/10.1007/978-3-642-59141-9)

\bibitem{milne-2020} J.~S.~Milne, \emph{Class Field Theory}, version 4.03, 2020. % [Author's text](https://www.jmilne.org/math/CourseNotes/CFT.pdf)

\bibitem{milne-2020a} J.~S.~Milne, \emph{Algebraic Number Theory}, version 3.08, 2020. % [Author's text](https://www.jmilne.org/math/CourseNotes/ANT.pdf)

\bibitem{kato-2004} K.~Kato, \(p\)-adic Hodge theory and values of zeta functions of modular forms, \emph{Ast\'erisque} 295 (2004), 117--290.

\bibitem{chenevier-harris-2013} G.~Chenevier and M.~Harris, Construction of automorphic Galois representations, II, \emph{Cambridge Journal of Mathematics} 1 (2013), 53--73. \url{https://doi.org/10.4310/CJM.2013.v1.n1.a2}.

\bibitem{green-1955} J.~A.~Green, The characters of the finite general linear groups, \emph{Transactions of the American Mathematical Society} 80 (1955), 402--447. \href{https://doi.org/10.1090/S0002-9947-1955-0072878-2}{doi:10.1090/S0002-9947-1955-0072878-2}.

\bibitem{lidl-niederreiter-1997} R.~Lidl and H.~Niederreiter, \emph{Finite Fields}, second edition, Encyclopedia of Mathematics and its Applications 20, Cambridge University Press, 1997.

\bibitem{zelevinsky-1981} A.~V.~Zelevinsky, \emph{Representations of Finite Classical Groups: A Hopf Algebra Approach}, Lecture Notes in Mathematics 869, Springer, 1981.

\bibitem{deriziotis-gotsis-1998} D.~I.~Deriziotis and C.~P.~Gotsis, The cuspidal modules of the finite general linear groups, \emph{LMS Journal of Computation and Mathematics} 1 (1998), 75--108.

\bibitem{gan-qiu-takeda-2014} W.~T.~Gan, Y.~Qiu and S.~Takeda, The regularized Siegel--Weil formula (the second term identity) and the Rallis inner product formula, \emph{Inventiones Mathematicae} 198 (2014), 739--831. \url{https://arxiv.org/abs/1207.4709}.

\bibitem{langlands-1976} R.~P.~Langlands, \emph{On the Functional Equations Satisfied by Eisenstein Series}, Lecture Notes in Mathematics 544, Springer-Verlag, Berlin--Heidelberg, 1976. \url{https://publications.ias.edu/node/39}.

\bibitem{weil-1965} A.~Weil, Sur la formule de Siegel dans la th\'eorie des groupes classiques, \emph{Acta Mathematica} 113 (1965), 1--87. \url{https://doi.org/10.1007/BF02391774}.

\bibitem{yamana-2013} S.~Yamana, The Siegel--Weil formula for unitary groups, \emph{Pacific Journal of Mathematics} 264 (2013), 235--256. \url{https://doi.org/10.2140/pjm.2013.264.235}.

\bibitem{scharlau-1985} W.~Scharlau, \emph{Quadratic and Hermitian Forms}, Grundlehren der mathematischen Wissenschaften 270, Springer-Verlag, Berlin--Heidelberg, 1985. \url{https://doi.org/10.1007/978-3-642-69971-9}.

\bibitem{moriwaki-2009} A.~Moriwaki, Continuity of volumes on arithmetic varieties, \emph{Journal of Algebraic Geometry} 18 (2009), 407--457. \href{https://arxiv.org/html/math/0612269v2}{Full text}.

\bibitem{zhang-1995} S.~Zhang, Positive line bundles on arithmetic varieties, \emph{J. Amer. Math. Soc.} 8 (1995), 187--221.

\bibitem{bernstein-2014} J.~Bernstein and B.~Kr\"otz, Smooth Fr\'echet globalizations of Harish-Chandra modules, \emph{Israel J. Math.} 199 (2014), 45--111. \url{https://www.math.tau.ac.il/~bernstei/Publication_list/publication_texts/Bern-Kroetz-2014.pdf}.

\bibitem{miller-2012} S.~D.~Miller and W.~Schmid, On the rapid decay of cuspidal automorphic forms, \emph{Adv. Math.} 231 (2012), 940--964. \url{https://arxiv.org/abs/1106.2149}.

\bibitem{neukirch-2008} J.~Neukirch, A.~Schmidt and K.~Wingberg, \emph{Cohomology of Number Fields}, second edition, Grundlehren der mathematischen Wissenschaften 323, Springer, Berlin, 2008, xv+825 pp., \url{https://link.springer.com/book/10.1007/978-3-540-37889-1}.

\bibitem{jacobowitz-1962} R.~Jacobowitz, Hermitian forms over local fields, \emph{American Journal of Mathematics} 84 (1962), 441--465. \url{https://www.jstor.org/stable/2372982}.

\bibitem{wallach-1992} N.~R.~Wallach, \emph{Real Reductive Groups II}, Pure and Applied Mathematics 132-II, Academic Press, Boston, 1992.

\bibitem{miller-schmid-2012} S.~D.~Miller and W.~Schmid, On the rapid decay of cuspidal automorphic forms, \emph{Advances in Mathematics} 231 (2012), 940--964. \url{https://doi.org/10.1016/j.aim.2012.06.001}.

\bibitem{shin-2011} S.~W.~Shin, Galois representations arising from some compact Shimura varieties, \emph{Annals of Mathematics} 173 (2011), 1645--1741.

\bibitem{caraiani-2012} A.~Caraiani, Local-global compatibility and the action of monodromy on nearby cycles, \emph{Duke Mathematical Journal} 161 (2012), 2311--2413.

\bibitem{caraiani-2014} A.~Caraiani, Monodromy and local-global compatibility for \(l=p\), \emph{Algebra \& Number Theory} 8 (2014), 1597--1646.

\bibitem{harris-taylor-2001} M.~Harris and R.~Taylor, \emph{The Geometry and Cohomology of Some Simple Shimura Varieties}, Annals of Mathematics Studies 151, Princeton University Press, 2001.

\bibitem{atiyah-bott-1968} M.~F.~Atiyah and R.~Bott, A Lefschetz fixed point formula for elliptic complexes, II: Applications, \emph{Annals of Mathematics} 88 (1968), 451--491.

\bibitem{milne-1980} J.~S.~Milne, \emph{Étale Cohomology}, Princeton Mathematical Series 33, Princeton University Press, 1980.

\bibitem{artin-grothendieck-verdier-1973} M.~Artin, A.~Grothendieck, and J.-L.~Verdier, \emph{Théorie des topos et cohomologie étale des schémas, Tome 3}, Lecture Notes in Mathematics 305, Springer, 1973.

\bibitem{bhatt-morrow-scholze-2018} B.~Bhatt, M.~Morrow, and P.~Scholze, Integral \(p\)-adic Hodge theory, \emph{Publications Mathématiques de l'IHÉS} 128 (2018), 219--397.

\bibitem{rouquier-1996} R.~Rouquier, Caractérisation des caractères et pseudo-caractères, \emph{Journal of Algebra} 180 (1996), 571--586.

\bibitem{taylor-1991} R.~Taylor, Galois representations associated to Siegel modular forms of low weight, \emph{Duke Mathematical Journal} 63 (1991), 281--332.

\bibitem{scholze-2013} P.~Scholze, The local Langlands correspondence for \(\GL_n\) over \(p\)-adic fields, \emph{Invent. Math.} 192 (2013), 663--715. \url{https://doi.org/10.1007/s00222-012-0420-5}.

\bibitem{silberger-zink-2008} A.~J.~Silberger and E.-W.~Zink, Explicit Shintani base change and the Macdonald correspondence for characters of \(\GL_n(k)\), \emph{J. Algebra} 319 (2008), 4147--4176. \url{https://doi.org/10.1016/j.jalgebra.2008.02.010}.

\bibitem{serre-1979} J.-P.~Serre, \emph{Local Fields}, Graduate Texts in Mathematics 67, Springer, 1979. \url{https://doi.org/10.1007/978-1-4757-5673-9}.

\bibitem{rodier-1982} F.~Rodier, Repr\'esentations de \(\GL(n,k)\) o\`u \(k\) est un corps \(p\)-adique, \emph{Ast\'erisque} 92--93 (1982), 201--218. \url{https://www.numdam.org/item/SB_1981-1982__24__201_0/}.

\bibitem{moy-prasad-1996} A.~Moy and G.~Prasad, Jacquet functors and unrefined minimal \(K\)-types, \emph{Comment. Math. Helv.} 71 (1996), 98--121. \url{https://doi.org/10.1007/BF02566411}.

\bibitem{zelevinsky-1980} A.~V.~Zelevinsky, Induced representations of reductive \(p\)-adic groups. II. On irreducible representations of \(\GL(n)\), \emph{Ann. Sci. \'Ecole Norm. Sup.} (4) 13 (1980), 165--210. \url{https://numdam.org/articles/10.24033/asens.1379/}.

\end{thebibliography}

\end{document}
